% Auto-generated from data/segments.json + layout.json (+ transforms) — do not edit by hand.
% volume: 27khu10
% mode: printing
% segments: 2353
% transform_rules: 158
% toc_marks: 96

% Volume layout defaults (layout.json ``layout``).
\csromanlayoutapply{1}{1.5}{21.6pt}{8pt}{8pt}{65pt}{2.5em}%

\setcsromanpitaka{สุตฺตนฺตปิฏก}
\setcsromannikaya{ขุทฺทกนิกาย}
\csromanheader{2}{\textbf{ขุทฺทกนิกาย}}
\csromanmatikaanchor{mtk.0}
\csromantocmarknopage{chapter}{เนตฺติปาฬิ}{mtk.0}
\bookmark[dest={mtk.0},level=0]{เนตฺติปาฬิ}
\gambhira{\textbf{เนตฺติปาฬิ}}
\namakkaram{นโม ตสฺส ภควโต อรหโต สมฺมา\-สมฺพุทฺธสฺส.}
\csromanmatikaanchor{mtk.1}
\csromantocmarkref{section}{1. สงฺคหวาร}{mtk.1}
\bookmark[dest={mtk.1},level=1]{1. สงฺคหวาร}
\csromanheader{1}{1. \textbf{สงฺคหวาร}}
\setlength{\csromangathapagewidth}{0pt}
\setlength{\csromangathawakwidth}{0pt}
\csromangathameasuregroup{ยํ โลโก ปูชยเต, \\ ตสฺเสต สาสนวรํ, \\ ทฺวาทส ปทานิ สุตฺตํ, \\ ตํ วิญฺเญยฺยํ อุภยํ, \\ โสฬสหารา เนตฺติ\textsuperscript{1}, \\ อฏฺฐารสมูลปทา, \\ หารา พฺยญฺชนวิจโย, \\ อุภยํ ปริคฺคหีตํ, \\ ยา เจว เทสนา ยญฺจ, \\ ตตฺรายมานุปุพฺพี,}{\csromangathabat{ยํ โลโก ปูชยเต,}{สโลกปาโล สทา นมสฺสติ จ.} \\ \csromangathabat{ตสฺเสต สาสนวรํ,}{วิทูหิ เญยฺยํ นรวรสฺส.} \\ \csromangathabat{ทฺวาทส ปทานิ สุตฺตํ,}{ตํ สพฺพํ พฺยญฺชนญฺจ อตฺโถ จ.} \\ \csromangathabat{ตํ วิญฺเญยฺยํ อุภยํ,}{โก อตฺโถ พฺยญฺชนํ กตมํ.} \\ \csromangathabat{โสฬสหารา เนตฺติ\textsuperscript{1},}{ปญฺจนยา สาสนสฺส ปริเยฏฺฐิ.} \\ \csromangathabat{อฏฺฐารสมูลปทา,}{มหกจฺจาเนน\textsuperscript{1} นิทฺทิฏฺฐา\textsuperscript{1}.} \\ \csromangathabat{หารา พฺยญฺชนวิจโย,}{สุตฺตสฺส นยา ตโย จ สุตฺตตฺโถ.} \\ \csromangathabat{อุภยํ ปริคฺคหีตํ,}{วุจฺจติ สุตฺตํ ยถาสุตฺตํ.} \\ \csromangathabat{ยา เจว เทสนา ยญฺจ,}{เทสิตํ อุภยเมว วิญฺเญยฺยํ.} \\ \csromangathabat{ตตฺรายมานุปุพฺพี,}{นววิธสุตฺตนฺตปริเยฏฺฐีติ.}}
\csromangathasetleft{ยํ โลโก ปูชยเต, \\ ตสฺเสต สาสนวรํ, \\ ทฺวาทส ปทานิ สุตฺตํ, \\ ตํ วิญฺเญยฺยํ อุภยํ, \\ โสฬสหารา เนตฺติ\textsuperscript{1}, \\ อฏฺฐารสมูลปทา, \\ หารา พฺยญฺชนวิจโย, \\ อุภยํ ปริคฺคหีตํ, \\ ยา เจว เทสนา ยญฺจ, \\ ตตฺรายมานุปุพฺพี,}
\csromangathagroup{\csromangathastanza{\csromanfolio{1}\csromangathabat{ยํ โลโก ปูชยเต,}{สโลกปาโล สทา นมสฺสติ จ.} \\ \csromangathabat{ตสฺเสต สาสนวรํ,}{วิทูหิ เญยฺยํ นรวรสฺส.}}\csromangathastanza{\csromangathabat{ทฺวาทส ปทานิ สุตฺตํ,}{ตํ สพฺพํ พฺยญฺชนญฺจ อตฺโถ จ.} \\ \csromangathabat{ตํ วิญฺเญยฺยํ อุภยํ,}{โก อตฺโถ พฺยญฺชนํ กตมํ.}}\csromangathastanza{\csromangathabat{โสฬสหารา เนตฺติ\csromansharedfootnote{3a5c5cc19cc401bf}{เนตฺตี (ก)},}{ปญฺจนยา สาสนสฺส ปริเยฏฺฐิ.} \\ \csromangathabat{อฏฺฐารสมูลปทา,}{มหกจฺจาเนน\csromansharedfootnote{3a5c5cc19cc401bf}{เนตฺตี (ก)} นิทฺทิฏฺฐา\csromansharedfootnote{3a5c5cc19cc401bf}{เนตฺตี (ก)}.}}\csromangathastanza{\csromangathabat{หารา พฺยญฺชนวิจโย,}{สุตฺตสฺส นยา ตโย จ สุตฺตตฺโถ.} \\ \csromangathabat{อุภยํ ปริคฺคหีตํ,}{วุจฺจติ สุตฺตํ ยถาสุตฺตํ.}}\csromangathastanza{\csromangathabat{ยา เจว เทสนา ยญฺจ,}{เทสิตํ อุภยเมว วิญฺเญยฺยํ.} \\ \csromangathabat{ตตฺรายมานุปุพฺพี,}{นววิธสุตฺตนฺตปริเยฏฺฐีติ.}}}

\vspace{\tipitakaparskip}
\csromancenter{สงฺคหวาโร.}
\csromanmatikaanchor{mtk.2}
\csromantocmarkref{section}{2. อุทฺเทสวาร}{mtk.2}
\bookmark[dest={mtk.2},level=1]{2. อุทฺเทสวาร}
\csromanheader{1}{2. \textbf{อุทฺเทสวาร}}
\proseitem{1}{ตตฺถ กตเม โสฬส หารา, เทสนา วิจโย ยุตฺติ ปทฏฺฐาโน ลกฺขโณ จตุพฺยูโห อาวฏฺโฏ วิภตฺติ ปริวตฺตโน เววจโน \csromanfolio{2}ปญฺญตฺติ โอตรโณ โสธโน อธิฏฺฐาโน ปริกฺขาโร สมาโรปโน อิติ.}

\csromanheader{2}{\textbf{ตสฺสานุคีติ}}
\prose{\csromansymbolfootnote{*}{อุปริ 169 ปิฏฺเฐ เปฏโกปเทเสปิ.}เทสนา วิจโย ยุตฺติ, ปทฏฺฐาโน จ ลกฺขโณ.}

\prose{จตุพฺยูโห จ อาวฏฺโฏ, วิภตฺติ ปริวตฺตโน.}

\prose{\csromansymbolmark{*}เววจโน จ ปญฺญตฺติ, โอตรโณ จ โสธโน.}

\prose{อธิฏฺฐาโน ปริกฺขาโร, สมาโรปโน โสฬโส\csromansharedfootnote{85a1fd0848f2171c}{โสฬส (สี)}.}

\prose{เอเต โสฬส หารา, ปกิตฺติตา อตฺถโต อสํกิณฺณา.}

\prose{เอเตสญฺเจว ภวติ, วิตฺถารตยา นยวิภตฺตีติ.}

\proseitem{1}{ตตฺถ กตเม ปญฺจ นยา, นนฺทิยาวฏฺโฏ ติปุกฺขโล สีหวิกฺกีฬิโต ทิสาโลจโน องฺกุโส อิติ. \textbf{ตสฺสานุคีติ}}

\setlength{\csromangathapagewidth}{0pt}
\setlength{\csromangathawakwidth}{0pt}
\csromangathameasuregroup{\textsuperscript{+}ปฐโม นนฺทิยาวฏฺโฏ, \\ สีหวิกฺกีฬิโต นาม, \\ \textsuperscript{+}ทิสาโลจนมาหํสุ, \\ ปญฺจโม องฺกุโส นาม,}{\csromangathabat{\textsuperscript{+}ปฐโม นนฺทิยาวฏฺโฏ,}{ทุติโย จ ติปุกฺขโล.} \\ \csromangathabat{สีหวิกฺกีฬิโต นาม,}{ตติโย นยลญฺชโก\textsuperscript{1}.} \\ \csromangathabat{\textsuperscript{+}ทิสาโลจนมาหํสุ,}{จตุตฺถํ นยมุตฺตมํ.} \\ \csromangathabat{ปญฺจโม องฺกุโส นาม,}{สพฺเพ ปญฺจ นยา คตาติ.}}
\csromangathasetleft{\textsuperscript{+}ปฐโม นนฺทิยาวฏฺโฏ, \\ สีหวิกฺกีฬิโต นาม, \\ \textsuperscript{+}ทิสาโลจนมาหํสุ, \\ ปญฺจโม องฺกุโส นาม,}
\csromangathagroup{\csromangathastanza{\csromangathaitembat{1}{\csromansymbolfootnote{+}{อุปริ 169 ปิฏฺเฐ เปฏโกปเทเส โถกํ วิสทิสํ.}ปฐโม นนฺทิยาวฏฺโฏ,}{ทุติโย จ ติปุกฺขโล.} \\ \csromangathacontbat{สีหวิกฺกีฬิโต นาม,}{ตติโย นยลญฺชโก\csromansharedfootnote{408e12090260cc39}{นยลญฺฉโก (สี)}.}}\csromangathastanza{\csromangathaitembat{1}{\csromansymbolmark{+}ทิสา\-โลจนมาหํสุ,}{จตุตฺถํ นยมุตฺตมํ.} \\ \csromangathacontbat{ปญฺจโม องฺกุโส นาม,}{สพฺเพ ปญฺจ นยา คตาติ.}}}

\prose{ตตฺถ กตมานิ อฏฺฐารส มูลปทานิ, นว ปทานิ กุสลานิ นว ปทานิ อกุสลานิ. ตตฺถ กตมานิ นว ปทานิ อกุสลานิ, ตณฺหา อวิชฺชา โลโภ โทโส โมโห สุภสญฺญา สุขสญฺญา นิจฺจสญฺญา อตฺตสญฺญาติ, อิมานิ นว ปทานิ อกุสลานิ, ยตฺถ สพฺโพ อกุสลปกฺโข สงฺคหํ สโมสรณํ คจฺฉติ.}

\prose{ตตฺถ กตมานิ นว ปทานิ กุสลานิ, สมโถ วิปสฺสนา อโลโภ อโทโส อโมโห อสุภสญฺญา ทุกฺขสญฺญา อนิจฺจสญฺญา อนตฺตสญฺญาติ, อิมานิ นว ปทานิ กุสลานิ, ยตฺถ สพฺโพ กุสลปกฺโข สงฺคหํ สโมสรณํ คจฺฉติ.}

\csromanheader{2}{3. นิทฺเทสวาร \textbf{ตตฺริทํ อุทฺทานํ}}
\setlength{\csromangathapagewidth}{0pt}
\setlength{\csromangathawakwidth}{0pt}
\csromangathameasuregroup{ตณฺหา จ อวิชฺชาปิ จ, \\ จตุโร จ วิปลฺลาสา, \\ สมโถ จ วิปสฺสนา จ, \\ จตุโร จติปฏฺฐานา, \\ นวหิ จ ปเทหิ กุสลา, \\ เอเต โข มูลปทา,}{\csromangathabat{ตณฺหา จ อวิชฺชาปิ จ,}{โลโภ โทโส ตเถว โมโห จ.} \\ \csromangathabat{จตุโร จ วิปลฺลาสา,}{กิเลสภูมี นว ปทานิ.} \\ \csromangathabat{สมโถ จ วิปสฺสนา จ,}{กุสลานิ จ ยานิ ตีณิ มูลานิ.} \\ \csromangathabat{จตุโร จติปฏฺฐานา,}{อินฺทฺริยภูมี นว ปทานิ.} \\ \csromangathabat{นวหิ จ ปเทหิ กุสลา,}{นวหิ จ ยุชฺชนฺติ อกุสลปกฺขา.} \\ \csromangathabat{เอเต โข มูลปทา,}{ภวนฺติ อฏฺฐารส ปทานีติ.} \\ อุทฺเทสวาโร.}
\csromangathasetleft{ตณฺหา จ อวิชฺชาปิ จ, \\ จตุโร จ วิปลฺลาสา, \\ สมโถ จ วิปสฺสนา จ, \\ จตุโร จติปฏฺฐานา, \\ นวหิ จ ปเทหิ กุสลา, \\ เอเต โข มูลปทา,}
\csromangathagroup{\csromangathastanza{\csromanfolio{3}\csromangathabat{ตณฺหา จ อวิชฺชาปิ จ,}{โลโภ โทโส ตเถว โมโห จ.} \\ \csromangathabat{จตุโร จ วิปลฺลาสา,}{กิเลสภูมี นว ปทานิ.}}\csromangathastanza{\csromangathabat{สมโถ จ วิปสฺสนา จ,}{กุสลานิ จ ยานิ ตีณิ มูลานิ.} \\ \csromangathabat{จตุโร จติปฏฺฐานา,}{อินฺทฺริยภูมี นว ปทานิ.}}\csromangathastanza{\csromangathabat{นวหิ จ ปเทหิ กุสลา,}{นวหิ จ ยุชฺชนฺติ อกุสลปกฺขา.} \\ \csromangathabat{เอเต โข มูลปทา,}{ภวนฺติ อฏฺฐารส ปทานีติ.} \\ อุทฺเทสวาโร.}}

\csromanmatikaanchor{mtk.3}
\csromantocmarkref{section}{3. นิทฺเทสวาร}{mtk.3}
\bookmark[dest={mtk.3},level=1]{3. นิทฺเทสวาร}
\csromanheader{1}{3. นิทฺเทสวาร}
\vspace{\tipitakaparskip}
\csromancenter{ตตฺถ สงฺเขปโต เนตฺติ กิตฺติตา.}
\csromanheader{2}{\textbf{หารสงฺเขป}}
\setlength{\csromangathapagewidth}{0pt}
\setlength{\csromangathawakwidth}{0pt}
\csromangathameasuregroup{อสฺสาทาทีนวตา, \\ อาณตฺตี จ ภควโต, \\ ยํ ปุจฺฉิตญฺจ วิสฺสชฺชิตญฺจ, \\ สุตฺตสฺส โย ป\textbf{วิจโย}, \\ สพฺเพสํ หารานํ, \\ ยุตฺตายุตฺตปริกฺขา, \\ ธมฺมํ เทเสติ ชิโน, \\ อิติ ยาว สพฺพธมฺมา, \\ วุตฺตมฺหิ เอกธมฺเม, \\ วุตฺตา ภวนฺติ สพฺเพ, \\ เนรุตฺตมธิปฺปาโย, \\ ปุพฺพาปรานุสนฺธี, \\ เอกมฺหิ ปทฏฺฐาเน, \\ อาวฏฺฏติ ปฏิปกฺเข, \\ กุสลากุสเล ธมฺเม, \\ ปริวตฺตติ ปฏิปกฺเข, \\ เววจนานิ พหูนิ ตุ, \\ โย ชานาติ สุตฺตวิทู, \\ เอกํ ภควา ธมฺมํ, \\ โส อากาโร เญยฺโย, \\ โย จ ปฏิจฺจุปฺปาโท, \\ เอเตหิ โอตรติ โย, \\ วิสฺสชฺชิตมฺหิ ปเญฺห, \\ สุทฺธาสุทฺธปริกฺขา, \\ เอกตฺตตาย ธมฺมา, \\ เต น วิกปฺปยิตพฺพา, \\ เย ธมฺมา ยํ ธมฺมํ, \\ เหตุมวกฑฺฒยิตฺวา, \\ เย ธมฺมา ยํมูลา, \\ เต สมโรปยิตพฺพา,}{\csromangathabat{อสฺสาทาทีนวตา,}{นิสฺสรณมฺปิ จ ผลํ อุปาโย จ.} \\ \csromangathabat{อาณตฺตี จ ภควโต,}{โยคีนํ \textbf{เทสนาหาโร.}} \\ \csromangathabat{ยํ ปุจฺฉิตญฺจ วิสฺสชฺชิตญฺจ,}{สุตฺตสฺส ยา จ อนุคีติ.} \\ \csromangathabat{สุตฺตสฺส โย ป\textbf{วิจโย},}{หาโร \textbf{วิจโย}ติ นิทฺทิฏฺโฐ.} \\ \csromangathabat{สพฺเพสํ หารานํ,}{ยา ภูมี โย จ โคจโร เตสํ.} \\ \csromangathabat{ยุตฺตายุตฺตปริกฺขา,}{หาโร \textbf{ยุตฺตี}ติ นิทฺทิฏฺโฐ.} \\ \csromangathabat{ธมฺมํ เทเสติ ชิโน,}{ตสฺส จ ธมฺมสฺส ยํ ปทฏฺฐานํ.} \\ \csromangathabat{อิติ ยาว สพฺพธมฺมา,}{เอโส หาโร \textbf{ปทฏฺฐาโน.}} \\ \csromangathabat{วุตฺตมฺหิ เอกธมฺเม,}{เย ธมฺมา เอกลกฺขณา เกจิ.} \\ \csromangathabat{วุตฺตา ภวนฺติ สพฺเพ,}{โส หาโร \textbf{ลกฺขโณ} นาม.} \\ \csromangathabat{เนรุตฺตมธิปฺปาโย,}{พฺยญฺชนมถ เทสนานิทานญฺจ.} \\ \csromangathabat{ปุพฺพาปรานุสนฺธี,}{เอโส หาโร \textbf{จตุพฺยูโห}.} \\ \csromangathabat{เอกมฺหิ ปทฏฺฐาเน,}{ปริเยสติ เสสกํ ปทฏฺฐานํ.} \\ \csromangathabat{อาวฏฺฏติ ปฏิปกฺเข,}{\textbf{อาวฏฺโฏ} นาม โส หาโร.} \\ ธมฺมญฺจ ปทฏฺฐานํ, ภูมิญฺจ วิภชฺชเต อยํ หาโร. \\ สาธารเณ อสาธารเณ จ เนยฺโย \textbf{วิภตฺตี}ติ. \\ \csromangathabat{กุสลากุสเล ธมฺเม,}{นิทฺทิฏฺเฐ ภาวิเต ปหีเน จ.} \\ \csromangathabat{ปริวตฺตติ ปฏิปกฺเข,}{หาโร \textbf{ปริวตฺตโน} นาม.} \\ \csromangathabat{เววจนานิ พหูนิ ตุ,}{สุตฺเต วุตฺตานิ เอกธมฺมสฺส.} \\ \csromangathabat{โย ชานาติ สุตฺตวิทู,}{\textbf{เววจโน} นาม โส หาโร.} \\ \csromangathabat{เอกํ ภควา ธมฺมํ,}{\textbf{ปญฺญตฺตี}หิ วิวิธาหิ เทเสติ.} \\ \csromangathabat{โส อากาโร เญยฺโย,}{\textbf{ปญฺญตฺตี} นาม หาโรติ.} \\ \csromangathabat{โย จ ปฏิจฺจุปฺปาโท,}{อินฺทฺริยขนฺธา จ ธาตุ อายตนา.} \\ \csromangathabat{เอเตหิ โอตรติ โย,}{\textbf{โอตรโณ} นาม โส หาโร.} \\ \csromangathabat{วิสฺสชฺชิตมฺหิ ปเญฺห,}{คาถายํ ปุจฺฉิตา ยมารพฺภ.} \\ \csromangathabat{สุทฺธาสุทฺธปริกฺขา,}{หาโร โส \textbf{โสธโน} นาม.} \\ \csromangathabat{เอกตฺตตาย ธมฺมา,}{เยปิ จ เวมตฺตตาย นิทฺทิฏฺฐา.} \\ \csromangathabat{เต น วิกปฺปยิตพฺพา,}{เอโส หาโร \textbf{อธิฏฺฐาโน.}} \\ \csromangathabat{เย ธมฺมา ยํ ธมฺมํ,}{ชนยนฺติปฺปจฺจยา ปรมฺปรโต.} \\ \csromangathabat{เหตุมวกฑฺฒยิตฺวา,}{เอโส หาโร \textbf{ปริกฺขาโร.}} \\ \csromangathabat{เย ธมฺมา ยํมูลา,}{เย เจกตฺถา ปกาสิตา มุนินา.} \\ \csromangathabat{เต สมโรปยิตพฺพา,}{เอส \textbf{สมาโรปโน} หาโร.}}
\csromangathameasurewak{ธมฺมญฺจ ปทฏฺฐานํ, ภูมิญฺจ วิภชฺชเต อยํ หาโร. \\ สาธารเณ อสาธารเณ จ เนยฺโย \textbf{วิภตฺตี}ติ.}
\csromangathasetleft{อสฺสาทาทีนวตา, \\ อาณตฺตี จ ภควโต, \\ ยํ ปุจฺฉิตญฺจ วิสฺสชฺชิตญฺจ, \\ สุตฺตสฺส โย ป\textbf{วิจโย}, \\ สพฺเพสํ หารานํ, \\ ยุตฺตายุตฺตปริกฺขา, \\ ธมฺมํ เทเสติ ชิโน, \\ อิติ ยาว สพฺพธมฺมา, \\ วุตฺตมฺหิ เอกธมฺเม, \\ วุตฺตา ภวนฺติ สพฺเพ, \\ เนรุตฺตมธิปฺปาโย, \\ ปุพฺพาปรานุสนฺธี, \\ เอกมฺหิ ปทฏฺฐาเน, \\ อาวฏฺฏติ ปฏิปกฺเข, \\ กุสลากุสเล ธมฺเม, \\ ปริวตฺตติ ปฏิปกฺเข, \\ เววจนานิ พหูนิ ตุ, \\ โย ชานาติ สุตฺตวิทู, \\ เอกํ ภควา ธมฺมํ, \\ โส อากาโร เญยฺโย, \\ โย จ ปฏิจฺจุปฺปาโท, \\ เอเตหิ โอตรติ โย, \\ วิสฺสชฺชิตมฺหิ ปเญฺห, \\ สุทฺธาสุทฺธปริกฺขา, \\ เอกตฺตตาย ธมฺมา, \\ เต น วิกปฺปยิตพฺพา, \\ เย ธมฺมา ยํ ธมฺมํ, \\ เหตุมวกฑฺฒยิตฺวา, \\ เย ธมฺมา ยํมูลา, \\ เต สมโรปยิตพฺพา,}
\csromangathagroup{\csromangathastanza{\csromangathaitembat{1}{อสฺสาทาทีนวตา,}{นิสฺสรณมฺปิ จ ผลํ อุปาโย จ.} \\ \csromangathacontbat{อาณตฺตี จ ภควโต,}{โยคีนํ \textbf{เทสนาหาโร.}}}\csromangathastanza{\csromangathaitembat{2}{ยํ ปุจฺฉิตญฺจ วิสฺสชฺชิตญฺจ,}{สุตฺตสฺส ยา จ อนุคีติ.} \\ \csromangathacontbat{สุตฺตสฺส โย ป\textbf{วิจโย},}{หาโร \textbf{วิจโย}ติ นิทฺทิฏฺโฐ.}}\csromangathastanza{\csromangathaitembat{3}{สพฺเพสํ หารานํ,}{ยา ภูมี โย จ โคจโร เตสํ.} \\ \csromangathacontbat{ยุตฺตายุตฺตปริกฺขา,}{หาโร \textbf{ยุตฺตี}ติ นิทฺทิฏฺโฐ.}}\csromangathastanza{\csromangathaitembat{4}{ธมฺมํ เทเสติ ชิโน,}{ตสฺส จ ธมฺมสฺส ยํ ปทฏฺฐานํ.} \\ \csromangathacontbat{อิติ ยาว สพฺพธมฺมา,}{เอโส หาโร \textbf{ปทฏฺฐาโน.}}}\csromangathastanza{\csromangathaitembat{5}{วุตฺตมฺหิ เอกธมฺเม,}{เย ธมฺมา เอกลกฺขณา เกจิ.} \\ \csromangathacontbat{วุตฺตา ภวนฺติ สพฺเพ,}{โส หาโร \textbf{ลกฺขโณ} นาม.}}\csromangathastanza{\csromangathaitembat{6}{เนรุ\-ตฺตม\-ธิปฺปาโย,}{พฺยญฺชนมถ เทสนา\-นิทานญฺจ.} \\ \csromangathacontbat{ปุพฺพาปรานุสนฺธี,}{เอโส หาโร \textbf{จตุพฺยูโห}.}}\csromangathastanza{\csromangathaitembat{7}{เอกมฺหิ ปทฏฺฐาเน,}{ปริเยสติ เสสกํ ปทฏฺฐานํ.} \\ \csromangathacontbat{อาวฏฺฏติ ปฏิปกฺเข,}{\textbf{อาวฏฺโฏ} นาม โส หาโร.}}\csromangathastanza{\csromangathawak{\csromanfolio{4}\csromangathaitemline{8}{ธมฺมญฺจ ปทฏฺฐานํ, ภูมิญฺจ วิภชฺชเต อยํ หาโร.} \\ \csromangathacontline{สาธารเณ อสาธารเณ จ เนยฺโย \textbf{วิภตฺตี}ติ.}}}\csromangathastanza{\csromangathaitembat{9}{กุสลากุสเล ธมฺเม,}{นิทฺทิฏฺเฐ ภาวิเต ปหีเน จ.} \\ \csromangathacontbat{ปริวตฺตติ ปฏิปกฺเข,}{หาโร \textbf{ปริวตฺตโน} นาม.}}\csromangathastanza{\csromangathaitembat{10}{เววจนานิ พหูนิ ตุ,}{สุตฺเต วุตฺตานิ เอกธมฺมสฺส.} \\ \csromangathacontbat{โย ชานาติ สุตฺตวิทู,}{\textbf{เววจโน} นาม โส หาโร.}}\csromangathastanza{\csromangathaitembat{11}{เอกํ ภควา ธมฺมํ,}{\textbf{ปญฺญตฺตี}หิ วิวิธาหิ เทเสติ.} \\ \csromangathacontbat{โส อากาโร เญยฺโย,}{\textbf{ปญฺญตฺตี} นาม หาโรติ.}}\csromangathastanza{\csromangathaitembat{12}{โย จ ปฏิจฺจุปฺปาโท,}{อินฺทฺริยขนฺธา จ ธาตุ อายตนา.} \\ \csromangathacontbat{เอเตหิ โอตรติ โย,}{\textbf{โอตรโณ} นาม โส หาโร.}}\csromangathastanza{\csromangathaitembat{13}{วิสฺสชฺชิตมฺหิ ปเญฺห,}{คาถายํ ปุจฺฉิตา ยมารพฺภ.} \\ \csromangathacontbat{สุทฺธาสุทฺธปริกฺขา,}{หาโร โส \textbf{โสธโน} นาม.}}\csromangathastanza{\csromangathaitembat{14}{เอกตฺตตาย ธมฺมา,}{เยปิ จ เวมตฺตตาย นิทฺทิฏฺฐา.} \\ \csromangathacontbat{เต น วิกปฺปยิตพฺพา,}{เอโส หาโร \textbf{อธิฏฺฐาโน.}}}\csromangathastanza{\csromangathaitembat{15}{เย ธมฺมา ยํ ธมฺมํ,}{ชนยนฺติ\-ปฺปจฺจยา ปรมฺปรโต.} \\ \csromangathacontbat{เหตุมวกฑฺฒยิตฺวา,}{เอโส หาโร \textbf{ปริกฺขาโร.}}}\csromangathastanza{\csromangathaitembat{16}{เย ธมฺมา ยํมูลา,}{เย เจกตฺถา ปกาสิตา มุนินา.} \\ \csromangathacontbat{เต สมโรปยิตพฺพา,}{เอส \textbf{สมาโรปโน} หาโร.}}}

\csromanheader{2}{\textbf{นยสงฺเขป}}
\csromanmatikaanchor{mtk.4}
\csromantocmarkref{section}{4. ปฏินิทฺเทสวาร}{mtk.4}
\bookmark[dest={mtk.4},level=1]{4. ปฏินิทฺเทสวาร}
\setlength{\csromangathapagewidth}{0pt}
\setlength{\csromangathawakwidth}{0pt}
\csromangathameasuregroup{ตณฺหญฺจ อวิชฺชมฺปิ จ, \\ สจฺเจหิ โยชยิตฺวา, \\ โย เนติ วิปลฺลาเสหิ, \\ เอตํ นยํ นยวิทู, \\ โอโลเกตฺวา ทิสโลจเนน, \\ สพฺเพ กุสลากุสเล, \\ โสฬส หารา ปฐมํ, \\ สงฺขิปิย องฺกุเสน หิ,}{\csromangathabat{ตณฺหญฺจ อวิชฺชมฺปิ จ,}{สมเถน วิปสฺสนาย โย เนติ.} \\ \csromangathabat{สจฺเจหิ โยชยิตฺวา,}{อยํ นโย \textbf{นนฺทิยาวฏฺโฏ.}} \\ โย อกุสเล สมูเลหิ, \\ เนติ กุสเล จ กุสลมูเลหิ. \\ ภูตํ ตถํ อวิตถํ, \textbf{ติปุกฺขลํ} ตํ นยํ อาหุ. \\ \csromangathabat{โย เนติ วิปลฺลาเสหิ,}{กิเลเส อินฺทฺริเยหิ สทฺธมฺเม.} \\ \csromangathabat{เอตํ นยํ นยวิทู,}{\textbf{สีหวิกฺกีฬิตํ} อาหุ.} \\ เวยฺยากรเณสุ หิ เย, \\ กุสลากุสลา ตหิํ ตหิํ วุตฺตา. \\ มนสา โวโลกยเต, ตํ ขุ \textbf{ทิสาโลจนํ} อาหุ. \\ \csromangathabat{โอโลเกตฺวา ทิสโลจเนน,}{อุกฺขิปิย ยํ สมาเนติ.} \\ \csromangathabat{สพฺเพ กุสลากุสเล,}{อยํ นโย \textbf{องฺกุโส} นาม.} \\ \csromangathabat{โสฬส หารา ปฐมํ,}{ทิสโลจนโต\textsuperscript{1} ทิสา วิโลเกตฺวา.} \\ \csromangathabat{สงฺขิปิย องฺกุเสน หิ,}{นเยหิ ตีหิ นิทฺทิเส สุตฺตํ.}}
\csromangathameasurewak{โย อกุสเล สมูเลหิ, \\ เนติ กุสเล จ กุสลมูเลหิ. \\ ภูตํ ตถํ อวิตถํ, \textbf{ติปุกฺขลํ} ตํ นยํ อาหุ. \\ เวยฺยากรเณสุ หิ เย, \\ กุสลากุสลา ตหิํ ตหิํ วุตฺตา. \\ มนสา โวโลกยเต, ตํ ขุ \textbf{ทิสาโลจนํ} อาหุ.}
\csromangathasetleft{ตณฺหญฺจ อวิชฺชมฺปิ จ, \\ สจฺเจหิ โยชยิตฺวา, \\ โย เนติ วิปลฺลาเสหิ, \\ เอตํ นยํ นยวิทู, \\ โอโลเกตฺวา ทิสโลจเนน, \\ สพฺเพ กุสลากุสเล, \\ โสฬส หารา ปฐมํ, \\ สงฺขิปิย องฺกุเสน หิ,}
\csromangathagroup{\csromangathastanza{\csromangathaitembat{17}{ตณฺหญฺจ อวิชฺชมฺปิ จ,}{สมเถน วิปสฺสนาย โย เนติ.} \\ \csromangathacontbat{สจฺเจหิ โยชยิตฺวา,}{อยํ นโย \textbf{นนฺทิยาวฏฺโฏ.}}}\csromangathastanza{\csromangathawak{\csromangathaitemline{18}{โย อกุสเล สมูเลหิ,} \\ \csromangathacontline{เนติ กุสเล จ กุสลมูเลหิ.} \\ \csromangathacontline{ภูตํ ตถํ อวิตถํ, \textbf{ติปุกฺขลํ} ตํ นยํ อาหุ.}}}\csromangathastanza{\csromangathaitembat{19}{โย เนติ วิปลฺลาเสหิ,}{กิเลเส อินฺทฺริเยหิ สทฺธมฺเม.} \\ \csromangathacontbat{เอตํ นยํ นยวิทู,}{\textbf{สีหวิกฺกีฬิตํ} อาหุ.}}\csromangathastanza{\csromangathawak{\csromanfolio{5}\csromangathaitemline{20}{เวยฺยากรเณสุ หิ เย,} \\ \csromangathacontline{กุสลากุสลา ตหิํ ตหิํ วุตฺตา.} \\ \csromangathacontline{มนสา โวโลกยเต, ตํ ขุ \textbf{ทิสาโลจนํ} อาหุ.}}}\csromangathastanza{\csromangathaitembat{21}{โอโลเกตฺวา ทิสโลจเนน,}{อุกฺขิปิย ยํ สมาเนติ.} \\ \csromangathacontbat{สพฺเพ กุสลากุสเล,}{อยํ นโย \textbf{องฺกุโส} นาม.}}\csromangathastanza{\csromangathaitembat{22}{โสฬส หารา ปฐมํ,}{ทิสโลจนโต\csromansharedfootnote{49434a2deb5c0c74}{ทิสโลจเนน (ก)} ทิสา วิโลเกตฺวา.} \\ \csromangathacontbat{สงฺขิปิย องฺกุเสน หิ,}{นเยหิ ตีหิ นิทฺทิเส สุตฺตํ.}}}

\csromanheader{2}{\textbf{ทฺวาทสปท}}
\setlength{\csromangathapagewidth}{0pt}
\setlength{\csromangathawakwidth}{0pt}
\csromangathameasuregroup{อกฺขรํ ปทํ พฺยญฺชนํ, \\ อาการฉฏฺฐวจนํ, \\ สงฺกาสนา ปกาสนา, \\ เอเตหิ ฉหิ ปเทหิ, \\ ตีณิ จ นยา อนูนา, \\ นวหิ ปเทหิ ภควโต, \\ อตฺถสฺส นวปฺปทานิ, \\ อุภยํ สงฺกลยิตฺวา\textsuperscript{1},}{\csromangathabat{อกฺขรํ ปทํ พฺยญฺชนํ,}{นิรุตฺติ ตเถว นิทฺเทโส.} \\ \csromangathabat{อาการฉฏฺฐวจนํ,}{เอตฺตาว พฺยญฺชนํ สพฺพํ.} \\ \csromangathabat{สงฺกาสนา ปกาสนา,}{วิวรณา วิภชนุตฺตานีกมฺมปญฺญตฺติ.} \\ \csromangathabat{เอเตหิ ฉหิ ปเทหิ,}{อตฺโถ กมฺมญฺจ นิทฺทิฏฺฐํ.} \\ \csromangathabat{ตีณิ จ นยา อนูนา,}{อตฺถสฺส จ ฉปฺปทานิ คณิตานิ.} \\ \csromangathabat{นวหิ ปเทหิ ภควโต,}{วจนสฺสตฺโถ สมายุตฺโต.} \\ \csromangathabat{อตฺถสฺส นวปฺปทานิ,}{พฺยญฺชนปริเยฏฺฐิยา จตุพฺพีส.} \\ \csromangathabat{อุภยํ สงฺกลยิตฺวา\textsuperscript{1},}{เตตฺติํสา เอตฺติกา เนตฺตีติ.}}
\csromangathasetleft{อกฺขรํ ปทํ พฺยญฺชนํ, \\ อาการฉฏฺฐวจนํ, \\ สงฺกาสนา ปกาสนา, \\ เอเตหิ ฉหิ ปเทหิ, \\ ตีณิ จ นยา อนูนา, \\ นวหิ ปเทหิ ภควโต, \\ อตฺถสฺส นวปฺปทานิ, \\ อุภยํ สงฺกลยิตฺวา\textsuperscript{1},}
\csromangathagroup{\csromangathastanza{\csromangathaitembat{23}{อกฺขรํ ปทํ พฺยญฺชนํ,}{นิรุตฺติ ตเถว นิทฺเทโส.} \\ \csromangathacontbat{อาการฉฏฺฐวจนํ,}{เอตฺตาว พฺยญฺชนํ สพฺพํ.}}\csromangathastanza{\csromangathaitembat{24}{สงฺกาสนา ปกาสนา,}{วิวรณา วิภชนุ\-ตฺตานีกมฺม\-ปญฺญตฺติ.} \\ \csromangathacontbat{เอเตหิ ฉหิ ปเทหิ,}{อตฺโถ กมฺมญฺจ นิทฺทิฏฺฐํ.}}\csromangathastanza{\csromangathaitembat{25}{ตีณิ จ นยา อนูนา,}{อตฺถสฺส จ ฉปฺปทานิ คณิตานิ.} \\ \csromangathacontbat{นวหิ ปเทหิ ภควโต,}{วจนสฺสตฺโถ สมายุตฺโต.}}\csromangathastanza{\csromangathaitembat{26}{อตฺถสฺส นวปฺปทานิ,}{พฺยญฺชน\-ปริเยฏฺฐิยา จตุพฺพีส.} \\ \csromangathacontbat{อุภยํ สงฺกลยิตฺวา\csromansharedfootnote{58e47349434fec55}{สงฺเขปยโต (ก)},}{เตตฺติํสา เอตฺติกา เนตฺตีติ.}}}

\vspace{\tipitakaparskip}
\csromancenter{นิทฺเทสวาโร.}
\csromanmatikaanchor{mtk.5}
\csromantocmarkref{section}{1. เทสนาหารวิภงฺค}{mtk.5}
\bookmark[dest={mtk.5},level=1]{1. เทสนาหารวิภงฺค}
\csromanheader{1}{1. \textbf{เทสนาหาร}\-\textbf{วิภงฺค}}
\proseitem{5}{ตตฺถ กตโม เทสนาหาโร, “อสฺสาทาทีนวตา”ติ คาถา อยํ เทสนาหาโร กิํ เทสยติ, อสฺสาทํ อาทีนวํ นิสฺสรณํ ผลํ อุปายํ อาณตฺติํ.\csromansymbolfootnote{*}{ม 3. 327 ปิฏฺเฐ.}ธมฺมํ โว ภิกฺขเว เทเสสฺสามิ อาทิกลฺยาณํ มชฺเฌกลฺยาณํ ปริโยสาน\-กลฺยาณํ สาตฺถํ สพฺยญฺชนํ เกวล\-ปริปุณฺณํ ปริสุทฺธํ พฺรหฺมจริยํ ปกาเสสฺสามีติ.}

\vspace{\tipitakaparskip}
\csromancenter{ตตฺถ กตโม อสฺสาโท\csromandash{}}
\setlength{\csromangathapagewidth}{0pt}
\setlength{\csromangathawakwidth}{0pt}
\csromangathameasuregroup{\textsuperscript{*}“กามํ\textsuperscript{1} กามยมานสฺส,}{\csromangathabat{\textsuperscript{*}“กามํ\textsuperscript{1} กามยมานสฺส,}{ตสฺส เจตํ สมิชฺฌติ.}}
\csromangathasetleft{\textsuperscript{*}“กามํ\textsuperscript{1} กามยมานสฺส,}
\csromangathagroup{\csromangathastanza{\csromanfolio{6}\csromangathabat{\csromansymbolfootnote{*}{อุปริ 59, 198 ปิฏฺเฐสุปิ.}“กามํ\csromansharedfootnote{738fc8a7b7bd28bd}{กามมาทิกา อิมา ฉ คาถา ขุ 1. 399 ปิฏฺเฐ ปสฺสิตพฺพา.} กามยมานสฺส,}{ตสฺส เจตํ สมิชฺฌติ.}}}

\vspace{\tipitakaparskip}
\csromancenter{อทฺธา ปีติมโน โหติ, ลทฺธา มจฺโจ ยทิจฺฉตี”ติ. อยํ อสฺสาโท. ตตฺถ กตโม อาทีนโว\csromandash{}}
\setlength{\csromangathapagewidth}{0pt}
\setlength{\csromangathawakwidth}{0pt}
\csromangathameasuregroup{\textsuperscript{+}“ตสฺส เจ กามยานสฺส,}{\csromangathabat{\textsuperscript{+}“ตสฺส เจ กามยานสฺส,}{ฉนฺทชาตสฺส ชนฺตุโน.}}
\csromangathasetleft{\textsuperscript{+}“ตสฺส เจ กามยานสฺส,}
\csromangathagroup{\csromangathastanza{\csromangathabat{\csromansymbolfootnote{+}{อุปริ 59 ปิฏฺเฐปิ.}“ตสฺส เจ กามยานสฺส,}{ฉนฺทชาตสฺส ชนฺตุโน.}}}

\prose{เต กามา ปริหายนฺติ, สลฺลวิทฺโธว รุปฺปตี”ติ. อยํ อาทีนโว. ตตฺถ กตมํ นิสฺสรณํ\csromandash{}}

\setlength{\csromangathapagewidth}{0pt}
\setlength{\csromangathawakwidth}{0pt}
\csromangathameasuregroup{\textsuperscript{*}“โย กาเม ปริวชฺเชติ,}{\csromangathabat{\textsuperscript{*}“โย กาเม ปริวชฺเชติ,}{สปฺปสฺเสว ปทา สิโร.}}
\csromangathasetleft{\textsuperscript{*}“โย กาเม ปริวชฺเชติ,}
\csromangathagroup{\csromangathastanza{\csromangathabat{\csromansymbolmark{*}“โย กาเม ปริวชฺเชติ,}{สปฺปสฺเสว ปทา สิโร.}}}

\prose{โสมํ วิสตฺติกํ โลเก, สโต สมติวตฺตตี”ติ. อิทํ นิสฺสรณํ. ตตฺถ กตโม อสฺสาโท\csromandash{}}

\setlength{\csromangathapagewidth}{0pt}
\setlength{\csromangathawakwidth}{0pt}
\csromangathameasuregroup{“เขตฺตํ วตฺถุํ หิรญฺญํ วา,}{\csromangathabat{“เขตฺตํ วตฺถุํ หิรญฺญํ วา,}{ควาสฺสํ ทาสโปริสํ.}}
\csromangathasetleft{“เขตฺตํ วตฺถุํ หิรญฺญํ วา,}
\csromangathagroup{\csromangathastanza{\csromangathabat{“เขตฺตํ วตฺถุํ หิรญฺญํ วา,}{ควาสฺสํ ทาสโปริสํ.}}}

\vspace{\tipitakaparskip}
\csromancenter{ถิโย พนฺธู ปุถู กาเม, โย นโร อนุคิชฺฌตี”ติ. อยํ อสฺสาโท. ตตฺถ กตโม อาทีนโว\csromandash{}}
\setlength{\csromangathapagewidth}{0pt}
\setlength{\csromangathawakwidth}{0pt}
\csromangathameasuregroup{“อพลา นํ พลียนฺติ,}{\csromangathabat{“อพลา นํ พลียนฺติ,}{มทฺทนฺเตนํ ปริสฺสยา.}}
\csromangathasetleft{“อพลา นํ พลียนฺติ,}
\csromangathagroup{\csromangathastanza{\csromangathabat{“อพลา นํ พลียนฺติ,}{มทฺทนฺเตนํ ปริสฺสยา.}}}

\vspace{\tipitakaparskip}
\csromancenter{ตโต นํ ทุกฺขมเนฺวติ, นาวํ ภินฺนมิโวทกนฺ”ติ. อยํ อาทีนโว. ตตฺถ กตมํ นิสฺสรณํ\csromandash{}}
\setlength{\csromangathapagewidth}{0pt}
\setlength{\csromangathawakwidth}{0pt}
\csromangathameasuregroup{“ตสฺมา ชนฺตุ สทา สโต, \\ เต ปหาย ตเร โอฆํ,}{\csromangathabat{“ตสฺมา ชนฺตุ สทา สโต,}{กามานิ ปริวชฺชเย.} \\ \csromangathabat{เต ปหาย ตเร โอฆํ,}{นาวํ สิตฺวาว ปารคู”ติ.} \\ อิทํ นิสฺสรณํ.}
\csromangathameasurewak{อิทํ นิสฺสรณํ.}
\csromangathasetleft{“ตสฺมา ชนฺตุ สทา สโต, \\ เต ปหาย ตเร โอฆํ,}
\csromangathagroup{\csromangathastanza{\csromangathabat{“ตสฺมา ชนฺตุ สทา สโต,}{กามานิ ปริวชฺชเย.} \\ \csromangathabat{เต ปหาย ตเร โอฆํ,}{นาวํ สิตฺวาว ปารคู”ติ.}}\csromangathastanza{\csromangathawak{อิทํ นิสฺสรณํ.}}}

\vspace{\tipitakaparskip}
\csromancenter{ตตฺถ กตมํ ผลํ\csromandash{}}
\setlength{\csromangathapagewidth}{0pt}
\setlength{\csromangathawakwidth}{0pt}
\csromangathameasuregroup{}{\textsuperscript{*}“ธมฺโม หเว รกฺขติ ธมฺมจาริํ, \\ ฉตฺตํ มหนฺตํ ยถ วสฺสกาเล. \\ เอสานิสํโส ธมฺเม สุจิณฺเณ,}
\csromangathameasurewak{\textsuperscript{*}“ธมฺโม หเว รกฺขติ ธมฺมจาริํ, \\ ฉตฺตํ มหนฺตํ ยถ วสฺสกาเล. \\ เอสานิสํโส ธมฺเม สุจิณฺเณ,}
\setlength{\csromangathaleftcol}{0pt}
\csromangathagroup{\csromangathastanza{\csromangathawak{\csromanfolio{7}\csromansymbolfootnote{*}{ขุ 5. 215 ปิฏฺเฐ โถกํ วิสทิสํ.}“ธมฺโม หเว รกฺขติ ธมฺมจาริํ, \\ ฉตฺตํ มหนฺตํ ยถ วสฺสกาเล. \\ เอสานิสํโส ธมฺเม สุจิณฺเณ,}}}

\prose{น ทุคฺคติํ คจฺฉติ ธมฺมจารี”ติ. อิทํ ผลํ. ตตฺถ กตโม อุปาโย\csromandash{}}

\prose{“สพฺเพ สงฺขารา อนิจฺจา”ติ ฯเปฯ.}

\prose{“สพฺเพ สงฺขารา ทุกฺขา”ติ ฯเปฯ.}

\prose{“สพฺเพ ธมฺมา อนตฺตา”ติ,}

\setlength{\csromangathapagewidth}{0pt}
\setlength{\csromangathawakwidth}{0pt}
\csromangathameasuregroup{}{ยทา ปญฺญาย ปสฺสติ.}
\csromangathameasurewak{ยทา ปญฺญาย ปสฺสติ.}
\setlength{\csromangathaleftcol}{0pt}
\csromangathagroup{\csromangathastanza{\csromangathawak{ยทา ปญฺญาย ปสฺสติ.}}}

\prose{อถ นิพฺพินฺทติ ทุกฺเข, เอส มคฺโค วิสุทฺธิยา”ติ\csromansharedfootnote{8e60e332cad3ccc3}{ขุ 1. 53 ปิฏฺเฐ ธมฺมปเท.}. อยํ อุปาโย. ตตฺถ กตมา อาóอตฺติ\csromandash{}}

\setlength{\csromangathapagewidth}{0pt}
\setlength{\csromangathawakwidth}{0pt}
\csromangathameasuregroup{“จกฺขุมา วิสมานีว,}{\csromangathabat{“จกฺขุมา วิสมานีว,}{วิชฺชมาเน ปรกฺกเม.} \\ อยํ อาณตฺติ.}
\csromangathameasurewak{อยํ อาณตฺติ.}
\csromangathasetleft{“จกฺขุมา วิสมานีว,}
\csromangathagroup{\csromangathastanza{\csromangathabat{“จกฺขุมา วิสมานีว,}{วิชฺชมาเน ปรกฺกเม.}}\csromangathastanza{\csromangathawak{อยํ อาณตฺติ.}}}

\prose{ปณฺฑิโต ชีวโลกสฺมิํ, ปาปานิ ปริวชฺชเย”ติ\csromansharedfootnote{9cd0225e86bdc70c}{ขุ 1. 137 ปิฏฺเฐ อุทาเน.}.}

\setlength{\csromangathapagewidth}{0pt}
\setlength{\csromangathawakwidth}{0pt}
\csromangathameasuregroup{}{“อตฺตานุทิฏฺฐิํ อูหจฺจ, \\ เอวํ มจฺจุตโร สิยา”\textsuperscript{1}. \\ อิทํ ผลํ.}
\csromangathameasurewak{“อตฺตานุทิฏฺฐิํ อูหจฺจ, \\ เอวํ มจฺจุตโร สิยา”\textsuperscript{1}. \\ อิทํ ผลํ.}
\setlength{\csromangathaleftcol}{0pt}
\csromangathagroup{\csromangathastanza{\csromangathawak{“อตฺตานุทิฏฺฐิํ อูหจฺจ, \\ เอวํ มจฺจุตโร สิยา”\csromansharedfootnote{338dd9d64fd7ee1e}{[มิสฺสินฺคฺ โนเต 0]}. \\ อิทํ ผลํ.}}}

\prose{“สุญฺญโต โลกํ อเวกฺขสฺสุ, โมฆราชา”ติ อาณตฺติ,}

\prose{“สทา สโต”ติ อุปาโย.}

\proseitem{6}{ตตฺต ภควา อุคฺฆฏิต\-ญฺญุสฺส ปุคฺคลสฺส นิสฺสรณํ เทสยติ, วิปญฺจิต\-ญฺญุสฺส\csromansharedfootnote{f4bf22b1f81616fb}{วิปจฺจิตญฺญุสฺส (สี)} ปุคฺคลสฺส อาทีนวญฺจ นิสฺสรณญฺจ เทสยติ, เนยฺยสฺส ปุคฺคลสฺส อสฺสาทญฺจ อาทีนวญฺจ นิสฺสรณญฺจ เทสยติ.}

\prose{ตตฺถ จตสฺโส ปฏิปทา, จตฺตาโร ปุคฺคลา. ตณฺหาจริโต มนฺโท สตินฺทฺริเยน ทุกฺขาย ปฏิปทาย ทนฺธาภิญฺญาย นิยฺยาติ สติปฏฺฐาเนหิ นิสฺสเยหิ.  ตณฺหาจริโต อุทตฺโถ\csromansharedfootnote{e657080d5d63a8f9}{อุทตฺโต (ก)} สมาธิ\-นฺทฺริเยน ทุกฺขาย ปฏิปทาย \csromanfolio{8}ขิปฺปาภิญฺญาย นิยฺยาติ ฌาเนหิ นิสฺสเยหิ.  ทิฏฺฐิจริโต มนฺโท วีริยินฺทฺริเยน สุขาย ปฏิปทาย ทนฺธาภิญฺญาย นิยฺยาติ สมฺม\-ปฺปธาเนหิ นิสฺสเยหิ.  ทิฏฺฐิจริโต อุทตฺโต ปญฺญินฺทฺริเยน สุขาย ปฏิปทาย ขิปฺปาภิญฺญาย นิยฺยาติ สจฺเจหิ นิสฺสเยหิ.}

\prose{อุโภ ตณฺหาจริตา สมถ\-ปุพฺพงฺคมาย วิปสฺสนาย นิยฺยนฺติ ราควิราคาย เจโตวิมุตฺติยา. อุโภ ทิฏฺฐิจริตา วิปสฺสนา\-ปุพฺพงฺคเมน สมเถน นิยฺยนฺติ อวิชฺชาวิราคาย ปญฺญาวิมุตฺติยา.}

\prose{ตตฺถ เย สมถ\-ปุพฺพงฺคมาหิ ปฏิปทาหิ นิยฺยนฺติ, เต นนฺทิยาวฏฺเฏน นเยน หาตพฺพา, เย วิปสฺสนา\-ปุพฺพงฺคมาหิ ปฏิปทาหิ นิยฺยนฺติ, เต สีหวิกฺกีฬิเตน นเยน หาตพฺพา.}

\prose{สฺวายํ หาโร กตฺถ สมฺภวติ, ยสฺส สตฺถา วา ธมฺมํ เทสยติ อญฺญตโร วา ครุฏฺฐานีโย สพฺรหฺมจารี, โส ตํ ธมฺมํ สุตฺวา สทฺธํ ปฏิลภติ. ตตฺถ ยา วีมํสา อุสฺสาหนา ตุลนา อุปปริกฺขา, อยํ สุตมยี ปญฺญา. ตถา สุเตน นิสฺสเยน ยา วีมํสา ตุลนา อุปปริกฺขา มนสานุเปกฺขณา, อยํ จินฺตามยี ปญฺญา. อิมาหิ ทฺวีหิ ปญฺญาหิ มนสิการ\-สมฺปยุตฺตสฺส ยํ ญาณํ อุปฺปชฺชติ ทสฺสน\-ภูมิยํ วา ภาวนา\-ภูมิยํ วา, อยํ ภาวนามยี ปญฺญา.}

\prose{ปรโตโฆสา สุตมยี ปญฺญา. ปจฺจตฺต\-สมุฏฺฐิตา โยนิโส มนสิการา จินฺตามยี ปญฺญา. ยํ ปรโต จ โฆเสน ปจฺจตฺต\-สมุฏฺฐิเตน จ โยนิโส\-มนสิกาเรน ญาณํ อุปฺปชฺชติ, อยํ ภาวนามยี ปญฺญา. ยสฺส อิมา เทฺว ปญฺญา อตฺถิ สุตมยี จินฺตามยี จ, อยํ อุคฺฆฏิตญฺญู. ยสฺส สุตมยี ปญฺญา อตฺถิ, จินฺตามยี นตฺถิ, อยํ วิปญฺจิตญฺญู. ยสฺส เนว สุตมยี ปญฺญา อตฺถิ น จินฺตามยี, อยํ เนยฺโย.}

\proseitem{9}{สายํ ธมฺมเทสนา กิํ เทสยติ, จตฺตาริ สจฺจานิ ทุกฺขํ สมุทยํ นิโรธํ มคฺคํ. อาทีนโว จ ผลญฺจ ทุกฺขํ, อสฺสาโท สมุทโย, นิสฺสรณํ นิโรโธ, อุปาโย อาณตฺติ จ มคฺโค. อิมานิ จตฺตาริ สจฺจานิ. อิทํ ธมฺมจกฺกํ.}

\prose{\csromanfolio{9}ยถาห ภควา\csromandash{}“อิทํ ทุกฺขนฺ”ติ เม ภิกฺขเว พาราณสิยํ อิสิปตเน มิคทาเย อนุตฺตรํ ธมฺมจกฺกํ ปวตฺติตํ อปฺปฏิวตฺติยํ สมเณน วา พฺราหฺมเณน วา เทเวน วา มาเรน วา พฺรหฺมุนา วา เกนจิ วา โลกสฺมิํ, สพฺพํ ธมฺมจกฺกํ.*}

\prose{ตตฺถ อปริมาณา ปทา, อปริมาณา อกฺขรา, อปริมาณา พฺยญฺชนา, อปริมาณา อาการา เนรุตฺตา นิทฺเทสา. เอตสฺเสว อตฺถสฺส สงฺกาสนา ปกาสนา วิวรณา วิภชนา อุตฺตานีกมฺมํ\csromansharedfootnote{93855b27a579489b}{อุตฺตานิกมฺมํ (ก)} ปญฺญตฺติ, อิติปิทํ ทุกฺขํ อริยสจฺจํ.}

\prose{“อยํ ทุกฺขสมุทโย”ติ เม ภิกฺขเว พาราณสิยํ อิสิปตเน มิคทาเย อนุตฺตรํ ธมฺมจกฺกํ ปวตฺติตํ ฯเปฯ “อยํ ทุกฺขนิโรโธ”ติ เม ภิกฺขเว ฯเปฯ “อยํ ทุกฺขนิโรธ\-คามินี ปฏิปทา”ติ เม ภิกฺขเว พาราณสิยํ อิสิปตเน มิคทาเย อนุตฺตรํ ธมฺมจกฺกํ ปวตฺติตํ อปฺปฏิวตฺติยํ สมเณน วา พฺราหฺมเณน วา เทเวน วา มาเรน วา พฺรหฺมุนา วา เกนจิ วา โลกสฺมิํ.}

\prose{ตตฺถ อปริมาณา ปทา, อปริมาณา อกฺขรา, อปริมาณา พฺยญฺชนา, อปริมาณา อาการา เนรุตฺตา นิทฺเทสา. เอตสฺเสว อตฺถสฺส สงฺกาสนา ปกาสนา วิวรณา วิภชนา อุตฺตานีกมฺมํ ปญฺญตฺติ อิติปิทํ ทุกฺขนิโรธ\-คามินี ปฏิปทา อริยสจฺจํ.}

\prose{ตตฺถ ภควา อกฺขเรหิ สงฺกาเสติ, ปเทหิ ปกาเสติ, พฺยญฺชเนหิ วิวรติ, อากาเรหิ วิภชติ, นิรุตฺตีหิ อุตฺตานีกโรติ\csromansharedfootnote{07c9a85d9dba38db}{อุตฺตานิํ กโรติ (ก)}, นิทฺเทเสหิ ปญฺญเปติ. ตตฺถ ภควา อกฺขเรหิ จ ปเทหิ จ อุคฺฆเฏติ\csromansharedfootnote{a44a1bd383235e3e}{อุคฺฆาเฏติ (สี)}, พฺยญฺชเนหิ จ อากาเรหิ จ วิปญฺจยติ, นิรุตฺตีหิ จ นิทฺเทเสหิ จ วิตฺถาเรติ. ตตฺถ อุคฺฆฏนา\csromansharedfootnote{65ea6c15c5b79b54}{อุคฺฆาฏนา (สี)} อาทิ, วิปญฺจนา มชฺเฌ, วิตฺถารณา ปริโยสานํ. โสยํ ธมฺมวินโย อุคฺฆฏียนฺโต อุคฺฆฏิตญฺญู\-ปุคฺคลํ วิเนติ, เตน นํ อาหุ “อาทิกลฺยาโณ”ติ. วิปญฺจียนฺโต วิปญฺจิตญฺญู\-ปุคฺคลํ วิเนติ, เตน นํ อาหุ “มชฺเฌกลฺยาโณ”ติ. วิตฺถารียนฺโต เนยฺยํ ปุคฺคลํ วิเนติ, เตน นํ อาหุ “ปริโยสาน\-กลฺยาโณ”ติ.}

\proseitem{10}{\csromanfolio{10}ตตฺถ ฉปฺปทานิ อตฺโถ สงฺกาสนา ปกาสนา วิวรณา วิภชนา อุตฺตานีกมฺมํ ปญฺญตฺติ, อิมานิ ฉปฺปทานิ อตฺโถ. ฉปฺปทานิ พฺยญฺชนํ อกฺขรํ ปทํ พฺยญฺชนํ อากาโร นิรุตฺติ นิทฺเทโส, อิมานิ ฉปฺปทานิ พฺยญฺชนํ. เตนาห ภควา “ธมฺมํ โว ภิกฺขเว เทเสสฺสามิ อาทิกลฺยาณํ มชฺเฌกลฺยาณํ ปริโยสาน\-กลฺยาณํ สาตฺถํ สพฺยญฺชนนฺ”ติ.}

\prose{\textbf{เกวลนฺ}ติ โลกุตฺตรํ น มิสฺสํ โลกิเยหิ ธมฺเมหิ. \textbf{ปริปุณฺณนฺ}ติ ปริปูรํ อนูนํ อนติเรกํ. \textbf{ปริสุทฺธนฺ}ติ นิมฺมลํ สพฺพ\-มลา\-ปคตํ ปริโยทาตํ อุปฏฺฐิตํ สพฺพวิเสสานํ, อิทํ วุจฺจติ ตถาคตปทํอิติปิ ตถาคตนิเสวิตํอิติปิ ตถาคตารญฺชิตํอิติปิ, อโต เจตํ พฺรหฺมจริยํ ปญฺญายติ. เตนาห ภควา “เกวล\-ปริปุณฺณํ ปริสุทฺธํ พฺรหฺมจริยํ ปกาเสสฺสามี”ติ.}

\prose{เกสํ อยํ ธมฺมเทสนา, โยคีนํ. เตนาห อายสฺมา มหากจฺจายโน\csromandash{}}

\setlength{\csromangathapagewidth}{0pt}
\setlength{\csromangathawakwidth}{0pt}
\csromangathameasuregroup{“อสฺสาทาทีนวตา, \\ อาณตฺตี จ ภควโต,}{\csromangathabat{“อสฺสาทาทีนวตา,}{นิสฺสรณมฺปิ จ ผลํ อุปาโย จ.} \\ \csromangathabat{อาณตฺตี จ ภควโต,}{โยคีนํ เทสนาหาโร”ติ.} \\ นิยุตฺโต เทสนาหาโร.}
\csromangathameasurewak{นิยุตฺโต เทสนาหาโร.}
\csromangathasetleft{“อสฺสาทาทีนวตา, \\ อาณตฺตี จ ภควโต,}
\csromangathagroup{\csromangathastanza{\csromangathabat{“อสฺสาทาทีนวตา,}{นิสฺสรณมฺปิ จ ผลํ อุปาโย จ.} \\ \csromangathabat{อาณตฺตี จ ภควโต,}{โยคีนํ เทสนาหาโร”ติ.}}\csromangathastanza{\csromangathawak{นิยุตฺโต เทสนาหาโร.}}}

\csromanmatikaanchor{mtk.6}
\csromantocmarkref{section}{2. วิจยหารวิภงฺค}{mtk.6}
\bookmark[dest={mtk.6},level=1]{2. วิจยหารวิภงฺค}
\csromanheader{1}{2. \textbf{วิจย}\-\textbf{หาร}\-\textbf{วิภงฺค}}
\proseitem{11}{ตตฺถ กตโม วิจโย หาโร, “ยํ ปุจฺฉิตญฺจ วิสฺสชฺชิตญฺจา”ติ คาถา, อยํ วิจโย หาโร.}

\prose{กิํ วิจินติ. ปทํ วิจินติ, ปญฺหํ วิจินติ, วิสชฺชนํ\csromansharedfootnote{50fa4a160f4a7550}{วิสฺสชฺชนํ (สี, ก)} วิจินติ, ปุพฺพาปรํ วิจินติ, อสฺสาทํ วิจินติ, อาทีนวํ วิจินติ, นิสฺสรณํ วิจินติ, ผลํ วิจินติ, อุปายํ วิจินติ, อาณตฺติํ วิจินติ, อนุคีติํ วิจินติ, สพฺเพ นว สุตฺตนฺเต วิจินติ. ยถา กิํ ภเว, ยถา อายสฺมา อชิโต ปารายเน ภควนฺตํ ปญฺหํ ปุจฺฉติ\csromandash{}}

\prose{“เกนสฺสุ นิวุโต โลโก, (อิจฺจายสฺมา อชิโต,)}

\prose{เกนสฺสุ นปฺปกาสติ.}

\setlength{\csromangathapagewidth}{0pt}
\setlength{\csromangathawakwidth}{0pt}
\csromangathameasuregroup{กิสฺสาภิเลปนํ พฺรูสิ,}{\csromangathabat{กิสฺสาภิเลปนํ พฺรูสิ,}{กิํ สุ ตสฺส มหพฺภยนฺ”ติ\textsuperscript{1}.}}
\csromangathasetleft{กิสฺสาภิเลปนํ พฺรูสิ,}
\csromangathagroup{\csromangathastanza{\csromangathabat{กิสฺสาภิเลปนํ พฺรูสิ,}{กิํ สุ ตสฺส มหพฺภยนฺ”ติ\csromansharedfootnote{f206042c5db8c83d}{ขุ 1. 434 ปิฏฺเฐ สุตฺตนิปาเต; อุปริ 60 ปิฏฺเฐปิ.}.}}}

\prose{\csromanfolio{11}อิมานิ จตฺตาริ ปทานิ ปุจฺฉิตานิ, โส เอโก ปโญฺห. กสฺมา, เอกวตฺถุ\-ปริคฺคหา, เอวญฺหิ อาห “เกนสฺสุ นิวุโต โลโก”ติ โลกาธิฏฺฐานํ ปุจฺฉติ, “เกนสฺสุ นปฺปกาสตี”ติ โลกสฺส อปฺปกาสนํ ปุจฺฉติ, “กิสฺสาภิเลปนํ พฺรูสี”ติ โลกสฺส อภิเลปนํ ปุจฺฉติ, “กิํสุ ตสฺส มหพฺภยนฺ”ติ ตสฺเสว โลกสฺส มหาภยํ ปุจฺฉติ. โลโก ติวิโธ กิเลสโลโก ภวโลโก อินฺทฺริยโลโก.}

\vspace{\tipitakaparskip}
\csromancenter{ตตฺถ วิสชฺชนา\csromandash{}}
\prose{\csromansymbolfootnote{*}{ขุ 1. 434 ปิฏฺเฐ สุตฺตนิปาเต; อุปริ 60, 177, 183 ปิฏฺเฐสุปิ.}“อวิชฺชาย นิวุโต โลโก, (อชิตาติ ภควา,)}

\prose{วิวิจฺฉา\csromansharedfootnote{c2a3da2a1980c56d}{เววิจฺฉา (ขุ 1. 434 ปิฏฺเฐ.)} ปมาทา นปฺปกาสติ.}

\setlength{\csromangathapagewidth}{0pt}
\setlength{\csromangathawakwidth}{0pt}
\csromangathameasuregroup{ชปฺปาภิเลปนํ พฺรูมิ,}{\csromangathabat{ชปฺปาภิเลปนํ พฺรูมิ,}{ทุกฺขมสฺส มหพฺภยนฺ”ติ.}}
\csromangathasetleft{ชปฺปาภิเลปนํ พฺรูมิ,}
\csromangathagroup{\csromangathastanza{\csromangathabat{ชปฺปา\-ภิเลปนํ พฺรูมิ,}{ทุกฺขมสฺส มหพฺภยนฺ”ติ.}}}

\prose{อิมานิ จตฺตาริ ปทานิ อิเมหิ จตูหิ ปเทหิ วิสชฺชิตานิ ปฐมํ ปฐเมน, ทุติยํ ทุติเยน, ตติยํ ตติเยน, จตุตฺถํ จตุตฺเถน.}

\prose{“เกนสฺสุ นิวุโต โลโก”ติ ปเญฺห “อวิชฺชาย นิวุโต โลโก”ติ วิสชฺชนา. นีวรเณหิ นิวุโต โลโก, อวิชฺชานีวรณา หิ สพฺเพ สตฺตา. ยถาห ภควา “สพฺพสตฺตานํ ภิกฺขเว สพฺพปาณานํ สพฺพภูตานํ ปริยายโต เอกเมว นีวรณํ วทามิ ยทิทํ อวิชฺชา, อวิชฺชานีวรณา หิ สพฺเพ สตฺตา. สพฺพโสว ภิกฺขเว อวิชฺชาย นิโรธา จาคา ปฏินิสฺสคฺคา นตฺถิ สตฺตานํ นีวรณนฺติ วทามี”ติ. เตน จ ปฐมสฺส ปทสฺส วิสชฺชนา ยุตฺตา.}

\prose{“เกนสฺสุ นปฺปกาสตี”ติ ปเญฺห “วิวิจฺฉา ปมาทา นปฺปกาสตี”ติ วิสชฺชนา. โย ปุคฺคโล นีวรเณหิ นิวุโต, โส วิวิจฺฉติ, วิวิจฺฉา นาม วุจฺจติ วิจิกิจฺฉา, โส วิจิกิจฺฉนฺโต นาภิสทฺทหติ, น อภิสทฺทหนฺโต วีริยํ นารภติ อกุสลานํ ธมฺมานํ ปหานาย กุสลานํ ธมฺมานํ สจฺฉิกิริยาย. โส อิธปฺปมาทม\-นุยุตฺโต วิหรติ ปมตฺโต, สุกฺเก ธมฺเม น อุปฺปาทิยติ, ตสฺส เต อนุปฺปาทิยมานา นปฺปกาสนฺติ, ยถาห ภควา\csromandash{}}

\setlength{\csromangathapagewidth}{0pt}
\setlength{\csromangathawakwidth}{0pt}
\csromangathameasuregroup{“ทูเร สนฺโต ปกาสนฺติ\textsuperscript{1}, \\ อสนฺเตตฺถ น ทิสฺสนฺติ, \\ เต คุเณหิ ปกาสนฺติ,}{\csromangathabat{“ทูเร สนฺโต ปกาสนฺติ\textsuperscript{1},}{หิมวนฺโตว ปพฺพโต.} \\ \csromangathabat{อสนฺเตตฺถ น ทิสฺสนฺติ,}{รตฺติํ ขิตฺตา\textsuperscript{1} ยถา สรา.} \\ \csromangathabat{เต คุเณหิ ปกาสนฺติ,}{กิตฺติยา จ ยเสน จา”ติ.}}
\csromangathasetleft{“ทูเร สนฺโต ปกาสนฺติ\textsuperscript{1}, \\ อสนฺเตตฺถ น ทิสฺสนฺติ, \\ เต คุเณหิ ปกาสนฺติ,}
\csromangathagroup{\csromangathastanza{\csromanfolio{12}\csromangathabat{“ทูเร สนฺโต ปกาสนฺติ\csromansharedfootnote{0a0770791602188f}{ปกาเสนฺติ (ขุ 1. 56 ปิฏฺเฐ.)},}{หิมวนฺโตว ปพฺพโต.} \\ \csromangathabat{อสนฺเตตฺถ น ทิสฺสนฺติ,}{รตฺติํ ขิตฺตา\csromansharedfootnote{0a0770791602188f}{ปกาเสนฺติ (ขุ 1. 56 ปิฏฺเฐ.)} ยถา สรา.} \\ \csromangathabat{เต คุเณหิ ปกาสนฺติ,}{กิตฺติยา จ ยเสน จา”ติ.}}}

\prose{เตน จ ทุติยสฺส ปทสฺส วิสชฺชนา ยุตฺตา.}

\prose{“กิสฺสาภิเลปนํ พฺรูสี”ติ ปเญฺห “ชปฺปา\-ภิเลปนํ พฺรูมี”ติ วิสชฺชนา. ชปฺปา นาม วุจฺจติ ตณฺหา, สา กถํ อภิลิมฺปติ. ยถาห ภควา\csromandash{}}

\setlength{\csromangathapagewidth}{0pt}
\setlength{\csromangathawakwidth}{0pt}
\csromangathameasuregroup{\textsuperscript{*}“รตฺโต อตฺถํ น ชานาติ, \\ อนฺธนฺตมํ\textsuperscript{1} ตทา โหติ,}{\csromangathabat{\textsuperscript{*}“รตฺโต อตฺถํ น ชานาติ,}{รตฺโต ธมฺมํ น ปสฺสติ.} \\ \csromangathabat{อนฺธนฺตมํ\textsuperscript{1} ตทา โหติ,}{ยํ ราโค สหเต นรนฺ”ติ.}}
\csromangathasetleft{\textsuperscript{*}“รตฺโต อตฺถํ น ชานาติ, \\ อนฺธนฺตมํ\textsuperscript{1} ตทา โหติ,}
\csromangathagroup{\csromangathastanza{\csromangathabat{\csromansymbolfootnote{*}{อุปริ 32 ปิฏฺเฐปิ.}“รตฺโต อตฺถํ น ชานาติ,}{รตฺโต ธมฺมํ น ปสฺสติ.} \\ \csromangathabat{อนฺธนฺตมํ\csromansharedfootnote{cdcc39fa49486f0b}{อนฺธตมํ (ก) ** ขุ 1. 434 ปิฏฺเฐ สุตฺตนิปาเต; อุปริ 61 ปิฏฺเฐปิ.} ตทา โหติ,}{ยํ ราโค สหเต นรนฺ”ติ.}}}

\prose{สายํ ตณฺหา อาสตฺติ\-พหุลสฺส ปุคฺคลสฺส “เอวํ อภิชปฺปา”ติ กริตฺวา ตตฺถ โลโก อภิลิตฺโต นาม ภวติ, เตน จ ตติยสฺส ปทสฺส วิสชฺชนา ยุตฺตา.}

\prose{“กิํ สุ ตสฺส มหพฺภยนฺ”ติ ปเญฺห “ทุกฺขมสฺส มหพฺภยนฺ”ติ วิสชฺชนา. ทุวิธํ ทุกฺขํ กายิกญฺจ เจตสิกญฺจ. ยํ กายิกํ อิทํ ทุกฺขํ, ยํ เจตสิกํ อิทํ โทมนสฺสํ. สพฺเพ สตฺตา หิ ทุกฺขสฺส อุพฺพิชฺชนฺติ, นตฺถิ ภยํ ทุกฺเขน สมสมํ, กุโต วา ปน ตสฺส อุตฺตริตรํ. ติสฺโส ทุกฺขตา ทุกฺขทุกฺขตา สงฺขาร\-ทุกฺขตา วิปริณาม\-ทุกฺขตา. ตตฺถ โลโก โอธโส กทาจิ กรหจิ ทุกฺข\-ทุกฺขตาย มุจฺจติ. ตถา วิปริณาม\-ทุกฺขตาย. ตํ กิสฺส เหตุ, โหนฺติ โลเก อปฺปาพาธาปิ ทีฆายุกาปิ. สงฺขาร\-ทุกฺขตาย ปน โลโก อนุปาทิเสสาย นิพฺพานธาตุยา มุจฺจติ, ตสฺมา สงฺขาร\-ทุกฺขตา ทุกฺขํ โลกสฺสาติ กตฺวา ทุกฺขมสฺส มหพฺภยนฺติ. เตน จ จตุตฺถสฺส ปทสฺส วิสชฺชนา ยุตฺตา. เตนาห ภควา “อวิชฺชาย นิวุโต โลโก”ติ.}

\prose{** สวนฺติ สพฺพธิ โสตา, (อิจฺจายสฺมา อชิโต,)}

\prose{โสตานํ กิํ นิวารณํ.}

\setlength{\csromangathapagewidth}{0pt}
\setlength{\csromangathawakwidth}{0pt}
\csromangathameasuregroup{โสตานํ สํวรํ พฺรูหิ,}{\csromangathabat{โสตานํ สํวรํ พฺรูหิ,}{เกน โสตา ปิธียเร\textsuperscript{1}.}}
\csromangathasetleft{โสตานํ สํวรํ พฺรูหิ,}
\csromangathagroup{\csromangathastanza{\csromangathabat{โสตานํ สํวรํ พฺรูหิ,}{เกน โสตา ปิธียเร\csromansharedfootnote{b026e62be20cf5e3}{ปิถียเร (สี), ปิธิยฺยเร (ก)}.}}}

\prose{อิมานิ จตฺตาริ ปทานิ ปุจฺฉิตานิ. เต เทฺว ปญฺหา. กสฺมา, อิเม หิ พตฺวาธิวจเนน ปุจฺฉิตา. เอวํ สมาปนฺนสฺส โลกสฺส เอวํ สํกิลิฏฺฐสฺส กิํ โลกสฺส โวทานํ วุฏฺฐานมิติ, เอวญฺหิ อาห.}

\prose{\csromanfolio{13}\textbf{สวนฺติ สพฺพธิ โสตา}ติ อสมาหิตสฺส สวนฺติ อภิชฺฌาพฺยาปาท\-ปฺปมาท\-พหุลสฺส. ตตฺถ ยา อภิชฺฌา อยํ โลโภ อกุสลมูลํ, โย พฺยาปาโท อยํ โทโส อกุสลมูลํ, โย ปมาโท อยํ โมโห อกุสลมูลํ. ตสฺเสวํ อสมาหิตสฺส ฉสุ อายตเนสุ ตณฺหา สวนฺติ รูปตณฺหา สทฺทตณฺหา คนฺธตณฺหา รสตณฺหา โผฏฺฐพฺพ\-ตณฺหา ธมฺมตณฺหา, ยถาห ภควา\csromandash{}}

\prose{“สวตี”ติ จ โข ภิกฺขเว ฉนฺเนตํ อชฺฌตฺติกานํ อายตนานํ อธิวจนํ. จกฺขุ สวติ มนาปิเกสุ รูเปสุ, อมนาปิเกสุ\csromansharedfootnote{1e9dcef71c3022b3}{อมนาปิเยสุ (ก)} ปฏิหญฺญตีติ. โสตํ ฯเปฯ. ฆานํ. ชิวฺหา. กาโย. มโน สวติ มนาปิเกสุ ธมฺเมสุ อมนาปิเกสุ ปฏิหญฺญตีติ. อิติ สพฺพา จ สวติ, สพฺพถา จ สวติ. เตนาห “สวนฺติ สพฺพธิ โสตา”ติ.}

\prose{“โสตานํ กิํ นิวารณนฺ”ติ ปริยุฏฺฐาน\-วิฆาตํ ปุจฺฉติ, อิทํ โวทานํ. “โสตานํ สํวรํ พฺรูหิ, เกน โสตา ปิธียเร”ติ อนุสย\-สมุคฺฆาตํ ปุจฺฉติ, อิทํ วุฏฺฐานํ. ตตฺถ วิสชฺชนา\csromandash{}}

\prose{\csromansymbolfootnote{*}{ขุ 1. 434 ปิฏฺเฐ สุตฺตนิปาเต; อุปริ 61, 180, 225 ปิฏฺเฐสุปิ. 2. ตสฺส (สี)}“ยานิ โสตานิ โลกสฺมิํ, (อชิตาติ ภควา,)}

\prose{สติ เตสํ นิวารณํ.}

\setlength{\csromangathapagewidth}{0pt}
\setlength{\csromangathawakwidth}{0pt}
\csromangathameasuregroup{โสตานํ สํวรํ พฺรูมิ,}{\csromangathabat{โสตานํ สํวรํ พฺรูมิ,}{ปญฺญาเยเต ปิธียเร”ติ.}}
\csromangathasetleft{โสตานํ สํวรํ พฺรูมิ,}
\csromangathagroup{\csromangathastanza{\csromangathabat{โสตานํ สํวรํ พฺรูมิ,}{ปญฺญาเยเต ปิธียเร”ติ.}}}

\prose{กายคตาย สติยา ภาวิตาย พหุลีกตาย จกฺขุ นาวิญฺฉติ มนาปิเกสุ รูเปสุ, อมนาปิเกสุ น ปฏิหญฺญติ, โสตํ ฯเปฯ ฆานํ. ชิวฺหา. กาโย. มโน นาวิญฺฉติ มนาปิเกสุ ธมฺเมสุ, อมนาปิเกสุ น ปฏิหญฺญติ. เกน การเณน, สํวุต\-นิวาริต\-ตฺตา อินฺทฺริยานํ. เกน เต สํวุต\-นิวาริตา, สติอารกฺเขน. เตนาห ภควา “สติ เตสํ นิวารณนฺ”ติ.}

\prose{ปญฺญาย อนุสยา ปหียนฺติ, อนุสเยสุ ปหีเนสุ ปริยุฏฺฐานา ปหียนฺติ. กิสฺส2, อนุสยสฺส ปหีนตฺตา. ตํ ยถา ขนฺธ\-วนฺตสฺส รุกฺขสฺส อนวเสส\-มูลุ\-ทฺธรเณ กเต ปุปฺผผลปลฺลว\-งฺกุร\-สนฺตติ สมุจฺฉินฺนา ภวติ.}

\prose{\csromanfolio{14}เอวํ อนุสเยสุ ปหีเนสุ ปริยุฏฺฐาน\-สนฺตติ สมุจฺฉินฺนา ภวติ ปิทหิตา ปฏิจฺฉนฺนา. เกน, ปญฺญาย. เตนาห ภควา “ปญฺญาเยเต ปิธียเร”ติ.}

\prose{\csromansymbolfootnote{*}{ขุ 1. 434 ปิฏฺเฐ สุตฺตนิปาเต; อุปริ 61 ปิฏฺเฐ.}“ปญฺญา เจว สติ จ, (อิจฺจายสฺมา อชิโต,)}

\prose{นามรูปญฺจ มาริส.}

\setlength{\csromangathapagewidth}{0pt}
\setlength{\csromangathawakwidth}{0pt}
\csromangathameasuregroup{เอตํ เม ปุฏฺโฐ ปพฺรูหิ, \\ \textsuperscript{*}“ยเมตํ ปญฺหํ อปุจฺฉิ\textsuperscript{1}, \\ ยตฺถ นามญฺจ รูปญฺจ, \\ วิญฺญาณสฺส นิโรเธน,}{\csromangathabat{เอตํ เม ปุฏฺโฐ ปพฺรูหิ,}{กตฺเถตํ อุปรุชฺฌตี”ติ.} \\ \csromangathabat{\textsuperscript{*}“ยเมตํ ปญฺหํ อปุจฺฉิ\textsuperscript{1},}{อชิต ตํ วทามิ เต.} \\ \csromangathabat{ยตฺถ นามญฺจ รูปญฺจ,}{อเสสํ อุปรุชฺฌติ.} \\ \csromangathabat{วิญฺญาณสฺส นิโรเธน,}{เอตฺเถตํ อุปรุชฺฌตี”ติ.}}
\csromangathasetleft{เอตํ เม ปุฏฺโฐ ปพฺรูหิ, \\ \textsuperscript{*}“ยเมตํ ปญฺหํ อปุจฺฉิ\textsuperscript{1}, \\ ยตฺถ นามญฺจ รูปญฺจ, \\ วิญฺญาณสฺส นิโรเธน,}
\csromangathagroup{\csromangathastanza{\csromangathabat{เอตํ เม ปุฏฺโฐ ปพฺรูหิ,}{กตฺเถตํ อุปรุชฺฌตี”ติ.} \\ \csromangathabat{\csromansymbolmark{*}“ยเมตํ ปญฺหํ อปุจฺฉิ\csromansharedfootnote{7503d1ec6e3cb507}{มํ ปุจฺฉิ (ก)},}{อชิต ตํ วทามิ เต.}}\csromangathastanza{\csromangathabat{ยตฺถ นามญฺจ รูปญฺจ,}{อเสสํ อุปรุชฺฌติ.} \\ \csromangathabat{วิญฺญาณสฺส นิโรเธน,}{เอตฺเถตํ อุปรุชฺฌตี”ติ.}}}

\prose{อยํ ปเญฺห\csromansharedfootnote{fb745e69dd94b661}{ปโญฺห (สี, ก) เนตฺติวิภาวนี ปสฺสิตพฺพา.} อนุสนฺธิํ ปุจฺฉติ. อนุสนฺธิํ ปุจฺฉนฺโต กิํ ปุจฺฉติ, อนุปาทิเสสํ นิพฺพานธาตุํ. ตีณิ จ สจฺจานิ สงฺขตานิ นิโรธ\-ธมฺมานิ ทุกฺขํ สมุทโย มคฺโค, นิโรโธ อสงฺขโต. ตตฺถ สมุทโย ทฺวีสุ ภูมีสุ ปหียติ ทสฺสนภูมิยา จ ภาวนาภูมิยา จ. ทสฺสเนน ตีณิ สํโยชนานิ ปหียนฺติ สกฺกายทิฏฺฐิ วิจิกิจฺฉา สีลพฺพต\-ปรามาโส, ภาวนาย สตฺต สํโยชนานิ ปหียนฺติ กามจฺฉนฺโท พฺยาปาโท รูปราโค อรูปราโค มาโน อุทฺธจฺจํ อวิชฺชาวเสสา\csromansharedfootnote{59c584b1ef2fd028}{อวิชฺชา จ นิรวเสสา (สี, ก)}. เตธาตุเก อิมานิ ทส สํโยชนานิ ปญฺโจ\-รมฺภาคิยานิ ปญฺจุ\-ทฺธมฺภาคิยานิ.}

\proseitem{12}{ตตฺถ ตีณิ สํโยชนานิ สกฺกายทิฏฺฐิ วิจิกิจฺฉา สีลพฺพต\-ปรามาโส อนญฺญาต\-ญฺญสฺสามี\-ตินฺทฺริยํ อธิฏฺฐาย นิรุชฺฌนฺติ. สตฺต สํโยชนานิ กามจฺฉนฺโท พฺยาปาโท รูปราโค อรูปราโค มาโน อุทฺธจฺจํ อวิชฺชา วเสสา อญฺญินฺทฺริยํ อธิฏฺฐาย นิรุชฺฌนฺติ. ยํ ปน เอวํ ชานาติ “ขีณา เม ชาตี”ติ, อิทํ ขเย ญาณํ. “นาปรํ อิตฺถตฺตายา”ติ ปชานาติ, อิทํ อนุปฺปาเท ญาณํ. อิมานิ เทฺว ญาณานิ อญฺญาตาวิ\-นฺทฺริยํ. ตตฺถ ยญฺจ อนญฺญาต\-ญฺญสฺสามี\-ตินฺทฺริยํ ยญฺจ อญฺญินฺทฺริยํ, อิมานิ อคฺคผลํ อรหตฺตํ ปาปุณนฺตสฺส นิรุชฺฌนฺติ, ตตฺถ ยญฺจ ขเย ญาณํ ยญฺจ อนุปฺปาเท ญาณํ, อิมานิ เทฺว ญาณานิ เอกปญฺญา.}

\prose{อปิ จ อารมฺมณ\-สงฺเกเตน เทฺว นามานิ ลพฺภนฺติ, “ขีณา เม ชาตี”ติ ปชานนฺตสฺส ขเย ญาณนฺติ นามํ ลภติ, “นาปรํ อิตฺถตฺตายา”ติ \csromanfolio{15}ปชานนฺตสฺส อนุปฺปาเท ญาณนฺติ นามํ ลภติ. สา ปชานนฏฺเฐน ปญฺญา, ยถาทิฏฺฐํ อปิลาป\-นฏฺเฐน สติ.}

\proseitem{13}{ตตฺถ เย ปญฺจุ\-ปาทาน\-กฺขนฺธา, อิทํ นามรูปํ. ตตฺถ เย ผสฺส\-ปญฺจมกา ธมฺมา, อิทํ นามํ. ยานิ ปญฺจินฺทฺริยานิ รูปานิ, อิทํ รูปํ. ตทุภยํ นามรูปํ วิญฺญาณ\-สมฺปยุตฺตํ ตสฺส นิโรธํ ภควนฺตํ ปุจฺฉนฺโต อายสฺมา อชิโต ปารายเน เอวมาห\csromandash{}}

\setlength{\csromangathapagewidth}{0pt}
\setlength{\csromangathawakwidth}{0pt}
\csromangathameasuregroup{“ปญฺญา เจว สติ จ, \\ เอตํ เม ปุฏฺโฐ ปพฺรูหิ,}{\csromangathabat{“ปญฺญา เจว สติ จ,}{นามรูปญฺจ มาริส.} \\ \csromangathabat{เอตํ เม ปุฏฺโฐ ปพฺรูหิ,}{กตฺเถตํ อุปรุชฺฌตี”ติ.}}
\csromangathasetleft{“ปญฺญา เจว สติ จ, \\ เอตํ เม ปุฏฺโฐ ปพฺรูหิ,}
\csromangathagroup{\csromangathastanza{\csromangathabat{“ปญฺญา เจว สติ จ,}{นามรูปญฺจ มาริส.} \\ \csromangathabat{เอตํ เม ปุฏฺโฐ ปพฺรูหิ,}{กตฺเถตํ อุปรุชฺฌตี”ติ.}}}

\prose{ตตฺถ สติ จ ปญฺญา จ จตฺตาริ อินฺทฺริยานิ, สติ เทฺว อินฺทฺริยานิ สตินฺทฺริยญฺจ สมาธินฺทฺริยญฺจ, ปญฺญา เทฺว อินฺทฺริยานิ ปญฺญินฺทฺริยญฺจ วีริยินฺทฺริยญฺจ. ยา อิเมสุ จตูสุ อินฺทฺริเยสุ สทฺทหนา โอกปฺปนา, อิทํ สทฺธินฺทฺริยํ. ตตฺถ ยา สทฺธา\-ธิปเตยฺยา จิตฺเตกคฺคตา, อยํ ฉนฺทสมาธิ. สมาหิเต จิตฺเต กิเลสานํ วิกฺขมฺภ\-นตาย ปฏิสงฺขาน\-พเลน วา ภาวนาพเลน วา, อิทํ ปหานํ. ตตฺถ เย อสฺสาสปสฺสาสา วิตกฺกวิจารา สญฺญาเวทยิตา สรสงฺกปฺปา, อิเม สงฺขารา. อิติ ปุริมโก จ ฉนฺทสมาธิ, กิเลส\-วิกฺขมฺภ\-นตาย จ ปหานํ อิเม จ สงฺขารา, ตทุภยํ ฉนฺท\-สมาธิ\-ปฺปธาน\-สงฺขาร\-สมนฺนาคตํ อิทฺธิปาทํ ภาเวติ วิเวกนิสฺสิตํ วิราคนิสฺสิตํ นิโรธ\-นิสฺสิตํ โวสฺสคฺค\-ปริณามิํ. ตตฺถ ยา วีริยา\-ธิปเตยฺยา จิตฺเตกคฺคตา, อยํ วีริยสมาธิ ฯเปฯ. ตตฺถ ยา จิตฺตา\-ธิปเตยฺยา จิตฺเตกคฺคตา, อยํ จิตฺตสมาธิ ฯเปฯ. ตตฺถ ยา วีมํสา\-ธิปเตยฺยา จิตฺเตกคฺคตา, อยํ วีมํสาสมาธิ. สมาหิเต จิตฺเต กิเลสานํ วิกฺขมฺภ\-นตาย ปฏิสงฺขาน\-พเลน วา ภาวนาพเลน วา, อิทํ ปหานํ. ตตฺถ เย อสฺสาสปสฺสาสา วิตกฺกวิจารา สญฺญาเวทยิตา สรสงฺกปฺปา, อิเม สงฺขารา. อิติ ปุริมโก จ วีมํสาสมาธิ, กิเลส\-วิกฺขมฺภ\-นตาย จ ปหานํ อิเม จ สงฺขารา, ตทุภยํ วีมํสา\-สมาธิ\-ปฺปธาน\-สงฺขาร\-สมนฺนาคตํ อิทฺธิปาทํ ภาเวติ วิเวกนิสฺสิตํ วิราคนิสฺสิตํ นิโรธ\-นิสฺสิตํ โวสฺสคฺค\-ปริณามิํ.}

\setlength{\csromangathapagewidth}{0pt}
\setlength{\csromangathawakwidth}{0pt}
\csromangathameasuregroup{}{สพฺโพ สมาธิ ญาณมูลโก ญาณปุพฺพงฺคโม ญาณานุปริวตฺติ, \\ ยถา ปุเร ตถา ปจฺฉา, ยถา ปจฺฉา ตถา ปุเร. \\ ยถา ทิวา ตถา รตฺติํ\textsuperscript{1}, ยถา รตฺติํ ตถา ทิวา.}
\csromangathameasurewak{สพฺโพ สมาธิ ญาณมูลโก ญาณปุพฺพงฺคโม ญาณานุปริวตฺติ, \\ ยถา ปุเร ตถา ปจฺฉา, ยถา ปจฺฉา ตถา ปุเร. \\ ยถา ทิวา ตถา รตฺติํ\textsuperscript{1}, ยถา รตฺติํ ตถา ทิวา.}
\setlength{\csromangathaleftcol}{0pt}
\csromangathagroup{\csromangathastanza{\csromangathawak{\csromangathaitemline{14}{สพฺโพ สมาธิ ญาณมูลโก ญาณปุพฺพงฺคโม ญาณา\-นุปริวตฺติ,} \\ \csromangathacontline{ยถา ปุเร ตถา ปจฺฉา, ยถา ปจฺฉา ตถา ปุเร.} \\ \csromangathacontline{ยถา ทิวา ตถา รตฺติํ\csromansharedfootnote{a5743f436274f122}{รตฺติ (ก) อยํ คาถา ขุ 2. 284 ปิฏฺเฐ อญฺญถา ทิสฺสติ.}, ยถา รตฺติํ ตถา ทิวา.}}}}

\prose{\csromanfolio{16}อิติ วิวเฏน เจตสา อปริโยนทฺเธน สปฺปภาสํ จิตฺตํ ภาเวติ ปญฺจินฺทฺริยานิ กุสลานิ จิตฺตสหภูนิ จิตฺเต อุปฺปชฺชมาเน อุปฺปชฺชนฺติ, จิตฺเต นิรุชฺฌมาเน นิรุชฺฌนฺติ. นามรูปญฺจ วิญฺญาณเหตุกํ วิญฺญาณปจฺจยา นิพฺพตฺตํ, ตสฺส มคฺเคน เหตุ อุปจฺฉินฺโน, วิญฺญาณํ อนาหารํ อนภินนฺทิตํ อปฺปฏิสนฺธิกํ ตํ นิรุชฺฌติ. นามรูปมปิ อเหตุ อปฺปจฺจยํ ปุนพฺภวํ น นิพฺพตฺตยติ\csromansharedfootnote{30604144677b47e5}{นิพฺพตฺติยติ (ก)}. เอวํ วิญฺญาณสฺส นิโรธา ปญฺญา จ สติ จ นามรูปญฺจ นิรุชฺฌติ. เตนาห ภควา\csromandash{}}

\setlength{\csromangathapagewidth}{0pt}
\setlength{\csromangathawakwidth}{0pt}
\csromangathameasuregroup{“ยเมตํ ปญฺหํ อปุจฺฉิ, \\ ยตฺถ นามญฺจ รูปญฺจ,}{\csromangathabat{“ยเมตํ ปญฺหํ อปุจฺฉิ,}{อชิต ตํ วทามิ เต.} \\ \csromangathabat{ยตฺถ นามญฺจ รูปญฺจ,}{อเสสํ อุปรุชฺฌติ.}}
\csromangathasetleft{“ยเมตํ ปญฺหํ อปุจฺฉิ, \\ ยตฺถ นามญฺจ รูปญฺจ,}
\csromangathagroup{\csromangathastanza{\csromangathabat{“ยเมตํ ปญฺหํ อปุจฺฉิ,}{อชิต ตํ วทามิ เต.} \\ \csromangathabat{ยตฺถ นามญฺจ รูปญฺจ,}{อเสสํ อุปรุชฺฌติ.}}}

\prose{วิญฺญาณสฺส นิโรเธน, เอตฺเถตํ อุปรุชฺฌตี”ติ\csromansharedfootnote{338dd9d64fd7ee1e}{[มิสฺสินฺคฺ โนเต 0]}.}

\proseitem{14}{“เย จ สงฺขาต\-ธมฺมาเส, (อิจฺจายสฺมา อชิโต,)}

\prose{เย จ เสกฺขา ปุถู อิธ.}

\setlength{\csromangathapagewidth}{0pt}
\setlength{\csromangathawakwidth}{0pt}
\csromangathameasuregroup{เตสํ เม นิปโก อิริยํ,}{\csromangathabat{เตสํ เม นิปโก อิริยํ,}{ปุฏฺโฐ ปพฺรูหิ มาริสา”ติ\textsuperscript{1}.}}
\csromangathasetleft{เตสํ เม นิปโก อิริยํ,}
\csromangathagroup{\csromangathastanza{\csromangathabat{เตสํ เม นิปโก อิริยํ,}{ปุฏฺโฐ ปพฺรูหิ มาริสา”ติ\csromansharedfootnote{9fe577add4a01715}{ขุ 1. 434 ปิฏฺเฐ สุตฺตนิปาเต; อุปริ 19, 226 ปิฏฺเฐสุปิ.}.}}}

\proseitem{15}{อิมานิ ตีณิ ปทานิ ปุจฺฉิตานิ, เต ตโย ปญฺหา. กิสฺส, เสขาเสข\-วิปสฺสนาปุพฺพงฺคม\-ปฺปหาน\-โยเคน, เอวญฺหิ อาห. “เย จ สงฺขาต\-ธมฺมาเส”ติ อรหตฺตํ ปุจฺฉติ, “เย จ เสกฺขา ปุถู อิธา”ติ เสขํ ปุจฺฉติ, “เตสํ เม นิปโก อิริยํ, ปุฏฺโฐ ปพฺรูหิ มาริสา”ติ วิปสฺสนา\-ปุพฺพงฺคมํ ปหานํ ปุจฺฉติ. ตตฺถ วิสชฺชนา\csromandash{}}

\prose{“กาเมสุ นาภิคิชฺเฌยฺย, (อชิตาติ ภควา,)}

\prose{มนสา’นาวิโล สิยา.}

\setlength{\csromangathapagewidth}{0pt}
\setlength{\csromangathawakwidth}{0pt}
\csromangathameasuregroup{กุสโล สพฺพธมฺมานํ,}{\csromangathabat{กุสโล สพฺพธมฺมานํ,}{สโต ภิกฺขุ ปริพฺพเช”ติ\textsuperscript{1}.}}
\csromangathasetleft{กุสโล สพฺพธมฺมานํ,}
\csromangathagroup{\csromangathastanza{\csromangathabat{กุสโล สพฺพธมฺมานํ,}{สโต ภิกฺขุ ปริพฺพเช”ติ\csromansharedfootnote{9fe577add4a01715}{ขุ 1. 434 ปิฏฺเฐ สุตฺตนิปาเต; อุปริ 19, 226 ปิฏฺเฐสุปิ.}.}}}

\prose{ภควโต สพฺพํ กายกมฺมํ ญาณ\-ปุพฺพงฺคมํ ญาณา\-นุปริวตฺติ, สพฺพํ วจีกมฺมํ ญาณ\-ปุพฺพงฺคมํ ญาณา\-นุปริวตฺติ, สพฺพํ มโนกมฺมํ ญาณ\-ปุพฺพงฺคมํ ญาณา\-นุปริวตฺติ.  อตีเต อํเส อปฺปฏิหต\-ญาณทสฺสนํ, อนาคเต อํเส อปฺปฏิหต\-ญาณทสฺสนํ, ปจฺจุปฺปนฺเน อํเส อปฺปฏิหต\-ญาณทสฺสนํ.}

\prose{\csromanfolio{17}โก จ ญณทสฺสนสฺส ปฏิฆาโต, ยํ อนิจฺเจ ทุกฺเข อนตฺตนิ จ อญฺญาณํ อทสฺสนํ, อยํ ญาณทสฺสนสฺส ปฏิฆาโต. ยถา อิธ ปุริโส ตารกรูปานิ ปสฺเสยฺย, โน จ คณน\-สงฺเกเตน ชาเนยฺย, อยํ ญาณทสฺสนสฺส ปฏิฆาโต.}

\prose{ภควโต ปน อปฺปฏิหต\-ญาณทสฺสนํ, อนาวรณ\-ญาณทสฺสนา หิ พุทฺธา ภควนฺโต. ตตฺถ เสเขน ทฺวีสุ ธมฺเมสุ จิตฺตํ รกฺขิตพฺพํ เคธา จ รชนีเยสุ ธมฺเมสุ, โทสา จ ปริยุฏฺฐานีเยสุ. ตตฺถ ยา อิจฺฉา มุจฺฉา ปตฺถนา ปิยายนา กีฬนา, ตํ ภควา นิวาเรนฺโต เอวมาห “กาเมสุ นาภิคิชฺเฌยฺยา”ติ.}

\prose{“มนสา’นาวิโล สิยา”ติ ปริยุฏฺฐาน\-วิฆาตํ อาห. ตถา หิ เสโข อภิคิชฺฌนฺโต อสมุปฺปนฺนญฺจ กิเลสํ อุปฺปาเทติ, อุปฺปนฺนญฺจ กิเลสํ ผาติํ กโรติ. โย ปน อนาวิล\-สงฺกปฺโป อนภิคิชฺฌนฺโต วายมติ, โส อนุปฺปนฺนานํ ปาปกานํ อกุสลานํ ธมฺมานํ อนุปฺปาทาย ฉนฺทํ ชเนติ วายมติ วีริยํ อารภติ จิตฺตํ ปคฺคณฺหาติ ปทหติ. โส อุปฺปนฺนานํ ปาปกานํ อกุสลานํ ธมฺมานํ ปหานาย ฉนฺทํ ชเนติ วายมติ วีริยํ อารภติ จิตฺตํ ปคฺคณฺหาติ ปทหติ. โส อนุปฺปนฺนานํ กุสลานํ ธมฺมานํ อุปฺปาทาย ฉนฺทํ ชเนติ วายมติ วีริยํ อารภติ จิตฺตํ ปคฺคณฺหาติ ปทหติ. โส อุปฺปนฺนานํ กุสลานํ ธมฺมานํ ฐิติยา อสมฺโมสาย ภิยฺโยภาวาย เวปุลฺลาย ภาวนาย ปาริปูริยา ฉนฺทํ ชเนติ วายมติ วีริยํ อารภติ จิตฺตํ ปคฺคณฺหาติ ปทหติ.}

\proseitem{16}{กตเม อนุปฺปนฺนา ปาปกา อกุสลา ธมฺมา, กามวิตกฺโก พฺยาปาทวิตกฺโก วิหิํสาวิตกฺโก, อิเม อนุปฺปนฺนา ปาปกา อกุสลา ธมฺมา. กตเม อุปฺปนฺนา ปาปกา อกุสลา ธมฺมา, อนุสยา อกุสลมูลานิ, อิเม อุปฺปนฺนา ปาปกา อกุสลา ธมฺมา. กตเม อนุปฺปนฺนา กุสลา ธมฺมา, ยานิ โสตาปนฺนสฺส อินฺทฺริยานิ, อิเม อนุปฺปนฺนา กุสลา ธมฺมา. กตเม อุปฺปนฺนา กุสลา ธมฺมา, ยานิ อฏฺฐมกสฺส อินฺทฺริยานิ, อิเม อุปฺปนฺนา กุสลา ธมฺมา.}

\prose{เยน กามวิตกฺกํ วาเรติ, อิทํ สตินฺทฺริยํ. เยน พฺยาปาท\-วิตกฺกํ วาเรติ, อิทํ สมาธินฺทฺริยํ. เยน วิหิํสา\-วิตกฺกํ วาเรติ, อิทํ วีริยินฺทฺริยํ. \csromanfolio{18}เยน อุปฺปนฺนุปฺปนฺเน ปาปเก อกุสเล ธมฺเม ปชหติ วิโนเทติ พฺยนฺตีกโรติ อนภาวํ คเมติ นาธิวาเสติ, อิทํ ปญฺญินฺทฺริยํ. ยา อิเมสุ จตูสุ อินฺทฺริเยสุ สทฺทหนา โอกปฺปนา, อิทํ สทฺธินฺทฺริยํ.}

\prose{ตตฺถ สทฺธินฺทฺริยํ กตฺถ ทฏฺฐพฺพํ, จตูสุ โสตาปตฺติยงฺเคสุ. วีริยินฺทฺริยํ กตฺถ ทฏฺฐพฺพํ, จตูสุ สมฺม\-ปฺปธาเนสุ. สตินฺทฺริยํ กตฺถ ทฏฺฐพฺพํ, จตูสุ สติปฏฺฐาเนสุ. สมาธินฺทฺริยํ กตฺถ ทฏฺฐพฺพํ, จตูสุ ฌาเนสุ. ปญฺญินฺทฺริยํ กตฺถ ทฏฺฐพฺพํ, จตูสุ อริยสจฺเจสุ. เอวํ เสโข สพฺเพหิ กุสเลหิ ธมฺเมหิ อปฺปมตฺโต วุตฺโต ภควตา อนาวิลตาย มนสา. เตนาห ภควา “มนสา’นาวิโลสิยา”ติ.}

\proseitem{17}{“กุสโล สพฺพธมฺมานนฺ”ติ โลโก นาม ติวิโธ กิเลสโลโก ภวโลโก อินฺทฺริยโลโก. ตตฺถ กิเลสโลเกน ภวโลโก สมุทาคจฺฉติ, โส อินฺทฺริยานิ นิพฺพตฺเตติ, อินฺทฺริเยสุ ภาวิยมาเนสุ เนยฺยสฺส ปริญฺญา ภวติ. สา ทุวิเธน อุปปริ\-กฺขิ\-ตพฺพา ทสฺสน\-ปริญฺญาย จ ภาวนา\-ปริญฺญาย จ. ยทา หิ เสโข เญยฺยํ ปริชานาติ, ตทา นิพฺพิทา\-สหคเตหิ สญฺญามนสิกาเรหิ เนยฺยํ ปริญฺญาตํ ภวติ. ตสฺส เทฺว ธมฺมา โกสลฺลํ คจฺฉนฺติ ทสฺสน\-โกสลฺลญฺจ ภาวนา\-โกสลฺลญฺจ.}

\prose{ตํ ญาณํ ปญฺจวิเธน เวทิตพฺพํ อภิญฺญา ปริญฺญา ปหานํ ภาวนา สจฺฉิกิริยา. ตตฺถ กตมา อภิญฺญา, ยํ ธมฺมานํ สลกฺขเณ ญาณํ ธมฺม\-ปฏิสมฺภิทา จ อตฺถ\-ปฏิสมฺภิทา จ, อยํ อภิญฺญา.}

\prose{ตตฺถ กตมา ปริญฺญา, เอวํ อภิชานิตฺวา ยา ปริชานนา “อิทํ กุสลํ, อิทํ อกุสลํ, อิทํ สาวชฺชํ, อิทํ อนวชฺชํ, อิทํ กณฺหํ, อิทํ สุกฺกํ, อิทํ เสวิตพฺพํ, อิทํ น เสวิตพฺพํ, อิเม ธมฺมา เอวํคหิตา, อิทํ ผลํ นิพฺพตฺเตนฺติ\csromansharedfootnote{6b9f884c5d5b5bfd}{นิพฺพตฺตาเปนฺติ (ก)}, เตสํ เอวํคหิตานํ อยํ อตฺโถ”ติ, อยํ ปริญฺญา.}

\prose{เอวํ ปริชานิตฺวา ตโย ธมฺมา อวสิฏฺฐา ภวนฺติ ปหาตพฺพา ภาเวตพฺพา สจฺฉิกาตพฺพา จ, ตตฺถ กตเม ธมฺมา ปหาตพฺพา, เย อกุสลา, ตตฺถ กตเม ธมฺมา ภาเวตพฺพา, เย กุสลา, ตตฺถ กตเม ธมฺมา สจฺฉิกาตพฺพา, ยํ อสงฺขตํ, โย เอวํ ชานาติ อยํ วุจฺจติ อตฺถกุสโล \csromanfolio{19}ธมฺมกุสโล กลฺยาณตา\-กุสโล ผลตากุสโล, อายกุสโล อปายกุสโล อุปายกุสโล มหตา โกสลฺเลน สมนฺนาคโตติ, เตนาห ภควา “กุสโล สพฺพธมฺมานนฺ”ติ.}

\prose{“สโต ภิกฺขุ ปริพฺพเช”ติ เตน ทิฏฺฐธมฺมสุขวิหารตฺถํ อภิกฺกนฺเต ปฏิกฺกนฺเต อาโลกิเต วิโลกิเต สมิญฺชิเต\csromansharedfootnote{c92294846cbf0bec}{สมฺมิญฺชิเต (สี)} ปสาริเต สงฺฆาฏิ\-ปตฺต\-จีวร\-ธารเณ อสิเต ปีเต ขายิเต สายิเต อุจฺจาร\-ปสฺสาว\-กมฺเม คเต ฐิเต นิสินฺเน สุตฺเต ชาคริเต ภาสิเต ตุณฺหิภาเว สเตน สมฺปชาเนน วิหาตพฺพํ. อิมา เทฺว จริยา อนุญฺญาตา ภควตา เอกา วิสุทฺธานํ, เอกา วิสุชฺฌนฺตานํ, เก วิสุทฺธา, อรหโต. เก วิสุชฺฌนฺตา, เสกฺขา. กตกิจฺจานิ หิ อรหโต อินฺทฺริยานิ. ยํ โพชฺฌํ, ตํ จตุพฺพิธํ ทุกฺขสฺส ปริญฺญา\-ภิสมเยน สมุทยสฺส ปหานา\-ภิสมเยน มคฺคสฺส ภาวนาภิ\-สมเยน นิโรธสฺส สจฺฉิกิริยาภิ\-สมเยน, อิทํ จตุพฺพิธํ โพชฺฌํ โย เอวํ ชานาติ, อยํ วุจฺจติ สโต อภิกฺกมติ สโต ปฏิกฺกมติ ขยา ราคสฺส ขยา โทสสฺส ขยา โมหสฺส. เตนาห ภควา “สโต ภิกฺขุ ปริพฺพเช”ติ, เตนาห\csromandash{}}

\prose{“กาเมสุ นาภิคิชฺเฌยฺย, (อชิตาติ ภควา.)}

\prose{มนสา’นาวิโล สิยา.}

\setlength{\csromangathapagewidth}{0pt}
\setlength{\csromangathawakwidth}{0pt}
\csromangathameasuregroup{กุสโล สพฺพธมฺมานํ,}{\csromangathabat{กุสโล สพฺพธมฺมานํ,}{สโต ภิกฺขุ ปริพฺพเช”ติ.}}
\csromangathasetleft{กุสโล สพฺพธมฺมานํ,}
\csromangathagroup{\csromangathastanza{\csromangathabat{กุสโล สพฺพธมฺมานํ,}{สโต ภิกฺขุ ปริพฺพเช”ติ.}}}

\prose{เอวํ ปุจฺฉิตพฺพํ, เอวํ วิสชฺชิตพฺพํ. สุตฺตสฺส จ อนุคีติ อตฺถโต จ พฺยญฺชนโต จ สมาเนตพฺพา\csromansharedfootnote{41a8672dce6135cc}{สมานยิตพฺพา. (สี, ก)}. อตฺถาปคตํ หิ พฺยญฺชนํ สมฺผปฺปลาปํ ภวติ. ทุนฺนิกฺขิตฺตสฺส ปทพฺยญฺชนสฺส อตฺโถปิ ทุนฺนโย ภวติ, ตสฺมา อตฺถ\-พฺยญฺชนู\-เปตํ สงฺคายิตพฺพํ. สุตฺตญฺจ ปวิจินิตพฺพํ. กิํ อิทํ สุตฺตํ อาหจฺจ วจนํ อนุสนฺธิ\-วจนํ นีตตฺถํ เนยฺยตฺถํ สํกิเลส\-ภาคิยํ นิพฺเพธ\-ภาคิยํ อเสกฺข\-ภาคิยํ. กุหิํ อิมสฺส สุตฺตสฺส สพฺพานิ สจฺจานิ ปสฺสิตพฺพานิ, อาทิ\-มชฺฌ\-ปริโยสาเนติ. เอวํ สุตฺตํ ปวิเจตพฺพํ. เตนาห อายสฺมา มหากจฺจายโน “ยํ ปุจฺฉิตญฺจ วิสฺสชฺชิตญฺจ, สุตฺตสฺส ยา จ อนุคีตี”ติ. นิยุตฺโต วิจโย หาโร.}

\csromanmatikaanchor{mtk.7}
\csromantocmarkref{section}{3. ยุตฺติหารวิภงฺค}{mtk.7}
\bookmark[dest={mtk.7},level=1]{3. ยุตฺติหารวิภงฺค}
\csromanheader{1}{3. \textbf{ยุตฺติ}\-\textbf{หาร}\-\textbf{วิภงฺค}}
\proseitem{18}{\csromanfolio{20}ตตฺถ กตโม ยุตฺติหาโร, “สพฺเพสํ หารานนฺ”ติ, อยํ ยุตฺติหาโร. กิํ โยชยติ\csromansharedfootnote{e18849c06894f52f}{โยเชติ (สี)}, จตฺตาโร มหาปเทสา พุทฺธาปเทโส สํฆาปเทโส สมฺพหุล\-ตฺเถรา\-ปเทโส\csromansharedfootnote{834575555955c4da}{สมฺปหุล... (ก)} เอก\-ตฺเถรา\-ปเทโส. อิเม จตฺตาโร มหาปเทสา, ตานิ ปทพฺยญฺชนานิ สุตฺเต โอตารยิตพฺพานิ, วินเย สนฺทสฺส\-ยิตพฺพานิ, ธมฺมตายํ อุปนิกฺขิปิตพฺพานิ.}

\prose{กตมสฺมิํ สุตฺเต โอตารยิตพฺพานิ, จตูสุ อริยสจฺเจสุ. กตมสฺมิํ วินเย สนฺทสฺส\-ยิตพฺพานิ, ราควินเย โทสวินเย โมหวินเย. กตมิสฺสํ\csromansharedfootnote{3b9b80829831bd3b}{กตมิยํ (สี) 4. สํ 1. 37 ปิฏฺเฐ.} ธมฺมตายํ อุปนิกฺขิปิตพฺพานิ, ปฏิจฺจ\-สมุปฺปาเท. ยทิ จตูสุ อริยสจฺเจสุ อวตรติ, กิเลสวินเย สนฺทิสฺสติ, ธมฺมตญฺจ น วิโลเมติ, เอวํ อาสเว น ชเนติ. จตูหิ มหาปเทเสหิ ยํ ยํ ยุชฺชติ, เยน เยน ยุชฺชติ, ยถา ยถา ยุชฺชติ, ตํ ตํ คเหตพฺพํ.}

\proseitem{19}{ปญฺหํ ปุจฺฉิเตน กติ ปทานิ ปเญฺหติ ปทโส ปริโย\-คาหิตพฺพํ วิเจตพฺพํ. ยทิ สพฺพานิ ปทานิ เอกํ อตฺถํ อภิวทนฺติ, เอโก ปโญฺห. อถ จตฺตาริ ปทานิ เอกํ อตฺถํ อภิวทนฺติ, เอโก ปโญฺห. อถ ตีณิ ปทานิ เอกํ อตฺถํ อภิวทนฺติ, เอโก ปโญฺห. อถ เทฺว ปทานิ เอกํ อตฺถํ อภิวทนฺติ, เอโก ปโญฺห. อถ เอกํ ปทํ เอกํ อตฺถํ อภิวทติ, เอโก ปโญฺห. ตํ อุปปริกฺขมาเนน อญฺญาตพฺพํ กิํ อิเม ธมฺมา นานตฺถา นานาพฺยญฺชนา, อุทาหุ อิเมสํ ธมฺมานํ เอโก อตฺโถ พฺยญฺชนเมว นานนฺติ. ยถา กิํ ภเว, ยถา สา เทวตา ภควนฺตํ ปญฺหํ ปุจฺฉติ.}

\setlength{\csromangathapagewidth}{0pt}
\setlength{\csromangathawakwidth}{0pt}
\csromangathameasuregroup{“เกนสฺสุพฺภาหโต โลโก, \\ เกน สลฺเลน โอติณฺโณ,}{\csromangathabat{“เกนสฺสุพฺภาหโต โลโก,}{เกนสฺสุ ปริวาริโต.} \\ \csromangathabat{เกน สลฺเลน โอติณฺโณ,}{กิสฺส ธูปายิโต สทา”ติ\textsuperscript{1}.}}
\csromangathasetleft{“เกนสฺสุพฺภาหโต โลโก, \\ เกน สลฺเลน โอติณฺโณ,}
\csromangathagroup{\csromangathastanza{\csromangathabat{“เกนสฺสุพฺภาหโต โลโก,}{เกนสฺสุ ปริวาริโต.} \\ \csromangathabat{เกน สลฺเลน โอติณฺโณ,}{กิสฺส ธูปายิโต สทา”ติ\csromansharedfootnote{f5a080d4d238a070}{สํ 1. 37 ปิฏฺเฐ.}.}}}

\prose{อิมานิ จตฺตาริ ปทานิ ปุจฺฉิตานิ. เต ตโย ปญฺหา กถํ ญายติ. ภควา หิ เทวตาย วิสชฺเชติ.}

\setlength{\csromangathapagewidth}{0pt}
\setlength{\csromangathawakwidth}{0pt}
\csromangathameasuregroup{“มจฺจุนาพฺภาหโต โลโก, \\ ตณฺหาสลฺเลน โอติณฺโณ,}{\csromangathabat{“มจฺจุนาพฺภาหโต โลโก,}{ชราย ปริวาริโต.} \\ \csromangathabat{ตณฺหาสลฺเลน โอติณฺโณ,}{อิจฺฉาธูปายิโต สทา”ติ\textsuperscript{1}.}}
\csromangathasetleft{“มจฺจุนาพฺภาหโต โลโก, \\ ตณฺหาสลฺเลน โอติณฺโณ,}
\csromangathagroup{\csromangathastanza{\csromangathabat{“มจฺจุนาพฺภาหโต โลโก,}{ชราย ปริวาริโต.} \\ \csromangathabat{ตณฺหาสลฺเลน โอติณฺโณ,}{อิจฺฉาธูปายิโต สทา”ติ\csromansharedfootnote{214bcfe6664d7a2f}{ขุ 2. 289 ปิฏฺเฐ; สํ 1. 37 ปิฏฺเฐ จ.}.}}}

\proseitem{20}{\csromanfolio{21}ตตฺถ ชรา จ มรณญฺจ อิมานิ เทฺว สงฺขตสฺส สงฺขต\-ลกฺขณานิ. ชรายํ ฐิตสฺส อญฺญถตฺตํ, มรณํ วโย. ตตฺถ ชราย จ มรณสฺส จ อตฺถโต นานตฺตํ. เกน การเณน, คพฺภคตาปิ หิ มียนฺติ, น จ เต ชิณฺณา ภวนฺติ. อตฺถิ จ เทวานํ มรณํ, น จ เตสํ สรีรานิ ชีรนฺติ. สกฺกเตว ชราย ปฏิกมฺมํ กาตุํ, น ปน สกฺกเต มรณสฺส ปฏิกมฺมํ กาตุํ อญฺญเตฺรว อิทฺธิมนฺตานํ อิทฺธิวิสยา. ยํ ปนาห ตณฺหาสลฺเลน โอติณฺโณติ ทิสฺสนฺติ วีตราคา ชีรนฺตาปิ มียนฺตาปิ. ยทิ จ ยถา ชรามรณํ, เอวํ ตณฺหาปิ สิยา. เอวํ สนฺเต สพฺเพ โยพฺพนฏฺฐาปิ วิคตตณฺหา สิยุํ. ยถา จ ตณฺหา ทุกฺขสฺส สมุทโย, เอวํ ชรามรณมฺปิ สิยา ทุกฺขสฺส สมุทโย, น จ สิยา ตณฺหา ทุกฺขสฺส สมุทโย, น หิ ชรามรณํ ทุกฺขสฺส สมุทโย, ตณฺหา ทุกฺขสฺส สมุทโย. ยถา จ ตณฺหา มคฺควชฺฌา, เอวํ ชรามรณมฺปิ สิยา มคฺควชฺฌํ. อิมาย ยุตฺติยา อญฺญมญฺเญหิ การเณหิ คเวสิตพฺพํ. ยทิ จ สนฺทิสฺสติ ยุตฺติ\-สมารูฬฺหํ อตฺถโต จ อญฺญตฺตํ, พฺยญฺชนโตปิ คเวสิตพฺพํ.}

\prose{สลฺโลติ วา ธูปายนนฺติ วา อิเมสํ ธมฺมานํ อตฺถโต เอกตฺตํ. น หิ ยุชฺชติ อิจฺฉาย จ ตณฺหาย จ อตฺถโต อญฺญตฺตํ. ตณฺหาย อธิปฺปาเย อปริปูรมาเน นวสุ อาฆาตวตฺถูสุ โกโธ จ อุปนาโห จ อุปฺปชฺชติ. อิมาย ยุตฺติยา ชราย จ มรณสฺส จ ตณฺหาย จ อตฺถโต อญฺญตฺตํ.}

\prose{ยํ ปนิทํ ภควตา ทฺวีหิ นาเมหิ\csromansharedfootnote{5a92e3f8be855aef}{ธมฺเมหิ (ฏฺฐ)} อภิลปิตํ อิจฺฉาติปิ ตณฺหาติปิ, อิทํ ภควตา พาหิรานํ วตฺถูนํ อารมฺมณวเสน ทฺวีหิ นาเมหิ อภิลปิตํ อิจฺฉาติปิ ตณฺหาติปิ, สพฺพา หิ ตณฺหา อชฺโฌสาน\-ลกฺขเณน เอกลกฺขณา. ยถา สพฺโพ อคฺคิ อุณฺหตฺต\-ลกฺขเณน เอกลกฺขโณ, อปิ จ อุปาทานวเสน อญฺญมญฺญานิ นามานิ ลภติ, กฏฺฐคฺคีติปิ ติณคฺคีติปิ สกลิกคฺคีติปิ โคมยคฺคีติปิ ถุสคฺคีติปิ สงฺการคฺคีติปิ, สพฺโพ หิ อคฺคิ อุณฺหตฺตลกฺขโณว. เอวํ สพฺพา ตณฺหา อชฺโฌสาน\-ลกฺขเณน เอกลกฺขณา, อปิ ตุ อารมฺมณอุปาทานวเสน อญฺญมญฺเญหิ นาเมหิ อภิลปิตา อิจฺฉาอิติปิ ตณฺหาอิติปิ สลฺโลอิติปิ ธูปายนาอิติปิ สริตาอิติปิ วิสตฺติกาอิติปิ สิเนโหอิติปิ กิลมโถอิติปิ ลตาอิติปิ มญฺญนาอิติปิ พนฺโธอิติปิ อาสาอิติปิ ปิปาสาอิติปิ อภินนฺทนาอิติปิ, อิติ สพฺพา ตณฺหา อชฺโฌสาน\-ลกฺขเณน เอกลกฺขณา. ยถา จ เววจเน วุตฺตา.}

\setlength{\csromangathapagewidth}{0pt}
\setlength{\csromangathawakwidth}{0pt}
\csromangathameasuregroup{}{“อาสา จ ปิหา อภินนฺทนา จ, \\ อเนกธาตูสุ สรา ปติฏฺฐิตา. \\ อญฺญาณมูลปฺปภวา ปชปฺปิตา, \\ สพฺพา มยา พฺยนฺติกตา สมูลิกา”ติ\textsuperscript{1}.}
\csromangathameasurewak{“อาสา จ ปิหา อภินนฺทนา จ, \\ อเนกธาตูสุ สรา ปติฏฺฐิตา. \\ อญฺญาณมูลปฺปภวา ปชปฺปิตา, \\ สพฺพา มยา พฺยนฺติกตา สมูลิกา”ติ\textsuperscript{1}.}
\setlength{\csromangathaleftcol}{0pt}
\csromangathagroup{\csromangathastanza{\csromangathawak{\csromanfolio{22}“อาสา จ ปิหา อภินนฺทนา จ, \\ อเนกธาตูสุ สรา ปติฏฺฐิตา. \\ อญฺญาณมูลปฺปภวา ปชปฺปิตา, \\ สพฺพา มยา พฺยนฺติกตา สมูลิกา”ติ\csromansharedfootnote{aa29bc1bef980e12}{สํ 1. 183 ปิฏฺเฐ โถกํ วิสทิสํ; อุปริ 45, 180, 183 ปิฏฺเฐสุปิ.}.}}}

\prose{ตณฺหาเยตํ เววจนํ. ยถาห ภควา\csromansymbolfootnote{*}{สํ 2. 87 ปิฏฺเฐ.}รูเป ติสฺส อวิคตราคสฺส อวิคตจฺฉนฺทสฺส อวิคตเปมสฺส อวิคตปิปาสสฺส อวิคตปริฬาหสฺส. เอวํ เวทนาย สญฺญาย สงฺขาเรสุ วิญฺญาเณ อวิคตราคสฺส อวิคตจฺฉนฺทสฺส อวิคตเปมสฺส อวิคตปิปาสสฺส อวิคตปริฬาหสฺส สพฺพํ สุตฺตํ วิตฺถาเรตพฺพํ. ตณฺหาเยตํ เววจนํ เอวํ ยุชฺชติ.}

\proseitem{21}{สพฺโพ ทุกฺขูปจาโร กามตณฺหา\-สงฺขาร\-มูลโก, น ปน ยุชฺชติ สพฺโพ นิพฺพิทูปจาโร กามตณฺหา\-ปริกฺขาร\-มูลโก. อิมาย ยุตฺติยา อญฺญมญฺเญหิ การเณหิ คเวสิตพฺพํ.}

\prose{ยถา หิ\csromansharedfootnote{d04b73c477052725}{ยถาห (สี)} ภควา ราคจริตสฺส ปุคฺคลสฺส อสุภํ เทสยติ, โทสจริตสฺส ภควา ปุคฺคลสฺส เมตฺตํ เทสยติ. โมหจริตสฺส ภควา ปุคฺคลสฺส ปฏิจฺจ\-สมุปฺปาทํ เทสยติ. ยทิ หิ ภควา ราคจริตสฺส ปุคฺคลสฺส เมตฺตํ เจโตวิมุตฺติํ เทเสยฺย. สุขํ วา ปฏิปทํ ทนฺธาภิญฺญํ สุขํ วา ปฏิปทํ ขิปฺปาภิญฺญํ วิปสฺสนา\-ปุพฺพงฺคมํ วา ปหานํ เทเสยฺย, น ยุชฺชติ เทสนา. เอวํ ยํ กิญฺจิ ราคสฺส อนุโลม\-ปฺปหานํ โทสสฺส อนุโลม\-ปฺปหานํ โมหสฺส อนุโลม\-ปฺปหานํ. สพฺพํ ตํ วิจเยน หาเรน วิจินิตฺวา ยุตฺติหาเรน โยเชตพฺพํ. ยาวติกา ญาณสฺส ภูมิ.}

\prose{เมตฺตาวิหาริสฺส สโต พฺยาปาโท จิตฺตํ ปริยาทาย ฐสฺสตีติ น ยุชฺชติ เทสนา, พฺยาปาโท ปหานํ อพฺภตฺถํ คจฺฉตีติ ยุชฺชติ เทสนา. กรุณา\-วิหาริสฺส สโต วิเหสา จิตฺตํ ปริยาทาย ฐสฺสตีติ น ยุชฺชติ เทสนา, วิเหสา ปหานํ อพฺภตฺถํ คจฺฉตีติ ยุชฺชติ เทสนา. มุทิตา\-วิหาริสฺส สโต อรติ จิตฺตํ ปริยาทาย ฐสฺสตีติ น ยุชฺชติ เทสนา, อรติ ปหานํ อพฺภตฺถํ คจฺฉตีติ ยุชฺชติ เทสนา. อุเปกฺขา\-วิหาริสฺส สโต ราโค จิตฺตํ ปริยาทาย ฐสฺสตีติ น ยุชฺชติ เทสนา, ราโค ปหานํ อพฺภตฺถํ คจฺฉตีติ ยุชฺชติ เทสนา. อนิมิตฺต\-วิหาริสฺส \csromanfolio{23}สโต นิมิตฺตานุสารี เตน เตเนว วิญฺญาณํ ปวตฺตตีติ น ยุชฺชติ เทสนา, นิมิตฺตํ ปหานํ อพฺภตฺถํ คจฺฉตีติ ยุชฺชติ เทสนา. อสฺมีติ วิคตํ อยมหมสฺมีติ น สมนุปสฺสามิ. อถ จ ปน เม กิสฺมีติ กถสฺมีติ วิจิกิจฺฉา กถํกถา\-สลฺลํ จิตฺตํ ปริยาทาย ฐสฺสตีติ น ยุชฺชตี เทสนา, วิจิกิจฺฉา กถํกถา\-สลฺลํ ปหานํ อพฺภตฺถํ คจฺฉตีติ ยุชฺชติ เทสนา.}

\prose{ยถา วา ปน ปฐมํ ฌานํ สมาปนฺนสฺส สโต กามราค\-พฺยาปาทา วิเสสาย สํวตฺตนฺตีติ น ยุชฺชติ เทสนา, หานาย สํวตฺตนฺตีติ ยุชฺชติ เทสนา. วิตกฺก\-สหคตา วา สญฺญา\-มนสิการา หานาย สํวตฺตนฺตีติ น ยุชฺชติ เทสนา, วิเสสาย สํวตฺตนฺตีติ ยุชฺชติ เทสนา.  ทุติยํ ฌานํ สมาปนฺนสฺส สโต วิตกฺกวิจาร\-สหคตา สญฺญา\-มนสิการา วิเสสาย สํวตฺตนฺตีติ น ยุชฺชติ เทสนา, หานาย สํวตฺตนฺตีติ ยุชฺชติ เทสนา. อุเปกฺขา\-สุขสหคตา วา สญฺญา\-มนสิการา หานาย สํวตฺตนฺตีติ น ยุชฺชติ เทสนา, วิเสสาย สํวตฺตนฺตีติ ยุชฺชติ เทสนา.  ตติยํ ฌานํ สมาปนฺนสฺส สโต ปีติสุขสหคตา สญฺญา\-มนสิการา วิเสสาย สํวตฺตนฺตีติ น ยุชฺชติ เทสนา, หานาย สํวตฺตนฺตีติ ยุชฺชติ เทสนา, อุเปกฺขา\-สติ\-ปาริสุทฺธิสหคตา วา สญฺญา\-มนสิการา หานาย สํวตฺตนฺตีติ น ยุชฺชติ เทสนา, วิเสสาย สํวตฺตนฺตีติ ยุชฺชติ เทสนา.  จตุตฺถํ ฌานํ สมาปนฺนสฺส สโต อุเปกฺขา\-สหคตา สญฺญา\-มนสิการา วิเสสาย สํวตฺตนฺตีติ น ยุชฺชติ เทสนา, หานาย สํวตฺตนฺตีติ ยุชฺชติ เทสนา.  อากาสา\-นญฺจา\-ยตนสหคตา วา สญฺญา\-มนสิการา หานาย สํวตฺตนฺตีติ น ยุชฺชติ เทสนา, วิเสสาย สํวตฺตนฺตีติ ยุชฺชติ เทสนา.}

\prose{อากาสา\-นญฺจา\-ยตนํ สมาปนฺนสฺส สโต รูปสหคตา สญฺญา\-มนสิการา วิเสสาย สํวตฺตนฺตีติ น ยุชฺชติ เทสนา, หานาย สํวตฺตนฺตีติ ยุชฺชติ เทสนา. วิญฺญาณญฺจา\-ยตนสหคตา วา สญฺญา\-มนสิการา หานาย สํวตฺตนฺตีติ น ยุชฺชติ เทสนา, วิเสสาย สํวตฺตนฺตีติ ยุชฺชติ เทสนา.  วิญฺญาณญฺจา\-ยตนํ สมาปนฺนสฺส สโต อากาสา\-นญฺจา\-ยตนสหคตา สญฺญา\-มนสิการา วิเสสาย สํวตฺตนฺตีติ น ยุชฺชติ เทสนา, หานาย สํวตฺตนฺตีติ ยุชฺชติ เทสนา. อากิญฺจญฺญายตน\-สหคตา วา สญฺญา\-มนสิการา หานาย สํวตฺตนฺตีติ น ยุชฺชติ เทสนา, วิเสสาย สํวตฺตนฺตีติ ยุชฺชติ เทสนา. \csromanfolio{24}อากิญฺจญฺญา\-ยตนํ สมาปนฺนสฺส สโต วิญฺญาณญฺจา\-ยตนสหคตา สญฺญา\-มนสิการา วิเสสาย สํวตฺตนฺตีติ น ยุชฺชติ เทสนา, หานาย สํวตฺตนฺตีติ ยุชฺชติ เทสนา.  เนว สญฺญานาสญฺญายตนสหคตา วา สญฺญา\-มนสิการา หานาย สํวตฺตนฺตีติ น ยุชฺชติ เทสนา, วิเสสาย สํวตฺตนฺตีติ ยุชฺชติ เทสนา. เนว\-สญฺญา\-นา\-สญฺญา\-ยตนํ สมาปนฺนสฺส สโต สญฺญูปจารา วิเสสาย สํวตฺตนฺตีติ น ยุชฺชติ เทสนา, หานาย สํวตฺตนฺตีติ ยุชฺชติ เทสนา.  สญฺญาเวทยิตนิโรธ\-สหคตา วา สญฺญา\-มนสิการา หานาย สํวตฺตนฺตีติ น ยุชฺชติ เทสนา, วิเสสาย สํวตฺตนฺตีติ ยุชฺชติ เทสนา.  กลฺลตา\-ปริจิตํ จิตฺตํ น จ อภินีหารํ ขมตีติ น ยุชฺชติ เทสนา, กลฺลตา\-ปริจิตํ จิตฺตํ อถ จ อภินีหารํ ขมตีติ ยุชฺชติ เทสนา.}

\prose{เอวํ สพฺเพ นวสุตฺตนฺตา ยถาธมฺมํ ยถาวินยํ ยถา\-สตฺถุสาสนํ สพฺพโต วิจเยน หาเรน วิจินิตฺวา ยุตฺติหาเรน โยเชตพฺพาติ. เตนาห อายสฺมา มหากจฺจายโน “สพฺเพสํ หารานํ ยา ภูมิโย จ โคจโร เตสนฺ”ติ.}

\vspace{\tipitakaparskip}
\csromancenter{นิยุตฺโต ยุตฺติหาโร.}
\csromanmatikaanchor{mtk.8}
\csromantocmarkref{section}{4. ปทฏฺฐานหารวิภงฺค}{mtk.8}
\bookmark[dest={mtk.8},level=1]{4. ปทฏฺฐานหารวิภงฺค}
\csromanheader{1}{4. \textbf{ปทฏฺฐาน}\-\textbf{หาร}\-\textbf{วิภงฺค}}
\proseitem{22}{ตตฺถ กตโม ปทฏฺฐาโน หาโร, “ธมฺมํ เทเสติ ชิโน”ติ, อยํ ปทฏฺฐาโน หาโร. กิํ เทเสติ, สพฺพธมฺมยาถาวอสมฺปฏิเวธลกฺขณา อวิชฺชา, ตสฺสา วิปลฺลาสา ปทฏฺฐานํ. อชฺโฌสาน\-ลกฺขณา ตณฺหา, ตสฺสา ปิยรูปํ สาตรูปํ ปทฏฺฐานํ. ปตฺถน\-ลกฺขโณ โลโภ, ตสฺส อทินฺนาทานํ ปทฏฺฐานํ. วณฺณ\-สณฺฐาน\-พฺยญฺชน\-คฺคหณ\-ลกฺขณา สุภสญฺญา, ตสฺสา อินฺทฺริยา สํวโร ปทฐานํ. สาสวผสฺส- อุปคมน\-ลกฺขณา สุขสญฺญา, ตสฺสา อสฺสาโท ปทฏฺฐานํ. สงฺขต\-ลกฺขณานํ ธมฺมานํ อสม\-นุปสฺสน\-ลกฺขณา นิจฺจสญฺญา, ตสฺสา วิญฺญาณํ ปทฏฺฐานํ, อนิจฺจสญฺญาทุกฺขสญฺญา- อสม\-นุปสฺสน\-ลกฺขณา อตฺตสญฺญา, ตสฺสา นามกาโย ปทฏฺฐานํ.  สพฺพธมฺมสมฺปฏิเวธลกฺขณา วิชฺชา, ตสฺสา สพฺพํ เนยฺยํ ปทฏฺฐานํ. จิตฺตวิกฺเขป\-ปฏิสํหรณ\-ลกฺขโณ สมโถ, ตสฺส อสุภา ปทฏฺฐานํ. อิจฺฉาวจร\-ปฏิสํหรณ\-ลกฺขโณ อโลโภ, ตสฺส อทินฺนาทานา เวรมณี\csromansharedfootnote{f96664fc1fbeb341}{เวรมณิ (ก)} ปทฏฺฐานํ \csromanfolio{25}อพฺยาปชฺช\-ลกฺขโณ อโทโส, ตสฺส ปาณาติปาตา เวรมณี ปทฏฺฐานํ. วตฺถุ- อวิปฺปฏิปตฺติลกฺขโณ\csromansharedfootnote{340269ab8d479b87}{วตฺถุอวิปฺปฏิปาทานลกฺขโณ (สี, ก)} อโมโห, ตสฺส สมฺมาปฏิปตฺติ ปทฏฺฐานํ. วินีลก\-วิปุพฺพก\-คหณ\-ลกฺขณา อสุภสญฺญา, ตสฺสา นิพฺพิทา ปทฏฺฐานํ. สาสว\-ผสฺส\-ปริชานน\-ลกฺขณา ทุกฺขสญฺญา, ตสฺสา เวทนา ปทฏฺฐานํ. สงฺขต\-ลกฺขณานํ ธมฺมานํ สมนุปสฺสน\-ลกฺขณา อนิจฺจสญฺญา, ตสฺสา อุปฺปาทวยา ปทฏฺฐานํ. สพฺพธมฺม\-อภินิเวส\-ลกฺขณา อนตฺตสญฺญา, ตสฺสา ธมฺมสญฺญา ปทฏฺฐานํ.}

\prose{ปญฺจ กามคุณา กามราคสฺส ปทฏฺฐานํ, ปญฺจินฺทฺริยานิ รูปีนิ รูปราคสฺส ปทฏฺฐานํ, ฉฏฺฐายตนํ ภวราคสฺส ปทฏฺฐานํ, นิพฺพตฺต\-ภวา\-นุปสฺสิตา ปญฺจนฺนํ อุปาทาน\-กฺขนฺธานํ ปทฏฺฐานํ, ปุพฺเพ\-นิวาสา\-นุสฺสติญาณทสฺสนสฺส ปทฏฺฐานํ. โอกปฺปน\-ลกฺขณา สทฺธา อธิมุตฺติ\-ปจฺจุปฏฺฐานา จ, อนาวิล\-ลกฺขโณ ปสาโท สมฺปสีทน\-ปจฺจุปฏฺฐาโน จ. อภิปตฺถิยนลกฺขณา สทฺธา, ตสฺสา อเวจฺจปสาโท ปทฏฺฐานํ. อนาวิล\-ลกฺขโณ ปสาโท, ตสฺส สทฺธา ปทฏฺฐานํ. อารมฺภ\-ลกฺขณํ วีริยํ, ตสฺส สมฺมปฺปธานํ ปทฏฺฐานํ. อปิลาปน\-ลกฺขณา สติ, ตสฺสา สติปฏฺฐานํ ปทฏฺฐานํ. เอกคฺคลกฺขโณ สมาธิ, ตสฺส ฌานานิ ปทฏฺฐานํ. ปชานน\-ลกฺขณา ปญฺญา, ตสฺสา สจฺจานิ ปทฏฺฐานํ.}

\prose{อปโร นโย, อสฺสาท\-มนสิการ\-ลกฺขโณ อโยนิโส\-มนสิกาโร, ตสฺส อวิชฺชา ปทฏฺฐานํ. สจฺจ\-สมฺโม\-หน\-ลกฺขณา อวิชฺชา, สา สงฺขารานํ ปทฏฺฐานํ. ปุนพฺภววิโรหณลกฺขณา สงฺขารา, เต\csromansharedfootnote{672b3cfb0804b7fc}{ตํ (ก)} วิญฺญาณสฺส ปทฏฺฐานํ. โอปปจฺจยิกนิพฺพตฺติลกฺขณํ วิญฺญาณํ, ตํ นามรูปสฺส ปทฏฺฐานํ. นามกาย\-รูปกาย\-สงฺฆาต\-ลกฺขณํ นามรูปํ, ตํ ฉฬายตนสฺส ปทฏฺฐานํ. อินฺทฺริย\-ววตฺถาน\-ลกฺขณํ ฉฬายตนํ, ตํ ผสฺสสฺส ปทฏฺฐานํ. จกฺขุ\-รูป\-วิญฺญาณ\-สนฺนิปาต\-ลกฺขโณ ผสฺโส, โส เวทนาย ปทฏฺฐานํ. อิฏฺฐานิฏฺฐ\-อนุภวน\-ลกฺขณา เวทนา, สา ตณฺหาย ปทฏฺฐานํ. อชฺโฌสาน\-ลกฺขณา ตณฺหา, สา อุปาทานสฺส ปทฏฺฐานํ. โอปปจฺจยิกํ อุปาทานํ, ตํ ภวสฺส ปทฏฺฐานํ. นามกาย\-รูปกาย\-สมฺภวน\-ลกฺขโณ ภโว, โส ชาติยา ปทฏฺฐานํ. ขนฺธ\-ปาตุ\-ภวน\-ลกฺขณา ชาติ, สา ชราย ปทฏฺฐานํ. อุปธิ\-ปริปาก\-ลกฺขณา ชรา, สา มรณสฺส ปทฏฺฐานํ. ชีวิตินฺทฺริยุ\-ปจฺเฉท\-ลกฺขณํ มรณํ, ตํ โสกสฺส ปทฏฺฐานํ. อุสฺสุกฺกการโก โสโก, โส ปริเทวสฺส \csromanfolio{26}ปทฏฺฐานํ. ลาลปฺปการโก ปริเทโว, โส ทุกฺขสฺส ปทฏฺฐานํ. กายสํปีฬนํ ทุกฺขํ, ตํ โทมนสฺสสฺส ปทฏฺฐานํ. จิตฺตสํปีฬนํ โทมนสฺสํ, ตํ อุปายาสสฺส ปทฏฺฐานํ. โอทหนการโก อุปายาโส, โส ภวสฺส ปทฏฺฐานํ. อิมานิ ภวงฺคานิ ยทา สมคฺคานิ นิพฺพตฺตานิ ภวนฺติ โส ภโว, ตํ สํสารสฺส ปทฏฺฐานํ. นิยฺยานิก\-ลกฺขโณ มคฺโค, โส นิโรธสฺส ปทฏฺฐานํ.}

\prose{ติตฺถญฺญุตา ปีตญฺญุตาย ปทฏฺฐานํ, ปีตญฺญุตา ปตฺตญฺญุตาย\csromansharedfootnote{a47024dc9d015804}{มตฺตญฺญุตาย (สี, ก)} ปทฏฺฐานํ, ปตฺตญฺญุตา อตฺตญฺญุตาย ปทฏฺฐานํ, อตฺตญฺญุตา ปุพฺเพ\-กตปุญฺญตาย ปทฏฺฐานํ, ปุพฺเพ\-กตปุญฺญตา ปติรูป\-เทส\-วาสสฺส ปทฏฺฐานํ, ปติรูป\-เทส\-วาโส สปฺปุริสู\-ปนิสฺสยสฺส ปทฏฺฐานํ, สปฺปุริสู\-ปนิสฺสโย อตฺตสมฺมาปณิธา\-นสฺส ปทฏฺฐานํ, อตฺต\-สมฺมา\-ปณิธานํ สีลานํ ปทฏฺฐานํ, สีลานิ อวิปฺปฏิสารสฺส ปทฏฺฐานํ, อวิปฺปฏิสาโร ปาโมชฺชสฺส ปทฏฺฐานํ, ปาโมชฺชํ ปีติยา ปทฏฺฐานํ, ปีติ ปสฺสทฺธิยา ปทฏฺฐานํ, ปสฺสทฺธิ สุขสฺส ปทฏฺฐานํ, สุขํ สมาธิสฺส ปทฏฺฐานํ, สมาธิ ยถา\-ภูต\-ญาณ\-ทสฺสนสฺส ปทฏฺฐานํ, ยถา\-ภูต\-ญาณ\-ทสฺสนํ นิพฺพิทาย ปทฏฺฐานํ, นิพฺพิทา วิราคสฺส ปทฏฺฐานํ, วิราโค วิมุตฺติยา ปทฏฺฐานํ. วิมุตฺติ วิมุตฺติ\-ญาณ\-ทสฺสนสฺส ปทฏฺฐานํ. เอวํ โย โกจิ อุปนิสฺสโย โย โกจิ ปจฺจโย, สพฺโพ โส ปทฏฺฐานํ. เตนาห อายสฺมา มหากจฺจายโน “ธมฺมํ เทเสติ ชิโน”ติ.}

\vspace{\tipitakaparskip}
\csromancenter{นิยุตฺโต ปทฏฺฐาโน หาโร.}
\csromanmatikaanchor{mtk.9}
\csromantocmarkref{section}{5. ลกฺขณหารวิภงฺค}{mtk.9}
\bookmark[dest={mtk.9},level=1]{5. ลกฺขณหารวิภงฺค}
\csromanheader{1}{5. \textbf{ลกฺขณ}\-\textbf{หาร}\-\textbf{วิภงฺค}}
\proseitem{23}{ตตฺถ กตโม ลกฺขโณ หาโร. “วุตฺตมฺหิ เอกธมฺเม”ติ, อยํ ลกฺขโณ หาโร, กิํ ลกฺขยติ, เย ธมฺมา เอกลกฺขณา, เตสํ ธมฺมานํ เอกสฺมิํ ธมฺเม วุตฺเต อวสิฏฺฐา ธมฺมา วุตฺตา ภวนฺติ. ยถา กิํ ภเว. ยถาห ภควา\csromandash{}}

\prose{จกฺขุํ ภิกฺขเว อนวฏฺฐิตํ อิตฺตรํ ปริตฺตํ ปภงฺคุ ปรโต ทุกฺขํ พฺยสนํ จลนํ\csromansharedfootnote{f3f1615c5d9f3ec8}{จลํ (สี)} กุกฺกุฬํ สงฺขารํ\csromansharedfootnote{4b183d3389b6ac6d}{สสงฺขารํ (ก)} วธกํ อมิตฺตมชฺเฌ. อิมสฺมิํ จกฺขุสฺมิํ วุตฺเต อวสิฏฺฐานิ อชฺฌตฺติกานิ อายตนานิ วุตฺตานิ ภวนฺติ, เกน การเณน, สพฺพานิ หิ ฉ อชฺฌตฺติกานิ อายตนานิ วธกฏฺเฐน เอกลกฺขณานิ. ยถา จาห ภควา\csromandash{}}

\prose{\csromanfolio{27}อตีเต ราธ รูเป อนเปกฺโข โหติ, อนาคตํ รูปํ มา อภินนฺทิ\csromansharedfootnote{c5b328cf4f4eb561}{อภินนฺท (ก)}, ปจฺจุปฺปนฺนสฺส รูปสฺส นิพฺพิทาย วิราคาย นิโรธาย จาคาย ปฏินิสฺสคฺคาย ปฏิปชฺช. อิมสฺมิํ รูปกฺขนฺเธ วุตฺเต อวสิฏฺฐา ขนฺธา วุตฺตา ภวนฺติ, เกน การเณน, สพฺเพ หิ ปญฺจกฺขนฺธา ยมโก\-วาทสุตฺเต\csromansharedfootnote{a360aeb723c93f57}{สํ 2. 89 ปิฏฺเฐ.} วธกฏฺเฐน เอกลกฺขณา วุตฺตา. ยถา จาห ภควา\csromandash{}}

\setlength{\csromangathapagewidth}{0pt}
\setlength{\csromangathawakwidth}{0pt}
\csromangathameasuregroup{“เยสญฺจ สุสมารทฺธา, \\ อกิจฺจํ เต น เสวนฺติ,}{\csromangathabat{“เยสญฺจ สุสมารทฺธา,}{นิจฺจํ กายคตาสติ.} \\ \csromangathabat{อกิจฺจํ เต น เสวนฺติ,}{กิจฺเจ สาตจฺจการิโน”\textsuperscript{1}.}}
\csromangathasetleft{“เยสญฺจ สุสมารทฺธา, \\ อกิจฺจํ เต น เสวนฺติ,}
\csromangathagroup{\csromangathastanza{\csromangathabat{“เยสญฺจ สุสมารทฺธา,}{นิจฺจํ กายคตาสติ.} \\ \csromangathabat{อกิจฺจํ เต น เสวนฺติ,}{กิจฺเจ สาตจฺจการิโน”\csromansharedfootnote{148d87c82e5012aa}{ขุ 1. 55 ปิฏฺเฐ ธมฺมปเท.}.}}}

\prose{อิติ กายคตาย สติยา วุตฺตาย วุตฺตา ภวนฺติ เวทนาคตา สติ จิตฺตคตา ธมฺมคตา จ. ตถา ยํ กิญฺจิ ทิฏฺฐํ วา สุตํ วา มุตํ วาติ วุตฺเต วุตฺตํ ภวติ วิญฺญาตํ. ยถา จาห ภควา\csromandash{}}

\prose{\csromansymbolfootnote{*}{สํ 3. 125, 165 ปิฏฺเฐสุ โถกํ วิสทิสํ.}ตสฺมาติห ตฺวํ ภิกฺขุ กาเย กายานุปสฺสี วิหราหิ อาตาปี สมฺปชาโน สติมา วิเนยฺย โลเก อภิชฺฌา\-โทมนสฺสํ. “อาตาปี”ติ วีริยินฺทฺริยํ, “สมฺปชาโน”ติ ปญฺญินฺทฺริยํ, “สติมา”ติ สตินฺทฺริยํ, “วิเนยฺย โลเก อภิชฺฌาโทมนสฺสนฺ”ติ สมาธินฺทฺริยํ, เอวํ กาเย กายานุปสฺสิโน วิหรโต จตฺตาโร สติปฏฺฐานา ภาวนา\-ปาริปูริํ คจฺฉนฺติ. เกน การเณน, เอกลกฺขณตฺตา จตุนฺนํ อินฺทฺริยานํ.}

\proseitem{24}{จตูสุ สติปฏฺฐาเนสุ ภาวิยมาเนสุ จตฺตาโร สมฺมปฺปธานา ภาวนา\-ปาริปูริํ คจฺฉนฺติ, จตูสุ สมฺม\-ปฺปธาเนสุ ภาวิยมาเนสุ จตฺตาโร อิทฺธิปาทา ภาวิยมาเนสุ ปญฺจินฺทฺริยานิ ภาวนา\-ปาริปูริํ คจฺฉนฺติ, จตูสุ อิทฺธิปาเทสุ ภาวนา\-ปาริปูริํ คจฺฉนฺติ, ปญฺจสุ อินฺทฺริเยสุ ภาวิยมาเนสุ ปญฺจ พลานิ ภาวนา\-ปาริปูริํ คจฺฉนฺติ, ปญฺจสุ พเลสุ ภาวิยมาเนสุ สตฺต โพชฺฌงฺคา ภาวนา\-ปาริปูริํ คจฺฉนฺติ, สตฺตสุ โพชฺฌงฺเคสุ ภาวิยมาเนสุ อริโย อฏฺฐงฺคิโก มคฺโค ภาวนา\-ปาริปูริํ คจฺฉติ, สพฺเพว\csromansharedfootnote{b4688ab7828edd5e}{สพฺเพ จ (สี, ก)} โพธงฺคมา ธมฺมา โพธิปกฺขิยา ภาวนา\-ปาริปูริํ คจฺฉนฺติ. เกน การเณน, สพฺเพ หิ โพธงฺคมา โพธิปกฺขิยา เนยฺยา\-นิก\-ลกฺขเณน เอกลกฺขณา, เต เอกลกฺขณตฺตา ภาวนา\-ปาริปูริํ คจฺฉนฺติ.}

\prose{เอวํ อกุสลาปิ ธมฺมา เอกลกฺขณตฺตา ปหานํ อพฺภตฺถํ คจฺฉนฺติ. จตูสุ สติปฏฺฐาเนสุ ภาวิยมาเนสุ วิปลฺลาสา ปหียนฺติ, \csromanfolio{28}อาหารา จสฺส ปริญฺญํ คจฺฉนฺติ, อุปาทาเนหิ อนุปาทาโน ภวติ, โยเคหิ จ วิสํยุตฺโต ภวติ, คนฺเถหิ จ วิปฺปยุตฺโต ภวติ, อาสเวหิ จ อนาสโว ภวติ, โอเฆหิ จ นิตฺถิณฺโณ ภวติ, สลฺเลหิ จ วิสลฺโล ภวติ, วิญฺญาณ\-ฏฺฐิติโย จสฺส ปริญฺญํ คจฺฉนฺติ, อคติคมเนหิ น อคติํ คจฺฉติ, เอวํ อกุสลาปิ ธมฺมา เอกลกฺขณตฺตา ปหานํ อพฺภตฺถํ คจฺฉนฺติ.}

\prose{ยตฺถ วา ปน รูปินฺทฺริยํ เทสิตํ, เทสิตา ตตฺเถว รูปธาตุ รูปกฺขนฺโธ รูปญฺจายตนํ. ยตฺถ วา ปน สุขา เวทนา เทสิตา, เทสิตํ ตตฺถ สุขินฺทฺริยํ โสมนสฺสิ\-นฺทฺริยํ ทุกฺขสมุทโย จ อริยสจฺจํ. ยตฺถ วา ปน ธุกฺขา เวทนา เทสิตา, เทสิตํ ตตฺถ ทุกฺขินฺทฺริยํ โทมนสฺสิ\-นฺทฺริยํ ทุกฺขญฺจ อริยสจฺจํ. ยตฺถ วา ปน อทุกฺขมสุขา เวทนา เทสิตา, เทสิตํ ตตฺถ อุเปกฺขินฺทฺริยํ สพฺโพ จ ปฏิจฺจ\-สมุปฺปาโท. เกน การเณน, อทุกฺขม\-สุขาย หิ เวทนาย อวิชฺชา อนุเสติ. อวิชฺชาปจฺจยา สงฺขารา, สงฺขาร\-ปจฺจยา วิญฺญาณํ, วิญฺญาณปจฺจยา นามรูปํ, นามรูป\-ปจฺจยา จฬายตนํ, สฬายตน\-ปจฺจยา ผสฺโส, ผสฺสปจฺจยา เวทนา, เวทนาปจฺจยา ตณฺหา, ตณฺหาปจฺจยา อุปาทานํ, อุปาทานปจฺจยา ภโว, ภวปจฺจยา ชาติ, ชาติปจฺจยา ชรามรณํ โสก\-ปริเทว\-ทุกฺข\-โทมนสฺสุ\-ปายาสา สมฺภวนฺติ, เอวเมตสฺส เกวลสฺส ทุกฺข\-กฺขนฺธสฺส สมุทโย โหติ. โส จ สราค\-สโทส\-สโมห\-สํกิเลส\-ปกฺเขน หาตพฺโพ, วีตราค\-วีตโทส\-วีตโมห\-อริยธมฺเมหิ หาตพฺโพ.}

\prose{เอวํ เย ธมฺมา เอกลกฺขณา กิจฺจโต จ ลกฺขณโต จ สามญฺญโต จ จุตูปปาตโต จ, เตสํ ธมฺมานํ เอกสฺมิํ ธมฺเม วุตฺเต อวสิฏฺฐา ธมฺมา วุตฺตา ภวนฺติ. เตนาห อายสฺมา มหากจฺจายโน “วุตฺตมฺหิ เอกธมฺเม”ติ.}

\vspace{\tipitakaparskip}
\csromancenter{นิยุตฺโต ลกฺขโณ หาโร.}
\csromanmatikaanchor{mtk.10}
\csromantocmarkref{section}{6. จตุพฺยูหหารวิภงฺค}{mtk.10}
\bookmark[dest={mtk.10},level=1]{6. จตุพฺยูหหารวิภงฺค}
\csromanheader{1}{6. \textbf{จตุพฺยูห}\-\textbf{หาร}\-\textbf{วิภงฺค}}
\proseitem{25}{ตตฺถ กตโม จตุพฺยูโห หาโร. “เนรุ\-ตฺตม\-ธิปฺปาโย”ติ อยํ. พฺยญฺชเนน สุตฺตสฺส เนรุตฺตญฺจ อธิปฺปาโย จ นิทานญฺจ ปุพฺพาปร\-สนฺธิ จ \csromanfolio{29}คเวสิตพฺโพ. ตตฺถ กตมํ เนรุตฺตํ, ยา นิรุตฺติ\-ปท\-สํหิตา, ยํ ธมฺมานํ นามโส ญาณํ. ยทา หิ ภิกฺขุ อตฺถสฺส จ นามํ ชานาติ, ธมฺมสฺส จ นามํ ชานาติ, ตถา ตถา นํ อภินิโรเปติ. อยญฺจ วุจฺจติ อตฺถกุสโล ธมฺมกุสโล พฺยญฺชนกุสโล นิรุตฺติกุสโล ปุพฺพาปร\-กุสโล เทสนากุสโล อตีตา\-ธิวจน\-กุสโล อนาคตา\-ธิวจน\-กุสโล ปจฺจุปฺปนฺนา\-ธิวจน\-กุสโล อิตฺถา\-ธิวจน\-กุสโล ปุริสา\-ธิวจน\-กุสโล นปุํสกา\-ธิวจน\-กุสโล เอกา\-ธิวจน\-กุสโล อเนกา\-ธิวจน\-กุสโล, เอวํ สพฺพานิ กาตพฺพานิ ชนปทนิรุตฺตานิ สพฺพา จ ชนปท\-นิรุตฺติโย. อยํ นิรุตฺติ\-ปท\-สํหิตา. (1)}

\proseitem{26}{ตตฺถ กตโม อธิปฺปาโย.}

\prose{\csromansymbolfootnote{+}{เหฏฺฐา 7 ปิฏฺเฐ; อุปริ 38, 197 ปิฏฺเฐสุปิ.}“ธมฺโม หเว รกฺขติ ธมฺมจาริํ,}

\prose{ฉตฺตํ มหนฺตํ ยถ วสฺสกาเล.}

\setlength{\csromangathapagewidth}{0pt}
\setlength{\csromangathawakwidth}{0pt}
\csromangathameasuregroup{}{เอสานิสํโส ธมฺเม สุจิณฺเณ, \\ น ทุคฺคติํ คจฺฉติ ธมฺมจารี”ติ.}
\csromangathameasurewak{เอสานิสํโส ธมฺเม สุจิณฺเณ, \\ น ทุคฺคติํ คจฺฉติ ธมฺมจารี”ติ.}
\setlength{\csromangathaleftcol}{0pt}
\csromangathagroup{\csromangathastanza{\csromangathawak{เอสานิสํโส ธมฺเม สุจิณฺเณ, \\ น ทุคฺคติํ คจฺฉติ ธมฺมจารี”ติ.}}}

\prose{อิธ ภควโต โก อธิปฺปาโย, เย อปาเยหิ ปริมุจฺจิตุกามา ภวิสฺสนฺติ, เต ธมฺมจาริโน ภวิสฺสนฺตีติ อยํ เอตฺถ ภควโต อธิปฺปาโย.}

\prose{\csromansymbolfootnote{*}{ขุ 2. 327 ปิฏฺเฐ ปสฺสิตพฺพํ; อุปริ 108 ปิฏฺเฐปิ.}“โจโร ยถา สนฺธิมุเข คหีโต,}

\prose{สกมฺมุนา หญฺญติ\csromansharedfootnote{c087ac6190faca7f}{หญฺญเต (สี)} พชฺฌเต จ.}

\setlength{\csromangathapagewidth}{0pt}
\setlength{\csromangathawakwidth}{0pt}
\csromangathameasuregroup{}{เอวํ อยํ เปจฺจ ปชา ปรตฺถ, \\ สกมฺมุนา หญฺญติ\textsuperscript{1} พชฺฌเต จา”ติ.}
\csromangathameasurewak{เอวํ อยํ เปจฺจ ปชา ปรตฺถ, \\ สกมฺมุนา หญฺญติ\textsuperscript{1} พชฺฌเต จา”ติ.}
\setlength{\csromangathaleftcol}{0pt}
\csromangathagroup{\csromangathastanza{\csromangathawak{เอวํ อยํ เปจฺจ ปชา ปรตฺถ, \\ สกมฺมุนา หญฺญติ\csromansharedfootnote{c087ac6190faca7f}{หญฺญเต (สี)} พชฺฌเต จา”ติ.}}}

\prose{อิธ ภควโต โก อธิปฺปาโย, สญฺเจตนิกานํ กตานํ กมฺมานํ อุปจิตานํ ทุกฺข\-เวทนียานํ อนิฏฺฐํ อสาตํ วิปากํ ปจฺจนุภวิสฺสตีติ อยํ เอตฺถ ภควโต อธิปฺปาโย.}

\setlength{\csromangathapagewidth}{0pt}
\setlength{\csromangathawakwidth}{0pt}
\csromangathameasuregroup{“สุขกามานิ ภูตานิ, \\ อตฺตโน สุขเมสาโน,}{\csromangathabat{“สุขกามานิ ภูตานิ,}{โย ทณฺเฑน วิหิํสติ.} \\ \csromangathabat{อตฺตโน สุขเมสาโน,}{เปจฺจ โส น ลภเต สุขนฺ”ติ\textsuperscript{1}.}}
\csromangathasetleft{“สุขกามานิ ภูตานิ, \\ อตฺตโน สุขเมสาโน,}
\csromangathagroup{\csromangathastanza{\csromangathabat{“สุขกามานิ ภูตานิ,}{โย ทณฺเฑน วิหิํสติ.} \\ \csromangathabat{อตฺตโน สุขเมสาโน,}{เปจฺจ โส น ลภเต สุขนฺ”ติ\csromansharedfootnote{dfcbd237dfa78591}{ม 3. 89; ขุ 1. 33 ปิฏฺเฐสุ; อุปริ 108, 113 ปิฏฺเฐสุปิ. 3. ปาปกํ กมฺมํ (ก)}.}}}

\prose{อิธ ภควโต โก อธิปฺปาโย, เย สุเขน อตฺถิกา ภวิสฺสนฺติ, เต ปาปกมฺมํ\csromansharedfootnote{94641755c255831e}{ปาปกํ กมฺมํ (ก)} น กริสฺสนฺตีติ อยํ เอตฺถ ภควโต อธิปฺปาโย.}

\setlength{\csromangathapagewidth}{0pt}
\setlength{\csromangathawakwidth}{0pt}
\csromangathameasuregroup{“มิทฺธี ยทา โหติ มหคฺฆโส จ,}{\csromangathabat{“มิทฺธี ยทา โหติ มหคฺฆโส จ,}{นิทฺทายิตา สมฺปริวตฺตสายี.} \\ มหาวราโหว นิวาปปุฏฺโฐ. \\ ปุนปฺปุนํ คพฺภมุเปติ มนฺโท”ติ\textsuperscript{1}.}
\csromangathasetleft{“มิทฺธี ยทา โหติ มหคฺฆโส จ,}
\csromangathagroup{\csromangathastanza{\csromanfolio{30}\csromangathabat{“มิทฺธี ยทา โหติ มหคฺฆโส จ,}{นิทฺทายิตา สมฺปริวตฺต\-สายี.} \\ มหาวราโหว นิวาปปุฏฺโฐ. \\ ปุนปฺปุนํ คพฺภมุเปติ มนฺโท”ติ\csromansharedfootnote{7b6e1321cab8e566}{ขุ 1. 59 ปิฏฺเฐ ธมฺมปเท; อุปริ 108 ปิฏฺเฐปิ.}.}}

\prose{อิธ ภควโต โก อธิปฺปาโย, เย ชรามรเณน อฏฺฏิยิตุกามา ภวิสฺสนฺติ, เต ภวิสฺสนฺติ โภชเน มตฺตญฺญุโน อินฺทฺริเยสุ คุตฺตทฺวารา ปุพฺพรตฺตาปรรตฺตํ ชาคริยา\-นุโยคม\-นุยุตฺตา วิปสฺสกา กุสเลสุ ธมฺเมสุ สคารวา จ สพฺรหฺมจารีสุ เถเรสุ นเวสุ มชฺฌิเมสูติ อยํ เอตฺถ ภควโต อธิปฺปาโย.}

\setlength{\csromangathapagewidth}{0pt}
\setlength{\csromangathawakwidth}{0pt}
\csromangathameasuregroup{“อปฺปมาโท อมตปทํ\textsuperscript{1}, \\ อปฺปมตฺตา น มียนฺติ,}{\csromangathabat{“อปฺปมาโท อมตปทํ\textsuperscript{1},}{ปมาโท มจฺจุโน ปทํ.} \\ \csromangathabat{อปฺปมตฺตา น มียนฺติ,}{เย ปมตฺตา ยถา มตา”ติ.}}
\csromangathasetleft{“อปฺปมาโท อมตปทํ\textsuperscript{1}, \\ อปฺปมตฺตา น มียนฺติ,}
\csromangathagroup{\csromangathastanza{\csromangathabat{“อปฺปมาโท อมตปทํ\csromansharedfootnote{09e23b6cb84fb5ba}{อมตํ ปทํ (ก) ขุ 1. 16 ปิฏฺเฐ ปสฺสิตพฺพํ.},}{ปมาโท มจฺจุโน ปทํ.} \\ \csromangathabat{อปฺปมตฺตา น มียนฺติ,}{เย ปมตฺตา ยถา มตา”ติ.}}}

\prose{อิธ ภควโต โก อธิปฺปาโย, เย อมต\-ปริเยสนํ ปริเยสิตุกามา ภวิสฺสนฺติ, เต อปฺปมตฺตา วิหริสฺสนฺตีติ อยํ เอตฺถ ภควโต อธิปฺปาโย. อยํ อธิปฺปาโย. (2)}

\proseitem{27}{ตตฺถ กตมํ นิทานํ, ยถา โส ธนิโย โคปาลโก ภควนฺตํ อาห\csromandash{}}

\prose{\csromansymbolfootnote{*}{สํ 1. 109; ขุ 1. 284 ปิฏฺเฐสุ; อุปริ 205 ปิฏฺเฐปิ.}“นนฺทติ ปุตฺเตหิ ปุตฺติมา,}

\prose{โคมา\csromansharedfootnote{8033b55626965b7f}{โคมิโก (สี), โคปิโก (ก) ขุ 1. 284 ปิฏฺเฐ; สํ 1. 109 ปิฏฺเฐ จ ปสฺสิตพฺพํ.} โคหิ ตเถว นนฺทติ.}

\setlength{\csromangathapagewidth}{0pt}
\setlength{\csromangathawakwidth}{0pt}
\csromangathameasuregroup{}{อุปธี หิ นรสฺส นนฺทนา, \\ น หิ โส นนฺทติ โย นิรูปธี”ติ.}
\csromangathameasurewak{อุปธี หิ นรสฺส นนฺทนา, \\ น หิ โส นนฺทติ โย นิรูปธี”ติ.}
\setlength{\csromangathaleftcol}{0pt}
\csromangathagroup{\csromangathastanza{\csromangathawak{อุปธี หิ นรสฺส นนฺทนา, \\ น หิ โส นนฺทติ โย นิรูปธี”ติ.}}}

\prose{ภควา อาห\csromandash{}}

\prose{\csromansymbolmark{*}“โสจติ ปุตฺเตหิ ปุตฺติมา,}

\prose{โคปิโก\csromansharedfootnote{535941215439a30c}{โคมิโก (สี)} โคหิ ตเถว โสจติ.}

\setlength{\csromangathapagewidth}{0pt}
\setlength{\csromangathawakwidth}{0pt}
\csromangathameasuregroup{}{อุปธี หิ นรสฺส โสจนา, \\ น หิ โส โสจติ โย นิรูปธี”ติ.}
\csromangathameasurewak{อุปธี หิ นรสฺส โสจนา, \\ น หิ โส โสจติ โย นิรูปธี”ติ.}
\setlength{\csromangathaleftcol}{0pt}
\csromangathagroup{\csromangathastanza{\csromangathawak{อุปธี หิ นรสฺส โสจนา, \\ น หิ โส โสจติ โย นิรูปธี”ติ.}}}

\prose{\csromanfolio{31}อิมินา วตฺถุนา อิมินา นิทาเนน เอวํ ญายติ “อิธ ภควา พาหิรํ ปริคฺคหํ อุปธิํ อาหา”ติ. ยถา จ มาโร ปาปิมา คิชฺฌกูฏา ปพฺพตา ปุถุสิลํ ปาเตสิ, ภควา อาห\csromandash{}}

\setlength{\csromangathapagewidth}{0pt}
\setlength{\csromangathawakwidth}{0pt}
\csromangathameasuregroup{“สเจปิ เกวลํ สพฺพํ, \\ เนว สมฺมาวิมุตฺตานํ,}{\csromangathabat{“สเจปิ เกวลํ สพฺพํ,}{คิชฺฌกูฏํ จเลสฺสสิ.} \\ \csromangathabat{เนว สมฺมาวิมุตฺตานํ,}{พุทฺธานํ อตฺถิ อิญฺชิตํ\textsuperscript{1}.}}
\csromangathasetleft{“สเจปิ เกวลํ สพฺพํ, \\ เนว สมฺมาวิมุตฺตานํ,}
\csromangathagroup{\csromangathastanza{\csromangathabat{“สเจปิ เกวลํ สพฺพํ,}{คิชฺฌกูฏํ จเลสฺสสิ.} \\ \csromangathabat{เนว สมฺมา\-วิมุตฺตานํ,}{พุทฺธานํ อตฺถิ อิญฺชิตํ\csromansharedfootnote{9495dd7b19c8898c}{สํ 1. 111 ปิฏฺเฐ.}.}}}

\prose{\csromansymbolfootnote{*}{สํ 1. 108 ปิฏฺเฐ.}นภํ ผเลยฺย\-ปฺปถวี จเลยฺย,}

\prose{สพฺเพว ปาณา อุท สนฺตเสยฺยุํ.}

\setlength{\csromangathapagewidth}{0pt}
\setlength{\csromangathawakwidth}{0pt}
\csromangathameasuregroup{}{สลฺลมฺปิ เจ อุรสิ ปกปฺปเยยฺยุํ\textsuperscript{1}, \\ อุปธีสุ ตาณํ น กโรนฺติ พุทฺธา”ติ.}
\csromangathameasurewak{สลฺลมฺปิ เจ อุรสิ ปกปฺปเยยฺยุํ\textsuperscript{1}, \\ อุปธีสุ ตาณํ น กโรนฺติ พุทฺธา”ติ.}
\setlength{\csromangathaleftcol}{0pt}
\csromangathagroup{\csromangathastanza{\csromangathawak{สลฺลมฺปิ เจ อุรสิ ปกปฺปเยยฺยุํ\csromansharedfootnote{29b839185e2f71dc}{ปกมฺปเยยฺยุํ (สี) ** ขุ 1. 63; ขุ 5. 53 ปิฏฺเฐสุ; อุปริ 132, 185 ปิฏฺเฐสุปิ.}, \\ อุปธีสุ ตาณํ น กโรนฺติ พุทฺธา”ติ.}}}

\prose{อิมินา วตฺถุนา อิมินา นิทาเนน เอวํ ญายติ “อิธ ภควา กายํ อุปธิํ อาหา”ติ. ยถา จาห\csromandash{}}

\prose{** “น ตํ ทฬฺหํ พนฺธนมาหุ ธีรา,}

\prose{ยทายสํ ทารุชปพฺพชญฺจ\csromansharedfootnote{9cc122aff2fe5e58}{ทารุชํ พพฺพชญฺจ (สี)}.}

\setlength{\csromangathapagewidth}{0pt}
\setlength{\csromangathawakwidth}{0pt}
\csromangathameasuregroup{}{สารตฺตรตฺตา มณิกุณฺฑเลสุ, \\ ปุตฺเตสุ ทาเรสุ จ ยา อเปกฺขา”ติ.}
\csromangathameasurewak{สารตฺตรตฺตา มณิกุณฺฑเลสุ, \\ ปุตฺเตสุ ทาเรสุ จ ยา อเปกฺขา”ติ.}
\setlength{\csromangathaleftcol}{0pt}
\csromangathagroup{\csromangathastanza{\csromangathawak{สารตฺตรตฺตา มณิกุณฺฑเลสุ, \\ ปุตฺเตสุ ทาเรสุ จ ยา อเปกฺขา”ติ.}}}

\prose{อิมินา วตฺถุนา อิมินา นิทาเนน เอวํ ญายติ “อิธ ภควา พาหิเรสุ วตฺถูสุ ตณฺหํ อาหา”ติ. ยถา จาห\csromandash{}}

\prose{** “เอตํ ทฬฺหํ พนฺธนมาหุ ธีรา,}

\prose{โอหารินํ สิถิลํ ทุปฺปมุญฺจํ.}

\setlength{\csromangathapagewidth}{0pt}
\setlength{\csromangathawakwidth}{0pt}
\csromangathameasuregroup{}{เอตมฺปิ เฉตฺวาน ปริพฺพชนฺติ, \\ อนเปกฺขิโน กามสุขํ ปหายา”ติ.}
\csromangathameasurewak{เอตมฺปิ เฉตฺวาน ปริพฺพชนฺติ, \\ อนเปกฺขิโน กามสุขํ ปหายา”ติ.}
\setlength{\csromangathaleftcol}{0pt}
\csromangathagroup{\csromangathastanza{\csromangathawak{เอตมฺปิ เฉตฺวาน ปริพฺพชนฺติ, \\ อนเปกฺขิโน กามสุขํ ปหายา”ติ.}}}

\prose{อิมินา วตฺถุนา อิมินา นิทาเนน เอวํ ญายติ “อิธ ภควา พาหิร\-วตฺถุกาย ตณฺหาย ปหานํ อาหา”ติ. ยถา จาห\csromandash{}}

\setlength{\csromangathapagewidth}{0pt}
\setlength{\csromangathawakwidth}{0pt}
\csromangathameasuregroup{“อาตุรํ อสุจิํ ปูติํ, \\ ปคฺฆรนฺตํ ทิวา รตฺติํ,}{\csromangathabat{“อาตุรํ อสุจิํ ปูติํ,}{ทุคฺคนฺธํ เทหนิสฺสิตํ.} \\ \csromangathabat{ปคฺฆรนฺตํ ทิวา รตฺติํ,}{พาลานํ อภินนฺทิตนฺ”ติ.}}
\csromangathasetleft{“อาตุรํ อสุจิํ ปูติํ, \\ ปคฺฆรนฺตํ ทิวา รตฺติํ,}
\csromangathagroup{\csromangathastanza{\csromangathabat{“อาตุรํ อสุจิํ ปูติํ,}{ทุคฺคนฺธํ เทหนิสฺสิตํ.} \\ \csromangathabat{ปคฺฆรนฺตํ ทิวา รตฺติํ,}{พาลานํ อภินนฺทิตนฺ”ติ.}}}

\prose{อิมินา วตฺถุนา อิมินา นิทาเนน เอวํ ญายติ “อิธ ภควา อชฺฌตฺติก\-วตฺถุกาย ตณฺหาย ปหานํ อาหา”ติ. ยถา จาห\csromandash{}}

\setlength{\csromangathapagewidth}{0pt}
\setlength{\csromangathawakwidth}{0pt}
\csromangathameasuregroup{“อุจฺฉินฺท สิเนหมตฺตโน, \\ สนฺติมคฺคเมว พฺรูหย,}{\csromangathabat{“อุจฺฉินฺท สิเนหมตฺตโน,}{กุมุทํ สารทิกํว ปาณินา.} \\ \csromangathabat{สนฺติมคฺคเมว พฺรูหย,}{นิพฺพานํ สุคเตน เทสิตนฺ”ติ\textsuperscript{1}.}}
\csromangathasetleft{“อุจฺฉินฺท สิเนหมตฺตโน, \\ สนฺติมคฺคเมว พฺรูหย,}
\csromangathagroup{\csromangathastanza{\csromanfolio{32}\csromangathabat{“อุจฺฉินฺท สิเนหมตฺตโน,}{กุมุทํ สารทิกํว ปาณินา.} \\ \csromangathabat{สนฺติมคฺคเมว พฺรูหย,}{นิพฺพานํ สุคเตน เทสิตนฺ”ติ\csromansharedfootnote{edc3de8dd2179eb9}{ขุ 1. 54 ปิฏฺเฐ ธมฺมปเท.}.}}}

\prose{อิมินา วตฺถุนา อิมินา นิทาเนน เอวํ ญายติ “อิธ ภควา อชฺฌตฺติก\-วตฺถุกาย ตณฺหาย ปหานํ อาหา”ติ. อิทํ นิทานํ. (3)}

\prose{ตตฺถ กตโม ปุพฺพาปร\-สนฺธิ. ยถาห\csromandash{}}

\prose{\csromansymbolfootnote{+}{ขุ 1. 172 ปิฏฺเฐ; อุปริ 107, 184, 196 ปิฏฺเฐสุปิ.}“กามนฺธา ชาลสญฺฉนฺนา, ตณฺหา\-ฉทน\-ฉาทิตา.}

\setlength{\csromangathapagewidth}{0pt}
\setlength{\csromangathawakwidth}{0pt}
\csromangathameasuregroup{ปมตฺตพนฺธนา\textsuperscript{1} พทฺธา\textsuperscript{1}, \\ ชรามรณมเนฺวนฺติ,}{\csromangathabat{ปมตฺตพนฺธนา\textsuperscript{1} พทฺธา\textsuperscript{1},}{มจฺฉาว กุมินามุเข.} \\ \csromangathabat{ชรามรณมเนฺวนฺติ,}{วจฺโฉ ขีรปโกว มาตรนฺ”ติ.}}
\csromangathasetleft{ปมตฺตพนฺธนา\textsuperscript{1} พทฺธา\textsuperscript{1}, \\ ชรามรณมเนฺวนฺติ,}
\csromangathagroup{\csromangathastanza{\csromangathabat{ปมตฺตพนฺธนา\csromansharedfootnote{69872777a1cd06d4}{ปมตฺตพนฺธุนา ขุ 1. 172 ปิฏฺเฐ. ตทฏฺฐกถาปิ ปสฺสิตพฺพา.} พทฺธา\csromansharedfootnote{d70ddbbe9aae4984}{พนฺธา (ก)},}{มจฺฉาว กุมินามุเข.} \\ \csromangathabat{ชรามรณมเนฺวนฺติ,}{วจฺโฉ ขีรปโกว มาตรนฺ”ติ.}}}

\prose{อยํ กามตณฺหา วุตฺตา. สา กตเมน ปุพฺพาปเรน ยุชฺชติ, ยถาห\csromandash{}}

\prose{\csromansymbolfootnote{*}{เหฏฺฐา 12 ปิฏฺเฐปิ.}“รตฺโต อตฺถํ น ชานาติ, รตฺโต ธมฺมํ น ปสฺสติ.}

\setlength{\csromangathapagewidth}{0pt}
\setlength{\csromangathawakwidth}{0pt}
\csromangathameasuregroup{อนฺธนฺตมํ ตทา โหติ,}{\csromangathabat{อนฺธนฺตมํ ตทา โหติ,}{ยํ ราโค สหเต นรนฺ”ติ.}}
\csromangathasetleft{อนฺธนฺตมํ ตทา โหติ,}
\csromangathagroup{\csromangathastanza{\csromangathabat{อนฺธนฺตมํ ตทา โหติ,}{ยํ ราโค สหเต นรนฺ”ติ.}}}

\prose{อิติ อนฺธตาย จ สญฺฉนฺนตาย จ สาเยว ตณฺหา อภิลปิตา. ยญฺจาห กามนฺธา ชาลสญฺฉนฺนา, ตณฺหา\-ฉทน\-ฉาทิตาติ. ยญฺจาห รตฺโต อตฺถํ น ชานาติ, รตฺโต ธมฺมํ น ปสฺสตีติ, อิเมหิ ปเทหิ ปริยุฏฺฐาเนหิ สาเยว ตณฺหา อภิลปิตา. ยํ อนฺธการํ, อยํ ทุกฺขสมุทโย, ยา จ ตณฺหา โปโนภวิกา, ยญฺจาห กามาติ อิเม กิเลสกามา. ยญฺจาห ชาลสญฺฉนฺนาติ เตสํ เยว\csromansharedfootnote{9a3af51fcf41c94d}{เตสํเยว นิคฺคหีตสนฺธิ \csromandash{} ม.พ.ป.} กามานํ ปโยเคน ปริยุฏฺฐานํ ทสฺเสติ, ตสฺมา กิเลสวเสน จ ปริยุฏฺฐาน\-วเสน จ ตณฺหาพนฺธนํ วุตฺตํ. เย เอทิสิกา, เต ชรามรณํ อเนฺวนฺติ, อยํ ภควตา ยถานิกฺขิตฺต\-คาถา\-พเลน ทสฺสิตา ชารามรณมเนฺวนฺตีติ.}

\setlength{\csromangathapagewidth}{0pt}
\setlength{\csromangathawakwidth}{0pt}
\csromangathameasuregroup{}{“ยสฺส ปปญฺจา ฐิตี จ นตฺถิ, \\ สนฺทานํ ปลิฆญฺจ\textsuperscript{1} วีติวตฺโต. \\ ตํ นิตฺตณฺหํ มุนิํ จรนฺตํ, \\ น วิชานาติ สเทวโกปิ โลโก”ติ.}
\csromangathameasurewak{“ยสฺส ปปญฺจา ฐิตี จ นตฺถิ, \\ สนฺทานํ ปลิฆญฺจ\textsuperscript{1} วีติวตฺโต. \\ ตํ นิตฺตณฺหํ มุนิํ จรนฺตํ, \\ น วิชานาติ สเทวโกปิ โลโก”ติ.}
\setlength{\csromangathaleftcol}{0pt}
\csromangathagroup{\csromangathastanza{\csromangathawak{“ยสฺส ปปญฺจา ฐิตี จ นตฺถิ, \\ สนฺทานํ ปลิฆญฺจ\csromansharedfootnote{554aea25e87686e6}{ปฬิฆญฺจ (สี) ขุ 1. 174 ปิฏฺเฐ ปสฺสิตพฺพํ.} วีติวตฺโต. \\ ตํ นิตฺตณฺหํ มุนิํ จรนฺตํ, \\ น วิชานาติ สเทวโกปิ โลโก”ติ.}}}

\prose{ปปญฺจา นาม ตณฺหาทิฏฺฐิ\-มานา, ตทภิสงฺขตา จ สงฺขารา. ฐิติ นาม อนุสยา. สนฺทานํ นาม ตณฺหาย ปริยุฏฺฐานํ, ยานิ ฉตฺติํ\-สตณฺหาย}

\prose{\csromanfolio{33}ชาลินิยา วิจริตานิ. ปลิโฆ นาม โมโห. เย จ ปปญฺจา สงฺขารา ยา จ ฐิติ ยํ สนฺทานญฺจ ยํ ปลิฆญฺจ โย เอตํ สพฺพํ สมติกฺกนฺโต, อยํ วุจฺจติ นิตฺตโณฺห อิติ.}

\proseitem{28}{ตตฺถ ปริยุฏฺฐาน\-สงฺขารา ทิฏฺฐทมฺมเวทนียา วา อุปปชฺช\-เวทนียา วา อปราปริย\-เวทนียา วา, เอวํ ตณฺหา ติวิธํ ผลํ เทติ ทิฏฺเฐ วา ธมฺเม อุปปชฺเช วา อปเร วา ปริยาเย, เอวํ ภควา อาห “ยํ โลภปกตํ กมฺมํ กโรติ กาเยน วา วาจาย วา มนสา วา, ตสฺส วิปากํ อนุโภติ ทิฏฺเฐ วา ธมฺเม อุปปชฺเช วา อปเร วา ปริยาเย”ติ. อิทํ ภควโต ปุพฺพาปเรน ยุชฺชติ. ตตฺถ ปริยุฏฺฐานํ ทิฏฺฐธมฺม\-เวทนียํ วา กมฺมํ อุปปชฺช\-เวทนียํ วา กมฺมํ อปราปริยาย\-เวทนียํ\csromansharedfootnote{685e1e0804adc9cc}{อปราปริยเวทนียํ (สี)} วา กมฺมํ, เอวํ กมฺมํ ติธา วิปจฺจติ ทิฏฺเฐ วา ธมฺเม อุปปชฺเช วา อปเร วา ปริยาเย. ยถาห\csromandash{}}

\prose{\csromansymbolfootnote{*}{ม 3. 257 ปิฏฺเฐ.}“ยญฺเจ พาโล อิธ ปาณาติปาตี โหติ ฯเปฯ มิจฺฉาทิฏฺฐิ โหติ, ตสฺส ทิฏฺเฐ วา ธมฺเม วิปากํ ปฏิสํเวเทติ อุปปชฺเช วา อปเร วา ปริยาเย”ติ. อิทํ ภควโต ปุพฺพาปเรน ยุชฺชติ. ตตฺถ ปริยุฏฺฐานํ ปฏิสงฺขาน\-พเลน ปหาตพฺพํ, สงฺขารา ทสฺสนพเลน, ฉตฺติํส ตณฺหา\-วิจริตานิ ภาวนาพเลน ปหาตพฺพานีติ เอวํ ตณฺหาปิ ติธา ปหียติ. ยา นิตฺตณฺหาตา อยํ สอุปาทิเสสา นิพฺพานธาตุ. เภทา กายสฺส อยํ อนุปาทิเสสา นิพฺพานธาตุ.}

\prose{ปปญฺโจ นาม วุจฺจติ อนุพนฺโธ. ยญฺจาห ภควา “ปปญฺเจติ อตีตา\-นาคต\-ปจฺจุปฺปนฺนํ จกฺขุ\-วิญฺเญยฺยํ รูปํ อารพฺภา”ติ. ยญฺจาห ภควา “อตีเต ราธ รูเป อนเปกฺโข โหติ, อนาคตํ รูปํ มา อภินนฺทิ, ปจฺจุปฺปนฺนสฺส รูปสฺส นิพฺพิทาย วิราคาย นิโรธาย ปฏินิสฺสคฺคาย ปฏิปชฺชา”ติ. อิทํ ภควโต ปุพฺพาปเรน ยุชฺชติ. โย จาปิ ปปญฺโจ เย จ สงฺขารา ยา จ อตีตา\-นาคต\-ปจฺจุปฺปนฺนสฺส อภินนฺทนา, อิทํ เอกตฺถํ. อปิ จ อญฺญมญฺเญหิ ปเทหิ อญฺญมญฺเญหิ อกฺขเรหิ อญฺญมญฺเญหิ พฺยญฺชเนหิ อปริมาณา ธมฺมเทสนา วุตฺตา ภควตา. เอวํ สุตฺเตน สุตฺตํ สํสนฺทยิตฺวา ปุพฺพาปเรน สทฺธิํ โยชยิตฺวา สุตฺตํ นิทฺทิฏฺฐํ ภวติ.}

\prose{โส จายํ\csromansharedfootnote{f3eb6e63bfbf92ee}{ส จายํ (สี)} ปุพฺพาปโร สนฺธิ จตุพฺพิโธ อตฺถสนฺธิ พฺยญฺชนสนฺธิ เทสนาสนฺธิ นิทฺเทส\-สนฺธีติ.}

\prose{\csromanfolio{34}ตตฺถ อตฺถสนฺธิ ฉปฺปทานิ สงฺกาสนา ปกาสนา วิวรณา วิภชนา อุตฺตานีกมฺมตา ปญฺญตฺตีติ.}

\prose{พฺยญฺชนสนฺธิ ฉปฺปทานิ อกฺขรํ ปทํ พฺยญฺชนํ อากาโร นิรุตฺติ นิทฺเทโสติ.}

\prose{เทสนาสนฺธิ น จ ปถวิํ นิสฺสาย ฌายติ ฌายี ฌายติ จ. น จ อาปํ นิสฺสาย ฌายติ ฌายี ฌายติ จ, น จ เตชํ นิสฺสาย ฌายติ ฌายี ฌายติ จ, น จ วายุํ นิสฺสาย ฌายติ ฌายี ฌายติ จ ฯเปฯ น จ อากาสา\-นญฺจา\-ยตนํ นิสฺสาย. น จ วิญฺญาณญฺจา\-ยตนํ นิสฺสาย. น จ อากิญฺจญฺญา\-ยตนํ นิสฺสาย. น จ เนว\-สญฺญา\-นา\-สญฺญา\-ยตนํ นิสฺสาย. น จ อิมํ โลกํ นิสฺสาย. น จ ปรโลกํ นิสฺสาย ฌายติ ฌายี ฌายติ จ. ยมิทํ อุภยม\-นฺตเรน ทิฏฺฐํ สุตํ มุตํ วิญฺญาตํ ปตฺตํ ปริเยสิตํ วิตกฺกิตํ วิจาริตํ มนสานุ\-จินฺติตํ, ตมฺปิ นิสฺสาย น ฌายติ ฌายี ฌายติ จ. อยํ สเทวเก โลเก สมารเก สพฺรหฺมเก สสฺสมณ\-พฺราหฺมณิยา ปชาย สเทว\-มนุสฺสาย อนิสฺสิเตน จิตฺเตน น ญายติ ฌายนฺโต.}

\prose{ยถา มาโร ปาปิมา โคธิกสฺส กุลปุตฺตสฺส\csromansharedfootnote{fc1d7ea555f20924}{สํ 1. 123 ปิฏฺเฐ.} วิญฺญาณํ สมเนฺวสนฺโต น ชานาติ น ปสฺสติ. โส หิ ปปญฺจาตีโต ตณฺหาปหาเนน ทิฏฺฐินิสฺสโยปิสฺส นตฺถิ. ยถา จ โคธิกสฺส, เอวํ วกฺกลิสฺส สเทวเกน โลเกน สมารเกน สพฺรหฺมเกน สสฺสมณ\-พฺราหฺมณิยา ปชาย สเทว\-มนุสฺสาย อนิสฺสิตจิตฺตา น ญายนฺติ ฌายมานา. อยํ เทสนาสนฺธิ.}

\prose{ตตฺถ กตมา นิทฺเทสสนฺธิ, นิสฺสิตจิตฺตา อกุสล\-ปกฺเขน นิทฺทิสิตพฺพา, อนิสฺสิตจิตฺตา กุสลปกฺเขน นิทฺทิสิตพฺพา. นิสฺสิตจิตฺตา กิเลเสน นิทฺทิสิตพฺพา, อนิสฺสิตจิตฺตา โวทาเนน นิทฺทิสิตพฺพา. นิสฺสิตจิตฺตา สํสาร\-ปฺปวตฺติยา นิทฺทิสิตพฺพา, อนิสฺสิตจิตฺตา สํสาร\-นิวตฺติยา นิทฺทิสิตพฺพา. นิสฺสิตจิตฺตา ตณฺหาย จ อวิชฺชาย จ นิทฺทิสิตพฺพา, อนิสฺสิตจิตฺตา สมเถน จ วิปสฺสนาย จ นิทฺทิสิตพฺพา. นิสฺสิตจิตฺตา อหิริเกน จ อโนตฺตปฺเปน จ นิทฺทิสิตพฺพา, อนิสฺสิตจิตฺตา หิริยา จ โอตฺตปฺเปน จ นิทฺทิสิตพฺพา. นิสฺสิตจิตฺตา อสติยา จ อสมฺปชญฺเญน จ นิทฺทิสิตพฺพา, อนิสฺสิตจิตฺตา สติยา จ สมฺปชญฺเญน จ นิทฺทิสิตพฺพา. นิสฺสิตจิตฺตา \csromanfolio{35}อโยนิยา จ อโยนิโส\-มนสิกาเรน จ นิทฺทิสิตพฺพา, อนิสฺสิตจิตฺตา โยนิยา จ โยนิโส\-มนสิกาเรน จ นิทฺทิสิตพฺพา. นิสฺสิตจิตฺตา โกสชฺเชน จ โทวจสฺเสน จ นิทฺทิสิตพฺพา, อนิสฺสิตจิตฺตา วีริยารมฺเภน จ โสวจสฺเสน จ นิทฺทิสิตพฺพา. นิสฺสิตจิตฺตา อสฺสทฺธิเยน จ ปมาเทน จ นิทฺทิสิตพฺพา, อนิสฺสิตจิตฺตา สทฺธาย จ อปฺปมาเทน จ นิทฺทิสิตพฺพา. นิสฺสิตจิตฺตา อสทฺธมฺม\-สฺสวเนน จ อสํวรเณน จ นิทฺทิสิตพฺพา, อนิสฺสิตจิตฺตา สทฺธมฺม\-สฺสวเนน จ สํวเรน จ นิทฺทิสิตพฺพา. นิสฺสิตจิตฺตา อภิชฺฌาย จ พฺยาปาเทน จ นิทฺทิสิตพฺพา, อนิสฺสิตจิตฺตา อนภิชฺฌาย จ อพฺยาปาเทน จ นิทฺทิสิตพฺพา. นิสฺสิตจิตฺตา นีวรเณหิ จ สญฺโญชนิเยหิ จ นิทฺทิสิตพฺพา, อนิสฺสิตจิตฺตา ราควิราคาย จ เจโตวิมุตฺติยา อวิชฺชาวิราคาย จ ปญฺญาวิมุตฺติยา นิทฺทิสิตพฺพา. นิสฺสิตจิตฺตา อุจฺเฉท\-ทิฏฺฐิยา จ สสฺสต\-ทิฏฺฐิยา จ นิทฺทิสิตพฺพา, อนิสฺสิตจิตฺตา สอุปาทิเสสาย จ อนุปาทิเสสาย จ นิพฺพานธาตุยา นิทฺทิสิตพฺพา. อยํ นิทฺเทสสนฺธิ. เตนาห อายสฺมา มหากจฺจายโน “เนรุ\-ตฺตม\-ธิปฺปาโย”ติ.}

\vspace{\tipitakaparskip}
\csromancenter{นิยุตฺโต จตุพฺยูโห หาโร.}
\csromanmatikaanchor{mtk.11}
\csromantocmarkref{section}{7. อาวฏฺฏหารวิภงฺค}{mtk.11}
\bookmark[dest={mtk.11},level=1]{7. อาวฏฺฏหารวิภงฺค}
\csromanheader{1}{7. \textbf{อาวฏฺฏ}\-\textbf{หาร}\-\textbf{วิภงฺค}}
\proseitem{29}{ตตฺถ กตโม อาวฏฺโฏ หาโร, “เอกมฺหิ ปทฏฺฐาเน”ติ อยํ.}

\setlength{\csromangathapagewidth}{0pt}
\setlength{\csromangathawakwidth}{0pt}
\csromangathameasuregroup{\textsuperscript{*}“อารมฺภถ\textsuperscript{1} นิกฺกมถ, \\ ธุนาถ มจฺจุโน เสนํ,}{\csromangathabat{\textsuperscript{*}“อารมฺภถ\textsuperscript{1} นิกฺกมถ,}{ยุญฺชถ พุทฺธสาสเน.} \\ \csromangathabat{ธุนาถ มจฺจุโน เสนํ,}{นฬาคารํว กุญฺชโร”ติ.}}
\csromangathasetleft{\textsuperscript{*}“อารมฺภถ\textsuperscript{1} นิกฺกมถ, \\ ธุนาถ มจฺจุโน เสนํ,}
\csromangathagroup{\csromangathastanza{\csromangathaitembat{29}{\csromansymbolfootnote{*}{สํ 1. 158; ขุ 2. 267 ปิฏฺเฐสุ; อุปริ 217 ปิฏฺเฐปิ.}“อารมฺภถ\csromansharedfootnote{5bec48cebfd3b614}{อารพฺภถ (สี)} นิกฺกมถ,}{ยุญฺชถ พุทฺธสาสเน.} \\ \csromangathacontbat{ธุนาถ มจฺจุโน เสนํ,}{นฬาคารํว กุญฺชโร”ติ.}}}

\prose{“อารมฺภถ นิกฺกมถา”ติ วีริยสฺส ปทฏฺฐานํ. “ยุญฺชถ พุทฺธสาสเน”ติ สมาธิสฺส ปทฏฺฐานํ. “ธุนาถ มจฺจุโน เสนํ, นฬาคารํว กุญฺชโร”ติ ปญฺญาย ปทฏฺฐานํ.  “อารมฺภถ นิกฺกมถา”ติ วีริยิ\-นฺทฺริยสฺส ปทฏฺฐานํ. “ยุญฺชถ พุทฺธสาสเน”ติ สมาธิ\-นฺทฺริยสฺส ปทฏฺฐานํ. “ธุนาถ มจฺจุโน เสนํ, นฬาคารํว กุญฺชโร”ติ ปญฺญินฺทฺริยสฺส ปทฏฺฐานํ. อิมานิ ปทฏฺฐานานิ เทสนา.}

\prose{อยุญฺชนฺตานํ วา สตฺตานํ โยเค, ยุญฺชนฺตานํ วา อารมฺโภ.}

\prose{ตตฺถ เย น ยุญฺชนฺติ, เต ปมาทมูลกา น ยุญฺชนฺติ. โส ปมาโท ทุวิโธ ตณฺหามูลโก อวิชฺชามูลโก จ. ตตฺถ อวิชฺชามูลโก \csromanfolio{36}เยน อญฺญาเณน นิวุโต เญยฺยฏฺฐานํ นปฺปชานาติ ปญฺจกฺขนฺธา อุปฺปาทวยธมฺมาติ, อยํ อวิชฺชามูลโก. โย ตณฺหามูลโก, โส ติวิโธ อนุปฺปนฺนานํ โภคานํ อุปฺปาทาย ปริเยสนฺโต ปมาทํ อาปชฺชติ, อุปฺปนฺนานํ โภคานํ อารกฺข\-นิมิตฺตํ ปริโภค\-นิมิตฺตญฺจ ปมาทํ อาปชฺชติ อยํ โลเก จตุพฺพิโธ ปมาโท เอกวิโธ อวิชฺชาย ติวิโธ ตณฺหาย. ตตฺถ อวิชฺชาย นามกาโย ปทฏฺฐานํ. ตณฺหาย รูปกาโย ปทฏฺฐานํ. ตํ กิสฺส เหตุ, รูปีสุ ภเวสุ อชฺโฌสานํ, อรูปีสุ สมฺโมโห. ตตฺถ รูปกาโย รูปกฺขนฺโธ นามกาโย จตฺตาโร อรูปิโน ขนฺธา. อิเม ปญฺจกฺขนฺธา กตเมน อุปาทาเนน สอุปาทานา, ตณฺหาย จ อวิชฺชาย จ. ตตฺถ ตณฺหา เทฺว อุปาทานานิ กามุปาทานญฺจ สีลพฺพตุปาทานญฺจ. อวิชฺชา เทฺว อุปาทานานิ ทิฏฺฐุปาทานญฺจ อตฺตวาทุปาทานญฺจ. อิเมหิ จตูหิ อุปาทาเนหิ เย สอุปาทานา ขนฺธา, อิทํ ทุกฺขํ. จตฺตาริ อุปาทานานิ, อยํ สมุทโย. ปญฺจกฺขนฺธา ทุกฺขํ. เตสํ ภควา ปริญฺญาย ปหานาย จ ธมฺมํ เทเสติ ทุกฺขสฺส ปริญฺญาย สมุทยสฺส ปหานาย.}

\proseitem{30}{ตตฺถ โย ติวิโธ ตณฺหามูลโก ปมาโท อนุปฺปนฺนานํ โภคานํ อุปฺปาทาย ปริเยสติ, อุปฺปนฺนานํ โภคานํ อารกฺขณญฺจ กโรติ ปริโภค\-นิมิตฺตญฺจ, ตสฺส สมฺปฏิเวเธน รกฺขณา ปฏิสํหรณา, อยํ สมโถ.}

\prose{โส กถํ ภวติ, ยทา ชานาติ กามานํ อสฺสาทญฺจ อสฺสาทโต อาทีนวญฺจ อาทีนวโต นิสฺสรณญฺจ นิสฺสรณโต โอการญฺจ สํกิเลสญฺจ โวทานญฺจ เนกฺขมฺเม จ อานิสํสํ. ตตฺถ ยา วีมํสา อุปปริกฺขา. อยํ วิปสฺสนา. อิเม เทฺว ธมฺมา ภาวนา\-ปาริปูริํ คจฺฉนฺติ สมโถ จ วิปสฺสนา จ. อิเมสุ ทฺวีสุ ธมฺเมสุ ภาวิยมาเนสุ เทฺว ธมฺมา ปหียนฺติ ตณฺหา จ อวิชฺชา จ, อิเมสุ ทฺวีสุ ธมฺเมสุ ปหีเนสุ จตฺตาริ อุปาทานานิ นิรุชฺฌนฺติ. อุปาทานนิโรธา ภวนิโรโธ. ภวนิโรธา ชาตินิโรโธ, ชาตินิโรธา ชรามรณํ โสก\-ปริเทว\-ทุกฺข\-โทมนสฺสุ\-ปายาสา นิรุชฺฌนฺติ. เอวเมตสฺส เกวลสฺส ทุกฺข\-กฺขนฺธสฺส นิโรโธ โหติ. อิติ ปุริมกานิ จ เทฺว สจฺจานิ ทุกฺขํ สมุทโย จ, สมโถ จ วิปสฺสนา จ มคฺโค. ภว นิโรโธ นิพฺพานํ อิมานิ จตฺตาริ สจฺจานิ. เตนาห ภควา “อารมฺภถ นิกฺกมถา”ติ.}

\prose{\csromanfolio{37}ยถาปิ มูเล อนุปทฺทเว ทเฬฺห.}

\prose{ฉินฺโนปิ รุกฺโข ปุนเรว\csromansharedfootnote{98277bfe3b7a2f66}{ปุนเทว (ก) ขุ 1. 62 ปิฏฺเฐ.} รูหติ.}

\setlength{\csromangathapagewidth}{0pt}
\setlength{\csromangathawakwidth}{0pt}
\csromangathameasuregroup{}{เอวมฺปิ ตณฺหานุสเย อนูหเต, \\ นิพฺพตฺตตี ทุกฺขมิทํ ปุนปฺปุนํ.}
\csromangathameasurewak{เอวมฺปิ ตณฺหานุสเย อนูหเต, \\ นิพฺพตฺตตี ทุกฺขมิทํ ปุนปฺปุนํ.}
\setlength{\csromangathaleftcol}{0pt}
\csromangathagroup{\csromangathastanza{\csromangathawak{เอวมฺปิ ตณฺหานุสเย อนูหเต, \\ นิพฺพตฺตตี ทุกฺขมิทํ ปุนปฺปุนํ.}}}

\prose{อยํ ตณฺหานุสโย กตมสฺสา ตณฺหาย, ภวตณฺหาย. โย เอตสฺส ธมฺมสฺส ปจฺจโย อยํ อวิชฺชา. อวิชฺชาปจฺจยา หิ ภวตณฺหา. อิเม เทฺว กิเลสา ตณฺหา จ อวิชฺชา จ. ตานิ จตฺตาริ อุปาทานานิ เตหิ จตูหิ อุปาทาเนหิ เย ส- อุปาทานา ขนฺธา, อิทํ ทุกฺขํ. จตฺตาริ อุปาทานานิ อยํ สมุทโย. ปญฺจกฺขนฺธา ทุกฺขํ. เตสํ ภควา ปริญฺญาย จ ปหานาย จ ธมฺมํ เทเสติ ทุกฺขสฺส ปริญฺญาย สมุทยสฺส ปหานาย.}

\prose{เยน ตณฺหานุสยํ สมูหนติ\csromansharedfootnote{b7f8b1aa77ffa0ab}{สมูหนฺติ (สี)}, อยํ สมโถ. เยน ตณฺหานุสยสฺส ปจฺจยํ อวิชฺชํ วารยติ, อยํ วิปสฺสนา. อิเม เทฺว ธมฺมา ภาวนา\-ปาริปูริํ คจฺฉนฺติ สมโถ จ วิปสฺสนา จ. ตตฺถ สมถสฺส ผลํ ราควิราคา เจโตวิมุตฺติ, วิปสฺสนาย ผลํ อวิชฺชาวิราคา ปญฺญาวิมุตฺติ. อิติ ปุริมกานิ จ เทฺว สจฺจานิ ทุกฺขํ สมุทโย จ, สมโถ วิปสฺสนา จ มคฺโค, เทฺว จ วิมุตฺติโย นิโรโธ. อิมานิ จตฺตาริ สจฺจานิ. เตนาห ภควา “ยถาปิ มูเล”ติ.}

\setlength{\csromangathapagewidth}{0pt}
\setlength{\csromangathawakwidth}{0pt}
\csromangathameasuregroup{\textsuperscript{*}“สพฺพปาปสฺส อกรณํ, \\ สจิตฺตปริโยทาปนํ\textsuperscript{1},}{\csromangathabat{\textsuperscript{*}“สพฺพปาปสฺส อกรณํ,}{กุสลสฺส อุปสมฺปทา.} \\ \csromangathabat{สจิตฺตปริโยทาปนํ\textsuperscript{1},}{เอตํ พุทฺธาน สาสนนฺ”ติ.}}
\csromangathasetleft{\textsuperscript{*}“สพฺพปาปสฺส อกรณํ, \\ สจิตฺตปริโยทาปนํ\textsuperscript{1},}
\csromangathagroup{\csromangathastanza{\csromangathabat{\csromansymbolfootnote{*}{ที 2. 42; ขุ 1. 41 ปิฏฺเฐสุ; อุปริ 68, 148, 161, 205, 208 ปิฏฺเฐสุปิ.}“สพฺพปาปสฺส อกรณํ,}{กุสลสฺส อุปสมฺปทา.} \\ \csromangathabat{สจิตฺต\-ปริโยทาปนํ\csromansharedfootnote{f72060a4dc0e967c}{ปริโยทปนํ (สี)},}{เอตํ พุทฺธาน สาสนนฺ”ติ.}}}

\prose{สพฺพปาปํ นาม ตีณิ ทุจฺจริตานิ กายทุจฺจริตํ วจีทุจฺจริตํ มโนทุจฺจริตํ, เต ทส อกุสล\-กมฺมปถา ปาณาติปาโต อทินฺนาทานํ กาเมสุ\-มิจฺฉาจาโร มุสาวาโท ปิสุณา วาจา ผรุสา วาจา สมฺผปฺปลาโป อภิชฺฌา พฺยาปาโท มิจฺฉาทิฏฺฐิ, ตานิ เทฺว กมฺมานิ เจตนา เจตสิกญฺจ. ตตฺถ โย จ ปาณาติปาโต ยา จ ปิสุณา วาจา ยา จ ผรุสา วาจา, อิทํ โทส\-สมุฏฺฐานํ. ยญฺจ อทินฺนาทานํ โย จ กาเมสุ\-มิจฺฉาจาโร โย จ มุสาวาโท, อิทํ โลภ\-สมุฏฺฐานํ, โย สมฺผปฺปลาโป, อิทํ โมห\-สมุฏฺฐานํ. อิมานิ สตฺต การณานิ เจตนากมฺมํ. ยา อภิชฺฌา, อยํ โลโภ อกุสลมูลํ. โย พฺยาปาโท, อยํ โทโส อกุสลมูลํ. \csromanfolio{38}ยา มิจฺฉาทิฏฺฐิ, อยํ มิจฺฉามคฺโค. อิมานิ ตีณิ การณานิ เจตสิกกมฺมํ. เตนาห “เจตนากมฺมํ เจตสิกกมฺมนฺ”ติ.}

\prose{อกุสลมูลํ ปโยคํ คจฺฉนฺตํ จตุพฺพิธํ อคติํ คจฺฉติ ฉนฺทา โทสา ภยา โมหา. ตตฺถ ยํ ฉนฺทา อคติํ คจฺฉติ, อิทํ โลภ\-สมุฏฺฐานํ. ยํ โทสา อคติํ คจฺฉติ, อิทํ โทส\-สมุฏฺฐานํ. ยํ ภยา จ โมหา จ อคติํ คจฺฉติ, อิทํ โมห\-สมุฏฺฐานํ. ตตฺถ โลโภ อสุภาย ปหียติ. โทโส เมตฺตาย. โมโห ปญฺญาย. ตถา โลโภ อุเปกฺขาย ปหียติ. โทโส เมตฺตาย จ กรุณาย จ. โมโห มุทิตาย ปหานํ อพฺภตฺถํ คจฺฉติ. เตนาห ภควา “สพฺพปาปสฺส อกรณนฺ”ติ.}

\proseitem{31}{สพฺพปาปํ นาม อฏฺฐ มิจฺฉตฺตานิ มิจฺฉาทิฏฺฐิ มิจฺฉาสงฺกปฺโป มิจฺฉาวาจา มิจฺฉากมฺมนฺโต มิจฺฉาอาชีโว มิจฺฉาวายาโม มิจฺฉาสติ มิจฺฉาสมาธิ, อิทํ วุจฺจติ สพฺพปาปํ. อิเมสํ อฏฺฐนฺนํ มิจฺฉตฺตานํ ยา อกิริยา อกรณํ อนชฺฌาจาโร, อิทํ วุจฺจติ สพฺพปาปสฺส อกรณํ.}

\prose{อฏฺฐสุ มิจฺฉตฺเตสุ ปหีเนสุ อฏฺฐ สมฺมตฺตานิ สมฺปชฺชนฺติ. อฏฺฐนฺนํ สมฺมตฺตานํ ยา กิริยา กรณํ สมฺปาทนํ, อยํ วุจฺจติ กุสลสฺส อุปสมฺปทา. สจิตฺต\-ปริโยทาปนนฺติ อตีตสฺส มคฺคสฺส ภาวนากิริยํ ทสฺสยติ, จิตฺเต ปริโยทาปิเต\csromansharedfootnote{f86e1486d01864d3}{ปริโยทปิเต (สี, ก)} ปญฺจกฺขนฺธา ปริโยทาปิตา ภวนฺติ, เอวญฺหิ ภควา อาห “เจโต\-วิสุทฺธตฺถํ ภิกฺขเว ตถาคเต พฺรหฺมจริยํ วุสฺสตี”ติ. ทุวิธา หิ ปริโยทาปนา นีวรณ\-ปฺปหานญฺจ อนุสย\-สมุคฺฆาโต จ. เทฺว ปริโยทาปน\-ภูมิโย ทสฺสนภูมิ จ, ภาวนาภูมิ จ, ตตฺถ ยํ ปฏิเวเธน ปริโยทาเปติ, อิทํ ทุกฺขํ. ยโต ปริโยทาเปติ, อยํ สมุทโย. เยน ปริโยทาเปติ, อยํ มคฺโค. ยํ ปริโยทาปิตํ, อยํ นิโรโธ. อิมานิ จตฺตาริ สจฺจานิ. เตนาห ภควา “สพฺพปาปสฺส อกรณนฺ”ติ.}

\setlength{\csromangathapagewidth}{0pt}
\setlength{\csromangathawakwidth}{0pt}
\csromangathameasuregroup{}{\textsuperscript{*}“ธมฺโม หเว รกฺขติ ธมฺมจาริํ, \\ ฉตฺตํ มหนฺตํ ยถ วสฺสกาเล. \\ เอสานิสํโส ธมฺเม สุจิณฺเณ, \\ น ทุคฺคติํ คจฺฉติ ธมฺมจารี”ติ.}
\csromangathameasurewak{\textsuperscript{*}“ธมฺโม หเว รกฺขติ ธมฺมจาริํ, \\ ฉตฺตํ มหนฺตํ ยถ วสฺสกาเล. \\ เอสานิสํโส ธมฺเม สุจิณฺเณ, \\ น ทุคฺคติํ คจฺฉติ ธมฺมจารี”ติ.}
\setlength{\csromangathaleftcol}{0pt}
\csromangathagroup{\csromangathastanza{\csromangathawak{\csromansymbolfootnote{*}{เหฏฺฐา 7. 29 ปิฏฺเฐสุ; อุปริ จ 197 ปิฏฺเฐปิ.}“ธมฺโม หเว รกฺขติ ธมฺมจาริํ, \\ ฉตฺตํ มหนฺตํ ยถ วสฺสกาเล. \\ เอสานิสํโส ธมฺเม สุจิณฺเณ, \\ น ทุคฺคติํ คจฺฉติ ธมฺมจารี”ติ.}}}

\prose{\csromanfolio{39}ธมฺโม นาม ทุวิโธ อินฺทฺริยสํวโร มคฺโค จ. ทุคฺคติ นาม ทุวิธา เทวมนุสฺเส วา อุปนิธาย อปายา ทุคฺคติ, นิพฺพานํ วา อุปนิธาย สพฺพา อุปปตฺติโย ทุคฺคติ. ตตฺถ ยา สํวรสีเล อขณฺฑการิตา, อยํ ธมฺโม สุจิณฺโณ อปาเยหิ รกฺขติ. เอวํ ภควา อาห\csromandash{}เทฺวมา ภิกฺขเว สีลวโต คติโย เทวา จ มนุสฺสา จ. เอวญฺจ นาฬนฺทายํ นิคเม อสิพนฺธก\-ปุตฺโต คามณิ ภควนฺตํ เอตทโวจ\csromandash{}}

\prose{\csromansymbolfootnote{*}{สํ 2. 498, 500 ปิฏฺเฐสุ.}“พฺราหฺมณา ภนฺเต ปจฺฉาภูมกา กามณฺฑลุกา เสวาลมาลิกา อุทโกโรหกา อคฺคิปริจารกา, เต มตํ กาลงฺกตํ อุยฺยาเปนฺติ นาม, สญฺญาเปนฺติ นาม, สคฺคํ นาม โอกฺกาเมนฺติ\csromansharedfootnote{9116a2814daec414}{อุคฺคเมนฺติ (สี)}. ภควา ปน ภนฺเต อรหํ สมฺมาสมฺพุทฺโธ ปโหติ ตถา กาตุํ, ยถา สพฺโพ โลโก กายสฺส เภทา ปรํ มรณา สุคติํ สคฺคํ โลกํ อุปปชฺเชยฺยา”ติ. เตน หิ คามณิ ตญฺเญเวตฺถ ปฏิปุจฺฉิสฺสามิ, ยถา เต ขเมยฺย, ตถา นํ พฺยากเรยฺยาสีติ.}

\prose{ตํ กิํ มญฺญสิ คามณิ, อิธสฺส ปุริโส ปาณาติปาตี อทินฺนาทายี กาเมสุ\-มิจฺฉาจารี มุสาวาที ปิสุณวาโจ ผรุสวาโจ สมฺผปฺปลาปี อภิชฺฌาลุ พฺยาปนฺนจิตฺโต มิจฺฉาทิฏฺฐิโก, ตเมนํ มหา ชนกาโย สงฺคมฺม สมาคมฺม อายาเจยฺย โถเมยฺย ปญฺชลิโก อนุปริสกฺเกยฺย “อยํ ปุริโส กายสฺส เภทา ปรํ มรณา สุคติํ สคฺคํ โลกํ อุปปชฺชตู”ติ. ตํ กิํ มญฺญสิ คามณิ, อปิ นุ โส ปุริโส มหโต ชนกายสฺส อายาจนเหตุ วา โถมนเหตุวา ปญฺชลิกํ\csromansharedfootnote{6a0a0d7e05f11a4a}{ปญฺชลิกา (สํ 2. 498 ปิฏฺเฐ.)} อนุปริสกฺกน\-เหตุ วา กายสฺส เภทา ปรํ มรณา สุคติํ สคฺคํ โลกํ อุปปชฺเชยฺยาติ. โน เหตํ ภนฺเต.}

\prose{เสยฺยถาปิ คามณิ ปุริโส มหติํ ปุถุสิลํ คมฺภีเร อุทกรหเท\csromansharedfootnote{5b40d0647ddfba1d}{อุทกทเห (ก)} ปกฺขิเปยฺย, ตเมนํ มหา ชนกาโย สงฺคมฺม สมาคมฺม อายาเจยฺย โถเมยฺย ปญฺชลิโก อนุปริสกฺเกยฺย “อุมฺมุชฺช โภ ปุถุสิเล, อุปฺลว โภ ปุถุสิเล, ถลมุปฺลว โภ ปุถุสิเล”ติ. ตํ กิํ มญฺญสิ คามณิ, อปิ นุ สา มหตี ปุถุสิลา มหโต ชนกายสฺส อายาจนเหตุ วา โถมนเหตุ วา ปญฺชลิกํ อนุปริสกฺกน\-เหตุ วา อุมฺมุชฺเชยฺย วา อุปฺลเวยฺย วา ถลํ วา อุปฺลเวยฺยาติ. โน \csromanfolio{40}เหตํ ภนฺเต. เอวเมว โข คามณิ โย โส ปุริโส ปาณาติปาตี ฯเปฯ มิจฺฉาทิฏฺฐิโก, กิญฺจาปิ นํ มหา ชนกาโย สงฺคมฺม สมาคมฺม อายาเจยฺย โถเมยฺย ปญฺชลิโก อนุปริสกฺเกยฺย “อยํ ปุริโส กายสฺส เภทา ปรํ มรณา สุคติํ สคฺคํ โลกํ อุปปชฺชตู”ติ. อถ โข โส ปุริโส กายสฺส เภทา ปรํ มรณา อปายํ ทุคฺคติํ วินิปาตํ นิรยํ อุปปชฺเชยฺย.}

\prose{ตํ กิํ มญฺญสิ คามณิ, อิธสฺส ปุริโส ปาณาติปาตา ปฏิวิรโต อทินฺนาทานา ปฏิวิรโต กาเมสุ\-มิจฺฉาจารา ปฏิวิรโต มุสาวาทา ปฏิวิรโต ปิสุณาย วาจาย ปฏิวิรโต ผรุสาย วาจาย ปฏิวิรโต สมฺผปฺปลาปา ปฏิวิรโต อนภิชฺฌาลุ อพฺยาปนฺนจิตฺโต สมฺมาทิฏฺฐิโก, ตเมนํ มหา ชนกาโย สงฺคมฺม สมาคมฺม อายาเจยฺย โถเมยฺย ปญฺชลิโก อนุปริสกฺเกยฺย “อยํ ปุริโส กายสฺส เภทา ปรํ มรณา อปายํ ทุคฺคติํ วินิปาตํ นิรยํ อุปปชฺชตู”ติ. ตํ กิํ มญฺญสิ คามณิ, อปิ นุ โส ปุริโส มหโต ชนกายสฺส อายาจนเหตุ วา โถมนเหตุ วา ปญฺชลิกํ อนุปริสกฺกน\-เหตุ วา กายสฺส เภทา ปรํ มรณา อปายํ ทุคฺคติํ วินิปาตํ นิรยํ อุปปชฺเชยฺยาติ. โน เหตํ ภนฺเต.}

\prose{เสยฺยถาปิ คามณิ ปุริโส สปฺปิกุมฺภํ วา เตลกุมฺภํ วา คมฺภีเร\csromansharedfootnote{c50eeb4bfdd8a333}{คมฺภีรํ (สี, ก)} อุทกรหเท โอคาเหตฺวา ภินฺเทยฺย. ตตฺร ยาสฺส สกฺขร\-กฐลํ\csromansharedfootnote{e617106a4bc8c5c0}{สกฺขรา วา กฐลา (สํ 2. 499 ปิฏฺเฐ)} วา, สา อโธคามี อสฺส. ยญฺจ ขฺวสฺส ตตฺร สปฺปิ วา เตลํ วา, ตํ อุทฺธํคามิ อสฺส. ตเมนํ มหา ชนกาโย สงฺคมฺม สมาคมฺม อายาเจยฺย โถเมยฺย ปญฺชลิโก อนุปริสกฺเกยฺย “โอสีท โภ สปฺปิเตล, สํสีท โภ สปฺปิเตล, อโธ คจฺฉ\csromansharedfootnote{08276c6c785b1fc5}{อวํคจฺฉ (สี, ก)} โภ สปฺปิเตลา”ติ. ตํ กิํ มญฺญสิ คามณิ, อปิ นุ ตํ สปฺปิเตลํ มหโต ชนกายสฺส อายาจนเหตุ วา โถมนเหตุ วา ปญฺชลิกํ อนุปริสกฺกน\-เหตุ วา “โอสีเทยฺย วา สํสีเทยฺย วา อโธ วา คจฺเฉยฺยา”ติ. โน เหตํ ภนฺเต.}

\prose{เอวเมว โข คามณิ โย โส ปุริโส ปาณาติปาตา ปฏิวิรโต ฯเปฯ สมฺมาทิฏฺฐิโก, กิญฺจาปิ นํ มหา ชนกาโย สงฺคมฺม สมาคมฺม อายาเจยฺย โถเมยฺย ปญฺชลิโก อนุปริสกฺเกยฺย “อยํ ปุริโส \csromanfolio{41}กายสฺส เภทา ปรํ มรณา อปายํ ทุคฺคติํ วินิปาตํ นิรยํ อุปปชฺชตู”ติ. อถ โข โส ปุริโส กายสฺส เภทา ปรํ มรณา สุคติํ สคฺคํ โลกํ อุปปชฺเชยฺย. อิติ ธมฺโม สุจิณฺโณ อปาเยหิ รกฺขติ. ตตฺถ ยา มคฺคสฺส ติกฺขตา อธิมตฺตตา, อยํ ธมฺโม สุจิณฺโณ สพฺพาหิ อุปปตฺตีหิ รกฺขติ. เอวํ ภควา อาห\csromandash{}}

\setlength{\csromangathapagewidth}{0pt}
\setlength{\csromangathawakwidth}{0pt}
\csromangathameasuregroup{“ตสฺมา รกฺขิตจิตฺตสฺส, \\ สมฺมาทิฏฺฐิปุเรกฺขาโร,}{\csromangathabat{“ตสฺมา รกฺขิตจิตฺตสฺส,}{สมฺมาสงฺกปฺปโคจโร.} \\ \csromangathabat{สมฺมาทิฏฺฐิปุเรกฺขาโร,}{ญตฺวาน อุทยพฺพยํ.}}
\csromangathasetleft{“ตสฺมา รกฺขิตจิตฺตสฺส, \\ สมฺมาทิฏฺฐิปุเรกฺขาโร,}
\csromangathagroup{\csromangathastanza{\csromangathabat{“ตสฺมา รกฺขิต\-จิตฺตสฺส,}{สมฺมาสงฺกปฺป\-โคจโร.} \\ \csromangathabat{สมฺมาทิฏฺฐิ\-ปุเรกฺขาโร,}{ญตฺวาน อุทยพฺพยํ.}}}

\prose{ถินมิทฺธาภิภู ภิกฺขุ, สพฺพา ทุคฺคติโย ชเห”ติ\csromansharedfootnote{338dd9d64fd7ee1e}{[มิสฺสินฺคฺ โนเต 0]}.}

\proseitem{32}{ตตฺถ ทุคฺคตีนํ เหตุ ตณฺหา จ อวิชฺชา จ, ตานิ จตฺตาริ อุปาทานานิ, เตหิ จตูหิ อุปาทาเนหิ เย สอุปาทานา ขนฺธา, อิทํ ทุกฺขํ. จตฺตาริ อุปาทานานิ, อยํ สมุทโย. ปญฺจกฺขนฺธา ทุกฺขํ, เตสํ ภควา ปริญฺญาย จ ปหานาย จ ธมฺมํ เทเสติ ทุกฺขสฺส ปริญฺญาย สมุทยสฺส ปหานาย. ตตฺถ ตณฺหาย ปญฺจินฺทฺริยานิ รูปีนิ ปทฏฺฐานํ. อวิชฺชาย มนินฺทฺริยํ ปทฏฺฐานํ. ปญฺจินฺทฺริยานิ รูปีนิ รกฺขนฺโต สมาธิํ ภาวยติ, ตณฺหญฺจ นิคฺคณฺหาติ. มนินฺทฺริยํ รกฺขนฺโต วิปสฺสนํ ภาวยติ, อวิชฺชญฺจ นิคฺคณฺหาติ, ตณฺหานิคฺคเหน เทฺว อุปาทานานิ ปหียนฺติ กามุปาทานญฺจ สีลพฺพตุปาทานญฺจ. อวิชฺชา\-นิคฺคเหน เทฺว อุปาทานานิ ปหียนฺติ ทิฏฺฐุปาทานญฺจ อตฺตวาทุปาทานญฺจ. จตูสุ อุปาทาเนสุ ปหีเนสุ เทฺว ธมฺมา ภาวนา\-ปาริปูริํ คจฺฉนฺติ สมโถ จ วิปสฺสนา จ. อิทํ วุจฺจติ พฺรหฺม\-จริยนฺติ.}

\prose{ตตฺถ พฺรหฺม\-จริยสฺส ผลํ จตฺตาริ สามญฺญผลานิ โสตาปตฺติผลํ สกทาคามิผลํ อนาคามิผลํ อรหตฺตํ\csromansharedfootnote{973a47a9ee2527cc}{อรหตฺตผลํ (ก)} อคฺคผลํ. อิมานิ จตฺตาริ พฺรหฺม\-จริยสฺส ผลานิ\csromansharedfootnote{f6b13c7630350e37}{พฺรหฺมจริยผลานีติ (สี)}. อิติ ปุริมกานิ จ เทฺว สจฺจานิ ทุกฺขํ สมุทโย จ. สมโถ จ วิปสฺสนา จ พฺรหฺมจริยญฺจ มคฺโค, พฺรหฺม\-จริยสฺส ผลานิ จ ตทารมฺมณา จ อสงฺขตาธาตุ นิโรโธ. อิมานิ จตฺตาริ สจฺจานิ. เตนาห ภควา “ธมฺโม หเว รกฺขตี”ติ.}

\prose{ตตฺถ ยํ ปฏิเวเธน รกฺขติ, อิทํ ทุกฺขํ. ยโต รกฺขติ, อยํ สมุทโย. เยน รกฺขติ, อยํ มคฺโค. ยํ รกฺขติ, อยํ นิโรโธ. อิมานิ จตฺตาริ สจฺจานิ. เตนาห อายสฺมา มหากจฺจายโน “เอกมฺหิ ปทฏฺฐาเน”ติ.}

\vspace{\tipitakaparskip}
\csromancenter{นิยุตฺโต อาวฏฺโฏ หาโร.}
\csromanmatikaanchor{mtk.12}
\csromantocmarkref{section}{8. วิภตฺติหารวิภงฺค}{mtk.12}
\bookmark[dest={mtk.12},level=1]{8. วิภตฺติหารวิภงฺค}
\csromanheader{1}{8. \textbf{วิภตฺติ}\-\textbf{หาร}\-\textbf{วิภงฺค}}
\proseitem{33}{\csromanfolio{42}ตตฺถ กตโม วิภตฺติหาโร, “ธมฺมญฺจ ปทฏฺฐานํ ภูมิญฺจา”ติ.}

\prose{เทฺว สุตฺตานิ วาสนาภาคิยญฺจ นิพฺเพธภาคิยญฺจ. เทฺว ปฏิปทา ปุญฺญภาคิยา จ ผลภาคิยา จ. เทฺว สีลานิ สํวรสีลญฺจ ปหานสีลญฺจ, ตตฺถ ภควา วาสนาภาคิยํ สุตฺตํ ปุญฺญภาคิยาย ปฏิปทาย เทสยติ, โส สํวรสีเล ฐิโต เตน พฺรหฺมจริเยน พฺรหฺมจารี ภวติ, ตตฺถ ภควา นิพฺเพธ\-ภาคิยํ สุตฺตํ ผลภาคิยาย ปฏิปทาย เทสยติ, โส ปหานสีเล ฐิโต เตน พฺรหฺมจริเยน พฺรหฺมจารี ภวติ.}

\prose{ตตฺถ กตมํ วาสนาภาคิยํ สุตฺตํ, วาสานาภาคิยํ นาม สุตฺตํ ทานกถา สีลกถา สคฺคกถา กามานํ อาทีนโว เนกฺขมฺเม อานิสํโสติ.}

\prose{ตตฺถ กตมํ นิพฺเพธ\-ภาคิยํ สุตฺตํ, นิพฺเพธ\-ภาคิยํ นาม สุตฺตํ ยา จตุสจฺจ\-ปฺปกาสนา, วาสนาภาคิเย สุตฺเต นตฺถิ ปชานนา, นตฺถิ มคฺโค, นตฺถิ ผลํ. นิพฺเพธ\-ภาคิเย สุตฺเต อตฺถิ ปชานนา, อตฺถิ มคฺโค, อตฺถิ ผลํ. อิมานิ จตฺตาริ สุตฺตานิ. อิเมสํ จตุนฺนํ สุตฺตานํ เทสนาย ผเลน สีเลน พฺรหฺมจริเยน สพฺพโต วิจเยน หาเรน วิจินิตฺวา ยุตฺติหาเรน โยชยิตพฺพา ยาวติกา ญาณสฺส ภูมิ.}

\proseitem{34}{ตตฺถ กตเม ธมฺมา สาธารณา, เทฺว ธมฺมา สาธารณา นามสาธารณา วตฺถุ\-สาธารณา จ. ยํ วา ปน กิญฺจิ อญฺญมฺปิ เอวํ ชาติยํ, มิจฺฉตฺต\-นิยตานํ สตฺตานํ อนิยตานญฺจ สตฺตานํ ทสฺสน\-ปฺปหาตพฺพา กิเลสา สาธารณา, ปุถุชฺชนสฺส โสตาปนฺนสฺส จ กามราค\-พฺยาปาทา สาธารณา, ปุถุชฺชนสฺส อนาคามิสฺส จ อุทฺธํภาคิยา สํโยชนา สาธารณา, ยํ กิญฺจิ อริยสาวโก โลกิยํ สมาปตฺติํ สมาปชฺชติ, สพฺพา สา อวีตราเคหิ\csromansharedfootnote{377ec4bd148fad78}{อวิคตราเคหิ (ก)} สาธารณา, สาธารณา หิ ธมฺมา เอวํ อญฺญมญฺญํ ปรํ ปรํ สกํ สกํ วิสยํ นาติวตฺตนฺติ. โยปิ อิเมหิ ธมฺเมหิ สมนฺนาคโต น โส ตํ ธมฺมํ อุปาติวตฺตติ. อิเม ธมฺมา สาธารณา.}

\prose{ตตฺถ กตเม ธมฺมา อสาธารณา ยาว เทสนํ อุปาทาย คเวสิตพฺพา เสกฺขาเสกฺขา ภพฺพาภพฺพาติ, อฏฺฐมกสฺส โสตาปนฺนสฺส จ กามราค\-พฺยาปาทา สาธารณา ธมฺมตา อสาธารณา, อฏฺฐมกสฺส อนาคามิสฺส จ \csromanfolio{43}อุทฺธมฺภาคิยา สํโยชนา สาธารณา ธมฺมตา อสาธารณา. สพฺเพสํ เสกฺขานํ นามํ สาธารณํ ธมฺมตา อสาธารณา. สพฺเพสํ ปฏิปนฺนกานํ นามํ สาธารณํ ธมฺมตา อสาธารณา. สพฺเพสํ เสกฺขานํ เสกฺขสีลํ สาธารณํ ธมฺมตา อสาธารณา. เอวํ วิเสสา\-นุปสฺสินา หีนุกฺกฏฺฐมชฺฌิมํ อุปาทาย คเวสิตพฺพํ.}

\prose{ทสฺสนภูมิ นิยามา\-วกฺกนฺติยา ปทฏฺฐานํ, ภาวนาภูมิ อุตฺตริกานํ ผลานํ ปตฺติยา ปทฏฺฐานํ, ทุกฺขา ปฏิปทา ทนฺธาภิญฺญา สมถสฺส ปทฏฺฐานํ, สุขา ปฏิปทา ขิปฺปาภิญฺญา วิปสฺสนาย ปทฏฺฐานํ, ทานมยํ ปุญฺญกิริย\-วตฺถุ ปรโต โฆสสฺส สาธารณํ ปทฏฺฐานํ, สีลมยํ ปุญฺญกิริย\-วตฺถุ จินฺตามยิยา ปญฺญาย สาธารณํ ปทฏฺฐานํ, ภาวนามยํ ปุญฺญกิริย\-วตฺถุ ภาวนามยิยา ปญฺญาย สาธารณํ ปทฏฺฐานํ. ทานมยํ ปุญฺญกิริย\-วตฺถุ ปรโต จ โฆสสฺส สุตมยิยา จ ปญฺญาย สาธารณํ ปทฏฺฐานํ สีลมยํ ปุญฺญกิริย\-วตฺถุ จินฺตามยิยา จ ปญฺญาย โยนิโส จ มนสิการสฺส สาธารณํ ปทฏฺฐานํ, ภาวนามยํ ปุญฺญกิริย\-วตฺถุ ภาวนามยิยา จ ปญฺญาย สมฺมาทิฏฺฐิยา จ สาธารณํ ปทฏฺฐานํ. ปติรูป\-เทส\-วาโส วิเวกสฺส จ สมาธิสฺส จ สาธารณํ ปทฏฺฐานํ, สปฺปุริสู\-ปนิสฺสโย ติณฺณญฺจ อเวจฺจ\-ปฺปสาทานํ สมถสฺส จ สาธารณํ ปทฏฺฐานํ, อตฺต\-สมฺมา\-ปณิธานํ หิริยา จ วิปสฺสนาย จ สาธารณํ ปทฏฺฐานํ, อกุสล\-ปริจฺจาโค กุสล\-วีมํสาย จ สมาธิ\-นฺทฺริยสฺส จ สาธารณํ ปทฏฺฐานํ, ธมฺม\-สฺวากฺขาตตา กุสลมูล\-โรปนาย จ ผลสมาปตฺติยา จ สาธารณํ ปทฏฺฐานํ, สํฆสุปฺปฏิปนฺนตา สํฆสุฏฺฐุตาย สาธารณํ ปทฏฺฐานํ, สตฺถุสมฺปทา อปฺปสนฺนานญฺจ ปสาทาย ปสนฺนานญฺจ ภิยฺโยภาวาย สาธารณํ ปทฏฺฐานํ, อปฺปฏิหต\-ปาติโมกฺขตา ทุมฺมงฺกูนญฺจ ปุคฺคลานํ นิคฺคหาย เปสลานญฺจ ปุคฺคลานํ ผาสุวิหาราย สาธารณํ ปทฏฺฐานํ. เตนาห อายสฺมา มหากจฺจายโน “ธมฺมญฺจ ปทฏฺฐานนฺ”ติ.}

\vspace{\tipitakaparskip}
\csromancenter{นิยุตฺโต วิภตฺติหาโร.}
\csromanmatikaanchor{mtk.13}
\csromantocmarkref{section}{9. ปริวตฺตนหารวิภงฺค}{mtk.13}
\bookmark[dest={mtk.13},level=1]{9. ปริวตฺตนหารวิภงฺค}
\csromanheader{1}{9. \textbf{ปริวตฺตน}\-\textbf{หาร}\-\textbf{วิภงฺค}}
\proseitem{35}{\csromanfolio{44}ตตฺถ กตโม ปริวตฺตโน หาโร, “กุสลากุสเล ธมฺเม”ติ.}

\prose{สมฺมาทิฏฺฐิสฺส ปุริส\-ปุคฺคลสฺส มิจฺจาทิฏฺฐิ นิชฺชิณฺณา ภวติ. เย จสฺส มิจฺฉาทิฏฺฐิ\-ปจฺจยา อุปฺปชฺเชยฺยุํ อเนเก\csromansharedfootnote{6fed1cf991a0e2cc}{อเนกา (ก)} ปาปกา อกุสลา ธมฺมา, เต จสฺส นิชฺชิณฺณา โหนฺติ. สมฺมาทิฏฺฐิ\-ปจฺจยา จสฺส อเนเก กุสลา ธมฺมา สมฺภวนฺติ, เต จสฺส ภาวนา\-ปาริปูริํ คจฺฉนฺติ. สมฺมา\-สงฺกปฺปสฺส ปุริส\-ปุคฺคลสฺส มิจฺฉาสงฺกปฺโป นิชฺชิณฺโณ ภวติ. เย จสฺส มิจฺฉาสงฺกปฺป\-ปจฺจยา อุปฺปชฺเชยฺยุํ อเนเก ปาปกา อกุสลา ธมฺมา, เต จสฺส นิชฺชิณฺณา โหนฺติ. สมฺมาสงฺกปฺป\-ปจฺจยา จสฺส อเนเก กุสลา ธมฺมา สมฺภวนฺติ. เต จสฺส ภาวนา\-ปาริปูริํ คจฺฉนฺติ. เอวํ สมฺมาวาจสฺส สมฺมา\-กมฺมนฺตสฺส สมฺมาอาชีวสฺส สมฺมาวายามสฺส สมฺมาสติสฺส สมฺมา\-สมาธิสฺส สมฺมา\-วิมุตฺตสฺส สมฺมา\-วิมุตฺติ\-ญาณ\-ทสฺสนสฺส ปุริส\-ปุคฺคลสฺส มิจฺฉา\-วิมุตฺติ\-ญาณ\-ทสฺสนํ นิชฺชิณฺณํ ภวติ. เย จสฺส มิจฺฉา\-วิมุตฺติญาณทสฺสน\-ปจฺจยา อุปฺปชฺเชยฺยุํ อเนเก ปาปกา อกุสลา ธมฺมา, เต จสฺส นิชฺชิณฺณา โหนฺติ. สมฺมา\-วิมุตฺติญาณทสฺสน\-ปจฺจยา จสฺส อเนเก กุสลา ธมฺมา สมฺภวนฺติ, เต จสฺส ภาวนา\-ปาริปูริํ คจฺฉนฺติ.}

\proseitem{36}{ยสฺส วา ปาณาติปาตา ปฏิวิรตสฺส ปาณาติปาโต ปหีโน โหติ. อทินฺนาทานา ปฏิวิรตสฺส อทินฺนาทานํ ปหีนํ โหติ. พฺรหฺมจาริสฺส อพฺรหฺมจริยํ ปหีนํ โหติ. สจฺจวาทิสฺส มุสาวาโท ปหีโน โหติ. อปิสุณวาจสฺส ปิสุณา วาจา ปหีนา โหติ. สณฺหวาจสฺส ผรุสา วาจา ปหีนา โหติ. กาลวาทิสฺส สมฺผปฺปลาโป ปหีโน โหติ. อนภิชฺฌาลุสฺส\csromansharedfootnote{7278bc0390edbf3d}{อนภิชฺฌามนสฺส (ก)} อภิชฺฌา ปหีนา โหติ. อพฺยาปนฺน\-จิตฺตสฺส พฺยาปาโท ปหีโน โหติ. สมฺมาทิฏฺฐิสฺส มิจฺฉาทิฏฺฐิ ปหีนา โหติ.}

\prose{เย จ โข เกจิ อริยํ อฏฺฐงฺคิกํ มคฺคํ ครหนฺติ, เนสํ สนฺทิฏฺฐิกา สหธมฺมิกา คารยฺหา วาทานุวาทา อาคจฺฉนฺติ. สมฺมาทิฏฺฐิญฺจ เต ภวนฺโต ธมฺมํ ครหนฺติ. เตน หิ เย มิจฺฉาทิฏฺฐิกา, เตสํ ภวนฺตานํ ปุชฺชา จ ปาสํสา จ. เอวํ สมฺมาสงฺกปฺปํ สมฺมาวาจํ สมฺมากมฺมนฺตํ สมฺมาอาชีวํ สมฺมาวายามํ สมฺมาสติํ สมฺมาสมาธิํ สมฺมาวิมุตฺติํ สมฺมา\-วิมุตฺติญาณทสฺสนญฺจ เต ภวนฺโต ธมฺมํ ครหนฺติ. เตน หิ เย มิจฺฉา\-วิมุตฺติญาณทสฺสนา, เตสํ ภวนฺตานํ ปุชฺชา จ ปาสํสา จ.}

\prose{\csromanfolio{45}เย จ โข เกจิ เอวมาหํสุ “ภุญฺชิตพฺพา กามา, ปริภุญฺชิ\-ตพฺพา กามา, อาเสวิตพฺพา กามา, นิเสวิตพฺพา กามา, ภาวยิตพฺพา กามา, พหุลีกาตพฺพา กามา”ติ. กาเมหิ เวรมณี เตสํ อธมฺโม.}

\prose{เย วา ปน เกจิ เอวมาหํสุ “อตฺต\-กิลมถา\-นุโยโค ธมฺโม”ติ. นิยฺยานิโก เตสํ ธมฺโม อธมฺโม. เย จ โข เกจิ เอวมาหํสุ “ทุกฺโข ธมฺโม”ติ. สุโข เตสํ ธมฺโม อธมฺโม. ยถา วา ปน ภิกฺขุโน สพฺพ\-สงฺขาเรสุ อสุภา\-นุปสฺสิโน วิหรโต สุภสญฺญา ปหียนฺติ. ทุกฺขา\-นุปสฺสิโน วิหรโต สุขสญฺญา ปหียนฺติ. อนิจฺจา\-นุปสฺสิโน วิหรโต นิจฺจสญฺญา ปหียนฺติ. อนตฺตา\-นุปสฺสิโน วิหรโต อตฺตสญฺญา ปหียนฺติ. ยํ ยํ วา ปน ธมฺมํ โรจยติ วา อุปคจฺฉติ วา, ตสฺส ตสฺส ธมฺมสฺส โย ปฏิปกฺโข, สฺวสฺส อนิฏฺฐโต อชฺฌาปนฺโน ภวติ. เตนาห อายสฺมา มหากจฺจายโน “กุสลา\-กุสลธมฺเม”ติ.}

\vspace{\tipitakaparskip}
\csromancenter{นิยุตฺโต ปริวตฺตโน หาโร.}
\csromanmatikaanchor{mtk.14}
\csromantocmarkref{section}{10. เววจนหารวิภงฺค}{mtk.14}
\bookmark[dest={mtk.14},level=1]{10. เววจนหารวิภงฺค}
\csromanheader{1}{10. \textbf{เววจน}\-\textbf{หาร}\-\textbf{วิภงฺค}}
\proseitem{37}{ตตฺถ กตโม เววจโน หาโร, “เววจนานิ พหูนี”ติ. ยถา เอกํ ภควา ธมฺมํ อญฺญมญฺเญหิ เววจเนหิ นิทฺทิสติ. ยถาห ภควา\csromandash{}}

\prose{\csromansymbolfootnote{*}{เหฏฺฐา 22 ปิฏฺเฐ; อุปริ 180, 183 ปิฏฺเฐสุปิ.}“อาสา จ ปิหา อภินนฺทนา จ,}

\prose{อเนกธาตูสุ สรา ปติฏฺฐิตา.}

\setlength{\csromangathapagewidth}{0pt}
\setlength{\csromangathawakwidth}{0pt}
\csromangathameasuregroup{}{อญฺญาณมูลปฺปภวา ปชปฺปิตา, \\ สพฺพา มยา พฺยนฺติกตา สมูลิกา”ติ.}
\csromangathameasurewak{อญฺญาณมูลปฺปภวา ปชปฺปิตา, \\ สพฺพา มยา พฺยนฺติกตา สมูลิกา”ติ.}
\setlength{\csromangathaleftcol}{0pt}
\csromangathagroup{\csromangathastanza{\csromangathawak{อญฺญาณมูลปฺปภวา ปชปฺปิตา, \\ สพฺพา มยา พฺยนฺติกตา สมูลิกา”ติ.}}}

\prose{\textbf{อาสา} นาม วุจฺจติ ยา ภวิสฺสสฺส อตฺถสฺส อาสีสนา\csromansharedfootnote{41826f8bc971e070}{อาสิํสนา (สี)} อวสฺสํ อาคมิสฺสตีติ อาสาสฺส อุปฺปชฺชติ. ปิหา นาม ยา วตฺตมานสฺส\csromansharedfootnote{ca94d1a9fc2621e1}{วตฺตมานกสฺส (สี)} อตฺถสฺส ปตฺถนา, เสยฺยตรํ วา ทิสฺวา “เอทิโส ภเวยฺยนฺ”ติ ปิหาสฺส อุปฺปชฺชติ. อตฺถนิปฺผตฺติปฏิปาลนา อภินนฺทนา นาม, ปิยํ วา ญาติํ อภินนฺทติ, ปิยํ วา ธมฺมํ อภินนฺทติ, อปฺปฏิกูลโต วา อภินนฺทติ.}

\prose{\csromanfolio{46}\textbf{อเนกธาตู}ติ จกฺขุธาตุ รูปธาตุ จกฺขุวิญฺญาณ\-ธาตุ, โสตธาตุ สทฺทธาตุ โสตวิญฺญาณ\-ธาตุ, ฆานธาตุ คนฺธธาตุ ฆานวิญฺญาณ\-ธาตุ, ชิวฺหาธาตุ รสธาตุ ชิวฺหาวิญฺญาณ\-ธาตุ, กายธาตุ โผฏฺฐพฺพ\-ธาตุ กายวิญฺญาณ\-ธาตุ, มโนธาตุ ธมฺมธาตุ มโนวิญฺญาณ\-ธาตุ.}

\prose{\textbf{สรา}ติ เกจิ รูปาธิมุตฺตา เกจิ สทฺทาธิมุตฺตา เกจิ คนฺธาธิมุตฺตา เกจิ รสาธิมุตฺตา เกจิ โผฏฺฐพฺพา\-ธิมุตฺตา เกจิ ธมฺมาธิมุตฺตา. ตตฺถ ยานิ ฉ เคหสิตานิ โทมนสฺสานิ ยานิ จ ฉ เคหสิตานิ โสมนสฺสานิ ยานิ จ ฉ เนกฺขมฺม\-สิตานิ โทมนสฺสานิ ยานิ จ ฉ เนกฺขมฺม\-สิตานิ โสมนสฺสานิ, อิมานิ จตุวีสปทานิ ตณฺหาปกฺโข, ตณฺหาย เอตํ เววจนํ. ยา ฉ อุเปกฺขา เคหสิตา, อยํ ทิฏฺฐิปกฺโข.}

\proseitem{38}{สาเยว ปตฺตนากาเรน ธมฺมนนฺที ธมฺมเปมํ ธมฺมชฺโฌสานนฺติ ตณฺหาย เอตํ เววจนํ. จิตฺตํ มโน วิญฺญาณนฺติ จิตฺตสฺส เอตํ เววจนํ. มนินฺทฺริยํ มโนธาตุ มนายตนํ วิชานนาติ มนสฺเสตํ เววจนํ. ปญฺญินฺทฺริยํ ปญฺญาพลํ อธิปญฺญา สิกฺขา ปญฺญา ปญฺญากฺขนฺโธ ธมฺมวิจย\-สมฺโพชฺฌงฺโค ญาณํ สมฺมาทิฏฺฐิ ตีรณา วิปสฺสนา ธมฺเม ญาณํ อตฺเถ ญาณํ อนฺวเย ญาณํ เขเย ญาณํ อนุปฺปาเท ญาณํ อนญฺญาต\-ญฺญสฺสามี\-ตินฺทฺริยํ อญฺญินฺทฺริยํ อญฺญาตาวิ\-นฺทฺริยํ จกฺขุ วิชฺชา พุทฺธิ ภูริ เมธา อาโลโก, ยํ วา ปน ยํ กิญฺจิ อญฺญํปิ เอวํชาติยํ, ปญฺญาย เอตํ เววจนํ. ปญฺจินฺทฺริยานิ โลกุตฺตรานิ, สพฺพา ปญฺญา. อปิ จ อาธิปเตยฺย\-ฏฺเฐน สทฺธา, อารมฺภฏฺเฐน วีริยํ, อปิลาป\-นฏฺเฐน สติ, อวิกฺเขป\-ฏฺเฐน สมาธิ, ปชานนฏฺเฐน ปญฺญา.}

\prose{ยถา จ พุทฺธา\-นุสฺสติยํ วุตฺตํ อิติปิ โส ภควา อรหํ สมฺมาสมฺพุทฺโธ วิชฺชา\-จรณ\-สมฺปนฺโน สุคโต โลกวิทู อนุตฺตโร ปุริส\-ทมฺม\-สารถิ สตฺถา เทวมนุสฺสานํ พุทฺโธ ภควา. พลนิปฺผตฺติ\-คโต เวสารชฺช\-ปฺปตฺโต อธิคตปฺปฏิสมฺภิโท จตุโยค\-วิปฺปหีโน อคติคมน\-วีติวตฺโต อุทฺธฏสลฺโล นิรูฬฺหวโณ มทฺทิตกณฺฑโก นิพฺพาปิต\-ปริยุฏฺฐาโน\csromansharedfootnote{52e8f95b5ce66752}{นิพฺพาหิต... (ก)} พนฺธนาตีโต คนฺถวินิเวฐโน อชฺฌาสย\-วีติวตฺโต ภินฺนนฺธกาโร จกฺขุมา โลกธมฺม\-สมติกฺกนฺโต อนุโรธวิโรธ\-วิปฺปยุตฺโต อิฏฺฐานิฏฺเฐสุ ธมฺเมสุ อสงฺเขปคโต พนฺธ\-นาติวตฺโต ฐปิตสงฺคาโม อภิกฺกนฺตตโร อุกฺกาธโร อาโลกกโร ปชฺโชตกโร ตโมนุโท รณญฺชโห อปริมาณวณฺโณ อปฺปเมยฺยวณฺโณ อสงฺเขยฺย\-วณฺโณ อาภํกโร ปภํกโร \csromanfolio{47}ธมฺโม\-ภาส\-ปชฺโชต\-กโรติ จ พุทฺธา ภควนฺโตติ จ พุทฺธา\-นุสฺสติยา เอตํ เววจนํ.}

\prose{ยถา จ ธมฺมา\-นุสฺสติยํ วุตฺตํ สฺวากฺขาโต ภควตา ธมฺโม สนฺทิฏฺฐิโก อกาลิโก เอหิปสฺสิโก โอปเนยฺยิโก\csromansharedfootnote{c60f85659d806f24}{โอปนยิโก (สี)} ปจฺจตฺตํ เวทิตพฺโพ วิญฺญูหิ *. ยทิทํ มทนิมฺมทโน ปิปาสวินโย อาลย\-สมุคฺฆาโฏ วฏฺฏูปจฺเฉโท สุญฺญโต อติทุลฺลโภ ตณฺหกฺขโย วิราโค นิโรโธ นิพฺพานํ.}

\prose{** “อสงฺขตํ’นนฺตมนาสวญฺจ,}

\prose{สจฺจญฺจ ปารํ นิปุณํ สุทุทฺทสํ.}

\setlength{\csromangathapagewidth}{0pt}
\setlength{\csromangathawakwidth}{0pt}
\csromangathameasuregroup{}{อชชฺชรํ ธุวํ อปโลกิตํ\textsuperscript{1}, \\ อนิทสฺสนํ นิปฺปปญฺจ สนฺตํ.}
\csromangathameasurewak{อชชฺชรํ ธุวํ อปโลกิตํ\textsuperscript{1}, \\ อนิทสฺสนํ นิปฺปปญฺจ สนฺตํ.}
\setlength{\csromangathaleftcol}{0pt}
\csromangathagroup{\csromangathastanza{\csromangathawak{อชชฺชรํ ธุวํ อปโลกิตํ\csromansharedfootnote{5957d09028db1fe7}{อปโลกิยํ (สี, ก)}, \\ อนิทสฺสนํ นิปฺปปญฺจ สนฺตํ.}}}

\prose{** อมตํ ปณีตญฺจ สิวญฺจ เขมํ,}

\prose{ตณฺหกฺขโย อจฺฉริยญฺจ อพฺภุตํ.}

\setlength{\csromangathapagewidth}{0pt}
\setlength{\csromangathawakwidth}{0pt}
\csromangathameasuregroup{}{อนีติกํ อนีติกธมฺมํ\textsuperscript{1}, \\ นิพฺพานเมตํ สุคเตน เทสิตํ. \\ อชาตํ อภูตํ อนุปทฺทวญฺจ, \\ อกตํ อโสกญฺจ อโถ วิโสกํ. \\ อนูปสคฺคํนุปสคฺคธมฺมํ, \\ นิพฺพานเมตํ สุคเตน เทสิตํ. \\ คมฺภีรญฺเจว ทุปฺปสฺสํ, \\ อุตฺตรญฺจ อนุตฺตรํ. \\ อสมํ อปฺปฏิสมํ, \\ เชฏฺฐํ เสฏฺฐนฺติ วุจฺจติ. \\ เลณญฺจ ตาณํ อรณํ อนงฺคณํ, \\ อกาจ เมตํ วิมลนฺติ วุจฺจติ. \\ ทีโป สุขํ อปฺปมาณํ ปติฏฺฐา, \\ อกิญฺจนํ อปฺปปญฺจนฺติ วุตฺตนฺ”ติ.}
\csromangathameasurewak{อนีติกํ อนีติกธมฺมํ\textsuperscript{1}, \\ นิพฺพานเมตํ สุคเตน เทสิตํ. \\ อชาตํ อภูตํ อนุปทฺทวญฺจ, \\ อกตํ อโสกญฺจ อโถ วิโสกํ. \\ อนูปสคฺคํนุปสคฺคธมฺมํ, \\ นิพฺพานเมตํ สุคเตน เทสิตํ. \\ คมฺภีรญฺเจว ทุปฺปสฺสํ, \\ อุตฺตรญฺจ อนุตฺตรํ. \\ อสมํ อปฺปฏิสมํ, \\ เชฏฺฐํ เสฏฺฐนฺติ วุจฺจติ. \\ เลณญฺจ ตาณํ อรณํ อนงฺคณํ, \\ อกาจ เมตํ วิมลนฺติ วุจฺจติ. \\ ทีโป สุขํ อปฺปมาณํ ปติฏฺฐา, \\ อกิญฺจนํ อปฺปปญฺจนฺติ วุตฺตนฺ”ติ.}
\setlength{\csromangathaleftcol}{0pt}
\csromangathagroup{\csromangathastanza{\csromangathawak{อนีติกํ อนีติกธมฺมํ\csromansharedfootnote{1d955c237476203f}{นีติกธมฺมเมว วา (สี, ก)}, \\ นิพฺพานเมตํ สุคเตน เทสิตํ. \\ อชาตํ อภูตํ อนุปทฺทวญฺจ, \\ อกตํ อโสกญฺจ อโถ วิโสกํ.}}\csromangathastanza{\csromangathawak{อนูปสคฺคํนุปสคฺคธมฺมํ, \\ นิพฺพานเมตํ สุคเตน เทสิตํ. \\ คมฺภีรญฺเจว ทุปฺปสฺสํ, \\ อุตฺตรญฺจ อนุตฺตรํ.}}\csromangathastanza{\csromangathawak{อสมํ อปฺปฏิสมํ, \\ เชฏฺฐํ เสฏฺฐนฺติ วุจฺจติ. \\ เลณญฺจ ตาณํ อรณํ อนงฺคณํ, \\ อกาจ เมตํ วิมลนฺติ วุจฺจติ. \\ ทีโป สุขํ อปฺปมาณํ ปติฏฺฐา, \\ อกิญฺจนํ อปฺปปญฺจนฺติ วุตฺตนฺ”ติ.}}}

\prose{ธมฺมา\-นุสฺสติยา เอตํ เววจนํ.}

\prose{\csromanfolio{48}ยถา จ สํฆานุสฺสติยํ วุตฺตํ สุปฺปฏิปนฺโน อุชุปฺปฏิปนฺโน ญายปฺปฏิปนฺโน สามีจิ\-ปฺปฏิปนฺโน ยทิทํ จตฺตาริ ปุริสยุคานิ อฏฺฐ ปุริสปุคฺคลา เอส ภควโต สาวกสํโฆ อาหุเนยฺโย ปาหุเนยฺโย ทกฺขิเณยฺโย อญฺชลิกรณีโย อนุตฺตรํ ปุญฺญกฺเขตฺตํ โลกสฺส, สีลสมฺปนฺโน สมาธิ\-สมฺปนฺโน ปญฺญาสมฺปนฺโน วิมุตฺติ\-สมฺปนฺโน วิมุตฺติ\-ญาณ\-ทสฺสน\-สมฺปนฺโน สตฺตานํ สาโร สตฺตานํ มณฺโฑ สตฺตานํ อุทฺธาโร สตฺตานํ เอสิกา\csromansharedfootnote{ceba3ca9d948fd58}{เอสิโก (ก)} สตฺตานํ สุรภิปสูนํ ปุชฺโช เทวานญฺจ มนุสฺสานญฺจาติ สํฆานุสฺสติยา เอตํ เววจนํ.}

\prose{ยถา จ สีลานุสฺสติยํ วุตฺตํ ยานิ ตานิ สีลานิ อขณฺฑานิ อจฺฉิทฺทานิ อสพลานิ อกมฺมาสานิ อริยานิ อริยกนฺตานิ ภุชิสฺสานิ วิญฺญุ\-ปฺปสตฺถานิ อปรามฏฺฐานิ สมาธิ\-สํวตฺตนิกานิ, อลงฺกาโร จ สีลํ อุตฺตม\-งฺโคป\-โสภณ\-ตาย, นิธานญฺจ สีลํ สพฺพ\-โทภคฺค\-สมติกฺกม\-นฏฺเฐน, สิปฺปญฺจ สีลํ อกฺขณ\-เวธิตาย, เวลา จ สีลํ อนติกฺกม\-นฏฺเฐน, ธญฺญญฺจ สีลํ ทลิทฺโท\-ปจฺเฉทน\-ฏฺเฐน\csromansharedfootnote{b9fc64cda8b15f6c}{ทฬิทฺโท... (สี)}, อาทาโส จ สีลํ ธมฺม\-โวโลกน\-ตาย, ปาสาโท จ สีลํ โวโลกนฏฺเฐน, สพฺพ\-ภูมา\-นุปริวตฺติ จ สีลํ อมต\-ปริโยสานนฺติ สีลานุสฺสติยา เอตํ เววจนํ.}

\prose{ยถา จ จาคานุสฺสติยํ วุตฺตํ ยสฺมิํ สมเย อริยสาวโก อคารํ อชฺฌาวสติ มุตฺตจาโค ปยตปาณิ โวสฺสคฺครโต ยาจโยโค ทาน\-สํวิภาค\-รโตติ จาคานุสฺสติยา เอตํ เววจนํ. เตนาห อายสฺมา มหากจฺจายโน “เววจนานิ พหูนี”ติ.}

\vspace{\tipitakaparskip}
\csromancenter{นิยุตฺโต เววจโน หาโร.}
\csromanmatikaanchor{mtk.15}
\csromantocmarkref{section}{11. ปญฺญตฺติหารวิภงฺค}{mtk.15}
\bookmark[dest={mtk.15},level=1]{11. ปญฺญตฺติหารวิภงฺค}
\csromanheader{1}{11. \textbf{ปญฺญตฺติ}\-\textbf{หาร}\-\textbf{วิภงฺค}}
\proseitem{39}{ตตฺถ กตโม ปญฺญตฺติหาโร, “เอกํ ภควา ธมฺมํ ปญฺญตฺตีหิ วิวิธาหิ เทเสตี”ติ.}

\prose{ยา ปกติกถาย เทสนา. อยํ นิกฺเขป\-ปญฺญตฺติ. กา จ ปกติกถาย เทสนา, จตฺตาริ สจฺจานิ. ยถา ภควา อาห “อิทํ ทุกฺขนฺ”ติ อยํ ปญฺญตฺติ ปญฺจนฺนํ ขนฺธานํ ฉนฺนํ ธาตูนํ อฏฺฐารสนฺนํ ธาตูนํ ทฺวาทสนฺนํ อายตนานํ ทสนฺนํ อินฺทฺริยานํ นิกฺเขป\-ปญฺญตฺติ.}

\prose{\csromanfolio{49}\csromansymbolfootnote{*}{สํ 1. 325, 326; อภิ 4. 114; ขุ 7. 19 ปิฏฺเฐสุ.}กพฬีกาเร เจ ภิกฺขเว อาหาเร อตฺถิ ราโค อตฺถิ นนฺที\csromansharedfootnote{ebd9b94714ee2be0}{นนฺทิ (สี)} อตฺถิ ตณฺหา, ปติฏฺฐิตํ ตตฺถ วิญฺญาณํ วิรูฬฺหํ. ยตฺถ ปติฏฺฐิตํ วิญฺญาณํ วิรูฬฺหํ, อตฺถิ ตตฺถ นามรูปสฺส อวกฺกนฺติ. ยตฺถ อตฺถิ นามรูปสฺส อวกฺกนฺติ, อตฺถิ ตตฺถ สงฺขารานํ วุทฺธิ\csromansharedfootnote{a95db8d921defae0}{พุทฺธิ (ก)}. ยตฺถ อตฺถิ สงฺขารานํ วุทฺธิ, อตฺถิ ตตฺถ อายติํ ปุนพฺภวา\-ภินิพฺพตฺติ. ยตฺถ อตฺถิ อายติํ ปุนพฺภวา\-ภินิพฺพตฺติ, อตฺถิ ตตฺถ อายติํ ชาติ\-ชรา\-มรณํ. ยตฺถ อตฺถิ อายติํ ชาติ\-ชรา\-มรณํ, สโสกํ ตํ ภิกฺขเว สทรํ สอุปายาสนฺติ วทามิ.}

\prose{ผสฺเส เจ ฯเปฯ. มโน\-สญฺเจตนาย เจ ภิกฺขเว อาหาเร. วิญฺญาเณ เจ ภิกฺขเว อาหาเร อตฺถิ ราโค อตฺถิ นนฺที อตฺถิ ตณฺหา, ปติฏฺฐิตํ ตตฺถ วิญฺญาณํ วิรูฬฺหํ. ยตฺถ ปติฏฺฐิตํ วิญฺญาณํ วิรูฬฺหํ, อตฺถิ ตตฺถ นามรูปสฺส อวกฺกนฺติ. ยตฺถ อตฺถิ นามรูปสฺส อวกฺกนฺติ, อตฺถิ ตตฺถ สงฺขารานํ วุทฺธิ. ยตฺถ อตฺถิ สงฺขารานํ วุทฺธิ, อตฺถิ ตตฺถ อายติํ ปุนพฺภวา\-ภินิพฺพตฺติ. ยตฺถ อตฺถิ อายติํ ปุนพฺภวา\-ภินิพฺพตฺติ, อตฺถิ ตตฺถ อายติํ ชาติ\-ชรา\-มรณํ. ยตฺถ อตฺถิ อายติํ ชาติ\-ชรา\-มรณํ, สโสกํ ตํ ภิกฺขเว สทรํ สอุปายาสนฺติ วทามิ. อยํ ปภว\-ปญฺญตฺติ ทุกฺขสฺส จ สมุทยสฺส จ.}

\prose{กพฬีกาเร เจ ภิกฺขเว อาหาเร นตฺถิ ราโค นตฺถิ นนฺที นตฺถิ ตณฺหา, อปฺปติฏฺฐิตํ ตตฺถ วิญฺญาณํ อวิรูฬฺหํ. ยตฺถ อปฺปติฏฺฐิตํ วิญฺญาณํ อวิรูฬฺหํ, นตฺถิ ตตฺถ นามรูปสฺส อวกฺกนฺติ. ยตฺถ นตฺถิ นามรูปสฺส อวกฺกนฺติ, นตฺถิ ตตฺถ สงฺขารานํ วุทฺธิ. ยตฺถ นตฺถิ สงฺขารานํ วุทฺธิ, นตฺถิ ตตฺถ อายติํ ปุนพฺภวา\-ภินิพฺพตฺติ. ยตฺถ นตฺถิ อายติํ ปุนพฺภวา\-ภินิพฺพตฺติ, นตฺถิ ตตฺถ อายติํ ชาติ\-ชรา\-มรณํ. ยตฺถ นตฺถิ อายติํ ชาติ\-ชรา\-มรณํ, อโสกํ ตํ ภิกฺขเว อทรํ อนุปายาสนฺติ วทามิ.}

\prose{ผสฺเส เจ ฯเปฯ. มโน\-สญฺเจตนาย เจ ภิกฺขเว อาหาเร. วิญฺญาเณ เจ ภิกฺขเว อาหาเร นตฺถิ ราโค นตฺถิ นนฺที นตฺถิ ตณฺหา, อปฺปติฏฺฐิตํ ตตฺถ วิญฺญาณํ อวิรูฬฺหํ. ยตฺถ อปฺปติฏฺฐิตํ วิญฺญาณํ อวิรูฬฺหํ, นตฺถิ ตตฺถ นามรูปสฺส อวกฺกนฺติ. ยตฺถ นตฺถิ นามรูปสฺส อวกฺกนฺติ, นตฺถิ ตตฺถ สงฺขารานํ วุทฺธิ. ยตฺถ นตฺถิ สงฺขารานํ วุทฺธิ, นตฺถิ ตตฺถ อายติํ ปุนพฺภวา\-ภินิพฺพตฺติ. ยตฺถ นตฺถิ อายติํ ปุนพฺภวา\-ภินิพฺพตฺติ, นตฺถิ ตตฺถ อายติํ ชาติ\-ชรา\-มรณํ. ยตฺถ นตฺถิ อายติํ ชาติ\-ชรา\-มรณํ, อโสกํ ตํ ภิกฺขเว อทรํ อนุปายาสนฺติ วทามิ.}

\prose{\csromanfolio{50}อยํ ปริญฺญา\-ปญฺญตฺติ ทุกฺขสฺส, ปหานปญฺญตฺติ สมุทยสฺส, ภาวนา\-ปญฺญตฺติ มคฺคสฺส, สจฺฉิกิริยา\-ปญฺญตฺติ นิโรธสฺส.}

\proseitem{40}{\csromansymbolfootnote{*}{สํ 2. 302 ปิฏฺเฐ.}สมาธิํ ภิกฺขเว ภาเวถ อปฺปมตฺโต นิปโก สโต, สมาหิโต ภิกฺขเว ภิกฺขุ ยถาภูตํ ปชานาติ. กิญฺจ ยถาภูตํ ปชานาติ “จกฺขุ\csromansharedfootnote{835c456fc5e0bea9}{จกฺขุํ (ก)} อนิจฺจนฺ”ติ ยถาภูตํ ปชานาติ. “รูปา อนิจฺจา”ติ ยถาภูตํ ปชานาติ “จกฺขุวิญฺญาณํ อนิจฺจนฺ”ติ ยถาภูตํ ปชานาติ. “จกฺขุ\-สมฺผสฺโส อนิจฺโจ”ติ ยถาภูตํ ปชานาติ. ยมฺปิทํ\csromansharedfootnote{894fb7b72c278ece}{ยมิทํ (สี, ก)} จกฺขุ\-สมฺผสฺส\-ปจฺจยา อุปฺปชฺชติ เวทยิตํ สุขํ วา ทุกฺขํ วา อทุกฺขมสุขํ วา, ตมฺปิ อนิจฺจนฺติ ยถาภูตํ ปชานาติ.}

\prose{\csromansymbolmark{*}โสตํ ฯเปฯ. ฆานํ. ชิวฺหา. กาโย. “มโน อนิจฺโจ”ติ\csromansharedfootnote{deca7f5de0d69a2f}{อนิจฺจนฺ”ติ (สํ 2. 302 ปิฏฺเฐ)} ยถาภูตํ ปชานาติ. “ธมฺมา อนิจฺจา”ติ ยถาภูตํ ปชานาติ. “มโนวิญฺญาณํ อนิจฺจนฺ”ติ ยถาภูตํ ปชานาติ. “มโนสมฺผสฺโส อนิจฺโจ”ติ ยถาภูตํ ปชานาติ. ยมฺปิทํ มโน\-สมฺผสฺส\-ปจฺจยา อุปฺปชฺชติ เวทยิตํ สุขํ วา ทุกฺขํ วา อทุกฺขมสุขํ วา, ตมฺปิ อนิจฺจนฺติ ยถาภูตํ ปชานาติ.}

\prose{อยํ ภาวนา\-ปญฺญตฺติ มคฺคสฺส, ปริญฺญา\-ปญฺญตฺติ ทุกฺขสฺส, ปหานปญฺญตฺติ สมุทยสฺส, สจฺฉิกิริยา\-ปญฺญตฺติ นิโรธสฺส.}

\prose{\csromansymbolfootnote{+}{สํ 2. 155 ปิฏฺเฐ.}รูปํ ราธ วิกิรถ วิธมถ วิทฺธํเสถ วิกีฬนิยํ\csromansharedfootnote{b7a81ded9f6b9208}{วิกีฬนิกํ (สี, ก)} กโรถ, ปญฺญาย ตณฺหกฺขยาย ปฏิปชฺชถ. ตณฺหกฺขยา ทุกฺขกฺขโย, ทุกฺขกฺขยา นิพฺพานํ. เวทนํ ฯเปฯ. สญฺญํ. สงฺขาเร วิญฺญาณํ วิกิรถ วิธมถ วิทฺธํเสถ วิกีฬนิยํ กโรถ, ปญฺญาย ตณฺหกฺขยาย ปฏิปชฺชถ. ตณฺหกฺขยา ทุกฺขกฺขโย, ทุกฺขกฺขยา นิพฺพานํ.}

\prose{อยํ นิโรธ\-ปญฺญตฺติ นิโรธสฺส, นิพฺพิทา\-ปญฺญตฺติ อสฺสาทสฺส, ปริญฺญา\-ปญฺญตฺติ ทุกฺขสฺส, ปหานปญฺญตฺติ สมุทยสฺส, ภาวนา\-ปญฺญตฺติ มคฺคสฺส, สจฺฉิกิริยา\-ปญฺญตฺติ นิโรธสฺส.}

\prose{“โส อิทํ ทุกฺขนฺ”ติ ยถาภูตํ ปชานาติ, “อยํ ทุกฺขสมุทโย”ติ ยถาภูตํ ปชานาติ, “อยํ ทุกฺขนิโรโธ”ติ ยถาภูตํ ปชานาติ “อยํ ทุกฺขนิโรธ\-คามินี ปฏิปทา”ติ ยถาภูตํ ปชานาติ.}

\prose{\csromanfolio{51}อยํ ปฏิเวธ\-ปญฺญตฺติ สจฺจานํ, นิกฺเขป\-ปญฺญตฺติ ทสฺสนภูมิยา, ภาวนา\-ปญฺญตฺติ มคฺคสฺส, สจฺฉิกิริยา\-ปญฺญตฺติ โสตาปตฺติ\-ผลสฺส. “โส อิเม อาสวา”ติ ยถาภูตํ ปชานาติ, “อยํ อาสวสมุทโย”ติ ยถาภูตํ ปชานาติ, “อยํ อาสวนิโรโธ”ติ ยถาภูตํ ปชานาติ. “อยํ อาสว\-นิโรธ\-คามินี ปฏิปทา”ติ ยถาภูตํ ปชานาติ. “อิเม อาสวา อเสสํ นิรุชฺฌนฺตี”ติ ยถาภูตํ ปชานาติ.}

\prose{อยํ อุปฺปาทปญฺญตฺติ ขเย ญาณสฺส, โอกาสปญฺญตฺติ อนุปฺปาเท ญาณสฺส, ภาวนา\-ปญฺญตฺติ มคฺคสฺส, ปริญฺญา\-ปญฺญตฺติ ทุกฺขสฺส, ปหานปญฺญตฺติ สมุทยสฺส, อารมฺภ\-ปญฺญตฺติ วีริยิ\-นฺทฺริยสฺส, อาหฏ\-นา\-ปญฺญตฺติ อาสาฏิกานํ, นิกฺเขป\-ปญฺญตฺติ ภาวนาภูมิยา, อภินิฆาตปญฺญตฺติ ปาปกานํ อกุสลานํ ธมฺมานํ.}

\proseitem{41}{อิทํ “ทุกฺขนฺ”ติ เม ภิกฺขเว ปุพฺเพ อนนุสฺสุเตสุ ธมฺเมสุ จกฺขุํ อุทปาทิ, ญาณํ อุทปาทิ, ปญฺญา อุทปาทิ, วิชฺชา อุทปาทิ, อาโลโก อุทปาทิ. อยํ “ทุกฺขสมุทโย”ติ เม ภิกฺขเว ฯเปฯ. อยํ “ทุกฺขนิโรโธ”ติ เม ภิกฺขเว. อยํ “ทุกฺขนิโรธ\-คามินี ปฏิปทา”ติ เม ภิกฺขเว ปุพฺเพ อนนุสฺสุเตสุ ธมฺเมสุ จกฺขุํ อุทปาทิ, ญาณํ อุทปาทิ, ปญฺญา อุทปาทิ, วิชฺชา อุทปาทิ, อาโลโก อุทปาทิ\csromansharedfootnote{78c505e641836356}{วิ 3. 15, 16; สํ 3. 369 ปิฏฺเฐสุ ปสฺสิตพฺพํ.}.}

\prose{อยํ เทสนาปญฺญตฺติ สจฺจานํ, นิกฺเขป\-ปญฺญตฺติ สุตมยิยา ปญฺญาย, สจฺฉิกิริยา\-ปญฺญตฺติ อนญฺญาต\-ญฺญสฺสามี\-ตินฺทฺริยสฺส, ปวตฺตนา\-ปญฺญตฺติ ธมฺมจกฺกสฺส.}

\prose{“ตํ โข ปนิทํ ทุกฺขํ ปริญฺเญยฺยนฺ”ติ เม ภิกฺขเว ปุพฺเพ อนนุสฺสุเตสุ ธมฺเมสุ จกฺขุํ อุทปาทิ, ญาณํ อุทปาทิ, ปญฺญา อุทปาทิ, วิชฺชา อุทปาทิ, อาโลโก อุทปาทิ. “โส โข ปนายํ ทุกฺขสมุทโย ปหาตพฺโพ”ติ เม ภิกฺขเว ฯเปฯ. “โส โข ปนายํ ทุกฺขนิโรโธ สจฺฉิกาตพฺโพ”ติ เม ภิกฺขเว ฯเปฯ. “สา โข ปนายํ ทุกฺขนิโรธ\-คามินี ปฏิปทา ภาเวตพฺพา”ติ เม ภิกฺขเว ปุพฺเพ อนนุสฺสุเตสุ ธมฺเมสุ จกฺขุํ อุทปาทิ, ญาณํ อุทปาทิ, ปญฺญา อุทปาทิ, วิชฺชา อุทปาทิ, อาโลโก อุทปาทิ\csromansharedfootnote{78c505e641836356}{วิ 3. 15, 16; สํ 3. 369 ปิฏฺเฐสุ ปสฺสิตพฺพํ.}.}

\prose{อยํ ภาวนา\-ปญฺญตฺติ มคฺคสฺส, นิกฺเขป\-ปญฺญตฺติ จินฺตามยิยา ปญฺญาย, สจฺฉิกิริยา\-ปญฺญตฺติ อญฺญินฺทฺริยสฺส.}

\prose{\csromanfolio{52}“ตํ โข ปนิทํ ทุกฺขํ ปริญฺญาตนฺ”ติ เม ภิกฺขเว ปุพฺเพ อนนุสฺสุเตสุ ธมฺเมสุ จกฺขุํ อุทปาทิ, ญาณํ อุทปาทิ, ปญฺญา อุทปาทิ, วิชฺชา อุทปาทิ, อาโลโก อุทปาทิ. “โส โข ปนายํ ทุกฺขสมุทโย ปหีโน”ติ เม ภิกฺขเว ฯเปฯ. “โส โข ปนายํ ทุกฺขนิโรโธ สจฺฉิกโต”ติ เม ภิกฺขเว ฯเปฯ. “สา โข ปนายํ ทุกฺขนิโรธ\-คามินี ปฏิปทา ภาวิตา”ติ เม ภิกฺขเว ปุพฺเพ อนนุสฺสุเตสุ ธมฺเมสุ จกฺขุํ อุทปาทิ, ญาณํ อุทปาทิ, ปญฺญา อุทปาทิ, วิชฺชา อุทปาทิ, อาโลโก อุทปาทิ.* อยํ ภาวนา\-ปญฺญตฺติ มคฺคสฺส, นิกฺเขป\-ปญฺญตฺติ ภาวนามยิยา ปญฺญาย, สจฺฉิกิริยา\-ปญฺญตฺติ อญฺญาตาวิโน อินฺทฺริยสฺส, ปวตฺตนา\-ปญฺญตฺติ ธมฺมจกฺกสฺส.}

\setlength{\csromangathapagewidth}{0pt}
\setlength{\csromangathawakwidth}{0pt}
\csromangathameasuregroup{\textsuperscript{+}“ตุลมตุลญฺจ สมฺภวํ, \\ อชฺฌตฺตรโต สมาหิโต,}{\csromangathabat{\textsuperscript{+}“ตุลมตุลญฺจ สมฺภวํ,}{ภวสงฺขารมวสฺสชิ มุนิ.} \\ \csromangathabat{อชฺฌตฺตรโต สมาหิโต,}{อภินฺทิ\textsuperscript{1} กวจมิวตฺตสมฺภวนฺ”ติ.}}
\csromangathasetleft{\textsuperscript{+}“ตุลมตุลญฺจ สมฺภวํ, \\ อชฺฌตฺตรโต สมาหิโต,}
\csromangathagroup{\csromangathastanza{\csromangathabat{\csromansymbolfootnote{+}{ที 2. 90; สํ 3. 230; ขุ 1. 155 ปิฏฺเฐสุ; อุปริ 215 ปิฏฺเฐปิ.}“ตุลมตุลญฺจ สมฺภวํ,}{ภวสงฺขารม\-วสฺสชิ มุนิ.} \\ \csromangathabat{อชฺฌตฺตรโต สมาหิโต,}{อภินฺทิ\csromansharedfootnote{bd59c5fa22730f14}{อภิทา (สี, ก) ที 2. 90 ปิฏฺเฐ. ** สํ 1. 118, 120 ปิฏฺเฐสุ ปสฺสิตพฺพํ; อุปริ 171 ปิฏฺเฐปิ.} กวจมิวตฺตสมฺภวนฺ”ติ.}}}

\prose{“ตุลนฺ”ติ สงฺขารธาตุ. “อตุลนฺ”ติ นิพฺพานธาตุ, “ตุลมตุลญฺจ สมฺภวนฺ”ติ อภิญฺญา\-ปญฺญตฺติ สพฺพธมฺมานํ. นิกฺเขป\-ปญฺญตฺติ ธมฺม\-ปฏิสมฺภิทาย. “ภวสงฺขารม\-วสฺสชิ มุนี”ติ ปริจฺจาค\-ปญฺญตฺติ สมุทยสฺส. ปริญฺญา\-ปญฺญตฺติ ทุกฺขสฺส. “อชฺฌตฺตรโต สมาหิโต”ติ ภาวนา\-ปญฺญตฺติ กายคตาย สติยา. ฐิติปญฺญตฺติ จิตฺเตกคฺคตาย. “อภินฺทิ กวจมิวตฺตสมฺภวนฺ”ติ อภินิพฺพิทา\-ปญฺญตฺติ จิตฺตสฺส, อุปาทาน\-ปญฺญตฺติ สพฺพญฺญุตาย, ปทาลนา\-ปญฺญตฺติ อวิชฺชณฺฑโกสานํ. เตนาห ภควา “ตุลมตุลญฺจ สมฺภวนฺ”ติ.}

\setlength{\csromangathapagewidth}{0pt}
\setlength{\csromangathawakwidth}{0pt}
\csromangathameasuregroup{}{** โย ทุกฺขมทฺทกฺขิ ยโตนิทานํ, \\ กาเมสุ โส ชนฺตุ กถํ นเมยฺย. \\ กามา หิ โลเก สงฺโคติ ญตฺวา, \\ เตสํ สตีมา วินยาย สิกฺเขติ.}
\csromangathameasurewak{** โย ทุกฺขมทฺทกฺขิ ยโตนิทานํ, \\ กาเมสุ โส ชนฺตุ กถํ นเมยฺย. \\ กามา หิ โลเก สงฺโคติ ญตฺวา, \\ เตสํ สตีมา วินยาย สิกฺเขติ.}
\setlength{\csromangathaleftcol}{0pt}
\csromangathagroup{\csromangathastanza{\csromangathawak{** โย ทุกฺขมทฺทกฺขิ ยโตนิทานํ, \\ กาเมสุ โส ชนฺตุ กถํ นเมยฺย. \\ กามา หิ โลเก สงฺโคติ ญตฺวา, \\ เตสํ สตีมา วินยาย สิกฺเขติ.}}}

\prose{“โย ทุกฺขนฺ”ติ เววจน\-ปญฺญตฺติ จ ทุกฺขสฺส ปริญฺญา\-ปญฺญตฺติ จ. “ยโตนิทานนฺ”ติ ปภว\-ปญฺญตฺติ จ สมุทยสฺส ปหานปญฺญตฺติ จ. “อทฺทกฺขี”ติ เววจน\-ปญฺญตฺติ จ ญาณจกฺขุสฺส ปฏิเวธ\-ปญฺญตฺติ จ. “กาเมสุ โส ชนฺตุ กถํ นเมยฺยา”ติ เววจน\-ปญฺญตฺติ จ กามตณฺหาย อภินิเวส\-ปญฺญตฺติ จ. “กามา หิ โลเก สงฺโคติ ญตฺวา”ติ ปจฺจตฺถิกโต ทสฺสนปญฺญตฺติกามานํ. กามา หิ องฺคารกาสูปมา มํสเปสูปมา ปาวกกปฺปา ปปาตอุรโคปมา จ. \csromanfolio{53}“เตสํ สตีมา”ติ อปจย\-ปญฺญตฺติ ปหานาย, นิกฺเขป\-ปญฺญตฺติ กายคตาย สติยา, ภาวนา\-ปญฺญตฺติ มคฺคสฺส. “วินยาย สิกฺเข”ติ ปฏิเวธ\-ปญฺญตฺติ ราควินยสฺส โทสวินยสฺส โมหวินยสฺส. “ชนฺตู”ติ เววจน\-ปญฺญตฺติ โยคิสฺส. ยาทา หิ โยคี กามา สงฺโคติ ปชานาติ. โส กามานํ อนุปฺปาทาย กุสเล ธมฺเม อุปฺปาทยติ, โส อนุปฺปนฺนานํ กุสลานํ ธมฺมานํ อุปฺปาทาย วายมติ. อยํ วายามปญฺญตฺติ อปฺปตฺตสฺส ปตฺติยา. นิกฺเขป\-ปญฺญตฺติ โอรมตฺติกาย อสนฺตุฏฺฐิยา. ตตฺถ โส อุปฺปนฺนานํ กุสลานํ ธมฺมานํ ฐิติยา วายมตีติ อยํ อปฺปมาท\-ปญฺญตฺติ ภาวนาย, นิกฺเขป\-ปญฺญตฺติ วีริยิ\-นฺทฺริยสฺส, อารกฺข\-ปญฺญตฺติ กุสลานํ ธมฺมานํ, ฐิติปญฺญตฺติ อธิจิตฺต\-สิกฺขาย. เตนาห ภควา “โย ทุกฺขมทฺทกฺขิ ยโตนิทานนฺ”ติ.}

\setlength{\csromangathapagewidth}{0pt}
\setlength{\csromangathawakwidth}{0pt}
\csromangathameasuregroup{\textsuperscript{*}“โมหสมฺพนฺธโน โลโก, \\ อุปธิพนฺธโน\textsuperscript{1} พาโล, \\ อสฺสิรี วิย\textsuperscript{1} ขายติ,}{\csromangathabat{\textsuperscript{*}“โมหสมฺพนฺธโน โลโก,}{ภพฺพรูโปว ทิสฺสติ.} \\ \csromangathabat{อุปธิพนฺธโน\textsuperscript{1} พาโล,}{ตมสา ปริวาริโต.} \\ \csromangathabat{อสฺสิรี วิย\textsuperscript{1} ขายติ,}{ปสฺสโต นตฺถิ กิญฺจนนฺ”ติ.}}
\csromangathasetleft{\textsuperscript{*}“โมหสมฺพนฺธโน โลโก, \\ อุปธิพนฺธโน\textsuperscript{1} พาโล, \\ อสฺสิรี วิย\textsuperscript{1} ขายติ,}
\csromangathagroup{\csromangathastanza{\csromangathabat{\csromansymbolfootnote{*}{ขุ 1. 176 ปิฏฺเฐ อุทาเน.}“โมห\-สมฺพนฺธโน โลโก,}{ภพฺพรูโปว ทิสฺสติ.} \\ \csromangathabat{อุปธิพนฺธโน\csromansharedfootnote{5066d32b4fb44822}{อุปธิสมฺพนฺธโน (สี)} พาโล,}{ตมสา ปริวาริโต.} \\ \csromangathabat{อสฺสิรี วิย\csromansharedfootnote{5066d32b4fb44822}{อุปธิสมฺพนฺธโน (สี)} ขายติ,}{ปสฺสโต นตฺถิ กิญฺจนนฺ”ติ.}}}

\prose{“โมห\-สมฺพนฺธโน โลโก”ติ เทสนาปญฺญตฺติ วิปลฺลาสานํ. “ภพฺพรูโปว ทิสฺสตี”ติ วิปรีต\-ปญฺญตฺติ โลกสฺส. “อุปธิพนฺธโน พาโล”ติ ปภว\-ปญฺญตฺติ ปาปกานํ อิจฺฉาวจรานํ, กิจฺจปญฺญตฺติ ปริยุฏฺฐานานํ. พลวปญฺญตฺติ กิเลสานํ. วิรูหนา\-ปญฺญตฺติ สงฺขารานํ. “ตมสา ปริวาริโต”ติ เทสนาปญฺญตฺติ อวิชฺช\-นฺธการสฺส เววจน\-ปญฺญตฺติ จ. “อสฺสิรี วิย ขายตี”ติ ทสฺสน\-ปญฺญตฺติ ทิพฺพจกฺขุสฺส, นิกฺเขป\-ปญฺญตฺติ ปญฺญาจกฺขุสฺส. “ปสฺสโต นตฺถิ กิญฺจนนฺ”ติ ปฏิเวธ\-ปญฺญตฺติ สตฺตานํ, ราโค กิญฺจนํ โทโส กิญฺจนํ โมโห กิญฺจนํ. เตนาห ภควา “โมห\-สมฺพนฺธโน โลโก”ติ.}

\prose{\csromansymbolfootnote{+}{ขุ 1. 178 ปิฏฺเฐ อุทาเน.}“อตฺถิ ภิกฺขเว อชาตํ อภูตํ อกตํ อสงฺขตํ, โน เจตํ ภิกฺขเว อภวิสฺส อชาตํ อภูตํ อกตํ อสงฺขตํ. นยิธ\csromansharedfootnote{1045bb5338e30543}{น อิข (สี, ก)} ชาตสฺส ภูตสฺส กตสฺส สงฺขตสฺส นิสฺสรณํ ปญฺญาเยถ. ยสฺมา จ โข ภิกฺขเว อตฺถิ อชาตํ อภูตํ อกตํ อสงฺขตํ, ตสฺมา ชาตสฺส ภูตสฺส กตสฺส สงฺขตสฺส นิสฺสรณํ ปญฺญายตี”ติ.}

\prose{“โน เจตํ ภิกฺขเว อภวิสฺส อชาตํ อภูตํ อกตํ อสงฺขตนฺ”ติ เทสนาปญฺญตฺติ นิพฺพานสฺส เววจน\-ปญฺญตฺติ จ. “นยิธ ชาตสฺส ภูตสฺส กตสฺส สงฺขตสฺส นิสฺสรณํ ปญฺญาเยถา”ติ เววจน\-ปญฺญตฺติ สงฺขตสฺส \csromanfolio{54}อุปนยน\-ปญฺญตฺติ จ. “ยสฺมา จ โข ภิกฺขเว อตฺถิ อชาตํ อภูตํ อกตํ อสงฺขตนฺ”ติ เววจน\-ปญฺญตฺติ นิพฺพานสฺส โชตนาปญฺญตฺติ จ. “ตสฺมา ชาตสฺส ภูตสฺส กตสฺส สงฺขตสฺส นิสฺสรณํ ปญฺญายตี”ติ อยํ เววจน\-ปญฺญตฺติ นิพฺพานสฺส, นิยฺยานิก\-ปญฺญตฺติ มคฺคสฺส, นิสฺสรณ\-ปญฺญตฺติ สํสารโต. เตนาห ภควา “โน เจตํ ภิกฺขเว อภวิสฺสา”ติ. เตนาห อายสฺมา มหากจฺจายโน “เอกํ ภควา ธมฺมํ, ปญฺญตฺตีหิ วิวิธาหิ เทเสตี”ติ.}

\vspace{\tipitakaparskip}
\csromancenter{นิยุตฺโต ปญฺญตฺติหาโร.}
\csromanmatikaanchor{mtk.16}
\csromantocmarkref{section}{12. โอตรณหารวิภงฺค}{mtk.16}
\bookmark[dest={mtk.16},level=1]{12. โอตรณหารวิภงฺค}
\csromanheader{1}{12. \textbf{โอตรณ}\-\textbf{หาร}\-\textbf{วิภงฺค}}
\proseitem{42}{ตตฺถ กตโม โอตรโณ หาโร, “โย จ ปฏิจฺจุปฺปาโท”ติ.}

\setlength{\csromangathapagewidth}{0pt}
\setlength{\csromangathawakwidth}{0pt}
\csromangathameasuregroup{}{\textsuperscript{*}“อุทฺธํ อโธ สพฺพธิ วิปฺปมุตฺโต, \\ อยํ อหสฺมีติ\textsuperscript{1} อนานุปสฺสี. \\ เอวํ วิมุตฺโต อุทตาริ โอฆํ, \\ อติณฺณปุพฺพํ อปุนพฺภวายา”ติ.}
\csromangathameasurewak{\textsuperscript{*}“อุทฺธํ อโธ สพฺพธิ วิปฺปมุตฺโต, \\ อยํ อหสฺมีติ\textsuperscript{1} อนานุปสฺสี. \\ เอวํ วิมุตฺโต อุทตาริ โอฆํ, \\ อติณฺณปุพฺพํ อปุนพฺภวายา”ติ.}
\setlength{\csromangathaleftcol}{0pt}
\csromangathagroup{\csromangathastanza{\csromangathawak{\csromansymbolfootnote{*}{ขุ 1. 170 ปิฏฺเฐ อุทาเน; อุปริ 123, 184 ปิฏฺเฐสุปิ.}“อุทฺธํ อโธ สพฺพธิ วิปฺปมุตฺโต, \\ อยํ อหสฺมีติ\csromansharedfootnote{16a75c793444baf6}{อยมหมสฺมีติ (สี)} อนานุปสฺสี. \\ เอวํ วิมุตฺโต อุทตาริ โอฆํ, \\ อติณฺณปุพฺพํ อปุนพฺภวายา”ติ.}}}

\prose{“อุทฺธนฺ”ติ รูปธาตุ จ อรูปธาตุ จ. “อโธ”ติ กามธาตุ. “สพฺพธิวิปฺปมุตฺโต”ติ เตธาตุเก อยํ อเสกฺขาวิมุตฺติ. ตานิเยว อเสกฺขานิ ปญฺจินฺทฺริยานิ, อยํ อินฺทฺริเยหิ โอตรณา.}

\prose{ตานิเยว อเสกฺขานิ ปญฺจินฺทฺริยานิ วิชฺชา, วิชฺชุปฺปาทา อวิชฺชานิโรโธ, อวิชฺชานิโรธา สงฺขาร\-นิโรโธ, สงฺขาร\-นิโรธา วิญฺญาณนิโรโธ, วิญฺญาณนิโรธา นามรูป\-นิโรโธ, นามรูป\-นิโรธา สฬายตน\-นิโรโธ, สฬายตน\-นิโรธา ผสฺสนิโรโธ, ผสฺสนิโรธา เวทนานิโรโธ, เวทนานิโรธา ตณฺหานิโรโธ, ตณฺหานิโรธา อุปาทานนิโรโธ, อุปาทานนิโรธา ภวนิโรโธ, ภวนิโรธา ชาตินิโรโธ, ชาตินิโรธา ชรามรณํ โสก\-ปริเทว\-ทุกฺข\-โทมนสฺสุ\-ปายาสา นิรุชฺฌนฺติ. เอวเมตสฺส เกวลสฺส ทุกฺข\-กฺขนฺธสฺส นิโรโธ โหติ. อยํ ปฏิจฺจ\-สมุปฺปาเทหิ โอตรณา.}

\prose{\csromanfolio{55}ตานิเยว อเสกฺขานิ ปญฺจินฺทฺริยานิ ตีหิ ขนฺเธหิ สงฺคหิตานิ สีลกฺขนฺเธน สมาธิ\-กฺขนฺเธน ปญฺญา\-กฺขนฺเธน, อยํ ขนฺเธหิ โอตรณา.}

\prose{ตานิเยว อเสกฺขานิ ปญฺจินฺทฺริยานิ สงฺขาร\-ปริยาปนฺนานิ เย สงฺขารา อนาสวา, โน จ ภวงฺคา, เต สงฺขารา ธมฺมธาตุ\-สงฺคหิตา. อยํ ธาตูหิ โอตรณา.}

\prose{สา ธมฺมธาตุ ธมฺมายตน\-ปริยาปนฺนา, ยํ อายตนํ อนาสวํ, โน จ ภวงฺคํ. อยํ อายตเนหิ โอตรณา.}

\prose{“อยํ อหสฺมีติ อนานุปสฺสี”ติ อยํ สกฺกาย\-ทิฏฺฐิยา สมุคฺฆาโต, สา เสกฺขาวิมุตฺติ, ตานิเยว เสกฺขานิ ปญฺจินฺทฺริยานิ. อยํ อินฺทฺริเยหิ โอตรณา.}

\prose{ตานิเยว เสกฺขานิ ปญฺจินฺทฺริยานิ วิชฺชา, วิชฺชุปฺปาทา อวิชฺชานิโรโธ, อวิชฺชานิโรธา สงฺขาร\-นิโรโธ, เอวํ สพฺโพ ปฏิจฺจ\-สมุปฺปาโท. อยํ ปฏิจฺจ\-สมุปฺปาเทหิ โอตรณา.}

\prose{สาเยว วิชฺชา ปญฺญากฺขนฺโธ. อยํ ขนฺเธหิ โอตรณา.}

\prose{สาเยว วิชฺชา สงฺขาร\-ปริยาปนฺนา, เย สงฺขารา อนาสวา, โน จ ภวงฺคา, เต สงฺขารา ธมฺมธาตุ\-สงฺคหิตา, อยํ ธาตูหิ โอตรณา.}

\prose{สา ธมฺมธาตุ ธมฺมายตน\-ปริยาปนฺนา, ยํ อายตนํ อนาสวํ, โน จ ภวงฺคํ, อยํ อายตเนหิ โอตรณา.}

\prose{เสกฺขาย จ วิมุตฺติยา อเสกฺขาย จ วิมุตฺติยา วิมุตฺโต อุทตาริ โอฆํ อติณฺณปุพฺพํ อปุนพฺภวาย. เตนาห ภควา “อุทฺธํ อโธ”ติ.}

\proseitem{43}{“นิสฺสิตสฺส จลิตํ, อนิสฺสิตสฺส จลิตํ นตฺถิ, จลิเต อสติ ปสฺสทฺธิ, ปสฺสทฺธิยา สติ นติ น โหติ, นติยา อสติ อาคติคติ น โหติ, อาคติคติยา อสติ จุตูปปาโต น โหติ, จุตูปปาเต อสติ เนวิธ น หุรํ น อุภยม\-นฺตเรน เอเสวนฺโต ทุกฺขสฺสา”ติ\csromansharedfootnote{def20f33cb795b40}{ขุ 1. 179 ปิฏฺเฐ อุทาเน.}.}

\prose{“นิสฺสิตสฺส จลิตนฺ”ติ \textbf{นิสฺสโย} นาม ทุวิโธ ตณฺหานิสฺสโย จ ทิฏฺฐินิสฺสโย จ, ตตฺถ ยา รตฺตสฺส เจตนา, อยํ ตณฺหานิสฺสโย. ยา มูฬฺหสฺส เจตนา, อยํ ทิฏฺฐินิสฺสโย. เจตนา ปน สงฺขารา สงฺขาร\-ปจฺจยา \csromanfolio{56}วิญฺญาณํ, วิญฺญาณปจฺจยา นามรูปํ, เอวํ สพฺโพ ปฏิจฺจ\-สมุปฺปาโท. อยํ ปฏิจฺจ\-สมุปฺปาเทหิ โอตรณา.}

\prose{ตตฺถ ยา รตฺตสฺส เวทนา, อยํ สุขา เวทนา. ยา สมฺมูฬฺหสฺส เวทนา, อยํ อทุกฺขมสุขา เวทนา, อิมา เทฺว เวทนา เวทนากฺขนฺโธ. อยํ ขนฺเธหิ โอตรณา.}

\prose{ตตฺถ สุขา เวทนา เทฺว อินฺทฺริยานิ สุขินฺทฺริยํ โสมนสฺสินฺทฺริยญฺจ, อทุกฺขมสุขา เวทนา อุเปกฺขินฺทฺริยํ. อยํ อินฺทฺริเยหิ โอตรณา.}

\prose{ตานิเยว อินฺทฺริยานิ สงฺขาร\-ปริยาปนฺนานิ, เย สงฺขารา สาสวา ภวงฺคา, เต สงฺขารา ธมฺมธาตุ\-สงฺคหิตา. อยํ ธาตูหิ โอตรณา.}

\prose{สา ธมฺมธาตุ ธมฺมายตน\-ปริยาปนฺนา, ยํ อายตนํ สาสวํ ภวงฺคํ, อยํ อายตเนหิ โอตรณา.}

\prose{“อนิสฺสิตสฺส จลิตํ นตฺถี”ติ สมถวเสน วา ตณฺหาย อนิสฺสิโต วิปสฺสนา\-วเสน วา ทิฏฺฐิยา อนิสฺสิโต. ยา วิปสฺสนา อยํ วิชฺชา, วิชฺชุปฺปาทา อวิชฺชานิโรโธ, อวิชฺชานิโรธา สงฺขาร\-นิโรโธ, สงฺขาร\-นิโรธา วิญฺญาณนิโรโธ, เอวํ สพฺโพ ปฏิจฺจ\-สมุปฺปาโท. อยํ ปฏิจฺจ\-สมุปฺปาเทหิ โอตรณา.}

\prose{สาเยว วิปสฺสนา ปญฺญากฺขนฺโธ. อยํ ขนฺเธหิ โอตรณา.}

\prose{สาเยว วิปสฺสนา เทฺว อินฺทฺริยานิ วีริยินฺทฺริยญฺจ ปญฺญินฺทฺริยญฺจ. อยํ อินฺทฺริเยหิ โอตรณา.}

\prose{สาเยว วิปสฺสนา สงฺขาร\-ปริยาปนฺนา, เย สงฺขารา อนาสวา, โน จ ภวงฺคา, เต สงฺขารา ธมฺมธาตุ\-สงฺคหิตา. อยํ ธาตูหิ โอตรณา.}

\prose{สา ธมฺมธาตุ ธมฺมายตน\-ปริยาปนฺนา, ยํ อายตนํ อนาสวํ, โน จ ภวงฺคํ. อยํ อายตเนหิ โอตรณา.}

\prose{“ปสฺสทฺธิยา สตี”ติ ทุวิธา \textbf{ปสฺสทฺธิ} กายิกา จ เจตสิกา จ. ยํ กายิกํ สุขํ, อยํ กายปสฺสทฺธิ. ยํ เจตสิกํ สุขํ, อยํ เจตสิกา ปสฺสทฺธิ. ปสฺสทฺธกาโย สุขํ เวทิยติ\csromansharedfootnote{00a1975d0a900c5a}{เวทยติ (ก)}, สุขิโน จิตฺตํ สมาธิยติ, สมาหิโต ยถาภูตํ ปชานาติ, ยถาภูตํ ปชานนฺโต นิพฺพินฺทติ, \csromanfolio{57}นิพฺพินฺทนฺโต วิรชฺชติ, วิราคา วิมุจฺจติ, วิมุตฺตสฺมิํ “วิมุตฺตมฺ”อิติ\csromansharedfootnote{73e2a9629da4d584}{วิมุตฺตมฺหีติ (สี, ก)} ญาณํ โหติ, “ขีณา ชาติ, วุสิตํ พฺรหฺมจริยํ, กตํ กรณียํ, นาปรํ อิตฺถตฺตายา”ติ ปชานาติ. โส น นมติ รูเปสุ, น สทฺเทสุ, น คนฺเธสุ, น รเสสุ, น โผฏฺฐพฺเพสุ, น ธมฺเมสุ ขยา ราคสฺส ขยา โทสสฺส ขยา โมหสฺส เยน รูเปน ตถาคตํ ตฏฺฐนฺตํ จรนฺตํ ปญฺญาปยมาโน ปญฺญาเปยฺย, ตสฺส รูปสฺส ขยา วิราคา นิโรธา จาคา ปฏินิสฺสคฺคา รูปสงฺขเย วิมุตฺโต, ตถาคโต อตฺถีติปิ น อุเปติ, นตฺถีติปิ น อุเปติ, อตฺถิ นตฺถีติปิ น อุเปติ, เนวตฺถิ โน นตฺถีติปิ น อุเปติ. อถ โข คมฺภีโร อปฺปเมยฺโย อสงฺเขยฺโย นิพฺพุโตติเยว สงฺขํ คจฺฉติ ขยา ราคสฺส, ขยา โทสสฺส, ขยา โมหสฺส.}

\prose{ยาย เวทนาย ฯเปฯ. ยาย สญฺญาย. เยหิ สงฺขาเรหิ. เยน วิญฺญาเณน ตถาคตํ ติฏฺฐนฺตํ จรนฺตํ ปญฺญาปยมาโน ปญฺญาเปยฺย, ตสฺส วิญฺญาณสฺส ขยา วิราคา นิโรธา จาคา ปฏินิสฺสคฺคา วิญฺญาณ\-สงฺขเย วิมุตฺโต, ตถาคโต อตฺถีติปิ น อุเปติ, นตฺถีติปิ น อุเปติ, อตฺถิ นตฺถีติปิ น อุเปติ, เนวตฺถิ โน นตฺถีติปิ น อุเปติ. อถ โข คมฺภีโร อปฺปเมยฺโย อสงฺเขยฺโย นิพฺพุโตติเยว สงฺขํ คจฺฉติ ขยา ราคสฺส, ขยา โทสสฺส, ขยา โมหสฺส. “อาคตี”ติ อิธาคติ. “คตี”ติ เปจฺจภโว. อาคติคตีปิ น ภวนฺติ, “เนวิธา”ติ ฉสุ อชฺฌตฺติเกสุ อายตเนสุ. “น หุรนฺ”ติ ฉสุ พาหิเรสุ อายตเนสุ. “น อุภยม\-นฺตเรนา”ติ ผสฺส\-สมุทิเตสุ ธมฺเมสุ อตฺตานํ น ปสฺสติ. “เอเสวนฺโต ทุกฺขสฺสา”ติ ปฏิจฺจ\-สมุปฺปาโท, โส ทุวิโธ โลกิโย จ โลกุตฺตโร จ. ตตฺถ โลกิโย อวิชฺชาปจฺจยา สงฺขารา, ยาว ชรามรณา. โลกุตฺตโร สีลวโต อวิปฺปฏิสาโร ชายติ, ยาว นาปรํ อิตฺถตฺตายาติ ปชานาติ. เตนาห ภควา “นิสฺสิตสฺส จลิตํ อนิสฺสิตสฺส จลิตํ นตฺถิ ฯเปฯ เอเสวนฺโต ทุกฺขสฺสา”ติ.}

\setlength{\csromangathapagewidth}{0pt}
\setlength{\csromangathawakwidth}{0pt}
\csromangathameasuregroup{}{“เย เกจิ โสกา ปริเทวิตา วา, \\ ทุกฺขา\textsuperscript{1} จ โลกสฺมิมเนกรูปา. \\ ปิยํ ปฏิจฺจปฺปภวนฺติ เอเต, \\ ปิเย อสนฺเต น ภวนฺติ เอเต\textsuperscript{1}. \\ ตสฺมา หิ เต สุขิโน วีตโสกา, \\ เยสํ ปิยํ นตฺถิ กุหิญฺจิ โลเก. \\ ตสฺมา อโสกํ วิรชํ ปตฺถยาโน, \\ ปิยํ น กยิราถ กุหิญฺจิ โลเก”ติ\textsuperscript{1}.}
\csromangathameasurewak{“เย เกจิ โสกา ปริเทวิตา วา, \\ ทุกฺขา\textsuperscript{1} จ โลกสฺมิมเนกรูปา. \\ ปิยํ ปฏิจฺจปฺปภวนฺติ เอเต, \\ ปิเย อสนฺเต น ภวนฺติ เอเต\textsuperscript{1}. \\ ตสฺมา หิ เต สุขิโน วีตโสกา, \\ เยสํ ปิยํ นตฺถิ กุหิญฺจิ โลเก. \\ ตสฺมา อโสกํ วิรชํ ปตฺถยาโน, \\ ปิยํ น กยิราถ กุหิญฺจิ โลเก”ติ\textsuperscript{1}.}
\setlength{\csromangathaleftcol}{0pt}
\csromangathagroup{\csromangathastanza{\csromangathawak{\csromangathaitemline{44}{“เย เกจิ โสกา ปริเทวิตา วา,} \\ \csromangathacontline{ทุกฺขา\csromansharedfootnote{bd9c7704c614c433}{ทุกฺขํ (สี, ก)} จ โลกสฺมิม\-เนกรูปา.} \\ \csromangathacontline{ปิยํ ปฏิจฺจ\-ปฺปภวนฺติ เอเต,} \\ \csromangathacontline{ปิเย อสนฺเต น ภวนฺติ เอเต\csromansharedfootnote{91e538fd1766ea0c}{ขุ 1. 191 ปิฏฺเฐ; อุปริ 178 ปิฏฺเฐปิ.}.}}}\csromangathastanza{\csromangathawak{\csromanfolio{58}ตสฺมา หิ เต สุขิโน วีตโสกา, \\ เยสํ ปิยํ นตฺถิ กุหิญฺจิ โลเก. \\ ตสฺมา อโสกํ วิรชํ ปตฺถยาโน, \\ ปิยํ น กยิราถ กุหิญฺจิ โลเก”ติ\csromansharedfootnote{ad3354cb2dd3e12e}{ขุ 1. 191 ปิฏฺเฐ อุทาเน.}.}}}

\prose{“เย เกจิ โสกา ปริเทวิตา วา, ทุกฺขา จ โลกสฺมิม\-เนกรูปา ปิยํ ปฏิจฺจ\-ปฺปภวนฺติ เอเต”ติ\csromandash{}อยํ ทุกฺขา เวทนา. “ปิเย อสนฺเต น ภวนฺติ เอเต”ติ\csromandash{}อยํ สุขา เวทนา. เวทนา เวทนากฺขนฺโธ. อยํ ขนฺเธหิ โอตรณา.}

\prose{เวทนาปจฺจยา ตณฺหา, ตณฺหาปจฺจยา อุปาทานํ, อุปาทานปจฺจยา ภโว, ภวปจฺจยา ชาติ, ชาติปจฺจยา ชรามรณํ, เอวํ สพฺพํ. อยํ ปฏิจฺจ\-สมุปฺปาเทหิ โอตรณา.}

\prose{ตตฺถ สุขา เวทนา เทฺว อินฺทฺริยานิ สุขินฺทฺริยํ โสมนสฺสินฺทฺริยญฺจ. ทุกฺขา เวทนา เทฺว อินฺทฺริยานิ ทุกฺขินฺทฺริยํ โทมนสฺสินฺทฺริยญฺจ. อยํ อินฺทฺริเยหิ โอตรณา.}

\prose{ตานิเยว อินฺทฺริยานิ สงฺขาร\-ปริยาปนฺนานิ, เย สงฺขารา สาสวา ภวงฺคา, เต สงฺขารา ธมฺมธาตุ\-สงฺคหิตา. อยํ ธาตูหิ โอตรณา.}

\prose{สา ธมฺมธาตุ ธมฺมายตน\-ปริยาปนฺนา, ยํ อายตนํ สาสวํ ภวงฺคํ. อยํ อายตเนหิ โอตรณา.}

\setlength{\csromangathapagewidth}{0pt}
\setlength{\csromangathawakwidth}{0pt}
\csromangathameasuregroup{}{ตสฺมา หิ เต สุขิโน วีตโสกา, \\ เยสํ ปิยํ นตฺถิ กุหิญฺจิ โลเก. \\ ตสฺมา อโสกํ วิรชํ ปตฺถยาโน, \\ ปิยํ น กยิราถ กุหิญฺจิ โลเกติ.}
\csromangathameasurewak{ตสฺมา หิ เต สุขิโน วีตโสกา, \\ เยสํ ปิยํ นตฺถิ กุหิญฺจิ โลเก. \\ ตสฺมา อโสกํ วิรชํ ปตฺถยาโน, \\ ปิยํ น กยิราถ กุหิญฺจิ โลเกติ.}
\setlength{\csromangathaleftcol}{0pt}
\csromangathagroup{\csromangathastanza{\csromangathawak{ตสฺมา หิ เต สุขิโน วีตโสกา, \\ เยสํ ปิยํ นตฺถิ กุหิญฺจิ โลเก. \\ ตสฺมา อโสกํ วิรชํ ปตฺถยาโน, \\ ปิยํ น กยิราถ กุหิญฺจิ โลเกติ.}}}

\prose{อิทํ ตณฺหาปหานํ. ตณฺหานิโรธา อุปาทานนิโรโธ, อุปาทานนิโรธา ภวนิโรโธ, เอวํ สพฺพํ. อยํ ปฏิจฺจ\-สมุปฺปาเทหิ โอตรณา.}

\prose{ตํเยว ตณฺหาปหานํ สมโถ. โส สมโถ เทฺว อินฺทฺริยานิ สตินฺทฺริยํ สมาธินฺทฺริยญฺจ. อยํ อินฺทฺริเยหิ โอตรณา.}

\prose{โสเยว สมโถ สมาธิ\-กฺขนฺโธ. อยํ ขนฺเธหิ โอตรณา.}

\prose{โสเยว สมโถ สงฺขาร\-ปริยาปนฺโน, เย สงฺขารา อนาสวา, โน จ ภวงฺคา, เต สงฺขารา ธมฺมธาตุ\-สงฺคหิตา. อยํ ธาตูหิ โอตรณา.}

\prose{\csromanfolio{59}สา ธมฺมธาตุ ธมฺมายตน\-ปริยาปนฺนา, ยํ อายตนํ อนาสวํ, โน จ ภวงฺคํ. อยํ อายตเนหิ โอตรณา. เตนาห ภควา “เย เกจิ โสกา”ติ.}

\setlength{\csromangathapagewidth}{0pt}
\setlength{\csromangathawakwidth}{0pt}
\csromangathameasuregroup{กามํ กามยมานสฺส, \\ อทฺธา ปีติมโน โหติ, \\ ตสฺส เจ กามยานสฺส, \\ เต กามา ปริหายนฺติ, \\ โย กาเม ปริวชฺเชติ, \\ โสมํ วิสตฺติกํ โลเก,}{\csromangathabat{กามํ กามยมานสฺส,}{ตสฺส เจ ตํ สมิชฺฌติ.} \\ \csromangathabat{อทฺธา ปีติมโน โหติ,}{ลทฺธา มจฺโจ ยทิจฺฉติ.} \\ \csromangathabat{ตสฺส เจ กามยานสฺส,}{ฉนฺทชาตสฺส ชนฺตุโน.} \\ \csromangathabat{เต กามา ปริหายนฺติ,}{สลฺลวิทฺโธว รุปฺปติ.} \\ \csromangathabat{โย กาเม ปริวชฺเชติ,}{สปฺปสฺเสว ปทา สิโร.} \\ \csromangathabat{โสมํ วิสตฺติกํ โลเก,}{สโต สมติวตฺตตีติ\textsuperscript{1}.}}
\csromangathasetleft{กามํ กามยมานสฺส, \\ อทฺธา ปีติมโน โหติ, \\ ตสฺส เจ กามยานสฺส, \\ เต กามา ปริหายนฺติ, \\ โย กาเม ปริวชฺเชติ, \\ โสมํ วิสตฺติกํ โลเก,}
\csromangathagroup{\csromangathastanza{\csromangathabat{กามํ กามยมานสฺส,}{ตสฺส เจ ตํ สมิชฺฌติ.} \\ \csromangathabat{อทฺธา ปีติมโน โหติ,}{ลทฺธา มจฺโจ ยทิจฺฉติ.}}\csromangathastanza{\csromangathabat{ตสฺส เจ กามยานสฺส,}{ฉนฺทชาตสฺส ชนฺตุโน.} \\ \csromangathabat{เต กามา ปริหายนฺติ,}{สลฺลวิทฺโธว รุปฺปติ.}}\csromangathastanza{\csromangathabat{โย กาเม ปริวชฺเชติ,}{สปฺปสฺเสว ปทา สิโร.} \\ \csromangathabat{โสมํ วิสตฺติกํ โลเก,}{สโต สมติ\-วตฺตตีติ\csromansharedfootnote{97ffbdfb2c34455a}{เหฏฺฐา 6 ปิฏฺเฐปิ.}.}}}

\prose{ตตฺถ ยา ปีติมนตา, อยํ อนุนโย. ยทาห สลฺลวิทฺโธว รุปฺปตีติ, อิทํ ปฏิฆํ. อนุนยํ ปฏิฆญฺจ ปน ตณฺหาปกฺโข, ตณฺหาย จ ปน ทสรูปีนิ อายตนานิ ปทฏฺฐานํ. อยํ อายตเนหิ โอตรณา.}

\prose{ตานิเยว ทส รูปีนิ รูปกาโย นามสมฺปยุตฺโต, ตทุภยํ นามรูปํ, นามรูป\-ปจฺจยา สฬายตนํ, สฬายตน\-ปจฺจยา ผสฺโส, ผสฺสปจฺจยา เวทนา, เวทนาปจฺจยา ตณฺหา, เอวํ สพฺพํ. อยํ ปฏิจฺจ\-สมุปฺปาเทหิ โอตรณา.}

\prose{ตเทว นามรูปํ ปญฺจกฺขนฺโธ. อยํ ขนฺเธหิ โอตรณา.}

\prose{ตเทว นามรูปํ อฏฺฐารส ธาตุโย. อยํ ธาตูหิ โอตรณา.}

\prose{ตตฺถ โย รูปกาโย อิมานิ ปญฺจ รูปีนิ อินฺทฺริยานิ, โย นามกาโย อิมานิ ปญฺจ อรูปีนิ อินฺทฺริยานิ, อิมานิ ทส อินฺทฺริยานิ. อยํ อินฺทฺริเยหิ โอตรณา.}

\prose{ตตฺถ ยทาห\csromandash{}}

\setlength{\csromangathapagewidth}{0pt}
\setlength{\csromangathawakwidth}{0pt}
\csromangathameasuregroup{“โย กาเม ปริวชฺเชติ, \\ โสมํ วิสตฺติกํ โลเก,}{\csromangathabat{“โย กาเม ปริวชฺเชติ,}{สปฺปสฺเสว ปทา สิโร.} \\ \csromangathabat{โสมํ วิสตฺติกํ โลเก,}{สโต สมติวตฺตตี”ติ.}}
\csromangathasetleft{“โย กาเม ปริวชฺเชติ, \\ โสมํ วิสตฺติกํ โลเก,}
\csromangathagroup{\csromangathastanza{\csromangathabat{“โย กาเม ปริวชฺเชติ,}{สปฺปสฺเสว ปทา สิโร.} \\ \csromangathabat{โสมํ วิสตฺติกํ โลเก,}{สโต สมติวตฺตตี”ติ.}}}

\prose{อยํ สอุปาทิเสสา นิพฺพานธาตุ, อยํ ธาตูหิ โอตรณา.}

\prose{สาเยว สอุปาทิเสสา นิพฺพานธาตุ วิชฺชา, วิชฺชุปฺปาทา อวิชฺชานิโรโธ, อวิชฺชานิโรธา สงฺขาร\-นิโรโธ, เอวํ สพฺพํ. อยํ ปฏิจฺจ\-สมุปฺปาเทหิ โอตรณา.}

\csromanmatikaanchor{mtk.17}
\csromantocmarkref{section}{13. โสทนหารวิภงฺค}{mtk.17}
\bookmark[dest={mtk.17},level=1]{13. โสทนหารวิภงฺค}
\prose{\csromanfolio{60}สาเยว วิชฺชา ปญฺญากฺขนฺโธ. อยํ ขนฺเธหิ โอตรณา.}

\prose{สาเยว วิชฺชา เทฺว อินฺทฺริยานิ วีริยินฺทฺริยํ ปญฺญินฺทฺริยญฺจ. อยํ อินฺทฺริเยหิ โอตรณา.}

\prose{สาเยว วิชฺชา สงฺขาร\-ปริยาปนฺนา, เย สงฺขารา อนาสวา, โน จ ภวงฺคา, เต สงฺขารา ธมฺมธาตุ\-สงฺคหิตา. อยํ ธาตูหิ โอตรณา.}

\prose{สา ธมฺมธาตุ ธมฺมายตน\-ปริยาปนฺนา, ยํ อายตนํ อนาสวํ, โน จ ภวงฺคํ. อยํ อายตเนหิ โอตรณา. เตนาห ภควา “กามํ กามยมานสฺสา”ติ.}

\prose{เอตฺตาวตา ปฏิจฺจ อินฺทฺริย\-ขนฺธธาตุอายตนานิ สโมสรโณ\-ตรณานิ ภวนฺติ. เอวํ ปฏิจฺจ อินฺทฺริย\-ขนฺธธาตุอายตนานิ โอตาเรตพฺพานิ. เตนาห อายสฺมา มหากจฺจายโน “โย จ ปฏิจฺจุปฺปาโท”ติ.}

\vspace{\tipitakaparskip}
\csromancenter{นิยุตฺโต โอตรโณ หาโร.}
\csromanheader{2}{13. \textbf{โสธน}\-\textbf{หาร}\-\textbf{วิภงฺค}}
\proseitem{45}{ตตฺถ กตโม โสธโน หาโร, “วิสฺสชฺชิตมฺหิ ปเญฺห”ติ คาถา. ยถา อายสฺมา อชิโต ปารายเน ภควนฺตํ ปญฺหํ ปุจฺฉติ\csromandash{}}

\setlength{\csromangathapagewidth}{0pt}
\setlength{\csromangathawakwidth}{0pt}
\csromangathameasuregroup{\textsuperscript{*}“เกนสฺสุ นิวุโต โลโก, \\ กิสฺสาภิเลปนํ พฺรูสิ,}{\csromangathabat{\textsuperscript{*}“เกนสฺสุ นิวุโต โลโก,}{เกนสฺสุ นปฺปกาสติ.} \\ \csromangathabat{กิสฺสาภิเลปนํ พฺรูสิ,}{กิํสุ ตสฺส มหพฺภยนฺ”ติ.}}
\csromangathasetleft{\textsuperscript{*}“เกนสฺสุ นิวุโต โลโก, \\ กิสฺสาภิเลปนํ พฺรูสิ,}
\csromangathagroup{\csromangathastanza{\csromangathaitembat{45}{\csromansymbolfootnote{*}{เหฏฺฐา 10, 11 ปิฏฺเฐสุ; อุปริ จ 177, 183 ปิฏฺเฐสุปิ.}“เกนสฺสุ นิวุโต โลโก,}{เกนสฺสุ นปฺปกาสติ.} \\ \csromangathacontbat{กิสฺสาภิเลปนํ พฺรูสิ,}{กิํสุ ตสฺส มหพฺภยนฺ”ติ.}}}

\prose{\csromansymbolmark{*}“อวิชฺชาย นิวุโต โลโก, (อชิตาติ ภควา,)}

\prose{วิวิจฺฉา ปมาทา นปฺปกาสติ.}

\setlength{\csromangathapagewidth}{0pt}
\setlength{\csromangathawakwidth}{0pt}
\csromangathameasuregroup{ชปฺปาภิเลปนํ พฺรูมิ,}{\csromangathabat{ชปฺปาภิเลปนํ พฺรูมิ,}{ทุกฺขมสฺส มหพฺภยนฺ”ติ.}}
\csromangathasetleft{ชปฺปาภิเลปนํ พฺรูมิ,}
\csromangathagroup{\csromangathastanza{\csromangathabat{ชปฺปา\-ภิเลปนํ พฺรูมิ,}{ทุกฺขมสฺส มหพฺภยนฺ”ติ.}}}

\prose{“เกนสฺสุ นิวุโต โลโก”ติ ปเญฺห “อวิชฺชาย นิวุโต โลโก”ติ ภควา ปทํ โสเธติ, โน จ อารมฺภํ. “เกนสฺสุ นปฺปกาสตี”ติ ปเญฺห “วิวิจฺฉา ปมาทา นปฺปกาสตี”ติ ภควา ปทํ โสเธติ, โน จ อารมฺภํ. “กิสฺสาภิเลปนํ พฺรูสี”ติ ปเญฺห “ชปฺปา\-ภิเลปนํ พฺรูมี”ติ ภควา ปทํ โสเธติ, โน จ อารมฺภํ. “กิํสุ ตสฺส มหพฺภยนฺ”ติ ปเญฺห “ทุกฺขมสฺส มหพฺภยนฺ”ติ สุทฺโธ อารมฺโภ. เตนาห ภควา “อวิชฺชาย นิวุโต โลโก”ติ.}

\prose{\csromanfolio{61}\csromansymbolfootnote{*}{เหฏฺฐา 12 ปิฏฺเฐปิ. ** เหฏฺฐา 13 ปิฏฺเฐ; อุปริ 180, 225 ปิฏฺเฐสุปิ.}“สวนฺติ สพฺพธิ โสตา, (อิจฺจายสฺมา อชิโต,)}

\prose{โสตานํ กิํ นิวารณํ.}

\setlength{\csromangathapagewidth}{0pt}
\setlength{\csromangathawakwidth}{0pt}
\csromangathameasuregroup{โสตานํ สํวรํ พฺรูหิ,}{\csromangathabat{โสตานํ สํวรํ พฺรูหิ,}{เกน โสตา ปิธียเร”ติ.}}
\csromangathasetleft{โสตานํ สํวรํ พฺรูหิ,}
\csromangathagroup{\csromangathastanza{\csromangathabat{โสตานํ สํวรํ พฺรูหิ,}{เกน โสตา ปิธียเร”ติ.}}}

\prose{** “ยานิ โสตานิ โลกสฺมิํ, (อชิตาติ ภควา,)}

\prose{สติ เตสํ นิวารณํ.}

\setlength{\csromangathapagewidth}{0pt}
\setlength{\csromangathawakwidth}{0pt}
\csromangathameasuregroup{โสตานํ สํวรํ พฺรูมิ,}{\csromangathabat{โสตานํ สํวรํ พฺรูมิ,}{ปญฺญาเยเต ปิธียเร”ติ.}}
\csromangathasetleft{โสตานํ สํวรํ พฺรูมิ,}
\csromangathagroup{\csromangathastanza{\csromangathabat{โสตานํ สํวรํ พฺรูมิ,}{ปญฺญาเยเต ปิธียเร”ติ.}}}

\prose{“สวนฺติ สพฺพธิ โสตา, โสตานํ กิํ นิวารณนฺ”ติ ปเญฺห “ยานิ โสตานิ โลกสฺมิํ, สติ เตสํ นิวารณนฺ”ติ ภควา ปทํ โสเธติ, โน จ อารมฺภํ. “โสตานํ สํวรํ พฺรูหิ, เกน โสตา ปิธียเร”ติ ปเญฺห โสตานํ สํวรํ พฺรูมิ, ปญฺญาเยเต ปิธียเร”ติ สุทฺโธ อารมฺโภ. เตนาห ภควา “ยานิ โสตานิ โลกสฺมินฺ”ติ.}

\prose{\csromansymbolfootnote{+}{เหฏฺฐา 14, 15 ปิฏฺเฐสุปิ.}“ปญฺญา เจว สติ จ, (อิจฺจายสฺมา อชิโต,)}

\prose{นามรูปญฺจ มาริส.}

\setlength{\csromangathapagewidth}{0pt}
\setlength{\csromangathawakwidth}{0pt}
\csromangathameasuregroup{เอตํ เม ปุฏฺโฐ ปพฺรูหิ,}{\csromangathabat{เอตํ เม ปุฏฺโฐ ปพฺรูหิ,}{กตฺเถตํ อุปรุชฺฌตี”ติ.}}
\csromangathasetleft{เอตํ เม ปุฏฺโฐ ปพฺรูหิ,}
\csromangathagroup{\csromangathastanza{\csromangathabat{เอตํ เม ปุฏฺโฐ ปพฺรูหิ,}{กตฺเถตํ อุปรุชฺฌตี”ติ.}}}

\prose{ปเญฺห\csromandash{}}

\setlength{\csromangathapagewidth}{0pt}
\setlength{\csromangathawakwidth}{0pt}
\csromangathameasuregroup{\textsuperscript{+}“ยเมตํ ปญฺหํ อปุจฺฉิ, \\ ยตฺถ นามญฺจ รูปญฺจ, \\ วิญฺญาณสฺส นิโรเธน,}{\csromangathabat{\textsuperscript{+}“ยเมตํ ปญฺหํ อปุจฺฉิ,}{อชิต ตํ วทามิ เต.} \\ \csromangathabat{ยตฺถ นามญฺจ รูปญฺจ,}{อเสสํ อุปรุชฺฌติ.} \\ \csromangathabat{วิญฺญาณสฺส นิโรเธน,}{เอตฺเถตํ อุปรุชฺฌตี”ติ.}}
\csromangathasetleft{\textsuperscript{+}“ยเมตํ ปญฺหํ อปุจฺฉิ, \\ ยตฺถ นามญฺจ รูปญฺจ, \\ วิญฺญาณสฺส นิโรเธน,}
\csromangathagroup{\csromangathastanza{\csromangathabat{\csromansymbolmark{+}“ยเมตํ ปญฺหํ อปุจฺฉิ,}{อชิต ตํ วทามิ เต.} \\ \csromangathabat{ยตฺถ นามญฺจ รูปญฺจ,}{อเสสํ อุปรุชฺฌติ.} \\ \csromangathabat{วิญฺญาณสฺส นิโรเธน,}{เอตฺเถตํ อุปรุชฺฌตี”ติ.}}}

\prose{สุทฺโธ อารมฺโภ. เตนาห ภควา “ยเมตํ ปญฺหํ อปุจฺฉี”ติ. ยตฺถ เอวํ สุทฺโธ อารมฺโภ, โส ปโญฺห วิสชฺชิโต ภวติ. ยตฺถ ปน อารมฺโภ อสุทฺโธ, น ตาว โส ปโญฺห วิสชฺชิโต ภวติ. เตนาห อายสฺมา มหากจฺจายโน “วิสฺสชฺชิตมฺหิ ปเญฺห”ติ.}

\vspace{\tipitakaparskip}
\csromancenter{นิยุตฺโต โสธโน หาโร.}
\csromanmatikaanchor{mtk.18}
\csromantocmarkref{section}{14. อธิฏฺฐานหารวิภงฺค}{mtk.18}
\bookmark[dest={mtk.18},level=1]{14. อธิฏฺฐานหารวิภงฺค}
\csromanheader{1}{14. \textbf{อธิฏฺฐาน}\-\textbf{หาร}\-\textbf{วิภงฺค}}
\proseitem{46}{ตตฺถ กตโม อธิฏฺฐาโน หาโร, “เอกตฺตตาย ธมฺมา, เยปิ จ เวมตฺตตาย นิทฺทิฏฺฐา”ติ.}

\prose{เย ตตฺถ นิทฺทิฏฺฐา, ตถา เต ธารยิตพฺพา.}

\prose{\csromanfolio{62}“ทุกฺขนฺ”ติ เอกตฺตตา. ตตฺถ กตมํ ทุกฺขํ, ชาติ ทุกฺขา, ชรา ทุกฺขา, พฺยาธิ ทุกฺโข, มรณํ ทุกฺขํ, อปฺปิเยหิ สมฺปโยโค ทุกฺโข, ปิเยหิ วิปฺปโยโค ทุกฺโข, ยมฺปิจฺฉํ น ลภติ ตมฺปิ ทุกฺขํ, สํขิตฺเตน ปญฺจุ\-ปาทาน\-กฺขนฺธา ทุกฺขา, รูปา ทุกฺขา, เวทนา ทุกฺขา, สญฺญา ทุกฺขา, สงฺขารา ทุกฺขา, วิญฺญาณํ ทุกฺขํ. อยํ เวมตฺตตา.}

\prose{“ทุกฺขสมุทโย”ติ เอกตฺตตา. ตตฺถ กตโม ทุกฺขสมุทโย, ยายํ ตณฺหา โปโนภวิกา\csromansharedfootnote{cb07fa3d6c2ab9a9}{โปโนพฺภวิกา (ก)} นนฺที\-ราค\-สหคตา ตตฺร\-ตตฺรา\-ภินนฺทินี. เสยฺยถิทํ, กามตณฺหา ภวตณฺหา วิภวตณฺหา. อยํ เวมตฺตตา.}

\prose{“ทุกฺขนิโรโธ”ติ เอกตฺตตา. ตตฺถ กตโม ทุกฺขนิโรโธ, โย ตสฺสาเยว ตณฺหาย อเสส\-วิราค\-นิโรโธ จาโค ปฏินิสฺสคฺโค มุตฺติ อนาลโย. อยํ เวมตฺตตา.}

\prose{“ทุกฺขนิโรธ\-คามินี ปฏิปทา”ติ เอกตฺตตา. ตตฺถ กตมา ทุกฺขนิโรธ\-คามินี ปฏิปทา, อยเมว อริโย อฏฺฐงฺคิโก มคฺโค. เสยฺยถิทํ, สมฺมาทิฏฺฐิ สมฺมาสงฺกปฺโป สมฺมาวาจา สมฺมากมฺมนฺโต สมฺมาอาชีโว สมฺมาวายาโม สมฺมาสติ สมฺมาสมาธิ. อยํ เวมตฺตตา.}

\prose{“มคฺโค”ติ เอกตฺตตา. ตตฺถ กตโม มคฺโค, นิรยคามี มคฺโค ติรจฺฉานโยนิ\-คามี มคฺโค เปตฺติวิสย\-คามี มคฺโค อสุรโยนิโย\csromansharedfootnote{5c146ad6315faec4}{อสุรโยนิคามิโย (สี), อสุรโยนิคามีนิโย (ก)} มคฺโค สคฺคคามิโย มคฺโค มนุสฺสคามี มคฺโค นิพฺพานคามี มคฺโค. อยํ เวมตฺตตา.}

\prose{“นิโรโธ”ติ เอกตฺตตา. ตตฺถ กตโม นิโรโธ, ปฏิสงฺขา\-นิโรโธ อปฺปฏิสงฺขา\-นิโรโธ อนุนยนิโรโธ ปฏิฆนิโรโธ มานนิโรโธ มกฺขนิโรโธ ปฬาสนิโรโธ อิสฺสานิโรโธ มจฺฉริย\-นิโรโธ สพฺพกิเลส\-นิโรโธ. อยํ เวมตฺตตา.}

\prose{“รูปนฺ”ติ เอกตฺตตา. ตตฺถ กตมํ รูปํ, จาตุมหาภูติกํ\csromansharedfootnote{9d8daf8048cc3e64}{จาตุมฺมหาภูติกํ (สี)} รูปํ จตุนฺนํ มหาภูตานํ อุปาทาย รูปสฺส ปญฺญตฺติ. ตตฺถ กตมานิ จตฺตาริ มหาภูตานิ, ปถวีธาตุ\csromansharedfootnote{10aff1a7188644ad}{ปฐวีธาตุ (สี)} อาโปธาตุ เตโชธาตุ วาโยธาตุ.}

\proseitem{47}{ทฺวีหิ อากาเรหิ ธาตุโย ปริคฺคณฺหาติ สงฺเขเปน จ วิตฺถาเรน จ. กถํ วิตฺถาเรน ธาตุโย ปริคฺคณฺหาติ, วีสติยา อากาเรหิ \csromanfolio{63}ปถวีธาตุํ วิตฺถาเรน ปริคฺคณฺหาติ, ทฺวาทสหิ อากาเรหิ อาโปธาตุํ วิตฺถาเรน ปริคฺคณฺหาติ, จตูหิ อากาเรหิ เตโชธาตุํ วิตฺถาเรน ปริคฺคณฺหาติ, ฉหิ อากาเรหิ วาโยธาตุํ วิตฺถาเรน ปริคฺคณฺหาติ.}

\prose{กตเมหิ วีสติยา อากาเรหิ ปถวีธาตุํ วิตฺถาเรน ปริคฺคณฺหาติ. อตฺถิ อิมสฺมิํ กาเย เกสา โลมา นขา ทนฺตา ตโจ, มํสํ นฺหารุ อฏฺฐิ อฏฺฐิมิญฺชํ\csromansharedfootnote{52fdc36c1772839f}{อฏฺฐิมิญฺชา (สี)} วกฺกํ, หทยํ ยกนํ กิโลมกํ ปิหกํ ปปฺผาสํ, อนฺตํ อนฺตคุณํ อุทริยํ กรีสํ มตฺถเก มตฺถลุงฺคนฺติ อิเมหิ วีสติยา อากาเรหิ ปถวีธาตุํ วิตฺถาเรน ปริคฺคณฺหาติ.}

\prose{กตเมหิ ทฺวาทสหิ อากาเรหิ อาโปธาตุํ วิตฺถาเรน ปริคฺคณฺหาติ. อตฺถิ อิมสฺมิํ กาเย ปิตฺตํ เสมฺหํ ปุพฺโพ โลหิตํ เสโท เมโท อสฺสุ วสา เขโฬ สิงฺฆาณิกา ลสิกา มุตฺตนฺติ อิเมหิ ทฺวาทสหิ อากาเรหิ อาโปธาตุํ วิตฺถาเรน ปริคฺคณฺหาติ.}

\prose{กตเมหิ จตูหิ อากาเรหิ เตโชธาตุํ วิตฺถาเรน ปริคฺคณฺหาติ. เยน จ สนฺตปฺปติ, เยน จ ชีรียติ\csromansharedfootnote{fb1080489dc28027}{ชีรติ (สี), ชีรยติ (ก); ม 3. 284 ปิฏฺเฐ ปสฺสิตพฺพํ.}, เยน จ ปริฑยฺหติ, เยน จ อสิต\-ปีต\-ขายิต\-สายิตํ สมฺมา ปริณามํ คจฺฉติ, อิเมหิ จตูหิ อากาเรหิ เตโชธาตุํ วิตฺถาเรน ปริคฺคณฺหาติ.}

\prose{กตเมหิ ฉหิ อากาเรหิ วาโยธาตุํ วิตฺถาเรน ปริคฺคณฺหาติ. อุทฺธงฺคมา วาตา, อโธคมา วาตา, กุจฺฉิสยา วาตา, โกฏฺฐาสยา\csromansharedfootnote{27cff9be9bb1b2ec}{โกฏฺฐสยา (สี)} วาตา, องฺคมงฺคา\-นุสาริโน วาตา, อสฺสาโส ปสฺสาโส อิติ, อิเมหิ ฉหิ อากาเรหิ วาโยธาตุํ วิตฺถาเรน ปริคฺคณฺหาติ.}

\prose{เอวํ อิเมหิ ทฺวาจตฺตาลีสาย อากาเรหิ วิตฺถาเรน ธาตุโย สภาวโต อุปลกฺขยนฺโต ตุลยนฺโต ปริวีมํสนฺโต ปริโยคาหนฺโต ปจฺจเวกฺขนฺโต น กิญฺจิ คยฺหูปคํ ปสฺสติ กายํ วา กายปเทสํ วา, ยถา จนฺทนิกํ ปวิจินนฺโต น กิญฺจิ คยฺหูปคํ ปสฺเสยฺย, ยถา สงฺการฏฺฐานํ ปวิจินนฺโต น กิญฺจิ คยฺหูปคํ ปสฺเสยฺย, ยถา วจฺจกุฏิํ ปวิจินนฺโต น กิญฺจิ คยฺหูปคํ ปสฺเสยฺย, ยถา สิวถิกํ\csromansharedfootnote{37e1b9824492e298}{สีวถิกํ (สี)} ปวิจินนฺโต น กิญฺจิ คยฺหูปคํ ปสฺเสยฺย. เอวเมว อิเมหิ ทฺวาจตฺตาลีสาย อากาเรหิ เอวํ วิตฺถาเรน ธาตุโย \csromanfolio{64}สภาวโต อุปลกฺขยนฺโต ตุลยนฺโต ปริวีมํสนฺโต ปริโยคาหนฺโต ปจฺจเวกฺขนฺโต น กิญฺจิ คยฺหูปคํ ปสฺสติ กายํ วา กายปเทสํ วา. เตนาห ภควา ยา เจว โข ปน อชฺฌตฺติกา ปถวีธาตุ\csromansharedfootnote{021772a977dab8a3}{เนเวสาหํ (สี, ก) ม 3. 283 ปิฏฺเฐ.}, ยา จ พาหิรา ปถวีธาตุ, ปถวีธาตุเรเวสา. ตํ “เนตํ มม, เนโสหมสฺมิ, น เมโส อตฺตา”ติ เอวเมตํ ยถาภูตํ สมฺมปฺปญฺญาย ทฏฺฐพฺพํ, เอวเมตํ ยถาภูตํ สมฺมปญฺญาย ทิสฺวา ปถวีธาตุยา นิพฺพินฺทติ, ปถวีธาตุยา จิตฺตํ วิราเชติ. ยา เจว โข ปน อชฺฌตฺติกา อาโปธาตุ, ยา จ พาหิรา อาโปธาตุ ฯเปฯ. ยา เจว โข ปน อชฺฌตฺติกา เตโชธาตุ, ยา จ พาหิรา เตโชธาตุ ฯเปฯ. ยา เจว โข ปน อชฺฌตฺติกา วาโยธาตุ, ยา จ พาหิรา วาโยธาตุ, วาโยธาตุเรเวสา. ตํ “เนตํ มม, เนโสหมสฺมิ, น เมโส อตฺตา”ติ เอวเมตํ ยถาภูตํ สมฺมปฺปญฺญาย ทฏฺฐพฺพํ, เอวเมตํ ยถาภูตํ สมฺมปฺปญฺญาย ทิสฺวา วาโยธาตุยา นิพฺพินฺทติ, วาโยธาตุยา จิตฺตํ วิราเชติ. อยํ เวมตฺตตา.}

\proseitem{48}{“อวิชฺชา”ติ เอกตฺตตา. ตตฺถ กตมา อวิชฺชา,\csromansymbolfootnote{*}{อภิ 1. 216 ปิฏฺฐาทีสุปิ.}ทุกฺเข อญฺญาณํ, ทุกฺขสมุทเย อญฺญาณํ, ทุกฺขนิโรเธ อญฺญาณํ, ทุกฺขนิโรธ\-คามินิยา ปฏิปทาย อญฺญาณํ ปุพฺพนฺเต อญฺญาณํ, อปรนฺเต อญฺญาณํ, ปุพฺพนฺตา\-ปรนฺเต อญฺญาณํ, อิทปฺปจฺจยตา\-ปฏิจฺจ\-สมุปฺปนฺเนสุ ธมฺเมสุ อญฺญาณํ, ยํ เอวรูปํ อญฺญาณํ อทสฺสนํ อนภิสมโย อนนุโพโธ อสมฺโพโธ อปฺปฏิเวโธ อสลฺลกฺขณา อนุปลกฺขณา อปจฺจุปลกฺขณา อสมเปกฺขณา\csromansharedfootnote{7651362809494e5e}{อสมเวกฺขณํ (ก)} อปจฺจกฺขกมฺมํ ทุมฺเมชฺฌํ พาลฺยํ อสมฺปชญฺญํ โมโห ปโมโห สมฺโมโห อวิชฺชา อวิชฺโชโฆ อวิชฺชาโยโค อวิชฺชานุสโย อวิชฺชา\-ปริยุฏฺฐานํ อวิชฺชาลงฺคี โมโห อกุสลมูลํ. อยํ เวมตฺตตา.}

\prose{“วิชฺชา”ติ เอกตฺตตา. ตตฺถ กตมา วิชฺชา,\csromansymbolfootnote{+}{อภิ 1. 217 ปิฏฺฐาทีสุปิ.}ทุกฺเข ญาณํ, ทุกฺขสมุทเย ญาณํ, ทุกฺขนิโรเธ ญาณํ, ทุกฺขนิโรธ\-คามินิยา ปฏิปทาย ญาณํ, ปุพฺพนฺเต ญาณํ, อปรนฺเต ญาณํ, ปุพฺพนฺตา\-ปรนฺเต ญาณํ, อิทปฺปจฺจยตา\-ปฏิจฺจ\-สมุปฺปนฺเนสุ ธมฺเมสุ ญาณํ, ยา เอวรูปา ปญฺญา ปชานนา วิจโย ปวิจโย ธมฺมวิจโย สลฺลกฺขณา อุปลกฺขณา ปจฺจุ\-ปลกฺขณา ปณฺฑิจฺจํ โกสลฺลํ เนปุญฺญํ เวภพฺยา\csromansharedfootnote{92f2c948ccb9bbd0}{เวภวฺยา (สี)} จินฺตา อุปปริกฺขา ภูรี เมธา ปริณายิกา วิปสฺสนา สมฺปชญฺญํ ปโตโท ปญฺญา ปญฺญินฺทฺริยํ ปญฺญาพลํ ปญฺญาสตฺถํ ปญฺญาปสาโท \csromanfolio{65}ปญฺญาอาโลโก ปญฺญาโอภาโส ปญฺญาปชฺโชโต ปญฺญารตนํ อโมโห ธมฺมวิจโย สมฺมาทิฏฺฐิ ธมฺมวิจย\-สมฺโพชฺฌงฺโค มคฺคงฺคํ มคฺค\-ปริยาปนฺนํ. อยํ เวมตฺตตา.}

\prose{“สมาปตฺตี”ติ เอกตฺตตา. ตตฺถ กตมา สมาปตฺติ, สญฺญาสมาปตฺติ อสญฺญา\-สมาปตฺติ, เนว\-สญฺญา\-นา\-สญฺญาสมาปตฺติ. วิภูต\-สญฺญาสมาปตฺติ นิโรธสมาปตฺตีติ. อยํ เวมตฺตตา.}

\prose{“ฌายี”ติ เอกตฺตตา. ตตฺถ กตโม ฌายี, อตฺถิ เสกฺโข ฌายี, อตฺถิ อเสกฺโข ฌายี, เนว\-เสกฺข\-นา\-เสกฺโข ฌายี, อาชานิโย ฌายี, อสฺสขลุงฺโก ฌายี, ทิฏฺฐุตฺตโร ฌายี, ตณฺหุตฺตโร ฌายี, ปญฺญุตฺตโร ฌายี. อยํ เวมตฺตตา.}

\prose{“สมาธี”ติ เอกตฺตตา. ตตฺถ กตโม สมาธิ, สรโณ สมาธิ, อรโณ สมาธิ, สเวโร สมาธิ, อเวโร สมาธิ, สพฺยาปชฺโช\csromansharedfootnote{fcad2d1e0a7dae69}{สพฺยาปชฺโฌ (สี)} สมาธิ, อพฺยาปชฺโช สมาธิ, สปฺปีติโก สมาธิ, นิปฺปีติโก สมาธิ, สามิโส สมาธิ, นิรามิโส สมาธิ, สสงฺขาโร สมาธิ, อสงฺขาโร สมาธิ, เอกํสภาวิโต สมาธิ, อุภยํส\-ภาวิโต สมาธิ, อุภยโต ภาวิตภาวโน สมาธิ, สวิตกฺก\-สวิจาโร สมาธิ, อวิตกฺก\-วิจารมตฺโต สมาธิ, อวิตกฺก\-อวิจาโร สมาธิ, หานภาคิโย สมาธิ, ฐิติภาคิโย สมาธิ, วิเสสภาคิโย สมาธิ, นิพฺเพธ\-ภาคิโย สมาธิ, โลกิโย สมาธิ, โลกุตฺตโร สมาธิ, มิจฺฉาสมาธิ, สมฺมาสมาธิ. อยํ เวมตฺตตา.}

\prose{“ปฏิปทา”ติ เอกตฺตตา. ตตฺถ กตมา ปฏิปทา, อาคาฬฺหปฏิปทา\csromansharedfootnote{eee52b591f2511c8}{อาคาฬฺหา ปฏิปทา (สี) อฏฺฐกถา โอโลเกตพฺพา.}, นิชฺฌาม\-ปฏิปทา, มชฺฌิม\-ปฏิปทา, อกฺขมา ปฏิปทา, ขมา ปฏิปทา, สมา ปฏิปทา, ทมา ปฏิปทา, ทุกฺขา ปฏิปทา ทนฺธาหิญฺญา, ทุกฺขา ปฏิปทา ขิปฺปาภิญฺญา, สุขา ปฏิปทา ทนฺธาภิญฺญา, สุขา ปฏิปทา ขิปฺปาภิญฺญาติ. อยํ เวมตฺตตา.}

\prose{“กาโย”ติ เอกตฺตตา. ตตฺถ กตโม กาโย, นามกาโย รูปกาโย จ. ตตฺถ กตโม รูปกาโย,\csromansymbolfootnote{*}{ที 2. 233-4; ม 1. 72; ขุ 1. 2; อภิ 2. 200 ปิฏฺเฐสุปิ.}เกสา โลมา นขา ทนฺตา ตโจ มํสํ นฺหารุ\csromansharedfootnote{4b9a2ce0c786a2d7}{นหารุ (สี)} อฏฺฐิ อฏฺฐิมิญฺชํ วกฺกํ หทยํ ยกนํ กิโลมกํ ปิหกํ ปปฺผาสํ อนฺตํ อนฺตคุณํ อุทริยํ กรีสํ ปิตฺตํ เสมฺหํ ปุพฺโพ โลหิตํ เสโท \csromanfolio{66}เมโท อสฺสุ วสา เขโฬ สิงฺฆาณิกา ลสิกา มุตฺตํ มตฺถลุงฺคนฺติ, อยํ รูปกาโย. นามกาโย นาม เวทนา สญฺญา เจตนา จิตฺตํ ผสฺโส มนสิกาโรติ, อยํ นามกาโยติ. อยํ เวมตฺตตา.}

\prose{เอวํ โย ธมฺโม ยสฺส ธมฺมสฺส สมานภาโว, โส ธมฺโม ตสฺส ธมฺมสฺส เอกตฺตตาย เอกี ภวติ. เยน เยน วา ปน วิลกฺขโณ, เตน เตน เวมตฺตํ คจฺฉติ. เอวํ สุตฺเต วา เวยฺยากรเณ วา คาถายํ วา ปุจฺฉิเตน วีมํส\-ยิตพฺพํ, กิํ เอกตฺตตาย ปุจฺฉติ, อุทาหุ เวมตฺตตายาติ. ยทิ เอกตฺตตาย ปุจฺฉิตํ, เอกตฺตตาย วิสชฺช\-ยิตพฺพํ. ยทิ เวมตฺตตาย ปุจฺฉิตํ, เวมตฺตตาย วิสชฺช\-ยิตพฺพํ. ยทิ สตฺตา\-ธิฏฺฐาเนน ปุจฺฉิตํ, สตฺตา\-ธิฏฺฐาเนน วิสชฺช\-ยิตพฺพํ. ยทิ ธมฺมา\-ธิฏฺฐาเนน ปุจฺฉิตํ, ธมฺมา\-ธิฏฺฐาเนน วิสชฺช\-ยิตพฺพํ. ยถา ยถา วา ปน ปุจฺฉิตํ, ตถา ตถา วิสชฺช\-ยิตพฺพํ. เตนาห อายสฺมา มหากจฺจายโน “เอกตฺตตาย ธมฺมา”ติ.}

\vspace{\tipitakaparskip}
\csromancenter{นิยุตฺโต อธิฏฺฐาโน หาโร.}
\csromanmatikaanchor{mtk.19}
\csromantocmarkref{section}{15. ปริกฺขารหารวิภงฺค}{mtk.19}
\bookmark[dest={mtk.19},level=1]{15. ปริกฺขารหารวิภงฺค}
\csromanheader{1}{15. \textbf{ปริกฺขาร}\-\textbf{หาร}\-\textbf{วิภงฺค}}
\proseitem{49}{ตตฺถ กตโม ปริกฺขาโร หาโร, “เย ธมฺมา ยํ ธมฺมํ ชนยนฺตี”ติ.}

\prose{โย ธมฺโม ยํ ธมฺมํ ชนยติ, ตสฺส โส ปริกฺขาโร. กิํลกฺขโณ ปริกฺขาโร, ชนกลกฺขโณ ปริกฺขาโร. เทฺว ธมฺมา ชนยนฺติ เหตุ จ ปจฺจโย จ. ตตฺถ กิํลกฺขโณ เหตุ, กิํลกฺขโณ ปจฺจโย. อสาธารณ\-ลกฺขโณ เหตุ, สาธารณ\-ลกฺขโณ ปจฺจโย. ยถา กิํ ภเว. ยถา องฺกุรสฺส นิพฺพตฺติยา พีชํ อสาธารณํ, ปถวี อาโป จ สาธารณา. องฺกุรสฺส หิ ปถวี อาโป จ ปจฺจโย สภาโว เหตุ. ยถา วา ปน ฆเฏ ทุทฺธํ ปกฺขิตฺตํ ทธิ ภวติ, น จตฺถิ เอกกา\-ลสม\-วธานํ ทุทฺธสฺส จ ทธิสฺส จ. เอวเมวํ นตฺถิ เอกกา\-ลสม\-วธานํ เหตุสฺส จ ปจฺจยสฺส จ.}

\prose{อยญฺหิ สํสาโร สเหตุ สปฺปจฺจโย นิพฺพตฺโต. วุตฺตํ หิ อวิชฺชาปจฺจยา สงฺขารา, สงฺขาร\-ปจฺจยา วิญฺญาณํ, เอวํ สพฺโพ ปฏิจฺจ\-สมุปฺปาโท. อิติ อวิชฺชา อวิชฺชาย เหตุ อโยนิโส มนสิกาโร ปจฺจโย. ปุริมิกา อวิชฺชา ปจฺฉิมิกาย อวิชฺชาย เหตุ. ตตฺถ ปุริมิกา อวิชฺชา \csromanfolio{67}อวิชฺชานุสโย ปจฺฉิมิกา อวิชฺชา อวิชฺชา\-ปริยุฏฺฐานํ, ปุริมิโก อวิชฺชานุสโย ปจฺฉิมิกสฺส อวิชฺชา\-ปริยุฏฺฐานสฺส เหตุภูโต ปริพฺรูหนาย, พีชงฺกุโร วิย สมนนฺตร\-เหตุตาย. ยํ ปน ยตฺถ ผลํ นิพฺพตฺตติ, อิทมสฺส ปรมฺปร\-เหตุตาย เหตุภูตํ. ทุวิโธ หิ เหตุ สมนนฺตร\-เหตุ ปรมฺปรเหตุ จ, เอวํ อวิชฺชายปิ ทุวิโธ เหตุ สมนนฺตร\-เหตุ ปรมฺปรเหตุ จ.}

\prose{ยถา วา ปน ถาลกญฺจ วฏฺฏิ จ เตลญฺจ ปทีปสฺส ปจฺจยภูตํ น สภาวเหตุ, น หิ สกฺกา ถาลกญฺจ วฏฺฏิญฺจ เตลญฺจ อนคฺคิกํ ทีเปตุํ ปทีสสฺส ปจฺจยภูตํ. ปทีโป วิย สภาโว เหตุ โหติ. อิติ สภาโว เหตุ, ปรภาโว ปจฺจโย. อชฺฌตฺติโก เหตุ, พาหิโร ปจฺจโย. ชนโก เหตุ, ปริคฺคาหโก ปจฺจโย. อสาธารโณ เหตุ, สาธารโณ ปจฺจโย.}

\prose{อวุปจฺเฉทตฺโถ สนฺตติ อตฺโถ, นิพฺพตฺติ อตฺโถ ผลตฺโถ, ปฏิสนฺธิ อตฺโถ ปุนพฺภวตฺโถ, ปลิโพธตฺโถ ปริยุฏฺฐา\-นตฺโถ, อสมุ\-คฺฆาต\-ตฺโถ อนุสยตฺโถ, อสมฺปฏิเวธตฺโถ อวิชฺชตฺโถ, อปริญฺญาตตฺโถ วิญฺญาณสฺส พีชตฺโถ. ยตฺถ อวุปจฺเฉโท ตตฺถ สนฺตติ, ยตฺถ สนฺตติ ตตฺถ นิพฺพตฺติ, ยตฺถ นิพฺพตฺติ ตตฺถ ผลํ, ยตฺถ ผลํ ตตฺถ ปฏิสนฺธิ, ยตฺถ ปฏิสนฺธิ ตตฺถ ปุนพฺภโว, ยตฺถ ปุนพฺภโว ตตฺถ ปลิโพโธ, ยตฺถ ปลิโพโธ ตตฺถ ปริยุฏฺฐานํ, ยตฺถ ปริยุฏฺฐานํ ตตฺถ อสมุคฺฆาโต. ยตฺถ อสมุคฺฆาโต ตตฺถ อนุสโย, ยตฺถ อนุสโย ตตฺถ อสมฺปฏิเวโธ, ยตฺถ อสมฺปฏิเวโธ ตตฺถ อวิชฺชา, ยตฺถ อวิชฺชา ตตฺถ สาสวํ วิญฺญาณํ อปริญฺญาตํ, ยตฺถ สาสวํ วิญฺญาณํ อปริญฺญาตํ ตตฺถ พีชตฺโถ.}

\prose{สีลกฺขนฺโธ สมาธิ\-กฺขนฺธสฺส ปจฺจโย, สมาธิ\-กฺขนฺโธ ปญฺญา\-กฺขนฺธสฺส ปจฺจโย, ปญฺญากฺขนฺโธ วิมุตฺติ\-กฺขนฺธสฺส ปจฺจโย, วิมุตฺติ\-กฺขนฺโธ วิมุตฺติ\-ญาณ\-ทสฺสนกฺขนฺธสฺส ปจฺจโย. ติตฺถญฺญุตา ปีตญฺญุตาย ปจฺจโย, ปีตญฺญุตา ปตฺตญฺญุตาย ปจฺจโย, ปตฺตญฺญุตา อตฺตญฺญุตาย ปจฺจโย.}

\prose{ยถา วา ปน จกฺขุญฺจ ปฏิจฺจ รูเป จ อุปฺปชฺชติ จกฺขุวิญฺญาณํ. ตตฺถ จกฺขุ อาธิปเตยฺย\-ปจฺจยตาย ปจฺจโย, รูปา อารมฺมณ\-ปจฺจยตาย ปจฺจโย. อาโลโก สินฺนิสฺสยตาย ปจฺจโย, มนสิกาโร สภาโว \csromanfolio{68}เหตุ. สงฺขารา วิญฺญาณสฺส ปจฺจโย สภาโว เหตุ. วิญฺญาณํ นามรูปสฺส ปจฺจโย สภาโว เหตุ. นามรูปํ สฬายตนสฺส ปจฺจโย สภาโว เหตุ. สฬายตนํ ผสฺสสฺส ปจฺจโย สภาโว เหตุ. ผสฺโส เวทนาย ปจฺจโย สภาโว เหตุ. เวทนา ตณฺหาย ปจฺจโย สภาโว เหตุ. ตณฺหา อุปาทานสฺส ปจฺจโย สภาโว เหตุ. อุปาทานํ ภวสฺส ปจฺจโย สภาโว เหตุ. ภโว ชาติยา ปจฺจโย สภาโว เหตุ. ชาติ ชารามรณสฺส ปจฺจโย สภาโว เหตุ. ชรามรณํ โสกสฺส ปจฺจโย สภาโว เหตุ. โสโก ปริเทวสฺส ปจฺจโย สภาโว เหตุ. ปริเทโว ทุกฺขสฺส ปจฺจโย สภาโว เหตุ. ทุกฺขํ โทมนสฺสสฺส ปจฺจโย สภาโว เหตุ. โทมนสฺสํ อุปายาสสฺส ปจฺจโย สภาโว เหตุ. เอวํ โย โกจิ อุปนิสฺสโย สพฺโพ โส ปริกฺขาโร. เตนาห อายสฺมา มหากจฺจายโน “เย ธมฺมา ยํ ธมฺมํ ชนยนฺตี”ติ.}

\vspace{\tipitakaparskip}
\csromancenter{นิยุตฺโต ปริกฺขาโร หาโร.}
\csromanmatikaanchor{mtk.20}
\csromantocmarkref{section}{16. สมาโรปนหารวิภงฺค}{mtk.20}
\bookmark[dest={mtk.20},level=1]{16. สมาโรปนหารวิภงฺค}
\csromanheader{1}{16. \textbf{สมาโรปน}\-\textbf{หาร}\-\textbf{วิภงฺค}}
\proseitem{50}{ตตฺถ กตโม สมาโรปโน หาโร, “เย ธมฺมา ยํมูลา, เย เจกตฺถา ปกาสิตา มุนินา”ติ.}

\prose{เอกสฺมิํ ปทฏฺฐาเน ยตฺตกานิ ปทฏฺฐานานิ โอตรนฺติ, สพฺพานิ ตานิ สมา\-โรป\-ยิตพฺพานิ. ยถา อาวฏฺเฏ หาเร พหุกานิ ปทฏฺฐานานิ โอตรนฺตีติ. ตตฺถ สมาโรปนา จตุพฺพิธา ปทฏฺฐานํ, เววจนํ, ภาวนา, ปหานมิติ.}

\prose{ตตฺถ กตมา ปทฏฺฐาเนน สมาโรปนา.}

\setlength{\csromangathapagewidth}{0pt}
\setlength{\csromangathawakwidth}{0pt}
\csromangathameasuregroup{\textsuperscript{*}“สพฺพปาปสฺส อกรณํ, \\ สจิตฺตปริโยทาปนํ,}{\csromangathabat{\textsuperscript{*}“สพฺพปาปสฺส อกรณํ,}{กุสลสฺส อุปสมฺปทา.} \\ \csromangathabat{สจิตฺตปริโยทาปนํ,}{เอตํ พุทฺธาน สาสนนฺ”ติ.}}
\csromangathasetleft{\textsuperscript{*}“สพฺพปาปสฺส อกรณํ, \\ สจิตฺตปริโยทาปนํ,}
\csromangathagroup{\csromangathastanza{\csromangathaitembat{50}{\csromansymbolfootnote{*}{เหฏฺฐา 37 ปิฏฺเฐ; อุปริ 148, 161, 205, 208 ปิฏฺเฐสุปิ.}“สพฺพปาปสฺส อกรณํ,}{กุสลสฺส อุปสมฺปทา.} \\ \csromangathacontbat{สจิตฺต\-ปริโยทาปนํ,}{เอตํ พุทฺธาน สาสนนฺ”ติ.}}}

\prose{ตสฺส กิํ ปทฏฺฐานํ, ตีณิ สุจริตานิ กายสุจริตํ วจีสุจริตํ มโนสุจริตํ, อิทํ ปทฏฺฐานํ, ตตฺถ ยํ กายิกญฺจ วาจสิกญฺจ สุจริตํ, อยํ สีลกฺขนฺโธ. มโนสุจริเต ยา อนภิชฺฌา อพฺยาปาโท จ, อยํ สมาธิ\-กฺขนฺโธ. ยา สมฺมาทิฏฺฐิ, อยํ ปญฺญากฺขนฺโธ. อิทํ ปทฏฺฐานํ, ตตฺถ \csromanfolio{69}สีลกฺขนฺโธ จ สมาธิ\-กฺขนฺโธ จ สมโถ, ปญฺญากฺขนฺโธ วิปสฺสนา. อิทํ ปทฏฺฐานํ, ตตฺถ สมถสฺส ผลํ ราควิราคา เจโตวิมุตฺติ, วิปสฺสนาย ผลํ อวิชฺชาวิราคา ปญฺญาวิมุตฺติ. อิทํ ปทฏฺฐานํ.}

\prose{วนํ วนถสฺส ปทฏฺฐานํ. กิญฺจ วนํ, โก จ วนโถ. วนํ นาม ปญฺจ กามคุณา, ตณฺหา วนโถ. อิทํ ปทฏฺฐานํ. วนํ นาม นิมิตฺตคฺคาโห “อิตฺถี”ติ วา “ปุริโส”ติ วา. วนโถ นาม เตสํ เตสํ องฺค\-ปจฺจงฺคานํ อนุพฺยญฺชน\-คฺคาโห “อโห จกฺขุ, อโห โสตํ, อโห ฆานํ, อโห ชิวฺหา, อโห กาโย, อิติ. อิทํ ปทฏฺฐานํ. วนํ นาม ฉ อชฺฌตฺติก\-พาหิรานิ อายตนานิ อปริญฺญาตานิ. ยํ ตทุภยํ ปฏิจฺจ อุปฺปชฺชติ สญฺโญชนํ, อยํ วนโถ. อิทํ ปทฏฺฐานํ. วนํ นาม อนุสโย. วนโถ นาม ปริยุฏฺฐานํ. อิทํ ปทฏฺฐานํ. เตนาห ภควา “เฉตฺวา วนญฺจ วนถญฺจา”ติ. อยํ ปทฏฺฐาเนน สมาโรปนา. (1)}

\proseitem{51}{ตตฺถ กตมา เววจเนน สมาโรปนา, ราควิราคา เจโตวิมุตฺติ เสกฺขผลํ. อวิชฺชาวิราคา ปญฺญาวิมุตฺติ อเสกฺขผลํ. อิทํ เววจนํ. ราควิราคา เจโตวิมุตฺติ อนาคามิผลํ. อวิชฺชาวิราคา ปญฺญาวิมุตฺติ อคฺคผลํ อรหตฺตํ. อิทํ เววจนํ. ราควิราคา เจโตวิมุตฺติ กามธาตุ\-สมติกฺกมนํ. อวิชฺชาวิราคา ปญฺญาวิมุตฺติ เตธาตุสมติกฺกมนํ. อิทํ เววจนํ. ปญฺญินฺทฺริยํ, ปญฺญาพลํ, อธิปญฺญา\-สิกฺขา, ปญฺญากฺขนฺโธ, ธมฺมวิจย\-สมฺโพชฺฌงฺโค, อุเปกฺขา\-สมฺโพชฺฌงฺโค, ญาณํ, สมฺมาทิฏฺฐิ, ตีรณา, สนฺตีรณา, หิรี, วิปสฺสนา, ธมฺเม ญาณํ, สพฺพํ, อิทํ เววจนํ. อยํ เววจเนน สมาโรปนา. (2)}

\prose{ตตฺถ กตมา ภาวนาย สมาโรปนา, ยถาห ภควา “ตสฺมาติห ตฺวํ ภิกฺขุ กาเย กายานุปสฺสี วิหราหิ, อาตาปี สมฺปชาโน สติมา วิเนยฺย โลเก อภิชฺฌา\-โทมนสฺสํ”. อาตาปีติ วีริยินฺทฺริยํ. สมฺปชาโนติ ปญฺญินฺทฺริยํ. สติมาติ สตินฺทฺริยํ. วิเนยฺย โลเก อภิชฺฌา\-โทมนสฺสนฺติ สมาธินฺทฺริยํ. เอวํ กาเย กายานุปสฺสิโน วิหรโต จตฺตาโร สติปฏฺฐานา ภาวนา\-ปาริปูริํ คจฺฉนฺติ. เกน การเณน, เอกลกฺขณตฺตา จตุนฺนํ อินฺทฺริยานํ. จตูสุ สติปฏฺฐาเนสุ ภาวิยมาเนสุ จตฺตาโร สมฺมปฺปธานา ภาวนา\-ปาริปูริํ คจฺฉนฺติ. จตูสุ สมฺม\-ปฺปธาเนสุ ภาวิยมาเนสุ จตฺตาโร อิทฺธิปาทา ภาวนา\-ปาริปูริํ คจฺฉนฺติ. จตูสุ อิทฺธิปาเทสุ ภาวิยมาเนสุ ปญฺจินฺทฺริยานิ ภาวนา\-ปาริปูริํ คจฺฉนฺติ. เอวํ \csromanfolio{70}สพฺเพ. เกน การเณน, สพฺเพ หิ โพธงฺคมา ธมฺมา โพธิปกฺขิยา นิยฺยานิก\-ลกฺขเณน เอกลกฺขณา, เต เอกลกฺขณตฺตา ภาวนา\-ปาริปูริํ คจฺฉนฺติ. อยํ ภาวนาย สมาโรปนา. (3)}

\prose{ตตฺถ กตมา ปหาเนน สมาโรปนา, กาเย กายานุปสฺสี วิหรนฺโต “อสุเภ สุภนฺ”ติ วิปลฺลาสํ ปชหติ, กพฬีกาโร จสฺส อาหาโร ปริญฺญํ คจฺฉติ, กามุปาทาเนน จ อนุปาทาโน ภวติ, กามโยเคน จ วิสํยุตฺโต ภวติ, อภิชฺฌา\-กายคนฺเถน จ วิปฺปยุชฺชติ, กามาสเวน จ อนาสโว ภวติ, กาโมฆญฺจ อุตฺติณฺโณ ภวติ, ราคสลฺเลน จ วิสลฺโล ภวติ, รูปูปิกา\csromansharedfootnote{eb2b031a3f8b5445}{รูปุปิกา (ก) เอวมุปริปิ.} จสฺส วิญฺญาณฏฺฐิติ ปริญฺญํ คจฺฉติ, รูปธาตุยํ จสฺส ราโค ปหีโน ภวติ, น จ ฉนฺทาคติํ คจฺฉติ.}

\prose{เวทนาสุ เวทนานุปสฺสี วิหรนฺโต “ทุกฺเข สุขนฺ”ติ วิปลฺลาสํ ปชหติ, ผสฺโส จสฺส อาหาโร ปริญฺญํ คจฺฉติ, ภวูปาทาเนน จ อนุปาทาโน ภวติ, ภวโยเคน จ วิสํยุตฺโต ภวติ, พฺยาปาท\-กายคนฺเถน จ วิปฺปยุชฺชติ, ภวาสเวน จ อนาสโว ภวติ, ภโวฆญฺจ อุตฺติณฺโณ ภวติ, โทสสลฺเลน จ วิสลฺโล ภวติ, เวทนูปิกา จสฺส วิญฺญาณฏฺฐิติ ปริญฺญํ คจฺฉติ, เวทนาธาตุยํ จสฺส ราโค ปหีโน ภวติ, น จ โทสาคติํ คจฺฉติ.}

\prose{จิตฺเต จิตฺตานุปสฺสี วิหรนฺโต “อนิจฺเจ นิจฺจนฺ”ติ วิปลฺลาสํ ปชหติ, วิญฺญาณํ จสฺส อาหาโร ปริญฺญํ คจฺฉติ, ทิฏฺฐุปาทาเนน จ อนุปาทาโน ภวติ, ทิฏฺฐิโยเคน จ วิสํยุตฺโต ภวติ, สีลพฺพตปรามาส\-กายคนฺเถน จ วิปฺปยุชฺชติ, ทิฏฺฐาสเวน จ อนาสโว ภวติ, ทิฏฺโฐฆญฺจ อุตฺติณฺโณ ภวติ, มานสลฺเลน จ วิสลฺโล ภวติ, สญฺญูปิกา จสฺส วิญฺญาณฏฺฐิติ ปริญฺญํ คจฺฉติ, สญฺญาธาตุยํ จสฺส ราโค ปหีโน ภวติ, น จ ภยาคติํ คจฺฉติ.}

\prose{ธมฺเมสุ ธมฺมานุปสฺสี วิหรนฺโต “อนตฺตนิ\csromansharedfootnote{82320957c33567f0}{อนตฺตนิเย (สี) ปสฺส อํ 1. 361 ปิฏฺเฐ.} อตฺตา”ติ วิปลฺลาสํ ปชหติ, มโนสญฺเจตนา จสฺส อาหาโร ปริญฺญํ คจฺฉติ, อตฺตวาทุ\-ปาทาเนน จ อนุปาทาโน ภวติ, อวิชฺชาโยเคน จ วิสํยุตฺโต ภวติ, อิทํสจฺจาภินิเวส\-กายคนฺเถน จ วิปฺปยุชฺชติ, อวิชฺชาสเวน จ อนาสโว ภวติ, อวิชฺโชฆญฺจ อุตฺติณฺโณ ภวติ, โมหสลฺเลน จ วิสลฺโล \csromanfolio{71}ภวติ, สงฺขารูปิกา จสฺส วิญฺญาณฏฺฐิติ ปริญฺญํ คจฺฉติ, สงฺขาร\-ธาตุยํ จสฺส ราโค ปหีโน ภวติ, น จ โมหาคติํ คจฺฉติ. อยํ ปหาเนน สมาโรปนา. (4)}

\prose{เตนาห อายสฺมา มหากจฺจายโน\csromandash{}}

\setlength{\csromangathapagewidth}{0pt}
\setlength{\csromangathawakwidth}{0pt}
\csromangathameasuregroup{“เย ธมฺมา ยํมูลา, \\ เต สมาโรปยิตพฺพา,}{\csromangathabat{“เย ธมฺมา ยํมูลา,}{เย เจกตฺถา ปกาสิตา มุนินา.} \\ \csromangathabat{เต สมาโรปยิตพฺพา,}{เอส สมาโรปโน หาโร”ติ.}}
\csromangathasetleft{“เย ธมฺมา ยํมูลา, \\ เต สมาโรปยิตพฺพา,}
\csromangathagroup{\csromangathastanza{\csromangathabat{“เย ธมฺมา ยํมูลา,}{เย เจกตฺถา ปกาสิตา มุนินา.} \\ \csromangathabat{เต สมา\-โรป\-ยิตพฺพา,}{เอส สมาโรปโน หาโร”ติ.}}}

\vspace{\tipitakaparskip}
\csromancenter{นิยุตฺโต สมาโรปโน หาโร.}
\nitthitam{นิฏฺฐิโต จ หารวิภงฺโค.}
\csromanmatikaanchor{mtk.21}
\csromantocmarkref{section}{1. เทสนาหารสมฺปาต}{mtk.21}
\bookmark[dest={mtk.21},level=1]{1. เทสนาหารสมฺปาต}
\csromanheader{1}{1. \textbf{เทสนาหาร}\-\textbf{สมฺปาต}}
\setlength{\csromangathapagewidth}{0pt}
\setlength{\csromangathawakwidth}{0pt}
\csromangathameasuregroup{}{“โสฬส หารา ปฐมํ, \\ ทิสโลจนโต ทิสา วิโลเกตฺวา. \\ สงฺขิปิย องฺกุเสน หิ, \\ นเยหิ ตีหิ นิทฺทิเส สุตฺตนฺ”ติ.}
\csromangathameasurewak{“โสฬส หารา ปฐมํ, \\ ทิสโลจนโต ทิสา วิโลเกตฺวา. \\ สงฺขิปิย องฺกุเสน หิ, \\ นเยหิ ตีหิ นิทฺทิเส สุตฺตนฺ”ติ.}
\setlength{\csromangathaleftcol}{0pt}
\csromangathagroup{\csromangathastanza{\csromangathawak{\csromangathaitemline{52}{“โสฬส หารา ปฐมํ,} \\ \csromangathacontline{ทิสโลจนโต ทิสา วิโลเกตฺวา.} \\ \csromangathacontline{สงฺขิปิย องฺกุเสน หิ,} \\ \csromangathacontline{นเยหิ ตีหิ นิทฺทิเส สุตฺตนฺ”ติ.}}}}

\prose{วุตฺตา, ตสฺสา นิทฺเทโส กุหิํ ทฏฺฐพฺโพ, หารสมฺปาเต. ตตฺถ กตโม เทสนาหาร\-สมฺปาโต.}

\setlength{\csromangathapagewidth}{0pt}
\setlength{\csromangathawakwidth}{0pt}
\csromangathameasuregroup{“อรกฺขิเตน จิตฺเตน\textsuperscript{1}, \\ ถินมิทฺธาภิภูเตน,}{\csromangathabat{“อรกฺขิเตน จิตฺเตน\textsuperscript{1},}{มิจฺฉาทิฏฺฐิหเตน จ.} \\ \csromangathabat{ถินมิทฺธาภิภูเตน,}{วสํ มารสฺส คจฺฉตี”ติ.}}
\csromangathasetleft{“อรกฺขิเตน จิตฺเตน\textsuperscript{1}, \\ ถินมิทฺธาภิภูเตน,}
\csromangathagroup{\csromangathastanza{\csromangathabat{“อรกฺขิเตน จิตฺเตน\csromansharedfootnote{c9ea673238c90054}{กาเยน (ขุ 1. 121 ปิฏฺเฐ.)},}{มิจฺฉา\-ทิฏฺฐิ\-หเตน จ.} \\ \csromangathabat{ถินมิทฺธา\-ภิภูเตน,}{วสํ มารสฺส คจฺฉตี”ติ.}}}

\prose{\textbf{อรกฺขิเตน จิตฺเตนา}ติ กิํ เทสยติ, ปมาทํ ตํ มจฺจุโน ปทํ. \textbf{มิจฺฉา}\-\textbf{ทิฏฺฐิ}\-\textbf{หเตน จา}ติ มิจฺฉา\-ทิฏฺฐิหตํ นาม วุจฺจติ ยทา “อนิจฺเจ นิจฺจนฺ”ติ ปสฺสติ, โส วิปลฺลาโส. โส ปน วิปลฺลาโส กิํลกฺขโณ, วิปรีต\-คฺคาห\-ลกฺขโณ วิปลฺลาโส. โส กิํ วิปลฺลาสยติ, ตโย ธมฺเม สญฺญํ จิตฺตํ ทิฏฺฐิมิติ. โส กุหิํ วิปลฺลาสยติ, จตูสุ อตฺตภาว\-วตฺถูสุ,\csromansymbolfootnote{*}{ม 1. 374; ม 3. 68; สํ 2. 36, 478; อํ 1. 535; อภิ 1. 208; อภิ 2. 379; ปฏิสํ 147 ปิฏฺเฐสุปิ.}รูปํ อตฺตโต สมนุปสฺสติ, รูปวนฺตํ วา อตฺตานํ, อตฺตนิ วา รูปํ, รูปสฺมิํ วา อตฺตานํ. เอวํ เวทนํ ฯเปฯ สญฺญํ. สงฺขาเร. วิญฺญาณํ อตฺตโต สมนุปสฺสติ, วิญฺญาณวนฺตํ วา อตฺตานํ, อตฺตนิ วา วิญฺญาณํ, วิญฺญาณสฺมิํ วา อตฺตานํ.}

\prose{\csromanfolio{72}ตตฺถ รูปํ ปฐมํ วิปลฺลาส\-วตฺถุ “อสุเภ สุภนฺ”ติ. เวทนา ทุติยํ วิปลฺลาส\-วตฺถุ “ทุกฺเข สุขนฺ”ติ. สญฺญา สงฺขารา จ ตติยํ วิปลฺลาส\-วตฺถุ “อนตฺตนิ อตฺตา”ติ. วิญฺญาณํ จตุตฺถํ วิปลฺลาส\-วตฺถุ “อนิจฺเจ นิจฺจนฺ”ติ. เทฺว ธมฺมา จิตฺตสฺส สํกิเลสา ตณฺหา จ อวิชฺชา จ. ตณฺหานิวุตํ จิตฺตํ ทฺวีหิ วิปลฺลาเสหิ วิปลฺลาสียติ “อสุเภ สุภนฺ”ติ “ทุกฺเข สุขนฺ”ติ. ทิฏฺฐินิวุตํ จิตฺตํ ทฺวีหิ วิปลฺลาเสหิ วิปลฺลาสียติ “อนิจฺเจ นิจฺจนฺ”ติ “อนตฺตนิ อตฺตา”ติ.}

\prose{ตตฺถ โย ทิฏฺฐิ\-วิปลฺลาโส, โส อตีตํ รูปํ อตฺตโต สมนุปสฺสติ, อตีตํ เวทนํ ฯเปฯ อตีตํ สญฺญํ, อตีเต สงฺขาเร. อตีตํ วิญฺญาณํ อตฺตโต สมนุปสฺสติ. ตตฺถ โย ตณฺหาวิปลฺลาโส, โส อนาคตํ รูปํ อภินนฺทติ, อนาคตํ เวทนํ ฯเปฯ อนาคตํ สญฺญํ, อนาคเต สงฺขาเร, อนาคตํ วิญฺญาณํ อภินนฺทติ. เทฺว ธมฺมา จิตฺตสฺส อุปกฺกิเลสา ตณฺหา จ อวิชฺชา จ. ตาหิ วิสุชฺฌนฺตํ จิตฺตํ วิสุชฺฌติ. เตสํ อวิชฺชา\-นีวรณานํ ตณฺหา\-สํโยชนานํ ปุพฺพา โกฏิ น ปญฺญายติ สนฺธาวนฺตานํ สํสรนฺตานํ สกิํ นิรยํ สกิํ ติรจฺฉาน\-โยนิํ สกิํ เปตฺติวิสยํ สกิํ อสุรกายํ สกิํ เทเว สกิํ มนุสฺเส.}

\prose{\textbf{ถินมิทฺธาภิภูเตนา}ติ ถินํ\csromansharedfootnote{69d3ef48af45fc1c}{ถีนํ (สี)} นาม ยา จิตฺตสฺส อกลฺลตา อกมฺมนิยตา. มิทฺธํ นาม ยํ กายสฺส ลีนตฺตํ. \textbf{วสํ มารสฺส คจฺฉตี}ติ กิเลสมารสฺส จ สตฺตมารสฺส จ วสํ คจฺฉติ, โส หิ นิวุโต สํสารา\-ภิมุโข โหติ. อิมานิ ภควตา เทฺว สจฺจานิ เทสิตานิ ทุกฺขํ สมุทโย จ. เตสํ ภควา ปริญฺญาย จ ปหานาย จ ธมฺมํ เทเสติ ทุกฺขสฺส ปริญฺญาย สมุทยสฺส ปหานาย. เยน จ ปริชานาติ เยน จ ปชหติ, อยํ มคฺโค. ยํ ตณฺหาย อวิชฺชาย จ ปหานํ, อยํ นิโรโธ. อิมานิ จตฺตาริ สจฺจานิ. เตนาห ภควา “อรกฺขิเตน จิตฺเตนา”ติ. เตนาหายสฺมา มหากจฺจายโน “อสฺสาทาทีนวตา”ติ.}

\vspace{\tipitakaparskip}
\csromancenter{นิยุตฺโต เทสนาหาร\-สมฺปาโต.}
\csromanmatikaanchor{mtk.22}
\csromantocmarkref{section}{2. วิจยหารสมฺปาต}{mtk.22}
\bookmark[dest={mtk.22},level=1]{2. วิจยหารสมฺปาต}
\csromanheader{1}{2. \textbf{วิจย}\-\textbf{หาร}\-\textbf{สมฺปาต}}
\proseitem{53}{\csromanfolio{73}ตตฺถ กตโม วิจโย หารสมฺปาโต. ตตฺถ ตณฺหา ทุวิธา กุสลาปิ อกุสลาปิ. อกุสลา สํสารคามินี, กุสลา อปจยคามินี ปหานตณฺหา. มาโนปิ ทุวิโธ กุสโลปิ อกุสโลปิ. ยํ มานํ นิสฺสาย มานํ ปชหติ, อยํ มาโน กุสโล. โย ปน มาโน ทุกฺขํ นิพฺพตฺตยติ, อยํ มาโน อกุสโล. ตตฺถ ยํ เนกฺขมฺมสิตํ โทมนสฺสํ กุทาสฺสุนามาหํ ตํ อายตนํ สจฺฉิกตฺวา อุปสมฺปชฺช วิหริสฺสํ ยํ อริยา สนฺตํ อายตนํ สจฺฉิกตฺวา อุปสมฺปชฺช วิหรนฺตีติ ตสฺส อุปฺปชฺชติ ปิหา, ปิหาปจฺจยา โทมนสฺสํ, อยํ ตณฺหา กุสลา ราควิราคา เจโตวิมุตฺติ, ตทารมฺมณา กุสลา อวิชฺชาวิราคา ปญฺญาวิมุตฺติ.}

\prose{ตสฺสา โก ปวิจโย, อฏฺฐ มคฺคงฺคานิ สมฺมาทิฏฺฐิ สมฺมาสงฺกปฺโป สมฺมาวาจา สมฺมากมฺมนฺโต สมฺมาอาชีโว สมฺมาวายาโม สมฺมาสติ สมฺมาสมาธิ. โส กตฺถ ทฏฺฐพฺโพ, จตุตฺเถ ฌาเน ปารมิตาย. จตุตฺเถ หิ ฌาเน อฏฺฐงฺค\-สมนฺนาคตํ จิตฺตํ ภาวยติ ปริสุทฺธํ ปริโยทาถํ อนงฺคณํ วิคตู\-ปกฺกิเลสํ มุทุ กมฺมนิยํ ฐิตํ อาเนญฺชปฺปตฺตํ. โส ตตฺถ อฏฺฐวิธํ อธิคจฺฉติ ฉ อภิญฺญา เทฺว จ วิเสเส, ตํ จิตฺตํ ยโต ปริสุทฺธํ, ตโต ปริโยทาตํ, ยโต ปริโยทาตํ, ตโต อนงฺคณํ, ยโต อนงฺคณํ, ตโต วิคตู\-ปกฺกิเลสํ, ยโต วิคตู\-ปกฺกิเลสํ, ตโต มุทุ, ยโต มุทุ, ตโต กมฺมนิยํ, ยโต กมฺมนิยํ, ตโต ฐิตํ, ยโต ฐิตํ, ตโต อาเนญฺชปฺปตฺตํ. ตตฺถ องฺคณา จ อุปกฺกิเลสา จ ตทุภยํ ตณฺหาปกฺโข. ยา จ อิญฺชนา ยา จ จิตฺตสฺส อฏฺฐิติ, อยํ ทิฏฺฐิปกฺโข.}

\prose{จตฺตาริ อินฺทฺริยานิ ทุกฺขินฺทฺริยํ โทมนสฺสิ\-นฺทฺริยํ สุขินฺทฺริยํ โสมนสฺสินฺทฺริยญฺจ จตุตฺถชฺฌาเน นิรุชฺฌนฺติ, ตสฺส อุเปกฺขินฺทฺริยํ อวสิฏฺฐํ ภวติ. โส อุปริมํ สมาปตฺติํ สนฺตโต มนสิกโรติ, ตสฺส อุปริมํ สมาปตฺติํ สนฺตโต มนสิกโรโต จตุตฺถชฺฌาเน โอฬาริกา สญฺญา สณฺฐหติ อุกฺกณฺฐา จ ปฏิฆสญฺญา, โส สพฺพโส รูปสญฺญานํ สมติกฺกมา ปฏิฆ\-สญฺญานํ อตฺถงฺคมา นานตฺต\-สญฺญานํ อมนสิการา “อนนฺตํ อากาสนฺ”ติ อากาสา\-นญฺจา\-ยตนสมาปตฺติํ สจฺฉิกตฺวา อุปสมฺปชฺช วิหรติ. อภิญฺญาภิ\-นีหาโร รูปสญฺญา โวกาโร นานตฺตสญฺญา สมติกฺกมติ ปฏิฆสญฺญา จสฺส อพฺภตฺถํ คจฺฉติ, เอวํ สมาธิ ตสฺส สมาหิตสฺส โอภาโส อนฺตรธายติ ทสฺสนญฺจ รูปานํ, โส สมาธิ ฉฬงฺค\-สมนฺนาคโต \csromanfolio{74}ปจฺจเวกฺขิตพฺโพ. อนภิชฺฌา\-สหคตํ เม มานสํ สพฺพโลเก, อพฺยาปนฺนํ เม จิตฺตํ สพฺพสตฺเตสุ, อารทฺธํ เม วีริยํ ปคฺคหิตํ, ปสฺสทฺโธ เม กาโย อสารทฺโธ, สมาหิตํ เม จิตฺตํ อวิกฺขิตฺตํ, อุปฏฺฐิตา เม สติ อสมฺมุฏฺฐา\csromansharedfootnote{9ac1b0acb3eb891f}{อปฺปมฺมุฏฺฐา (สี)}, ตตฺถ ยญฺจ อนภิชฺฌา\-สหคตํ มานสํ สพฺพโลเก ยญฺจ อพฺยาปนฺนํ จิตฺตํ สพฺพสตฺเตสุ ยญฺจ อารทฺธํ วีริยํ ปคฺคหิตํ ยญฺจ สมาหิตํ จิตฺตํ อวิกฺขิตฺตํ, อยํ สมโถ. โย ปสฺสทฺโธ กาโย อสารทฺโธ, อยํ สมาธิ\-ปริกฺขาโร. ยา อุปฏฺฐิตา สติ อสมฺมุฏฺฐา อยํ วิปสฺสนา.}

\proseitem{54}{โส สมาธิ ปญฺจวิเธน เวทิตพฺโพ อยํ สมาธิ “ปจฺจุปฺปนฺน\-สุโข”ติ อิติสฺส ปจฺจตฺตเมว ญาณทสฺสนํ ปจฺจุปฏฺฐิตํ ภวติ, อยํ สมาธิ “อายติํ สุขวิปาโก”ติ อิติสฺส ปจฺจตฺตเมว ญาณทสฺสนํ ปจฺจุปฏฺฐิตํ ภวติ, อยํ สมาธิ “อริโย นิรามิโส”ติ อิติสฺส ปจฺจตฺตเมว ญาณทสฺสนํ ปจฺจุปฏฺฐิตํ ภวติ, อยํ สมาธิ “อกาปุริสเสวิโต”ติ อิติสฺส ปจฺจตฺตเมว ญาณทสฺสนํ ปจฺจุปฏฺฐิตํ ภวติ, อยํ สมาธิ “สนฺโต เจว ปณีโต จ ปฏิปฺปสฺสทฺธิ\-ลทฺโธ จ เอโกทิภาวา\-ธิคโต จ น สสงฺขาร\-นิคฺคยฺห\-วาริต\-คโต\csromansharedfootnote{d426ffa22b53de19}{สสงฺขารนิคฺคยฺหวาริตวโต (สี), สสงฺขารนิคฺคยฺหวาริวาวโฏ (ก)} จา”ติ อิติสฺส ปจฺจตฺตเมว ญาณทสฺสนํ ปจฺจุปฏฺฐิตํ ภวติ. ตํ โข ปนิมํ สมาธิํ “สโต สมาปชฺชามิ สโต วุฏฺฐหามี”ติ อิติสฺส ปจฺจตฺตเมว ญาณทสฺสนํ ปจฺจุปฏฺฐิตํ ภวติ. ตตฺถ โย จ สมาธิ ปจฺจุปฺปนฺน\-สุโข โย จ สมาธิ อายติํ สุขวิปาโก อยํ สมโถ. โย จ สมาธิ อริโย นิรามิโส, โย จ สมาธิ อกาปุริสเสวิโต, โย จ สมาธิ สนฺโต เจว ปณีโต จ ปฏิปฺปสฺสทฺธิ\-ลทฺโธ จ เอโกทิภาวา\-ธิคโต จ น สสงฺขาร\-นิคฺคยฺห\-วาริต\-คโต จ ยญฺจาหํ ตํ โข ปนิมํ สมาธิํ สโต สมาปชฺชามิ สโต วุฏฺฐหามีติ. อยํ วิปสฺสนา.}

\prose{โส สมาธิ ปญฺจวิเธน เวทิตพฺโพ ปีติผรณตา สุขผรณตา เจโตผรณตา อาโลกผรณตา ปจฺจเวกฺขณา\-นิมิตฺตํ. ตตฺถ โย จ ปีติผรโณ โย จ สุขผรโณ โย จ เจโตผรโณ, อยํ สมโถ. โย จ อาโลกผรโณ ยญฺจ ปจฺจเวกฺขณา\-นิมิตฺตํ. อยํ วิปสฺสนา.}

\proseitem{55}{\csromanfolio{75}ทส กสิณายตนานิ ปถวีกสิณํ อาโปกสิณํ เตโชกสิณํ วาโยกสิณํ นีลกสิณํ ปีตกสิณํ โลหิตกสิณํ โอทาตกสิณํ อากาสกสิณํ วิญฺญาณกสิณํ. ตตฺถ ยญฺจ ปถวีกสิณํ ยญฺจ อาโปกสิณํ เอวํ สพฺพํ, ยญฺจ โอทาตกสิณํ. อิมานิ อฏฺฐ กสิณานิ สมโถ. ยญฺจ อากาสกสิณํ ยญฺจ วิญฺญาณกสิณํ, อยํ วิปสฺสนา. เอวํ สพฺโพ อริโย มคฺโค เยน เยน อากาเรน วุตฺโต, เตน เตน สมถ\-วิปสฺสเนน โยชยิตพฺโพ. เต ตีหิ ธมฺเมหิ สงฺคหิตา อนิจฺจตาย ทุกฺขตาย อนตฺตตาย. โส สมถ\-วิปสฺสนํ ภาวยมาโน ตีณิ วิโมกฺขมุขานิ ภาวยติ. ตีณิ วิโมกฺขมุขานิ ภาวยนฺโต ตโย ขนฺเธ ภาวยติ. ตโย ขนฺเธ ภาวยนฺโต อริยํ อฏฺฐงฺคิกํ มคฺคํ ภาวยติ.}

\prose{ราคจริโต ปุคฺคโล อนิมิตฺเตน วิโมกฺขมุเขน นิยฺยาติ\csromansharedfootnote{e27e06a70fc5326d}{นียาติ (สี)} อธิจิตฺต\-สิกฺขาย สิกฺขนฺโต โลภํ อกุสลมูลํ ปชหนฺโต สุขเวทนียํ ผสฺสํ อนุปคจฺฉนฺโต สุขํ เวทนํ ปริชานนฺโต ราคมลํ ปวาเหนฺโต ราครชํ นิทฺธุนนฺโต ราควิสํ วเมนฺโต ราคคฺคิํ นิพฺพาเปนฺโต ราคสลฺลํ อุปฺปาเฏนฺโต ราคชฏํ วิชเฏนฺโต. โทสจริโต ปุคฺคโล อปฺปณิหิเตน วิโมกฺขมุเขน นิยฺยาติ อธิสีล\-สิกฺขาย สิกฺขนฺโต โทสํ อกุสลมูลํ ปชหนฺโต ทุกฺข\-เวทนียํ ผสฺสํ อนุปคจฺฉนฺโต ทุกฺขเวทนํ ปริชานนฺโต โทสมลํ ปวาเหนฺโต โทสรชํ นิทฺธุนนฺโต โทสวิสํ วเมนฺโต โทสคฺคิํ นิพฺพาเปนฺโต โทสสลฺลํ อุปฺปาเฏนฺโต โทสชฏํ วิชเฏนฺโต. โมหจริโต ปุคฺคโล สุญฺญต\-วิโมกฺขมุเขน นิยฺยาติ อธิปญฺญา\-สิกฺขาย สิกฺขนฺโต โมหํ อกุสลมูลํ ปชหนฺโต อทุกฺขม\-สุข\-เวทนียํ ผสฺสํ อนุปคจฺฉนฺโต อทุกฺขมสุขํ เวทนํ ปริชานนฺโต โมหมลํ ปวาเหนฺโต โมหรชํ นิทฺธุนนฺโต โมหวิสํ วเมนฺโต โมหคฺคิํ นิพฺพาเปนฺโต โมหสลฺลํ อุปฺปาเฏนฺโต โมหชฏํ วิชเฏนฺโต.}

\prose{ตตฺถ สุญฺญต\-วิโมกฺขมุขํ ปญฺญากฺขนฺโธ, อนิมิตฺต\-วิโมกฺขมุขํ สมาธิ\-กฺขนฺโธ, อปฺปณิหิต\-วิโมกฺขมุขํ สีลกฺขนฺโธ, โส ตีณิ วิโมกฺขมุขานิ ภาวยนฺโต ตโย ขนฺเธ ภาวยติ, ตโย ขนฺเธ ภาวยนฺโต อริยํ อฏฺฐงฺคิกํ มคฺคํ ภาวยติ, ตตฺถ ยา จ สมฺมาวาจา โย จ สมฺมากมฺมนฺโต โย จ สมฺมา- อาชีโว, อยํ สีลกฺขนฺโธ, โย จ \csromanfolio{76}สมฺมาวายาโม ยา จ สมฺมาสติ โย จ สมฺมาสมาธิ, อยํ สมาธิ\-กฺขนฺโธ, ยา จ สมฺมาทิฏฺฐิ โย จ สมฺมาสงฺกปฺโป, อยํ ปญฺญากฺขนฺโธ.}

\prose{ตตฺถ สีลกฺขนฺโธ จ สมาธิ\-กฺขนฺโธ จ สมโถ, ปญฺญากฺขนฺโธ วิปสฺสนา, โย สมถ\-วิปสฺสนํ ภาเวติ, ตสฺส เทฺว ภวงฺคานิ ภาวนํ คจฺฉนฺติ กาโย จิตฺตญฺจ, ภวนิโรธ\-คามินี ปฏิปทา เทฺว ปทานิ สีลํ สมาธิ จ. โส โหติ ภิกฺขุ ภาวิตกาโย ภาวิตสีโล ภาวิตจิตฺโต ภาวิตปญฺโญ, กาเย ภาวิยมาเน เทฺว ธมฺมา ภาวนํ คจฺฉนฺติ สมฺมากมฺมนฺโต สมฺมาวายาโม จ, สีเล ภาวิยมาเน เทฺว ธมฺมา ภาวนํ คจฺฉนฺติ สมฺมาวาจา สมฺมาอาชีโว จ, จิตฺเต ภาวิยมาเน เทฺว ธมฺมา ภาวนํ คจฺฉนฺติ สมฺมาสติ สมฺมาสมาธิ จ, ปญฺญาย ภาวิยมานาย เทฺว ธมฺมา ภาวนํ คจฺฉนฺติ สมฺมาทิฏฺฐิ สมฺมาสงฺกปฺโป จ.}

\prose{ตตฺถ โย จ สมฺมากมฺมนฺโต โย จ สมฺมาวายาโม สิยา กายิโก สิยา เจตสิโก, ตตฺถ โย กายสงฺคโห, โส กาเย ภาวิเต ภาวนํ คจฺฉติ, โย จิตฺตสงฺคโห, โส จิตฺเต ภาวิเต ภาวนํ คจฺฉติ, โส สมถ\-วิปสฺสนํ ภาวยนฺโต ปญฺจวิธํ อธิคมํ คจฺฉติ\csromansharedfootnote{50fd70cd1daa79d6}{อธิคจฺฉติ (สี)} ขิปฺปาธิคโม จ โหติ, วิมุตฺตาธิคโม จ โหติ, มหาธิคโม จ โหติ, วิปุลาธิคโม จ โหติ, อนวเสสา\-ธิคโม จ โหติ, ตตฺถ สมเถน ขิปฺปาธิคโม จ มหาธิคโม จ วิปุลาธิคโม จ โหติ, วิปสฺสนาย วิมุตฺตาธิคโม จ อนวเสสา\-ธิคโม จ โหติ.}

\proseitem{56}{ตตฺถ โย เทสยติ, โส ทสพล\-สมนฺนาคโต สตฺถา โอวาเทน สาวเก น วิสํวาทยติ, โส ติวิธํ อิทํ กโรถ อิมินา อุปาเยน กโรถ อิทํ โว กุรุมานานํ หิตาย สุขาย ภวิสฺสติ, โส ตถา โอวทิโต ตถานุสิฏฺโฐ ตถากโรนฺโต ตถา\-ปฏิปชฺชนฺโต ตํ ภูมิํ น ปาปุณิสฺสตีติ เนตํ ฐานํ วิชฺชติ, โส ตถา โอวทิโต ตถานุสิฏฺโฐ สีลกฺขนฺธํ อปริปูรยนฺโต ตํ ภูมิํ อนุปาปุณิสฺสตีติ เนตํ ฐานํ วิชฺชติ, โส ตถา โอวทิโต ตถานุสิฏฺโฐ สีลกฺขนฺธํ ปริปูรยนฺโต ตํ ภูมิํ อนุปาปุณิสฺสตีติ ฐานเมตํ วิชฺชติ.}

\prose{2 สมฺมา\-สมฺพุทฺธสฺส เต สโต อิเม ธมฺมา อนภิสมฺพุทฺธาติ เนตํ ฐานํ วิชฺชติ, สพฺพาสว\-ปริกฺขีณสฺส เต สโต อิเม อาสวา อปริกฺขีณาติ \csromanfolio{77}เนตํ ฐานํ วิชฺชติ, ยสฺส เต อตฺถาย ธมฺโม เทสิโต, โส น นิยฺยาติ ตกฺกรสฺส สมฺมา ทุกฺขกฺขยายาติ เนตํ ฐานํ วิชฺชติ, สาวโก โข ปน เต ธมฺมา\-นุธมฺม\-ปฺปฏิปนฺโน สามีจิ\-ปฺปฏิปนฺโน อนุธมฺมจารี โส ปุพฺเพน อปรํ อุฬารํ วิเสสาธิคมํ น สจฺฉิกริสฺสตีติ เนตํ ฐานํ วิชฺชติ.}

\prose{เย โข ปน ธมฺมา อนฺตรายิกา, เต ปฏิเสวโต นาลํ อนฺตรายายาติ เนตํ ฐานํ วิชฺชติ. เย โข ปน ธมฺมา อนิยฺยานิกา, เต นิยฺยนฺติ ตกฺกรสฺส สมฺมา ทุกฺขกฺขยายาติ เนตํ ฐานํ วิชฺชติ. เย โข ปน ธมฺมา นิยฺยานิกา, เต นิยฺยนฺติ ตกฺกรสฺส สมฺมา ทุกฺขกฺขยายาติ ฐานเมตํ วิชฺชติ, สาวโก โข ปน เต สอุปาทิเสโส อนุปาทิเสสํ นิพฺพานธาตุํ อนุปาปุณิสฺสตีติ เนตํ ฐานํ วิชฺชติ.}

\prose{1 ทิฏฺฐิสมฺปนฺโน มาตรํ ชีวิตา โวโรเปยฺย หตฺเถหิ วา ปาเทหิ วา สุหตํ กเรยฺยาติ เนตํ ฐานํ วิชฺชติ, ปุถุชฺชโน มาตรํ ชีวิตา โวโรเปยฺย หตฺเถหิ วา ปาเทหิ วา สุหตํ กเรยฺยาติ ฐานเมตํ วิชฺชติ, เอวํ ปิตรํ, อรหนฺตํ ภิกฺขุํ, ทิฏฺฐิสมฺปนฺโน ปุคฺคโล สํฆํ ภินฺเทยฺย สํเฆ วา สํฆราชิํ ชเนยฺยาติ เนตํ ฐานํ วิชฺชติ, ปุถุชฺชโน สํฆํ ภินฺเทยฺย สํเฆ วา สํฆราชิํ ชเนยฺยาติ ฐานเมตํ วิชฺชติ, ทิฏฺฐิสมฺปนฺโน ตถาคตสฺส ทุฏฺฐจิตฺโต โลหิตํ อุปฺปาเทยฺย, ปรินิพฺพุตสฺส วา ตถาคตสฺส ทุฏฺฐจิตฺโต ถูปํ ภินฺเทยฺยาติ เนตํ ฐานํ วิชฺชติ, ปุถุชฺชโน ตถาคตสฺส ทุฏฺฐจิตฺโต โลหิตํ อุปฺปาเทยฺย, ปรินิพฺพุตสฺส วา ตถาคตสฺส ทุฏฺฐจิตฺโต ถูปํ ภินฺเทยฺยาติ ฐานเมตํ วิชฺชติ. ทิฏฺฐิสมฺปนฺโน อญฺญํ สตฺถารํ อปทิเสยฺย อปิ ชีวิตเหตูติ เนตํ ฐานํ วิชฺชติ, ปุถุชฺชโน อญฺญํ สตฺถารํ อปทิเสยฺยาติ ฐานเมตํ วิชฺชติ, ทิฏฺฐิสมฺปนฺโน อิโต พหิทฺธา อญฺญํ ทกฺขิเณยฺยํ ปริเยเสยฺยาติ เนตํ ฐานํ วิชฺชติ, ปุถุชฺชโน อิโต พหิทฺธา อญฺญํ ทกฺขิเณยฺยํ ปริเยเสยฺยาติ ฐานเมตํ วิชฺชติ, ทิฏฺฐิสมฺปนฺโน กุตูหลมงฺคเลน สุทฺธิํ ปจฺเจยฺยาติ เนตํ ฐานํ วิชฺชติ, ปุถุชฺชโน กุตูหลมงฺคเลน สุทฺธิํ ปจฺเจยฺยาติ ฐานเมตํ วิชฺชติ.}

\proseitem{57}{1 อิตฺถี ราชา จกฺกวตฺตี สิยาติ เนตํ ฐานํ วิชฺชติ, ปุริโส ราชา จกฺกวตฺตี สิยาติ ฐานเมตํ วิชฺชติ, อิตฺถี สกฺโก เทวานมินฺโท สิยาติ เนตํ ฐานํ วิชฺชติ, ปุริโส สกฺโก เทวานมินฺโท สิยาติ ฐานเมตํ วิชฺชติ, อิตฺถี มาโร ปาปิมา สิยาติ เนตํ ฐานํ วิชฺชติ, ปุริโส มาโร \csromanfolio{78}ปาปิมา สิยาติ ฐานเมตํ วิชฺชติ, อิตฺถี มหาพฺรหฺมา สิยาติ เนตํ ฐานํ วิชฺชติ, ปุริโส มหาพฺรหฺมา สิยาติ ฐานเมตํ วิชฺชติ, อิตฺถี ตถาคโต อรหํ สมฺมาสมฺพุทฺโธ สิยาติ เนตํ ฐานํ วิชฺชติ, ปุริโส ตถาคโต อรหํ สมฺมาสมฺพุทฺโธ สิยาติ ฐานเมตํ วิชฺชติ.}

\prose{เทฺว ตถาคตา อรหนฺโต สมฺมาสมฺพุทฺธา อปุพฺพํ อจริมํ เอกิสฺสา โลกธาตุยา อุปฺปชฺเชยฺยุํ วา ธมฺมํ วา เทเสยฺยุนฺติ เนตํ ฐานํ วิชฺชติ, เอโกว ตถาคโต อรหํ สมฺมาสมฺพุทฺโธ เอกิสฺสา โลกธาตุยา อุปฺปชฺชิสฺสติ วา ธมฺมํ วา เทเสสฺสตีติ ฐานเมตํ วิชฺชติ.}

\prose{ติณฺณํ ทุจฺจริตานํ อิฏฺโฐ กนฺโต ปิโย มนาโป วิปาโก ภวิสฺสตีติ เนตํ ฐานํ วิชฺชติ, ติณฺณํ ทุจฺจริตานํ อนิฏฺโฐ อกนฺโต อปฺปิโย อมนาโป วิปาโก ภวิสฺสตีติ ฐานเมตํ วิชฺชติ, ติณฺณํ สุจริตานํ อนิฏฺโฐ อกนฺโต อปฺปิโย อมนาโป วิปาโก ภวสฺสตีติ เนตํ ฐานํ วิชฺชติ, ติณฺณํ สุจริตานํ อิฏฺโฐ กนฺโต ปิโย มนาโป วิปาโก ภวิสฺสตีติ ฐานเมตํ วิชฺชติ.}

\prose{อญฺญตโร สมโณ วา พฺราหฺมโณ วา กุหโก ลปโก เนมิตฺตโก กุหนลปน\-เนมิตฺต\-กตฺตํ ปุพฺพงฺคมํ กตฺวา ปญฺจ นีวรเณ อปฺปหาย เจตโส อุปกฺกิเลเส ปญฺญาย ทุพฺพลีกรเณ จตูสุ สติปฏฺฐาเนสุ อนุปฏฺฐิต\-สฺสติ\csromansharedfootnote{c4d69e08e7fc5b0b}{อนุปฏฺฐิตสติ (สี)} วิหรนฺโต สตฺต โพชฺฌงฺเค อภาวยิตฺวา อนุตฺตรํ สมฺมาสมฺโพธิํ อภิสมฺพุชฺฌิสฺสตีติ เนตํ ฐานํ วิชฺชติ, อญฺญตโร สมโณ วา พฺราหฺมโณ วา สพฺพ\-โทสาปคโต ปญฺจ นีวรเณ ปหาย เจตโส อุปกฺกิเลเส ปญฺญาย ทุพฺพลีกรเณ จตูสุ สติปฏฺฐาเนสุ อุปฏฺฐิตสฺสติ วิหรนฺโต สตฺต โพชฺฌงฺเค ภาวยิตฺวา อนุตฺตรํ สมฺมาสมฺโพธิํ อภิสมฺพุชฺฌิสฺสตีติ ฐานเมตํ วิชฺชติ. ยํ เอตฺถ ญาณํ เหตุโส ฐานโส อโนธิโส อิทํ วุจฺจติ ฐานาฏฺฐาน\-ญาณํ ปฐมํ ตถาคตพลํ. (1)}

\prose{อิติ ฐานาฏฺฐาน\-คตา สพฺเพ ขยธมฺมา วยธมฺมา วิราคธมฺมา นิโรธธมฺมา เกจิ สคฺคูปคา เกจิ อปายูปคา เกจิ นิพฺพานูปคา, เอวํ ภควา อาห\csromandash{}}

\setlength{\csromangathapagewidth}{0pt}
\setlength{\csromangathawakwidth}{0pt}
\csromangathameasuregroup{1 สพฺเพ สตฺตา มริสฺสนฺติ, \\ ยถากมฺมํ คมิสฺสนฺติ, \\ นิรยํ ปาปกมฺมนฺตา, \\ อปเร จ มคฺคํ ภาเวตฺวา,}{\csromangathabat{1 สพฺเพ สตฺตา มริสฺสนฺติ,}{มรณนฺตํ หิ ชีวิตํ.} \\ \csromangathabat{ยถากมฺมํ คมิสฺสนฺติ,}{ปุญฺญปาปผลูปคา.} \\ \csromangathabat{นิรยํ ปาปกมฺมนฺตา,}{ปุญฺญกมฺมา จ สุคฺคติํ1.} \\ \csromangathabat{อปเร จ มคฺคํ ภาเวตฺวา,}{ปรินิพฺพนฺติ’นาสวาติ.}}
\csromangathasetleft{1 สพฺเพ สตฺตา มริสฺสนฺติ, \\ ยถากมฺมํ คมิสฺสนฺติ, \\ นิรยํ ปาปกมฺมนฺตา, \\ อปเร จ มคฺคํ ภาเวตฺวา,}
\csromangathagroup{\csromangathastanza{\csromanfolio{79}\csromangathaitembat{58}{1 สพฺเพ สตฺตา มริสฺสนฺติ,}{มรณนฺตํ หิ ชีวิตํ.} \\ \csromangathacontbat{ยถากมฺมํ คมิสฺสนฺติ,}{ปุญฺญ\-ปาป\-ผลู\-ปคา.} \\ \csromangathacontbat{นิรยํ ปาปกมฺมนฺตา,}{ปุญฺญกมฺมา จ สุคฺคติํ1.} \\ \csromangathacontbat{อปเร จ มคฺคํ ภาเวตฺวา,}{ปรินิพฺพนฺติ’นาสวาติ.}}}

\prose{\textbf{สพฺเพ สตฺตา}ติ อริยา จ อนริยา จ สกฺกาย\-ปริยาปนฺนา จ สกฺกาย\-วีติวตฺตา จ. \textbf{มริสฺสนฺตี}ติ ทฺวีหิ มรเณหิ ทนฺธมรเณน จ อทนฺธ\-มรเณน จ, สกฺกาย\-ปริยาปนฺนานํ อทนฺธมรณํ สกฺกาย\-วีติวตฺตานํ ทนฺธมรณํ. \textbf{มรณนฺตํ หิ ชีวิตนฺ}ติ ขยา อายุสฺส อินฺทฺริยานํ อุปโรธา ชีวิต\-ปริยนฺโต มรณ\-ปริยนฺโต. \textbf{ยถากมฺมํ คมิสฺสนฺตี}ติ กมฺมสฺสกตา. \textbf{ปุญฺญ}\-\textbf{ปาป}\-\textbf{ผลู}\-\textbf{ปคา}ติ กมฺมานํ ผลทสฺสาวิตา จ อวิปฺปวาโส จ.}

\prose{\textbf{นิรยํ ปาปกมฺมนฺตา}ติ อปุญฺญสงฺขารา. \textbf{ปุญฺญกมฺมา จ สุคฺคตินฺ}ติ ปุญฺญสงฺขารา สุคติํ คมิสฺสนฺติ. \textbf{อปเร จ มคฺคํ ภาเวตฺวา, ปรินิพฺพนฺติ’นาสวา}ติ สพฺพ\-สงฺขารานํ สมติกฺกมนํ. เตนาห ภควา\csromandash{}“สพฺเพ ฯเปฯ นาสวา”ติ.}

\prose{“สพฺเพ สตฺตา มริสฺสนฺติ, มรณนฺตํ หิ ชีวิตํ. ยถากมฺมํ คมิสฺสนฺติ, ปุญฺญ\-ปาป\-ผลู\-ปคา. นิรยํ ปาปกมฺมนฺตา”ติ อาคาฬฺหา จ นิชฺฌามา จ ปฏิปทา. “อปเร จ มคฺคํ ภาเวตฺวา, ปรินิพฺพนฺติ’นาสวา”ติ มชฺฌิมา ปฏิปทา. “สพฺเพ สตฺตา มริสฺสนฺติ, มรณนฺตํ หิ ชีวิตํ, ยถากมฺมํ คมิสฺสนฺติ, ปุญฺญ\-ปาป\-ผลู\-ปคา, นิรยํ ปาปกมฺมนฺตา”ติ อยํ สํกิเลโส. เอวํ สํสารํ นิพฺพตฺตยติ. “สพฺเพ สตฺตา มริสฺสนฺติ ฯเปฯ นิรยํ ปาปกมฺมนฺตา”ติ อิเม ตโย วฏฺฏา ทุกฺขวฏฺโฏ กมฺมวฏฺโฏ กิเลสวฏฺโฏ. “อปเร จ มคฺคํ ภาเวตฺวา, ปรินิพฺพนฺติ’นาสวา”ติ ติณฺณํ วฏฺฏานํ วิวฏฺฏนา. “สพฺเพ สตฺตา มริสฺสนฺติ ฯเปฯ นิรยํ ปาปกมฺมนฺตา”ติ อาทีนโว, “ปุญฺญกมฺมา จ สุคฺคตินฺ”ติ อสฺสาโท, “อปเร จ มคฺคํ ภาเวตฺวา, ปรินิพฺพนฺติ’นาสวา”ติ นิสฺสรณํ. “สพฺเพ สตฺตา มริสฺสนฺติ ฯเปฯ นิรยํ ปาปกมฺมนฺตา”ติ เหตุ จ ผลญฺจ, ปญฺจกฺขนฺธา ผลํ, ตณฺหา เหตุ, “อปเร จ มคฺคํ ภาเวตฺวา, ปรินิพฺพนฺติ’นาสวา”ติ มคฺโค จ ผลญฺจ. “สพฺเพ สตฺตา มริสฺสนฺติ, มรณนฺตํ หิ ชีวิตํ. ยถากมฺมํ คมิสฺสนฺติ, ปุญฺญ\-ปาป\-ผลู\-ปคา, นิรยํ ปาปกมฺมนฺตา”ติ อยํ สํกิเลโส, โส สํกิเลโส ติวิโธ ตณฺหาสํกิเลโส ทิฏฺฐิ\-สํกิเลโส ทุจฺจริตสํกิเลโสติ.}

\proseitem{59}{\csromanfolio{80}ตตฺถ ตณฺหาสํกิเลโส ตีหิ ตณฺหาหิ นิทฺทิสิตพฺโพ กามตณฺหาย ภวตณฺหาย วิภวตณฺหาย, เยน เยน วา ปน วตฺถุนา อชฺโฌสิโต, เตน เตเนว นิทฺทิสิตพฺโพ, ตสฺสา วิตฺถาโร ฉตฺติํสาย ตณฺหาย ชาลินิยา วิจริตานิ. ตตฺถ ทิฏฺฐิ\-สํกิเลโส อุจฺเฉท\-สสฺสเตน นิทฺทิสิตพฺโพ, เยน เยน วา ปน วตฺถุนา ทิฏฺฐิวเสน อภินิวิสติ, “อิทเมว สจฺจํ โมฆมญฺญนฺ”ติ เตน เตเนว นิทฺทิสิตพฺโพ, ตสฺสา วิตฺถาโร ทฺวาสฏฺฐิ ทิฏฺฐิคตานิ. ตตฺถ ทุจฺจริต\-สํกิเลโส เจตนา เจตสิก\-กมฺเมน นิทฺทิสิตพฺโพ, ตีหิ ทุจฺจริเตหิ กาย\-ทุจฺจริเตน วจีทุจฺจริเตน มโน\-ทุจฺจริเตน, ตสฺสา วิตฺถาโร ทส อกุสล\-กมฺมปถา. อปเร จ มคฺคํ ภาเวตฺวา, ปรินิพฺพนฺติ’นาสวาติ อิทํ โวทานํ.}

\prose{ตยิทํ โวทานํ ติวิธํ, ตณฺหาสํกิเลโส สมเถน วิสุชฺฌติ, โส สมโถ สมาธิ\-กฺขนฺโธ, ทิฏฺฐิ\-สํกิเลโส วิปสฺสนาย วิสุชฺฌติ, สา วิปสฺสนา ปญฺญากฺขนฺโธ, ทุจฺจริต\-สํกิเลโส สุจริเตน วิสุชฺฌติ, ตํ สุจริตํ สีลกฺขนฺโธ.}

\prose{“สพฺเพ สตฺตา มริสฺสนฺติ, มรณนฺตํ หิ ชีวิตํ, ยถากมฺมํ คมิสฺสนฺติ, ปุญฺญ\-ปาป\-ผลู\-ปคา, นิรยํ ปาปกมฺมนฺตา”ติ อปุญฺญ\-ปฺปฏิปทา, “ปุญฺญกมฺมา จ สุคฺคตินฺ”ติ ปุญฺญปฺปฏิปทา, “อปเร จ มคฺคํ ภาเวตฺวา, ปรินิพฺพนฺติ’นาสวา”ติ ปุญฺญปาป\-สมติกฺกม\-ปฺปฏิปทา, ตตฺถ ยา จ ปุญฺญปฺปฏิปทา ยา จ อปุญฺญ\-ปฺปฏิปทา, อยํ เอกา ปฏิปทา สพฺพตฺถ\-คามินี เอกา อปาเยสุ, เอกา เทเวสุ, ยา จ ปุญฺญปาป\-สมติกฺกมา ปฏิปทา อยํ ตตฺถ ตตฺถ คามินี ปฏิปทา.}

\prose{ตโย ราสี มิจฺฉตฺต\-นิยโต ราสิ สมฺมตฺตนิยโต ราสิ อนิยโต ราสิ, ตตฺถ โย จ มิจฺฉตฺต\-นิยโต ราสิ โย จ สมฺมตฺตนิยโต ราสิ เอกา ปฏิปทา ตตฺถ ตตฺถ คามินี, ตตฺถ โย อนิยโต ราสิ, อยํ สพฺพตฺถ\-คามินี ปฏิปทา, เกน การเณน ปจฺจยํ ลภนฺโต นิรเย อุปปชฺเชยฺย, ปจฺจยํ ลภนฺโต ติรจฺฉาน\-โยนีสุ อุปปชฺเชยฺย, ปจฺจยํ ลภนฺโต เปตฺติวิสเยสุ อุปปชฺเชยฺย, ปจฺจยํ ลภนฺโต อสุเรสุ อุปปชฺเชยฺย, ปจฺจยํ ลภนฺโต เทเวสุ อุปปชฺเชยฺย, ปจฺจยํ ลภนฺโต มนุสฺเสสุ อุปปชฺเชยฺย, ปจฺจยํ ลภนฺโต ปรินิพฺพาเยยฺย, ตสฺมายํ สพฺพตฺถ\-คามินี ปฏิปทา, ยํ เอตฺถ ญาณํ เหตุโส ฐานโส อโนธิโส, อิทํ วุจฺจติ สพฺพตฺถ\-คามินี ปฏิปทา ญาณํ ทุติยํ ตถาคตพลํ. (2)}

\prose{\csromanfolio{81}อิติ สพฺพตฺถ\-คามินี ปฏิปทา อเนกธาตุโลโก, ตตฺถ ตตฺถ คามินี ปฏิปทา นานาธาตุโลโก. ตตฺถ กตโม อเนกธาตุโลโก, จกฺขุธาตุ รูปธาตุ จกฺขุวิญฺญาณ\-ธาตุ, โสตธาตุ สทฺทธาตุ โสตวิญฺญาณ\-ธาตุ, ฆานธาตุ คนฺธธาตุ ฆานวิญฺญาณ\-ธาตุ, ชิวฺหาธาตุ รสธาตุ ชิวฺหาวิญฺญาณ\-ธาตุ, กายธาตุ โผฏฺฐพฺพ\-ธาตุ กายวิญฺญาณ\-ธาตุ, มโนธาตุ ธมฺมธาตุ มโนวิญฺญาณ\-ธาตุ, ปถวีธาตุ, อาโปธาตุ, เตโชธาตุ, วาโยธาตุ, อากาสธาตุ, วิญฺญาณธาตุ, กามธาตุ, พฺยาปาทธาตุ, วิหิํสาธาตุ, เนกฺขมฺมธาตุ, อพฺยาปาทธาตุ, อวิหิํสาธาตุ, ทุกฺขธาตุ, โทมนสฺสธาตุ, อวิชฺชาธาตุ, สุขธาตุ, โสมนสฺสธาตุ, อุเปกฺขาธาตุ, รูปธาตุ, อรูปธาตุ, นิโรธธาตุ, สงฺขารธาตุ, นิพฺพานธาตุ, อยํ อเนกธาตุโลโก.}

\prose{ตตฺถ กตโม นานาธาตุโลโก, อญฺญา จกฺขุธาตุ, อญฺญา รูปธาตุ, อญฺญา จกฺขุวิญฺญาณ\-ธาตุ, เอวํ สพฺพา, อญฺญา นิพฺพานธาตุ. ยํ เอตฺถ ญาณํ เหตุโส ฐานโส อโนธิโส, อิทํ วุจฺจติ อเนกธาตุ นานาธาตุ ญาณํ ตติยํ ตถาคตพลํ. (3)}

\proseitem{60}{อิติ อเนกธาตุ นานาธาตุกสฺส โลกสฺส ยํ ยเทว ธาตุํ สตฺตา อธิมุจฺจนฺติ, ตํ ตเทว อธิฏฺฐหนฺติ อภินิวิสนฺติ, เกจิ รูปาธิมุตฺตา, เกจิ สทฺทาธิมุตฺตา, เกจิ คนฺธาธิมุตฺตา, เกจิ รสาธิมุตฺตา, เกจิ โผฏฺฐพฺพา\-ธิมุตฺตา, เกจิ ธมฺมาธิมุตฺตา, เกจิ อิตฺถาธิมุตฺตา, เกจิ ปุริสาธิมุตฺตา, เกจิ จาคาธิมุตฺตา, เกจิ หีนาธิมุตฺตา, เกจิ ปณีตาธิมุตฺตา, เกจิ เทวาธิมุตฺตา, เกจิ มนุสฺสา\-ธิมุตฺตา, เกจิ นิพฺพานา\-ธิมุตฺตา. ยํ เอตฺถ ญาณํ เหตุโส ฐานโส อโนธิโส, อยํ เวเนยฺโย, อยํ น เวเนยฺโย, อยํ สคฺคคามี, อยํ ทุคฺคติคามีติ, อิทํ วุจฺจติ สตฺตานํ นานา\-ธิมุตฺติกตา ญาณํ จตุตฺถํ ตถาคตพลํ. (4)}

\prose{อิติ เต ยถาธิมุตฺตา จ ภวนฺติ, ตํ ตํ กมฺมสมาทานํ สมาทิยนฺติ, เต ฉพฺพิธํ กมฺมํ สมาทิยนฺติ, เกจิ โลภวเสน, เกจิ โทสวเสน, เกจิ โมหวเสน, เกจิ สทฺธาวเสน, เกจิ วีริยวเสน, เกจิ ปญฺญาวเสน, ตํ วิภชฺชมานํ ทุวิธํ สํสารคามิ จ นิพฺพานคามิ จ.}

\prose{ตตฺถ ยํ โลภวเสน โทสวเสน โมหวเสน จ กมฺมํ กโรติ, อิทํ กมฺมํ กณฺหํ กณฺหวิปากํ, ตตฺถ ยํ สทฺธาวเสน กมฺมํ กโรติ, อิทํ กมฺมํ \csromanfolio{82}สุกฺกํ สุกฺกวิปากํ. ตตฺถ ยํ โลภวเสน โทสวเสน โมหวเสน สทฺธาวเสน จ กมฺมํ กโรติ, อิทํ กมฺมํ กณฺหสุกฺกํ กณฺห\-สุกฺก\-วิปากํ. ตตฺถ ยํ วีริยวเสน ปญฺญาวเสน จ กมฺมํ กโรติ, อิทํ กมฺมํ อกณฺหํ อสุกฺกํ อกณฺห\-อสุกฺก\-วิปากํ กมฺมุตฺตมํ กมฺมเสฏฺฐํ กมฺมกฺขยาย สํวตฺตติ.}

\prose{จตฺตาริ กมฺม\-สมาทานานิ อตฺถิ กมฺมสมาทานํ ปจฺจุปฺปนฺนสุขํ อายติํ ทุกฺขวิปากํ, อตฺถิ กมฺมสมาทานํ ปจฺจุปฺปนฺน\-ทุกฺขํ อายติํ สุขวิปากํ, อตฺถิ กมฺมสมาทานํ ปจฺจุปฺปนฺนทุกฺข\-ญฺเจว อายติํ จ ทุกฺขวิปากํ, อตฺถิ กมฺมสมาทานํ ปจฺจุปฺปนฺนสุข\-ญฺเจว อายติํ จ สุขวิปากํ, ยํ เอวํ ชาติยํ กมฺมสมาทานํ, อิมินา ปุคฺคเลน อกุสล\-กมฺมสมาทานํ อุปจิตํ อวิปกฺกํ วิปากาย ปจฺจุปฏฺฐิตํ น จ ภพฺโพ อภินิพฺพิธา คนฺตุนฺติ ตํ ภควา น โอวทติ. ยถา เทวทตฺตํ โกกาลิกํ สุนกฺขตฺตํ ลิจฺฉวิปุตฺตํ, เย วา ปนญฺเญปิ สตฺตา มิจฺฉตฺต\-นิยตา อิเมสญฺจ ปุคฺคลานํ อุปจิตํ อกุสลํ น จ ตาว ปาริปูริํ คตํ, ปุรา ปาริปูริํ คจฺฉติ. ปุรา ผลํ นิพฺพตฺตยติ, ปุรา มคฺค\-มา\-วารยติ, ปุรา เวเนยฺยตฺตํ สมติ\-กฺกมตีติ เต ภควา อสมตฺเต โอวทติ. ยถา ปุณฺณญฺจ โควติกํ อเจลญฺจ กุกฺกุรวติกํ.}

\proseitem{61}{อิมสฺส จ ปุคฺคลสฺส อกุสล\-กมฺมสมาทานํ ปริปูรมานํ มคฺคํ อาวารยิสฺสติ ปุรา ปาริปูริํ คจฺฉติ, ปุรา ผลํ นิพฺพตฺตยติ, ปุรา มคฺค\-มา\-วารยติ, ปุรา เวเนยฺยตฺตํ สมติ\-กฺกมตีติ ตํ ภควา อสมตฺตํ โอวทติ. ยถา อายสฺมนฺตํ องฺคุลิมาลํ.}

\prose{สพฺเพสํ มุทุมชฺฌา\-ธิมตฺตตา, ตตฺถ มุทุ อาเนญฺชา\-ภิสงฺขารา มชฺฌํ อวเสส\-กุสล\-สงฺขารา, อธิมตฺตํ อกุสล\-สงฺขารา, ยํ เอตฺถ ญาณํ เหตุโส ฐานโส อโนธิโส, อิทํ ทิฏฺฐธมฺม\-เวทนียํ, อิทํ อุปปชฺช\-เวทนียํ, อิทํ อปราปริย\-เวทนียํ, อิทํ นิรย\-เวทนียํ, อิทํ ติรจฺฉาน\-เวทนียํ, อิทํ เปตฺติวิสยเวทนียํ, อิทํ อสุรเวทนียํ, อิทํ เทวเวทนียํ, อิทํ มนุสฺสเวทนียนฺติ, อิทํ วุจฺจติ อตีตา\-นาคต\-ปจฺจุปฺปนฺนานํ กมฺม\-สมาทานานํ เหตุโส ฐานโส อโนธิโส วิปาก\-เวมตฺตตา ญาณํ ปญฺจมํ ตถาคตพลํ. (5)}

\proseitem{62}{อิติ ตถา สมาทินฺนานํ กมฺมานํ สมาทินฺนานํ ฌานานํ วิโมกฺขานํ สมาธีนํ สมาปตฺตีนํ อยํ สํกิเลโส, อิทํ โวทานํ, อิทํ วุฏฺฐานํ, เอวํ สํกิลิสฺสติ, เอวํ โวทายติ, เอวํ วุฏฺฐหตีติ ญาณํ อนาวรณํ.}

\prose{\csromanfolio{83}ตตฺถ กติ ฌานานิ, จตฺตาริ ฌานานิ. กติ วิโมกฺขา, เอกาทส จ อฏฺฐ จ สตฺต จ ตโย จ เทฺว จ. กติ สมาธี, ตโย สมาธี สวิตกฺโก สวิจาโร สมาธิ, อวิตกฺโก วิจารมตฺโต สมาธิ, อวิตกฺโก อวิจาโร สมาธิ. กติ สมาปตฺติโย, ปญฺจ สมาปตฺติโย สญฺญาสมาปตฺติ อสญฺญา\-สมาปตฺติ เนว\-สญฺญา\-นา\-สญฺญาสมาปตฺติ วิภูต\-สญฺญาสมาปตฺติ\csromansharedfootnote{2bdf0fea25d81940}{วิภูตสมาปตฺติ (สี, ก) 59 ปิฏฺเฐ นตฺถิ ปาฐนานตฺตํ.} นิโรธ\-สมาปตฺติ.}

\prose{ตตฺถ กตโม สํกิเลโส, ปฐม\-ชฺฌานสฺส กามราค\-พฺยาปาทา สํกิเลโส เย จ กุกฺกุฏฌายี เทฺว ปฐมกา โย วา ปน โกจิ หานภาคิโย สมาธิ, อยํ สํกิเลโส, ตตฺถ กตมํ โวทานํ, นีวรณ\-ปาริสุทฺธิ, ปฐมสฺส ฌานสฺส เย จ กุกฺกุฏฌายี เทฺว ปจฺฉิมกา โย วา ปน โกจิ วิเสสภาคิโย สมาธิ, อิทํ โวทานํ, ตตฺถ กตมํ วุฏฺฐานํ, ยํ สมาปตฺติ\-วุฏฺฐาน\-โกสลฺลํ, อิทํ วุฏฺฐานํ, ยํ เอตฺถ ญาณํ เหตุโส ฐานโส อโนธิโส, อิทํ วุจฺจติ สพฺเพสํ ฌาน\-วิโมกฺข\-สมาธิ\-สมาปตฺตีนํ สํกิเลส\-โวทาน\-วุฏฺฐาน\-ญาณํ ฉฏฺฐํ ตาถาคตพลํ. (6)}

\proseitem{63}{อิติ ตสฺเสว สมาธิสฺส ตโย ธมฺมา ปริวารา อินฺทฺริยานิ พลานิ วีริยมิติ, ตานิเยว อินฺทฺริยานิ วีริยวเสน พลานิ ภวนฺติ, อาธิปเตยฺย\-ฏฺเฐน อินฺทฺริยานิ, อกมฺปิยฏฺเฐน พลานิ, อิติ เตสํ มุทุมชฺฌา\-ธิมตฺตตา อยํ มุทินฺทฺริโย อยํ มชฺฌินฺทฺริโย อยํ ติกฺขินฺทฺริโยติ. ตตฺถ ภควา ติกฺขินฺทฺริยํ สํขิตฺเตน โอวาเทน โอวทติ, มชฺฌินฺทฺริยํ ภควา สํขิตฺต\-วิตฺถาเรน โอวทติ, มุทินฺทฺริยํ ภควา วิตฺถาเรน โอวทติ. ตตฺถ ภควา ติกฺขิ\-นฺทฺริยสฺส มุทุกํ ธมฺมเทสนํ อุปทิสติ, มชฺฌิ\-นฺทฺริยสฺส ภควา มุทุติกฺข\-ธมฺมเทสนํ อุปทิสติ, มุทินฺทฺริยสฺส ภควา ติกฺขํ ธมฺมเทสนํ อุปทิสติ. ตตฺถ ภควา ติกฺขิ\-นฺทฺริยสฺส สมถํ อุปทิสติ, มชฺฌิ\-นฺทฺริยสฺส ภควา สมถ\-วิปสฺสนํ อุปทิสติ, มุทินฺทฺริยสฺส ภควา วิปสฺสนํ อุปทิสติ. ตตฺถ ภควา ติกฺขิ\-นฺทฺริยสฺส นิสฺสรณํ อุปทิสติ, มชฺฌิ\-นฺทฺริยสฺส ภควา อาทีนวญฺจ นิสฺสรณญฺจ อุปทิสติ, มุทินฺทฺริยสฺส ภควา อสฺสาทญฺจ อาทีนวญฺจ นิสฺสรณญฺจ อุปทิสติ. ตตฺถ ภควา ติกฺขิ\-นฺทฺริยสฺส อธิปญฺญา\-สิกฺขาย ปญฺญาปยติ, มชฺฌิ\-นฺทฺริยสฺส ภควา อธิจิตฺต\-สิกฺขาย ปญฺญาปยติ, มุทินฺทฺริยสฺส ภควา อธิสีล\-สิกฺขาย ปญฺญาปยติ.}

\prose{\csromanfolio{84}ยํ เอตฺถ ญาณํ เหตุโส ฐานโส อโนธิโส อยํ อิมํ ภูมิํ ภาวนญฺจ คโต, อิมาย เวลาย อิมาย อนุสาสนิยา เอวํ ธาตุโก จายํ อยํ จสฺส อาสโย อยญฺจ อนุสโย อิติ, อิทํ วุจฺจติ ปรสตฺตานํ ปรปุคฺคลานํ อินฺทฺริยปโรปริยตฺต\-เวมตฺตตา ญาณํ สตฺตมํ ตถาคตพลํ. (7)}

\prose{อิติ ตตฺถ ยํ อเนกวิหิตํ ปุพฺเพนิวาสํ อนุสฺสรติ. เสยฺยถิทํ, เอกมฺปิ ชาติํ เทฺวปิ ชาติโย ติสฺโสปิ ชาติโย จตสฺโสปิ ชาติโย ปญฺจปิ ชาติโย ทสปิ ชาติโย วีสมฺปิ ชาติโย ติํสมฺปิ ชาติโย จตฺตารีสมฺปิ ชาติโย ปญฺญาสมฺปิ ชาติโย ชาติสตมฺปิ ชาติสหสฺสมฺปิ ชาติ\-สต\-สหสฺสมฺปิ อเนกานิปิ ชาติสตานิ อเนกานิปิ ชาติสหสฺสานิ อเนกานิปิ ชาติ\-สต\-สหสฺสานิ อเนเกปิ สํวฏฺฏกปฺเป อเนเกปิ วิวฏฺฏกปฺเป อเนเกปิ สํวฏฺฏ\-วิวฏฺฏ\-กปฺเป. อมุตฺราสิํ เอวํนาโม เอวํโคตฺโต เอวํวณฺโณ เอวมาหาโร เอวํ\-สุขทุกฺข\-ปฺปฏิสํเวที เอวมา\-ยุปริยนฺโต, โส ตโต จุโต อมุตฺร อุทปาทิํ. ตตฺราปาสิํ เอวํนาโม เอวํโคตฺโต เอวํวณฺโณ เอวมาหาโร เอวํ\-สุขทุกฺข\-ปฺปฏิสํเวที เอวมา\-ยุปริยนฺโต, โส ตโต จุโต อิธูปปนฺโนติ, อิติ สาการํ สอุทฺเทสํ อเนกวิหิตํ ปุพฺเพนิวาสํ อนุสฺสรติ.}

\proseitem{64}{ตตฺถ สคฺคูปเคสุ จ สตฺเตสุ มนุสฺสูปเคสุ จ สตฺเตสุ อปายูปเคสุ จ สตฺเตสุ อิมสฺส ปุคฺคลสฺส โลภาทโย อุสฺสนฺนา อโลภาทโย มนฺทา, อิมสฺส ปุคฺคลสฺส อโลภาทโย อุสฺสนฺนา โลภาทโย มนฺทา, เย วา ปน อุสฺสนฺนา เย วา ปน มนฺทา อิมสฺส ปุคฺคลสฺส อิมานิ อินฺทฺริยานิ อุปจิตานิ อิมสฺส ปุคฺคลสฺส อิมานิ อินฺทฺริยานิ อนุปจิตานิ อมุกาย วา กปฺปโกฏิยํ กปฺป\-สตสหสฺเส วา กปฺปสหสฺเส วา กปฺปสเต วา กปฺเป วา อนฺตรกปฺเป วา อุปฑฺฒกปฺเป วา สํวจฺฉเร วา อุปฑฺฒ\-สํวจฺฉเร วา มาเส วา ปกฺเข วา ทิวเส วา มุหุตฺเต วา อิมินา ปมาเทน วา ปสาเทน วาติ. ตํ ตํ ภวํ ภควา อนุสฺสรนฺโต อเสสํ ชานาติ, ตตฺถ ยํ ทิพฺเพน จกฺขุนา วิสุทฺเธน อติกฺกนฺต\-มานุสเกน สตฺเต ปสฺสติ จวมาเน อุปปชฺชมาเน หีเน ปณีเต สุวณฺเณ ทุพฺพณฺเณ สุคเต ทุคฺคเต ยถากมฺมูปเค สตฺเต ปชานาติ อิเม วต โภนฺโต สตฺตา กาย\-ทุจฺจริเตน สมนฺนาคตา วจีทุจฺจริเตน สมนฺนาคตา มโน\-ทุจฺจริเตน สมนฺนาคตา อริยานํ อุปวาทกา \csromanfolio{85}มิจฺฉาทิฏฺฐิกา มิจฺฉาทิฏฺฐิ\-กมฺม\-สมาทานา, เต กายสฺส เภทา ปรํ มรณา อปายํ ทุคฺคติํ วินิปาตํ นิรยํ อุปปนฺนา. อิเม วา ปน โภนฺโต สตฺตา กายสุจริเตน สมนฺนาคตา วจีสุจริเตน สมนฺนาคตา.}

\prose{มโนสุจริเตน สมนฺนาคตา อริยานํ อนุปวาทกา สมฺมาทิฏฺฐิกา สมฺมา\-ทิฏฺฐิ\-กมฺม\-สมาทานา, เต กายสฺส เภทา ปรํ มรณา สุคติํ สคฺคํ โลกํ อุปปนฺนา, ตตฺถ สคฺคูปเคสุ จ สตฺเตสุ มนุสฺสูปเคสุ จ สตฺเตสุ อปายูปเคสุ จ สตฺเตสุ อิมินา ปุคฺคเลน เอวรูปํ กมฺมํ อมุกาย กปฺปโกฏิยํ อุปจิตํ กปฺป\-สตสหสฺเส วา กปฺปสหสฺเส วา กปฺปสเต วา กปฺเป วา อนฺตรกปฺเป วา อุปฑฺฒกปฺเป วา สํวจฺฉเร วา อุปฑฺฒ\-สํวจฺฉเร วา มาเส วา ปกฺเข วา ทิวเส วา มุหุตฺเต วา อิมินา ปมาเทน วา ปสาเทน วาติ. อิมานิ ภควโต เทฺว ญาณานิ ปุพฺเพ\-นิวาสา\-นุสฺสติ\-ญาณญฺจ ทิพฺพจกฺขุ จ อฏฺฐมํ นวมํ ตถาคตพลํ. \mbox{(8-9)}}

\prose{อิติ ตตฺถ ยํ สพฺพญฺญุตา ปตฺตา วิทิตา สพฺพธมฺมา วิรชํ วีตมลํ อุปฺปนฺนํ สพฺพญฺญุต\-ญาณํ นิหโต มาโร โพธิมูเล, อิทํ ภควโต ทสมํ พลํ สพฺพาสว\-ปริกฺขยํ ญาณํ. ทสพล\-สมนฺนาคตา หิ พุทฺธา ภควนฺโตติ. (10)}

\vspace{\tipitakaparskip}
\csromancenter{นิยุตฺโต วิจโย หารสมฺปาโต.}
\csromanmatikaanchor{mtk.23}
\csromantocmarkref{section}{3. ยุตฺติหารสมฺปาต}{mtk.23}
\bookmark[dest={mtk.23},level=1]{3. ยุตฺติหารสมฺปาต}
\csromanheader{1}{3. \textbf{ยุตฺติ}\-\textbf{หาร}\-\textbf{สมฺปาต}}
\proseitem{65}{ตตฺถ กตโม ยุตฺติ\-หาร\-สมฺปาโต.}

\prose{\csromansymbolfootnote{*}{เหฏฺฐา 41 ปิฏฺเฐ; อุปริ 90 ปิฏฺเฐปิ.}“ตสฺมา รกฺขิต\-จิตฺตสฺส, สมฺมาสงฺกปฺป\-โคจโร.}

\setlength{\csromangathapagewidth}{0pt}
\setlength{\csromangathawakwidth}{0pt}
\csromangathameasuregroup{สมฺมาทิฏฺฐิปุเรกฺขาโร, \\ ถินมิทฺธาภิภู ภิกฺขุ,}{\csromangathabat{สมฺมาทิฏฺฐิปุเรกฺขาโร,}{ญตฺวาน อุทยพฺพยํ.} \\ \csromangathabat{ถินมิทฺธาภิภู ภิกฺขุ,}{สพฺพา ทุคฺคติโย ชเห”ติ.}}
\csromangathasetleft{สมฺมาทิฏฺฐิปุเรกฺขาโร, \\ ถินมิทฺธาภิภู ภิกฺขุ,}
\csromangathagroup{\csromangathastanza{\csromangathabat{สมฺมาทิฏฺฐิ\-ปุเรกฺขาโร,}{ญตฺวาน อุทยพฺพยํ.} \\ \csromangathabat{ถินมิทฺธาภิภู ภิกฺขุ,}{สพฺพา ทุคฺคติโย ชเห”ติ.}}}

\prose{“ตสฺมา รกฺขิต\-จิตฺตสฺส, สมฺมาสงฺกปฺป\-โคจโร”ติ รกฺขิต\-จิตฺตสฺส สมฺมาสงฺกปฺป\-โคจโร ภวิสฺสตีติ ยุชฺชติ, สมฺมาสงฺกปฺป\-โคจโร สมฺมาทิฏฺฐิ ภวิสฺสตีติ ยุชฺชติ, สมฺมาทิฏฺฐิ\-ปุเรกฺขาโร วิหรนฺโต อุทยพฺพยํ ปฏิวิชฺฌิสฺสตีติ ยุชฺชติ, อุทยพฺพยํ ปฏิวิชฺฌนฺโต สพฺพา ทุคฺคติโย ชหิสฺสตีติ ยุชฺชติ. สพฺพา ทุคฺคติโย ชหนฺโต สพฺพานิ ทุคฺคติ\-วินิปาต\-ภยานิ สมติกฺกมิสฺสตีติ ยุชฺชตีติ.}

\vspace{\tipitakaparskip}
\csromancenter{นิยุตฺโต ยุตฺติ\-หาร\-สมฺปาโต.}
\csromanmatikaanchor{mtk.24}
\csromantocmarkref{section}{4. ปทฏฺฐานหารสมฺปาต}{mtk.24}
\bookmark[dest={mtk.24},level=1]{4. ปทฏฺฐานหารสมฺปาต}
\csromanheader{1}{4. \textbf{ปทฏฺฐาน}\-\textbf{หาร}\-\textbf{สมฺปาต}}
\proseitem{66}{\csromanfolio{86}ตตฺถ กตโม ปทฏฺฐาโน หารสมฺปาโต.}

\prose{“ตสฺมา รกฺขิต\-จิตฺตสฺส, สมฺมาสงฺกปฺป\-โคจโร”ติ คาถา. “ตสฺมา รกฺขิต\-จิตฺตสฺสา”ติ ติณฺณํ สุจริตานํ ปทฏฺฐานํ. “สมฺมาสงฺกปฺป\-โคจโร”ติ สมถสฺส ปทฏฺฐานํ. “สมฺมาทิฏฺฐิ\-ปุเรกฺขาโร”ติ วิปสฺสนาย ปทฏฺฐานํ. “ญตฺวาน อุทยพฺพยนฺ”ติ ทสฺสนภูมิยา ปทฏฺฐานํ. “ถินมิทฺธาภิภู ภิกฺขู”ติ วีริยสฺส ปทฏฺฐานํ. “สพฺพา ทุคฺคติโย ชเห”ติ ภาวนาย ปทฏฺฐานํ. นิยุตฺโต ปทฏฺฐาโน หารสมฺปาโต.}

\csromanmatikaanchor{mtk.25}
\csromantocmarkref{section}{5. ลกฺขณหารสมฺปาต}{mtk.25}
\bookmark[dest={mtk.25},level=1]{5. ลกฺขณหารสมฺปาต}
\csromanheader{1}{5. \textbf{ลกฺขณ}\-\textbf{หาร}\-\textbf{สมฺปาต}}
\proseitem{67}{ตตฺถ กตโม ลกฺขโณ หารสมฺปาโต.}

\prose{“ตสฺมา รกฺขิต\-จิตฺตสฺส, สมฺมาสงฺกปฺป\-โคจโร”ติ คาถา. “ตสฺมา รกฺขิต\-จิตฺตสฺส, สมฺมาสงฺกปฺป\-โคจโร”ติ อิทํ สตินฺทฺริยํ, สตินฺทฺริเย คหิเต คหิตานิ ภวนฺติ ปญฺจินฺทฺริยานิ. “สมฺมาทิฏฺฐิ\-ปุเรกฺขาโร”ติ สมฺมาทิฏฺฐิยา คหิตาย คหิโต ภวติ อริโย อฏฺฐงฺคิโก มคฺโค. ตํ กิสฺส เหตุ, สมฺมาทิฏฺฐิโต หิ สมฺมาสงฺกปฺโป ปภวติ, สมฺมา\-สงฺกปฺปโต สมฺมาวาจา ปภวติ, สมฺมาวาจาโต สมฺมากมฺมนฺโต ปภวติ, สมฺมา\-กมฺมนฺตโต สมฺมาอาชีโว ปภวติ, สมฺมาอาชีวโต สมฺมาวายาโม ปภวติ, สมฺมาวายามโต สมฺมาสติ ปภวติ, สมฺมาสติโต สมฺมาสมาธิ ปภวติ, สมฺมาสมาธิโต สมฺมาวิมุตฺติ ปภวติ, สมฺมาวิมุตฺติโต สมฺมา\-วิมุตฺติ\-ญาณ\-ทสฺสนํ ปภวติ. นิยุตฺโต ลกฺขโณ หารสมฺปาโต.}

\csromanmatikaanchor{mtk.26}
\csromantocmarkref{section}{6. จตุพฺยูหหารสมฺปาต}{mtk.26}
\bookmark[dest={mtk.26},level=1]{6. จตุพฺยูหหารสมฺปาต}
\csromanheader{1}{6. \textbf{จตุพฺยูห}\-\textbf{หาร}\-\textbf{สมฺปาต}}
\proseitem{68}{ตตฺถ กตโม จตุพฺยูโห หารสมฺปาโต.}

\prose{“ตสฺมา รกฺขิต\-จิตฺตสฺส, สมฺมาสงฺกปฺป\-โคจโร”ติ คาถา. “ตสฺมา รกฺขิต\-จิตฺตสฺสา”ติ รกฺขิตํ ปริปาลียตีติ เอสา นิรุตฺติ. อิธ ภควโต โก อธิปฺปาโย, เย ทุคฺคตีติ ปริมุจฺจิตุกามา ภวิสฺสนฺติ, เต ธมฺมจาริโน ภวิสฺสนฺตีติ อยํ เอตฺถ ภควโต อธิปฺปาโย, โกกาลิโก หิ \csromanfolio{87}สาริปุตฺต\-โมคฺคลฺลาเนสุ เถเรสุ จิตฺตํ ปโทสยิตฺวา มหา\-ปทุมนิรเย อุปปนฺโน. ภควา จ สติอารกฺเขน เจตสา สมนฺนาคโต, สุตฺตมฺหิ วุตฺตํ “สติยา จิตฺตํ รกฺขิตพฺพนฺ”ติ. นิยุตฺโต จตุพฺยูโห หารสมฺปาโต.}

\csromanmatikaanchor{mtk.27}
\csromantocmarkref{section}{7. อาวฏฺฏหารสมฺปาต}{mtk.27}
\bookmark[dest={mtk.27},level=1]{7. อาวฏฺฏหารสมฺปาต}
\csromanheader{1}{7. \textbf{อาวฏฺฏ}\-\textbf{หาร}\-\textbf{สมฺปาต}}
\proseitem{69}{ตตฺถ กตโม อาวฏฺโฏ หารสมฺปาโต.}

\prose{“ตสฺมา รกฺขิต\-จิตฺตสฺส, สมฺมาสงฺกปฺป\-โคจโร”ติ คาถา. “ตสฺมา รกฺขิต\-จิตฺตสฺส, สมฺมาสงฺกปฺป\-โคจโร”ติ สมโถ\csromansharedfootnote{37ec514ed28d82fd}{อยํ สมมถา (สี, ก)}. “สมฺมาทิฏฺฐิ\-ปุเรกฺขาโร”ติ วิปสฺสนา. “ญตฺวาน อุทยพฺพยนฺ”ติ ทุกฺขปริญฺญา. “ถินมิทฺธาภิภู ภิกฺขู”ติ สมุทย\-ปหานํ. “สพฺพา ทุคฺคติโย ชเห”ติ นิโรโธ\csromansharedfootnote{ac3a3d3f7c2650f4}{อยํ นิโรโธ (สี, ก)}. อิมานิ จตฺตาริ สจฺจานิ. นิยุตฺโต อาวฏฺโฏ หารสมฺปาโต.}

\csromanmatikaanchor{mtk.28}
\csromantocmarkref{section}{8. วิภตฺติหารสมฺปาต}{mtk.28}
\bookmark[dest={mtk.28},level=1]{8. วิภตฺติหารสมฺปาต}
\csromanheader{1}{8. \textbf{วิภตฺติ}\-\textbf{หาร}\-\textbf{สมฺปาต}}
\proseitem{70}{ตตฺถ กตโม วิภตฺติ\-หาร\-สมฺปาโต.}

\prose{“ตสฺมา รกฺขิต\-จิตฺตสฺส, สมฺมาสงฺกปฺป\-โคจโร”ติ คาถา. กุสลปกฺโข กุสลปกฺเขน นิทฺทิสิตพฺโพ. อกุสลปกฺโข อกุสล\-ปกฺเขน นิทฺทิสิตพฺโพ. นิยุตฺโต วิภตฺติ\-หาร\-สมฺปาโต.}

\csromanmatikaanchor{mtk.29}
\csromantocmarkref{section}{9. ปริวตฺตนหารสมฺปาต}{mtk.29}
\bookmark[dest={mtk.29},level=1]{9. ปริวตฺตนหารสมฺปาต}
\csromanheader{1}{9. \textbf{ปริวตฺตน}\-\textbf{หาร}\-\textbf{สมฺปาต}}
\proseitem{71}{ตตฺถ กตฺตโม ปริวตฺตโน หารสมฺปาโต.}

\prose{“ตสฺมา รกฺขิต\-จิตฺตสฺส, สมฺมาสงฺกปฺป\-โคจโร”ติ คาถา. สมถ\-วิปสฺสนาย ภาวิตาย นิโรโธ ผลํ, ปริญฺญาตํ ทุกฺขํ, สมุทโย ปหีโน, มคฺโค ภาวิโต ปฏิปกฺเขน. นิยุตฺโต ปริวตฺตโน หารสมฺปาโต.}

\csromanmatikaanchor{mtk.30}
\csromantocmarkref{section}{10. เววจนหารสมฺปาต}{mtk.30}
\bookmark[dest={mtk.30},level=1]{10. เววจนหารสมฺปาต}
\csromanheader{1}{10. \textbf{เววจน}\-\textbf{หาร}\-\textbf{สมฺปาต}}
\proseitem{72}{\csromanfolio{88}ตตฺถ กตโม เววจโน หารสมฺปาโต.}

\prose{“ตสฺมา รกฺขิต\-จิตฺตสฺส, สมฺมาสงฺกปฺป\-โคจโร”ติ คาถา. “ตสฺมา รกฺขิต\-จิตฺตสฺสา”ติ จิตฺตํ มโน วิญฺญาณํ มนินฺทฺริยํ มนายตนํ วิชานนา วิชานิตตฺตํ, อิทํ เววจนํ. “สมฺมาสงฺกปฺป\-โคจโร”ติ เนกฺขมฺม\-สงฺกปฺโป อพฺยาปาทสงฺกปฺโป อวิหิํสาสงฺกปฺโป, อิทํ เววจนํ. “สมฺมาทิฏฺฐิ\-ปุเรกฺขาโร”ติ สมฺมาทิฏฺฐิ นาม ปญฺญาสตฺถํ ปญฺญาขคฺโค ปญฺญารตนํ ปญฺญาปชฺโชโต ปญฺญาปโตโท ปญฺญาปาสาโท, อิทํ เววจนํ. นิยุตฺโต เววจโน หารสมฺปาโต.}

\csromanmatikaanchor{mtk.31}
\csromantocmarkref{section}{11. ปญฺญตฺติหารสมฺปาต}{mtk.31}
\bookmark[dest={mtk.31},level=1]{11. ปญฺญตฺติหารสมฺปาต}
\csromanheader{1}{11. \textbf{ปญฺญตฺติ}\-\textbf{หาร}\-\textbf{สมฺปาต}}
\proseitem{73}{ตตฺถ กตโม ปญฺญตฺติ\-หาร\-สมฺปาโต.}

\prose{“ตสฺมา รกฺขิต\-จิตฺตสฺส, สมฺมาสงฺกปฺป\-โคจโร”ติ คาถา. “ตสฺมา รกฺขิต\-จิตฺตสฺสา”ติ ปทฏฺฐาน\-ปญฺญตฺติ สติยา. “สมฺมาสงฺกปฺป\-โคจโร”ติ ภาวนา\-ปญฺญตฺติ สมถสฺส. “สมฺมาทิฏฺฐิ\-ปุเรกฺขาโร, ญตฺวาน อุทยพฺพยนฺ”ติ ทสฺสนภูมิยา นิกฺเขป\-ปญฺญตฺติ. “ถินมิทฺธาภิภู ภิกฺขู”ติ สมุทยสฺส อนวเสส\-ปฺปหาน\-ปญฺญตฺติ, “สพฺพา ทุคฺคติโย ชเห”ติ ภาวนา\-ปญฺญตฺติ มคฺคสฺส. นิยุตฺโต ปญฺญตฺติ\-หาร\-สมฺปาโต.}

\csromanmatikaanchor{mtk.32}
\csromantocmarkref{section}{12. โอตรณหารสมฺปาต}{mtk.32}
\bookmark[dest={mtk.32},level=1]{12. โอตรณหารสมฺปาต}
\csromanheader{1}{12. \textbf{โอตรณ}\-\textbf{หาร}\-\textbf{สมฺปาต}}
\proseitem{74}{ตตฺถ กตโม โอตรโณ หารสมฺปาโต.}

\prose{“ตสฺมา รกฺขิต\-จิตฺตสฺส, สมฺมาสงฺกปฺป\-โคจโร”ติ คาถา. “ตสฺมา รกฺขิต\-จิตฺตสฺส, สมฺมาสงฺกปฺป\-โคจโร. สมฺมาทิฏฺฐิ\-ปุเรกฺขาโร”ติ สมฺมาทิฏฺฐิยา คหิตาย คหิตานิ ภวนฺติ ปญฺจินฺทฺริยานิ, อยํ อินฺทฺริเยหิ โอตรณา.}

\prose{ตานิเยว อินฺทฺริยานิ วิชฺชา, วิชฺชุปฺปาทา อวิชฺชานิโรโธ, อวิชฺชานิโรธา สงฺขาร\-นิโรโธ, สงฺขาร\-นิโรธา วิญฺญาณนิโรโธ, เอวํ สพฺพํ, อยํ ปฏิจฺจ\-สมุปฺปาเทน โอตรณา.}

\prose{ตานิเยว ปญฺจินฺทฺริยานิ ตีหิ ขนฺเธหิ สงฺคหิตานิ สีลกฺขนฺเธน สมาธิ\-กฺขนฺเธน ปญฺญา\-กฺขนฺเธน, อยํ ขนฺเธหิ โอตรณา.}

\prose{\csromanfolio{89}ตานิ เยว ปญฺจินฺทฺริยานิ สงฺขาร\-ปริยาปนฺนานิ. เย สงฺขารา อนาสวา โน จ ภวงฺคา, เต สงฺขารา ธมฺมธาตุ\-สงฺคหิตา, อยํ ธาตูหิ โอตรณา.}

\prose{สา ธมฺมธาตุ ธมฺมายตน\-ปริยาปนฺนา, ยํ อายตนํ อนาสวํ โน จ ภวงฺคํ, อยํ อายตเนหิ โอตรณา. นิยุตฺโต โอตรโณ หารสมฺปาโต.}

\csromanmatikaanchor{mtk.33}
\csromantocmarkref{section}{13. โสธนหารสมฺปาต}{mtk.33}
\bookmark[dest={mtk.33},level=1]{13. โสธนหารสมฺปาต}
\csromanheader{1}{13. \textbf{โสธน}\-\textbf{หาร}\-\textbf{สมฺปาต}}
\proseitem{75}{ตตฺถ กตโม โสธโน หารสมฺปาโต.}

\prose{“ตสฺมา รกฺขิต\-จิตฺตสฺส, สมฺมาสงฺกปฺป\-โคจโร”ติ คาถา. ยตฺถ อารมฺโภ สุทฺโธ, โส ปโญฺห วิสชฺชิโต ภวติ. ยตฺถ ปน อารมฺโภ น สุทฺโธ, น ตาว โส ปโญฺห วิสชฺชิโต ภวติ. นิยุตฺโต โสธโน หารสมฺปาโต.}

\csromanmatikaanchor{mtk.34}
\csromantocmarkref{section}{14. อธิฏฺฐานหารสมฺปาต}{mtk.34}
\bookmark[dest={mtk.34},level=1]{14. อธิฏฺฐานหารสมฺปาต}
\csromanheader{1}{14. \textbf{อธิฏฺฐาน}\-\textbf{หาร}\-\textbf{สมฺปาต}}
\proseitem{76}{ตตฺถ กตโม อธิฏฺฐาโน หารสมฺปาโต.}

\prose{ตสฺมา รกฺขิต\-จิตฺตสฺส, สมฺมาสงฺกปฺป\-โคจโรติ คาถา. ตสฺมา รกฺขิต\-จิตฺตสฺสาติ เอกตฺตตา. จิตฺตํ มโน วิญฺญาณํ, อยํ เวมตฺตตา. สมฺมาสงฺกปฺป\-โคจโรติ เอกตฺตตา. เนกฺขมฺม\-สงฺกปฺโป อพฺยาปาทสงฺกปฺโป อวิหิํสาสงฺกปฺโป อยํ เวมตฺตตา. สมฺมาทิฏฺฐิ\-ปุเรกฺขาโรติ เอกตฺตตา. สมฺมาทิฏฺฐิ นาม ยํ ทุกฺเข ญาณํ ทุกฺขสมุทเย ญาณํ ทุกฺขนิโรเธ ญาณํ ทุกฺขนิโรธ\-คามินิยา ปฏิปทาย ญาณํ มคฺเค ญาณํ เหตุมฺหิ ญาณํ เหตุ\-สมุปฺปนฺเนสุ ธมฺเมสุ ญาณํ ปจฺจเย ญาณํ ปจฺจยสมุปฺปนฺเนสุ ธมฺเมสุ ญาณํ, ยํ ตตฺถ ตตฺถ ยถาภูตํ ญาณทสฺสนํ อภิสมโย สมฺปฏิเวโธ สจฺจาคมนํ, อยํ เวมตฺตตา. ญตฺวาน อุทยพฺพยนฺติ เอกตฺตตา, อุทเยน อวิชฺชาปจฺจยา สงฺขารา, สงฺขาร\-ปจฺจยา วิญฺญาณํ, เอวํ สพฺพํ สมุทโย ภวติ. วเยน อวิชฺชานิโรธา สงฺขาร\-นิโรโธ, เอวํ สพฺพํ นิโรโธ โหติ, อยํ เวมตฺตตา. ถินมิทฺธาภิภู ภิกฺขูติ เอกตฺตตา, ถินํ นาม ยา จิตฺตสฺส อกลฺลตา อกมฺมนิยตา, มิทฺธํ นาม ยํ กายสฺส ลีนตฺตํ, อยํ เวมตฺตตา. สพฺพา ทุคฺคติโย ชเหติ เอกตฺตตา, เทวมนุสฺส วา \csromanfolio{90}อุปนิธาย อปายา ทุคฺคติ, นิพฺพานํ วา อุปนิธาย สพฺพา อุปปตฺติโย ทุคฺคติ, อยํ เวมตฺตตา. นิยุตฺโต อธิฏฺฐาโน หารสมฺปาโต.}

\csromanmatikaanchor{mtk.35}
\csromantocmarkref{section}{15. ปริกฺขารหารสมฺปาต}{mtk.35}
\bookmark[dest={mtk.35},level=1]{15. ปริกฺขารหารสมฺปาต}
\csromanheader{1}{15. \textbf{ปริกฺขาร}\-\textbf{หาร}\-\textbf{สมฺปาต}}
\proseitem{77}{ตตฺถ กตโม ปริกฺขาโร หารสมฺปาโต.}

\prose{ตสฺมา รกฺขิต\-จิตฺตสฺส, สมฺมาสงฺกปฺป\-โคจโรติ คาถา. อยํ สมถ\-วิปสฺสนาย ปริกฺขาโร. นิยุตฺโต ปริกฺขาโร หารสมฺปาโต.}

\csromanmatikaanchor{mtk.36}
\csromantocmarkref{section}{16. สมาโรปนหารสมฺปาต}{mtk.36}
\bookmark[dest={mtk.36},level=1]{16. สมาโรปนหารสมฺปาต}
\csromanheader{1}{16. \textbf{สมาโรปน}\-\textbf{หาร}\-\textbf{สมฺปาต}}
\proseitem{78}{ตตฺถ กตโม สมาโรปโน หารสมฺปาโต.}

\prose{\csromansymbolfootnote{*}{เหฏฺฐา 41, 85 ปิฏฺเฐสุปิ.}“ตสฺมา รกฺขิต\-จิตฺตสฺส, สมฺมาสงฺกปฺป\-โคจโร.}

\setlength{\csromangathapagewidth}{0pt}
\setlength{\csromangathawakwidth}{0pt}
\csromangathameasuregroup{สมฺมาทิฏฺฐิปุเรกฺขาโร, \\ ถินมิทฺธาภิภู ภิกฺขุ,}{\csromangathabat{สมฺมาทิฏฺฐิปุเรกฺขาโร,}{ญตฺวาน อุทยพฺพยํ.} \\ \csromangathabat{ถินมิทฺธาภิภู ภิกฺขุ,}{สพฺพา ทุคฺคติโย ชเห”ติ.}}
\csromangathasetleft{สมฺมาทิฏฺฐิปุเรกฺขาโร, \\ ถินมิทฺธาภิภู ภิกฺขุ,}
\csromangathagroup{\csromangathastanza{\csromangathabat{สมฺมาทิฏฺฐิ\-ปุเรกฺขาโร,}{ญตฺวาน อุทยพฺพยํ.} \\ \csromangathabat{ถินมิทฺธาภิภู ภิกฺขุ,}{สพฺพา ทุคฺคติโย ชเห”ติ.}}}

\prose{ตสฺมา รกฺขิตจิตฺตสฺสาติ ติณฺณํ สุจริตานํ ปทฏฺฐานํ, จิตฺเต รกฺขิเต ตํ รกฺขิตํ ภวติ กายกมฺมํ วจีกมฺมํ มโนกมฺมํ. สมฺมาทิฏฺฐิ\-ปุเรกฺขาโรติ สมฺมาทิฏฺฐิยา ภาวิตาย ภาวิโต ภวติ อริโย อฏฺฐงฺคิโก มคฺโค. เกน การเณน, สมฺมาทิฏฺฐิโต หิ สมฺมาสงฺกปฺโป ปภวติ, สมฺมา\-สงฺกปฺปโต สมฺมาวาจา ปภวติ, สมฺมาวาจาโต สมฺมากมฺมนฺโต ปภวติ, สมฺมา\-กมฺมนฺตโต สมฺมาอาชีโว ปภวติ, สมฺมาอาชีวโต สมฺมาวายาโม ปภวติ, สมฺมาวายามโต สมฺมาสติ ปภวติ, สมฺมาสติโต สมฺมาสมาธิ ปภวติ, สมฺมาสมาธิโต สมฺมาวิมุตฺติ ปภวติ, สมฺมาวิมุตฺติโต สมฺมา\-วิมุตฺติ\-ญาณ\-ทสฺสนํ ปภวติ. อยํ อนุปาทิเสโส ปุคฺคโล อนุปาทิเสสา จ นิพฺพานธาตุ.}

\vspace{\tipitakaparskip}
\csromancenter{นิยุตฺโต สมาโรปโน หารสมฺปาโต.}
\prose{\csromanfolio{91}เตนาห อายสฺมา มหากจฺจายโน\csromandash{}}

\setlength{\csromangathapagewidth}{0pt}
\setlength{\csromangathawakwidth}{0pt}
\csromangathameasuregroup{“โสฬส หารา ปฐมํ, \\ สงฺขิปิย องฺกุเสน หิ,}{\csromangathabat{“โสฬส หารา ปฐมํ,}{ทิสโลจนโต ทิสา วิโลเกตฺวา.} \\ \csromangathabat{สงฺขิปิย องฺกุเสน หิ,}{นเยหิ ตีหิ นิทฺทิเส สุตฺตนฺ”ติ.} \\ นิยุตฺโต หารสมฺปาโต.}
\csromangathameasurewak{นิยุตฺโต หารสมฺปาโต.}
\csromangathasetleft{“โสฬส หารา ปฐมํ, \\ สงฺขิปิย องฺกุเสน หิ,}
\csromangathagroup{\csromangathastanza{\csromangathabat{“โสฬส หารา ปฐมํ,}{ทิสโลจนโต ทิสา วิโลเกตฺวา.} \\ \csromangathabat{สงฺขิปิย องฺกุเสน หิ,}{นเยหิ ตีหิ นิทฺทิเส สุตฺตนฺ”ติ.}}\csromangathastanza{\csromangathawak{นิยุตฺโต หารสมฺปาโต.}}}

\csromanmatikaanchor{mtk.37}
\csromantocmarkref{subsection}{นยสมุฏฺฐาน}{mtk.37}
\bookmark[dest={mtk.37},level=2]{นยสมุฏฺฐาน}
\csromanheader{2}{\textbf{นยสมุฏฺฐาน}}
\proseitem{79}{ตตฺถ กตมํ นยสมุฏฺฐานํ. ปุพฺพา โกฏิ น ปญฺญายติ อวิชฺชาย จ ภวตณฺหาย จ, ตตฺถ อวิชฺชา\-นีวรณํ ตณฺหา\-สํโยชนํ. อวิชฺชานีวรณา สตฺตา อวิชฺชาสํยุตฺตา\csromansharedfootnote{a9fad056866a9993}{อวิชฺชาย สํยุตฺตา (สี, ก)} อวิชฺชาปกฺเขน วิจรนฺติ, เต วุจฺจนฺติ ทิฏฺฐิจริตาติ. ตณฺหาสํโยชนา สตฺตา ตณฺหาสํยุตฺตา ตณฺหาปกฺเขน วิจรนฺติ, เต วุจฺจนฺติ ตณฺหาจริตาติ. ทิฏฺฐิจริตา อิโต พหิทฺธา ปพฺพชิตา อตฺตกิลมถา\-นุโยคมนุยุตฺตา วิหรนฺติ. ตณฺหาจริตา อิโต พหิทฺธา ปพฺพชิตา กาเมสุ กามสุขลฺลิกา\-นุโยคมนุยุตฺตา วิหรนฺติ.}

\prose{ตตฺถ กิํการณํ ยํ ทิฏฺฐิจริตา อิโต พหิทฺธา ปพฺพชิตา อตฺตกิลมถา\-นุโยคมนุยุตฺตา วิหรนฺติ. ตณฺหาจริตา อิโต พหิทฺธา ปพฺพชิตา กาเมสุ กามสุขลฺลิกา\-นุโยคมนุยุตฺตา วิหรนฺติ. อิโต พหิทฺธา นตฺถิ สจฺจ\-ววตฺถานํ, กุโต จตุสจฺจ\-ปฺปกาสนา วา สมถ\-วิปสฺสนา\-โกสลฺลํ วา อุปสมสุข\-ปฺปตฺติ วา. เต อุปสม\-สุขสฺส อนภิญฺญา วิปรีตเจตา เอวมาหํสุ “นตฺถิ สุเขน สุขํ, ทุกฺเขน นาม สุขํ อธิคนฺตพฺพนฺ”ติ. โย กาเม ปฏิเสวติ, โส โลกํ วฑฺฒยติ, โย โลกํ วฑฺฒยติ, โส พหุํ ปุญฺญํ ปสวตีติ เต เอวํสญฺญี เอวํทิฏฺฐี ทุกฺเขน สุขํ ปตฺถยมานา กาเมสุ ปุญฺญสญฺญี อตฺตกิลมถา\-นุโยคมนุยุตฺตา จ วิหรนฺติ กามสุขลฺลิกา\-นุโยคมนุยุตฺตา จ, เต ตทภิญฺญา สนฺตา โรคเมว วฑฺฒยนฺติ, คณฺฑเมว วฑฺฒยนฺติ, สลฺลเมว วฑฺฒยนฺติ, เต โรคาภิตุนฺนา คณฺฑ\-ปฏิปีฬิตา สลฺลานุวิทฺธา นิรย\-ติรจฺฉานโยนิ\-เปตา\-สุเรสุ อุมฺมุชฺชนิมุชฺชานิ กโรนฺตา อุคฺฆาต\-นิคฺฆาตํ ปจฺจนุโภนฺตา โรคคณฺฑสลฺลเภสชฺฌํ น วินฺทนฺติ. ตตฺถ อตฺต\-กิลมถา\-นุโยโค กามสุขลฺลิกา\-นุโยโค จ สํกิเลโส, สมถ\-วิปสฺสนา โวทานํ. อตฺต\-กิลมถา\-นุโยโค กามสุขลฺลิกา\-นุโยโค จ โรโค, สมถ\-วิปสฺสนา โรค\-นิคฺฆาตก\-เภสชฺชํ. อตฺต\-กิลมถา\-นุโยโค กามสุขลฺลิกา\-นุโยโค จ คณฺโฑ, สมถ\-วิปสฺสนา \csromanfolio{92}คณฺฑ\-นิคฺฆาตก\-เภสชฺชํ. อตฺต\-กิลมถา\-นุโยโค กามสุขลฺลิกา\-นุโยโค จ สลฺโล, สมถ\-วิปสฺสนา สลฺลุ\-ทฺธรณ\-เภสชฺชํ.}

\prose{ตตฺถ สํกิเลโส ทุกฺขํ, ตทภิสงฺโค ตณฺหา สมุทโย, ตณฺหานิโรโธ ทุกฺขนิโรโธ, สมถ\-วิปสฺสนา ทุกฺขนิโรธ\-คามินี ปฏิปทา, อิมานิ จตฺตาริ จจฺจานิ. ทุกฺขํ ปริญฺเญยฺยํ, สมุทโย ปหาตพฺโพ, มคฺโค ภาเวตพฺโพ, นิโรโธ สจฺฉิกาตพฺโพ.}

\proseitem{80}{ตตฺถ ทิฏฺฐิจริตา รูปํ อตฺตโต อุปคจฺฉนฺติ. เวทนํ ฯเปฯ สญฺญํ. สงฺขาเร. วิญฺญาณํ อตฺตโต อุปคจฺฉนฺติ. ตณฺหาจริตา รูปวนฺตํ อตฺตานํ อุปคจฺฉนฺติ. อตฺตนิ วา รูปํ, รูปสฺมิํ วา อตฺตานํ, เวทนาวนฺตํ ฯเปฯ สญฺญาวนฺตํ. สงฺขารวนฺตํ. วิญฺญาณวนฺตํ อตฺตานํ อุปคจฺฉนฺติ, อตฺตนิ วา วิญฺญาณํ, วิญฺญาณสฺมิํ วา อตฺตานํ, อยํ วุจฺจติ วีสติวตฺถุกา สกฺกายทิฏฺฐิ.}

\prose{ตสฺสา ปฏิปกฺโข โลกุตฺตรา สมฺมาทิฏฺฐิ, อนฺวายิกา สมฺมาสงฺกปฺโป สมฺมาวาจา สมฺมากมฺมนฺโต สมฺมาอาชีโว สมฺมาวายาโม สมฺมาสติ สมฺมาสมาธิ, อยํ อริโย อฏฺฐงฺคิโก มคฺโค, เต ตโย ขนฺธา สีลกฺขนฺโธ สมาธิ\-กฺขนฺโธ ปญฺญากฺขนฺโธ. สีลกฺขนฺโธ สมาธิ\-กฺขนฺโธ จ สมโถ, ปญฺญากฺขนฺโธ วิปสฺสนา. ตตฺถ สกฺกาโย ทุกฺขํ, สกฺกาย\-สมุทโย ทุกฺขสมุทโย, สกฺกายนิโรโธ ทุกฺขนิโรโธ, อริโยอฏฺฐงฺคิโก มคฺโค ทุกฺขนิโรธ\-คามินี ปฏิปทา, อิมานิ จตฺตาริ สจฺจานิ. ทุกฺขํ ปริญฺเญยฺยํ, สมุทโย ปหาตพฺโพ, มคฺโค ภาเวตพฺโพ นิโรโธ สจฺฉิกาตพฺโพ.}

\prose{ตตฺถ เย รูปํ อตฺตโต อุปคจฺฉนฺติ. เวทนํ ฯเปฯ สญฺญํ. สงฺขาเร. วิญฺญาณํ อตฺตโต อุปคจฺฉนฺติ. อิเม วุจฺจนฺติ “อุจฺเฉทวาทิโน”ติ. เย รูปวนฺตํ อตฺตานํ อุปคจฺฉนฺติ. อตฺตนิ วา รูปํ, รูปสฺมิํ วา อตฺตานํ. เย เวทนาวนฺตํ ฯเปฯ. เย สญฺญาวนฺตํ, เย สงฺขารวนฺตํ, เย วิญฺญาณวนฺตํ อตฺตานํ อุปคจฺฉนฺติ, อตฺตนิ วา วิญฺญาณํ, วิญฺญาณสฺมิํ วา อตฺตานํ. อิเม วุจฺจนฺติ “สสฺสตวาทิโน”ติ, ตตฺถ อุจฺเฉท\-สสฺสตวาทา อุโภ อนฺตา, อยํ สํสารปวตฺติ. ตสฺส ปฏิปกฺโข มชฺฌิมา ปฏิปทา อริโย อฏฺฐงฺคิโก มคฺโค, อยํ สํสารนิวตฺติ. ตตฺถ ปวตฺติ ทุกฺขํ, ตทภิสงฺโค ตณฺหา สมุทโย, ตณฺหานิโรโธ ทุกฺขนิโรโธ, อริโย อฏฺฐงฺคิโก มคฺโค \csromanfolio{93}ทุกฺขนิโรธ\-คามินี ปฏิปทา, อิมานิ จิตฺตาริ สจฺจานิ. ทุกฺขํ ปริญฺเญยฺยํ, สมุทโย ปหาตพฺโพ, มคฺโค ภาเวตพฺโพ, นิโรโธ สจฺฉิกาตพฺโพ.}

\prose{ตตฺถ อุจฺเฉท\-สสฺสตํ สมาสโต วีสติวุตฺถุกา สกฺกายทิฏฺฐิ, วิตฺถารโต ทฺวาสฏฺฐิ ทิฏฺฐิคตานิ, เตสํ ปฏิปกฺโข เตจตฺตาลีสํ โพธิปกฺขิยา ธมฺมา อฏฺฐ วิโมกฺขา ทส กสิณายตนานิ. ทฺวาสฏฺฐิ ทิฏฺฐิคตานิ โมหชาลํ อนาทิอนิธนปฺปวตฺตํ. เตจตฺตาลีสํ\csromansharedfootnote{3265c5dda93a9f8c}{เตตาลีสํ (สี)} โพธิปกฺขิยา ธมฺมา ญาณวชิรํ โมหชาล\-ปฺปทาลนํ. ตตฺถ โมโห อวิชฺชา, ชาลํ ภวตณฺหา, เตน วุจฺจติ “ปุพฺพา โกฏิ น ปญฺญายติ อวิชฺชาย จ ภวตณฺหาย จา”ติ.}

\proseitem{81}{ตตฺถ ทิฏฺฐิจริโต อสฺมิํ สาสเน ปพฺพชิโต สลฺเลขา\-นุสนฺตตวุตฺติ ภวติ สลฺเลเข ติพฺพคารโว. ตณฺหาจริโต อสฺมิํ สาสเน ปพฺพชิโต สิกฺขา\-นุสนฺตตวุตฺติ ภวติ สิกฺขาย ติพฺพคารโว. ทิฏฺฐิจริโต สมฺมตฺต\-นิยามํ โอกฺกมนฺโต ธมฺมานุสารี ภวติ. ตณฺหาจริโต สมฺมตฺต\-นิยามํ โอกฺกมนฺโต สทฺธานุสารี ภวติ, ทิฏฺฐิจริโต สุขาย ปฏิปทาย ทนฺธาภิญฺญาย ขิปฺปาภิญฺญาย จ นิยฺยาติ. ตณฺหาจริโต ทุกฺขาย ปฏิปทาย ทนฺธาภิญฺญาย ขิปฺปาภิญฺญาย จ นิยฺยาติ.}

\prose{ตตฺถ กิํ การณํ, ยํ ตณฺหาจริโต ทุกฺขาย ปฏิปทาย ทนฺธาภิญฺญาย ขิปฺปาภิญฺญาย จ นิยฺยาติ, ตสฺส หิ กามา อปริจฺจตฺตา ภวนฺติ, โส กาเมหิ วิเวจิยมาโน ทุกฺเขน ปฏินิสฺสรติ ทนฺธญฺจ ธมฺมํ อาชานาติ. โย ปนายํ ทิฏฺฐิจริโต อยํ อาทิโตเยว กาเมหิ อนตฺถิโก ภวติ. โส ตโต วิเวจิยมาโน ขิปฺปญฺจ ปฏินิสฺสรติ, ขิปฺปญฺจ ธมฺมํ อาชานาติ. ทุกฺขาปิ ปฏิปทา ทุวิธา ทนฺธาภิญฺญา จ ขิปฺปาภิญฺญา จ. สุขาปิ ปฏิปทา ทุวิธา ทนฺธาภิญฺญา จ ขิปฺปาภิญฺญา จ. สตฺตาปิ ทุวิธา มุทินฺทฺริยาปิ ติกฺขินฺทฺริยาปิ. เย มุทินฺทฺริยา, เต ทนฺธญฺจ ปฏินิสฺสรนฺติ ทนฺธญฺจ ธมฺมํ อาชานนฺติ. เย ติกฺขินฺทฺริยา, เต ขิปฺปญฺจ ปฏินิสฺสรนฺติ, ขิปฺปญฺจ ธมฺมํ อาชานนฺติ, อิมา จตสฺโส ปฏิปทา. เย หิ เกจิ นิยฺยิํสุ วา นิยฺยนฺติ วา นิยฺยิสฺสนฺติ วา, เต อิมาหิ เอว จตูหิ ปฏิปทาหิ. เอวํ อริยา จตุกฺกมคฺคํ ปญฺญาเปนฺติ อพุธ\-ชน\-เสวิตาย พาลกนฺตาย รตฺตวาสินิยา นนฺทิยา ภวตณฺหาย อวฏฺฏนตฺถํ\csromansharedfootnote{38dae69e8017d9d4}{อาวฏฺฏนตฺถํ (สี, ก)}. อยํ วุจฺจติ นนฺทิยาวฏฺฏสฺสนฺ นยสฺส ภูมีติ, เตนาห “ตณฺหญฺจ อวิชฺชมฺปิ จ สมเถนา”ติ.}

\proseitem{82}{\csromanfolio{94}เวยฺยากรเณสุ หิ เย กุสลากุสลาติ เต ทุวิธา อุปปริ\-กฺขิ\-ตพฺพา โลก\-วฏฺฏานุสารี จ โลก\-วิวฏฺฏา\-นุ\-สารี จ. วฏฺฏํ นาม สํสาโร. วิวฏฺฏํ นิพฺพานํ. กมฺมกิเลสา เหตุ สํสารสฺส. ตตฺถ กมฺมํ เจตนา เจตสิกญฺจ นิทฺทิสิตพฺพํ. ตํ กถํ ทฏฺฐพฺพํ, อุปจเยน สพฺเพปิ กิเลสา จตูหิ วิปลฺลาเสหิ นิทฺทิสิตพฺพา, เต กตฺถ ทฏฺฐพฺพา, ทส วตฺถุเก กิเลสปุญฺเช. กตมานิ ทส วตฺถูนิ, จตฺตาโร อาหารา, จตฺตาโร วิปลฺลาสา, จตฺตาริ อุปาทานานิ, จตฺตาโร โยคา, จตฺตาโร คนฺถา, จตฺตาโร อาสวา, จตฺตาโร โอฆา, จตฺตาโร สลฺลา, จตสฺโส วิญฺญาณ\-ฏฺฐิติโย จตฺตาริ อคติคมนานิ. ปฐเม อาหาเร ปฐโม วิปลฺลาโส, ทุติเย อาหาเร ทุติโย วิปลฺลาโส, ตติเย อาหาเร ตติโย วิปลฺลาโส, จตุตฺเถ อาหาเร จตุตฺโถ วิปลฺลาโส. ปฐเม วิปลฺลาเส ปฐมํ อุปาทานํ. ทุติเย วิปลฺลาเส ทุติยํ อุปาทานํ, ตติเย วิปลฺลาเส ตติยํ อุปาทานํ, จตุตฺเถ วิปลฺลาเส จตุตฺถํ อุปาทานํ. ปฐเม อุปาทาเน ปฐโม โยโค, ทุติเย อุปาทาเน ทุติโย โยโค, ตติเย อุปาทาเน ตติโย โยโค, จตุตฺเถ อุปาทาเน จตุตฺโถ โยโค. ปฐเม โยเค ปฐโม คนฺโถ, ทุติเย โยเค ทุติโย คนฺโถ, ตติเย โยเค ตติโย คนฺโถ, จตุตฺเถ โยเค จตุตฺโถ คนฺโถ, ปฐเม คนฺเถ ปฐโม อาสโว, ทุติเย คนฺเถ ทุติโย อาสโว, ตติเย คนฺเถ ตติโย อาสโว, จตุตฺเถ คนฺเถ จตุตฺโถ อาสโว. ปฐเม อาสเว ปฐโม โอโฆ, ทุติเย อาสเว ทุติโย โอโฆ, ตติเย อาสเว ตติโย โอโฆ, จตุตฺเถ อาสเว จตุตฺโถ โอโฆ. ปฐเม โอเฆ ปฐโม สลฺโล, ทุติเย โอเฆ ทุติโย สลฺโล, ตติเย โอเฆ ตติโย สลฺโล, จตุตฺเถ โอเฆ จตุตฺโถ สลฺโล. ปฐเม สลฺเล ปฐมา วิญฺญาณฏฺฐิติ, ทุติเย สลฺเล ทุติยา วิญฺญาณฏฺฐิติ, ตติเย สลฺเล ตติยา วิญฺญาณฏฺฐิติ, จตุตฺเถ สลฺเล จตุตฺถี\csromansharedfootnote{8e64e3ea2300b5fa}{จตุตฺถา (สี)} วิญฺญาณฏฺฐิติ, ปฐมายํ วิญฺญาณ\-ฏฺฐิติยํ ปฐมํ อคติคมนํ. ทุติยายํ วิญฺญาณ\-ฏฺฐิติยํ ทุติยํ อคติคมนํ. ตติยายํ วิญฺญาณ\-ฏฺฐิติยํ ตติยํ อคติคมนํ, จตุตฺถิยํ\csromansharedfootnote{b30425a6fbf2a049}{จตุตฺถายํ (สี)} วิญฺญาณ\-ฏฺฐิติยํ จตุตฺถํ อคติคมนํ.}

\proseitem{83}{\csromanfolio{95}ตตฺถ โย จ กพฬีกาโร อาหาโร ผสฺโส อาหาโร, อิเม ตณฺหาจริตสฺส ปุคฺคลสฺส อุปกฺกิเลสา. โย จ มโนสญฺเจตนา\-หาโร โย จ วิญฺญาณาหาโร, อิเม ทิฏฺฐิ\-จริตสฺส ปุคฺคลสฺส อุปกฺกิเลสา. ตตฺถ โย จ “อสุเภ สุภนฺ”ติ วิปลฺลาโส, โย จ “ทุกฺเข สุขนฺ”ติ วิปลฺลาโส, อิเม ตณฺหาจริตสฺส ปุคฺคลสฺส อุปกฺกิเลสา. โย จ “อนิจฺเจ นิจฺจนฺ”ติ วิปลฺลาโส, โย จ “อนตฺตนิ อตฺตา”ติ วิปลฺลาโส, อิเม ทิฏฺฐิ\-จริตสฺส ปุคฺคลสฺส อุปกฺกิเลสา.  ตตฺถ ยญฺจ กามุปาทานํ ยญฺจ ภวุปาทานํ, อิเม ตณฺหาจริตสฺส ปุคฺคลสฺส อุปกฺกิเลสา. ยญฺจ ทิฏฺฐุปาทานํ ยญฺจ อตฺตวาทุ\-ปาทานํ, อิเม ทิฏฺฐิ\-จริตสฺส ปุคฺคลสฺส อุปกฺกิเลสา.  ตตฺถ โย จ กามโยโค, โย จ ภวโยโค, อิเม ตณฺหาจริตสฺส ปุคฺคลสฺส อุปกฺกิเลสา. โย จ ทิฏฺฐิโยโค, โย จ อวิชฺชาโยโค, อิเม ทิฏฺฐิ\-จริตสฺส ปุคฺคลสฺส อุปกฺกิเลสา.  ตตฺถ โย จ อภิชฺฌา\-กายคนฺโถ, โย จ พฺยาปาโท กายคนฺโถ, อิเม ตณฺหาจริตสฺส ปุคฺคลสฺส อุปกฺกิเลสา. โย จ ปรามาส\-กายคนฺโถ, โย จ อิทํสจฺจาภินิเวส\-กายคนฺโถ, อิเม ทิฏฺฐิ\-จริตสฺส ปุคฺคลสฺส อุปกฺกิเลสา.  ตตฺถ โย จ กามาสโว, โย จ ภวาสโว, อิเม ตณฺหาจริตสฺส ปุคฺคลสฺส อุปกฺกิเลสา. โย จ ทิฏฺฐาสโว, โย จ อวิชฺชาสโว, อิเม ทิฏฺฐิ\-จริตสฺส ปุคฺคลสฺส อุปกฺกิเลสา.  ตตฺถ โย จ กาโมโฆ, โย จ ภโวโฆ, อิเม ตณฺหาจริตสฺส ปุคฺคลสฺส อุปกฺกิเลสา. โย จ ทิฏฺโฐโฆ, โย จ อวิชฺโชโฆ, อิเม ทิฏฺฐิ\-จริตสฺส ปุคฺคลสฺส อุปกฺกิเลสา.  ตตฺถ โย จ ราคสลฺโล, โย จ โทสสลฺโล, อิเม ตณฺหาจริตสฺส ปุคฺคลสฺส อุปกฺกิเลสา. โย จ มานสลฺโล, โย จ โมหสลฺโล, อิเม ทิฏฺฐิ\-จริตสฺส ปุคฺคลสฺส อุปกฺกิเลสา. ตตฺถ ยา จ รูปูปคา วิญฺญาณฏฺฐิติ, ยา จ เวทนูปคา วิญฺญาณฏฺฐิติ, อิเม ตณฺหาจริตสฺส ปุคฺคลสฺส อุปกฺกิเลสา. ยา จ สญฺญูปคา วิญฺญาณฏฺฐิติ, ยา จ สงฺขารูปคา วิญฺญาณฏฺฐิติ, อิเม ทิฏฺฐิ\-จริตสฺส ปุคฺคลสฺส อุปกฺกิเลสา. ตตฺถ ยญฺจ ฉนฺทา อคติคมนํ ยญฺจ โทสา อคติคมนํ, อิเม ตณฺหาจริตสฺส ปุคฺคลสฺส อุปกฺกิเลสา. ยญฺจ ภยา อคติคมนํ, ยญฺจ โมหา อคติคมนํ, อิเม ทิฏฺฐิ\-จริตสฺส ปุคฺคลสฺส อุปกฺกิเลสา.}

\proseitem{84}{ตตฺถ กพฬีกาเร อาหาเร “อสุเภ สุภนฺ”ติ วิปลฺลาโส. ผสฺเส อาหาเร “ทุกฺเข สุขนฺ”ติ วิปลฺลาโส. วิญฺญาเณ อาหาเร \csromanfolio{96}“อนิจฺเจ นิจฺจนฺ”ติ วิปลฺลาโส. มโน\-สญฺเจตนาย อาหาเร “อนตฺตนิ อตฺตา”ติ วิปลฺลาโส, ปฐเม วิปลฺลาเส ฐิโต กาเม อุปาทิยติ, อิทํ วุจฺจติ กามุปาทานํ. ทุติเย วิปลฺลาเส ฐิโต อนาคตํ ภวํ อุปาทิยติ, อิทํ วุจฺจติ ภวุปาทานํ. ตติเย วิปลฺลาเส ฐิโต สํสารา\-ภินนฺทินิํ ทิฏฺฐิํ อุปาทิยติ, อิทํ วุจฺจติ ทิฏฺฐุปาทานํ, จตุตฺเถ วิปลฺลาเส ฐิโต อตฺตานํ กปฺปิยํ อุปาทิยติ, อิทํ วุจฺจติ อตฺตวาทุ\-ปาทานํ.}

\prose{กามุปาทาเนน กาเมหิ สํยุชฺชติ, อยํ วุจฺจติ กามโยโค. ภวุปาทาเนน ภเวหิ สํยุชฺชติ, อยํ วุจฺจติ ภวโยโค. ทิฏฺฐุปาทาเนน ปาปิกาย ทิฏฺฐิยา สํยุชฺชติ, อยํ วุจฺจติ ทิฏฺฐิโยโค. อตฺตวาทุ\-ปาทาเนน อวิชฺชาย สํยุชฺชติ, อยํ วุจฺจติ อวิชฺชาโยโค.}

\prose{ปฐเม โยเค ฐิโต อภิชฺฌาย กายํ คนฺถติ, อยํ วุจฺจติ อภิชฺฌา\-กายคนฺโถ. ทุติเย โยเค ฐิโต พฺยาปาเทน กายํ คนฺถติ, อยํ วุจฺจติ พฺยาปาท\-กายคนฺโถ. ตติเย โยเค ฐิโต ปรามาเสน กายํ คนฺถติ, อยํ วุจฺจติ ปรามาส\-กายคนฺโถ. จตุตฺเถ โยเค ฐิโต อิทํสจฺจา\-ภินิเวเสน กายํ คนฺถติ, อยํ วุจฺจติ อิทํสจฺจาภินิเวส\-กายคนฺโถ.}

\prose{ตสฺส เอวํคนฺถิตา กิเลสา อาสวนฺติ. กุโต จ วุจฺจติ อาสวนฺตีติ, อนุสยโต วา ปริยุฏฺฐานโต วา. ตตฺถ อภิชฺฌา\-กายคนฺเถน กามาสโว, พฺยาปาท\-กายคนฺเถน ภวาสโว, ปรามาส\-กายคนฺเถน ทิฏฺฐาสโว, อิทํสจฺจาภินิเวส\-กายคนฺเถน อวิชฺชาสโว.}

\prose{ตสฺส อิเม จตฺตาโร อาสวา เวปุลฺลํ คตา โอฆา ภวนฺติ. อิติ อาสวเวปุลฺลา โอฆเวปุลฺลํ. ตตฺถ กามาสเวน กาโมโฆ, ภวาสเวน ภโวโฆ, ทิฏฺฐาสเวน ทิฏฺโฐโฆ, อวิชฺชาสเวน อวิชฺโชโฆ.}

\prose{ตสฺส อิเม จตฺตาโร โอฆา อนุสย\-สหคตา อชฺฌาสยํ อนุปวิฏฺฐา หทยํ อาหจฺจ ติฏฺฐนฺติ, เตน วุจฺจนฺติ สลฺลาอิติ. ตตฺถ กาโมเฆน ราคสลฺโล, ภโวเฆน โทสสลฺโล, ทิฏฺโฐเฆน มานสลฺโล, อวิชฺโชเฆน โมหสลฺโล.}

\prose{\csromanfolio{97}ตสฺส อิเมหิ จตูหิ สลฺเลหิ ปริยาทินฺนํ\csromansharedfootnote{6a9099f7dbcedaf4}{ปริยาทิณฺณํ (ก)} วิญฺญาณํ จตูสุ ธมฺเมสุ สณฺฐหติ รูเป เวทนาย สญฺญาย สงฺขาเรสุ. ตตฺถ ราคสลฺเลน นนฺทู\-ปเสจเนน วิญฺญาเณน รูปูปคา วิญฺญาณฏฺฐิติ, โทสสลฺเลน นนฺทู\-ปเสจเนน วิญฺญาเณน เวทนูปคา วิญฺญาณฏฺฐิติ, มานสลฺเลน นนฺทู\-ปเสจเนน วิญฺญาเณน สญฺญูปคา วิญฺญาณฏฺฐิติ, โมหสลฺเลน นนฺทู\-ปเสจเนน วิญฺญาเณน สงฺขารูปคา วิญฺญาณฏฺฐิติ.}

\prose{ตสฺส อิมาหิ จตูหิ วิญฺญาณ\-ฏฺฐิตีหิ อุปตฺถทฺธํ วิญฺญาณํ จตูหิ ธมฺเมหิ อคติํ คจฺฉติ ฉนฺทา โทสา ภยา โมหา. ตตฺถ ราเคน ฉนฺทาคติํ คจฺฉติ, โทเสน โทสาคติํ คจฺฉติ, ภเยน ภยาคติํ คจฺฉติ, โมเหน โมหาคติํ คจฺฉติ. อิติ โข ตญฺจ กมฺมํ อิเม จ กิเลสา, เอส เหตุ สํสารสฺส, เอวํ สพฺเพ กิเลสา จตูหิ วิปลฺลาเสหิ นิทฺทิสิตพฺพา.}

\proseitem{85}{ตตฺถ อิมา จตสฺโส ทิสา กพฬีกาโร อาหาโร “อสุเภ สุภนฺ”ติ วิปลฺลาโส, กามุปาทานํ, กามโยโค, อภิชฺฌา\-กายคนฺโถ, กามาสโว, กาโมโฆ, ราคสลฺโล, รูปูปคา วิญฺญาณฏฺฐิติ, ฉนฺทา อคติคมนนฺติ ปฐมา ทิสา.}

\prose{ผสฺโส อาหาโร, “ทุกฺเข สุขนฺ”ติ วิปลฺลาโส, ภวุปาทานํ, ภวโยโค, พฺยาปาท\-กายคนฺโถ, ภวาสโว, ภโวโฆ, โทสสลฺโล, เวทนูปคา วิญฺญาณฏฺฐิติ, โทสา อคติคมนนฺติ ทุติยา ทิสา.}

\prose{วิญฺญาณาหาโร “อนิจฺเจ นิจฺจนฺ”ติ วิปลฺลาโส, ทิฏฺฐุปาทานํ, ทิฏฺฐิโยโค ปรามาส\-กายคนฺโถ, ทิฏฺฐาสโว, ทิฏฺโฐโฆ, มานสลฺโล, สญฺญูปคา วิญฺญาณฏฺฐิติ, ภยา อคติคมนนฺติ ตติยา ทิสา.}

\prose{มโนสญฺเจตนา\-หาโร “อนตฺตนิ อตฺตา”ติ วิปลฺลาโส, อตฺตวาทุ\-ปาทานํ, อวิชฺชาโยโค, อิทํสจฺจาภินิเวส\-กายคนฺโถ, อวิชฺชาสโว, อวิชฺโชโฆ, โมหสลฺโล, สงฺขารูปคา วิญฺญาณฏฺฐิติ, โมหา อคติคมนนฺติ จตุตฺถี ทิสา.}

\prose{ตตฺถ โย จ กพฬีกาโร อาหาโร โย จ “อสุเภ สุภนฺ”ติ วิปลฺลาโส, กามุปาทานํ, กามโยโค, อภิชฺฌา\-กายคนฺโถ, \csromanfolio{98}กามาสโว, กาโมโฆ, ราคสลฺโล, รูปูปคา วิญฺญาณฏฺฐิติ ฉนฺทา อคติคมนนฺติ, อิเมสํ ทสนฺนํ สุตฺตานํ เอโก อตฺโถ, พฺยญฺชนเมว นานํ. อิเม ราคจริตสฺส ปุคฺคลสฺส อุปกฺกิเลสา.}

\prose{ตตฺถ โย จ ผสฺโส อาหาโร โย จ “ทุกฺเข สุขนฺ”ติ วิปลฺลาโส ภวุปาทานํ, ภวโยโค, พฺยาปาท\-กายคนฺโถ, ภวาสโว, ภโวโฆ, โทสสลฺโล, เวทนูปคา วิญฺญาณฏฺฐิติ, โทสา อคติคมนนฺติ อิเมสํ ทสนฺนํ สุตฺตานํ เอโก อตฺโถ พฺยญฺชนเมว นานํ, อิเม โทสจริตสฺส ปุคฺคลสฺส อุปกฺกิเลสา.}

\prose{ตตฺถ โย จ วิญฺญาณาหาโร โย จ “อนิจฺเจ นิจฺจนฺ”ติ วิปลฺลาโส, ทิฏฺฐุปาทานํ, ทิฏฺฐิโยโค, ปรามาส\-กายคนฺโถ, ทิฏฺฐาสโว, ทิฏฺโฐโฆ, มานสลฺโล, สญฺญูปคา วิญฺญาณฏฺฐิติ, ภยา อคติคมนนฺติ อิเมสํ ทสนฺนํ สุตฺตานํ เอโก อตฺโถ, พฺยญฺชนเมว นานํ. อิเม ทิฏฺฐิ\-จริตสฺส มนฺทสฺส อุปกฺกิเลสา.}

\prose{ตตฺถ โย จ มโนสญฺเจตนา\-หาโร โย จ “อนตฺตนิ อตฺตา”ติ วิปลฺลาโส, อตฺตวาทุ\-ปาทานํ, อวิชฺชาโยโค, อิทํสจฺจาภินิเวส\-กายคนฺโถ, อวิชฺชาสโว, อวิชฺโชโฆ, โมหสลฺโล, สงฺขารูปคา วิญฺญาณฏฺฐิติ, โมหา อคติคมนนฺติ, อิเมสํ ทสนฺนํ สุตฺตานํ เอโก อตฺโถ, พฺยญฺชนเมว นานํ. อิเม ทิฏฺฐิ\-จริตสฺส อุทตฺตสฺส\csromansharedfootnote{390b77ff2df70690}{อุทตฺถสฺส (สี, ก)} อุปกฺกิเลสา.}

\prose{ตตฺถ โย จ กพฬีกาโร อาหาโร โย จ ผสฺโส อาหาโร, อิเม อปฺปณิหิเตน วิโมกฺขมุเขน ปริญฺญํ คจฺฉนฺติ, วิญฺญาณาหาโร สุญฺญตาย, มโนสญฺเจตนา\-หาโร อนิมิตฺเตน, ตตฺถ โย จ “อสุเภ สุภนฺ”ติ วิปลฺลาโส, โย จ “ทุกฺเข สุขนฺ”ติ วิปลฺลาโส, อิเม อปฺปณิหิเตน วิโมกฺขมุเขน ปหานํ อพฺภตฺถํ คจฺฉนฺติ. “อนิจฺเจ นิจฺจนฺ”ติ วิปลฺลาโส สุญฺญตาย, “อนตฺตนิ อตฺตา”ติ วิปลฺลาโส อนิมิตฺเตน.  ตตฺถ กามุปาทานญฺจ ภวุปาทานญฺจ อปฺปณิหิเตน วิโมกฺขมุเขน ปหานํ คจฺฉนฺติ. ทิฏฺฐุปาทานํ สุญฺญตาย, อตฺตวาทุ\-ปาทานํ อนิมิตฺเตน.  ตตฺถ กามโยโค จ ภวโยโค จ อปฺปณิหิเตน วิโมกฺขมุเขน ปหานํ คจฺฉนฺติ, ทิฏฺฐิโยโค สุญฺญตาย, อวิชฺชาโยโค อนิมิตฺเตน.  ตตฺถ อภิชฺฌา\-กายคนฺโถ จ พฺยาปาท\-กายคนฺโถ จ \csromanfolio{99}อปฺปณิหิเตน วิโมกฺขมุเขน ปหานํ คจฺฉนฺติ, ปรามาส\-กายคนฺโถ สุญฺญตาย, อิทํสจฺจาภินิเวส\-กายคนฺโถ อนิมิตฺเตน.}

\prose{ตตฺถ กามาสโว จ ภวาสโว จ อปฺปณิหิเตน วิโมกฺขมุเขน ปหานํ คจฺฉนฺติ, ทิฏฺฐาสโว สุญฺญตาย, อวิชฺชาสโว อนิมิตฺเตน.  ตตฺถ กาโมโฆ จ ภโวโฆ จ อปฺปณิหิเตน วิโมกฺขมุเขน ปหานํ คจฺฉนฺติ, ทิฏฺโฐโฆ สุญฺญตาย, อวิชฺโชโฆ อนิมิตฺเตน.  ตตฺถ ราคสลฺโล จ โทสสลฺโล จ อปฺปณิหิเตน วิโมกฺขมุเขน ปหานํ คจฺฉนฺติ, มานสลฺโล สุญฺญตาย, โมหสลฺโล อนิมิตฺเตน.  ตตฺถ รูปูปคา จ วิญฺญาณฏฺฐิติ เวทนูปคา จ วิญฺญาณฏฺฐิติ อปฺปณิหิเตน วิโมกฺขมุเขน ปริญฺญํ คจฺฉนฺติ, สญฺญูปคา สุญฺญตาย, สงฺขารูปคา อนิมิตฺเตน.}

\prose{ตตฺถ ฉนฺทา จ อคติคมนํ โทสา จ อคติคมนํ อปฺปณิหิเตน วิโมกฺขมุเขน ปหานํ คจฺฉนฺติ, ภยา อคติคมนํ สุญฺญตาย, โมหา อคติคมนํ อนิมิตฺเตน วิโมกฺขมุเขน ปหานํ คจฺฉนฺติ. อิติ สพฺเพ โลก\-วฏฺฏานุสาริโน ธมฺมา นิยฺยนฺติ. เต โลกา ตีหิ วิโมกฺขมุเขหิ.}

\prose{86. ตตฺริทํ นิยฺยานํ\csromandash{}}

\prose{จตสฺโส ปฏิปทา, จตฺตาโร สติปฏฺฐานา, จตฺตาริ ฌานานิ, จตฺตาโร วิหารา, จตฺตาโร สมฺมปฺปธานา, จตฺตาโร อจฺฉริยา อพฺภุตา ธมฺมา, จตฺตาริ อธิฏฺฐานานิ, จตสฺโส สมาธิภาวนา, จตฺตาโร สุขภาคิยา ธมฺมา, จตสฺโส อปฺปมาณา.}

\prose{ปฐมา ปฏิปทา ปฐมํ สติปฏฺฐานํ, ทุติยา ปฏิปทา ทุติยํ สติปฏฺฐานํ, ตติยา ปฏิปทา ตติยํ สติปฏฺฐานํ, จตุตฺถี ปฏิปทา จตุตฺถํ สติปฏฺฐานํ. ปฐมํ สติปฏฺฐานํ ปฐมํ ฌานํ, ทุติยํ สติปฏฺฐานํ ทุติยํ ฌานํ, ตติยํ สติปฏฺฐานํ ตติยํ ฌานํ. จตุตฺถํ สติปฏฺฐานํ จตุตฺถํ ฌานํ. ปฐมํ ฌานํ ปฐโม วิหาโร, ทุติยํ ฌานํ ทุติโย วิหาโร, ตติยํ ฌานํ ตติโย วิหาโร, จตุตฺถํ ฌานํ จตุตฺโถ วิหาโร.  ปฐโม วิหาโร ปฐมํ สมฺมปฺปธานํ, ทุติโย วิหาโร ทุติยํ สมฺมปฺปธานํ, ตติโย วิหาโร ตติยํ สมฺมปฺปธานํ, จตุตฺโถ วิหาโร จตุตฺถํ สมฺมปฺปธานํ.  ปฐมํ สมฺมปฺปธานํ ปฐโม อจฺฉริโย อพฺภุโต ธมฺโม, ทุติยํ ทุติโย. ตติยํ ตติโย. จตุตฺถํ สมฺมปฺปธานํ จตุตฺโถ อจฺฉริโย อพฺภุโต ธมฺโม.  ปฐโม อจฺฉริโย อพฺภุโต ธมฺโม}

\prose{\csromanfolio{100}ปฐมํ อธิฏฺฐานํ, ทุติโย อจฺฉริโย อพฺภุโต ธมฺโม ทุติยํ อธิฏฺฐานํ, ตติโย อจฺฉริโย อพฺภุโต ธมฺโม ตติยํ อธิฏฺฐานํ, จตุตฺโถ อจฺฉริโย อพฺภุโต ธมฺโม จตุตฺถํ อธิฏฺฐานํ.  ปฐมํ อธิฏฺฐานํ ปฐมา สมาธิภาวนา, ทุติยํ อธิฏฺฐานํ ทุติยา สมาธิภาวนา, ตติยํ อธิฏฺฐานํ ตติยา สมาธิภาวนา, จตุตฺถํ อธิฏฺฐานํ จตุตฺถี สมาธิภาวนา.  ปฐมา สมาธิภาวนา ปฐโม สุขภาคิโย ธมฺโม, ทุติยา สมาธิภาวนา ทุติโย สุขภาคิโย ธมฺโม, ตติยา สมาธิภาวนา ตติโย สุขภาคิโย ธมฺโม, จตุตฺถี สมาธิภาวนา จตุตฺโถ สุขภาคิโย ธมฺโม.  ปฐโม สุขภาคิโย ธมฺโม ปฐมํ อปฺปมาณํ, ทุติโย สุขภาคิโย ธมฺโม ทุติยํ อปฺปมาณํ, ตติโย สุขภาคิโย ธมฺโม ตติยํ อปฺปมาณํ, จตุตฺโถ สุขภาคิโย ธมฺโม จตุตฺถํ อปฺปมาณํ.  ปฐมา ปฏิปทา ภาวิตา พหุลีกตา\csromansharedfootnote{c82028f8d7ae6d4f}{พหุลิกตา (ก)} ปฐมํ สติปฏฺฐานํ ปริปูเรติ, ทุติยา ปฏิปทา ภาวิตา พหุลีกตา ทุติยํ สติปฏฺฐานํ ปริปูเรติ, ตติยา ปฏิปทา ภาวิตา พหุลีกตา ตติยํ สติปฏฺฐานํ ปริปูเรติ, จตุตฺถี ปฏิปทา ภาวิตา พหุลีกตา จตุตฺถํ สติปฏฺฐานํ ปริปูเรติ.  ปฐโม สติปฏฺฐาโน ภาวิโต พหุลีกโต ปฐมํ ฌานํ ปริปูเรติ, ทุติโย สติปฏฺฐาโน ภาวิโต พหุลีกโต ทุติยํ ฌานํ ปริปูเรติ, ตติโย สติปฏฺฐาโน ภาวิโต พหุลีกโต ตติยํ ฌานํ ปริปูเรติ, จตุตฺโถ สติปฏฺฐาโน ภาวิโต พหุลีกโต จตุตฺถํ ฌานํ ปริปูเรติ.}

\prose{ปฐมํ ฌานํ ภาวิตํ พหุลีกตํ ปฐมํ วิหารํ ปริปูเรติ, ทุติยํ ฌานํ ภาวิตํ พหุลีกตํ ทุติยํ วิหารํ ปริปูเรติ, ตติยํ ฌานํ ภาวิตํ พหุลีกตํ ตติยํ วิหารํ ปริปูเรติ, จตุตฺถํ ฌานํ ภาวิตํ พหุลีกตํ จตุตฺถํ วิหารํ ปริปูเรติ.  ปฐโม วิหาโร ภาวิโต พหุลีกโต อนุปฺปนฺนานํ ปาปกานํ อกุสลานํ ธมฺมานํ อนุปฺปาทํ ปริปูเรติ, ทุติโย วิหาโร ภาวิโต พหุลีกโต อุปฺปนฺนานํ ปาปกานํ อกุสลานํ ธมฺมานํ ปหานํ ปริปูเรติ, ตติโย วิหาโร ภาวิโต พหุลีกโต อนุปฺปนฺนานํ กุสลานํ ธมฺมานํ อุปฺปาทํ ปริปูเรติ, จตุตฺโถ วิหาโร ภาวิโต พหุลีกโต อุปฺปนฺนานํ กุสลานํ ธมฺมานํ ฐิติํ อสมฺโมสํ ภิยฺโยภาวํ ปริปูเรติ.  ปฐมํ สมฺมปฺปธานํ ภาวิตํ พหุลีกตํ}

\prose{\csromanfolio{101}มานปฺปหานํ ปริปูเรติ, ทุติยํ สมฺมปฺปธานํ ภาวิตํ พหุลีกตํ อาลย\-สมุคฺฆาตํ ปริปูเรติ, ตติยํ สมฺมปฺปธานํ ภาวิตํ พหุลีกตํ อวิชฺชาปหานํ ปริปูเรติ, จตุตฺถํ สมฺมปฺปธานํ ภาวิตํ พหุลีกตํ ภวูปสมํ ปริปูเรติ.  มานปฺปหานํ ภาวิตํ พหุลีกตํ สจฺจา\-ธิฏฺฐานํ ปริปูเรติ, อาลย\-สมุคฺฆาโต ภาวิโต พหุลีกโต จาคาธิฏฺฐานํ ปริปูเรติ, อวิชฺชาปหานํ ภาวิตํ พหุลีกตํ ปญฺญา\-ธิฏฺฐานํ ปริปูเรติ, ภวูปสโม ภาวิโต พหุลีกโต อุปสมา\-ธิฏฺฐานํ ปริปูเรติ.  สจฺจา\-ธิฏฺฐานํ ภาวิตํ พหุลีกตํ ฉนฺทสมาธิํ ปริปูเรติ, จาคาธิฏฺฐานํ ภาวิตํ พหุลีกตํ วีริยสมาธิํ ปริปูเรติ, ปญฺญา\-ธิฏฺฐานํ ภาวิตํ พหุลีกตํ จิตฺตสมาธิํ ปริปูเรติ, อุปสมา\-ธิฏฺฐานํ ภาวิตํ พหุลีกตํ วีมํสา\-สมาธิํ ปริปูเรติ. ฉนฺทสมาธิ ภาวิโต พหุลีกโต อินฺทฺริย\-สํวรํ ปริปูเรติ, วีริยสมาธิ ภาวิโต พหุลีกโต ตปํ ปริปูเรติ, จิตฺตสมาธิ ภาวิโต พหุลีกโต พุทฺธิํ ปริปูเรติ, วีมํสาสมาธิ ภาวิโต พหุลีกโต สพฺพู\-ปธิ\-ปฏินิสฺสคฺคํ ปริปูเรติ.  อินฺทฺริยสํวโร ภาวิโต พหุลีกโต เมตฺตํ ปริปูเรติ, ตโป ภาวิโต พหุลีกโต กรุณํ ปริปูเรติ, พุทฺธิ ภาวิตา พหุลีกตา มุทิตํ ปริปูเรติ, สพฺพู\-ปธิ\-ปฏินิสฺสคฺโค ภาวิโต พหุลีกโต อุเปกฺขํ ปริปูเรติ.}

\proseitem{87}{ตตฺถ อิมา จตสฺโส ทิสา ปฐมา ปฏิปทา ปฐโม สติปฏฺฐาโน ปฐมํ ฌานํ ปฐโม วิหาโร ปฐโม สมฺมปฺปธาโน ปฐโม อจฺฉริโย อพฺภุโต ธมฺโม สจฺจา\-ธิฏฺฐานํ ฉนฺทสมาธิ อินฺทฺริยสํวโร เมตฺตา อิติ ปฐมา ทิสา.}

\prose{ทุติยา ปฏิปทา ทุติโย สติปฏฺฐาโน ทุติโย วิหาโร ทุติโย สมฺมปฺปธาโน ทุติโย อจฺฉริโย อพฺภุโต ธมฺโม ภวา\-ธิฏฺฐานํ วีริยสมาธิ ตโป กรุณา อิติ ทุติยา ทิสา.}

\prose{ตติยา ปฏิปทา ตติโย สติปฏฺฐาโน ตติยํ ฌานํ ตติโย วิหาโร ตติโย สมฺมปฺปธาโน ตติโย อจฺฉริโย อพฺภุโต ธมฺโม ปญฺญา\-ธิฏฺฐานํ จิตฺตสมาธิ พุทฺธิ มุทิตา อิติ ตติยา ทิสา.}

\prose{จตุตฺถี ปฏิปทา จตุตฺโถ สติปฏฺฐาโน จตุตฺถํ ฌานํ จตุตฺโถ วิหาโร จตุตฺโถ สมฺมปฺปธาโน จตุตฺโถ อจฺฉริโย อพฺภุโต ธมฺโม อุปสมา\-ธิฏฺฐานํ วีมํสาสมาธิ สพฺพู\-ปธิ\-ปฏินิสฺสคฺโค อุเปกฺขา อิติ จตุตฺถี ทิสา.}

\prose{\csromanfolio{102}ตตฺถ ปฐมา ปฏิปทา ปฐโม สติปฏฺฐาโน ปฐมํ ฌานํ ปฐโม วิหาโร ปฐโม สมฺมปฺปธาโน ปฐโม อจฺฉริโย อพฺภุโต ธมฺโม สจฺจา\-ธิฏฺฐานํ ฉนฺทสมาธิ อินฺทฺริยสํวโร, เมตฺตา อิติ อิเมสํ ทสนฺนํ สุตฺตานํ เอโก อตฺโถ, พฺยญฺชนเมว นานํ. อิทํ ราคจริตสฺส ปุคฺคลสฺส เภสชฺชํ.}

\prose{ทุติยา ปฏิปทา ทุติโย สติปฏฺฐาโน ทุติยํ ฌานํ ทุติโย วิหาโร ทุติโย สมฺมปฺปธาโน ทุติโย อจฺฉริโย อพฺภุโต ธมฺโม จาคาธิฏฺฐานํ วีริยสมาธิ ตโป กรุณา อิติ อิเมสํ ทสนฺนํ สุตฺตานํ เอโก อตฺโถ, พฺยญฺชนเมว นานํ. อิทํ โทสจริตสฺส ปุคฺคลสฺส เภสชฺชํ.}

\prose{ตติยา ปฏิปทา ตติโย สติปฏฺฐาโน ตติยํ ฌานํ ตติโย วิหาโร ตติโย สมฺมปฺปธาโน ตติโย อจฺฉริโย อพฺภุโต ธมฺโม ปญฺญา\-ธิฏฺฐานํ จิตฺตสมาธิ พุทฺธิ มุทิตา อิติ อิเมสํ ทสนฺนํ สุตฺตานํ เอโก อตฺโถ, พฺยญฺชนเมว นานํ. อิทํ ทิฏฺฐิ\-จริตสฺส มนฺทสฺส เภสชฺชํ.}

\prose{จตุตฺถี ปฏิปทา จตุตฺโถ สติปฏฺฐาโน จตุตฺถํ ฌานํ จตุตฺโถ วิหาโร จตุตฺโถ สมฺมปฺปธาโน จตุตฺโถ อจฺฉริโย อพฺภุโต ธมฺโม อุปสมา\-ธิฏฺฐานํ วีมํสาสมาธิ สพฺพู\-ปธิ\-ปฏินิสฺสคฺโค อุเปกฺขา อิติ อิเมสํ ทสนฺนํ สุตฺตานํ เอโก อตฺโถ, พฺยญฺชนเมว นานํ. อิทํ ทิฏฺฐิ\-จริตสฺส อุทตฺตสฺส เภสชฺชํ.}

\prose{ตตฺถ ทุกฺขา จ ปฏิปทา ทนฺธาภิญฺญา ทุกฺขา จ ปฏิปทา ขิปฺปาภิญฺญา อปฺปณิหิตํ วิโมกฺขมุขํ, สุขา ปฏิปทา ทนฺธาภิญฺญา สุญฺญตํ วิโมกฺขมุขํ, สุขา ปฏิปทา ขิปฺปาภิญฺญา อนิมิตฺตํ วิโมกฺขมุขํ.}

\prose{ตตฺถ กาเย กายานุปสฺสิตา สติปฏฺฐานญฺจ เวทนาสุ เวทนานุปสฺสิตา ปติปฏฺฐานญฺจ อปฺปณิหิตํ วิโมกฺขมุขํ, จิตฺเต จิตฺตานุปสฺสิตา สุญฺญตํ วิโมกฺขมุขํ. ธมฺเมสุ ธมฺมานุปสฺสิตา อนิมิตฺตํ วิโมกฺขมุขํ.}

\prose{ตตฺถ ปฐมญฺจ ฌานํ ทุติยญฺจ ฌานํ อปฺปณิหิตํ วิโมกฺขมุขํ, ตติยํ ฌานํ สุญฺญตํ วิโมกฺขมุขํ, จตุตฺถํ ฌานํ อนิมิตฺตํ วิโมกฺขมุขํ.}

\prose{ตตฺถ ปฐโม จ วิหาโร ทุติโย จ วิหาโร อปฺปณิหิตํ วิโมกฺขมุขํ, ตติโย วิหาโร สุญฺญตํ วิโมกฺขมุขํ, จตุตฺโถ วิหาโร อนิมิตฺตํ วิโมกฺขมุขํ.}

\prose{\csromanfolio{103}ยตฺถ ปฐมญฺจ สมฺมปฺปธานํ ทุติยญฺจ สมฺมปฺปธานํ อปฺปณิหิตํ วิโมกฺขมุขํ, ตติยํ สมฺมปฺปธานํ สุญฺญตํ วิโมกฺขมุขํ, จตุตฺถํ สมฺมปฺปธานํ อนิมิตฺตํ วิโมกฺขมุขํ.}

\prose{ตตฺถ มานปฺปหานญฺจ อาลย\-สมุคฺฆาโต จ อปฺปณิหิตํ วิโมกฺขมุขํ, อวิชฺชาปหานํ สุญฺญตํ วิโมกฺขมุขํ, ภวูปสโม อนิมิตฺตํ วิโมกฺขมุขํ.}

\prose{ตตฺถ สจฺจาธิฏฺฐานญฺจ จาคาธิฏฺฐานญฺจ อปฺปณิหิตํ วิโมกฺขมุขํ, ปญฺญา\-ธิฏฺฐานํ สุญฺญตํ วิโมกฺขมุขํ, อุปสมา\-ธิฏฺฐานํ อนิมิตฺตํ วิโมกฺขมุขํ.}

\prose{ตตฺถ ฉนฺทสมาธิ จ วีริยสมาธิ จ อปฺปณิหิตํ วิโมกฺขมุขํ, จิตฺตสมาธิ สุญฺญตํ วิโมกฺขมุขํ, วีมํสาสมาธิ อนิมิตฺตํ วิโมกฺขมุขํ.}

\prose{ตตฺถ อินฺทฺริยสํวโร จ ตโป จ อปฺปณิหิตํ วิโมกฺขมุขํ, พุทฺธิ สุญฺญตํ วิโมกฺขมุขํ สพฺพู\-ปธิ\-ปฏินิสฺสคฺโค อนิมิตฺตํ วิโมกฺขมุขํ.}

\prose{ตตฺถ เมตฺตา จ กรุณา จ อปฺปณิหิตํ วิโมกฺขมุขํ, มุทิตา สุญฺญตํ วิโมกฺขมุขํ อุเปกฺขา อนิมิตฺตํ วิโมกฺขมุขํ.}

\prose{เตสํ วิกฺกีฬิตํ. จตฺตาโร อาหารา เตสํ ปฏิปกฺโข จตสฺโส ปฏิปทา ฯเปฯ จตฺตาโร วิปลฺลาสา เตสํ ปฏิปกฺโข จตฺตาโร สติปฏฺฐานา. จตฺตาริ อุปาทานานิ เตสํ ปฏิปกฺโข จตฺตาริ ฌานานิ. จตฺตาโร โยคา เตสํ ปฏิปกฺโข จตฺตาโร วิหารา. จตฺตาโร คนฺถา เตสํ ปฏิปกฺโข จตฺตาโร สมฺมปฺปธานา. จตฺตาโร อาสวา เตสํ ปฏิปกฺโข จตฺตาโร อจฺฉริยา อพฺภุตา ธมฺมา. จตฺตาโร โอฆา เตสํ ปฏิปกฺโข จตฺตาริ อธิฏฺฐานานิ. จตฺตาโร สลฺลา เตสํ ปฏิปกฺโข จตสฺโส สมาธิภาวนา. จตสฺโส วิญฺญาณ\-ฏฺฐิติโย ตาสํ ปฏิปกฺโข จตฺตาโร สุขภาคิยา ธมฺมา. จตฺตาริ อคติคมนานิ เตสํ ปฏิปกฺโข จตสฺโส อปฺปมาณา.}

\prose{สีหา พุทฺธา ปจฺเจกพุทฺธา สาวกา จ หตราคโทสโมหา, เตสํ วิกฺกีฬิตํ ภาวนา สจฺฉิกิริยา พฺยนฺตีกิริยา จ. วิกฺกีฬิตํ อินฺทฺริยา\-ธิฏฺฐานํ วิกฺกีฬิตํ วิปริยาสา\-นธิฏฺฐานญฺจ. อินฺทฺริยานิ สทฺธมฺม\-โคจโร วิปริยาสา กิเลสโคจโร. อยํ วุจฺจติ สีหวิกฺกีฬิตสฺส จ นยสฺส ทิสาโลจนสฺส จ นยสฺส ภูมีติ. เตนาห “โย เนติ วิปลฺลาเสหิ สํกิเลเส”ติ. เวยฺยากรเณสุ หิ เย “กุสลากุสลา”ติ จ.}

\prose{\csromanfolio{104}ตตฺถ เย ทุกฺขาย ปฏิปทาย ทนฺธาภิญฺญาย ขิปฺปาภิญฺญาย จ นิยฺยนฺติ, อิเม เทฺว ปุคฺคลา, เย สุขาย ปฏิปทาย ทนฺธาภิญฺญาย ขิปฺปาภิญฺญาย จ นิยฺยนฺติ, อิเม เทฺว ปุคฺคลา. เตสํ จตุนฺนํ ปุคฺคลานํ อยํ สํกิเลโส, จตฺตาโร อาหารา, จตฺตาโร วิปลฺลาสา, จตฺตาริ อุปาทานานิ, จตฺตาโร โยคา, จตฺตาโร คนฺถา, จตฺตาโร อาสวา, จตฺตาโร โอฆา, จตฺตาโร สลฺลา, จตสฺโส วิญฺญาณ\-ฏฺฐิติโย, จตฺตาริ อคติคมนานีติ. เตสํ จตุนฺนํ ปุคฺคลานํ อิทํ โวทานํ, จตสฺโส ปฏิปทา, จตฺตาโร สติปฏฺฐานา, จตฺตาริ ฌานานิ, จตฺตาโร วิหารา, จตฺตาโร สมฺมปฺปธานา, จตฺตาโร อจฺฉริยา อพฺภุตา ธมฺมา, จตฺตาริ อธิฏฺฐานานิ, จตสฺโส สมาธิภาวนา, จตฺตาโร สุขภาคิยา ธมฺมา, จตสฺโส อปฺปมาณา อิติ.}

\proseitem{88}{ตตฺถ เย ทุกฺขาย ปฏิปทาย ทนฺธาภิญฺญาย ขิปฺปาภิญฺญาย จ นิยฺยนฺติ อิเม เทฺว ปุคฺคลา. เย สุขาย ปฏิปทาย ทนฺธาภิญฺญาย ขิปฺปาภิญฺญาย จ นิยฺยนฺติ, อิเม เทฺว ปุคฺคลา. ตตฺถ โย สุขาย ปฏิปทาย ขิปฺปาภิญฺญาย นิยฺยาติ, อยํ อุคฺฆฏิตญฺญู. โย สาธารณาย, อยํ วิปญฺจิตญฺญู. โย ทุกฺขาย ปฏิปทาย ทนฺธาภิญฺญาย นิยฺยาติ, อยํ เนยฺโย.}

\prose{ตตฺถ ภควา อุคฺฆฏิต\-ญฺญุสฺส ปุคฺคลสฺส สมถํ อุปทิสติ, เนยฺยสฺส วิปสฺสนํ, สมถ\-วิปสฺสนํ วิปญฺจิต\-ญฺญุสฺส. ตตฺถ ภควา อุคฺฆฏิต\-ญฺญุสฺส ปุคฺคลสฺส มุทุกํ ธมฺมเทสนํ อุปทิสติ, ติกฺขํ เนยฺยสฺส, มุทุติกฺขํ วิปญฺจิต\-ญฺญุสฺส, ตตฺถ ภควา อุคฺฆฏิต\-ญฺญุสฺส ปุคฺคลสฺส สํขิตฺเตน ธมฺมํ เทสยติ, สํขิตฺต\-วิตฺถาเรน วิปญฺจิต\-ญฺญุสฺส, วิตฺถาเรน เนยฺยสฺส. ตตฺถ ภควา อุคฺฆฏิต\-ญฺญุสฺส ปุคฺคลสฺส นิสฺสรณํ อุปทิสติ, วิปญฺจิต\-ญฺญุสฺส อาทีนวญฺจ นิสฺสรณญฺจ อุปทิสติ, เนยฺยสฺส อสฺสาทญฺจ อาทีนวญฺจ นิสฺสรณญฺจ อุปทิสติ. ตตฺถ ภควา อุคฺฆฏิต\-ญฺญุสฺส อธิปญฺญา\-สิกฺขํ ปญฺญาปยติ, อธิจิตฺตํ วิปญฺจิต\-ญฺญุสฺส อธิสีลํ เนยฺยสฺส.}

\prose{ตตฺถ เย ทุกฺขาย ปฏิปทาย ทนฺธาภิญฺญาย ขิปฺปาภิญฺญาย จ นิยฺยนฺติ, อิเม เทฺว ปุคฺคลา. เย สุขาย ปฏิปทาย ทนฺธาภิญฺญาย ขิปฺปาภิญฺญาย จ นิยฺยนฺติ. อิเม เทฺว ปุคฺคลา. อิติ โข จตฺตาริ หุตฺวา ตีณิ ภวนฺติ อุคฺฆฏิตญฺญู วิปญฺจิตญฺญู เนยฺโยติ.}

\prose{เตสํ ติณฺณํ ปุคฺคลานํ อยํ สํกิเลโส, ตีณิ อกุสลมูลานิ โลโภ อกุสลมูลํ โทโส อกุสลมูลํ โมโห อกุสลมูลํ, \csromanfolio{105}ตีณิ ทุจฺจริตานิ กายทุจฺจริตํ วจีทุจฺจริตํ มโนทุจฺจริตํ, ตโย อกุสลวิตกฺกา กามวิตกฺโก พฺยาปาทวิตกฺโก วิหิํสาวิตกฺโก, ติสฺโส อกุสลสญฺญา กามสญฺญา พฺยาปาทสญฺญา วิหิํสาสญฺญา, ติสฺโส วิปรีตสญฺญา นิจฺจสญฺญา สุขสญฺญา อตฺตสญฺญา, ติสฺโส เวทนา สุขา เวทนา ทุกฺขา เวทนา อทุกฺขมสุขา เวทนา, ติสฺโส ทุกฺขตา ทุกฺขทุกฺขตา สงฺขาร\-ทุกฺขตา วิปริณาม\-ทุกฺขตา, ตโย อคฺคี ราคคฺคิ โทสคฺคิ โมหคฺคิ, ตโย สลฺลา ราคสลฺโล โทสสลฺโล โมหสลฺโล, ติสฺโส ชฏา ราคชฏา โทสชฏา โมหชฏา, ติสฺโส อกุสลู\-ปปริกฺขา อกุสลํ กายกมฺมํ อกุสลํ วจีกมฺมํ อกุสลํ มโนกมฺมํ. ติสฺโส วิปตฺติโย สีลวิปตฺติ ทิฏฺฐิวิปตฺติ อาจารวิปตฺตีติ. เตสํ ติณฺณํ ปุคฺคลานํ อิทํ โวทานํ. ตีณิ กุสลมูลานิ อโลโภ กุสลมูลํ อโทโส กุสลมูลํ อโมโห กุสลมูลํ. ตีณิ สุจริตานิ กายสุจริตํ วจีสุจริตํ มโนสุจริตํ. ตโย กุสลวิตกฺกา เนกฺขมฺม\-วิตกฺโก อพฺยาปาทวิตกฺโก อวิหิํสา\-วิตกฺโก. ตโย สมาธี สวิตกฺโก สวิจาโร สมาธิ อวิตกฺโก วิจารมตฺโต สมาธิ อวิตกฺโก อวิจาโร สมาธิ, ติสฺโส กุสลสญฺญา เนกฺขมฺมสญฺญา อพฺยาปาทสญฺญา อวิหิํสาสญฺญา, ติสฺโส อวิปรีตสญฺญา อนิจฺจสญฺญา ทุกฺขสญฺญา อนตฺตสญฺญา, ติสฺโส กุสลู\-ปปริกฺขา กุสลํ กายกมฺมํ กุสลํ วจีกมฺมํ กุสลํ มโนกมฺมํ, ตีณิ โสเจยฺยานิ กายโสเจยฺยํ วจีโสเจยฺยํ มโนโสเจยฺยํ, ติสฺโส สมฺปตฺติโย สีลสมฺปตฺติ สมาธิ\-สมฺปตฺติ ปญฺญาสมฺปตฺติ, ติสฺโส สิกฺขา อธิสีลสิกฺขา อธิจิตฺต\-สิกฺขา อธิปญฺญา\-สิกฺขา, ตโย ขนฺธา สีลกฺขนฺโธ สมาธิ\-กฺขนฺโธ ปญฺญากฺขนฺโธ, ตีณิ วิโมกฺขมุขานิ สุญฺญตํ อนิมิตฺตํ อปฺปณิหิตนฺติ.}

\prose{อิติ โข จตฺตาริ หุตฺวา ตีณิ ภวนฺติ, ตีณิ หุตฺวา เทฺว ภวนฺติ ตณฺหาจริโต จ ทิฏฺฐิจริโต จ.}

\prose{เตสํ ทฺวินฺนํ ปุคฺคลานํ อยํ สํกิเลโส, ตณฺหา จ อวิชฺชา จ อหิริกญฺจ อโนตฺตปฺปญฺจ อสฺสติ จ อสมฺปชญฺญญฺจ อโยนิโส มนสิกาโร จ โกสชฺชญฺจ โทวจสฺสญฺจ อหํกาโร จ มมํกาโร จ อสฺสทฺธา จ ปมาโท จ อสทฺธมฺมสฺสวนญฺจ อสํวโร จ อภิชฺฌา จ พฺยาปาโท จ นีวรณญฺจ สํโยชนญฺจ โกโธ จ อุปนาโห จ มกฺโข จ ปลาโส จ อิสฺสา จ มจฺเฉรญฺจ มายา จ สาเฐยฺยญฺจ สสฺสตทิฏฺฐิ จ อุจฺเฉททิฏฺฐิ จาติ.}

\prose{\csromanfolio{106}เตสํ ทฺวินฺนํ ปุคฺคลานํ อิทํ โวทานํ, สมโถ จ วิปสฺสนา จ หิรี จ โอตฺตปฺปญฺจ สติ จ สมฺปชญฺญญฺจ โยนิโส มนสิกาโร จ วีริยารมฺโภ จ โสวจสฺสญฺจ ธมฺเม ญาณญฺจ อนฺวเย ญาณญฺจ ขเย ญาณญฺจ อนุปฺปาเท ญาณญฺจ สทฺธา จ อปฺปมาโท จ สทฺธมฺมสฺสวนญฺจ สํวโร จ อนภิชฺฌา จ อพฺยาปาโท จ ราควิราคา จ เจโตวิมุตฺติ อวิชฺชาวิราคา จ ปญฺญาวิมุตฺติ อภิสมโย จ อปฺปิจฺฉตา จ สนฺตุฏฺฐิ จ อกฺโกโธ จ อนุปนาโห จ อมกฺโข จ อปลาโส จ อิสฺสาปหานญฺจ มจฺฉริย\-ปฺปหานญฺจ วิชฺชา จ วิมุตฺติ จ สงฺขตา\-รมฺมโณ จ วิโมกฺโข อสงฺขตา\-รมฺมโณ จ วิโมกฺโข สอุปาทิเสสา จ นิพฺพานธาตุ อนุปาทิเสสา จ นิพฺพานธาตูติ.}

\prose{อยํ วุจฺจติ ติปุกฺขลสฺส จ นยสฺส องฺกุสสฺส จ นยสฺส ภูมีติ. เตนาห “โย อกุสเล สมูเลหิ เนตี”ติ “โอโลเกตฺวา ทิสโลจเนนา”ติ จ.}

\vspace{\tipitakaparskip}
\csromancenter{นิยุตฺตํ นยสมุฏฺฐานํ.}
\csromanmatikaanchor{mtk.38}
\csromantocmarkref{subsection}{สาสนปฏฺฐาน}{mtk.38}
\bookmark[dest={mtk.38},level=2]{สาสนปฏฺฐาน}
\csromanheader{2}{\textbf{สาสนปฏฺฐาน}}
\proseitem{89}{ตตฺถ อฏฺฐารส มูลปทา กุหิํ ทฏฺฐพฺพา, สาสนปฏฺฐาเน. ตตฺถ กตมํ สาสน\-ปฏฺฐานํ\csromandash{}สํกิเลส\-ภาคิยํ สุตฺตํ, วาสนาภาคิยํ สุตฺตํ, นิพฺเพธ\-ภาคิยํ สุตฺตํ, อเสกฺข\-ภาคิยํ สุตฺตํ, สํกิเลสภาคิยญฺจ วาสนาภาคิยญฺจ สุตฺตํ, สํกิเลสภาคิยญฺจ นิพฺเพธภาคิยญฺจ สุตฺตํ, สํกิเลสภาคิยญฺจ อเสกฺขภาคิยญฺจ สุตฺตํ, สํกิเลสภาคิยญฺจ นิพฺเพธภาคิยญฺจ อเสกฺขภาคิยญฺจ สุตฺตํ, สํกิเลสภาคิยญฺจ วาสนาภาคิยญฺจ นิพฺเพธภาคิยญฺจ สุตฺตํ, วาสนาภาคิยญฺจ นิพฺเพธภาคิยญฺจ สุตฺตํ, ตณฺหา\-สํกิเลส\-ภาคิยํ สุตฺตํ, ทิฏฺฐิ\-สํกิเลส\-ภาคิยํ สุตฺตํ, ทุจฺจริต\-สํกิเลส\-ภาคิยํ สุตฺตํ, ตณฺหา\-โวทาน\-ภาคิยํ สุตฺตํ, ทิฏฺฐิ\-โวทาน\-ภาคิยํ สุตฺตํ, ทุจฺจริต\-โวทาน\-ภาคิยํ สุตฺตํ.}

\prose{ตตฺถ สํกิเลโส ติวิโธ ตณฺหาสํกิเลโส ทิฏฺฐิ\-สํกิเลโส ทุจฺจริต\-สํกิเลโส. ตตฺถ ตณฺหาสํกิเลโส สมเถน วิสุชฺฌติ, โส สมโถ สมาธิ\-กฺขนฺโธ. ทิฏฺฐิ\-สํกิเลโส วิปสฺสนาย วิสุชฺฌติ, \csromanfolio{107}สา วิปสฺสนา ปญฺญากฺขนฺโธ. ทุจฺจริต\-สํกิเลโส สุจริเตน วิสุชฺฌติ, ตํ สุจริตํ สีลกฺขนฺโธ. ตสฺส สีเล ปติฏฺฐิตสฺส ยทิ อาสตฺติ อุปฺปชฺชติ ภเวสุ, เอวํ สายํ สมถ\-วิปสฺสนา ภาวนามยํ ปุญฺญ\-กฺริย\-วตฺถุ ภวติ ตตฺรูปปตฺติยา สํวตฺตติ. อิมานิ จตฺตาริ สุตฺตานิ, สาธารณานิ กตานิ อฏฺฐ ภวนฺติ, ตานิเยว อฏฺฐ สุตฺตานิ สาธารณานิ กตานิ โสฬส ภวนฺติ.}

\prose{อิเมหิ โสฬสหิ สุตฺเตหิ ภินฺเนหิ นววิธํ สุตฺตํ ภินฺนํ ภวติ. คาถาย คาถา อนุมินิตพฺพา, เวยฺยากรเณน เวยฺยากรณํ อนุมินิตพฺพํ. สุตฺเตน สุตฺตํ อนุมินิตพฺพํ.}

\proseitem{90}{ตตฺถ กตมํ สํกิเลส\-ภาคิยํ สุตฺตํ.}

\setlength{\csromangathapagewidth}{0pt}
\setlength{\csromangathawakwidth}{0pt}
\csromangathameasuregroup{“กามนฺธา ชาลสญฺฉนฺนา, \\ ปมตฺตพนฺธนา พทฺธา,}{\csromangathabat{“กามนฺธา ชาลสญฺฉนฺนา,}{ตณฺหาฉทนฉาทิตา.} \\ \csromangathabat{ปมตฺตพนฺธนา พทฺธา,}{มจฺฉาว กุมินามุเข.} \\ อิทํ สํกิเลสภาคิยํ สุตฺตํ.}
\csromangathameasurewak{อิทํ สํกิเลสภาคิยํ สุตฺตํ.}
\csromangathasetleft{“กามนฺธา ชาลสญฺฉนฺนา, \\ ปมตฺตพนฺธนา พทฺธา,}
\csromangathagroup{\csromangathastanza{\csromangathabat{“กามนฺธา ชาลสญฺฉนฺนา,}{ตณฺหา\-ฉทน\-ฉาทิตา.} \\ \csromangathabat{ปมตฺตพนฺธนา พทฺธา,}{มจฺฉาว กุมินามุเข.}}\csromangathastanza{\csromangathawak{อิทํ สํกิเลส\-ภาคิยํ สุตฺตํ.}}}

\prose{ชรามรณมเนฺวนฺติ, วจฺโฉ ขีรปโกว มาตรนฺ”ติ\csromansharedfootnote{8c1256d3c50b9e9c}{เหฏฺฐา 32 ปิฏฺเฐ; อุปริ 184, 196 ปิฏฺเฐสุปิ.}.}

\proseitem{90}{จตฺตาริมานิ ภิกฺขเว อคติคมนานิ. กตมานิ จตฺตาริ, ฉนฺทาคติํ\csromansharedfootnote{2aed9cc3abdef971}{ฉนฺทา อคติํ (สี, ก)} คจฺฉติ, โทสาคติํ คจฺฉติ, โมหาคติํ คจฺฉติ, ภยาคติํ คจฺฉติ. อิมานิ โข ภิกฺขเว จตฺตาริ อคติคมนานิ. อิทมโวจ ภควา, อิทํ วตฺวาน สุคโต, อถาปรํ เอตทโวจ สตฺถา\csromandash{}}

\prose{\csromansymbolfootnote{*}{อุปริ 212 ปิฏฺเฐปิ.}“ฉนฺทา โทสา ภยา โมหา, โย ธมฺมํ อติวตฺตติ.}

\setlength{\csromangathapagewidth}{0pt}
\setlength{\csromangathawakwidth}{0pt}
\csromangathameasuregroup{}{อิทํ สํกิเลสภาคิยํ สุตฺตํ. \\ “มโนปุพฺพงฺคมา ธมฺมา, \\ มโนเสฏฺฐา มโนมยา.}
\csromangathameasurewak{อิทํ สํกิเลสภาคิยํ สุตฺตํ. \\ “มโนปุพฺพงฺคมา ธมฺมา, \\ มโนเสฏฺฐา มโนมยา.}
\setlength{\csromangathaleftcol}{0pt}
\csromangathagroup{\csromangathastanza{\csromangathawak{อิทํ สํกิเลส\-ภาคิยํ สุตฺตํ. \\ “มโนปุพฺพงฺคมา ธมฺมา, \\ มโนเสฏฺฐา มโนมยา.}}}

\prose{นิหียติ ตสฺส ยโส, กาฬปกฺเขว จนฺทิมา”ติ\csromansharedfootnote{6b71b4ee81ab5e3f}{อํ 1. 325 ปิฏฺเฐ; อุปริ 140 ปิฏฺเฐปิ เปยฺยาลมุเขน.}.}

\setlength{\csromangathapagewidth}{0pt}
\setlength{\csromangathawakwidth}{0pt}
\csromangathameasuregroup{มนสา เจ ปุทุฏฺเฐน,}{\csromangathabat{มนสา เจ ปุทุฏฺเฐน,}{ภาสติ วา กโรติ วา.} \\ อิทํ สํกิเลสภาคิยํ สุตฺตํ.}
\csromangathameasurewak{อิทํ สํกิเลสภาคิยํ สุตฺตํ.}
\csromangathasetleft{มนสา เจ ปุทุฏฺเฐน,}
\csromangathagroup{\csromangathastanza{\csromangathabat{มนสา เจ ปุทุฏฺเฐน,}{ภาสติ วา กโรติ วา.}}\csromangathastanza{\csromangathawak{อิทํ สํกิเลส\-ภาคิยํ สุตฺตํ.}}}

\prose{ตโต นํ ทุกฺขมเนฺวติ, จกฺกํว วหโต ปทนฺ”ติ\csromansharedfootnote{4bd8970f16d89103}{ขุ 1. 13 ปิฏฺเฐ.}.}

\setlength{\csromangathapagewidth}{0pt}
\setlength{\csromangathawakwidth}{0pt}
\csromangathameasuregroup{}{“มิทฺธี ยทา โหติ มหคฺฆโส จ, \\ นิทฺทายิตา สมฺปริวตฺตสายี. \\ มหาวราโหว นิวาปปุฏฺโฐ, \\ อิทํ สํกิเลสภาคิยํ สุตฺตํ.}
\csromangathameasurewak{“มิทฺธี ยทา โหติ มหคฺฆโส จ, \\ นิทฺทายิตา สมฺปริวตฺตสายี. \\ มหาวราโหว นิวาปปุฏฺโฐ, \\ อิทํ สํกิเลสภาคิยํ สุตฺตํ.}
\setlength{\csromangathaleftcol}{0pt}
\csromangathagroup{\csromangathastanza{\csromangathawak{\csromanfolio{108}“มิทฺธี ยทา โหติ มหคฺฆโส จ, \\ นิทฺทายิตา สมฺปริวตฺต\-สายี. \\ มหาวราโหว นิวาปปุฏฺโฐ,}}\csromangathastanza{\csromangathawak{อิทํ สํกิเลส\-ภาคิยํ สุตฺตํ.}}}

\prose{ปุนปฺปุนํ คพฺภมุเปติ มนฺโท”ติ\csromansharedfootnote{23bf4b0bb8e8a91a}{เหฏฺฐา 30 ปิฏฺเฐปิ. * ขุ 1. 48 ปิฏฺเฐ; อุปริ 173, 201 ปิฏฺเฐสุปิ.}.}

\setlength{\csromangathapagewidth}{0pt}
\setlength{\csromangathawakwidth}{0pt}
\csromangathameasuregroup{\textsuperscript{*}“อยสาว มลํ สมุฏฺฐิตํ, \\ เอวํ อติโธนจารินํ, \\ อตฺตโน สุขเมสาโน,}{\csromangathabat{\textsuperscript{*}“อยสาว มลํ สมุฏฺฐิตํ,}{ตตุฏฺฐาย\textsuperscript{1} ตเมว ขาทติ.} \\ \csromangathabat{เอวํ อติโธนจารินํ,}{สานิ\textsuperscript{1} กมฺมานิ นยนฺติ ทุคฺคตินฺ”ติ.} \\ อิทํ สํกิเลสภาคิยํ สุตฺตํ. \\ \textsuperscript{+}“โจโร ยถา สนฺธิมุเข คหีโต, \\ สกมฺมุนา หญฺญติ พชฺฌเต จ. \\ เอวํ อยํ เปจฺจ ปชา ปรตฺถ, \\ สกมฺมุนา หญฺญติ พชฺฌเต จา”ติ. \\ อิทํ สํกิเลสภาคิยํ สุตฺตํ. \\ \csromangathabat{อตฺตโน สุขเมสาโน,}{เปจฺจ โส น ลภเต\textsuperscript{1} สุขนฺ”ติ.} \\ อิทํ สํกิเลสภาคิยํ สุตฺตํ.}
\csromangathameasurewak{อิทํ สํกิเลสภาคิยํ สุตฺตํ. \\ \textsuperscript{+}“โจโร ยถา สนฺธิมุเข คหีโต, \\ สกมฺมุนา หญฺญติ พชฺฌเต จ. \\ เอวํ อยํ เปจฺจ ปชา ปรตฺถ, \\ สกมฺมุนา หญฺญติ พชฺฌเต จา”ติ. \\ อิทํ สํกิเลสภาคิยํ สุตฺตํ.}
\csromangathasetleft{\textsuperscript{*}“อยสาว มลํ สมุฏฺฐิตํ, \\ เอวํ อติโธนจารินํ, \\ อตฺตโน สุขเมสาโน,}
\csromangathagroup{\csromangathastanza{\csromangathabat{\csromansymbolmark{*}“อยสาว มลํ สมุฏฺฐิตํ,}{ตตุฏฺฐาย\csromansharedfootnote{338dd9d64fd7ee1e}{[มิสฺสินฺคฺ โนเต 0]} ตเมว ขาทติ.} \\ \csromangathabat{เอวํ อติโธนจารินํ,}{สานิ\csromansharedfootnote{338dd9d64fd7ee1e}{[มิสฺสินฺคฺ โนเต 0]} กมฺมานิ นยนฺติ ทุคฺคตินฺ”ติ.}}\csromangathastanza{\csromangathawak{อิทํ สํกิเลส\-ภาคิยํ สุตฺตํ.}}\csromangathastanza{\csromangathawak{\csromansymbolmark{+}“โจโร ยถา สนฺธิมุเข คหีโต, \\ สกมฺมุนา หญฺญติ พชฺฌเต จ. \\ เอวํ อยํ เปจฺจ ปชา ปรตฺถ, \\ สกมฺมุนา หญฺญติ พชฺฌเต จา”ติ.}}\csromangathastanza{\csromangathawak{อิทํ สํกิเลส\-ภาคิยํ สุตฺตํ.}}\csromangathastanza{\csromangathabat{อตฺตโน สุขเมสาโน,}{เปจฺจ โส น ลภเต\csromansharedfootnote{338dd9d64fd7ee1e}{[มิสฺสินฺคฺ โนเต 0]} สุขนฺ”ติ.} \\ อิทํ สํกิเลส\-ภาคิยํ สุตฺตํ.}}

\prose{** “สุขกามานิ ภูตานิ, โย ทณฺเฑน วิหิํสติ.}

\prose{++ “คุนฺนํ เจ ตรมานานํ,}

\setlength{\csromangathapagewidth}{0pt}
\setlength{\csromangathawakwidth}{0pt}
\csromangathameasuregroup{เอวเมว มนุสฺเสสุ, \\ โส เจ อธมฺมํ จรติ, \\ สพฺพํ รฏฺฐํ ทุกฺขํ เสติ, \\ คจฺฉนฺติ เต พหุชนสนฺนิวาสํ,}{ชิมฺหํ คจฺฉติ ปุงฺคโว. \\ สพฺพา ตา ชิมฺหํ คจฺฉนฺติ, \\ เนตฺเต ชิมฺหํ คเต\textsuperscript{1} สติ. \\ \csromangathabat{เอวเมว มนุสฺเสสุ,}{โย โหติ เสฏฺฐสมฺมโต.} \\ \csromangathabat{โส เจ อธมฺมํ จรติ,}{ปเคว อิตรา ปชา.} \\ \csromangathabat{สพฺพํ รฏฺฐํ ทุกฺขํ เสติ,}{ราชา เจ โหติ อธมฺมิโก”ติ.} \\ อิทํ สํกิเลสภาคิยํ สุตฺตํ. \\ “สุกิจฺฉรูปาวติเม มนุสฺสา, \\ กโรนฺติ ปาปํ อุปธีสุ รตฺตา. \\ \csromangathabat{คจฺฉนฺติ เต พหุชนสนฺนิวาสํ,}{นิรยํ อวีจิํ กฏุกํ ภยานกนฺ”ติ.} \\ อิทํ สํกิเลสภาคิยํ สุตฺตํ. \\ “ผลํ เว กทลิํ หนฺติ, \\ ผลํ เวฬุํ ผลํ นฬํ. \\ อิทํ สํกิเลสภาคิยํ สุตฺตํ. \\ \textsuperscript{*}“โกธมกฺขครุ ภิกฺขุ, \\ ลาภสกฺการคารโว\textsuperscript{1}. \\ สุเขตฺเต ปูติพีชํว, \\ สทฺธมฺเม น วิรูหตี”ติ. \\ อิทํ สํกิเลสภาคิยํ สุตฺตํ.}
\csromangathameasurewak{ชิมฺหํ คจฺฉติ ปุงฺคโว. \\ สพฺพา ตา ชิมฺหํ คจฺฉนฺติ, \\ เนตฺเต ชิมฺหํ คเต\textsuperscript{1} สติ. \\ กโรนฺติ ปาปํ อุปธีสุ รตฺตา. \\ ผลํ เวฬุํ ผลํ นฬํ. \\ อิทํ สํกิเลสภาคิยํ สุตฺตํ. \\ \textsuperscript{*}“โกธมกฺขครุ ภิกฺขุ, \\ ลาภสกฺการคารโว\textsuperscript{1}. \\ สุเขตฺเต ปูติพีชํว, \\ สทฺธมฺเม น วิรูหตี”ติ. \\ อิทํ สํกิเลสภาคิยํ สุตฺตํ.}
\csromangathasetleft{เอวเมว มนุสฺเสสุ, \\ โส เจ อธมฺมํ จรติ, \\ สพฺพํ รฏฺฐํ ทุกฺขํ เสติ, \\ คจฺฉนฺติ เต พหุชนสนฺนิวาสํ,}
\csromangathagroup{\csromangathastanza{\csromangathawak{ชิมฺหํ คจฺฉติ ปุงฺคโว. \\ สพฺพา ตา ชิมฺหํ คจฺฉนฺติ, \\ เนตฺเต ชิมฺหํ คเต\csromansharedfootnote{bcf0c8dd6cfed689}{ชิมฺหคเต (ก)} สติ.}}\csromangathastanza{\csromangathabat{เอวเมว มนุสฺเสสุ,}{โย โหติ เสฏฺฐสมฺมโต.} \\ \csromangathabat{โส เจ อธมฺมํ จรติ,}{ปเคว อิตรา ปชา.}}\csromangathastanza{\csromangathabat{สพฺพํ รฏฺฐํ ทุกฺขํ เสติ,}{ราชา เจ โหติ อธมฺมิโก”ติ.} \\ อิทํ สํกิเลส\-ภาคิยํ สุตฺตํ. \\ “สุกิจฺฉรูปาวติเม มนุสฺสา,}\csromangathastanza{\csromangathawak{กโรนฺติ ปาปํ อุปธีสุ รตฺตา.}}\csromangathastanza{\csromangathabat{คจฺฉนฺติ เต พหุชนสนฺนิวาสํ,}{นิรยํ อวีจิํ กฏุกํ ภยานกนฺ”ติ.} \\ อิทํ สํกิเลส\-ภาคิยํ สุตฺตํ. \\ “ผลํ เว กทลิํ หนฺติ,}\csromangathastanza{\csromangathawak{ผลํ เวฬุํ ผลํ นฬํ.}}\csromangathastanza{\csromangathawak{\csromanfolio{109}อิทํ สํกิเลส\-ภาคิยํ สุตฺตํ. \\ \csromansymbolmark{*}“โกธมกฺขครุ ภิกฺขุ, \\ ลาภสกฺการคารโว\csromansharedfootnote{338dd9d64fd7ee1e}{[มิสฺสินฺคฺ โนเต 0]}. \\ สุเขตฺเต ปูติพีชํว, \\ สทฺธมฺเม น วิรูหตี”ติ. \\ อิทํ สํกิเลส\-ภาคิยํ สุตฺตํ.}}}

\prose{สกฺกาโร กาปุริสํ หนฺติ, คพฺโภ อสฺสตริํ ยถา”ติ\csromansharedfootnote{d00e7c338853800b}{อํ 1. 385 ปิฏฺเฐ.}.}

\proseitem{91}{“อิธาหํ ภิกฺขเว เอกจฺจํ ปุคฺคลํ ปทุฏฺฐจิตฺตํ เอวํ เจตสา เจโต ปริจฺจ ปชานามิ, (ยถา โข อยํ ปุคฺคโล อิริยติ, ยญฺจ ปฏิปทํ ปฏิปนฺโน, ยญฺจ มคฺคํ สมารูโฬฺห)\csromansharedfootnote{2a52f91ed03f44a0}{( ) นตฺถิ อํ 1. 8 ปิฏฺเฐ; ขุ 1. 203 ปิฏฺเฐ จ.} อิมมฺหิ จายํ สมเย กาลํ กเรยฺย, ยถาภตํ นิกฺขิตฺโต, เอวํ นิรเย. ตํ กิสฺส เหตุ, จิตฺตํ หิสฺส ภิกฺขเว ปทุฏฺฐํ\csromansharedfootnote{84f6dfedfdfc6346}{ปโทสิตํ (สี, ก) อํ 1. 8 ปิฏฺเฐ; ขุ 1. 203 ปิฏฺเฐ จ ปสฺสิตพฺพํ.}, เจโตปโทสเหตุ\csromansharedfootnote{a6665fb6f9b82ee9}{จิตฺตปโทสเหตุ (สี, ก)} โข ปน ภิกฺขเว เอวมิเธกจฺเจ สตฺตา กายสฺส เภทา ปรํ มรณา อปายํ ทุคฺคติํ วินิปาตํ นิรยํ อุปปชฺชนฺตี”ติ. เอตมตฺถํ ภควา อโวจ, ตตฺเถตํ อิติ วุจฺจติ\csromandash{}}

\setlength{\csromangathapagewidth}{0pt}
\setlength{\csromangathawakwidth}{0pt}
\csromangathameasuregroup{“ปทุฏฺฐจิตฺตํ ญตฺวาน, \\ เอตมตฺถญฺจ พฺยากาสิ, \\ อิมมฺหิ จายํ สมเย, \\ นิรยํ อุปปชฺเชยฺย, \\ เจโตปโทสเหตุ หิ, \\ ยถาภตํ นิกฺขิเปยฺย, \\ กายสฺส เภทา ทุปฺปญฺโญ, \\ \textsuperscript{+}“สเจ ภายถ ทุกฺขสฺส, \\ มากตฺถ ปาปกํ กมฺมํ, \\ \textsuperscript{*}สเจ จ ปาปกํ กมฺมํ, \\ น โว ทุกฺขา ปมุตฺยตฺถิ, \\ มเมติ พาลา มญฺญนฺติ, \\ อนฺตรายา สุ ภวิสฺสนฺติ, \\ มตา สคฺคํ น คจฺฉนฺติ, \\ กถํ วิวฏฺฏเต ธมฺมา, \\ โลภา ขณติ อตฺตานํ, \\ โลภา วิวฏฺฏเต ธมฺมา, \\ กโรนฺตา ปาปกํ กมฺมํ, \\ น ตํ กมฺมํ กตํ สาธุ, \\ ยสฺส อสฺสุมุโข โรทํ,}{\csromangathabat{“ปทุฏฺฐจิตฺตํ ญตฺวาน,}{เอกจฺจํ อิธ ปุคฺคลํ.} \\ \csromangathabat{เอตมตฺถญฺจ พฺยากาสิ,}{พุทฺโธ\textsuperscript{1} ภิกฺขูน สนฺติเก.} \\ \csromangathabat{อิมมฺหิ จายํ สมเย,}{กาลํ กยิราถ ปุคฺคโล.} \\ \csromangathabat{นิรยํ อุปปชฺเชยฺย,}{จิตฺตํ หิสฺส ปทูสิตํ.} \\ \csromangathabat{เจโตปโทสเหตุ หิ,}{สตฺตา คจฺฉนฺติ ทุคฺคติํ.} \\ \csromangathabat{ยถาภตํ นิกฺขิเปยฺย,}{เอวเมว ตถาวิโธ.} \\ \csromangathabat{กายสฺส เภทา ทุปฺปญฺโญ,}{นิรยํ โส’ปปชฺชตี”ติ.} \\ อยมฺปิ อตฺโถ วุตฺโต ภควตา อิติ เม สุตนฺติ. \\ อิทํ สํกิเลสภาคิยํ สุตฺตํ. \\ \csromangathabat{\textsuperscript{+}“สเจ ภายถ ทุกฺขสฺส,}{สเจ โว ทุกฺขมปฺปิยํ.} \\ \csromangathabat{มากตฺถ ปาปกํ กมฺมํ,}{อาวิ\textsuperscript{1} วา ยทิ วา รโห.} \\ \csromangathabat{\textsuperscript{*}สเจ จ ปาปกํ กมฺมํ,}{กริสฺสถ กโรถ วา.} \\ \csromangathabat{น โว ทุกฺขา ปมุตฺยตฺถิ,}{อุเปจฺจปิ ปลายตนฺ”ติ.} \\ อิทํ สํกิเลสภาคิยํ สุตฺตํ. \\ “อธมฺเมน ธนํ ลทฺธา, \\ มุสาวาเทน จูภยํ. \\ \csromangathabat{มเมติ พาลา มญฺญนฺติ,}{ตํ กถํ นุ ภวิสฺสติ.} \\ \csromangathabat{อนฺตรายา สุ ภวิสฺสนฺติ,}{สมฺภต’สฺส วินสฺสติ.} \\ \csromangathabat{มตา สคฺคํ น คจฺฉนฺติ,}{นนุ เอตฺตาวตา หตา”ติ.} \\ อิทํ สํกิเลสภาคิยํ สุตฺตํ. \\ “กถํ ขณติ อตฺตานํ, \\ กถํ มิตฺเตหิ ชีรติ. \\ \csromangathabat{กถํ วิวฏฺฏเต ธมฺมา,}{กถํ สคฺคํ น คจฺฉติ.} \\ \csromangathabat{โลภา ขณติ อตฺตานํ,}{ลุทฺโธ มิตฺเตหิ ชีรติ.} \\ \csromangathabat{โลภา วิวฏฺฏเต ธมฺมา,}{โลภา สคฺคํ น คจฺฉตี”ติ.} \\ อิทํ สํกิเลสภาคิยํ สุตฺตํ. \\ \textsuperscript{+}“จรนฺติ พาลา ทุมฺเมธา, \\ อมิตฺเตเนว อตฺตนา. \\ \csromangathabat{กโรนฺตา ปาปกํ กมฺมํ,}{ยํ โหติ กฏุกปฺผลํ\textsuperscript{1}.} \\ \csromangathabat{น ตํ กมฺมํ กตํ สาธุ,}{ยํ กตฺวา อนุตปฺปติ.} \\ \csromangathabat{ยสฺส อสฺสุมุโข โรทํ,}{วิปากํ ปฏิเสวตี”ติ.}}
\csromangathasetleft{“ปทุฏฺฐจิตฺตํ ญตฺวาน, \\ เอตมตฺถญฺจ พฺยากาสิ, \\ อิมมฺหิ จายํ สมเย, \\ นิรยํ อุปปชฺเชยฺย, \\ เจโตปโทสเหตุ หิ, \\ ยถาภตํ นิกฺขิเปยฺย, \\ กายสฺส เภทา ทุปฺปญฺโญ, \\ \textsuperscript{+}“สเจ ภายถ ทุกฺขสฺส, \\ มากตฺถ ปาปกํ กมฺมํ, \\ \textsuperscript{*}สเจ จ ปาปกํ กมฺมํ, \\ น โว ทุกฺขา ปมุตฺยตฺถิ, \\ มเมติ พาลา มญฺญนฺติ, \\ อนฺตรายา สุ ภวิสฺสนฺติ, \\ มตา สคฺคํ น คจฺฉนฺติ, \\ กถํ วิวฏฺฏเต ธมฺมา, \\ โลภา ขณติ อตฺตานํ, \\ โลภา วิวฏฺฏเต ธมฺมา, \\ กโรนฺตา ปาปกํ กมฺมํ, \\ น ตํ กมฺมํ กตํ สาธุ, \\ ยสฺส อสฺสุมุโข โรทํ,}
\csromangathagroup{\csromangathastanza{\csromangathabat{“ปทุฏฺฐจิตฺตํ ญตฺวาน,}{เอกจฺจํ อิธ ปุคฺคลํ.} \\ \csromangathabat{เอตมตฺถญฺจ พฺยากาสิ,}{พุทฺโธ\csromansharedfootnote{8fe3171211a0ea00}{สตฺถา (สี, ก)} ภิกฺขูน สนฺติเก.}}\csromangathastanza{\csromangathabat{อิมมฺหิ จายํ สมเย,}{กาลํ กยิราถ ปุคฺคโล.} \\ \csromangathabat{นิรยํ อุปปชฺเชยฺย,}{จิตฺตํ หิสฺส ปทูสิตํ.}}\csromangathastanza{\csromangathabat{เจโตปโทสเหตุ หิ,}{สตฺตา คจฺฉนฺติ ทุคฺคติํ.} \\ \csromangathabat{ยถาภตํ นิกฺขิเปยฺย,}{เอวเมว ตถาวิโธ.}}\csromangathastanza{\csromangathabat{กายสฺส เภทา ทุปฺปญฺโญ,}{นิรยํ โส’ปปชฺชตี”ติ.} \\ อยมฺปิ อตฺโถ วุตฺโต ภควตา อิติ เม สุตนฺติ. \\ อิทํ สํกิเลส\-ภาคิยํ สุตฺตํ. \\ \csromangathabat{\csromansymbolmark{+}“สเจ ภายถ ทุกฺขสฺส,}{สเจ โว ทุกฺขมปฺปิยํ.}}\csromangathastanza{\csromangathabat{มากตฺถ ปาปกํ กมฺมํ,}{อาวิ\csromansharedfootnote{53ea7028c38c929b}{อาวี (สี)} วา ยทิ วา รโห.} \\ \csromangathabat{\csromansymbolfootnote{*}{ขุ 1. 137 ปิฏฺเฐ.}สเจ จ ปาปกํ กมฺมํ,}{กริสฺสถ กโรถ วา.}}\csromangathastanza{\csromanfolio{110}\csromangathabat{น โว ทุกฺขา ปมุตฺยตฺถิ,}{อุเปจฺจปิ ปลายตนฺ”ติ.} \\ อิทํ สํกิเลส\-ภาคิยํ สุตฺตํ. \\ “อธมฺเมน ธนํ ลทฺธา,}\csromangathastanza{มุสาวาเทน จูภยํ.}\csromangathastanza{\csromangathabat{มเมติ พาลา มญฺญนฺติ,}{ตํ กถํ นุ ภวิสฺสติ.} \\ \csromangathabat{อนฺตรายา สุ ภวิสฺสนฺติ,}{สมฺภต’สฺส วินสฺสติ.}}\csromangathastanza{\csromangathabat{มตา สคฺคํ น คจฺฉนฺติ,}{นนุ เอตฺตาวตา หตา”ติ.} \\ อิทํ สํกิเลส\-ภาคิยํ สุตฺตํ. \\ “กถํ ขณติ อตฺตานํ,}\csromangathastanza{กถํ มิตฺเตหิ ชีรติ.}\csromangathastanza{\csromangathabat{กถํ วิวฏฺฏเต ธมฺมา,}{กถํ สคฺคํ น คจฺฉติ.} \\ \csromangathabat{โลภา ขณติ อตฺตานํ,}{ลุทฺโธ มิตฺเตหิ ชีรติ.}}\csromangathastanza{\csromangathabat{โลภา วิวฏฺฏเต ธมฺมา,}{โลภา สคฺคํ น คจฺฉตี”ติ.} \\ อิทํ สํกิเลส\-ภาคิยํ สุตฺตํ. \\ \csromansymbolmark{+}“จรนฺติ พาลา ทุมฺเมธา,}\csromangathastanza{อมิตฺเตเนว อตฺตนา.}\csromangathastanza{\csromangathabat{กโรนฺตา ปาปกํ กมฺมํ,}{ยํ โหติ กฏุกปฺผลํ\csromansharedfootnote{e133bdc91f1ab9f9}{กฏกํ ผลํ (ก)}.} \\ \csromangathabat{น ตํ กมฺมํ กตํ สาธุ,}{ยํ กตฺวา อนุตปฺปติ.} \\ \csromangathabat{ยสฺส อสฺสุมุโข โรทํ,}{วิปากํ ปฏิเสวตี”ติ.}}}

\vspace{\tipitakaparskip}
\csromancenter{อิทํ สํกิเลส\-ภาคิยํ สุตฺตํ.}
\setlength{\csromangathapagewidth}{0pt}
\setlength{\csromangathawakwidth}{0pt}
\csromangathameasuregroup{“ทุกฺกรํ ทุตฺติติกฺขญฺจ, \\ พหู หิ ตตฺถ สมฺพาธา, \\ โย หิ อตฺถญฺจ ธมฺมญฺจ, \\ มนํ ปโทสเย พาโล, \\ \textsuperscript{*}“อปฺปเมยฺยํ ปมินนฺโต, \\ อปฺปเมยฺยํ ปมายินํ\textsuperscript{1}, \\ \textsuperscript{+}“ปุริสสฺส หิ ชาตสฺส, \\ ยาย ฉินฺทติ อตฺตานํ, \\ น หิ สตฺถํ สุนิสิตํ, \\ เอวํ วิรทฺธํ ปาเตติ,}{\csromangathabat{“ทุกฺกรํ ทุตฺติติกฺขญฺจ,}{อพฺยตฺเตน จ\textsuperscript{1} สามญฺญํ.} \\ \csromangathabat{พหู หิ ตตฺถ สมฺพาธา,}{ยตฺถ พาโล วิสีทติ\textsuperscript{1}.} \\ \csromangathabat{โย หิ อตฺถญฺจ ธมฺมญฺจ,}{ภาสมาเน ตถาคเต.} \\ \csromangathabat{มนํ ปโทสเย พาโล,}{โมฆํ โข ตสฺส ชีวิตํ.} \\ เอตญฺจาหํ อรหามิ, \\ ทุกฺขญฺจ อิโต จ ปาปิยตรํ ภนฺเต. \\ โย อปฺปเมยฺเยสุ ตถาคเตสุ, \\ จิตฺตํ ปโทเสมิ อวีตราโค”ติ. \\ อิทํ สํกิเลสภาคิยํ สุตฺตํ. \\ \csromangathabat{\textsuperscript{*}“อปฺปเมยฺยํ ปมินนฺโต,}{โกธ วิทฺวา วิกปฺปเย.} \\ \csromangathabat{อปฺปเมยฺยํ ปมายินํ\textsuperscript{1},}{นิวุตํ ตํ มญฺเญ อกิสฺสวนฺ”ติ.} \\ อิทํ สํกิเลสภาคิยํ สุตฺตํ. \\ \csromangathabat{\textsuperscript{+}“ปุริสสฺส หิ ชาตสฺส,}{กุฐารี\textsuperscript{1} ชายเต มุเข.} \\ \csromangathabat{ยาย ฉินฺทติ อตฺตานํ,}{พาโล ทุพฺภาสิตํ ภณํ.} \\ \csromangathabat{น หิ สตฺถํ สุนิสิตํ,}{วิสํ หลาหลํ อิว.} \\ \csromangathabat{เอวํ วิรทฺธํ ปาเตติ,}{วาจา ทุพฺภาสิตา ยถา”ติ.} \\ อิทํ สํกิเลสภาคิยํ สุตฺตํ.}
\csromangathameasurewak{เอตญฺจาหํ อรหามิ, \\ ทุกฺขญฺจ อิโต จ ปาปิยตรํ ภนฺเต. \\ โย อปฺปเมยฺเยสุ ตถาคเตสุ, \\ จิตฺตํ ปโทเสมิ อวีตราโค”ติ. \\ อิทํ สํกิเลสภาคิยํ สุตฺตํ. \\ อิทํ สํกิเลสภาคิยํ สุตฺตํ. \\ อิทํ สํกิเลสภาคิยํ สุตฺตํ.}
\csromangathasetleft{“ทุกฺกรํ ทุตฺติติกฺขญฺจ, \\ พหู หิ ตตฺถ สมฺพาธา, \\ โย หิ อตฺถญฺจ ธมฺมญฺจ, \\ มนํ ปโทสเย พาโล, \\ \textsuperscript{*}“อปฺปเมยฺยํ ปมินนฺโต, \\ อปฺปเมยฺยํ ปมายินํ\textsuperscript{1}, \\ \textsuperscript{+}“ปุริสสฺส หิ ชาตสฺส, \\ ยาย ฉินฺทติ อตฺตานํ, \\ น หิ สตฺถํ สุนิสิตํ, \\ เอวํ วิรทฺธํ ปาเตติ,}
\csromangathagroup{\csromangathastanza{\csromangathabat{“ทุกฺกรํ ทุตฺติติกฺขญฺจ,}{อพฺยตฺเตน จ\csromansharedfootnote{c162c8ac730aea0e}{อวิยตฺเตน (สี, ก)} สามญฺญํ.} \\ \csromangathabat{พหู หิ ตตฺถ สมฺพาธา,}{ยตฺถ พาโล วิสีทติ\csromansharedfootnote{c162c8ac730aea0e}{อวิยตฺเตน (สี, ก)}.}}\csromangathastanza{\csromangathabat{โย หิ อตฺถญฺจ ธมฺมญฺจ,}{ภาสมาเน ตถาคเต.} \\ \csromangathabat{มนํ ปโทสเย พาโล,}{โมฆํ โข ตสฺส ชีวิตํ.}}\csromangathastanza{\csromangathawak{เอตญฺจาหํ อรหามิ, \\ ทุกฺขญฺจ อิโต จ ปาปิยตรํ ภนฺเต. \\ โย อปฺปเมยฺเยสุ ตถาคเตสุ, \\ จิตฺตํ ปโทเสมิ อวีตราโค”ติ.}}\csromangathastanza{\csromangathawak{อิทํ สํกิเลส\-ภาคิยํ สุตฺตํ.}}\csromangathastanza{\csromanfolio{111}\csromangathabat{\csromansymbolmark{*}“อปฺปเมยฺยํ ปมินนฺโต,}{โกธ วิทฺวา วิกปฺปเย.} \\ \csromangathabat{อปฺปเมยฺยํ ปมายินํ\csromansharedfootnote{338dd9d64fd7ee1e}{[มิสฺสินฺคฺ โนเต 0]},}{นิวุตํ ตํ มญฺเญ อกิสฺสวนฺ”ติ.}}\csromangathastanza{\csromangathawak{อิทํ สํกิเลส\-ภาคิยํ สุตฺตํ.}}\csromangathastanza{\csromangathabat{\csromansymbolfootnote{+}{สํ 1. 151; อํ 3. 396; ขุ 1. 381 ปิฏฺเฐสุ.}“ปุริสสฺส หิ ชาตสฺส,}{กุฐารี\csromansharedfootnote{e11bd696de7cb918}{กุธารี (ก) ** สํ 1. 151; อํ 1. 309, 310; อํ 3. 396; ขุ 1. 381, 2 ปิฏฺเฐสุ.} ชายเต มุเข.} \\ \csromangathabat{ยาย ฉินฺทติ อตฺตานํ,}{พาโล ทุพฺภาสิตํ ภณํ.}}\csromangathastanza{\csromangathabat{น หิ สตฺถํ สุนิสิตํ,}{วิสํ หลาหลํ อิว.} \\ \csromangathabat{เอวํ วิรทฺธํ ปาเตติ,}{วาจา ทุพฺภาสิตา ยถา”ติ.}}\csromangathastanza{\csromangathawak{อิทํ สํกิเลส\-ภาคิยํ สุตฺตํ.}}}

\proseitem{92}{\csromansymbolfootnote{*}{สํ 1. 151 ปิฏฺเฐ.}* “โย นินฺทิยํ ปสํสติ,}

\prose{ตํ วา นินฺทติ โย ปสํสิโย.}

\setlength{\csromangathapagewidth}{0pt}
\setlength{\csromangathawakwidth}{0pt}
\csromangathameasuregroup{สพฺพสฺสาปิ สหาปิ อตฺตนา,}{วิจินาติ มุเขน โส กลิํ, \\ กลินา เตน สุขํ น วินฺทติ. \\ ** อปฺปมตฺโต อยํ กลิ, \\ โย อกฺเขสุ ธนปราชโย. \\ \csromangathabat{สพฺพสฺสาปิ สหาปิ อตฺตนา,}{อยเมว มหนฺตตโร\textsuperscript{1} กลิ.} \\ โย สุคเตสุ มนํ ปโทสเย. \\ ** สตํ สหสฺสานํ นิรพฺพุทานํ, \\ ฉตฺติํสตี ปญฺจ จ อพฺพุทานิ. \\ ยมริยครหี นิรยํ อุเปติ, \\ วาจํ มนญฺจ ปณิธาย ปาปกนฺ”ติ. \\ อิทํ สํกิเลสภาคิยํ สุตฺตํ.}
\csromangathameasurewak{วิจินาติ มุเขน โส กลิํ, \\ กลินา เตน สุขํ น วินฺทติ. \\ ** อปฺปมตฺโต อยํ กลิ, \\ โย อกฺเขสุ ธนปราชโย. \\ ฉตฺติํสตี ปญฺจ จ อพฺพุทานิ. \\ ยมริยครหี นิรยํ อุเปติ, \\ วาจํ มนญฺจ ปณิธาย ปาปกนฺ”ติ. \\ อิทํ สํกิเลสภาคิยํ สุตฺตํ.}
\csromangathasetleft{สพฺพสฺสาปิ สหาปิ อตฺตนา,}
\csromangathagroup{\csromangathastanza{\csromangathawak{วิจินาติ มุเขน โส กลิํ, \\ กลินา เตน สุขํ น วินฺทติ. \\ ** อปฺปมตฺโต อยํ กลิ, \\ โย อกฺเขสุ ธนปราชโย.}}\csromangathastanza{\csromangathabat{สพฺพสฺสาปิ สหาปิ อตฺตนา,}{อยเมว มหนฺตตโร\csromansharedfootnote{a897b92ddee5bcb2}{มหตฺตโร (ก) ++ ขุ 1. 382 ปิฏฺเฐ.} กลิ.} \\ โย สุคเตสุ มนํ ปโทสเย. \\ ** สตํ สหสฺสานํ นิรพฺพุทานํ,}\csromangathastanza{\csromangathawak{ฉตฺติํสตี ปญฺจ จ อพฺพุทานิ. \\ ยมริยครหี นิรยํ อุเปติ, \\ วาจํ มนญฺจ ปณิธาย ปาปกนฺ”ติ. \\ อิทํ สํกิเลส\-ภาคิยํ สุตฺตํ.}}}

\prose{++ “โย โลภคุเณ อนุยุตฺโต,}

\setlength{\csromangathapagewidth}{0pt}
\setlength{\csromangathawakwidth}{0pt}
\csromangathameasuregroup{}{โส วจสา\textsuperscript{1} ปริภาสติ อญฺเญ. \\ อสฺสทฺโธ กทริโย\textsuperscript{1} อวทญฺญู, \\ มจฺฉริ เปสุณิยํ อนุยุตฺโต. \\ \textsuperscript{*}มุขทุคฺค วิภูต อนริย, \\ ภูนหุ ปาปก ทุกฺกฏการิ. \\ ปุริสนฺต กลี อวชาตปุตฺต\textsuperscript{1}, \\ มา พหุภาณิธ เนรยิโกสิ. \\ ** รชมากิรสี อหิตาย, \\ สนฺเต ครหสิ กิพฺพิสการี. \\ พหูนิ ทุจฺจริตานิ จริตฺวา, \\ คจฺฉสิ โข ปปตํ จิรรตฺตนฺ”ติ. \\ อิทํ สํกิเลสภาคิยํ สุตฺตํ. \\ ตตฺถ กตมํ วาสนาภาคิยํ สุตฺตํ. \\ \textsuperscript{+}“มโนปุพฺพงฺคมา ธมฺมา, \\ มโนเสฏฺฐา มโนมยา. \\ มนสา เจ ปสนฺเนน, \\ ภาสติ วา กโรติ วา. \\ ตโต นํ สุขมเนฺวติ, \\ อิทํ วาสนาภาคิยํ สุตฺตํ.}
\csromangathameasurewak{โส วจสา\textsuperscript{1} ปริภาสติ อญฺเญ. \\ อสฺสทฺโธ กทริโย\textsuperscript{1} อวทญฺญู, \\ มจฺฉริ เปสุณิยํ อนุยุตฺโต. \\ \textsuperscript{*}มุขทุคฺค วิภูต อนริย, \\ ภูนหุ ปาปก ทุกฺกฏการิ. \\ ปุริสนฺต กลี อวชาตปุตฺต\textsuperscript{1}, \\ มา พหุภาณิธ เนรยิโกสิ. \\ ** รชมากิรสี อหิตาย, \\ สนฺเต ครหสิ กิพฺพิสการี. \\ พหูนิ ทุจฺจริตานิ จริตฺวา, \\ คจฺฉสิ โข ปปตํ จิรรตฺตนฺ”ติ. \\ อิทํ สํกิเลสภาคิยํ สุตฺตํ. \\ ตตฺถ กตมํ วาสนาภาคิยํ สุตฺตํ. \\ \textsuperscript{+}“มโนปุพฺพงฺคมา ธมฺมา, \\ มโนเสฏฺฐา มโนมยา. \\ มนสา เจ ปสนฺเนน, \\ ภาสติ วา กโรติ วา. \\ ตโต นํ สุขมเนฺวติ, \\ อิทํ วาสนาภาคิยํ สุตฺตํ.}
\setlength{\csromangathaleftcol}{0pt}
\csromangathagroup{\csromangathastanza{\csromangathawak{โส วจสา\csromansharedfootnote{7c8efbd222de26ce}{วจสา จ (ก)} ปริภาสติ อญฺเญ. \\ อสฺสทฺโธ กทริโย\csromansharedfootnote{052b6c36dd31528a}{อนริโย (สี, ก)} อวทญฺญู, \\ มจฺฉริ เปสุณิยํ อนุยุตฺโต. \\ \csromansymbolmark{*}มุขทุคฺค วิภูต อนริย,}}\csromangathastanza{\csromangathawak{\csromanfolio{112}ภูนหุ ปาปก ทุกฺกฏการิ. \\ ปุริสนฺต กลี อวชาตปุตฺต\csromansharedfootnote{03e69b852f4e92e1}{อวชาตกปุตฺต (สี, ก) ** ขุ 1. 383 ปิฏฺเฐ.}, \\ มา พหุภาณิธ เนรยิโกสิ. \\ ** รชมากิรสี อหิตาย,}}\csromangathastanza{\csromangathawak{สนฺเต ครหสิ กิพฺพิสการี. \\ พหูนิ ทุจฺจริตานิ จริตฺวา, \\ คจฺฉสิ โข ปปตํ จิรรตฺตนฺ”ติ. \\ อิทํ สํกิเลส\-ภาคิยํ สุตฺตํ. \\ ตตฺถ กตมํ วาสนาภาคิยํ สุตฺตํ. \\ \csromansymbolmark{+}“มโนปุพฺพงฺคมา ธมฺมา,}}\csromangathastanza{\csromangathawak{มโนเสฏฺฐา มโนมยา. \\ มนสา เจ ปสนฺเนน, \\ ภาสติ วา กโรติ วา. \\ ตโต นํ สุขมเนฺวติ,}}\csromangathastanza{\csromangathawak{อิทํ วาสนาภาคิยํ สุตฺตํ.}}}

\prose{ฉายาว อนปายินี”ติ\csromansharedfootnote{78518dbcd57ff8d5}{อนุปายินีติ (ก) ++ สํ 3. 323 ปิฏฺเฐ.}.}

\proseitem{93}{\csromansymbolfootnote{+}{ขุ 1. 13 ปิฏฺเฐ.}+ “มหานาโม สกฺโก ภควนฺตํ เอตทโวจ\csromandash{}อิทํ ภนฺเต กปิลวตฺถุ อิทฺธญฺเจว ผีตญฺจ พาหุชญฺญํ\csromansharedfootnote{d41f71506521d039}{พหุชนํ (สี, ก)} อากิณฺณมนุสฺสํ สมฺพาธพฺยูตํ, โส โข อหํ ภนฺเต ภควนฺตํ วา ปยิรุปาสิตฺวา มโนภาวนีเย วา ภิกฺขู สายนฺหสมยํ กปิลวตฺถุํ ปวิสนฺโต ภนฺเตนปิ หตฺถินา สมาคจฺฉามิ, ภนฺเตนปิ อสฺเสน สมาคจฺฉามิ, ภนฺเตนปิ รเถน สมาคจฺฉามิ, ภนฺเตนปิ สกเฏน สมาคจฺฉามิ, ภนฺเตนปิ ปุริเสน สมาคจฺฉามิ, ตสฺส มยฺหํ ภนฺเต ตสฺมิํ สมเย มุสฺสเตว ภควนฺตํ อารพฺภ สติ, มุสฺสติ ธมฺมํ อารพฺภ สติ, มุสฺสติ สํฆํ อารพฺภ สติ. ตสฺส มยฺหํ ภนฺเต เอวํ โหติ “อิมมฺหิ จาหํ สายนฺหสมเย กาลํ กเรยฺยํ, กา มยฺหํ\csromansharedfootnote{e56d078e42d217a7}{มมสฺส (สี, ก)} คติ, โก อภิสมฺปราโย”ติ.}

\prose{มา ภายิ มหานาม, มา ภายิ มหานาม, อปาปกํ เต มรณํ ภวิสฺสติ, อปาปิกา\csromansharedfootnote{4f3b77235ff89279}{อปาปีกา เต (สี)} กาลํกิริยา. จตูหิ โข มหานาม ธมฺเมหิ \csromanfolio{113}สมนฺนาคโต อริยสาวโก นิพฺพานนินฺโน โหติ นิพฺพานโปโณ นิพฺพาน\-ปพฺภาโร. กตเมหิ จตูหิ, อิธ มหานาม อริยสาวโก พุทฺเธ อเวจฺจ\-ปฺปสาเทน สมนฺนาคโต โหติ, อิติปิ โส ภควา อรหํ ฯเปฯ พุทฺโธ ภควาติ. ธมฺเม ฯเปฯ. สํเฆ ฯเปฯ. อริยกนฺเตหิ สีเลหิ สมนฺนาคโต โหติ อขณฺเฑหิ ฯเปฯ สมาธิ\-สํวตฺตนิเกหิ. เสยฺยถาปิ มหานาม รุกฺโข ปาจีนนินฺโน ปาจีนโปโณ ปาจีนปพฺภาโร, โส มูลจฺฉินฺโน\csromansharedfootnote{cd6b511d19345d2f}{มูเลหิ ฉินฺโน (สี, ก)} กตเมน ปปเตยฺยาติ. เยน ภนฺเต นินฺโน เยน โปโณ เยน ปพฺภาโรติ. เอวเมว โข มหานาม อิเมหิ จตูหิ ธมฺเมหิ สมนฺนาคโต อริยสาวโก นิพฺพานนินฺโน โหติ นิพฺพานโปโณ นิพฺพาน\-ปพฺภาโร. มา ภายิ มหานาม, มา ภายิ มหานาม, อปาปกํ เต มรณํ ภวิสฺสติ, อปาปิกา กาลํกิริยา”ติ. อิทํ วาสนาภาคิยํ สุตฺตํ.}

\setlength{\csromangathapagewidth}{0pt}
\setlength{\csromangathawakwidth}{0pt}
\csromangathameasuregroup{\textsuperscript{*}“สุขกามานิ ภูตานิ, \\ อตฺตโน สุขเมสาโน,}{\csromangathabat{\textsuperscript{*}“สุขกามานิ ภูตานิ,}{โย ทณฺเฑน น หิํสติ.} \\ \csromangathabat{อตฺตโน สุขเมสาโน,}{เปจฺจ โส ลภเต สุขนฺ”ติ.} \\ อิทํ วาสนาภาคิยํ สุตฺตํ.}
\csromangathameasurewak{อิทํ วาสนาภาคิยํ สุตฺตํ.}
\csromangathasetleft{\textsuperscript{*}“สุขกามานิ ภูตานิ, \\ อตฺตโน สุขเมสาโน,}
\csromangathagroup{\csromangathastanza{\csromangathabat{\csromansymbolmark{*}“สุขกามานิ ภูตานิ,}{โย ทณฺเฑน น หิํสติ.} \\ \csromangathabat{อตฺตโน สุขเมสาโน,}{เปจฺจ โส ลภเต สุขนฺ”ติ.}}\csromangathastanza{\csromangathawak{อิทํ วาสนาภาคิยํ สุตฺตํ.}}}

\prose{** “คุนฺนญฺเจ ตรมานานํ,}

\setlength{\csromangathapagewidth}{0pt}
\setlength{\csromangathawakwidth}{0pt}
\csromangathameasuregroup{เอวเมว มนุสฺเสสุ, \\ โส สเจ\textsuperscript{1} ธมฺมํ จรติ, \\ สพฺพํ รฏฺฐํ สุขํ เสติ,}{อุชุํ คจฺฉติ ปุงฺคโว. \\ สพฺพา ตา อุชุํ คจฺฉนฺติ, \\ เนตฺเต อุชุํ คเต สติ. \\ \csromangathabat{เอวเมว มนุสฺเสสุ,}{โย โหติ เสฏฺฐสมฺมโต.} \\ \csromangathabat{โส สเจ\textsuperscript{1} ธมฺมํ จรติ,}{ปเคว อิตรา ปชา.} \\ \csromangathabat{สพฺพํ รฏฺฐํ สุขํ เสติ,}{ราชา เจ โหติ ธมฺมิโก”ติ.} \\ อิทํ วาสนาภาคิยํ สุตฺตํ.}
\csromangathameasurewak{อุชุํ คจฺฉติ ปุงฺคโว. \\ สพฺพา ตา อุชุํ คจฺฉนฺติ, \\ เนตฺเต อุชุํ คเต สติ.}
\csromangathasetleft{เอวเมว มนุสฺเสสุ, \\ โส สเจ\textsuperscript{1} ธมฺมํ จรติ, \\ สพฺพํ รฏฺฐํ สุขํ เสติ,}
\csromangathagroup{\csromangathastanza{\csromangathawak{อุชุํ คจฺฉติ ปุงฺคโว. \\ สพฺพา ตา อุชุํ คจฺฉนฺติ, \\ เนตฺเต อุชุํ คเต สติ.}}\csromangathastanza{\csromangathabat{เอวเมว มนุสฺเสสุ,}{โย โหติ เสฏฺฐสมฺมโต.} \\ \csromangathabat{โส สเจ\csromansharedfootnote{338dd9d64fd7ee1e}{[มิสฺสินฺคฺ โนเต 0]} ธมฺมํ จรติ,}{ปเคว อิตรา ปชา.}}\csromangathastanza{\csromangathabat{สพฺพํ รฏฺฐํ สุขํ เสติ,}{ราชา เจ โหติ ธมฺมิโก”ติ.} \\ อิทํ วาสนาภาคิยํ สุตฺตํ.}}

\proseitem{94}{ภควา สาวตฺถิยํ วิหรติ เชตวเน อนาถ\-ปิณฺฑิกสฺส อาราเม.\csromansymbolfootnote{+}{สํ 3. 303 ปิฏฺฐาทีสุ.}เตน โข ปน สมเยน สมฺพหุลา ภิกฺขู ภควโต จีวรกมฺมํ กโรนฺติ “นิฏฺฐิตจีวโร ภควา เตมาสจฺจเยน จาริกํ ปกฺกมิสฺสตี”ติ. เตน โข ปน สมเยน อิสิทตฺตปุราณา ถปตโย สาเกเต\csromansharedfootnote{2361fee5c3874d16}{สาธุเก (สํ 3. 303 ปิฏฺเฐ)} ปฏิวสนฺติ เกนจิ เทว กรณีเยน. อสฺโสสุํ โข อิสิทตฺตปุราณา ถปตโย “สมฺพหุลา กิร ภิกฺขู ภควโต จีวรกมฺมํ กโรนฺติ “นิฏฺฐิตจีวโร ภควา เตมาสจฺจเยน จาริกํ ปกฺกมิสฺสตี”ติ.}

\prose{\csromanfolio{114}อถ โข อิสิทตฺตปุราณา ถปตโย มคฺเค ปุริสํ ฐเปสุํ “ยทา ตฺวํ อมฺโภ ปุริส ปสฺเสยฺยาสิ ภควนฺตํ อาคจฺฉนฺตํ อรหนฺตํ สมฺมา\-สมฺพุทฺธํ, อถ อมฺหากํ อาโรเจยฺยาสี”ติ. ทฺวีหตีหํ ฐิโต โข โส ปุริโส อทฺทส ภควนฺตํ ทูรโตว อาคจฺฉนฺตํ, ทิสฺวาน เยน อิสิทตฺตปุราณา ถปตโย เตนุปสงฺกมิ, อุปสงฺกมิตฺวา อิสิทตฺตปุราเณ ถปตโย เอตทโวจ “อยํ โส ภนฺเต\csromansharedfootnote{b78825393138f395}{อยํ ภนฺเต (สี, ก)} ภควา อาคจฺฉติ อรหํ สมฺมาสมฺพุทฺโธ, ยสฺสทานิ กาลํ มญฺญถา”ติ.}

\prose{อถ โข อิสิทตฺตปุราณา ถปตโย เยน ภควา เตนุ\-ปสงฺกมิํสุ, อุปสงฺกมิตฺวา ภควนฺตํ อภิวาเทตฺวา ภควนฺตํ ปิฏฺฐิโต ปิฏฺฐิโต อนุพนฺธิํสุ. อถ โข ภควา มคฺคา โอกฺกมฺม เยน อญฺญตรํ รุกฺขมูลํ เตนุปสงฺกมิ, อุปสงฺกมิตฺวา ปญฺญตฺเต อาสเน นิสีทิ. อิสิทตฺตปุราณา ถปตโย ภควนฺตํ อภิวาเทตฺวา เอกมนฺตํ นิสีทิํสุ, เอกมนฺตํ นิสินฺนา โข อิสิทตฺตปุราณา ถปตโย ภควนฺตํ เอตทโวจุํ\csromandash{}}

\prose{ยทา มยํ ภนฺเต ภควนฺตํ สุโณม “สาวตฺถิยา โกสเลสุ จาริกํ ปกฺกมิสฺสตี”ติ, โหติ โน ตสฺมิํ สมเย อนตฺตมนตา โหติ โทมนสฺสํ “ทูเร โน ภควา ภวิสฺสตี”ติ. ยทา ปน มยํ ภนฺเต ภควนฺตํ สุโณม “สาวตฺถิยา โกสเลสุ จาริกํ ปกฺกนฺโต”ติ, โหติ โน ตสฺมิํ สมเย อนตฺตมนตา โหติ โทมนสฺสํ “ทูเร โน ภควา”ติ ฯเปฯ.}

\prose{ยทา ปน มยํ ภนฺเต ภควนฺตํ สุโณม “กาสีสุ มคเธสุ\csromansharedfootnote{737cbd828e1cbd40}{กาสีหิ มาคเธ (สํ 3. 304 ปิฏฺเฐ)} จาริกํ ปกฺกมิสฺสตี”ติ, โหติ โน ตสฺมิํ สมเย อนตฺตมนตา โหติ โทมนสฺสํ “ทูเร โน ภควา ภวิสฺสตี”ติ. ยทา ปน มยํ ภนฺเต ภควนฺตํ สุโณม “กาสีสุ มคเธสุ จาริกํ ปกฺกนฺโต”ติ, อนปฺปกา โน ตสฺมิํ สมเย อนตฺตมนตา โหติ อนปฺปกํ โทมนสฺสํ “ทูเร โน ภควา”ติ.}

\prose{ยทา ปน มยํ ภนฺเต ภควนฺตํ สุโณม “มคเธสุ กาสีสุ\csromansharedfootnote{69554af8e84ec1ec}{มคเธหิ กาสีสุ (สํ 3. 304 ปิฏฺเฐ.)} จาริกํ ปกฺกมิสฺสตี”ติ, โหติ โน ตสฺมิํ สมเย อตฺตมนตา โหติ โสมนสฺสํ “อาสนฺเน โน ภควา ภวิสฺสตี”ติ. ยทา ปน มยํ ภนฺเต ภควนฺตํ สุโณม “มคเธสุ กาสีสุ จาริกํ ปกฺกนฺโต”ติ, โหติ \csromanfolio{115}โน ตสฺมิํ สมเย อตฺตมนตา โหติ โสมนสฺสํ “อาสนฺเน โน ภควา”ติ ฯเปฯ.}

\prose{ยทา ปน มยํ ภนฺเต ภควนฺตํ สุโณม “โกสเลสุ สาวตฺถิํ\csromansharedfootnote{5c6fb3e5488adea7}{โกสเลสุ สาวตฺถิยํ (สี, ก), โกสเลหิ สาวตฺถิํ (สํ 3. 305 ปิฏฺเฐ.)} จาริกํ ปกฺกมิสฺสตี”ติ. โหติ โน ตสฺมิํ สมเย อตฺตมนตา โหติ โสมนสฺสํ “อาสนฺเน โน ภควา ภวิสฺสตี”ติ.}

\prose{ยทา ปน มยํ ภนฺเต ภควนฺตํ สุโณม “สาวตฺถิยํ วิหรติ เชตวเน อนาถ\-ปิณฺฑิกสฺส อาราเม”ติ. โหติ อนปฺปกา โน ตสฺมิํ สมเย อตฺตมนตา, โหติ อนปฺปกํ โสมนสฺสํ “อาสนฺเน โน ภควา”ติ.}

\prose{ตสฺมาติห ถปตโย สมฺพาโธ ฆราวาโส รชาปโถ, อพฺโภกาโส ปพฺพชฺชา, อลญฺจ ปน โว ถปตโย อปฺปมาทายาติ. อตฺถิ โข โน ภนฺเต เอตมฺหา สมฺพาธา อญฺโญ สมฺพาโธ สมฺพาธตโร เจว สมฺพาธ\-สงฺขาต\-ตโร จาติ. กตโม ปน โว ถปตโย เอตมฺหา สมฺพาธา อญฺโญ สมฺพาโธ สมฺพาธตโร เจว สมฺพาธ\-สงฺขาต\-ตโร จาติ.}

\prose{อิธ มยํ ภนฺเต ยทา ราชา ปเสนทิ โกสโล อุยฺยานภูมิํ นิยฺยาตุกาโม\csromansharedfootnote{107450b87f969bff}{คนฺตุกาโม (สี, ก)} โหติ, เย เต รญฺโญ ปเสนทิสฺส โกสลสฺส นาคา โอปวยฺหา, เต กปฺเปตฺวา ยา ตา รญฺโญ ปเสนทิสฺส โกสลสฺส ปชาปติโย ปิยา มนาปา, ตา\csromansharedfootnote{0bc25b2d1c60ffc8}{ตาสํ (สี, ก)} เอกํ ปุรโต เอกํ ปจฺฉโต นิสีทาเปม, ตาสํ โข ปน ภนฺเต ภคินีนํ เอวรูโป คนฺโธ โหติ. เสยฺยถาปิ นาม คนฺธ\-กรณฺฑกสฺส ตาวเทว วิวริย\-มานสฺส, ยถา ตํ ราชกญฺญานํ\csromansharedfootnote{72d13cba044d6ceb}{ราชารเหน (สี, ก)} คนฺเธน วิภูสิตานํ. ตาสํ โข ปน ภนฺเต ภคินีนํ เอวรูโป กายสมฺผสฺโส โหติ, เสยฺยถาปิ นาม ตูลปิจุโน วา กปฺปาสปิจุโน วา, ยถา ตํ ราชกญฺญานํ สุเขธิตานํ. ตสฺมิํ โข ปน ภนฺเต สมเย นาโคปิ รกฺขิตพฺโพ โหติ, ตาปิ ภคินิโย รกฺขิตพฺพา โหนฺติ, อตฺตาปิ รกฺขิตพฺโพ โหติ. น โข ปน มยํ ภนฺเต อภิชานาม ตาสุ ภคินีสุ ปาปกํ จิตฺตํ อุปฺปาเทนฺตา, อยํ โข โน ภนฺเต เอตมฺหา สมฺพาธา อญฺโญ สมฺพาโธ สมฺพาธตโร เจว สมฺพาธ\-สงฺขาต\-ตโร จาติ.}

\prose{\csromanfolio{116}ตสฺมาติห ถปตโย สมฺพาโธ ฆราวาโส รชาปโถ, อพฺโภกาโส ปพฺพชฺชา. อลญฺจ ปน โว ถปตโย อปฺปมาทาย. จตูหิ โข ถปตโย ธมฺเมหิ สมนฺนาคโต อริยสาวโก โสตาปนฺโน โหติ อวินิปาต\-ธมฺโม นิยโต สมฺโพธิ\-ปรายโณ.}

\prose{กตเมหิ จตูหิ, อิธ ถปตโย สุตวา อริยสาวโก พุทฺเธ อเวจฺจ\-ปฺปสาเทน สมนฺนาคโต โหติ อิติปิ โส ภควา อรหํ ฯเปฯ พุทฺโธ ภควาติ. ธมฺเม ฯเปฯ. สํเฆ ฯเปฯ. วิคต\-มล\-มจฺเฉเรน เจตสา อคารํ อชฺฌาวสติ, มุตฺตจาโค ปยตปาณิ โวสฺสคฺครโต ยาจโยโค ทาน\-สํวิภาค\-รโต อปฺปฏิวิภตฺตํ. อิเมหิ โข ถปตโย จตูหิ ธมฺเมหิ สมนฺนาคโต อริยสาวโก โสตาปนฺโน โหติ อวินิปาต\-ธมฺโม นิยโต สมฺโพธิ\-ปรายโณ.}

\prose{ตุเมฺห โข ถปตโย พุทฺเธ อเวจฺจ\-ปฺปสาเทน สมนฺนาคตา อิติปิ โส ภควา อรหํ ฯเปฯ พุทฺโธ ภควาติ. ธมฺเม ฯเปฯ. สํเฆ ฯเปฯ. ยํ โข ปน กิญฺจิ กุเล เทยฺยธมฺมํ, สพฺพํ ตํ อปฺปฏิวิภตฺตํ สีลวนฺเตหิ ลลฺยาณธมฺเมหิ, ตํ กิํ มญฺญถ ถปตโย กติวิธา เต โกสเลสุ มนุสฺสา เย ตุมฺหากํ สมสมา ยทิทํ ทานสํวิภาเคหีติ. ลาภา โน ภนฺเต, สุลทฺธํ โน ภนฺเต, เยสํ โน ภควา เอวํ ปชานาตีติ.}

\prose{อิทํ วาสนาภาคิยํ สุตฺตํ.}

\proseitem{95}{“เอกปุปฺผํ จชิตฺวาน\csromansharedfootnote{52bebb2eff8f712f}{ยชิตฺวาน (ก)}, สหสฺสํ กปฺปโกฏิโย.}

\setlength{\csromangathapagewidth}{0pt}
\setlength{\csromangathawakwidth}{0pt}
\csromangathameasuregroup{เทเว เจว มนุสฺเส จ, \\ เอกํ พุทฺธคตํ สญฺญํ, \\ อชฺช ติํสํ ตโต กปฺปา, \\ ติสฺโส วิชฺชา สจฺฉิกตา, \\ อนุกมฺปโก ปุเรภตฺตํ, \\ ปุริสสฺส วฏํสโก หตฺเถ, \\ โส อทฺทสาสิ สมฺพุทฺธํ, \\ ปวิสนฺตํ ราชมคฺเคน, \\ หฏฺโฐ จิตฺตํ ปสาเทตฺวา, \\ โส ตํ วฏํสกํ สุรภิํ, \\ สมฺพุทฺธสฺสุปนาเมสิ, \\ ตโต อคฺคิสิขา วณฺณา, \\ สหสฺสรํสิ วิชฺชุริว, \\ ปทกฺขิณํ กริตฺวาน, \\ ติกฺขตฺตุํ ปริวฏฺเฏตฺวา, \\ อิทํ ทิสฺวา อจฺฉริยํ, \\ เอกํสํ จีวรํ กตฺวา, \\ โก เหตุ สิตกมฺมสฺส, \\ ธมฺมาโลโก ภวิสฺสติ, \\ ยสฺส ตํ สพฺพธมฺเมสุ, \\ กงฺขิํ เวมติกํ เถรํ, \\ โย โส อานนฺท ปุริโส, \\ จตุราสีติกปฺปานิ, \\ เทเวสุ เทวโสภคฺคํ, \\ มนุเชสุ มนุชินฺโท, \\ โส จริมํ ปพฺพชิตฺวา, \\ ปจฺเจกพุทฺโธ ธุตราโค, \\ นตฺถิ จิตฺเต ปสนฺนมฺหิ, \\ ตถาคเต วา สมฺพุทฺเธ, \\ เอวํ อจินฺติยา พุทฺธา,}{\csromangathabat{เทเว เจว มนุสฺเส จ,}{เสเสน ปรินิพฺพุโต”ติ.} \\ อิทํ วาสนาภาคิยํ สุตฺตํ. \\ “อสฺสตฺเถ หริโตภาเส, \\ สํวิรูฬฺหมฺหิ ปาทเป. \\ \csromangathabat{เอกํ พุทฺธคตํ สญฺญํ,}{อลภิํ’หํ\textsuperscript{1} ปติสฺสโต\textsuperscript{1}.} \\ \csromangathabat{อชฺช ติํสํ ตโต กปฺปา,}{นาภิชานามิ ทุคฺคติํ.} \\ \csromangathabat{ติสฺโส วิชฺชา สจฺฉิกตา,}{ตสฺสา สญฺญาย วาสนา”ติ.} \\ อิทํ วาสนาภาคิยํ สุตฺตํ. \\ “ปิณฺฑาย โกสลํ ปุรํ, \\ ปาวิสิ อคฺคปุคฺคโล. \\ \csromangathabat{อนุกมฺปโก ปุเรภตฺตํ,}{ตณฺหานิฆาตโก มุนิ.} \\ \csromangathabat{ปุริสสฺส วฏํสโก หตฺเถ,}{สพฺพปุปฺเผหิ’ลงฺกโต.} \\ \csromangathabat{โส อทฺทสาสิ สมฺพุทฺธํ,}{ภิกฺขุสํฆปุรกฺขตํ.} \\ \csromangathabat{ปวิสนฺตํ ราชมคฺเคน,}{เทวมานุสปูชิตํ.} \\ \csromangathabat{หฏฺโฐ จิตฺตํ ปสาเทตฺวา,}{สมฺพุทฺธมุปสงฺกมิ.} \\ \csromangathabat{โส ตํ วฏํสกํ สุรภิํ,}{วณฺณวนฺตํ มโนรมํ.} \\ \csromangathabat{สมฺพุทฺธสฺสุปนาเมสิ,}{ปสนฺโน เสหิ ปาณิภิ.} \\ \csromangathabat{ตโต อคฺคิสิขา วณฺณา,}{พุทฺธสฺส ลปนนฺตรา.} \\ \csromangathabat{สหสฺสรํสิ วิชฺชุริว,}{โอกฺกา นิกฺขมิ อานนา.} \\ \csromangathabat{ปทกฺขิณํ กริตฺวาน,}{สีเส อาทิจฺจพนฺธุโน.} \\ \csromangathabat{ติกฺขตฺตุํ ปริวฏฺเฏตฺวา,}{มุทฺธนนฺตรธายถ.} \\ \csromangathabat{อิทํ ทิสฺวา อจฺฉริยํ,}{อพฺภุตํ โลมหํสนํ.} \\ \csromangathabat{เอกํสํ จีวรํ กตฺวา,}{อานนฺโท เอตทพฺรวิ.} \\ \csromangathabat{โก เหตุ สิตกมฺมสฺส,}{พฺยากโรหิ มหามุเน.} \\ \csromangathabat{ธมฺมาโลโก ภวิสฺสติ,}{กงฺขํ วิตร โน มุเน.} \\ \csromangathabat{ยสฺส ตํ สพฺพธมฺเมสุ,}{สทา ญาณํ ปวตฺตติ.} \\ \csromangathabat{กงฺขิํ เวมติกํ เถรํ,}{อานนฺทํ เอตทพฺรวิ.} \\ \csromangathabat{โย โส อานนฺท ปุริโส,}{มยิ จิตฺตํ ปสาทยิ.} \\ \csromangathabat{จตุราสีติกปฺปานิ,}{ทุคฺคติํ น คมิสฺสติ.} \\ \csromangathabat{เทเวสุ เทวโสภคฺคํ,}{ทิพฺพํ รชฺชํ ปสาสิย.} \\ \csromangathabat{มนุเชสุ มนุชินฺโท,}{ราชา รฏฺเฐ ภวิสฺสติ.} \\ \csromangathabat{โส จริมํ ปพฺพชิตฺวา,}{สจฺฉิกตฺวา จ\textsuperscript{1} ธมฺมตํ.} \\ \csromangathabat{ปจฺเจกพุทฺโธ ธุตราโค,}{วฏํสโก นาม ภวิสฺสติ.} \\ \csromangathabat{นตฺถิ จิตฺเต ปสนฺนมฺหิ,}{อปฺปกา นาม ทกฺขิณา.} \\ \csromangathabat{ตถาคเต วา สมฺพุทฺเธ,}{อถ วา ตสฺส สาวเก\textsuperscript{1}.} \\ \csromangathabat{เอวํ อจินฺติยา พุทฺธา,}{พุทฺธธมฺมา อจินฺติยา.} \\ อิทํ วาสนาภาคิยํ สุตฺตํ.}
\csromangathameasurewak{อิทํ วาสนาภาคิยํ สุตฺตํ.}
\csromangathasetleft{เทเว เจว มนุสฺเส จ, \\ เอกํ พุทฺธคตํ สญฺญํ, \\ อชฺช ติํสํ ตโต กปฺปา, \\ ติสฺโส วิชฺชา สจฺฉิกตา, \\ อนุกมฺปโก ปุเรภตฺตํ, \\ ปุริสสฺส วฏํสโก หตฺเถ, \\ โส อทฺทสาสิ สมฺพุทฺธํ, \\ ปวิสนฺตํ ราชมคฺเคน, \\ หฏฺโฐ จิตฺตํ ปสาเทตฺวา, \\ โส ตํ วฏํสกํ สุรภิํ, \\ สมฺพุทฺธสฺสุปนาเมสิ, \\ ตโต อคฺคิสิขา วณฺณา, \\ สหสฺสรํสิ วิชฺชุริว, \\ ปทกฺขิณํ กริตฺวาน, \\ ติกฺขตฺตุํ ปริวฏฺเฏตฺวา, \\ อิทํ ทิสฺวา อจฺฉริยํ, \\ เอกํสํ จีวรํ กตฺวา, \\ โก เหตุ สิตกมฺมสฺส, \\ ธมฺมาโลโก ภวิสฺสติ, \\ ยสฺส ตํ สพฺพธมฺเมสุ, \\ กงฺขิํ เวมติกํ เถรํ, \\ โย โส อานนฺท ปุริโส, \\ จตุราสีติกปฺปานิ, \\ เทเวสุ เทวโสภคฺคํ, \\ มนุเชสุ มนุชินฺโท, \\ โส จริมํ ปพฺพชิตฺวา, \\ ปจฺเจกพุทฺโธ ธุตราโค, \\ นตฺถิ จิตฺเต ปสนฺนมฺหิ, \\ ตถาคเต วา สมฺพุทฺเธ, \\ เอวํ อจินฺติยา พุทฺธา,}
\csromangathagroup{\csromangathastanza{\csromangathabat{เทเว เจว มนุสฺเส จ,}{เสเสน ปรินิพฺพุโต”ติ.} \\ อิทํ วาสนาภาคิยํ สุตฺตํ. \\ “อสฺสตฺเถ หริโตภาเส,}\csromangathastanza{สํวิรูฬฺหมฺหิ ปาทเป.}\csromangathastanza{\csromangathabat{เอกํ พุทฺธคตํ สญฺญํ,}{อลภิํ’หํ\csromansharedfootnote{1eb227b28db631fb}{อลภิตฺถํ (ก)} ปติสฺสโต\csromansharedfootnote{2ac2b612a01e821c}{ขุ 2. 261 ปิฏฺเฐ.}.} \\ \csromangathabat{อชฺช ติํสํ ตโต กปฺปา,}{นาภิชานามิ ทุคฺคติํ.}}\csromangathastanza{\csromangathabat{ติสฺโส วิชฺชา สจฺฉิกตา,}{ตสฺสา สญฺญาย วาสนา”ติ.} \\ อิทํ วาสนาภาคิยํ สุตฺตํ. \\ “ปิณฺฑาย โกสลํ ปุรํ,}\csromangathastanza{ปาวิสิ อคฺคปุคฺคโล.}\csromangathastanza{\csromanfolio{117}\csromangathabat{อนุกมฺปโก ปุเรภตฺตํ,}{ตณฺหานิฆาตโก มุนิ.} \\ \csromangathabat{ปุริสสฺส วฏํสโก หตฺเถ,}{สพฺพปุปฺเผหิ’ลงฺกโต.}}\csromangathastanza{\csromangathabat{โส อทฺทสาสิ สมฺพุทฺธํ,}{ภิกฺขุสํฆ\-ปุรกฺขตํ.} \\ \csromangathabat{ปวิสนฺตํ ราชมคฺเคน,}{เทวมานุสปูชิตํ.}}\csromangathastanza{\csromangathabat{หฏฺโฐ จิตฺตํ ปสาเทตฺวา,}{สมฺพุทฺธมุปสงฺกมิ.} \\ \csromangathabat{โส ตํ วฏํสกํ สุรภิํ,}{วณฺณวนฺตํ มโนรมํ.}}\csromangathastanza{\csromangathabat{สมฺพุทฺธสฺสุปนาเมสิ,}{ปสนฺโน เสหิ ปาณิภิ.} \\ \csromangathabat{ตโต อคฺคิสิขา วณฺณา,}{พุทฺธสฺส ลปนนฺตรา.}}\csromangathastanza{\csromangathabat{สหสฺสรํสิ วิชฺชุริว,}{โอกฺกา นิกฺขมิ อานนา.} \\ \csromangathabat{ปทกฺขิณํ กริตฺวาน,}{สีเส อาทิจฺจพนฺธุโน.}}\csromangathastanza{\csromangathabat{ติกฺขตฺตุํ ปริวฏฺเฏตฺวา,}{มุทฺธนนฺตรธายถ.} \\ \csromangathabat{อิทํ ทิสฺวา อจฺฉริยํ,}{อพฺภุตํ โลมหํสนํ.}}\csromangathastanza{\csromangathabat{เอกํสํ จีวรํ กตฺวา,}{อานนฺโท เอตทพฺรวิ.} \\ \csromangathabat{โก เหตุ สิตกมฺมสฺส,}{พฺยากโรหิ มหามุเน.}}\csromangathastanza{\csromangathabat{ธมฺมาโลโก ภวิสฺสติ,}{กงฺขํ วิตร โน มุเน.} \\ \csromangathabat{ยสฺส ตํ สพฺพธมฺเมสุ,}{สทา ญาณํ ปวตฺตติ.}}\csromangathastanza{\csromangathabat{กงฺขิํ เวมติกํ เถรํ,}{อานนฺทํ เอตทพฺรวิ.} \\ \csromangathabat{โย โส อานนฺท ปุริโส,}{มยิ จิตฺตํ ปสาทยิ.}}\csromangathastanza{\csromangathabat{จตุราสีติกปฺปานิ,}{ทุคฺคติํ น คมิสฺสติ.} \\ \csromangathabat{เทเวสุ เทวโสภคฺคํ,}{ทิพฺพํ รชฺชํ ปสาสิย.}}\csromangathastanza{\csromangathabat{มนุเชสุ มนุชินฺโท,}{ราชา รฏฺเฐ ภวิสฺสติ.} \\ \csromangathabat{โส จริมํ ปพฺพชิตฺวา,}{สจฺฉิกตฺวา จ\csromansharedfootnote{3a44bcd740210156}{สจฺฉิกตฺวาน (ก)} ธมฺมตํ.}}\csromangathastanza{\csromangathabat{ปจฺเจกพุทฺโธ ธุตราโค,}{วฏํสโก นาม ภวิสฺสติ.} \\ \csromangathabat{นตฺถิ จิตฺเต ปสนฺนมฺหิ,}{อปฺปกา นาม ทกฺขิณา.}}\csromangathastanza{\csromangathabat{ตถาคเต วา สมฺพุทฺเธ,}{อถ วา ตสฺส สาวเก\csromansharedfootnote{338dd9d64fd7ee1e}{[มิสฺสินฺคฺ โนเต 0]}.} \\ \csromangathabat{เอวํ อจินฺติยา พุทฺธา,}{พุทฺธธมฺมา อจินฺติยา.}}\csromangathastanza{\csromangathawak{อิทํ วาสนาภาคิยํ สุตฺตํ.}}}

\prose{อจินฺติเย ปสนฺนานํ, วิปาโก โหติ อจินฺติโย”ติ\csromansharedfootnote{3952bd01c43886e8}{ขุ 2. 66 ปิฏฺเฐ.}.}

\proseitem{96}{\csromanfolio{118}“อิธาหํ ภิกฺขเว เอกจฺจํ ปุคฺคลํ ปสนฺนจิตฺตํ เอวํ เจตสา เจโต ปริจฺจ ปชานามิ “(ยถา โข อยํ ปุคฺคโล อิริยติ, ยญฺจ ปฏิปทํ ปฏิปนฺโน, ยญฺจ มคฺคํ สมารูโฬฺห.)\csromansharedfootnote{0dee15b25e1f98fd}{( ) นตฺถิ อํ 1. 8 ปิฏฺเฐ; ขุ 1. 204 ปิฏฺเฐ จ.} อิมมฺหิ จายํ สมเย กาลํ กเรยฺย, ยถาภตํ นิกฺขิตฺโต เอวํ สคฺเค. ตํ กิสฺส เหตุ, จิตฺตํ หิสฺส ภิกฺขเว ปสนฺนํ, เจโตปสาทเหตุ\csromansharedfootnote{33a931832fb4cb9c}{จิตฺตปฺปสาทเหตุ (สี, ก)} โข ปน เอวมิเธกจฺเจ สตฺตา กายสฺส เภทา ปรํ มรณา สุคติํ สคฺคํ โลกํ อุปปชฺชนฺตี”ติ. เอตมตฺถํ ภควา อโวจ, ตตฺเถตํ อิติ วุจฺจติ\csromandash{}}

\setlength{\csromangathapagewidth}{0pt}
\setlength{\csromangathawakwidth}{0pt}
\csromangathameasuregroup{“ปสนฺนจิตฺตํ ญตฺวาน, \\ เอตมตฺถญฺจ พฺยากาสิ, \\ อิมมฺหิ จายํ สมเย, \\ สคฺคมฺหิ อุปปชฺเชยฺย, \\ เจโตปสาทเหตุ หิ,}{\csromangathabat{“ปสนฺนจิตฺตํ ญตฺวาน,}{เอกจฺจํ อิธ ปุคฺคลํ.} \\ \csromangathabat{เอตมตฺถญฺจ พฺยากาสิ,}{พุทฺโธ\textsuperscript{1} ภิกฺขูน สนฺติเก.} \\ \csromangathabat{อิมมฺหิ จายํ สมเย,}{กาลํ กยิราถ ปุคฺคโล.} \\ \csromangathabat{สคฺคมฺหิ อุปปชฺเชยฺย,}{จิตฺตํ หิสฺส ปสาทิตํ.} \\ \csromangathabat{เจโตปสาทเหตุ หิ,}{สตฺตา คจฺฉนฺติ สุคฺคติํ.}}
\csromangathasetleft{“ปสนฺนจิตฺตํ ญตฺวาน, \\ เอตมตฺถญฺจ พฺยากาสิ, \\ อิมมฺหิ จายํ สมเย, \\ สคฺคมฺหิ อุปปชฺเชยฺย, \\ เจโตปสาทเหตุ หิ,}
\csromangathagroup{\csromangathastanza{\csromangathabat{“ปสนฺนจิตฺตํ ญตฺวาน,}{เอกจฺจํ อิธ ปุคฺคลํ.} \\ \csromangathabat{เอตมตฺถญฺจ พฺยากาสิ,}{พุทฺโธ\csromansharedfootnote{8fe3171211a0ea00}{สตฺถา (สี, ก)} ภิกฺขูน สนฺติเก.}}\csromangathastanza{\csromangathabat{อิมมฺหิ จายํ สมเย,}{กาลํ กยิราถ ปุคฺคโล.} \\ \csromangathabat{สคฺคมฺหิ อุปปชฺเชยฺย,}{จิตฺตํ หิสฺส ปสาทิตํ.} \\ \csromangathabat{เจโตปสาทเหตุ หิ,}{สตฺตา คจฺฉนฺติ สุคฺคติํ.}}}

\prose{ยถาภตํ นิกฺขิเปยฺย, เอวเมวํ ตถาวิโธ}

\setlength{\csromangathapagewidth}{0pt}
\setlength{\csromangathawakwidth}{0pt}
\csromangathameasuregroup{กายสฺส เภทา สปฺปญฺโญ,}{\csromangathabat{กายสฺส เภทา สปฺปญฺโญ,}{สคฺคํ โส อุปปชฺชตีติ.}}
\csromangathasetleft{กายสฺส เภทา สปฺปญฺโญ,}
\csromangathagroup{\csromangathastanza{\csromangathabat{กายสฺส เภทา สปฺปญฺโญ,}{สคฺคํ โส อุปปชฺชตีติ.}}}

\prose{อยมฺปิ อตฺโถ วุตฺโต ภควตา อิติ เม สุตนฺ”ติ. อิทํ วาสนาภาคิยํ สุตฺตํ.}

\setlength{\csromangathapagewidth}{0pt}
\setlength{\csromangathawakwidth}{0pt}
\csromangathameasuregroup{“สุวณฺณจฺฉทนํ นาวํ, \\ โอคาหสิ\textsuperscript{1} โปกฺขรณิํ, \\ เกน เต ตาทิโส วณฺโณ, \\ อุปฺปชฺชนฺติ จ เต โภคา, \\ ปุจฺฉิตา เทวเต สํส, \\ สา เทวตา อตฺตมนา, \\ ปญฺห ปุฏฺฐา วิยากาสิ, \\ อทฺธานํ ปฏิปนฺนาหํ, \\ ตตฺถ จิตฺตํ ปสาเทสิํ, \\ ปทฺมปุปฺเผหิ ปูเชสิํ, \\ กมฺมสฺส ผลํ วิปาโก,}{\csromangathabat{“สุวณฺณจฺฉทนํ นาวํ,}{นาริ อารุยฺห ติฏฺฐสิ.} \\ \csromangathabat{โอคาหสิ\textsuperscript{1} โปกฺขรณิํ,}{ปทฺมํ ฉินฺทสิ ปาณินา\textsuperscript{1}.} \\ \csromangathabat{เกน เต ตาทิโส วณฺโณ,}{อานุภาโว ชุติ จ เต.} \\ \csromangathabat{อุปฺปชฺชนฺติ จ เต โภคา,}{เย เกจิ มนสิจฺฉิตา.} \\ \csromangathabat{ปุจฺฉิตา เทวเต สํส,}{กิสฺส กมฺมสฺสิทํ ผลํ.} \\ \csromangathabat{สา เทวตา อตฺตมนา,}{เทวราเชน ปุจฺฉิตา.} \\ \csromangathabat{ปญฺห ปุฏฺฐา วิยากาสิ,}{สกฺกสฺส อิติ เม สุตํ.} \\ \csromangathabat{อทฺธานํ ปฏิปนฺนาหํ,}{ทิสฺวา ถูปํ มโนรมํ.} \\ \csromangathabat{ตตฺถ จิตฺตํ ปสาเทสิํ,}{กสฺสปสฺส ยสสฺสิโน.} \\ \csromangathabat{ปทฺมปุปฺเผหิ ปูเชสิํ,}{ปสนฺนา เสหิ ตสฺเสว.} \\ \csromangathabat{กมฺมสฺส ผลํ วิปาโก,}{เอตาทิสํ กตปุญฺญา ลภนฺตี”ติ.} \\ อิทํ วาสนาภาคิยํ สุตฺตํ.}
\csromangathasetleft{“สุวณฺณจฺฉทนํ นาวํ, \\ โอคาหสิ\textsuperscript{1} โปกฺขรณิํ, \\ เกน เต ตาทิโส วณฺโณ, \\ อุปฺปชฺชนฺติ จ เต โภคา, \\ ปุจฺฉิตา เทวเต สํส, \\ สา เทวตา อตฺตมนา, \\ ปญฺห ปุฏฺฐา วิยากาสิ, \\ อทฺธานํ ปฏิปนฺนาหํ, \\ ตตฺถ จิตฺตํ ปสาเทสิํ, \\ ปทฺมปุปฺเผหิ ปูเชสิํ, \\ กมฺมสฺส ผลํ วิปาโก,}
\csromangathagroup{\csromangathastanza{\csromangathabat{“สุวณฺณ\-จฺฉทนํ นาวํ,}{นาริ อารุยฺห ติฏฺฐสิ.} \\ \csromangathabat{โอคาหสิ\csromansharedfootnote{68c5d0db6afc5511}{โอคาหเส (สี, ก)} โปกฺขรณิํ,}{ปทฺมํ ฉินฺทสิ ปาณินา\csromansharedfootnote{3c60fe90fb2ece67}{ขุ 2. 6, 7 ปิฏฺเฐสุปิ.}.}}\csromangathastanza{\csromangathabat{เกน เต ตาทิโส วณฺโณ,}{อานุภาโว ชุติ จ เต.} \\ \csromangathabat{อุปฺปชฺชนฺติ จ เต โภคา,}{เย เกจิ มนสิจฺฉิตา.}}\csromangathastanza{\csromangathabat{ปุจฺฉิตา เทวเต สํส,}{กิสฺส กมฺมสฺสิทํ ผลํ.} \\ \csromangathabat{สา เทวตา อตฺตมนา,}{เทวราเชน ปุจฺฉิตา.}}\csromangathastanza{\csromanfolio{119}\csromangathabat{ปญฺห ปุฏฺฐา วิยากาสิ,}{สกฺกสฺส อิติ เม สุตํ.} \\ \csromangathabat{อทฺธานํ ปฏิปนฺนาหํ,}{ทิสฺวา ถูปํ มโนรมํ.}}\csromangathastanza{\csromangathabat{ตตฺถ จิตฺตํ ปสาเทสิํ,}{กสฺสปสฺส ยสสฺสิโน.} \\ \csromangathabat{ปทฺมปุปฺเผหิ ปูเชสิํ,}{ปสนฺนา เสหิ ตสฺเสว.}}\csromangathastanza{\csromangathabat{กมฺมสฺส ผลํ วิปาโก,}{เอตาทิสํ กตปุญฺญา ลภนฺตี”ติ.} \\ อิทํ วาสนาภาคิยํ สุตฺตํ.}}

\prose{“ทานกถา สีลกถา สคฺคกถา ปุญฺญกถา ปุญฺญวิปาก\-กถา”ติ. อิทํ วาสนาภาคิยํ สุตฺตํ.}

\prose{“อปิจาปิ ปํสุถูเปสุ อุทฺทิสฺสกเตสุ ทสพล\-ธรานํ ตตฺถปิ การํ กตฺวา สคฺเคสุ นรา ปโมทนฺตี”ติ. อิทํ วาสนาภาคิยํ สุตฺตํ.}

\proseitem{97}{“เทวปุตฺต\-สรีรวณฺณา, สพฺเพ สุภค\-สณฺฐิตี.}

\setlength{\csromangathapagewidth}{0pt}
\setlength{\csromangathawakwidth}{0pt}
\csromangathameasuregroup{อุปฺปลานิ จ จตฺตาริ, \\ อชฺช ติํสํ ตโต กปฺปา, \\ วินิปาตํ น คจฺฉามิ, \\ อชฺช ติํสํ ตโต กปฺปา, \\ วินิปาตญฺจ น คจฺฉามิ, \\ “กปิลํ นาม นครํ, \\ อากิณฺณมิทฺธํ ผีตญฺจ, \\ กุมฺมาสํ วิกฺกิณิํ ตตฺถ, \\ โสหํ อทฺทสิํ สมฺพุทฺธํ, \\ หฏฺโฐ จิตฺตํ ปสาเทตฺวา, \\ อริฏฺฐํ ธุวภตฺเตน, \\ ตโต จ กตฺติโก ปุณฺโณ\textsuperscript{1}, \\ นวํ ทุสฺสยุคํ คยฺห, \\ ปสนฺนจิตฺตํ ญตฺวาน, \\ อนุกมฺปโก การุณิโก, \\ ตาหํ กมฺมํ กริตฺวาน, \\ เทเว เจว มนุสฺเส จ, \\ พาราณสิยํ นคเร, \\ อฑฺฒ กุลสฺมิํ อุปฺปชฺชิํ, \\ ตโต จ วิญฺญุตํ ปตฺโต, \\ ปาสาทา โอรูหิตฺวาน, \\ โส เม ธมฺมมเทสยิ, \\ ทุกฺขํ ทุกฺขสมุปฺปาทํ, \\ อริยํ อฏฺฐงฺคิกํ มคฺคํ, \\ จตฺตาริ อริยสจฺจานิ, \\ ตสฺสาหํ วจนํ สุตฺวา, \\ สมถํ ปฏิวิชฺฌาหํ, \\ อชฺฌตฺตญฺจ พหิทฺธา จ, \\ สพฺเพ อาสุํ สมุจฺฉินฺนา, \\ ปริยนฺตกตํ ทุกฺขํ, \\ ชาติมรณสํสาโร,}{อุทเกน ปํสุํ เตเมตฺวา, \\ ถูปํ วฑฺเฒถ กสฺสปํ. \\ อยํ สุคตฺเต สุคตสฺส ถูโป, \\ มเหสิโน ทสพลธมฺมธาริโน. \\ ตสฺมิํ\textsuperscript{1} อิเม เทวมนุชา ปสนฺนา, \\ การํ กโรนฺตา ชรามรณา ปมุจฺจเร”ติ. \\ อิทํ วาสนาภาคิยํ สุตฺตํ. \\ “อุฬารํ วต ตํ อาสิ, \\ ยาหํ ถูปํ มเหสิโน. \\ \csromangathabat{อุปฺปลานิ จ จตฺตาริ,}{มาลญฺจ อภิโรปยิํ.} \\ \csromangathabat{อชฺช ติํสํ ตโต กปฺปา,}{นาภิชานามิ ทุคฺคติํ.} \\ \csromangathabat{วินิปาตํ น คจฺฉามิ,}{ถูปํ ปูเชตฺว\textsuperscript{1} สตฺถุโน”ติ.} \\ อิทํ วาสนาภาคิยํ สุตฺตํ. \\ “พาตฺติํสลกฺขณธรสฺส, \\ วิชิตวิชยสฺส โลกนาถสฺส. \\ สตสหสฺสํ กปฺเป, \\ มุทิโต ถูปํ อปูเชสิ. \\ ยํ มยา ปสุตํ ปุญฺญํ, \\ เตน จ ปุญฺเญน เทว โสภคฺคํ. \\ รชฺชานิ จ การิตานิ, \\ อนาคนฺตุน วินิปาตํ. \\ ยํ จกฺขุ อทนฺตทมกสฺส, \\ สาสเน ปณิหิตํ ตถา. \\ จิตฺตํ ตํ เม สพฺพํ, \\ ลทฺธํ วิมุตฺตจิตฺตมฺหิ วิธูตลโต”ติ. \\ อิทํ วาสนาภาคิยํ สุตฺตํ. \\ “สามากปตฺโถทนมตฺตเมว หิ, \\ ปจฺเจกพุทฺธมฺหิ อทาสิ ทกฺขิณํ. \\ วิมุตฺตจิตฺเต อขิเล อนาสเว, \\ อรณวิหาริมฺหิ อสงฺคมานเส. \\ ตสฺมิํ จ โอกปฺปยิ ธมฺมมุตฺตมํ, \\ ตสฺมิํ จ ธมฺเม ปณิเธสิํ มานสํ. \\ เอวํวิหารีหิ เม สงฺคโม สิยา, \\ ภเว กุทาสุปิ จ มา อเปกฺขวา. \\ ตสฺเสว กมฺมสฺส วิปากโต อหํ, \\ สหสฺสกฺขตฺตุํ กุรุสูปปชฺชถ\textsuperscript{1}. \\ ทีฆายุเกสุ อมเมสุ ปาณิสุ, \\ วิเสสคามีสุ อหีนคามิสุ. \\ ตสฺเสว กมฺมสฺส วิปากโต อหํ, \\ สหสฺสกฺขตฺตุํ ติทโสปปชฺชถ. \\ วิจิตฺรมาลาภรณานุเลปิสุ, \\ วิสิฏฺฐกายูปคโต ยสสฺสิสุ. \\ ตสฺเสว กมฺมสฺส วิปากโต อหํ, \\ วิมุตฺตจิตฺโต อขิโล อนาสโว. \\ อิเมหิ เม อนฺติมเทหธาริภิ, \\ สมาคโม อาสิ หิตาหิตาสิหิ. \\ ปจฺจกฺขํ ขฺวิมํ อวจ ตถาคโต ชิโน, \\ สมิชฺฌเต สีลวโต ยทิจฺฉติ. \\ ยถา ยถา เม มนสา วิจินฺติตํ, \\ ตถา สมิทฺธํ อยมนฺติโม ภโว”ติ. \\ อิทํ วาสนาภาคิยํ สุตฺตํ. \\ “เอกติํสมฺหิ กปฺปมฺหิ ชิโน อเนโช, \\ อนนฺตทสฺสี ภควา สิขีติ. \\ ตสฺสาปิ ราชา ภาตา สิขิทฺเธ\textsuperscript{1}, \\ พุทฺเธ จ ธมฺเม จ อภิปฺปสนฺโน. \\ ปรินิพฺพุเต โลกวินายกมฺหิ, \\ ถูปํ สกาสิ วิปุลํ มหนฺตํ. \\ สมนฺตโต คาวุติกํ มเหสิโน, \\ เทวาติเทวสฺส นรุตฺตมสฺส. \\ ตสฺมิํ มนุสฺโส พลิมาภิหารี, \\ ปคฺคยฺห ชาติสุมนํ ปหฏฺโฐ. \\ วาเตน ปุปฺผํ ปติตสฺส เอกํ, \\ ตาหํ คเหตฺวาน ตสฺเสว ทาสิ. \\ โส มํ อโวจาภิปสนฺนจิตฺโต, \\ ตุยฺหเมว เอตํ ปุปฺผํ ททามิ. \\ ตาหํ คเหตฺวา อภิโรปเยสิํ, \\ ปุนปฺปุนํ พุทฺธมนุสฺสรนฺโต. \\ \csromangathabat{อชฺช ติํสํ ตโต กปฺปา,}{นาภิชานามิ ทุคฺคติํ.} \\ \csromangathabat{วินิปาตญฺจ น คจฺฉามิ,}{ถูปปูชายิทํ ผลนฺ”ติ.} \\ อิทํ วาสนาภาคิยํ สุตฺตํ. \\ \csromangathabat{“กปิลํ นาม นครํ,}{สุวิภตฺตํ มหาปถํ.} \\ \csromangathabat{อากิณฺณมิทฺธํ ผีตญฺจ,}{พฺรหฺมทตฺตสฺส ราชิโน.} \\ \csromangathabat{กุมฺมาสํ วิกฺกิณิํ ตตฺถ,}{ปญฺจาลานํ ปุรุตฺตเม.} \\ \csromangathabat{โสหํ อทฺทสิํ สมฺพุทฺธํ,}{อุปริฏฺฐํ ยสสฺสินํ.} \\ \csromangathabat{หฏฺโฐ จิตฺตํ ปสาเทตฺวา,}{นิมนฺเตสิํ นรุตฺตมํ.} \\ \csromangathabat{อริฏฺฐํ ธุวภตฺเตน,}{ยํ เม เคหมฺหิ วิชฺชถ.} \\ \csromangathabat{ตโต จ กตฺติโก ปุณฺโณ\textsuperscript{1},}{ปุณฺณมาสี อุปฏฺฐิตา.} \\ \csromangathabat{นวํ ทุสฺสยุคํ คยฺห,}{อริฏฺฐสฺโสปนามยิํ.} \\ \csromangathabat{ปสนฺนจิตฺตํ ญตฺวาน,}{ปฏิคฺคณฺหิ นรุตฺตโม.} \\ \csromangathabat{อนุกมฺปโก การุณิโก,}{ตณฺหานิฆาตโก มุนิ.} \\ \csromangathabat{ตาหํ กมฺมํ กริตฺวาน,}{กลฺยาณํ พุทฺธวณฺณิตํ.} \\ \csromangathabat{เทเว เจว มนุสฺเส จ,}{สนฺธาวิตฺวา ตโต จุโต.} \\ \csromangathabat{พาราณสิยํ นคเร,}{เสฏฺฐิสฺส เอกปุตฺตโก.} \\ \csromangathabat{อฑฺฒ กุลสฺมิํ อุปฺปชฺชิํ,}{ปาเณหิ จ ปิยตโร.} \\ \csromangathabat{ตโต จ วิญฺญุตํ ปตฺโต,}{เทวปุตฺเตน โจทิโต.} \\ \csromangathabat{ปาสาทา โอรูหิตฺวาน,}{สมฺพุทฺธมุปสงฺกมิํ.} \\ \csromangathabat{โส เม ธมฺมมเทสยิ,}{อนุกมฺปาย โคตโม.} \\ \csromangathabat{ทุกฺขํ ทุกฺขสมุปฺปาทํ,}{ทุกฺขสฺส จ อติกฺกมํ.} \\ \csromangathabat{อริยํ อฏฺฐงฺคิกํ มคฺคํ,}{ทุกฺขูปสมคามินํ.} \\ \csromangathabat{จตฺตาริ อริยสจฺจานิ,}{มุนิ ธมฺมมเทสยิ.} \\ \csromangathabat{ตสฺสาหํ วจนํ สุตฺวา,}{วิหริํ สาสเน รโต.} \\ \csromangathabat{สมถํ ปฏิวิชฺฌาหํ,}{รตฺตินฺทิวมตนฺทิโต.} \\ \csromangathabat{อชฺฌตฺตญฺจ พหิทฺธา จ,}{เย เม วิชฺชิํสุ\textsuperscript{1} อาสวา.} \\ \csromangathabat{สพฺเพ อาสุํ สมุจฺฉินฺนา,}{น จ อุปฺปชฺชเร ปุน.} \\ \csromangathabat{ปริยนฺตกตํ ทุกฺขํ,}{จริโมยํ สมุสฺสโย.} \\ \csromangathabat{ชาติมรณสํสาโร,}{นตฺถิทานิ ปุนพฺภโว”ติ.} \\ อิทํ วาสนาภาคิย สุตฺตํ.}
\csromangathameasurewak{อุทเกน ปํสุํ เตเมตฺวา, \\ ถูปํ วฑฺเฒถ กสฺสปํ. \\ อยํ สุคตฺเต สุคตสฺส ถูโป, \\ มเหสิโน ทสพลธมฺมธาริโน. \\ ยาหํ ถูปํ มเหสิโน. \\ วิชิตวิชยสฺส โลกนาถสฺส. \\ สตสหสฺสํ กปฺเป, \\ มุทิโต ถูปํ อปูเชสิ. \\ ยํ มยา ปสุตํ ปุญฺญํ, \\ เตน จ ปุญฺเญน เทว โสภคฺคํ. \\ รชฺชานิ จ การิตานิ, \\ อนาคนฺตุน วินิปาตํ. \\ ยํ จกฺขุ อทนฺตทมกสฺส, \\ สาสเน ปณิหิตํ ตถา. \\ จิตฺตํ ตํ เม สพฺพํ, \\ ลทฺธํ วิมุตฺตจิตฺตมฺหิ วิธูตลโต”ติ. \\ อิทํ วาสนาภาคิยํ สุตฺตํ. \\ “สามากปตฺโถทนมตฺตเมว หิ, \\ ปจฺเจกพุทฺธมฺหิ อทาสิ ทกฺขิณํ. \\ วิมุตฺตจิตฺเต อขิเล อนาสเว, \\ อรณวิหาริมฺหิ อสงฺคมานเส. \\ ตสฺมิํ จ โอกปฺปยิ ธมฺมมุตฺตมํ, \\ ตสฺมิํ จ ธมฺเม ปณิเธสิํ มานสํ. \\ เอวํวิหารีหิ เม สงฺคโม สิยา, \\ ภเว กุทาสุปิ จ มา อเปกฺขวา. \\ ตสฺเสว กมฺมสฺส วิปากโต อหํ, \\ สหสฺสกฺขตฺตุํ กุรุสูปปชฺชถ\textsuperscript{1}. \\ ทีฆายุเกสุ อมเมสุ ปาณิสุ, \\ วิเสสคามีสุ อหีนคามิสุ. \\ ตสฺเสว กมฺมสฺส วิปากโต อหํ, \\ สหสฺสกฺขตฺตุํ ติทโสปปชฺชถ. \\ วิจิตฺรมาลาภรณานุเลปิสุ, \\ วิสิฏฺฐกายูปคโต ยสสฺสิสุ. \\ ตสฺเสว กมฺมสฺส วิปากโต อหํ, \\ วิมุตฺตจิตฺโต อขิโล อนาสโว. \\ อิเมหิ เม อนฺติมเทหธาริภิ, \\ สมาคโม อาสิ หิตาหิตาสิหิ. \\ ปจฺจกฺขํ ขฺวิมํ อวจ ตถาคโต ชิโน, \\ สมิชฺฌเต สีลวโต ยทิจฺฉติ. \\ ยถา ยถา เม มนสา วิจินฺติตํ, \\ ตถา สมิทฺธํ อยมนฺติโม ภโว”ติ. \\ อิทํ วาสนาภาคิยํ สุตฺตํ. \\ “เอกติํสมฺหิ กปฺปมฺหิ ชิโน อเนโช, \\ อนนฺตทสฺสี ภควา สิขีติ. \\ ตสฺสาปิ ราชา ภาตา สิขิทฺเธ\textsuperscript{1}, \\ พุทฺเธ จ ธมฺเม จ อภิปฺปสนฺโน. \\ ปรินิพฺพุเต โลกวินายกมฺหิ, \\ ถูปํ สกาสิ วิปุลํ มหนฺตํ. \\ สมนฺตโต คาวุติกํ มเหสิโน, \\ เทวาติเทวสฺส นรุตฺตมสฺส. \\ ตสฺมิํ มนุสฺโส พลิมาภิหารี, \\ ปคฺคยฺห ชาติสุมนํ ปหฏฺโฐ. \\ วาเตน ปุปฺผํ ปติตสฺส เอกํ, \\ ตาหํ คเหตฺวาน ตสฺเสว ทาสิ. \\ โส มํ อโวจาภิปสนฺนจิตฺโต, \\ ตุยฺหเมว เอตํ ปุปฺผํ ททามิ. \\ ตาหํ คเหตฺวา อภิโรปเยสิํ, \\ ปุนปฺปุนํ พุทฺธมนุสฺสรนฺโต. \\ อิทํ วาสนาภาคิยํ สุตฺตํ. \\ อิทํ วาสนาภาคิย สุตฺตํ.}
\csromangathasetleft{อุปฺปลานิ จ จตฺตาริ, \\ อชฺช ติํสํ ตโต กปฺปา, \\ วินิปาตํ น คจฺฉามิ, \\ อชฺช ติํสํ ตโต กปฺปา, \\ วินิปาตญฺจ น คจฺฉามิ, \\ “กปิลํ นาม นครํ, \\ อากิณฺณมิทฺธํ ผีตญฺจ, \\ กุมฺมาสํ วิกฺกิณิํ ตตฺถ, \\ โสหํ อทฺทสิํ สมฺพุทฺธํ, \\ หฏฺโฐ จิตฺตํ ปสาเทตฺวา, \\ อริฏฺฐํ ธุวภตฺเตน, \\ ตโต จ กตฺติโก ปุณฺโณ\textsuperscript{1}, \\ นวํ ทุสฺสยุคํ คยฺห, \\ ปสนฺนจิตฺตํ ญตฺวาน, \\ อนุกมฺปโก การุณิโก, \\ ตาหํ กมฺมํ กริตฺวาน, \\ เทเว เจว มนุสฺเส จ, \\ พาราณสิยํ นคเร, \\ อฑฺฒ กุลสฺมิํ อุปฺปชฺชิํ, \\ ตโต จ วิญฺญุตํ ปตฺโต, \\ ปาสาทา โอรูหิตฺวาน, \\ โส เม ธมฺมมเทสยิ, \\ ทุกฺขํ ทุกฺขสมุปฺปาทํ, \\ อริยํ อฏฺฐงฺคิกํ มคฺคํ, \\ จตฺตาริ อริยสจฺจานิ, \\ ตสฺสาหํ วจนํ สุตฺวา, \\ สมถํ ปฏิวิชฺฌาหํ, \\ อชฺฌตฺตญฺจ พหิทฺธา จ, \\ สพฺเพ อาสุํ สมุจฺฉินฺนา, \\ ปริยนฺตกตํ ทุกฺขํ, \\ ชาติมรณสํสาโร,}
\csromangathagroup{\csromangathastanza{\csromangathawak{อุทเกน ปํสุํ เตเมตฺวา, \\ ถูปํ วฑฺเฒถ กสฺสปํ. \\ อยํ สุคตฺเต สุคตสฺส ถูโป, \\ มเหสิโน ทสพลธมฺมธาริโน.}}\csromangathastanza{ตสฺมิํ\csromansharedfootnote{338dd9d64fd7ee1e}{[มิสฺสินฺคฺ โนเต 0]} อิเม เทวมนุชา ปสนฺนา, \\ การํ กโรนฺตา ชรามรณา ปมุจฺจเร”ติ. \\ อิทํ วาสนาภาคิยํ สุตฺตํ. \\ “อุฬารํ วต ตํ อาสิ,}\csromangathastanza{\csromangathawak{ยาหํ ถูปํ มเหสิโน.}}\csromangathastanza{\csromangathabat{อุปฺปลานิ จ จตฺตาริ,}{มาลญฺจ อภิโรปยิํ.} \\ \csromangathabat{อชฺช ติํสํ ตโต กปฺปา,}{นาภิชานามิ ทุคฺคติํ.}}\csromangathastanza{\csromangathabat{วินิปาตํ น คจฺฉามิ,}{ถูปํ ปูเชตฺว\csromansharedfootnote{338dd9d64fd7ee1e}{[มิสฺสินฺคฺ โนเต 0]} สตฺถุโน”ติ.} \\ อิทํ วาสนาภาคิยํ สุตฺตํ. \\ “พาตฺติํสลกฺขณธรสฺส,}\csromangathastanza{\csromangathawak{วิชิต\-วิชยสฺส โลกนาถสฺส.}}\csromangathastanza{\csromangathawak{สตสหสฺสํ กปฺเป, \\ มุทิโต ถูปํ อปูเชสิ. \\ ยํ มยา ปสุตํ ปุญฺญํ, \\ เตน จ ปุญฺเญน เทว โสภคฺคํ.}}\csromangathastanza{\csromangathawak{\csromanfolio{120}รชฺชานิ จ การิตานิ, \\ อนาคนฺตุน วินิปาตํ. \\ ยํ จกฺขุ อทนฺตทมกสฺส, \\ สาสเน ปณิหิตํ ตถา.}}\csromangathastanza{\csromangathawak{จิตฺตํ ตํ เม สพฺพํ, \\ ลทฺธํ วิมุตฺตจิตฺตมฺหิ วิธูตลโต”ติ. \\ อิทํ วาสนาภาคิยํ สุตฺตํ.}}\csromangathastanza{\csromangathawak{\csromangathaitemline{98}{“สามาก\-ปตฺโถทน\-มตฺตเมว หิ,} \\ \csromangathacontline{ปจฺเจกพุทฺธมฺหิ อทาสิ ทกฺขิณํ.} \\ \csromangathacontline{วิมุตฺตจิตฺเต อขิเล อนาสเว,} \\ \csromangathacontline{อรณวิหาริมฺหิ อสงฺคมานเส.} \\ \csromangathacontline{ตสฺมิํ จ โอกปฺปยิ ธมฺมมุตฺตมํ,} \\ \csromangathacontline{ตสฺมิํ จ ธมฺเม ปณิเธสิํ มานสํ.} \\ \csromangathacontline{เอวํวิหารีหิ เม สงฺคโม สิยา,} \\ \csromangathacontline{ภเว กุทาสุปิ จ มา อเปกฺขวา.} \\ \csromangathacontline{ตสฺเสว กมฺมสฺส วิปากโต อหํ,} \\ \csromangathacontline{สหสฺส\-กฺขตฺตุํ กุรุสูปปชฺชถ\csromansharedfootnote{13bab22be1b3ccb4}{กุรูสูปปชฺชถ (สี)}.} \\ \csromangathacontline{ทีฆายุเกสุ อมเมสุ ปาณิสุ,} \\ \csromangathacontline{วิเสสคามีสุ อหีนคามิสุ.} \\ \csromangathacontline{ตสฺเสว กมฺมสฺส วิปากโต อหํ,} \\ \csromangathacontline{สหสฺส\-กฺขตฺตุํ ติทโสปปชฺชถ.} \\ \csromangathacontline{วิจิตฺรมาลาภรณานุเลปิสุ,} \\ \csromangathacontline{วิสิฏฺฐกายูปคโต ยสสฺสิสุ.} \\ \csromangathacontline{ตสฺเสว กมฺมสฺส วิปากโต อหํ,} \\ \csromangathacontline{วิมุตฺตจิตฺโต อขิโล อนาสโว.} \\ \csromangathacontline{อิเมหิ เม อนฺติมเทหธาริภิ,} \\ \csromangathacontline{สมาคโม อาสิ หิตาหิตาสิหิ.}}}\csromangathastanza{\csromangathawak{\csromanfolio{121}ปจฺจกฺขํ ขฺวิมํ อวจ ตถาคโต ชิโน, \\ สมิชฺฌเต สีลวโต ยทิจฺฉติ. \\ ยถา ยถา เม มนสา วิจินฺติตํ, \\ ตถา สมิทฺธํ อยมนฺติโม ภโว”ติ.}}\csromangathastanza{\csromangathawak{อิทํ วาสนาภาคิยํ สุตฺตํ.}}\csromangathastanza{\csromangathawak{“เอกติํสมฺหิ กปฺปมฺหิ ชิโน อเนโช, \\ อนนฺตทสฺสี ภควา สิขีติ. \\ ตสฺสาปิ ราชา ภาตา สิขิทฺเธ\csromansharedfootnote{3624e2d9d47aa531}{สิขณฺฑิ (สี)}, \\ พุทฺเธ จ ธมฺเม จ อภิปฺปสนฺโน.}}\csromangathastanza{\csromangathawak{ปรินิพฺพุเต โลกวินายกมฺหิ, \\ ถูปํ สกาสิ วิปุลํ มหนฺตํ. \\ สมนฺตโต คาวุติกํ มเหสิโน, \\ เทวาติเทวสฺส นรุตฺตมสฺส.}}\csromangathastanza{\csromangathawak{ตสฺมิํ มนุสฺโส พลิมาภิหารี, \\ ปคฺคยฺห ชาติสุมนํ ปหฏฺโฐ. \\ วาเตน ปุปฺผํ ปติตสฺส เอกํ, \\ ตาหํ คเหตฺวาน ตสฺเสว ทาสิ.}}\csromangathastanza{\csromangathawak{โส มํ อโวจาภิปสนฺนจิตฺโต, \\ ตุยฺหเมว เอตํ ปุปฺผํ ททามิ. \\ ตาหํ คเหตฺวา อภิโรปเยสิํ, \\ ปุนปฺปุนํ พุทฺธมนุสฺสรนฺโต.}}\csromangathastanza{\csromangathabat{อชฺช ติํสํ ตโต กปฺปา,}{นาภิชานามิ ทุคฺคติํ.} \\ \csromangathabat{วินิปาตญฺจ น คจฺฉามิ,}{ถูปปูชายิทํ ผลนฺ”ติ.}}\csromangathastanza{\csromangathawak{อิทํ วาสนาภาคิยํ สุตฺตํ.}}\csromangathastanza{\csromangathabat{“กปิลํ นาม นครํ,}{สุวิภตฺตํ มหาปถํ.} \\ \csromangathabat{อากิณฺณมิทฺธํ ผีตญฺจ,}{พฺรหฺมทตฺตสฺส ราชิโน.}}\csromangathastanza{\csromanfolio{122}\csromangathabat{กุมฺมาสํ วิกฺกิณิํ ตตฺถ,}{ปญฺจาลานํ ปุรุตฺตเม.} \\ \csromangathabat{โสหํ อทฺทสิํ สมฺพุทฺธํ,}{อุปริฏฺฐํ ยสสฺสินํ.}}\csromangathastanza{\csromangathabat{หฏฺโฐ จิตฺตํ ปสาเทตฺวา,}{นิมนฺเตสิํ นรุตฺตมํ.} \\ \csromangathabat{อริฏฺฐํ ธุวภตฺเตน,}{ยํ เม เคหมฺหิ วิชฺชถ.}}\csromangathastanza{\csromangathabat{ตโต จ กตฺติโก ปุณฺโณ\csromansharedfootnote{e94dc28b39e530b1}{กตฺติกา ปุณฺณา (ก)},}{ปุณฺณมาสี อุปฏฺฐิตา.} \\ \csromangathabat{นวํ ทุสฺสยุคํ คยฺห,}{อริฏฺฐสฺโสปนามยิํ.}}\csromangathastanza{\csromangathabat{ปสนฺนจิตฺตํ ญตฺวาน,}{ปฏิคฺคณฺหิ นรุตฺตโม.} \\ \csromangathabat{อนุกมฺปโก การุณิโก,}{ตณฺหานิฆาตโก มุนิ.}}\csromangathastanza{\csromangathabat{ตาหํ กมฺมํ กริตฺวาน,}{กลฺยาณํ พุทฺธวณฺณิตํ.} \\ \csromangathabat{เทเว เจว มนุสฺเส จ,}{สนฺธาวิตฺวา ตโต จุโต.}}\csromangathastanza{\csromangathabat{พาราณสิยํ นคเร,}{เสฏฺฐิสฺส เอกปุตฺตโก.} \\ \csromangathabat{อฑฺฒ กุลสฺมิํ อุปฺปชฺชิํ,}{ปาเณหิ จ ปิยตโร.}}\csromangathastanza{\csromangathabat{ตโต จ วิญฺญุตํ ปตฺโต,}{เทวปุตฺเตน โจทิโต.} \\ \csromangathabat{ปาสาทา โอรูหิตฺวาน,}{สมฺพุทฺธมุปสงฺกมิํ.}}\csromangathastanza{\csromangathabat{โส เม ธมฺมมเทสยิ,}{อนุกมฺปาย โคตโม.} \\ \csromangathabat{ทุกฺขํ ทุกฺข\-สมุปฺปาทํ,}{ทุกฺขสฺส จ อติกฺกมํ.}}\csromangathastanza{\csromangathabat{อริยํ อฏฺฐงฺคิกํ มคฺคํ,}{ทุกฺขู\-ปสม\-คามินํ.} \\ \csromangathabat{จตฺตาริ อริยสจฺจานิ,}{มุนิ ธมฺมมเทสยิ.}}\csromangathastanza{\csromangathabat{ตสฺสาหํ วจนํ สุตฺวา,}{วิหริํ สาสเน รโต.} \\ \csromangathabat{สมถํ ปฏิวิชฺฌาหํ,}{รตฺตินฺทิวม\-ตนฺทิโต.}}\csromangathastanza{\csromangathabat{อชฺฌตฺตญฺจ พหิทฺธา จ,}{เย เม วิชฺชิํสุ\csromansharedfootnote{9750e3edeceaaadf}{วิชฺฌิํสุ (สี)} อาสวา.} \\ \csromangathabat{สพฺเพ อาสุํ สมุจฺฉินฺนา,}{น จ อุปฺปชฺชเร ปุน.}}\csromangathastanza{\csromangathabat{ปริยนฺตกตํ ทุกฺขํ,}{จริโมยํ สมุสฺสโย.} \\ \csromangathabat{ชาติ\-มรณ\-สํสาโร,}{นตฺถิทานิ ปุนพฺภโว”ติ.}}\csromangathastanza{\csromangathawak{อิทํ วาสนาภาคิย สุตฺตํ.}}}

\proseitem{99}{\csromanfolio{123}ตตฺถ กตมํ นิพฺเพธ\-ภาคิยํ สุตฺตํ.}

\prose{\csromansymbolfootnote{*}{เหฏฺฐา 54 ปิฏฺเฐ; อุปริ จ 184 ปิฏฺเฐปิ.}“อุทฺธํ อโธ สพฺพธิ วิปฺปมุตฺโต,}

\prose{อยํ อหสฺมีติ อนานุปสฺสี.}

\setlength{\csromangathapagewidth}{0pt}
\setlength{\csromangathawakwidth}{0pt}
\csromangathameasuregroup{}{เอวํ วิมุตฺโต อุทตาริ โอฆํ, \\ อติณฺณปุพฺพํ อปุนพฺภวายา”ติ. \\ อิทํ นิพฺเพธภาคิยํ สุตฺตํ.}
\csromangathameasurewak{เอวํ วิมุตฺโต อุทตาริ โอฆํ, \\ อติณฺณปุพฺพํ อปุนพฺภวายา”ติ. \\ อิทํ นิพฺเพธภาคิยํ สุตฺตํ.}
\setlength{\csromangathaleftcol}{0pt}
\csromangathagroup{\csromangathastanza{\csromangathawak{เอวํ วิมุตฺโต อุทตาริ โอฆํ, \\ อติณฺณปุพฺพํ อปุนพฺภวายา”ติ. \\ อิทํ นิพฺเพธ\-ภาคิยํ สุตฺตํ.}}}

\prose{“สีลวโต อานนฺท น เจตนา กรณียา “กินฺติ เม อวิปฺปฏิสาโร ชาเยยฺยา”ติ. ธมฺมตา เอสา อานนฺท ยํ สีลวโต อวิปฺปฏิสาโร ชาเยยฺย.  อวิปฺปฏิสารินา อานนฺท น เจตนา กรณียา “กินฺติ เม ปาโมชฺชํ ชาเยยฺยา”ติ. ธมฺมตา เอสา อานนฺท ยํ อวิปฺปฏิสาริโน ปาโมชฺชํ ชาเยยฺย.  ปมุทิเตน อานนฺท น เจตนา กรณียา “กินฺติ เม ปีติ ชาเยยฺยา”ติ. ธมฺมตา เอสา อานนฺท ยํ ปมุทิตสฺส ปีติ ชาเยยฺย.  ปีติมนสฺส อานนฺท น เจตนา กรณียา “กินฺติ เม กาโย ปสฺสมฺเภยฺยา”ติ. ธมฺมตา เอสา อานนฺท ยํ ปีติมนสฺส กาโย ปสฺสมฺเภยฺย.  ปสฺสทฺธ\-กายสฺส อานนฺท น เจตนา กรณียา “กินฺตาหํ สุขํ เวทิเยยฺยนฺ”ติ. ธมฺมตา เอสา อานนฺท ยํ ปสฺสทฺธกาโย สุขํ เวทิเยยฺย.  สุขิโน อานนฺท น เจตนา กรณียา “กินฺติ เม สมาธิ ชาเยยฺยา”ติ. ธมฺมตา เอสา อานนฺท ยํ สุขิโน สมาธิ ชาเยยฺย.  สมาหิตสฺส อานนฺท น เจตนา กรณียา “กินฺตาหํ ยถาภูตํ ปชาเนยฺยนฺ”ติ. ธมฺมตา เอสา อานนฺท ยํ สมาหิโต ยถาภูตํ ปชาเนยฺย.  ยถาภูตํ ปชานตา อานนฺท น เจตนา กรณียา “กินฺติ เม นิพฺพิทา ชาเยยฺยา”ติ. ธมฺมตา เอสา อานนฺท ยํ ยถาภูตํ ปชานนฺโต นิพฺพินฺเตยฺย.  นิพฺพินฺทนฺเตน อานนฺท น เจตนา กรณียา “กินฺติ เม วิราโค ชาเยยฺยา”ติ. ธมฺมตา เอสา อานนฺท ยํ นิพฺพินฺทนฺโต วิรชฺเชยฺย.  วิรชฺชนฺเตน อานนฺท น เจตนา กรณียา “กินฺติ เม วิมุตฺติ ชาเยยฺยา”ติ. ธมฺมตา เอสา อานนฺท ยํ วิรชฺชนฺโต วิมุจฺเจยฺย.  วิมุตฺเตน อานนฺท น เจตนา กรณียา “กินฺติ เม วิมุตฺติ\-ญาณ\-ทสฺสนํ อุปฺปชฺเชยฺยา”ติ. ธมฺมตา เอสา อานนฺท ยํ วิมุตฺตสฺส วิมุตฺติ\-ญาณ\-ทสฺสนํ อุปฺปชฺเชยฺยา”ติ\csromansharedfootnote{14649d1460b1b062}{อํ 3. 516, 7 ปิฏฺเฐสุ ปสฺสิตพฺพํ.}. อิทํ นิพฺเพธ\-ภาคิยํ สุตฺตํ.}

\setlength{\csromangathapagewidth}{0pt}
\setlength{\csromangathawakwidth}{0pt}
\csromangathameasuregroup{\textsuperscript{*}“ยทา หเว ปาตุภวนฺติ ธมฺมา,}{\csromangathabat{\textsuperscript{*}“ยทา หเว ปาตุภวนฺติ ธมฺมา,}{อาตาปิโน ฌายโต พฺราหฺมณสฺส.} \\ อถสฺส กงฺขา วปยนฺติ สพฺพา. \\ ยโต ปชานาติ สเหตุธมฺมนฺ”ติ. \\ อิทํ นิพฺเพธภาคิยํ สุตฺตํ. \\ \textsuperscript{*}“ยทา หเว ปาตุภวนฺติ ธมฺมา, \\ อาตาปิโน ฌายโต พฺราหฺมณสฺส. \\ อถสฺส กงฺขา วปยนฺติ สพฺพา, \\ ยโต ขยํ ปจฺจยานํ อเวที”ติ. \\ อิทํ นิพฺเพธภาคิยํ สุตฺตํ. \\ “กิํนุ กุชฺฌสิ มา กุชฺฌิ, \\ อกฺโกโธ ติสฺส เต วรํ. \\ โกธมานมกฺขวินยตฺถํ หิ, \\ อิทํ นิพฺเพธภาคิยํ สุตฺตํ.}
\csromangathameasurewak{อิทํ นิพฺเพธภาคิยํ สุตฺตํ. \\ \textsuperscript{*}“ยทา หเว ปาตุภวนฺติ ธมฺมา, \\ อาตาปิโน ฌายโต พฺราหฺมณสฺส. \\ อถสฺส กงฺขา วปยนฺติ สพฺพา, \\ ยโต ขยํ ปจฺจยานํ อเวที”ติ. \\ อิทํ นิพฺเพธภาคิยํ สุตฺตํ. \\ “กิํนุ กุชฺฌสิ มา กุชฺฌิ, \\ อกฺโกโธ ติสฺส เต วรํ. \\ โกธมานมกฺขวินยตฺถํ หิ, \\ อิทํ นิพฺเพธภาคิยํ สุตฺตํ.}
\csromangathasetleft{\textsuperscript{*}“ยทา หเว ปาตุภวนฺติ ธมฺมา,}
\csromangathagroup{\csromangathastanza{\csromanfolio{124}\csromangathabat{\csromansymbolmark{*}“ยทา หเว ปาตุภวนฺติ ธมฺมา,}{อาตาปิโน ฌายโต พฺราหฺมณสฺส.} \\ อถสฺส กงฺขา วปยนฺติ สพฺพา. \\ ยโต ปชานาติ สเหตุธมฺมนฺ”ติ.}\csromangathastanza{\csromangathawak{อิทํ นิพฺเพธ\-ภาคิยํ สุตฺตํ.}}\csromangathastanza{\csromangathawak{\csromansymbolmark{*}“ยทา หเว ปาตุภวนฺติ ธมฺมา, \\ อาตาปิโน ฌายโต พฺราหฺมณสฺส. \\ อถสฺส กงฺขา วปยนฺติ สพฺพา, \\ ยโต ขยํ ปจฺจยานํ อเวที”ติ.}}\csromangathastanza{\csromangathawak{อิทํ นิพฺเพธ\-ภาคิยํ สุตฺตํ.}}\csromangathastanza{\csromangathawak{“กิํนุ กุชฺฌสิ มา กุชฺฌิ, \\ อกฺโกโธ ติสฺส เต วรํ. \\ โกธมานมกฺขวินยตฺถํ หิ,}}\csromangathastanza{\csromangathawak{อิทํ นิพฺเพธ\-ภาคิยํ สุตฺตํ.}}}

\prose{ติสฺส พฺรหฺมจริยํ วุสฺสตี”ติ\csromansharedfootnote{26c8b7e2a2eadf85}{สํ 1. 469 ปิฏฺเฐ.}.}

\setlength{\csromangathapagewidth}{0pt}
\setlength{\csromangathawakwidth}{0pt}
\csromangathameasuregroup{“กทาหํ นนฺทํ ปสฺเสยฺยํ,}{\csromangathabat{“กทาหํ นนฺทํ ปสฺเสยฺยํ,}{อารญฺญํ ปํสุกูลิกํ.} \\ อิทํ นิพฺเพธภาคิยํ สุตฺตํ.}
\csromangathameasurewak{อิทํ นิพฺเพธภาคิยํ สุตฺตํ.}
\csromangathasetleft{“กทาหํ นนฺทํ ปสฺเสยฺยํ,}
\csromangathagroup{\csromangathastanza{\csromangathabat{“กทาหํ นนฺทํ ปสฺเสยฺยํ,}{อารญฺญํ ปํสุกูลิกํ.}}\csromangathastanza{\csromangathawak{อิทํ นิพฺเพธ\-ภาคิยํ สุตฺตํ.}}}

\prose{อญฺญาตุญฺเฉน ยาเปนฺตํ, กาเมสุ อนเปกฺขินนฺ”ติ\csromansharedfootnote{0a4880e63afaadab}{สํ 1. 468 ปิฏฺเฐ.}.}

\setlength{\csromangathapagewidth}{0pt}
\setlength{\csromangathawakwidth}{0pt}
\csromangathameasuregroup{“กิํสุ เฉตฺวา สุขํ เสติ, \\ กิสฺสสฺสุ\textsuperscript{1} เอกธมฺมสฺส, \\ โกธํ เฉตฺวา สุขํ เสติ, \\ โกธสฺส วิสมูลสฺส,}{\csromangathabat{“กิํสุ เฉตฺวา สุขํ เสติ,}{กิํสุ เฉตฺวา น โสจติ.} \\ \csromangathabat{กิสฺสสฺสุ\textsuperscript{1} เอกธมฺมสฺส,}{วธํ โรเจสิ โคตมาติ.} \\ \csromangathabat{โกธํ เฉตฺวา สุขํ เสติ,}{โกธํ เฉตฺวา น โสจติ.} \\ \csromangathabat{โกธสฺส วิสมูลสฺส,}{มธุรคฺคสฺส พฺราหฺมณ.} \\ อิทํ นิพฺเพธภาคิยํ สุตฺตํ. \\ “กิํสุ หเน อุปฺปติตํ, \\ กิํสุ ชาตํ วิโนทเย.}
\csromangathameasurewak{อิทํ นิพฺเพธภาคิยํ สุตฺตํ. \\ “กิํสุ หเน อุปฺปติตํ, \\ กิํสุ ชาตํ วิโนทเย.}
\csromangathasetleft{“กิํสุ เฉตฺวา สุขํ เสติ, \\ กิสฺสสฺสุ\textsuperscript{1} เอกธมฺมสฺส, \\ โกธํ เฉตฺวา สุขํ เสติ, \\ โกธสฺส วิสมูลสฺส,}
\csromangathagroup{\csromangathastanza{\csromangathabat{“กิํสุ เฉตฺวา สุขํ เสติ,}{กิํสุ เฉตฺวา น โสจติ.} \\ \csromangathabat{กิสฺสสฺสุ\csromansharedfootnote{45065ab5b6868253}{กิสฺสสฺส (สี, ก)} เอกธมฺมสฺส,}{วธํ โรเจสิ โคตมาติ.}}\csromangathastanza{\csromangathabat{โกธํ เฉตฺวา สุขํ เสติ,}{โกธํ เฉตฺวา น โสจติ.} \\ \csromangathabat{โกธสฺส วิสมูลสฺส,}{มธุรคฺคสฺส พฺราหฺมณ.}}\csromangathastanza{\csromangathawak{อิทํ นิพฺเพธ\-ภาคิยํ สุตฺตํ. \\ “กิํสุ หเน อุปฺปติตํ, \\ กิํสุ ชาตํ วิโนทเย.}}}

\prose{วธํ อริยา ปสํสนฺติ, ตํ หิ เฉตฺวา น โสจตี”ติ\csromansharedfootnote{9d04f10a4815bc3f}{สํ 1. 162 ปิฏฺเฐ.}.}

\setlength{\csromangathapagewidth}{0pt}
\setlength{\csromangathawakwidth}{0pt}
\csromangathameasuregroup{กิญฺจสฺสุ ปชเห ธีโร, \\ โกธํ หเน อุปฺปติตํ, \\ อวิชฺชํ ปชเห ธีโร, \\ กามราคปฺปหานาย, \\ สตฺติยา วิย โอมฏฺโฐ, \\ สกฺกายทิฏฺฐิปฺปหานาย,}{\csromangathabat{กิญฺจสฺสุ ปชเห ธีโร,}{กิสฺสาภิสมโย สุโข.} \\ \csromangathabat{โกธํ หเน อุปฺปติตํ,}{ราคํ ชาตํ วิโนทเย.} \\ \csromangathabat{อวิชฺชํ ปชเห ธีโร,}{สจฺจาภิสมโย สุโข”ติ.} \\ อิทํ นิพฺเพธภาคิยํ สุตฺตํ. \\ \textsuperscript{*}“สตฺติยา วิย โอมฏฺโฐ, \\ ฑยฺหมาโนว\textsuperscript{1} มตฺถเก. \\ \csromangathabat{กามราคปฺปหานาย,}{สโต ภิกฺขุ ปริพฺพเช.} \\ \csromangathabat{สตฺติยา วิย โอมฏฺโฐ,}{ฑยฺหมาโนว มตฺถเก.} \\ \csromangathabat{สกฺกายทิฏฺฐิปฺปหานาย,}{สโต ภิกฺขุ ปริพฺพเช”ติ.} \\ อิทํ นิพฺเพธภาคิยํ สุตฺตํ.}
\csromangathasetleft{กิญฺจสฺสุ ปชเห ธีโร, \\ โกธํ หเน อุปฺปติตํ, \\ อวิชฺชํ ปชเห ธีโร, \\ กามราคปฺปหานาย, \\ สตฺติยา วิย โอมฏฺโฐ, \\ สกฺกายทิฏฺฐิปฺปหานาย,}
\csromangathagroup{\csromangathastanza{\csromangathabat{กิญฺจสฺสุ ปชเห ธีโร,}{กิสฺสาภิสมโย สุโข.} \\ \csromangathabat{โกธํ หเน อุปฺปติตํ,}{ราคํ ชาตํ วิโนทเย.}}\csromangathastanza{\csromanfolio{125}\csromangathabat{อวิชฺชํ ปชเห ธีโร,}{สจฺจาภิสมโย สุโข”ติ.} \\ อิทํ นิพฺเพธ\-ภาคิยํ สุตฺตํ. \\ \csromansymbolmark{*}“สตฺติยา วิย โอมฏฺโฐ,}\csromangathastanza{ฑยฺหมาโนว\csromansharedfootnote{338dd9d64fd7ee1e}{[มิสฺสินฺคฺ โนเต 0]} มตฺถเก.}\csromangathastanza{\csromangathabat{กาม\-ราค\-ปฺปหานาย,}{สโต ภิกฺขุ ปริพฺพเช.} \\ \csromangathabat{สตฺติยา วิย โอมฏฺโฐ,}{ฑยฺหมาโนว มตฺถเก.}}\csromangathastanza{\csromangathabat{สกฺกายทิฏฺฐิ\-ปฺปหานาย,}{สโต ภิกฺขุ ปริพฺพเช”ติ.} \\ อิทํ นิพฺเพธ\-ภาคิยํ สุตฺตํ.}}

\prose{** “สพฺเพ ขยนฺตา นิจยา, ปตนนฺตา สมุสฺสยา.}

\setlength{\csromangathapagewidth}{0pt}
\setlength{\csromangathawakwidth}{0pt}
\csromangathameasuregroup{สพฺเพ ขยนฺตา นิจยา, \\ สพฺเพสํ มรณมาคมฺม, \\ \textsuperscript{+}เอตํ ภยํ มรเณ เปกฺขมาโน, \\ เยสํ ฌานรตํ จิตฺตํ, \\ อารทฺธวีริโย ปหิตตฺโต, \\ วิรโต กามสญฺญาย,}{สพฺเพสํ มรณมาคมฺม, \\ สพฺเพสํ ชีวิตมทฺธุวํ. \\ \textsuperscript{+}เอตํ ภยํ มรเณ\textsuperscript{1} เปกฺขมาโน, \\ ปุญฺญานิ กยิราถ สุขาวหานิ. \\ \csromangathabat{สพฺเพ ขยนฺตา นิจยา,}{ปตนนฺตา สมุสฺสยา.} \\ \csromangathabat{สพฺเพสํ มรณมาคมฺม,}{สพฺเพสํ ชีวิตมทฺธุวํ.} \\ \csromangathabat{\textsuperscript{+}เอตํ ภยํ มรเณ เปกฺขมาโน,}{โลกามิสํ ปชเห สนฺติเปกฺโข”ติ.} \\ อิทํ นิพฺเพธภาคิยํ สุตฺตํ. \\ “สุขํ สยนฺติ มุนโย, \\ น เต โสจนฺติ มาวิธ. \\ \csromangathabat{เยสํ ฌานรตํ จิตฺตํ,}{ปญฺญวา สุสมาหิโต.} \\ \csromangathabat{อารทฺธวีริโย ปหิตตฺโต,}{โอฆํ ตรติ ทุตฺตรํ.} \\ \csromangathabat{วิรโต กามสญฺญาย,}{สพฺพสํโยชนาตีโต\textsuperscript{1}.} \\ อิทํ นิพฺเพธภาคิยํ สุตฺตํ.}
\csromangathameasurewak{สพฺเพสํ มรณมาคมฺม, \\ สพฺเพสํ ชีวิตมทฺธุวํ. \\ \textsuperscript{+}เอตํ ภยํ มรเณ\textsuperscript{1} เปกฺขมาโน, \\ ปุญฺญานิ กยิราถ สุขาวหานิ. \\ น เต โสจนฺติ มาวิธ. \\ อิทํ นิพฺเพธภาคิยํ สุตฺตํ.}
\csromangathasetleft{สพฺเพ ขยนฺตา นิจยา, \\ สพฺเพสํ มรณมาคมฺม, \\ \textsuperscript{+}เอตํ ภยํ มรเณ เปกฺขมาโน, \\ เยสํ ฌานรตํ จิตฺตํ, \\ อารทฺธวีริโย ปหิตตฺโต, \\ วิรโต กามสญฺญาย,}
\csromangathagroup{\csromangathastanza{\csromangathawak{สพฺเพสํ มรณมาคมฺม, \\ สพฺเพสํ ชีวิตมทฺธุวํ. \\ \csromansymbolfootnote{+}{สํ 1. 53 ปิฏฺเฐ.}เอตํ ภยํ มรเณ\csromansharedfootnote{688fefc3c1776ed3}{มรณํ (ก)} เปกฺขมาโน, \\ ปุญฺญานิ กยิราถ สุขาวหานิ.}}\csromangathastanza{\csromangathabat{สพฺเพ ขยนฺตา นิจยา,}{ปตนนฺตา สมุสฺสยา.} \\ \csromangathabat{สพฺเพสํ มรณมาคมฺม,}{สพฺเพสํ ชีวิตมทฺธุวํ.}}\csromangathastanza{\csromangathabat{\csromansymbolmark{+}เอตํ ภยํ มรเณ เปกฺขมาโน,}{โลกามิสํ ปชเห สนฺติเปกฺโข”ติ.} \\ อิทํ นิพฺเพธ\-ภาคิยํ สุตฺตํ. \\ “สุขํ สยนฺติ มุนโย,}\csromangathastanza{\csromangathawak{น เต โสจนฺติ มาวิธ.}}\csromangathastanza{\csromangathabat{เยสํ ฌานรตํ จิตฺตํ,}{ปญฺญวา สุสมาหิโต.} \\ \csromangathabat{อารทฺธวีริโย ปหิตตฺโต,}{โอฆํ ตรติ ทุตฺตรํ.} \\ \csromangathabat{วิรโต กามสญฺญาย,}{สพฺพสํโยชนาตีโต\csromansharedfootnote{338dd9d64fd7ee1e}{[มิสฺสินฺคฺ โนเต 0]}.}}\csromangathastanza{\csromangathawak{อิทํ นิพฺเพธ\-ภาคิยํ สุตฺตํ.}}}

\prose{นนฺทิภวปริกฺขีโณ\csromansharedfootnote{9e2591f223fe423c}{สพฺพสํโยชนาติโค (สี) ** สํ 1. 12 ปิฏฺเฐ.}, โส คมฺภีเร น สีทตี”ติ\csromansharedfootnote{9e2591f223fe423c}{สพฺพสํโยชนาติโค (สี) ** สํ 1. 12 ปิฏฺเฐ.}.}

\setlength{\csromangathapagewidth}{0pt}
\setlength{\csromangathawakwidth}{0pt}
\csromangathameasuregroup{“สทฺทหาโน อรหตํ, \\ สุสฺสูสํ ลภเต ปญฺญํ, \\ \textsuperscript{*}ปติรูปการี ธุรวา, \\ สจฺเจน กิตฺติํ ปปฺโปติ, \\ อสฺมา โลกา ปรํ โลกํ,}{\csromangathabat{“สทฺทหาโน อรหตํ,}{ธมฺมํ นิพฺพานปตฺติยา.} \\ \csromangathabat{สุสฺสูสํ ลภเต ปญฺญํ,}{อปฺปมตฺโต วิจกฺขโณ\textsuperscript{1}.} \\ \csromangathabat{\textsuperscript{*}ปติรูปการี ธุรวา,}{อุฏฺฐาตา วินฺทเต ธนํ.} \\ \csromangathabat{สจฺเจน กิตฺติํ ปปฺโปติ,}{ททํ มิตฺตานิ คนฺถติ.} \\ \csromangathabat{อสฺมา โลกา ปรํ โลกํ,}{เอวํ\textsuperscript{1} เปจฺจ น โสจตี”ติ.} \\ อิทํ นิพฺเพธภาคิยํ สุตฺตํ.}
\csromangathasetleft{“สทฺทหาโน อรหตํ, \\ สุสฺสูสํ ลภเต ปญฺญํ, \\ \textsuperscript{*}ปติรูปการี ธุรวา, \\ สจฺเจน กิตฺติํ ปปฺโปติ, \\ อสฺมา โลกา ปรํ โลกํ,}
\csromangathagroup{\csromangathastanza{\csromangathabat{“สทฺทหาโน อรหตํ,}{ธมฺมํ นิพฺพานปตฺติยา.} \\ \csromangathabat{สุสฺสูสํ ลภเต ปญฺญํ,}{อปฺปมตฺโต วิจกฺขโณ\csromansharedfootnote{731e0c7c41dc9270}{สํ 1. 217 ปิฏฺเฐ.}.}}\csromangathastanza{\csromanfolio{126}\csromangathabat{\csromansymbolmark{*}ปติรูปการี ธุรวา,}{อุฏฺฐาตา วินฺทเต ธนํ.} \\ \csromangathabat{สจฺเจน กิตฺติํ ปปฺโปติ,}{ททํ มิตฺตานิ คนฺถติ.}}\csromangathastanza{\csromangathabat{อสฺมา โลกา ปรํ โลกํ,}{เอวํ\csromansharedfootnote{338dd9d64fd7ee1e}{[มิสฺสินฺคฺ โนเต 0]} เปจฺจ น โสจตี”ติ.} \\ อิทํ นิพฺเพธ\-ภาคิยํ สุตฺตํ.}}

\proseitem{101}{** “สพฺพคนฺถ\-ปหีนสฺส, วิปฺปมุตฺตสฺส เต สโต.}

\setlength{\csromangathapagewidth}{0pt}
\setlength{\csromangathawakwidth}{0pt}
\csromangathameasuregroup{สมณสฺส น ตํ สาธุ, \\ เยน เกนจิ วณฺเณน, \\ น ตํ อรหติ สปฺปญฺโญ, \\ มนสา เจ ปสนฺเนน, \\ น เตน โหติ สํยุตฺโต,}{\csromangathabat{สมณสฺส น ตํ สาธุ,}{ยทญฺญมนุสาสสีติ.} \\ \csromangathabat{เยน เกนจิ วณฺเณน,}{สํวาโส สกฺก ชายติ.} \\ \csromangathabat{น ตํ อรหติ สปฺปญฺโญ,}{มนสา อนุกมฺปิตุํ\textsuperscript{1}.} \\ \csromangathabat{มนสา เจ ปสนฺเนน,}{ยทญฺญมนุสาสติ.} \\ \csromangathabat{น เตน โหติ สํยุตฺโต,}{ยานุกมฺปา อนุทฺทยา”ติ.} \\ อิทํ นิพฺเพธภาคิยํ สุตฺตํ.}
\csromangathasetleft{สมณสฺส น ตํ สาธุ, \\ เยน เกนจิ วณฺเณน, \\ น ตํ อรหติ สปฺปญฺโญ, \\ มนสา เจ ปสนฺเนน, \\ น เตน โหติ สํยุตฺโต,}
\csromangathagroup{\csromangathastanza{\csromangathabat{สมณสฺส น ตํ สาธุ,}{ยทญฺญมนุสาสสีติ.} \\ \csromangathabat{เยน เกนจิ วณฺเณน,}{สํวาโส สกฺก ชายติ.}}\csromangathastanza{\csromangathabat{น ตํ อรหติ สปฺปญฺโญ,}{มนสา อนุกมฺปิตุํ\csromansharedfootnote{338dd9d64fd7ee1e}{[มิสฺสินฺคฺ โนเต 0]}.} \\ \csromangathabat{มนสา เจ ปสนฺเนน,}{ยทญฺญมนุสาสติ.}}\csromangathastanza{\csromangathabat{น เตน โหติ สํยุตฺโต,}{ยานุกมฺปา อนุทฺทยา”ติ.} \\ อิทํ นิพฺเพธ\-ภาคิยํ สุตฺตํ.}}

\proseitem{102}{“ราโค จ โทโส จ กุโตนิทานา,}

\prose{อรตี รตี\csromansharedfootnote{9502569c4dc6a6ae}{อรติ รติ (ก)} โลมหํโส กุโตชา.}

\setlength{\csromangathapagewidth}{0pt}
\setlength{\csromangathawakwidth}{0pt}
\csromangathameasuregroup{}{กุโต สมุฏฺฐาย มโนวิตกฺกา, \\ กุมารกา ธงฺกมิโวสชนฺติ. \\ ราโค จ โทโส จ อิโตนิทานา, \\ อรตี รตี โลมหํโส อิโตชา. \\ อิโต สมุฏฺฐาย มโนวิตกฺกา, \\ กุมารกา ธงฺกมิโวสชนฺติ. \\ สฺเนหชา อตฺตสมฺภูตา, \\ นิโคฺรธสฺเสว ขนฺธชา. \\ ปุถุ วิสตฺตา กาเมสุ, \\ มาลุวาว วิตตา วเน. \\ เย นํ ปชานนฺติ ยโตนิทานํ, \\ เต นํ วิโนเทนฺติ สุโณหิ ยกฺข. \\ เต ทุตฺตรํ โอฆมิมํ ตรนฺติ, \\ อิทํ นิพฺเพธภาคิยํ สุตฺตํ.}
\csromangathameasurewak{กุโต สมุฏฺฐาย มโนวิตกฺกา, \\ กุมารกา ธงฺกมิโวสชนฺติ. \\ ราโค จ โทโส จ อิโตนิทานา, \\ อรตี รตี โลมหํโส อิโตชา. \\ อิโต สมุฏฺฐาย มโนวิตกฺกา, \\ กุมารกา ธงฺกมิโวสชนฺติ. \\ สฺเนหชา อตฺตสมฺภูตา, \\ นิโคฺรธสฺเสว ขนฺธชา. \\ ปุถุ วิสตฺตา กาเมสุ, \\ มาลุวาว วิตตา วเน. \\ เย นํ ปชานนฺติ ยโตนิทานํ, \\ เต นํ วิโนเทนฺติ สุโณหิ ยกฺข. \\ เต ทุตฺตรํ โอฆมิมํ ตรนฺติ, \\ อิทํ นิพฺเพธภาคิยํ สุตฺตํ.}
\setlength{\csromangathaleftcol}{0pt}
\csromangathagroup{\csromangathastanza{\csromangathawak{กุโต สมุฏฺฐาย มโนวิตกฺกา, \\ กุมารกา ธงฺกมิโวสชนฺติ. \\ ราโค จ โทโส จ อิโตนิทานา, \\ อรตี รตี โลมหํโส อิโตชา.}}\csromangathastanza{\csromangathawak{อิโต สมุฏฺฐาย มโนวิตกฺกา, \\ กุมารกา ธงฺกมิโวสชนฺติ. \\ สฺเนหชา อตฺตสมฺภูตา, \\ นิโคฺรธสฺเสว ขนฺธชา.}}\csromangathastanza{\csromangathawak{ปุถุ วิสตฺตา กาเมสุ, \\ มาลุวาว วิตตา วเน. \\ เย นํ ปชานนฺติ ยโตนิทานํ, \\ เต นํ วิโนเทนฺติ สุโณหิ ยกฺข.}}\csromangathastanza{\csromangathawak{เต ทุตฺตรํ โอฆมิมํ ตรนฺติ,}}\csromangathastanza{\csromangathawak{อิทํ นิพฺเพธ\-ภาคิยํ สุตฺตํ.}}}

\prose{อติณฺณปุพฺพํ อปุนพฺภวายา”ติ\csromansharedfootnote{1d63869ddb1a76a5}{สํ 1. 209 ปิฏฺเฐ; ขุ 1. 320 ปิฏฺเฐ จ.}.}

\proseitem{102}{\csromanfolio{127}\csromansymbolfootnote{*}{สํ 1. 45, 46 ปิฏฺเฐสุ.}“ทุกฺกรํ ภควา สุทุกฺกรํ ภควา”ติ.}

\prose{“ทุกฺกรํ วาปิ กโรนฺติ, (กามทาติ ภควา)}

\prose{เสกฺขา สีลสมาหิตา.}

\setlength{\csromangathapagewidth}{0pt}
\setlength{\csromangathawakwidth}{0pt}
\csromangathameasuregroup{}{ฐิตตฺตา อนคาริยุเปตสฺส, \\ ตุฏฺฐิ โหติ สุขาวหา”ติ.}
\csromangathameasurewak{ฐิตตฺตา อนคาริยุเปตสฺส, \\ ตุฏฺฐิ โหติ สุขาวหา”ติ.}
\setlength{\csromangathaleftcol}{0pt}
\csromangathagroup{\csromangathastanza{\csromangathawak{ฐิตตฺตา อนคาริยุ\-เปตสฺส, \\ ตุฏฺฐิ โหติ สุขาวหา”ติ.}}}

\prose{“ทุลฺลภา\csromansharedfootnote{7e3a75dfe4f0e155}{ทุลฺลภํ (สี, ก)} ภควา ยทิทํ ตุฏฺฐี”ติ.}

\prose{“ทุลฺลภํ วาปิ ลภนฺติ, (กามทาติ ภควา)}

\prose{จิตฺตวูปสเม รตา.}

\setlength{\csromangathapagewidth}{0pt}
\setlength{\csromangathawakwidth}{0pt}
\csromangathameasuregroup{}{เยสํ ทิวา จ รตฺโต จ, \\ ภาวนาย รโต มโน”ติ.}
\csromangathameasurewak{เยสํ ทิวา จ รตฺโต จ, \\ ภาวนาย รโต มโน”ติ.}
\setlength{\csromangathaleftcol}{0pt}
\csromangathagroup{\csromangathastanza{\csromangathawak{เยสํ ทิวา จ รตฺโต จ, \\ ภาวนาย รโต มโน”ติ.}}}

\prose{“ทุสฺสมาทหํ ภควา ยทิทํ จิตฺตนฺ”ติ.}

\vspace{\tipitakaparskip}
\csromancenter{“ทุสฺสมาทหํ วาปิ สมาทหนฺติ, (กามทาติ ภควา)}
\prose{อินฺทฺริยูปสเม รตา.}

\setlength{\csromangathapagewidth}{0pt}
\setlength{\csromangathawakwidth}{0pt}
\csromangathameasuregroup{}{เต เฉตฺวา มจฺจุโน ชาลํ, \\ อริยา คจฺฉนฺติ กามทา”ติ.}
\csromangathameasurewak{เต เฉตฺวา มจฺจุโน ชาลํ, \\ อริยา คจฺฉนฺติ กามทา”ติ.}
\setlength{\csromangathaleftcol}{0pt}
\csromangathagroup{\csromangathastanza{\csromangathawak{เต เฉตฺวา มจฺจุโน ชาลํ, \\ อริยา คจฺฉนฺติ กามทา”ติ.}}}

\prose{“ทุคฺคโม ภควา วิสโม มคฺโค”ติ.}

\setlength{\csromangathapagewidth}{0pt}
\setlength{\csromangathawakwidth}{0pt}
\csromangathameasuregroup{“ทุคฺคเม วิสเม วาปิ, \\ อนริยา วิสเม มคฺเค, \\ อริยานํ สโม มคฺโค,}{\csromangathabat{“ทุคฺคเม วิสเม วาปิ,}{อริยา คจฺฉนฺติ กามท\textsuperscript{1}.} \\ \csromangathabat{อนริยา วิสเม มคฺเค,}{ปปตนฺติ อวํสิรา.} \\ \csromangathabat{อริยานํ สโม มคฺโค,}{อริยา หิ วิสเม สมา”ติ.} \\ อิทํ นิพฺเพธภาคิยํ สุตฺตํ.}
\csromangathasetleft{“ทุคฺคเม วิสเม วาปิ, \\ อนริยา วิสเม มคฺเค, \\ อริยานํ สโม มคฺโค,}
\csromangathagroup{\csromangathastanza{\csromangathabat{“ทุคฺคเม วิสเม วาปิ,}{อริยา คจฺฉนฺติ กามท\csromansharedfootnote{338dd9d64fd7ee1e}{[มิสฺสินฺคฺ โนเต 0]}.} \\ \csromangathabat{อนริยา วิสเม มคฺเค,}{ปปตนฺติ อวํสิรา.}}\csromangathastanza{\csromangathabat{อริยานํ สโม มคฺโค,}{อริยา หิ วิสเม สมา”ติ.} \\ อิทํ นิพฺเพธ\-ภาคิยํ สุตฺตํ.}}

\proseitem{103}{\csromansymbolmark{*}* “อิทํ หิ ตํ เชตวนํ, อิสิสํฆ\-นิเสวิตํ.}

\setlength{\csromangathapagewidth}{0pt}
\setlength{\csromangathawakwidth}{0pt}
\csromangathameasuregroup{อาวุตฺถํ ธมฺมราเชน, \\ กมฺมํ วิชฺชา จ ธมฺโม จ, \\ เอเตน มจฺจา สุชฺฌนฺติ, \\ ตสฺมา หิ ปณฺฑิโต โปโส, \\ โยนิโส วิจิเน ธมฺมํ, \\ สาริปุตฺโตว ปญฺญาย, \\ โยปิ ปารงฺคโต ภิกฺขุ, \\ ยทตีตํ ปหีนํ\textsuperscript{1} ตํ, \\ ปจฺจุปฺปนฺนญฺจ โย ธมฺมํ\textsuperscript{1}, \\ อสํหีรํ อสํกุปฺปํ, \\ อชฺเชว กิจฺจมาตปฺปํ\textsuperscript{1}, \\ น หิ โน สงฺครํ เตน, \\ เอวํ วิหาริํ อาตาปิํ,}{\csromangathabat{อาวุตฺถํ ธมฺมราเชน,}{ปีติสญฺชนนํ มม.} \\ \csromangathabat{กมฺมํ วิชฺชา จ ธมฺโม จ,}{สีลํ ชีวิตมุตฺตมํ.} \\ \csromangathabat{เอเตน มจฺจา สุชฺฌนฺติ,}{น โคตฺเตน ธเนน วา.} \\ \csromangathabat{ตสฺมา หิ ปณฺฑิโต โปโส,}{สมฺปสฺสํ อตฺถมตฺตโน.} \\ \csromangathabat{โยนิโส วิจิเน ธมฺมํ,}{เอวํ ตตฺถ วิสุชฺฌติ.} \\ \csromangathabat{สาริปุตฺโตว ปญฺญาย,}{สีเลน อุปสเมน จ.} \\ \csromangathabat{โยปิ ปารงฺคโต ภิกฺขุ,}{เอตาวปรโม สิยา”ติ.} \\ อิทํ นิพฺเพธภาคิยํ สุตฺตํ. \\ \textsuperscript{*}“อตีตํ นานฺวาคเมยฺย, \\ นปฺปฏิกงฺเข อนาคตํ. \\ \csromangathabat{ยทตีตํ ปหีนํ\textsuperscript{1} ตํ,}{อปฺปตฺตญฺจ อนาคตํ.} \\ \csromangathabat{ปจฺจุปฺปนฺนญฺจ โย ธมฺมํ\textsuperscript{1},}{ตตฺถ ตตฺถ วิปสฺสติ.} \\ \csromangathabat{อสํหีรํ อสํกุปฺปํ,}{ตํ วิทฺวา มนุพฺรูหเย.} \\ \csromangathabat{อชฺเชว กิจฺจมาตปฺปํ\textsuperscript{1},}{โก ชญฺญา มรณํ สุเว.} \\ \csromangathabat{น หิ โน สงฺครํ เตน,}{มหาเสเนน มจฺจุนา.} \\ \csromangathabat{เอวํ วิหาริํ อาตาปิํ,}{อโหรตฺตมตนฺทิตํ.} \\ อิทํ นิพฺเพธภาคิยํ สุตฺตํ.}
\csromangathameasurewak{อิทํ นิพฺเพธภาคิยํ สุตฺตํ.}
\csromangathasetleft{อาวุตฺถํ ธมฺมราเชน, \\ กมฺมํ วิชฺชา จ ธมฺโม จ, \\ เอเตน มจฺจา สุชฺฌนฺติ, \\ ตสฺมา หิ ปณฺฑิโต โปโส, \\ โยนิโส วิจิเน ธมฺมํ, \\ สาริปุตฺโตว ปญฺญาย, \\ โยปิ ปารงฺคโต ภิกฺขุ, \\ ยทตีตํ ปหีนํ\textsuperscript{1} ตํ, \\ ปจฺจุปฺปนฺนญฺจ โย ธมฺมํ\textsuperscript{1}, \\ อสํหีรํ อสํกุปฺปํ, \\ อชฺเชว กิจฺจมาตปฺปํ\textsuperscript{1}, \\ น หิ โน สงฺครํ เตน, \\ เอวํ วิหาริํ อาตาปิํ,}
\csromangathagroup{\csromangathastanza{\csromangathabat{อาวุตฺถํ ธมฺมราเชน,}{ปีติสญฺชนนํ มม.} \\ \csromangathabat{กมฺมํ วิชฺชา จ ธมฺโม จ,}{สีลํ ชีวิตมุตฺตมํ.}}\csromangathastanza{\csromangathabat{เอเตน มจฺจา สุชฺฌนฺติ,}{น โคตฺเตน ธเนน วา.} \\ \csromangathabat{ตสฺมา หิ ปณฺฑิโต โปโส,}{สมฺปสฺสํ อตฺถมตฺตโน.}}\csromangathastanza{\csromangathabat{โยนิโส วิจิเน ธมฺมํ,}{เอวํ ตตฺถ วิสุชฺฌติ.} \\ \csromangathabat{สาริปุตฺโตว ปญฺญาย,}{สีเลน อุปสเมน จ.}}\csromangathastanza{\csromanfolio{128}\csromangathabat{โยปิ ปารงฺคโต ภิกฺขุ,}{เอตาวปรโม สิยา”ติ.} \\ อิทํ นิพฺเพธ\-ภาคิยํ สุตฺตํ. \\ \csromansymbolmark{*}“อตีตํ นานฺวาคเมยฺย,}\csromangathastanza{นปฺปฏิกงฺเข อนาคตํ.}\csromangathastanza{\csromangathabat{ยทตีตํ ปหีนํ\csromansharedfootnote{81d5cf101c3e9b8e}{ปหีณํ (สี)} ตํ,}{อปฺปตฺตญฺจ อนาคตํ.} \\ \csromangathabat{ปจฺจุปฺปนฺนญฺจ โย ธมฺมํ\csromansharedfootnote{81d5cf101c3e9b8e}{ปหีณํ (สี)},}{ตตฺถ ตตฺถ วิปสฺสติ.}}\csromangathastanza{\csromangathabat{อสํหีรํ อสํกุปฺปํ,}{ตํ วิทฺวา มนุพฺรูหเย.} \\ \csromangathabat{อชฺเชว กิจฺจมาตปฺปํ\csromansharedfootnote{64548209108f29b0}{กิจฺจํ อาตปฺปํ (สี)},}{โก ชญฺญา มรณํ สุเว.}}\csromangathastanza{\csromangathabat{น หิ โน สงฺครํ เตน,}{มหาเสเนน มจฺจุนา.} \\ \csromangathabat{เอวํ วิหาริํ อาตาปิํ,}{อโหรตฺตม\-ตนฺทิตํ.}}\csromangathastanza{\csromangathawak{อิทํ นิพฺเพธ\-ภาคิยํ สุตฺตํ.}}}

\prose{ตํ เว “ภทฺเทกรตฺโต”ติ, สนฺโต อาจิกฺขเต มุนี”ติ.}

\proseitem{103}{“จตฺตาริมานิภิกฺขเว สจฺฉิ\-กาตพฺพานิ. กตมานิ จตฺตาริ, อตฺถิ ภิกฺขเว ธมฺมา จกฺขุนา ปญฺญาย จ สจฺฉิกาตพฺพา, อตฺถิ ธมฺมา สติยา ปญฺญาย จ สจฺฉิกาตพฺพา, อตฺถิ ธมฺมา กาเยน ปญฺญาย จ สจฺฉิกาตพฺพา, อตฺถิ ธมฺมา ปญฺญาย เวทิตพฺพา, ปญฺญาย จ สจฺฉิกาตพฺพา.}

\prose{กตเม จ ภิกฺขเว ธมฺมา จกฺขุนา ปญฺญาย จ สจฺฉิกาตพฺพา, ทิพฺพจกฺขุ สุวิสุทฺธํ อติกฺกนฺต\-มานุสกํ จกฺขุนา ปญฺญาย จ สจฺฉิกาตพฺพํ.}

\prose{กตเม จ ภิกฺขเว ธมฺมา สติยา ปญฺญาย จ สจฺฉิกาตพฺพา, ปุพฺเพ\-นิวาสา\-นุสฺสติ สติยา ปญฺญาย จ สจฺฉิกาตพฺพา.}

\prose{กตเม จ ภิกฺขเว ธมฺมา กาเยน ปญฺญาย จ สจฺฉิกาตพฺพา, อิทฺธิวิธา นิโรธา กาเยน ปญฺญาย จ สจฺฉิกาตพฺพา.}

\prose{กตเม จ ภิกฺขเว ธมฺมา ปญฺญาย เวทิตพฺพา, ปญฺญาย สจฺฉิกาตพฺพา, อาสวานํ ขเย ญาณํ ปญฺญาย เวทิตพฺพํ, ปญฺญาย จ สจฺฉิกาตพฺพนฺ”ติ. อิทํ นิพฺเพธ\-ภาคิยํ สุตฺตํ.}

\proseitem{104}{\csromanfolio{129}ตตฺถ กตมํ อเสกฺข\-ภาคิยํ สุตฺตํ.}

\setlength{\csromangathapagewidth}{0pt}
\setlength{\csromangathawakwidth}{0pt}
\csromangathameasuregroup{“ยสฺส เสลูปมํ จิตฺตํ, \\ วิรตฺตํ รชนีเยสุ,}{\csromangathabat{“ยสฺส เสลูปมํ จิตฺตํ,}{ฐิตํ นานุปกมฺปติ.} \\ \csromangathabat{วิรตฺตํ รชนีเยสุ,}{โกปเนยฺเย น กุปฺปติ.} \\ อิทํ อเสกฺขภาคิยํ สุตฺตํ.}
\csromangathameasurewak{อิทํ อเสกฺขภาคิยํ สุตฺตํ.}
\csromangathasetleft{“ยสฺส เสลูปมํ จิตฺตํ, \\ วิรตฺตํ รชนีเยสุ,}
\csromangathagroup{\csromangathastanza{\csromangathabat{“ยสฺส เสลูปมํ จิตฺตํ,}{ฐิตํ นานุปกมฺปติ.} \\ \csromangathabat{วิรตฺตํ รชนีเยสุ,}{โกปเนยฺเย น กุปฺปติ.}}\csromangathastanza{\csromangathawak{อิทํ อเสกฺข\-ภาคิยํ สุตฺตํ.}}}

\prose{ยสฺเสวํ ภาวิตํ จิตฺตํ, กุโต นํ ทุกฺขเมสฺสตี”ติ\csromansharedfootnote{79861c79659f7ecc}{ขุ 1. 152 ปิฏฺเฐ; อุปริ 185 ปิฏฺเฐปิ.}.}

\proseitem{104}{อายสฺมโต จ สาริปุตฺตสฺส จาริกาทสมํ เวยฺยากรณํ กาตพฺพนฺติ. อิทํ อเสกฺข\-ภาคิยํ สุตฺตํ.}

\prose{\csromansymbolfootnote{*}{วิ 3. 3; ขุ 1. 80 ปิฏฺเฐสุ.}“โย พฺราหฺมโณ พาหิต\-ปาปธมฺโม,}

\prose{นิหุํหุงฺโก\csromansharedfootnote{f9613c0b42d18b95}{นิหุหุงฺโก (สี) ** ขุ 1. 81 ปิฏฺเฐ.} นิกฺกสาโว ยตตฺโต.}

\setlength{\csromangathapagewidth}{0pt}
\setlength{\csromangathawakwidth}{0pt}
\csromangathameasuregroup{เวทนฺตคู วูสิตพฺรหฺมจริโย,}{\csromangathabat{เวทนฺตคู วูสิตพฺรหฺมจริโย,}{ธมฺเมน โส พฺรหฺมวาทํ วเทยฺย.} \\ ยสฺสุ’สฺสทา นตฺถิ กุหิญฺจิ โลเก”ติ. \\ อิทํ อเสกฺขภาคิยํ สุตฺตํ.}
\csromangathasetleft{เวทนฺตคู วูสิตพฺรหฺมจริโย,}
\csromangathagroup{\csromangathastanza{\csromangathabat{เวทนฺตคู วูสิต\-พฺรหฺมจริโย,}{ธมฺเมน โส พฺรหฺมวาทํ วเทยฺย.} \\ ยสฺสุ’สฺสทา นตฺถิ กุหิญฺจิ โลเก”ติ. \\ อิทํ อเสกฺข\-ภาคิยํ สุตฺตํ.}}

\prose{** “พาหิตฺวา ปาปเก ธมฺเม, เย จรนฺติ สทา สตา.}

\setlength{\csromangathapagewidth}{0pt}
\setlength{\csromangathawakwidth}{0pt}
\csromangathameasuregroup{ขีณสํโยชนา พุทฺธา,}{\csromangathabat{ขีณสํโยชนา พุทฺธา,}{เต เว โลกสฺมิ\textsuperscript{1} พฺราหฺมณา”ติ.} \\ อิทํ อเสกฺขภาคิยํ สุตฺตํ.}
\csromangathasetleft{ขีณสํโยชนา พุทฺธา,}
\csromangathagroup{\csromangathastanza{\csromangathabat{ขีณสํโยชนา พุทฺธา,}{เต เว โลกสฺมิ\csromansharedfootnote{338dd9d64fd7ee1e}{[มิสฺสินฺคฺ โนเต 0]} พฺราหฺมณา”ติ.} \\ อิทํ อเสกฺข\-ภาคิยํ สุตฺตํ.}}

\prose{\csromansymbolfootnote{+}{ขุ 1. 86, 87 ปิฏฺเฐสุ.}“ยตฺถ อาโป จ ปถวี, เตโช วาโย น คาธติ.}

\setlength{\csromangathapagewidth}{0pt}
\setlength{\csromangathawakwidth}{0pt}
\csromangathameasuregroup{น ตตฺถ สุกฺกา โชตนฺติ, \\ น ตตฺถ จนฺทิมา ภาติ, \\ ยทา จ อตฺตนา’เวทิ\textsuperscript{1}, \\ อถ รูปา อรูปา จ, \\ “ยทา สเกสุ ธมฺเมสุ,}{\csromangathabat{น ตตฺถ สุกฺกา โชตนฺติ,}{อาทิจฺโจ นปฺปกาสติ.} \\ \csromangathabat{น ตตฺถ จนฺทิมา ภาติ,}{ตโม ตตฺถ น วิชฺชติ.} \\ \csromangathabat{ยทา จ อตฺตนา’เวทิ\textsuperscript{1},}{มุนิ โมเนน พฺรหฺมโณ.} \\ \csromangathabat{อถ รูปา อรูปา จ,}{สุขทุกฺขา ปมุจฺจตี”ติ.} \\ อิทํ อเสกฺขภาคิยํ สุตฺตํ. \\ \csromangathabat{“ยทา สเกสุ ธมฺเมสุ,}{ปารคู โหติ พฺราหฺมโณ.} \\ อิทํ อเสกฺขภาคิยํ สุตฺตํ.}
\csromangathameasurewak{อิทํ อเสกฺขภาคิยํ สุตฺตํ. \\ อิทํ อเสกฺขภาคิยํ สุตฺตํ.}
\csromangathasetleft{น ตตฺถ สุกฺกา โชตนฺติ, \\ น ตตฺถ จนฺทิมา ภาติ, \\ ยทา จ อตฺตนา’เวทิ\textsuperscript{1}, \\ อถ รูปา อรูปา จ, \\ “ยทา สเกสุ ธมฺเมสุ,}
\csromangathagroup{\csromangathastanza{\csromangathabat{น ตตฺถ สุกฺกา โชตนฺติ,}{อาทิจฺโจ นปฺปกาสติ.} \\ \csromangathabat{น ตตฺถ จนฺทิมา ภาติ,}{ตโม ตตฺถ น วิชฺชติ.}}\csromangathastanza{\csromangathabat{ยทา จ อตฺตนา’เวทิ\csromansharedfootnote{338dd9d64fd7ee1e}{[มิสฺสินฺคฺ โนเต 0]},}{มุนิ โมเนน พฺรหฺมโณ.} \\ \csromangathabat{อถ รูปา อรูปา จ,}{สุขทุกฺขา ปมุจฺจตี”ติ.}}\csromangathastanza{\csromangathawak{อิทํ อเสกฺข\-ภาคิยํ สุตฺตํ.}}\csromangathastanza{\csromangathabat{“ยทา สเกสุ ธมฺเมสุ,}{ปารคู โหติ พฺราหฺมโณ.}}\csromangathastanza{\csromangathawak{อิทํ อเสกฺข\-ภาคิยํ สุตฺตํ.}}}

\prose{อถ เอตํ ปิสาจญฺจ, ปกฺกุลญฺจาติวตฺตตี”ติ\csromansharedfootnote{b054ae76de616a8a}{ขุ 1. 82 ปิฏฺเฐ.}.}

\setlength{\csromangathapagewidth}{0pt}
\setlength{\csromangathawakwidth}{0pt}
\csromangathameasuregroup{นาภินนฺทติ อายนฺติํ\textsuperscript{1}, \\ สงฺคา สงฺคามชิํ มุตฺตํ, \\ “น อุทเกน สุจี\textsuperscript{1} โหติ,}{\csromangathabat{นาภินนฺทติ อายนฺติํ\textsuperscript{1},}{ปกฺกมนฺติํ น โสจติ.} \\ \csromangathabat{สงฺคา สงฺคามชิํ มุตฺตํ,}{ตมหํ พฺรูมิ พฺราหฺมณนฺ”ติ.} \\ อิทํ อเสกฺขภาคิยํ สุตฺตํ. \\ \csromangathabat{“น อุทเกน สุจี\textsuperscript{1} โหติ,}{พเหฺวตฺถ นฺหายตี\textsuperscript{1} ชโน.} \\ อิทํ อเสกฺขภาคิยํ สุตฺตํ.}
\csromangathameasurewak{อิทํ อเสกฺขภาคิยํ สุตฺตํ. \\ อิทํ อเสกฺขภาคิยํ สุตฺตํ.}
\csromangathasetleft{นาภินนฺทติ อายนฺติํ\textsuperscript{1}, \\ สงฺคา สงฺคามชิํ มุตฺตํ, \\ “น อุทเกน สุจี\textsuperscript{1} โหติ,}
\csromangathagroup{\csromangathastanza{\csromanfolio{130}\csromangathabat{นาภินนฺทติ อายนฺติํ\csromansharedfootnote{338dd9d64fd7ee1e}{[มิสฺสินฺคฺ โนเต 0]},}{ปกฺกมนฺติํ น โสจติ.} \\ \csromangathabat{สงฺคา สงฺคามชิํ มุตฺตํ,}{ตมหํ พฺรูมิ พฺราหฺมณนฺ”ติ.}}\csromangathastanza{\csromangathawak{อิทํ อเสกฺข\-ภาคิยํ สุตฺตํ.}}\csromangathastanza{\csromangathabat{“น อุทเกน สุจี\csromansharedfootnote{338dd9d64fd7ee1e}{[มิสฺสินฺคฺ โนเต 0]} โหติ,}{พเหฺวตฺถ นฺหายตี\csromansharedfootnote{61785d03c494c2ae}{[มิสฺสินฺคฺ โนเต 1]} ชโน.}}\csromangathastanza{\csromangathawak{อิทํ อเสกฺข\-ภาคิยํ สุตฺตํ.}}}

\prose{ยมฺหิ สจฺจญฺจ ธมฺโม จ, โส สุจี โส จ พฺราหฺมโณ”ติ\csromansharedfootnote{0e9a37ecb09cc368}{สุจิ (สี, ก) ขุ 1. 84 ปิฏฺเฐ.}.}

\setlength{\csromangathapagewidth}{0pt}
\setlength{\csromangathawakwidth}{0pt}
\csromangathameasuregroup{นโม เต ปุริสาชญฺญ, \\ ยสฺส เต นาภิชานาม, \\ “สหายา วติเม ภิกฺขู, \\ สเมติ เนสํ สทฺธมฺโม, \\ สุวินีตา กปฺปิเนน,}{\textsuperscript{*}“ยทา หเว ปาตุภวนฺติ ธมฺมา, \\ อาตาปิโน ฌายโต พฺราหฺมณสฺส. \\ วิธูปยํ ติฏฺฐติ มารเสนํ, \\ สูริโยว โอภาสยมนฺตลิกฺขนฺ”ติ. \\ อิทํ อเสกฺขภาคิยํ สุตฺตํ. \\ “สนฺตินฺทฺริยํ ปสฺสถ อิริยมานํ, \\ เตวิชฺชปตฺตํ อปหานธมฺมํ. \\ สพฺพานิ โยคานิ อุปาติวตฺโต, \\ อกิญฺจโน อิริยติ ปํสุกูลิโก. \\ ตํ เทวตา สมฺพหุลา อุฬารา, \\ พฺรหฺมวิมานํ อุปสงฺกมิตฺวา. \\ อาชานิยํ ชาติพลํ นิเสธํ, \\ นีตํ\textsuperscript{1} นมสฺสนฺติ ปสนฺนจิตฺตา. \\ \csromangathabat{นโม เต ปุริสาชญฺญ,}{นโม เต ปุริสุตฺตม.} \\ \csromangathabat{ยสฺส เต นาภิชานาม,}{กิํ ตฺวํ นิสฺสาย ฌายสี”ติ.} \\ อิทํ อเสกฺขภาคิยํ สุตฺตํ. \\ \csromangathabat{“สหายา วติเม ภิกฺขู,}{จิรรตฺตสเมติกา.} \\ \csromangathabat{สเมติ เนสํ สทฺธมฺโม,}{ธมฺเม พุทฺธปฺปเวทิเต.} \\ \csromangathabat{สุวินีตา กปฺปิเนน,}{ธมฺเม อริยปฺปเวทิเต.} \\ อิทํ อเสกฺขภาคิยํ สุตฺตํ.}
\csromangathameasurewak{\textsuperscript{*}“ยทา หเว ปาตุภวนฺติ ธมฺมา, \\ อาตาปิโน ฌายโต พฺราหฺมณสฺส. \\ วิธูปยํ ติฏฺฐติ มารเสนํ, \\ สูริโยว โอภาสยมนฺตลิกฺขนฺ”ติ. \\ อิทํ อเสกฺขภาคิยํ สุตฺตํ. \\ “สนฺตินฺทฺริยํ ปสฺสถ อิริยมานํ, \\ เตวิชฺชปตฺตํ อปหานธมฺมํ. \\ สพฺพานิ โยคานิ อุปาติวตฺโต, \\ อกิญฺจโน อิริยติ ปํสุกูลิโก. \\ ตํ เทวตา สมฺพหุลา อุฬารา, \\ พฺรหฺมวิมานํ อุปสงฺกมิตฺวา. \\ อาชานิยํ ชาติพลํ นิเสธํ, \\ นีตํ\textsuperscript{1} นมสฺสนฺติ ปสนฺนจิตฺตา. \\ อิทํ อเสกฺขภาคิยํ สุตฺตํ. \\ อิทํ อเสกฺขภาคิยํ สุตฺตํ.}
\csromangathasetleft{นโม เต ปุริสาชญฺญ, \\ ยสฺส เต นาภิชานาม, \\ “สหายา วติเม ภิกฺขู, \\ สเมติ เนสํ สทฺธมฺโม, \\ สุวินีตา กปฺปิเนน,}
\csromangathagroup{\csromangathastanza{\csromangathawak{\csromansymbolmark{*}“ยทา หเว ปาตุภวนฺติ ธมฺมา, \\ อาตาปิโน ฌายโต พฺราหฺมณสฺส. \\ วิธูปยํ ติฏฺฐติ มารเสนํ, \\ สูริโยว โอภาสยมนฺตลิกฺขนฺ”ติ.}}\csromangathastanza{\csromangathawak{อิทํ อเสกฺข\-ภาคิยํ สุตฺตํ.}}\csromangathastanza{\csromangathawak{“สนฺตินฺทฺริยํ ปสฺสถ อิริยมานํ, \\ เตวิชฺชปตฺตํ อปหานธมฺมํ. \\ สพฺพานิ โยคานิ อุปาติวตฺโต, \\ อกิญฺจโน อิริยติ ปํสุกูลิโก.}}\csromangathastanza{\csromangathawak{ตํ เทวตา สมฺพหุลา อุฬารา, \\ พฺรหฺมวิมานํ อุปสงฺกมิตฺวา. \\ อาชานิยํ ชาติพลํ นิเสธํ, \\ นีตํ\csromansharedfootnote{c30948a354f55383}{นิธ (ก); สํ ฏฺฐ 2. 237 ปิฏฺเฐ ปสฺสิตพฺพํ.} นมสฺสนฺติ ปสนฺนจิตฺตา.}}\csromangathastanza{\csromangathabat{นโม เต ปุริสาชญฺญ,}{นโม เต ปุริสุตฺตม.} \\ \csromangathabat{ยสฺส เต นาภิชานาม,}{กิํ ตฺวํ นิสฺสาย ฌายสี”ติ.}}\csromangathastanza{\csromangathawak{อิทํ อเสกฺข\-ภาคิยํ สุตฺตํ.}}\csromangathastanza{\csromangathabat{“สหายา วติเม ภิกฺขู,}{จิรรตฺตสเมติกา.} \\ \csromangathabat{สเมติ เนสํ สทฺธมฺโม,}{ธมฺเม พุทฺธ\-ปฺปเวทิเต.} \\ \csromanfolio{131}\csromangathabat{สุวินีตา กปฺปิเนน,}{ธมฺเม อริยปฺปเวทิเต.}}\csromangathastanza{\csromangathawak{อิทํ อเสกฺข\-ภาคิยํ สุตฺตํ.}}}

\prose{ธาเรนฺติ อนฺติมํ เทตํ, เชตฺวา มารํ สวาหินินฺ”ติ\csromansharedfootnote{55c564c2c84f0ccf}{สวาหนนฺ”ติ (ก) สํ 1. 472 ปิฏฺเฐ ปสฺสิตพฺพํ.}.}

\setlength{\csromangathapagewidth}{0pt}
\setlength{\csromangathawakwidth}{0pt}
\csromangathameasuregroup{\textsuperscript{*}“นยิทํ สิถิลมารพฺภ, \\ นิพฺพานํ อธิคนฺตพฺพํ, \\ อยญฺจ ทหโร ภิกฺขุ, \\ ธาเรติ อนฺติมํ เทหํ, \\ “ทุพฺพณฺณโก ลูขจีวโร, \\ ขีณาสโว วิสํยุตฺโต, \\ เตวิชฺโช อิทฺธิปฺปตฺโต จ, \\ ธาเรติ อนฺติมํ เทตํ,}{\csromangathabat{\textsuperscript{*}“นยิทํ สิถิลมารพฺภ,}{นยิทํ อปฺเปน ถามสา.} \\ \csromangathabat{นิพฺพานํ อธิคนฺตพฺพํ,}{สพฺพทุกฺขปฺปโมจนํ\textsuperscript{1}.} \\ \csromangathabat{อยญฺจ ทหโร ภิกฺขุ,}{อยมุตฺตมโปริโส.} \\ \csromangathabat{ธาเรติ อนฺติมํ เทหํ,}{เชตฺวา มารํ สวาหินินฺ”ติ.} \\ อิทํ อเสกฺขภาคิยํ สุตฺตํ. \\ \csromangathabat{“ทุพฺพณฺณโก ลูขจีวโร,}{โมฆราชา สทา สโต.} \\ \csromangathabat{ขีณาสโว วิสํยุตฺโต,}{กตกิจฺโจ อนาสโว.} \\ \csromangathabat{เตวิชฺโช อิทฺธิปฺปตฺโต จ,}{เจโตปริยายโกวิโท.} \\ \csromangathabat{ธาเรติ อนฺติมํ เทตํ,}{เชตฺวา มารํ สวาหินินฺ”ติ.} \\ อิทํ อเสกฺขภาคิยํ สุตฺตํ.}
\csromangathameasurewak{อิทํ อเสกฺขภาคิยํ สุตฺตํ. \\ อิทํ อเสกฺขภาคิยํ สุตฺตํ.}
\csromangathasetleft{\textsuperscript{*}“นยิทํ สิถิลมารพฺภ, \\ นิพฺพานํ อธิคนฺตพฺพํ, \\ อยญฺจ ทหโร ภิกฺขุ, \\ ธาเรติ อนฺติมํ เทหํ, \\ “ทุพฺพณฺณโก ลูขจีวโร, \\ ขีณาสโว วิสํยุตฺโต, \\ เตวิชฺโช อิทฺธิปฺปตฺโต จ, \\ ธาเรติ อนฺติมํ เทตํ,}
\csromangathagroup{\csromangathastanza{\csromangathabat{\csromansymbolfootnote{*}{สํ 1. 466 ปิฏฺเฐ.}“นยิทํ สิถิลมา\-รพฺภ,}{นยิทํ อปฺเปน ถามสา.} \\ \csromangathabat{นิพฺพานํ อธิคนฺตพฺพํ,}{สพฺพทุกฺขปฺปโมจนํ\csromansharedfootnote{dbe186fd019a039e}{สพฺพคนฺถปโมจนํ (ก) ** สํ 2. 54 ปิฏฺเฐ.}.}}\csromangathastanza{\csromangathabat{อยญฺจ ทหโร ภิกฺขุ,}{อยมุ\-ตฺตมโปริโส.} \\ \csromangathabat{ธาเรติ อนฺติมํ เทหํ,}{เชตฺวา มารํ สวาหินินฺ”ติ.}}\csromangathastanza{\csromangathawak{อิทํ อเสกฺข\-ภาคิยํ สุตฺตํ.}}\csromangathastanza{\csromangathabat{“ทุพฺพณฺณโก ลูขจีวโร,}{โมฆราชา สทา สโต.} \\ \csromangathabat{ขีณาสโว วิสํยุตฺโต,}{กตกิจฺโจ อนาสโว.}}\csromangathastanza{\csromangathabat{เตวิชฺโช อิทฺธิปฺปตฺโต จ,}{เจโตปริยายโกวิโท.} \\ \csromangathabat{ธาเรติ อนฺติมํ เทตํ,}{เชตฺวา มารํ สวาหินินฺ”ติ.}}\csromangathastanza{\csromangathawak{อิทํ อเสกฺข\-ภาคิยํ สุตฺตํ.}}}

\proseitem{105}{\csromansymbolmark{*}* ตถาคโต ภิกฺขเว อรหํ สมฺมาสมฺพุทฺโธ รูปสฺส นิพฺพิทา วิราคา นิโรธา อนุปาทา วิมุตฺโต สมฺมา\-สมฺพุทฺโธติ วุจฺจติ. ภิกฺขุปิ ภิกฺขเว ปญฺญาวิมุตฺโต รูปสฺส นิพฺพิทา วิราคา นิโรธา อนุปาทา วิมุตฺโต ปญฺญาวิมุตฺโตติ วุจฺจติ.}

\prose{ตถาคโต ภิกฺขเว อรหํ สมฺมาสมฺพุทฺโธ เวทนาย ฯเปฯ สญฺญาย. สงฺขารานํ. วิญฺญาณสฺส นิพฺพิทา วิราคา นิโรธา อนุปาทา วิมุตฺโต สมฺมา\-สมฺพุทฺโธติ วุจฺจติ. ภิกฺขุปิ ภิกฺขเว ปญฺญาวิมุตฺโต วิญฺญาณสฺส นิพฺพิทา วิราคา นิโรธา อนุปาทา วิมุตฺโต ปญฺญาวิมุตฺโตติ วุจฺจติ.}

\prose{ตตฺร โข ภิกฺขเว โก วิเสโส โก อธิปฺปยาโส\csromansharedfootnote{60a1a8ac9cfbc2bd}{อธิปฺปาโย (สี, ก)} กิํ นานากรณํ ตถาคตสฺส อรหโต สมฺมา\-สมฺพุทฺธสฺส ปญฺญาวิมุตฺเตน ภิกฺขุนาติ. ภควํมูลกา โน ภนฺเต ธมฺมา ฯเปฯ.}

\prose{ตถาคโต ภิกฺขเว อรหํ สมฺมาสมฺพุทฺโธ อนุปฺปนฺนสฺส มคฺคสฺส อุปฺปาเทตา, อสญฺชาตสฺส มคฺคสฺส สญฺชเนตา, อนกฺขาตสฺส มคฺคสฺส \csromanfolio{132}อกฺขาตา, มคฺคญฺญู มคฺควิทู มคฺคโกวิโท, มคฺคานุคา จ ภิกฺขเว เอตรหิ สาวกา วิหรนฺติ ปจฺฉา\-สมนฺนาคตา. อยํ โข ภิกฺขเว วิเสโส, อยํ อธิปฺปยาโส, อิทํ นานากรณํ ตถาคตสฺส อรหโต สมฺมาส\-มฺพุ\-ทฺทสฺส ปญฺญาวิมุตฺเตน ภิกฺขุนา”ติ. อิทํ อเสกฺข\-ภาคิยํ สุตฺตํ.}

\setlength{\csromangathapagewidth}{0pt}
\setlength{\csromangathawakwidth}{0pt}
\csromangathameasuregroup{}{ตตฺถ กตมํ สํกิเลสภาคิยญฺจ วาสนาภาคิยญฺจ สุตฺตํ. \\ “ฉนฺนมติวสฺสติ, วิวฏํ นาติวสฺสติ. \\ ตสฺมา ฉนฺนํ วิวเรถ, เอวํ ตํ นาติวสฺสตี”ติ\textsuperscript{1}.}
\csromangathameasurewak{ตตฺถ กตมํ สํกิเลสภาคิยญฺจ วาสนาภาคิยญฺจ สุตฺตํ. \\ “ฉนฺนมติวสฺสติ, วิวฏํ นาติวสฺสติ. \\ ตสฺมา ฉนฺนํ วิวเรถ, เอวํ ตํ นาติวสฺสตี”ติ\textsuperscript{1}.}
\setlength{\csromangathaleftcol}{0pt}
\csromangathagroup{\csromangathastanza{\csromangathawak{\csromangathaitemline{106}{ตตฺถ กตมํ สํกิเลสภาคิยญฺจ วาสนาภาคิยญฺจ สุตฺตํ.} \\ \csromangathacontline{“ฉนฺนม\-ติวสฺสติ, วิวฏํ นาติวสฺสติ.} \\ \csromangathacontline{ตสฺมา ฉนฺนํ วิวเรถ, เอวํ ตํ นาติวสฺสตี”ติ\csromansharedfootnote{6bb8d7c7d71501d4}{ขุ 1. 144 ปิฏฺเฐ; อุปริ 185 ปิฏฺเฐ จ.}.}}}}

\prose{“ฉนฺน\-มติ\-วสฺสตี”ติ สํกิเลโส, “วิวฏํ นาติวสฺสตี”ติ วาสนา, “ตสฺมา ฉนฺนํ วิวเรถ, เอวํ ตํ นาติวสฺสตี”ติ อยํ สํกิเลโส จ วาสนา จ. อิทํ สํกิเลสภาคิยญฺจ วาสนาภาคิยญฺจ สุตฺตํ.}

\prose{“จตฺตาโรเม มหาราช\csromansharedfootnote{75fe9d4dc1e5cb72}{ภิกฺขเว (อํ 1. 397 ปิฏฺเฐ.)} ปุคฺคลา สนฺโต สํวิชฺชมานา โลกสฺมิํ, กตเม จตฺตาโร, ตโม ตมปรายโณ ตโม โชติปรายโณ โชติ ตมปรายโณ โชติ โชติปรายโณ”ติ. ตตฺถ โย จ ปุคฺคโล โชติ ตมปรายโณ โย จ ปุคฺคโล ตโม ตมปรายโณ, อิเม เทฺว ปุคฺคลา สํกิเลส\-ภาคิยา, โย จ ปุคฺคโล ตโม โชติปรายโณ โย จ ปุคฺคโล โชติ โชติปรายโณ, อิเม เทฺว ปุคฺคลา วาสนาภาคิยา. อิทํ สํกิเลสภาคิยญฺจ วาสนาภาคิยญฺจ สุตฺตํ.}

\prose{ตตฺถ กตมํ สํกิเลสภาคิยญฺจ นิพฺเพธภาคิยญฺจ สุตฺตํ.}

\prose{\csromansymbolfootnote{*}{เหฏฺฐา 31 ปิฏฺเฐ; อุปริ 185 ปิฏฺเฐปิ. ** เหฏฺฐา 31 ปิฏฺเฐปิ.}“น ตํ ทฬฺหํ พนฺธนมาหุ ธีรา,}

\prose{“ยทายสํ ทารุชปพฺพชญฺจ.}

\setlength{\csromangathapagewidth}{0pt}
\setlength{\csromangathawakwidth}{0pt}
\csromangathameasuregroup{}{สารตฺตรตฺตา มณิกุณฺฑเลสุ, \\ ปุตฺเตสุ ทาเรสุ จ ยา อเปกฺขา”ติ.อยํ สํกิเลโส.}
\csromangathameasurewak{สารตฺตรตฺตา มณิกุณฺฑเลสุ, \\ ปุตฺเตสุ ทาเรสุ จ ยา อเปกฺขา”ติ.อยํ สํกิเลโส.}
\setlength{\csromangathaleftcol}{0pt}
\csromangathagroup{\csromangathastanza{\csromangathawak{สารตฺตรตฺตา มณิกุณฺฑเลสุ, \\ ปุตฺเตสุ ทาเรสุ จ ยา อเปกฺขา”ติ.  อยํ สํกิเลโส.}}}

\prose{** “เอตํ ทฬฺหํ พนฺธนมาหุ ธีรา,}

\prose{โอหารินํ สิถิลํ ทุปฺปมุญฺจํ.}

\setlength{\csromangathapagewidth}{0pt}
\setlength{\csromangathawakwidth}{0pt}
\csromangathameasuregroup{}{เอตมฺปิ เฉตฺวาน ปริพฺพชนฺติ, \\ อนเปกฺขิโน กามสุขํ ปหายา”ติ.}
\csromangathameasurewak{เอตมฺปิ เฉตฺวาน ปริพฺพชนฺติ, \\ อนเปกฺขิโน กามสุขํ ปหายา”ติ.}
\setlength{\csromangathaleftcol}{0pt}
\csromangathagroup{\csromangathastanza{\csromangathawak{เอตมฺปิ เฉตฺวาน ปริพฺพชนฺติ, \\ อนเปกฺขิโน กามสุขํ ปหายา”ติ.}}}

\prose{\csromanfolio{133}อยํ นิพฺเพโธ.  อิทํ สํกิเลสภาคิยญฺจ นิพฺเพธภาคิยญฺจ สุตฺตํ.}

\proseitem{107}{\csromansymbolfootnote{*}{สํ 1. 295 ปิฏฺเฐ.}“ยญฺจ ภิกฺขเว เจเตติ, ยญฺจ ปกปฺเปติ, ยญฺจ อนุเสติ. อารมฺมณเมตํ โหติ วิญฺญาณสฺส ฐิติยา, อารมฺมเณ สติ ปติฏฺฐา วิญฺญาณสฺส โหติ, ตสฺมิํ ปติฏฺฐิเต วิญฺญาเณ วิรูเฬฺห อายติํ\csromansharedfootnote{4030749e407aebb8}{อายติ (สี, ก)} ปุนพฺภวา\-ภินิพฺพตฺติ โหติ, อายติํ ปุนพฺภวา\-ภินิพฺพตฺติยา สติ อายติํ\csromansharedfootnote{4030749e407aebb8}{อายติ (สี, ก)} ชาติ\-ชรา\-มรณํ โสก\-ปริเทว\-ทุกฺข\-โทมนสฺสุ\-ปายาสา สมฺภวนฺติ, เอวเมตสฺส เกวลสฺส ทุกฺข\-กฺขนฺธสฺส สมุทโย โหติ.}

\prose{โน เจ ภิกฺขเว เจเตติ, โน เจ ปกปฺเปติ, อถ เจ อนุเสติ. อารมฺมณเมตํ โหติ วิญฺญาณสฺส ฐิติยา, อารมฺมเณ สติ ปติฏฺฐา วิญฺญาณสฺส\csromansharedfootnote{9d5f96ab3107e8ad}{ตสฺส วิญฺญาณสฺส (สี, ก)} โหติ, ตสฺมิํ ปติฏฺฐิเต วิญฺญาเณ วิรูเฬฺห อายติํ ปุนพฺภวา\-ภินิพฺพตฺติ โหติ, อายติํ ปุนพฺภวา\-ภินิพฺพตฺติยา สติ อายติํ ชาติ\-ชรา\-มรณํ โสก\-ปริเทว\-ทุกฺข\-โทมนสฺสุ\-ปายาสา สมฺภวนฺติ, เอวเมตสฺส เกวลสฺส ทุกฺข\-กฺขนฺธสฺส สมุทโย โหตี”ติ. อยํ สํกิเลโส.}

\prose{“ยโต จ โข ภิกฺขเว โน เจว\csromansharedfootnote{b4e15e34a3ecdd8e}{จ (สี, ก) ** สํ 2. 367 ปิฏฺเฐ.} เจเตติ, โน จ ปกปฺเปติ, โน จ อนุเสติ. อารมฺมณเมตํ น โหติ วิญฺญาณสฺส ฐิติยา, อารมฺมเณ อสติ ปติฏฺฐา วิญฺญาณสฺส น โหติ, ตสฺมิํ อปฺปติฏฺฐิเต วิญฺญาเณ อวิรูเฬฺห อายติํ ปุนพฺภวา\-ภินิพฺพตฺติ น โหติ, อายติํ ปุนพฺภวา\-ภินิพฺพตฺติยา อสติ อายติํ ชาติ\-ชรา\-มรณํ โสก\-ปริเทว\-ทุกฺข\-โทมนสฺสุ\-ปายาสา นิรุชฺฌนฺติ. เอวเมตสฺส เกวลสฺส ทุกฺข\-กฺขนฺธสฺส นิโรโธ โหตี”ติ, อยํ นิพฺเพโธ. อิทํ สํกิเลสภาคิยญฺจ นิพฺเพธภาคิยญฺจ สุตฺตํ.}

\proseitem{108}{ตตฺถ กตมํ สํกิเลสภาคิยญฺจ อเสกฺขภาคิยญฺจ สุตฺตํ.}

\prose{** “สมุทฺโท สมุทฺโท”ติ โข ภิกฺขเว อสฺสุตวา ปุถุชฺชโน ภาสติ, เนโส ภิกฺขเว อริยสฺส วินเย สมุทฺโท, มหา เอโส ภิกฺขเว อุทกราสิ มหา อุทกณฺณโว. จกฺขุ ภิกฺขเว ปุริสสฺส สมุทฺโท, ตสฺส รูปมโย เวโค. อยํ สํกิเลโส.}

\prose{\csromanfolio{134}\csromansymbolfootnote{*}{สํ 2. 367 ปิฏฺเฐ.}“โย ตํ รูปมยํ เวคํ สหติ. อยํ วุจฺจติ ภิกฺขเว อตริ\csromansharedfootnote{9b585011d2e91002}{อตาริ (สี, ก)} จกฺขุสมุทฺทํ สอูมิํ สาวฏฺฏํ สคหํ\csromansharedfootnote{c38c9b017bcdc050}{สคาหํ (สํ 2. 367 ปิฏฺเฐ.) ** สํ 2. 368 ปิฏฺเฐ.} สรกฺขสํ ติณฺโณ ปารงฺคโต ถเล ติฏฺฐติ พฺราหฺมโณ”ติ. อยํ อเสกฺโข.}

\prose{\csromansymbolmark{*}“โสตํ ภิกฺขเว ฯเปฯ. ฆานํ. ชิวฺหา. กาโย. มโน ภิกฺขเว ปุริสสฺส สมุทฺโท ตสฺส ธมฺมมโย เวโค”ติ.  อยํ สํกิเลโส.}

\prose{\csromansymbolmark{*}“โย ตํ ธมฺมมยํ เวคํ สหติ, อยํ วุจฺจติ ภิกฺขเว อตริ มโนสมุทฺทํ สอูมิํ สาวฏฺฏํ สคหํ สรกฺขสํ ติณฺโณ ปารงฺคโต ถเล ติฏฺฐติ พฺราหฺมโณ”ติ. อยํ อเสกฺโข. อิทมโวจ ภควา, อิทํ วตฺวาน สุคโต, อถาปรํ เอตทโวจ สตฺถา\csromandash{}}

\setlength{\csromangathapagewidth}{0pt}
\setlength{\csromangathawakwidth}{0pt}
\csromangathameasuregroup{}{โย อิมํ สมุทฺทํ สคหํ\textsuperscript{1} สรกฺขสํ, \\ สอูมิํ สาวฏฺฏํ สภยํ ทุตฺตรํ อจฺจตริ. \\ ส เวทนฺตคู วุสิตพฺรหฺมจริโย, \\ โลกนฺตคู ปารคโตติ วุจฺจตี”ติ.}
\csromangathameasurewak{โย อิมํ สมุทฺทํ สคหํ\textsuperscript{1} สรกฺขสํ, \\ สอูมิํ สาวฏฺฏํ สภยํ ทุตฺตรํ อจฺจตริ. \\ ส เวทนฺตคู วุสิตพฺรหฺมจริโย, \\ โลกนฺตคู ปารคโตติ วุจฺจตี”ติ.}
\setlength{\csromangathaleftcol}{0pt}
\csromangathagroup{\csromangathastanza{\csromangathawak{โย อิมํ สมุทฺทํ สคหํ\csromansharedfootnote{c38c9b017bcdc050}{สคาหํ (สํ 2. 367 ปิฏฺเฐ.) ** สํ 2. 368 ปิฏฺเฐ.} สรกฺขสํ, \\ สอูมิํ สาวฏฺฏํ สภยํ ทุตฺตรํ อจฺจตริ. \\ ส เวทนฺตคู วุสิต\-พฺรหฺมจริโย, \\ โลกนฺตคู ปารคโตติ วุจฺจตี”ติ.}}}

\prose{อยํ อเสกฺโข. อิทํ สํกิเลสภาคิยญฺจ อเสกฺขภาคิยญฺจ สุตฺตํ.}

\prose{** “ฉยิเม ภิกฺขเว พฬิสา โลกสฺมิํ อนยาย สตฺตานํ พฺยาพาธาย\csromansharedfootnote{6dff10e831f2b413}{วธาย (สํ 2. 368 ปิฏฺเฐ.)} ปาณีนํ. กตเม ฉ, สนฺติ ภิกฺขเว จกฺขุวิญฺเญยฺยา รูปา อิฏฺฐา กนฺตา มนาปา ปิยรูปา กามูปสํหิตา รชนียา, ตญฺเจ ภิกฺขุ อภินนฺทติ อภิวทติ อชฺโฌสาย ติฏฺฐติ, อยํ วุจฺจติ ภิกฺขเว ภิกฺขุ คิลิตพฬิโส\csromansharedfootnote{a036577b3192b8d1}{คิลพฬิโส (สี, ก) สํ 1. 182 ปิฏฺเฐปิ.} มารสฺส อนยํ อาปนฺโน, พฺยสนํ อาปนฺโน, ยถา\-กาม\-กรณีโย\csromansharedfootnote{90d62a4ef5840c6a}{ยถากามํ กรณีโย (ก)} ปาปิมโต.}

\prose{สนฺติ ภิกฺขเว โสตวิญฺเญยฺยา สทฺทา ฯเปฯ ฆานวิญฺเญยฺยา คนฺธา. ชิวฺหาวิญฺเญยฺยา รสา. กายวิญฺเญยฺยา โผฏฺฐพฺพา. มโนวิญฺเญยฺยา ธมฺมา อิฏฺฐา กนฺตา มนาปา ปิยรูปา กามูปสํหิตา รชนียา, ตญฺเจ ภิกฺขุ อภินนฺทติ อภิวทติ อชฺโฌสาย ติฏฺฐติ. อยํ วุจฺจติ ภิกฺขเว ภิกฺขุ คิฬิตพลิโส มารสฺส อนยํ อาปนฺโน, พฺยสนํ อาปนฺโน, ยถา\-กาม\-กรณีโย ปาปิมโต”ติ.  อยํ สํกิเลโส.}

\prose{\csromanfolio{135}“สนฺติ จ ภิกฺขเว จกฺขุวิญฺเญยฺยา รูปา อิฏฺฐา กนฺตา มนาปา ปิยรูปา กามูปสํหิตา รชนียา, ตญฺเจ ภิกฺขุ นาภินนฺทติ นาภิวทติ นาชฺโฌสาย ติฏฺฐติ, อยํ วุจฺจติ ภิกฺขเว ภิกฺขุ น คิลิตพฬิโส มารสฺส, อเภทิ พฬิสํ, ปริเภทิ พฬิสํ, น อนยํ อาปนฺโน, น พฺยสนํ อาปนฺโน, น ยถา\-กาม\-กรณีโย ปาปิมโต.}

\prose{สนฺติ จ ภิกฺขเว โสตวิญฺเญยฺยา สทฺทา ฯเปฯ มโนวิญฺเญยฺยา ธมฺมา อิฏฺฐา กนฺตา มนาปา ปิยรูปา กามุปสํหิตา รชนียา, ตญฺเจ ภิกฺขุ นาภินนฺทติ นาภิวทติ, นาชฺโฌสาย ติฏฺฐติ. อยํ วุจฺจติ ภิกฺขเว ภิกฺขุ น คิลิตพฬิโส มารสฺส, อเภทิ พฬิสํ, ปริเภทิ พฬิสํ, น อนยํ อาปนฺโน, น พฺยสนํ อาปนฺโน, น ยถา\-กาม\-กรณีโย ปาปิมโต”ติ.  อยํ อเสกฺโข. อิทํ สํกิเลสภาคิยญฺจ อเสกฺขภาคิยญฺจ สุตฺตํ.}

\proseitem{109}{ตตฺถ กตมํ สํกิเลสภาคิยญฺจ นิพฺเพธภาคิยญฺจ อเสกฺขภาคิยญฺจ สุตฺตํ.}

\setlength{\csromangathapagewidth}{0pt}
\setlength{\csromangathawakwidth}{0pt}
\csromangathameasuregroup{อญฺญถาภาวี ภวสตฺโต โลโก,}{\textsuperscript{*}อยํ โลโก สนฺตาปชาโต, \\ ผสฺสปเรโต โรคํ วทติ อตฺตโต\textsuperscript{1}. \\ เยน เยน หิ มญฺญนฺติ\textsuperscript{1}, \\ ตโต ตํ โหติ อญฺญถา. \\ \csromangathabat{อญฺญถาภาวี ภวสตฺโต โลโก,}{ภวปเรโต ภวเมวาภินนฺทติ.} \\ ยทภินนฺทติ ตํ ภยํ. \\ ยสฺส ภายติ ตํ ทุกฺขนฺ”ติ. \\ อยํ สํกิเลโส.}
\csromangathameasurewak{\textsuperscript{*}อยํ โลโก สนฺตาปชาโต, \\ ผสฺสปเรโต โรคํ วทติ อตฺตโต\textsuperscript{1}. \\ เยน เยน หิ มญฺญนฺติ\textsuperscript{1}, \\ ตโต ตํ โหติ อญฺญถา. \\ อยํ สํกิเลโส.}
\csromangathasetleft{อญฺญถาภาวี ภวสตฺโต โลโก,}
\csromangathagroup{\csromangathastanza{\csromangathawak{\csromansymbolfootnote{*}{ขุ 1. 115 ปิฏฺเฐ อุทาเน.}อยํ โลโก สนฺตาปชาโต, \\ ผสฺสปเรโต โรคํ วทติ อตฺตโต\csromansharedfootnote{bbac3e1be509c3ba}{อตฺตโน (สี, ก)}. \\ เยน เยน หิ มญฺญนฺติ\csromansharedfootnote{6f1dcc58f238b7a5}{มญฺญติ (ขุ 1. 115 ปิฏฺเฐ.)}, \\ ตโต ตํ โหติ อญฺญถา.}}\csromangathastanza{\csromangathabat{อญฺญถาภาวี ภวสตฺโต โลโก,}{ภวปเรโต ภวเม\-วา\-ภินนฺทติ.} \\ ยทภินนฺทติ ตํ ภยํ. \\ ยสฺส ภายติ ตํ ทุกฺขนฺ”ติ.}\csromangathastanza{\csromangathawak{อยํ สํกิเลโส.}}}

\prose{\csromansymbolmark{*}“ภววิปฺปหานาย โข ปนิทํ พฺรหฺมจริยํ วุสฺสตี”ติ. อยํ นิพฺเพโธ.}

\prose{\csromansymbolmark{*}“เย หิ เกจิ สมณา วา พฺราหฺมณา วา ภเวน ภวสฺส วิปฺปโมกฺขมาหํสุ, สพฺเพ เต “อวิปฺปมุตฺตา ภวสฺมา”ติ วทามิ. เย วา ปน เกจิ สมณา วา พฺราหฺมณา วา วิภเวน ภวสฺส นิสฺสรณมาหํสุ, สพฺเพ เต “อนิสฺสฏา ภวสฺมา”ติ วทามิ. อุปธิํ\csromansharedfootnote{7e2d84e02d91217e}{อุปธี (สี, ก)} หิ ปฏิจฺจ ทุกฺขมิทํ สมฺโภตี”ติ. อยํ สํกิเลโส.}

\prose{\csromanfolio{136}\csromansymbolfootnote{*}{ขุ 1. 115 ปิฏฺเฐ อุทาเน.}“สพฺพุ\-ปาทาน\-กฺขยา นตฺถิ ทุกฺขสฺส สมฺภโว”ติ. อยํ นิพฺเพโธ.}

\prose{\csromansymbolmark{*}“โลกมิมํ ปสฺส, ปุถู อวิชฺชาย ปเรตา ภูตา ภูตรตา\csromansharedfootnote{7782b45637b27abe}{ปเรตํ ภูตํ ภูตรตํ. (ฏฺฐ) ** ขุ 1. 116 ปิฏฺเฐ อุทาเน; อุปริ 186 ปิฏฺเฐปิ.}, ภวา อปริมุตฺตา, เย หิ เกจิ ภวา สพฺพธิ สพฺพตฺถตาย, สพฺเพ เต ภวา อนิจฺจา ทุกฺขา วิปริณาม\-ธมฺมา”ติ. อยํ สํกิเลโส.}

\prose{** “เอวเมตํ ยถาภูตํ, สมฺมปฺปญฺญาย ปสฺสโต.}

\setlength{\csromangathapagewidth}{0pt}
\setlength{\csromangathawakwidth}{0pt}
\csromangathameasuregroup{}{ภวตณฺหา ปหียติ, \\ วิภวํ นาภินนฺทติ. \\ สพฺพโส ตณฺหานํ ขยา, \\ อเสสวิราคนิโรโธ นิพฺพานนฺ”ติ. \\ อยํ นิพฺเพโธ.}
\csromangathameasurewak{ภวตณฺหา ปหียติ, \\ วิภวํ นาภินนฺทติ. \\ สพฺพโส ตณฺหานํ ขยา, \\ อเสสวิราคนิโรโธ นิพฺพานนฺ”ติ. \\ อยํ นิพฺเพโธ.}
\setlength{\csromangathaleftcol}{0pt}
\csromangathagroup{\csromangathastanza{\csromangathawak{ภวตณฺหา ปหียติ, \\ วิภวํ นาภินนฺทติ. \\ สพฺพโส ตณฺหานํ ขยา, \\ อเสส\-วิราค\-นิโรโธ นิพฺพานนฺ”ติ.}}\csromangathastanza{\csromangathawak{อยํ นิพฺเพโธ.}}}

\prose{** “ตสฺส นิพฺพุตสฺส ภิกฺขุโน,}

\prose{อนุปาทา ปุนพฺภโว น โหติ.}

\setlength{\csromangathapagewidth}{0pt}
\setlength{\csromangathawakwidth}{0pt}
\csromangathameasuregroup{}{อภิภูโต มาโร วิชิตสงฺคาโม, \\ อุปจฺจคา สพฺพภาวานิ ตาที”ติ.}
\csromangathameasurewak{อภิภูโต มาโร วิชิตสงฺคาโม, \\ อุปจฺจคา สพฺพภาวานิ ตาที”ติ.}
\setlength{\csromangathaleftcol}{0pt}
\csromangathagroup{\csromangathastanza{\csromangathawak{อภิภูโต มาโร วิชิตสงฺคาโม, \\ อุปจฺจคา สพฺพภาวานิ ตาที”ติ.}}}

\prose{อยํ อเสกฺโข. อิทํ สํกิเลสภาคิยญฺจ นิพฺเพธภาคิยญฺจ อเสกฺขภาคิยญฺจ สุตฺตํ.}

\prose{\csromansymbolfootnote{+}{อํ 1. 311 ปิฏฺเฐ.}จตฺตาโรเม ภิกฺขเว ปุคฺคลา. กตเม จตฺตาโร, อนุโสตคามี ปฏิโสตคามี ฐิตตฺโต ติณฺโณ ปารงฺคโต ถเล ติฏฺฐติ พฺราหฺมโณติ. ตตฺถ โยยํ ปุคฺคโล อนุโสตคามี, อยํ ปุคฺคโล สํกิเลส\-ภาคิโย. ตตฺถ โยยํ ปุคฺคโล ปฏิโสตคามี โย จ ฐิตตฺโต, อิเม เทฺว ปุคฺคลา นิพฺเพธ\-ภาคิยา. ตตฺถ โยยํ ปุคฺคโล ติณฺโณ ปารงฺคโต ถเล ติฏฺฐติ พฺราหฺมโณ, อยํ อเสกฺโข. อิทํ สํกิเลสภาคิยญฺจ นิพฺเพธภาคิยญฺจ อเสกฺขภาคิยญฺจ สุตฺตํ.}

\proseitem{110}{ตตฺถ กตมํ สํกิเลสภาคิยญฺจ วาสนาภาคิยญฺจ นิพฺเพธภาคิยญฺจ สุตฺตํ.}

\prose{2 ฉฬาภิชาติโก อตฺถิ ปุคฺคโล กโณฺห กณฺหาภิชาติโก กณฺหํ ธมฺมํ อภิชายติ, อตฺถิ ปุคฺคโล กโณฺห กณฺหาภิชาติโก สุกฺกํ ธมฺมํ อภิชายติ, อตฺถิ ปุคฺคโล กโณฺห กณฺหาภิชาติโก อกณฺหํ อสุกฺกํ อกณฺห\-อสุกฺก\-วิปากํ อจฺจนฺตทิฏฺฐํ\csromansharedfootnote{fd68a29404a41ac8}{อนฺตํ นิฏฺฐํ (สี)} นิพฺพานํ อาราเธติ, \csromanfolio{137}อตฺถิ ปุคฺคโล สุกฺโก สุกฺกาภิชาติโก กณฺหํ ธมฺมํ อภิชายติ, อตฺถิ ปุคฺคโล สุกฺโก สุกฺกาภิชาติโก สุกฺกํ ธมฺมํ อภิชายติ, อตฺถิ ปุคฺคโล สุกฺโก สุกฺกาภิชาติโก อกณฺหํ อสุกฺกํ อกณฺห\-อสุกฺก\-วิปากํ อจฺจนฺตทิฏฺฐํ นิพฺพานํ อาราเธติ.}

\prose{ตตฺถ โย จ ปุคฺคโล กโณฺห กณฺหาภิชาติโก กณฺหํ ธมฺมํ อภิชายติ, โย จ ปุคฺคโล สุกฺโก สุกฺกาภิชาติโก กณฺหํ ธมฺมํ อภิชายติ, อิเม เทฺว ปุคฺคลา สํกิเลส\-ภาคิยา.}

\prose{ตตฺถ โย จ ปุคฺคโล กโณฺห กณฺหาภิชาติโก สุกฺกํ ธมฺมํ อภิชายติ, โย จ ปุคฺคโล สุกฺโก สุกฺกาภิชาติโก สุกฺกํ ธมฺมํ อภิชายติ, อิเม เทฺว ปุคฺคลา วาสนาภาคิยา.}

\prose{ตตฺถ โย จ ปุคฺคโล กโณฺห กณฺหาภิชาติโก อกณฺหํ อสุกฺกํ อกณฺห- อสุกฺกวิปากํ อจฺจนฺตทิฏฺฐํ นิพฺพานํ อาราเธติ, โย จ ปุคฺคโล สุกฺโก สุกฺกาภิชาติโก อกณฺหํ อสุกฺกํ อกณฺห\-อสุกฺก\-วิปากํ อจฺจนฺตทิฏฺฐํ นิพฺพานํ อาราเธติ, อิเม เทฺว ปุคฺคลา นิพฺเพธ\-ภาคิยา, อิทํ สํกิเลสภาคิยญฺจ วาสนาภาคิยญฺจ นิพฺเพธภาคิยญฺจ สุตฺตํ.}

\prose{จตฺตาริมานิ ภิกฺขเว กมฺมานิ, กตมานิ จตฺตาริ, อตฺถิ กมฺมํ กณฺหํ กณฺหวิปากํ, อตฺถิ กมฺมํ สุกฺกํ สุกฺกวิปากํ, อตฺถิ กมฺมํ กณฺหสุกฺกํ กณฺห\-สุกฺก\-วิปากํ, อตฺถิ กมฺมํ อกณฺหํ อสุกฺกํ อกณฺห\-อสุกฺก\-วิปากํ กมฺมุตฺตมํ กมฺมเสฏฺฐํ กมฺมกฺขยาย สํวตฺตติ\csromansharedfootnote{5bcd42ef2368153a}{อํ 1. 553 ปิฏฺเฐ.}.}

\prose{ตตฺถ ยญฺจ กมฺมํ กณฺหํ กณฺหวิปากํ, ยญฺจ กมฺมํ กณฺหสุกฺกํ กณฺห\-สุกฺก\-วิปากํ, อยํ สํกิเลโส. ยญฺจ กมฺมํ สุกฺกํ สุกฺกวิปากํ, อยํ วาสนา. ยญฺจ กมฺมํ อกณฺหํ อสุกฺกํ อกณฺห\-อสุกฺก\-วิปากํ กมฺมุตฺตมํ กมฺมเสฏฺฐํ กมฺมกฺขยาย สํวตฺตติ, อยํ นิพฺเพโธ. อิทํ สํกิเลสภาคิยญฺจ วาสนาภาคิยญฺจ นิพฺเพธภาคิยญฺจ สุตฺตํ.}

\setlength{\csromangathapagewidth}{0pt}
\setlength{\csromangathawakwidth}{0pt}
\csromangathameasuregroup{ตตฺถ กตมํ วาสนาภาคิยญฺจ, \\ ลทฺธาน มานุสตฺตํ เทฺว, \\ สุกิจฺจํ เจว ปุญฺญานิ,}{\csromangathabat{ตตฺถ กตมํ วาสนาภาคิยญฺจ,}{นิพฺเพธภาคิยญฺจ สุตฺตํ.} \\ \csromangathabat{ลทฺธาน มานุสตฺตํ เทฺว,}{กิจฺจํ อกิจฺจเมว จ.} \\ \csromangathabat{สุกิจฺจํ เจว ปุญฺญานิ,}{สํโยชนวิปฺปหานํ วาติ.}}
\csromangathasetleft{ตตฺถ กตมํ วาสนาภาคิยญฺจ, \\ ลทฺธาน มานุสตฺตํ เทฺว, \\ สุกิจฺจํ เจว ปุญฺญานิ,}
\csromangathagroup{\csromangathastanza{\csromangathaitembat{111}{ตตฺถ กตมํ วาสนาภาคิยญฺจ,}{นิพฺเพธภาคิยญฺจ สุตฺตํ.} \\ \csromangathacontbat{ลทฺธาน มานุสตฺตํ เทฺว,}{กิจฺจํ อกิจฺจเมว จ.} \\ \csromangathacontbat{สุกิจฺจํ เจว ปุญฺญานิ,}{สํโยชน\-วิปฺปหานํ วาติ.}}}

\prose{\csromanfolio{138}“สุกิจฺจํ เจว ปุญฺญานี”ติ วาสนา. “สํโยชน\-วิปฺปหานํ วา”ติ นิพฺเพโธ.}

\setlength{\csromangathapagewidth}{0pt}
\setlength{\csromangathawakwidth}{0pt}
\csromangathameasuregroup{ปุญฺญานิ กริตฺวาน, \\ สํโยชนปฺปหานา,}{\csromangathabat{ปุญฺญานิ กริตฺวาน,}{สคฺคา สคฺคํ วชนฺติ กตปุญฺญา.} \\ \csromangathabat{สํโยชนปฺปหานา,}{ชรามรณา วิปฺปมุจฺจนฺตีติ.}}
\csromangathasetleft{ปุญฺญานิ กริตฺวาน, \\ สํโยชนปฺปหานา,}
\csromangathagroup{\csromangathastanza{\csromangathabat{ปุญฺญานิ กริตฺวาน,}{สคฺคา สคฺคํ วชนฺติ กตปุญฺญา.} \\ \csromangathabat{สํโยชน\-ปฺปหานา,}{ชรามรณา วิปฺปมุจฺจนฺตีติ.}}}

\prose{“ปุญฺญานิ กริตฺวาน, สคฺคา สคฺคํ วชนฺติ กตปุญฺญา”ติ วาสนา. “สํโยชน\-ปฺปหานา ชรามรณา วิปฺปมุจฺจนฺตี”ติ นิพฺเพโธ. อิทํ วาสนาภาคิยญฺจ นิพฺเพธภาคิยญฺจ สุตฺตํ.}

\prose{เทฺวมานิ ภิกฺขเว ปธานานิ\csromansharedfootnote{e2aff0a308b1752a}{อํ 1. 51 ปิฏฺเฐปิ ปสฺสิตพฺพํ.}, กตมานิ เทฺว, โย จ อคารสฺมา อนคาริยํ ปพฺพชิเตสุ จีวร\-ปิณฺฑปาต\-เสนาสนคิลานปจฺจย- เภสชฺช\-ปริกฺขารํ ปริจฺจชติ, โย จ อคารสฺมา อนคาริยํ ปพฺพชิเตสุ สพฺพู\-ปธิ\-ปฏินิสฺสคฺโค ตณฺหกฺขโย วิราโค นิโรโธ นิพฺพานนฺติ, ตตฺถ โย อคารสฺมา อนคาริยํ ปพฺพชิเตสุ จีวร\-ปิณฺฑปาต\-เสนาสนคิลานปจฺจย- เภสชฺช\-ปริกฺขารํ ปริจฺจชติ, อยํ วาสนา.}

\prose{โย อคารสฺมา อนคาริยํ ปพฺพชิเตสุ สพฺพู\-ปธิ\-ปฏินิสฺสคฺโค ตณฺหกฺขโย วิราโค นิโรโธ นิพฺพานํ, อยํ นิพฺเพโธ. อิทํ วาสนาภาคิยญฺจ นิพฺเพธภาคิยญฺจ สุตฺตํ.}

\prose{ตตฺถ ตณฺหา\-สํกิเลส\-ภาคิยํ สุตฺตํ ตณฺหา\-ปกฺเขเนว นิทฺทิสิตพฺพํ ตีหิ ตณฺหาหิ กามตณฺหาย ภวตณฺหาย วิภวตณฺหาย. เยน เยน วา ปน วตฺถุนา อชฺโฌสิตา, เตน เตเนว นิทฺทิสิตพฺพํ, ตสฺสา วิตฺถาโร ฉตฺติํ\-สตณฺหา\-ชาลินิยา วิจริตานิ.}

\prose{ตตฺถ ทิฏฺฐิ\-สํกิเลส\-ภาคิยํ สุตฺตํ ทิฏฺฐิ\-ปกฺเขเนว นิทฺทิสิตพฺพํ อุจฺเฉท\-สสฺสเตน, เยน เยน วา ปน วตฺถุนา ทิฏฺฐิวเสน อภินิวิสติ “อิทเมว สจฺจํ โมฆมญฺญนฺ”ติ, เตน เตเนว นิทฺทิสิตพฺพํ, ตสฺสา วิตฺถาโร ทฺวาสฏฺฐิ\-ทิฏฺฐิคตานิ.}

\prose{ตตฺถ ทุจฺจริต\-สํกิเลส\-ภาคิยํ สุตฺตํ เจตนาย เจตสิก\-กมฺเมน นิทฺทิสิตพฺพํ ตีหิ ทุจฺจริเตหิ กาย\-ทุจฺจริเตน วจีทุจฺจริเตน มโน\-ทุจฺจริเตน, ตสฺส วิตฺถาโร ทสอกุสล\-กมฺมปถา.}

\prose{\csromanfolio{139}ตตฺถ ตณฺหา\-โวทาน\-ภาคิยํ สุตฺตํ สมเถน นิทฺทิสิตพฺพํ, ทิฏฺฐิ\-โวทาน\-ภาคิยํ สุตฺตํ วิปสฺสนาย นิทฺทิสิตพฺพํ, ทุจฺจริต\-โวทาน\-ภาคิยํ สุตฺตํ สุจริเตน นิทฺทิสิตพฺพํ, ตีณิ อกุสลมูลานิ, ตํ กิสฺส เหตุ, สํสารสฺส นิพฺพตฺติยา, ตถา นิพฺพตฺเต สํสาเร กายทุจฺจริตํ กายสุจริตํ วจีทุจฺจริตํ วจีสุจริตํ มโนทุจฺจริตํ มโนสุจริตํ อิมินา อสุเภน กมฺมวิปาเกน อิทํ พาลลกฺขณํ นิพฺพตฺตตีติ. อิทํ สํกิเลส\-ภาคิยํ สุตฺตํ.}

\prose{อิมินา สุเภน กมฺมวิปาเกน อิทํ มหา\-ปุริส\-ลกฺขณํ นิพฺพตฺตตีติ อิทํ วาสนาภาคิยํ สุตฺตํ.}

\prose{ตตฺถ สํกิเลส\-ภาคิยํ สุตฺตํ จตูหิ กิเลสภูมีหิ นิทฺทิสิตพฺพํ อนุสยภูมิยา ปริยุฏฺฐาน\-ภูมิยา สํโยชน\-ภูมิยา อุปาทานภูมิยา, สานุสยสฺส ปริยุฏฺฐานํ ชายติ, ปริยุฏฺฐิโต สํยุชฺชติ, สํยุชฺชนฺโต อุปาทิยติ, อุปาทานปจฺจยา ภโว, ภวปจฺจยา ชาติ, ชาติปจฺจยา ชรามรณํ โสก\-ปริเทว\-ทุกฺข\-โทมนสฺสุ\-ปายาสา สมฺภวนฺติ, เอวเมตสฺส เกวลสฺส ทุกฺข\-กฺขนฺธสฺส สมุทโย โหติ. อิมาหิ จตูหิ กิเลสภูมีหิ สพฺเพ กิเลสา สงฺคหํ สโมสรณํ คจฺฉนฺติ, อิทํ สํกิเลส\-ภาคิยํ สุตฺตํ.}

\prose{วาสนาภาคิยํ สุตฺตํ ตีหิ สุจริเตหิ นิทฺทิสิตพฺพํ, นิพฺเพธ\-ภาคิยํ สุตฺตํ จตูหิ สจฺเจหิ นิทฺทิสิตพฺพํ, อเสกฺข\-ภาคิยํ สุตฺตํ ตีหิ ธมฺเมหิ นิทฺทิสิตพฺพํ พุทฺธธมฺเมหิ ปจฺเจก\-พุทฺธธมฺเมหิ สาวกภูมิยา, ฌายิวิสเย นิทฺทิสิตพฺพนฺติ.}

\proseitem{112}{ตตฺถ กตเม อฏฺฐารส มูลปทา, โลกิยํ โลกุตฺตรํ โลกิยญฺจ โลกุตฺตรญฺจ, สตฺตา\-ธิฏฺฐานํ ธมฺมา\-ธิฏฺฐานํ สตฺตา\-ธิฏฺฐานญฺจ ธมฺมาธิฏฺฐานญฺจ, ญาณํ เญยฺยํ ญาณญฺจ เญยฺยญฺจ, ทสฺสนํ ภาวนา ทสฺสนญฺจ ภาวนา จ, สกวจนํ ปรวจนํ สกวจนญฺจ ปรวจนญฺจ, วิสชฺชนียํ อวิสชฺชนียํ วิสชฺชนียญฺจ อวิสชฺชนียญฺจ, กมฺมํ วิปาโก กมฺมญฺจ วิปาโก จ, กุสลํ อกุสลํ กุสลญฺจ อกุสลญฺจ, อนุญฺญาตํ ปฏิกฺขิตฺตํ อนุญฺญาตญฺจ ปฏิกฺขิตฺตญฺจ, ถโว จาติ.}

\prose{\csromanfolio{140}ตตฺถ กตมํ โลกิยํ.}

\setlength{\csromangathapagewidth}{0pt}
\setlength{\csromangathawakwidth}{0pt}
\csromangathameasuregroup{\textsuperscript{*}“น หิ ปาปํ กตํ กมฺมํ, \\ qอหนฺตํ\textsuperscript{1} พาลมเนฺวติ,}{\csromangathabat{\textsuperscript{*}“น หิ ปาปํ กตํ กมฺมํ,}{สชฺชุขีรํว มุจฺจติ.} \\ \csromangathabat{qอหนฺตํ\textsuperscript{1} พาลมเนฺวติ,}{ภสฺมจฺฉนฺโนว\textsuperscript{1} ปาวโกติ.} \\ อิทํ โลกิยํ.}
\csromangathameasurewak{อิทํ โลกิยํ.}
\csromangathasetleft{\textsuperscript{*}“น หิ ปาปํ กตํ กมฺมํ, \\ qอหนฺตํ\textsuperscript{1} พาลมเนฺวติ,}
\csromangathagroup{\csromangathastanza{\csromangathabat{\csromansymbolmark{*}“น หิ ปาปํ กตํ กมฺมํ,}{สชฺชุขีรํว มุจฺจติ.} \\ \csromangathabat{qอหนฺตํ\csromansharedfootnote{7b43f88d7341dbad}{ภสฺมาฉนฺโนว (ก) ** เหฏฺฐา 107 ปิฏฺเฐปิ.} พาลมเนฺวติ,}{ภสฺม\-จฺฉนฺโนว\csromansharedfootnote{7b43f88d7341dbad}{ภสฺมาฉนฺโนว (ก) ** เหฏฺฐา 107 ปิฏฺเฐปิ.} ปาวโกติ.}}\csromangathastanza{\csromangathawak{อิทํ โลกิยํ.}}}

\prose{** “จตฺตาริมานิ ภิกฺขเว อคติคมนานิ สพฺพํ ฯเปฯ. นิหียเต ตสฺส ยโส กาฬปกฺเขว จนฺทิมา”ติ. อิทํ โลกิยํ.}

\prose{อฏฺฐิเม ภิกฺขเว โลกธมฺมา, กตเม อฏฺฐ, ลาโภ อลาโภ, ยโส อยโส, นินฺทา ปสํสา, สุขํ ทุกฺขํ. อิเม โข ภิกฺขเว อฏฺฐ โลกธมฺมาติ\csromansharedfootnote{7feb633210ce3efc}{ที 3. 215; อํ 1. 7 ปิฏฺเฐสุ ปสฺสิตพฺพํ.}. อิทํ โลกิยํ.}

\setlength{\csromangathapagewidth}{0pt}
\setlength{\csromangathawakwidth}{0pt}
\csromangathameasuregroup{}{ตตฺถ กตมํ โลกุตฺตรํ. \\ “ยสฺสินฺทฺริยานิ สมถงฺคตานิ.}
\csromangathameasurewak{ตตฺถ กตมํ โลกุตฺตรํ. \\ “ยสฺสินฺทฺริยานิ สมถงฺคตานิ.}
\setlength{\csromangathaleftcol}{0pt}
\csromangathagroup{\csromangathastanza{\csromangathawak{\csromangathaitemline{112}{ตตฺถ กตมํ โลกุตฺตรํ.} \\ \csromangathacontline{“ยสฺสินฺทฺริยานิ สมถงฺคตานิ.}}}}

\prose{อสฺสา ยถา สารถินา สุทนฺตา.}

\setlength{\csromangathapagewidth}{0pt}
\setlength{\csromangathawakwidth}{0pt}
\csromangathameasuregroup{}{ปหีนมานสฺส อนาสวสฺส, \\ เทวาปิ ตสฺส ปิหยนฺติ ตาทิโน”ติ\textsuperscript{1}. \\ อิทํ โลกุตฺตรํ.}
\csromangathameasurewak{ปหีนมานสฺส อนาสวสฺส, \\ เทวาปิ ตสฺส ปิหยนฺติ ตาทิโน”ติ\textsuperscript{1}. \\ อิทํ โลกุตฺตรํ.}
\setlength{\csromangathaleftcol}{0pt}
\csromangathagroup{\csromangathastanza{\csromangathawak{ปหีนมานสฺส อนาสวสฺส, \\ เทวาปิ ตสฺส ปิหยนฺติ ตาทิโน”ติ\csromansharedfootnote{b9a372f868a73b03}{ขุ 1. 27 ปิฏฺเฐ ธมฺมปเท; อุปริ 200 ปิฏฺเฐปิ.}. \\ อิทํ โลกุตฺตรํ.}}}

\prose{“ปญฺจิมานิ ภิกฺขเว อินฺทฺริยานิ โลกุตฺตรานิ, กตมานิ ปญฺจ, สทฺธินฺทฺริยํ วีริยินฺทฺริยํ สตินฺทฺริยํ สมาธินฺทฺริยํ ปญฺญินฺทฺริยํ. อิมานิ โข ภิกฺขเว ปญฺจินฺทฺริยานิ โลกุตฺตรานี”ติ. อิทํ โลกุตฺตรํ. (1)}

\prose{ตตฺถ กตมํ โลกิยญฺจ โลกุตฺตรญฺจ.}

\prose{“ลทฺธาน มานุสตฺตํ เทฺว, กิจฺจํ อกิจฺจเมว จา”ติ\csromansharedfootnote{f8284b1bb98de70a}{เหฏฺฐา 137, 138 ปิฏฺเฐสุ.} เทฺว คาถา. ยํ อิห “สุกิจฺจํ เจว ปุญฺญานี”ติ จ “ปุญฺญานิ กริตฺวาน, สคฺคา สคฺคํ วชนฺติ กตปุญฺญา”ติ จ. อิทํ โลกิยํ.}

\prose{ยํ อิห “สํโยชน\-วิปฺปหานํ วา”ติ จ “สํโยชน\-ปฺปหานา, ชรามรณา วิปฺปมุจฺจนฺตี”ติ จ, อิทํ โลกุตฺตรํ. อิทํ โลกิยญฺจ โลกุตฺตรญฺจ.}

\prose{\csromanfolio{141}วิญฺญาเณ เจ ภิกฺขเว อาหาเร สติ นามรูปสฺส อวกฺกนฺติ โหติ, นามรูปสฺส อวกฺกนฺติยา สติ ปุนพฺภโว โหติ, ปุนพฺภเว สติ ชาติ โหติ, ชาติยา สติ ชรามรณํ โสก\-ปริเทว\-ทุกฺข\-โทมนสฺสุ\-ปายาสา สมฺภวนฺติ. เอวเมตสฺส เกวลสฺส ทุกฺข\-กฺขนฺธสฺส สมุทโย โหติ.\csromansharedfootnote{0de00cc0b8e0ccbf}{สํ 1. 314 ปิฏฺเฐ.} เสยฺยถาปิ ภิกฺขเว มหารุกฺโข, ตสฺส ยานิ เจว มูลานิ อโธคมานิ ยานิ จ ติริยงฺคมานิ, สพฺพานิ ตานิ อุทฺธํ โอชํ อภิหรนฺติ. เอวํ หิ โส ภิกฺขเว มหารุกฺโข ตทาหาโร ตทุปาทาโน จิรํ ทีฆมทฺธานํ ติฏฺเฐยฺย\csromansharedfootnote{0de00cc0b8e0ccbf}{สํ 1. 314 ปิฏฺเฐ.}. เอวเมว โข ภิกฺขเว วิญฺญาเณ อาหาเร สติ นามรูปสฺส อวกฺกนฺติ โหติ สพฺพํ ฯเปฯ. เอวเมตสฺส เกวลสฺส ทุกฺข\-กฺขนฺธสฺส สมุทโย โหตีติ. อิทํ โลกิยํ.}

\prose{วิญฺญาเณ เจ ภิกฺขเว อาหาเร อสติ นามรูปสฺส อวกฺกนฺติ น โหติ, นามรูปสฺส อวกฺกนฺติยา อสติ ปุนพฺภโว น โหติ, ปุนพฺภเว อสติ ชาติ น โหติ, ชาติยา อสติ ชรามรณํ โสก\-ปริเทว\-ทุกฺข\-โทมนสฺสุ\-ปายาสา นิรุชฺฌนฺติ. เอวเมตสฺส เกวลสฺส ทุกฺข\-กฺขนฺธสฺส นิโรโธ โหติ\csromansharedfootnote{cf14146ee4f34943}{กทาลปิฏกํ (ก)}. เสยฺยถาปิ ภิกฺขเว มหารุกฺโข อถ ปุริโส อาคจฺเฉยฺย กุทฺทาลปิฏกํ\csromansharedfootnote{5a1ff6e23f69e345}{มูลํ (สํ 1. 314 ปิฏฺเฐ.)} อาทาย, โส ตํ รุกฺขํ มูเล ฉินฺเทยฺย, มูเล\csromansharedfootnote{e948e927b56d7dc1}{เฉตฺวา (สี, ก)} เฉตฺวา ปลิขเณยฺย, ปลิขณิตฺวา มูลานิ อุทฺธเรยฺย อนฺตมโส อุสีร\-นาฬิ\-มตฺตานิปิ. โส ตํ รุกฺขํ ขณฺฑาขณฺฑิกํ ฉินฺเทยฺย, ขณฺฑาขณฺฑิกํ ฉินฺทิตฺวา\csromansharedfootnote{d0a5ff33140bed22}{อนภาวํคโต (สี)} ผาเลยฺย, ผาเลตฺวา สกลิกํ สกลิกํ กเรยฺย, สกลิกํ สกลิกํ กริตฺวา วาตาตเป วิโสเสยฺย, วาตาตเป วิโสเสตฺวา อคฺคินา ฑเหยฺย, อคฺคินา ฑเหตฺวา มสิํ กเรยฺย, มสิํ กริตฺวา มหาวาเต วา โอผุเณยฺย, นทิยา วา สีฆโสตาย ปวาเหยฺย, เอวํ หิ โส ภิกฺขเว มหาริกฺโข อุจฺฉินฺนมูโล อสฺส ตาลวตฺถุกโต อนภาวํกโต\csromansharedfootnote{0de00cc0b8e0ccbf}{สํ 1. 314 ปิฏฺเฐ.} อายติํ อนุปฺปาทธมฺโม\csromansharedfootnote{0de00cc0b8e0ccbf}{สํ 1. 314 ปิฏฺเฐ.}. เอวเมว โข ภิกฺขเว วิญฺญาเณ อาหาเร อสติ นามรูปสฺส อวกฺกนฺติ น โหติ, นามรูปสฺส อวกฺกนฺติยา อสติ สพฺพํ ฯเปฯ. เอวเมตสฺส เกวลสฺส ทุกฺข\-กฺขนฺธสฺส นิโรโธ โหตีติ อิทํ โลกุตฺตรํ. อิทํ โลกิยญฺจ โลกุตฺตรญฺจ. (2)}

\proseitem{113}{\csromanfolio{142}ตตฺถ กตมํ สตฺตา\-ธิฏฺฐานํ.}

\setlength{\csromangathapagewidth}{0pt}
\setlength{\csromangathawakwidth}{0pt}
\csromangathameasuregroup{}{“สพฺพา ทิสา อนุปริคมฺม เจตสา, \\ เนวชฺฌคา ปิยตรมตฺตนา กฺวจิ. \\ เอวํ ปิโย ปุถุ อตฺตา ปเรสํ, \\ อิทํ สตฺตาธิฏฺฐานํ.}
\csromangathameasurewak{“สพฺพา ทิสา อนุปริคมฺม เจตสา, \\ เนวชฺฌคา ปิยตรมตฺตนา กฺวจิ. \\ เอวํ ปิโย ปุถุ อตฺตา ปเรสํ, \\ อิทํ สตฺตาธิฏฺฐานํ.}
\setlength{\csromangathaleftcol}{0pt}
\csromangathagroup{\csromangathastanza{\csromangathawak{“สพฺพา ทิสา อนุปริคมฺม เจตสา, \\ เนวชฺฌคา ปิยตรมตฺตนา กฺวจิ. \\ เอวํ ปิโย ปุถุ อตฺตา ปเรสํ,}}\csromangathastanza{\csromangathawak{อิทํ สตฺตา\-ธิฏฺฐานํ.}}}

\prose{ตสฺมา น หิํเส ปรมตฺตกาโม”ติ\csromansharedfootnote{b4a4ddfa8f86fd0f}{ปรํ อตฺตกาโมติ (สี) สํ 1. 75 ปิฏฺเฐ; ขุ 1. 133 ปิฏฺเฐ จ ปสฺสิตพฺพํ.}.}

\proseitem{113}{\csromansymbolfootnote{*}{ขุ 1. 134 ปิฏฺเฐ.}“เย เกจิ ภูตา ภวิสฺสนฺติ เย วาปิ\csromansharedfootnote{5ee0fcf1723bb47e}{จ (สี, ก)},}

\prose{สพฺเพ คมิสฺสนฺติ ปหาย เทตํ.}

\setlength{\csromangathapagewidth}{0pt}
\setlength{\csromangathawakwidth}{0pt}
\csromangathameasuregroup{}{ตํ สพฺพชานิํ กุสโล วิทิตฺวา, \\ อาตาปิโย\textsuperscript{1} พฺรหฺมจริยํ จเรยฺยา”ติ. \\ อิทํ สตฺตาธิฏฺฐานํ.}
\csromangathameasurewak{ตํ สพฺพชานิํ กุสโล วิทิตฺวา, \\ อาตาปิโย\textsuperscript{1} พฺรหฺมจริยํ จเรยฺยา”ติ. \\ อิทํ สตฺตาธิฏฺฐานํ.}
\setlength{\csromangathaleftcol}{0pt}
\csromangathagroup{\csromangathastanza{\csromangathawak{ตํ สพฺพชานิํ กุสโล วิทิตฺวา, \\ อาตาปิโย\csromansharedfootnote{338dd9d64fd7ee1e}{[มิสฺสินฺคฺ โนเต 0]} พฺรหฺมจริยํ จเรยฺยา”ติ. \\ อิทํ สตฺตา\-ธิฏฺฐานํ.}}}

\prose{** สตฺตหิ ภิกฺขเว องฺเคหิ สมนฺนาคตํ กลฺยาณมิตฺตํ อปิ วิเวจิยมาเนน ปณามิยมาเนน คเล ปิสน\-มชฺช\-มาเนน\csromansharedfootnote{1bdde453dac1d84c}{คเลปิ ปมชฺชมาเนน (สี)} ยาวชีวํ น วิชหิตพฺพํ. กตเมหิ สตฺตหิ, ปิโย จ โหติ มนาโป จ ครุ จ ภาวนีโย จ วตฺตา จ วจนกฺขโม จ คมฺภีรญฺจ กถํ กตฺตา โหติ, โน จ อฏฺฐาเน\csromansharedfootnote{70dac8025889ffa6}{น จ อฏฺฐาเน (สี, ก)} นิโยเชติ. อิเมหิ โข ภิกฺขเว สตฺตหิ ฯเปฯ น วิชหิตพฺพํ. อิทมโวจ ภควา, อิทํ วตฺวาน สุคโต. อถาปรํ เอตทโวจ สตฺถา\csromandash{}}

\setlength{\csromangathapagewidth}{0pt}
\setlength{\csromangathawakwidth}{0pt}
\csromangathameasuregroup{“ปิโย ครุ ภาวนีโย, \\ คมฺภีรญฺจ กถํ กตฺตา, \\ ตํ มิตฺตํ มิตฺตกาเมน,}{\csromangathabat{“ปิโย ครุ ภาวนีโย,}{วตฺตา จ วจนกฺขโม.} \\ \csromangathabat{คมฺภีรญฺจ กถํ กตฺตา,}{น จฏฺฐาเน นิโยชโก.} \\ \csromangathabat{ตํ มิตฺตํ มิตฺตกาเมน,}{ยาวชีวมฺปิ เสวิยนฺ”ติ.} \\ อิทํ สตฺตาธิฏฺฐานํ.}
\csromangathasetleft{“ปิโย ครุ ภาวนีโย, \\ คมฺภีรญฺจ กถํ กตฺตา, \\ ตํ มิตฺตํ มิตฺตกาเมน,}
\csromangathagroup{\csromangathastanza{\csromangathabat{“ปิโย ครุ ภาวนีโย,}{วตฺตา จ วจนกฺขโม.} \\ \csromangathabat{คมฺภีรญฺจ กถํ กตฺตา,}{น จฏฺฐาเน นิโยชโก.}}\csromangathastanza{\csromangathabat{ตํ มิตฺตํ มิตฺตกาเมน,}{ยาวชีวมฺปิ เสวิยนฺ”ติ.} \\ อิทํ สตฺตา\-ธิฏฺฐานํ.}}

\prose{ตตฺถ กตมํ ธมฺมา\-ธิฏฺฐานํ.}

\setlength{\csromangathapagewidth}{0pt}
\setlength{\csromangathawakwidth}{0pt}
\csromangathameasuregroup{“ยญฺจ กามสุขํ โลเก,}{\csromangathabat{“ยญฺจ กามสุขํ โลเก,}{ยญฺจิทํ ทิวิยํ สุขํ.} \\ อิทํ ธมฺมาธิฏฺฐานํ. \\ “สุสุขํ วต นิพฺพานํ, \\ สมฺมาสมฺพุทฺธเทสิตํ.}
\csromangathameasurewak{อิทํ ธมฺมาธิฏฺฐานํ. \\ “สุสุขํ วต นิพฺพานํ, \\ สมฺมาสมฺพุทฺธเทสิตํ.}
\csromangathasetleft{“ยญฺจ กามสุขํ โลเก,}
\csromangathagroup{\csromangathastanza{\csromangathabat{“ยญฺจ กามสุขํ โลเก,}{ยญฺจิทํ ทิวิยํ สุขํ.}}\csromangathastanza{\csromangathawak{อิทํ ธมฺมา\-ธิฏฺฐานํ. \\ “สุสุขํ วต นิพฺพานํ, \\ สมฺมาสมฺพุทฺธ\-เทสิตํ.}}}

\prose{ตณฺหกฺขยสุขสฺเสเต, กลํ นาคฺฆนฺติ โสฬสินฺ”ติ\csromansharedfootnote{64f7ccc33cd8e93f}{ขุ 1. 89 ปิฏฺเฐ; อุปริ 206, 209 ปิฏฺเฐสุปิ.}.}

\proseitem{113}{\csromanfolio{143}อโสกํ วิรชํ เขมํ, ยตฺถ ทุกฺขํ นิรุชฺฌตี”ติ\csromansharedfootnote{7b5a4e276102d698}{ขุ 2. 264 ปิฏฺเฐ.}. อิทํ ธมฺมา\-ธิฏฺฐานํ. (3)}

\prose{ตตฺถ กตมํ สตฺตา\-ธิฏฺฐานญฺจ ธมฺมาธิฏฺฐานญฺจ.}

\prose{\csromansymbolfootnote{*}{ขุ 1. 55 ปิฏฺเฐ ธมฺมปเท.}“มาตรํ ปิตรํ หนฺตฺวา, ราชาโน เทฺว จ ขตฺติเย.}

\prose{รฏฺฐํ สานุจรํ หนฺตฺวา”ติ.  อิทํ ธมฺมา\-ธิฏฺฐานํ.}

\prose{“อนีโฆ ยาติ พฺราหฺมโณ”ติ.  อิทํ สตฺตา\-ธิฏฺฐานํ.}

\prose{อิทํ สตฺตา\-ธิฏฺฐานญฺจ ธมฺมาธิฏฺฐานญฺจ.}

\prose{จตฺตาโรเม ภิกฺขเว อิทฺธิปาทา\csromansharedfootnote{693a669d1904515e}{สํ 3. 222 ปิฏฺฐาทีสุ อิทฺธิปาทสํยุตฺเต ปสฺสิตพฺพํ.}. กตเม จตฺตาโร, ฉนฺท\-สมาธิ\-ปธาน\-สงฺขาร\-สมนฺนาคโต อิทฺธิปาโท, วีริย ฯเปฯ จิตฺต. วีมํสา\-สมาธิ\-ปธาน\-สงฺขาร\-สมนฺนาคโต อิทฺธิปาโทติ, อิทํ ธมฺมา\-ธิฏฺฐานํ.}

\prose{โส กาเยปิ จิตฺตํ สโมทหติ, จิตฺเตปิ กายํ สโมทหติ, กาเย สุขสญฺญญฺจ ลหุสญฺญญฺจ โอกฺกมิตฺวา อุปสมฺปชฺช วิหรติ. อิทํ สตฺตา\-ธิฏฺฐานํ, อิทํ สตฺตา\-ธิฏฺฐานญฺจ ธมฺมาธิฏฺฐานญฺจ. (4)}

\proseitem{114}{ตตฺถ กตมํ ญาณํ.}

\setlength{\csromangathapagewidth}{0pt}
\setlength{\csromangathawakwidth}{0pt}
\csromangathameasuregroup{“ยํ ตํ โลกุตฺตรํ ญาณํ, \\ น ตสฺส ปริหาน’ตฺถิ, \\ “ปญฺญา หิ เสฏฺฐา โลกสฺมิํ, \\ ยาย สมฺมา ปชานาติ,}{\csromangathabat{“ยํ ตํ โลกุตฺตรํ ญาณํ,}{สพฺพญฺญู เยน วุจฺจติ.} \\ \csromangathabat{น ตสฺส ปริหาน’ตฺถิ,}{สพฺพกาเล ปวตฺตตี”ติ.} \\ อิทํ ญาณํ. \\ \csromangathabat{“ปญฺญา หิ เสฏฺฐา โลกสฺมิํ,}{ยายํ นิพฺพานคามินี\textsuperscript{1}.} \\ \csromangathabat{ยาย สมฺมา ปชานาติ,}{ชาติมรณสงฺขยนฺ”ติ.} \\ อิทํ ญาณํ.}
\csromangathameasurewak{อิทํ ญาณํ. \\ อิทํ ญาณํ.}
\csromangathasetleft{“ยํ ตํ โลกุตฺตรํ ญาณํ, \\ น ตสฺส ปริหาน’ตฺถิ, \\ “ปญฺญา หิ เสฏฺฐา โลกสฺมิํ, \\ ยาย สมฺมา ปชานาติ,}
\csromangathagroup{\csromangathastanza{\csromangathabat{“ยํ ตํ โลกุตฺตรํ ญาณํ,}{สพฺพญฺญู เยน วุจฺจติ.} \\ \csromangathabat{น ตสฺส ปริหาน’ตฺถิ,}{สพฺพกาเล ปวตฺตตี”ติ.}}\csromangathastanza{\csromangathawak{อิทํ ญาณํ.}}\csromangathastanza{\csromangathabat{“ปญฺญา หิ เสฏฺฐา โลกสฺมิํ,}{ยายํ นิพฺพานคามินี\csromansharedfootnote{338dd9d64fd7ee1e}{[มิสฺสินฺคฺ โนเต 0]}.} \\ \csromangathabat{ยาย สมฺมา ปชานาติ,}{ชาติมรณสงฺขยนฺ”ติ.}}\csromangathastanza{\csromangathawak{อิทํ ญาณํ.}}}

\prose{ตตฺถ กตมํ เญยฺยํ.}

\prose{** “กิตฺตยิสฺสามิ เต\csromansharedfootnote{40ced95b06cc6e51}{โว (สี, ก)} สนฺติํ, (โธตกาติ ภควา,)}

\prose{ทิฏฺเฐ ธมฺเม อนีติหํ.}

\setlength{\csromangathapagewidth}{0pt}
\setlength{\csromangathawakwidth}{0pt}
\csromangathameasuregroup{ยํ วิทิตฺวา สโต จรํ, \\ ตญฺจาหํ อภินนฺทามิ, \\ ยํ วิทิตฺวา สโต จรํ,}{\csromangathabat{ยํ วิทิตฺวา สโต จรํ,}{ตเร โลเก วิสตฺติกํ.} \\ \csromangathabat{ตญฺจาหํ อภินนฺทามิ,}{มเหสิ สนฺติมุตฺตมํ.} \\ \csromangathabat{ยํ วิทิตฺวา สโต จรํ,}{ตเร โลเก วิสตฺติกํ.}}
\csromangathasetleft{ยํ วิทิตฺวา สโต จรํ, \\ ตญฺจาหํ อภินนฺทามิ, \\ ยํ วิทิตฺวา สโต จรํ,}
\csromangathagroup{\csromangathastanza{\csromangathabat{ยํ วิทิตฺวา สโต จรํ,}{ตเร โลเก วิสตฺติกํ.} \\ \csromangathabat{ตญฺจาหํ อภินนฺทามิ,}{มเหสิ สนฺติมุตฺตมํ.} \\ \csromangathabat{ยํ วิทิตฺวา สโต จรํ,}{ตเร โลเก วิสตฺติกํ.}}}

\prose{\csromanfolio{144}ยํ กิญฺจิ สมฺปชานาสิ, (โธตกาติ ภควา,)}

\prose{อุทฺธํ อโธ ติริยญฺจาปิ มชฺเฌ.}

\setlength{\csromangathapagewidth}{0pt}
\setlength{\csromangathawakwidth}{0pt}
\csromangathameasuregroup{}{เอตํ วิทิตฺวา สงฺโคติ โลเก, \\ ภวาภวาย มากาสิ ตณฺหนฺ”ติ. \\ อิทํ เญยฺยํ.}
\csromangathameasurewak{เอตํ วิทิตฺวา สงฺโคติ โลเก, \\ ภวาภวาย มากาสิ ตณฺหนฺ”ติ. \\ อิทํ เญยฺยํ.}
\setlength{\csromangathaleftcol}{0pt}
\csromangathagroup{\csromangathastanza{\csromangathawak{เอตํ วิทิตฺวา สงฺโคติ โลเก, \\ ภวาภวาย มากาสิ ตณฺหนฺ”ติ. \\ อิทํ เญยฺยํ.}}}

\prose{“จตุนฺนํ ภิกฺขเว อริยสจฺจานํ อนนุโพธา อปฺปฏิเวธา เอวมิทํ ทีฆมทฺธานํ สนฺธาวิตํ สํสริตํ มมญฺเจว ตุมฺหากญฺจ ฯเปฯ. ตยิทํ ภิกฺขเว ทุกฺขํ อริยสจฺจํ อนุพุทฺธํ ปฏิวิทฺธํ, ทุกฺข\-สมุทยํ\csromansharedfootnote{055fd63993cc1971}{ทุกฺขสมุทโย (สี, ก)} อริยสจฺจํ อนุพุทฺธํ ปฏิวิทฺธํ, ทุกฺขนิโรธํ\csromansharedfootnote{f713f9b89b9d6039}{ทุกฺขนิโรโธ (สี, ก)} อริยสจฺจํ อนุพุทฺธํ ปฏิวิทฺธํ, ทุกฺขนิโรธ\-คามินี ปฏิปทา อริยสจฺจํ อนุพุทฺธํ ปฏิวิทฺธํ. อุจฺฉินฺนา ภวตณฺหา, ขีณา ภวเนตฺติ, นตฺถิทานิ ปุนพฺภโว”ติ. อิทมโวจ ภควา, อิทํ วตฺวาน สุคโต, อถาปรํ เอตทโวจ สตฺถา\csromandash{}}

\setlength{\csromangathapagewidth}{0pt}
\setlength{\csromangathawakwidth}{0pt}
\csromangathameasuregroup{จตุนฺนํ อริยสจฺจานํ, \\ สํสิตํ\textsuperscript{1} ทีฆมทฺธานํ, \\ ตานิ เอตานิ ทิฏฺฐานิ,}{\csromangathabat{จตุนฺนํ อริยสจฺจานํ,}{ยถาภูตํ อทสฺสนา.} \\ \csromangathabat{สํสิตํ\textsuperscript{1} ทีฆมทฺธานํ,}{ตาสุ ตาเสฺวว ชาติสุ.} \\ \csromangathabat{ตานิ เอตานิ ทิฏฺฐานิ,}{ภวเนตฺติ สมูหตา.}}
\csromangathasetleft{จตุนฺนํ อริยสจฺจานํ, \\ สํสิตํ\textsuperscript{1} ทีฆมทฺธานํ, \\ ตานิ เอตานิ ทิฏฺฐานิ,}
\csromangathagroup{\csromangathastanza{\csromangathabat{จตุนฺนํ อริยสจฺจานํ,}{ยถาภูตํ อทสฺสนา.} \\ \csromangathabat{สํสิตํ\csromansharedfootnote{3ca51b90b9bee6c8}{สํสริตํ (สี)} ทีฆมทฺธานํ,}{ตาสุ ตาเสฺวว ชาติสุ.} \\ \csromangathabat{ตานิ เอตานิ ทิฏฺฐานิ,}{ภวเนตฺติ สมูหตา.}}}

\prose{อุจฺฉินฺนํ มูลํ ทุกฺขสฺส, นตฺถิทานิ ปุนพฺภ โว”ติ\csromansharedfootnote{12fe2c7e35132ca9}{วิ 3. 325, 326; ที 277; สํ 3. 377, 378 ปิฏฺเฐสุ.}. อิทํ เญยฺยํ. (5)}

\prose{ตตฺถ กหมํ ญาณญฺจ เญยฺยญฺจ. รูปํ อนิจฺจํ, เวทนา อนิจฺจา, สญฺญา อนิจฺจา, สงฺขารา อนิจฺจา, วิญฺญาณํ อนิจฺจนฺติ. อิทํ เญยฺยํ.}

\prose{เอวํ ชานํ เอวํ ปสฺสํ อริยสาวโก “รูปํ อนิจฺจนฺ”ติ ปสฺสติ, “เวทนา อนิจฺจา”ติ ปสฺสติ, “สญฺญํ. สงฺขาเร. วิญฺญาณํ อนิจฺจนฺ”ติ ปสฺสตีติ. อิทํ ญาณํ.}

\prose{โส ปริมุจฺจติ รูเปน, ปริมุจฺจติ เวทนาย, ปริมุจฺจติ สญฺญาย, ปริมุจฺจติ สงฺขาเรหิ, ปริมุจฺจติ วิญฺญาณมฺหา, ปริมุจฺจติ ทุกฺขสฺมาติ วาทามีติ. อิทํ ญาณญฺจ เญยฺยญฺจ.}

\prose{\csromanfolio{145}“สพฺเพ สงฺขารา อนิจฺจา”ติ อิทํ เญยฺยํ. “ยทา ปญฺญาย ปสฺสตี”ติ อิทํ ญาณํ. “อถ นิพฺพินฺทติ ทุกฺเข เอส มคฺโค วิสุทฺธิยา”ติ อิทํ ญาณญฺจ เญยฺยญฺจ.}

\prose{“สพฺเพ สงฺขารา ทุกฺขา”ติ อิทํ เญยฺยํ. “ยทา ปญฺญาย ปสฺสตี”ติ อิทํ ญาณํ. “อถ นิพฺพินฺทติ ทุกฺเข เอส มคฺโค วิสุทฺธิยา”ติ อิทํ ญาณญฺจ เญยฺยญฺจ.}

\prose{“สพฺเพ ธมฺมา อนตฺตา”ติ อิทํ เญยฺยํ. “ยทา ปญฺญาย ปสฺสตี”ติ อิทํ ญาณํ. “อถ นิพฺพินฺทติ ทุกฺเข เอส มคฺโค วิสุทฺธิยา”ติ อิทํ ญาณญฺจ เญยฺยญฺจ.}

\prose{เย หิ เกจิ โสณ\csromansharedfootnote{1218570f1e275c27}{สํ 2. 40 ปิฏฺเฐ โสณสุตฺเต.} สมณา วา พฺราหฺมณา วา อนิจฺเจน รูเปน ทุกฺเขน วิปริณาม\-ธมฺเมน “เสยฺโยหมสฺมี”ติ วา สมนุปสฺสนฺติ, “สทิโสหมสฺมี”ติ วา สมนุปสฺสนฺติ, “หีโนหมสฺมี”ติ วา สมนุปสฺสนฺติ. กิมญฺญตฺร ยถาภูตสฺส อทสฺสนา. อนิจฺจาย เวทนาย ฯเปฯ อนิจฺจาย สญฺญาย. อนิจฺเจหิ สงฺขาเรหิ. อนิจฺเจน วิญฺญาเณน ทุกฺเขน วิปริณาม\-ธมฺเมน “เสยฺโยหมสฺมี”ติ วา สมนุปสฺสนฺติ, “สทิโสหมสฺมี”ติ วา สมนุปสฺสนฺติ, “หีโนหมสฺมี”ติ วา สมนุปสฺสนฺติ, กิมญฺญตฺร ยถาภูตสฺส อทสฺสนาติ. อิทํ เญยฺยํ.}

\prose{เย จ โข เกจิ โสณ สมณา วา พฺราหฺมณา วา อนิจฺเจน รูเปน ทุกฺเขน วิปริณาม\-ธมฺเมน “เสยฺโยหมสฺมี”ติปิ น สมนุปสฺสนฺติ, “สทิโสหมสฺมี”ติปิ น สมนุปสฺสนฺติ, “หีโนหมสฺมี”ติปิ น สมนุปสฺสนฺติ, กิมญฺญตฺร ยถาภูตสฺส ทสฺสนา. อนิจฺจาย เวทนาย ฯเปฯ อนิจฺจาย สญฺญาย. อนิจฺเจหิ สงฺขาเรหิ. อนิจฺเจน วิญฺญาเณน ทุกฺเขน วิปริณาม\-ธมฺเมน “เสยฺโยหมสฺมี”ติปิ น สมนุปสฺสนฺติ, “สทิโสหมสฺมี”ติปิ น สมนุปสฺสนฺติ, “หีโนหมสฺมี”ติปิ น สมนุปสฺสนฺติ, กิมญฺญตฺร ยถาภูตสฺส ทสฺสนาติ. อิทํ ญาณํ. อิทํ ญาณญฺจ เญยฺยญฺจ. (6)}

\prose{ตตฺถ กตมํ ทสฺสนํ.}

\proseitem{115}{“เย อริยสจฺจานิ วิภาวยนฺติ,}

\prose{คมฺภีร\-ปญฺเญน สุเทสิตานิ.}

\setlength{\csromangathapagewidth}{0pt}
\setlength{\csromangathawakwidth}{0pt}
\csromangathameasuregroup{}{กิญฺจาปิ เต โหนฺติ ภุสํ ปมตฺตา\textsuperscript{1}, \\ อิทํ ทสฺสนํ. \\ \textsuperscript{*}“ยถินฺทขีโล ปถวิสฺสิโต สิยา, \\ จตุพฺภิ วาเตหิ อสมฺปกมฺปิโย. \\ ตถูปมํ สปฺปุริสํ วทามิ, \\ โย อริยสจฺจานิ อเวจฺจ ปสฺสตี”ติ. \\ อิทํ ทสฺสนํ.}
\csromangathameasurewak{กิญฺจาปิ เต โหนฺติ ภุสํ ปมตฺตา\textsuperscript{1}, \\ อิทํ ทสฺสนํ. \\ \textsuperscript{*}“ยถินฺทขีโล ปถวิสฺสิโต สิยา, \\ จตุพฺภิ วาเตหิ อสมฺปกมฺปิโย. \\ ตถูปมํ สปฺปุริสํ วทามิ, \\ โย อริยสจฺจานิ อเวจฺจ ปสฺสตี”ติ. \\ อิทํ ทสฺสนํ.}
\setlength{\csromangathaleftcol}{0pt}
\csromangathagroup{\csromangathastanza{\csromangathawak{กิญฺจาปิ เต โหนฺติ ภุสํ ปมตฺตา\csromansharedfootnote{338dd9d64fd7ee1e}{[มิสฺสินฺคฺ โนเต 0]},}}\csromangathastanza{\csromangathawak{อิทํ ทสฺสนํ. \\ \csromansymbolmark{*}“ยถินฺทขีโล ปถวิสฺสิโต สิยา, \\ จตุพฺภิ วาเตหิ อสมฺปกมฺปิโย. \\ ตถูปมํ สปฺปุริสํ วทามิ, \\ โย อริยสจฺจานิ อเวจฺจ ปสฺสตี”ติ. \\ อิทํ ทสฺสนํ.}}}

\prose{น เต ภวํ อฏฺฐมมาทิยนฺตี”ติ\csromansharedfootnote{7cb0c83b51fadf10}{ภุสปฺปมตฺตา (สี)}.}

\proseitem{115}{\csromanfolio{146}จตูหิ ภิกฺขเว โสตาปตฺติยงฺเคหิ สมนฺนาคโต อริยสาวโก อากงฺขมาโน อตฺตนาว อตฺตานํ พฺยากเรยฺย “ขีณนิรโยมฺหิ, ขีณติรจฺฉานโยนิ, ขีณเปตฺติวิสโย, ขีณาปายเทคฺคติวินิปาโต, โสตาปนฺโนหมสฺมิ อวินิปาต\-ธมฺโม นิยโต สมฺโพธิ\-ปรายโณ, สตฺตกฺขตฺตุปรมํ\csromansharedfootnote{4a61fc517eaef33d}{สตฺตกฺขตฺตุปรโม (สี)} เทเว จ มนุสฺเส จ สนฺธาวิตฺวา สํสริตฺวา ทุกฺขสฺสนฺตํ กริสฺสามี”ติ. กตเมหิ จตูหิ, อิธ ภิกฺขเว อริยสาวกสฺส ตถาคเต สทฺธา นิวิฏฺฐา ปติฏฺฐิตา วิรูฬฺหา มูลชาตา อสํหาริยา สมเณน วา พฺราหฺมเณน วา เทเวน วา มาเรน วา พฺรหฺมุนา วา เกนจิ วา โลกสฺมิํ สห ธมฺเมน, ธมฺเม โข ปน นิฏฺฐํ คโต โหติ, สฺวากฺขาโต ภควตา ธมฺโม สนฺทิฏฺฐิโก อกาลิโก เอหิปสฺสิโก โอปเนยฺยิโก ปจฺจตฺตํ เวทิตพฺโพ วิญฺญูหิ, ยทิทํ มทนิมฺมทโน ฯเปฯ นิโรโธ นิพฺพานํ, สห ธมฺมิยา โข ปนสฺส โหนฺติ อิฏฺฐา กนฺตา ปิยา มนาปา คิหี เจว ปพฺพชิตา จ. อริยกนฺเตหิ โข ปน สีเลหิ สมนฺนาคโต โหติ อขณฺเฑหิ อจฺฉิทฺเทหิ อสพเลหิ อกมฺมาเสหิ ภุชิสฺเสหิ วิญฺญุ\-ปฺปสฏฺเฐหิ อปรามฏฺเฐหิ สมาธิ\-สํวตฺตนิเกหิ. อิเมหิ โข ภิกฺขเว จตูหิ โสตาปตฺติยงฺเคหิ สมนฺนาคโต อริยสาวโก อากงฺขมาโน อตฺตนาว อตฺตานํ พฺยากเรยฺย “ขีณนิรโยมฺหิ, ขีณติรจฺฉานโยนิ, ขีณเปตฺติวิสโย, ขีณา\-ปายทุคฺคติวินิปาโต, โสตาปนฺโนหมสฺมิ อวินิปาต\-ธมฺโม นิยโต สมฺโพธิ\-ปรายโณ, สตฺตกฺขตฺตุปรมํ เทเว จ มนุสฺเส จ สนฺธาวิตฺวา สํสริตฺวา ทุกฺขสฺสนฺตํ กริสฺสามี”ติ. อิทํ ทสฺสนํ.}

\prose{\csromanfolio{147}ตตฺถ กตมา ภาวนา.}

\setlength{\csromangathapagewidth}{0pt}
\setlength{\csromangathawakwidth}{0pt}
\csromangathameasuregroup{}{\textsuperscript{*}“ยสฺสินฺทฺริยานิ ภาวิตานิ\textsuperscript{1}, \\ อชฺฌตฺตํ พหิทฺธา จ สพฺพโลเก. \\ นิพฺพิชฺฌ อิมํ ปรญฺจ โลกํ, \\ กาลํ กงฺขติ ภาวิโต ส ทนฺโต”ติ. \\ อยํ ภาวนา.}
\csromangathameasurewak{\textsuperscript{*}“ยสฺสินฺทฺริยานิ ภาวิตานิ\textsuperscript{1}, \\ อชฺฌตฺตํ พหิทฺธา จ สพฺพโลเก. \\ นิพฺพิชฺฌ อิมํ ปรญฺจ โลกํ, \\ กาลํ กงฺขติ ภาวิโต ส ทนฺโต”ติ. \\ อยํ ภาวนา.}
\setlength{\csromangathaleftcol}{0pt}
\csromangathagroup{\csromangathastanza{\csromangathawak{\csromansymbolmark{*}“ยสฺสินฺทฺริยานิ ภาวิตานิ\csromansharedfootnote{338dd9d64fd7ee1e}{[มิสฺสินฺคฺ โนเต 0]}, \\ อชฺฌตฺตํ พหิทฺธา จ สพฺพโลเก. \\ นิพฺพิชฺฌ อิมํ ปรญฺจ โลกํ, \\ กาลํ กงฺขติ ภาวิโต ส ทนฺโต”ติ.}}\csromangathastanza{\csromangathawak{อยํ ภาวนา.}}}

\prose{“จตฺตาริมานิ ภิกฺขเว ธมฺมปทานิ. กตมานิ จตฺตาริ, อนภิชฺฌา ธมฺมปทํ, อพฺยาปาโท ธมฺมปทํ, สมฺมาสติ ธมฺมปทํ, สมฺมาสมาธิ ธมฺมปทํ, อิมานิ โข ภิกฺขเว จตฺตาริ ธมฺมปทานี”ติ\csromansharedfootnote{2d7a8131550f4057}{อํ 1. 337 ปิฏฺเฐ ปสฺสิตพฺพํ.}. อยํ ภาวนา. (7)}

\prose{ตตฺถ กตมํ ทสฺสนญฺจ ภาวนา จ “ปญฺจ ฉินฺเท ปญฺจ ชเห”ติ อิทํ ทสฺสนํ. “ปญฺจ จุตฺตริ ภาวเย. ปญฺจ สงฺคาติโค ภิกฺขุ, โอฆติณฺโณติ วุจฺจตี”ติ\csromansharedfootnote{65c81cb5613c5556}{ขุ 1. 66 ปิฏฺเฐ. 4-4. สํ 3. 180 ปิฏฺเฐ.}. อยํ ภาวนา. อิทํ ทสฺสนญฺจ ภาวนา จ.}

\prose{4 ตีณิมานิ ภิกฺขเว อินฺทฺริยานิ, กตมานิ ตีณิ, อนญฺญาต\-ญฺญสฺสามี\-ตินฺทฺริยํ อญฺญินฺทฺริยํ อญฺญาตาวิ\-นฺทฺริยํ4. กตมญฺจ ภิกฺขเว อนญฺญาต\-ญฺญสฺสามี\-ตินฺทฺริยํ, อิธ ภิกฺขเว ภิกฺขุ อนภิสเมตสฺส ทุกฺขสฺส อริยสจฺจสฺส อภิสมยาย ฉนฺทํ ชเนติ วายมติ วีริยํ อารภติ จิตฺตํ ปคฺคณฺหาติ ปทหติ, อนภิสเมตสฺส ทุกฺข\-สมุทยสฺส อริยสจฺจสฺส ฯเปฯ ทุกฺข\-นิโรธสฺส ฯเปฯ ทุกฺขนิโรธ\-คามินิยา ปฏิปทาย อริยสจฺจสฺส อภิสมายาย ฉนฺทํ ชเนติ วายมติ วีริยํ อารภติ จิตฺตํ ปคฺคณฺหาติ ปทหติ. อิทํ ภิกฺขเว อนญฺญาต\-ญฺญสฺสามี\-ตินฺทฺริยนฺติ. อิทํ ทสฺสนํ.}

\prose{กตมญฺจ ภิกฺขเว อญฺญินฺทฺริยํ, อิธ ภิกฺขเว ภิกฺขุ “อิทํ ทุกฺขนฺ”ติ ยถาภูตํ ปชานาติ, “อยํ ทุกฺขสมุทโย”ติ ยถาภูตํ ปชานาติ, “อยํ ทุกฺขนิโรโธ”ติ ฯเปฯ “อยํ ทุกฺขนิโรธ\-คามินี ปฏิปทา”ติ ยถาภูตํ ปชานาติ. อิทํ ภิกฺขเว อญฺญินฺทฺริยํ.}

\prose{กตมญฺจ ภิกฺขเว อญฺญาตาวิ\-นฺทฺริยํ, อิธ ภิกฺขเว ภิกฺขุ อาสวานํ ขยา อนาสวํ เจโตวิมุตฺติํ ปญฺญาวิมุตฺติํ ทิฏฺเฐว ธมฺเม สยํ อภิญฺญา สจฺฉิกตฺวา อุปสมฺปชฺช วิหรติ, ขีณา ชาติ, วุสิตํ พฺรหฺมจริยํ, กตํ \csromanfolio{148}กรณียํ นาปรํ อิตฺถตฺตายาติ ปชานาติ. อิทํ ภิกฺขเว อญฺญาตาวิ\-นฺทฺริยนฺติ. อยํ ภาวนา. อิทํ ทสฺสนญฺจ ภาวนา จ. (8)}

\proseitem{116}{ตตฺถ กตมํ สกวจนํ.}

\prose{\csromansymbolfootnote{*}{เหฏฺฐา 37, 68 ปิฏฺเฐสุ; อุปริ 161, 205, 208 ปิฏฺเฐสุปิ.}“สพฺพปาปสฺส อกรณํ, กุสลสฺส อุปสมฺปทา.}

\setlength{\csromangathapagewidth}{0pt}
\setlength{\csromangathawakwidth}{0pt}
\csromangathameasuregroup{สจิตฺตปริโยทาปนํ,}{\csromangathabat{สจิตฺตปริโยทาปนํ,}{เอตํ พุทฺธาน สาสนนฺ”ติ.} \\ อิทํ สกวจนํ.}
\csromangathasetleft{สจิตฺตปริโยทาปนํ,}
\csromangathagroup{\csromangathastanza{\csromangathabat{สจิตฺต\-ปริโยทาปนํ,}{เอตํ พุทฺธาน สาสนนฺ”ติ.} \\ อิทํ สกวจนํ.}}

\prose{\csromansymbolfootnote{+}{ม 3. 203; อํ 1. 100 ปิฏฺเฐสุ ปสฺสิตพฺพํ.}ตีณิมานิ ภิกฺขเว พาลสฺส พาลลกฺขณานิ พาลนิมิตฺตานิ พาลาปทานานิ, เยหิ พาลํ พาโลติ ปเร สญฺชานนฺติ. กตมานิ ตีณิ, พาโล ภิกฺขเว ทุจฺจินฺติต\-จินฺตี จ โหติ, ทุพฺภาสิต\-ภาสี จ โหติ, ทุกฺกฏ\-กมฺม\-การี\csromansharedfootnote{8ac7796405f9903c}{ทุกฺกตกมฺมการี (สี) ม 3. 201 ปิฏฺเฐ; อํ 1. 100 ปิฏฺเฐ จ ปสฺสิตพฺพํ. ++ ม 3. 208; อํ 1. 101 ปิฏฺเฐสุ ปสฺสิตพฺพํ.} จ โหติ. อิมานิ โข ภิกฺขเว ตีณิ พาลสฺส พาลลกฺขณานิ พาลนิมิตฺตานิ พาลาปทานานิ.}

\prose{++ ตีณิมานิ ภิกฺขเว ปณฺฑิตสฺส ปณฺฑิต\-ลกฺขณานิ ปณฺฑิต\-นิมิตฺตานิ ปณฺฑิตา\-ปทานานิ, เยหิ ปณฺฑิตํ ปณฺฑิโตติ ปเร สญฺชานนฺติ. กตมานิ ตีณิ, ปณฺฑิโต ภิกฺขเว สุจินฺติตจินฺตี จ โหติ, สุภาสิตภาสี จ โหติ, สุกต\-กมฺม\-การี จ โหติ. อิมานิ โข ภิกฺขเว ตีณิ ปณฺฑิตสฺส ปณฺฑิต\-ลกฺขณานิ ปณฺฑิต\-นิมิตฺตานิ ปณฺฑิตาปทานานีติ. อิทํ สกวจนํ.}

\prose{ตตฺถ กตมํ ปรวจนํ.}

\setlength{\csromangathapagewidth}{0pt}
\setlength{\csromangathawakwidth}{0pt}
\csromangathameasuregroup{}{“ปถวีสโม นตฺถิ วิตฺถโต, \\ นินฺโน ปาตาลสโม น วิชฺชติ. \\ เมรุสโม นตฺถิ อุนฺนโต, \\ จกฺกวตฺติสทิโส นตฺถิ โปริโส”ติ. \\ อิทํ ปรวจนํ.}
\csromangathameasurewak{“ปถวีสโม นตฺถิ วิตฺถโต, \\ นินฺโน ปาตาลสโม น วิชฺชติ. \\ เมรุสโม นตฺถิ อุนฺนโต, \\ จกฺกวตฺติสทิโส นตฺถิ โปริโส”ติ. \\ อิทํ ปรวจนํ.}
\setlength{\csromangathaleftcol}{0pt}
\csromangathagroup{\csromangathastanza{\csromangathawak{“ปถวีสโม นตฺถิ วิตฺถโต, \\ นินฺโน ปาตาลสโม น วิชฺชติ. \\ เมรุสโม นตฺถิ อุนฺนโต, \\ จกฺกวตฺติสทิโส นตฺถิ โปริโส”ติ.}}\csromangathastanza{\csromangathawak{อิทํ ปรวจนํ.}}}

\prose{\csromanfolio{149}** โหตุ เทวานมินฺท สุภาสิเตน ชโยติ. โหตุ เวปจตฺติ สุภาสิเตน ชโยติ. ภณ เวปจิตฺติ คาถนฺติ. อถ โข ภิกฺขเว เวปจิตฺติ อสุรินฺโท อิมํ คาถํ อภาสิ\csromandash{}}

\setlength{\csromangathapagewidth}{0pt}
\setlength{\csromangathawakwidth}{0pt}
\csromangathameasuregroup{“ภิยฺโย พาลา ปกุชฺเฌยฺยุํ\textsuperscript{1}, \\ ตสฺมา ภุเสน ทณฺเฑน,}{\csromangathabat{“ภิยฺโย พาลา ปกุชฺเฌยฺยุํ\textsuperscript{1},}{โน จสฺส ปฏิเสธโก.} \\ \csromangathabat{ตสฺมา ภุเสน ทณฺเฑน,}{ธีโร พาลํ นิเสธเย”ติ.}}
\csromangathasetleft{“ภิยฺโย พาลา ปกุชฺเฌยฺยุํ\textsuperscript{1}, \\ ตสฺมา ภุเสน ทณฺเฑน,}
\csromangathagroup{\csromangathastanza{\csromangathabat{“ภิยฺโย พาลา ปกุชฺเฌยฺยุํ\csromansharedfootnote{7a48634cb5ca9244}{ปภิชฺเชยฺยุํ (ก)},}{โน จสฺส ปฏิเสธโก.} \\ \csromangathabat{ตสฺมา ภุเสน ทณฺเฑน,}{ธีโร พาลํ นิเสธเย”ติ.}}}

\prose{ภาสิตาย โข ปน ภิกฺขเว เวปจิตฺตินา อสุรินฺเทน คาถาย อสุรา อนุโมทิํสุ, เทวา ตุณฺหี อเหสุํ. อถ โข ภิกฺขเว เวปจิตฺติ อสุรินฺโท สกฺกํ เทวานมินฺทํ เอตทโวจ “ภณ เทวานมินฺท คาถนฺ”ติ. เอวํ วุตฺเต ภิกฺขเว สกฺโก เทวานมินฺโท อิมํ คาถํ อภาสิ\csromandash{}}

\setlength{\csromangathapagewidth}{0pt}
\setlength{\csromangathawakwidth}{0pt}
\csromangathameasuregroup{“เอตเทว อหํ มญฺเญ, \\ ปรํ สงฺกุปิตํ ญตฺวา,}{\csromangathabat{“เอตเทว อหํ มญฺเญ,}{พาลสฺส ปฏิเสธนํ.} \\ \csromangathabat{ปรํ สงฺกุปิตํ ญตฺวา,}{โย สโต อุปสมฺมตี”ติ.}}
\csromangathasetleft{“เอตเทว อหํ มญฺเญ, \\ ปรํ สงฺกุปิตํ ญตฺวา,}
\csromangathagroup{\csromangathastanza{\csromangathabat{“เอตเทว อหํ มญฺเญ,}{พาลสฺส ปฏิเสธนํ.} \\ \csromangathabat{ปรํ สงฺกุปิตํ ญตฺวา,}{โย สโต อุปสมฺมตี”ติ.}}}

\prose{ภาสิตาย โข ปน ภิกฺขเว สกฺเกน เทวานมินฺเทน คาถาย เทวา อนุโมทิํสุ, อสุรา ตุณฺหี อเหสุํ. อถ โข ภิกฺขเว สกฺโก เทวานมินฺโท เวปจิตฺติํ อสุรินฺทํ เอตทโวจ “ภณ เวปจิตฺติ คาถนฺ”ติ. เอวํ วุตฺเต ภิกฺขเว เวปจิตฺติ อสุรินฺโท อิมํ คาถํ อภาสิ\csromandash{}}

\setlength{\csromangathapagewidth}{0pt}
\setlength{\csromangathawakwidth}{0pt}
\csromangathameasuregroup{“เอตเทว ติติกฺขาย, \\ ยทา นํ มญฺญติ\textsuperscript{1} พาโล, \\ อชฺฌารุหติ ทุมฺเมโธ,}{\csromangathabat{“เอตเทว ติติกฺขาย,}{วชฺชํ ปสฺสามิ วาสว.} \\ \csromangathabat{ยทา นํ มญฺญติ\textsuperscript{1} พาโล,}{ภยา มฺยายํ ติติกฺขติ.} \\ \csromangathabat{อชฺฌารุหติ ทุมฺเมโธ,}{โคว ภิยฺโย ปลายินนฺ”ติ.}}
\csromangathasetleft{“เอตเทว ติติกฺขาย, \\ ยทา นํ มญฺญติ\textsuperscript{1} พาโล, \\ อชฺฌารุหติ ทุมฺเมโธ,}
\csromangathagroup{\csromangathastanza{\csromangathabat{“เอตเทว ติติกฺขาย,}{วชฺชํ ปสฺสามิ วาสว.} \\ \csromangathabat{ยทา นํ มญฺญติ\csromansharedfootnote{3fe51462d28e2f75}{มญฺญตี (สี)} พาโล,}{ภยา มฺยายํ ติติกฺขติ.} \\ \csromangathabat{อชฺฌารุหติ ทุมฺเมโธ,}{โคว ภิยฺโย ปลายินนฺ”ติ.}}}

\prose{ภาสิตาย โข ปน ภิกฺขเว เวปจิตฺตินา อสุรินฺเทน คาถาย อสุรา อนุโมทิํสุ, เทวา ตุณฺหี อเหสุํ. อถ โข เวปจิตฺติ อสุรินฺโท สกฺกํ เทวานมินฺทํ เอตทโวจ “ภณ เทวานมินฺท คาถนฺ”ติ. เอวํ วุตฺเต ภิกฺขเว สกฺโก เทวานมินฺโต อิมา คาถาโย อภาสิ\csromandash{}}

\setlength{\csromangathapagewidth}{0pt}
\setlength{\csromangathawakwidth}{0pt}
\csromangathameasuregroup{“กามํ มญฺญตุ วา มา วา, \\ สทตฺถปรมา อตฺถา, \\ โย หเว พลวา สนฺโต, \\ ตมาหุ ปรมํ ขนฺติํ, \\ อพลํ ตํ พลํ อาหุ, \\ พลสฺส ธมฺมคุตฺตสฺส, \\ ตสฺเสว เตน ปาปิโย, \\ กุทฺธํ อปฺปฏิกุชฺฌนฺโต, \\ อุภินฺนมตฺถํ จรติ, \\ ปรํ สงฺกุปิตํ ญตฺวา, \\ อุภินฺนํ ติกิจฺฉนฺตานํ, \\ ชนา มญฺญนฺติ พาโลติ,}{\csromangathabat{“กามํ มญฺญตุ วา มา วา,}{ภยา มฺยายํ ติติกฺขติ.} \\ \csromangathabat{สทตฺถปรมา อตฺถา,}{ขนฺตฺยา ภิยฺโย น วิชฺชติ.} \\ \csromangathabat{โย หเว พลวา สนฺโต,}{ทุพฺพลสฺส ติติกฺขติ.} \\ \csromangathabat{ตมาหุ ปรมํ ขนฺติํ,}{นิจฺจํ ขมติ ทุพฺพโล.} \\ \csromangathabat{อพลํ ตํ พลํ อาหุ,}{ยสฺส พาลพลํ พลํ.} \\ \csromangathabat{พลสฺส ธมฺมคุตฺตสฺส,}{ปฏิวตฺตา น วิชฺชติ.} \\ \csromangathabat{ตสฺเสว เตน ปาปิโย,}{โย กุทฺธํ ปฏิกุชฺฌติ.} \\ \csromangathabat{กุทฺธํ อปฺปฏิกุชฺฌนฺโต,}{สงฺคามํ เชติ ทุชฺชยํ.} \\ \csromangathabat{อุภินฺนมตฺถํ จรติ,}{อตฺตโน จ ปรสฺส จ.} \\ \csromangathabat{ปรํ สงฺกุปิตํ ญตฺวา,}{โย สโต อุปสมฺมติ.} \\ \csromangathabat{อุภินฺนํ ติกิจฺฉนฺตานํ,}{อตฺตโน จ ปรสฺส จ.} \\ \csromangathabat{ชนา มญฺญนฺติ พาโลติ,}{เย ธมฺมสฺส อโกวิทา”ติ.}}
\csromangathasetleft{“กามํ มญฺญตุ วา มา วา, \\ สทตฺถปรมา อตฺถา, \\ โย หเว พลวา สนฺโต, \\ ตมาหุ ปรมํ ขนฺติํ, \\ อพลํ ตํ พลํ อาหุ, \\ พลสฺส ธมฺมคุตฺตสฺส, \\ ตสฺเสว เตน ปาปิโย, \\ กุทฺธํ อปฺปฏิกุชฺฌนฺโต, \\ อุภินฺนมตฺถํ จรติ, \\ ปรํ สงฺกุปิตํ ญตฺวา, \\ อุภินฺนํ ติกิจฺฉนฺตานํ, \\ ชนา มญฺญนฺติ พาโลติ,}
\csromangathagroup{\csromangathastanza{\csromangathabat{“กามํ มญฺญตุ วา มา วา,}{ภยา มฺยายํ ติติกฺขติ.} \\ \csromangathabat{สทตฺถปรมา อตฺถา,}{ขนฺตฺยา ภิยฺโย น วิชฺชติ.}}\csromangathastanza{\csromangathabat{โย หเว พลวา สนฺโต,}{ทุพฺพลสฺส ติติกฺขติ.} \\ \csromangathabat{ตมาหุ ปรมํ ขนฺติํ,}{นิจฺจํ ขมติ ทุพฺพโล.}}\csromangathastanza{\csromangathabat{อพลํ ตํ พลํ อาหุ,}{ยสฺส พาลพลํ พลํ.} \\ \csromangathabat{พลสฺส ธมฺมคุตฺตสฺส,}{ปฏิวตฺตา น วิชฺชติ.}}\csromangathastanza{\csromanfolio{150}\csromangathabat{ตสฺเสว เตน ปาปิโย,}{โย กุทฺธํ ปฏิกุชฺฌติ.} \\ \csromangathabat{กุทฺธํ อปฺปฏิกุชฺฌนฺโต,}{สงฺคามํ เชติ ทุชฺชยํ.}}\csromangathastanza{\csromangathabat{อุภินฺนมตฺถํ จรติ,}{อตฺตโน จ ปรสฺส จ.} \\ \csromangathabat{ปรํ สงฺกุปิตํ ญตฺวา,}{โย สโต อุปสมฺมติ.}}\csromangathastanza{\csromangathabat{อุภินฺนํ ติกิจฺฉนฺตานํ,}{อตฺตโน จ ปรสฺส จ.} \\ \csromangathabat{ชนา มญฺญนฺติ พาโลติ,}{เย ธมฺมสฺส อโกวิทา”ติ.}}}

\prose{ภาสิตาสุ โข ปน ภิกฺขเว สกฺเกน เทวานมินฺเทน คาถาสุ เทวา อนุโมทิํสุ, อสุรา ตุณฺหี อเหสุนฺติ. อิทํ ปรวจนํ. (9)}

\proseitem{117}{ตตฺถ กตมํ สกวจนญฺจ ปรวจนญฺจ.}

\prose{ยญฺจ ปตฺตํ ยญฺจ ปตฺตพฺพํ อุภยเมตํ รชานุกิณฺณํ อาตุรสฺสา\-นุสิกฺขโต. เย จ สิกฺขาสารา สีลํ วตํ ชีวิตํ พฺรหฺมจริยํ อุปฏฺฐานสารา, อยเมโก อนฺโต. เย จ เอวํวาทิโน เอวํทิฏฺฐิโน “นตฺถิ กาเมสุ โทโส”ติ, อยํ ทุติโย อนฺโต. อิจฺเจเต อุโภ อนฺตา กฏสิวฑฺฒนา กฏสิโย ทิฏฺฐิํ วฑฺเฒนฺติ. เอเต อุโภ อนฺเต อนภิญฺญาย โอลียนฺติ เอเก อติธาวนฺติ เอเกติ. อิทํ ปรวจนํ.}

\prose{เย จ โข เต อุโภ อนฺเต อภิญฺญาย ตตฺร จ น อเหสุํ, เต น จ อมญฺญิํสุ, วฏฺฏํ เตสํ นตฺถิ ปญฺญาปนายาติ. อิทํ สกวจนํ. อยํ อุทาโน สกวจนญฺจ ปรวจนญฺจ.}

\prose{\csromansymbolfootnote{*}{สํ 1. 71, 72 ปิฏฺเฐสุ.}ราชา ปเสนทิ\csromansharedfootnote{fd1ef24072062670}{ปสฺเสนทิ (ก) สํ 1. 71 ปิฏฺเฐ.} โกสโล ภควนฺตํ เอตทโวจ\csromandash{}อิธ มยฺหํ ภนฺเต รโหคตสฺส ปฏิสลฺลีนสฺส เอวํ เจตโส ปริวิตกฺโก อุทปาทิ “เกสํ นุ โข ปิโย อตฺตา, เกสํ อปฺปิโย อตฺตา”ติ. ตสฺส มยฺหํ ภนฺเต เอตทโหสิ “เย จ โข เกจิ กาเยน ทุจฺจริตํ จรนฺติ, วาจาย ทุจฺจริตํ จรนฺติ, มนสา ทุจฺจริตํ จรนฺติ, เตสํ อปฺปิโย อตฺตา. กิญฺจาปิ เต เอวํ วเทยฺยุํ ‘ปิโย โน อตฺตา’ติ, อถ โข เตสํ อปฺปิโย อตฺตา. ตํ กิสฺส เหตุ, ยํ หิ อปฺปิโย อปฺปิยสฺส กเรยฺย, ตํ เต อตฺตนาว อตฺตโน กโรนฺติ, ตสฺมา เตสํ อปฺปิโย อตฺตา. เย จ โข เกจิ กาเยน สุจริตํ จรนฺติ, วาจาย สุจริตํ จรนฺติ, มนสา สุจริตํ จรนฺติ, เตสํ ปิโย อตฺตา. กิญฺจาปิ เต เอวํ วเทยฺยุํ ‘อปฺปิโย โน อตฺตา’ติ, อถ โข เตสํ ปิโย อตฺตา. ตํ กิสฺส เหตุ, ยํ หิ ปิโย ปิยสฺส \csromanfolio{151}กเรยฺย. ตํ เต อตฺตนาว อตฺตโน กโรนฺติ. ตสฺมา เตสํ ปิโย อตฺตา”ติ.}

\prose{เอวเมตํ มหาราช, เอวเมตํ มหาราช, เย หิ เกจิ มหาราช กาเยน ทุจฺจริตํ จรนฺติ, วาจาย ทุจฺจริตํ จรนฺติ, มนสา ทุจฺจริตํ จรนฺติ ตสฺมา เตสํ อปฺปิโย อตฺตา. กิญฺจาปิ เต เอวํ วเทยฺยุํ “ปิโย โน อตฺตา”ติ, อถ โข เตสํ อปฺปิโย อตฺตา. ตํ กิสฺส เหตุ, ยํ หิ มหาราช อปฺปิโย อปฺปิยสฺส กเรยฺย, ตํ เต อตฺตนาว อตฺตโน กโรนฺติ, ตสฺมา เตสํ อปฺปิโย อตฺตา. เย จ โข เกจิ มหาราช กาเยน สุจริตํ จรนฺติ, วาจาย สุจริตํ จรนฺติ, มนสา สุจริตํ จรนฺติ, เตสํ ปิโย อตฺตา. กิญฺจาปิ เต เอวํ วเทยฺยุํ “อปฺปิโย โน อตฺตา”ติ, อถ โข เตสํ ปิโย อตฺตา. ตํ กิสฺส เหตุ, ยํ หิ มหาราช ปิโย ปิยสฺส กเรยฺย, ตํ เต อตฺตนาว อตฺตโน กโรนฺติ, ตสฺมา เตสํ ปิโย อตฺตาติ. อิทมโวจ ภควา ฯเปฯ สตฺถา\csromandash{}}

\setlength{\csromangathapagewidth}{0pt}
\setlength{\csromangathawakwidth}{0pt}
\csromangathameasuregroup{“อตฺตานญฺเจ ปิยํ ชญฺญา, \\ น หิ ตํ สุลภํ โหติ, \\ \textsuperscript{*}อนฺตเกนาธิปนฺนสฺส\textsuperscript{1}, \\ กิํ หิ ตสฺส สกํ โหติ, \\ กิญฺจสฺส อนุคํ โหติ, \\ อุโภ ปุญฺญญฺจ ปาปญฺจ, \\ ตญฺหิ ตสฺส สกํ โหติ, \\ ตํวสฺส อนุคํ โหติ, \\ ตสฺมา กเรยฺย กลฺยาณํ, \\ ปุญฺญานิ ปรโลกสฺมิํ,}{\csromangathabat{“อตฺตานญฺเจ ปิยํ ชญฺญา,}{น นํ ปาเปน สํยุเช.} \\ \csromangathabat{น หิ ตํ สุลภํ โหติ,}{สุขํ ทุกฺกฏการินา.} \\ \csromangathabat{\textsuperscript{*}อนฺตเกนาธิปนฺนสฺส\textsuperscript{1},}{ชหโต มานุสํ ภวํ.} \\ \csromangathabat{กิํ หิ ตสฺส สกํ โหติ,}{กิญฺจ อาทาย คจฺฉติ.} \\ \csromangathabat{กิญฺจสฺส อนุคํ โหติ,}{ฉายาว อนปายินี.} \\ \csromangathabat{อุโภ ปุญฺญญฺจ ปาปญฺจ,}{ยํ มจฺโจ กุรุเต อิธ.} \\ \csromangathabat{ตญฺหิ ตสฺส สกํ โหติ,}{ตํว อาทาย คจฺฉติ.} \\ \csromangathabat{ตํวสฺส อนุคํ โหติ,}{ฉายาว อนปายินี.} \\ \csromangathabat{ตสฺมา กเรยฺย กลฺยาณํ,}{นิจยํ สมฺปรายิกํ.} \\ \csromangathabat{ปุญฺญานิ ปรโลกสฺมิํ,}{ปติฏฺฐา โหนฺติ ปาณินนฺ”ติ.}}
\csromangathasetleft{“อตฺตานญฺเจ ปิยํ ชญฺญา, \\ น หิ ตํ สุลภํ โหติ, \\ \textsuperscript{*}อนฺตเกนาธิปนฺนสฺส\textsuperscript{1}, \\ กิํ หิ ตสฺส สกํ โหติ, \\ กิญฺจสฺส อนุคํ โหติ, \\ อุโภ ปุญฺญญฺจ ปาปญฺจ, \\ ตญฺหิ ตสฺส สกํ โหติ, \\ ตํวสฺส อนุคํ โหติ, \\ ตสฺมา กเรยฺย กลฺยาณํ, \\ ปุญฺญานิ ปรโลกสฺมิํ,}
\csromangathagroup{\csromangathastanza{\csromangathabat{“อตฺตานญฺเจ ปิยํ ชญฺญา,}{น นํ ปาเปน สํยุเช.} \\ \csromangathabat{น หิ ตํ สุลภํ โหติ,}{สุขํ ทุกฺกฏการินา.}}\csromangathastanza{\csromangathabat{\csromansymbolfootnote{*}{อุปริ 154 ปิฏฺเฐปิ.}อนฺตเกนา\-ธิปนฺนสฺส\csromansharedfootnote{f4b7ba5aaca940dd}{มรเณนาภิภูตสฺส (ก)},}{ชหโต มานุสํ ภวํ.} \\ \csromangathabat{กิํ หิ ตสฺส สกํ โหติ,}{กิญฺจ อาทาย คจฺฉติ.}}\csromangathastanza{\csromangathabat{กิญฺจสฺส อนุคํ โหติ,}{ฉายาว อนปายินี.} \\ \csromangathabat{อุโภ ปุญฺญญฺจ ปาปญฺจ,}{ยํ มจฺโจ กุรุเต อิธ.}}\csromangathastanza{\csromangathabat{ตญฺหิ ตสฺส สกํ โหติ,}{ตํว อาทาย คจฺฉติ.} \\ \csromangathabat{ตํวสฺส อนุคํ โหติ,}{ฉายาว อนปายินี.}}\csromangathastanza{\csromangathabat{ตสฺมา กเรยฺย กลฺยาณํ,}{นิจยํ สมฺปรายิกํ.} \\ \csromangathabat{ปุญฺญานิ ปรโลกสฺมิํ,}{ปติฏฺฐา โหนฺติ ปาณินนฺ”ติ.}}}

\prose{อิทํ สุตฺตํ ปรวจนํ. อนุคีติ สกวจนํ. อิทํ สกวจนญฺจ ปรวจนญฺจ. (10)}

\proseitem{118}{ตตฺถ กตมํ วิสชฺชนียํ.}

\prose{ปเญฺห ปจฺฉิเต อิทํ อภิญฺเญยฺยํ, อิทํ ปริญฺเญยฺยํ, อิทํ ปหาตพฺพํ, อิทํ ภาเวตพฺพํ, อิทํ สจฺฉิกาตพฺพํ, อิเม ธมฺมา เอวํคหิตา อิทํ ผลํ นิพฺพตฺตยนฺติ. \csromanfolio{152}เตสํ เอวํคหิตานํ อยมตฺโถ อิติ อิทํ วิสชฺชนียํ. “อุฬฺหาโร พุทฺโธ ภควา”ติ พุทฺธอุฬารตํ ธมฺม\-สฺวากฺขาตตํ สํฆสุปฺปฏิปตฺติญฺจ เอกํเสเนว นิทฺทิเส. “สพฺเพ สงฺขารา อนิจฺจา”ติ “สพฺเพ สงฺขารา ทุกฺขา”ติ “สพฺเพ ธมฺมา อนตฺตา”ติ เอกํเสเนว นิทฺทิเส. ยํ วา ปนญฺญมฺปิ เอวํ ชาติยํ. อิทํ วิสชฺชนียํ.}

\prose{ตตฺถ กตมํ อวิสชฺชนียํ.}

\setlength{\csromangathapagewidth}{0pt}
\setlength{\csromangathawakwidth}{0pt}
\csromangathameasuregroup{}{“อากงฺขโต เต นรทมฺมสารถิ\textsuperscript{1}, \\ เทวา มนุสฺสา มนสา วิจินฺติตํ. \\ สพฺเพ น ชญฺญา กสิณาปิ ปาณิโน, \\ สนฺตํ สมาธิํ อรณํ นิเสวโต. \\ กินฺตํ ภควา อากงฺขตี”ติ. \\ อิทํ อวิสชฺชนียํ.}
\csromangathameasurewak{“อากงฺขโต เต นรทมฺมสารถิ\textsuperscript{1}, \\ เทวา มนุสฺสา มนสา วิจินฺติตํ. \\ สพฺเพ น ชญฺญา กสิณาปิ ปาณิโน, \\ สนฺตํ สมาธิํ อรณํ นิเสวโต. \\ กินฺตํ ภควา อากงฺขตี”ติ. \\ อิทํ อวิสชฺชนียํ.}
\setlength{\csromangathaleftcol}{0pt}
\csromangathagroup{\csromangathastanza{\csromangathawak{“อากงฺขโต เต นรทมฺมสารถิ\csromansharedfootnote{8f9f50060a644ab5}{นรทมฺมสารถี (สี)}, \\ เทวา มนุสฺสา มนสา วิจินฺติตํ. \\ สพฺเพ น ชญฺญา กสิณาปิ ปาณิโน, \\ สนฺตํ สมาธิํ อรณํ นิเสวโต. \\ กินฺตํ ภควา อากงฺขตี”ติ. \\ อิทํ อวิสชฺชนียํ.}}}

\prose{เอตฺตโก ภควา สีลกฺขนฺเธ สมาธิ\-กฺขนฺเธ ปญฺญากฺขนฺเธ วิมุตฺติ\-กฺขนฺเธ วิมุตฺติ\-ญาณ\-ทสฺสนกฺขนฺเธ อิริยายํ ปภาเว หิเตสิตายํ กรุณายํ อิทฺธิยนฺติ. อิทํ อวิสชฺชนียํ.}

\prose{ตถาคตสฺส ภิกฺขเว อรหโต สมฺมา\-สมฺพุทฺธสฺส โลเก อุปฺปาทา ติณฺณํ รตนานํ อุปฺปาโท\csromansharedfootnote{d49ab0f03e28ca4a}{อุปฺปาทา (ก)} พุทฺธ\-รตนสฺส ธมฺม\-รตนสฺส สํฆรตนสฺส. กิํปมาณานิ ตีณิ รตนานีติ. อิทํ อวิสชฺชนียํ.}

\prose{พุทฺธวิสโย อวิสชฺชนีโย. ปุคฺคลปโรปรญฺญุตา อวิสชฺชนียา. ปุพฺพา ภิกฺขเว โกฏิ น ปญฺญายติ อวิชฺชา\-นีวรณานํ สตฺตานํ ตณฺหา\-สํโยชนานํ สกิํ นิรยํ สกิํ ติรจฺฉาน\-โยนิํ สกิํ เปตฺติวิสยํ สกิํ อสุรโยนิํ สกิํ เทเว สกิํ มนุสฺเส สนฺธาวิตํ สํสริตํ. กตมา ปุพฺพา โกฏีติ อวิสชฺชนียํ. น ปญฺญายตีติ สาวกานํ ญาณเวกลฺเลน. ทุวิธา พุทฺธานํ ภควนฺตานํ เทสนา อตฺตูปนายิกา จ ปรูปนายิกา จ. น ปญฺญายตีติ ปรูปนายิกา. นตฺถิ พุทฺธานํ ภควนฺตานํ อวิชานนาติ\csromansharedfootnote{435d36e8f6988194}{อปฺปชานนาติ (สี)} อตฺตูปนายิกา. ยถา ภควา โกกาลิกํ ภิกฺขุํ อารพฺภ อญฺญตรํ ภิกฺขุํ เอวมาห\csromandash{}}

\prose{\csromanfolio{153}เสยฺยถาปิ ภิกฺขุ วีสติขาริโก โกสลโก ติลวาโห ฯเปฯ. น เตฺวว เอโก อพฺพุโท นิรโย. เสยฺยถาปิ ภิกฺขุ วีสติ อพฺพุทา นิรยา, เอวเมโก นิรพฺพุโท นิรโย. เสยฺยถาปิ ภิกฺขุ วีสติ นิรพฺพุทา นิรยา, เอวเมโก อพโพ นิรโย. เสยฺยถาปิ ภิกฺขุ วีสติ อพพา นิรยา, เอวเมโก อฏโฏ นิรโย. เสยฺยถาปิ ภิกฺขุ วีสติ อฏฏา นิรยา, เอวเมโก อหโห นิรโย. เสยฺยถาปิ ภิกฺขุ วีสติ อหหา นิรยา, เอวเมโก กุมุโท นิรโย. เสยฺยถาปิ ภิกฺขุ วีสติ กุมุทา นิรยา, เอวเมโก โสคนฺธิโก นิรโย. เสยฺยถาปิ ภิกฺขุ วีสติ โสคนฺธิกา นิรยา, เอวเมโก อุปฺปลโก นิรโย. เสยฺยถาปิ ภิกฺขุ วีสติ อุปฺปลกา นิรยา, เอวเมโก ปุณฺฑรีโก นิรโย. เสยฺยถาปิ ภิกฺขุ วีสติ ปุณฺฑรีกา นิรยา, เอวเมโก ปทุโม นิรโย. ปทุเม ปน ภิกฺขุ นิรเย โกกาลิโก ภิกฺขุ อุปปนฺโน สาริปุตฺต\-โมคฺคลฺลาเนสุ จิตฺตํ อาฆาเตตฺวาติ\csromansharedfootnote{fa5db0b334d3bfa3}{สํ 1. 153, 154; อํ 3. 395, 396; ขุ 1. 381 ปิฏฺเฐสุ.}. ยํ วา ปน กิญฺจิ ภควา อาห “อยํ อปฺปเมยฺโย อสงฺเขฺยโย”ติ. สพฺพํ ตํ อวิสชฺชนียํ. อิทํ อวิสชฺชนียํ. (11)}

\proseitem{119}{ตตฺถ กตมํ วิสชฺชนียญฺจ อวิสชฺชนียญฺจ, ยทา โส อุปโก อาชีวโก ภควนฺตํ อาห “กุหิํ อาวุโส โคตม คมิสฺสสี”ติ. ภควา อาห\csromandash{}}

\setlength{\csromangathapagewidth}{0pt}
\setlength{\csromangathawakwidth}{0pt}
\csromangathameasuregroup{“พาราณสิํ คมิสฺสามิ, \\ ธมฺมจกฺกํ ปวตฺเตตุํ,}{\csromangathabat{“พาราณสิํ คมิสฺสามิ,}{อหํ ตํ อมตทุนฺทุภิํ.} \\ \csromangathabat{ธมฺมจกฺกํ ปวตฺเตตุํ,}{โลเก อปฺปฏิวตฺติยนฺ”ติ.}}
\csromangathasetleft{“พาราณสิํ คมิสฺสามิ, \\ ธมฺมจกฺกํ ปวตฺเตตุํ,}
\csromangathagroup{\csromangathastanza{\csromangathabat{“พาราณสิํ คมิสฺสามิ,}{อหํ ตํ อมตทุนฺทุภิํ.} \\ \csromangathabat{ธมฺมจกฺกํ ปวตฺเตตุํ,}{โลเก อปฺปฏิวตฺติยนฺ”ติ.}}}

\prose{อุปโก อาชีวโก อาห “‘ชิโน’ติ โข อาวุโส โภ โคตม ปฏิชานาสี”ติ. ภควา อาห\csromandash{}}

\setlength{\csromangathapagewidth}{0pt}
\setlength{\csromangathawakwidth}{0pt}
\csromangathameasuregroup{“มาทิสา เว ชินา\textsuperscript{1} โหนฺติ, \\ ชิตา เม ปาปกา ธมฺมา,}{\csromangathabat{“มาทิสา เว ชินา\textsuperscript{1} โหนฺติ,}{เย ปตฺตา อาสวกฺขยํ.} \\ \csromangathabat{ชิตา เม ปาปกา ธมฺมา,}{ตสฺมาหํ อุปกา ชิโน”ติ\textsuperscript{1}.}}
\csromangathasetleft{“มาทิสา เว ชินา\textsuperscript{1} โหนฺติ, \\ ชิตา เม ปาปกา ธมฺมา,}
\csromangathagroup{\csromangathastanza{\csromangathabat{“มาทิสา เว ชินา\csromansharedfootnote{7971a5d80fae218e}{ชินา เว มาทิสา (สี, ก)} โหนฺติ,}{เย ปตฺตา อาสวกฺขยํ.} \\ \csromangathabat{ชิตา เม ปาปกา ธมฺมา,}{ตสฺมาหํ อุปกา ชิโน”ติ\csromansharedfootnote{7971a5d80fae218e}{ชินา เว มาทิสา (สี, ก)}.}}}

\prose{กถํ ชิโน เกน ชิโนติ วิสชฺชนียํ. กตโม ชิโนติ อวิสชฺชนียํ. กตโม อาสวกฺขโย, ราคกฺขโย โทสกฺขโย โมหกฺขโยติ วิสชฺชนียํ. กิตฺตโก อาสวกฺขโยติ อวิสชฺชนียํ. อิทํ วิสชฺชนียญฺจ อวิสชฺชนียญฺจ.}

\prose{\csromanfolio{154}อตฺถิ ตถาคโตติ วิสชฺชนียํ. อตฺถิ รูปนฺติ วิสชฺชนียํ. รูปํ ตถาคโตติ อวิสชฺชนียํ. รูปวา ตถาคโตติ อวิสชฺชนียํ. รูเป ตถาคโตติ อวิสชฺชนียํ. ตถาคเต รูปนฺติ อวิสชฺชนียํ. เอวํ อตฺถิ เวทนา ฯเปฯ สญฺญา. สงฺขารา. อตฺถิ วิญฺญาณนฺติ วิสชฺชนียํ. วิญฺญาณํ ตถาคโตติ อวิสชฺชนียํ. วิญฺญาณวา ตถาคโตติ อวิสชฺชนียํ. วิญฺญาเณ ตถาคโตติ อวิสชฺชนียํ. ตถาคเต วิญฺญาณนฺติ อวิสชฺชนียํ. อญฺญตฺร รูเปน ตถาคโตติ อวิสชฺชนียํ. อญฺญตฺร เวทนาย ฯเปฯ สญฺญาย. สงฺขาเรหิ. วิญฺญาเณน ตถาคโตติ อวิสชฺชนียํ. อยํ โส ตถาคโต อรูปโก. อเวทนโก. อสญฺญโก. อสงฺขารโก. อวิญฺญาณโกติ อวิสชฺชนียํ. อิทํ วิสชฺชนียญฺจ อวิสชฺชนียญฺจ.}

\prose{ปสฺสติ ภควา ทิพฺเพน จกฺขุนา วิสุทฺเธน อติกฺกนฺต\-มานุสเกน สตฺเต จวมาเน อุปปชฺชมาเน เอวํ สพฺพํ ฯเปฯ ยถากมฺมูปเค สตฺเต ปชานาตีติ วิสชฺชนียํ. กตเม สตฺตา, กตโม ตถาคโตติ อวิสชฺชนียํ. อิทํ วิสชฺชนียญฺจ อวิสชฺชนียญฺจ.}

\prose{อตฺถิ ตถาคโตติ วิสชฺชนียํ. อตฺถิ ตถาคโต ปรํ มรณาติ อวิสชฺชนียํ. อิทํ วิสชฺชนียญฺจ อวิสชฺชนียญฺจ. (12)}

\proseitem{120}{ตตฺถ กตมํ กมฺมํ.}

\prose{\csromansymbolfootnote{*}{สํ 1. 72 ปิฏฺเฐ; เหฏฺฐา 151 ปิฏฺเฐปิ.}“อนฺตเกนา\-ธิปนฺนสฺส, ชหโต มานุสํ ภวํ.}

\setlength{\csromangathapagewidth}{0pt}
\setlength{\csromangathawakwidth}{0pt}
\csromangathameasuregroup{กิํ หิ ตสฺส สกํ โหติ, \\ กิญฺจสฺส อนุคํ โหติ, \\ อุโภ ปุญฺญญฺจ ปาปญฺจ, \\ ตญฺหิ ตสฺส สกํ โหติ, \\ ตํวสฺส อนุคํ โหติ,}{\csromangathabat{กิํ หิ ตสฺส สกํ โหติ,}{กิญฺจ อาทาย คจฺฉติ.} \\ \csromangathabat{กิญฺจสฺส อนุคํ โหติ,}{ฉายาว อนปายินี.} \\ \csromangathabat{อุโภ ปุญฺญญฺจ ปาปญฺจ,}{ยํ มจฺโจ กุรุเต อิธ.} \\ \csromangathabat{ตญฺหิ ตสฺส สกํ โหติ,}{ตํว\textsuperscript{1} อาทาย คจฺฉติ.} \\ \csromangathabat{ตํวสฺส อนุคํ โหติ,}{ฉายาว อนปายินี”ติ.} \\ อิทํ กมฺมํ.}
\csromangathasetleft{กิํ หิ ตสฺส สกํ โหติ, \\ กิญฺจสฺส อนุคํ โหติ, \\ อุโภ ปุญฺญญฺจ ปาปญฺจ, \\ ตญฺหิ ตสฺส สกํ โหติ, \\ ตํวสฺส อนุคํ โหติ,}
\csromangathagroup{\csromangathastanza{\csromangathabat{กิํ หิ ตสฺส สกํ โหติ,}{กิญฺจ อาทาย คจฺฉติ.} \\ \csromangathabat{กิญฺจสฺส อนุคํ โหติ,}{ฉายาว อนปายินี.}}\csromangathastanza{\csromangathabat{อุโภ ปุญฺญญฺจ ปาปญฺจ,}{ยํ มจฺโจ กุรุเต อิธ.} \\ \csromangathabat{ตญฺหิ ตสฺส สกํ โหติ,}{ตํว\csromansharedfootnote{338dd9d64fd7ee1e}{[มิสฺสินฺคฺ โนเต 0]} อาทาย คจฺฉติ.}}\csromangathastanza{\csromangathabat{ตํวสฺส อนุคํ โหติ,}{ฉายาว อนปายินี”ติ.} \\ อิทํ กมฺมํ.}}

\prose{\csromansymbolfootnote{+}{ม 3. 203 ปิฏฺเฐ.}“ปุน จปรํ ภิกฺขเว พาลํ ปีฐสมารูฬฺหํ วา มญฺจ\-สมารูฬฺหํ วา ฉมายํ\csromansharedfootnote{7e443653b22d1558}{ฉมาย (สี, ก)} วา เสมานํ ยานิสฺส ปุพฺเพ ปาปกานิ กมฺมานิ กตานิ กาเยน ทุจฺจริตานิ วาจาย ทุจฺจริตานิ มนสา ทุจฺจริตานิ, ตานิสฺส ตมฺหิ สมเย โอลมฺพนฺติ \csromanfolio{155}อชฺโฌลมฺพนฺติ อภิปฺป\-ลมฺพนฺติ. เสยฺยถาปิ ภิกฺขเว มหตํ ปพฺพตกูฏานํ ฉายา สายนฺหสมยํ ปถวิยํ โอลมฺพนฺติ อชฺโฌลมฺพนฺติ อภิปฺป\-ลมฺพนฺติ. เอวเมว โข ภิกฺขเว พาลํ ปีฐสมารูฬฺหํ วา มญฺจ\-สมารูฬฺหํ วา ฉมายํ วา เสมานํ ยานิสฺส ปุพฺเพ ปาปกานิ กมฺมานิ กตานิ กาเยน ทุจฺจริตานิ วาจาย ทุจฺจริตานิ มนสา ทุจฺจริตานิ, ตานิสฺส ตมฺหิ สมเย โอลมฺพนฺติ อชฺโฌลมฺพนฺติ อภิปฺป\-ลมฺพนฺติ. ตตฺร ภิกฺขเว พาลสฺส เอวํ โหติ “อกตํ วต เม กลฺยาณํ, อกตํ กุสลํ, อกตํ ภีรุตฺตาณํ. กตํ ปาปํ, กตํ ลุทฺทํ, กตํ กิพฺพิสํ, ยาวตา โภ อกต\-กลฺยาณานํ อกตกุสลานํ อกต\-ภีรุตฺตาณานํ กตปาปานํ กตลุทฺทานํ กตกิพฺพิสานํ คติ, ตํ คติํ เปจฺจ คจฺฉามี”ติ, โส โสจติ กิลมติ ปริเทวติ อุรตฺตาฬิํ กนฺทติ สมฺโมหํ อาปชฺชตี”ติ.}

\prose{\csromansymbolfootnote{*}{ม 3. 210 ปิฏฺเฐ.}ปุน จปรํ ภิกฺขเว ปณฺฑิตํ ปีฐสมารูฬฺหํ วา มญฺจ\-สมารูฬฺหํ วา ฉมายํ วา เสมานํ ยานิสฺส ปุพฺเพ กลฺยาณานิ กมฺมานิ กตานิ กาเยน สุจริตานิ วาจาย สุจริตานิ มนสา สุจริตานิ, ตานิสฺส ตมฺหิ สมเย โอลมฺพนฺติ อชฺโฌลมฺพนฺติ อภิปฺป\-ลมฺพนฺติ. เสยฺยถาปิ ภิกฺขเว มหตํ ปพฺพตกูฏานํ ฉายา สายนฺหสมยํ ปถวิยํ โอลมฺพนฺติ อชฺโฌลมฺพนฺติ อภิปฺป\-ลมฺพนฺติ. เอวเมว โข ภิกฺขเว ปณฺฑิตํ ปีฐสมารูฬฺหํ วา มญฺจ\-สมารูฬฺหํ วา ฉมายํ วา เสมานํ ยานิสฺส ปุพฺเพ อลฺยาณานิ กมฺมานิ กตานิ กาเยน สุจริตานิ วาจาย สุจริตานิ มนสา สุจริตานิ, ตานิสฺส ตมฺหิ สมเย โอลมฺพนฺติ อชฺโฌลมฺพนฺติ อภิปฺป\-ลมฺพนฺติ. ตตฺร ภิกฺขเว ปณฺฑิตสฺส “เอวํ โหติ อกตํ วต เม ปาปํ, อกตํ ลุทฺทํ, อกตํ กิพฺพิสํ. กตํ กลฺยาณํ, กตํ กุสลํ, กตํ ภีรุตฺตาณํ, ยาวตา โภ อกตปาปานํ อกตลุทฺทานํ อกตกิพฺพิสานํ กตกลฺยาณานํ กตกุสลานํ กตภีรุตฺตาณานํ คติ, ตํ คติํ เปจฺจ คจฺฉามี”ติ, โส น โสจติ น กิลมติ น ปริเทวติ น อุรตฺตาฬิํ กนฺทติ น สมฺโมหํ อาปชฺชติ, “กตํ เม ปุญฺญํ, อกตํ ปาปํ, ยา ภวิสฺสติ คติ อกตปาปสฺส อกตลุทฺทสฺส อกตกิพฺพิสสฺส กตปุญฺญสฺส กตกุสลสฺส กตภีรุตฺตาณสฺส, ตํ เปจฺจ ภเว คติํ ปจฺจนุภวิสฺสามี”ติ วิปฺปฏิสาโร น ชายติ. อวิปฺปฏิสาริโน โข ภิกฺขเว อิตฺถิยา วา ปุริสสฺส วา คิหิโน วา ปพฺพชิตสฺส วา ภทฺทกํ มรณํ ภทฺทิกา กาลํกิริยาติ วทามี”ติ. อิทํ กมฺมํ.}

\prose{\csromanfolio{156}ตีณิมานิ ภิกฺขเว ทุจฺจริตานิ. กตมานิ ตีณิ, กายทุจฺจริตํ วจีทุจฺจริตํ มโนทุจฺจริตํ, อิมานิ โข ภิกฺขเว ตีณิ ทุจฺจริตานิ. ตีณิมานิ ภิกฺขเว สุจริตานิ. กตมานิ ตีณิ, กายสุจริตํ วจีสุจริตํ มโนสุจริตํ. อิมานิ โข ภิกฺขเว ตีณิ สุจริตานิ. อิทํ กมฺมํ.}

\prose{ตตฺถ กตโม วิปาโก.}

\prose{\csromansymbolfootnote{*}{สํ 2. 341 ปิฏฺเฐ.}ลาภา โว ภิกฺขเว, สุลทฺธํ โว ภิกฺขเว, ขโณ โว ภิกฺขเว ปฏิลทฺโธ พฺรหฺมจริย\-วาสาย. ทิฏฺฐา มยา ภิกฺขเว ฉผสฺสา\-ยตนิกา นาม นิรยา. ตตฺถ ยํ กิญฺจิ จกฺขุนา รูปํ ปสฺสติ อนิฏฺฐ\-รูปํเยว ปสฺสติ, โน อิฏฺฐรูปํ. อกนฺต\-รูปํเยว ปสฺสติ, โน กนฺตรูปํ. อมนาป\-รูปํเยว ปสฺสติ, โน มนาปรูปํ.}

\prose{ยํ กิญฺจิ โสเตน สทฺทํ สุณาติ ฯเปฯ ฆาเนน. ชิวฺหาย. กาเยน. ยํ กิญฺจิ มนสา ธมฺมํ วิชานาติ อนิฏฺฐธมฺมํเยว\csromansharedfootnote{0040da3766a603ba}{อนิฏฺฐรูปํเยว (สํ 2. 341 ปิฏฺเฐ, ตถา ปเรสุ โน อิฏฺฐธมฺมนฺติอาทิปเทสุปิ.)} วิชานาติ, โน อิฏฺฐธมฺมํ. อกนฺตธมฺมํเยว วิชานาติ, โน กนฺตธมฺมํ. อมนาปธมฺมํเยว วิชานาติ, โน มนาปธมฺมํ. ลาภา โว ภิกฺขเว, สุลทฺธํ โว ภิกฺขเว, ขโณ โว ภิกฺขเว ปฏิลทฺโธ พฺรหฺมจริย\-วาสาย.}

\prose{ทิฏฺฐา มยา ภิกฺขเว ฉผสฺสา\-ยตนิกา นาม สคฺคา. ตตฺถ ยํกิญฺจิ จกฺขุนา รูปํ ปสฺสติ อิฏฺฐรูปํเยว ปสฺสติ, โน อนิฏฺฐรูปํ. กนฺตรูปํเยว ปสฺสติ, โน อกนฺตรูปํ. มนาป\-รูปํเยว ปสฺสติ, โน อมนาปรูปํ. ยํ กิญฺจิ โสเตน สทฺทํ สุณาติ ฯเปฯ ฆาเนน. ชิวฺหาย. กาเยน. มนสา ธมฺมํ วิชานาติ อิฏฺฐธมฺมํเยว วิชานาติ, โน อนิฏฺฐธมฺมํ. กนฺตธมฺมํเยว วิชานาติ, โน อกนฺตธมฺมํ. มนาปธมฺมํเยว วิชานาติ, โน อมนาปธมฺมํ. ลาภา โว ภิกฺขเว, สุลทฺธํ โว ภิกฺขเว, ขโณ โว ภิกฺขเว ปฏิลทฺโธ พฺรหฺมจริยวาสายา”ติ. อยํ วิปาโก.}

\setlength{\csromangathapagewidth}{0pt}
\setlength{\csromangathawakwidth}{0pt}
\csromangathameasuregroup{สฏฺฐิวสฺสสหสฺสานิ, \\ นิรเย ปจฺจมานานํ\textsuperscript{1}, \\ นตฺถิ อนฺโต กุโต อนฺโต,}{\csromangathabat{สฏฺฐิวสฺสสหสฺสานิ,}{ปริปุณฺณานิ สพฺพโส.} \\ \csromangathabat{นิรเย ปจฺจมานานํ\textsuperscript{1},}{กทา อนฺโต ภวิสฺสติ.} \\ \csromangathabat{นตฺถิ อนฺโต กุโต อนฺโต,}{น อนฺโต ปฏิทิสฺสติ\textsuperscript{1}.}}
\csromangathasetleft{สฏฺฐิวสฺสสหสฺสานิ, \\ นิรเย ปจฺจมานานํ\textsuperscript{1}, \\ นตฺถิ อนฺโต กุโต อนฺโต,}
\csromangathagroup{\csromangathastanza{\csromangathabat{สฏฺฐิ\-วสฺสสหสฺสานิ,}{ปริปุณฺณานิ สพฺพโส.} \\ \csromangathabat{นิรเย ปจฺจมานานํ\csromansharedfootnote{22153d9c7b2bbb65}{ปจฺจมานสฺส (ก)},}{กทา อนฺโต ภวิสฺสติ.} \\ \csromangathabat{นตฺถิ อนฺโต กุโต อนฺโต,}{น อนฺโต ปฏิทิสฺสติ\csromansharedfootnote{22153d9c7b2bbb65}{ปจฺจมานสฺส (ก)}.}}}

\vspace{\tipitakaparskip}
\csromancenter{ตทา หิ ปกตํ ปาปํ, ตุยฺหํ มยฺหญฺจ มาริสาติ\csromansharedfootnote{48a7cb49aa2ed023}{ขุ 2. 216; ขุ 5. 100 ปิฏฺเฐสุ; อุปริ 201, 208 ปิฏฺเฐสุปิ.}. อยํ วิปาโก. (13)}
\proseitem{121}{\csromanfolio{157}ตตฺถ กตมํ กมฺมญฺจ วิปาโก จ.}

\setlength{\csromangathapagewidth}{0pt}
\setlength{\csromangathawakwidth}{0pt}
\csromangathameasuregroup{น หิ ธมฺโม อธมฺโม จ, \\ อธมฺโม นิรยํ เนติ,}{อธมฺมจารี หิ นโร ปมตฺโต, \\ ยหิํ ยหิํ คจฺฉติ ทุคฺคติํ โย. \\ โส นํ อธมฺโม จริโต หนาติ, \\ สยํ คหีโต ยถา กณฺหสปฺโป. \\ \csromangathabat{น หิ ธมฺโม อธมฺโม จ,}{อุโภ สมวิปากิโน.} \\ \csromangathabat{อธมฺโม นิรยํ เนติ,}{ธมฺโม ปาเปติ สุคฺคตินฺติ\textsuperscript{1}.} \\ อิทํ กมฺมญฺจ วิปาโก จ.}
\csromangathameasurewak{อธมฺมจารี หิ นโร ปมตฺโต, \\ ยหิํ ยหิํ คจฺฉติ ทุคฺคติํ โย. \\ โส นํ อธมฺโม จริโต หนาติ, \\ สยํ คหีโต ยถา กณฺหสปฺโป. \\ อิทํ กมฺมญฺจ วิปาโก จ.}
\csromangathasetleft{น หิ ธมฺโม อธมฺโม จ, \\ อธมฺโม นิรยํ เนติ,}
\csromangathagroup{\csromangathastanza{\csromangathawak{อธมฺมจารี หิ นโร ปมตฺโต, \\ ยหิํ ยหิํ คจฺฉติ ทุคฺคติํ โย. \\ โส นํ อธมฺโม จริโต หนาติ, \\ สยํ คหีโต ยถา กณฺหสปฺโป.}}\csromangathastanza{\csromangathabat{น หิ ธมฺโม อธมฺโม จ,}{อุโภ สมวิปากิโน.} \\ \csromangathabat{อธมฺโม นิรยํ เนติ,}{ธมฺโม ปาเปติ สุคฺคตินฺติ\csromansharedfootnote{338dd9d64fd7ee1e}{[มิสฺสินฺคฺ โนเต 0]}.}}\csromangathastanza{\csromangathawak{อิทํ กมฺมญฺจ วิปาโก จ.}}}

\prose{\csromansymbolfootnote{*}{ขุ 1. 205 ปิฏฺเฐ; อํ 1. 465 ปิฏฺเฐปิ ปสฺสิตพฺพํ.}“มา ภิกฺขเว ปุญฺญานํ ภายิตฺถ, สุขสฺเสตํ ภิกฺขเว อธิวจนํ อิฏฺฐสฺส กนฺตสฺส ปิยสฺส มนาปสฺส ยทิทํ ปุญฺญานิ. อภิชานามิ โข ปนาหํ ภิกฺขเว ทีฆรตฺตํ กตานํ ปุญฺญานํ อิฏฺฐํ\csromansharedfootnote{0d5b036ef1db5ef7}{ทีฆรตฺตํ อิฏฺฐํ (สี, ก) 3. น ยิมํ (ขุ 1. 205 ปิฏฺเฐ.)} กนฺตํ ปิยํ มนาปํ วิปากํ ปจฺจนุภูตํ, สตฺต วสฺสานิ เมตฺตจิตฺตํ ภาเวตฺวา สตฺต สํวฏฺฏ\-วิวฏฺฏ\-กปฺเป น อิมํ\csromansharedfootnote{1a9c92200200e8fb}{สุทํ (ขุ 1. 205 ปิฏฺเฐ.)} โลกํ ปุนราคมาสิํ. สํวฏฺฏมาเน สุทาหํ\csromansharedfootnote{e1e7561d4f6d11eb}{จกฺกวตฺติ (ก)} ภิกฺขเว กปฺเป อาภสฺสรูปโค โหมิ. วิวฏฺฏมาเน กปฺเป สุญฺญํ พฺรหฺมวิมานํ อุปปชฺชามิ. ตตฺร สุทาหํ\csromansharedfootnote{f0a418e340053646}{ม 3. 243 ปิฏฺฐาทีสุ.} ภิกฺขเว พฺรหฺมา โหมิ มหาพฺรหฺมา อภิภู อนภิภูโต อญฺญทตฺถุทโส วสวตฺตี. ฉตฺติํส\-กฺขตฺตุํ โข ปนาหํ ภิกฺขเว สกฺโก อโหสิํ เทวานมินฺโท, อเนก\-สต\-กฺขตฺตุํ ราชา อโหสิํ จกฺกวตฺตี\csromansharedfootnote{5e2b06c21b65d0fa}{น ยิมํ (ขุ 1. 205 ปิฏฺเฐ.)} ธมฺมิโก ธมฺมราชา จาตุรนฺโต วิชิตาวี ชนปท\-ตฺถาวริย\-ปฺปตฺโต สตฺต\-รตน\-สมนฺนาคโต, โก ปน วาโท ปเทสรชฺชสฺส. ตสฺส มยฺหํ ภิกฺขเว เอตทโหสิ “กิสฺส นุ โข เม อิทํ กมฺมสฺส ผลํ, กิสฺส กมฺมสฺส วิปาโก, เยนาหํ เอตรหิ เอวํมหิทฺธิโก เอวํ\-มหา\-นุภาโว”ติ. ตสฺส มยฺหํ ภิกฺขเว เอตทโหสิ “ติณฺณํ โข เม อิทํ กมฺมานํ ผลํ, ติณฺณํ กมฺมานํ วิปาโก. เยนาหํ เอตรหิ เอวํมหิทฺธิโก เอวํ\-มหา\-นุภาโว”ติ. เสยฺยถิทํ, ทานสฺส ทมสฺส สํยมสฺสา”ติ, ตตฺถ ยญฺจ ทานํ โย จ ทโม โย จ สํยโม, อิทํ กมฺมํ. โย ตปฺปจฺจยา วิปาโก ปจฺจนุภูโต, อยํ วิปาโก. ตถา จูฬกมฺมวิภงฺโค\csromansharedfootnote{1a9c92200200e8fb}{สุทํ (ขุ 1. 205 ปิฏฺเฐ.)} วตฺตพฺโพ.}

\prose{\csromanfolio{158}ยํ สุภสฺส มาณวสฺส โตเทยฺย\-ปุตฺตสฺส เทสิตํ. ตตฺถ เย ธมฺมา อปฺปายุก\-ทีฆายุกตาย สํวตฺตนฺติ พหฺวาพาธ\-อปฺปาพาธตาย อปฺเปสกฺข\-มเหสกฺข\-ตาย ทุพฺพณฺณ\-สุวณฺณตาย นีจกุลิก- อุจฺจกุลิกตาย อปฺปโภค\-มหาโภคตาย ทุปฺปญฺญ\-ปญฺญวนฺต\-ตาย จ สํวตฺตนฺติ, อิทํ กมฺมํ. ยา ตตฺถ อปฺปายุก\-ทีฆายุกตา ฯเปฯ ทุปฺปญฺญ\-ปญฺญวนฺตตา, อยํ วิปาโก. อิทํ กมฺมญฺจ วิปาโก จ. (14)}

\proseitem{122}{ตตฺถ กตมํ กุสลํ.}

\prose{\csromansymbolfootnote{*}{ขุ 1. 53 ปิฏฺเฐ ธมฺมปเท; อุปริ 216 ปิฏฺเฐปิ.}“วาจานุรกฺขี มนสา สุสํวุโต,}

\prose{กาเยน จ นากุสลํ กยิรา\csromansharedfootnote{9a5f6de8edfd7aed}{อกุสลํ น กยิรา (สี) ** ขุ 1. 69 ปิฏฺเฐ ธมฺมปเท.}.}

\setlength{\csromangathapagewidth}{0pt}
\setlength{\csromangathawakwidth}{0pt}
\csromangathameasuregroup{}{เอเต ตโย กมฺมปเถ วิโสธเย, \\ อาราธเย มคฺค’มิสิปฺปเวทิตนฺ”ติ. \\ อิทํ กุสลํ.}
\csromangathameasurewak{เอเต ตโย กมฺมปเถ วิโสธเย, \\ อาราธเย มคฺค’มิสิปฺปเวทิตนฺ”ติ. \\ อิทํ กุสลํ.}
\setlength{\csromangathaleftcol}{0pt}
\csromangathagroup{\csromangathastanza{\csromangathawak{เอเต ตโย กมฺมปเถ วิโสธเย, \\ อาราธเย มคฺค’มิสิปฺปเวทิตนฺ”ติ. \\ อิทํ กุสลํ.}}}

\prose{** “ยสฺส กาเยน วาจาย, มนสา นตฺถิ ทุกฺกฏํ.}

\setlength{\csromangathapagewidth}{0pt}
\setlength{\csromangathawakwidth}{0pt}
\csromangathameasuregroup{สํวุตํ ตีหิ ฐาเนหิ,}{\csromangathabat{สํวุตํ ตีหิ ฐาเนหิ,}{ตมหํ พฺรูมิ พฺราหฺมณนฺ”ติ.} \\ อิทํ กุสลํ.}
\csromangathasetleft{สํวุตํ ตีหิ ฐาเนหิ,}
\csromangathagroup{\csromangathastanza{\csromangathabat{สํวุตํ ตีหิ ฐาเนหิ,}{ตมหํ พฺรูมิ พฺราหฺมณนฺ”ติ.} \\ อิทํ กุสลํ.}}

\prose{\csromansymbolfootnote{+}{อํ 1. 203 ปิฏฺเฐ. ++ สํ 3. 2 ปิฏฺเฐ.}ตีณิมานิ ภิกฺขเว กุสลมูลานิ. กตมานิ ตีณิ, อโลโภ กุสลมูลํ, อโทโส กุสลมูลํ, อโมโห กุสลมูลํ. อิมานิ โข ภิกฺขเว ตีณิ กุสลมูลานิ. อิทํ กุสลํ. ++ วิชฺชา ภิกฺขเว\csromansharedfootnote{7440759bfb676ff8}{วิชฺชา จ โข ภิกฺขเว (สํ 3. 2 ปิฏฺเฐ.)} ปุพฺพงฺคมา กุสลานํ ธมฺมานํ สมาปตฺติยา อนุเทว\csromansharedfootnote{33b9d9391bbbbfc5}{อนฺวเทว (สี, ก, สํ 3. 2 ปิฏฺเฐ จ), สฺยาทิกณฺเฑ (โมคฺคลฺลาเน) 11 สุตฺตํ ปสฺสิตพฺพํ.} หิรี\csromansharedfootnote{5f44079e15dd5122}{หิริํ (?) ++ ขุ 1. 37 ปิฏฺเฐ ธมฺมปเท.} โอตฺตปฺปญฺจาติ. อิทํ กุสลํ.}

\prose{ตตฺถ กตมํ อกุสลํ.}

\prose{++ “ยสฺส อจฺจนฺต\-ทุสฺสีลฺยํ, มาลุวา สาลมิ\-โว\-ตฺถตํ.}

\setlength{\csromangathapagewidth}{0pt}
\setlength{\csromangathawakwidth}{0pt}
\csromangathameasuregroup{กโรติ โส ตถ’ตฺตานํ,}{\csromangathabat{กโรติ โส ตถ’ตฺตานํ,}{ยถา นํ อิจฺฉตี ทิโส”ติ.} \\ อิทํ อกุสลํ.}
\csromangathasetleft{กโรติ โส ตถ’ตฺตานํ,}
\csromangathagroup{\csromangathastanza{\csromangathabat{กโรติ โส ตถ’ตฺตานํ,}{ยถา นํ อิจฺฉตี ทิโส”ติ.} \\ อิทํ อกุสลํ.}}

\prose{++ “อตฺตนา หิ กตํ ปาปํ, อตฺตชํ อตฺตสมฺภวํ.}

\setlength{\csromangathapagewidth}{0pt}
\setlength{\csromangathawakwidth}{0pt}
\csromangathameasuregroup{อภิมตฺถติ\textsuperscript{1} ทุมฺเมธํ,}{\csromangathabat{อภิมตฺถติ\textsuperscript{1} ทุมฺเมธํ,}{วชิรํว’สฺมมยํ มณินฺ”ติ.} \\ อิทํ อกุสลํ. \\ “ทส กมฺมปเถ นิเสวิย, \\ อกุสลากุสเลหิ วิวชฺชิตา. \\ ครหา จ ภวนฺติ เทวเต, \\ พาลมตี นิรเยสุ ปจฺจเร”ติ. \\ อิทํ อกุสลํ.}
\csromangathameasurewak{อกุสลากุสเลหิ วิวชฺชิตา. \\ ครหา จ ภวนฺติ เทวเต, \\ พาลมตี นิรเยสุ ปจฺจเร”ติ. \\ อิทํ อกุสลํ.}
\csromangathasetleft{อภิมตฺถติ\textsuperscript{1} ทุมฺเมธํ,}
\csromangathagroup{\csromangathastanza{\csromangathabat{อภิมตฺถติ\csromansharedfootnote{338dd9d64fd7ee1e}{[มิสฺสินฺคฺ โนเต 0]} ทุมฺเมธํ,}{วชิรํว’สฺมมยํ มณินฺ”ติ.} \\ อิทํ อกุสลํ. \\ “ทส กมฺมปเถ นิเสวิย,}\csromangathastanza{\csromangathawak{อกุสลากุสเลหิ วิวชฺชิตา. \\ ครหา จ ภวนฺติ เทวเต, \\ พาลมตี นิรเยสุ ปจฺจเร”ติ. \\ อิทํ อกุสลํ.}}}

\prose{\csromanfolio{159}\csromansymbolfootnote{*}{อํ 1. 202 ปิฏฺเฐ.}ตีณิมานิ ภิกฺขเว อกุสลมูลานิ, กตมานิ ตีณิ, โลโภ อกุสลมูลํ, โทโส อกุสลมูลํ, โมโห อกุสลมูลํ. อิมานิ โข ภิกฺขเว ตีณิ อกุสลมูลานิ. อิทํ อกุสลํ. (15)}

\setlength{\csromangathapagewidth}{0pt}
\setlength{\csromangathawakwidth}{0pt}
\csromangathameasuregroup{}{ตตฺถ กตมํ กุสลญฺจ อกุสลญฺจ. \\ “ยาทิสํ วปเต พีชํ, ตาทิสํ หรเต ผลํ. \\ กลฺยาณการี กลฺยาณํ, ปาปการี จ ปาปกนฺ”ติ\textsuperscript{1}.}
\csromangathameasurewak{ตตฺถ กตมํ กุสลญฺจ อกุสลญฺจ. \\ “ยาทิสํ วปเต พีชํ, ตาทิสํ หรเต ผลํ. \\ กลฺยาณการี กลฺยาณํ, ปาปการี จ ปาปกนฺ”ติ\textsuperscript{1}.}
\setlength{\csromangathaleftcol}{0pt}
\csromangathagroup{\csromangathastanza{\csromangathawak{\csromangathaitemline{122}{ตตฺถ กตมํ กุสลญฺจ อกุสลญฺจ.} \\ \csromangathacontline{“ยาทิสํ วปเต พีชํ, ตาทิสํ หรเต ผลํ.} \\ \csromangathacontline{กลฺยาณการี กลฺยาณํ, ปาปการี จ ปาปกนฺ”ติ\csromansharedfootnote{06bdb60366bd7e4f}{สํ 1. 229 ปิฏฺเฐ.}.}}}}

\prose{ตตฺถ ยํ อาห “กลฺยาณการี กลฺยาณนฺ”ติ, อิทํ กุสลํ.  ยํ อาห “ปาปการี จ ปาปกนฺ”ติ, อิทํ อกุสลํ. อิทํ กุสลญฺจ อกุสลญฺจ.}

\setlength{\csromangathapagewidth}{0pt}
\setlength{\csromangathawakwidth}{0pt}
\csromangathameasuregroup{}{สุเภน กมฺเมน วชนฺติ สุคฺคติํ, \\ อปายภูมิํ อสุเภน กมฺมุนา. \\ ขยา จ กมฺมสฺส วิมุตฺตเจตโส, \\ นิพฺพนฺติ เต โชติริวินฺธนกฺขยา.}
\csromangathameasurewak{สุเภน กมฺเมน วชนฺติ สุคฺคติํ, \\ อปายภูมิํ อสุเภน กมฺมุนา. \\ ขยา จ กมฺมสฺส วิมุตฺตเจตโส, \\ นิพฺพนฺติ เต โชติริวินฺธนกฺขยา.}
\setlength{\csromangathaleftcol}{0pt}
\csromangathagroup{\csromangathastanza{\csromangathawak{สุเภน กมฺเมน วชนฺติ สุคฺคติํ, \\ อปายภูมิํ อสุเภน กมฺมุนา. \\ ขยา จ กมฺมสฺส วิมุตฺตเจตโส, \\ นิพฺพนฺติ เต โชติริวินฺธนกฺขยา.}}}

\prose{ตตฺถ ยํ อาห “สุเภน กมฺเมน วชนฺติ สุคฺคตินฺ”ติ, อิทํ กุสลํ.  ยํ อาห “อปายภูมิํ อสุเภน กมฺมุนา”ติ, อิทํ อกุสลํ. อิทํ กุสลญฺจ อกุสลญฺจ. (16)}

\setlength{\csromangathapagewidth}{0pt}
\setlength{\csromangathawakwidth}{0pt}
\csromangathameasuregroup{}{ตตฺถ กตมํ อนุญฺญาตํ. \\ “ยถาปิ ภมโร ปุปฺผํ, วณฺณคนฺธมเหฐยํ. \\ ปเลติ\textsuperscript{1} รสมาทาย, เอวํ คาเม มุนี จเร”ติ\textsuperscript{1}.}
\csromangathameasurewak{ตตฺถ กตมํ อนุญฺญาตํ. \\ “ยถาปิ ภมโร ปุปฺผํ, วณฺณคนฺธมเหฐยํ. \\ ปเลติ\textsuperscript{1} รสมาทาย, เอวํ คาเม มุนี จเร”ติ\textsuperscript{1}.}
\setlength{\csromangathaleftcol}{0pt}
\csromangathagroup{\csromangathastanza{\csromangathawak{\csromangathaitemline{123}{ตตฺถ กตมํ อนุญฺญาตํ.} \\ \csromangathacontline{“ยถาปิ ภมโร ปุปฺผํ, วณฺณคนฺธม\-เหฐยํ.} \\ \csromangathacontline{ปเลติ\csromansharedfootnote{6f82779f9476916a}{ปเฬติ (ก)} รสมาทาย, เอวํ คาเม มุนี จเร”ติ\csromansharedfootnote{b782edd3d6bafc16}{ขุ 1. 20 ปิฏฺเฐ ธมฺมปเท.}.}}}}

\prose{อิทํ อนุญฺญาตํ.}

\prose{\csromanfolio{160}ตีณิมานิ ภิกฺขเว ภิกฺขูนํ กรณียานิ. กตมานิ ตีณิ, อิธ ภิกฺขเว ภิกฺขุ ปาติโมกฺข\-สํวร\-สํวุโต วิหรติ อาจารโคจร\-สมฺปนฺโน อณุมตฺเตสุ วชฺเชสุ ภยทสฺสาวี, สมาทาย สีกฺขติ สิกฺขาปเทสุ, กายกมฺม\-วจีกมฺเมน สมนฺนาคโต กุสเลน ปริสุทฺธาชีโว. อารทฺธวีริโย โข ปน โหติ ถามวา ทฬฺหปรกฺกโม อนิกฺขิตฺตขุโร อกุสลานํ ธมฺมานํ ปหานาย กุสลานํ ธมฺมานํ ภาวนาย สจฺฉิกิริยาย. ปญฺญวา โข ปน โหติ อุทย\-ตฺถคามินิยา ปญฺญาย สมนฺนาคโต อริยาย นิพฺเพธิกาย สมฺมา ทุกฺข\-กฺขย\-คามินิยา. อิมานิ โข ภิกฺขเว ภิกฺขูนํ ตีณิ กรณียานีติ. อิทํ อนุญฺญาตํ.}

\prose{\csromansymbolfootnote{*}{อํ 3. 325 ปิฏฺเฐ}ทสยิเม\csromansharedfootnote{39931d5581b3f106}{ทสอิเม (สี, ก)} ภิกฺขเว ธมฺมา ปพฺพชิเตน อภิณฺหํ ปจฺจเวกฺขิตพฺพา. กตเม ทส, “เววณฺณิยมฺหิ อชฺฌุปคโต”ติ ปพฺพชิเตน อภิณฺหํ ปจฺจเวกฺขิตพฺพํ ฯเปฯ. อิเม โข ภิกฺขเว ทส ธมฺมา ปพฺพชิเตน อภิณฺหํ ปจฺจเวกฺขิตพฺพาติ. อิทํ อนุญฺญาตํ.}

\prose{ตีณิมานิ ภิกฺขเว กรณียานิ. กตมานิ ตีณิ, กายสุจริตํ วจีสุจริตํ มโน\-สุจริตนฺติ. อิมานิ โข ภิกฺขเว ตีณิ กรณียานีติ. อิทํ อนุญฺญาตํ.}

\prose{ตตฺถ กตมํ ปฏิกฺขิตฺตํ.}

\prose{** “นตฺถิ ปุตฺตสมํ เปมํ,}

\setlength{\csromangathapagewidth}{0pt}
\setlength{\csromangathawakwidth}{0pt}
\csromangathameasuregroup{}{นตฺถิ โคสมิตํ\textsuperscript{1} ธนํ. \\ นตฺถิ สุริยสมา\textsuperscript{1} อาภา, \\ สมุทฺทปรมา สรา”ติ.}
\csromangathameasurewak{นตฺถิ โคสมิตํ\textsuperscript{1} ธนํ. \\ นตฺถิ สุริยสมา\textsuperscript{1} อาภา, \\ สมุทฺทปรมา สรา”ติ.}
\setlength{\csromangathaleftcol}{0pt}
\csromangathagroup{\csromangathastanza{\csromangathawak{นตฺถิ โคสมิตํ\csromansharedfootnote{93e7bec8090f8599}{โคณสมํ (ก)} ธนํ. \\ นตฺถิ สุริยสมา\csromansharedfootnote{93e7bec8090f8599}{โคณสมํ (ก)} อาภา, \\ สมุทฺทปรมา สรา”ติ.}}}

\prose{ภควา อาห\csromandash{}}

\setlength{\csromangathapagewidth}{0pt}
\setlength{\csromangathawakwidth}{0pt}
\csromangathameasuregroup{\textsuperscript{+}“นตฺถิ อตฺตสมํ เปมํ, \\ นตฺถิ ปญฺญาสมา อาภา, \\ เอตฺถ ยํ ปุริมกํ,}{\csromangathabat{\textsuperscript{+}“นตฺถิ อตฺตสมํ เปมํ,}{นตฺถิ ธญฺญสมํ ธนํ.} \\ \csromangathabat{นตฺถิ ปญฺญาสมา อาภา,}{วุฏฺฐิเวปรมา สรา”ติ.} \\ \csromangathabat{เอตฺถ ยํ ปุริมกํ,}{อิทํ ปฏิกฺขิตฺตํ.}}
\csromangathasetleft{\textsuperscript{+}“นตฺถิ อตฺตสมํ เปมํ, \\ นตฺถิ ปญฺญาสมา อาภา, \\ เอตฺถ ยํ ปุริมกํ,}
\csromangathagroup{\csromangathastanza{\csromangathabat{\csromansymbolmark{+}“นตฺถิ อตฺตสมํ เปมํ,}{นตฺถิ ธญฺญสมํ ธนํ.} \\ \csromangathabat{นตฺถิ ปญฺญาสมา อาภา,}{วุฏฺฐิเวปรมา สรา”ติ.} \\ \csromangathabat{เอตฺถ ยํ ปุริมกํ,}{อิทํ ปฏิกฺขิตฺตํ.}}}

\prose{ตีณิมานิ ภิกฺขเว อกรณียานิ. กตมานิ ตีณิ, กายทุจฺจริตํ วจีทุจฺจริตํ มโน\-ทุจฺจริตนฺติ. อิมานิ โข ภิกฺขเว ตีณิ อกรณียานีติ. อิทํ ปฏิกฺขิตฺตํ. (17)}

\setlength{\csromangathapagewidth}{0pt}
\setlength{\csromangathawakwidth}{0pt}
\csromangathameasuregroup{}{ตตฺถ กตมํ อนุญฺญาตญฺจ ปฏิกฺขิตฺตญฺจ. \\ “กิํ สูธ ภีตา ชนตา อเนกา, \\ มคฺโค จ’เนกายตโน ปวุตฺโต\textsuperscript{1}. \\ ปุจฺฉามิ ตํ โคตม ภูริปญฺญ, \\ กิสฺมิํ ฐิโต ปรโลกิ น ภาเยติ. \\ วาจํ มนญฺจ ปณิธาย สมฺมา, \\ กาเยน ปาปานิ อกุพฺพมาโน. \\ พหฺวนฺนปานํ ฆรมาวสนฺโต, \\ สทฺโธ มุทู สํวิภาคี วทญฺญู. \\ เอเตสุ ธมฺเมสุ ฐิโต จตูสุ, \\ ธมฺเม ฐิโต ปรโลกํ น ภาเย”ติ\textsuperscript{1}.}
\csromangathameasurewak{ตตฺถ กตมํ อนุญฺญาตญฺจ ปฏิกฺขิตฺตญฺจ. \\ “กิํ สูธ ภีตา ชนตา อเนกา, \\ มคฺโค จ’เนกายตโน ปวุตฺโต\textsuperscript{1}. \\ ปุจฺฉามิ ตํ โคตม ภูริปญฺญ, \\ กิสฺมิํ ฐิโต ปรโลกิ น ภาเยติ. \\ วาจํ มนญฺจ ปณิธาย สมฺมา, \\ กาเยน ปาปานิ อกุพฺพมาโน. \\ พหฺวนฺนปานํ ฆรมาวสนฺโต, \\ สทฺโธ มุทู สํวิภาคี วทญฺญู. \\ เอเตสุ ธมฺเมสุ ฐิโต จตูสุ, \\ ธมฺเม ฐิโต ปรโลกํ น ภาเย”ติ\textsuperscript{1}.}
\setlength{\csromangathaleftcol}{0pt}
\csromangathagroup{\csromangathastanza{\csromangathawak{\csromanfolio{161}\csromangathaitemline{124}{ตตฺถ กตมํ อนุญฺญาตญฺจ ปฏิกฺขิตฺตญฺจ.} \\ \csromangathacontline{“กิํ สูธ ภีตา ชนตา อเนกา,} \\ \csromangathacontline{มคฺโค จ’เนกายตโน ปวุตฺโต\csromansharedfootnote{6ad32d0e001a5a5a}{จเนกายตนปฺปวุตฺโต (สํ 1. 39 ปิฏฺเฐ)}.} \\ \csromangathacontline{ปุจฺฉามิ ตํ โคตม ภูริปญฺญ,} \\ \csromangathacontline{กิสฺมิํ ฐิโต ปรโลกิ น ภาเยติ.} \\ \csromangathacontline{วาจํ มนญฺจ ปณิธาย สมฺมา,} \\ \csromangathacontline{กาเยน ปาปานิ อกุพฺพมาโน.} \\ \csromangathacontline{พหฺวนฺนปานํ ฆรมาวสนฺโต,} \\ \csromangathacontline{สทฺโธ มุทู สํวิภาคี วทญฺญู.} \\ \csromangathacontline{เอเตสุ ธมฺเมสุ ฐิโต จตูสุ,} \\ \csromangathacontline{ธมฺเม ฐิโต ปรโลกํ น ภาเย”ติ\csromansharedfootnote{c9a13e384cec98d4}{สํ 1. 39, 40 ปิฏฺเฐสุ.}.}}}}

\prose{ตตฺถ ยํ อาห “วาจํ มนญฺจ ปณิธาย สมฺมา”ติ, อิทํ อนุญฺญาตํ. “กาเยน ปาปานิ อกุพฺพมาโน”ติ, อิทํ ปฏิกฺขิตฺตํ. “พหฺวนฺนปานํ ฆรมาวสนฺโต, สทฺโธ มุทู สํวิภาคี วทญฺญู. เอเตสุ ธมฺเมสุ ฐิโต จตูสุ, ธมฺเม ฐิโต ปรโลกํ นภาเย”ติ, อิทํ อนุญฺญาตํ. อิทํ อนุญฺญาตญฺจ ปฏิกฺขิตฺตญฺจ.}

\setlength{\csromangathapagewidth}{0pt}
\setlength{\csromangathawakwidth}{0pt}
\csromangathameasuregroup{\textsuperscript{*}สพฺพปาปสฺส อกรณํ, \\ สจิตฺตปริโยทาปนํ,}{\csromangathabat{\textsuperscript{*}สพฺพปาปสฺส อกรณํ,}{กุสลสฺส อุปสมฺปทา.} \\ \csromangathabat{สจิตฺตปริโยทาปนํ,}{เอตํ พุทฺธานสาสนํ.}}
\csromangathasetleft{\textsuperscript{*}สพฺพปาปสฺส อกรณํ, \\ สจิตฺตปริโยทาปนํ,}
\csromangathagroup{\csromangathastanza{\csromangathabat{\csromansymbolfootnote{*}{เหฏฺฐา 37, 68, 148 ปิฏฺเฐสุ; อุปริ 205, 208 ปิฏฺเฐสุปิ. ** ที 2. 223 ปิฏฺฐาทีสุ.}สพฺพปาปสฺส อกรณํ,}{กุสลสฺส อุปสมฺปทา.} \\ \csromangathabat{สจิตฺต\-ปริโยทาปนํ,}{เอตํ พุทฺธานสาสนํ.}}}

\prose{ตตฺถ ยํ อาห “สพฺพปาปสฺส อกรณนฺ”ติ, อิทํ ปฏิกฺขิตฺตํ, ยํ อาห “กุสลสฺส อุปสมฺปทา”ติ, อิทํ อนุญฺญาตํ. อิทํ อนุญฺญาตญฺจ ปฏิกฺขิตฺตญฺจ.}

\prose{** กายสมาจารมฺปาหํ เทวานมินฺท ทุวิเธน วทามิ เสวิตพฺพมฺปิ อเสวิตพฺพมฺปิ. วจีสามาจารมฺปาหํ เทวานมินฺท ทุวิเธน วทามิ เสวิตพฺพมฺปิ อเสวิตพฺพมฺปิ. มโนสมาจารมฺปาหํ เทวานมินฺท ทุวิเธน วทามิ ฯเปฯ. ปริเยสนมฺปาหํ เทวานมินฺท ทุวิเธน วทามิ เสวิตพฺพมฺปิ อเสวิตพฺพมฺปิ.}

\prose{“กายสมาจารมฺปาหํ เทวานมินฺท ทุวิเธน วทามิ เสวิตพฺพมฺปิ อเสวิตพฺพมฺปี”ติ อิติ โข ปเนตํ วุตฺตํ, กิญฺเจตํ ปฏิจฺจ วุตฺตํ. ยถารูปญฺจ โข กายสมาจารํ เสวโต อกุสลา ธมฺมา อภิวฑฺฒนฺติ, กุสลา ธมฺมา ปริหายนฺติ, เอวรูโป กายสมาจาโร น เสวิตพฺโพ. ตตฺถ ยํ ชญฺญา \csromanfolio{162}กายสมาจารํ “อิมํ โข เม กายสมาจารํ เสวโต อกุสลา ธมฺมา ปริหายนฺติ, กุสลา ธมฺมา อภิวฑฺฒนฺตี”ติ, เอวรูโป กายสมาจาโร เสวิตพฺโพ. “กายสมาจารมฺปาหํ เทวานมินฺท ทุวิเธน วทามิ เสวิตพฺพมฺปิ อเสวิตพฺพมฺปี”ติ อิติ ยํ ตํ วุตฺตํ, อิทเมติ ปฏิจฺจ วุตฺตํ. “วจีสมาจารํ ฯเปฯ. “ปริเยสนมฺปาหํ เทวานมินฺท ทุวิเธน วทามิ เสวิตพฺพมฺปิ อเสวิตพฺพมฺปี”ติ อิติ โข ปเนตํ วุตฺตํ, กิญฺเจตํ ปฏิจฺจ วุตฺตํ. ยถารูปญฺจ โข ปริเยสนํ เสวโต อกุสลา ธมฺมา อภิวฑฺฒนฺติ, กุสลา ธมฺมา ปริหายนฺติ, เอวรูปา ปริเยสนา น เสวิตพฺพา. ตตฺถ ยํ ชญฺญา ปริเยสนํ “อิมํ โข เม ปริเยสนํ เสวโต อกุสลา ธมฺมา ปริหายนฺติ, กุสลา ธมฺมา อภิวฑฺฒนฺตี”ติ, เอวรูปา ปริเยสนา เสวิตพฺพา. “ปริเยสนมฺปาหํ เทวานมินฺท ทุวิเธน วทามิ เสวิตพฺพมฺปิ อเสวิตพฺพมฺปี”ติ อิติ ยํ ตํ วุตฺตํ, อิทเมตํ ปฏิจฺจ วุตฺตํ.}

\prose{ตตฺถ ยํ อาห “เสวิตพฺพมฺปี”ติ, อิทํ อนุญฺญาตํ. ยํ อาห “น เสวิตพฺพมฺปี”ติ, อิทํ ปฏิกฺขิตฺตํ. อิทํ อนุญฺญาตญฺจ ปฏิกฺขิตฺตญฺจ. (18)}

\setlength{\csromangathapagewidth}{0pt}
\setlength{\csromangathawakwidth}{0pt}
\csromangathameasuregroup{}{ตตฺถ กตโม ถโว. \\ “มคฺคานฏฺฐงฺคิโก เสฏฺโฐ, สจฺจานํ จตุโร ปทา. \\ วิราโค เสฏฺโฐ ธมฺมานํ, ทฺวิปทานญฺจ จกฺขุมา”ติ\textsuperscript{1}.}
\csromangathameasurewak{ตตฺถ กตโม ถโว. \\ “มคฺคานฏฺฐงฺคิโก เสฏฺโฐ, สจฺจานํ จตุโร ปทา. \\ วิราโค เสฏฺโฐ ธมฺมานํ, ทฺวิปทานญฺจ จกฺขุมา”ติ\textsuperscript{1}.}
\setlength{\csromangathaleftcol}{0pt}
\csromangathagroup{\csromangathastanza{\csromangathawak{\csromangathaitemline{170}{ตตฺถ กตโม ถโว.} \\ \csromangathacontline{“มคฺคานฏฺฐงฺคิโก เสฏฺโฐ, สจฺจานํ จตุโร ปทา.} \\ \csromangathacontline{วิราโค เสฏฺโฐ ธมฺมานํ, ทฺวิปทานญฺจ จกฺขุมา”ติ\csromansharedfootnote{8688dde4a6623752}{ขุ 1. 52 ปิฏฺเฐ; อุปริ 206 ปิฏฺเฐปิ.}.}}}}

\prose{อยํ ถโว.}

\prose{ตีณิมานิ ภิกฺขเว อคฺคานิ. กตมานิ ตีณิ, ยาวตา ภิกฺขเว สตฺตา อปทา วา ทฺวิปทา วา จตุปฺปทา วา พหุปฺปทา วา รูปิโน วา อรูปิโน วา สญฺญิโน วา อสญฺญิโน วา เนว\-สญฺญี\-นา\-สญฺญิโน วา, ตถาคโต เตสํ อคฺคมกฺขายติ เสฏฺฐม\-กฺขายติ ปวรม\-กฺขายติ, ยทิทํ อรหํ สมฺมาสมฺพุทฺโธ. ยาวตา ภิกฺขเว ธมฺมานํ ปณฺณตฺติ\-สงฺขตานํ วา อสงฺขตานํ วา, วิราโค เตสํ ธมฺมานํ อคฺคมกฺขายติ เสฏฺฐม\-กฺขายติ ปวรม\-กฺขายติ, ยทิทํ มทนิมฺมทโน ฯเปฯ นิโรโธ นิพฺพานํ. ยาวตา ภิกฺขเว สํฆานํ ปณฺณตฺติ คณานํ ปณฺณตฺติ มหาชน\-สนฺนิปาตานํ ปณฺณตฺติ, ตถาคต\-สาวกสํโฆ เตสํ อคฺคมกฺขายติ เสฏฺฐม\-กฺขายติ ปวรม\-กฺขายติ, ยทิทํ จตฺตาริ ปุริสยุคานิ อฏฺฐ ปุริสปุคฺคลา ฯเปฯ ปุญฺญกฺเขตฺตํ โลกสฺสาติ\csromansharedfootnote{6a85e7188c227d4d}{อํ 1. 343 ปิฏฺเฐ; ขุ 1. 254 ปิฏฺเฐ จ ปสฺสิตพฺพํ.}.}

\setlength{\csromangathapagewidth}{0pt}
\setlength{\csromangathawakwidth}{0pt}
\csromangathameasuregroup{สพฺพโลกุตฺตโร สตฺถา, \\ คโณ จ นรสีหสฺส,}{\csromangathabat{สพฺพโลกุตฺตโร สตฺถา,}{ธมฺโม จ กุสลกฺขโต\textsuperscript{1}.} \\ \csromangathabat{คโณ จ นรสีหสฺส,}{ตานิ ตีณิ วิสฺสิสฺสเร.} \\ สมณปทุมสญฺจโย คโณ, \\ ธมฺมวโร จ วิทูนํ สกฺกโต. \\ นรวรทมโก จ จกฺขุมา, \\ ตานิ ตีณิ โลกสฺส อุตฺตริ. \\ สตฺถา จ อปฺปฏิสโม, \\ ธมฺโม จ สพฺโพ นิรุปทาโห. \\ อริโย จ คณวโร, \\ ตานิ ขลุ วิสฺสิสฺสเร ตีณิ. \\ สจฺจนาโม ชิโน เขโม สพฺพาภิภู, \\ สจฺจธมฺโม นตฺถญฺโญ ตสฺส อุตฺตริ. \\ อริยสํโฆ นิจฺจํ วิญฺญูนํ ปูชิโต, \\ ตานิ ตีณิ โลกสฺส อุตฺตริ. \\ \textsuperscript{*}เอกายนํ ชาติขยนฺตทสฺสี, \\ มคฺคํ ปชานาติ หิตานุกมฺปี. \\ เอเตน มคฺเคน ตริํสุ ปุพฺเพ, \\ ตริสฺสนฺติ เย จ\textsuperscript{1} ตรนฺติ โอฆํ. \\ ตํ ตาทิสํ เทวมนุสฺสเสฏฺฐํ, \\ สตฺตา นมสฺสนฺติ วิสุทฺธิปิกฺขา”ติ.}
\csromangathameasurewak{สมณปทุมสญฺจโย คโณ, \\ ธมฺมวโร จ วิทูนํ สกฺกโต. \\ นรวรทมโก จ จกฺขุมา, \\ ตานิ ตีณิ โลกสฺส อุตฺตริ. \\ สตฺถา จ อปฺปฏิสโม, \\ ธมฺโม จ สพฺโพ นิรุปทาโห. \\ อริโย จ คณวโร, \\ ตานิ ขลุ วิสฺสิสฺสเร ตีณิ. \\ สจฺจนาโม ชิโน เขโม สพฺพาภิภู, \\ สจฺจธมฺโม นตฺถญฺโญ ตสฺส อุตฺตริ. \\ อริยสํโฆ นิจฺจํ วิญฺญูนํ ปูชิโต, \\ ตานิ ตีณิ โลกสฺส อุตฺตริ. \\ \textsuperscript{*}เอกายนํ ชาติขยนฺตทสฺสี, \\ มคฺคํ ปชานาติ หิตานุกมฺปี. \\ เอเตน มคฺเคน ตริํสุ ปุพฺเพ, \\ ตริสฺสนฺติ เย จ\textsuperscript{1} ตรนฺติ โอฆํ. \\ ตํ ตาทิสํ เทวมนุสฺสเสฏฺฐํ, \\ สตฺตา นมสฺสนฺติ วิสุทฺธิปิกฺขา”ติ.}
\csromangathasetleft{สพฺพโลกุตฺตโร สตฺถา, \\ คโณ จ นรสีหสฺส,}
\csromangathagroup{\csromangathastanza{\csromanfolio{163}\csromangathabat{สพฺพโลกุตฺตโร สตฺถา,}{ธมฺโม จ กุสลกฺขโต\csromansharedfootnote{e7f8c5337bc76717}{กุสลมกฺขโต (ก)}.} \\ \csromangathabat{คโณ จ นรสีหสฺส,}{ตานิ ตีณิ วิสฺสิสฺสเร.}}\csromangathastanza{\csromangathawak{สมณปทุมสญฺจโย คโณ, \\ ธมฺมวโร จ วิทูนํ สกฺกโต. \\ นรวรทมโก จ จกฺขุมา, \\ ตานิ ตีณิ โลกสฺส อุตฺตริ.}}\csromangathastanza{\csromangathawak{สตฺถา จ อปฺปฏิสโม, \\ ธมฺโม จ สพฺโพ นิรุปทาโห. \\ อริโย จ คณวโร, \\ ตานิ ขลุ วิสฺสิสฺสเร ตีณิ.}}\csromangathastanza{\csromangathawak{สจฺจนาโม ชิโน เขโม สพฺพาภิภู, \\ สจฺจธมฺโม นตฺถญฺโญ ตสฺส อุตฺตริ. \\ อริยสํโฆ นิจฺจํ วิญฺญูนํ ปูชิโต, \\ ตานิ ตีณิ โลกสฺส อุตฺตริ.}}\csromangathastanza{\csromangathawak{\csromansymbolfootnote{*}{สํ 3. 146; ขุ 7. 362; ขุ 8. 212 ปิฏฺเฐสุ.}เอกายนํ ชาติ\-ขยนฺต\-ทสฺสี, \\ มคฺคํ ปชานาติ หิตานุกมฺปี. \\ เอเตน มคฺเคน ตริํสุ ปุพฺเพ, \\ ตริสฺสนฺติ เย จ\csromansharedfootnote{e58e4f42c1ca597a}{เย จาปิ (สี, ก)} ตรนฺติ โอฆํ. \\ ตํ ตาทิสํ เทวมนุสฺสเสฏฺฐํ, \\ สตฺตา นมสฺสนฺติ วิสุทฺธิปิกฺขา”ติ.}}}

\prose{อยํ ถโวติ.}

\prose{ตตฺถ โลกิยํ สุตฺตํ ทฺวีหิ สุตฺเตหิ นิทฺทิสิตพฺพํ สํกิเลส\-ภาคิเยน จ วาสนา\-ภาคิเยน จ. โลกุตฺตรํปิ สุตฺตํ ตีหิ สุตฺเตหิ นิทฺทิสิตพฺพํ ทสฺสน\-ภาคิเยน จ ภาวนา\-ภาคิเยน จ อเสกฺขภาคิเย น จ. โลกิยญฺจ โลกุตฺตรญฺจ. ยสฺมิํ สุตฺเต ยํ ยํ ปทํ ทิสฺสติ สํกิเลส\-ภาคิยํ วา วาสนาภาคิยํ วา, เตน เตน โลกิยนฺติ นิทฺทิสิตพฺพํ, ทสฺสน\-ภาคิยํ วา ภาวนา\-ภาคิยํ วา อเสกฺข\-ภาคิยํ วา ยํ ยํ ปทํ ทิสฺสติ เตน เตน โลกุตฺตรนฺติ นิทฺทิสิตพฺพํ.}

\prose{\csromanfolio{164}วาสนาภาคิยํ สุตฺตํ สํกิเลส\-ภาคิยสฺส สุตฺตสฺส นิคฺฆาตาย, ทสฺสน\-ภาคิยํ สุตฺตํ วาสนา\-ภาคิยสฺส สุตฺตสฺส นิคฺฆาตาย, ภาวนา\-ภาคิยํ สุตฺตํ ทสฺสน\-ภาคิยสฺส สุตฺตสฺส ปฏินิสฺสคฺคาย, อเสกฺข\-ภาคิยํ สุตฺตํ ภาวนา\-ภาคิยสฺส สุตฺตสฺส ปฏินิสฺสคฺคาย, อเสกฺข\-ภาคิยํ สุตฺตํ ทิฏฺฐธมฺมสุขวิหารตฺถํ.}

\prose{โลกุตฺตรํ สุตฺตํ สตฺตา\-ธิฏฺฐานํ ฉพฺพีสติยา ปุคฺคเลหิ นิทฺทิสิตพฺพํ, เต ตีหิ สุตฺเตหิ สมเนฺวสิตพฺพา ทสฺสน\-ภาคิเยน ภาวนา\-ภาคิเยน อเสกฺข\-ภาคิเยน จาติ.}

\prose{ตตฺถ ทสฺสน\-ภาคิยํ สุตฺตํ ปญฺจหิ ปุคฺคเลหิ นิทฺทิสิตพฺพํ เอกพีชินา โกลํโกเลน สตฺตกฺขตฺตุ\-ปรเมน สทฺธานุสารินา ธมฺมานุสารินา จาติ, ทสฺสน\-ภาคิยํ สุตฺตํ อิเมหิ ปญฺจหิ ปุคฺคเลหิ นิทฺทิสิตพฺพํ.  ภาวนา\-ภาคิยํ สุตฺตํ ทฺวาทสหิ ปุคฺคเลหิ นิทฺทิสิตพฺพํ สกทาคามิ\-ผล\-สจฺฉิกิริยาย ปฏิปนฺเนน, สกทาคามินา, อนาคามิ\-ผล\-สจฺฉิกิริยาย ปฏิปนฺเนน, อนาคามินา, อนฺตรา ปรินิพฺพายินา, อุปหจฺจ ปรินิพฺพายินา, อสงฺขารปรินิพฺพายินา, สสงฺขาร\-ปรินิพฺพายินา, อุทฺธํโสเตน อกนิฏฺฐิคมินา, สทฺธา\-วิมุตฺเตน, ทิฏฺฐิปฺปตฺเตน, กายสกฺขินา จาติ, ภาวนา\-ภาคิยํ สุตฺตํ อิเมหิ ทฺวาทสหิ ปุคฺคเลหิ นิทฺทิสิตพฺพํ.  อเสกฺข\-ภาคิยํ สุตฺตํ นวหิ ปุคฺคเลหิ นิทฺทิสิตพฺพํ สทฺธา\-วิมุตฺเตน, ปญฺญาวิมุตฺเตน, สุญฺญต\-วิมุตฺเตน, อนิมิตฺต\-วิมุตฺเตน, อปฺปณิหิต\-วิมุตฺเตน, อุภโตภาค\-วิมุตฺเตน สมสีสินา ปจฺเจกพุทฺธ\-สมฺมา\-สมฺพุทฺเธหิ จาติ, อเสกฺข\-ภาคิยํ สุตฺตํ อิเมหิ นวหิ ปุคฺคเลหิ นิทฺทิสิตพฺพํ. เอวํ โลกุตฺตรํ สุตฺตํ สตฺตา\-ธิฏฺฐานํ อิเมหิ ฉพฺพีสติยา ปุคฺคเลหิ นิทฺทิสิตพฺพํ.}

\prose{โลกิยํ สุตฺตํ สตฺตา\-ธิฏฺฐานํ เอกูนวีสติยา ปุคฺคเลหิ นิทฺทิสิตพฺพํ, เต จริเตหิ นิทฺทิฏฺฐา สมเนฺวสิตพฺพา เกจิ ราคจริตา, เกจิ โทสจริตา, เกจิ โมหจริตา, เกจิ ราคจริตา จ โทสจริตา จ, เกจิ ราคจริตา จ โมหจริตา จ, เกจิ โทสจริตา จ โมหจริตา จ, เกจิ ราคจริตา จ โทสจริตา จ โมหจริตา จ, ราคมุเข ฐิโต ราคจริโต, ราคมุเข ฐิโต โทสจริโต, ราคมุเข ฐิโต โมหจริโต, ราคมุเข ฐิโต ราคจริโต จ โทสจริโต จ โมหจริโต จ, โทสมุเข ฐิโต โทสจริโต, โทสมุเข ฐิโต โมหจริโต, โทสมุเข ฐิโต ราคจริโต, โทสมุเข ฐิโต ราคจริโต จ \csromanfolio{165}โทสจริโต จ โมหจริโต จ, โมหมุเข ฐิโต โมหจริโต, โมหมุเข ฐิโต ราคจริโต โมหมุเข ฐิโต โทสจริโต, โมหมุเข ฐิโต ราคจริโต จ โทสจริโต จ โมหจริโต จาติ, โลกิยํ สุตฺตํ สตฺตา\-ธิฏฺฐานํ อิเมหิ เอกูนวีสติยา ปุคฺคเลหิ นิทฺทิสิตพฺพํ.}

\prose{วาสนาภาคิยํ สุตฺตํ สีลวนฺเตหิ นิทฺทิสิตพฺพํ, เต สีลวนฺโต ปญฺจ ปุคฺคลา ปกติสีลํ สมาทานสีลํ จิตฺตปฺปสาโท สมโถ วิปสฺสนา จาติ, วาสนาภาคิยํ สุตฺตํ อิเมหิ ปญฺจหิ ปุคฺคเลหิ นิทฺทิสิตพฺพํ.}

\prose{โลกุตฺตรํ สุตฺตํ ธมฺมา\-ธิฏฺฐานํ ตีหิ สุตฺเตหิ นิทฺทิสิตพฺพํ ทสฺสน\-ภาคิเยน ภาวนา\-ภาคิเยน อเสกฺข\-ภาคิเยน จ.}

\prose{โลกิยญฺจ โลกุตฺตรญฺจ สตฺตา\-ธิฏฺฐานญฺจ ธมฺมาธิฏฺฐานญฺจ อุภเยน นิทฺทิสิตพฺพํ, ญาณํ ปญฺญาย นิทฺทิสิตพฺพํ ปญฺญินฺทฺริเยน ปญฺญาพเลน อธิปญฺญา\-สิกฺขาย ธมฺมวิจย\-สมฺโพชฺฌงฺเคน สมฺมาทิฏฺฐิยา ตีรณาย สนฺตีรณาย ธมฺเม ญาเณน อนฺวเย ญาเณน ขเย ญาเณน อนุปฺปาเท ญาเณน อนญฺญาต\-ญฺญสฺสามี\-ตินฺทฺริเยน อญฺญินฺเทฺรเยน อญฺญาตาวิ\-นฺทฺริเยน จกฺขุนา วิชฺชาย พุทฺธิยา ภูริยา เมธาย, ยํ ยํ วา ปน ลพฺภติ, เตน เตน ปญฺญา\-ธิวจเนน นิทฺทิสิตพฺพํ.}

\prose{เญยฺยํ อตีตา\-นาคต\-ปจฺจุปฺปนฺเนหิ อชฺฌตฺติก\-พาหิเรหิ หีนปฺปณีเตหิ ทูรสนฺติเกหิ สงฺขตา\-สงฺขเตหิ กุสลากุสลา\-พฺยากเตหิ สงฺเขปโต วา ฉหิ อารมฺมเณหิ นิทฺทิสิตพฺพํ.  ญาณญฺจ เญ ยฺยญฺจ ตทุภเยน นิทฺทิสิตพฺพํ, ปญฺญาปิ อารมฺมณภูตา เญยฺยํ, ยํ กิญฺจิ อารมฺมณภูตํ อชฺฌตฺติกํ วา พาหิรํ วา, สพฺพํ ตํ สงฺขเตน อสงฺขเตน จ นิทฺทิสิตพฺพํ.}

\prose{ทสฺสนํ ภาวนา\csromansharedfootnote{2bbc3d7fa7673c05}{ทสฺสนา ภาวนา (สี)} สกวจนํ ปรวจนํ วิสชฺชนียํ อวิสชฺชนียํ กมฺมํ วิปาโกติ สพฺพตฺถ ตทุภยํ สุตฺเต ยถา นิทฺทิฏฺฐํ, ตถา อุปธารยิตฺวา ลพฺภมานโต นิทฺทิสิตพฺพํ, ยํ วา ปน กิญฺจิ ภควา อญฺญตร\-วจนํ ภาสติ, สพฺพํ ตํ ยถา\-นิทฺทิฏฺฐํ ธารยิตพฺพํ.}

\prose{ทุวิโธ เหตุ ยญฺจ กมฺมํ เย จ กิเลสา, สมุทโย กิเลสา. ตตฺถ กิเลสา สํกิเลส\-ภาคิเยน สุตฺเตน นิทฺทิสิตพฺพา, สมุทโย \csromanfolio{166}สํกิเลส\-ภาคิเยน จ วาสนา\-ภาคิเยน จ สุตฺเตน นิทฺทิสิตพฺโพ. ตตฺถ กุสลํ จตูหิ สุตฺเตหิ นิทฺทิสิตพฺพํ วาสนา\-ภาคิเยน ทสฺสน\-ภาคิเยน \textbf{ภาวนา}\-\textbf{ภาคิเยน อเสกฺข}\-\textbf{ภาคิเยน จ. อกุสลํ สํ กิเลสภาคิเยน} สุตฺเตน นิทฺทิสิตพฺพํ. กุสลญฺจ อกุสลญฺจ ตทุภเยน\csromansharedfootnote{d4820aecc669a54b}{ตทุภเยหิ (สี)} นิทฺทิสิตฺตพฺพํ. \textbf{อนุญฺญาตํ ภควโต อนุญฺญาตาย นิทฺทิสิตพฺพํ, ตํ ปญฺจวิธํ} สํวโร ปหานํ ภาวนา สจฺฉิกิริยา กปฺปิยานุโลโมติ, ยํ ทิสฺสติ ตาสุ ตาสุ ภูมีสุ, ตํ กปฺปิยา\-นุโลเมน นิทฺทิสิตพฺพํ. ปฏิกฺขิตฺตํ ภควตา ปฏิกฺขิตฺต\-การเณน นิทฺทิสิตพฺพํ. อนุญฺญาตญฺจ ปฏิกฺขิตฺตญฺจ ตทุภเยน นิทฺทิสิตพฺพํ. ถโว ปสํสาย นิทฺทิสิตพฺโพ, โส ปญฺจวิเธน เวทิตพฺโพ \textbf{ภควโต ธมฺมสฺส อริยสํฆสฺส อริยธมฺมานํ สิกฺขาย} \textbf{โลกิย}\-\textbf{คุณสมฺปตฺติยาติ. เอวํ ถโว ปญฺจวิเธน นิทฺทิสิตพฺโพ.}}

\prose{อินฺทฺริยภูมิ นวหิ ปเทหิ นิทฺทิสิตพฺพา, กิเลสภูมิ นวหิ ปเทหิ นิทฺทิสิตพฺพา, เอวเมตานิ อฏฺฐารส ปทานิ โหนฺติ นว ปทานิ กุสลานิ นวปทานิ อกุสลานีติ, ตถาหิ วุตฺตํ “อฏฺฐรส มูลปทา กุหิํ ทฏฺฐพฺพา, สาสน\-ปฺปฏฺฐาเน”ติ\csromansharedfootnote{e1d78de9b5e58064}{เหฏฺฐา 106 ปิฏฺเฐ.}. เตนาห อายสฺมา มหากจฺจายโน\csromandash{}}

\setlength{\csromangathapagewidth}{0pt}
\setlength{\csromangathawakwidth}{0pt}
\csromangathameasuregroup{}{“นวหิ จ ปเทหิ กุสลา, \\ นวหิ จ ยุชฺชนฺติ อกุสลปกฺขา. \\ เอเต โข\textsuperscript{1} มูลปทา, \\ ภวนฺติ อฏฺฐารส ปทานี”ติ\textsuperscript{1}.}
\csromangathameasurewak{“นวหิ จ ปเทหิ กุสลา, \\ นวหิ จ ยุชฺชนฺติ อกุสลปกฺขา. \\ เอเต โข\textsuperscript{1} มูลปทา, \\ ภวนฺติ อฏฺฐารส ปทานี”ติ\textsuperscript{1}.}
\setlength{\csromangathaleftcol}{0pt}
\csromangathagroup{\csromangathastanza{\csromangathawak{“นวหิ จ ปเทหิ กุสลา, \\ นวหิ จ ยุชฺชนฺติ อกุสลปกฺขา. \\ เอเต โข\csromansharedfootnote{e316391c38bd1cfa}{ขลุ (ก)} มูลปทา, \\ ภวนฺติ อฏฺฐารส ปทานี”ติ\csromansharedfootnote{967ff8f112953c71}{เหฏฺฐา 3 ปิฏฺเฐ.}.}}}

\vspace{\tipitakaparskip}
\csromancenter{นิยุตฺตํ สาสน\-ปฺปฏฺฐานํ.}
\prosewithcloser{เอตฺตาวตาสมตฺตา เนตฺติ ยา อายสฺมตา มหา\-กจฺจายเนน ภาสิตา ภควตา อนุโมทิตา มูลสงฺคีติยํ สงฺคีตาติ.}{\textbf{เนตฺติปกรณํ นิฏฺฐิตํ.}}
\csromanpalirecto
\csromanheader{2}{\textbf{ขุทฺทกนิกาย}}
\csromanmatikaanchor{mtk.39}
\csromantocmarknopage{chapter}{เปฏโกปเทสปาฬิ}{mtk.39}
\bookmark[dest={mtk.39},level=0]{เปฏโกปเทสปาฬิ}
\gambhira{\textbf{เปฏโกปเทส}\-\textbf{ปาฬิ}}
\namakkaram{นโม ตสฺส ภควโต อรหโต สมฺมา\-สมฺพุทฺธสฺส.}
\csromanmatikaanchor{mtk.40}
\csromantocmarkref{section}{1. อริยสจฺจปฺปกาสนปฐมภูมิ}{mtk.40}
\bookmark[dest={mtk.40},level=1]{1. อริยสจฺจปฺปกาสนปฐมภูมิ}
\prose{\csromanfolio{167}นโม สมฺมา\-สมฺพุทฺธานํ ปรมตฺถ\-ทสฺสีนํ สีลาทิคุณ\-ปารมิปฺปตฺตานํ.}

\proseitem{1}{ทุเว เหตู ทุเว ปจฺจยา สาวกสฺส สมฺมาทิฏฺฐิยา อุปฺปาทาย ปรโต จ โฆโส สจฺจานุสนฺธิ อชฺฌตฺตญฺจ โยนิโส มนสิกาโร. ตตฺถ กตโม ปรโต โฆโส, ยา ปรโต เทสนา โอวาโท อนุสาสนี สจฺจกถา สจฺจานุโลโม. จตฺตาริ สจฺจานิ ทุกฺขํ สมุทโย นิโรโธ มคฺโค. อิเมสํ จตุนฺนํ สจฺจานํ ยา เทสนา สนฺทสฺสนา วิวรณา วิภชนา อุตฺตานีกิริยา\csromansharedfootnote{96d03426ef9ce515}{อุตฺตานิกิริยา (ก)} ปกาสนา. อยํ วุจฺจติ สจฺจานุโลโม โฆโสติ.}

\proseitem{2}{ตตฺถ กตโม อชฺฌตฺตํ โยนิโส มนสิกาโร. อชฺฌตฺตํ โยนิโส มนสิกาโร นาม โย ยถาเทสิเต ธมฺเม พหิทฺธา อารมฺมณํ อนภินีหริตฺวา โยนิโส มนสิกาโร. อยํ วุจฺจติ โยนิโส มนสิกาโร. ตํอากาโร โยนิโส ทฺวาโร วิธิ อุปาโย. ยถา ปุริโส สุกฺเข กฏฺเฐ วิคตสฺเนเห สุกฺขาย อุตฺตรารณิยา ถเล อภิมนฺถ\-มานํ ภพฺโพ โชติสฺส อธิคมาย. ตํ กิสฺส เหตุ. โยนิโส อคฺคิสฺส อธิคมาย. เอวเมวสฺส ยมิทํ ทุกฺขสมุทย\-นิโรธ\-มคฺคานํ อวิปรีตธมฺม\-เทสนํ มนสิกโรติ. อยํ วุจฺจติ โยนิโส มนสิกาโร.}

\prose{\csromanfolio{168}ยถา ติสฺโส อุปมา ปุพฺเพ อสฺสุตา จ อสฺสุตปุพฺพา จ ปฏิภนฺติ. โย หิ โก จิ กาเมสุ อวีตราโคติ ฯเปฯ ทุเว อุปมา อโยนิโส กาตพฺพา ปจฺฉิเมสุ วุตฺตํ. ตตฺถ โย จ ปรโต โฆโสโย จ อชฺฌตฺตํ โยนิโส มนสิกาโร. อิเม เทฺว ปจฺจยา ปรโต โฆเสน ยา อุปฺปชฺชติ ปญฺญา. อยํ วุจฺจติ สุตมยี ปญฺญา. ยา อชฺฌตฺตํ โยนิโส มนสิกาเรน อุปฺปชฺชติ ปญฺญา. อยํ วุจฺจติ จินฺตามยี ปญฺญาติ อิมา เทฺว ปญฺญา เวทิตพฺพา. ปุริมกา จ เทฺว ปจฺจยา อิเม เทฺว เหตู เทฺว ปจฺจยา สาวกสฺส สมฺมาทิฏฺฐิยา อุปฺปาทาย.}

\proseitem{3}{ตตฺถ ปรโต โฆสสฺส สจฺจา\-นุ\-สนฺธิสฺส เทสิตสฺส อตฺถํ อวิชานนฺโต อตฺถ\-ปฺปฏิสํเวที ภวิสฺสตีติ เนตํ ฐานํ วิชฺชติ. น จ อตฺถ\-ปฺปฏิสํเวที โยนิโส มนสิกริสฺสตีติ เนตํ ฐานํ วิชฺชติ. ปรโต โฆสสฺส สจฺจา\-นุ\-สนฺธิสฺส เทสิตสฺส อตฺถํ วิชานนฺโต อตฺถ\-ปฺปฏิสํเวที ภวิสฺสตีติ ฐานเมตํ วิชฺชติ. เอส เหตุ เอตํ อารมฺมณํ เอโส อุปาโย สาวกสฺส นิยฺยานสฺส นตฺถญฺโญ. โสยํ น จ สุตฺตสฺส อตฺถ\-วิชานนาย สห ยุตฺโต นาปิ โฆสานุโยเคน ปรโต โฆสสฺส อตฺถํ อวิชานนฺเตน สกฺกา อุตฺตริ\-มนุสฺสธมฺมํ อลมริย\-ญาณทสฺสนํ อธิคนฺตุํ, ตสฺมา นิพฺพายิตุ\-กาเมน สุตมเยน อตฺถา ปริเยสิตพฺพา. ตตฺถ ปริเยสนาย อยํ อนุปุพฺพี ภวติ โสฬส หารา, ปญฺจ นยา, อฏฺฐารส มูลปทานิ.}

\csromanheader{2}{\textbf{ตตฺถายํ อุทฺทานคาถา}}
\setlength{\csromangathapagewidth}{0pt}
\setlength{\csromangathawakwidth}{0pt}
\csromangathameasuregroup{โสฬสหารา เนตฺตี, \\ อฏฺฐารสมูลปทา,}{\csromangathabat{โสฬสหารา เนตฺตี,}{ปญฺจนยา สาสนสฺส ปริเยฏฺฐิ.} \\ \csromangathabat{อฏฺฐารสมูลปทา,}{กจฺจายนโคตฺตนิทฺทิฏฺฐา.}}
\csromangathasetleft{โสฬสหารา เนตฺตี, \\ อฏฺฐารสมูลปทา,}
\csromangathagroup{\csromangathastanza{\csromangathabat{โสฬสหารา เนตฺตี,}{ปญฺจนยา สาสนสฺส ปริเยฏฺฐิ.} \\ \csromangathabat{อฏฺฐารสมูลปทา,}{กจฺจายนโคตฺตนิทฺทิฏฺฐา.}}}

\proseitem{4}{ตตฺถ กตเม โสฬสหารา, เทสนา วิจโย ยุตฺติปทฏฺฐานํ ลกฺขณํ จตุพฺยูโห อาวฏฺโฏ วิภตฺติ ปริวตฺตโน เววจโน ปญฺญตฺติ โอตรโณ โสธโน อธิฏฺฐาโน ปริกฺขาโร สมาโรปโน, อิเม โสฬส หารา.}

\csromanheader{2}{1. อริย\-สจฺจ\-ปฺปกาสน\-ปฐมภูมิ \textbf{ตตฺถ อุทฺทานคาถา}}
\setlength{\csromangathapagewidth}{0pt}
\setlength{\csromangathawakwidth}{0pt}
\csromangathameasuregroup{\textsuperscript{*}เทสนา วิจโย ยุตฺติ, \\ จตุพฺยูโห จ อาวฏฺโฏ, \\ เววจโน จ ปญฺญตฺติ, \\ อธิฏฺฐาโน ปริกฺขาโร,}{\csromangathabat{\textsuperscript{*}เทสนา วิจโย ยุตฺติ,}{ปทฏฺฐาโน จ ลกฺขโณ\textsuperscript{1}.} \\ \csromangathabat{จตุพฺยูโห จ อาวฏฺโฏ,}{วิภตฺติ ปริวตฺตโน.} \\ \csromangathabat{เววจโน จ ปญฺญตฺติ,}{โอตรโณ จ โสธโน.} \\ \csromangathabat{อธิฏฺฐาโน ปริกฺขาโร,}{สมาโรปโน โสฬโส\csromandash{}\textsuperscript{1}.}}
\csromangathasetleft{\textsuperscript{*}เทสนา วิจโย ยุตฺติ, \\ จตุพฺยูโห จ อาวฏฺโฏ, \\ เววจโน จ ปญฺญตฺติ, \\ อธิฏฺฐาโน ปริกฺขาโร,}
\csromangathagroup{\csromangathastanza{\csromanfolio{169}\csromangathabat{\csromansymbolfootnote{*}{เหฏฺฐา 2 ปิฏฺเฐ.}เทสนา วิจโย ยุตฺติ,}{ปทฏฺฐาโน จ ลกฺขโณ\csromansharedfootnote{79d9b5fd77ccb58d}{ปทฏฺฐานญฺจ ลกฺขณํ (อิ)}.} \\ \csromangathabat{จตุพฺยูโห จ อาวฏฺโฏ,}{วิภตฺติ ปริวตฺตโน.}}\csromangathastanza{\csromangathabat{เววจโน จ ปญฺญตฺติ,}{โอตรโณ จ โสธโน.} \\ \csromangathabat{อธิฏฺฐาโน ปริกฺขาโร,}{สมาโรปโน โสฬโส\csromandash{}\csromansharedfootnote{4bc181d58ee06d6e}{โสฬส หารา (อิ, ก)}.}}}

\proseitem{5}{ตตฺถ กตเม ปญฺจนยา, นนฺทิยาวฏฺโฏ ติปุกฺขโล สีหวิกฺกีฬิโต ทิสาโลจโน องฺกุโสติ.}

\csromanheader{2}{\textbf{ตตฺถ อุทฺทานคาถา}}
\setlength{\csromangathapagewidth}{0pt}
\setlength{\csromangathawakwidth}{0pt}
\csromangathameasuregroup{ปฐโม นนฺทิยาวฏฺโฏ, \\ สีหวิกฺกีฬิโต นาม, \\ ทิสาโลจนมาหํสุ, \\ ปญฺจโม องฺกุโส นาม\textsuperscript{1},}{\csromangathabat{ปฐโม นนฺทิยาวฏฺโฏ,}{ทุติโย จ ติปุกฺขโล.} \\ \csromangathabat{สีหวิกฺกีฬิโต นาม,}{ตติโย โหติ โส นโย.} \\ \csromangathabat{ทิสาโลจนมาหํสุ,}{จตุตฺโถ นยลญฺชโก.} \\ \csromangathabat{ปญฺจโม องฺกุโส นาม\textsuperscript{1},}{สพฺเพ ปญฺจ นยา คตา.}}
\csromangathasetleft{ปฐโม นนฺทิยาวฏฺโฏ, \\ สีหวิกฺกีฬิโต นาม, \\ ทิสาโลจนมาหํสุ, \\ ปญฺจโม องฺกุโส นาม\textsuperscript{1},}
\csromangathagroup{\csromangathastanza{\csromangathabat{ปฐโม นนฺทิยาวฏฺโฏ,}{ทุติโย จ ติปุกฺขโล.} \\ \csromangathabat{สีหวิกฺกีฬิโต นาม,}{ตติโย โหติ โส นโย.}}\csromangathastanza{\csromangathabat{ทิสา\-โลจนมาหํสุ,}{จตุตฺโถ นยลญฺชโก.} \\ \csromangathabat{ปญฺจโม องฺกุโส นาม\csromansharedfootnote{4c771757706f5e18}{ปญฺจมํ องฺกุสํ อาหุ (อิ, ก)},}{สพฺเพ ปญฺจ นยา คตา.}}}

\proseitem{6}{ตตฺถ กตมานิ อฏฺฐารส มูลปทานิ อวิชฺชา ตณฺหา โลโภ โทโส โมโห สุภสญฺญา สุขสญฺญา นิจฺจสญฺญา อตฺตสญฺญา สมโถ วิปสฺสนา อโลโภ อโทโส อโมโห อสุภสญฺญา ทุกฺขสญฺญา อนิจฺจสญฺญา อนตฺตสญฺญา, อิมานิ อฏฺฐารส มูลปทานิ. ตตฺถ นว ปทานิ อกุสลานิ ยตฺถ สพฺพํ อกุสลํ สโมสรติ. นว ปทานิ กุสลานิ ยตฺถ สพฺพํ กุสลํ สโมสรติ.}

\prose{กตมานิ นว ปทานิ อกุสลานิ ยตฺถ สพฺพํ อกุสลํ สโมสรโต, อวิชฺชา ยาว อตฺตสญฺญา, อิมานิ นว ปทานิ อกุสลานิ, ยตฺถ สพฺพํ อกุสลํ สโมสรติ. กตมานิ นว ปทานิ กุสลานิ ยตฺถ สพฺพํ กุสลํ สโมสรติ. สมโถ ยาว อนตฺตสญฺญา, อิมานิ นว ปทานิ กุสลานิ ยตฺถ สพฺพํ กุสลํ สโมสรติ. อิมานิ อฏฺฐารส มูลปทานิ.}

\csromanpalirecto
\gambhira{เปฏโกปเทส\-ปาฬิ \textbf{ตตฺถ อิมา อุทฺทานคาถา}}
\setlength{\csromangathapagewidth}{0pt}
\setlength{\csromangathawakwidth}{0pt}
\csromangathameasuregroup{}{ตณฺหา จ อวิชฺชา โลโภ, \\ โทโส ตเถว โมโห จ. \\ จตฺตาโร จ วิปลฺลาสา, \\ กิเลสภูมิ นว ปทานิ. \\ เย จ สติปฏฺฐานา สมโถ, \\ วิปสฺสนา กุสลมูลํ. \\ เอตํ สพฺพํ กุสลํ, \\ อินฺทฺริยภูมิ นวปทานิ. \\ สพฺพํ กุสลํ นวหิ ปเทหิ ยุชฺชติ, \\ นวหิ เจว อกุสลํ. \\ เอกเก นว มูลปทานิ, \\ อุภยโต อฏฺฐารส มูลปทานิ.}
\csromangathameasurewak{ตณฺหา จ อวิชฺชา โลโภ, \\ โทโส ตเถว โมโห จ. \\ จตฺตาโร จ วิปลฺลาสา, \\ กิเลสภูมิ นว ปทานิ. \\ เย จ สติปฏฺฐานา สมโถ, \\ วิปสฺสนา กุสลมูลํ. \\ เอตํ สพฺพํ กุสลํ, \\ อินฺทฺริยภูมิ นวปทานิ. \\ สพฺพํ กุสลํ นวหิ ปเทหิ ยุชฺชติ, \\ นวหิ เจว อกุสลํ. \\ เอกเก นว มูลปทานิ, \\ อุภยโต อฏฺฐารส มูลปทานิ.}
\setlength{\csromangathaleftcol}{0pt}
\csromangathagroup{\csromangathastanza{\csromangathawak{\csromanfolio{170}ตณฺหา จ อวิชฺชา โลโภ, \\ โทโส ตเถว โมโห จ. \\ จตฺตาโร จ วิปลฺลาสา, \\ กิเลสภูมิ นว ปทานิ.}}\csromangathastanza{\csromangathawak{เย จ สติปฏฺฐานา สมโถ, \\ วิปสฺสนา กุสลมูลํ. \\ เอตํ สพฺพํ กุสลํ, \\ อินฺทฺริยภูมิ นวปทานิ.}}\csromangathastanza{\csromangathawak{สพฺพํ กุสลํ นวหิ ปเทหิ ยุชฺชติ, \\ นวหิ เจว อกุสลํ. \\ เอกเก นว มูลปทานิ, \\ อุภยโต อฏฺฐารส มูลปทานิ.}}}

\prose{อิเมสํ อฏฺฐารสนฺนํ มูลปทานํ ยานิ นว ปทานิ อกุสลานิ, อยํ ทุกฺขสมุทโย. ยานิ นว ปทานิ กุสลานิ, อยํ ทุกฺขนิโรธ\-คามินี ปฏิปทา. อิติ สมุทยสฺส ทุกฺขํ ผลํ, ทุกฺขนิโรธ\-คามินิยา ปฏิปทาย นิโรธํ ผลํ. อิมานิ จตฺตาริ อริยสจฺจานิ ภควตา พาราณสิยํ เทสิตานิ.}

\proseitem{7}{ตตฺถ ทุกฺขสฺส อริยสจฺจสฺส อปริมาณานิ อกฺขรานิ ปทานิ พฺยญฺจนานิ อาการานิ นิรุตฺติโย นิทฺเทสา เทสิตา เอตสฺเส\-ว\-ตฺถสฺส สงฺกาสนาย ปกาสนาย วิวรณาย วิภชนาย อุตฺตานีกมฺม\-ตาย ปญฺญาปนายาติ. ยา เอวํ สพฺเพสํ สจฺจานํ อิติ เอกเมกํ สจฺจํ อปริมาเณหิ อกฺขร\-ปทพฺยญฺชน\-อาการ\-นิรุตฺติ\-นิทฺเทเสหิ ปริเยสิตพฺพํ, ตญฺจ พฺยญฺชนํ อตฺถ\-ปุถุตฺเตน ปน อตฺเถว พฺยญฺชน\-ปุถุตฺเตน.}

\prose{โย หิ โกจิ สมโณ วา พฺราหฺมโณ วา เอวํ วเทยฺย “อหํ อิทํ ทุกฺขํ ปจฺจกฺขาย อญฺญํ ทุกฺขํ ปญฺญเปสฺสามี”ติ. ตสฺส ตํ วาจาวตฺถุกเม\-วสฺส, ปุจฺฉิโต จ น สมฺปายิสฺสติ. เอวํ สจฺจานิ ยญฺจ รตฺติํ ภควา อภิสมฺพุทฺโธ, ยญฺจ รตฺติํ อนุปาทาย ปรินิพฺพุโต, เอตฺถนฺตเร ยํ กิญฺจิ ภควตา ภาสิตํ สุตฺตํ เคยฺยํ เวยฺยากรณํ คาถา อุทานํ อิติวุตฺตกํ ชาตกํ อพฺภุตธมฺมํ เวทลฺลํ, สพฺพํ ตํ ธมฺมจกฺกํ ปวตฺติตํ. น กิญฺจิ พุทฺธานํ \csromanfolio{171}ภควนฺตานํ ธมฺมเทสนาย ธมฺมจกฺกโต พหิทฺธา, ตสฺส สพฺพํ สุตฺตํ อริยธมฺเมสุ ปริเยสิตพฺพํ. ตตฺถ ปริคฺคณฺหนาย อาโลกสภานิ จตฺตาริ อริยสจฺจานิ ถาวรานิ อิมานิ.}

\prose{ตตฺถ กตมํ ทุกฺขํ, ชาติ ชรา พฺยาธิ มรณํ สงฺขิตฺเตน ปญฺจุ\-ปาทาน\-กฺขนฺธา ทุกฺขา. ตตฺถายํ ลกฺขณ\-นิทฺเทโส, ปาตุภาว\-ลกฺขณา ชาติ, ปริปาก\-ลกฺขณา ชรา, ทุกฺขทุกฺขตา\-ลกฺขโณ พฺยาธิ, จุติลกฺขณํ มรณํ, ปิยวิปฺปโยควิปริณามปริตาปน- ลกฺขโณ โสโก, ลาลปฺปน\-ลกฺขโณ ปริเทโว, กาย\-สมฺปีฬน\-ลกฺขณํ ทุกฺขํ, จิตฺต\-สมฺปีฬน\-ลกฺขณํ โทมนสฺสํ, กิเลส\-ปริทหน\-ลกฺขโณ อุปายาโส, อมนาป\-สโมธาน\-ลกฺขโณ อปฺปิย\-สมฺปโยโค, มนาป\-วินาภาว\-ลกฺขโณ ปิยวิปฺปโยโค, อธิปฺปาย\-วิวตฺต\-นล\-กฺขโณ อลาโภ, อปริญฺญาลกฺขณา ปญฺจุ\-ปาทาน\-กฺขนฺธา, ปริปาก\-จุติ\-ลกฺขณํ ชรามรณํ, ปาตุภาว\-จุติ\-ลกฺขณํ จุโตปปตฺติ, ปฏิสนฺธิ\-นิพฺพตฺตน\-ลกฺขโณ สมุทโย, สมุทยปริชหนลกฺขโณ นิโรโธ, อนุสย\-สมุจฺเฉท\-ลกฺขโณ มคฺโค. พฺยาธิ\-ลกฺขณํ ทุกฺขํ, สญฺชานน\-ลกฺขโณ สมุทโย, นิยฺยานิก\-ลกฺขโณ มคฺโค, สนฺติลกฺขโณ นิโรโธ. อปฺปฏิสนฺธิ\-ภาว\-นิโรธ\-ลกฺขณา อนุปาทิเสสา นิพฺพานธาตุ, ทุกฺขญฺจ สมุทโย จ, ทุกฺขญฺจ นิโรโธ จ, ทุกฺขญฺจ มคฺโค จ, สมุทโย จ ทุกฺขญฺจ, สมุทโย จ นิโรโธ จ, สมุทโย จ มคฺโค จ, นิโรโธ จ สมุทโย จ, นิโรโธ จ ทุกฺขญฺจ, นิโรโธ จ มคฺโค จ, มคฺโค จ นิโรโธ จ, มคฺโค จ สมุทโย จ, มคฺโค จ ทุกฺขญฺจ.}

\proseitem{8}{ตตฺถิมานิ สุตฺตานิ.}

\setlength{\csromangathapagewidth}{0pt}
\setlength{\csromangathawakwidth}{0pt}
\csromangathameasuregroup{“ยเมกรตฺติํ ปฐมํ, \\ อพฺภุฏฺฐิโตว โส ยาติ,}{\csromangathabat{“ยเมกรตฺติํ ปฐมํ,}{คพฺเภ วสติ มาณโว.} \\ \csromangathabat{อพฺภุฏฺฐิโตว โส ยาติ,}{ส คจฺฉํ น นิวตฺตตี”ติ\textsuperscript{1}.}}
\csromangathasetleft{“ยเมกรตฺติํ ปฐมํ, \\ อพฺภุฏฺฐิโตว โส ยาติ,}
\csromangathagroup{\csromangathastanza{\csromangathaitembat{8}{“ยเมกรตฺติํ ปฐมํ,}{คพฺเภ วสติ มาณโว.} \\ \csromangathacontbat{อพฺภุฏฺฐิโตว โส ยาติ,}{ส คจฺฉํ น นิวตฺตตี”ติ\csromansharedfootnote{6e59fdedac9d9830}{ขุ 5 วีสตินิปาเต อโยฆรชาตเก.}.}}}

\prose{อฏฺฐิมา อานนฺท ทานุปปตฺติโย เอกุตฺตริเก สุตฺตํ, อยํ ชาติ.}

\prose{ตตฺถ กตมา ชรา.}

\setlength{\csromangathapagewidth}{0pt}
\setlength{\csromangathawakwidth}{0pt}
\csromangathameasuregroup{อจริตฺวา พฺรหฺมจริยํ, \\ ชิณฺณโกญฺจาว ฌายนฺติ, \\ ปญฺจ ปุพฺพนิมิตฺตานิ เทเวสุ,}{\csromangathabat{อจริตฺวา พฺรหฺมจริยํ,}{อลทฺธา โยพฺพเน ธนํ.} \\ \csromangathabat{ชิณฺณโกญฺจาว ฌายนฺติ,}{ขีณมจฺเฉว ปลฺลเล\textsuperscript{1}.} \\ \csromangathabat{ปญฺจ ปุพฺพนิมิตฺตานิ เทเวสุ,}{อยํ ชรา.}}
\csromangathasetleft{อจริตฺวา พฺรหฺมจริยํ, \\ ชิณฺณโกญฺจาว ฌายนฺติ, \\ ปญฺจ ปุพฺพนิมิตฺตานิ เทเวสุ,}
\csromangathagroup{\csromangathastanza{\csromangathabat{อจริตฺวา พฺรหฺมจริยํ,}{อลทฺธา โยพฺพเน ธนํ.} \\ \csromangathabat{ชิณฺณโกญฺจาว ฌายนฺติ,}{ขีณมจฺเฉว ปลฺลเล\csromansharedfootnote{338dd9d64fd7ee1e}{[มิสฺสินฺคฺ โนเต 0]}.} \\ \csromangathabat{ปญฺจ ปุพฺพนิมิตฺตานิ เทเวสุ,}{อยํ ชรา.}}}

\csromanhangingbegin
\hangingheadline{\csromanfolio{172}ตตฺถ กตโม พฺยาธิ.}
\hangingbody{สามํ เตน กุโต ราช, ตุวมฺปิ ชรายนฺติ เวเทสิ.}
\hangingbody{ขตฺติย กมฺมสฺส ผโล, โลโก น หิ กมฺมํ ปนยติ.}
\csromanhangingend

\setlength{\csromangathapagewidth}{0pt}
\setlength{\csromangathawakwidth}{0pt}
\csromangathameasuregroup{ตโย คิลานา,}{\csromangathabat{ตโย คิลานา,}{อยํ พฺยาธิ.}}
\csromangathasetleft{ตโย คิลานา,}
\csromangathagroup{\csromangathastanza{\csromangathabat{ตโย คิลานา,}{อยํ พฺยาธิ.}}}

\csromanhangingbegin
\hangingheadline{ตตฺถ กตมํ มรณํ.}
\hangingbody{ยถาปิ กุมฺภการสฺส, กตํ มตฺติกภาชนํ.}
\hangingbody{ขุทฺทกญฺจ มหนฺตญฺจ, ยํ ปกฺกํ ยญฺจ อามกํ.}
\hangingbody{สพฺพํ เภทนปริยนฺตํ, เอวํ มจฺจาน ชีวิตํ1.}
\hangingbody{มมายิเต ปสฺสถ ผนฺทมาเน2,}
\hangingbody{มจฺเฉว อปฺโปทเก ขีณโสเต.}
\hangingbody{เอตมฺปิ ทิสฺวา อมโม จเรยฺย,}
\hangingbody{ภเวสุ อาสตฺติมกุพฺพมาโน.}
\csromanhangingend

\setlength{\csromangathapagewidth}{0pt}
\setlength{\csromangathawakwidth}{0pt}
\csromangathameasuregroup{อุทกปฺปนสุตฺตํ,}{\csromangathabat{อุทกปฺปนสุตฺตํ,}{อิทํ มรณํ.}}
\csromangathasetleft{อุทกปฺปนสุตฺตํ,}
\csromangathagroup{\csromangathastanza{\csromangathabat{อุทกปฺปนสุตฺตํ,}{อิทํ มรณํ.}}}

\prose{ตตฺถ กตโม โสโก.}

\setlength{\csromangathapagewidth}{0pt}
\setlength{\csromangathawakwidth}{0pt}
\csromangathameasuregroup{}{อิธ โสจติ เปจฺจ โสจติ, \\ ปาปการี อุภยตฺถ โสจติ. \\ โส โสจติ โส วิหญฺญติ, \\ ทิสฺวา กมฺมกิลิฏฺฐมตฺตโม\textsuperscript{1}. \\ ตีณิ ทุจฺจริตานิ. \\ อยํ โสโก.}
\csromangathameasurewak{อิธ โสจติ เปจฺจ โสจติ, \\ ปาปการี อุภยตฺถ โสจติ. \\ โส โสจติ โส วิหญฺญติ, \\ ทิสฺวา กมฺมกิลิฏฺฐมตฺตโม\textsuperscript{1}. \\ ตีณิ ทุจฺจริตานิ. \\ อยํ โสโก.}
\setlength{\csromangathaleftcol}{0pt}
\csromangathagroup{\csromangathastanza{\csromangathawak{อิธ โสจติ เปจฺจ โสจติ, \\ ปาปการี อุภยตฺถ โสจติ. \\ โส โสจติ โส วิหญฺญติ, \\ ทิสฺวา กมฺมกิลิฏฺฐมตฺตโม\csromansharedfootnote{338dd9d64fd7ee1e}{[มิสฺสินฺคฺ โนเต 0]}. \\ ตีณิ ทุจฺจริตานิ. \\ อยํ โสโก.}}}

\prose{ตตฺถ กตโม ปริเทโว.}

\setlength{\csromangathapagewidth}{0pt}
\setlength{\csromangathawakwidth}{0pt}
\csromangathameasuregroup{}{กาเมสุ คิทฺธา ปสุตา ปมูฬฺหา, \\ อวทานิยา เต วิสเม นิวิฏฺฐา. \\ ทุกฺขูปนีตา ปริเทวยนฺติ, \\ กิํ สุ ภวิสฺสาม อิโต จุตาเส\textsuperscript{1}. \\ ติสฺโส วิปตฺติโย, \\ อยํ ปริเทโว.}
\csromangathameasurewak{กาเมสุ คิทฺธา ปสุตา ปมูฬฺหา, \\ อวทานิยา เต วิสเม นิวิฏฺฐา. \\ ทุกฺขูปนีตา ปริเทวยนฺติ, \\ กิํ สุ ภวิสฺสาม อิโต จุตาเส\textsuperscript{1}. \\ ติสฺโส วิปตฺติโย,}
\setlength{\csromangathaleftcol}{0pt}
\csromangathagroup{\csromangathastanza{\csromangathawak{กาเมสุ คิทฺธา ปสุตา ปมูฬฺหา, \\ อวทานิยา เต วิสเม นิวิฏฺฐา. \\ ทุกฺขูปนีตา ปริเทวยนฺติ, \\ กิํ สุ ภวิสฺสาม อิโต จุตาเส\csromansharedfootnote{338dd9d64fd7ee1e}{[มิสฺสินฺคฺ โนเต 0]}.}}\csromangathastanza{\csromangathawak{ติสฺโส วิปตฺติโย,}}\csromangathastanza{อยํ ปริเทโว.}}

\csromanhangingbegin
\hangingheadline{\csromanfolio{173}ตตฺถ กตมํ ทุกฺขํ.}
\hangingbody{สตํ อาสิ อโยสงฺกู1, สพฺเพ ปจฺจตฺตเวทนา.}
\hangingbody{ชลิตา ชาตเวทาว, อจฺจิสํฆสมากุลา.}
\csromanhangingend

\prose{มหา วต โส ปริฬาโห\csromansharedfootnote{0b7e3e1710df08a4}{ปริทาโฆ (อิ, ก) สํ 3. 393 ปิฏฺเฐ.} สํยุตฺตเก สุตฺตํ สจฺจ\-สํยุตฺเตสุ. อิทํ ทุกฺขํ.}

\prose{ตตฺถ กตมํ โทมนสฺสํ.}

\setlength{\csromangathapagewidth}{0pt}
\setlength{\csromangathawakwidth}{0pt}
\csromangathameasuregroup{สงฺกปฺเปหิ ปเรโต\textsuperscript{1} โส, \\ สุตฺวา ปเรสํ นิคฺโฆสํ,}{\csromangathabat{สงฺกปฺเปหิ ปเรโต\textsuperscript{1} โส,}{กปโณ วิย ฌายติ.} \\ \csromangathabat{สุตฺวา ปเรสํ นิคฺโฆสํ,}{มงฺกุ โหติ ตถาวิโธ.} \\ เทฺวเม ตปนียา ธมฺมา. \\ อิทํ โทมนสฺสํ.}
\csromangathameasurewak{เทฺวเม ตปนียา ธมฺมา. \\ อิทํ โทมนสฺสํ.}
\csromangathasetleft{สงฺกปฺเปหิ ปเรโต\textsuperscript{1} โส, \\ สุตฺวา ปเรสํ นิคฺโฆสํ,}
\csromangathagroup{\csromangathastanza{\csromangathabat{สงฺกปฺเปหิ ปเรโต\csromansharedfootnote{338dd9d64fd7ee1e}{[มิสฺสินฺคฺ โนเต 0]} โส,}{กปโณ วิย ฌายติ.} \\ \csromangathabat{สุตฺวา ปเรสํ นิคฺโฆสํ,}{มงฺกุ โหติ ตถาวิโธ.}}\csromangathastanza{\csromangathawak{เทฺวเม ตปนียา ธมฺมา. \\ อิทํ โทมนสฺสํ.}}}

\prose{ตตฺถ กตโม อุปายาโส.}

\setlength{\csromangathapagewidth}{0pt}
\setlength{\csromangathawakwidth}{0pt}
\csromangathameasuregroup{กมฺมารานํ ยถา อุกฺกา, \\ เอวํ ฑยฺหติ เม หทยํ,}{\csromangathabat{กมฺมารานํ ยถา อุกฺกา,}{อนฺโต ฑยฺหติ โน พหิ.} \\ \csromangathabat{เอวํ ฑยฺหติ เม หทยํ,}{สุตฺวา นิพฺพตฺตมมฺพุชํ.} \\ ตโย อคฺคี. \\ อยํ อุปายาโส.}
\csromangathameasurewak{ตโย อคฺคี. \\ อยํ อุปายาโส.}
\csromangathasetleft{กมฺมารานํ ยถา อุกฺกา, \\ เอวํ ฑยฺหติ เม หทยํ,}
\csromangathagroup{\csromangathastanza{\csromangathabat{กมฺมารานํ ยถา อุกฺกา,}{อนฺโต ฑยฺหติ โน พหิ.} \\ \csromangathabat{เอวํ ฑยฺหติ เม หทยํ,}{สุตฺวา นิพฺพตฺตมมฺพุชํ.}}\csromangathastanza{\csromangathawak{ตโย อคฺคี. \\ อยํ อุปายาโส.}}}

\csromanhangingbegin
\hangingheadline{ตตฺถ กตโม อปฺปิย\-สมฺปโยโค.}
\hangingbody{อยสาว มลํ สมุฏฺฐิตํ, ตตุฏฺฐาย ตเมว ขาทติ.}
\hangingbody{เอวํ อติโธนจารินํ, สานิ กมฺมานิ นยนฺติ ทุคฺคติํ4.}
\csromanhangingend

\prose{เทฺวเม ตถาคตํ อพฺภาจิกฺขนฺติ. เอกุตฺตริเก สุตฺตํ ทุเกสุ. อยํ อปฺปิย\-สมฺปโยโค.}

\csromanhangingbegin
\hangingheadline{ตตฺถ กตโม ปิยวิปฺปโยโค.}
\hangingbody{สุปิเนน ยถาปิ สงฺคตํ, ปฏิพุทฺโธ ปุริโส น ปสฺสติ.}
\hangingbody{เอวมฺปิ ปิยายิตํ5 ชนํ, เปตํ กาลงฺกตํ6 น ปสฺสติ.}
\csromanhangingend

\prose{เต เทวา จวนธมฺมํ วิทิตฺวา ตีหิ วาจาหิ อนุสาสนฺติ. อยํ ปิยวิปฺปโยโค.}

\csromanheader{2}{เปฏโกปเทส\-ปาฬิ ยมฺปิจฺฉํ น ลภติ, ติสฺโส มารธีตโร.}
\prose{\csromanfolio{174}\csromansymbolfootnote{*}{เหฏฺฐา 6, 57 ปิฏฺเฐสุ.}ตสฺส เจ กามยานสฺส\csromansharedfootnote{6389d285ed0ec3e2}{กามยมานสฺส (ก) ขุ 1. 399 ปิฏฺเฐ.}, ฉนฺทชาตสฺส ชนฺตุโน.}

\setlength{\csromangathapagewidth}{0pt}
\setlength{\csromangathawakwidth}{0pt}
\csromangathameasuregroup{เต กามา ปริหายนฺติ,}{\csromangathabat{เต กามา ปริหายนฺติ,}{สลฺลวิทฺโธว รุปฺปติ.} \\ สํขิตฺเตน ปญฺจุปาทานกฺขนฺธา ทุกฺขา. \\ จกฺขุ โสตญฺจ ฆานญฺจ, \\ ชิวฺหา กาโย ตโต มนํ. \\ เอเต โลกามิสา โฆรา, \\ ยตฺถ สตฺตา ปุถุชฺชนา. \\ ปญฺจิเม ภิกฺขเว ขนฺธา. \\ อิทํ ทุกฺขํ.}
\csromangathameasurewak{ชิวฺหา กาโย ตโต มนํ. \\ เอเต โลกามิสา โฆรา, \\ ยตฺถ สตฺตา ปุถุชฺชนา. \\ ปญฺจิเม ภิกฺขเว ขนฺธา. \\ อิทํ ทุกฺขํ.}
\csromangathasetleft{เต กามา ปริหายนฺติ,}
\csromangathagroup{\csromangathastanza{\csromangathabat{เต กามา ปริหายนฺติ,}{สลฺลวิทฺโธว รุปฺปติ.} \\ สํขิตฺเตน ปญฺจุ\-ปาทาน\-กฺขนฺธา ทุกฺขา. \\ จกฺขุ โสตญฺจ ฆานญฺจ,}\csromangathastanza{\csromangathawak{ชิวฺหา กาโย ตโต มนํ. \\ เอเต โลกามิสา โฆรา, \\ ยตฺถ สตฺตา ปุถุชฺชนา. \\ ปญฺจิเม ภิกฺขเว ขนฺธา.}}\csromangathastanza{\csromangathawak{อิทํ ทุกฺขํ.}}}

\prose{ตตฺถ กตมา ชรา จ มรณญฺจ.}

\setlength{\csromangathapagewidth}{0pt}
\setlength{\csromangathawakwidth}{0pt}
\csromangathameasuregroup{อปฺปํ วต ชีวิตํ อิทํ, \\ อถ วาปิ อกิจฺฉํ ชีวิตํ,}{\csromangathabat{อปฺปํ วต ชีวิตํ อิทํ,}{โอรํ วสฺสสตาปิ มียเต\textsuperscript{1}.} \\ \csromangathabat{อถ วาปิ อกิจฺฉํ ชีวิตํ,}{อถ โข โส ชรสาปิ มียเต.}}
\csromangathasetleft{อปฺปํ วต ชีวิตํ อิทํ, \\ อถ วาปิ อกิจฺฉํ ชีวิตํ,}
\csromangathagroup{\csromangathastanza{\csromangathabat{อปฺปํ วต ชีวิตํ อิทํ,}{โอรํ วสฺสสตาปิ มียเต\csromansharedfootnote{c8eb941a735e4fa9}{มียติ (ขุ 1. 405 ปิฏฺเฐ)}.} \\ \csromangathabat{อถ วาปิ อกิจฺฉํ ชีวิตํ,}{อถ โข โส ชรสาปิ มียเต.}}}

\prose{สํยุตฺตเก ปเสนทิ\-สํยุตฺตเก สุตฺตํ อยฺยิกา เม กาลงฺกตา. อยํ ชรา จ มรณญฺจ.}

\prose{ตตฺถ กตมา จุติ จ อุปปตฺติ จ.}

\prose{** สพฺเพ สตฺตา มริสฺสนฺติ, มรณนฺตํ หิ ชีวิตํ.}

\setlength{\csromangathapagewidth}{0pt}
\setlength{\csromangathawakwidth}{0pt}
\csromangathameasuregroup{ยถากมฺมํ คมิสฺสนฺติ,}{\csromangathabat{ยถากมฺมํ คมิสฺสนฺติ,}{อตฺตกมฺมผลูปคาติ\textsuperscript{1}.} \\ อยํ จุติ จ อุปปตฺติ จ.}
\csromangathasetleft{ยถากมฺมํ คมิสฺสนฺติ,}
\csromangathagroup{\csromangathastanza{\csromangathabat{ยถากมฺมํ คมิสฺสนฺติ,}{อตฺตกมฺมผลูปคาติ\csromansharedfootnote{338dd9d64fd7ee1e}{[มิสฺสินฺคฺ โนเต 0]}.} \\ อยํ จุติ จ อุปปตฺติ จ.}}

\prose{อิเมหิ สุตฺเตหิ เอกสทิเสหิ จ อญฺเญหิ นววิธํ สุตฺตํ ตํ อนุปวิฏฺเฐหิ ลกฺขณโต ทุกฺขํ ญตฺวา สาธารณญฺจ อสาธารณญฺจ ทุกฺขํ อริยสจฺจํ นิทฺทิสิตพฺพํ. คาถาหิ คาถา อนุมินิตพฺพา, พฺยากรเณหิ วา พฺยากรณํ. อิทํ ทุกฺขํ.}

\proseitem{9}{ตตฺถ กตโม ทุกฺขสมุทโย.}

\prose{\csromansymbolfootnote{+}{อุปริ 203 ปิฏฺเฐ.}กาเมสุ สตฺตา กามสงฺคสตฺตา\csromansharedfootnote{103906e7309b6307}{กามปสงฺคสตฺตา (อิ) ขุ 1. 171 ปิฏฺเฐ. ** เหฏฺฐา 79 ปิฏฺเฐ.},}

\prose{สํโยชเน วชฺชม\-ปสฺสมานา.}

\prose{น หิ ชาตุ สํโยชนสงฺค\-สตฺตา,}

\prose{โอฆํ ตเรยฺยุํ วิปุลํ มหนฺตํ.}

\prose{จตฺตาโร อาสวา สุตฺตํ. อยํ ทุกฺขสมุทโย.}

\csromanhangingbegin
\hangingheadline{\csromanfolio{175}ตตฺถ กตโม ทุกฺขนิโรโธ.}
\hangingbody{ยมฺหิ น มายา วสตี น มาโน,}
\hangingbody{โย วีตโลโภ อมโม นิราโส,}
\hangingbody{ปนุณฺณโกโธ1 อภินิพฺพุตตฺโต.}
\hangingbody{โส พฺราหฺมโณ โส สมโณ ส ภิกฺขุ.}
\csromanhangingend

\prose{เทฺวมา วิมุตฺติโย, ราควิราคา จ เจโตวิมุตฺติ, อวิชฺชาวิราคา จ ปญฺญาวิมุตฺติ, อยํ นิโรโธ.}

\prose{ตตฺถ กตโม มคฺโค.}

\setlength{\csromangathapagewidth}{0pt}
\setlength{\csromangathawakwidth}{0pt}
\csromangathameasuregroup{เอเสว มคฺโค นตฺถญฺโญ, \\ อริโย อฏฺฐงฺคิโก มคฺโค, \\ สตฺติเม ภิกฺขเว โพชฺฌงฺคา,}{\csromangathabat{เอเสว มคฺโค นตฺถญฺโญ,}{ทสฺสนสฺส วิสุทฺธิยา.} \\ \csromangathabat{อริโย อฏฺฐงฺคิโก มคฺโค,}{มารสฺเสตํ ปโมหนํ.} \\ \csromangathabat{สตฺติเม ภิกฺขเว โพชฺฌงฺคา,}{อยํ มคฺโค.}}
\csromangathasetleft{เอเสว มคฺโค นตฺถญฺโญ, \\ อริโย อฏฺฐงฺคิโก มคฺโค, \\ สตฺติเม ภิกฺขเว โพชฺฌงฺคา,}
\csromangathagroup{\csromangathastanza{\csromangathabat{เอเสว มคฺโค นตฺถญฺโญ,}{ทสฺสนสฺส วิสุทฺธิยา.} \\ \csromangathabat{อริโย อฏฺฐงฺคิโก มคฺโค,}{มารสฺเสตํ ปโมหนํ.} \\ \csromangathabat{สตฺติเม ภิกฺขเว โพชฺฌงฺคา,}{อยํ มคฺโค.}}}

\csromanhangingbegin
\hangingheadline{ตตฺถ กตมานิ จตฺตาริ อริยสจฺจานิ.}
\hangingbody{“เย ธมฺมา2 เหตุปฺปภวา, เตสํ เหตุํ ตถาคโต อาห.}
\hangingbody{เตสญฺจ โย นิโรโธ, เอวํวาที มหาสมโณ”ติ.}
\csromanhangingend

\prose{เหตุปฺปภวา ธมฺมา ทุกฺขํ, เหตุสมุทโย, ยํ ภควโต วจนํ. อยํ ธมฺโม โย นิโรโธ, เย หิ เกจิ สํโยชนิเยสุ ธมฺเมสุ อสฺสทานุปสฺสิโน วิหรนฺติ. กิเลสา ตณฺหา ปวฑฺฒติ, ตณฺหาปจฺจยา อุปาทานํ ฯเปฯ เอวเมตสฺส เกวลสฺส ทุกฺข\-กฺขนฺธสฺส สมุทโย โหติ. ตตฺถ ยํ สํโยชนํ, อยํ สมุทโย. เย สํโยชนิยา ธมฺมา เย จ โสก\-ปริเทว\-ทุกฺข\-โทมนสฺสุ\-ปายาสา สมฺภวนฺติ, อิทํ ทุกฺขํ. ยา สํโยชนิเยสุ ธมฺเมสุ อาทีนวา\-นุปสฺสนา, อยํ มคฺโค. ปริมุจฺจติ ชาติยา ชราย พฺยาธีหิ มรเณหิ โสเกหิ ปริเทเวหิ ยาว อุปายาเสหิ, อิทํ นิพฺพานํ. อิมานิ จตฺตาริ สจฺจานิ.}

\prose{ตตฺถ กตมา อนุปาทิเสสา นิพฺพานธาตุ.}

\setlength{\csromangathapagewidth}{0pt}
\setlength{\csromangathawakwidth}{0pt}
\csromangathameasuregroup{}{อตฺถงฺคตสฺส นปมาณมตฺถิ, \\ ตํ หิ วา นตฺถิ เยน นํ ปญฺญเปยฺย. \\ สพฺพสงฺคานํ สมูหตตฺตา วิทู, \\ สิตา วาทสต’สฺสุ\textsuperscript{1} สพฺเพ. \\ สํยุตฺตเก โคธิกสํยุตฺตํ.}
\csromangathameasurewak{อตฺถงฺคตสฺส นปมาณมตฺถิ, \\ ตํ หิ วา นตฺถิ เยน นํ ปญฺญเปยฺย. \\ สพฺพสงฺคานํ สมูหตตฺตา วิทู, \\ สิตา วาทสต’สฺสุ\textsuperscript{1} สพฺเพ. \\ สํยุตฺตเก โคธิกสํยุตฺตํ.}
\setlength{\csromangathaleftcol}{0pt}
\csromangathagroup{\csromangathastanza{\csromangathawak{อตฺถงฺคตสฺส นปมาณมตฺถิ, \\ ตํ หิ วา นตฺถิ เยน นํ ปญฺญเปยฺย. \\ สพฺพสงฺคานํ สมูหตตฺตา วิทู, \\ สิตา วาทสต’สฺสุ\csromansharedfootnote{338dd9d64fd7ee1e}{[มิสฺสินฺคฺ โนเต 0]} สพฺเพ.}}\csromangathastanza{\csromangathawak{สํยุตฺตเก โคธิกสํยุตฺตํ.}}}

\prose{\csromanfolio{176}อิมานิ อสาธารณานิ สุตฺตานิ. ยหิํ ยหิํ สจฺจานิ นิทฺทิฏฺฐานิ, ตหิํ ตหิํ สจฺจ\-ลกฺขณโต โอตาเรตฺวา\csromansharedfootnote{c24c76a4a05cfe92}{โอหาเรตฺวา (อิ, ก)} อปริมาเณหิ พฺยญฺชเนหิ โส อตฺโถ ปริเยสิตพฺโพ. ตตฺถ อตฺถา\-นุปริวตฺติ พฺยญฺชเนน ปุน พฺยญฺชนา\-นุปริวตฺติ อตฺเถน ตสฺส เอกเมกสฺส อปริมาณานิ พฺยญฺชนานิ อิเมหิ สุตฺเตหิ ยถา\-นิกฺขิตฺเตหิ จตฺตาริ อริยสจฺจานิ นิทฺทิสิตพฺพานิ. ปญฺจนิกาเย อนุปวิฏฺฐาหิ คาถาหิ คาถา อนุมินิตพฺพา, พฺยากรเณน พฺยากรณํ. อิมานิ อสาธารณานิ สุตฺตานิ.}

\csromanheader{2}{\textbf{เตสํ อิมา อุทฺทานคาถา}}
\setlength{\csromangathapagewidth}{0pt}
\setlength{\csromangathawakwidth}{0pt}
\csromangathameasuregroup{ยเมกรตฺติํ ปฐมํ, \\ ปญฺจ ปุพฺพนิมิตฺตานิ, \\ สามํ เตน กุโต ราช, \\ ยถาปิ กุมฺภการสฺส, \\ อิธ โสจติ เปจฺจ โสจติ, \\ กาเมสุ คิทฺธา ปสุตา, \\ สตํ อาสิ\textsuperscript{1} อโยสงฺกู, \\ สงฺกปฺเปหิ ปเรโต โส, \\ กมฺมารานํ ยถา อุกฺกา, \\ อยโต มลมุปฺปนฺนํ, \\ ติวิธํ เทวานุสาสนฺติ, \\ ติสฺโส เจว มารธีตา, \\ จกฺขุ โสตญฺจ ฆานญฺจ, \\ อปฺปํ วต ชีวิตํ อิทํ, \\ สพฺเพ สตฺตา มริสฺสนฺติ, \\ กาเมสุ สตฺตา ปสุตา, \\ ยมฺหิ น มายา วสติ, \\ เอเสว มคฺโค นตฺถญฺโญ, \\ อตฺถงฺคตสฺส น ปมาณมตฺถิ, \\ เย ธมฺมา เหตุปฺปภวา,}{\csromangathabat{ยเมกรตฺติํ ปฐมํ,}{อฏฺฐ ทานูปปตฺติโย.} \\ \csromangathabat{ปญฺจ ปุพฺพนิมิตฺตานิ,}{ขีณมจฺฉํว ปลฺลลํ.} \\ \csromangathabat{สามํ เตน กุโต ราช,}{ตโย เทวา คิลานกา.} \\ \csromangathabat{ยถาปิ กุมฺภการสฺส,}{ยถา นทิทกปฺปนํ.} \\ \csromangathabat{อิธ โสจติ เปจฺจ โสจติ,}{ตีณิ ทุจฺจริตานิ จ.} \\ \csromangathabat{กาเมสุ คิทฺธา ปสุตา,}{ยาว ติสฺโส วิปตฺติโย.} \\ \csromangathabat{สตํ อาสิ\textsuperscript{1} อโยสงฺกู,}{ปริฬาโห มหตฺตโร.} \\ \csromangathabat{สงฺกปฺเปหิ ปเรโต โส,}{ตตฺถ ตปนิเยหิ จ.} \\ \csromangathabat{กมฺมารานํ ยถา อุกฺกา,}{ตโย อคฺคี ปกาสิตา.} \\ \csromangathabat{อยโต มลมุปฺปนฺนํ,}{อพฺภกฺขานํ ตถาคเต.} \\ \csromangathabat{ติวิธํ เทวานุสาสนฺติ,}{สุปิเนน สงฺคโม ยถา.} \\ \csromangathabat{ติสฺโส เจว มารธีตา,}{สลฺลวิทฺโธว รุปฺปติ.} \\ \csromangathabat{จกฺขุ โสตญฺจ ฆานญฺจ,}{ปญฺจกฺขนฺธา ปกาสิตา.} \\ \csromangathabat{อปฺปํ วต ชีวิตํ อิทํ,}{อยฺยิกา เม มหลฺลิกา.} \\ \csromangathabat{สพฺเพ สตฺตา มริสฺสนฺติ,}{อุปปตฺติ จุติจยํ.} \\ \csromangathabat{กาเมสุ สตฺตา ปสุตา,}{อาสเวหิ จตูหิ จ.} \\ \csromangathabat{ยมฺหิ น มายา วสติ,}{เทฺวมา เจโตวิมุตฺติโย.} \\ \csromangathabat{เอเสว มคฺโค นตฺถญฺโญ,}{โพชฺฌงฺคา จ สุเทสิตา.} \\ \csromangathabat{อตฺถงฺคตสฺส น ปมาณมตฺถิ,}{โคธิโก ปรินิพฺพุโต.} \\ \csromangathabat{เย ธมฺมา เหตุปฺปภวา,}{สํโยชนานุปสฺสิโน.}}
\csromangathasetleft{ยเมกรตฺติํ ปฐมํ, \\ ปญฺจ ปุพฺพนิมิตฺตานิ, \\ สามํ เตน กุโต ราช, \\ ยถาปิ กุมฺภการสฺส, \\ อิธ โสจติ เปจฺจ โสจติ, \\ กาเมสุ คิทฺธา ปสุตา, \\ สตํ อาสิ\textsuperscript{1} อโยสงฺกู, \\ สงฺกปฺเปหิ ปเรโต โส, \\ กมฺมารานํ ยถา อุกฺกา, \\ อยโต มลมุปฺปนฺนํ, \\ ติวิธํ เทวานุสาสนฺติ, \\ ติสฺโส เจว มารธีตา, \\ จกฺขุ โสตญฺจ ฆานญฺจ, \\ อปฺปํ วต ชีวิตํ อิทํ, \\ สพฺเพ สตฺตา มริสฺสนฺติ, \\ กาเมสุ สตฺตา ปสุตา, \\ ยมฺหิ น มายา วสติ, \\ เอเสว มคฺโค นตฺถญฺโญ, \\ อตฺถงฺคตสฺส น ปมาณมตฺถิ, \\ เย ธมฺมา เหตุปฺปภวา,}
\csromangathagroup{\csromangathastanza{\csromangathabat{ยเมกรตฺติํ ปฐมํ,}{อฏฺฐ ทานูปปตฺติโย.} \\ \csromangathabat{ปญฺจ ปุพฺพนิมิตฺตานิ,}{ขีณมจฺฉํว ปลฺลลํ.}}\csromangathastanza{\csromangathabat{สามํ เตน กุโต ราช,}{ตโย เทวา คิลานกา.} \\ \csromangathabat{ยถาปิ กุมฺภการสฺส,}{ยถา นทิทกปฺปนํ.}}\csromangathastanza{\csromangathabat{อิธ โสจติ เปจฺจ โสจติ,}{ตีณิ ทุจฺจริตานิ จ.} \\ \csromangathabat{กาเมสุ คิทฺธา ปสุตา,}{ยาว ติสฺโส วิปตฺติโย.}}\csromangathastanza{\csromangathabat{สตํ อาสิ\csromansharedfootnote{fad696375e303f51}{สตมายุ (สี), สตธาตุ (อิ)} อโยสงฺกู,}{ปริฬาโห มหตฺตโร.} \\ \csromangathabat{สงฺกปฺเปหิ ปเรโต โส,}{ตตฺถ ตปนิเยหิ จ.}}\csromangathastanza{\csromangathabat{กมฺมารานํ ยถา อุกฺกา,}{ตโย อคฺคี ปกาสิตา.} \\ \csromangathabat{อยโต มลมุปฺปนฺนํ,}{อพฺภกฺขานํ ตถาคเต.}}\csromangathastanza{\csromangathabat{ติวิธํ เทวานุสาสนฺติ,}{สุปิเนน สงฺคโม ยถา.} \\ \csromangathabat{ติสฺโส เจว มารธีตา,}{สลฺลวิทฺโธว รุปฺปติ.}}\csromangathastanza{\csromangathabat{จกฺขุ โสตญฺจ ฆานญฺจ,}{ปญฺจกฺขนฺธา ปกาสิตา.} \\ \csromangathabat{อปฺปํ วต ชีวิตํ อิทํ,}{อยฺยิกา เม มหลฺลิกา.}}\csromangathastanza{\csromangathabat{สพฺเพ สตฺตา มริสฺสนฺติ,}{อุปปตฺติ จุติจยํ.} \\ \csromangathabat{กาเมสุ สตฺตา ปสุตา,}{อาสเวหิ จตูหิ จ.}}\csromangathastanza{\csromangathabat{ยมฺหิ น มายา วสติ,}{เทฺวมา เจโตวิมุตฺติโย.} \\ \csromangathabat{เอเสว มคฺโค นตฺถญฺโญ,}{โพชฺฌงฺคา จ สุเทสิตา.}}\csromangathastanza{\csromanfolio{177}\csromangathabat{อตฺถงฺคตสฺส น ปมาณมตฺถิ,}{โคธิโก ปรินิพฺพุโต.} \\ \csromangathabat{เย ธมฺมา เหตุปฺปภวา,}{สํโยชนานุปสฺสิโน.}}}

\csromanheader{2}{อิมา ทส เตสํ อุทฺทานคาถา.}
\proseitem{10}{ตตฺถิมานิ สาธารณานิ สุตฺตานิ เยสุ สุตฺเตสุ สาธารณานิ สจฺจานิ เทสิตานิ อนุโลมมฺปิ ปฏิโลมมฺปิ โวมิสฺสกมฺปิ. ตตฺถ อยํ อาทิ.}

\prose{อวิชฺชาย นิวุโต โลโก, (อชิตาติ ภควา,)}

\prose{วิวิจฺฉา ปมาทา นปฺปกาสติ.}

\setlength{\csromangathapagewidth}{0pt}
\setlength{\csromangathawakwidth}{0pt}
\csromangathameasuregroup{ชปฺปาภิเลปนํ\textsuperscript{1} พฺรูมิ,}{\csromangathabat{ชปฺปาภิเลปนํ\textsuperscript{1} พฺรูมิ,}{ทุกฺขมสฺส มหพฺภยํ.}}
\csromangathasetleft{ชปฺปาภิเลปนํ\textsuperscript{1} พฺรูมิ,}
\csromangathagroup{\csromangathastanza{\csromangathabat{ชปฺปา\-ภิเลปนํ\csromansharedfootnote{a3a51af6b426dd57}{ชปฺปานุเลปนํ (ข); ขุ 1. 434; เหฏฺฐา 11, 60 ปิฏฺเฐสุ.} พฺรูมิ,}{ทุกฺขมสฺส มหพฺภยํ.}}}

\prose{ตตฺถ ยา อวิชฺชา จ วิวิจฺฉา จ, อยํ สมุทโย. ยํ มหพฺภยํ, อิทํ ทุกฺขํ. อิมานิ เทฺว สจฺจานิ ทุกฺขญฺจ สมุทโย จ. “สํโยชนํ สํโยชนิยา จ ธมฺมา”ติ สํยุตฺตเก จิตฺต\-สํยุตฺตเกสุ พฺยากรณํ. ตตฺถ ยํ สํโยชนํ, อยํ สมุทโย. เย สํโยชนิยา ธมฺมา, อิทํ ทุกฺขํ. อิมานิ เทฺว สจฺจานิ ทุกฺขญฺจ สมุทโย จ.}

\setlength{\csromangathapagewidth}{0pt}
\setlength{\csromangathawakwidth}{0pt}
\csromangathameasuregroup{}{ตตฺถ กตมํ ทุกฺขญฺจ นิโรโธ จ. \\ อุจฺฉินฺนภวตณฺหสฺส, เนตฺติจฺฉินฺนสฺส\textsuperscript{1} ภิกฺขุโน. \\ วิกฺขีโณ ชาติสํสาโร, นตฺถิ ทานิ ปุนพฺภโว.}
\csromangathameasurewak{ตตฺถ กตมํ ทุกฺขญฺจ นิโรโธ จ. \\ อุจฺฉินฺนภวตณฺหสฺส, เนตฺติจฺฉินฺนสฺส\textsuperscript{1} ภิกฺขุโน. \\ วิกฺขีโณ ชาติสํสาโร, นตฺถิ ทานิ ปุนพฺภโว.}
\setlength{\csromangathaleftcol}{0pt}
\csromangathagroup{\csromangathastanza{\csromangathawak{\csromangathaitemline{10}{ตตฺถ กตมํ ทุกฺขญฺจ นิโรโธ จ.} \\ \csromangathacontline{อุจฺฉินฺน\-ภวตณฺหสฺส, เนตฺติจฺฉินฺนสฺส\csromansharedfootnote{6cf37fa05a019d2a}{สนฺตจิตฺตสฺส (ขุ 1. 394 ปิฏฺเฐ.)} ภิกฺขุโน.} \\ \csromangathacontline{วิกฺขีโณ ชาติสํสาโร, นตฺถิ ทานิ ปุนพฺภโว.}}}}

\prose{ยํ จิตฺตํ, อิทํ ทุกฺขํ. โย ภวตณฺหาย อุปจฺเฉโท, อยํ ทุกฺขนิโรโธ. วิกฺขีโณ ชาติสํสาโร, นตฺถิ ทานิ ปุนพฺภโวติ นิทฺเทโส. อิมานิ เทฺว สจฺจานิ ทุกฺขญฺจ นิโรโธ จ. เทฺวมา ภิกฺขเว วิมุตฺติโย ราควิราคา จ เจโตวิมุตฺติ, อวิชฺชาวิราคา จ ปญฺญาวิมุตฺติ. ยํ จิตฺตํ, อิทํ ทุกฺขํ. ยา วิมุตฺติ, อยํ นิโรโธ. อิมานิ เทฺว สจฺจานิ ทุกฺขญฺจ นิโรโธ จ.}

\setlength{\csromangathapagewidth}{0pt}
\setlength{\csromangathawakwidth}{0pt}
\csromangathameasuregroup{}{ตตฺถ กตมํ ทุกฺขญฺจ มคฺโค จ. \\ กุมฺภูปมํ\textsuperscript{1} กายมิมํ วิทิตฺวา, \\ นครูปมํ จิตฺตมิทํ ฐเปตฺวา. \\ โยเธถ มารํ ปญฺญาวุเธน, \\ ชิตญฺจ รกฺเข อนิเวสโน สิยา.}
\csromangathameasurewak{ตตฺถ กตมํ ทุกฺขญฺจ มคฺโค จ. \\ กุมฺภูปมํ\textsuperscript{1} กายมิมํ วิทิตฺวา, \\ นครูปมํ จิตฺตมิทํ ฐเปตฺวา. \\ โยเธถ มารํ ปญฺญาวุเธน, \\ ชิตญฺจ รกฺเข อนิเวสโน สิยา.}
\setlength{\csromangathaleftcol}{0pt}
\csromangathagroup{\csromangathastanza{\csromangathawak{\csromangathaitemline{10}{ตตฺถ กตมํ ทุกฺขญฺจ มคฺโค จ.} \\ \csromangathacontline{กุมฺภูปมํ\csromansharedfootnote{ccb32d36fd533087}{ขุ 1. 18 ปิฏฺเฐ.} กายมิมํ วิทิตฺวา,} \\ \csromangathacontline{นครูปมํ จิตฺตมิทํ ฐเปตฺวา.} \\ \csromangathacontline{โยเธถ มารํ ปญฺญาวุเธน,} \\ \csromangathacontline{ชิตญฺจ รกฺเข อนิเวสโน สิยา.}}}}

\prose{\csromanfolio{178}ตตฺถ ยญฺจ กุมฺภูปโม กาโย ยญฺจ นครูปมํ จิตฺตํ, อิทํ ทุกฺขํ. ยํ ปญฺญาวุเธน มารํ โยเธถาติ อยํ มคฺโค. อิมานิ เทฺว สจฺจานิ. ยํ ภิกฺขเว น ตุมฺหากํ, ตํ ปชหิตพฺพํ. ยา สํโยชนา, อยํ มคฺโค. เย เต ธมฺมา อนตฺตนิยา ปหาตพฺพา, รูปํ ยาว วิญฺญาณํ, อิทํ ทุกฺขญฺจ มคฺโค จ.}

\prose{ตตฺถ กตมํ ทุกฺขญฺจ สมุทโย จ นิโรโธ จ.}

\prose{\csromansymbolfootnote{*}{เหฏฺฐา 57 ปิฏฺเฐ.}เย เกจิ โสกา ปริเทวิตา วา,}

\prose{ทุกฺขา จ\csromansharedfootnote{92ffbf33eedfd58f}{ทุกฺขญฺจ (อิ, ก) ขุ 1. 191 ปิฏฺเฐ.} โลกสฺมิม\-เนกรูปา.}

\setlength{\csromangathapagewidth}{0pt}
\setlength{\csromangathawakwidth}{0pt}
\csromangathameasuregroup{}{ปิยํ ปฏิจฺจปฺปภวนฺติ เอเต, \\ ปิเย อสนฺเต น ภวนฺติ เอเต.}
\csromangathameasurewak{ปิยํ ปฏิจฺจปฺปภวนฺติ เอเต, \\ ปิเย อสนฺเต น ภวนฺติ เอเต.}
\setlength{\csromangathaleftcol}{0pt}
\csromangathagroup{\csromangathastanza{\csromangathawak{ปิยํ ปฏิจฺจ\-ปฺปภวนฺติ เอเต, \\ ปิเย อสนฺเต น ภวนฺติ เอเต.}}}

\prose{เย โสกปริเทวา, ยํ จ อเนกรูปํ ทุกฺขํ, ยํ เปมโต ภวติ, อิทํ ทุกฺขํ. ยํ เปมํ, อยํ สมุทโย. โย ตตฺถ ฉนฺทราค\-วินโย ปิยสฺส อกิริยา, อยํ นิโรโธ. อิมานิ ตีณิ สจฺจานิ. ติมฺพรุโก ปริพฺพาชโก ปจฺเจติ “สยํกตํ ปรํกตนฺ”ติ. ยเถสา วีมํสา, อิทํ ทุกฺขํ. ยา เอเต เทฺว อนฺเต อนุปคมฺม มชฺฌิมา ปฏิปทา อวิชฺชาปจฺจยา สงฺขารา ยาว ชาติปจฺจยา ชรามรณํ, อิทมฺปิ ทุกฺขญฺจ สมุทโย จ. วิญฺญาณํ นามรูปํ สฬายตนํ ผสฺโส เวทนา ภโว ชาติ ชารามรณํ, อิทํ ทุกฺขํ. อวิชฺชา สงฺขารา ตณฺหา อุปาทานํ, อยํ สมุทโย. อิติ อิทํ สยํกตํ วีมํเสยฺยาติ\csromansharedfootnote{24c0110e3d77ec98}{วีมํสียติ (อิ, ก) ** เหฏฺฐา 52 ปิฏฺเฐ.} ยญฺจ ปฏิจฺจ\-สมุปฺปาเท ทุกฺขํ, อิทํ เอโส สมุทโย นิทฺทิฏฺโฐ. อวิชฺชานิโรธา สงฺขาร\-นิโรโธ จ ยาว จ ชรามรณ\-นิโรโธติ อยํ นิโรโธ. อิมานิ ตีณิ สจฺจานิ ทุกฺขญฺจ สมุทโย จ นิโรโธ จ.}

\proseitem{11}{ตตฺถ กตมํ ทุกฺขญฺจ สมุทโย จ มคฺโค จ.}

\prose{** “โย ทุกฺขมทฺทกฺขิ ยโตนิทานํ,}

\prose{กาเมสุ โส ชนฺตุ กถํ นเมยฺย.}

\setlength{\csromangathapagewidth}{0pt}
\setlength{\csromangathawakwidth}{0pt}
\csromangathameasuregroup{}{กามา หิ โลเก สงฺคาติ ญตฺวา, \\ เตสํ สตีมา วินยาย สิกฺเข”ติ\textsuperscript{1}.}
\csromangathameasurewak{กามา หิ โลเก สงฺคาติ ญตฺวา, \\ เตสํ สตีมา วินยาย สิกฺเข”ติ\textsuperscript{1}.}
\setlength{\csromangathaleftcol}{0pt}
\csromangathagroup{\csromangathastanza{\csromangathawak{กามา หิ โลเก สงฺคาติ ญตฺวา, \\ เตสํ สตีมา วินยาย สิกฺเข”ติ\csromansharedfootnote{45bda9c7bcfa3be7}{สํ 1. 120 ปิฏฺเฐ.}.}}}

\prose{\csromanfolio{179}โย ทุกฺขมทฺทกฺขิ, อิทํ ทุกฺขํ. ยโต ภวติ, อยํ สมุทโย. สนฺทิฏฺฐํ ยโต ภวติ ยาว ตสฺส วินยาย สิกฺขา, อยํ มคฺโค. อิมานิ ตีณิ สจฺจานิ. เอกาทส\-งฺคุตฺตเรสุ โคปาลโก\-ปม\-สุตฺตํ.}

\prose{ตตฺถ ยาว รูปสญฺญุตฺตา ยญฺจ สฬายตนํ ยถา วณํ ปฏิจฺฉาเทติ ยญฺจ ติตฺถํ ยถา จ ลภติ ธมฺมู\-ปสญฺหิตํ อุฬารํ ปีติปาโมชฺชํ จตุพฺพิธํ จ อตฺตภาวโต จ วตฺถุ, อิทํ ทุกฺขํ. ยาว อาสาฏิกํ หาเรตา\csromansharedfootnote{75670da3a128227a}{สาเฏตา (สี, อิ) อํ 3. 546 ปิฏฺเฐ.} โหติ, อยํ สมุทโย. รูปสญฺญุตฺตา อาสาฏกหรณํ\csromansharedfootnote{9ec8fe8f2b297794}{อาสาฏิกสาฏนา (อิ)} วณปฏิจฺฉาทนํ วีถิญฺญุตา โคจรกุสลญฺจ, อยํ มคฺโค. อวเสสา ธมฺมา อตฺถิ เหตู อตฺถิ ปจฺจยา อตฺถิ นิสฺสยา สวเสสโทหิตา อเนกปูชา จ กลฺยาณมิตฺต\-ตปฺปจฺจยา ธมฺมา วีถิญฺญุตา จ เหตุ, อิมานิ ตีณิ สจฺจานิ.}

\prose{ตตฺถ กตมํ ทุกฺขญฺจ มคฺโค จ นิโรโธ จ.}

\setlength{\csromangathapagewidth}{0pt}
\setlength{\csromangathawakwidth}{0pt}
\csromangathameasuregroup{}{สติ กายคตา อุปฏฺฐิตา, \\ ฉสุ ผสฺสายตเนสุ สํวุโต\textsuperscript{1}. \\ สตตํ ภิกฺขุ สมาหิโต, \\ ชญฺญา\textsuperscript{1} นิพฺพานมตฺตโน.}
\csromangathameasurewak{สติ กายคตา อุปฏฺฐิตา, \\ ฉสุ ผสฺสายตเนสุ สํวุโต\textsuperscript{1}. \\ สตตํ ภิกฺขุ สมาหิโต, \\ ชญฺญา\textsuperscript{1} นิพฺพานมตฺตโน.}
\setlength{\csromangathaleftcol}{0pt}
\csromangathagroup{\csromangathastanza{\csromangathawak{สติ กายคตา อุปฏฺฐิตา, \\ ฉสุ ผสฺสายตเนสุ สํวุโต\csromansharedfootnote{6b896b42decf9bd0}{สํวโร (อิ, ก) ขุ 1. 108 ปิฏฺเฐ.}. \\ สตตํ ภิกฺขุ สมาหิโต, \\ ชญฺญา\csromansharedfootnote{c2e10618c78f0a05}{ชาเนยฺย (อิ, ก)} นิพฺพานมตฺตโน.}}}

\prose{ตตฺถ ยา จ กายคตา สติ ยญฺจ สฬายตนํ ยตฺถ สพฺพญฺเจตํ ทุกฺขํ. ยา จ กายคตา สติ โย จ สีลสํวโร โย จ สมาธิ ยตฺถ ยา สติ, อยํ ปญฺญากฺขนฺโธ. สพฺพมฺปิ สีลกฺขนฺโธ สมาธิ\-กฺขนฺโธ, อยํ มคฺโค. เอวํวิหารินา ญาตพฺพํ นิพฺพานํ. อยํ นิโรโธ, อิมานิ ตีณิ สจฺจานิ.  สีเล ปติฏฺฐาย เทฺว ธมฺมา ภาเวตพฺพา สมโถ จ วิปสฺสนา จ. ตตฺถ ยํ จิตฺตสหชาตา ธมฺมา, อิทํ ทุกฺขํ. โย จ สมโถ ยา จ วิปสฺสนา, อยํ มคฺโค. ราควิราคา จ เจโตวิมุตฺติ, อวิชฺชาวิราคา จ ปญฺญาวิมุตฺติ, อยํ นิโรโธ. อิมานิ ตีณิ สจฺจานิ.}

\prose{\csromanfolio{180}ตตฺถ กตโม สมุทโย จ นิโรโธ จ.}

\setlength{\csromangathapagewidth}{0pt}
\setlength{\csromangathawakwidth}{0pt}
\csromangathameasuregroup{}{อาสา จ ปีหา อภินนฺทนา จ, \\ อเนกธาตูสุ สรา ปติฏฺฐิตา. \\ อญฺญาณมูลปฺปภวา ปชปฺปิตา, \\ สพฺพา มยา พฺยนฺติกตา สมูลิกา.}
\csromangathameasurewak{อาสา จ ปีหา อภินนฺทนา จ, \\ อเนกธาตูสุ สรา ปติฏฺฐิตา. \\ อญฺญาณมูลปฺปภวา ปชปฺปิตา, \\ สพฺพา มยา พฺยนฺติกตา สมูลิกา.}
\setlength{\csromangathaleftcol}{0pt}
\csromangathagroup{\csromangathastanza{\csromangathawak{อาสา จ ปีหา อภินนฺทนา จ, \\ อเนกธาตูสุ สรา ปติฏฺฐิตา. \\ อญฺญาณมูลปฺปภวา ปชปฺปิตา, \\ สพฺพา มยา พฺยนฺติกตา สมูลิกา.}}}

\prose{อญฺญาณ\-มูล\-ปฺปภวาติ ปุริมเกหิ สมุทโย. สพฺพา มยา พฺยนฺติกตา สมูลิกาติ นิโรโธ. อิมานิ เทฺว สจฺจานิ. จตุนฺนํ ธมฺมานํ อนนุโพธา อปฺปฏิเวธา วิตฺถาเรน กาตพฺพํ. อริยสฺส สีลสฺส สมาธิโน ปญฺญาย วิมุตฺติยา. ตตฺถ โย อิเมสํ จตุนฺนํ ธมฺมานํ อนนุโพธา อปฺปฏิเวธา, อยํ สมุทโย. ปฏิเวโธ ภวเนตฺติยา, อยํ นิโรโธ. อยํ สมุทโย จ นิโรโธ จ.}

\prose{ตตฺถ กตโม สมุทโย จ มคฺโค จ.}

\prose{ยานิ โสตานิ โลกสฺมิํ, (อชิตาติ ภควา,)}

\prose{สติ เตสํ นิวารณํ.}

\setlength{\csromangathapagewidth}{0pt}
\setlength{\csromangathawakwidth}{0pt}
\csromangathameasuregroup{โสตานํ สํวรํ พฺรูมิ,}{\csromangathabat{โสตานํ สํวรํ พฺรูมิ,}{ปญฺญาเยเต ปิธียเร\textsuperscript{1}.}}
\csromangathasetleft{โสตานํ สํวรํ พฺรูมิ,}
\csromangathagroup{\csromangathastanza{\csromangathabat{โสตานํ สํวรํ พฺรูมิ,}{ปญฺญาเยเต ปิธียเร\csromansharedfootnote{eaf9c113afe6366a}{ขุ 1. 434 ปิฏฺเฐ.}.}}}

\prose{ยานิ โสตานีติ อยํ สมุทโย. ยา จ ปญฺญา ยา จ สติ นิวารณํ ปิธานญฺจ, อยํ มคฺโค. อิมานิ เทฺว สจฺจานิ.  สญฺเจตนิยํ สุตฺตํ ทฬฺหเนมิ\-ยานา\-กาโร ฉหิ มาเสหิ นิทฺทิฏฺโฐ. ตตฺถ ยํ กายํ กายกมฺมํ สวงฺกํ สโทสํ สกสาวํ ยา สวงฺกตา สโทสตา สกสาวตา, อยํ สมุทโย. เอวํ วจีกมฺมํ มโนกมฺมํ อวงฺกํ อโทสํ อกสาวํ, ยา อวงฺกตา อโทสตา อกสาวตา, อยํ มคฺโค. เอวํ วจีกมฺมํ มโนกมฺมํ. อิมานิ เทฺว สจฺจานิ สมุทโย จ มคฺโค จ.}

\prose{ตตฺถ กตโม สมุทโย จนิโรโธ จ มคฺโค จ.}

\prose{“นิสฺสิตสฺส จลิตํ, อนิสฺสิตสฺส จลิตํ นตฺถิ, จลิเต อสติ ปสฺสทฺธิ, ปสฺสทฺธิยา สติ นติ น โหติ, นติยา อสติ\csromansharedfootnote{4580259bbf8556c3}{อสติยา (อิ) ขุ 1. 179; สํ 2. 284 ปิฏฺเฐสุ ปสฺสิตพฺพํ.} อาคติคติ น โหติ, อาคติคติยา อสติ จุตูปปาโต น โหติ, จุตูปปาเต อสติ เนวิธ น หุรํ น อุภยม\-นฺตเรน. เอเสวนฺโต ทุกฺขสฺสา”ติ.}

\prose{\csromanfolio{181}ตตฺถ เทฺว นิสฺสยา, อยํ สมุทโย. โย จ อนิสฺสโย, ยา จ อนติ, อยํ มคฺโค. ยา อาคติคติ น โหติ จุตูปปาโต จ โย เอเสวนฺโต ทุกฺขสฺสาติ, อยํ นิโรโธ. อิมานิ ตีณิ สจฺจานิ.  อนุปฏฺฐิต\-กายคตา สติ ฯเปฯ ยํ วิมุตฺติ\-ญาณ\-ทสฺสนํ, อยํ สมุทโย. เอการสอุปนิสฺสยา วิมุตฺติโย ยาว อุปนิสฺสยอุปสมฺปทา อุปฏฺฐิต\-กายคตา\-สติสฺส วิหรติ. สีลสํวโร โสสานิโย โหติ, ยญฺจ วิมุตฺติ\-ญาณ\-ทสฺสนํ, อยํ มคฺโค. ยา จ วิมุตฺติ, อยํ นิโรโธ. อิมานิ ตีณิ สจฺจานิ. สมุทโย จ นิโรโธ จ มคฺโค จ.}

\setlength{\csromangathapagewidth}{0pt}
\setlength{\csromangathawakwidth}{0pt}
\csromangathameasuregroup{}{ตตฺถ กตโม นิโรโธ จ มคฺโค จ. \\ สยํ กเตน สจฺเจน, \\ เตน อตฺตนา อภินิพฺพานคโต วิติณฺณกงฺโข. \\ วิภวญฺจ ญตฺวา โลกสฺมิํ, \\ ตาว ขีณปุนพฺภโว ส ภิกฺขุ.}
\csromangathameasurewak{ตตฺถ กตโม นิโรโธ จ มคฺโค จ. \\ สยํ กเตน สจฺเจน, \\ เตน อตฺตนา อภินิพฺพานคโต วิติณฺณกงฺโข. \\ วิภวญฺจ ญตฺวา โลกสฺมิํ, \\ ตาว ขีณปุนพฺภโว ส ภิกฺขุ.}
\setlength{\csromangathaleftcol}{0pt}
\csromangathagroup{\csromangathastanza{\csromangathawak{\csromangathaitemline{12}{ตตฺถ กตโม นิโรโธ จ มคฺโค จ.} \\ \csromangathacontline{สยํ กเตน สจฺเจน,} \\ \csromangathacontline{เตน อตฺตนา อภินิพฺพานคโต วิติณฺณกงฺโข.} \\ \csromangathacontline{วิภวญฺจ ญตฺวา โลกสฺมิํ,} \\ \csromangathacontline{ตาว ขีณปุนพฺภโว ส ภิกฺขุ.}}}}

\prose{ยํ สจฺเจน, อยํ มคฺโค. ยํ ขีณปุนพฺภโว, อยํ นิโรโธ. อิมานิ เทฺว สจฺจานิ.  ปญฺจ วิมุตฺตา\-ยตนานิ สตฺถา วา ธมฺมํ เทเสสิ อญฺญตโร วา วิญฺญู สพฺราหฺมจารีติ วิตฺถาเรน กาตพฺพา. ตสฺส อตฺถ\-ปฺปฏิสํเวทิสฺส ปาโมชฺชํ ชายติ, ปมุทิตสฺส ปีติ ชายติ, ยาว นิพฺพินฺทนฺโต วิรชฺชติ, อยํ มคฺโค. ยา วิมุตฺติ, อยํ นิโรโธ. เอวํ ปญฺจ วิมุตฺตา\-ยตนานิ วิตฺถาเรน, อิมานิ เทฺว สจฺจานิ นิโรโธ จ มคฺโค จ.}

\prose{อิมานิ สาธารณานิ สุตฺตานิ. อิเมหิ สาธารเณหิ สุตฺเตหิ ยถา\-นิกฺขิตฺเตหิ ปฏิเวธโต จ ลกฺขณโต จ โอตาเรตฺวา อญฺญานิ สุตฺตานิ นิทฺทิสิตพฺพานิ อปริหายนฺเตน. คาถาหิ คาถา อนุมินิตพฺพา, พฺยากรเณหิ พฺยากรณํ. อิเม จ สาธารณา ทส ปริวฑฺฒกา เอโก จ จตุกฺโก นิทฺเทโส สาธารโณ. อยญฺจ ปกิณฺณก\-นิทฺเทโส. เอกํ ปญฺจ ฉ จ สเวกเทโส สพฺพํ. อิเม เทฺว ปริวชฺชนา ปุริมกา จ ทส. อิเม ทฺวาทส ปริวฑฺฒกา สจฺจานิ.  เอตฺตาวตา สพฺพํ สุตฺตํ นตฺถิ, ตํ พฺยากรณํ วา คาถา วิย. อิเมหิ ทฺวาทสหิ ปริวฑฺฒเกหิ น โอตริตุํ อปฺปมตฺเตน ปริเยสิตฺวา นิทฺทิสิตพฺพา.}

\prose{ตตฺถายํ สงฺเขโป. สพฺพํ ทุกฺขํ สตฺตหิ ปเทหิ สโมสรณํ คจฺฉติ. กตเรหิ สตฺตหิ, อปฺปิย\-สมฺปโยโค จ ปิยวิปฺปโยโค จ, อิเมหิ \csromanfolio{182}ทฺวีหิ ปเทหิ สพฺพํ ทุกฺขํ นิทฺทิสิตพฺพํ. ตสฺส เทฺว นิสฺสยา กาโย จ จิตฺตญฺจ, เตน วุจฺจติ “กายิกํ ทุกฺขํ เจตสิกญฺเจ”ติ, นตฺถิ ตํ ทุกฺขํ น กายิกํ วา น เจตสิกํ, สพฺพํ ทุกฺขํ ทฺวีหิ ทุกฺเขหิ นิทฺทิสิตพฺพํ กายิเกน จ เจตสิเกน จ. ตีหิ ทุกฺขตาหิ สงฺคหิตํ ทุกฺข\-ทุกฺขตาย สงฺขาร\-ทุกฺขตาย วิปริณาม\-ทุกฺขตาย. อิติ ตํ สพฺพํ ทุกฺขํ ตีหิ ทุกฺขตาหิ สงฺคหิตํ, อิติ อิทญฺจ ทุกฺขํ ติวิธํ, ทุวิธํ ทุกฺขํ กายิกญฺจ เจตสิกญฺจ. ทุวิธํ อปฺปิย\-สมฺปโยโค จ ปิยวิปฺปโยโค จ, อิทํ สตฺตวิธํ ทุกฺขํ.}

\prose{ตตฺถ ติวิโธ สมุทโย อจตุตฺโถ อปญฺจโม. กตโม ติวิโธ, ตณฺหา จ ทิฏฺฐิ จ กมฺมํ, ตตฺถ ตณฺหา จ ภวสมุทโย กมฺมํ. ตถา\csromansharedfootnote{3558827ed63cec44}{ตตฺถ (อิ)} นิพฺพตฺตสฺส หีนปณีตตา\csromansharedfootnote{a3df5326abe354fb}{หีนปณีตตาย (อิ)}, อยํ สมุทโย. อิติ ยาปิ ภวคตีสุ หีนตา จ ปณีตตา จ, ยาปิ ตีหิ ทุกฺขตาหิ สงฺคหิตา, โยปิ ทฺวีหิ มูเลหิ สมุทานีโต อวิชฺชาย นิวุตสฺส ภวตณฺหา\-สํยุตฺตสฺส สวิญฺญาณโก กาโย, โสปิ ตีหิ ทุกฺขตาหิ สงฺคหิโต.}

\prose{ตถา วิปลฺลาสโต ทิฏฺฐิ อาคนฺตพฺพา\csromansharedfootnote{da7f3e573d6d8fe5}{มวคนฺตพฺพา (ก)}. สา สตฺตวิธา นิทฺทิสิตพฺพา. เอโก วิปลฺลาโส ตีณิ นิทฺทิสียติ จตฺตาริ วิปลฺลาส\-วตฺถูนิ. ตตฺถ กตโม เอโก วิปลฺลาโส, โย วิปรีตคฺโคโห ปฏิกฺเขเปน, โอตรณํ ยถา “อนิจฺเจ นิจฺจมฺ”อิติ วิปรีตํ คณฺหาติ. เอวํ จตฺตาโร วิปลฺลาสา. อยเมโก วิปลฺลาสียติ สญฺญา จิตฺตํ ทิฏฺฐิ. กตมานิ จตฺตาริ วิปลฺลาส\-วตฺถูนิ, กาโย เวทนา จิตฺตํ ธมฺมา. เอวํ วิปลฺลาส\-คตสฺส อกุสลญฺจ ปวฑฺเฒติ. ตตฺถ สญฺญาวิปลฺลาโส โทสํ อกุสลมูลํ ปวฑฺเฒติ. จิตฺตวิปลฺลาโส โลภํ อกุสลมูลํ ปวฑฺเฒติ. ทิฏฺฐิ\-วิปลฺลาโส โมหํ อกุสลมูลํ ปวฑฺเฒติ. ตตฺถ โทสสฺส อกุสลมูลสฺส ตีณิ มิจฺฉตฺตานิ ผลํ มิจฺฉาวาจา มิจฺฉากมฺมนฺโต มิจฺฉาอาชีโว, โลภสฺส อกุสลมูลสฺส ตีณิ มิจฺฉตฺตานิ ผลํ มิจฺฉาสงฺกปฺโป มิจฺฉาวายาโม มิจฺฉาสมาธิ, โมหสฺส อกุสลมูลสฺส เทฺว มิจฺฉตฺตานิ ผลํ มิจฺฉาทิฏฺฐิ จ มิจฺฉาสติ จ. เอวํ อกุสลํ สเหตุ สปฺปจฺจยํ วิปลฺลาสา จ ปจฺจโย, อกุสลมูลานิ สเหตู เอเตเยว ปฏิปกฺเขน อนูนา อนธิกา ทฺวีหิ ปจฺจเยหิ นิทฺทิสิตพฺพา. นิโรเธ จ มคฺเค จ วิปลฺลาสมุ\-ปาทาย ปรโต\csromansharedfootnote{f4d468ac769035f1}{ปริโต (อิ)} ปฏิปกฺเขน จตสฺโส.}

\csromanmatikaanchor{mtk.41}
\csromantocmarkref{section}{2. สาสนปฏฺฐานทุติยภูมิ}{mtk.41}
\bookmark[dest={mtk.41},level=1]{2. สาสนปฏฺฐานทุติยภูมิ}
\csromanheader{1}{2. สาสน\-ปฏฺฐาน\-ทุติยภูมิ \textbf{ตตฺถิมา อุทฺทานคาถา}}
\setlength{\csromangathapagewidth}{0pt}
\setlength{\csromangathawakwidth}{0pt}
\csromangathameasuregroup{อวิชฺชาย นิวุโต โลโก, \\ สา ปจฺฉินฺนภวตณฺหา, \\ กุมฺภูปมํ กายมิมํ, \\ เย เกจิ โสกปริเทวา, \\ ทุกฺขํ ทิฏฺฐิ จ อุปฺปนฺนํ, \\ สติ กายคตา มาหุ, \\ อาสา ปิหา จ อภินนฺทนา จ, \\ ยานิ โสตานิ โลกสฺมิํ, \\ ยํ นิสฺสิตสฺส จลิตํ, \\ สยํ กเตน สจฺเจน,}{\csromangathabat{อวิชฺชาย นิวุโต โลโก,}{จิตฺตํ สํโยชนมฺปิ.} \\ \csromangathabat{สา ปจฺฉินฺนภวตณฺหา,}{เทฺวมา เจว วิมุตฺติโย.} \\ \csromangathabat{กุมฺภูปมํ กายมิมํ,}{ยํ น ตุมฺหากํ ตํ ปชห\textsuperscript{1}.} \\ \csromangathabat{เย เกจิ โสกปริเทวา,}{ติมฺพรุโก จ สยํกตํ.} \\ \csromangathabat{ทุกฺขํ ทิฏฺฐิ จ อุปฺปนฺนํ,}{ยญฺจ โคปาลโกปมํ.} \\ \csromangathabat{สติ กายคตา มาหุ,}{สมโถ จ วิปสฺสนา.} \\ \csromangathabat{อาสา ปิหา จ อภินนฺทนา จ,}{จตุนฺนมนนุโพธนา.} \\ \csromangathabat{ยานิ โสตานิ โลกสฺมิํ,}{ทฬฺหํ เนมิยานากาโร.} \\ \csromangathabat{ยํ นิสฺสิตสฺส จลิตํ,}{อนุปฏฺฐิตกายคตาสติ.} \\ \csromangathabat{สยํ กเตน สจฺเจน,}{วิมุตฺตายตเนหิ จ.}}
\csromangathasetleft{อวิชฺชาย นิวุโต โลโก, \\ สา ปจฺฉินฺนภวตณฺหา, \\ กุมฺภูปมํ กายมิมํ, \\ เย เกจิ โสกปริเทวา, \\ ทุกฺขํ ทิฏฺฐิ จ อุปฺปนฺนํ, \\ สติ กายคตา มาหุ, \\ อาสา ปิหา จ อภินนฺทนา จ, \\ ยานิ โสตานิ โลกสฺมิํ, \\ ยํ นิสฺสิตสฺส จลิตํ, \\ สยํ กเตน สจฺเจน,}
\csromangathagroup{\csromangathastanza{\csromanfolio{183}\csromangathabat{อวิชฺชาย นิวุโต โลโก,}{จิตฺตํ สํโยชนมฺปิ.} \\ \csromangathabat{สา ปจฺฉินฺนภวตณฺหา,}{เทฺวมา เจว วิมุตฺติโย.}}\csromangathastanza{\csromangathabat{กุมฺภูปมํ กายมิมํ,}{ยํ น ตุมฺหากํ ตํ ปชห\csromansharedfootnote{2a13fbcf832b034a}{ชหา (อิ, ก)}.} \\ \csromangathabat{เย เกจิ โสกปริเทวา,}{ติมฺพรุโก จ สยํกตํ.}}\csromangathastanza{\csromangathabat{ทุกฺขํ ทิฏฺฐิ จ อุปฺปนฺนํ,}{ยญฺจ โคปาลโกปมํ.} \\ \csromangathabat{สติ กายคตา มาหุ,}{สมโถ จ วิปสฺสนา.}}\csromangathastanza{\csromangathabat{อาสา ปิหา จ อภินนฺทนา จ,}{จตุนฺนมนนุโพธนา.} \\ \csromangathabat{ยานิ โสตานิ โลกสฺมิํ,}{ทฬฺหํ เนมิยานากาโร.}}\csromangathastanza{\csromangathabat{ยํ นิสฺสิตสฺส จลิตํ,}{อนุปฏฺฐิตกายคตาสติ.} \\ \csromangathabat{สยํ กเตน สจฺเจน,}{วิมุตฺตายตเนหิ จ.}}}

\prose{เปฏโกปเทเส มหา\-กจฺจายเนน ภาสิเต ปฐมภูมิ อริยสจฺจปฺปกาสนา นาม ตํ ชีวตา ภควตา มาทิเสน สมุทฺทเนน ตถาคเตนาติ.}

\proseitem{13}{ตตฺถ กตมํ สาสน\-ปฺปฏฺฐานํ, สํกิเลส\-ภาคิยํ สุตฺตํ, วาสนา ภาคิยํ สุตฺตํ, นิพฺเพธ\-ภาคิยํ สุตฺตํ, อเสกฺข\-ภาคิยํ สุตฺตํ, สํกิเลสภาคิยญฺจ วาสนาภาคิยญฺจ, สํกิเลสภาคิยญฺจ นิพฺเพธภาคิยญฺจ, สํกิเลสภาคิยญฺจ นิพฺเพธภาคิยญฺจ อเสกฺขภาคิยญฺจ, วาสนาภาคิยญฺจ นิพฺเพธภาคิยญฺจ. อาณตฺติ, ผลํ, อุปาโย, อาณตฺติ จ ผลญฺจ, ผลญฺจ อุปาโย จ, อาณตฺติ จ ผลญฺจ อุปาโย จ. อสฺสาโท, อาทีนโว, นิสฺสรณํ, อสฺสาโท จ อาทีนโว จ, อสฺสาโท จ นิสฺสรณญฺจ, อาทีนโว จ นิสฺสรณญฺจ, อสฺสาโท จ อาทีนโว จ นิสฺสรณญฺจ. โลกิกํ, โลกุตฺตรํ, โลกิกญฺจ โลกุตฺตรญฺจ. กมฺมํ, วิปาโก, กมฺมญฺจ วิปาโก จ. นิทฺทิฏฺฐํ, อนิทฺทิฐํ, นิทฺทิฏฺฐญฺจ อนิทฺทิฏฺฐญฺจ. ญาณํ, เญยฺยํ, ญาณญฺจ เญยฺยญฺจ. ทสฺสนํ, ภาวนา, ทสฺสนญฺจ ภาวนา จ. วิปากกมฺมํ, น วิปากกมฺมํ, \csromanfolio{184}เนววิปากนวิปากกมฺมํ. สกวจนํ, ปรวจนํ, สกวจนญฺจ ปรวจนญฺจ. สตฺตา\-ธิฏฺฐานํ, ธมฺมา\-ธิฏฺฐานํ, สตฺตา\-ธิฏฺฐานญฺจ ธมฺมาธิฏฺฐานญฺจ. ถโว, สกวจนา\-ธิฏฺฐานํ, ปรวจนา\-ธิฏฺฐานํ, สกวจนา\-ธิฏฺฐานญฺจ ปรวจนา\-ธิฏฺฐานญฺจ. กิริยํ, ผลํ, กิริยญฺจ ผลญฺจ. อนุญฺญาตํ, ปฏิกฺขิตฺตํ, อนุญฺญาตญฺจ ปฏิกฺขิตฺตญฺจ. อิมานิ ฉ ปฏิกฺขิตฺตานิ.}

\proseitem{14}{ตตฺถ กตมํ สํกิเลส\-ภาคิยํ สุตฺตํ.}

\setlength{\csromangathapagewidth}{0pt}
\setlength{\csromangathawakwidth}{0pt}
\csromangathameasuregroup{กามนฺธา ชาลสญฺฉนฺนา, \\ ปมตฺตพนฺธุนา พทฺธา, \\ ชรามรณมเนฺวนฺติ,}{\csromangathabat{กามนฺธา ชาลสญฺฉนฺนา,}{ตณฺหาฉทนฉาทิตา.} \\ \csromangathabat{ปมตฺตพนฺธุนา พทฺธา,}{มจฺฉาว กุมินามุเข.} \\ \csromangathabat{ชรามรณมเนฺวนฺติ,}{วิจฺโฉ ขีรปโกว\textsuperscript{1} มาตรํ.} \\ ปญฺจิเม ภิกฺขเว นีวรณา.}
\csromangathasetleft{กามนฺธา ชาลสญฺฉนฺนา, \\ ปมตฺตพนฺธุนา พทฺธา, \\ ชรามรณมเนฺวนฺติ,}
\csromangathagroup{\csromangathastanza{\csromangathabat{กามนฺธา ชาลสญฺฉนฺนา,}{ตณฺหา\-ฉทน\-ฉาทิตา.} \\ \csromangathabat{ปมตฺต\-พนฺธุนา พทฺธา,}{มจฺฉาว กุมินามุเข.}}\csromangathastanza{\csromangathabat{ชรามรณมเนฺวนฺติ,}{วิจฺโฉ ขีรปโกว\csromansharedfootnote{338dd9d64fd7ee1e}{[มิสฺสินฺคฺ โนเต 0]} มาตรํ.} \\ ปญฺจิเม ภิกฺขเว นีวรณา.}}

\prose{ตตฺถ กตมํ วาสนาภาคิยํ สุตฺตํ.}

\prose{\csromansymbolfootnote{*}{เหฏฺฐา 112 ปิฏฺเฐ; อุปริ 277 ปิฏฺเฐปิ. ** เหฏฺฐา 54, 123 ปิฏฺเฐสุปิ.}มโนปุพฺพงฺคมา ธมฺมา, มโนเสฏฺฐา มโนมยา.}

\setlength{\csromangathapagewidth}{0pt}
\setlength{\csromangathawakwidth}{0pt}
\csromangathameasuregroup{มนสา เจ ปสนฺเนน, \\ ตโต นํ สุขมเนฺวติ,}{\csromangathabat{มนสา เจ ปสนฺเนน,}{ภาสติ วา กโรติ วา.} \\ \csromangathabat{ตโต นํ สุขมเนฺวติ,}{ฉายาว อนปายินี.} \\ สํยุตฺตเก สุตฺตํ.}
\csromangathameasurewak{สํยุตฺตเก สุตฺตํ.}
\csromangathasetleft{มนสา เจ ปสนฺเนน, \\ ตโต นํ สุขมเนฺวติ,}
\csromangathagroup{\csromangathastanza{\csromangathabat{มนสา เจ ปสนฺเนน,}{ภาสติ วา กโรติ วา.} \\ \csromangathabat{ตโต นํ สุขมเนฺวติ,}{ฉายาว อนปายินี.}}\csromangathastanza{\csromangathawak{สํยุตฺตเก สุตฺตํ.}}}

\prose{มหานามสฺส สกฺกสฺส อิทํ ภควา สกฺยานํ กปิล\-วตฺถุมฺหิ นคเร นยวิตฺถาเรน สทฺธา\-สีล\-ปริภาวิตํ สุตฺตํ ภาวญฺเญน ปริภาวิตํ ตํ นาม ปจฺฉิเม กาเล.}

\prose{ตตฺถ กตมํ นิพฺเพธ\-ภาคิยํ สุตฺตํ.}

\prose{** อุทฺธํ อโธ สพฺพธิ วิปฺปมุตฺโต,}

\prose{อยํ อหสฺมีติ อนานุปสฺสี.}

\setlength{\csromangathapagewidth}{0pt}
\setlength{\csromangathawakwidth}{0pt}
\csromangathameasuregroup{}{เอวํ วิมุตฺโต อุทตาริ โอฆํ, \\ อติณฺณปุพฺพํ อปุนพฺภวาย\textsuperscript{1}.}
\csromangathameasurewak{เอวํ วิมุตฺโต อุทตาริ โอฆํ, \\ อติณฺณปุพฺพํ อปุนพฺภวาย\textsuperscript{1}.}
\setlength{\csromangathaleftcol}{0pt}
\csromangathagroup{\csromangathastanza{\csromangathawak{เอวํ วิมุตฺโต อุทตาริ โอฆํ, \\ อติณฺณปุพฺพํ อปุนพฺภวาย\csromansharedfootnote{e1d5ea17ce7e006e}{ขุ 1. 170 ปิฏฺเฐ.}.}}}

\prose{สีลานิ นุ โข ภวนฺติ กิมตฺถิยานิ อานนฺโท ปุจฺฉติ สตฺถารํ.}

\prose{\csromanfolio{185}ตตฺถ กตมํ อเสกฺข\-ภาคิยํ สุตฺตํ.}

\prose{\csromansymbolfootnote{*}{เหฏฺฐา 129 ปิฏฺเฐ.}ยสฺส เสลูปมํ จิตฺตํ, ฐิตํ นานุปกมฺปติ.}

\setlength{\csromangathapagewidth}{0pt}
\setlength{\csromangathawakwidth}{0pt}
\csromangathameasuregroup{วิรตฺตํ รชนีเยสุ, \\ ยสฺเสวํ ภาวิตํ จิตฺตํ,}{\csromangathabat{วิรตฺตํ รชนีเยสุ,}{โกปเนยฺเย\textsuperscript{1} น กุปฺปติ.} \\ \csromangathabat{ยสฺเสวํ ภาวิตํ จิตฺตํ,}{กุโต ตํ ทุกฺขเมสฺสตีติ.}}
\csromangathasetleft{วิรตฺตํ รชนีเยสุ, \\ ยสฺเสวํ ภาวิตํ จิตฺตํ,}
\csromangathagroup{\csromangathastanza{\csromangathabat{วิรตฺตํ รชนีเยสุ,}{โกปเนยฺเย\csromansharedfootnote{90dbf0ddbe4bf4f8}{โกปนีเย (ก) ขุ 1. 125 ปิฏฺเฐ. ** เหฏฺฐา 132 ปิฏฺเฐ.} น กุปฺปติ.} \\ \csromangathabat{ยสฺเสวํ ภาวิตํ จิตฺตํ,}{กุโต ตํ ทุกฺขเมสฺสตีติ.}}}

\prose{สาริปุตฺโต นาม ภควา เถรญฺญตโร โส มํ อาสชฺช อปฺปฏินิสฺสชฺช จาริกํ ปกฺกมติ, สาริปุตฺตสฺส พฺยากรณํ กาตพฺพํ. ยสฺส นูน ภควา กายคตา สติ อภาวิตา อสฺส อพหุลีกตา วิตฺถาเรน กาตพฺพํ.}

\proseitem{15}{ตตฺถ กตมํ สํกิเลสภาคิยญฺจ วาสนาภาคิยญฺจ.}

\setlength{\csromangathapagewidth}{0pt}
\setlength{\csromangathawakwidth}{0pt}
\csromangathameasuregroup{** ฉนฺนมติวสฺสติ, \\ ตสฺมา ฉนฺนํ วิวเรถ,}{\csromangathabat{** ฉนฺนมติวสฺสติ,}{วิวฏํ นาติวสฺสติ.} \\ \csromangathabat{ตสฺมา ฉนฺนํ วิวเรถ,}{เอวํ ตํ นาติวสฺสติ\textsuperscript{1}.}}
\csromangathasetleft{** ฉนฺนมติวสฺสติ, \\ ตสฺมา ฉนฺนํ วิวเรถ,}
\csromangathagroup{\csromangathastanza{\csromangathaitembat{15}{** ฉนฺนม\-ติวสฺสติ,}{วิวฏํ นาติวสฺสติ.} \\ \csromangathacontbat{ตสฺมา ฉนฺนํ วิวเรถ,}{เอวํ ตํ นาติวสฺสติ\csromansharedfootnote{591252522f51efb5}{ขุ 1. 144 ปิฏฺเฐ. + เหฏฺฐา 31, 132 ปิฏฺเฐสุปิ.}.}}}

\prose{ฉนฺนม\-ติวสฺสตีติ สํกิเลโส. วิวฏํ นาติวสฺสตีติ วาสนา. ตโม ตมปรายโนติ วิตฺถาเรน. ตตฺถ โย จ ตโม โย จ ตมปรายโน, อยํ สํกิเลโส. โย จ โชติ โย จ โชติปรายโน, อยํ วาสนา.}

\prose{ตตฺถ กตมํ สํกิเลสภาคิยญฺจ นิพฺเพธภาคิยญฺจ สุตฺตํ.}

\prose{\csromansymbolmark{+}น ตํ ทฬฺหํ พนฺธนมาหุ ธีรา,}

\prose{ยทายสํ ทารุชปพฺพชญฺจ\csromansharedfootnote{321063897b501b0d}{ทารุชํ พพฺพชญฺจ (อิ) ขุ 1. 63 ปิฏฺเฐ; สํ 1. 77 ปิฏฺเฐ จ.}.}

\setlength{\csromangathapagewidth}{0pt}
\setlength{\csromangathawakwidth}{0pt}
\csromangathameasuregroup{}{สารตฺตรตฺตา มณิกุณฺฑเลสุ, \\ ปุตฺเตสุ ทาเรสุ จ ยา อเปกฺขา.}
\csromangathameasurewak{สารตฺตรตฺตา มณิกุณฺฑเลสุ, \\ ปุตฺเตสุ ทาเรสุ จ ยา อเปกฺขา.}
\setlength{\csromangathaleftcol}{0pt}
\csromangathagroup{\csromangathastanza{\csromangathawak{สารตฺตรตฺตา มณิกุณฺฑเลสุ, \\ ปุตฺเตสุ ทาเรสุ จ ยา อเปกฺขา.}}}

\prose{น ตํ ทฬํ พนฺธนมาหุ ธีรา, ยทา ปุตฺเตสุ ทาเรสุ จ ยา อเปกฺขา, อยํ สํกิเลโส. เอตมฺปิ เฉตฺวา ปริพฺพชนฺติ ธีรา อนเปกฺขิโน สพฺพกาเม ปหายาติ, อยํ นิพฺเพโธ.  ยํ เจตยิตํ ปกปฺปิตํ ยา จ นามรูปสฺส อวกฺกนฺติ โหติ. อิเมหิ จตูหิ ปเทหิ สํกิเลโส. ปจฺฉิมเกหิ จตูหิ นิพฺเพโธ.}

\prose{ตตฺถ กตมํ สํกิเลสภาคิยญฺจ นิพฺเพธภาคิยญฺจ อเสกฺขภาคิยญฺจ สุตฺตํ.}

\setlength{\csromangathapagewidth}{0pt}
\setlength{\csromangathawakwidth}{0pt}
\csromangathameasuregroup{}{\textsuperscript{*}อยํ โลโก สนฺตาปชาโต, \\ ผสฺสปเรโต โรคํ\textsuperscript{1} วทติ อตฺตโต. \\ เยน เยน หิ มญฺญนฺติ, \\ ตโต ตํ โหติ อญฺญถา. \\ \textsuperscript{*}อญฺญถาภาวี ภวสตฺโต โลโก, \\ ภวปเรโต ภวเมวาภินนฺทติ. \\ ยทภินนฺทติ ตํ ภยํ, \\ ยสฺส ภายติ ตํ ทุกฺขํ. \\ ภวิปฺปหานาย โข ปนิทํ พฺรหฺมจริยํ วุสฺสติ.}
\csromangathameasurewak{\textsuperscript{*}อยํ โลโก สนฺตาปชาโต, \\ ผสฺสปเรโต โรคํ\textsuperscript{1} วทติ อตฺตโต. \\ เยน เยน หิ มญฺญนฺติ, \\ ตโต ตํ โหติ อญฺญถา. \\ \textsuperscript{*}อญฺญถาภาวี ภวสตฺโต โลโก, \\ ภวปเรโต ภวเมวาภินนฺทติ. \\ ยทภินนฺทติ ตํ ภยํ, \\ ยสฺส ภายติ ตํ ทุกฺขํ.}
\setlength{\csromangathaleftcol}{0pt}
\csromangathagroup{\csromangathastanza{\csromangathawak{\csromanfolio{186}\csromansymbolfootnote{*}{เหฏฺฐา 135 ปิฏฺเฐ.}อยํ โลโก สนฺตาปชาโต, \\ ผสฺสปเรโต โรคํ\csromansharedfootnote{9eb4f6ad6cd926db}{โรทํ (อิ) ขุ 1. 115 ปิฏฺเฐ. ** เหฏฺฐา 136 ปิฏฺเฐ.} วทติ อตฺตโต. \\ เยน เยน หิ มญฺญนฺติ, \\ ตโต ตํ โหติ อญฺญถา.}}\csromangathastanza{\csromangathawak{\csromansymbolmark{*}อญฺญถาภาวี ภวสตฺโต โลโก, \\ ภวปเรโต ภวเม\-วา\-ภินนฺทติ. \\ ยทภินนฺทติ ตํ ภยํ, \\ ยสฺส ภายติ ตํ ทุกฺขํ.}}\csromangathastanza{ภวิปฺปหานาย โข ปนิทํ พฺรหฺมจริยํ วุสฺสติ.}}

\prose{เย หิ เกจิ สมณา วา พฺราหฺมณา วา ภเวน ภวสฺส วิปฺปโมกฺขมาหํสุ, สพฺเพเต “อวิปฺปมุตฺตา ภวสฺมา”ติ วทามิ. เย วา ปน เกจิ สมณา วา พฺราหฺมณา วา วิภเวน ภวสฺส นิสฺสรณมาหํสุ, สพฺเพเต “อนิสฺสฏา ภวสฺมา”ติ วทามิ. อุปธิํ หิ ปฏิจฺจ ทุกฺขมิทํ สมฺโภติ, สพฺพุ\-ปาทาน\-กฺขยา นตฺถิ ทุกฺขสฺส สมฺภโว, โลกมิมํ ปสฺส, ปุถู อวิชฺชาย ปเรตา ภูตา ภูตรตา ภวา อปริมุตฺตา. เย หิ เกจิ ภวา สพฺพธิ สพฺพตฺถตาย, สพฺเพเต ภวา อนิจฺจา ทุกฺขา วิปริณาม\-ธมฺมาติ.}

\prose{** “เอวเมตํ ยถาภูตํ,}

\setlength{\csromangathapagewidth}{0pt}
\setlength{\csromangathawakwidth}{0pt}
\csromangathameasuregroup{}{สมฺมปฺปญฺญาย ปสฺสโต. \\ ภวตณฺหา ปหียติ, \\ วิภวํ นาภินนฺทติ. \\ สพฺพโส ตณฺหานํ ขยา, \\ อเสสวิราคนิโรโธ นิพฺพานํ. \\ ** ตสฺส นิพฺพุตสฺส ภิกฺขุโน, \\ อนุปาทา ปุนพฺภโว น โหติ. \\ อภิภูโต มาโร วิชิตสงฺคาโม, \\ อุเปจฺจคา สพฺพภวานิ ตาที”ติ.}
\csromangathameasurewak{สมฺมปฺปญฺญาย ปสฺสโต. \\ ภวตณฺหา ปหียติ, \\ วิภวํ นาภินนฺทติ. \\ สพฺพโส ตณฺหานํ ขยา, \\ อเสสวิราคนิโรโธ นิพฺพานํ. \\ ** ตสฺส นิพฺพุตสฺส ภิกฺขุโน, \\ อนุปาทา ปุนพฺภโว น โหติ. \\ อภิภูโต มาโร วิชิตสงฺคาโม, \\ อุเปจฺจคา สพฺพภวานิ ตาที”ติ.}
\setlength{\csromangathaleftcol}{0pt}
\csromangathagroup{\csromangathastanza{\csromangathawak{สมฺมปฺปญฺญาย ปสฺสโต. \\ ภวตณฺหา ปหียติ, \\ วิภวํ นาภินนฺทติ.}}\csromangathastanza{\csromangathawak{สพฺพโส ตณฺหานํ ขยา, \\ อเสส\-วิราค\-นิโรโธ นิพฺพานํ. \\ ** ตสฺส นิพฺพุตสฺส ภิกฺขุโน, \\ อนุปาทา ปุนพฺภโว น โหติ. \\ อภิภูโต มาโร วิชิตสงฺคาโม, \\ อุเปจฺจคา สพฺพภวานิ ตาที”ติ.}}}

\prose{อยํ โลโก สนฺตาปชาโต ยาว ทุกฺขนฺติ ยํ ตณฺหา สํกิเลโส.}

\prose{ยํ ปุนคฺคหณํ เย หิ เกจิ สมณา วา พฺราหฺมณา วา ภเวน ภวสฺส วิโมกฺขมาหํสุ, สพฺเพเต “อวิมุตฺตา ภวสฺมา”ติ วทามิ. เย วา ปน เกจิ สมณา วา พฺราหฺมณา วา วิภเวน ภวสฺส นิสฺสรณมาหํสุ “อนิสฺสฏา ภวสฺมา”ติ วทามิ. อยํ ทิฏฺฐิ\-สํกิเลโส, ตํ \csromanfolio{187}ทิฏฺฐิ\-สํกิเลโส จ ตณฺหาสํกิเลโส จ, อุภยเมตํ สํกิเลโส. ยํ ปุนคฺคหณํ ภววิปฺปหานาย พฺรหฺมจริยํ วุสฺสติ, ยาว สพฺพโส อุปาทานกฺขยา สมฺภวา, อิทํ นิพฺเพธ\-ภาคิยํ. ตสฺส นิพฺพุตสฺส ภิกฺขุโน ยาว อุปจฺจคา สพฺพภวานิ ตาทีติ อิทํ อเสกฺข\-ภาคิยํ. จตฺตาโร ปุคฺคลา อนุโสตคามี สํกิเลโส ฐิตตฺโต จ ปฏิโสตคามี จ นิพฺเพโธ. ถเล ติฏฺฐตีติ อเสกฺขภูมิ.}

\proseitem{16}{ตตฺถ กตมํ วาสนาภาคิยญฺจ นิพฺเพธภาคิยญฺจ สุตฺตํ.}

\setlength{\csromangathapagewidth}{0pt}
\setlength{\csromangathawakwidth}{0pt}
\csromangathameasuregroup{}{“ททโต ปุญฺญํ ปวฑฺฒติ, \\ สํยมโต เวรํ น จียติ. \\ กุสโล จ ชหาติ ปาปกํ, \\ ราคโทสโมหกฺขยา สนิพฺพุโต”ติ\textsuperscript{1}.}
\csromangathameasurewak{“ททโต ปุญฺญํ ปวฑฺฒติ, \\ สํยมโต เวรํ น จียติ. \\ กุสโล จ ชหาติ ปาปกํ, \\ ราคโทสโมหกฺขยา สนิพฺพุโต”ติ\textsuperscript{1}.}
\setlength{\csromangathaleftcol}{0pt}
\csromangathagroup{\csromangathastanza{\csromangathawak{“ททโต ปุญฺญํ ปวฑฺฒติ, \\ สํยมโต เวรํ น จียติ. \\ กุสโล จ ชหาติ ปาปกํ, \\ ราคโทสโมห\-กฺขยา สนิพฺพุโต”ติ\csromansharedfootnote{715f40a7348ff1ad}{ขุ 1. 183 ปิฏฺเฐ; ที 2. 113 ปิฏฺเฐ จ.}.}}}

\prose{ททโต ปุญฺญํ ปวฑฺฒติ, สํยมโต เวรํ น จียตีติ วาสนา.  กุสโล จ ชหาติ ปาปกํ, ราคโทสโมห\-กฺขยา สนิพฺพุโตติ นิพฺเพโธ.}

\proseitem{16}{โสตานุคเตสุ ธมฺเมสุ วจสา ปริจิเตสุ มนสา\-นุเปกฺขิเตสุ ทิฏฺฐิยา สุปฺปฏิวิทฺเธสุ ปญฺจานิสํสา ปาฏิกงฺขา. อิเธกจฺจสฺส พหุสฺสุตา ธมฺมา โหนฺติ ธาตา อปมุฏฺฐา วจสา ปริจิตา มนสา\-นุเปกฺขิตา ทิฏฺฐิยา สุปฺปฏิวิทฺธา, โส ยุญฺชนฺโต ฆเฏนฺโต วายามนฺโต ทิฏฺเฐว ธมฺเม วิเสสํ ปปฺโปติ. โน เจ ทิฏฺเฐว ธมฺเม วิเสสํ ปปฺโปติ, คิลาโน ปปฺโปติ. โน เจ คิลาโน ปปฺโปติ, มรณกาล\-สมเย ปปฺโปติ. โน เจ มรณกาล\-สมเย ปปฺโปติ, เทวภูโต ปาปุณาติ. โน เจ เทวภูโต ปาปุณาติ, เตน ธมฺมราเคน ตาย ธมฺมนนฺทิยา ปจฺเจกโพธิํ ปาปุณาติ.}

\prose{ตตฺถายํ ทิฏฺเฐว ธมฺเม ปาปุณาติ, อยํ นิพฺเพโธ. ยํ สมฺปราเย ปจฺเจกโพธิํ ปาปุณาติ, อยํ วาสนา. อิมานิ โสฬส สุตฺตานิ สพฺพสาสนํ อติคฺคณฺหนฺโต ติฏฺฐนฺติ. อิเมหิ โสฬสหิ สุตฺเตหิ นววิโธ สุตฺตนฺโต วิภตฺโต ภวติ. โส จ ปญฺญวโต โน ทุปฺปญฺญสฺส, ยุตฺตสฺส โน อยุตฺตสฺส, อกมฺมสฺส วิหาริสฺส ปกติยา โลเก \csromanfolio{188}สํกิเลโส จรติ. โส สํกิเลโส ติวิโธ ตณฺหาสํกิเลโส ทิฏฺฐิ\-สํกิเลโส ทุจฺจริต\-สํกิเลโส. ตโต สํกิเลสโต อุฏฺฐหนฺโต สํกิเลโส ธมฺเมสุ ปติฏฺฐหติ, โลกิเยสุ ปติฏฺฐหตีติ. ตตฺถากุสโล ทิฏฺฐโต สเจ ตํ สีลญฺจ ทิฏฺฐิญฺจ ปรามสติ, ตสฺส โส ตณฺหาสํกิเลโส โหติ. สเจ ปนสฺส เอวํ โหติ “อิมินาหํ สีเลน วา วเตน วา พฺรหฺมจริเยน วา เทโว วา ภวิสฺสํ\csromansharedfootnote{1c59f891ee27e5da}{ภวิสฺสามิ (อิ)} เทวญฺญตโร วา”ติ ยสฺส โหติ มิจฺฉาทิฏฺฐิ, เอตสฺส มิจฺฉาทิฏฺฐิ\-สํกิเลโส ภวติ. สเจ ปน สีเล ปติฏฺฐิโต อปรามฏฺฐสฺส หิ สีลวตํ โหติ, ตสฺส ตํ สีลวโต โยนิโส คหิตํ อวิปฺปฏิสารํ ชเนติ ยาว วิมุตฺติ\-ญาณ\-ทสฺสนํ, ตญฺจ ตสฺส ทิฏฺเฐว ธมฺเม กาลงฺกตสฺส วา ตมฺหิเยว วา ปน อปรา\-ปริยาเยน วา, อญฺเญสุ ขนฺเธสุ เอวํ สุตํ “สุจริตํ วาสนาย สํวตฺตตี”ติ วาสนาภาคิยํ สุตฺตํ วุจฺจติ. ตตฺถ สีเลสุ ฐิตสฺส วินีวรณํ จิตฺตํ, ตํ ตโต สกฺกายทิฏฺฐิ\-ปฺปหานาย ภควา ธมฺมํ เทเสติ. โส อจฺจนฺตนิฏฺฐํ นิพฺพานํ ปาปุณาติ, ยทิ วา สาสนนฺตเร อจฺจนฺตํ นิพฺพานํ ปาปุณาติ, ยทิ วา เอกาสเน ฉ อภิญฺเญ. ตตฺถ เทฺว ปุคฺคลา อริยธมฺเม ปาปุณนฺติ สทฺธานุสารี จ ธมฺมานุสารี จ. ตตฺถ ธมฺมานุสารี อุคฺฆฏิตญฺญู, สทฺธานุสารี เนยฺโย. ตตฺถ อุคฺฆฏิตญฺญู ทุวิโธ โกจิ ติกฺขินฺทฺริโย โกจิ มุทินฺทฺริโย. ตตฺถ เนยฺโยปิ ทุวิโธ โกจิ ติกฺขินฺทฺริโย โกจิ มุทินฺทฺริโย. ตตฺถ โย จ อุคฺฆฏิตญฺญู มุทินฺทฺริโย, โย จ เนยฺโย ติกฺขินฺทฺริโย, อิเม ปุคฺคลา อสมินฺทฺริยา โหนฺติ. ตตฺถ อิเม ปุคฺคลา สมินฺทฺริยา ปริหายนฺติ จ อุคฺฆฏิตญฺญุโต, วิปญฺจิตญฺญู เนยฺยโต, อิเม มชฺฌิมา ภูมิคตา วิปญฺจิตญฺญู โหติ. อิเม ตโย ปุคฺคลา.}

\proseitem{17}{ตตฺถ จตุตฺถา ปน ปญฺจมา อุคฺฆฏิตญฺญู วิปญฺจิตญฺญู เนยฺโย จ, ตตฺถ อุคฺฆฏิตญฺญู ปุคฺคโล อินฺทฺริยานิ ปฏิลภิตฺวา ทสฺสน\-ภูมิยํ ฐิโต โสตาปตฺติผลญฺจ ปาปุณาติ, เอกพีชี โหติ ปฐโม โสตาปนฺโน. ตตฺถ วิปญฺจิตญฺญู ปุคฺคโล อินฺทฺริยานิ ปฏิลภิตฺวา ทสฺสน\-ภูมิยํ ฐิโต โสตาปตฺติผลญฺจ ปาปุณาติ, โกลํโกโล จ โหติ ทุติโย โสตาปนฺโน. ตตฺถ เนยฺโย ปุคฺคโล อินฺทฺริยานิ ปฏิลภิตฺวา ทสฺสน\-ภูมิยํ ฐิโต โสตาปตฺติผลญฺจ ปาปุณาติ, สตฺตกฺขตฺตุ\-ปรโม จ \csromanfolio{189}โหติ, อยํ ตติโย โสตาปนฺโน. อิเม ตโย ปุคฺคลา อินฺทฺริย\-เวมตฺตตาย โสตาปตฺติผเล ฐิตา.}

\prose{อุคฺฆฏิตญฺญู เอกพีชี โหติ, วิปญฺจิตญฺญู โกลํโกโล โหติ, เนยฺโย สตฺตกฺขตฺตุ\-ปรโม โหติ. อิทํ นิพฺเพธ\-ภาคิยํ สุตฺตํ. สเจ ปน ตทุตฺตริ วายมติ อจฺจนฺตนิฏฺฐํ นิพฺพานํ ปาปุณาติ. ตตฺถ อุคฺฆฏิตญฺญู ปุคฺคโล โย ติกฺขินฺทฺริโย, เต เทฺว ปุคฺคลา โหนฺติ อนาคามิผลํ ปาปุณิตฺวา อนฺตรา\-ปรินิพฺพายี จ อุปหจฺจ\-ปรินิพฺพายี จ. ตตฺถ วิปญฺจิตญฺญู ปุคฺคโล โย ติกฺขินฺทฺริโย, เต เทฺว ปุคฺคลา โหนฺติ อนาคามิผลํ ปาปุณนฺติ อสงฺขารปรินิพฺพายี จ สสงฺขาร\-ปรินิพฺพายี จ. ตตฺถ เนยฺโย อนาคามิผลํ ปาปุณนฺโต อุทฺธํโสโต อกนิฏฺฐคามี โหติ. อุคฺฆฏิตญฺญู จ วิปญฺจิตญฺญู จ, อินฺทฺริย\-นานตฺเตน อุคฺฆฏิตญฺญู ปุคฺคโล ติกฺขินฺทฺริโย อนฺตรา\-ปรินิพฺพายี โหติ, อุคฺคฏิตญฺญู มุทินฺทฺริโย อุทฺธํโสโต อกนิฏฺฐคามี โหติ. อุคฺฆฏิตญฺญู จ วิปญฺจิตญฺญู จ อินฺทฺริย\-นานตฺเตน อุคฺฆฏิตญฺญู ปุคฺคโล ติกฺขินฺทฺริโย สสงฺขาร\-ปรินิพฺพายี โหติ, ติกฺขินฺทฺริโย อนฺตรา\-ปรินิพฺพายี โหติ, อุคฺฆฏิตญฺญู มุทินฺทฺริโย อุปหจฺจ\-ปรินิพฺพายี โหติ. วิปญฺจิตญฺญู ติกฺขินฺทฺริโย อสงฺขารปรินิพฺพายี โหติ, วิปญฺจิตญฺญู มุทินฺทฺริโย สสงฺขาร\-ปรินิพฺพายี โหติ, เนยฺโย อุปหจฺจ\-ปรินิพฺพายี โหติ, วิปญฺจิตญฺญู ติกฺขินฺทฺริโย อสงฺขารปรินิพฺพายี โหติ. วิปญฺจิตญฺญู มุทินฺทฺริโย สสงฺขาร\-ปรินิพฺพายี โหติ, เนยฺโย อุทฺธํโสโต อกนิฏฺฐคามี โหติ. อิติ ปญฺจ อนาคามิโน, ฉฏฺโฐ สกทาคามี, ตโย จ โสตาปนฺนาติ อิเม นว เสกฺขา.}

\prose{ตตฺถ อุคฺฆฏิตญฺญู ปุคฺคโล ติกฺขินฺทฺริโย อรหตฺตํ ปาปุณนฺโต เทฺว ปุคฺคลา โหนฺติ อุภโตภาค\-วิมุตฺโต ปญฺญาวิมุตฺโต จ. ตตฺถ อุคฺฆฏิตญฺญู ปุคฺคโล มุทินฺทฺริโย อรหตฺตํ ปาปุณนฺโต เทฺว ปุคฺคลา โหนฺติ, ฐิตกปฺปี\csromansharedfootnote{8be9add65ffc78f2}{ฐิตกปฺปิ (อิ, ก) อภิ 3. 116 ปิฏฺเฐ.} จ ปฏิเวธ\-นภาโว จ. ตตฺถ วิปญฺจิตญฺญู ปุคฺคโล จ ติกฺขินฺทฺริโย โส อรหตฺตํ ปาปุณนฺโต เทฺว ปุคฺคลา โหนฺติ เจตนาภพฺโพ จ รกฺขณาภพฺโพ จ. ตตฺถ วิปญฺจิตญฺญู มุทินฺทฺริโย อรหตฺตํ ปาปุณนฺโต เทฺว ปุคฺคลา โหนฺติ, สเจ เจเตติ น ปรินิพฺพายี, โน เจ เจเตติ ปรินิพฺพายีติ. สเจ อนุรกฺขติ น ปรินิพฺพายี, โน เจ อนุรกฺขติ ปรินิพฺพายีติ. ตตฺถ เนยฺโย ปุคฺคโล ภาวนา\-นุโยคม\-นุยุตฺโต \csromanfolio{190}ปริหานธมฺโม โหติ กมฺมนิยโต วา สมสีสิ วา, อิเม นว อรหนฺโต อิทํ จตุพฺพิธํ สุตฺตํ สํกิเลส\-ภาคิยํ อเสกฺข\-ภาคิยํ. อิเมสุ ปุคฺคเลสุ ตถาคตสฺส ทสวิธํ พลํ ปวตฺตติ.}

\proseitem{18}{กตมํ ทสวิธํ, อิธ พุทฺธานํ ภควนฺตานํ อปฺปวตฺติเต ธมฺมจกฺเก มเหสกฺขา เทวปุตฺตา ยาจนาย อภิยาตา\csromansharedfootnote{f52ad2065eae035d}{อติยาตา (อิ, ก)} โหนฺติ “เทเสตุ สุคโต ธมฺมนฺ”ติ. โส อนุตฺตเรน พุทฺธจกฺขุนา โวโลเกนฺโต อทฺทสาสิ สตฺตานํ ตโย ราสีนํ สมฺมตฺตนิยโต มิจฺฉตฺต\-นิยโต อนิยโต\csromansharedfootnote{9694b7540d7564f7}{ที 3. 182 ปิฏฺเฐ.}. ตตฺถ สมฺมตฺตนิยโต ราสิ มิจฺฉาสติํ อาปชฺเชยฺยาติ เนตํ ฐานํ วิชฺชติ, อสตฺถุโก ปรินิพฺพาเยยฺยาติ เนตํ ฐานํ วิชฺชติ, สมาปตฺติํ อาปชฺเชยฺยาติ ฐานเมตํ วิชฺชติ. ตตฺถ มิจฺฉตฺต\-นิยโต ราสิ อริย\-สมาปตฺติํ ปฏิปชฺชิสฺสตีติ เนตํ ฐานํ วิชฺชติ, อนริย\-มิจฺฉา\-ปฏิปตฺติํ ปฏิปชฺชิสฺสตีติ ฐานเมตํ วิชฺชติ. ตตฺถ อนิยโต ราสิ สมฺมา\-ปฏิปชฺชมานํ สมฺมตฺตนิยต\-ราสิํ คมิสฺสตีติ ฐานเมตํ วิชฺชติ, มิจฺฉา\-ปฏิปชฺชมาโน สมฺมตฺตนิยต\-ราสิํ คมิสฺสตีติ เนตํ ฐานํ วิชฺชติ. สมฺมา\-ปฏิปชฺชมานํ สมฺมตฺตนิยต\-ราสิํ คมิสฺสตีติ ฐานเมตํ วิชฺชติ, มิจฺฉา\-ปฏิปชฺชมานํ มิจฺฉตฺตนิยต\-ราสิํ คมิสฺสตีติ ฐานเมตํ วิชฺชติ. อิเม ตโย อนุตฺตเรน พุทฺธจกฺขุนา โวโลเกนฺตสฺส สมฺมา\-สมฺพุทฺธสฺส เม สโต อิเม ธมฺมา อนภิสมฺพุทฺธาติ เอตฺตวตา มํ โกจิ สหธมฺเมน ปฏิโจทิสฺสตีติ เนตํ ฐานํ วิชฺชติ, วีตราคสฺส เต ปฏิชานโต อขีณาสวตาย สหธมฺเมน โกจิ ปฏิโจทิสฺสตีติ เนตํ ฐานํ วิชฺชติ. ยโต ปน อิมสฺส อนิยตสฺส ราสิสฺส ธมฺมเทสนา, สา น ทิสฺสติ ตกฺกรสฺส สมฺมา\-ทุกฺข\-กฺขยายาติ เนตํ ฐานํ วิชฺชติ, ตถา โอวทิโต ยํ ปน เม อนิยตราสิ สาวโก ปุพฺเพนาปรํ วิเสสํ น สจฺฉิกริสฺสตีติ เนตํ ฐานํ วิชฺชติ.}

\proseitem{19}{ยํ โข มุนิ นานปฺปการสฺส นานานิรุตฺติโย เทว\-นาค\-ยกฺขานํ ทเมติ ธมฺเม ววตฺถาเนน วตฺวา การณโต อญฺญํ ปารํ คมิสฺสตีติ เนตํ ฐานํ วิชฺชติ. ธมฺม\-ปฏิสมฺภิทา. ยโต ปนิมา นิรุตฺติโต สตฺต สตฺต นิรุตฺติโย นาภิสมฺภุเนยฺยาติ เนตํ ฐานํ วิชฺชติ. นิรุตฺติ\-ปฏิสมฺภิทา. นิรุตฺติ โข ปน อภิสม\-คฺค\-รตานํ สาวกานํ ตมตฺถมวิญฺญาปเยติ เนตํ ฐานํ วิชฺชติ. อตฺถ\-ปฏิสมฺภิทา. มเหสกฺขา เทวปุตฺตา อุปสงฺกมิตฺวา ปเญฺห \csromanfolio{191}ปุจฺฉิํสุ. กายิเกน วา มานสิเกน วา ปริปีฬิตสฺส หตฺถกุณีติ วา ปาเท วา ขญฺเช ทนฺธสฺส\csromansharedfootnote{9ce7846810038007}{ทนฺตสฺส (อิ, ก)} โส อตฺโถ น ปริภา\-ชิ\-ยตีติ เนตํ ฐานํ วิชฺชติ. ปฏิภาน\-ปฏิสมฺภิทา. ยมฺหิ ตํ เตสํ โหติ ตมฺหิ อสนฺตํ ภวตีติ เนตํ ฐานํ วิชฺชติ. ยํ หิ นาสํ เตสํ น ภวติ, ตมฺหิ นาสํ เตสํ ภวิสฺสตีติ เนตํ ฐานํ วิชฺชติ. เอวํ สมุทยสฺส นิโรธาย ทส อกุสล\-กมฺมปถา. มาโร วา อินฺโท วา พฺรหฺมา วา ตถาคโต วา จกฺกวตฺตี วา โส วต นาม มาตุคาโม ภวิสฺสตีติ เนตํ ฐานํ วิชฺชติ, ปุริโส อสฺส ราชา จกฺกวตฺตี สกฺโก เทวานมินฺโท ภวิสฺสตีติ ฐานเมตํ วิชฺชติ. อิติสฺส เอวรูปํ พลํ เอวรูปํ ญาณํ, อิทํ วุจฺจติ ฐานาฏฺฐาน\-ญาณํ ปฐมํ ตถาคตพลํ ตํ นิทฺทิสิตพฺพํ. ตีหิ ราสีหิ จตูหิ เวสารชฺเชหิ จตูหิ ปฏิสมฺพิทาหิ ปฏิจฺจ\-สมุปฺปาทสฺส ปวตฺติยํ นิวตฺติยํ ภาคิยญฺจ. กุสลํ กุสลวิปาเกสุ จ อุปปชฺชติ ยญฺจ อิตฺถิปุริสานํ. อิทํ ปฐมํ พลํ ตถาคโต เอวํ ชานาติ.}

\prose{เยสํ ปน สมฺมตฺตนิยโต ราสิ, นายํ สพฺพตฺถ\-คามินี ปฏิปทา, นิพฺพานคามินีเยวายํ ปฏิปทา. ตตฺถ สิยา มิจฺฉตฺตินิยโต ราสิ, เอสาปิ น สพฺพตฺถ\-คามินี ปฏิปทา. สกฺกายสมุทยคามินีเยวายํ ปฏิปทา โหตุ, อยํ ตตฺถ ตตฺถ ปฏิปตฺติยา ฐิโต คจฺฉติ นิพฺพานํ. คจฺฉติ อปายํ. คจฺฉติ เทวมนุสฺสสฺส. ยํ ยํ วา ปฏิปทํ ปฏิปชฺเชยฺย สพฺพตฺถ คจฺเฉยฺย, อยํ สพฺพตฺถ\-คามินี ปฏิปทา. ยํ เอตฺถ ญาณํ ยถาภูตํ, อิทํ วุจฺจติ สพฺพตฺถ\-คามินี ปฏิปทาญาณํ ทุติยํ ตถาคตพลํ.}

\prose{สา โข ปนายํ สพฺพตฺถ\-คามินี ปฏิปาทา นานาธิมุตฺตา เกจิ กาเมสุ เกจิ ทุกฺกรการิยํ เกจิ อตฺตกิลมถา\-นุโยคมนุยุตฺตา เกจิ สํสาเรน สุทฺธิํ ปจฺเจนฺติ เกจิ อนชฺชาภาวนาติ. เตน เตน จริเตน วินิพนฺธานํ สตฺตานํ ยํ ญาณํ ยถาภูตํ นานาคตํ โลกสฺส อเนกา\-ธิ\-มุตฺตคตํ ยถาภูตํ ปชานาติ. อิทํ ตติยํ ตถาคตพลํ.}

\prose{ตตฺถ สตฺตานํ อธิมุตฺตา ภวนฺติ อาเสวนฺติ ภาเวนฺติ พหุลีกโรนฺติ. เตสํ กมฺมุปสยานํ ตทา\-ธิมุตฺตานํ. สา เจว ธาตุ สํวหติ. กตรา ปเนสา ธาตุ เนกฺขมฺมธาตุ พลธาตุ กาจิ สมฺปตฺติ กาจิ มิจฺฉตฺตญฺจ ธาตุ อธิมุตฺตา ภวนฺติ. อญฺญตรา อุตฺตริ น สมนุปสฺสนฺติ. เต ตเทวฏฺฐานํ \csromanfolio{192}มหา ชรามรณสฺส อภินิวิสฺส โวหรนฺติ “อิทเมว สจฺจํ โมฆมญฺญนฺ”ติ. ยถา ภควา สกฺกสฺส เทวานมินฺทสฺส ภาสิตํ. ยํ ตตฺถ ยถาภูตํ ญาณํ. อิทํ วุจฺจติ จตุตฺถํ ตถาคตพลํ.}

\proseitem{2}{ตตฺถ ยํเยว ธาตุ\csromansharedfootnote{ec0830f2150fd8a0}{ยํ ยเทว ธาตุํ (ก)} เสฏฺฐนฺติ ตํ ตํ กาเยน จ วาจาย จ อารมฺภนฺติ เจตสิโก. อารมฺโภ เจตนา กมฺมํ กายิกา วาจสิกา อารมฺโภ เสตสิกตฺตา กมฺมนฺตรํ ตถาคโต เอวํ ปชานาติ “อิมินา สตฺเตน เอวํ ธาตุเกน เอวรูปํ กมฺมํ กตํ, ตํ อตีตมทฺธานํ อิมินา เหตุนา ตสฺส เอวรูโป วิปาโก วิปจฺจติ เอตรหิ วิปจฺจิสฺสติ วา อนาคตมทฺธานนฺติ. เอวํ ปจฺจุปฺปนฺน\-ม\-ทฺธานํ ปชานาติ “อยํ ปุคฺคโล เอวํธาตุโก อิทํ กมฺมํ กโรติ. ตณฺหาย จ ทิฏฺฐิยา จ อิมินา เหตุนา น ตสฺส วิปาโก ทิฏฺเฐเยว ธมฺเม นิพฺพตฺติสฺสติ, อุปปชฺเช วา”ติ อปรมฺหิ วา ปริยาเย เอวํ ปชานาติ อยํ ปุคฺคโล เอวรูปํ กมฺมํ กริสฺสติ อนาคตมทฺธานํ, อิมินา เหตุนา ตสฺส เอวรูโป วิปาโก นิพฺพตฺติสฺสติ, อิมินา เหตุนา ยานิ จตฺตาริ กมฺมฏฺฐานานิ อิทํ กมฺมฏฺฐานํ ปจฺจุปฺปนฺนสุขํ อายติํ จ สุขวิปากํ ฯเปฯ. อิติ อยํ อตีตา\-นาคต\-ปจฺจุปฺปนฺนานํ กมฺม\-สมาทานานํ เหตุโส ฐานโส วิปาก\-เวมตฺตตํ ปชานาติ อุจฺจาวจา หีนปณีตตา, อิทํ วุจฺจติ กมฺมวิปาก\-ญาณํ ปญฺจมํ ตถาคตพลํ.}

\prose{ตถา สตฺตา ยํ วา กมฺมสมาทานํ สมาทิยนฺตา ตตฺถ เอวํ ปชานาติ อิมสฺส ปุคฺคลสฺส กมฺมา\-ธิมุตฺตสฺส ราคจริตสฺส เนกฺขมฺม\-ธาตูนํ ปาริปูริํ คจฺฉนฺติ, ตสฺส ราคานุคเต สุญฺญมานสฺส ปฐมํ ฌานํ สํกิลิสฺสติ, สเจ ปุน อุตฺตริ วายามโต ฌาน\-โวทาน\-คเต มานเส วิเสสภาคิยํ ปฏิปทํ อนุยุญฺชิยติ. ตสฺส หิ ฌานภาคิยํเยว ปฐมชฺฌาเน ฐิตสฺส ทุติยํ ฌานํ โวทานํ คจฺฉติ, ตติยญฺจ ฌานํ สมาปชฺชิตุ\-กามสฺส โสมนสฺสิ\-นฺทฺริยํ จิตฺตํ ปริยาทาย ติฏฺฐติ, ตสฺส สา ปีติ อวิเสส\-ภาคิยํ ตติยํ ฌานํ อาทิสฺส ติฏฺฐติ. สเจ ตสฺส นิสฺสรณํ ยถาภูตํ ปชานาติ. ตถาคตสฺส จตุตฺถ\-ชฺฌานํ โวทานํ คจฺฉติเยว, จตุตฺถสฺส ฌานสฺส หานภาคิยา ธมฺมา, เต จ ธมฺมา ยตฺถ ปชายนฺติ เยหิ จตุตฺถ\-ชฺฌานํ โวทานํ ทิสฺสติ. เอวํ อชฺฌาสย\-สมาปตฺติยา ยา จตสฺโส สมาปตฺติโย ตีณิ วิโมกฺขมุขานิ อฏฺฐ วิโมกฺข\-ฌานานีติ จตฺตาริ ฌานานิ วิโมกฺขาติ. อฏฺฐ จ วิโมกฺขา ตีณิ จ วิโมกฺขมุขานิ. สมาธีติ \csromanfolio{193}จตฺตาโร สมาธี ฉนฺทสมาธิ วีริยสมาธิ จิตฺตสมาธิ วีมํสาสมาธีติ, สมาปตฺติโย จตสฺโส อชฺฌาสย\-สมาปตฺติโย อิติ อิเมสํ ฌานานํ วิโมกฺข\-สมาปตฺตีติ เอวรูโป สํกิเลโส ราคจริตสฺส ปุคฺคลสฺส. เอวํ โทสจริตสฺส. โมหจริตสฺส. ราคจริตสฺส ปุคฺคลสฺส เอวรูปํ โวทานํ อิติ ยํ เอตฺถ ญาณํ ยถาภูตํ อสาธารณํ สพฺพสตฺเตหิ. อิทํ วุจฺจติ ฉฏฺฐํ ตถาคตพลํ.}

\prose{ตตฺถ ตถาคโต เอวํ ปชานาติ โลกิกา ธมฺมา โลกุตฺตรา ธมฺมา ภาวนา\-ภาคิยํ อินฺทฺริยํ นามํ ลภนฺติ. อาธิปเตยฺยภูมิํ อุปาทาย พลํ นามํ ลภนฺติ ถามคตํ มโน มนินฺทฺริยํ ตํ อุปาทาย. วีริยํ นามํ ลภนฺติ อารมฺภธาตุํ อุปาทาย. อิติสฺส เทว เอวรูปํ ญาณํ อิเมหิ จ ธมฺเมหิ อิเม ปุคฺคลา สมนฺนาคตาติปิ ธมฺมเทสนํ อกาสิ. อาการโต จ โวการโต จ อาสย\-ชฺฌาสยสฺส อธิมุตฺติ\-สมนฺนาคตานํ. อิทํ วุจฺจติ ปรสตฺตานํ ปรปุคฺคลานํ อินฺทฺริย\-พลวีริย\-เวมตฺตตํ ญาณํ สตฺตมํ ตถาคตพลํ.}

\prose{ตตฺถ จ ตถาคโต โลกาทีสุ จ ภูมีสุ สํโยชนานญฺจ เสกฺขานํ ทฺวีหิ พเลหิ คติํ ปชานาติ, ปุพฺเพ\-นิวาสา\-นุสฺสติยา อตีเต สํสาเร เอตรหิ จ ปจฺจุปฺปนฺเน ทิพฺพจกฺขุนา จุตูปปาตํ อิติ อิมานิ เทฺว พลานิ ทิพฺพจกฺขุโต อภินีหิตานิ. โส อตีตมทฺธานํ ทิพฺพสฺส จกฺขุโน โคจโร โส เอตรหิ สติ โคจโร อิติ อตฺตโน จ ปเรสํ จ ปุพฺเพนิวาสญาณํ อเนกวิธํ นานปฺปการกํ ปจฺจุปฺปนฺน\-ม\-ทฺธานํ ทิพฺเพน จกฺขุนา อิมานิ เทฺว ตถาคต\-พลานิ, อฏฺฐมํ ปุพฺเพนิวาโส, นวมํ ทิพฺพจกฺขุ.}

\prose{ปุน จปรํ ตถาคโต อริย\-ปุคฺคลานํ ฌานํ โวทานํ นิพฺเพธ\-ภาคิยํ ปชานาติ อยํ ปุคฺคโล อิมินา มคฺเคน อิมาย ปฏิปทาย อาสวานํ ขยา อนาสวํ เจโตวิมุตฺติํ ปญฺญาวิมุตฺติํ ทิฏฺเฐว ธมฺเม สจฺฉิกตฺวา อุปสมฺปชฺช วิหรตีติ อิติ อตฺตโน จ อาสวานํ ขยํ ญาณํ ทิฏฺเฐ\-กฏฺฐานํ จตุภูมิมุ\-ปาทาย ยาว นวนฺนํ อรหนฺตานํ อาสวกฺขโย โอธิโส เสกฺขานํ อโนธิโส อรหนฺตานํ. ตตฺถ เจโตวิมุตฺติ ทฺวีหิ อาสเวหิ อนาสวา กามาสเวน จ ภวาสเวน จ, ปญฺญาวิมุตฺติ ทฺวีหิ อาสเวหิ อนาสวา ทิฏฺฐาสเวน จ อวิชฺชาสเวน จ, อิมาสํ ทฺวินฺนํ วิมุตฺตีนํ ยถาภูตํ ญาณํ, อิทํ วุจฺจติ อาสวกฺขเย ญาณํ. ทสมํ ตถาคตพลํ.}

\proseitem{20}{\csromanfolio{194}อิเมสุ ทสสุ พเลสุ ฐิโต ตถาคโต ปญฺจวิธํ สาสนํ เทเสติ สํกิเลส\-ภาคิยํ วาสนาภาคิยํ ทสฺสน\-ภาคิยํ ภาวนา\-ภาคิยํ อเสกฺข\-ภาคิยํ. ตตฺถ โย ตณฺหาสํกิเลโส, อิมสฺส อโลโภ นิสฺสรณํ. โย ทิฏฺฐิ\-สํกิเลโส, อิมสฺส อโมโห นิสฺสรณํ. โย ทุจฺจริต\-สํกิเลโส, อิมสฺส ตีณิ กุสลานิ นิสฺสรณํ. กิํ นิทานํ, ตีณิ อิมานิ\csromansharedfootnote{4c1c83275f090d5d}{ตีณิ หิ อิมานิ (อิ)} มโน\-ทุจฺจริตานิ อภิชฺฌา พฺยาปาโท มิจฺฉาทิฏฺฐิ. ตตฺถ อภิชฺฌฺยา มโนทุจฺจริตํ กายกมฺมํ อุปฏฺฐเปติ, อทินฺนาทานํ สพฺพญฺจ ตทุปนิพฺพทฺธํ วาจากมฺมํ อุปฏฺฐเปติ, มุสาวาทญฺจ สพฺพวิตถํ สพฺพํ วาจมภาวํ สพฺพมกฺขํ ปลาสํ อภิชฺฌา อกุสลมูลนฺติ, สุจริเต สุจริตํ มุสาวาทา อทินฺนาทานา อภิชฺฌาย เจตนา, ตตฺถ พฺยาปาโท มโนทุจฺจริตํ กายกมฺมํ อุปฏฺฐเปติ, ปาณาติปาตํ สพฺพญฺจ เมตํ อากฑฺฒนํ ปริกฑฺฒนํ นิพฺพทฺธํ โรจนํ วาจากมฺมํ อุปฏฺฐเปติ, ปิสุณวาจํ ผรุสวาจํ มิจฺฉาทิฏฺฐิ มโนทุจฺจริตญฺจ อภิชฺฌํ พฺยาปาทํ มิจฺฉาทิฏฺฐิํ ปโยเชติ, ตสฺส โย โกจิ มิจฺฉาทิฏฺฐิ จาโค ราคโช วา โทสโช วา สพฺพโส มิจฺฉาทิฏฺฐิ สมฺภูโต อิมินา การเณน มิจฺฉาทิฏฺฐิํ อุปฏฺฐเปติ, กาเมสุ\-มิจฺฉาจารํ วจีกมฺมํ อุปฏฺฐเปติ สมฺผปฺปลาปํ. อิมานิ ตีณิ ทุจฺจริตานิ อกุสลมูลานิ.}

\prose{ยา อภิชฺฌา, โส โลโภ. โย พฺยาปาโท, โสโทโส. ยา มิจฺฉาทิฏฺฐิ, โส โมโห. ตานิ อฏฺฐ มิจฺฉตฺตานิ อุปฏฺฐเปนฺติ. เตสุ คหิเตสุ ตีสุ อกุสลมูเลสุ ทสวิธํ อกุสลมูลํ ปาริปูริํ คจฺฉติ, ตสฺส ติวิธสฺส ทุจฺจริต\-สํกิเลสสฺส วาสนาภาคิยญฺจ สุตฺตํ นิสฺสรณํ. ตตฺถ โย พหุสิโต นิทฺเทโส ยถา โลโภ โทโส โมโหปิ, ตตฺถ อสิตุํ เอตฺถ โลโภ อุสฺสโท เตน การเณน เตสุ วา ธมฺเมสุ โลโภ ปญฺญปิยติ. ตตฺถายํ โมโห อกุสลํ โมโห อยํ อวิชฺชา, สา จตุพฺพิธา รูเป อภินิวิฏฺฐา, รูปํ อตฺตโต สมนุปสฺสติ, อวิชฺชาคโต รูปวนฺตํ อตฺตานํ, อตฺตนิ วา รูปํ, รูปสฺมิํ วา อตฺตานํ. ตตฺถ กตมํ ปทํ สกฺกาย\-ทิฏฺฐิยา อุจฺเฉทํ วทติ “ตํ ชีวํ ตํ สรีรนฺ”ติ นตฺถิกทิฏฺฐิ อธิจฺจสมุปฺปนฺน\-ทิฏฺฐิ จ อญฺโญ จ กโรติ, อญฺโญ ปฏิสํเวทิยติ. ปจฺฉิม\-สฏฺฐิ\-กปฺปานํ ตีณิ ปทานิ สกฺกาย\-ทิฏฺฐิยา สสฺสตํ ภชนฺติ “อญฺญํ ชีวํ อญฺญํ สรีรนฺ”ติ อกิริยญฺจ ตํ ทุกฺขมิจฺฉโต อเหตุกา จ ปตนฺติ อนชฺฌาภาโว จ กมฺมานํ สพฺพญฺจ มานยิ\csromansharedfootnote{befbeca378407b5b}{มานติ (อิ)}. ตตฺถ “อิทเมว สจฺจํ โมฆมญฺญนฺ”ติ \csromanfolio{195}สํสาเรน สุทฺธิ อาชีวกา ฉฬาสีติ ปญฺญเปนฺติ. ยถารูเป สกฺกาย\-ทิฏฺฐิยา จตุวตฺถุกา, เอวํ ปญฺจสุ ขนฺเธสุ วีสติวตฺถุกา สกฺกาย\-ทิฏฺฐิยา สสฺสตํ ภชติ. อญฺญาชีวกา จ สสฺสตวาทิเก จ สีลพฺพตํ ภชนฺติ ปรามสนฺติ อิมินา ภวิสฺสามิ เทโว วา เทวญฺญตโร วา, อยํ สีลพฺพต\-ปรามาโส. ตตฺถ สกฺกาย\-ทิฏฺฐิยา โส รูปํ อตฺตโต สมนุปสฺสติ, “ตํ ชีวํ ตํ สรีรมฺ”อิติ ตํ กงฺขติ วิจิกิจฺฉติ นาธิมุจฺจติ นาภิปฺปสีทติ ปุพฺพนฺเต อปรนฺเต ปุพฺพนฺตา\-ปรนฺเต ฯเปฯ อิติ วาสนา\-ภาคิเยสุ ฐิตสฺส อยํ อุปกฺกิเลโส.}

\proseitem{21}{ตตฺถ สทฺธินฺทฺริเยน สพฺพํ วิจิกิจฺฉิตํ ปชหติ, ปญฺญินฺทฺริเยน อุทยพฺพยํ ปสฺสติ, สมาธิ\-นฺทฺริเยน จิตฺตํ เอโกทิ กโรติ วีริยินฺทฺริเยน อารภติ. โส อิเมหิ ปญฺจหิ อินฺทฺริเยหิ สทฺธานุสารี อเวจฺจปฺปสาเท นิรโต อนนฺตริยํ สมาธิํ อุปฺปาเทติ. อินฺทฺริเยหิ สุทฺเธหิ ธมฺมานุสารี อปฺปจฺจยตาย อนนฺตริยํ สมาธิํ อุปฺปาเทติ. โส “อิทํ ทุกฺขนฺ”ติ ยถาภูตํ ปชานาติ. สจฺจานิ อิทํ ทสฺสน\-ภาคิยํ สุตฺตํ. ตสฺส ปญฺจนฺนํ โอรมฺภาคิยานํ สํโยชนานํ ตีณิ สํโยชนานิ ทสฺสน\-ปหาตพฺพานิ สพฺเพน สพฺพํ ปหีนานิ เทฺว ปุคฺคลกตานิ. ตตฺถ ตีณิ อกุสลมูลานิ. ภาวนา\-ปหาตพฺพานิ อุปริกฺขิตฺตานิ ฉ ภเว นิพฺพตฺเตนฺติ. ตตฺถ เตสุ อภิชฺฌาย จ พฺยาปาเทสุ ตนุกเตสุ ฉ ภวา ปริกฺขยา มริยาทํ คจฺฉนฺติ, เทฺว ภวา อวสิฏฺฐา. ตสฺส อภิชฺฌา จ พฺยาปาโท จ สพฺเพน สพฺพํ ปริกฺขีณา โหนฺติ. เอโก ภโว อวสิฏฺโฐ โหติ. โส จ มานวเสน นิพฺพตฺเตติ. กิญฺจาปิ เอตฺถ อญฺเญปิ จตฺตาโร กิเลสา รูปราโค ภวราโค อวิชฺชา อุทฺธจฺจํ เกตุสฺมิมา\-น\-ภูตา นปฺปฏิพลา อสฺมิมานํ วินิวตฺเตตุํ, สพฺเพปิ เต อสฺมิมานสฺส ปหานํ อารภเต. ขีเณสุ น จ เตสุ อิทมุ\-ตฺตริ\-ทสฺสน\-ภูมิยํ ปญฺจสุ เสกฺข\-ปุคฺคเลสุ ตีสุ จ ปฏิปฺปนฺนเกสุ ทฺวีสุ จ ผลฏฺเฐสุ ภาวนา\-ภาคิยํ สุตฺตํ. ตทุตฺตริ อเสกฺขภาคิย\-สุตฺตํ, กตฺถจิ ภูมิ นิปีฬิยติ. อิทญฺจ ปญฺจมํ สุตฺตํ. ติณฺณํ ปุคฺคลานํ เทสิตํ ปุถุชฺชนสฺส เสกฺขสฺส อเสกฺขสฺส สํกิเลส\-ภาคิยํ วาสนาภาคิยํ. ปุถุชฺชนสฺส ทสฺสน\-ภาคิยํ. ภาวนา\-ภาคิยํ ปญฺจนฺนํ เสกฺขานํ. ยํ ปฐม\-นิทฺทิฏฺฐํ อเสกฺข\-ภาคิยํ สพฺเพสํ อรหนฺตานํ. สา ปน ปญฺจวิธา สตฺตวีส- อากาเร\csromansharedfootnote{a9f4cf252a2e537b}{สตฺตวีสํ อากาเร (อิ)} ปริเยสิตพฺพํ. เอเตสุ ตสฺส คตีนํ ตโต อุตฺตริ. ตญฺจ โข สงฺเขเปน ปญฺญาสาย \csromanfolio{196}อากาเรหิ สมฺปตติ, เย ปญฺญาส อาการา สาสเน นิทฺทิฏฺฐา, เต สงฺขิปิยนฺตา ทสหิ อากาเรหิ ปตนฺติ. เย อริยสจฺจํ นิกฺเขเปน ฐิเต สงฺขิปิยตฺตา อฏฺฐสุ อากาเรสุ ปตนฺติ. จตูสุ จ สาธารเณสุ สุตฺเตสุ ยา หารสมฺปาตสฺส ภูมิ, เต สงฺขิปิยนฺตา ปญฺจสุ สุตฺเตสุ ปตนฺติ. สํกิเลส\-ภาคิเย วาสนาภาคิเย ภาวนาภาคิเย นิพฺเพธ\-ภาคิเย อเสกฺขภาคิเย จ. เต สงฺขิปิยนฺตา จตูสุ สุตฺเตสุ ปตนฺติ. สํกิเลส\-ภาคิเย วาสนาภาคิเย นิพฺเพธ\-ภาคิเย อเสกฺขภาคิเย จ. เต สงฺขิปิยมานา ตีสุ สุตฺเตสุ ปตนฺติ, ปุถุชฺชน\-ภาคิเย เสกฺขภาคิเย อเสกฺขภาคิเย จ. เต สงฺขิปิยนฺตา ทฺวีสุ สุตฺเตสุ ปตนฺติ นิพฺเพธ\-ภาคิเย จ ปุพฺพโยค\-ภาคิเย จ. ยถา วุตฺตํ ภควตา เทฺว อตฺถวเส ยมฺปสฺสมานา ตถาคตา อรหนฺโต สมฺมาสมฺพุทฺธา ธมฺมํ เทเสนฺติ สุตฺตํ เคยฺยํ ฯเปฯ สตฺถา ปุพฺพโยค\-สมนฺนาคเต อปฺปกสิเรน มญฺญมานา วสิยนฺติ ปุพฺพโยคา จ ภวิสฺสนฺติ สนฺตานํ มญฺญมานา\-ธราย. ตตฺถ ปญฺญา\-เวมตฺตตํ อตฺตโน สมนุ\-ปสฺส\-มาเนน อฏฺฐวิเธ สุตฺตสงฺเขเป, ยตฺถ ยตฺถ สกฺโกติ, ตตฺถ ตตฺถ โยเชตพฺพํ. ตตฺถ ตตฺถ โยเชตฺวา สุตฺตสฺส อตฺโถ นิทฺทิสิตพฺโพ. น หิ สติ เวทนา มโน ธาเรตฺวา สกฺกา เยน เกนจิ สุตฺตสฺส อตฺโถ ยถาภูตํ นิทฺทิสิตุํ.}

\csromanheader{2}{\textbf{ตตฺถ ปุริมกานํ สุตฺตานํ อิมา อุทฺทานคาถา}}
\setlength{\csromangathapagewidth}{0pt}
\setlength{\csromangathawakwidth}{0pt}
\csromangathameasuregroup{กามนฺธา ชาลสญฺฉนฺนา, \\ มโนปุพฺพงฺคมา ธมฺมา, \\ อุทฺธํ อโธ วิปฺปมุตฺโต, \\ ยสฺส เสลูปมํ จิตฺตํ, \\ ยสฺส กายคตาสติ, \\ น ตํ ทฬฺหํ เจตสิกํ, \\ จตฺตาโร เจว ปุคฺคลา,}{\csromangathabat{กามนฺธา ชาลสญฺฉนฺนา,}{ปญฺจ นีวรณานิ จ.} \\ \csromangathabat{มโนปุพฺพงฺคมา ธมฺมา,}{มหานาโม จ สากิโย.} \\ \csromangathabat{อุทฺธํ อโธ วิปฺปมุตฺโต,}{ยญฺจ สีลกิมตฺถิยา.} \\ \csromangathabat{ยสฺส เสลูปมํ จิตฺตํ,}{อุปติสฺส ปุจฺฉาทิกา.} \\ \csromangathabat{ยสฺส กายคตาสติ,}{ฉนฺนํ ตโมปรายโณ.} \\ \csromangathabat{น ตํ ทฬฺหํ เจตสิกํ,}{อยํ โลโกติอาทิกํ.} \\ \csromangathabat{จตฺตาโร เจว ปุคฺคลา,}{ททโต ปุญฺญํ ปวฑฺฒิตํ.}}
\csromangathasetleft{กามนฺธา ชาลสญฺฉนฺนา, \\ มโนปุพฺพงฺคมา ธมฺมา, \\ อุทฺธํ อโธ วิปฺปมุตฺโต, \\ ยสฺส เสลูปมํ จิตฺตํ, \\ ยสฺส กายคตาสติ, \\ น ตํ ทฬฺหํ เจตสิกํ, \\ จตฺตาโร เจว ปุคฺคลา,}
\csromangathagroup{\csromangathastanza{\csromangathabat{กามนฺธา ชาลสญฺฉนฺนา,}{ปญฺจ นีวรณานิ จ.} \\ \csromangathabat{มโนปุพฺพงฺคมา ธมฺมา,}{มหานาโม จ สากิโย.}}\csromangathastanza{\csromangathabat{อุทฺธํ อโธ วิปฺปมุตฺโต,}{ยญฺจ สีลกิมตฺถิยา.} \\ \csromangathabat{ยสฺส เสลูปมํ จิตฺตํ,}{อุปติสฺส ปุจฺฉาทิกา.}}\csromangathastanza{\csromangathabat{ยสฺส กายคตาสติ,}{ฉนฺนํ ตโมปรายโณ.} \\ \csromangathabat{น ตํ ทฬฺหํ เจตสิกํ,}{อยํ โลโกติอาทิกํ.} \\ \csromangathabat{จตฺตาโร เจว ปุคฺคลา,}{ททโต ปุญฺญํ ปวฑฺฒิตํ.}}}

\csromanheader{2}{โสตานุคต\-ธมฺเมสุ, อิมา เตสํ อุทฺทานคาถา.}
\setlength{\csromangathapagewidth}{0pt}
\setlength{\csromangathawakwidth}{0pt}
\csromangathameasuregroup{}{ตตฺถ กตมา อาณตฺติ. \\ สเจ ภายถ ทุกฺขสฺส, สเจ โว ทุกฺขมปฺปิยํ. \\ มากตฺถ ปาปกํ กมฺมํ, อาวิ วา ยทิ วา รโห\textsuperscript{1}.}
\csromangathameasurewak{ตตฺถ กตมา อาณตฺติ. \\ สเจ ภายถ ทุกฺขสฺส, สเจ โว ทุกฺขมปฺปิยํ. \\ มากตฺถ ปาปกํ กมฺมํ, อาวิ วา ยทิ วา รโห\textsuperscript{1}.}
\setlength{\csromangathaleftcol}{0pt}
\csromangathagroup{\csromangathastanza{\csromangathawak{\csromangathaitemline{22}{ตตฺถ กตมา อาณตฺติ.} \\ \csromangathacontline{สเจ ภายถ ทุกฺขสฺส, สเจ โว ทุกฺขมปฺปิยํ.} \\ \csromangathacontline{มากตฺถ ปาปกํ กมฺมํ, อาวิ วา ยทิ วา รโห\csromansharedfootnote{7e24d6a1754ec759}{ขุ 1. 137; เหฏฺฐา 109 ปิฏฺเฐสุปิ.}.}}}}

\prose{\csromanfolio{197}อตีเต ราธ รูเป อนเปกฺโข โหหีติ วิตฺถาเรน กาตพฺพา.  สีลวนฺเตน อานนฺท ปุคฺคเลน สทา กรณียา กินฺติเม อวิปฺปฏิสาโร อสฺสาติ. อยํ วุจฺจติ อาณตฺติ.}

\prose{ตตฺถ กตมํ ผลํ.}

\setlength{\csromangathapagewidth}{0pt}
\setlength{\csromangathawakwidth}{0pt}
\csromangathameasuregroup{}{\textsuperscript{*}ธมฺโม หเว รกฺขติ ธมฺมจาริํ, \\ ฉตฺตํ มหนฺตํ ยถ วสฺสกาเล. \\ เอสานิสํโส ธมฺเม สุจิณฺเณ, \\ น ทุคฺคติํ คจฺฉติ ธมฺมจารี.}
\csromangathameasurewak{\textsuperscript{*}ธมฺโม หเว รกฺขติ ธมฺมจาริํ, \\ ฉตฺตํ มหนฺตํ ยถ วสฺสกาเล. \\ เอสานิสํโส ธมฺเม สุจิณฺเณ, \\ น ทุคฺคติํ คจฺฉติ ธมฺมจารี.}
\setlength{\csromangathaleftcol}{0pt}
\csromangathagroup{\csromangathastanza{\csromangathawak{\csromansymbolfootnote{*}{เหฏฺฐา 7, 29, 38 ปิฏฺเฐสุปิ.}ธมฺโม หเว รกฺขติ ธมฺมจาริํ, \\ ฉตฺตํ มหนฺตํ ยถ วสฺสกาเล. \\ เอสานิสํโส ธมฺเม สุจิณฺเณ, \\ น ทุคฺคติํ คจฺฉติ ธมฺมจารี.}}}

\prose{อิทํ ผลํ.}

\csromanhangingbegin
\hangingheadline{ตตฺถ กตโม อุปาโย.}
\hangingbody{“สพฺเพ ธมฺมา อนตฺตา”ติ, ยทา ปญฺญาย ปสฺสติ.}
\hangingbody{อถ นิพฺพินฺทติ ทุกฺเข, เอส มคฺโค วิสุทฺธิยา1.}
\csromanhangingend

\prose{สตฺตหงฺเคหิ สมนฺนาคโต โข ภิกฺขุ อปิ หิมวนฺตํ ปพฺพตราชานํ จาเลยฺย, โก ปน วาโท ฉวํ อวิชฺชํ สตฺตเกสุ เวยฺยากรณํ กาตพฺพํ. อยํ อุปาโย.}

\csromanhangingbegin
\hangingheadline{ตตฺถ กตมา อาณตฺติ จ ผลญฺจ.}
\hangingbody{สเจ ภายถ ทุกฺขสฺส, สเจ โว ทุกฺขมปฺปิยํ.}
\hangingbody{มากตฺถ ปาปกํ กมฺมํ, อาวิ วา ยทิ วา รโห.}
\hangingbody{สเจ หิ ปาปกํ กมฺมํ, กโรถ วา กริสฺสถ.}
\hangingbody{น โว ทุกฺขา ปโมกฺขาตฺถิ, อุปจฺจาปิ ปลายตํ2.}
\csromanhangingend

\prose{ปุริมิกาย คาถาย อาณตฺติ ปจฺฉิมิกาย ผลํ.  สีเล ปติฏฺฐาย เทฺว ธมฺมา ภาเวตพฺพา ยา จ จิตฺตภาวนา ยา จ ปญฺญาภาวนา ยา จ อาณตฺติ ราควิราคา จ ผลํ.}

\csromanhangingbegin
\hangingheadline{ตตฺถ กตมํ ผลญฺจ อุปาโย จ.}
\hangingbody{สีเล ปติฏฺฐาย นโร สปญฺโญ, จิตฺตํ ปญฺญญฺจ ภาวยํ.}
\hangingbody{อาตาปี นิปโก ภิกฺขุ, โส อิมํ วิชฏเย ชฏํ3.}
\csromanhangingend

\prose{ปุริมิกาย อฑฺฒคาถาย อุปาโย, ปจฺฉิมิกาย อฑฺฒคาถาย ผลํ. นนฺทิโย\csromansharedfootnote{ebc145709c390d7a}{นนฺทิโก (อิ, ก)} สกฺโก อิสิวุตฺถปุริริกามเอกรกฺเข\csromansharedfootnote{5f55808c1aa96efe}{อิสิวุตฺต... (อิ)} สุตฺตํ มูลโต อุปาทาย}

\prose{\csromanfolio{198}ยาว ฉสุ ธมฺเมสุ. อุตฺตริ ปญฺจสุ ธมฺเมสุ ยาจโยโค\csromansharedfootnote{00f240d92a81a21e}{โย จ โยโค (อิ)} กรณีโย, อยํ อุปาโย. อสหคตสฺส กามาสวาปิ จิตฺตํ มุจฺจตีติ. สพฺพาสุ ฉสุ ตีสุ. อยํ อุปาโย จ ผลญฺจ.}

\csromanhangingbegin
\hangingheadline{ตตฺถ กตมา อาณตฺติ จ ผลญฺจ อุปาโย จ.}
\hangingbody{สุญฺญโต โลกํ อเวกฺขสฺสุ, โมฆราช สทา สโต.}
\hangingbody{อตฺตานุทิฏฺฐิํ อุหจฺจ2, เอวํ มจฺจุตโร สิยา.}
\csromanhangingend

\prose{สุญฺญโต โลกํ อเวกฺขสฺสุ, โมฆราชาติ อาณตฺติ. สทา สโตติ อุปาโย. อตฺตานุทิฏฺฐิํ อุหจฺจ, เอวํ มจฺจุตโร สิยาติ ผลํ.  สมาธิํ ภิกฺขเว ภาเวถ, สมาหิโต ภิกฺขเว ภิกฺขุ รูปํ อนิจฺจนฺติ ปชานาติ. เอวํ ปสฺสํ อริยสาวโก ปริมุจฺจติ ชาติยาปิ ฯเปฯ อุปายาเสหิปิ อิธ ตีณิปิ.}

\proseitem{23}{ตตฺถ กตโม อสฺสาโท.}

\prose{กามํ กามยมานสฺส, ตสฺส เจตํ สมิชฺฌติ. อยํ อสฺสาโท.}

\prose{ธมฺมจริยา สมจริยา กุสลจริยา เหตูหิ พฺราหฺมณ เอวมิเธกจฺเจ สตฺตา กายสฺส เภทา สุคติํ สคฺคํ โลกํ อุปปชฺชนฺติ. อยํ อสฺสาโท.}

\prose{ตตฺถ กตโม อาทีนโว.}

\prose{กาเมสุ เว หญฺญเต สพฺพา มุจฺเจว. อยํ อาทีนโว. ปเสนทิ\-สํยุตฺตเก สุตฺเต ปพฺพโตปมา. อยํ อาทีนโว.}

\prose{ตตฺถ กตมํ นิสฺสรณํ.}

\setlength{\csromangathapagewidth}{0pt}
\setlength{\csromangathawakwidth}{0pt}
\csromangathameasuregroup{\textsuperscript{*}โย กาเม ปริวชฺเชติ, \\ โสมํ วิสตฺติกํ โลเก,}{\csromangathabat{\textsuperscript{*}โย กาเม ปริวชฺเชติ,}{สปฺปสฺเสว ปทา สิโร.} \\ \csromangathabat{โสมํ วิสตฺติกํ โลเก,}{สโต สมติวตฺตติ.}}
\csromangathasetleft{\textsuperscript{*}โย กาเม ปริวชฺเชติ, \\ โสมํ วิสตฺติกํ โลเก,}
\csromangathagroup{\csromangathastanza{\csromangathaitembat{23}{\csromansymbolfootnote{*}{เหฏฺฐา 6, 59 ปิฏฺเฐสุ.}โย กาเม ปริวชฺเชติ,}{สปฺปสฺเสว ปทา สิโร.} \\ \csromangathacontbat{โสมํ วิสตฺติกํ โลเก,}{สโต สมติวตฺตติ.}}}

\prose{สํยุตฺตเก สุตฺตํ ปาริจฺฉตฺตโก ปณฺฑุปลาโส สนฺนิปลาโส. อิทํ นิสฺสรณํ.}

\prose{\csromanfolio{199}ตตฺถ กตโม อสฺสาโท จ อาทีนโว จ.}

\setlength{\csromangathapagewidth}{0pt}
\setlength{\csromangathawakwidth}{0pt}
\csromangathameasuregroup{ยานิ กโรติ ปุริโส, \\ กลฺยาณการี กลฺยาณํ,}{\csromangathabat{ยานิ กโรติ ปุริโส,}{ตานิ อตฺตนิ ปสฺสติ.} \\ \csromangathabat{กลฺยาณการี กลฺยาณํ,}{ปาปการี จ ปาปกํ\textsuperscript{1}.}}
\csromangathasetleft{ยานิ กโรติ ปุริโส, \\ กลฺยาณการี กลฺยาณํ,}
\csromangathagroup{\csromangathastanza{\csromangathabat{ยานิ กโรติ ปุริโส,}{ตานิ อตฺตนิ ปสฺสติ.} \\ \csromangathabat{กลฺยาณการี กลฺยาณํ,}{ปาปการี จ ปาปกํ\csromansharedfootnote{6ee0db9176153c98}{ขุ 5. 60 ทุกนิปาเต.}.}}}

\prose{ตตฺถ ยํ ปาปการี ปจฺจนุโภติ อยํ อสฺสาโท. ลาภาลาภ\-อฏฺฐเกสุ พฺยากรณํ, ตตฺถ อลาโภ อยโส นินฺทา ทุกฺขํ, อยํ อาทีนโว. ลาโภ ยโส สุขํ ปสํสา, อยํ อสฺสาโท.}

\prose{ตตฺถ กตมํ อสฺสาโท จ นิสฺสรณญฺจ.}

\setlength{\csromangathapagewidth}{0pt}
\setlength{\csromangathawakwidth}{0pt}
\csromangathameasuregroup{สุโข วิปาโล ปุญฺญานํ, \\ ขิปฺปญฺจ ปรมํ สนฺติํ,}{\csromangathabat{สุโข วิปาโล ปุญฺญานํ,}{อธิปฺปาโย จ อิชฺฌติ.} \\ \csromangathabat{ขิปฺปญฺจ ปรมํ สนฺติํ,}{นิพฺพานมธิคจฺฉตีติ.}}
\csromangathasetleft{สุโข วิปาโล ปุญฺญานํ, \\ ขิปฺปญฺจ ปรมํ สนฺติํ,}
\csromangathagroup{\csromangathastanza{\csromangathabat{สุโข วิปาโล ปุญฺญานํ,}{อธิปฺปาโย จ อิชฺฌติ.} \\ \csromangathabat{ขิปฺปญฺจ ปรมํ สนฺติํ,}{นิพฺพานมธิคจฺฉตีติ.}}}

\prose{โย จ วิปาโก ปุญฺญานํ ยา จ อธิปฺปายสฺส อิชฺฌนา, อยํ อสฺสาโท. ยํ ขิปฺปญฺจ ปรมํ สนฺติํ นิพฺพานม\-ธิคจฺฉติ, อิทํ นิสฺสรณํ.}

\prose{พาตฺติํสาย เจว มหา\-ปุริส\-ลกฺขเณหิ สมนฺนาคตสฺส มหาปุริสสฺส เทฺวเยว คติโย โหนฺติ, สเจ อคารํ อชฺฌาวสติ, ราชา โหติ จกฺกวตฺตี ยาว อภิวิชินิตฺวา อชฺฌาวสติ อยํ อสฺสาโท. สเจ อคารสฺมา อนคาริยํ ปพฺพชติ สพฺเพน โอเฆน\csromansharedfootnote{af3a1e360a3190c4}{โอสเธน (อิ, ก)} นิสฺสรณํ. อยํ อสฺสาโท จ นิสฺสรณญฺจ.}

\prose{ตตฺถ กตโม อาทีนโว จ นิสฺสรณญฺจ.}

\setlength{\csromangathapagewidth}{0pt}
\setlength{\csromangathawakwidth}{0pt}
\csromangathameasuregroup{อาทานสฺส\textsuperscript{1} ภยํ ญตฺวา, \\ อนาทาตุํ นิพฺพตฺตติ,}{\csromangathabat{อาทานสฺส\textsuperscript{1} ภยํ ญตฺวา,}{ชาติมรณสมฺภวํ.} \\ \csromangathabat{อนาทาตุํ นิพฺพตฺตติ,}{ชาติมรณสงฺขยา.}}
\csromangathasetleft{อาทานสฺส\textsuperscript{1} ภยํ ญตฺวา, \\ อนาทาตุํ นิพฺพตฺตติ,}
\csromangathagroup{\csromangathastanza{\csromangathabat{อาทานสฺส\csromansharedfootnote{e7927954d99038c5}{อาทินฺนสฺส (อิ)} ภยํ ญตฺวา,}{ชาติมรณสมฺภวํ.} \\ \csromangathabat{อนาทาตุํ นิพฺพตฺตติ,}{ชาติมรณสงฺขยา.}}}

\prose{ปุริมิกาย อฑฺฒคาถาย ชาติมรณ\-สมฺภโว อาทีนโว. อนาทาตุํ นิพฺพตฺตติ ชาติ\-มรณ\-สงฺขยาติ นิสฺสรณํ.}

\prose{กิจฺฉํ วตายํ โลโก อาปนฺโน ยมิทํ ชายเต จ มียเต จ. ยาว กุทสฺสุนามสฺส ทุกฺขสฺส อนฺโต ภวิสฺสติ ปรโต วาติ เอตฺถ ยา อุปริกฺขา, อยํ อาทีนโว. โย เคธํ ญตฺวา อภินิกฺขมติ ยาว ปุราณกาย ราชธานิยา, อิทํ นิสฺสรณํ. อยํ อาทีนโว จ นิสฺสรณญฺจ.}

\csromanhangingbegin
\hangingheadline{\csromanfolio{200}ตตฺถ กตโม อสฺสาโท จ อาทีนโว จ นิสฺสรณญฺจ.}
\hangingbody{กามา หิ จิตฺรา วิวิธา1 มโนรมา,}
\hangingbody{วิรูปรูเปหิ มเถนฺติ จิตฺตํ.}
\hangingbody{ตสฺมา อหํ2 ปพฺพชิโตมฺหิ ราช,}
\hangingbody{อปณฺณกํ สามญฺญเมว เสยฺโย.}
\csromanhangingend

\prose{ยํ กามา หิ จิตฺรา วิวิธา มโนรมาติ อยํ อสฺสาโท. ยํ วิรูปรูเปหิ มเถนฺติ จิตฺตนฺติ อยํ อาทีนโว. ยํ อหํ อคารสฺมา ปพฺพชิโตมฺหิ ราช อปณฺณกํ สามญฺญเมว เสยฺโยติ อิทํ นิสฺสรณํ.}

\prose{พลวํ พาโลปมสุตฺตํ ยํ อาสาย วา เวทนียํ กมฺมํ คาหติ, ตถา เจปิ ยํ ยํ ปาปกมฺมํ อนุโภติ, ตตฺถ ทุกฺข\-เวทนีเยน กมฺเมน อภาวิตกาเยน จ ยาว ปริตฺตเจตโส จ อาทีนวํ ทสฺเสติ สุขเวทนีเยน กมฺเมน อสฺสาเทติ. ยํ ปุราสทิโส โหติ. ภาวิตจิตฺโต ภาวิตกาโย ภาวิตปญฺโญ มหานาโม อปริตฺตเจตโส, อิทํ นิสฺสรณํ.}

\proseitem{24}{ตตฺถ กตมํ โลกิกํ สุตฺตํ.}

\setlength{\csromangathapagewidth}{0pt}
\setlength{\csromangathawakwidth}{0pt}
\csromangathameasuregroup{\textsuperscript{*}น หิ ปาปํ กตํ กมฺมํ, \\ qอหนฺตํ พาลมเนฺวติ,}{\csromangathabat{\textsuperscript{*}น หิ ปาปํ กตํ กมฺมํ,}{สชฺชุขีรํว มุจฺจติ.} \\ \csromangathabat{qอหนฺตํ พาลมเนฺวติ,}{ภสฺมจฺฉนฺโนว\textsuperscript{1} ปาวโก.}}
\csromangathasetleft{\textsuperscript{*}น หิ ปาปํ กตํ กมฺมํ, \\ qอหนฺตํ พาลมเนฺวติ,}
\csromangathagroup{\csromangathastanza{\csromangathaitembat{24}{\csromansymbolfootnote{*}{เหฏฺฐา 140 ปิฏฺเฐ.}น หิ ปาปํ กตํ กมฺมํ,}{สชฺชุขีรํว มุจฺจติ.} \\ \csromangathacontbat{qอหนฺตํ พาลมเนฺวติ,}{ภสฺม\-จฺฉนฺโนว\csromansharedfootnote{31e5449a98b9c0a4}{ภสฺมาฉนฺโนว (ก) ขุ 1. 23 ปิฏฺเฐ.} ปาวโก.}}}

\prose{จตฺตาริ อคติคมนานิ, อิทํ โลกิกํ สุตฺตํ.0}

\prose{ตตฺถ กตมํ โลกุตฺตรํ สุตฺตํ.}

\prose{\csromansymbolmark{*}ยสฺสินฺทฺริยานิ สมถงฺคตานิ\csromansharedfootnote{4236cc4a935b2610}{สมถํ คตานิ (อิ) ขุ 1. 27 ปิฏฺเฐ.},}

\prose{อสฺสา ยถา สารถินา สุทนฺตา.}

\setlength{\csromangathapagewidth}{0pt}
\setlength{\csromangathawakwidth}{0pt}
\csromangathameasuregroup{}{ปหีนมานสฺส อนาสวสฺส, \\ เทวาปิ ตสฺส ปิหยนฺติ ตาทิโนติ.}
\csromangathameasurewak{ปหีนมานสฺส อนาสวสฺส, \\ เทวาปิ ตสฺส ปิหยนฺติ ตาทิโนติ.}
\setlength{\csromangathaleftcol}{0pt}
\csromangathagroup{\csromangathastanza{\csromangathawak{ปหีนมานสฺส อนาสวสฺส, \\ เทวาปิ ตสฺส ปิหยนฺติ ตาทิโนติ.}}}

\prose{อริยํ โว ภิกฺขเว สมฺมาสมาธิํ เทเสสฺสามีติ อิทํ โลกุตฺตรํ สุตฺตํ.}

\prose{\csromanfolio{201}ตตฺถ กตมํ โลกิกํ โลกุตฺตรญฺจ สุตฺตํ.}

\setlength{\csromangathapagewidth}{0pt}
\setlength{\csromangathawakwidth}{0pt}
\csromangathameasuregroup{\textsuperscript{*}สตฺติยา วิย โอมฏฺโฐ, \\ กามราคปฺปหานาย, \\ สตฺติยา วิย โอมฏฺโฐ, \\ กามราคปฺปหานาย,}{\csromangathabat{\textsuperscript{*}สตฺติยา วิย โอมฏฺโฐ,}{ทยฺหมาโนว มตฺถเก.} \\ \csromangathabat{กามราคปฺปหานาย,}{สโต ภิกฺขุ ปริพฺพเช.} \\ \csromangathabat{สตฺติยา วิย โอมฏฺโฐ,}{ทยฺหมาโนว มตฺถเกติ โลกิกํ.} \\ \csromangathabat{กามราคปฺปหานาย,}{สโต ภิกฺขุ ปริพฺพเชติ โลกุตฺตรํ.}}
\csromangathasetleft{\textsuperscript{*}สตฺติยา วิย โอมฏฺโฐ, \\ กามราคปฺปหานาย, \\ สตฺติยา วิย โอมฏฺโฐ, \\ กามราคปฺปหานาย,}
\csromangathagroup{\csromangathastanza{\csromangathabat{\csromansymbolfootnote{*}{สํ 1. 12; เหฏฺฐา 125 ปิฏฺเฐสุปิ.}สตฺติยา วิย โอมฏฺโฐ,}{ทยฺหมาโนว มตฺถเก.} \\ \csromangathabat{กาม\-ราค\-ปฺปหานาย,}{สโต ภิกฺขุ ปริพฺพเช.}}\csromangathastanza{\csromangathabat{สตฺติยา วิย โอมฏฺโฐ,}{ทยฺหมาโนว มตฺถเกติ โลกิกํ.} \\ \csromangathabat{กาม\-ราค\-ปฺปหานาย,}{สโต ภิกฺขุ ปริพฺพเชติ โลกุตฺตรํ.}}}

\prose{กพฬีกาเร อาหาเร อตฺถิ ฉนฺโทติ โลกิกํ. นตฺถิ ฉนฺโทติ โลกุตฺตรํ สุตฺตํ.}

\prose{ตตฺถ กตมํ กมฺมํ.}

\setlength{\csromangathapagewidth}{0pt}
\setlength{\csromangathawakwidth}{0pt}
\csromangathameasuregroup{โย ปาณมติปาเตติ, \\ โลเก อทินฺนํ อาทิยติ\textsuperscript{1}, \\ สุราเมรยปานญฺจ, \\ อปฺปหาย ปญฺจ เวรานิ,}{\csromangathabat{โย ปาณมติปาเตติ,}{มุสาวาทญฺจ ภาสติ.} \\ \csromangathabat{โลเก อทินฺนํ อาทิยติ\textsuperscript{1},}{ปรทารญฺจ คจฺฉติ.} \\ \csromangathabat{สุราเมรยปานญฺจ,}{โย นโร อนุยุญฺชติ.} \\ \csromangathabat{อปฺปหาย ปญฺจ เวรานิ,}{ทุสฺสีโล อิติ วุจฺจติ.}}
\csromangathasetleft{โย ปาณมติปาเตติ, \\ โลเก อทินฺนํ อาทิยติ\textsuperscript{1}, \\ สุราเมรยปานญฺจ, \\ อปฺปหาย ปญฺจ เวรานิ,}
\csromangathagroup{\csromangathastanza{\csromangathabat{โย ปาณมติปาเตติ,}{มุสาวาทญฺจ ภาสติ.} \\ \csromangathabat{โลเก อทินฺนํ อาทิยติ\csromansharedfootnote{9c5209d417a51d4d}{อาทิยิ (ก) อํ 2. 181; อุปริ 210 ปิฏฺเฐสุปิ. 2. สํ 2. 341 ปิฏฺเฐ. ** ขุ 1. 48; เหฏฺฐา 108, 173 ปิฏฺเฐสุปิ.},}{ปรทารญฺจ คจฺฉติ.}}\csromangathastanza{\csromangathabat{สุรา\-เมรย\-ปานญฺจ,}{โย นโร อนุยุญฺชติ.} \\ \csromangathabat{อปฺปหาย ปญฺจ เวรานิ,}{ทุสฺสีโล อิติ วุจฺจติ.}}}

\vspace{\tipitakaparskip}
\csromancenter{ตีณิมานิ ภิกฺขเว ทุจฺจริตานิ. อิทํ กมฺมํ.}
\prose{ตตฺถ กตโม วิปาโก.}

\prose{สฏฺฐิ\-วสฺสสหสฺสานิ, ยถารูปี วิปจฺจคา.}

\prose{ทิฏฺฐา มยา ภิกฺขเว ฉ ผสฺสายตนิกา นาม นิรยา. ทิฏฺฐา มยา ภิกฺขเว ฉ ผสฺสายตนิกา นาม สคฺคา\csromansharedfootnote{b57a64e15fcad444}{สํ 2. 341 ปิฏฺเฐ. ** ขุ 1. 48; เหฏฺฐา 108, 173 ปิฏฺเฐสุปิ.}. อยํ วิปาโก.}

\prose{ตตฺถ กตมํ กมฺมญฺจ วิปาโก จ.}

\setlength{\csromangathapagewidth}{0pt}
\setlength{\csromangathawakwidth}{0pt}
\csromangathameasuregroup{** อยสาว มลํ สมุฏฺฐิตํ, \\ เอวํ อติโธนจารินํ,}{\csromangathabat{** อยสาว มลํ สมุฏฺฐิตํ,}{ตตุฏฺฐาย ตเมว ขาทติ.} \\ \csromangathabat{เอวํ อติโธนจารินํ,}{สานิ กมฺมานิ นยนฺติ ทุคฺคติํ.}}
\csromangathasetleft{** อยสาว มลํ สมุฏฺฐิตํ, \\ เอวํ อติโธนจารินํ,}
\csromangathagroup{\csromangathastanza{\csromangathabat{** อยสาว มลํ สมุฏฺฐิตํ,}{ตตุฏฺฐาย ตเมว ขาทติ.} \\ \csromangathabat{เอวํ อติโธนจารินํ,}{สานิ กมฺมานิ นยนฺติ ทุคฺคติํ.}}}

\prose{อยสาว มลํ สมุฏฺฐิตํ, ยาว สานิ กมฺมานีติ อิทํ กมฺมํ. นยนฺติ ทุคฺคตินฺติ วิปาโก.}

\prose{จตูสุ สมฺมา\-ปฏิปชฺชมาโน มาตริ ปิตริ ตถาคเต ตถาคต\-สาวเก ยา สมฺมาปฏิปตฺติ, อิทํ กมฺมํ. ยํ เทเวสุ อุปปชฺชติ, อยํ วิปาโก. อิทํ กมฺมญฺจ วิปาโก จ.}

\setlength{\csromangathapagewidth}{0pt}
\setlength{\csromangathawakwidth}{0pt}
\csromangathameasuregroup{}{ตตฺถ กตมํ นิทฺทิฏฺฐํ สุตฺตํ. \\ เนลงฺโค เสตปจฺฉาโท, เอกาโร วตฺตตี\textsuperscript{1} รโถ. \\ อนีฆํ ปสฺส อายนฺตํ, ฉินฺนโสตํ อพนฺธนํ. \\ ยํ วา จิตฺตํ สมเณสุ, จิตฺตาคหปติ ทิสฺสติ.}
\csromangathameasurewak{ตตฺถ กตมํ นิทฺทิฏฺฐํ สุตฺตํ. \\ เนลงฺโค เสตปจฺฉาโท, เอกาโร วตฺตตี\textsuperscript{1} รโถ. \\ อนีฆํ ปสฺส อายนฺตํ, ฉินฺนโสตํ อพนฺธนํ. \\ ยํ วา จิตฺตํ สมเณสุ, จิตฺตาคหปติ ทิสฺสติ.}
\setlength{\csromangathaleftcol}{0pt}
\csromangathagroup{\csromangathastanza{\csromangathawak{\csromanfolio{202}\csromangathaitemline{25}{ตตฺถ กตมํ นิทฺทิฏฺฐํ สุตฺตํ.} \\ \csromangathacontline{เนลงฺโค เสตปจฺฉาโท, เอกาโร วตฺตตี\csromansharedfootnote{07021f5e5e5dd272}{วตฺตเต (ก) ขุ 1. 174 ปิฏฺเฐ.} รโถ.} \\ \csromangathacontline{อนีฆํ ปสฺส อายนฺตํ, ฉินฺนโสตํ อพนฺธนํ.} \\ \csromangathacontline{ยํ วา จิตฺตํ สมเณสุ, จิตฺตาคหปติ ทิสฺสติ.}}}}

\prose{เอวํ อิมาย คาถาย นิทฺทิฏฺโฐ อตฺโถ.}

\prose{โคปาลโกปเม เอกาทส ปทานิ. เอวํ โข ภิกฺขเว ภิกฺขุ รูปญฺญู โหติ. ยา จ อติเรกปูชาย ปูเชตา โหตีติ. อิมานิ เอกาทส ปทานิ ยถาภาสิตานิ นิทฺทิฏฺโฐ อตฺโถ.}

\setlength{\csromangathapagewidth}{0pt}
\setlength{\csromangathawakwidth}{0pt}
\csromangathameasuregroup{}{ตตฺถ กตโม อนิทฺทิฏฺโฐ อตฺโถ. \\ สุโข วิเวโก ตุฏฺฐสฺส, สุตธมฺมสฺส ปสฺสโต. \\ อพฺยาปชฺชํ\textsuperscript{1} สุขํ โลเก, ปาณภูเตสุ สํยโมติ. \\ สุขา วิราคตา โลเก, กามานํ สมติกฺกโม. \\ อสฺมิมานสฺส โย วินโย, เอตํ เว ปรมํ สุขนฺติ.}
\csromangathameasurewak{ตตฺถ กตโม อนิทฺทิฏฺโฐ อตฺโถ. \\ สุโข วิเวโก ตุฏฺฐสฺส, สุตธมฺมสฺส ปสฺสโต. \\ อพฺยาปชฺชํ\textsuperscript{1} สุขํ โลเก, ปาณภูเตสุ สํยโมติ. \\ สุขา วิราคตา โลเก, กามานํ สมติกฺกโม. \\ อสฺมิมานสฺส โย วินโย, เอตํ เว ปรมํ สุขนฺติ.}
\setlength{\csromangathaleftcol}{0pt}
\csromangathagroup{\csromangathastanza{\csromangathawak{\csromangathaitemline{25}{ตตฺถ กตโม อนิทฺทิฏฺโฐ อตฺโถ.} \\ \csromangathacontline{สุโข วิเวโก ตุฏฺฐสฺส, สุตธมฺมสฺส ปสฺสโต.} \\ \csromangathacontline{อพฺยาปชฺชํ\csromansharedfootnote{aa8eca7d9acdaf43}{อพฺยาปชฺฌํ (อิ, ก) ขุ 1. 88 ปิฏฺเฐ.} สุขํ โลเก, ปาณภูเตสุ สํยโมติ.} \\ \csromangathacontline{สุขา วิราคตา โลเก, กามานํ สมติกฺกโม.} \\ \csromangathacontline{อสฺมิมานสฺส โย วินโย, เอตํ เว ปรมํ สุขนฺติ.}}}}

\prose{อิทํ อนิทฺทิฏฺฐํ, อฏฺฐ มหาปุริส\-วิตกฺกา. อิทํ อนิทฺทิฏฺฐํ.}

\setlength{\csromangathapagewidth}{0pt}
\setlength{\csromangathawakwidth}{0pt}
\csromangathameasuregroup{}{ตตฺถ กตมํ นิทฺทิฏฺฐญฺจ อนิทฺทิฏฺฐญฺจ. \\ ปสนฺนเนตฺโต สุมุโข, พฺรหา อุชุ ปตาปวา. \\ มชฺเฌ สมณสํฆสฺส, อาทิจฺโจว วิโรจสิ\textsuperscript{1}.}
\csromangathameasurewak{ตตฺถ กตมํ นิทฺทิฏฺฐญฺจ อนิทฺทิฏฺฐญฺจ. \\ ปสนฺนเนตฺโต สุมุโข, พฺรหา อุชุ ปตาปวา. \\ มชฺเฌ สมณสํฆสฺส, อาทิจฺโจว วิโรจสิ\textsuperscript{1}.}
\setlength{\csromangathaleftcol}{0pt}
\csromangathagroup{\csromangathastanza{\csromangathawak{\csromangathaitemline{25}{ตตฺถ กตมํ นิทฺทิฏฺฐญฺจ อนิทฺทิฏฺฐญฺจ.} \\ \csromangathacontline{ปสนฺนเนตฺโต สุมุโข, พฺรหา อุชุ ปตาปวา.} \\ \csromangathacontline{มชฺเฌ สมณ\-สํฆสฺส, อาทิจฺโจว วิโรจสิ\csromansharedfootnote{e585b61f9570f254}{ขุ 1. 367 สุตฺตนิปาเต.}.}}}}

\prose{ปสนฺนเนตฺโต ยาว อาทิจฺโจว วิโรจสีติ นิทฺทิฏฺโฐ. ปสนฺนเนตฺโต โย ภควา กถญฺจ ปน ปสนฺนเนตฺตตา, กถํ สุมุขตา, กถํ พฺรหกายตา, กถํ อุชุกตา, กถํ ปตาปวตา, กถํ วิโรจตาติ อนิทฺทิฏฺโฐ. เผณปิณฺโฑปมํ เวยฺยากรณํ ยถา เผณปิณฺโฑ เอวํ รูปํ ยถา ปุพฺพุโฬ เอวํ เวทนา มายา วิญฺญาณํ ปญฺจกฺขนฺธา ปญฺจหิ อุปมาหิ นิทฺทิฏฺฐา. เกน การเณน เผณปิณฺโฑปมํ รูปํ สพฺพญฺจ จกฺขุ\-วิญฺเญยฺยํ ยํ วา จตูหิ อายตเนหิ. กถํ เวทนา ปุพฺพุฬูปมา. กตรา จ สา เวทนา สุขา ทุกฺขา อทุกฺขมสุขา, เอวเมสา อนิทฺทิฏฺฐา. เอวํ นิทฺทิฏฺฐญฺจ อนิทฺทิฏฺฐญฺจ.}

\setlength{\csromangathapagewidth}{0pt}
\setlength{\csromangathawakwidth}{0pt}
\csromangathameasuregroup{}{ตตฺถ กตมํ ญาณํ. \\ ปญฺญา หิ เสฏฺฐา โลกสฺมิํ, ยายํ นิพฺเพธคามินี. \\ ยาย\textsuperscript{1} สมฺมา ปชานาติ, ชาติมรณสงฺขยํ.}
\csromangathameasurewak{ตตฺถ กตมํ ญาณํ. \\ ปญฺญา หิ เสฏฺฐา โลกสฺมิํ, ยายํ นิพฺเพธคามินี. \\ ยาย\textsuperscript{1} สมฺมา ปชานาติ, ชาติมรณสงฺขยํ.}
\setlength{\csromangathaleftcol}{0pt}
\csromangathagroup{\csromangathastanza{\csromangathawak{\csromanfolio{203}\csromangathaitemline{26}{ตตฺถ กตมํ ญาณํ.} \\ \csromangathacontline{ปญฺญา หิ เสฏฺฐา โลกสฺมิํ, ยายํ นิพฺเพธคามินี.} \\ \csromangathacontline{ยาย\csromansharedfootnote{86d7cd69fa4fea94}{ยายํ (ก) ขุ 1. 219; เหฏฺฐา 143 ปิฏฺเฐสุปิ.} สมฺมา ปชานาติ, ชาติมรณสงฺขยํ.}}}}

\prose{ตีณิมานิ อินฺทฺริยานิ อนญฺญาต\-ญฺญสฺสามี\-ตินฺทฺริยํ อญฺญินฺทฺริยํ อญฺญาตาวิ\-นฺทฺริยํ, อิทํ ญาณํ.}

\setlength{\csromangathapagewidth}{0pt}
\setlength{\csromangathawakwidth}{0pt}
\csromangathameasuregroup{}{ตตฺถ กตมํ เนยฺยํ. \\ กาเมสุ สตฺตา กามสงฺคสตฺตา, \\ สํโยชเน วชฺชมปสฺสมานา. \\ น หิ ชาตุ สํโยชนสงฺคสตฺตา, \\ โอฆํ ตเรยฺยุํ วิปุลํ มหนฺตํ\textsuperscript{1}.}
\csromangathameasurewak{ตตฺถ กตมํ เนยฺยํ. \\ กาเมสุ สตฺตา กามสงฺคสตฺตา, \\ สํโยชเน วชฺชมปสฺสมานา. \\ น หิ ชาตุ สํโยชนสงฺคสตฺตา, \\ โอฆํ ตเรยฺยุํ วิปุลํ มหนฺตํ\textsuperscript{1}.}
\setlength{\csromangathaleftcol}{0pt}
\csromangathagroup{\csromangathastanza{\csromangathawak{\csromangathaitemline{26}{ตตฺถ กตมํ เนยฺยํ.} \\ \csromangathacontline{กาเมสุ สตฺตา กามสงฺคสตฺตา,} \\ \csromangathacontline{สํโยชเน วชฺชม\-ปสฺสมานา.} \\ \csromangathacontline{น หิ ชาตุ สํโยชนสงฺค\-สตฺตา,} \\ \csromangathacontline{โอฆํ ตเรยฺยุํ วิปุลํ มหนฺตํ\csromansharedfootnote{cc4ff3fd1f454f6a}{ขุ 1. 171 อุทาเน; เหฏฺฐา 174 ปิฏฺเฐสุปิ.}.}}}}

\prose{จตูหิ องฺเคหิ สมนฺนาคตา กายสฺส เภทา เทเวสุ อุปฺปชฺชนฺติ. อุทาเน กาปิยํ สุตฺตํ อปณฺณก\-ปสาทนียํ. อิทํ เนยฺยํ.}

\prose{ตตฺถ กตมํ ญาณญฺจ เนยฺยญฺจ.}

\setlength{\csromangathapagewidth}{0pt}
\setlength{\csromangathawakwidth}{0pt}
\csromangathameasuregroup{“สพฺเพ ธมฺมา อนตฺตาติ, \\ อถ นิพฺพินฺทติ ทุกฺเข,}{\csromangathabat{“สพฺเพ ธมฺมา อนตฺตาติ,}{ยทา ปญฺญาย ปสฺสติ.} \\ \csromangathabat{อถ นิพฺพินฺทติ ทุกฺเข,}{เอส มคฺโค วิสุทฺธิยา.}}
\csromangathasetleft{“สพฺเพ ธมฺมา อนตฺตาติ, \\ อถ นิพฺพินฺทติ ทุกฺเข,}
\csromangathagroup{\csromangathastanza{\csromangathabat{“สพฺเพ ธมฺมา อนตฺตาติ,}{ยทา ปญฺญาย ปสฺสติ.} \\ \csromangathabat{อถ นิพฺพินฺทติ ทุกฺเข,}{เอส มคฺโค วิสุทฺธิยา.}}}

\prose{ยทา ปสฺสตีติ ญาณํ. โย สพฺพธมฺเม อนตฺตากาเรน อุปฏฺฐเปติ. อิทํ เนยฺยํ.}

\prose{จตฺตาริ อริยสจฺจานิ, ตตฺถ ตีณิ เนยฺยานิ มคฺคสจฺจํ สีลกฺขนฺโธ จ ปญฺญากฺขนฺโธ จ, อิทํ ญาณญฺจ เนยฺยญฺจ.}

\setlength{\csromangathapagewidth}{0pt}
\setlength{\csromangathawakwidth}{0pt}
\csromangathameasuregroup{}{ตตฺถ กตมํ ทสฺสนํ. \\ เอเสว มคฺโค นตฺถญฺโญ, ทสฺสนสฺส วิสุทฺธิยา. \\ เอตญฺหิ ตุเมฺห ปฏิปชฺชถ, มารสฺเสตํ ปโมหนํ\textsuperscript{1}.}
\csromangathameasurewak{ตตฺถ กตมํ ทสฺสนํ. \\ เอเสว มคฺโค นตฺถญฺโญ, ทสฺสนสฺส วิสุทฺธิยา. \\ เอตญฺหิ ตุเมฺห ปฏิปชฺชถ, มารสฺเสตํ ปโมหนํ\textsuperscript{1}.}
\setlength{\csromangathaleftcol}{0pt}
\csromangathagroup{\csromangathastanza{\csromangathawak{\csromangathaitemline{27}{ตตฺถ กตมํ ทสฺสนํ.} \\ \csromangathacontline{เอเสว มคฺโค นตฺถญฺโญ, ทสฺสนสฺส วิสุทฺธิยา.} \\ \csromangathacontline{เอตญฺหิ ตุเมฺห ปฏิปชฺชถ, มารสฺเสตํ ปโมหนํ\csromansharedfootnote{cf43595373364069}{ขุ 1. 52 ธมฺมปเท.}.}}}}

\prose{จตูหิ องฺเคหิ สมนฺนาคโต อริยสาวโก อตฺตนาว\csromansharedfootnote{6b63b5dce6b78086}{อตฺตนาเยว (ก) สํ 3. 310 ปิฏฺเฐ.} อตฺตานํ พฺยากเรยฺย ‘ขีณนิรโยมฺหิ ยาว โสตาปนฺโนหมสฺมิ อวินิปาต\-ธมฺโม นิยโต สมฺโมธิปรายโณ”ติ. อิทํ ทสฺสนํ.}

\csromanhangingbegin
\hangingheadline{\csromanfolio{204}ตตฺถ กตมา ภาวนา.}
\hangingbody{ยสฺสินฺทฺริยานิ สุภาวิตานิ,}
\hangingbody{อชฺฌตฺตํ พหิทฺธา จ สพฺพโลเก.}
\hangingbody{โส ปุคฺคโล มติ จ รูปสญฺญี,}
\hangingbody{สุโมหคตา น ชานาติ1.}
\csromanhangingend

\prose{จตฺตาริ ธมฺมปทานิ อนภิชฺฌา อพฺยาปาโท สมฺมาสติ สมฺมาสมาธิ. อยํ ภาวนา.}

\csromanhangingbegin
\hangingheadline{ตตฺถ กตมํ ทสฺสนญฺจ ภาวนา จ.}
\hangingbody{วจสา มนสาถ กมฺมุนา จ,}
\hangingbody{อวิรุทฺโธ สมฺมา วิทิตฺวา2 ธมฺมํ.}
\hangingbody{นิพฺพานปทาภิปตฺถยาโน,}
\hangingbody{สมฺมา โส โลเก ปริพฺพเชยฺย.}
\csromanhangingend

\prose{โสตาปตฺติผลํ สจฺฉิ\-กาตุกาเมน กตเม ธมฺมา มนสิกาตพฺพา, ภควา อาห ปญฺจุ\-ปาทาน\-กฺขนฺธา. อิทํ ทสฺสนญฺจ ภาวนา จ.}

\proseitem{28}{ตตฺถ กตเม วิปาก\-ธมฺม\-ธมฺมา.}

\prose{ยานิ กโรติ ปุริโสติ วิตฺถาโร.  ตีณิมานิ ภิกฺขเว สุจริตานิ. อิเม วิปาก\-ธมฺม\-ธมฺมา.}

\setlength{\csromangathapagewidth}{0pt}
\setlength{\csromangathawakwidth}{0pt}
\csromangathameasuregroup{}{ตตฺถ กตเม นวิปากธมฺมธมฺมา. \\ รูปํ เวทยิตํ สญฺญา, วิญฺญาณํ ยา เจว เจตนา. \\ เนโสหมสฺมิ น เมโส อตฺตา, อิติ ทิฏฺโฐ วิรชฺชติ.}
\csromangathameasurewak{ตตฺถ กตเม นวิปากธมฺมธมฺมา. \\ รูปํ เวทยิตํ สญฺญา, วิญฺญาณํ ยา เจว เจตนา. \\ เนโสหมสฺมิ น เมโส อตฺตา, อิติ ทิฏฺโฐ วิรชฺชติ.}
\setlength{\csromangathaleftcol}{0pt}
\csromangathagroup{\csromangathastanza{\csromangathawak{\csromangathaitemline{28}{ตตฺถ กตเม นวิปาก\-ธมฺม\-ธมฺมา.} \\ \csromangathacontline{รูปํ เวทยิตํ สญฺญา, วิญฺญาณํ ยา เจว เจตนา.} \\ \csromangathacontline{เนโสหมสฺมิ น เมโส อตฺตา, อิติ ทิฏฺโฐ วิรชฺชติ.}}}}

\prose{ปญฺจิเม ภิกฺขเว ขนฺธา, อิเม นวิปาก\-ธมฺม\-ธมฺมา.}

\setlength{\csromangathapagewidth}{0pt}
\setlength{\csromangathawakwidth}{0pt}
\csromangathameasuregroup{}{ตตฺถ กตโม เนววิปาโก นวิปากธมฺมธมฺโม. \\ เย เอวํ ปฏิปชฺชนฺติ, นยํ พุทฺเธน เทสิตํ. \\ เต ทุกฺขสฺสนฺตํ กริสฺสนฺติ, สตฺถุสาสนการกาติ.}
\csromangathameasurewak{ตตฺถ กตโม เนววิปาโก นวิปากธมฺมธมฺโม. \\ เย เอวํ ปฏิปชฺชนฺติ, นยํ พุทฺเธน เทสิตํ. \\ เต ทุกฺขสฺสนฺตํ กริสฺสนฺติ, สตฺถุสาสนการกาติ.}
\setlength{\csromangathaleftcol}{0pt}
\csromangathagroup{\csromangathastanza{\csromangathawak{\csromangathaitemline{28}{ตตฺถ กตโม เนววิปาโก นวิปาก\-ธมฺม\-ธมฺโม.} \\ \csromangathacontline{เย เอวํ ปฏิปชฺชนฺติ, นยํ พุทฺเธน เทสิตํ.} \\ \csromangathacontline{เต ทุกฺขสฺสนฺตํ กริสฺสนฺติ, สตฺถุสาสนการกาติ.}}}}

\prose{อิติ ยา จ สมฺมาปฏิปตฺติ โย จ นิโรโธ, อุภยเมตํ เนววิปาโก นวิปากธมฺโม. พฺรหฺมจริยํ โว ภิกฺขเว เทเสสฺสามิ, พฺรหฺมจริย\-ผลานิ จ พฺรหฺมจริยญฺจ อริโย อฏฺฐงฺคิโก มคฺโค พฺรหฺมจริย\-ผลานิ โสตาปตฺติผลํ ยาว อรหตฺตํ.}

\proseitem{29}{\csromanfolio{205}ตตฺถ กตมํ สกวจนํ.}

\setlength{\csromangathapagewidth}{0pt}
\setlength{\csromangathawakwidth}{0pt}
\csromangathameasuregroup{\textsuperscript{*}สพฺพปาปสฺส อกรณํ, \\ สจิตฺตปริโยทาปนํ,}{\csromangathabat{\textsuperscript{*}สพฺพปาปสฺส อกรณํ,}{กุสลสฺส อุปสมฺปทา.} \\ \csromangathabat{สจิตฺตปริโยทาปนํ,}{เอตํ พุทฺธาน สาสนํ.}}
\csromangathasetleft{\textsuperscript{*}สพฺพปาปสฺส อกรณํ, \\ สจิตฺตปริโยทาปนํ,}
\csromangathagroup{\csromangathastanza{\csromangathaitembat{29}{\csromansymbolfootnote{*}{ขุ 1. 41; ที 2. 42; เหฏฺฐา 141, 161 ปิฏฺเฐสุปิ. ** สํ 1. 6; เหฏฺฐา 160 ปิฏฺเฐสุปิ.}สพฺพปาปสฺส อกรณํ,}{กุสลสฺส อุปสมฺปทา.} \\ \csromangathacontbat{สจิตฺต\-ปริโยทาปนํ,}{เอตํ พุทฺธาน สาสนํ.}}}

\prose{ตีณิมานิ ภิกฺขเว วิโมกฺขมุขานิ. อิทํ สกวจนํ.}

\prose{ตตฺถ กตมํ ปรวจนํ.}

\setlength{\csromangathapagewidth}{0pt}
\setlength{\csromangathawakwidth}{0pt}
\csromangathameasuregroup{** นตฺถิ ปุตฺตสมํ เปมํ, \\ นตฺถิ สูริยสมา อาภา,}{\csromangathabat{** นตฺถิ ปุตฺตสมํ เปมํ,}{นตฺถิ โคณสมิตํ ธนํ.} \\ \csromangathabat{นตฺถิ สูริยสมา อาภา,}{สมุทฺทปรมา สรา.}}
\csromangathasetleft{** นตฺถิ ปุตฺตสมํ เปมํ, \\ นตฺถิ สูริยสมา อาภา,}
\csromangathagroup{\csromangathastanza{\csromangathaitembat{29}{** นตฺถิ ปุตฺตสมํ เปมํ,}{นตฺถิ โคณสมิตํ ธนํ.} \\ \csromangathacontbat{นตฺถิ สูริยสมา อาภา,}{สมุทฺทปรมา สรา.}}}

\prose{เหตุนา มาริสา โกสิยา สุภาสิเตน สงฺคามวิชโย โสปิ นาม ภิกฺขเว สกฺโก เทวานมินฺโท สกํ ผลํ ปริภุญฺช\-มาโนติ วิตฺถาเรน กาตพฺพํ. อิทํ ปรวจนํ.}

\prose{ตตฺถ กตมํ สกวจนญฺจ ปรวจนญฺจ.}

\setlength{\csromangathapagewidth}{0pt}
\setlength{\csromangathawakwidth}{0pt}
\csromangathameasuregroup{“ยํ ปตฺตํ ยญฺจ ปตฺตพฺพํ, \\ เย เอวํวาทิโน นตฺถิ,}{\csromangathabat{“ยํ ปตฺตํ ยญฺจ ปตฺตพฺพํ,}{อุภยเมตํ รชานุกิณฺณํ.} \\ \csromangathabat{เย เอวํวาทิโน นตฺถิ,}{เตสํ กาเมสุ โทโส”ติ.}}
\csromangathasetleft{“ยํ ปตฺตํ ยญฺจ ปตฺตพฺพํ, \\ เย เอวํวาทิโน นตฺถิ,}
\csromangathagroup{\csromangathastanza{\csromangathabat{“ยํ ปตฺตํ ยญฺจ ปตฺตพฺพํ,}{อุภยเมตํ รชานุกิณฺณํ.} \\ \csromangathabat{เย เอวํวาทิโน นตฺถิ,}{เตสํ กาเมสุ โทโส”ติ.}}}

\prose{อิทํ ปรวจนํ. เย จ โข เต อุโภ อนฺเต อนุปคมฺม วฏฺฏํ เตสํ นตฺถิ ปญฺญาปนาย. อิทํ สกวจนํ.}

\prose{\csromansymbolfootnote{+}{เหฏฺฐา 30 ปิฏฺเฐ.}“นนฺทติ ปุตฺเตหิ ปุตฺติมา,}

\prose{โคมา โคหิ\csromansharedfootnote{5877e2ed0a4985d1}{โภคิโก โภเคหิ (อิ) สํ 1. 6 ปิฏฺเฐ.} ตเถว นนฺทติ.}

\setlength{\csromangathapagewidth}{0pt}
\setlength{\csromangathawakwidth}{0pt}
\csromangathameasuregroup{}{อุปธี หิ นรสฺส นนฺทนา, \\ น หิ โส นนฺทติ โย นิรูปธี”ติ. \\ ปรวจนํ.}
\csromangathameasurewak{อุปธี หิ นรสฺส นนฺทนา, \\ น หิ โส นนฺทติ โย นิรูปธี”ติ. \\ ปรวจนํ.}
\setlength{\csromangathaleftcol}{0pt}
\csromangathagroup{\csromangathastanza{\csromangathawak{อุปธี หิ นรสฺส นนฺทนา, \\ น หิ โส นนฺทติ โย นิรูปธี”ติ. \\ ปรวจนํ.}}}

\prose{\csromansymbolmark{+}“โสจติ ปุตฺเตหิ ปุตฺติมา,}

\prose{โคมา โคหิ ตเถว โสจติ.}

\setlength{\csromangathapagewidth}{0pt}
\setlength{\csromangathawakwidth}{0pt}
\csromangathameasuregroup{อุปธี หิ นรสฺส โสจนา,}{\csromangathabat{อุปธี หิ นรสฺส โสจนา,}{น หิ โส โสจติ โย นิรูปธี”ติ.} \\ สกวจนํ. \\ อิทํ สกวจนํ ปรวจนญฺจ. \\ ตตฺถ กตมํ สตฺตาธิฏฺฐานํ. \\ เย เกจิ ภูตา ภวิสฺสนฺติ เย วาปิ, \\ สพฺเพ คมิสฺสนฺติ ปหาย เทหํ. \\ ตํ สพฺพชานิํ กุสโล วิทิตฺวา, \\ ธมฺเม ฐิโต\textsuperscript{1} พฺรหฺมจริยํ จเรยฺย.}
\csromangathameasurewak{ตตฺถ กตมํ สตฺตาธิฏฺฐานํ. \\ เย เกจิ ภูตา ภวิสฺสนฺติ เย วาปิ, \\ สพฺเพ คมิสฺสนฺติ ปหาย เทหํ. \\ ตํ สพฺพชานิํ กุสโล วิทิตฺวา, \\ ธมฺเม ฐิโต\textsuperscript{1} พฺรหฺมจริยํ จเรยฺย.}
\csromangathasetleft{อุปธี หิ นรสฺส โสจนา,}
\csromangathagroup{\csromangathastanza{\csromangathabat{อุปธี หิ นรสฺส โสจนา,}{น หิ โส โสจติ โย นิรูปธี”ติ.} \\ สกวจนํ. \\ อิทํ สกวจนํ ปรวจนญฺจ.}\csromangathastanza{\csromangathawak{\csromanfolio{206}\csromangathaitemline{30}{ตตฺถ กตมํ สตฺตา\-ธิฏฺฐานํ.} \\ \csromangathacontline{เย เกจิ ภูตา ภวิสฺสนฺติ เย วาปิ,} \\ \csromangathacontline{สพฺเพ คมิสฺสนฺติ ปหาย เทหํ.} \\ \csromangathacontline{ตํ สพฺพชานิํ กุสโล วิทิตฺวา,} \\ \csromangathacontline{ธมฺเม ฐิโต\csromansharedfootnote{6fccd45a21d55f88}{อาตาปิโย (ขุ 1. 134 อุทาเน)} พฺรหฺมจริยํ จเรยฺย.}}}}

\prose{ตโยเม ภิกฺขเว สตฺถาโร, ตถาคโต อรหํ เสกฺโข ปฏิปโท. อิทํ สตฺตา\-ธิฏฺฐานํ.}

\setlength{\csromangathapagewidth}{0pt}
\setlength{\csromangathawakwidth}{0pt}
\csromangathameasuregroup{}{ตตฺถ กตมํ ธมฺมาธิฏฺฐานํ. \\ ยญฺจ กามสุขํ\textsuperscript{1} โลเก, ยญฺจิทํ ทิวิยํ สุขํ. \\ ตณฺหกฺขยสุขสฺเสเต, กลํ นาคฺฆนฺติ โสฬสิํ.}
\csromangathameasurewak{ตตฺถ กตมํ ธมฺมาธิฏฺฐานํ. \\ ยญฺจ กามสุขํ\textsuperscript{1} โลเก, ยญฺจิทํ ทิวิยํ สุขํ. \\ ตณฺหกฺขยสุขสฺเสเต, กลํ นาคฺฆนฺติ โสฬสิํ.}
\setlength{\csromangathaleftcol}{0pt}
\csromangathagroup{\csromangathastanza{\csromangathawak{\csromangathaitemline{30}{ตตฺถ กตมํ ธมฺมา\-ธิฏฺฐานํ.} \\ \csromangathacontline{ยญฺจ กามสุขํ\csromansharedfootnote{cc1c8025d4f17658}{(ขุ 1. 89 อุทาเน; เหฏฺฐา 142 ปิฏฺเฐสุปิ.)} โลเก, ยญฺจิทํ ทิวิยํ สุขํ.} \\ \csromangathacontline{ตณฺหกฺขยสุขสฺเสเต, กลํ นาคฺฆนฺติ โสฬสิํ.}}}}

\prose{สตฺติเม ภิกฺขเว โพชฺฌงฺคา, อิทํ ธมฺมา\-ธิฏฺฐานํ.}

\prose{ตตฺถ กตมํ สตฺตา\-ธิฏฺฐานญฺจ ธมฺมาธิฏฺฐานญฺจ. ทุทฺทสมนฺตํ สจฺจํ ทุทฺทโส ปฏิเวโธ พาเลหิ, ชานโต ปสฺสโต นตฺถิ นนฺทีติ วทามิ. ทุทฺทสมนฺตํ สจฺจํ ทุทฺทโส ปฏิเวโธ พาเลหีติ ธมฺมา\-ธิฏฺฐานํ. ชานโต ปสฺสโต นตฺถิ นนฺทีติ สตฺตา\-ธิฏฺฐานํ. ทารุกฺขนฺโธปมํ คงฺคาย ตีริยา โอริมญฺจ ตีรํ ปาริมญฺจ ตีรํ ถเล วา\csromansharedfootnote{7a2d2751bcfeebff}{ถเลว จ (ก) สํ 2. 386 ปิฏฺเฐ.} น จ อุสฺสีทนํ, มชฺเฌ จ น สํสีทนํ มนุสฺสคฺคาโห จ อมนุสฺสคฺคาโห จ อนฺโตปูติภาโว จ, อิทํ ธมฺมา\-ธิฏฺฐานํ. เอวํ ปน ภิกฺขุ นิพฺพานนินฺโน ภวิสฺสติ นิพฺพาน\-ปรายโณติ สตฺตา\-ธิฏฺฐานํ. อิทํ สตฺตา\-ธิฏฺฐานญฺจ ธมฺมาธิฏฺฐานญฺจ.}

\prose{ตตฺถ กตโม ถโว.}

\setlength{\csromangathapagewidth}{0pt}
\setlength{\csromangathawakwidth}{0pt}
\csromangathameasuregroup{\textsuperscript{*}มคฺคานฏฺฐงฺคิโก เสฏฺโฐ, \\ วิราโค เสฏฺโฐ ธมฺมานํ,}{\csromangathabat{\textsuperscript{*}มคฺคานฏฺฐงฺคิโก เสฏฺโฐ,}{สจฺจานํ จตุโร ปทา.} \\ \csromangathabat{วิราโค เสฏฺโฐ ธมฺมานํ,}{ทฺวิปทานญฺจ จกฺขุมา.}}
\csromangathasetleft{\textsuperscript{*}มคฺคานฏฺฐงฺคิโก เสฏฺโฐ, \\ วิราโค เสฏฺโฐ ธมฺมานํ,}
\csromangathagroup{\csromangathastanza{\csromangathaitembat{30}{\csromansymbolfootnote{*}{เหฏฺฐา 162 ปิฏฺเฐ.}มคฺคานฏฺฐงฺคิโก เสฏฺโฐ,}{สจฺจานํ จตุโร ปทา.} \\ \csromangathacontbat{วิราโค เสฏฺโฐ ธมฺมานํ,}{ทฺวิปทานญฺจ จกฺขุมา.}}}

\prose{ตีณิมานิ ภิกฺขเว อคฺคานิ, พุทฺโธ สตฺตานํ, วิราโค ธมฺมานํ, สํโฆ คณานํ. อยํ ถโว.}

\proseitem{31}{\csromanfolio{207}ตตฺถ กตมํ อนุญฺญาตํ.}

\setlength{\csromangathapagewidth}{0pt}
\setlength{\csromangathawakwidth}{0pt}
\csromangathameasuregroup{กาเยน สํวโร สาธุ, \\ มนสา สํวโร สาธุ, \\ สพฺพตฺถ สํวุโต ภิกฺขุ,}{\csromangathabat{กาเยน สํวโร สาธุ,}{สาธุ วาจาย สํวโร.} \\ \csromangathabat{มนสา สํวโร สาธุ,}{สาธุ สพฺพตฺถ สํวุโต.} \\ \csromangathabat{สพฺพตฺถ สํวุโต ภิกฺขุ,}{สพฺพทุกฺขา ปมุจฺจติ\textsuperscript{1}.} \\ อิทํ ภควตา อนุญฺญาตํ.}
\csromangathasetleft{กาเยน สํวโร สาธุ, \\ มนสา สํวโร สาธุ, \\ สพฺพตฺถ สํวุโต ภิกฺขุ,}
\csromangathagroup{\csromangathastanza{\csromangathabat{กาเยน สํวโร สาธุ,}{สาธุ วาจาย สํวโร.} \\ \csromangathabat{มนสา สํวโร สาธุ,}{สาธุ สพฺพตฺถ สํวุโต.}}\csromangathastanza{\csromangathabat{สพฺพตฺถ สํวุโต ภิกฺขุ,}{สพฺพทุกฺขา ปมุจฺจติ\csromansharedfootnote{338dd9d64fd7ee1e}{[มิสฺสินฺคฺ โนเต 0]}.} \\ อิทํ ภควตา อนุญฺญาตํ.}}

\prose{ตีณิมานิ ภิกฺขเว กรณียานิ กายสุจริตํ วจีสุจริตํ มโนสุจริตํ. อิทํ อนุญฺญาตํ.}

\prose{ตตฺถ กตมํ ปฏิกฺขิตฺตํ.}

\prose{นตฺถิ ปุตฺตสมํ เปมํ. วิตฺถาโร อิทํ ปฏิกฺขิตฺตํ.}

\prose{ตีณิมานิ ภิกฺขเว อกรณียานิ สยํ อภิญฺญาย เทสิตานิ. กตมานิ ตีณิ, กายทุจฺจริตํ วจีทุจฺจริตํ มโนทุจฺจริตํ. อิทํ ปฏิกฺขิตฺตํ.}

\prose{ตตฺถ กตมํ อนุญฺญาตญฺจ ปฏิกฺขิตฺตญฺจ.}

\setlength{\csromangathapagewidth}{0pt}
\setlength{\csromangathawakwidth}{0pt}
\csromangathameasuregroup{\textsuperscript{*}กาเยน กุสลํ กเร, \\ กายทุจฺจริตํ หิตฺวา,}{\csromangathabat{\textsuperscript{*}กาเยน กุสลํ กเร,}{อสฺส กาเยน สํวุโต.} \\ \csromangathabat{กายทุจฺจริตํ หิตฺวา,}{กายสุจริตํ จเร.}}
\csromangathasetleft{\textsuperscript{*}กาเยน กุสลํ กเร, \\ กายทุจฺจริตํ หิตฺวา,}
\csromangathagroup{\csromangathastanza{\csromangathaitembat{31}{\csromansymbolfootnote{*}{อุปริ 216 ปิฏฺเฐ.}กาเยน กุสลํ กเร,}{อสฺส กาเยน สํวุโต.} \\ \csromangathacontbat{กายทุจฺจริตํ หิตฺวา,}{กายสุจริตํ จเร.}}}

\prose{ทฺวีหิ ปฐมปเทหิ จตุตฺเถน จ ปเทน อนุชานาติ. กายทุจฺจริตํ หิตฺวาติ ตติเยน ปเทน ปฏิกฺขิตฺตนฺติ. มหาวิภงฺโค อจิรตปานาโท.}

\csromanheader{2}{\textbf{ตตฺถิมา อุทฺทานคาถา}}
\csromanmatikaanchor{mtk.42}
\csromantocmarkref{section}{3. สุตฺตาธิฏฺฐานตติยภูมิ}{mtk.42}
\bookmark[dest={mtk.42},level=1]{3. สุตฺตาธิฏฺฐานตติยภูมิ}
\setlength{\csromangathapagewidth}{0pt}
\setlength{\csromangathawakwidth}{0pt}
\csromangathameasuregroup{สเจ ภายสิ ทุกฺขสฺส, \\ วสฺสกาเล ยถา ฉตฺตํ, \\ สพฺเพ ธมฺมา อนตฺตาติ, \\ น โว ทุกฺขา ปโมกฺขาตฺถิ, \\ กามจฺฉนฺทํ อุปาทาย, \\ สุภาวิตตฺเต โพชฺฌงฺเค, \\ สุญฺญโต โลกํ อเวกฺขสฺสุ, \\ กามํ กามยมานสฺส, \\ หญฺญเต สพฺพา มุจฺเจว, \\ โย กาเม ปริวชฺเชติ, \\ ยานิ กโรติ ปุริโส, \\ สุโข วิปาโก ปุญฺญานํ, \\ อาทานสฺส ภยํ ญตฺวา, \\ กามา หิ จิตฺรา วิวิธา, \\ น หิ ปาปํ กตํ กมฺมํ, \\ ยสฺสินฺทฺริยานิ สมถงฺคตานิ, \\ สตฺติยา วิย โอมฏฺโฐ, \\ โย ปาณมติปาเตติ, \\ สฏฺฐิวสฺสสหสฺสานิ, \\ อยสาว มลํ สมุฏฺฐิตํ, \\ เนลงฺโค เสตปจฺฉาโท, \\ สุโข วิเวโก ตุฏฺฐสฺส, \\ เผณปิณฺโฑปมํ รูปํ, \\ ปญฺญา หิ เสฏฺฐา โลกสฺมิํ, \\ กาเมสุ สตฺตา กามสงฺคสตฺตา, \\ สพฺเพ ธมฺมา อนตฺตาติ, \\ เอเสว มคฺโค นตฺถญฺโญ, \\ ยสฺสินฺทฺริยานิ สุภาวิตานิ, \\ วจสา มนสา เจว, \\ ยานิ กโรติ ปุริโส, \\ รูปํ เวทยิตํ สญฺญา, \\ โย เอวํ ปฏิปชฺชติ, \\ สพฺพปาปสฺส อกรณํ, \\ นตฺถิ ปุตฺตสมํ เปมํ, \\ ยํ ปตฺตํ ยญฺจ ปตฺตพฺพํ, \\ เย เกจิ ภูตา ภวิสฺสนฺติ, \\ ยญฺจ กามสุขํ โลเก, \\ มคฺคานฏฺฐงฺคิโก เสฏฺโฐ, \\ กาเยน สํวโร สาธุ, \\ นตฺถิ อตฺตสมํ เปมํ, \\ กาเยน กุสลํ อภิรโต, \\ โพชฺฌงฺคา จ สุเทสิตา,}{\csromangathabat{สเจ ภายสิ ทุกฺขสฺส,}{มาภินนฺทิ อนาคตํ.} \\ \csromangathabat{วสฺสกาเล ยถา ฉตฺตํ,}{กุสลานิ กมตฺถเก.} \\ \csromangathabat{สพฺเพ ธมฺมา อนตฺตาติ,}{สมาคตํ วิจาลเย.} \\ \csromangathabat{น โว ทุกฺขา ปโมกฺขาตฺถิ,}{สมโถ จ วิปสฺสนา.} \\ \csromangathabat{กามจฺฉนฺทํ อุปาทาย,}{โย โส วิตกฺเกหิ ขชฺชติ.} \\ \csromangathabat{สุภาวิตตฺเต โพชฺฌงฺเค,}{โส อิมํ วิชฏเย ชฏํ.} \\ \csromangathabat{สุญฺญโต โลกํ อเวกฺขสฺสุ,}{สมาธิภาวิ ภาวเส.} \\ \csromangathabat{กามํ กามยมานสฺส,}{ธมฺมจริยาย สุคติํ.} \\ \csromangathabat{หญฺญเต สพฺพา มุจฺเจว,}{นิปฺโปเฐนฺโต จตุทฺทิสา.} \\ \csromangathabat{โย กาเม ปริวชฺเชติ,}{ปาริฉตฺโตปเมว จ.} \\ \csromangathabat{ยานิ กโรติ ปุริโส,}{โลกธมฺมา ปกาสิตา.} \\ \csromangathabat{สุโข วิปาโก ปุญฺญานํ,}{ตติยํ อญฺญํ น วิชฺชติ.} \\ \csromangathabat{อาทานสฺส ภยํ ญตฺวา,}{ชายเต ชียเตปิ จ.} \\ \csromangathabat{กามา หิ จิตฺรา วิวิธา,}{อถ โลณสลฺโลปมํ.} \\ \csromangathabat{น หิ ปาปํ กตํ กมฺมํ,}{อคตีหิ จ คจฺฉติ.} \\ \csromangathabat{ยสฺสินฺทฺริยานิ สมถงฺคตานิ,}{ตเถว ปญฺจญาณิโก.} \\ \csromangathabat{สตฺติยา วิย โอมฏฺโฐ,}{วิญฺญาณญฺจ ปติฏฺฐิตา.} \\ \csromangathabat{โย ปาณมติปาเตติ,}{ตีณิ ทุจฺจริตานิ จ.} \\ \csromangathabat{สฏฺฐิวสฺสสหสฺสานิ,}{ขณํ ลทฺธาน ทุลฺลภํ.} \\ \csromangathabat{อยสาว มลํ สมุฏฺฐิตํ,}{จตูสุ ปฏิปตฺติสุ.} \\ \csromangathabat{เนลงฺโค เสตปจฺฉาโท,}{อถ โคปาลโกปมํ.} \\ \csromangathabat{สุโข วิเวโก ตุฏฺฐสฺส,}{วิตกฺกา จ สุเทสิตา.} \\ \csromangathabat{เผณปิณฺโฑปมํ รูปํ,}{พฺรหา อุชุ ปตาปวา.} \\ \csromangathabat{ปญฺญา หิ เสฏฺฐา โลกสฺมิํ,}{อนญฺญา ตีณิ อินฺทฺริยานิ.} \\ \csromangathabat{กาเมสุ สตฺตา กามสงฺคสตฺตา,}{อถ วณฺโณ รหสฺสวา.} \\ \csromangathabat{สพฺเพ ธมฺมา อนตฺตาติ,}{อริยสจฺจญฺจ เทสิตํ.} \\ \csromangathabat{เอเสว มคฺโค นตฺถญฺโญ,}{โสตาปนฺโนติ พฺยากเร.} \\ \csromangathabat{ยสฺสินฺทฺริยานิ สุภาวิตานิ,}{อถ ธมฺมปเทหิ จ.} \\ \csromangathabat{วจสา มนสา เจว,}{ปญฺจกฺขนฺธา อนิจฺจโต.} \\ \csromangathabat{ยานิ กโรติ ปุริโส,}{ตีณิ สุจริตานิ จ.} \\ \csromangathabat{รูปํ เวทยิตํ สญฺญา,}{ปญฺจกฺขนฺธา ปกาสิตา.} \\ \csromangathabat{โย เอวํ ปฏิปชฺชติ,}{พฺรหฺมา เจว ผลานิ จ.} \\ \csromangathabat{สพฺพปาปสฺส อกรณํ,}{วิโมกฺขา ตํ หิ เทสิภา.} \\ \csromangathabat{นตฺถิ ปุตฺตสมํ เปมํ,}{เทวานํ อสุราน จ.} \\ \csromangathabat{ยํ ปตฺตํ ยญฺจ ปตฺตพฺพํ,}{นนฺทติ โสจติ นิจฺจํ.} \\ \csromangathabat{เย เกจิ ภูตา ภวิสฺสนฺติ,}{สตฺถาโร จ ปกาสิตา.} \\ \csromangathabat{ยญฺจ กามสุขํ โลเก,}{โพชฺฌงฺคา จ สุเทสิตา.} \\ \csromangathabat{มคฺคานฏฺฐงฺคิโก เสฏฺโฐ,}{ตโย จ อคฺคปตฺติโย.} \\ \csromangathabat{กาเยน สํวโร สาธุ,}{กรณียญฺจ เทสิตํ.} \\ \csromangathabat{นตฺถิ อตฺตสมํ เปมํ,}{อริยา ตีณิ จ เทสิตา.} \\ \csromangathabat{กาเยน กุสลํ อภิรโต,}{วินยญฺจ กามสุขํ โลเก.} \\ \csromangathabat{โพชฺฌงฺคา จ สุเทสิตา,}{ทุทฺทสํ อนตํ เจว ปราปรํ จ.}}
\csromangathasetleft{สเจ ภายสิ ทุกฺขสฺส, \\ วสฺสกาเล ยถา ฉตฺตํ, \\ สพฺเพ ธมฺมา อนตฺตาติ, \\ น โว ทุกฺขา ปโมกฺขาตฺถิ, \\ กามจฺฉนฺทํ อุปาทาย, \\ สุภาวิตตฺเต โพชฺฌงฺเค, \\ สุญฺญโต โลกํ อเวกฺขสฺสุ, \\ กามํ กามยมานสฺส, \\ หญฺญเต สพฺพา มุจฺเจว, \\ โย กาเม ปริวชฺเชติ, \\ ยานิ กโรติ ปุริโส, \\ สุโข วิปาโก ปุญฺญานํ, \\ อาทานสฺส ภยํ ญตฺวา, \\ กามา หิ จิตฺรา วิวิธา, \\ น หิ ปาปํ กตํ กมฺมํ, \\ ยสฺสินฺทฺริยานิ สมถงฺคตานิ, \\ สตฺติยา วิย โอมฏฺโฐ, \\ โย ปาณมติปาเตติ, \\ สฏฺฐิวสฺสสหสฺสานิ, \\ อยสาว มลํ สมุฏฺฐิตํ, \\ เนลงฺโค เสตปจฺฉาโท, \\ สุโข วิเวโก ตุฏฺฐสฺส, \\ เผณปิณฺโฑปมํ รูปํ, \\ ปญฺญา หิ เสฏฺฐา โลกสฺมิํ, \\ กาเมสุ สตฺตา กามสงฺคสตฺตา, \\ สพฺเพ ธมฺมา อนตฺตาติ, \\ เอเสว มคฺโค นตฺถญฺโญ, \\ ยสฺสินฺทฺริยานิ สุภาวิตานิ, \\ วจสา มนสา เจว, \\ ยานิ กโรติ ปุริโส, \\ รูปํ เวทยิตํ สญฺญา, \\ โย เอวํ ปฏิปชฺชติ, \\ สพฺพปาปสฺส อกรณํ, \\ นตฺถิ ปุตฺตสมํ เปมํ, \\ ยํ ปตฺตํ ยญฺจ ปตฺตพฺพํ, \\ เย เกจิ ภูตา ภวิสฺสนฺติ, \\ ยญฺจ กามสุขํ โลเก, \\ มคฺคานฏฺฐงฺคิโก เสฏฺโฐ, \\ กาเยน สํวโร สาธุ, \\ นตฺถิ อตฺตสมํ เปมํ, \\ กาเยน กุสลํ อภิรโต, \\ โพชฺฌงฺคา จ สุเทสิตา,}
\csromangathagroupwithcloser{\csromangathastanza{\csromangathabat{สเจ ภายสิ ทุกฺขสฺส,}{มาภินนฺทิ อนาคตํ.} \\ \csromangathabat{วสฺสกาเล ยถา ฉตฺตํ,}{กุสลานิ กมตฺถเก.}}\csromangathastanza{\csromangathabat{สพฺเพ ธมฺมา อนตฺตาติ,}{สมาคตํ วิจาลเย.} \\ \csromangathabat{น โว ทุกฺขา ปโมกฺขาตฺถิ,}{สมโถ จ วิปสฺสนา.}}\csromangathastanza{\csromangathabat{กามจฺฉนฺทํ อุปาทาย,}{โย โส วิตกฺเกหิ ขชฺชติ.} \\ \csromangathabat{สุภาวิตตฺเต โพชฺฌงฺเค,}{โส อิมํ วิชฏเย ชฏํ.}}\csromangathastanza{\csromangathabat{สุญฺญโต โลกํ อเวกฺขสฺสุ,}{สมาธิภาวิ ภาวเส.} \\ \csromangathabat{กามํ กามยมานสฺส,}{ธมฺมจริยาย สุคติํ.}}\csromangathastanza{\csromangathabat{หญฺญเต สพฺพา มุจฺเจว,}{นิปฺโปเฐนฺโต จตุทฺทิสา.} \\ \csromangathabat{โย กาเม ปริวชฺเชติ,}{ปาริฉตฺโตปเมว จ.}}\csromangathastanza{\csromanfolio{208}\csromangathabat{ยานิ กโรติ ปุริโส,}{โลกธมฺมา ปกาสิตา.} \\ \csromangathabat{สุโข วิปาโก ปุญฺญานํ,}{ตติยํ อญฺญํ น วิชฺชติ.}}\csromangathastanza{\csromangathabat{อาทานสฺส ภยํ ญตฺวา,}{ชายเต ชียเตปิ จ.} \\ \csromangathabat{กามา หิ จิตฺรา วิวิธา,}{อถ โลณสลฺโลปมํ.}}\csromangathastanza{\csromangathabat{น หิ ปาปํ กตํ กมฺมํ,}{อคตีหิ จ คจฺฉติ.} \\ \csromangathabat{ยสฺสินฺทฺริยานิ สมถงฺคตานิ,}{ตเถว ปญฺจญาณิโก.}}\csromangathastanza{\csromangathabat{สตฺติยา วิย โอมฏฺโฐ,}{วิญฺญาณญฺจ ปติฏฺฐิตา.} \\ \csromangathabat{โย ปาณมติปาเตติ,}{ตีณิ ทุจฺจริตานิ จ.}}\csromangathastanza{\csromangathabat{สฏฺฐิ\-วสฺสสหสฺสานิ,}{ขณํ ลทฺธาน ทุลฺลภํ.} \\ \csromangathabat{อยสาว มลํ สมุฏฺฐิตํ,}{จตูสุ ปฏิปตฺติสุ.}}\csromangathastanza{\csromangathabat{เนลงฺโค เสตปจฺฉาโท,}{อถ โคปาลโกปมํ.} \\ \csromangathabat{สุโข วิเวโก ตุฏฺฐสฺส,}{วิตกฺกา จ สุเทสิตา.}}\csromangathastanza{\csromangathabat{เผณปิณฺโฑปมํ รูปํ,}{พฺรหา อุชุ ปตาปวา.} \\ \csromangathabat{ปญฺญา หิ เสฏฺฐา โลกสฺมิํ,}{อนญฺญา ตีณิ อินฺทฺริยานิ.}}\csromangathastanza{\csromangathabat{กาเมสุ สตฺตา กามสงฺคสตฺตา,}{อถ วณฺโณ รหสฺสวา.} \\ \csromangathabat{สพฺเพ ธมฺมา อนตฺตาติ,}{อริยสจฺจญฺจ เทสิตํ.}}\csromangathastanza{\csromangathabat{เอเสว มคฺโค นตฺถญฺโญ,}{โสตาปนฺโนติ พฺยากเร.} \\ \csromangathabat{ยสฺสินฺทฺริยานิ สุภาวิตานิ,}{อถ ธมฺมปเทหิ จ.}}\csromangathastanza{\csromangathabat{วจสา มนสา เจว,}{ปญฺจกฺขนฺธา อนิจฺจโต.} \\ \csromangathabat{ยานิ กโรติ ปุริโส,}{ตีณิ สุจริตานิ จ.}}\csromangathastanza{\csromangathabat{รูปํ เวทยิตํ สญฺญา,}{ปญฺจกฺขนฺธา ปกาสิตา.} \\ \csromangathabat{โย เอวํ ปฏิปชฺชติ,}{พฺรหฺมา เจว ผลานิ จ.}}\csromangathastanza{\csromangathabat{สพฺพปาปสฺส อกรณํ,}{วิโมกฺขา ตํ หิ เทสิภา.} \\ \csromangathabat{นตฺถิ ปุตฺตสมํ เปมํ,}{เทวานํ อสุราน จ.}}\csromangathastanza{\csromangathabat{ยํ ปตฺตํ ยญฺจ ปตฺตพฺพํ,}{นนฺทติ โสจติ นิจฺจํ.} \\ \csromangathabat{เย เกจิ ภูตา ภวิสฺสนฺติ,}{สตฺถาโร จ ปกาสิตา.}}\csromangathastanza{\csromanfolio{209}\csromangathabat{ยญฺจ กามสุขํ โลเก,}{โพชฺฌงฺคา จ สุเทสิตา.} \\ \csromangathabat{มคฺคานฏฺฐงฺคิโก เสฏฺโฐ,}{ตโย จ อคฺคปตฺติโย.}}\csromangathastanza{\csromangathabat{กาเยน สํวโร สาธุ,}{กรณียญฺจ เทสิตํ.} \\ \csromangathabat{นตฺถิ อตฺตสมํ เปมํ,}{อริยา ตีณิ จ เทสิตา.}}\csromangathastanza{\csromangathabat{กาเยน กุสลํ อภิรโต,}{วินยญฺจ กามสุขํ โลเก.} \\ \csromangathabat{โพชฺฌงฺคา จ สุเทสิตา,}{ทุทฺทสํ อนตํ เจว ปราปรํ จ.}}}{เปฏโกปเทเส สาสน\-ปฺปฏฺฐานํ นาม ทุติยภูมิ สมตฺตา.}
\proseitem{32}{ตตฺถ กตมํ สุตฺตา\-ธิฏฺฐานํ.}

\prose{โลภา\-ธิฏฺฐานํ โทสาธิฏฺฐานํ โมหาธิฏฺฐานํ อโลภา\-ธิฏฺฐานํ อโทสา\-ธิฏฺฐานํ อโมหา\-ธิฏฺฐานํ กายกมฺมา\-ธิฏฺฐานํ วาจากมฺมา\-ธิฏฺฐานํ มโนกมฺมา\-ธิฏฺฐานํ สทฺธินฺทฺริยา\-ธิฏฺฐานํ วีริยินฺทฺริยา\-ธิฏฺฐานํ สตินฺทฺริยา\-ธิฏฺฐานํ สมาธินฺทฺริยา\-ธิฏฺฐานํ ปญฺญินฺทฺริยา\-ธิฏฺฐานํ.}

\setlength{\csromangathapagewidth}{0pt}
\setlength{\csromangathawakwidth}{0pt}
\csromangathameasuregroup{}{ตตฺถ กตมํ โลภาธิฏฺฐานํ. \\ วิตกฺกมถิตสฺส\textsuperscript{1} ชนฺตุโน, \\ ติพฺพราคสฺส สุภานุปสฺสิโน. \\ ภิยฺโย ตณฺหา ปวฑฺฒติ, \\ เอส โข คาฬฺหํ กโรติ พนฺธนํ.}
\csromangathameasurewak{ตตฺถ กตมํ โลภาธิฏฺฐานํ. \\ วิตกฺกมถิตสฺส\textsuperscript{1} ชนฺตุโน, \\ ติพฺพราคสฺส สุภานุปสฺสิโน. \\ ภิยฺโย ตณฺหา ปวฑฺฒติ, \\ เอส โข คาฬฺหํ กโรติ พนฺธนํ.}
\setlength{\csromangathaleftcol}{0pt}
\csromangathagroup{\csromangathastanza{\csromangathawak{\csromangathaitemline{32}{ตตฺถ กตมํ โลภา\-ธิฏฺฐานํ.} \\ \csromangathacontline{วิตกฺก\-มถิตสฺส\csromansharedfootnote{b9c72dbd36f83ba4}{วิตกฺกนิมฺมถิตสฺส (ก) ขุ 1. 63 ธมฺมปเท.} ชนฺตุโน,} \\ \csromangathacontline{ติพฺพราคสฺส สุภานุปสฺสิโน.} \\ \csromangathacontline{ภิยฺโย ตณฺหา ปวฑฺฒติ,} \\ \csromangathacontline{เอส โข คาฬฺหํ กโรติ พนฺธนํ.}}}}

\prose{วิตกฺกมถิตสฺสาติ กามราโค. สุภา\-นุปสฺสิโนติ กามราควตฺถุ. ภิยฺโย ตณฺหา ปวฑฺฒตีติ กามตณฺหา. เอส คาฬฺหํ กโรติ พนฺธนนฺติ ราคํ, อิติ โย โย ธมฺโม มูลนิกฺขิตฺโต, โส เยเวตฺถ ธมฺโม อุคฺคาวหิตพฺโพ\csromansharedfootnote{8e56e870d96f9934}{อุคฺคาปยิตพฺโพ (อิ, ก)}. น ภควา เอกํ ธมฺมํ อารพฺภ อญฺญํ ธมฺมํ เทเสติ. ยสฺส วิตกฺเกติ กามวิตกฺโก ตเมว วิตกฺกํ กามวิตกฺเกน นิทฺทิสียติ. ติพฺพราคสฺสาติ ตสฺเสว วิตกฺกสฺส วตฺถุํ นิทฺทิสติ. สุภานุปสฺสิโน ภิยฺโย ตณฺหา ปวฑฺฒตีติ ตเมว ราคํ กามตณฺหาติ นิทฺทิสติ. เอส คาฬฺหํ กโรติ พนฺธนนฺติ ตเมว ตณฺหา\-สํโยชนํ นิทฺทิสติ. เอวํ คาถาสุ อนุมินิตพฺพํ. เอวํ สเวยฺยากรเณสุ.}

\prose{\csromanfolio{210}ตตฺถ ภควา เอกํ ธมฺมํ ติวิธํ นิทฺทิสติ, นิสฺสนฺทโต เหตุโต ผลโต.}

\setlength{\csromangathapagewidth}{0pt}
\setlength{\csromangathawakwidth}{0pt}
\csromangathameasuregroup{}{ททํ ปิโย โหติ ภชนฺติ นํ พหู, \\ กิตฺติญฺจ ปปฺโปติ ยโส จ วฑฺฒติ. \\ อมงฺกุภูโต ปริสํ วิคาหติ, \\ วิสารโท โหติ นโร อมจฺฉรี\textsuperscript{1}.}
\csromangathameasurewak{ททํ ปิโย โหติ ภชนฺติ นํ พหู, \\ กิตฺติญฺจ ปปฺโปติ ยโส จ วฑฺฒติ. \\ อมงฺกุภูโต ปริสํ วิคาหติ, \\ วิสารโท โหติ นโร อมจฺฉรี\textsuperscript{1}.}
\setlength{\csromangathaleftcol}{0pt}
\csromangathagroup{\csromangathastanza{\csromangathawak{ททํ ปิโย โหติ ภชนฺติ นํ พหู, \\ กิตฺติญฺจ ปปฺโปติ ยโส จ วฑฺฒติ. \\ อมงฺกุภูโต ปริสํ วิคาหติ, \\ วิสารโท โหติ นโร อมจฺฉรี\csromansharedfootnote{e9191470b17e9eb4}{สํ 2. 34 ปิฏฺเฐ.}.}}}

\prose{ททนฺติ ยํ ยํ ทานํ, อิทํ ทานมยิกํ ปุญฺญกฺริยํ. ตตฺถ เหตุ. ยํ เจตํ. ภชนฺติ นํ พหู, กิตฺตินฺติ โย จ กลฺยาโณ กิตฺติสทฺโท โลเก อพฺภุคฺคจฺฉติ, ยํ พหุกสฺส ชนสฺส ปิโย ภวติ มนาโป จ. ยญฺจ อวิปฺปฏิสารี กาลงฺกโรติ อยํ นิสฺสนฺโท. ยํ กายสฺส เภทา เทเวสุ อุปปชฺชตีติ อิทํ ผลํ. อิทํ โลภา\-ธิฏฺฐานํ.}

\proseitem{33}{ตตฺถ กตมํ โทสาธิฏฺฐานํ.}

\prose{\csromansymbolfootnote{*}{เหฏฺฐา 201 ปิฏฺเฐ.}โย ปาณมติปาเตติ, มุสาวาทญฺจ ภาสติ.}

\setlength{\csromangathapagewidth}{0pt}
\setlength{\csromangathawakwidth}{0pt}
\csromangathameasuregroup{โลเก อทินฺนํ อาทิยติ, \\ สุราเมรยปานญฺจ, \\ อปฺปหาย ปญฺจ เวทฺรานิ, \\ กายสฺส เภทา ทุปฺปญฺโญ,}{\csromangathabat{โลเก อทินฺนํ อาทิยติ,}{ปรทารญฺจ คจฺฉติ.} \\ \csromangathabat{สุราเมรยปานญฺจ,}{โย นโร อนุยุญฺชติ\textsuperscript{1}.} \\ \csromangathabat{อปฺปหาย ปญฺจ เวทฺรานิ,}{ทุสฺสีโล อิติ วุจฺจติ.} \\ \csromangathabat{กายสฺส เภทา ทุปฺปญฺโญ,}{นิรยํ โสปปชฺชติ.}}
\csromangathasetleft{โลเก อทินฺนํ อาทิยติ, \\ สุราเมรยปานญฺจ, \\ อปฺปหาย ปญฺจ เวทฺรานิ, \\ กายสฺส เภทา ทุปฺปญฺโญ,}
\csromangathagroup{\csromangathastanza{\csromangathabat{โลเก อทินฺนํ อาทิยติ,}{ปรทารญฺจ คจฺฉติ.} \\ \csromangathabat{สุรา\-เมรย\-ปานญฺจ,}{โย นโร อนุยุญฺชติ\csromansharedfootnote{06942e5887bce717}{อภิคิชฺฌติ (อิ, ก) อํ 2. 181 ปิฏฺเฐ.}.}}\csromangathastanza{\csromangathabat{อปฺปหาย ปญฺจ เวทฺรานิ,}{ทุสฺสีโล อิติ วุจฺจติ.} \\ \csromangathabat{กายสฺส เภทา ทุปฺปญฺโญ,}{นิรยํ โสปปชฺชติ.}}}

\prose{โย ปาณมติปาเตตีติ ทุฏฺโฐ ปาณมติปาเตติ. มุสาวาทญฺจ ภาสตีติ โทโสปฆาตาย มุสาวาทญฺจ ภาสติ. สุรา\-เมรย\-ปานญฺจ, โย นโร อนุยุญฺชตีติ โทโส นิทานํ. โย จ สุรา\-เมรย\-ปานํ อนุยุญฺชติ ยถา\-ปรทาร\-วิหารี\csromansharedfootnote{1f1673bbf7e3cb60}{ยถาปมุทิตวิหารี (ก)} อมิตฺตา ชนยนฺติ.}

\prose{ปญฺจ เวรานิ อปฺปหายาติ ปญฺจนฺนํ สิกฺขาปทานํ สมติกฺกมนํ สพฺเพสํ โทสชานํ สา ปณฺณตฺติ, เตเนว โทสชนิเตน กมฺเมน ทุสฺสีโล อิติ วุจฺจติ โสปิ ธมฺโม เหตุนา นิทฺทิสิตพฺโพ, นิสฺสนฺเทน ผเลน จ.}

\prose{ตีณิ พาลสฺส พาลลกฺขณานิ ทุพฺภาสิต\-ภาสี\csromansharedfootnote{0f111363c0510176}{ทุพฺภาสิตภาสิตา (อิ, ก) อํ 1. 100 ปิฏฺเฐ.} จ โหติ, ทุจฺจินฺติต\-จินฺตี จ ทุกฺกฏ\-กมฺม\-การี จ. ตตฺถ ยํ กาเยน จ วาจาย จ ปรกฺกมติ, อิทมสฺส ทุกฺกฏ\-กมฺม\-การี. ตายํ ยถา จ มุสาวาทํ ภาสติ ยถา \csromanfolio{211}ปุพฺพ\-นิทฺทิฏฺฐํ, อิทมสฺส ทุพฺภาสิตา. ยญฺจ สงฺกปฺเปติ มโนทุจฺจริตํ พฺยาปาทํ. อิทมสฺส ทุจฺจินฺติต\-จินฺติตา, ยํ โส อิเมหิ ตีหิ พาลลกฺขเณหิ สมนฺนาคโต ตีณิ ตชฺชานิ ทุกฺขานิ โทมนสฺสานิ อนุภวติ, โส จ โหติ สภคฺคโต วา ปริสคฺคโต วา ตชฺชํ กถํ กถนฺติ. ยทา ภวติ โส จ ปาณาติปาตา\-ทิ\-ทส\-อกุสลกมฺมปโถ, โส ตโตนิทานํ ทุกฺขํ โทมนสฺสํ ปฏิสํเวเทตีติ. ปุน จปรํ ยทา ปสฺสติ โจรํ ราชาปราธิกํ รญฺญา คหิตํ ชีวิตา โวโรเปตํ, ตสฺเสวํ ภวติ สเจ มมมฺปิ ราชา ชาเนยฺย มมมฺปิ ราชา คาหาเปตฺวา ชีวิตา โวโรเปยฺยาติ, โส ตโตนิทานํ ทุกฺขํ โทมนสฺสํ ปฏิสํเวเทติ. ปุน จปรํ พาโล ยทา ภวติ อาสนา สมารูโฬฺห ยาว ยา เม คติ ภวิสฺสติ อิโต เปจฺจ ปรํ มรณาติ โส ตโตนิทานํ ทุกฺขํ โทมนสฺสํ ปฏิสํเวเทติ อิติ พาลลกฺขณํ เหตุ. ตีณิ ตชฺชานิ ทุกฺขานิ นิสฺสนฺโท. กายสฺส เภทา นิรเยสุ อุปปชฺชติ, อิทํ ผลํ.  อิทํ โทสาธิฏฺฐานํ.}

\setlength{\csromangathapagewidth}{0pt}
\setlength{\csromangathawakwidth}{0pt}
\csromangathameasuregroup{}{ตตฺถ กตมํ โมหาธิฏฺฐานํ. \\ สตญฺเจว สหสฺสานํ, กปฺปานํ สํสริสฺสติ. \\ อถวาปิ ตโต ภิยฺโย, คพฺภา คพฺภํ คมิสฺสถ. \\ อนุปาทาย พุทฺธวจนํ, สงฺขาเร อตฺตโต อุปาทาย. \\ ทุกฺขสฺสนฺตํ กริสฺสนฺติ, ฐานเมตํ น วิชฺชติ.}
\csromangathameasurewak{ตตฺถ กตมํ โมหาธิฏฺฐานํ. \\ สตญฺเจว สหสฺสานํ, กปฺปานํ สํสริสฺสติ. \\ อถวาปิ ตโต ภิยฺโย, คพฺภา คพฺภํ คมิสฺสถ. \\ อนุปาทาย พุทฺธวจนํ, สงฺขาเร อตฺตโต อุปาทาย. \\ ทุกฺขสฺสนฺตํ กริสฺสนฺติ, ฐานเมตํ น วิชฺชติ.}
\setlength{\csromangathaleftcol}{0pt}
\csromangathagroup{\csromangathastanza{\csromangathawak{\csromangathaitemline{34}{ตตฺถ กตมํ โมหาธิฏฺฐานํ.} \\ \csromangathacontline{สตญฺเจว สหสฺสานํ, กปฺปานํ สํสริสฺสติ.} \\ \csromangathacontline{อถวาปิ ตโต ภิยฺโย, คพฺภา คพฺภํ คมิสฺสถ.} \\ \csromangathacontline{อนุปาทาย พุทฺธวจนํ, สงฺขาเร อตฺตโต อุปาทาย.} \\ \csromangathacontline{ทุกฺขสฺสนฺตํ กริสฺสนฺติ, ฐานเมตํ น วิชฺชติ.}}}}

\prose{โย ยํ อนมตคฺค\-สํสารํ สมาปนฺโน ชายเต จ มียเต จ, อยํ อวิชฺชาเหตุกา. ยานิปิ จ สงฺขารานํ ปโยชนานิ, ตานิปิ อวิชฺชา\-ปจฺจยานิ, ยํ อทสฺสนํ พุทฺธ\-วจนสฺส, อยํ อวิชฺชา\-สุตฺเตเยว นิทฺทิฏฺฐํ. โย จ สงฺขาเร อตฺตโต หรติ ปญฺจกฺขนฺเธ ปญฺจ ทิฏฺฐิโย อุปคจฺฉติ. “เอตํ มม, เอโสหมสฺมิ, เอโส เม อตฺตา”ติ อิทํ สุตฺตํ อวิชฺชาย นิกฺขิตฺตํ, อวิชฺชาย นิกฺขิปิตํ. เอวํ สตฺถา สุตฺเต นเยน\csromansharedfootnote{c4fb3579ee9ac8c2}{สุตนเยน (อิ)} ธมฺเมน นิทฺทิสติ. อสาธารเณน ตํเยว ตตฺถ นิทฺทิสิตพฺพํ. น อญฺญํ.}

\prose{เย หิ เกจิ ภิกฺขเว สมณา วา พฺราหฺมณา วา “อิทํ ทุกฺขนฺ”ติ นปฺปชานนฺติ จตฺตาริ สจฺจานิ วิตฺถาเรน, ยํ ตตฺถ อปฺปชานนา, อิทํ ทุกฺขํ, อยํ เหตุ. อปฺปชานนฺโต วิวิเธ สงฺขาเร อภิสงฺขโรติ, อยํ นิสฺสนฺโท. ยญฺจ ทิฏฺฐิคตานิ ปรามสติ “อิทเมว สจฺจํ โมฆมญฺญนฺ”ติ อยํ นิสฺสนฺโท. \csromanfolio{212}ยํ ปุนพฺภวํ นิพฺพตฺเตติ, อิทํ ผลํ. อยมฺปิ ธมฺโม สนิทฺทิฏฺโฐ เหตุโต จ ผลโต จ นิสฺสนฺทโต จ.}

\prose{เอตฺถ ปน เกจิ ธมฺมา สาธารณา ภวนฺติ. เหตุ ขลุ อาทิโตเยว สุตฺเต นิกฺขิปิสฺสนฺติ. ยถา กิํ ภเว จตฺตาริมานิ ภิกฺขเว อคติคมนานิ ตตฺถ ยญฺจ ฉนฺทาคติํ คจฺฉติ ยญฺจ ภยาคติํ คจฺฉติ, อยํ โลโภ อกุสลมูลํ. ยํ โทสา, อยํ โทโสเยว. ยํ โมหา, อยํ โมโหเยว. เอวํ อิมานิ ตีณิ อกุสลมูลานิ อาทิโตเยว อุปปริ\-กฺขิ\-ตพฺพานิ. ยตฺถ เอกํ นิทฺทิสิตพฺพํ, ตตฺถ เอกํ นิทฺทิสียติ. ตถา เทฺว ยถา ตีณิ, น หิ อาทีหิ อนิกฺขิตฺเต เหตุ วา นิสฺสนฺโท วา ผลํ วา นิทฺทิสิตพฺพํ. อยญฺเจตฺถ คาถา.}

\setlength{\csromangathapagewidth}{0pt}
\setlength{\csromangathawakwidth}{0pt}
\csromangathameasuregroup{\textsuperscript{*}ฉนฺทา โทสา ภยา โมหา, \\ นิหียติ\textsuperscript{1} ตสฺส ยโส,}{\csromangathabat{\textsuperscript{*}ฉนฺทา โทสา ภยา โมหา,}{โย ธมฺมํ อติวตฺตติ.} \\ \csromangathabat{นิหียติ\textsuperscript{1} ตสฺส ยโส,}{กาฬปกฺเขว จนฺทิมา.}}
\csromangathasetleft{\textsuperscript{*}ฉนฺทา โทสา ภยา โมหา, \\ นิหียติ\textsuperscript{1} ตสฺส ยโส,}
\csromangathagroup{\csromangathastanza{\csromangathaitembat{34}{\csromansymbolfootnote{*}{เหฏฺฐา 107 ปิฏฺเฐ.}ฉนฺทา โทสา ภยา โมหา,}{โย ธมฺมํ อติวตฺตติ.} \\ \csromangathacontbat{นิหียติ\csromansharedfootnote{1cf84155d040f3cf}{นิหียเต (อิ, ก) อํ 1. 352 ปิฏฺเฐ.} ตสฺส ยโส,}{กาฬปกฺเขว จนฺทิมา.}}}

\prose{กตฺถ ฉนฺทา จ อยํ โลโภ ยถา นิทฺทิฏฺฐํ ปุพฺเพ. อิทํ โมหาธิฏฺฐานํ.}

\proseitem{35}{ตตฺถ กตมํ อโลภา\-ธิฏฺฐานํ.}

\setlength{\csromangathapagewidth}{0pt}
\setlength{\csromangathawakwidth}{0pt}
\csromangathameasuregroup{“อสุภานุปสฺสิํ\textsuperscript{1} วิหรนฺตํ, \\ โภชนมฺหิ จ มตฺตญฺญุํ, \\ ตํ เว นปฺปสหติ มาโร,}{\csromangathabat{“อสุภานุปสฺสิํ\textsuperscript{1} วิหรนฺตํ,}{อินฺทฺริเยสุ สุสํวุตํ.} \\ \csromangathabat{โภชนมฺหิ จ มตฺตญฺญุํ,}{สทฺธํ อารทฺธวีริยํ.} \\ \csromangathabat{ตํ เว นปฺปสหติ มาโร,}{วาโต เสลํว ปพฺพตนฺ”ติ.}}
\csromangathasetleft{“อสุภานุปสฺสิํ\textsuperscript{1} วิหรนฺตํ, \\ โภชนมฺหิ จ มตฺตญฺญุํ, \\ ตํ เว นปฺปสหติ มาโร,}
\csromangathagroup{\csromangathastanza{\csromangathabat{“อสุภานุปสฺสิํ\csromansharedfootnote{f13966cc3451e2bd}{อสุภานุปสฺสี (อิ) ขุ 1. 14 ธมฺมปเท.} วิหรนฺตํ,}{อินฺทฺริเยสุ สุสํวุตํ.} \\ \csromangathabat{โภชนมฺหิ จ มตฺตญฺญุํ,}{สทฺธํ อารทฺธวีริยํ.} \\ \csromangathabat{ตํ เว นปฺปสหติ มาโร,}{วาโต เสลํว ปพฺพตนฺ”ติ.}}}

\prose{ตตฺถ ยา อสุภาย อุปปริกฺขา, อยํ กาเมสุ อาทีนว\-ทสฺสเนน ปริจฺจาโค. อินฺทฺริเยสุ สุสํวุโต ตสฺเสว อโลภสฺส ปาริปูริยํ มม อายตนโสจิตํ อนุปาทาย. โภชนมฺหิ จ มตฺตญฺญุนฺติ รสตณฺหา\-ปหานํ. อิติ อยํ อโลโภ อสุภานุปสฺสิ\-ตาย วตฺถุโต ธารยติ, โส อโลโภ เหตุ. อินฺทฺริเยสุ คุตฺตทฺวารตาย โคจรโต ธารยติ, โภชเน\-มตฺตญฺญุตาย ปรโต ธารยติ, อยํ นิสฺสนฺโท. ตํ เว นปฺปสหติ มาโร, วาโต เสลํ ว ปพฺพตนฺติ, อิทํ ผลํ. อิติ โยเยว ธมฺโม อาทิมฺหิ นิกฺขิตฺโต, โสเยว มชฺเฌ เจว อวสาเน จ.}

\prose{นาหํ ภิกฺเขเว อญฺญํ เอกธมฺมมฺปิ สมนุปสฺสามิ อสมุปฺปนฺนสฺส กามจฺฉนฺทสฺส อนุปฺปาทาย อุปฺปนฺนสฺส วา ปหานาย, ยถยิทํ\csromansharedfootnote{dc0871dbc0b28398}{ยทิทํ (อิ, ก) อํ 1. 3 ปิฏฺเฐ.} อสุภนิมิตฺตํ. ตตฺถ \csromanfolio{213}อสุภนิมิตฺตํ มนสิ\-กโรนฺตสฺส อนุปฺปนฺโน เจว กามจฺฉนฺโท น อุปฺปชฺชติ, อุปฺปนฺโน จ กามจฺฉนฺโท ปหียติ. อิทํ อโลภสฺส วตฺถุ. ยํ ปุน อนุปฺปนฺโน กามราโค ปริยาทิยติ รูปราคํ อรูปราคํ, อิติ ผลํ. อิติ อยมฺปิ จ ธมฺโม นิทฺทิฏฺโฐ เหตุโต จ นิสฺสนฺทโต จ ผลโต จ. อิทํ อโลภา\-ธิฏฺฐานํ.}

\setlength{\csromangathapagewidth}{0pt}
\setlength{\csromangathawakwidth}{0pt}
\csromangathameasuregroup{}{ตตฺถ กตมํ อโทสาธิฏฺฐานํ. \\ เอกมฺปิ เจ ปาณมทุฏฺฐจิตฺโต, \\ เมตฺตายติ กุสโล\textsuperscript{1} เตน โหติ. \\ สพฺเพ จ ปาเณ มนสานุกมฺปํ\textsuperscript{1}, \\ ปหูตมริโย ปกโรติ ปุญฺญํ.}
\csromangathameasurewak{ตตฺถ กตมํ อโทสาธิฏฺฐานํ. \\ เอกมฺปิ เจ ปาณมทุฏฺฐจิตฺโต, \\ เมตฺตายติ กุสโล\textsuperscript{1} เตน โหติ. \\ สพฺเพ จ ปาเณ มนสานุกมฺปํ\textsuperscript{1}, \\ ปหูตมริโย ปกโรติ ปุญฺญํ.}
\setlength{\csromangathaleftcol}{0pt}
\csromangathagroup{\csromangathastanza{\csromangathawak{\csromangathaitemline{35}{ตตฺถ กตมํ อโทสา\-ธิฏฺฐานํ.} \\ \csromangathacontline{เอกมฺปิ เจ ปาณม\-ทุฏฺฐจิตฺโต,} \\ \csromangathacontline{เมตฺตายติ กุสโล\csromansharedfootnote{b3f5f76d7f700426}{กุสลี (ก) ขุ 1. 209 ปิฏฺเฐ.} เตน โหติ.} \\ \csromangathacontline{สพฺเพ จ ปาเณ มนสานุกมฺปํ\csromansharedfootnote{e08d5d3c2cd54323}{อนุกมฺปมาโน (อิ)},} \\ \csromangathacontline{ปหูตมริโย ปกโรติ ปุญฺญํ.}}}}

\prose{เอกมฺปิ เจ ปาณม\-ทุฏฺฐจิตฺโต เมตฺตายตีติ อยํ อโทโส. นิคฺฆาเตน อสฺสาโท, กุสโล เตน โหตีติ เตน กุสเลน ธมฺเมน สํยุตฺโต ธมฺม\-ปญฺญตฺติํ คจฺฉติ. กุสโลติ ยถา ปญฺญาย ปญฺโญ ปณฺฑิจฺเจน ปณฺฑิโต. ปหูตมริโย ปกโรติ ปุญฺญนฺติ ตสฺสาเยว วิปาโก อยํ โลกิยสฺส, น หิ โลกุตฺตรสฺส. ตตฺถ ยา เมตฺตายนา, อยํ เหตุ. ยํ กุสโล ภวติ อยํ นิสฺสนฺโท. ยาว อพฺยาปชฺโช ภูมิยํ พหุปุญฺญํ ปสวติ, อิทํ ผลํ. อิติ อโทโส นิทฺทิฏฺโฐ เหตุโต จ นิสฺสนฺทโต จ ผลโต จ.}

\prose{เอกาทสานิสํสา เมตฺตาย เจโตวิมุตฺติยา. ตตฺถ ยา เมตฺตา\-เจโตวิมุตฺติ, อยํ อริยธมฺเมสุ ราควิราคา เจโตวิมุตฺติ, โลกิกาย ภูมิกา เหตุ, ยํ สุขํ อายติํ มนาโป โหติ มนุสฺสานํ, อิเม เอกาทส ธมฺมา นิสฺสนฺโท. ยญฺจ อกตาวี พฺรหฺมกาเย อุปปชฺชติ. อิทํ ผลํ. อิทํ อโทสา\-ธิฏฺฐานํ.}

\setlength{\csromangathapagewidth}{0pt}
\setlength{\csromangathawakwidth}{0pt}
\csromangathameasuregroup{}{ตตฺถ กตมํ อโมหาธิฏฺฐานํ. \\ ปญฺญา หิ เสฏฺฐา โลกสฺมิํ, ยายํ นิพฺเพธคามินี\textsuperscript{1}. \\ ยาย สมฺมา ปชานาติ, ชาติมรณสงฺขยํ.}
\csromangathameasurewak{ตตฺถ กตมํ อโมหาธิฏฺฐานํ. \\ ปญฺญา หิ เสฏฺฐา โลกสฺมิํ, ยายํ นิพฺเพธคามินี\textsuperscript{1}. \\ ยาย สมฺมา ปชานาติ, ชาติมรณสงฺขยํ.}
\setlength{\csromangathaleftcol}{0pt}
\csromangathagroup{\csromangathastanza{\csromangathawak{\csromangathaitemline{36}{ตตฺถ กตมํ อโมหา\-ธิฏฺฐานํ.} \\ \csromangathacontline{ปญฺญา หิ เสฏฺฐา โลกสฺมิํ, ยายํ นิพฺเพธคามินี\csromansharedfootnote{fecfce6ecd1ca705}{นิพฺเพธภาคินี (อิ, ก) ขุ 1. 219; เหฏฺฐา 143, 203 ปิฏฺเฐสุปิ.}.} \\ \csromangathacontline{ยาย สมฺมา ปชานาติ, ชาติมรณสงฺขยํ.}}}}

\prose{\csromanfolio{214}ปญฺญา หิ เสฏฺฐาติ วตฺถุํ. นิพฺเพธคามินีติ นิพฺพานคามินิยํ ยถาภูตํ ปฏิวิชฺฌติ. สมฺมา ปชานาติ, ชาติ\-มรณ\-สงฺขยนฺติ อโมโห. ปญฺญาติ เหตุ. ยํ ปชานาติ อยํ นิสฺสนฺโท. โย ชาติ\-มรณ\-สงฺขโย, อิทํ ปลํ. อิติ อโมโห นิทฺทิฏฺโฐ เหตุนา จ นิสฺสนฺเทน จ ผเลน จ.}

\prose{ตีณิมานิ ภิกฺขเว อินฺทฺริยานิ อนญฺญาต\-ญฺญสฺสามี\-ตินฺทฺริยํ อญฺญินฺทฺริยํ อญฺญาตาวิ\-นฺทฺริยํ\csromansharedfootnote{33b600b98abe4f45}{ขุ 1. 231 ติกนิปาเต.}. ตตฺถ กตมํ อนญฺญาต\-ญฺญสฺสามี\-ตินฺทฺริยํ, อิธ ภิกฺขเว ภิกฺขุ อนภิสเมตสฺส ทุกฺขสฺส อริยสจฺจสฺส อภิสมยาย ฉนฺทํ ชเนติ วายมติ, วีริยํ อารภติ, จิตฺตํ ปคฺคณฺหาติ ปทหติ. เอวํ จตุนฺนํ อริยสจฺจานํ กาตพฺพํ. ตตฺถ กตมํ อญฺญินฺทฺริยํ, อิธ ภิกฺขเว ภิกฺขุ “อิทํ ทุกฺขํ อริยสจฺจนฺ”ติ ยถาภูตํ ปชานาติ, ยา จ มคฺโค, อิทํ อญฺญินฺทฺริยํ. อาสวกฺขยา อนาสโว โหติ, อิทํ วุจฺจติ อญฺญาตาวิ\-นฺทฺริยํ. ตถายํ ปญฺญา, อยํ เหตุ. ยํ ฉนฺทํ ชเนติ วายมติ, ยา ปชานาติ, อยํ นิสฺสนฺโท. เยน สพฺพโส อาสวานํ ขยา เหตุ, ยํ ขเย ญาณมุปฺปชฺชติ, อนุปฺปาเท ญาณญฺจ, อยํ นิสฺสนฺโท. ยํ อรหตฺตํ, อิทํ ผลํ. ตตฺถ ขีณา เม ชาติ, วุสุตํ พฺรหฺมจริยํ, กตํ กรณียนฺติ, อิทํ ขเย ญาณํ. นาปรํ อิตฺถตฺตายาติ ปชานามีติ อิทํ อนุปฺปาเท ญาณํ. อิติ อิมานิ อินฺทฺริยานิ อโมโห นิทฺทิฏฺโฐ เหตุนา จ นิสฺสนฺเทน จ ผเลน จ. อิมานิ อสาธารณานิ นิทฺทิฏฺฐานิ.}

\prose{ตตฺถ กตมานิ กุสลมูลานิ สาธารณานิ. กุสลญฺจ โว ภิกฺขเว เทเสสฺสามิ กุสลมูล\-ญฺเจว. ตตฺถ กตมํ กุสลมูลํ, อโลโภ อโทโส อโมโห. ตตฺถ กตมํ กุสลํ, อฏฺฐ สมฺมตฺตานิ สมฺมาทิฏฺฐิ ยาว สมฺมาสมาธิ. ตตฺถ ยานิ กุสลมูลานิ, อยํ เหตุ. ยญฺจ อโลโภ ตีณิ กมฺมานิ สมุฏฺฐาเปติ สงฺกปฺปํ วายามํ สมาธิญฺจ, อยํ อโลภสฺส นิสฺสนฺโท. ตตฺถ โย อโทโส, อยํ เหตุ. ยํ ตโย ธมฺเม ปฏฺฐเปติ สมฺมาวาจํ สมฺมากมฺมนฺตํ สมฺมาอาชีวญฺจ, อยํ นิสฺสนฺโท. ตตฺถ โย อโมโห เหตุ, ยํ เทฺว ธมฺเม อุปฏฺฐเปติ อวิปรีต\-ทสฺสนมฺปิ จ อนภิลาปนํ, อยํ นิสฺสนฺโท. อิมสฺส พฺรหฺม\-จริยสฺส ยํ ผลํ, ตา เทฺว วิมุตฺติโย ราควิราคา เจโตวิมุตฺติ อวิชฺชา วิราคา จ ปญฺญาวิมุตฺติ, อิทํ ผลํ. อิติ อิมานิ ตีณิ กุสลมูลานิ นิทฺทิฏฺฐานิ \csromanfolio{215}เหตุโต จ นิสฺสนฺทโต จ ผลโต จ. เอวํ สาธารณานิ กุสลานิ ปฏิวิชฺฌิตพฺพานิ.}

\csromanheader{2}{ยตฺถ ทุเว ยตฺถ ตีณิ. อยญฺเจตฺถ คาถา.}
\setlength{\csromangathapagewidth}{0pt}
\setlength{\csromangathawakwidth}{0pt}
\csromangathameasuregroup{\textsuperscript{*}“ตุลมตุลญฺจ สมฺภวํ, \\ อชฺฌตฺตรโต สมาหิโต,}{\csromangathabat{\textsuperscript{*}“ตุลมตุลญฺจ สมฺภวํ,}{ภวสงฺขารมวสฺสชิ มุนิ.} \\ \csromangathabat{อชฺฌตฺตรโต สมาหิโต,}{อภินฺทิ กวจมิวตฺตสมฺภวนฺ”ติ.}}
\csromangathasetleft{\textsuperscript{*}“ตุลมตุลญฺจ สมฺภวํ, \\ อชฺฌตฺตรโต สมาหิโต,}
\csromangathagroup{\csromangathastanza{\csromangathabat{\csromansymbolfootnote{*}{เหฏฺฐา 52 ปิฏฺเฐ.}“ตุลมตุลญฺจ สมฺภวํ,}{ภวสงฺขารม\-วสฺสชิ มุนิ.} \\ \csromangathabat{อชฺฌตฺตรโต สมาหิโต,}{อภินฺทิ กวจมิวตฺตสมฺภวนฺ”ติ.}}}

\prose{ตุลมตุลญฺจ สมฺภวนฺติ ตุลสงฺขตํ อตุลสงฺขตํ. ตตฺถ เย สงฺขตาตุลํ, เต เทฺว ธมฺมา อสฺสาโท จ อาทีนโว จ ตุลิตา ภวนฺติ. เอตฺตโก กาเมสุ อสฺสาโท. เอตฺตโก อาทีนโว อิมสฺส, อิทํ นิสฺสรณนฺติ อิติ นิพฺพานํ ปชานาติ. ทฺวีหิ การเณหิ อตุลํ น จ สกฺกา ตุลยิตุํ. เอตฺตกํ เอตํ เนตํ ปรมตฺถีติ เตน อตุลํ. อถ ปาปุณา รตนํ กริตฺวา อจฺฉริย\-ภาเวน อตุลํ. ตตฺถ กุสลสฺส จ อภิสมฺภวา ชานนา ปสฺสนา, อยํ อโมโห. ยํ ตตฺถ ญาตา โอสิรณา ภวสงฺขารานํ. อยํ อโลโภ. ยํ อชฺฌตฺตรโต สมาหิโตติ วิกฺเขป\-ปฏิสํหรณา. อยํ อโทโส อิติ อิมานิ ตีณิ กุสลมูลานิ. ตุลมตุล\-สมฺภวนฺติ อยํ อโมโห. โย ภวสงฺขารานํ สโมสรณํ โลโภ สมฺมา\-สมาธีนํ อสฺสาโท, อยํ เหตุ, ยํ อชฺฌตฺตรโต อวิชฺชณฺฑโกสํ สมฺเภโท, อยํ นิสฺสนฺโท. สา ปวตฺติ อิมานิ ตีณิ นิทฺทิฏฺฐานิ กุสลมูลานิ เหตุโต จ นิสฺสนฺทโต จ ผลโต จ.}

\prose{เอตฺตาวตา เอสา ปวตฺติ จ นิวตฺติ จ กุสลมูเลหิ ปวตฺตติ, กุสลมูเลหิ นิวตฺตตีติ อิเมหิ จ ตีหิ สพฺพํ อกุสลมูลํ สโมสรณํ คจฺฉติ. โส ธมฺเม วา วจนโต นิทฺทิฏฺโฐ ตณฺหาติ วา โกโธติ วา อสมฺปชญฺญนฺติ วา อนุสโยติ วา มกฺโขติ วา ปฬาโสติ วา อสฺสตีติ วา อิสฺสาติ วา มจฺฉริยนฺติ วา อญฺญาณนฺติ วา, เตหิ เย จ วตฺถูหิ นิทฺทิสิตพฺพํ. ยสฺสิมานิ เทฺว วจนานิ ธมฺมปทานิ นิทฺทิฏฺฐานิ น โส อตฺถิ กิเลสา, โย อิเมสุ นวสุ ปเทสุ สโมธานํ สโมสรณํ คจฺฉติ, อยํ กิเลโส, น จ โลโภ, น จ โทโส, น จ โมโห. ยถา อกุสลมูลานิ, เอวํ กุสลานิ ปฏิกฺเขเปน นิทฺทิสิตพฺพานิ. อิทํ อโมหา\-ธิฏฺฐานํ.}

\proseitem{37}{\csromanfolio{216}ตตฺถ กตมํ กายกมฺมา\-ธิฏฺฐานํ.}

\setlength{\csromangathapagewidth}{0pt}
\setlength{\csromangathawakwidth}{0pt}
\csromangathameasuregroup{\textsuperscript{*}กาเยน กุสลํ กเร, \\ กายทุจฺจริตํ หิตฺวา,}{\csromangathabat{\textsuperscript{*}กาเยน กุสลํ กเร,}{อสฺส กาเยน สํวุโต.} \\ \csromangathabat{กายทุจฺจริตํ หิตฺวา,}{กาเยน สุจริตํ จเร.}}
\csromangathasetleft{\textsuperscript{*}กาเยน กุสลํ กเร, \\ กายทุจฺจริตํ หิตฺวา,}
\csromangathagroup{\csromangathastanza{\csromangathaitembat{37}{\csromansymbolfootnote{*}{เหฏฺฐา 207 ปิฏฺเฐ.}กาเยน กุสลํ กเร,}{อสฺส กาเยน สํวุโต.} \\ \csromangathacontbat{กายทุจฺจริตํ หิตฺวา,}{กาเยน สุจริตํ จเร.}}}

\prose{ตีณิมานิ ภิกฺขเว สุจริตานิ\csromansharedfootnote{baabe175f999360f}{ขุ 1. 233 สุจริตสุตฺเต.}. ปาณาติปาตา เวรมณี, อทินฺนาทานา เวรมณี, กาเมสุ\-มิจฺฉาจารา เวรมณี, อิทํ กายกมฺมา\-ธิฏฺฐานํ.}

\setlength{\csromangathapagewidth}{0pt}
\setlength{\csromangathawakwidth}{0pt}
\csromangathameasuregroup{}{ตตฺถ กตมํ วาจากมฺมาธิฏฺฐานํ. \\ สุภาสิตํ อุตฺตมมาหุ สนฺโต, \\ ธมฺมํ ภเณ นาธมฺมํ ตํ ทุติยํ. \\ ปิยํ ภเณ นาปฺปิยํ ตํ ตติยํ, \\ สจฺจํ ภเณ นาลิกํ ตํ จตุตฺถํ\textsuperscript{1}.}
\csromangathameasurewak{ตตฺถ กตมํ วาจากมฺมาธิฏฺฐานํ. \\ สุภาสิตํ อุตฺตมมาหุ สนฺโต, \\ ธมฺมํ ภเณ นาธมฺมํ ตํ ทุติยํ. \\ ปิยํ ภเณ นาปฺปิยํ ตํ ตติยํ, \\ สจฺจํ ภเณ นาลิกํ ตํ จตุตฺถํ\textsuperscript{1}.}
\setlength{\csromangathaleftcol}{0pt}
\csromangathagroup{\csromangathastanza{\csromangathawak{\csromangathaitemline{37}{ตตฺถ กตมํ วาจากมฺมา\-ธิฏฺฐานํ.} \\ \csromangathacontline{สุภาสิตํ อุตฺตมมาหุ สนฺโต,} \\ \csromangathacontline{ธมฺมํ ภเณ นาธมฺมํ ตํ ทุติยํ.} \\ \csromangathacontline{ปิยํ ภเณ นาปฺปิยํ ตํ ตติยํ,} \\ \csromangathacontline{สจฺจํ ภเณ นาลิกํ ตํ จตุตฺถํ\csromansharedfootnote{05b3aafc697701d2}{ขุ 1. 344 สุตฺตนิปาเต.}.}}}}

\prose{จตฺตาริมานิ จ วจีสุจริตานิ อิทํ วาจากมฺมา\-ธิฏฺฐานํ.}

\setlength{\csromangathapagewidth}{0pt}
\setlength{\csromangathawakwidth}{0pt}
\csromangathameasuregroup{}{ตตฺถ กตมํ มโนกมฺมาธิฏฺฐานํ. \\ มเนน กุสลํ กมฺมํ, มนสา สํวุโต ภเว. \\ มโนทุจฺจริตํ หิตฺวา, มนสา สุจริตํ จเร.}
\csromangathameasurewak{ตตฺถ กตมํ มโนกมฺมาธิฏฺฐานํ. \\ มเนน กุสลํ กมฺมํ, มนสา สํวุโต ภเว. \\ มโนทุจฺจริตํ หิตฺวา, มนสา สุจริตํ จเร.}
\setlength{\csromangathaleftcol}{0pt}
\csromangathagroup{\csromangathastanza{\csromangathawak{\csromangathaitemline{37}{ตตฺถ กตมํ มโนกมฺมา\-ธิฏฺฐานํ.} \\ \csromangathacontline{มเนน กุสลํ กมฺมํ, มนสา สํวุโต ภเว.} \\ \csromangathacontline{มโนทุจฺจริตํ หิตฺวา, มนสา สุจริตํ จเร.}}}}

\prose{ตีณิมานิ มโนสุจริตานิ, อนภิชฺฌา, อพฺยาปาโท, สมฺมาทิฏฺฐิ, อิทํ มโนกมฺมา\-ธิฏฺฐานํ. อิมานิ อสาธารณานิ สุตฺตานิ.}

\setlength{\csromangathapagewidth}{0pt}
\setlength{\csromangathawakwidth}{0pt}
\csromangathameasuregroup{}{ตตฺถ กตมานิ สาธารณานิ สุตฺตานิ. \\ วาจานุรกฺขี มนสา สุสํวุโต, \\ กาเยน จ นากุสลํ กยิรา\textsuperscript{1}. \\ เอเต ตโย กมฺมปเถ วิโสธเย, \\ อาราธเย มคฺคมิสิปฺปเวทิตํ.}
\csromangathameasurewak{ตตฺถ กตมานิ สาธารณานิ สุตฺตานิ. \\ วาจานุรกฺขี มนสา สุสํวุโต, \\ กาเยน จ นากุสลํ กยิรา\textsuperscript{1}. \\ เอเต ตโย กมฺมปเถ วิโสธเย, \\ อาราธเย มคฺคมิสิปฺปเวทิตํ.}
\setlength{\csromangathaleftcol}{0pt}
\csromangathagroup{\csromangathastanza{\csromangathawak{\csromangathaitemline{37}{ตตฺถ กตมานิ สาธารณานิ สุตฺตานิ.} \\ \csromangathacontline{วาจานุรกฺขี มนสา สุสํวุโต,} \\ \csromangathacontline{กาเยน จ นากุสลํ กยิรา\csromansharedfootnote{11f62785663408be}{อกุสลํ น กยิรา (อิ, ก) ขุ 1. 53 ธมฺมปเท; เหฏฺฐา 158 ปิฏฺเฐปิ.}.} \\ \csromangathacontline{เอเต ตโย กมฺมปเถ วิโสธเย,} \\ \csromangathacontline{อาราธเย มคฺคมิสิปฺปเวทิตํ.}}}}

\prose{ติสฺโส อิมา ภิกฺขเว ปาริสุทฺธิโย กายกมฺม\-ปาริสุทฺธิ, วาจากมฺม\-ปาริสุทฺธิ, มโนกมฺม\-ปาริสุทฺธิ.}

\prose{ตตฺถ กตมา กายกมฺม\-ปาริสุทฺธิ, ปาณาติปาตา เวรมณี, อทินฺนาทานา เวรมณี, กาเมสุ\-มิจฺฉาจารา เวรมณี. ตตฺถ กตมา วจีกมฺม\-ปาริสุทฺธิ, มุสาวาทา เวรมณี ฯเปฯ สมฺผปฺปลาปา เวรมณี. ตตฺถ กตมา มโนกมฺม\-ปาริสุทฺธิ, อนภิชฺฌา อพฺยาปาโท สมฺมาทิฏฺฐิ. อิทํ สาธารณ\-สุตฺตํ.}

\prose{\csromanfolio{217}อิติ สาธารณานิ จ สุตฺตานิ อสาธารณานิ จ สุตฺตานิ ปฏิวิชฺฌิตพฺพานิ. ปฏิวิชฺฌิตฺวา วาจาย กาเยน จ สุตฺตสฺส อตฺโถ นิทฺทิสิตพฺโพ.}

\proseitem{38}{ตตฺถ กตมํ สทฺธินฺทฺริยา\-ธิฏฺฐานํ.}

\setlength{\csromangathapagewidth}{0pt}
\setlength{\csromangathawakwidth}{0pt}
\csromangathameasuregroup{ยสฺส สทฺธา ตถาคเต, \\ สีลญฺจ ยสฺส กลฺยาณํ, \\ สํเฆ ปสาโท ยสฺสตฺถิ, \\ อทลิทฺโทติ ตํ อาหุ, \\ สทฺธา เว นนฺทิกา อาราธิโก, \\ สพฺพํ สิยาติ ภควนฺตํ,}{\csromangathabat{ยสฺส สทฺธา ตถาคเต,}{อจลา สุปฺปติฏฺฐิตา.} \\ \csromangathabat{สีลญฺจ ยสฺส กลฺยาณํ,}{อริยกนฺตํ ปสํสิตํ.} \\ \csromangathabat{สํเฆ ปสาโท ยสฺสตฺถิ,}{อุชุภูตญฺจ ทสฺสนํ.} \\ \csromangathabat{อทลิทฺโทติ ตํ อาหุ,}{อโมฆํ ตสฺส ชีวิตํ\textsuperscript{1}.} \\ \csromangathabat{สทฺธา เว นนฺทิกา อาราธิโก,}{โน ตสฺส สทฺโธติ.} \\ \csromangathabat{สพฺพํ สิยาติ ภควนฺตํ,}{ตถารูโป ธมฺมสมฺปสาโท.} \\ อิทํ สทฺธินฺทฺริยาธิฏฺฐานํ. \\ ตตฺถ กตมํ วีริยาธิฏฺฐานํ. \\ อารมฺภถ\textsuperscript{1} นิกฺกมถ, ยุญฺชถ พุทฺธสาสเน. \\ ธุนาถ มจฺจุโน เสนํ, นฬาคารํว กุญฺชโร.}
\csromangathameasurewak{อิทํ สทฺธินฺทฺริยาธิฏฺฐานํ. \\ ตตฺถ กตมํ วีริยาธิฏฺฐานํ. \\ อารมฺภถ\textsuperscript{1} นิกฺกมถ, ยุญฺชถ พุทฺธสาสเน. \\ ธุนาถ มจฺจุโน เสนํ, นฬาคารํว กุญฺชโร.}
\csromangathasetleft{ยสฺส สทฺธา ตถาคเต, \\ สีลญฺจ ยสฺส กลฺยาณํ, \\ สํเฆ ปสาโท ยสฺสตฺถิ, \\ อทลิทฺโทติ ตํ อาหุ, \\ สทฺธา เว นนฺทิกา อาราธิโก, \\ สพฺพํ สิยาติ ภควนฺตํ,}
\csromangathagroup{\csromangathastanza{\csromangathabat{ยสฺส สทฺธา ตถาคเต,}{อจลา สุปฺปติฏฺฐิตา.} \\ \csromangathabat{สีลญฺจ ยสฺส กลฺยาณํ,}{อริยกนฺตํ ปสํสิตํ.}}\csromangathastanza{\csromangathabat{สํเฆ ปสาโท ยสฺสตฺถิ,}{อุชุภูตญฺจ ทสฺสนํ.} \\ \csromangathabat{อทลิทฺโทติ ตํ อาหุ,}{อโมฆํ ตสฺส ชีวิตํ\csromansharedfootnote{e027d2bc7f1d8fb2}{สํ 1. 234 ปิฏฺเฐ; ขุ 2. 296 อฏฺฐกนิปาเต จ ปสฺสิตพฺพํ.}.}}\csromangathastanza{\csromangathabat{สทฺธา เว นนฺทิกา อาราธิโก,}{โน ตสฺส สทฺโธติ.} \\ \csromangathabat{สพฺพํ สิยาติ ภควนฺตํ,}{ตถารูโป ธมฺมสมฺปสาโท.}}\csromangathastanza{\csromangathawak{อิทํ สทฺธินฺทฺริยา\-ธิฏฺฐานํ.}}\csromangathastanza{\csromangathawak{\csromangathaitemline{38}{ตตฺถ กตมํ วีริยา\-ธิฏฺฐานํ.} \\ \csromangathacontline{อารมฺภถ\csromansharedfootnote{fb8f7af740e13c6e}{อารภถ (อิ) สํ 1. 159 ปิฏฺเฐ; เหฏฺฐา 35 ปิฏฺเฐปิ.} นิกฺกมถ, ยุญฺชถ พุทฺธสาสเน.} \\ \csromangathacontline{ธุนาถ มจฺจุโน เสนํ, นฬาคารํว กุญฺชโร.}}}}

\prose{จตฺตาโรเม ภิกฺขเว สมฺมปฺปธานา, อิทํ วีริยา\-ธิฏฺฐานํ.}

\setlength{\csromangathapagewidth}{0pt}
\setlength{\csromangathawakwidth}{0pt}
\csromangathameasuregroup{}{ตตฺถ กตมํ สตินฺทฺริยาธิฏฺฐานํ. \\ สตีมโต สทา ภิทฺธํ, ภทฺทมตฺถุ สตีมโต. \\ สตีมโต สทา\textsuperscript{1} เสยฺโย, สตีมา สุขเมธติ.}
\csromangathameasurewak{ตตฺถ กตมํ สตินฺทฺริยาธิฏฺฐานํ. \\ สตีมโต สทา ภิทฺธํ, ภทฺทมตฺถุ สตีมโต. \\ สตีมโต สทา\textsuperscript{1} เสยฺโย, สตีมา สุขเมธติ.}
\setlength{\csromangathaleftcol}{0pt}
\csromangathagroup{\csromangathastanza{\csromangathawak{\csromangathaitemline{38}{ตตฺถ กตมํ สตินฺทฺริยา\-ธิฏฺฐานํ.} \\ \csromangathacontline{สตีมโต สทา ภิทฺธํ, ภทฺทมตฺถุ สตีมโต.} \\ \csromangathacontline{สตีมโต สทา\csromansharedfootnote{71b8b9991d22b968}{สุเว (สํ 1. 210 ปิฏฺเฐ)} เสยฺโย, สตีมา สุขเมธติ.}}}}

\prose{จตฺตาโร สติปฏฺฐานา วิตฺถาเรน กาตพฺพา, อิทํ สตินฺทฺริยา\-ธิฏฺฐานํ.}

\prose{ตตฺถ กตมํ สมาธินฺทฺริยา\-ธิฏฺฐานํ.}

\setlength{\csromangathapagewidth}{0pt}
\setlength{\csromangathawakwidth}{0pt}
\csromangathameasuregroup{}{\textsuperscript{*}อากงฺขโต เต นรทมฺมสารถิ, \\ เทวา มนุสฺสา มนสา วิจินฺติตํ. \\ สพฺเพน ชญฺญา กสิณาปิ ปาณิโน, \\ สนฺตํ สมาธิํ อรณํ นิเสวโต.}
\csromangathameasurewak{\textsuperscript{*}อากงฺขโต เต นรทมฺมสารถิ, \\ เทวา มนุสฺสา มนสา วิจินฺติตํ. \\ สพฺเพน ชญฺญา กสิณาปิ ปาณิโน, \\ สนฺตํ สมาธิํ อรณํ นิเสวโต.}
\setlength{\csromangathaleftcol}{0pt}
\csromangathagroup{\csromangathastanza{\csromangathawak{\csromansymbolfootnote{*}{เหฏฺฐา 52 ปิฏฺเฐ.}อากงฺขโต เต นรทมฺมสารถิ, \\ เทวา มนุสฺสา มนสา วิจินฺติตํ. \\ สพฺเพน ชญฺญา กสิณาปิ ปาณิโน, \\ สนฺตํ สมาธิํ อรณํ นิเสวโต.}}}

\prose{ตโยเม ภิกฺขเว สมาธี\csromandash{}สวิตกฺโก สวิจาโร, อวิตกฺโก วิจารมตฺโต, อวิตกฺโก อวิจาโร, อิทํ สมาธินฺทฺริยา\-ธิฏฺฐานํ.}

\prose{\csromanfolio{218}ตตฺถ กตมํ ปญฺญินฺทฺริยา\-ธิฏฺฐานํ.}

\prose{ปญฺญา หิ เสฏฺฐา โลกสฺมินฺติ วิตฺถาเรน.}

\prose{ติสฺโส อิมา ภิกฺขเว ปญฺญา, สุตมยี, จินฺตามยี, ภาวนามยี, อิทํ ปญฺญินฺทฺริยา\-ธิฏฺฐานํ สุตฺตํ, อิมานิ อินฺทฺริยา\-ธิฏฺฐานานิ อสาธารณานิ สุตฺตานิ.}

\setlength{\csromangathapagewidth}{0pt}
\setlength{\csromangathawakwidth}{0pt}
\csromangathameasuregroup{}{ตตฺถ กตมานิ สาธารณานิ อินฺทฺริยาธิฏฺฐานานิ สุตฺตานิ. \\ อวีตราโค กาเมสุ, ยสฺส ปญฺจินฺทฺริยา มุทู. \\ สทฺธา สติ จ วีริยํ, สมโถ จ วิปสฺสนา. \\ ตาทิสํ ภิกฺขุมาสชฺช, ปุพฺเพว อุปหญฺญติ\textsuperscript{1}.}
\csromangathameasurewak{ตตฺถ กตมานิ สาธารณานิ อินฺทฺริยาธิฏฺฐานานิ สุตฺตานิ. \\ อวีตราโค กาเมสุ, ยสฺส ปญฺจินฺทฺริยา มุทู. \\ สทฺธา สติ จ วีริยํ, สมโถ จ วิปสฺสนา. \\ ตาทิสํ ภิกฺขุมาสชฺช, ปุพฺเพว อุปหญฺญติ\textsuperscript{1}.}
\setlength{\csromangathaleftcol}{0pt}
\csromangathagroup{\csromangathastanza{\csromangathawak{\csromangathaitemline{39}{ตตฺถ กตมานิ สาธารณานิ อินฺทฺริยา\-ธิฏฺฐานานิ สุตฺตานิ.} \\ \csromangathacontline{อวีตราโค กาเมสุ, ยสฺส ปญฺจินฺทฺริยา มุทู.} \\ \csromangathacontline{สทฺธา สติ จ วีริยํ, สมโถ จ วิปสฺสนา.} \\ \csromangathacontline{ตาทิสํ ภิกฺขุมาสชฺช, ปุพฺเพว อุปหญฺญติ\csromansharedfootnote{2b74806e03ebcc8b}{อํ 2. 327 ปิฏฺเฐ.}.}}}}

\prose{ปญฺจิมานิ อินฺทฺริยานิ, สทฺธินฺทฺริยา\-ทิ\-อินฺทฺริยํ ทฏฺฐพฺพํ, ตีสุ อเวจฺจปฺปสาเท วิตฺถาเรน สุตฺตํ กาตพฺพํ, อิมานิ สาธารณานิ อินฺทฺริยํธิฏฺฏานานิ สุตฺตานิ. ยํ ยสฺส สมฺพนฺธํ กุสลสฺส วา อกุสลสฺส วา เตน เตน อธิฏฺฐาเนน ตํ สุตฺตํ นิทฺทิสิตพฺพํ, นตฺถญฺโญ ธมฺโม นิทฺทิสิตพฺโพ, ตตฺถ สาธารณํ กุสลํ นาปิ กุสลํ อกุสลํ ยถา สาธารณานิ จ กุสลมูลานิ สาธารณานิ จ อกุสลมูลานิ อุปฺปนฺนํ กามวิตกฺกํ ปชหติ ฯเปฯ จตฺถาโร สมฺมปฺปธานา กุสลํ อกุสลญฺจ.}

\csromanheader{2}{\textbf{ตตฺถิมา อุทฺทานคาถา}}
\setlength{\csromangathapagewidth}{0pt}
\setlength{\csromangathawakwidth}{0pt}
\csromangathameasuregroup{วิตกฺโก หิ มมตฺถิโก\textsuperscript{1}, \\ โย ปาณมติปาเตติ, \\ สตญฺเจว สหสฺสานํ, \\ ฉนฺทา โทสา ภยา โมหา, \\ อสุภานุปสฺสิํ วิหรนฺตํ, \\ เอกมฺปิ เจ ปิยํ ปาณํ, \\ ปญฺญา หิ เสฏฺฐา โลกสฺมิํ, \\ กุสลากุสลมูลานิ จ, \\ กาเยน กุสลํ กเร, \\ สุภาสิตํ อุตฺตมมาหุ,}{\csromangathabat{วิตกฺโก หิ มมตฺถิโก\textsuperscript{1},}{ททํ ปิโย นโร อิติ.} \\ \csromangathabat{โย ปาณมติปาเตติ,}{ตีณิ ตสฺส พาลลกฺขณํ.} \\ \csromangathabat{สตญฺเจว สหสฺสานํ,}{เย จ สมณพฺราหฺมณา.} \\ \csromangathabat{ฉนฺทา โทสา ภยา โมหา,}{จตูหิ อคตีหิ จ.} \\ \csromangathabat{อสุภานุปสฺสิํ วิหรนฺตํ,}{นิมิตฺเตสุ อสุภา จ.} \\ \csromangathabat{เอกมฺปิ เจ ปิยํ ปาณํ,}{มิตฺตา สเจ สุภาสิตา.} \\ \csromangathabat{ปญฺญา หิ เสฏฺฐา โลกสฺมิํ,}{อนุญฺญา ตีณิ อินฺทฺริยานิ.} \\ \csromangathabat{กุสลากุสลมูลานิ จ,}{ตุลมตุลญฺจ สมฺภวํ.} \\ \csromangathabat{กาเยน กุสลํ กเร,}{ตีณิ สุจริตานิ จ.} \\ \csromangathabat{สุภาสิตํ อุตฺตมมาหุ,}{สนฺโต วจีสุจริตานิ จ.}}
\csromangathasetleft{วิตกฺโก หิ มมตฺถิโก\textsuperscript{1}, \\ โย ปาณมติปาเตติ, \\ สตญฺเจว สหสฺสานํ, \\ ฉนฺทา โทสา ภยา โมหา, \\ อสุภานุปสฺสิํ วิหรนฺตํ, \\ เอกมฺปิ เจ ปิยํ ปาณํ, \\ ปญฺญา หิ เสฏฺฐา โลกสฺมิํ, \\ กุสลากุสลมูลานิ จ, \\ กาเยน กุสลํ กเร, \\ สุภาสิตํ อุตฺตมมาหุ,}
\csromangathagroup{\csromangathastanza{\csromangathabat{วิตกฺโก หิ มมตฺถิโก\csromansharedfootnote{1f12aea5ceea3653}{ปมตฺถิโก (อิ)},}{ททํ ปิโย นโร อิติ.} \\ \csromangathabat{โย ปาณมติปาเตติ,}{ตีณิ ตสฺส พาลลกฺขณํ.}}\csromangathastanza{\csromangathabat{สตญฺเจว สหสฺสานํ,}{เย จ สมณพฺราหฺมณา.} \\ \csromangathabat{ฉนฺทา โทสา ภยา โมหา,}{จตูหิ อคตีหิ จ.}}\csromangathastanza{\csromangathabat{อสุภานุปสฺสิํ วิหรนฺตํ,}{นิมิตฺเตสุ อสุภา จ.} \\ \csromangathabat{เอกมฺปิ เจ ปิยํ ปาณํ,}{มิตฺตา สเจ สุภาสิตา.}}\csromangathastanza{\csromangathabat{ปญฺญา หิ เสฏฺฐา โลกสฺมิํ,}{อนุญฺญา ตีณิ อินฺทฺริยานิ.} \\ \csromangathabat{กุสลากุสลมูลานิ จ,}{ตุลมตุลญฺจ สมฺภวํ.}}\csromangathastanza{\csromangathabat{กาเยน กุสลํ กเร,}{ตีณิ สุจริตานิ จ.} \\ \csromangathabat{สุภาสิตํ อุตฺตมมาหุ,}{สนฺโต วจีสุจริตานิ จ.}}}

\csromanmatikaanchor{mtk.43}
\csromantocmarkref{section}{4. สุตฺตวิจยจตุตฺถภูมิ}{mtk.43}
\bookmark[dest={mtk.43},level=1]{4. สุตฺตวิจยจตุตฺถภูมิ}
\csromanheader{1}{4. สุตฺต\-วิจย\-จตุตฺถภูมิ}
\setlength{\csromangathapagewidth}{0pt}
\setlength{\csromangathawakwidth}{0pt}
\csromangathameasuregroup{กาเยน จ กุสลํ กยิรา, \\ กายานุรกฺขี จ สทา, \\ ยสฺส สทฺธา ตถาคเต, \\ อารมฺภถ นิกฺกมถ, \\ สตีมโต สทา ภทฺทํ, \\ อากงฺขโต จ อนญฺญาณํ, \\ ปญฺญา หิ เสฏฺฐา โลกสฺมิํ, \\ อวีตราโค กาเมสุ,}{\csromangathabat{กาเยน จ กุสลํ กยิรา,}{มโนทุจฺจริตานิ จ.} \\ \csromangathabat{กายานุรกฺขี จ สทา,}{ติสฺโส จ ปาริสุทฺธิโย.} \\ \csromangathabat{ยสฺส สทฺธา ตถาคเต,}{สมุปฺปาเท จ เทสิโต.} \\ \csromangathabat{อารมฺภถ นิกฺกมถ,}{ยา จ สมฺมปฺปธานตา.} \\ \csromangathabat{สตีมโต สทา ภทฺทํ,}{สติปฏฺฐานภาวนา.} \\ \csromangathabat{อากงฺขโต จ อนญฺญาณํ,}{เย จ ตีณิ สมาธโย.} \\ \csromangathabat{ปญฺญา หิ เสฏฺฐา โลกสฺมิํ,}{ติสฺโส ปญฺญา ปกาสิตา.} \\ \csromangathabat{อวีตราโค กาเมสุ,}{ตเถว ปญฺจินฺทฺริยา.}}
\csromangathasetleft{กาเยน จ กุสลํ กยิรา, \\ กายานุรกฺขี จ สทา, \\ ยสฺส สทฺธา ตถาคเต, \\ อารมฺภถ นิกฺกมถ, \\ สตีมโต สทา ภทฺทํ, \\ อากงฺขโต จ อนญฺญาณํ, \\ ปญฺญา หิ เสฏฺฐา โลกสฺมิํ, \\ อวีตราโค กาเมสุ,}
\csromangathagroup{\csromangathastanza{\csromanfolio{219}\csromangathabat{กาเยน จ กุสลํ กยิรา,}{มโน\-ทุจฺจริตานิ จ.} \\ \csromangathabat{กายานุรกฺขี จ สทา,}{ติสฺโส จ ปาริสุทฺธิโย.}}\csromangathastanza{\csromangathabat{ยสฺส สทฺธา ตถาคเต,}{สมุปฺปาเท จ เทสิโต.} \\ \csromangathabat{อารมฺภถ นิกฺกมถ,}{ยา จ สมฺมปฺปธานตา.}}\csromangathastanza{\csromangathabat{สตีมโต สทา ภทฺทํ,}{สติปฏฺฐาน\-ภาวนา.} \\ \csromangathabat{อากงฺขโต จ อนญฺญาณํ,}{เย จ ตีณิ สมาธโย.}}\csromangathastanza{\csromangathabat{ปญฺญา หิ เสฏฺฐา โลกสฺมิํ,}{ติสฺโส ปญฺญา ปกาสิตา.} \\ \csromangathabat{อวีตราโค กาเมสุ,}{ตเถว ปญฺจินฺทฺริยา.}}}

\prose{อิติ เถรสฺส มหา\-กจฺจายนสฺส ชมฺพุ\-วนวาสิโน เปฏโกปเทเส ตติยภูมิ สุตฺตา\-ธิฏฺฐานํ นาม.}

\csromanheader{2}{4. สุตฺต\-วิจย\-จตุตฺถภูมิ}
\proseitem{40}{ตตฺถ กตโม สุตฺตวิจโย, ตตฺถ กุสเลหิ ธมฺเมหิ อกุสเลหิ ธมฺเมหิ ปุพฺพาปรโส สาธุกํ อุปปริกฺขิยติ. กิํ นุ โข อิทํ สุตฺตํ อารภิ ฯเปฯ เตหิ สุตฺเตหิ สห อธิสนฺน\-ฏฺเฐหิ ยุชฺชติ. อุทาหุ น ยุชฺชตีติ.}

\prose{ยถา ภควา กิเลเส อาทิมฺหิ ตตฺถ เทเสติ. กิํ เทสิตํ เตสํ กิเลสานํ ปหานํ อุทาหุ โน เทสิตนฺติ อุปปริ\-กฺขิ\-ตพฺพํ. ยทิ น เทสิตํ ภควติ เตสํ กิเลสานํ ปหานํ กุสลา ธมฺมา ปริเยสิตพฺพา. ยตฺถ เต อกุสลา ปหานํ คจฺฉนฺติ. สเจ สมนฺเนหมาโน น ลภติ. ตตฺถ อกุสลา ธมฺมา อปกฑฺฒิตพฺพา วีมํสิตพฺพา, สํกิเลสภาคิย\-สุตฺตํ, ยทิ กิเลสา อปกฑฺฒิยนฺตา. เย วา น เทนฺติ. ตตฺถ อุปปริ\-กฺขิ\-ตพฺพา อริยมคฺค\-ธมฺมา ตาสุ ภูมีสุ กิเลสา ปหานํ คจฺฉนฺติ, อุทาหุ น คจฺฉนฺตีติ. ยตฺตกา ปน กิเลสา เทสิตา. น ตตฺตกา อริยธมฺมา เทสิตา. ยตฺถ กิเลสา ปหานํ คจฺฉนฺติ, ตตฺถ เย กิเลสา อริยธมฺมานํ ปฏิปกฺเขน น ยุชฺชนฺติ, เต อปกฑฺฒิตพฺพา, สเจ อปกฑฺฒิยนฺตา โยชนํ เทติ. ตตฺถ เอวํ วีมํสิตพฺพํ. เทฺว ตีณิ วา ตทุตฺตริ วา กิเลสา เอเกน อริยมคฺเคน ปหานํ คจฺฉนฺตีติ. สเจ เอวํ วีมํสิยนฺตา โยชนํ เทติ, ตตฺถ อุปปริ\-กฺขิ\-ตพฺพํ. ปรมฺปราย วา ปิฏก\-สมฺปทาเนน วา สุตฺตสฺส อตฺโถ จ นตฺโถ จ. ยํ วา น สกฺกา สุตฺตํ นิทฺทิสิตุํ เนว สุตฺตํ วิจิกิจฺฉิตพฺพํ. เอวํ ยถา อาทิมฺหิ กุสลา ธมฺมา โหนฺติ. เย \csromanfolio{220}กิเลสา เต ปหีเนยฺยาติ. เต อุปปริ\-กฺขิ\-ตพฺพา. ปุโร วา กุสโล ปฏิปกฺเขน วา ปุโร เทสนา, อนูนา อนธิกา อุคฺคเหตพฺพา. ยถา ปฐโม อุตฺติโล เยสมิทานิ กิเลสานํ เย อริยธมฺมา เทสิตา อิเม กิเลสา อิเมหิ อริยธมฺเมหิ ปหียนฺติ, อุทาหุ นปฺปหียนฺตีติ วิจินิตพฺพา. ยทิ อุปปริกฺขิยมานา ยุชฺชนฺติ, คเหตพฺพา. อถ น ยุชฺชนฺติ, เย กิเลสา อปฏิปกฺขา โหนฺติ, เต กิเลสา อปริปกฺขิตพฺพา. เย จ อริยธมฺมา ปฏิปกฺขา โหนฺติ, เต อริยธมฺมา อปกฑฺฒิตพฺพา. น หิ อริยธมฺมา อนาคามิ\-กิเลสปฺปหานํ คจฺฉนฺติ, นาปิ อริยธมฺมา สพฺพกิเลสานํ ปหานาย สํวตฺตนฺติ. ยถา กุสลา เมตฺตา อกุสโล ราโค น ตุ กุสลา เมตฺตาติ กาเรตฺวา อกุสลสฺส ราคสฺส ปหานาย สมฺภวติ พฺยาปาโท เมตฺตาย ปหานํ คจฺฉติ. ตสฺมา อุโภ กิเลสา อุปปริ\-กฺขิ\-ตพฺพา. โย โย จ ธมฺโม อุปทิสิยติ. กุสโล วา อกุสโล วา โส อปกฑฺฒิตพฺโพ. สเจ เต ยุชฺชนฺติ อปกฑฺฒิยมาโน นตฺถิ อุปปริ\-กฺขิ\-ตพฺพํ. เทฺว วา กิเลสา เอเกน อริยธมฺเมน ปหีเนยฺยาติ ทฺวีหิ วา อริยธมฺเมหิ เอโก วา กิเลโส ปหียตีติ.}

\prose{อถ วา เอวมฺปิ อุปปริกฺขิ\-ยมานํ ยุชฺชติ, ตตฺถ วีมํสิตพฺพํ วา ยถา ยุชฺชติ ตตฺถ วีมํสิตพฺพํ วา, ยถา นนุ สกฺกา สุตฺตํ นิทฺทิสิตุํ, น หิ สุตฺเต วิจิกิจฺฉิตพฺพํ. กิเลโส มํ อริยธมฺเมสุ เทสิเตสุ อุภยโต อุปปริ\-กฺขิ\-ตพฺพํ. กิร เย วา อิเม กิเลสา เทสิตา เย จ อริยธมฺมา เทสิตา คาถาย วา พฺยากรเณน วา, กิํ นุ โข อิเม กิเลสา อิเมหิ อริยธมฺเมหิ ปหียนฺติ, อุทาหุ นปฺปหียนฺติ. อิเม วา อริยธมฺมา อิเมสํ กิเลสานํ ปหานาย สํวตฺตนฺตีติ. กิญฺจาปิ กุสเลหิ ธมฺเมหิ อกุสลา ธมฺมา ปหานํ คจฺฉนฺติ. น ตุ สพฺเพหิ อริยธมฺเมหิ สพฺพากุสลา ปหานํ คจฺฉนฺติ. ยถา เมตฺตา กุสโล อกุสโล จ ราโค น ตุ กุสลา เมตฺตา อกุสโล ราโคติ กาเรตฺวา เมตฺตาย ราโค ปหานํ, พฺยาปาโท เมตฺตาย ปหานํ คจฺฉนฺติ. เอวํ กิเลโสติ กาเรตฺวา สุตฺเตน ปหานํ คจฺฉติ. น สุตฺโต ธมฺโมติ กาเรตฺวา สพฺพํ กิเลสสฺส ปหานาย สํวตฺตติ. ยํ ตุ สุตฺตสฺส อริยธมฺโม สํกิเลส\-ปฏิปกฺโข, โส เตน ปหานํ คจฺฉตีติ.}

\proseitem{41}{ตตฺถ กุสเล เทสิเต สุตฺเต พฺยากรเณ วา สํกิเลสา น ยุชฺชนฺติ อริยธมฺมา วา, เต มหาปเทเส นิทฺทิสิตพฺพา\-วยเวน}

\vspace{\tipitakaparskip}
\csromancenter{สุตฺต\-วิจย\-จตุตฺถภูมิ อปกฑฺฒิตพฺพา.  ตตฺถ กิเลเสหิ จ เทสิเตหิ อริยธมฺเมสุ จ ยทิปิ เตน อริยธมฺเมน เต กิเลสา ปหานํ คจฺฉนฺติ. ตตฺถปิ อุตฺตริ อุปปริ\-กฺขิ\-ตพฺพํ. เกน การเณน เอเต กิเลสา ปชหิตพฺพา, เกน การเณน อริยธมฺมา เทสิตาติ เยน เยน วา อากาเรน อริยธมฺมา เทสิตา, เตน เตน ปกาเรน อยํ กิเลโส ฐิโต. อตฺถิ หิ เอโก กิเลโส, เตน วา อริยธมฺมา น อญฺญถา อญฺญถา ปหาตพฺโพ, ยถา ทิฏฺฐิ ราโค อวิชฺชา จ ทสฺสเนน ปหาตพฺพา. สา เจ เอวญฺจ อวิชฺชา ภาวนาย ภูมิ วา ธมฺมา ภาวนาย ปหาตพฺพา. สาเยว อุทฺธํภาคิยํ อสงฺขต\-ทสฺสนาย วิมุตฺติยา อนิมิตฺเตน เจโตสมาธินา อมนสิกาเรน ปหียติ. เอวํ สาตฺถํ สพฺยญฺชนํ อุปปริ\-กฺขิ\-ตพฺพํ. เย ทสฺสเนน ปหาตพฺพา กิเลสา ทสฺสนากาเรน อริยธมฺโม เทสิโต, ภาวนาย ปหาตพฺพา ภาวนากาเรน อริยธมฺโม เทสิโต, ปติเสวนา ปหาตพฺพา ปติเสวนา\-กาเรน อริยธมฺโม เทสิโต, เอวํ วิโนทน\-ปหาตพฺพา ยาว สตฺต อาสวา กาตพฺพา, ยาวญฺญถา. อญฺญถา เหส ธมฺโม ปหาตพฺโพ อญฺเญนากาเรน อริยธมฺโม เทสิโต, โส อริยธมฺโม อญฺญถา ปริเยสิตพฺโพ. ยทิ อยํ ธมฺโม ปริเยสโต โย จ เทเสติ เยน เยนากาเรน, โส อริยธมฺโม ปริเยสิตพฺโพ, เตนากาเรน กิเลโส ปหียติ. โส ตตฺถ อุปปริ\-กฺขิ\-ตพฺโพ. อถ น ยุชฺชติ ยทิ หิ เตน สุตฺเตน วิหิตํ สุตฺตํ วีมํสิตพฺพํ. ยถา ยุชฺชติ, ตถา คเหตพฺพํ. ยถา น ยุชฺชติ, ตถา น คเหตพฺพํ, อทฺธา เอตํ ภควตา น ภาสิตํ, อายสฺมตา วา ทุคฺคหิตํ, ยถา มหาปเทเส นิทฺทิสิตพฺพํ, ภควตา ยถาภูตํ เทสิตํ, โย จ ธมฺโม เทสิโต กุสโล จ อกุสโล จ ตสฺส ธมฺมสฺส ปจฺจโย ปริเยสิตพฺโพ. น หิ ปจฺจยา วินา ธมฺโม อปฺปจฺจโย อุปฺปชฺชติ. ตตฺถ โก อากาโร ปริเยสนาย.}
\prose{\csromanfolio{221}ตตฺถ ตถารูปํ สเหตุ สปฺปจฺจยํ โสยํ ธมฺโม วุตฺโตติ อิทํ วีมํสิตพฺพํ. โส จ ปจฺจโย ติวิโธ มุทุ มชฺโฌ อธิมตฺโต. ตตฺถ มุทุมฺหิ ปจฺจเย มุทุธมฺโม คเหตพฺโพ, เอวํ สเตฺยส ปจฺจโย ทุวิโธ ปรํปรา\-ปจฺจโย จ สมนนฺตร\-ปจฺจโย จ. โส ปจฺจโย มุทุเตน พฺยาธิมตฺตํ ปริเยสิตพฺพํ. กิํ การณํ, อญฺญตโรปิ ปจฺจโย อญฺเญหิ ปจฺจเยหิ ปริยตฺติํ วา ปาริปูริํ วา คจฺฉติ. ตตฺถ โย ธมฺโม เทสิโต, ตสฺส \csromanfolio{222}ธมฺมสฺส เอเตน วา การเณน วา เหตุ ปริเยสิตพฺโพ. ยถา ปจฺจโย เหตุนา ปจฺจเยน จ, โส ตสฺส ธมฺมสฺส นิสฺสนฺโท ปริเยสิตพฺโพ. ยถา นิทฺทิฏฺโฐ อธิฏฺฐาเน ปธานํ ปริเยสติ, โส ปจฺจโย ปริเยสิตพฺโพ. น หิ มุทุสฺส ธมฺมสฺส อธิมตฺโต นิสฺสนฺโท อธิมตฺตสฺส วา นิสฺสนฺทสฺส มุทุธมฺโม, อถ มุทุสฺส มุทุ มชฺฌาย มชฺโฌ อธิมตฺตสฺส อธิมตฺโต ยุชฺชติ, ตํ คเหตพฺพํ, อถ น ยุชฺชติ น คเหตพฺพํ. ยญฺจ ภควา อารภติ ธมฺมํ เทเสตุํ, ตํเยว ธมฺมํ มชฺฌ\-นฺต\-ปริโยสานํ เทเสติ, ยถา สุตฺตาธิฏฺฐาเน ธมฺมา อาทิมฺหิ นิทฺทิสติ, ตํเยว พหุ ตสฺส สุตฺตสฺส ปริโยสานํ. ตสฺส หิ ธมฺมสฺส วเสน ตํ สุตฺตํ โหติ คาถา วา พฺยากรณํ ขุทฺทกํ มหนฺตํ วา, ยถา ปน ทุวิธา อนุรูปนฺติ วา ถปนา จ เทสนาถปนา. รูปนฺติปิ ธมฺมสฺส ปริเยสิตพฺพา. ยถา จ ภควตา ปญฺจนฺนํ อินฺทฺริยานํ สํวรณํ เทสิตํ ตณฺหาย นิคฺคหณตฺถํ อิจฺฉาว โหติ. เทเสติ ยถา โคปาลโกปเม สุตฺเต อญฺเญหิปิ สุตฺเตหิ ภควา ภาสติ อิจฺฉาว โหติ มชฺฌิมนิกาเย วิตกฺโก อยํ ภควโต เทสนานุรูปนฺติ อิติ โส ธมฺโม อญฺเญสุปิ เวยฺยากรเณสุ ปริเยสิตพฺโพ. น หิ เอกํ หิ สุตฺเต ทฏฺฐพฺโพ. ยุชฺชนํ ตํ คเหตพฺพํ.}

\csromanmatikaanchor{mtk.44}
\csromantocmarkref{section}{5. ปญฺจมภูมิ}{mtk.44}
\bookmark[dest={mtk.44},level=1]{5. ปญฺจมภูมิ}
\proseitem{42}{ตตฺถ กตมํ อนุญฺญาตํ. ยํ กิญฺจิ สุตฺตํ ภควตา น ภาสิตํ ตญฺจ สุตฺเตสุเยว ทิสฺสติ, เอวเมตํ ธาเรตพฺพํ. ยถา อสุเกน ภาสิตนฺติ, ตํ สุตฺตํ วีมํสิตพฺพํ. กิํ นุ โข อิมํ สุตฺตํ อนุญฺญาตํ ขมํ ภควโต อุทาหุ นานุญฺญาตํ ขมํ, กิญฺจิ รูปญฺจ สุตฺตํ ภควโต อนุญฺญาตํ ขมํ กิญฺจิ รูปญฺจ นานุญฺญาตํ ขมํ. ยํ สพฺพโส อโนตาเรตฺวา ทสพโล โคจรํ เทเสติ, ตํ สพฺพํ สุตฺตํ ภควโต นานุญฺญาตํ ขมํ. อตฺถิปิ โส สาวโก ทสพลานํ โคจรํ ชานาติ โอธิโส อโนธิโส, ตํ ปน พลํ สพฺพโส น ชานาติ อญฺญถา นาม สวเนน, ยถา อายสฺมตา สาริปุตฺเตน เยน พฺราหฺมโณ โอวทิโต, ตสฺส อายสฺมโต นตฺถิ อินฺทฺริย\-พล\-เวมตฺต\-ญาณํ, เตน ปุคฺคล\csromansharedfootnote{1cc7c9381d943c78}{ปุคฺคโล ปโรปรญฺจ (อิ)} ปโร ปรญฺจ ตํ อชานนฺโต สติ อุตฺตริกรณีเย อุปฺปาทิโต, โส ภควตา อปสาทิโต. ยถาว อายสฺมา มหากสฺสโป ภาคิเนยฺยํ โอวทติ อนนฺต\-ริย\-สมนฺนาคโต อิทฺธิปฏิหีเรน องฺคุลิโย \csromanfolio{223}อทีเปตฺวา ยํ สพฺเพสํ ธมฺมานํ กมฺม\-สมาทานานํ เหตุโส ฐานโส ยถาภูตํ ญาณํ, ตสฺส อายสฺมโต สํวิชฺชเต, เตน นํ โอวทติ, ตํ ภควา กโรติ.}

\setlength{\csromangathapagewidth}{0pt}
\setlength{\csromangathawakwidth}{0pt}
\csromangathameasuregroup{สเจปิ ทส ปชฺโชเต, \\ เนว ทกฺขติ รูปานิ,}{\csromangathabat{สเจปิ ทส ปชฺโชเต,}{ธารยิสฺสสิ กสฺสป.} \\ \csromangathabat{เนว ทกฺขติ รูปานิ,}{จกฺขุ ตสฺส น วิชฺชตีติ.}}
\csromangathasetleft{สเจปิ ทส ปชฺโชเต, \\ เนว ทกฺขติ รูปานิ,}
\csromangathagroup{\csromangathastanza{\csromangathabat{สเจปิ ทส ปชฺโชเต,}{ธารยิสฺสสิ กสฺสป.} \\ \csromangathabat{เนว ทกฺขติ รูปานิ,}{จกฺขุ ตสฺส น วิชฺชตีติ.}}}

\prose{อปิ จ โข ยถา ทูโต ราชวจเนน สตฺตม\-นุสาสติ, เอวํ เสสานุโค อญฺญาตกํ โฆสํ ปเรสํ เทเสติ. อนุญฺญาต\-ขม\-สุตฺตํ คเหตพฺพํ. อนนุญฺญาตขมํ น คเหตพฺพํ.}

\prose{ตตฺถ กตโม สุตฺตสงฺกโร. ปญฺจวิธํ สุตฺตํ สํกิเลส\-ภาคิยํ วาสนาภาคิยํ ทสฺสน\-ภาคิยํ ภาวนา\-ภาคิยํ อเสกฺข\-ภาคิยํ. อญฺญํ อาราเธยฺย อญฺญํ เทเสติ อญฺญสฺส จ สุตฺตสฺส อตฺถํ อญฺญมฺหิ สุตฺเต นิทฺทิสติ. สุตฺตสฺส วา หิ อเนกาการํ อตฺถํ นิทฺทิสติ. อริยธมฺม\-สาธเน อตฺถํ วิวรติ. วาสนา\-ภาคิยสฺส อตฺถํ ทสฺสน\-ภาคิเยสุ นิทฺทิสติ. โอรมฺภาคิยานํ สํโยชนานํ อตฺถํ อุทฺธํภาคิเยสุ นิทฺทิสติ. มุทุมชฺฌานํ อินฺทฺริยานํ อธิมตฺเตสุ สุตฺเตสุ นิทฺทิสติ. อิติ อยํ สุตฺตํ สมฺเภทํ เหตุนา จ นิสฺสนฺเทน จ ผเลน จ นิทฺเทเสน จ มุทุมชฺฌา\-ธิมตฺตตายปิ จ อตฺเถน จ พฺยญฺชเนน จ โย สมฺเภโท, อยํ วุจฺจติ สุตฺตสงฺกโร. โย อสมฺเภโท, อยํ วุจฺจติ สุตฺตวิจโย.}

\csromanheader{2}{\textbf{ตตฺถายํ อุทฺทานคาถา}}
\setlength{\csromangathapagewidth}{0pt}
\setlength{\csromangathawakwidth}{0pt}
\csromangathameasuregroup{ปุริมานํ อกฺขณฺฑํ, \\ นิสฺสนฺโท วาสนาสทฺธิ,}{\csromangathabat{ปุริมานํ อกฺขณฺฑํ,}{ยถาภูตสฺส ปจฺจโย.} \\ \csromangathabat{นิสฺสนฺโท วาสนาสทฺธิ,}{อนุญฺญา สุตฺตสงฺกโร.} \\ เถรสฺส มหากจฺจายนสฺส สุตฺตวิจโย นาม จตุตฺถภูมิ.}
\csromangathameasurewak{เถรสฺส มหากจฺจายนสฺส สุตฺตวิจโย นาม จตุตฺถภูมิ.}
\csromangathasetleft{ปุริมานํ อกฺขณฺฑํ, \\ นิสฺสนฺโท วาสนาสทฺธิ,}
\csromangathagroup{\csromangathastanza{\csromangathabat{ปุริมานํ อกฺขณฺฑํ,}{ยถาภูตสฺส ปจฺจโย.} \\ \csromangathabat{นิสฺสนฺโท วาสนาสทฺธิ,}{อนุญฺญา สุตฺตสงฺกโร.}}\csromangathastanza{\csromangathawak{เถรสฺส มหา\-กจฺจายนสฺส สุตฺตวิจโย นาม จตุตฺถภูมิ.}}}

\proseitem{43}{ตตฺถ กตโม หารวิภงฺโค. ยตฺถ โสฬส หารา อกฺขรโส เภทํ คจฺฉนฺติ. ตตฺถ อาทิมฺหิ เทสนาหาโร. ตตฺถ อยํ คาถา กุสลา วา อกุสลา วา สจฺจานิ วา สจฺเจกเทโส วา. กิํ เทสิตนฺติ สุตฺเต วีมํสา เทสนาหาโร. ยถา อริยสจฺจานิ นิกฺเขโป จตฺตาริ \csromanfolio{224}สจฺจานิ สาธารณานิ อสาธารณานิ จ. ยานิ จ อฏฺฐารส ปทานิ ทุกฺขโต สตฺต ปทานิ สงฺเขเปน กายิเกน เจตสิเกน ทุกฺเขน, อปฺปิย\-สมฺปโยเคน ปิยวิปฺปโยเคน จ ตีหิ จ สงฺขตาหิ. ตตฺถ ตีณิ สงฺขต\-ลกฺขณานิ ติสฺโส ทุกฺขตา อุปฺปาโท สงฺขต\-ลกฺขณํ, สงฺขาร\-ทุกฺขตาย ทุกฺขตา จ สงฺขต\-ลกฺขณํ, วิปริณาม\-ทุกฺขตาย ทุกฺขตาหิ อญฺญถตฺตํ จ สงฺขต\-ลกฺขณํ, ทุกฺข\-ทุกฺขตาย จ ทุกฺขตา, อิเมสํ ติณฺณํ สงฺขต\-ลกฺขณานํ ตีสุ เวทนาภูมีสุ อทุกฺขมสุขา เวทนา อุปฺปาโท สงฺขต\-ลกฺขณํ, สงฺขาร\-ทุกฺขตาย จ ทุกฺขตา ตโย สงฺขต\-ลกฺขณํ, สุขา เวทนาย จ วิปริณาม\-ทุกฺขตาย จ ทุกฺขตาติ อญฺญถตฺตํ สงฺขต\-ลกฺขณํ, ทุกฺขาเวทนา ทุกฺขทุกฺขตา จ ทุกฺขตา อิมมฺหิ อิเมสุ นวปเทสุ ปฐมเกสุ สตฺตสุ ปเทสุ โสฬสสุ ปเทสุ ธุกฺขา ปริเยสิตพฺพา, เอกาทส ทุกฺขตาย จ ลกฺขณํ นิทฺเทเส นิทฺทิฏฺฐํ. ปาตุภาว\-ลกฺขณา ชาติยา จ ปาตุภาว\-จุติ\-ลกฺขโณ จุโตติ วิตฺถาเรน ปนฺนรส\-ปทานิ กตฺตพฺพานิ, เอวํ สาธารณานิ อสาธารณานิ จ สตฺตสุ ทสสุ ปเทสุ สญฺญาส ติวิเธ จ สาสน\-ปฺปฏฺฐาเน อฏฺฐารส\-วิเธสุ จ สุตฺตา\-ธิฏฺฐาเนสุ ทสวิเธสุ จ สุตฺต\-วิเธยฺเยสุ โสฬสวิเธสุ จ หาเรสุ เอก\-วีสติวิธาย จ ปวิจย\-วีมํสายาติ อิทํ เทสิตํ. ยถาภูตญฺจ เทสิตนฺติ, อยํ วุจฺจติ เทสนาหาโร.}

\setlength{\csromangathapagewidth}{0pt}
\setlength{\csromangathawakwidth}{0pt}
\csromangathameasuregroup{}{ตตฺถ กตโม วิจโย หาโร. \\ ปทํ ปญฺหา จ ปุจฺฉา จ, กิํ ปุพฺพํ กิญฺจ ปจฺฉิมํ. \\ อนุคีติ สา จ วิจโย, หาโร วิจโยติ นิทฺทิฏฺโฐ.}
\csromangathameasurewak{ตตฺถ กตโม วิจโย หาโร. \\ ปทํ ปญฺหา จ ปุจฺฉา จ, กิํ ปุพฺพํ กิญฺจ ปจฺฉิมํ. \\ อนุคีติ สา จ วิจโย, หาโร วิจโยติ นิทฺทิฏฺโฐ.}
\setlength{\csromangathaleftcol}{0pt}
\csromangathagroup{\csromangathastanza{\csromangathawak{\csromangathaitemline{44}{ตตฺถ กตโม วิจโย หาโร.} \\ \csromangathacontline{ปทํ ปญฺหา จ ปุจฺฉา จ, กิํ ปุพฺพํ กิญฺจ ปจฺฉิมํ.} \\ \csromangathacontline{อนุคีติ สา จ วิจโย, หาโร วิจโยติ นิทฺทิฏฺโฐ.}}}}

\prose{ปทนฺติ ปฐมํ ปทํ, ตสฺส โก อตฺโถ, ยํ ภควา ปุฏฺโฐ อายสฺมตา อชิเตน ตํ คเหตพฺพํ, กติปทานิ ปุฏฺฐานิ ยถากิํ เกนสฺสุ นิวุโต โลโกติ คาถา, อิมานิ กติปทานิ จตฺตาริ อิติ วิสชฺชนาย ปุจฺฉา. ยตฺตเกหิ ปเทหิ ภควตา วิสชฺชิตานิ ปทานิ อิติ ปุจฺฉาย จ ยา ปทานํ สงฺกาสนา, อิทํ วุจฺจติ ปทนฺติ.}

\prose{ปญฺหาติ อิมานิ จตฺตาริ ปทานิ, กติ ปญฺหา, เอโก วา เทฺว วา ตทุตฺตริ วา อิมานิ จตฺตาริ ปทานิ เอโก ปโญฺห, อตฺถา\-นุปริวตฺติ พฺยญฺชนํ โหติ, สมฺพหุลานิปิ ปทานิ เอกเมวตฺถํ ปุจฺฉติ. อิมานิ จตฺตาริ ปทานิ อนุปริวตฺตีนิ ตํ พฺยญฺชเนน เอโก ปโญฺหว โหติ. เกนสฺสุ นิวุโต \csromanfolio{225}โลโกติ โลกํ สนฺธาย ปุจฺฉติ, เกนสฺสุ นปฺปกาสติ กิสฺสาภิเลปนํ พฺรูสีติ ตํเยว ปุจฺฉติ. กิํสุ ตสฺส มหพฺภยนฺติ ตํเยว ปุจฺฉติ. เอวํ อตฺถา\-นุปริวตฺติ พฺยญฺชนํ เอโก ปโญฺห โหติ, โส ปโญฺห จตุพฺพิโธ เอกํส\-พฺยากรณีโย วิภชฺช\-พฺยากรณีโย ปฏิปุจฺฉา\-พฺยากรณีโย ฐปนิโยติ, ตตฺถ จกฺขุ อนิจฺจนฺติ เอกํส\-พฺยากรณีโย, ยํ อนิจฺจํ ตํ ทุกฺขนฺติ วิภชฺช\-พฺยากรณีโย, สิยา อนิจฺจํ น จกฺขุ, ยานิปิ อายตนานิ จ น จกฺขุ, ตานิปิ อนิจฺจนฺติ น จกฺขุเยว, อยํ วิภชฺช\-พฺยากรณีโย, ยํ จกฺขุ ตํ จกฺขุนฺทฺริยํ เนติ ปฏิปุจฺฉา\-พฺยากรณีโย, ตํ จกฺขุ ตถาคโตติ ฐปนิโย. อญฺญตฺร จกฺขุนาติ ฐปนิโย ปโญฺห. อิทํ ปญฺหํ ภควา กิํ ปุจฺฉิโต, โลกสฺส สํกิเลโส ปุจฺฉิโต. กิํ การณํ, ติวิโธ หิ สํกิเลโส ตณฺหาสํกิเลโส จ ทิฏฺฐิ\-สํกิเลโส จ ทุจฺจริต\-สํกิเลโส จ. ตตฺถ อวิชฺชาย นิวุโตติ อวิชฺชํ ทสฺเสติ, ชปฺปาติ ตณฺหํ ทสฺเสติ, มหพฺภยนฺติ อกุสลสฺส กมฺมสฺส วิปากํ ทสฺเสติ, โสตํ นาม สุขเวทนียสฺส กมฺมสฺส ทุกฺขเวทนีโย วิปาโก ภวิสฺสตีติ เนตํ ฐานํ วิชฺชตีติ ภควา วิสชฺเชติ, จตูหิ โย ปเทหิ อวิชฺชาย นิวุโต โลโกติ ฯเปฯ เอวํ วุจฺจติ.}

\proseitem{45}{ตทุตฺตริ ปฏิปุจฺฉติ, สวนฺติ สพฺพธิ โสตาติ คาถา, จตฺตาริ ปทานิ ปุจฺฉติ ตํ ภควา ทฺวีหิ ปเทหิ วิสชฺเชติ.}

\setlength{\csromangathapagewidth}{0pt}
\setlength{\csromangathawakwidth}{0pt}
\csromangathameasuregroup{\textsuperscript{*}ยานิ โสตานิ โลกสฺมิํ, \\ โสตานํ สํวรํ พฺรูมิ,}{\csromangathabat{\textsuperscript{*}ยานิ โสตานิ โลกสฺมิํ,}{สติ เตสํ นิวารณํ.} \\ \csromangathabat{โสตานํ สํวรํ พฺรูมิ,}{ปญฺญาเยเต ปิมียเร.}}
\csromangathasetleft{\textsuperscript{*}ยานิ โสตานิ โลกสฺมิํ, \\ โสตานํ สํวรํ พฺรูมิ,}
\csromangathagroup{\csromangathastanza{\csromangathaitembat{45}{\csromansymbolfootnote{*}{เหฏฺฐา 13, 180 ปิฏฺเฐสุปิ.}ยานิ โสตานิ โลกสฺมิํ,}{สติ เตสํ นิวารณํ.} \\ \csromangathacontbat{โสตานํ สํวรํ พฺรูมิ,}{ปญฺญาเยเต ปิมียเร.}}}

\prose{อิมานิ จตฺตาริ ปทานิ ทฺวีหิ ปเทหิ วิสชฺเชติ. อิทํ ปทนฺติ ปุจฺฉิโต, ตสฺส สํกิลิฏฺฐสฺส โลกสฺส โวทานํ ปุจฺฉิโต, โสตานิ ฉ ตณฺหากายา พหุลา\-ธิวจเนน นิทฺทิฏฺฐา ภวนฺติ สพฺเพหิ อายตเนหิ. ตานิ โสตานิ เกน นิวาริยนฺตีติ ปริยุฏฺฐาน\-ปหานํ ปุจฺฉติ, เกน โสตา ปิธียเรติ อนุสย\-สมุคฺฆาตํ ปุจฺฉติ. ตตฺถ ภควา ฉสุ ทฺวาเรสุ สติยา เทเสติ, โย หิ สมฺปชาโน วิหรติ สติโทวาริเกน จ ตสฺส อินฺทฺริยานิ คุตฺตานิ สมฺภวนฺติ. ตตฺถ คุตฺเตสุ อินฺทฺริเยสุ ยา ยา วิปสฺสนา, สา สา เตสํ เตสํ โสตานํ ตสฺสา จ อวิชฺชาย โย โลโก นิวุโต อจฺจนฺต\-ปหานาย สํวตฺตติ. เอวํ โสตานิ ปิหิตานิปิ ภวนฺติ ตโต อุตฺตริ ปุจฺฉติ.}

\prose{\csromanfolio{226}ปญฺญา จ สติ จ นามรูปสฺส โข ตสฺส ภควนฺตํ ปุฏฺฐุมาคมฺม กตฺเถตํ อุปสมฺมติ อิมานิ จตฺตาริ ปทานิ ภควา เอเกน ปเทน วิสชฺเชติ.}

\setlength{\csromangathapagewidth}{0pt}
\setlength{\csromangathawakwidth}{0pt}
\csromangathameasuregroup{ยเมตํ ปญฺหํ อปุจฺฉิ\textsuperscript{1}, \\ วิญฺญาณสฺส นิโรเธน,}{\csromangathabat{ยเมตํ ปญฺหํ อปุจฺฉิ\textsuperscript{1},}{อชิต ตํ วทามิ เต ฯเปฯ.} \\ \csromangathabat{วิญฺญาณสฺส นิโรเธน,}{เอตฺเถตํ อุปสมฺมติ.}}
\csromangathasetleft{ยเมตํ ปญฺหํ อปุจฺฉิ\textsuperscript{1}, \\ วิญฺญาณสฺส นิโรเธน,}
\csromangathagroup{\csromangathastanza{\csromangathabat{ยเมตํ ปญฺหํ อปุจฺฉิ\csromansharedfootnote{7b35b3e3a4bbb846}{ปุจฺฉเส ปญฺหํ (อิ, ก) ขุ 1. 434 ปิฏฺเฐ.},}{อชิต ตํ วทามิ เต ฯเปฯ.} \\ \csromangathabat{วิญฺญาณสฺส นิโรเธน,}{เอตฺเถตํ อุปสมฺมติ.}}}

\prose{อิมินา ปเญฺหน กิํ ปุจฺฉติ, อนุปาทิเสส\-นิพฺพานธาตุํ ปุจฺฉติ, ตํ ภควา อนุปาทิเสสาย นิพฺพานธาตุยา วิสชฺเชติ. ตตฺถ ปฐเมน ปเญฺหน สํกิเลสํ ปุจฺฉติ. ทุติเยน ปเญฺหน โวทานํ ปุจฺฉติ. ตติเยน ปเญฺหน โสปาทิเสสนิพฺพานธาตุํ ปุจฺฉติ. จตุตฺเถน ปเญฺหน อนุปาทิเสส\-นิพฺพานธาตุํ ปฏิปุจฺฉติ ตโต อุตฺตริ ปฏิปุจฺฉติ.}

\setlength{\csromangathapagewidth}{0pt}
\setlength{\csromangathawakwidth}{0pt}
\csromangathameasuregroup{เย จ สงฺขาตธมฺมาเส, \\ เตสํ เม นิปโก อิริยํ,}{\csromangathabat{เย จ สงฺขาตธมฺมาเส,}{เย จ เสขา\textsuperscript{1} ปุถู อิธ.} \\ \csromangathabat{เตสํ เม นิปโก อิริยํ,}{ปุฏฺโถ ปพฺรูหิ มาริส.}}
\csromangathasetleft{เย จ สงฺขาตธมฺมาเส, \\ เตสํ เม นิปโก อิริยํ,}
\csromangathagroup{\csromangathastanza{\csromangathabat{เย จ สงฺขาต\-ธมฺมาเส,}{เย จ เสขา\csromansharedfootnote{956ef50c91d43a25}{เสกฺขา (ก) ขุ 1. 434; เหฏฺฐา 16 ปิฏฺเฐสุปิ.} ปุถู อิธ.} \\ \csromangathabat{เตสํ เม นิปโก อิริยํ,}{ปุฏฺโถ ปพฺรูหิ มาริส.}}}

\prose{อิมานิ จตฺตาริ ปทานิ ปุจฺฉติ. กติ จ ปน เต ปเญฺห สงฺขาตธมฺมา จ อรหนฺตา เสกฺขา จ. กิํ ปุพฺพํ กิญฺจ ปจฺฉิมนฺติ อยมตฺโถ. ตตฺถ กตรํ ปฐมํ ปุจฺฉติ, กตรํ ปุจฺฉา. อรํหนฺตํ ปฐมํ ปุจฺฉติ. เสกฺขธมฺเม ตตฺถ เกน ปเทน สงฺขาตธมฺมาติ อรหนฺโต คหิตา. ปุถูติ เสกฺขา คหิตา. เตสํ เม นิปโกติ สาธารณํ ปทํ ภควนฺตํ ปุจฺฉติ. ตสฺส สาธารณานิ จ อสาธารณานิ จ ปเญฺหสุ ปุจฺฉิตพฺพานิ. ตํ ภควา วิสชฺเชติ. น ตถา ปุฏฺฐํ, ปฐมํ ปุฏฺฐํ, ตํ ปุจฺฉา วิสชฺเชติ. ยํ ปจฺฉา ปุจฺฉิตํ ปฐมํ วิสชฺเชติ. กิญฺจ อิทํ ปุจฺฉิตํ วิสุทฺธานํ วิสุชฺฌ\-นฺตานญฺจ กา อิริยาติ อิทํ ปุจฺฉิ, ตํ กาเมสุ นาภิคิชฺเฌยฺย. มนสานาวิโล สิยาติ ปริยุฏฺฐานานิ วิตกฺเกน จ ภควา นิวาเรติ, เทฺว ปน วิตกฺก\-อนาวิล\-ตาย ปริยุฏฺฐานํ, ยถา นีวรเณสุ นิทฺทิฏฺฐํ. กุสลา สพฺพธมฺเมสูติ อรหนฺตํ วิสชฺเชติ.}

\prose{เกนสฺสุ ตรติ โอฆนฺติ คาถา, อิมานิ จตฺตาริ ปทานิ. จตฺตาโรเยว ปญฺหา. กิํ การณํ, น หิ เอตฺถ อตฺถา\-นุปริวตฺติ พฺยญฺชนํ\csromansharedfootnote{1ccb3a9c32ab70dc}{ยถานุปริวตฺถิวชฺชํ (อิ, ก)} ยถา ปฐมํ อชิตปเญฺหสุ, ตสฺส น เอกํเสน พหูนิ วิสชฺชนานิ, พหุกา ปญฺหา, เอโกว น จาปิ, สพฺเพ ปุจฺฉติ, ปุพฺเพ วิสชฺชิโต, ยถา จตุตฺโถ อชิโตปเญฺห, \csromanfolio{227}ยํ เอตฺถ ยถาภูตํ ปริเยสนา\-ปท\-พนฺเธน วิสชฺชนาโย เอวํ ยถาภูตํ ปริเยสติ. โย ปุน เอตฺถ ยํ เอวํ ปุจฺฉติ ตตฺถ อยมากาโร ปุจฺฉนายํ อนฺโตชฏา พหิชฏาติ คาถา\csromansharedfootnote{772b02f5b21a4af8}{สํ 1. 13 ปิฏฺเฐ.} ปุจฺฉิต\-วิสชฺช\-นาย มคฺคิตพฺพา. กถํ วิสชฺชิตาติ ภควาติ วิสชฺเชติ. สีเล ปติฏฺฐาย นโร สปญฺโญติ คาถา. ตตฺถ จิตฺตภาวนาย สมถา, ปญฺญาภาวนาย วิปสฺสนา. ตตฺถ เอวํ อนุมียติ, เย ธมฺมา สมเถน จ วิปสฺสนาย จ ปหียนฺติ, เต อิเม อนฺโตชฏา พหิชฏา. ตตฺถ วิสชฺชนํ สมเถน ราโค ปหียติ, วิปสฺสนาย อวิชฺชา. อชฺฌตฺต\-วตฺถุโก ราโค อนฺโตชฏา, พาหิรวตฺถุโก ราโค พหิชฏา. อชฺฌตฺต\-วตฺถุกา สกฺกายทิฏฺฐิ, อยํ อนฺโตชฏา. เอกสฏฺฐิ ทิฏฺฐิคตานิ จ พาหิร\-วตฺถุกานิ พหิชฏา, ยา หิ อชฺฌตฺต\-วตฺถุกา ยา ทิฏฺฐิ\-ภาคิเยน ภวิสฺสติ, อยํ ชฏา. ตถา สํขิตฺเตน ยา กาจิ อชฺฌตฺต\-วตฺถุกา ตณฺหา จ ทิฏฺฐิ จ, อยํ อนฺโตชฏา. ยา กาจิ พาหิรวตฺถุกา ตณฺหา จ ทิฏฺฐิ จ, อยํ พหิชฏา.}

\proseitem{45}{ยถา เทวตา ภควนฺตํ ปุจฺฉติ “จตุจกฺกํ นวทฺวารํ”ติ คาถา\csromansharedfootnote{5f3920b5ab910035}{สํ 1. 15 ปิฏฺเฐ.}. ตตฺถ ภควา วิสชฺเชติ “เฉตฺวา นทฺธิํ วรตฺตํ จา”ติ คาถา, อิทํ ภควา ทุกฺขนิโรธ\-คามินิํ ปฏิปทํ วิสชฺเชติ. อิมาย วิสชฺชนาย ภควา อนุมียติ กิเลเส เอตฺถ ปุริมาย คาถาย นิทฺทิสิตพฺเพน. ตํ หิ จตุจกฺกนฺติ จตฺตาโร วา หตฺถปาทา. นวทฺวารนฺติ นว วณมุขานิ. ยถา จตุจกฺกนฺติ จตฺตาโร อุปาทานา, อุปาทาน\-ปฺปจฺจยา ภโว, อุปาทานนิโรธา ภวนิโรโธ. นวทฺวารนฺติ นว มานวิธา, มานชาติกาย หิ ทุกฺขํ เสยฺเยนมฺหิ ปรโส ตีณิ ติกานิ ปุณฺณํ. ติเกน สํยุตฺตํ หิ ปญฺจ\-กาม\-คุณิโก ราโค. ตตฺถ นทฺธีติ ตณฺหา วิสชฺชียติ. วรตฺตนฺติ มานํ วิสชฺเชติ, อิจฺฉา โลโภ จ ปาปโกติ ปญฺจ\-กาม\-คุณิโก ราโค. ตตฺถ วิสมโลโภ ปาปโกติ นิทฺทิสิยติ สมูลตณฺหนฺติ. อญฺญาณมูลกา ตณฺหาติ อญฺญาณมูลกา ตณฺหา, ตณฺหาย จ ทิฏฺฐิยา จ ปหานํ. เย จ ปุน อญฺเญปิ เกจิ จตุจกฺก\-โยเคน เตเนว การเณน จ ยุชฺชนฺติ, สํสารคามิโน ธมฺมา สพฺเพ นิทฺทิสิตพฺพา. ตตฺถายํ คาถา วิสชฺชนา ปุจฺฉาย จ วิสชฺชนาย สเมติ\csromansharedfootnote{5a524000a4783071}{สมํติ (อิ)}. ยํ ยทิ สนฺเทน อถ สห พฺยากรเณน อนุคีติยํ จ โส วิจโยติ ภควา ยตฺตกานิ ปทานิ นิกฺขิปติ, ตตฺถเกหิ อนุคายติ.}

\proseitem{46}{\csromanfolio{228}อฏฺฐหิ ภิกฺขเว องฺเคหิ สมนฺนาคโต ภิกฺขุ ทูเตยฺยํ คนฺตุมรหติ\csromansharedfootnote{933970d6dcc690e1}{กาตุมรหติ (อิ, ก) อํ 3. 37 ปิฏฺเฐ.}. อิมานิ อฏฺฐ ปทานิ นิกฺขิตฺตานิ. ฉหิ ปเทหิ ภควา อนุคายติ.}

\setlength{\csromangathapagewidth}{0pt}
\setlength{\csromangathawakwidth}{0pt}
\csromangathameasuregroup{“โย เว น พฺยถติ\textsuperscript{1} ปตฺวา, \\ น จ หาเปติ วจนํ, \\ อสนฺทิทฺธิํ จ ภณติ, \\ ส เว ตาทิสโก ภิกฺขุ,}{\csromangathabat{“โย เว น พฺยถติ\textsuperscript{1} ปตฺวา,}{ปริสํ อุคฺควาทินิํ.} \\ \csromangathabat{น จ หาเปติ วจนํ,}{น จ ฉาเทติ สาสนํ.} \\ \csromangathabat{อสนฺทิทฺธิํ จ ภณติ,}{ปุจฺฉิโต น จ กุปฺปติ.} \\ \csromangathabat{ส เว ตาทิสโก ภิกฺขุ,}{ทูเตยฺยํ คนฺตุมรหตี”ติ.}}
\csromangathasetleft{“โย เว น พฺยถติ\textsuperscript{1} ปตฺวา, \\ น จ หาเปติ วจนํ, \\ อสนฺทิทฺธิํ จ ภณติ, \\ ส เว ตาทิสโก ภิกฺขุ,}
\csromangathagroup{\csromangathastanza{\csromangathabat{“โย เว น พฺยถติ\csromansharedfootnote{252e5374af655ee4}{พฺยาถติ (ก)} ปตฺวา,}{ปริสํ อุคฺควาทินิํ.} \\ \csromangathabat{น จ หาเปติ วจนํ,}{น จ ฉาเทติ สาสนํ.}}\csromangathastanza{\csromangathabat{อสนฺทิทฺธิํ จ ภณติ,}{ปุจฺฉิโต น จ กุปฺปติ.} \\ \csromangathabat{ส เว ตาทิสโก ภิกฺขุ,}{ทูเตยฺยํ คนฺตุมรหตี”ติ.}}}

\prose{ตตฺถ ปน ภควา ยตฺตกานิ ปทานิ นิกฺขิปติ, ตตฺถเกหิ อนุคายติ. สตฺตหิ ภิกฺขเว องฺเคหิ สมนฺนาคโต กลฺยาณมิตฺโต ปิโย ครุภาวนีโยติ วิตฺถาเรน, อิทํ ภควา สตฺตหิ ปเทหิ อนุคายติ. อิติ พหุสฺสุตวา อนุคายติ, อปฺปตรกถํ ปทํ วา นิกฺเขโป, พหุสฺสุตวา นว ปทานิ นิกฺเขโป, อปฺปตริกา อนุคีติยา พหุตริกา อนุคายติ. อยํ วุจฺจติ เต อนุคีติ จ วิจโย, อยํ วิจโย นาม หาโร.}

\prose{ตตฺถ กตโม ยุตฺติหาโร.}

\setlength{\csromangathapagewidth}{0pt}
\setlength{\csromangathawakwidth}{0pt}
\csromangathameasuregroup{\textsuperscript{*}สพฺเพสํ หารานํ, \\ ยุตฺตายุตฺติ ปริกฺขา,}{\csromangathabat{\textsuperscript{*}สพฺเพสํ หารานํ,}{ยา ภูมี โย จ โคจโร เตสํ.} \\ \csromangathabat{ยุตฺตายุตฺติ ปริกฺขา,}{หาโร ยุตฺตีติ นิทฺทิฏฺโฐ.}}
\csromangathasetleft{\textsuperscript{*}สพฺเพสํ หารานํ, \\ ยุตฺตายุตฺติ ปริกฺขา,}
\csromangathagroup{\csromangathastanza{\csromangathabat{\csromansymbolfootnote{*}{เหฏฺฐา 3 ปิฏฺเฐ.}สพฺเพสํ หารานํ,}{ยา ภูมี โย จ โคจโร เตสํ.} \\ \csromangathabat{ยุตฺตายุตฺติ ปริกฺขา,}{หาโร ยุตฺตีติ นิทฺทิฏฺโฐ.}}}

\prose{หารานํ โสฬสนฺนํ ยถา เทสนา ยถา วิจโย โย จ นิทฺทิสิยติ, อยํ นิทฺเทโส. อยํ ปุจฺฉา สุตฺเตสุ น ยุชฺชตีติ ยา ตตฺถ วีมํสา, อยํ ยุตฺติ.}

\prose{ยถา หิ สเหตู สปฺปจฺจยา สตฺตา สํกิลิสฺสนฺติ, อตฺถิ เหตุ อตฺถิ ปจฺจโย สตฺตานํ สํกิเลสาย, สเหตู สปฺปจฺจยา สตฺตา วิสุชฺฌนฺติ, อตฺถิ เหตุ อตฺถิ ปจฺจโย สตฺตานํ วิสุทฺธิยา. สีลวตา อานนฺท ปุคฺคเลน น เวยฺยากรณิยา กินฺติ เม วิปฺปฏิสาโร อุปฺปาเทยฺย ฯเปฯ อพฺยากรณํ กตฺตพฺพํ, อยํ วิสุทฺธิยา มคฺโค. ตสฺส เหตุ โก ปจฺจโย, สีลกฺขนฺธสฺส จตฺตาริ จตฺตาริ เหตุ จ ปจฺจโย จ. สปฺปุริส\-สํเสโว โย จ ปติรูป\-เทส\-วาโส จ, อยํ อุปาทาปจฺจยตา สปฺปจฺจโย. ยํ โปราณกมฺมํ อสฺส วิปาโก ปจฺจโย, ตาย ปจฺจยาย อตฺต\-สมฺมา\-ปณิธิ, \csromanfolio{229}อยํ เหตุ. อิติ สีลกฺขนฺโธ สเหตุ สปฺปจฺจโยติ อิทํ โลกิกํ สีลํ.}

\prose{ยํ ปน โลกุตฺตรํ สีลํ, ตสฺส ตีณิ อินฺทฺริยานิ ปจฺจโย สทฺธินฺทฺริยํ วีริยินฺทฺริยํ สมาธินฺทฺริยํ, อยํ ปจฺจโย. สตินฺทฺริยญฺจ ปญฺญินฺทฺริยญฺจ เหตุ. ปญฺญาย นิพฺเพธ\-คามินิยา, ยํ สีลํ ชายติ. โสตาปนฺนสฺส จ สีลํ เตนายํ เหตุ อยํ ปจฺจโย. ยํ ปุน สมาธิโน ปสฺสทฺธิ จ ปีติ จ ปาโมชฺชํ ปจฺจโย. ยํ สุขํ เหตุ เตน สมาธิ\-กฺขนฺโธ สเหตุ สปฺปจฺจโย. ยํ สมาหิโต ยถาภูตํ ปชานาติ, อยํ ปญฺญา. ตสฺส ปรโตโฆโส อชฺฌตฺตํ จ โยนิโส มนสิกาโร เหตุ จ ปจฺจโย จ, อิติ อิเม ตโย ขนฺธา สเหตู สปฺปจฺจยา เอวํ สตฺต ปญฺญา. สตฺตพฺยากรณีสุ จ สุตฺเตสุ น ยุชฺชติ. อยํ ยุตฺติหาโร. โส จตูสุ มหาปเทเสสุ ทฏฺฐพฺโพ.}

\proseitem{47}{ตตฺถ กตมํ ปทฏฺฐานํ.}

\setlength{\csromangathapagewidth}{0pt}
\setlength{\csromangathawakwidth}{0pt}
\csromangathameasuregroup{\textsuperscript{*}ธมฺมํ เทเสติ ชิโน, \\ อิติ ยาว สพฺพธมฺมา,}{\csromangathabat{\textsuperscript{*}ธมฺมํ เทเสติ ชิโน,}{ตสฺส จ ธมฺมสฺส ยํ ปทฏฺฐานํ.} \\ \csromangathabat{อิติ ยาว สพฺพธมฺมา,}{เอโส หาโร ปทฏฺฐาโน.}}
\csromangathasetleft{\textsuperscript{*}ธมฺมํ เทเสติ ชิโน, \\ อิติ ยาว สพฺพธมฺมา,}
\csromangathagroup{\csromangathastanza{\csromangathaitembat{47}{\csromansymbolfootnote{*}{เหฏฺฐา 3 ปิฏฺเฐ.}ธมฺมํ เทเสติ ชิโน,}{ตสฺส จ ธมฺมสฺส ยํ ปทฏฺฐานํ.} \\ \csromangathacontbat{อิติ ยาว สพฺพธมฺมา,}{เอโส หาโร ปทฏฺฐาโน.}}}

\prose{ตตฺถ ปญฺจ กามคุณา กามราคสฺส ปทฏฺฐานํ. เยสํ เกสญฺจิ กามราโค อุปฺปชฺชติ อุปฺปนฺโน วา อุปฺปชฺชิสฺสติ วา, เอเตสุ เยปิ ปญฺจสุ รูเปสุ อายตเนสุ นาญฺญตฺร เอเตหิ กามราคสฺส ปทฏฺฐานนฺติ. วุจฺจเต, เตน ปญฺจ กามคุณา กามราคสฺส ปทฏฺฐานํ. ปญฺจินฺทฺริยานิ รูปราคสฺส ปทฏฺฐานํ. มนินฺทฺริยํ ภวราคสฺส ปทฏฺฐานํ. ปญฺจกฺขนฺธา สกฺกาย\-ทิฏฺฐิยา ปทฏฺฐานํ. เอกสฏฺฐิ ทิฏฺฐิคตานิ ทิฏฺฐิราคสฺส ปทฏฺฐานํ. กามธาตุ กามราคสฺส ปทฏฺฐานํ. อรูปธาตุ อรูปราคสฺส ปทฏฺฐานํ. สุขสญฺญา กามราคสฺส ปทฏฺฐานํ. พฺยาปาทสญฺญา พฺยาปาทสฺส ปทฏฺฐานํ. อสมฺปชญฺญตา สมฺโมหสฺส ปทฏฺฐานํ. นว อาฆาตวตฺถูนิ พฺยาปาทสฺส ปทฏฺฐานํ. นววิธํ มานํ มานสฺส\csromansharedfootnote{4591d5f9869ae750}{นวมานํ วิธมานสฺส (อิ, ก)} ปทฏฺฐานํ. สุขา เวทนา ราคานุสยสฺส ปทฏฺฐานํ. ทุกฺขา เวทนา ปฏิฆา\-นุสยสฺส ปทฏฺฐานํ. อทุกฺขมสุขา เวทนา อวิชฺชา\-นุสยสฺส ปทฏฺฐานํ. อตฺตวาทุปาทานญฺจ มุสาวาโท จ โลภสฺส ปทฏฺฐานํ. ปาณาติปาโต จ ปิสุณวาจา จ ผรุสวาจา จ พฺยาปาทสฺส ปทฏฺฐานํ. มิจฺฉตฺตญฺจ สมฺผปฺปลาโป จ โมหสฺส ปทฏฺฐานํ. ภวํ โภคญฺจ โวกาโร อหํการสฺส ปทฏฺฐานํ. พาหิรานํ \csromanfolio{230}ปริคฺคโห มมํการสฺส ปทฏฺฐานํ. กายสฺส สงฺคํ\csromansharedfootnote{f0ff574c85dca1ce}{กายวงฺกํ (อิ)} ทิฏฺฐิยา ปทฏฺฐานํ. กายิกโทโส โทสสฺส ปทฏฺฐานํ. กายิกกาสาโว โลภสฺส ปทฏฺฐานํ. โย โย วา ปน ธมฺโม เยน เยน อารมฺมเณน อุปฺปชฺชติ สจฺจา\-ธิฏฺฐาเนน วา ธมฺมา\-ธิฏฺฐาเนน วา อนุสยเนน วา, โส ธมฺโม ตสฺส ปทฏฺฐานํ. เตน สารมฺมเณน โส ธมฺโม อุปฺปชฺชติ.}

\prose{ยถา มนุสฺโส ปุริมสฺส ปทสฺส ปทฏฺฐานํ อลภนฺโต ทุติยํ ปทํ อุทฺธรติ, โส ปจฺฉานุปทํ สํหรติ. ยทิ ปน โย น ทุติยปทสฺส ปทฏฺฐานํ ลภติ, อปรํ ปทํ อุทฺธรติ. ตสฺส โย เจโส ปจฺจโย ภวติ. เอวํ ธมฺโม กุสโล วา อกุสโล วา อพฺยากโต วา ปทฏฺฐานํ อลภนฺโต น ปวตฺตติ. ยถา ปยุตฺตสฺส ธมฺมสฺส โยนิลาโภ\csromansharedfootnote{6c0fa8604d1c434c}{โยนิโส ลาโภ (อิ)}, อยํ วุจฺจติ ปทฏฺฐาโน หาโร.}

\proseitem{48}{ตตฺถ กตโม ลกฺขโณ หาโร.}

\setlength{\csromangathapagewidth}{0pt}
\setlength{\csromangathawakwidth}{0pt}
\csromangathameasuregroup{\textsuperscript{*}วุตฺตมฺหิ เอกธมฺเม, \\ สพฺเพ ภวนฺติ วุตฺตา,}{\csromangathabat{\textsuperscript{*}วุตฺตมฺหิ เอกธมฺเม,}{เย ธมฺมา เอกลกฺขณา เตน.} \\ \csromangathabat{สพฺเพ ภวนฺติ วุตฺตา,}{โส หาโร ลกฺขโณ นาม.}}
\csromangathasetleft{\textsuperscript{*}วุตฺตมฺหิ เอกธมฺเม, \\ สพฺเพ ภวนฺติ วุตฺตา,}
\csromangathagroup{\csromangathastanza{\csromangathaitembat{48}{\csromansymbolfootnote{*}{เหฏฺฐา 3 ปิฏฺเฐ.}วุตฺตมฺหิ เอกธมฺเม,}{เย ธมฺมา เอกลกฺขณา เตน.} \\ \csromangathacontbat{สพฺเพ ภวนฺติ วุตฺตา,}{โส หาโร ลกฺขโณ นาม.}}}

\prose{เยสญฺจ สุสมารทฺธา, นิจฺจํ กายคตาสตีติ คาถาย วุตฺตาย กายคตาสติยา วุตฺตา เวทนาคตา จิตฺตคตา ธมฺมคตา จ สติ จตุนฺนํ สติปฏฺฐานานํ เอเกน สติปฏฺฐาเนน. น หิ จิตฺตํ เอกสฺมิํ วิญฺญาณ\-ฏฺฐิติยา ปวตฺตติ, นานาสุ คตีสุ ปวตฺตติ, กายคตาสติยา วุตฺตาย วุตฺตา เวทนาคตา จิตฺต\-ธมฺมคตา จ. น หิ กายคตาสติยา ภาวิตาย สติปฏฺฐานา จตฺตาโร ภาวนา\-ปาริปูริํ น คจฺฉนฺติ. เอวํ ตสฺสทิเสสุ ธมฺเมสุ วุตฺเตสุ สพฺพธมฺมา วุตฺตา จ ภวนฺติ.}

\prose{สจิตฺต\-ปริโยทาปนํ, เอตํ พุทฺธาน สาสนนฺติ คาถา เจตสิกา ธมฺมา วุตฺตา, จิตฺเต รูปํ วุตฺตํ. อิทํ นามรูปํ ทุกฺขํ อริยสจฺจํ. ตโต สจิตฺต\-ปริโยทาปนา ยํ ยํ โอทเปติ, ตํ ทุกฺขํ. เยน โอทเปติ, โส มคฺโค. ยโต โอทปนา, โส นิโรโธ. จกฺขุํ จ ปฏิจฺจ รูเป จ อุปฺปชฺชติ จกฺขุวิญฺญาณํ, ตตฺถ สหชาตา เวทนา สญฺญา เจตนา ผสฺโส มนสิกาโร เอเตเต ธมฺมา เอกลกฺขณา อุปฺปาท\-ลกฺขเณน. โย จ รูเป นิพฺพินฺทติ, เวทนาย โส นิพฺพินฺทติ, สญฺญาสงฺขารวิญฺญาเณสุปิ โส นิพฺพินฺทติ. \csromanfolio{231}อิติ เย เอกลกฺขณา ธมฺมา, เตสํ เอกมฺหิ ธมฺเม นิทฺทิฏฺเฐ สพฺเพ ธมฺมา นิทฺทิฏฺฐา โหนฺติ, อยํ วุจฺจติ ลกฺขโณ หาโร.}

\prose{ตตฺถ กตโม จตุพฺยูโห หาโร.}

\setlength{\csromangathapagewidth}{0pt}
\setlength{\csromangathawakwidth}{0pt}
\csromangathameasuregroup{\textsuperscript{*}นิรุตฺติ อธิปฺปาโย จ, \\ สุตฺตตฺโถ ปุพฺพาปรสนฺธิ,}{\csromangathabat{\textsuperscript{*}นิรุตฺติ อธิปฺปาโย จ,}{พฺยญฺชนา เทสนาย จ.} \\ \csromangathabat{สุตฺตตฺโถ ปุพฺพาปรสนฺธิ,}{เอโส หาโร จตุพฺยูโห.}}
\csromangathasetleft{\textsuperscript{*}นิรุตฺติ อธิปฺปาโย จ, \\ สุตฺตตฺโถ ปุพฺพาปรสนฺธิ,}
\csromangathagroup{\csromangathastanza{\csromangathaitembat{48}{\csromansymbolfootnote{*}{เหฏฺฐา 3 ปิฏฺเฐ ปสฺสิตพฺพํ.}นิรุตฺติ อธิปฺปาโย จ,}{พฺยญฺชนา เทสนาย จ.} \\ \csromangathacontbat{สุตฺตตฺโถ ปุพฺพาปร\-สนฺธิ,}{เอโส หาโร จตุพฺยูโห.}}}

\proseitem{5}{ตตฺถ กตมา นิรุตฺติ, สา กถํ ปริเยสิตพฺพา\csromansharedfootnote{053b86a324002b14}{ปสฺสิตพฺพา (อิ, ก)}. ยถา วุตฺตํ ภควตา เอกาทสหิ องฺเคหิ สมนฺนาคโต ภิกฺขุ ขิปฺปํ ธมฺเมสุ มหตฺตํ ปาปุณาติ, อตฺถกุสโล จ โหติ, ธมฺมกุสโล จ โหติ, นิรุตฺติกุสโล จ โหติ, อิตฺถา\-ธิวจน\-กุสโล จ โหติ, ปุริสา\-ธิวจน\-กุสโล จ, วิปุริสา\-ธิวจน\-กุสโล จ, อตีตา\-ธิวจน\-กุสโล จ, อนาคตา\-ธิวจน\-กุสโล จ, ปจฺจุปฺปนฺนา\-ธิวจน\-กุสโล จ. เอกาธิปฺปาเยน กุสโล นานาธิปฺปาเยน กุสโล. กิมฺหิ เทสิตํ, อตีตา\-นาคต\-ปจฺจุปฺปนฺนํ. อิตฺถา\-ธิวจเนน ปุริสา\-ธิวจเนน วิปุริสา\-ธิวจเนน สพฺพํ ยถาสุตฺตํ นิทฺทิฏฺฐํ. ตํ พฺยญฺชนโต นิรุตฺติ\-โกสลฺลโต โย ยํ สุตฺตสฺส สุนิรุตฺติ\-ทุ\-นฺนิรุตฺติตํ อเวกฺขติ, อิทํ เอวํ นิโรปยิตพฺพํ. อิทมฺปิ น นิโรปยิตพฺพํ. อิทํ วุจฺจเต นิรุตฺติ\-โกสลฺลํ.}

\proseitem{49}{ตตฺถ กตมํ อธิปฺปาย\-โกสลฺลํ. ยถาเทสิตสฺส สุตฺตสฺส สพฺพสฺส วารํ คจฺฉติ อิเมน ภควตา เทสิตพฺพนฺติ. ยถา กิํ อปฺปมาโท อมตํ ปทํ, ปมาโท มจฺจุโน ปทนฺติ คาถา. เอตฺถ ภควโต โก อธิปฺปาโย เย อสีติเมว อากงฺขนฺติ เต อปฺปมตฺตา วิหริสฺสนฺติ, อยํ อธิปฺปาโย.}

\setlength{\csromangathapagewidth}{0pt}
\setlength{\csromangathawakwidth}{0pt}
\csromangathameasuregroup{}{โยคสฺส กาลํ น นิวตฺตติ ยา จ, \\ โส น ตตฺถ ปาปินฺตเว ภวนฺติ. \\ เวทนามคฺคอิสินา\textsuperscript{1} ปเวทิตํ, \\ ธุตรชาสวา ทุกฺขา ปโมกฺขาตา.}
\csromangathameasurewak{โยคสฺส กาลํ น นิวตฺตติ ยา จ, \\ โส น ตตฺถ ปาปินฺตเว ภวนฺติ. \\ เวทนามคฺคอิสินา\textsuperscript{1} ปเวทิตํ, \\ ธุตรชาสวา ทุกฺขา ปโมกฺขาตา.}
\setlength{\csromangathaleftcol}{0pt}
\csromangathagroup{\csromangathastanza{\csromangathawak{โยคสฺส กาลํ น นิวตฺตติ ยา จ, \\ โส น ตตฺถ ปาปินฺตเว ภวนฺติ. \\ เวทนามคฺคอิสินา\csromansharedfootnote{c40e6729ac7738a0}{เวทนามคฺคํ อิสินา (อิ)} ปเวทิตํ, \\ ธุตรชาสวา ทุกฺขา ปโมกฺขาตา.}}}

\prose{เอตฺถ ภควโต โก อธิปฺปาโย. เย ทุกฺเข นาสฺสาทกา\csromansharedfootnote{5bb2ddc8e1a1a9be}{ทุกฺเขน สาธกา (อิ)}, เต วีริยมา\-รภิสฺสนฺติ ทุกฺขกฺขยายาติ. อยํ ตตฺถ ภควโต อธิปฺปาโย. อิติ คาถาย วา พฺยากรเณน วา เทสิเต อิมินา สุตฺเตน สาธกา, \csromanfolio{232}โย เอวํ ธมฺมา\-นุธมฺมํ ปฏิปชฺชตีติ โส อธิปฺปาโย, อยํ วุจฺจติ เทสนาธิปฺปาโย.}

\prose{ตตฺถ กตโม ปุพฺพาปร\-สนฺธิ. ยํ คาถายํ วา สุตฺเตสุ วา ปทานิ อสีติ ตานิ ภวนฺติ เอวํ วา เอวเมติ ตสฺสา คาถาย สุตฺตสฺส วา ยานิ ปุริมานิ ปทานิ ยานิ จ ปจฺฉิมกานิ, ตานิ สโม\-สาเรตพฺพานิ. เอวํ โส ปุพฺพาปเรน สนฺธิ ญายติ. ยา เอกา สมารทฺธา คาถา เทฺว ตีณิ วา ตสฺส เมกเทเส ภาสิตานํ อภาสิตาหิ คาถาหิ อนิทฺทิฏฺโฐ อตฺโถ ภวติ ตทุปธาริตพฺพํ. ยํว สพฺพา\csromansharedfootnote{2faf625d6debb303}{ยํ วตฺตพฺพํ (อิ)} อิติสฺส ปริเยสมานสฺส ปริเยสนา กงฺขา, ตสฺส วา ปุคฺคลสฺส ปญฺญตฺตีนํ อปเร ปริเยสิตพฺพํ. อิทํ วุจฺจเต ปุพฺพาปเรน สนฺธิ. โกสลฺลนฺติ วตฺถุโต นิทานโกสลฺลํ. พฺยญฺชนโต นิรุตฺติ\-โกสลฺลํ. เทสนา\-ธิปฺปาย\-โกสลฺลํ. ปุพฺพาปเรน สนฺธิโกสลฺลํ. ตตฺถ ตสฺส คาถา ปริเยสิตา นิทานํ วา. อุปลพฺภิตุํ น อตฺโถ นิทฺทิสิตพฺโพ วตฺถุโต นิทานโกสลฺลํ อตฺถโกสลฺลํ อิเมหิ จตูหิ ปเทหิ อตฺโถ ปริเยสิยนฺโต ยถาภูตํ ปริยิฏฺโฐ โหติ. อถ จ สพฺโพ วตฺถุโต วา นิทาเนน วา โย อธิปฺปาโย พฺยญฺชโน นิรุตฺติ สนฺธิ จ อนุตฺตโร เอโส ปุพฺพาปเรน เอวํ สุตฺตตฺเถน เทสิตพฺพํ. อยํ จตุพฺยูโห หาโร.}

\proseitem{50}{ตตฺถ กตโม อาวฏฺโฏ หาโร.}

\setlength{\csromangathapagewidth}{0pt}
\setlength{\csromangathawakwidth}{0pt}
\csromangathameasuregroup{\textsuperscript{*}เอกมฺหิ ปทฏฺฐาเน, \\ อาวฏฺฏติ ปฏิปกฺเข,}{\csromangathabat{\textsuperscript{*}เอกมฺหิ ปทฏฺฐาเน,}{ปริเยสติ เสสกํ ปทฏฺฐานํ.} \\ \csromangathabat{อาวฏฺฏติ ปฏิปกฺเข,}{อาวฏฺโฏ นาม โส หาโร.}}
\csromangathasetleft{\textsuperscript{*}เอกมฺหิ ปทฏฺฐาเน, \\ อาวฏฺฏติ ปฏิปกฺเข,}
\csromangathagroup{\csromangathastanza{\csromangathaitembat{50}{\csromansymbolfootnote{*}{เหฏฺฐา 3 ปิฏฺเฐ.}เอกมฺหิ ปทฏฺฐาเน,}{ปริเยสติ เสสกํ ปทฏฺฐานํ.} \\ \csromangathacontbat{อาวฏฺฏติ ปฏิปกฺเข,}{อาวฏฺโฏ นาม โส หาโร.}}}

\prose{ยถา กิํ อุนฺนฬานํ ปมตฺตานนฺติ คาถาโย. ยํ ปมาโท, อิทํ กิสฺส ปทฏฺฐานํ, กุสลานํ ธมฺมานํ โอสคฺคสฺส. กุสลธมฺโม\-สคฺโค ปน กิสฺส ปทฏฺฐานํ, อกุสลธมฺม\-ปฏิเสวนาย. กิสฺส ปทฏฺฐานํ, กุสลธมฺม\-ปฏิเสวนาย. กิสฺส ปทฏฺฐานํ, กิเลสวตฺถุ\-ปฏิเสวนาย. อิติ ปมาเทน โมหปกฺขิยา ทิฏฺฐิ อวิชฺชา ฉนฺทราค\-ปกฺขิยา. ตตฺถ ตณฺหา จ ทิฏฺฐิ จตฺตาโร อาสวา ตณฺหา กามาสโว จ ภวาสโว จ ทิฏฺฐาสโว จ อวิชฺชาสโว จ. ตตฺถ จิตฺเต อตฺถีติ ทิฏฺฐิ เจตสิเกสุ นิจฺจนฺติ ปญฺจสุ กามคุเณสุ อชฺฌาวหเนน กามาสโว, อุปปตฺตีสุ อาสตฺติ ภวาสโว. ตตฺถ \csromanfolio{233}รูปกาโย กามาสวสฺส ภวาสวสฺส จ ปทฏฺฐานํ. นามกาโย ทิฏฺฐาสวสฺส อวิชฺชาสวสฺส จ ปทฏฺฐานํ.}

\prose{ตตฺถ อลฺลิยนาย อชฺฌตฺต\-วาหนํ กามาสวสฺส ลกฺขณํ. ปตฺถนคนฺถ- นอภิสงฺขารกายสงฺขารณํ ภวาสวสฺส ลกฺขณํ, อภินิเวโส จ ปรามาโส จ ทิฏฺฐาสวสฺส ลกฺขณํ. อปฺปฏิเวโธ ธมฺเมสุ อสมฺปชญฺญา จ อวิชฺชาสวสฺส ลกฺขณํ. อิเม จตฺตาโร อาสวา จตฺตาริ อุปาทานานิ. กามาสโว กามุปาทานํ, ภวาสโว ภวุปาทานํ, ทิฏฺฐาสโว ทิฏฺฐุปาทานํ, อวิชฺชาสโว อตฺตวาทุ\-ปาทานํ, อิเมหิ จตูหิ อุปาทาเนหิ ปญฺจกฺขนฺธา. ตตฺถ อวิชฺชาสโว จิตฺเต ปหาตพฺโพ, โส จิตฺเต จิตฺตา\-นุปสฺสิสฺส ปหียติ. ทิฏฺฐาสโว ธมฺเมสุ ปหาตพฺโพ, โส ธมฺเมสุ ธมฺมา\-นุปสฺสิสฺส ปหียติ. ภวาสโว อาสตฺติยา ปหาตพฺโพ, โส เวทนาสุ เวทนา\-นุปสฺสิสฺส ปหียติ. กามาสโว ปญฺจสุ กามคุเณสุ ปหาตพฺโพ, โส กาเย กายานุปสฺสิสฺส ปหียติ. ตตฺถ กายานุปสฺสนา ทุกฺขม\-ริยสจฺจํ ภชติ. เวทนา\-นุปสฺสนา ปญฺจนฺนํ อินฺทฺริยานํ ปจฺจโย สุขินฺทฺริยสฺส ทุกฺขิ\-นฺทฺริยสฺส โสมนสฺสิ\-นฺทฺริยสฺส โทมนสฺสิ\-นฺทฺริยสฺส อุเปกฺขิ\-นฺทฺริยสฺส, สตฺต\-กิเลโส\-ปจาโร เตน สมุทยํ ภชติ. จิตฺเต จิตฺตานุปสฺสนา นิโรธํ ภชติ. ธมฺเมสุ ธมฺมา\-นุปสฺสนา มคฺคํ ภชติ. เตนสฺส จตูสุ จ ทสฺสเนน ตสฺเสว สพฺเพ ปหียนฺติ, เยน นิทฺทิฏฺฐา ปฐมํ อุนฺนฬานํ ปมตฺตานํ เตสํ วฑฺฒนฺติ อาสวา. ชานโต โห ปสฺสโต อาสวานํ ขโย ทุกฺขํ สมุทโย นิโรโธ มคฺโค หิ อกุสลา ธมฺมา. เอวํ ปริเยสิตพฺพา. ยาว ตสฺส อกุสลสฺส คติ ตโต ปฏิปกฺเขน อกุสเล ธมฺเม ปริเยสติ เตสํ กิเลสานํ หาเรน อาวฏฺฏติ. อยํ วุจฺจเต อาวฏฺโฏ หาโร. เอวํ สุกฺกาปิ ธมฺมา ปริเยสิตพฺพา. อกุสลธมฺเม อาคมิสฺส.}

\prose{ตตฺถ อาวฏฺฏสฺส หารสฺส อยํ ภูมิ สติ อุปฏฺฐานา จ วิปลฺลาสา จ จตฺตาริ ญาณานิ สกฺกาย\-สมุปฺปาทาย\-คามินี จ ปฏิปทา สกฺกาย\-นิโรธ\-คามินี ปฏิปทา.}

\proseitem{51}{ตตฺถ กตโม วิภตฺติ หาโร. ยํ กิญฺจิ วิภชฺช\-พฺยากรณียํ วุจฺจติ วิภตฺติ หาโร. ยถา กิํ อาคนฺตฺวา จ ปุน ปุคฺคโล โหติ, โน \csromanfolio{234}วาคตํ น ปริภาสติ\csromansharedfootnote{762d81b06ff57f44}{โน วา น ปริภาสติ (อิ), น ตาวายํ ปริภาสิ (ก)} ปริปุจฺฉตาย ปญฺหาย อติยนํ เอกสฺส กิญฺจิ, อยํ วุจฺจเต วิภตฺติ หาโร.}

\prose{ตตฺถ กตโม ปริวตฺตโน หาโร. ยํ กิญฺจิ ปฏิปกฺข\-นิทฺเทโส, อยํ วุจฺจติ ปริวตฺตโน หาโร. ยถา วุตฺตํ ภควตา สมฺมา\-ทิฏฺฐิกสฺส ปุริส\-ปุคฺคลสฺส มิจฺฉาทิฏฺฐิ นิชฺชิณฺณา โหตีติ วิตฺถาเรน สพฺพานิ มคฺคงฺคานิ. อยํ วุจฺจเต ปริวตฺตโน หาโร.}

\prose{ตตฺถ กตโม เววจโน หาโร.}

\setlength{\csromangathapagewidth}{0pt}
\setlength{\csromangathawakwidth}{0pt}
\csromangathameasuregroup{}{เววจเนหิ อเนเกหิ, \\ เอกํ ธมฺมํ ปกาสิตํ. \\ สุตฺเต โย ชานาติ สุตฺตวิทู, \\ เววจโน นาม โส หาโร.}
\csromangathameasurewak{เววจเนหิ อเนเกหิ, \\ เอกํ ธมฺมํ ปกาสิตํ. \\ สุตฺเต โย ชานาติ สุตฺตวิทู, \\ เววจโน นาม โส หาโร.}
\setlength{\csromangathaleftcol}{0pt}
\csromangathagroup{\csromangathastanza{\csromangathawak{เววจเนหิ อเนเกหิ, \\ เอกํ ธมฺมํ ปกาสิตํ. \\ สุตฺเต โย ชานาติ สุตฺตวิทู, \\ เววจโน นาม โส หาโร.}}}

\prose{ยถา อายสฺมา สาริปุตฺโต เอกมฺหิ วตฺถุมฺหิ เววจเนน นานาวุตฺเตน ภควตา ปสํสิโต “มหาปญฺโญ สาริปุตฺโต หาสปญฺโญ ชวนปญฺโญ”ติ อิทํ ปญฺญาย เววจนํ. ยถา จ มคฺควิภงฺเค นิยฺยานตฺโถ เอกเมกํ มคฺคงฺคํ เววจเนหิ นิทฺทิฏฺฐํ. เอวํ อวิชฺชาย เววจนา. เอกํ อกุสลมูลํ ตเทว สนฺตํ เตสุ เตสุ ชนปเทสุ เตน เตน ปชานนฺติ. น หิ อเนน ตเทวปิ อาลปิยนฺติ อญฺญํ ภชติ. สพฺพกาม\-ชหสฺส ภิกฺขุโนติ กามา อาลปิตา. ยสฺส นิตฺถิณฺโณ สงฺโกติ เตเยว กาเม สงฺกาติ อาลปติ. สุณมานสฺส ปุเรตรํ รชฺชนฺติ เตเยว กาเม รชฺชนฺติ อาลปติ. เอวํ สุตฺตมฺหิ โย ธมฺโม เทสิยติ ตสฺส ปริเยฏฺฐิ “กตมสฺส ธมฺมสฺส อิทํ นามํ กตมสฺส อิทํ เววจนนฺ”ติ. สพฺพญฺญู หิ เยสํ เยสํ ยา นิรุตฺติ โหติ, ยถา คามิ เตน เตน เทเสตีติ ตสฺส เววจนํ ปริเยสิตพฺพํ. อยํ เววจโน หาโร.}

\proseitem{52}{ตตฺถ กตโม ปญฺญตฺติ หาโร. จตฺตาริ อริยสจฺจานีติ สุตฺตํ นิทฺทิสติ, นิกฺเขป\-ปญฺญตฺติ. ยา สมุทย\-ปญฺญตฺติ. กพฬีกาเร อาหาเร อตฺถิ ฉนฺโท อตฺถิ ราโค ยาว ปติฏฺฐิตํ. ตตฺถ วิญฺญาณํ ปภว\-ปญฺญตฺติํ ปญฺญเปติ. กพฬีกาเร อาหาเร นตฺถิ ฉนฺโท ฯเปฯ สมุคฺฆาติ ปญฺญตฺติ.}

\prose{ตสฺส กามาสวาปิ จิตฺตํ วิมุจฺจติ, ภวาสวาปิ จิตฺตํ วิมุจฺจติ, อวิชฺชาสวาปิ จิตฺตํ วิมุจฺจตีติ ปหาน\-ปญฺญตฺติํ ปญฺญเปติ. ตญฺหา ยสฺส ปุรกฺขตา \csromanfolio{235}ปญฺญา ปริวตฺตติ คาถา มนาป\-ปญฺญตฺติํ ปญฺญเปติ. เอวํ ปน มนาป\-ปญฺญตฺตีติ เอกธมฺมํ ภควา ปญฺญเปติ. น หิ ตณฺหา ทุกฺข\-สมุทโยติ กาเรตฺวา สพฺพตฺถ ตณฺหาสมุทโย นิทฺทิสิตพฺโพ. ยถา อุปฺปนฺนํ กามวิตกฺกํ นาธิวาเสติ วิโนเทติ ปชหตีติ ปฏิกฺเขป\-ปญฺญตฺติ. เอวํ สพฺเพสํ ธมฺมานํ กุสลานญฺจ อกุสลานญฺจ ยญฺจสฺส ธมฺมกฺเขตฺตํ ภวติ, โส เจว ธมฺโม ตตฺถ ปวตฺตติ. ตทวสิฏฺฐา ธมฺมา ตสฺสานุวตฺตกา โหนฺติ. สา ทุวิธา ปญฺญตฺติ ปราธีน\-ปญฺญตฺติ จ สาธีนปญฺญตฺติ จ. กตมา สาธีนปญฺญตฺติ. สมาธิํ ภิกฺเขเว\csromansharedfootnote{b4e06d5e35b790ec}{ปสฺส สํ 2. 12 ปิฏฺเฐ.} ภาเวถ, สมาหิโต ภิกฺขเว ภิกฺขุ ยถาภูตํ ปชานาติ. “รูปํ อนิจฺจนฺ”ติ ยถาภูตํ ปชานาติ, อยํ สาธีนปญฺญตฺติ ปราธีน\-ปญฺญตฺติ จ, สา ปญฺญตฺติ ปญฺญาย จ สีลสฺส จ, ยถา จตฺตาริ ฌานานิ ภาเวถ. ตสฺส อตฺถิ สมาธินฺทฺริยํ มุทูนิ จตฺตาริ อินฺทฺริยานิ ตานิ จตุปราธีนานิ, ตีณิ อเวจฺจ\-ปฺปสาเทติ ปราธีนํ สมาธินฺทฺริยํ จตฺตาริ อินฺทฺริยานิ ปราธีนาติ จตูสุ อริยสจฺเจสุ อปราธีนํ ปญฺญินฺทฺริยํ สติปฏฺฐาเนสุ สมฺม\-ปฺปธาเนสุ วีริยินฺทฺริยํ. อิติ สเก ปทฏฺฐาเน สเก เขตฺตสาธีโน โส ธมฺโม, โส จ ตตฺถ ปญฺญาเปตพฺโพ. ตสฺส ปฏิปกฺขา นิฆาโต นิทฺทิสิตพฺโพ. เอตฺถายํ อเนกาการ\-ปญฺญตฺติ เกน การเณน อยํ ธมฺโม ปญฺญตฺโตติ. อยํ วุจฺจเต ปญฺญตฺติ.}

\proseitem{53}{ตตฺถ กตโม โอตรโณ หาโร. ฉสุ ธมฺเมสุ โอตาเรตพฺพํ. กตเมสุ ฉสุ, ขนฺเธสุ ธาตูสุ อายตเนสุ อินฺทฺริเยสุ สจฺเจสุ ปฏิจฺจ\-สมุปฺปาเทสุ. นตฺถิ ตํ สุตฺตํ วา คาถา วา พฺยากรณํ วา. อิเมสุ ฉนฺนํ ธมฺมานํ อญฺญตรสฺมิํ น สนฺทิสฺสติ. เอตฺตาวตา เอส สพฺพา เทสนา ยา ตา ขนฺธา วา ธาตุโย วา อายตนานิ วา สจฺจานิ วา ปฏิจฺจ\-สมุปฺปาโท วา, ตตฺถ ปญฺจนฺนํ ขนฺธานํ เวทนากฺขนฺโธ ราค\-โทส\-โมหานํ ปทฏฺฐานํ. ตตฺถ ติสฺโส เวทนาโย ตสฺส สุขาย เวทนาย โสมนสฺโส สวิจาโร, ทุกฺขาย เวทนาย โทมนสฺโส สวิจาโร, อทุกฺขม\-สุขาย เวทนาย อุเปกฺโข สวิจาโร. ยํ ปุน ตตฺถ เวทยิตํ อิทํ ทุกฺขสจฺจํ, ขนฺเธสุ สงฺขาร\-กฺขนฺโธ ตตฺถ กาโย ปมตฺตํ สอุปวตฺตติ, ตญฺจ สงฺขารคโต ทฺวิธา จ ภวงฺโค\-ตรณํ กมฺมํ ตีณิ จ สงฺขารานิ ปุญฺญา\-ภิสงฺขารา วา อปุญฺญา วา อาเนญฺชา วา เหตุ สพฺพสราคสฺส โน วีตราคสฺส, โทสสฺส อภิสงฺขารานิ จ อวีตราโค เจเตติ จ ปกปฺเปติ จ, \csromanfolio{236}วีตราโค ปน เจเตติ จ โน อภิสงฺขโรติ, ยํ อุณฺหํ วชิรํ กฏฺเฐ วา รุกฺเข วา อญฺญตฺถ วา ปตนฺตํ ภินฺทติ จ ฑหติ จ, เอวํ สราคเจตนา เจเตติ จ อภิสงฺขโรติ จ. ยถา สตํ วชิรํ น ภินฺทติ น จ ฑหติ, เอวํ วีตราคเจตนา เจเตติ น จ อภิสงฺขโรติ. ตตฺถ ปญฺจนฺนํ ขนฺธานํ เอโก ขนฺโธ อนินฺทฺริยสรีรํ สญฺญากฺขนฺโธ.}

\prose{ตตฺถ ธาตูนํ อฏฺฐารส ธาตุโย. ตตฺถ ยา รูปี ทส ธาตุโย, ตาสุ เทสิยมานาสุ รูปกฺขนฺโธ นิทฺทิสิตพฺโพ, ทุกฺขํ อริยสจฺจํ. เยปิ จ ฉ วิญฺญาณกายา มโนธาตุ\-สตฺตมา, ตตฺถ วิญฺญาณ\-กฺขนฺโธ จ นิทฺทิสิตพฺโพ, ทุกฺขํ อริยสจฺจํ. ธมฺมธาตุ ปน ธมฺม\-สโมสรณา, โส ธมฺโม เหตุนา จ นิสฺสนฺเทน จ ผเลน จ กิจฺเจน จ เววจเนน จ เยน เยน อุปลพฺภติ, เตน เตน นิทฺทิสิตพฺโพ. ยทิ วา กุสลา ยทิ วา อกุสลา ยทิ วา อพฺยากตา ยทิ วา อสงฺขตา. ทฺวาทสนฺนํ อายตนานํ ทส อายตนานิ รูปานิ ตํ ทุกฺขํ อริยสจฺจํ นิทฺทิสิตพฺพํ. รูปกฺขนฺโธ จ มนายตนญฺจ วิญฺญาณ\-กฺขนฺเธน นิทฺทิสิตพฺพํ, ทุกฺขํ อริยสจฺจํ. ธมฺมายตนํ นานา\-ธมฺม\-สโมสรณํ. ตตฺถ เย ธมฺมา อินฺทฺริยานํ อินฺทฺริเยสุ นิทฺทิสิตพฺพา. เย อนินฺทฺริยานํ อนินฺทฺริเยสุ นิทฺทิสิตพฺพา. ปริยายโต จ โอตาเรตพฺพา. ยถา สา ธมฺมธาตุ ตถา ธมฺมายตนํ ปริเยสิตพฺพํ. ยาเยว หิ ธมฺมธาตุ ตเทว ธมฺมายตนํ อนูนํ อนธิกํ.}

\prose{ตตฺถ ปฏิจฺจ\-สมุปฺปาโท อตฺถิ ติวิโธ, อตฺถิ จตุพฺพิโธ, อตฺถิ ทุวิโธ. ตตฺถ ติวิโธ ปฏิจฺจ\-สมุปฺปาโท เหตุ\-ผล\-นิสฺสนฺโท. อวิชฺชา สงฺขารา ตณฺหา อุปาทานํ จ อยํ เหตุ, วิญฺญาณํ นามรูปํ สฬายตนํ ผสฺโส เวทนา จ อยํ ปจฺจโย, โย ภโว อยํ วิปาโก. ยา ชาติ มรณํ, อยํ นิสฺสนฺโท.}

\prose{กถํ จตุพฺพิโธ เหตุ ปจฺจโย วิปาโก นิสฺสนฺโท จ. อวิชฺชา จ ตณฺหาสงฺขารา จ อุปาทานํ จ อยํ เหตุ. วิญฺญาณํ นามรูปสฺส ปจฺจโย. นามรูปํ อุปปชฺชติ. ตถา อุปปนฺนสฺส สฬายตนํ ผสฺโส เวทานา จ อยํ ปจฺจโย. โย ภโว, อยํ วิปาโก. ยา ชาติ ยา จ ชรามรณํ, อยํ นิสฺสนฺโท.}

\prose{กถํ ทุวิโธ ปฏิจฺจ\-สมุปฺปาโท. อวิชฺชา สงฺขารา ตณฺหา อุปาทานํ, อยํ สมุทโย, วิญฺญาณํ นามรูปํ สฬายตนํ ผสฺโส เวทนา ภโว ชาติ \csromanfolio{237}มรณํ จ, อิทํ ทุกฺขํ. ยํ ปน อวิชฺชานิโรธา สงฺขาร\-นิโรโธ อิมานิ ตปฺปฏิปกฺเขน เทฺว สจฺจานิ. ตสฺมา ปฏิจฺจ\-สมุปฺปาโท เยน อากาเรน นิทฺทิฏฺโฐ, เตน เตน นิทฺทิสิตพฺโพ.}

\prose{ตถา พาวีสติ อินฺทฺริยานิ, ทฺวาทส อินฺทฺริยานิ จกฺขุนฺทฺริยานิ จกฺขุนฺทฺริยํ เยน โทมนสฺสิ\-นฺทฺริยํ, อิทํ ทุกฺขํ. ปุริสินฺทฺริยํ จ ทิฏฺฐิยา จ ตณฺหา\-ปทฏฺฐานํ. ยโต ปุริโส ปุริสกานํ ตํ เอวํ กาตพฺพตา. อถ อชฺฌตฺตํ สารชฺชติ. อยํ อหํกาโร ตํ ยสา สารตฺโต พหิทฺธา ปริเยสติ, อยํ มมํกาโร เอวํ อิตฺถี, ตตฺถ สุขินฺทฺริยํ จ โสมนสฺสิ\-นฺทฺริยํ จ ปุริสินฺทฺริยสฺสา\-นุวตฺตกา โหนฺติ. ตสฺส อธิปฺปาย\-ปริปุณฺณา โลภธมฺมา กุสลมูเล ปวฑฺเฒนฺติ. ตสฺส เจ อยมธิปฺปาโย น ปาริปูริํ คจฺฉติ. ตสฺส ทุกฺขินฺทฺริยํ จ โทมนสฺสิ\-นฺทฺริยํ จ วตฺตติ. โทโส จ อกุสลมูลํ ปวฑฺฒติ. สเจ ปน อุเปกฺขา ภาเวติ อุเปกฺขิ\-นฺทฺริยสฺส อนุวตฺตกามา ภวติ. อโมโห จ กุสลมูลํ ปวฑฺฒติ. อิติ สตฺต อินฺทฺริยานิ กิเลสวตฺถุมุ\-ปาทาย อนเนฺวมานิ อวมานิ สพฺพสฺส เวทนา อิตฺถินฺทฺริยํ ปุริสินฺทฺริยํ. ตตฺถ อฏฺฐ อินฺทฺริยานิ สทฺธินฺทฺริยํ ยาว อญฺญาตาวิโน อินฺทฺริยํ, อยํ ทุกฺขนิโรธ\-คามินี ปฏิปทา, ทสนฺนํ ปญฺญินฺทฺริยานํ กามราคสฺส ปทฏฺฐานํ. มนินฺทฺริยํ ภวราคสฺส ปทฏฺฐานํ. ปญฺญินฺทฺริยานิ รูปราคสฺส ปทฏฺฐานํ. อิตฺถินฺทฺริยํ จ ปุริสินฺทฺริยํ จ สตฺต ปญฺญตฺติยา ปทฏฺฐานํ. ตตฺถ เยน เยน อินฺทฺริเยน ยุตฺตํ วา คาถาย โอตาเรตุํ สกฺโกติ, เตน เตน นิทฺทิสิตพฺโพ. เอวํ ขนฺเธสุ ธาตูสุ อายตเนสุ สจฺเจสุ ปฏิจฺจ\-สมุปฺปาเทสุ, อยํ โอตรโณ หาโร.}

\proseitem{54}{ตตฺถ กตโม โสธโน หาโร. โย คาถา เอเกน อารมฺโภ ภาสิสฺสนฺติ. ตตฺถ เอกิสฺสา ภาสิตาย อวสิฏฺฐาสุ ภาสิตาสุ โส อตฺโถ น นิทฺทิสิตพฺโพ. กิํ การณํ, น หิ ตาว โส อตฺโถ ภาสิโต, โส อภาสิโต น สกฺกา นิทฺทิสิตุํ. ยถา กิํ อปฺปมาโท อมตํ ปทนฺติ คาถา อยเมกา คาถา นิทฺทิสิตพฺพา. กิํ การณํ, อตฺถิกฺขาตาว อิมสฺส อารมฺภสฺส อนภาสิตํ.}

\setlength{\csromangathapagewidth}{0pt}
\setlength{\csromangathawakwidth}{0pt}
\csromangathameasuregroup{เอวํ\textsuperscript{1} วิเสสโต ญตฺวา, \\ อปฺปมาเท ปโมทนฺติ,}{\csromangathabat{เอวํ\textsuperscript{1} วิเสสโต ญตฺวา,}{อปฺปมาทิมฺหิ ปณฺฑิตา.} \\ \csromangathabat{อปฺปมาเท ปโมทนฺติ,}{อริยานํ โคจเร รตาติ.}}
\csromangathasetleft{เอวํ\textsuperscript{1} วิเสสโต ญตฺวา, \\ อปฺปมาเท ปโมทนฺติ,}
\csromangathagroup{\csromangathastanza{\csromangathabat{เอวํ\csromansharedfootnote{ac0a872b8c44730c}{เอตํ (อิ) ขุ 1. 16 ธมฺมปเท.} วิเสสโต ญตฺวา,}{อปฺปมาทิมฺหิ ปณฺฑิตา.} \\ \csromangathabat{อปฺปมาเท ปโมทนฺติ,}{อริยานํ โคจเร รตาติ.}}}

\prose{\csromanfolio{238}อิทํ อภาสิตํ. อิมิสฺสาปิ คาถาย ภาสิตาย อตฺโถ นิทฺทิสิตพฺโพ. กิํ การณํ, อตฺถิ ตตฺถ อวสิฏฺฐํ. เต ฌายิโน\csromansharedfootnote{8f6bced035981d8f}{ขุ 1. 16 ธมฺมปเท.} สาตติกา, นิจฺจํ ทฬฺหปรกฺกมาติ คาถา, เอวํ อิมา คาถาโย อุปธาริตา ยทา ภวนฺติ, ตทา อตฺโถ นิทฺทิสิตพฺโพ. เอวํ อสฺสุตปุพฺเพสุ สุตฺเตสุ พฺยากรเณสุ วา เอกุทฺเทโส ภาสิโต. ยา วีมํสา ตุลนา อิทํ อตฺถิ กิจฺจํ, อิทํ สุตฺตํ ภาสิตํ ตสฺส เววจนํ นิทฺทิฏฺฐํ วา น วาติ. ตตฺถ ยา วีมํสา, อยํ วุจฺจเต โสธโน หาโร.}

\proseitem{55}{ตตฺถ กตโม อธิฏฺฐาโน หาโร. เอกตฺตตา จ เวมตฺตตา จ. ตตฺถ กิตปญฺญตฺติ จ กิจฺจปญฺญตฺติ จ. สา เอกตฺตตา จ เวมตฺตตา จ ยถา ปญฺญตฺติ เอกเววจเนน เวมตฺตตา ปชานาตีติ ปญฺญา, สา จ อาธิปเตยฺย\-ฏฺเฐน ปญฺญตฺติ. ยํ อโนมตฺติยฏฺเฐน ปญฺญตฺตนฺติ. ตํ อโนมตฺติยฏฺเฐน ปญฺญาพลํ. ตนุภูตา โคจรตฺตวสา เสวสติ ตีสุ รตเนสุ อนุสฺสติ พุทฺธานุสฺสติ ธมฺมานุสฺสติ สํฆานุสฺสติ อวิปรีตา\-นุสฺสรณ\-ตาย. สมฺมาทิฏฺฐิ ธมฺมานํ ปวิจเยน ธมฺมวิจย\-สมฺโพชฺฌงฺโค อภินีหารโต อภิญฺญาติ. สงฺเขเปน มคฺคา กา วตฺถุ อวิโกปนตาย เอกตฺตา, ยถา อุเณฺหน สํสฏฺฐํ อุโณฺหทกํ, สีเตน สํสฏฺฐํ สีโตทกํ ขาโรทกํ คุโฬฺหทกนฺติ, อิทํ เอกตฺตตา เวมตฺตตา จ.}

\prose{อตฺถิ ปุน ธมฺโม นานาธมฺมสํฆโต เอกโต ยถารูปํ จตฺตาโร วาเรตพฺพา, ตญฺจ รูปนฺติ เอกตฺตตา. ปถวีธาตุ อาโป เตโช วาโยธาตูติ เวมตฺตตา. เอวํ สพฺพา จตสฺโส ธาตุโย รูปนฺติ เอกตฺตตา, ปถวีธาตุ อาโป เตโช วาโยธาตูติ เวมตฺตตา. ปถวีธาตูติ ลกฺขณโต เอกตฺตตา, สํกิณฺณ\-วตฺถุโต เวมตฺตตา. ยํ กิญฺจิ กกฺขฬ\-ลกฺขณํ, สพฺพํ ตํ ปถวีธาตูติ เอกตฺตตา. เกสา โลมา นขา ทนฺตา ฉวิ จมฺมนฺติ เวมตฺตตา. เอวํ สพฺพํ จตสฺโส ธาตุโย รูปนฺติ เอกตฺตํ. สทฺทา คนฺธา รสา โผฏฺฐาพฺพาติ เวมตฺตตา.}

\prose{อตฺถิ ปุน ธมฺโม เวมตฺตตา อญฺโญ นามํ ลภติ. ยถา กายา\-นุปสฺสนาย นวสญฺญา วินีลกสญฺญา อุทฺธุมาตก\-สญฺญา, อยํ อสุภสญฺญา, ยา เอกตฺตตา อารมฺมณโต เวมตฺตโต, สา เอวํ สญฺญาเวทนาสุ อาทีนวํ สมนุปสฺสโต ตถา\-ธิฏฺฐานํ สมาธินฺทฺริยํ จ \csromanfolio{239}สาเยว ธมฺเมสุ ตตฺถ สญฺญาภาวนา วีริยินฺทฺริยํ จ ธมฺเมสุ ธมฺมา\-นุปสฺสนา จิตฺเต อตฺตสญฺญํ ปชหโต ปญฺญินฺทฺริยํ จ จิตฺเต จิตฺตานุปสฺสนา. (อิติ)\csromansharedfootnote{3dd9dfe504bbaa2e}{( ) นตฺถิ อิโปตฺถเก.} โย โกจิ ญาณปจาโร สพฺพโส ปญฺญาย โคจโร ปญฺญา, อยํ เวมตฺตตา, ยถา กามราโค ภวราโค ทิฏฺฐิราโคติ เวมตฺตตา ตณฺหาย. อิติ ยํ เอกตฺตตาย จ เวมตฺตตาย จ ญาณํ วีมํสนา ตุลนา. อยํ อธิฏฺฐาโน หาโร.}

\proseitem{56}{ตตฺถ กตโม ปริกฺขาโร หาโร. สเหตุ สปฺปจฺจยํ โวทานญฺจ สํกิเลโส จ, ยํ ตทุภยํ ปริเยฏฺฐิ, สปริกฺขาโร หาโร. อิติ ธมฺมานํ สเหตุกานํ เหตุ ปริเยสิตพฺโพ, สปฺปจฺจยานํ ปจฺจโย ปริเยสิตพฺโพ.}

\prose{ตตฺถ กิํ นานากรณํ, เหตุสฺส จ ปจฺจยสฺส จ. สภาโว เหตุ, ปรภาโว ปจฺจโย. ปรภาวสฺส ปจฺจโย เหตุปิ, สภาวสฺส เหตุยา ปรภาวสฺส กสฺสจิ ปจฺจโย อวุตฺโต เหตุ, วุตฺโต ปจฺจโย. อชฺฌตฺติโก เหตุ, พาหิโร ปจฺจโย. สภาโว เหตุ, ปรภาโว ปจฺจโย. นิพฺพตฺตโก เหตุ, ปฏิคฺคาหโก\csromansharedfootnote{1768d6be7768af00}{ปริคฺคาหโก (ก)} ปจฺจโย. เนวาสิโก เหตุ, อาคนฺตุโก ปจฺจโย. อสาธารโณ เหตุ, สาธารโณ ปจฺจโย. เอโกเยว เหตุ, อปราปโร ปจฺจโย.}

\prose{เหตุสฺส อุปกรณํ สมุทาเนตพฺโพ. สมุทานํ เหตุ, ตตฺถ ทุวิโธ เหตุ, ทุวิโธ ปจฺจโย สมนนฺตร\-ปจฺจโย จ ปรมฺปร\-ปจฺจโย จ. เหตุปิ ทุวิโธ สมนนฺตร\-เหตุ จ ปรมฺปรเหตุ จ. ตตฺถ กตโม ปรมฺปร\-ปจฺจโย, อวิชฺชา นามรูปสฺส ปรมฺปร\-ปจฺจโย, วิญฺญาณํ สมนนฺตร\-ปจฺจยตาย ปจฺจโย. ยทิ อาทิมฺหิ อวิชฺชานิโรโธ ภวติ นามรูปสฺส นิโรโธปิ. ตตฺถ สมนนฺตรํ กิํ การณํ ปรมฺปร\-ปจฺจโย สมนนฺตร\-ปจฺจโย สมุทฺทานิโต, อยํ ปจฺจยโต. ตตฺถ กตโม ปรมฺปรเหตุ, วิชานนฺตสฺส ปรมฺปร\-เหตุตาย เหตุ, อญฺญากาโร สมนนฺตร\-เหตุตาย เหตุ. ยสฺส หิ ยํ สมนนฺตรํ นิพฺพตฺตติ, โส ตสฺส เหตุปิ ชาตินิโรธา พหิ อาการนิโรโธ, อาการนิโรธา ทณฺฑนิโรโธ, ทณฺฑนิโรธา ขณฺฑนิโรโธ. เอวํ เหตุปิ ทฺวิธา โส ตาหิ ปสฺสิตพฺโพ.}

\prose{\csromanfolio{240}ปฏิจฺจ\-สมุปฺปาโท ยถา อวิชฺชาปจฺจโย ตสฺส ปุน กิํ ปจฺจโย, อโยนิโส มนสิกาโร. โส กสฺส ปจฺจโย สงฺขารานํ, อิติ ปจฺจโย จ สมุปฺปนฺนํ จ ตสฺส โก เหตุ อวิชฺชาเยว, ตถา หิ ปุริมา โกฏิ น ปญฺญายติ. ตตฺถ อวิชฺชานุสโย อวิชฺชา\-ปริยุฏฺฐานสฺส เหตุ ปุริมา เหตุ ปจฺฉา ปจฺจโย, สาปิ อวิชฺชา\-สงฺขารานํ ปจฺจโย จตูหิ การเณหิ สหชาต\-ปจฺจยตาย สมนนฺตร\-ปจฺจยตาย อภิสนฺทน\-ปจฺจยตาย ปติฏฺฐาน\-ปจฺจยตาย.}

\proseitem{57}{กถํ สหชาต\-ปจฺจยตาย อวิชฺชา\-สงฺขารานํ ปจฺจโย, ยํ จิตฺตํ ราค\-ปริยุฏฺฐํ, ตตฺถ อวิชฺชา\-ปริยุฏฺฐาเนน สพฺพํ ปญฺญาย โคจรํ หนฺติ. ตตฺถ สงฺขารา ติปจฺจยฏฺฐิกา อทฺธา\-ภูมิ\-การ\-มหตฺตสฺส\csromansharedfootnote{52b887983e7749ff}{ลทฺธา ภูมิกรมหตฺตสฺส (อิ, ก)} อยํ อวิชฺชา\-สห\-สมุปฺปนฺนํ วุทฺธิํ วิรูฬฺหิํ เวปุลฺลตมา\-ปชฺชนฺตี จตูหิ การเณหิ ปญฺญา ปหียติ. กตเมหิ จตูหิ อนุสโย ปริยุฏฺฐานํ สํโยชนํ อุปาทานํ. ตตฺถ อนุสโย ปริยุฏฺฐานํ ชาติ ปริยุฏฺฐิตา สํยุชฺชติ สํยุตฺตา อุปาทิยติ อุปาทานปจฺจยา ภโว. เอวํ เต สงฺขารา ติวิธา อุปฺปนฺนา ภูมิคตา นาส’ญฺญตฺถ อยํ มคฺเคน วินีตตฺตายาติ\csromansharedfootnote{f4d8c17a13e7a5f7}{วินิภตฺตาย (อิ), วินิภตฺตตาย (ก)} เต ถามคตา อปติวินีตาติปิ เต สงฺขาราติ วุจฺจติ, เอวํ สเหตุสมุปฺปนฺน\-ฏฺเฐน อตฺถิ เมว ปจฺจยา สงฺขารานํ ปจฺจโย นิทฺทิฏฺฐํ อปเนตฺวา กุสลํ อกุสลํ กุสโล จ อกุสโล จ ปกฺขิปิตพฺโพ, วิปากธมฺมา อปเนตฺวา วจนียํ อวจนียํ วจนียญฺจ อวจนียญฺจ ปกฺขิปิตพฺพํ, ภวอเปวิริตฺตา, สพฺพสุตฺตํ ปริกฺขิปิตพฺพํ.}

\prose{ทส ตถาคต\-พลานิ จตฺตาริ เวสารชฺชานิ ปุญฺญานิ อนญฺญากตํ อวิชฺชา สมนนฺตร\-ปจฺจยตาย สงฺขารานํ ปจฺจโย เยน จิตฺเตน สห สมุปฺปนฺนา อวิชฺชา ตสฺส จิตฺตสฺส สมนนฺตร\-จิตฺตํ สมุปฺปนฺนนฺติ, ตสฺส ยํ สมนนฺตร\-จิตฺตํ สมุปฺปนฺนนฺติ, ตสฺส ปจฺฉิมสฺส จิตฺตสฺส ปุริมจิตฺตํ เหตุ\-ปจฺจยตาย ปจฺจโย, เตน อวิชฺชา เหตุ เตน จิตฺเตน อุปาทานํ อโนกาสกตา ญาณํ น อุปฺปชฺชนฺติ. ยา ตสฺส อปฺปมาทา ธาตุ อภิชฺฌาภิ\-สนฺทิตา ตหิํ วิปลฺลาสา อุปฺปชฺชนฺติ “อสุเภ สุภนฺ”ติ “ธุกฺเข สุขนฺ”ติ, ตตฺถ สงฺขารา อุปฺปชฺชนฺติ รตฺตา ทุฏฺฐา มูลสฺส เจตนา ราค\-ปริยุฏฺฐาเนน พฺยาปาท\-ปริยุฏฺฐาเนน อวิชฺชา\-ปริยุฏฺฐาเนน ทิฏฺฐิ\-วิปลฺลาโส วตฺถุนิทฺเทเส \csromanfolio{241}นิทฺทิสิตพฺโพ, ยํ วิปรีตจิตฺโต วิชานาติ อยํ จิตฺตวิปลฺลาโส, ยา วิปรีตสญฺญา อุปคฺคณฺหาติ อยํ สญฺญาวิปลฺลาโส. ยํ วิปรีตทิฏฺฐิ อภินิวิสติ อยํ ทิฏฺฐิ\-วิปลฺลาโส. อฏฺฐ มิจฺฉตฺตานิ วฑฺฒนฺติ, ตีณิ อกุสลานิ อโยนิโส มนสิกาเร อุปฺปนฺนํ วิญฺญาณญฺจ วิชฺชญฺจ กโรนฺติ. อิติ ปุพฺพาปรนฺเต อกุสลา\-นาตริ\-ตโร สงฺขารา วุทฺธิํ เวปุลฺลตํ คจฺฉนฺติ. เต จ มหตา จ อปฺปฏิวิทิตา โปโนภวิกา\csromansharedfootnote{cb07fa3d6c2ab9a9}{โปโนพฺภวิกา (ก)} สงฺขารา ภวนฺติ. อิติ เอวํ อวิชฺชา สหชาต\-ปจฺจยตาย สงฺขารานํ ปจฺจโย สมนนฺตร\-ปจฺจยตาย จ.}

\proseitem{58}{กถํ อภิสนฺทนา\-กาเรน อวิชฺชา สงฺขารานํ ปจฺจโย, สา อวิชฺชา เต สงฺขาเร อภิสนฺเนติ ปริปฺผรติ. เสยฺยถาปิ นาม อุปฺปลํ วา ปทุมํ วา ตํ อุทเก วฑฺฒํ อสฺส, สีเตน วารินา อภิสนฺนํ ปริสนฺทนํ วุทฺธิํ วิรูฬฺหิํ เวปุลฺลตํ อาปชฺชติ. เอวํ อภิสนฺท\-นฏฺเฐน อวิชฺชา สงฺขารานํ ปจฺจโย.}

\prose{กถํ ปติฏฺฐห\-นฏฺเฐน อวิชฺชา สงฺขารานํ ปจฺจโย, เต สงฺขารา อวิชฺชายํ นิสฺสาย วุทฺธิํ วิรูฬฺหิํ เวปุลฺลตํ อาปชฺชนฺติ. เสยฺยถาปิ นาม อุปฺปลํ วา ปทุมํ วา ปถวิํ นิสฺสาย ปถวิํ ปติฏฺฐาย วุทฺธิํ วิรูฬฺหิํ เวปุลฺลตํ อาปชฺชติ. เอเต สงฺขารา อวิชฺชายํ ปติฏฺฐิตา อวิชฺชายํ นิสฺสาย วุทฺธิํ วิรูฬฺหิํ เวปุลฺลตํ คจฺฉนฺติ. เอวํ ปติฏฺฐห\-นฏฺเฐน อวิชฺชา สงฺขารานํ ปจฺจโย.}

\prose{ปุน ราค\-สหคตสฺส กมฺมสฺส วิปาเกน ปฏิสนฺธิมฺหิ ภโว นิพฺพตฺตติ, ตํ กมฺมสฺส\csromansharedfootnote{3b90f5942fb93bda}{กามสฺส (อิ)} สพฺพํ อภินิวิฏฺฐํ อญฺญาณวเสน โปโนภวิกา สงฺขาราติ วุจฺจนฺติ, เอวมฺปิ อวิชฺชาปจฺจยา สงฺขารา อตฺถิ. ปุน ปญฺจสุ เย จ เสกฺขา ปุคฺคลา, เย จ อสญฺญิ\-สมาปตฺติํ สมาปนฺนา, เย จ ภวคตา, เย จ อนฺโตคตาเยว สํเสทชา, เย จ วา ปน อญฺโญ หิ โกจิ อนาคามิภูตา น เจเตนฺติ น จ ปตฺเถนฺติ, เตสํ กิํ ปจฺจยา สงฺขารา. ปุน ราคา อตฺถิ เตสํ สงฺขารานิ อุปาทานานิ จิตฺตม\-นุสฺสรนฺติเยว อวิปกฺกวิปาก\-สมูหตา อสมุจฺฉินฺน\-ปจฺจยา เตสํ ปุน จ คโต ภวติ. เอวมฺปิ หิ อวิชฺชาปจฺจยา สงฺขารา. ปุน สา เต น อุปาทานา นปิ สงฺขารา อตฺถิ, ปุน เตสํ สตฺต อนุสยา อสมูหตา อสมุจฺฉินฺนา ตทารมฺมณํ ภวติ. วิญฺญาณสฺส ปติฏฺฐาย วิญฺญาณปจฺจยา นามรูปํ. เอวมฺปิ อวิชฺชาปจฺจยา สงฺขารา. ปุน สา ยํ กิญฺจิ กมฺมํ อาจยคามิ. สพฺพํ ตํ อวิชฺชาวเสน \csromanfolio{242}อภิสงฺขริยติ ตณฺหาวเสน จ อลฺลียติ. อญฺญณวเสน จ ตตฺถ อาทีนวมฺปิ น ชานาติ. ตเทว วิญฺญาณพีชํ ภวติ, สาเยว ตณฺหาสิเนโห ภวติ. สาเยว อวิชฺชา สมฺโมโหติ. เอวมฺปิ อวิชฺชาปจฺจยา สงฺขารา วตฺตพฺพา. อิติ อิเมหิ อากาเรหิ อวิชฺชา สงฺขารานํ ปจฺจโย.}

\prose{ตตฺถ อวิชฺชาย เหตุ อโยนิโส มนสิกาโร ปจฺจโย โหติ. ตตฺถ อภิจฺเฉโท อยํ ตตฺถ ตติยํ พลํ\csromansharedfootnote{e1b2d60a7bacdb50}{ผลํ (อิ)} นิวตฺติ, อยํ ปฏิสนฺธิ. ตตฺถ ปุนพฺภโว โย อเวจฺเฉโท อสมุ\-คฺฆาต\-นฏฺเฐน อยํ อนุสโย. ยถา ปฏากํ วา สาฏกํ วา เทฺว ชนา ปีเฬสุ จ เอกา วา พลํ วา อสฺส นิวาฏสฺเสสุ, น ปน ปีเฬสุ โสเสยฺย. ตตฺถ ยํ สิเนหา อาโปธาตุ อนุปุลฺลนา โสเสตพฺพา. อุณฺห\-ธาตุมา\-คมฺม สเจ ปุน ตํ อากาเส นิกฺขิเปยฺย ตํ อุสฺสาเวน เยภุยฺยตรํ สิเนหมา\-ปชฺเชยฺย, น หิ อนาคมฺม เตโชธาตุํ ปริเสสํ คจฺเฉยฺย. เอวเมว ภวคฺค\-ปรมาปิ สมาปตฺติ น อนุรูปสฺส สมุคฺคาตาย สํวตฺตติ. เต หิ อาลยนฺติ สมฺมสนฺติ, น จ ตณฺหาย ตณฺหาปหานํ คจฺฉนฺติ. ตตฺถ โส อสมุคฺฆาโต. อวิชฺชาย อนุสโย จ จิตฺตสฺส สมฺปลิโพโธ, อิทํ ปริยุฏฺฐานํ. ยถาภูตํ วิญฺญาณสฺส อปฺปฏิเวโธ อยํ อวิชฺชาอาสโว อวิชฺชา\-วิญฺญาณพีชํ ภวติ. ยํ พีชํ โส เหตุ น สมุจฺฉิชฺชติ, อสมุ\-จฺฉิชฺช\-นฺโต ปฏิสนฺเทหติ. ปฏิสนฺทหนฺโต น สมุคฺฆาตํ คจฺฉติ. อสมุคฺฆาตํ จิตฺตํ ปริโยนหติ, ปริโยนทฺธ\-จิตฺโต ยถาภูตํ นปฺปชานาติ, อิติ สญฺญาณสฺส สาสวตฺโถ, อวิชฺชตฺโถ, เหตุอตฺโถ, อวจฺเฉทตฺโถ, อนิวตฺติอตฺโถ, ผลตฺโถ. ปฏิสนฺธิ- อตฺโถ, ปุนพฺภวตฺโถ, อสมุ\-คฺฆาต\-ตฺโถ, อนุสยตฺโถ, ปริยุฏฺฐา\-นตฺโถ, อปฏิเวธนตฺโถ. เอตฺตาวตา อวิชฺชาย เขตฺตํ นิทฺทิฏฺฐํ ภวติ. อยํ วุจฺจเต ปริกฺขาโร นาม หาโร.}

\proseitem{59}{ตตฺถ กตโม สมาโรปโน หาโร. อุคฺฆฏิตมฺหิ ตมฺหิ สนฺตญฺเจว จ นํ วิตฺถารํ ปน วตฺตพฺพํ. วิตฺถารวิธํ จิตฺตญฺญา อยํ สมาโรปโน หาโร.  ตตฺถ นามนิทฺเทโส อุปฆฏกา\csromansharedfootnote{566613768b29ab07}{อุคฺฆฏกา (อิ)} วตฺถุนิทฺเทโส เววจนํ วตฺถุภูโต วิตฺถาโร. ยถา กิํ, ยา ภิกฺขูนํ วตฺตโต\csromansharedfootnote{4bedc620d006abf1}{นิวตฺตโต (อิ)} ปหาตพฺโพ, อยํ อุปฆฏนา.}

\prose{\csromanfolio{243}ตตฺถ กตโม สมาโรปโน, กิญฺจิ น วตฺตพฺพํ, รูปราคํ วา นาม\-วนฺต\-ปหาตพฺพํ\csromansharedfootnote{6aca5287c6d8429e}{นามมนฺติปหาตพฺพํ (ก)}. ยาว วิญฺญาณนฺติ วิตฺถาเรน กาตพฺพานิ. อวิชฺชา ตา โอปมฺเมน ปญฺญาเปตพฺพา, อยํ สมาโรปโน. นิสฺสิต\-จิตฺตสฺส จ มตฺติโก จ นิสฺสโย ตณฺหา จ ทิฏฺฐิ จ. ตตฺถ ทิฏฺฐิ อวิชฺชา ตณฺหา สงฺขารา. ตตฺถ ทิฏฺฐิปจฺจยา ตณฺหา อิเม อวิชฺชาปจฺจยา สงฺขารา. ตตฺถ นิสฺสิตํ วิญฺญาณํ อิทํ สงฺขาร\-ปจฺจยา วิญฺญาณํ ยาว ชรามรณํ, อิทํ สํขิตฺเตน ภาสิเต อวสิฏฺฐํ ปโรปยติ.}

\prose{อนิสฺสิตสฺส\csromansharedfootnote{566826cb96cf50e1}{ปสฺส ขุ 1. 179 อุทาเน; สํ 2. 284 ปิฏฺเฐสุปิ.} จลิตํ นตฺถีติ ตสฺส เอวํ ทิฏฺฐิยา ตณฺหาย จ ปหานํ ตตฺถ ทิฏฺฐิ\-อวิชฺชานิโรธาย ภูตํ วิญฺญาณํ สราคฏฺฐานิเยสุ ธมฺเมสุ ตํ ตํ ธมฺมํ อุเปจฺจ อญฺญํ ธมฺมํ ธาวติ มกฺกโฏ\-ปม\-ตาย, อถ ขฺวสฺส ปริตฺเตสุ ธมฺเมสุ สรคฏฺฐานิเยสุ ฉนฺทราโค นตฺถิ กุโต ตโต จลนา, อธิมตฺเตสุ สตฺเตสุ จิตฺตํ นิเวสฺสยติ ตํ อปติฏฺฐิตํ วิญฺญาณํ อนาหารํ นิรุชฺฌติ วิญฺญาณนิโรธา นามรูป\-นิโรโธ ยาว ชรามรณ\-นิโรโธ. อยํ สมาโรปโน.}

\prose{ตตฺถ ราควเสน วิญฺญาณสฺส จลิตํ สปริคฺคโห, ตสฺมิํ จลิเต อสติ โย ปริกิเลโส\-ปจาโร ติวิโธ อคฺคิ ปฏิปฺปสฺสทฺโธ ภวติ. เตนาห จลิเต อสนฺเต ปสฺสทฺธิ โหติ. ตตฺถ ยํ สมาโรปนา ปสฺสทฺธกาโย สุขํ เวเทติ, สุขิโน จิตฺตํ สมาธิยติ. ยาว วิมุตฺติตมิติ ญาณทสฺสนํ ภวติ. โส อาสวานํ ขยา จ วิมุตฺติ โน อุปปชฺชติ. ตสฺส อุปปตฺติสฺส อาคติคติยา อสนฺติยา เนวิธ น หุรํ น อุภยม\-นฺตเรน. เอเสวนฺโต ธุกฺขสฺสาติ อนุปาทิเสสา นิพฺพานธาตุ. อิทมสฺส สุตฺตสฺส มชฺเฌ สมาโรปิตํ ปฏิจฺจ\-สมุปฺปาเท จ วิมุตฺติยํ จ โยโค น จ เอตํ ตสฺส สํขิตฺเตน ภาสิตสฺส วิตฺถาเรน อตฺถํ วิภชฺชนฺติ. อยํ วุจฺจเต สมาโรปโน หาโร. น จ สํกิเลส\-ภาคิเยน สุตฺเตน สํกิเลส\-ภาคิโย เย จ ธมฺมา สมา\-โรป\-ยิตพฺพา นาญฺเญ. เอวํ วาสนาภาคิเย นิพฺเพธ\-ภาคิเย, อยํ สมาโรปโน หาโร.  อิเม โสฬส หารา.  สุวีรสฺส มหา\-กจฺจายนสฺส ชมฺพุ\-วนวาสิโน เปฏโกปเทเส ปญฺจมา ภูมิ.}

\csromanmatikaanchor{mtk.45}
\csromantocmarkref{section}{6. สุตฺตตฺถสมุจฺจยภูมิ}{mtk.45}
\bookmark[dest={mtk.45},level=1]{6. สุตฺตตฺถสมุจฺจยภูมิ}
\csromanmarkh{เปฏโกปเทส\-ปาฬิ}\csromanmarkhii{6. สุตฺตตฺถ\-สมุจฺจยภูมิ}\csromanheadernomark{1}{เปฏโกปเทส\-ปาฬิ \textbf{6. สุตฺตตฺถ}\-\textbf{สมุจฺจยภูมิ}}
\proseitem{60}{\csromanfolio{244}พุทฺธานํ ภควนฺตานํ สาสนํ ติวิเธน สงฺคหํ คจฺฉติ, ขนฺเธสุ ธาตูสุ อายตเนสุ จ. ตตฺถ ปญฺจกฺขนฺธา รูปกฺขนฺโธ ยาว วิญฺญาณ\-กฺขนฺโธ. ทส รูปอายตนานิ จกฺขู รูปา จ ยาว กาโย โผฏฺฐพฺพา จ, อยํ รูปกฺขนฺโธ. ตตฺถ ฉ เวทนากายา เวทนากฺขนฺโธ จกฺขุ\-สมฺผสฺสชา เวทนา ยาว มโน\-สมฺผสฺสชา เวทนา, อยํ เวทนากฺขนฺโธ. ตตฺถ ฉ สญฺญากายา สญฺญากฺขนฺโธ, รูปสญฺญา ยาว ธมฺมสญฺญา อิเม ฉ สญฺญากายา, อยํ สญฺญากฺขนฺโธ. ตตฺถ ฉ เจตนากายา สงฺขาร\-กฺขนฺโธ, รูปสญฺเจตนา ยาว ธมฺม\-สญฺเจตนา อิเม ฉ เจตนากายา, อยํ สงฺขาร\-กฺขนฺโธ. ตตฺถ ฉ วิญฺญาณกายา วิญฺญาณ\-กฺขนฺโธ, จกฺขุวิญฺญาณํ ยาว มโนวิญฺญาณํ อิเม ฉ วิญฺญาณกายา, อยํ วิญฺญาณ\-กฺขนฺโธ. อิเม ปญฺจกฺขนฺธา.}

\prose{เตสํ กา ปริญฺญา, อนิจฺจํ ทุกฺขํ สญฺญา อนตฺตาติ เอสา เอเตสํ ปริญฺญา. ตตฺถ กตโม ขนฺธตฺโถ, สมูหตฺโถ ขนฺธตฺโถ, ปุญฺชตฺโถ ขนฺธตฺโถ, ราสตฺโถ ขนฺธตฺโถ. ตํ ยถา ทพฺพกฺขนฺโธ วนกฺขนฺโธ ทารุกฺขนฺโธ อคฺคิกฺขนฺโธ อุทกกฺขนฺโธ วายุกฺขนฺโธ อิติ เอวํ ขนฺเธสุ สพฺพสงฺคโหว เอวํ ขนฺธตฺโถ.}

\prose{ตตฺถ อฏฺฐารส ธาตุโย จกฺขุธาตุ รูปธาตุ จกฺขุวิญฺญาณ\-ธาตุ ฯเปฯ มโนธาตุ ธมฺมธาตุ มโนวิญฺญาณ\-ธาตุ. เอตาโย อฏฺฐารส ธาตุโย. ตาสํ ปริญฺญา อนิจฺจํ ทุกฺขํ สญฺญา อนตฺตาติ เอสา เอตาสํ ปริญฺญา. ตตฺถ โก ธาตุอตฺโถ, วุจฺจเต อวยวตฺโถ ธาตุอตฺโถ. อวยโวติ จกฺขุ โน ปสาโท จกฺขุธาตุ. เอวํ ปญฺจสุ ธาตูสุ ปุน ราคววจฺเฉทตฺโถ ธาตุอตฺโถ. ววจฺฉินฺนา หิ จกฺขุธาตุ. เอวํ ปญฺจสุ ปุนราห เอกนฺติปกตฺยตฺเถน ธาตุอตฺโถติ วุจฺจเต. ตํ ยถา, ปกติยา อยํ ปุริโส ปิตฺติโก เสมฺหิโก วาติโก สนฺนิปาติโกติ เอวํ ปกติ\-จกฺขุธาตุ ทสนฺนํ ปิยา จ สพฺเพสุ อินฺทฺริเยสุ ฯเปฯ วิสภาคตฺโถ ธาตุอตฺโถ.}

\prose{ตตฺถ ทฺวาทสา\-ยตนานิ กตมานิ, ฉ อชฺฌตฺติกานิ ฉ พาหิรานิ. จกฺขายตนํ ยาว มนายตนนฺติ อชฺฌตฺติกํ, รูปายตนํ ยาว ธมฺมา\-ยตนนฺติ พาหิรํ. เอตานิ ทฺวาทส อายตนานิ. เอเตสํ กา ปริญฺญา, อนิจฺจํ ทุกฺขํ สญฺญา อนตฺตาติ, เอสา เอเตสํ ปริญฺญา. \csromanfolio{245}อปิ จ ทฺวิธา ปริญฺญา ญาตปริญฺญา จ ปหานปริญฺญา จ. ตตฺถ ญาตปริญฺญา นาม อนิจฺจํ ทุกฺขํ สญฺญา อนตฺตาติ, เอสา ญาตปริญฺญา. ปหานปริญฺญา ปน ฉนฺท\-ราค\-ปฺปหานา, เอสา ปหานปริญฺญา. ตตฺถ กตโม อายตนตฺโถ, วุจฺจเต อาการตฺโถ อายตนตฺโถ. ยถา สุวณฺณากโร ทุพฺพณฺณากโร, ยถา ทฺวีหิ เตหิ อากาเรหิ เต เต คาวา อุตฺติฏฺฐนฺติ. เอวํ เอเต หิ จิตฺตเจตสิกา คาวา อุตฺติฏฺฐนฺติ กมฺมกิเลสา ทุกฺขธมฺมา จ. ปุนราห อายทานตฺโถ อายตนตฺโถ. ยถา รญฺโญ อายทาเนหิ อาโย ภวติ. เอวํ อายทานตฺโถ อายตนตฺโถ.}

\proseitem{61}{จตฺตาริ อริยสจฺจานิ ทุกฺขํ สมุทโย นิโรโธ มคฺโค จ. ทุกฺขํ ยถา สมาเสน ธมฺมาจริยํ มานสญฺจ, สมุทโย สมาเสน อวิชฺชา จ ตณฺหา จ, นิโรโธ สมาเสน วิชฺชา จ วิมุตฺติ จ, มคฺโค สมาเสน สมโถ จ วิปสฺสนา จ.}

\prose{ตตฺถ สตฺตติํส โพธิปกฺขิกา ธมฺมา. กตเม จตฺตาโร สติปฏฺฐานา ยาว อริโย อฏฺฐงฺคิโก มคฺโค, เอวเมเต สตฺตติํส โพธิปกฺขิกา ธมฺมา. เย ธมฺมา อตีตา\-นาคต\-ปจฺจุปฺปนฺนานํ พุทฺธานํ ภควนฺตานํ ปจฺเจก\-พุทฺธานํ สาวกานํ จ นิพฺพานาย สํวตฺตนฺตีติ, โส มคฺโค จตฺตาโร สติปฏฺฐานา. กตเม จตฺตาโร, อิธ ภิกฺขุ กาเย กายานุปสฺสี วิหรติ, สมฺมปฺปธานํ. อิทฺธิปาทํ. อินฺทฺริยานิ. พลานิ. ตตฺถ โก อินฺทฺริยตฺโถ, อินฺทตฺโถ อินฺทฺริยตฺโถ, อาธิปเตยฺย\-ตฺโถ อินฺทฺริยตฺโถ, ปสาทตฺโถ อินฺทฺริยตฺโถ, อสาธารณํ กสฺส กิริยตฺโถ อินฺทฺริยตฺโถ. อนวปริยตฺโถ พลตฺโถ, ถามตฺโถ พลตฺโถ, อุปาทายตฺโถ พลตฺโถ, อุปตฺถมฺพนตฺโถ พลตฺโถ.}

\prose{ตตฺถ กตเม สตฺต โพชฺฌงฺคา, สติสมฺโพชฺฌงฺโค ยาว อุเปกฺขา\-สมฺโพชฺฌงฺโค. ตตฺถ กตโม อฏฺฐงฺคิโก มคฺโค, สมฺมาทิฏฺฐิ ยาว สมฺมาสมาธิ. ตตฺถ อฏฺฐงฺคิโก มคฺโคติ ขนฺโธ สีลกฺขนฺโธ จ สมาธิ\-กฺขนฺโธ จ ปญฺญากฺขนฺโธ จ. ตตฺถ ยา จ สมฺมาวาจา โย จ สมฺมากมฺมนฺโต โย จ สมฺมาอาชีโว, อยํ สีลกฺขนฺโธ. ยา จ สมฺมาสติ โย จ สมฺมาวายาโม โย จ สมฺมาสมาธิ, อยํ สมาธิ\-กฺขนฺโธ. โย จ สมฺมาสงฺกปฺโป ยา จ สมฺมาทิฏฺฐิ, อยํ ปญฺญากฺขนฺโธ. เอวํ ตาโย ติสฺโส สิกฺขา. เอวํ ตีหากาเรหิ ทส ปทานิ ฯเปฯ.}

\prose{\csromanfolio{246}ตตฺถ โยคาวจโร สีลกฺขนฺเธ ฐิโต โทสํ อกุสลํ น อุปาทิยติ, โทสานุสยํ สมูหนติ, โทสสลฺลํ อุทฺธรติ, ทุกฺขเวทนํ ปริชานาติ, กามธาตุํ สมติกฺกมติ. สมาธิ\-กฺขนฺเธ ฐิโต โลภํ อกุสลํ น อุปาทิยติ, ราคานุสยํ สมูหนติ, โลภสลฺลํ อุทฺธรติ, สุขเวทนํ ปริชานาติ, รูปธาตุํ สมติกฺกมติ. ปญฺญากฺขนฺเธ ฐิโต โมหํ อกุสลํ น อุปาทิยติ, อวิชฺชานุสยํ สมูหนติ, โมหสลฺลํ ทิฏฺฐิสลฺลญฺจ อุทฺธรติ, อทุกฺขมสุข\-เวทนํ ปริชานาติ, อรูปธาตุํ สมติกฺกมติ. อิติ ตีหิ ขนฺเธหิ ตีณิ อกุสลมูลานิ น อุปาทิยติ, จตฺตาริ สลฺลานิ อุทฺธรติ, ติสฺโส เวทนา ปริชานาติ, เตธาตุกํ สมติกฺกมติ.}

\proseitem{62}{ตตฺถ กตมา อวิชฺชา, ยํ จตูสุ อริยสจฺเจสุ อญฺญาณนฺติ วิตฺถาเรน ยถา โส ปาณสชฺเชสุ กถํกถา กาตพฺพํ. ตตฺถ กตมํ วิญฺญาณํ, ฉ วิญฺญาณกายา เวทนา สญฺญา เจตนา ผสฺโส มนสิกาโร, อิทํ นามํ. ยตฺถ กตมํ รูปํ, จาตุมหาภูติกํ จตุนฺนํ มหาภูตานํ อุปาทายรูปสฺส ปญฺญตฺติํ. อิติ ปุริมกญฺจ นามํ อิทญฺจ รูปํ ตทุภยํ นามรูปนฺติ วุจฺจติ. ตตฺถ ฉฬายตนนฺติ ฉ อชฺฌตฺติกานิ อายตนานิ, จกฺขุ อชฺฌตฺติกํ อายตนํ ยาว มโน อชฺฌตฺติกํ อายตนํ. ผสฺโสติ ฉ ผสฺสกายา จกฺขุ\-สมฺผสฺโส ยาว มโน\-สมฺผสฺโสติ ผสฺโส. ฉ เวทนากายา เวทนา. ตณฺหาติ ฉ ตณฺหากายา ตณฺหา. อุปาทานนฺติ จตฺตาริ อุปาทานานิ กามุปาทานํ ทิฏฺฐุปาทานํ สีลพฺพตุ\-ปาทานํ อตฺตวาทุ\-ปาทานนฺติ อุปาทานํ. ภโวติ ตโย ภวา กามภโว รูปภโว อรูปภโว. ตตฺถ กตมา ชาติ, ยา ปฐมํ ขนฺธานํ ปฐมํ ธาตูนํ ปฐมํ อายตนานํ อุปฺปตฺติ ชาติ สญฺชาติ โอกฺกนฺติ อภินิพฺพตฺติ ขนฺธานํ ปาตุภาโว, อยํ ชาติ. ตตฺถ กตมา ชรา, ชรา นาม ยํ ตํ ขณฺฑิจฺจํ ปาลิจฺจํ วลิตฺตจตา ปวิวิตฺตํ จตุนฺนํ มหาภูตานํ วิวณฺณตํ ภคฺโค ตํ ชรา หียนา ปหียนา อายุโน หานิ สํหานิ อินฺทฺริยานํ ปริเภโท อุปนาโห ปริปาโก, อยํ ชรา. ตตฺถ กตมํ มรณํ, มรณํ นาม ยํ ตสฺมิํ ตสฺมิํ สตฺตนิกาเย เตสํ เตสํ สตฺตานํ จุติ จวนตา มรณํ กาลํกิริยา อุทฺธุมาตกานํ เภโท กายสฺส ชีวิติ\-นฺทฺริยสฺส อุปจฺเฉโท, อิทํ มรณํ. อิติ ปุริมิกา จ ชรา อิทญฺจ มรณํ ตทุภยํ ชรามรณํ.}

\prose{\csromanfolio{247}ตตฺถ อนฺธ\-การ\-ติมิสา ยถาภูตํ อปฺป\-ชานน\-ลกฺขณา อวิชฺชา สงฺขารานํ ปทฏฺฐานํ, อภิสงฺขรณ\-ลกฺขณา สงฺขารา, อุปจยปุนพฺภวาภิโรปนปจฺจุปฏฺฐานา. เต วิญฺญาณสฺส ปทฏฺฐานํ, วตฺถุ สวิญฺญตฺติ\-ลกฺขณํ วิญฺญาณํ, ตํ นามรูปสฺส ปทฏฺฐานํ. อเนก\-สนฺนิสฺสย\-ลกฺขณํ นามรูปํ, ตํ สฬายตนสฺส ปทฏฺฐานํ. อินฺทฺริย\-ววตฺถาปน\-ลกฺขณํ สฬายตนํ, ตํ ผสฺสสฺส ปทฏฺฐานํ. สนฺนิปาต\-ลกฺขโณ ผสฺโส, โส เวทนาย ปทฏฺฐานํ. อนุภวน\-ลกฺขณา เวทนา, สา ตณฺหาย ปทฏฺฐานํ. อชฺโฌสาน\-ลกฺขณา ตณฺหา, สา อุปาทานสฺส ปทฏฺฐานํ. อาทานปริหนนลกฺขณํ อุปาทานํ, ตํ ภวสฺส ปทฏฺฐานํ. นานา\-คติ\-วิกฺเขป\-ลกฺขโณ ภโว, โส ชาติยา ปทฏฺฐานํ. ขนฺธานํ ปาตุภาว\-ลกฺขณา ชาติ, สา ชราย ปทฏฺฐานํ. อุปนย\-ปริปาก\-ลกฺขณา ชรา, สา มรณสฺส ปทฏฺฐานํ. อายุกฺขยชีวิตอุปโรธลกฺขณํ มรณํ, ตํ ทุกฺขสฺส ปทฏฺฐานํ. กาย\-สมฺปีฬน\-ลกฺขณํ ทุกฺขํ, ตํ โทมนสฺสสฺส ปทฏฺฐานํ. จิตฺต\-สมฺปีฬน\-ลกฺขณํ โทมนสฺสํ, ตํ โสกสฺส ปทฏฺฐานํ. โสจนลกฺขโณ โสโก, โส ปริเทวสฺส ปทฏฺฐานํ. วจีนิจฺฉารณ\-ลกฺขโณ ปริเทโว, โส อุปายาสสฺส ปทฏฺฐานํ. เย อายาสา เต อุปายาสา.}

\prose{นว ปทานิ ยตฺถ สพฺโพ อกุสลปกฺโข สงฺคหํ สโมสรณํ คจฺฉติ. กตมานิ นว ปทานิ, เทฺว มูลกิเลสา, ตีณิ อกุสลมูลานิ, จตฺตาโร วิปลฺลาสา. ตตฺถ เทฺว มูลกิเลสา อวิชฺชา จ ภวตณฺหา จ, ตีณิ อกุสลมูลานิ โลโภ โทโส โมโห จ. จตฺตาโร วิปลฺลาสา\csromansharedfootnote{953a1a4de1b0fcf7}{อํ 1. 361 ปิฏฺเฐ.} “อนิจฺเจ นิจฺจนฺ”ติ สญฺญาวิปลฺลาโส จิตฺตวิปลฺลาโส ทิฏฺฐิ\-วิปลฺลาโส. “ทุกฺเข สุขนฺ”ติ สญฺญาวิปลฺลาโส จิตฺตวิปลฺลาโส ทิฏฺฐิ\-วิปลฺลาโส. “อนตฺตนิ อตฺตา”ติ สญฺญาวิปลฺลาโส จิตฺตวิปลฺลาโส ทิฏฺฐิ\-วิปลฺลาโส. “อสุเภ สุภนฺ”ติ สญฺญาวิปลฺลาโส จิตฺตวิปลฺลาโส ทิฏฺฐิ\-วิปลฺลาโส.}

\proseitem{63}{ตตฺถ อวิชฺชา นาม จตูสุ อริยสจฺเจสุ ยถาภูตํ อญฺญาณํ, อยํ อวิชฺชา. ภวตณฺหา นาม โย ภเวสุ ราโค สาราโค อิจฺฉา มุจฺฉา ปตฺถนา นนฺที อชฺโฌสานํ อปริจฺจาโค, อยํ ภวตณฺหา.}

\prose{ตตฺถ กตโม โลโภ อกุสลมูลํ, โลโภ นาม โส เตสุ เตสุ ปรวตฺถูสุ ปรทพฺเพสุ ปรฏฺฐาเนสุ ปรสาปเตยฺเยสุ \csromanfolio{248}ปรปริคฺคหิเตสุ โลโภ ลุพฺภนา อิจฺฉา มุจฺฉา ปตฺถนา นนฺที อชฺโฌสานํ อปริจฺจาโค, อยํ โลโภ อกุสลมูลํ. กสฺเสตํ มูลํ, โลโภ โลภชสฺส อกุสลสฺส กายกมฺมสฺส วจีกมฺมสฺส มโนกมฺมสฺส จ, ตถา ยถา ตํสมฺปยุตฺตานํ จิตฺต\-เจตสิกานํ ธมฺมานํ มูลํ.}

\prose{ตตฺถ กตโม โทโส อกุสลมูลํ, โส สตฺเตสุ อาฆาโต อกฺขนฺติ อปฺปจฺจโย พฺยาปาโท ปโทโส อนตฺถกามตา เจตโส ปฏิฆาโต, อยํ โทโส อกุสลมูลํ. กสฺเสตํ มูลํ, โทสชสฺส กายกมฺมสฺส วจีกมฺมสฺส มโนกมฺมสฺส สมฺปยุตฺตานญฺจ จิตฺต\-เจตสิกานํ ธมฺมานํ มูลํ.}

\prose{ตตฺถ กตโม โมโห อกุสลมูลํ, ยํ จตูสุ อริยสจฺเจสุ อนภิสมโย อสมฺปชฺชคฺคาโห อปฺปฏิเวโธ โมโห มุยฺหนา สมฺโมโห สมฺมุยฺหนา อวิชฺชา ตโม อนฺธกาโร อาวรณํ นีวรณํ ฉทนํ อจฺฉทนํ\csromansharedfootnote{243aabd79c70e29f}{อเวจฺฉทนํ (อิ, ก)} อปสจฺฉาคมนํ กุสลานํ ธมฺมานํ, อยํ โมโห อกุสลมูลํ. กสฺเสตํ มูลํ, โมหชสฺส อกุสลสฺส กายกมฺมสฺส วจีกมฺมสฺส มโนกมฺมสฺส จ ตํสมฺปยุตฺตกานญฺจ จิตฺต\-เจตสิกานํ ธมฺมานํ มูลํ.}

\prose{ตตฺถ วิปลฺลาสา ชานิตพฺพา, วิปลฺลาสานํ วตฺถุ ชานิตพฺพํ. ยํ วิปลฺลาสํ สิยา, ตํ ชานิตพฺพํ. ตตฺถ เอโก วิปลฺลาโส ตีณิ วิปลฺลาสานิ จตฺตาริ วิปลฺลาส\-วตฺถูนิ. กตโม เอโก วิปลฺลาโส จ, เยน ปฏิปกฺเขน วิปลฺลาสิตํ คณฺหาติ “อนิจฺเจ นิจฺจนฺ”ติ, “ทุกฺเข สุขนฺ”ติ, “อนตฺตนิ อตฺตา”ติ, “อสุเภ สุภนฺ”ติ, อยํ เอโก วิปลฺลาโส. กตมานิ จตฺตาริ วิปลฺลาส\-วตฺถูนิ, กาโย เวทนา จิตฺตํ ธมฺมา จ. อิมานิ จตฺตาริ วิปลฺลาส\-วตฺถูนิ, กตมานิ ตีณิ วิปลฺลาสานิ, สญฺญา จิตฺตํ ทิฏฺฐิ จ. อิมานิ ตีณิ วิปลฺลาสานิ.}

\prose{ตตฺถ มนาปิเก วตฺถุมฺหิ อินฺทฺริยวตฺเถ วณฺณายตเน วา โย นิมิตฺตสฺส อุคฺคาโห, อยํ สญฺญาวิปลฺลาโส. ตตฺถ วิปรีต\-จิตฺตสฺส วตฺถุมฺหิ สติ วิญฺญตฺติ, อยํ จิตฺตวิปลฺลาโส. ตตฺถ วิปรีต\-จิตฺตสฺส ตมฺหิ รูเป “อสุเภ สุภนฺ”ติ ยา ขนฺติ รุจิ อุเปกฺขนา นิจฺฉโย ทิฏฺฐิ นิทสฺสนํ สนฺตีรณา, อยํ ทิฏฺฐิ\-วิปลฺลาโส. ตตฺถ วตฺถุเภเทน กาเยสุ ทฺวาทส วิปลฺลาสา ภวนฺติ. \csromanfolio{249}ตโยกาเย ตโย เวทนาย ตโย จิตฺเต ตโย ธมฺเม, จตฺตาโร สญฺญาวิปลฺลาสา จตฺตาโร จิตฺตวิปลฺลาสา จตฺตาโร ทิฏฺฐิ\-วิปลฺลาสา, อายตนู\-ปจยโต จกฺขุวิญฺญาณ\-สญฺญา\-สมงฺคิสฺส รูเปสุ ทฺวาทส วิปลฺลาสา ยาว มโน สญฺญา\-สมงฺคิสฺส, ธมฺเมสุ ทฺวาทส วิปลฺลาสา ฉ ทฺวาทสกา จตฺตาริ วิปลฺลาสา ภวนฺติ. อารมฺมณ\-นานตฺตโต หิ อปริมิต\-สงฺเขยฺยานํ สตฺตานํ\csromansharedfootnote{2d392b2c90d82830}{อตฺตานํ (ก)} อปริมิตม\-สงฺเขยฺยา วิปลฺลาสา ภวนฺติ หีนุกฺกฏฺฐมชฺฌิม\-ตาย.}

\proseitem{64}{ตตฺถ ปญฺจกฺขนฺธา จตฺตาริ อตฺตภาว\-วตฺถูนิ ภวนฺติ. โย รูปกฺขนฺโธ, โส กาโย อตฺตภาว\-วตฺถุ. โย เวทนากฺขนฺโธ, โส เวทนา อตฺตภาว\-วตฺถุ. โย สญฺญากฺขนฺโธ จ สงฺขาร\-กฺขนฺโธ จ, เต ธมฺมา อตฺตภาว\-วตฺถุ. โย วิญฺญาณ\-กฺขนฺโธ, โส จิตฺตํ อตฺตภาว\-วตฺถุ. อิติ ปญฺจกฺขนฺธา จตฺตาริ อตฺตภาว\-วตฺถูนิ. ตตฺถ กาเย “อสุเภ สุภนฺ”ติ วิปลฺลาโส ภวติ. เอวํ เวทนาสุ. จิตฺเต. ธมฺเมสุ จ อตฺตวิปลฺลาโส ภวติ. ตตฺถ จตุนฺนํ วิปลฺลาสานํ สมุคฺฆาต\-นตฺถํ ภควา จตฺตาโร สติปฏฺฐาเน เทเสติ ปญฺญเปติ กาเย กายานุปสฺสี วิหรโต “อสุเภ สุภนฺ”ติ. วิปลฺลาสํ สมุคฺฆาเตติ, เอวํ เวทนาสุ. จิตฺเต. ธมฺเมสุ จ กาตพฺพํ.}

\prose{ตตฺถ อนฺธ\-การ\-ติมิสา อปฺปฏิเวธ\-ลกฺขณา อวิชฺชา, ตสฺสา วิปลฺลาส\-ปทฏฺฐานํ. อชฺโฌสาน\-ลกฺขณา ตณฺหา, ตสฺสา ปิยรูป\-สาตรูปํ ปทฏฺฐานํ. อตฺตา\-สย\-วญฺจนา\-ลกฺขโณ โลโภ, ตสฺส อทินฺนาทานํ ปทฏฺฐานํ. อิธ วิวาทลกฺขโณ โทโส, ตสฺส ปาณาติปาโต ปทฏฺฐานํ. วตฺถุ\-วิปฺปฏิปตฺติ\-ลกฺขโณ โมโห, ตสฺส มิจฺฉา\-ปฏิปตฺติ ปทฏฺฐานํ. สงฺขตานํ ธมฺมานํ อวินา\-ส\-คฺคหณ\-ลกฺขณา นิจฺจสญฺญา, ตสฺสา สพฺพสงฺขารา ปทฏฺฐานํ. สาสว\-ผสฺโส\-ปคมน\-ลกฺขณา สุขสญฺญา, ตสฺสา มมงฺกาโร ปทฏฺฐานํ. ธมฺเมสุ อุปคมน\-ลกฺขณา อตฺตสญฺญา, ตสฺสา อหงฺกาโร ปทฏฺฐานํ. วณฺณ\-สงฺคหณ\-ลกฺขณา สุภสญฺญา, ตสฺสา อินฺทฺริย\-อสํวโร ปทฏฺฐานํ. เอเตหิ นวหิ ปเทหิ อุทฺทิฏฺเฐหิ สพฺโพ อกุสลปกฺโข นิทฺทิฏฺโฐ ภวติ, โส จ โข พหุสฺสุเตน สกฺกา ชานิตุํ โน อปฺปสฺสุเตน, ปญฺญวตา โน ทุปฺปญฺเญน, ยุตฺเตน โน อยุตฺเตน.}

\prose{\csromanfolio{250}นว ปทานิ กุสลานิ ยตฺถ สพฺโพ กุสลปกฺโข สงฺคโห สโมสรณํ คจฺฉนฺติ. กตมานิ นว ปทานิ, สมโถ วิปสฺสนา อโลโภ อโทโส อโมโห อนิจฺจสญฺญา ทุกฺขสญฺญา อนตฺตสญฺญา อสุภสญฺญา จ.}

\prose{ตตฺถ กตโม สมโถ, ยา จิตฺตสฺส ฐิติ สณฺฐิติ อวฏฺฐิติ ฐานํ ปฏฺฐานํ อุปฏฺฐานํ สมาธิ สมาธานํ อวิกฺเขโป อวิปฺปฏิสาโร วูปสโม มานโส เอกคฺคํ จิตฺตสฺส, อยํ สมโถ.}

\prose{ตตฺถ กตมา วิปสฺสนา, ขนฺเธสุ วา ธาตูสุ วา อายตเนสุ วา นามรูเปสุ วา ปฏิจฺจ\-สมุปฺปาเทสุ วา ปฏิจฺจ\-สมุปฺปนฺเนสุ วา ธมฺเมสุ ทุกฺเขสุ วา สมุทเยสุ วา นิโรเธ วา มคฺเค วา กุสลากุสเลสุ วา ธมฺเมสุ สาวชฺช\-อนวชฺเชสุ วา กณฺหสุกฺเกสุ วา เสวิตพฺพ\-อเสวิตพฺเพสุ วา โส ยถาภูตํ วิจโย ปวิจโย วีมํสา ปรวีมํสา คาหนา อคฺคาหนา ปริคฺคาหนา จิตฺเตน ปริจิตนา ตุลนา อุปปริกฺขา ญาณํ วิชฺชา วา จกฺขุ พุทฺธิ เมธา ปญฺญา โอภาโส อาโลโก อาภา ปภา ขคฺโค นาราโจ\csromansharedfootnote{a3ef47c45556f09d}{นารชฺโช (อิ, ก)} ธมฺมวิจย\-สมฺโพชฺฌงฺโค สมฺมาทิฏฺฐิ มคฺคงฺคํ, อยํ วิปสฺสนา. เตเนสา วิปสฺสนา อิติ วุจฺจติ วิวิธา วา เอสา วิปสฺสนาติ, ตสฺมา เอสา วิปสฺสนาติ วุจฺจติ. ทฺวิธา เจสา หิ วิปสฺสนา ธมฺมวิปสฺสนาติ วุจฺจติ, ทฺวิธา อิมาย ปสฺสติ สุภญฺจ อสุภญฺจ กณฺหญฺจ สุกฺกญฺจ เสวิตพฺพญฺจ อเสวิตพฺพญฺจ กมฺมญฺจ วิปากญฺจ พนฺธญฺจ วิโมกฺขญฺจ อาจยญฺจ อปจยญฺจ ปวตฺติญฺจ นิวตฺติญฺจ สํกิเลสญฺจ โวทานญฺจ, เอวํ วิปสฺสนาติ วุจฺจติ. อถ วา วิอิติ อุปสคฺโค ปสฺสนาติ อตฺโถ ตสฺมา วิปสฺสนาติ วุจฺจเต, อยํ วิปสฺสนา.}

\proseitem{65}{ตตฺถ เทฺว โรคา สตฺตานํ อวิชฺชา จ ภวตณฺหา จ, เอเตสํ ทฺวินฺนํ โรคานํ นิฆาตาย ภควตา เทฺว เภสชฺชานิ วุตฺตานิ สมโถ จ วิปสฺสนา จ. อิมานิ เทฺว เภสชฺชานิ ปฏิเสเวนฺโต เทฺว อโรเค สจฺฉิกโรติ ราควิราคํ เจโตวิมุตฺติํ อวิชฺชาวิราคญฺจ ปญฺญาวิมุตฺติํ. ตตฺถ ตณฺหาโรคสฺส สมโถ เภสชฺชํ, ราควิราคา เจโตวิมุตฺติ อโรคํ. อวิชฺชาโรคสฺส วิปสฺสนา\-เภสชฺชํ อวิชฺชาวิราคา ปญฺญาวิมุตฺติ อโรคํ. เอวญฺหิ ภควา จาห, “เทฺว ธมฺมา ปริญฺเญยฺยา\csromansharedfootnote{96634675ca383466}{ทิ 3. 228 ปิฏฺเฐ.} นามญฺจ รูปญฺจ, เทฺว ธมฺมา ปหาตพฺพา อวิชฺชา จ ภวตณฺหา จ, เทฺว ธมฺมา ภาเวตพฺพา สมโถ จ \csromanfolio{251}วิปสฺสนา จ, เทฺว ธมฺมา สจฺฉิกาตพฺพา วิชฺชา จ วิมุตฺติ จา”ติ. ตตฺถ สมถํ ภาเวนฺโต รูปํ ปริชานาติ, รูปํ ปริชานนฺโต ตณฺหํ ปชหติ, ตณฺหํ ปชหนฺโต ราควิราคา เจโตวิมุตฺติํ สจฺฉิกโรติ, วิปสฺสนํ ภาเวนฺโต นามํ ปริชานาติ, นามํ ปริชานนฺโต อวิชฺชํ ปชหติ, อวิชฺชํ ปชหนฺโต อวิชฺชาวิราคา ปญฺญาวิมุตฺติํ สจฺฉิกโรติ. ยทา ภิกฺขุโน เทฺว ธมฺมา ปริญฺญาตา ภวนฺติ นามญฺจ รูปญฺจ, ตถาสฺส เทฺว ธมฺมา ปหีนา ภวนฺติ อวิชฺชา จ ภวตณฺหา จ. เทฺว ธมฺมา ภาวิตา ภวนฺติ สมโถ จ วิปสฺสนา จ, เทฺว ธมฺมา สจฺฉิกาตพฺพา ภวนฺติ วิชฺชา จ วิมุตฺติ จ. เอตฺตาวตา ภิกฺขุ กตกิจฺโจ ภวติ. เอสา โสปาทิเสสา นิพฺพานธาตุ. ตสฺส อายุปริยาทานา ชีวิติ\-นฺทฺริยสฺส อุปโรธา อิทญฺจ ทุกฺขํ นิรุชฺฌติ, อญฺญญฺจ ทุกฺขํ น อุปฺปชฺชติ. ตตฺถ โย อิเมสํ ขนฺธานํ ธาตุ\-อายตนานํ นิโรโธ วูปสโม อญฺเญสญฺจ ขนฺธ\-ธาตุ\-อายตนานํ อปฺปฏิสนฺธิ อปาตุภาโว, อยํ อนุปาทิเสสา นิพฺพานธาตุ.}

\prose{ตตฺถ กตมํ อโลโภ กุสลมูลํ, ยํธาตุโก อโลโภ อลุพฺภนา อลุพฺภิตตฺตํ อนิจฺฉา อปตฺถนา อกนฺตา อนชฺโฌสานํ. อยํ อโลโภ กุสลมูลํ. กสฺเสตํ มูลํ, อโลภชสฺส กุสลสฺส กายกมฺมสฺส วจีกมฺมสฺส มโนกมฺมสฺส ตํสมฺปยุตฺตานญฺจ จิตฺต\-เจตสิกานํ ธมฺมานํ มูลํ. อถ วา อริโย อฏฺฐงฺคิโก มคฺโค กุสลนฺติ วุจฺจติ, โส ติณฺณํ มคฺคงฺคานํ มูลํ. กตเมสํ ติณฺณํ, สมฺมา\-สงฺกปฺปสฺส สมฺมาวายามสฺส สมฺมา\-สมาธิสฺส จ อิเมสํ มูลนฺติ, ตสฺมา กุสลมูลนฺติ วุจฺจติ.}

\prose{ตตฺถ กตมํ อโทโส กุสลมูลํ, ยา สตฺเตสุ วา สงฺขาเรสุ วา อนฆาโต อปฺปฏิฆาโต อพฺยาปตฺติ อพฺยาปาโท อโทโส เมตฺตา เมตฺตายนา อตฺถกามตา หิตกามตา เจตโส ปสาโท, อยํ อโทโส กุสลมูลํ. กสฺเสตํ มูลํ, อโทสชสฺส กุสลสฺส กายกมฺมสฺส วจีกมฺมสฺส มโนกมฺมสฺส ตํสมฺปยุตฺตานญฺจ จิตฺต\-เจตสิกานํ ธมฺมานํ มูลํ. อถ วา ติณฺณํ มคฺคงฺคานํ มูลํ. กตเมสํ ติณฺณํ, สมฺมาวาจาย สมฺมา\-กมฺมนฺตสฺส สมฺมาอาชีวสฺส จ อิเมสํ ติณฺณํ มคฺคงฺคานํ มูลํ, ตสฺมา กุสลมูลนฺติ วุจฺจติ.}

\prose{ตตฺถ กตมํ อโมโห กุสลมูลํ, ยํ จตูสุ อริยสจฺเจสุ ยถาภูตํ ญาณทสฺสนํ อภิสมโย สมฺมา จ ปจฺจาคโม ปฏิเวโธ อโมโห อสมฺมุยฺหนา อสมฺโมโห วิชฺชาปกาโส อาโลโก \csromanfolio{252}อนาวรณํ เสกฺขานํ กุสลานํ ธมฺมานํ, อยํ อโมโห กุสลมูลํ. กสฺเสตํ มูลํ, อโมหชสฺส กุสลสฺส กายกมฺมสฺส วจีกมฺมสฺส มโนกมฺมสฺส ตํสมฺปยุตฺตานญฺจ จิตฺต\-เจตสิกานํ ธมฺมานํ มูลํ. อถ วา ทฺวินฺนํ มคฺคงฺคานํ เอตํ มูลํ. กตเมสํ ทฺวินฺน, สมฺมาทิฏฺฐิยา จ สมฺมาสติยา จ อิเมสํ ทฺวินฺนํ มคฺคงฺคานํ มูลํ, ตสฺมา กุสลมูลนฺติ วุจฺจติ. เอวํ อิเมสํ ตีหิ กุสลมูเลหิ อฏฺฐงฺคิโก มคฺโค โยเชตพฺโพ.}

\proseitem{66}{ตตฺถ กตมา อนิจฺจสญฺญา, “สพฺเพ สงฺขารา อุปฺปาท\-วย\-ธมฺมิโน”ติ จ ยา สญฺญา สญฺชานนา ววตฺถปนา อุคฺคาโห, อยํ อนิจฺจสญฺญา. ตสฺสา โก นิสฺสนฺโท, อนิจฺจสญฺญาย ภาวิตาย พหุลีกตาย อฏฺฐสุ โลกธมฺเมสุ จิตฺตํ นานุสนฺธติ น สนฺธติ น สณฺฐหติ, อุเปกฺขา วา ปฏิกฺกูลตา วา สณฺฐหติ, อยมสฺสา นิสฺสนฺโท.}

\prose{ตตฺถ กตมา ทุกฺขสญฺญา, “สพฺเพ สงฺขารา ทุกฺขา”ติ ยา สญฺญา สญฺชานนา ววตฺถปนา อุคฺคาโห, อยํ ทุกฺขสญฺญา. ตสฺสา โก นิสฺสนฺโท, ทุกฺขสญฺญาย ภาวิตาย พหุลีกตาย อาลสฺเส สํปมาเท วิมฺหเย จ จิตฺตํ นานุสนฺธติ น สนฺธติ น สณฺฐหติ, อุเปกฺขา วา ปฏิกฺกูลตา วา สณฺฐหติ, อยมสฺสา นิสฺสนฺโท.}

\prose{ตตฺถ กตมา อนตฺตสญฺญา, “สพฺเพสุ ธมฺเมสุ อนตฺตา”ติ ยา สญฺญา สญฺชานนา ววตฺถปนา อุคฺคาโห, อยํ อนตฺตสญฺญา. ตสฺสา โก นิสฺสนฺโท, อนตฺตสญฺญาย ภาวิตาย พหุลีกตาย อหงฺกาโร จิตฺตํ นานุสนฺธติ น สนฺธติ, มมงฺกาโร น สณฺฐหติ, อุเปกฺขา วา ปฏิกฺกูลตา วา สณฺฐหติ, อยมสฺสา นิสฺสนฺโท.}

\prose{ตตฺถ กตมา อสุภสญฺญา, “สตฺต สงฺขารา อสุภา”ติ ยา สญฺญา สญฺชานนา ววตฺถปนา อุคฺคาโห, อยํ อสุภสญฺญา. ตสฺสา โก นิสฺสนฺโท, อสุภสญฺญาย ภาวิตาย พหุลีกตาย สุภนิมิตฺเต จิตฺตํ นานุสนฺธติ น สนฺธติ น สณฺฐหติ, อุเปกฺขา วา ปฏิกฺกูลตา วา สณฺฐหติ, อยมสฺสา นิสฺสนฺโท.}

\prose{ตตฺถ ปญฺจนฺนํ ขนฺธานํ ปริญฺญา ภควตา เทสิตา, โย ตตฺถ อสุภสญฺญา รูปกฺขนฺธสฺส ปริญฺญตฺตํ, ทุกฺขสญฺญา เวทนา\-กฺขนฺธสฺส ปริญฺญตฺตํ, อนตฺตสญฺญา สญฺญา\-กฺขนฺธสฺส สงฺขาร\-กฺขนฺธสฺส ปริญฺญตฺตํ, อนิจฺจสญฺญา วิญฺญาณ\-กฺขนฺธสฺส ปริญฺญตฺตํ. ตตฺถ สมเถน ตณฺหํ สมุคฺฆาเตติ, วิปสฺสนาย \csromanfolio{253}อวิชฺชํ สมุคฺฆาเตติ, อโทเสน โทสํ สมุคฺฆาเตติ, อโมเหน โมหํ สมุคฺฆาเตติ, อนิจฺจสญฺญาย นิจฺจสญฺญํ สมุคฺฆาเตติ, ทุกฺขสญฺญาย สุขสญฺญํ สมุคฺฆาเตติ, อนตฺตสญฺญาย อตฺตสญฺญํ สมุคฺฆาเตติ, อสุภสญฺญาย สุภสญฺญํ สมุคฺฆาเตติ.}

\prose{จิตฺตวิกฺเขป\-ปฏิสํหรณ\-ลกฺขโณ สมโถ, ตสฺส ฌานานิ ปทฏฺฐานํ. สพฺพธมฺมํ ยถาภูตํ ปฏิเวธ\-ลกฺขณา วิปสฺสนา, ตสฺสา สพฺพเนยฺยํ ปทฏฺฐานํ. อิจฺฉา\-ปฏิสํหรณ\-ลกฺขโณ อโลโภ, ตสฺส อทินฺนาทานา เวรมณี ปทฏฺฐานํ. อพฺยาปาท\-ลกฺขโณ อโทโส, ตสฺส ปาณาติปาตา เวรมณี ปทฏฺฐานํ. วตฺถุ\-อปฺปฏิหต\-ลกฺขโณ อโมโห, ตสฺส สมฺมาปฏิปตฺติ ปทฏฺฐานํ. สงฺขตานํ ธมฺมานํ วินาส\-คฺคหณ\-ลกฺขณา อนิจฺจสญฺญา, ตสฺสา อุทยพฺพโย ปทฏฺฐานํ. สาสว\-ผสฺส\-สญฺชานนลกฺขณา ทุกฺขสญฺญา, ตสฺสา เวทนา ปทฏฺฐานํ. สพฺพธมฺม\-อนุปคมน\-ลกฺขณา อนตฺตสญฺญา, ตสฺสา ธมฺมสญฺญา ปทฏฺฐานํ. วินีลกวิปุพฺพกอุทฺธุมาตกสมุคฺคหณลกฺขณา อสุภสญฺญา, ตสฺสา นิพฺพิทา ปทฏฺฐานํ. อิเมสุ นวสุ ปเทสุ อุปทิฏฺเฐสุ สพฺโพ กุสลปกฺโข อุปทิฏฺโฐ ภวติ, โส จ พหุสฺสุเตน สกฺกา ชานิตุํ โน อปฺปสฺสุเตน, ปญฺญวตา โน ทุปฺปญฺเญน, ยุตฺเตน โน อยุตฺเตนาติ.}

\proseitem{67}{ตตฺถ นิจฺจสญฺญา\-ธิมุตฺตสฺส อปราปรํ จิตฺตํ ปณาเมนฺโต สติม\-ปจฺจเวกฺขโต อนิจฺจสญฺญา น อุปฏฺฐาติ, ปญฺจสุ กามคุเณสุ สุขสฺสาทา\-ธิมุตฺตสฺส อิริยาปถสฺส อคติม\-ปจฺจเวกฺขโต ทุกฺขสญฺญา น อุปฏฺฐาติ, ขนฺธ\-ธาตุ\-อายตเนสุ อตฺตา\-ธิมุตฺตสฺส นานาธาตุ- อเนกธาตุวินิพฺโภคมปจฺจเวกฺขโต อนตฺตสญฺญา น อุปฏฺฐาติ, วณฺณ\-สณฺฐานาภิ\-รตสฺส กาเย สุภา\-ธิมุตฺตสฺส จ วิปฺปฏิจฺฉนฺนา อสุภสญฺญา น อุปฏฺฐาติ.}

\prose{อวิปฺปฏิสาร\-ลกฺขณา สทฺธา, สทฺทหนา ปจฺจุปฏฺฐานํ. ตสฺส จตฺตาริ โสตาปตฺติยงฺคานิ ปทฏฺฐานํ. เอวญฺหิ วุตฺตํ ภควตา\csromansharedfootnote{1c14b16660711759}{สํ 3. 172 ปิฏฺเฐ.} สทฺธินฺทฺริยํ ภิกฺขเว กุหิํ ทฏฺฐพฺพํ, จตูสุ โสตาปตฺติยงฺเคสุ กุสเลสุ ธมฺเมสุ.}

\prose{สูราอปฏิกฺเขปนลกฺขณํ วีริยินฺทฺริยํ, วีริยินฺทฺริยา\-รมฺโภ ปจฺจุปฏฺฐานํ. ตสฺส อตีตา จตฺตาโร สมฺมปฺปธานา ปทฏฺฐานํ. ยถา วุตฺตํ ภควตา\csromansharedfootnote{1c14b16660711759}{สํ 3. 172 ปิฏฺเฐ.} วีริยินฺทฺริยํ ภิกฺขเว กุหิํ ทฏฺฐพฺพํ, จตูสุ สมฺม\-ปฺปธาเนสุ.}

\prose{\csromanfolio{254}สติ สรณลกฺขณา, อสมฺโมห\-ปจฺจุปฏฺฐานา. ตสฺส อตีตา จตฺตาโร สติปฏฺฐานา ปทฏฺฐานํ. ยถา วุตฺตํ ภควตา\csromansharedfootnote{1c14b16660711759}{สํ 3. 172 ปิฏฺเฐ.} สตินฺทฺริยํ ภิกฺขเว กุหิํ ทฏฺฐพฺพํ, จตูสุ สติปฏฺฐาเนสุ.}

\prose{เอกคฺคลกฺขโณ สมาธิ, อวิกฺเขป\-ปจฺจุปฏฺฐาโน, ตสฺส จตฺตาริ ญาณานิ ปทฏฺฐานํ. ยถา วุตฺตํ ภควตา\csromansharedfootnote{1c14b16660711759}{สํ 3. 172 ปิฏฺเฐ.} สมาธินฺทฺริยํ ภิกฺขเว กุหิํ ทฏฺฐพฺพํ, จตูสุ ฌาเนสุ.}

\prose{ปชานน\-ลกฺขณา ปญฺญา, ภูตตฺถ\-สนฺตีรณา ปจฺจุปฏฺฐานา, ตสฺส จตฺตาริ อริยสจฺจานิ ปทฏฺฐานํ. ยถา วุตฺตํ ภควตา\csromansharedfootnote{1c14b16660711759}{สํ 3. 172 ปิฏฺเฐ.} ปญฺญินฺทฺริยํ ภิกฺขเว กุหิํ ทฏฺฐพฺพํ, จตูสุ อริยสจฺเจสุ.}

\prose{จตฺตาริ จกฺกานิ\csromansharedfootnote{888a040ded07730b}{อํ 1. 341 ปิฏฺเฐ.} ปติรูป\-เทส\-วาโส จกฺกํ, สปฺปุริสู\-ปนิสฺสโย จกฺกํ, อตฺต\-สมฺมา\-ปณิธานํ จกฺกํ, ปุพฺเพ กตปุญฺญตา จกฺกํ. ตตฺถ อริย\-สนฺนิสฺสย\-ลกฺขโณ ปติรูป\-เทส\-วาโส, โส สปฺปุริสู\-ปนิสฺสยสฺส ปทฏฺฐานํ. อริย\-สนฺนิสฺสย\-ลกฺขโณ สปฺปุริ\-สฺสู\-ปนิสฺสโย, โส อตฺตสมฺมาปณิธา\-นสฺส ปทฏฺฐานํ. สมฺมาปฏิปตฺติ\-ลกฺขณํ อตฺต\-สมฺมา\-ปณิธานํ, ตํ ปุญฺญานํ ปทฏฺฐานํ. กุสลธมฺโม\-ปจย\-ลกฺขณํ ปุญฺญํ, ตํ สพฺพ\-สมฺปตฺตีนํ ปทฏฺฐานํ.}

\prose{เอกา\-ทสสีล\-มูลกา ธมฺมา สีลวโต อวิปฺปฏิสาโร ภวติ ฯเปฯ โส วิมุตฺติ\-ญาณ\-ทสฺสนํ “นาปรํ อิตฺถตฺตายา”ติ ปชานนา. ตตฺถ เวร\-มณิลกฺขณํ สีลํ, ตํ อวิปฺปฏิสารสฺส ปทฏฺฐานํ. น อตฺตานุวาท\-ลกฺขโณ อวิปฺปฏิสาโร, โส ปาโมชฺชสฺส ปทฏฺฐานํ. อภิปฺปโมท\-นล\-กฺขณํ ปาโมชฺชํ, ตํ ปีติยา ปทฏฺฐานํ. อตฺตมน\-ลกฺขณา ปีติ, สา ปสฺสทฺธิยา ปทฏฺฐานํ. กมฺมนิย\-ลกฺขณา ปสฺสทฺธิ, สา สุขสฺส ปทฏฺฐานํ. อพฺยาปาท\-ลกฺขณํ สุขํ, ตํ สมาธิโน ปทฏฺฐานํ. อวิกฺเขป\-นล\-กฺขโณ สมาธิ, โส ยถา\-ภูต\-ญาณ\-ทสฺสนสฺส ปทฏฺฐานํ. อวิปรีต\-สนฺตีรณ\-ลกฺขณา ปญฺญา, สา นิพฺพิทาย ปทฏฺฐานํ อนาลย\-นล\-กฺขณา นิพฺพิทา, สา วิราคสฺส ปทฏฺฐานํ. อสํกิเลสลกฺขโณ วิราโค, โส วิมุตฺติยา ปทฏฺฐานํ. อกุสลธมฺม\-วิเวก\-ลกฺขณา วิมุตฺติ, สา วิมุตฺติโน โวทานสฺส ปทฏฺฐานํ.}

\proseitem{68}{จตสฺโส อริยภูมิโย จตฺตาริ สามญฺญผลานิ. ตตฺถ โย ยถาภูตํ ปชานาติ, เอสา ทสฺสนภูมิ. โสตาปตฺติผลญฺจ โส \csromanfolio{255}ยถาภูตํ ปชานิตฺวา นิพฺพินฺทติ, อิทํ ตนุกามราคสฺส ปทฏฺฐานํ พฺยาปาทานํ. สกทาคามิ\-ผลญฺจ สณฺหํ วิรชฺชติ, อยํ ราควิราคา เจโตวิมุตฺติ. อนาคามิผลญฺจ ยํ อวิชฺชาวิราคา วิมุจฺจติ, อยํ กตาภูมิ. อรหตฺตญฺจ สามญฺญผลานีติ โก วจนตฺโถ, อริโย อฏฺฐงฺคิโก มคฺโค สามญฺญํ, ตสฺเสตานิ ผลานิ สามญฺญผลานีติ วุจฺจติ. กิสฺส พฺรหฺมญฺญผลานีติ วุจฺจนฺเต, พฺรหฺมญฺญอริโย อฏฺฐงฺคิโก มคฺโค, ตสฺส ตานิ ผลานีติ พฺรหฺมญฺญผลานีติ วุจฺจนฺเต.}

\prose{ตตฺถ โสตาปนฺโน กถํ โหติ, สห สจฺจาภิสมยา อริยสาวกสฺส ตีณิ สํโยชนานิ ปหียนฺติ สกฺกายทิฏฺฐิ วิจิกิจฺฉา สีลพฺพต\-ปรามาโส จ, อิเมสํ ติณฺณํ สํโยชนานํ ปหานา ปริกฺขยา อริยสาวโก โหติ โสตาปนฺโน อวินิปาต\-ธมฺโม ยาว ทุกฺขสฺสนฺตํ กโรติ.}

\prose{ตตฺถ กตมา สกฺกายทิฏฺฐิ, อสฺสุตวา พาโล ปุถุชฺชโน ยาว อริยธมฺเม อโกวิโท, โส รูปํ อตฺตโต สมนุปสฺสติ ยาว วิญฺญาณสฺมิํ อตฺตานํ, โส อิเมสุ ปญฺจสุ ขนฺเธสุ อตฺตคฺคาโห วา อตฺตนิยคฺคาโห วา เอโสหมสฺมิ เอกสฺมิํ วสวตฺติโก\csromansharedfootnote{612657f23b64c938}{อวตฺติโต (อิ, ก)} ปกฺขิตฺโต อนุคฺคโห อนุสยนฺโต องฺคมงฺคนฺติ ปรติ. ยา ตถาภูตสฺส ขนฺติ รุจิ เปกฺขนา อาการ\-ปริวิตกฺโก ทิฏฺฐินิ\-ชฺฌายนา อภิปฺปสนฺนา, อยํ วุจฺจเต สกฺกายทิฏฺฐีติ.}

\prose{ตตฺถ ปญฺจ ทิฏฺฐิโย อุจฺเฉทํ ภชนฺติ. กตมาโย ปญฺจ, รูปํ อตฺตโต สมนุปสฺสติ, ยาว วิญฺญาณํ อตฺตโต สมนุปสฺสติ, อิมาโย ปญฺจ อุจฺเฉทํ ภชนฺติ, อวเสสาโย ปนฺนรส สสฺสตํ ภชนฺติ. อิติ สกฺกายทิฏฺฐิ\-ปหานา ทฺวาสฏฺฐิ\-ทิฏฺฐิคตานิ ปหียนฺติ. ปหานา อุจฺเฉทํ สสฺสตญฺจ น ภชติ. อิติ อุจฺเฉท\-สสฺสต\-ปฺปหานา อริยสาวกสฺส น กิญฺจิ ทิฏฺฐิคตํ ภวติ, อญฺญา วา โลกุตฺตราย สมฺมาทิฏฺฐิยา. กถํ ปน สกฺกายทิฏฺฐิ น ภวติ, อิธ อริยสาวโก สุตวา โหติ, สพฺโพ สุกฺกปกฺโข กาตพฺโพ, ยาว อริยธมฺเมสุ โกวิโท รูปํ อนตฺตโต สมนุปสฺสติ, ยาว วิญฺญาณํ ฯเปฯ. เอวมสฺส สมนุ\-ปสฺส\-นฺตสฺส สกฺกายทิฏฺฐิ น ภวติ.}

\prose{\csromanfolio{256}กถํ วิจิกิจฺฉา น ภวติ, อิธ อริยสาวโก พุทฺเธ น กงฺขติ, น วิจิกิจฺฉติ อภิปฺปสีทติ, อิติปิ โส ภควาติ สพฺพํ. ธมฺเม น กงฺขติ น วิจิกิจฺฉติ สพฺพํ. ยาว ตณฺหกฺขโย วิราโค นิโรโธ นิพฺพานนฺติ, อิมินา ทุติเยน อากงฺขิเยน ธมฺเมน สมนฺนาคโต โหติ. สํเฆ น กงฺขติ ฯเปฯ. ยาว ปูชา เทวานญฺจ มนุสฺสานญฺจาติ, อิมินา ตติเยน อากงฺขิเยน ธมฺเมน สมนฺนาคโต โหติ.}

\prose{สพฺเพ สงฺขารา ทุกฺขาติ น กงฺขติ น วิจิกิจฺฉติ อธิมุจฺจติ อภิปฺปสีทติ. ตณฺหา ทุกฺข\-สมุทโยติ น กงฺขติ น วิจิกิจฺฉติ. ตณฺหานิโรธา ทุกฺข\-นิโรโธติ น กงฺขติ น วิจิกิจฺฉติ. อริโย อฏฺฐงฺคิโก มคฺโค ทุกฺขนิโรธ\-คามินี ปฏิปทาติ น กงฺขติ น วิจิกิจฺฉติ อธิมุจฺจติ อภิปฺปสีทติ. ยาว พุทฺเธ วา ธมฺเม วา สํเฆ วา ทุกฺเข วา สมุทเย วา นิโรเธ วา มคฺเค วา กงฺขายนา วิมติ วิจิกิจฺฉา เทฺวธาปถา อาสปฺปนา\csromansharedfootnote{3748f6b80440f1a7}{อปฺปนา (อิ, ก) อภิ 1. 208 นิกฺเขปกณฺเฑ ปสฺสิตพฺพํ.} ปริสปฺปนา อนวฏฺฐานํ อธิฏฺฐา\-คมนํ\csromansharedfootnote{d6858cd4e73514c9}{อนิฏฺฐาคมนํ (ก)} อเนกํโส อเนกํสิกตา, เต ตสฺส ปหีนา ภวนฺติ ปณุนฺนา อุจฺฉินฺนมูลา ตาลาวตฺถุกตา อนภาวํกตา อายติํ อนุปฺปาทธมฺมา.}

\proseitem{69}{ตตฺถ สีลพฺพต\-ปรามาโส ทฺวิธา สีลสฺส วา สุทฺธสฺส วา, ตตฺถ สีลสฺส สีลพฺพต\-ปรามาโส อิมินาหํ สีเลน วา วเตน วา ตเปน วา พฺรหฺมจริเยน วา เทโว วา ภวิสฺสามิ เทวญฺญตโร วา ตตฺถ กโปตปาทาหิ อจฺฉราหิ สทฺธิํ กีฬิสฺสามิ รมิสฺสามิ ปริจริสฺสามีติ. ยถาภูตทสฺสนนฺติ รุจิวิมุตฺติ ราโค ราค\-ปริวตฺตกา ทิฏฺฐิรูปนา ปสฺสนา อสนฺตุสฺสิตสฺส สีลพฺพต\-ปรามาโส. ตตฺถ กตโม สุทฺธสฺส สีลพฺพต\-ปรามาโส, อิเธกจฺโจ สีลํ ปรามสติ, สีเลน สุชฺฌติ, สีเลน นียติ, สีเลน มุจฺจติ, สุขํ วีติกฺกมติ, ทุกฺขํ วีติกฺกมติ, สุขทุกฺขํ วีติกฺกมติ อนุปาปุณาติ อุปริเมน. ตทุภยํ สีลวตํ ปรามสติ ตทุภเยน สีลวเตน สุชฺฌนฺติ มุจฺจนฺติ นียนฺติ, สุขํ วีติกฺกมนฺติ, ทุกฺขํ วีติกฺกมนฺติ, สุขทุกฺขํ วีติกฺกมนฺติ, อนุปาปุณนฺตีติ อวิสุจิกรํ ธมฺมํ อวิมุตฺติกรํ ธมฺมํ วิสุจิโต วิมุตฺติโต ปจฺจาคจฺฉนฺตสฺส ยา ตถาภูตสฺส ขนฺติ รุจิ มุตฺติ เปกฺขนา อาการ\-ปริวิตกฺโก ทิฏฺฐินิ\-ชฺฌายนา ปสฺสนา, อยํ สุทฺธสฺส สีลพฺพต\-ปรามาโส, เอเต อุโภ ปรามาสา อริยสาวกสฺส \csromanfolio{257}ปหีนา ภวนฺติ ยาว อายติํ อนุปฺปาทธมฺมา, โส สีลวา ภวติ อริยกนฺเตหิ สีเลหิ สมนฺนาคโต อกฺขณฺเฑหิ ยาว อุปสม\-สํวตฺตนิเกหิ. อิเมสํ ติณฺณํ สํโยชนานํ ปหานา สุตวา อริยสาวโก ภวติ โสตาปนฺโน อวินิปาต\-ธมฺโม, สพฺพํ.}

\prose{สหสจฺจาภิสมยา, อิติ โก วจนตฺโถ. จตฺตาโร อภิสมยา, ปริญฺญา\-ภิสมโย ปหานา\-ภิสมโย สจฺฉิกิริยาภิ\-สมโย ภาวนาภิ\-สมโย.}

\prose{ตตฺถ อริยสาวโก ทุกฺขํ ปริญฺญา\-ภิสมเยน อภิสเมติ, สมุทยํ ปหานา\-ภิสมเยน อภิสเมติ, นิโรธํ สจฺฉิกิริยาภิ\-สมเยน อภิสเมติ, มคฺคํ ภาวนาภิ\-สมเยน อภิสเมติ. กิํ การณํ, ทุกฺขสฺส ปริญฺญา\-ภิสมโย, สมุทยสฺส ปหานา\-ภิสมโย, นิโรธสฺส สจฺฉิกิริยาภิ\-สมโย, มคฺคสฺส ภาวนาภิ\-สมโย. สมถ\-วิปสฺสนาย กถํ อภิสเมติ, อารมฺมเณ จิตฺตํ อุปนิพนฺเธตฺวา ปญฺจกฺขนฺเธ ทุกฺขโต ปสฺสติ. ตตฺถ โย อุปนิพนฺโธ, อยํ สมโถ. ยา ปริโยคาหนา, อยํ วิปสฺสนา. ปญฺจกฺขนฺเธ ทุกฺขาติ ปสฺสโต โย ปญฺจ\-กฺขนฺเธสุ อาลโย นิกนฺติ อุปคมนํ อชฺโฌสานา อิจฺฉา มุจฺฉา ปณิธิ ปตฺถนา ปหียติ. ตตฺถ ปญฺจกฺขนฺธา ทุกฺขํ. โย ตตฺถ อาลโย นิกนฺติ อุปคมนํ อชฺโฌสานํ อิจฺฉา มุจฺฉา ปณิธิ ปตฺถนา, อยํ สมุทโย. ยํ ตสฺส ปหานํ, โส นิโรโธ สมโถ วิปสฺสนา จ มคฺโค, เอวํ เตสํ จตุนฺนํ อริยสจฺจานํ เอกกาเล เอกกฺขเณ เอกจิตฺเต อปุพฺพํ อจริมํ อภิสมโย ภวติ. เตนาห ภควา “สหสจฺจาภิสมยา อริยสาวกสฺส ตีณิ สํโยชนานิ ปหียนฺตี”ติ.}

\proseitem{70}{ตตฺถ สมถ\-วิปสฺสนา ยุคนทฺธา วตฺตมานา เอกกาเล เอกกฺขเณ เอกจิตฺเต จตฺตาริ กิจฺจานิ กโรติ, ทุกฺขํ ปริญฺญา\-ภิสมเยน อภิสเมติ, ยาว มคฺคํ ภาวนาภิ\-สมเยน อภิสเมติ. กิํ การณา, ทุกฺขํ ปริญฺญา\-ภิสมโย, ยาว มคฺคํ ภาวนาภิ\-สมโย. เอวํ ทิฏฺฐนฺโต ยถา นาวา ชลํ คจฺฉนฺตี จตฺตาริ กิจฺจานิ กโรติ, ปาริมํ ตีรํ ปาเปติ, โอริมํ ตีรํ ชหติ, ภารํ วหติ, โสตํ ฉินฺทติ. เอวเมว สมถ\-วิปสฺสนา ยุคนทฺธา วตฺตมานา เอกกาเล เอกกฺขเณ เอกจิตฺเต จตฺตาริ กิจฺจานิ \csromanfolio{258}กโรติ, ทุกฺขํ ปริญฺญา\-ภิสมเยน อภิสเมติ, ยาว มคฺคํ ภาวนาภิ\-สมเยน อภิสเมติ. ยถา วา สูริโย อุทยนฺโต เอกกาเล อปุพฺพํ อจริมํ จตฺตาริ กิจฺจานิ กโรติ, อนฺธการํ วิธมติ, อาโลกํ ปาตุกโรติ, รูปํ นิทสฺสียติ, สีตํ ปริยาทิยติ. เอวเมว สมถ\-วิปสฺสนา ยุคนทฺธา วตฺตมานา เอกกาเล ฯเปฯ. ยถา ปทีโป ชลนฺโต เอกกาเล อปุพฺพํ อจริมํ จตฺตาริ กิจฺจานิ กโรติ, อนฺธการํ วิธมติ, อาโลกํ ปาตุกโรติ, รูปํ นิทสฺสียติ, อุปาทานํ ปริยาทิยติ. เอวเมว สมถ\-วิปสฺสนา ยุคนทฺธา วตฺตมานา เอกกาเล ฯเปฯ.}

\prose{ยทา อริยสาวโก โสตาปนฺโน ภวติ อวินิปาต\-ธมฺโม นิยโต ยาว ทุกฺขสฺสนฺตํ กโรติ, อยํ ทสฺสนภูมิ. โสตาปตฺติผลญฺจ โสตาปตฺติผเล ฐิโต อุตฺตริ สมถ\-วิปสฺสนํ ภาเวนฺโต ยุคนทฺธา วตฺตมานา กามราค\-พฺยาปาทานํ เยภุยฺเยน ปหานา อริยสาวโก โหติ. สกทาคามิ ปรินิฏฺฐิตตฺตา สกิเทว อิมํ โลกํ อาคนฺตฺวา ทุกฺขสฺสนฺตํ กาโรติ, อยํ ตนุภูมิ.}

\prose{สกทาคามิ\-ผลญฺจ โย สกทาคามิ\-ผเล ฐิโต วิปสฺสนํ ภาเวนฺโต กามราค\-พฺยาปาเท สานุสเย อนวเสสํ ปชหติ, กามราค\-พฺยาปาเทสุ อนวเสสํ ปหีเนสุ ปญฺโจ\-รมฺภาคิยานิ สํโยชนานิ ปหีนานิ ภวนฺติ สกฺกายทิฏฺฐิ สีลพฺพต\-ปรามาโส วิจิกิจฺฉา กามจฺฉนฺโท พฺยาปาโท จ, อิเมสํ ปญฺจนฺนํ โอรมฺภาคิยานํ สํโยชนานํ ปหานา\csromansharedfootnote{21da85ee56fc8bf6}{ปหานาย (อิ, ก)} อริยสาวโก โหติ อนาคามี ตตฺถ ปรินิพฺพายี อนาวตฺติธมฺโม ตสฺมา โลกา, อยํ วีตราคภูมิ.}

\prose{อนาคามิผลญฺจ อนาคามิผเล ฐิโต อุตฺตริ สมถ\-วิปสฺสนํ ภาเวนฺโต ปญฺจ อุทฺธมฺ\-ภาคิยานิ สํโยชนานิ ปชหติ รูปราค- อรูปราคมานอุทฺธจฺจอวิชฺชญฺจ. อิเมสํ ปญฺจนฺนํ อุทฺธมฺ\-ภาคิยานํ สํโยชนานํ ปหานา อริยสาวโก อรหา ภวติ, ขีณาสโว วุสิตวา สมฺมทญฺญา\csromansharedfootnote{2b28179a7dfee4c5}{สมฺปชญฺโญ (อิ, ก)} วิมุตฺโต ปริกฺขีณ\-ภวสํโยชโน อนุปฺปตฺต\-สทตฺโถ, อยํ กตาภูมิ.}

\prose{อรหนฺโตว อยํ โสปาทิเสสา นิพฺพานธาตุ. ตสฺส อายุกฺขยา ชีวิตินฺทฺริยาปโรธา อิทญฺจ ทุกฺขํ นิรุชฺฌติ, อญฺญญฺจ ทุกฺขํ น อุปฺปชฺชติ. โย อิมสฺส ทุกฺขสฺส นิโรโธ วูปสโม, อญฺญสฺส จ อปาตุภาโว. \csromanfolio{259}อยํ อนุปาทิเสสา นิพฺพานธาตุ. อิมา เทฺว นิพฺพานธาตุโย. อิติ สจฺจานิ วุตฺตานิ. สจฺจาภิสมโย วุตฺโต, กิเลส\-ววตฺถานํ วุตฺตํ, ปหานํ วุตฺตํ, ภูมิโย วุตฺตา, ผลานิ วุตฺตานิ, นิพฺพานธาตุโย วุตฺตา. เอวมิเมสุ วุตฺเตสุ สพฺพโพธิ วุตฺตา ภวติ. เอตฺถ โยโค กรณีโย.}

\proseitem{71}{ตตฺถ กตมาโย นว อนุปุพฺพ\-สมาปตฺติโย, จตฺตาริ ฌานานิ จตสฺโส จ อรูป\-สมาปตฺติโย นิโรธ\-สมาปตฺติ จ. ตตฺถ จตฺตาริ ฌานานิ กตมานิ, อิธ ภิกฺขเว\csromansharedfootnote{fcc498ad2cc35e18}{ที 1. 69 ปิฏฺเฐ.} ภิกฺขุ วิวิจฺเจว กาเมหีติ วิตฺถาเรน กาตพฺพานิ. ตตฺถ กตมา จตฺตาโร อรูป\-สมาปตฺติโย, วิราคิโน วต วตฺตพฺโพ, ยาว นิโรธ\-สมาปตฺติ วิตฺถาเรน กาตพฺพา. อิมาโย นว อนุปุพฺพ\-สมาปตฺติโย.}

\prose{ตตฺถ กตมํ ปฐมํ ฌานํ, ปญฺจงฺค\-วิปฺปยุตฺตํ ปญฺจ\-งฺค\-สมนฺนาคตํ, กตเมหิ ปญฺจหิ องฺเคหิ วิปฺปยุตฺตํ ปญฺจหิ นีวรเณหิ. ตตฺถ กตมานิ ปญฺจ นีวรณานิ, กามจฺฉนฺโทติ วิตฺถาเรตพฺโพ. ตตฺถ กตโม กามจฺฉนฺโท, โย ปญฺจสุ กามคุเณสุ ฉนฺทราโค เปมํ นิกนฺติ อชฺโฌสานํ อิจฺฉา มุจฺฉา ปตฺถนา อปริจฺจาโค อนุสโย ปริยุฏฺฐานํ, อยํ กามจฺฉนฺท\-นีวรณํ. ตตฺถ กตมํ พฺยาปาท\-นีวรณํ, โย สตฺเตสุ สงฺขาเรสุ จ อาฆาโต ฯเปฯ ยถา โทเส ตถา นิ- โอฏฺฐานา, อยํ พฺยาปาโท นีวรณํ. ตตฺถ กตมํ มิทฺธํ, ยา จิตฺตสฺส ชฬตา จิตฺตสฺส ครุตฺตํ จิตฺตสฺส อกมฺมนียตา จิตฺตสฺส นิกฺเขโป นิทฺทายนา ปจลิกตา ปจลายนา ปจลายนํ, อิทํ มิทฺธํ. ตตฺถ กตมํ ถินํ\csromansharedfootnote{bdaf47f2a24a8aba}{ถีนํ (อิ)}, ยา กายสฺส ถินตา ชฬตา กายสฺส ครุตฺตา กายสฺส อปฺปสฺสทฺธิ, อิทํ ถินํ. อิติ อิทญฺจ ถินํ ปุริมกญฺจ มิทฺธํ ตทุภยํ ถินมิทฺธ\-นีวรณนฺติ วุจฺจติ. ตตฺถ กตมํ อุทฺธจฺจํ, โย อวูปสโม จิตฺตสฺส, อิทํ อุทฺธจฺจํ. ตตฺถ กตมํ กุกฺกุจฺจํ, โย เจตโส วิเลโข อลญฺจนา วิลญฺจนา หทยเลโข วิปฺปฏิสาโร, อิทํ กุกฺกุจฺจํ. อิติ อิทญฺจ กุกฺกุจฺจํ ปุริมกญฺจ อุทฺธจฺจํ ตทุภยํ อุทฺธจฺจกุกฺกุจฺจ\-นีวรณนฺติ วุจฺจติ. ตตฺถ กตมํ วิจิกิจฺฉา\-นีวรณํ, โย พุทฺเธ วา ธมฺเม วา สํเฆ วา ฯเปฯ อยํ วิจิกิจฺฉา. อปิ จ โข ปน ปญฺจ วิจิกิจฺฉาโย สมนนฺตรายิกา เทสนฺตรายิกา สมาปตฺต\-นฺตรายิกา มคฺคนฺตรายิกา สคฺคนฺตรายิกา, อิมาโย ปญฺจ \csromanfolio{260}วิจิกิจฺฉาโย. อิธ ปน สมาปตฺต\-นฺตรายิกา วิจิกิจฺฉา อธิปฺเปตา. อิเม ปญฺจ นีวรณา.}

\prose{ตตฺถ นีวรณานีติ โก วจนตฺโถ, กุโต นิวารยนฺตีติ, สพฺพโต กุสลปกฺขิกา นิวารยนฺติ. กถํ\csromansharedfootnote{cc63ea560f1e8073}{กิํ ตํ (อิ, ก)} นิวารยนฺติ, กามจฺฉนฺโท อสุภโต นิวารยติ, พฺยาปาโท เมตฺตาย\csromansharedfootnote{3bda306813569e8b}{เมตฺตโต (อิ)} นิวารยติ, ถินํ ปสฺสทฺธิโต นิวารยติ, มิทฺธํ วีริยารมฺภโต นิวารยติ, อุทฺธจฺจํ สมถโต นิวารยติ, กุกฺกุจฺจํ อวิปฺปฏิสารโต นิวารยติ, วิจิกิจฺฉา ปญฺญาโต ปฏิจฺจ\-สมุปฺปาทโต นิวารยติ.}

\prose{อปโร ปริยาโย, กามจฺฉนฺโท อโลภโต กุสลมูลโต นิวารยติ, พฺยาปาโท อโทสโต นิวารยติ, ถินมิทฺธํ สมาธิโต นิวารยติ, อุทฺธจฺจ\-กุกฺกุจฺจํ สติปฏฺฐาเนหิ นิวารยติ, วิจิกิจฺฉา อโมหโต กุสลมูลโต นิวารยติ.}

\prose{อปโร ปริยาโย, ตโย วิหารา ทิพฺพวิหาโร พฺรหฺมวิหาโร อริยวิหาโร. ทิพฺพวิหาโร จตฺตาริ ฌานานิ, พฺรหฺมวิหาโร จตฺตาริ อปฺปมาณานิ, อริยวิหาโร สตฺตติํส โพธิปกฺขิยา ธมฺมา. ตตฺถ กามจฺฉนฺโธ อุทฺธจฺจํ กุกฺกุจฺจญฺจ ทิพฺพวิหารํ นิวารยติ, พฺยาปาโท พฺรหฺมวิหารํ นิวารยติ, ถินมิทฺธํ วิจิกิจฺฉา จ อริยวิหารํ นิวารยติ.}

\prose{อปโร ปริยาโย, กามจฺฉนฺโท พฺยาปาโท อุทฺธจฺจ\-กุกฺกุจฺจญฺจ สมถํ นิวารยนฺติ, ถินมิทฺธํ วิจิกิจฺฉา จ วิปสฺสนํ นิวารยนฺติ, อโต นีวรณนฺติ วุจฺจนฺเต. อิเมหิ ปญฺจหิ องฺเคหิ วิปฺปยุตฺตํ ปฐมํ ฌานํ.}

\prose{กตเมหิ ปญฺจหิ องฺเคหิ สมฺปยุตฺตํ ปฐมํ ฌานํ, วิตกฺกวิจาเรหิ ปีติยา สุเขน จ จิตฺเตกคฺคตาย จ. อิเมสํ ปญฺจนฺนํ องฺคานํ อุปฺปาท\-ปฏิลาภ\-สมนฺนาคโม สจฺฉิกิริยํ ปฐมํ ฌานํ ปฏิลทฺธนฺติ วุจฺจติ. อิมานิ ปญฺจ องฺคานิ อุปฺปาเทตฺวา วิหรตีติ, เตน วุจฺจเต ปฐมํ ฌานํ อุปสมฺปชฺช วิหรตีติ ทิพฺเพน วิหาเรน. (1)}

\prose{ตตฺถ ทุติยํ ฌานํ จตุรงฺคสมนฺนาคตํ ปีติสุเขน จิตฺเตกคฺคตาย อชฺฌตฺตํ สมฺปสาทเนน อิมานิ จตฺตาริ องฺคานิ อุปฺปาเทตฺวา สมฺปาเทตฺวา วิหรติ, เตน วุจฺจติ ทุติยํ ฌานํ อุปสมฺปชฺช วิหรตีติ. (2)}

\csromanmatikaanchor{mtk.46}
\csromantocmarkref{section}{7. หารสมฺปาตภูมิ}{mtk.46}
\bookmark[dest={mtk.46},level=1]{7. หารสมฺปาตภูมิ}
\prose{\csromanfolio{261}ตตฺถ ปญฺจ\-งฺค\-สมนฺนาคตํ ตติยํ ฌานํ สติยา สมฺปชญฺเญน สุเขน จิตฺเตกคฺคตาย อุเปกฺขาย อิมานิ ปญฺจงฺคานิ อุปฺปาเทตฺวา สมฺปาเทตฺวา วิหรติ, เตน วุจฺจติ ตติยํ ฌานํ อุปสมฺปชฺช วิหรตีติ. (3)}

\prose{ตตฺถ จตุตฺถํ ฌานํ จตุรงฺคสมนฺนาคตํ อุเปกฺขาย สติปาริสุทฺธิยา อทุกฺขม\-สุขาย เวทนาย จิตฺเตกคฺคตา จ, อิเมหิ จตูหงฺเคหิ สมนฺนาคตํ จตุตฺถํ ฌานํ. อิติ อิเมสํ จตุนฺนํ องฺคานํ อุปฺปาโท ปฏิลาโภ สมนฺนาคโม สจฺฉิกิริยา จตุตฺถํ ฌานํ ปฏิลทฺธนฺติ วุจฺจติ. อิมานิ จตฺตาริ ฌานานิ อุปฺปาเทตฺวา สมฺปาเทตฺวา อุปสมฺปชฺช วิหรติ, เตน วุจฺจติ ทิพฺเพน วิหาเรน วิหรตีติ. (4)}

\prose{ตตฺถ กตโม อนิจฺจฏฺโฐ, ปีฬนฏฺโฐ อนิจฺจฏฺโฐ ปภงฺคฏฺโฐ สมฺปาปนฏฺโฐ วิเวกฏฺโฐ อนิจฺจฏฺโฐ, อยํ อนิจฺจฏฺโฐ.}

\prose{ตตฺถ กตโม ทุกฺขฏฺโฐ, ปีฬนฏฺโฐ ทุกฺขฏฺโฐ สมฺปีฬนฏฺโฐ สํเวคฏฺโฐ พฺยาธินฏฺโฐ, อยํ ทุกฺขฏฺโฐ.}

\prose{ตตฺถ กตโม สุญฺญฏฺโฐ, อนุปลิตฺโต สุญฺญฏฺโฐ, อสมฺภชนฏฺโฐ คตปฏฺโฐ\csromansharedfootnote{22bd7b8ba2502367}{อปฺปฏฺโถ (อิ)} วิวฏฺฏฏฺโฐ, อยํ สุญฺญฏฺโฐ.}

\prosewithcloser{ตตฺถ กตโม อนตฺตฏฺโฐ, อนิสฺสริยฏฺโฐ อนตฺตฏฺโฐ, อวสวตฺต\-นฏฺโฐ, อกามการิฏฺโฐ ปริวิทฏฺโฐ, อยํ อนตฺตฏฺโฐติ.}{สุตฺตตฺถ\-สมุจฺจโย นาม สํวตฺติ\-สนฺติกา เปฏกภูมิ สมตฺตา.}
\proseitem{72}{ฌานํ วิราโค. จตฺตาริ ฌานานิ วิตฺถาเรน กาตพฺพานิ. ตานิ ทุวิธานิ โพชฺฌงฺค\-วิปฺปยุตฺตานิ จ โพชฺฌงฺค\-สมฺปยุตฺตานิ จ. ตตฺถ โพชฺฌงฺค\-วิปฺปยุตฺตานิ พาหิรกานิ, โพชฺฌงฺค\-สมฺปยุตฺตานิ อริยปุคฺคลานิ. ตตฺถ เยน ฉ ปุคฺคลมูลานิ เตสํ นิกฺขิเปตฺวา ราคจริโต, โทสจริโต, โมหจริโต, ราคโทสจริโต, ราคโมหจริโต, โทสโมหจริโต, สมภาค\-จริโต, อิติ อิเมสํ ปุคฺคลานํ ฌานํ สมาปชฺชิตานํ ปญฺจ นีวรณานิ ปฏิปกฺโข เตสํ ปฏิฆาตาย ยถา อสมตฺโถ ตีณิ \csromanfolio{262}อกุสลมูลานิ นิคฺคณฺหาติ. โลเภน อกุสลมูเลน อภิชฺฌา จ อุทฺธจฺจญฺจ อุปฺปิลวตํ อโลเภน กุสลมูเลน นิคฺคณฺหาติ, กุกฺกุจฺจญฺจ วิจิกิจฺฉา จ โมหปกฺโข, ตํ อโมเหน นิคฺคณฺหาติ. โทโส จ ถินมิทฺธญฺจ โทสปกฺโข, ตํ อโทเสน นิคฺคณฺหาติ.}

\prose{ตตฺถ อโลภสฺส ปาริปูริยา เนกฺขมฺม\-วิตกฺกํ วิตกฺเกติ. ตตฺถ อโทสสฺส ปาริปูริยา อพฺยาปาทวิตกฺกํ วิตกฺเกติ. ตตฺถ อโมหสฺส ปาริปูริยา อวิหิํสา\-วิตกฺกํ วิตกฺเกติ. ตตฺถ อโลภสฺส ปาริปูริยา วิวิตฺโต โหติ กาเมหิ. ตตฺถ อโทสสฺส ปาริปูริยา อโมหสฺส ปาริปูริยา จ วิวิตฺโต โหติ ปาปเกหิ อกุสเลหิ ธมฺเมหิ, สวิตกฺกํ สวิจารํ วิเวกชํ ปีติสุขํ ปฐมํ ฌานํ อุปสมฺปชฺช วิหรติ.}

\prose{วิตกฺกาติ ตโย วิตกฺกา เนกฺขมฺม\-วิตกฺโก อพฺยาปาทวิตกฺโก อวิหิํสา\-วิตกฺโก. ตตฺถ ปฐมภินิปาโต วิตกฺโก, ปฏิลทฺธสฺส วิจรณํ วิจาโร. ยถา ปุริโส ทูรโต ปุริสํ ปสฺสติ อาคจฺฉนฺตํ, น จ ตาว ชานาติ เอโส อิตฺถีติ วา ปุริโสติ วา ยทา ตุ ปฏิลภติ อิตฺถีติ วา ปุริโสติ วา เอวํ วณฺโณติ วา เอวํ สณฺฐาโนติ วา อิเม วิตกฺกยนฺโต อุตฺตริ อุปปริกฺขนฺติ กิํ นุ โข อยํ สีลวา อุทาหุ ทุสฺสีโล อฑฺโฒ วา ทุคฺคโตติ วา. เอวํ วิจาโร วิตกฺเก อปฺเปติ, วิจาโร จริยติ จ อนุวตฺตติ จ. ยถา ปกฺขี ปุพฺพํ อายูหติ ปจฺฉา นายูหติ ยถา อายูหนา เอวํ วิตกฺโก, ยถา ปกฺขานํ ปสารณํ เอวํ วิจาโร อนุปาลติ วิตกฺเกติ วิจรติ วิจาเรติ. วิตกฺกยติ วิตกฺเกติ, อนุวิจรติ วิจาเรติ. กามสญฺญาย ปฏิปกฺโข วิตกฺโก, พฺยาปาทสญฺญาย วิหิํสสญฺญาย จ ปฏิปกฺโข วิจาโร. วิตกฺกานํ กมฺมํ อกุสลสฺส อมนสิกาโร, วิจารานํ กมฺมํ เชฏฺฐานํ สํวารณา. ยถา ปลิโก ตุณฺหิโก สชฺฌายํ กโรติ เอวํ วิตกฺโก, ยถา ตํเยว อนุปสฺสติ เอวํ วิจาโร. ยถา อปริญฺญา เอวํ วิตกฺโก. ยถา ปริญฺญา เอวํ วิจาโร. นิรุตฺติ\-ปฏิสมฺภิทายญฺจ ปฏิภาน\-ปฏิสมฺภิทายญฺจ วิตกฺโก, ธมฺม\-ปฏิสมฺภิทายญฺจ อตฺถ\-ปฏิสมฺภิทายญฺจ วิจาโร. กลฺลิตา โกสลฺลตฺตํ จิตฺตสฺส วิตกฺโก. อภินีหาร\-โกสลฺลํ จิตฺตสฺส วิจาโร. อิทํ กุสลํ อิทํ อกุสลํ อิทํ ภาเวตพฺพํ อิทํ ปหาตพฺพํ อิทํ สจฺฉิ\-กาตพฺพนฺติ วิตกฺโก, ยถา ปหานญฺจ ภาวนา จ สจฺฉิกิริยา จ เอวํ วิจาโร. อิเมสุ วิตกฺกวิจาเรสุ ฐิตสฺส ทุวิธํ ทุกฺขํ น อุปฺปชฺชติ กายิกญฺจ เจตสิกญฺจ. ทุวิธํ สุขํ อุปฺปชฺชติ \csromanfolio{263}กายิกญฺจ เจตสิกญฺจ. อิติ วิตกฺกชนิตํ เจตสิกํ สุขํ ปีติ กายิกํ สุขํ กายิโกเยว. ยา ตตฺถ จิตฺตสฺส เอกคฺคตา, อยํ สมาธิ. อิติ ปฐมํ ฌานํ ปญฺจ\-งฺค\-วิปฺปหีนํ ปญฺจ\-งฺค\-สมนฺนาคตํ.}

\prose{เตสํเยว วิตกฺก\-วิจารานํ อภิกฺขณํ อาเสวนาย ตสฺส ตปฺโปณมานสํ โหติ. ตสฺส วิตกฺกวิจารา โอฬาริกา ขายนฺติ. ยญฺจ ปีติสุขญฺจ เนกฺขมฺมญฺจ โอฬาริกํ ภวติ. อปิ จ สมาธิชา ปีติ รติ จ ชายติ. ตสฺส วิจารารมฺมณํ. เตสํ วูปสมา อชฺฌตฺตํ เจโต สมฺปสีทติ. เย วิตกฺกวิจารา เทฺว ธมฺมา\-นุสฺสริตพฺพา. ปจฺจุปฺปนฺนา ทรณิตพฺพํ. เตสํ วูปสมา เอโกทิภาวํ จิตฺเตกคฺคตํ โหติ. ตสฺส เอโกทิภาเวน ปีติ ปาริปูริํ คจฺฉติ. ยา ปีติ, ตํ โสมนสฺสิ\-นฺทฺริยํ, ยํ สุขํ, ตํ สุขินฺทฺริยํ. ยา จิตฺเตกคฺคตา, อยํ สมาธิ. ตํ ทุติยํ ฌานํ จตุรงฺคสมนฺนาคตํ. โส ปีติยา วิราคา ยาติ โอชหิ ชลฺลสหคตํ.}

\proseitem{73}{ตตฺถ โสมนสฺสจิตฺตมุปาทานนฺติ จ โส ตํ วิจินนฺโต อุเปกฺขเมว มนสิกโรติ. โส ปีติยา วิราคา อุเปกฺขโก วิหรติ. ยถา จ ปีติยา สุขมานิตํ, ตํ กาเยน ปฏิสํเวเทติ สมฺปชาโน วิหรติ. เยน สติสมฺปชญฺเญน อุเปกฺขา\-ปาริปูริํ คจฺฉติ. อิทํ ตติยํ ฌานํ จตุรงฺคสมนฺนาคตํ.}

\prose{ตถา กายิกสฺส สุขสฺส ปหานาย ปฐเม ฌาเน โสมนสฺสิ\-นฺทฺริยํ นิรุชฺฌติ. ทุติเย ฌาเน ทุกฺขินฺทฺริยํ นิรุชฺฌติ. โส สุขสฺส จ ปหานา ทุกฺขสฺส จ ปหานา ปุพฺเพว โสมนสฺส\-โทมนสฺสานํ อตฺถงฺคมา อทุกฺขมสุขํ อุเปกฺขา\-สติ\-ปาริสุทฺธิํ จตุตฺถํ ฌานํ อุปสมฺปชฺช วิหรติ. ตตฺถ จตูหิ อินฺทฺริเยหิ อุเปกฺขา ปสาทา โหติ, ทุกฺขินฺทฺริเยน โทมนสฺสิ\-นฺทฺริเยน สุขินฺทฺริเยน โสมนสฺสิ\-นฺทฺริเยน จ. เตสํ นิโรธา อุเปกฺขา\-สมฺปชญฺญํ โหติ, ตตฺถ สุขินฺทฺริเยน โสมนสฺสิ\-นฺทฺริเยน จ อสติ โหติ, เตสํ นิโรธา สติมา โหติ, ทุกฺขินฺทฺริเยน โทมนสฺสิ\-นฺทฺริเยน จ อสมฺปชญฺญํ, เตสํ นิโรธา สมฺปชญฺญํ โหติ, อิติ อุเปกฺขาย จ สญฺญา, สโต สมฺปชาโน จิตฺเตกคฺคตา จ อิทํ วุจฺจเต จ จตุตฺถํ ฌานํ.}

\prose{ตตฺถ โย ราคจริโต ปุคฺคโล ตสฺส สุขินฺทฺริยญฺจ โสมนสฺสินฺทฺริยญฺจ. โย โทสจริโต ปุคฺคโล ตสฺส ทุกฺขินฺทฺริยญฺจ โทมนสฺสินฺทฺริยญฺจ. โย โมหจริโต ปุคฺคโล ตสฺส อสติ จ อสมฺปชญฺญญฺจ.}

\prose{\csromanfolio{264}ตตฺถ ราคจริตสฺส ปุคฺคลสฺส ตติเย ฌาเน จตุตฺเถ จ อนุนโย นิรุชฺฌติ, โทสจริตสฺส ปฐเม ฌาเน ทุติเย จ ปฏิฆํ นิรุชฺฌติ, โมหจริตสฺส ปุคฺคลสฺส ปฐเม ฌาเน ทุติเย จ อสมฺปชญฺญํ นิรุชฺฌติ. ตติเย ฌาเน จตุตฺเถ จ อสติ นิรุชฺฌติ, เอวเมว เตสํ ติณฺณํ ปุคฺคลานํ จตฺตาริ ฌานานิ โวทานํ คมิสฺสนฺติ.}

\prose{ตตฺถ ราค\-โทสจริตสฺส ปุคฺคลสฺส อสมฺปชญฺญญฺจ อนุนโย จ ปฏิฆญฺจ, เตน หานภาคิยํ\csromansharedfootnote{deca906aeb4d99ff}{ปหานภาคิยํ (อิ, ก)} ฌานํ โหติ. ตตฺถ ราค\-โมหจริตสฺส ปุคฺคลสฺส อนุนยตฺตํ จ อาทีนวํ ทสฺสิตา, ตํ ตสฺส หานภาคิยํ ฌานํ โหติ. ตตฺถ โทส\-โมหจริตสฺส ปุคฺคลสฺส ปฏิโฆ จ อสติ จ อสมฺปชญฺญญฺจ อาทีนวํ ทสฺสิตา เตน ตสฺส หานภาคิยํ ฌานํ โหติ.}

\prose{ตตฺถ ราคโทสโมห\-สมภาค\-จริตสฺส ปุคฺคลสฺส วิเสสภาคิยํ ฌานํ โหติ, อิมานิ จตฺตาริ ฌานานิ สตฺตสุ ปุคฺคเลสุ นิทฺทิสิตพฺพานิ. จตูสุ จ สมาธีสุ ฉนฺท\-สมาธินา ปฐมํ ฌานํ, วีริย\-สมาธินา ทุติยํ ฌานํ, จิตฺตสมาธินา ตติยํ ฌานํ, วีมํสา\-สมาธินา จตุตฺถํ ฌานํ. อปฺปณิหิเตน ปฐมํ ฌานํ, สุญฺญตาย ทุติยํ ฌานํ, อนิมิตฺเตน ตติยํ ฌานํ, อานาปานสฺสติยา จตุตฺถํ ฌานํ. กามวิตกฺก\-พฺยาปาทานญฺจ ตํ ตํ วูปสเมน ปฐมํ ฌานํ โหติ, วิตกฺก\-วิจารานํ วูปสเมน ทุติยํ ฌานํ, สุขินฺทฺริย\-โสมนสฺสินฺทฺริยานํ วูปสเมน ตติยํ ฌานํ, กาย\-สงฺขารานํ วูปสเมน จตุตฺถํ ฌานญฺจ. จาคา\-ธิฏฺฐาเนน ปฐมํ ฌานํ, สจฺจา\-ธิฏฺฐาเนน ทุติยํ ฌานํ, ปญฺญา\-ธิฏฺฐาเนน ตติยํ ฌานํ, อุปสมา\-ธิฏฺฐาเนน จตุตฺถํ ฌานํ. อิมานิ จตฺตาริ ฌานานิ สงฺเขป\-นิทฺเทเสน นิทฺทิฏฺฐานิ, ตตฺถ สมาธินฺทฺริยํ ปาริปูริํ คจฺฉติ. อนุวตฺตนกานิ จตฺตาริ, ตตฺถ โย ปฐมํ ฌานํ นิสฺสาย อาสวกฺขยํ ปาปุณาติ, โส สุขาย ปฏิปทาย ทนฺธาภิญฺญาย โทมนสฺสินฺทฺริย\-ปฏิปกฺเขน. โย ทุติยํ ฌานํ นิสฺสาย อาสวานํ ขยํ ปาปุณาติ, โส สุขาย ปฏิปทาย ขิปฺปาภิญฺญาย ทุกฺขินฺทฺริย\-ปฏิปกฺเขน. โย ตติยํ ฌานํ นิสฺสาย อาสวานํ ขยํ ปาปุณาติ, โส สุขาย ปฏิปทาย ทนฺธาภิญฺญาย โสมนสฺสินฺทฺริย\-ปฏิปกฺเขน. โย จตุตฺถํ ฌานํ นิสฺสาย อาสวานํ ขยํ ปาปุณาติ, โส สุขาย ปฏิปทาย ขิปฺปาภิญฺญาย สุขินฺทฺริย\-ปฏิปกฺเขน คโต.}

\vspace{\tipitakaparskip}
\csromancenter{ปกิณฺณก\-นิทฺเทโส.}
\proseitem{74}{\csromanfolio{265}ยานิ จตฺตาริ ฌานานิ, เตสํ ฌานานํ อิมานิ องฺคานิ, เตสํ องฺคานํ สมูโห\csromansharedfootnote{a4ce800aa7f56501}{สมฺโมโห (อิ, ก)} อสฺส องฺคา, อยํ ฌานภูมิ โก วิเสโสติ อสฺส วิเสโส. อิเม สมฺภารา เตหิ อยํ สมุทาคโม, ตสฺส สมุทาคมสฺส อยํ อุปนิสา, ตาย อุปนิสาย อยํ ภาวนา. ตสฺสา ภาวนาย อยํ อาทีนโว. เตน อยํ ปริหานิ. กสฺส ปริหานีติ ตทุปค\-ชฺฌายิโน\csromansharedfootnote{354b44103df3712a}{ตทุปกชฺฌายิโน (อิ, ก)}. ตํ ยถา ภณิตํ ปจฺจเวกฺขนฺโต อยํ วิเสโส. เตน วิเสเสน อยํ อสฺสาโท, โส กสฺส อสฺสาโท อฌานิยา ฌายิโน, ตสฺสา อฌานิยา ฌายิโน, อิทํ กลฺลิตา โกสลฺเล ฐิตชฺฌานํ อโนมทฺทิยตํ คจฺฉติ ฌานพลํ, ฌานพเล ฐิตสฺส อยํ ปารมิ\-ปฺปตฺตสฺส อิมานิ ฌานงฺคานิ อนาวิล\-สงฺกปฺโป ปฐเม ฌาเน ฌานงฺคานิ ภาวี. โส ปีติ ตทนุสา\-ริตฺตาว ปฐเม ฌาเน ฌานงฺคํ ตสฺสงฺคุโน จ ธมฺมา ตทภิสนฺนิ\-ตาย จ. ปีติ ทุติเย ฌาเน ฌานงฺค\-ธมฺมตา โข ปน ตถา ปวตฺตสฺส สหคตํ ฌานงฺค\-ธมฺมํ สสุขตาย อชฺฌตฺตํ สมฺปสาโท ทุติเย ฌาเน ฌานงฺคํ มโน\-สมฺปสาท\-นตาย ตทภิสนฺนิ\-ตาย จ. ปีติ ทุติเย ฌาเน ฌานงฺคํ อชฺฌาตฺตํ สมฺปสาทนํ สมาธิตา\csromansharedfootnote{b9bdcab15e780d03}{สมาธิกา (อิ)} ปีติ ทุติเยฌาเนฌานงฺคํ, เจตโส เอโกทิภาโว ทุติเย ฌาเน ฌานงฺคํ, อุเปกฺขา ผสฺสตา ตติเย ฌาเน ฌานงฺคํ, สุขํ ตสฺส องฺคนฺติ จ. เจตโส เอโกทิภาโว จตุตฺเถ ฌาเน ฌานงฺคํ, อุเปกฺขา อทุกฺขมสุขา จตุตฺเถ ฌาเน ฌานงฺคํ, อภินิสาภูมิ อุเปกฺขา\-สติ\-ปาริสุทฺธิ จตุตฺเถ ฌาเน ฌานงฺคํ. สติปาริสุทฺธิ จ อเนก\-ชฺฌา\-ภูมีสุ ฌานงฺค\-สมา\-ยุตฺตา ปีติ เจตโส เอโกทิภาโว จตุตฺเถ ฌาเน ฌานงฺคํ.}

\prose{ตตฺถ กตมา ฌานภูมิ. สวิตกฺเก สวิจาเร วิเวกา อนุคตา ปฐเม ฌาเน ฌานภูมิ. อวิตกฺเก อวิจาเร อชฺฌตฺตํ สมฺปสาทนํ ชนิตํ ปีติมนุคตา ทุติเย ฌาเน ฌานภูมิ. สุขสาต\-สโมหิตา สปฺปีติกา ตติเย ฌาเน ฌานภูมิ. ตสฺส สุขทุกฺข\-สหคตา อภินีหาร\-สหคตา จตุตฺเถ ฌาเน ฌานภูมิ. อปฺปมาณ\-สหคตา สตฺตารมฺมณา ปฐเม ฌาเน ฌานภูมิ. อภิภูมิอายตนสหคตา รูปสญฺญีสุ ทุติเย ฌาเน ฌานภูมิ. วิโมกฺข\-สหคตานํ วิโมกฺเขสุ ตติเย ฌาเน ฌานภูมิ. อนุปสฺสนา\-สหคตา กายสงฺขารา สมฺมา จตุตฺถสฺส ฌานสฺส ภูมิ.}

\proseitem{75}{\csromanfolio{266}ตตฺถ กตเม ฌานวิเสสา. วิวิจฺเจว กาเมหิ วิวิจฺจ ปาปเกหิ อกุสเลหิ ธมฺเมหิ จิตฺตเจตสิก\-สหคตา กามธาตุ\-สมติกฺกมนตาปิ, อยํ ฌานวิเสโส. อวิตกฺกา เจว อวิจารา จ สปฺปีติกาย สติสหคตาย ปีติสหคตา สญฺญา\-มนสิการา สมุทาจรนฺติ. อยํ ฌานวิเสโส. อวิตกฺกาย ภูมิยา อวิจาเรเยว สติ อนุคตา อุเปกฺขา\-สหคตา มนสิการา สมุทาจรนฺติ. ตทนุธมฺมตาย จ สติ สณฺฑหติ\csromansharedfootnote{f7cc80a392b8498c}{สนฺทหติ (อิ)}. ตญฺจ ภูมิํ อุปสมฺปชฺช วิหรติ, อยํ ฌานวิเสโส. สติปาริสุทฺธิ\-สหคตา สญฺญา\-มนสิการา สมุทาจรนฺติ, ตญฺจ ภูมิํ อุปสมฺปชฺช วิหรติ, อยํ ฌานวิเสโส. วิญฺญาณญฺจา\-ยตนสหคตาย ภูมิยํ อากิญฺจญฺญายตน\-สหคตา สญฺญา\-มนสิการา สมุทาจรนฺติ, ตญฺจ ภูมิํ อุปสมฺปชฺช วิหรติ, อยํ ฌานวิเสโส.}

\prose{ฌานสมฺภารา เนกฺขมฺม\-วิตกฺโก สมฺภาโร กามวิตกฺก\-วิโนทนา\-ธิปฺปายตา. อพฺยาปาทวิตกฺโก สมฺภาโร พฺยาปาทวิตกฺก\-ปฏิวิโนทนา\-ธิปฺปายตา. อวิหิํสา\-วิตกฺโก สมฺภาโร วิหิํสาวิตกฺก\-ปฏิวิโนทนา\-ธิปฺปายตา. อินฺทฺริเยสุ คุตฺตทฺวารตา อปฺปิจฺฉตา สมฺภาโร ปริสุทฺธาชีโว จตุนฺนํ สมาปตฺตีนํ สมฺภาโร อกมฺมสฺส วิหาริตา. มคฺคสมฺภาโร สมาปตฺติ\-ปชฺช\-นตา. ผลสมฺภาโร ฌาน\-นิพฺพตฺติตาย ฌาน\-สมุทาคโม. กุสลเหตุ ยํ ฌานํ สมุทยํ คจฺฉนฺติ โก จ\csromansharedfootnote{2e64b5c0bcc173a5}{โกจิ (ก)} น กุโตจิ เนกฺขมฺม\-ปฺปตฺตา สมุทา\-ค\-จฺฉนฺติ. อาลมฺพ\-นิโรธ\-สมาธิ สนฺโต สมุทา\-ค\-จฺฉนฺติ. อวีติกฺกนฺตา สมุทา\-ค\-จฺฉนฺติ. สุขินฺทฺริยํ โสมนสฺสิ\-นฺทฺริยํ ปหานาย เต จ อพฺยาปชฺชตาย สมุทา\-ค\-จฺฉนฺติ. ตํ ปน สนฺธาย สมุทา\-ค\-จฺฉนฺติ. อปริทาหนาย สมุทา\-ค\-จฺฉนฺติ. อยํ ญาณสมุทาคโม.}

\proseitem{76}{ตตฺถ กตมา อุปนิสา. กลฺยาณมิตฺตตา ฌานสฺส อุปนิสา. กลฺยาณ\-สมฺปวงฺกตา ฌานสฺส อุปนิสา. อินฺทฺริเยสุ คุตฺตทฺวารตา ฌานสฺส อุปนิสา. อสนฺตุฏฺฐิตา กุสเลสุ ธมฺเมสุ ฌานสฺส อุปนิสา. สทฺธมฺม\-สฺสวนํ ฌานสฺส อุปนิสา. สํเวชนิเย ฐาเน สํวิคฺคสฺส โยนิโส ปธานํ. อยํ ฌาโนปนิสา.}

\prose{\csromanfolio{267}ตตฺถ กตมา ภาวนา, เมตฺตาเสวนา อพฺยาปาทวิตกฺก\-ภาวนา. กรุณาเสวนา อวิหิํสาวิตกฺก\-ภาวนา. มุทิตาภาวนา ปีติสุข\-สมฺปชญฺญา การิตา. อุเปกฺขาภาวนา ปสฺสวตา อุเปกฺขาภาวนา อปสฺสวตา อุเปกฺขา จ อชฺฌุเปกฺขา จ, อสุตสญฺญาภาวนา ทุกฺขาปฏิปทา ทนฺธาภิญฺญา ภวสนฺธา\-ภิญฺญา ภวสนฺธานํ, สา ฉพฺพิธา ภาวนา ภาวิตา พหุลีกตา อนุฏฺฐิตา วตฺถุกตา ยานีกตา ปริจิตา สุสมารทฺธา. อยํ ภาวนา.}

\prose{เอวํ ภาวยนฺตสฺส อยํ อาทีนโว. ปฐเม ฌาเน สงฺขาร\-สมนฺนาคโต เอโส ธมฺโม อสฺสุโต สาสโว. สเจ เอส ธมฺโม อยํ สีโล อาสนฺน\-ปฏิปกฺโข จ เอส ธมฺโม กาโม ปติจาโร ปติวิจาโร สมาปตฺตีนํ จ สพฺโพฬาริโก เอส ธมฺโม วิตกฺกวิจาโร จ. ตตฺต จิตฺตํ โขเภนฺติ, กาโย เจตฺถ กิลมติ, กายมฺหิ เจตฺถ กิลนฺเต จิตฺตํ วิหญฺญติ. อนภินีหารกฺขโมว อภิญฺญานํ อิเม อาทีนวา ปฐเม ฌาเน.}

\prose{ทุติเย ฌาเน อิเม อาทีนวา ปีติผรณ\-สหคโต จ เอโส ธมฺโม, น สมุทาจารสฺเสติ จิตฺตํ. อโสธยํ อุปคโม เจส ธมฺโม อุปคมิ\-ปริสฺสโย\csromansharedfootnote{983dc31be1c9536d}{อุปคมิปริจโย (อิ)} โทมนสฺส\-ปจฺจตฺถิโก เจส ธมฺโม. ตตฺถ ตตฺถ ยุตฺตีนํ ปีติ ปรชฺชโต เจส ธมฺโม ทุกฺกรํ โหติ, อวตฺต\-สนฺตาส\-ภูมิ\-ปริวชฺชยนฺโต จตูสุ ทุกฺขตาสุ เอส ธมฺโม อนุวิทฺธา\-ปน\-สทฺธาย\csromansharedfootnote{720382340e9360b5}{อนุวิทฺธา ปสฺสติยา (อิ)} ทุกฺขตาย จ น ปลิโพธ\-ทุกฺขตาย จ อภิญฺญา\-ทุกฺขตาย จ โรคทุกฺขตาย จ, อิเม อาทีนวา ทุติเย ฌาเน.}

\prose{ตตฺถ กตเม อาทีนวา ตติเย ฌาเน. อุเปกฺขา\-สุข\-สหคตาย ตตฺถ สาตาวีนํ ปญฺจนฺนํ อุเปกฺขาสุขํ ปริวตฺติโต เอส ธมฺโม เตน นิจฺจ\-สญฺญิตานญฺจ ยํ โหติ. ทุกฺโขปนิยํ สุขํ จิตฺตสฺส สงฺโขภตํ อุปาทาย สุขทุกฺขาย คโต สวติ. สุขทุกฺขา\-นุกตญฺจ อุปาทาย อนภิหาร\-กฺขมํ จิตฺตํ โหติ. อภิญฺญาย สจฺฉิกิริยาสุ สพฺเพปิ เจเต ธมฺมา ตีสุ ฌานสมาปตฺตีสุ จตูหิ จ ทุกฺขตาหิ อนุวิทฺธานํ สา ภยา ทุกฺขตาย ปลิโพธ\-ทุกฺขตาย จ อภิญฺญาย ทุกฺขตาย จ อิเม อาทีนวา ตติเย ฌาเน.}

\prose{\csromanfolio{268}ตตฺถ กตเม อาทีนวา จตุตฺเถ ฌาเน, อากิญฺจญฺญา\-สมาปตฺติกา เต ธมฺมา\-นุ\-สมาปตฺติกา เอติสฺสา จ ภูมิยํ สาตานํ พาล\-ปุถุชฺชนานํ อเนกวิธานิ ทิฏฺฐิคตานิ อุปฺปชฺชนฺติ. โอฬาริกา สุขุเมหิ จ รูปสญฺญาหิ อนุวิธานิ เอตานิ ฌานานิ สทา อนุทยเมตฺตาฌานกลานุทนุกลาย สาธารณา, ทุกฺกรา จ สพฺเพ จตฺตาโร มหาสมฺภารา สมุทาคตานิ จ เอตานิ ฌานานิ อญฺญมญฺญํ นิสฺสาย สมุทา\-ค\-จฺฉนฺติ. เอตฺถ สมุทาคตา จ เอเต ธมฺมา น สมตฺตา โหนฺติ. อสมุ\-คฺคหิต\-นิมิตฺตา จ เอเต ธมฺมา ปริหายนฺติ. นิรุชฺฌนฺติ จ เอเต ธมฺมา น อุปาทิยนฺติ นิรุชฺฌงฺคานิ จ, เอเตสํ ธมฺมานํ ฌานานิ นิมิตฺตานิ น ฌาน\-นิมิตฺต\-สญฺญา โวกิรติ. อปฺปฏิลทฺธปุพฺพา จ ฌายีวเสน จ ภวติ\csromansharedfootnote{9dfd576d866de293}{ฌายี จ วเสน จ ภวติ (อิ, ก)}. อิเมหิ อาทีนเวหิ อยํ ฌานปริหานิ.}

\proseitem{77}{นิโรธ\-สมาปตฺติยา อปฏิสงฺขาย อวเสสสญฺญิโน อากิญฺจญฺญายตน\-สหคตา สญฺญา\-มนสิการา สมุทาจรนฺติ, โส นิโรธ\-สมาปตฺติโต ปริหายติ. อาเนญฺชสญฺญิโน อสญฺญายตนํ สมาปนฺนสฺส อากิญฺจญฺญายตน\-สหคตา มนสิการา สมุทาจรนฺติ, ตญฺจ ภูมิํ น ปชานาติ, โส ตโต ปริหายติ. อากิญฺจญฺญา\-ยตนํ สมาปนฺนสฺส วิญฺญาณญฺจา\-ยตนสญฺญา มนสิการา สมุทาจรนฺติ, ตญฺจ ภูมิํ น ปชานาติ, โส ตโต ปริหายติ. วิญฺญาณญฺจา\-ยตนํ สมาปนฺนสฺส รูปสญฺญา\-สหคตา. วิตฺถาเรน ฯเปฯ. ยาว ปฐเม ฌาเน กามสญฺญา\-สหคตา กาตพฺพา. สกสฺส\csromansharedfootnote{bd05cfe189b6ee58}{สา กสฺส (อิ, ก)} ปริหายติ, กลงฺกชฺฌาเน กลงฺกํ ฌายติ, ปริสมนฺตโต ฌายติ, ภินฺทนฺโต ฌายติ, น สชฺฌายติ, อายูหนฺโต ฌายติ, กิญฺจิ จ นิปริจิโต ฌายติ. อติวิธาวนฺโต ฌายติ, อติมญฺญนฺโต ฌายติ, กายสงฺขาเร อปฺปฏิสมฺภาเร ฌายติ, ปริยุฏฺฐานสฺส นิสฺสรณํ อชานนฺโต ฌายติ, นีวรณา\-ภิภูโต ฌายติ, อสฺสา\-ปตฺติ\-มนสิกโรนฺโต ฌานสฺส อสฺสาโท กามราค\-ปริยุฏฺฐานํ ปหานํ ฌานสฺส อสฺสาโท กามราค\-เหตูนํ ธมฺมานํ อุทยนฺติ, นิรุชฺฌงฺคานิ เอเตสํ ธมฺมานํ ฌานานิ อุปริมา สุขุเปกฺขา กาม\-กมฺมกิเลสานํ ปหานํ อสฺสาโท, เอวํ โข ปุน ฌานสฺส อสฺสาโท มหาสํวาสมปฺปีฬิเต โลกสํนิวาเส อสมฺโพโธกาสา วิคเมสฺสมิทํ ฌานปฺปหานา. อยํ ปลิโรธมปฺปลิโรธโลกสนฺนิวาเส เอสนิธมิทํ ฌานํ อนมตคฺคสํสาร\-สมาปนฺนานํ สตฺตานํ สํสาร\-ปฺปหา\-นนา อานิสํโส, ยมิทํ \csromanfolio{269}ฌานสฺส อสฺสาโท กายสฺส อฌานิยฌายิโน ภวติ. อฌานิยฌานิยฌายีหิ อปรามสนฺโต อฌานิยฌายิตํ ฌายติ, ยานิ กลงฺก\-ชฺฌายิโน ปทานิ, ตานิ อนุธิตานิ ปฏิปกฺเข.}

\proseitem{78}{ตตฺถ กตมํ ฌานโกสลฺลํ, สมาปตฺติ\-โกสลฺลํ ฌานโกสลฺลํ, ฌาน\-วิเสส\-โกสลฺลํ ฌานโกสลฺลํ, ฌานนฺตริก\-โกสลฺลํ ฌานโกสลฺลํ, สมาปตฺติ\-วุฏฺฐาน\-โกสลฺลํ ฌานโกสลฺลํ, ฌาเน สภาว\-โกสลฺลํ ฌานโกสลฺลํ, ฌาเน อาทีนว\-โกสลฺลํ ฌานโกสลฺลํ, ฌาเน นิสฺสรณ\-โกสลฺลํ ฌานโกสลฺลํ, ฌานผเลน อุปาทาย โกสลฺลํ, ฌานผเลน ปฏิสงฺขาน\-ผเล อปริหาน\-ธมฺมตา นิพฺพตฺติฌาเน จ กีฬิตาปิ วิเสสภาคิยํ ฌานํ ปฏิลพฺภติ. อิทํ ปนสฺสาติ ภวหาริตา จ อารมฺมณา\-นิมิตฺตคฺคาโห อนภินีหารพลํ, จิตฺเตกคฺคตา นิมิตฺตาสุ คติสหิตา สมถพเลน อสํสีทนญฺจ ฌาเน มคฺคผลํ สมถํ ปวตฺเต สมาธิโน อุเปกฺขาปลิปุพฺพาปรนิมิตฺตาสโย ปคฺคาหิโน\csromansharedfootnote{624dd9da7d330bd1}{มคฺคาหิโน (อิ)} สติพลํ ตํ ปวตฺติตานญฺจ วิปสฺสนานํ สมญฺญาพเล.}

\prose{ตตฺถ กตมา ฌานปารมิตา, สุปารมิตา เมตฺตา กาเมสุ สตฺตา กาม\-สงฺค\-สตฺตาติ\csromansharedfootnote{0757f760b67f67a2}{ขุ 1. 171 อุทาเน ปสฺสิตพฺพํ.} ยมฺหิ สุตฺเต เทสนาย โวหาเรน เทฺว สจฺจานิ นิทฺทิฏฺฐานิ, ทุกฺขญฺจ สมุทโย จ, วิจเยน หาเรน เย สํโยชนีเยสุ ธมฺเมสุ วชฺชํ น ปสฺสนฺติ, เต โอฆํ ตริสฺสนฺตีติ เนตํ ฐานํ วิชฺชติ. น ตริสฺสนฺตีติ อตฺถิ เอสา ยุตฺติ จ วิจโย จ อิทํ นุ กิสฺส ปทฏฺฐานํ, กาเมสุ สตฺตาติ ปญฺจ กามคุณา, ตํ กามตณฺหาย ปทฏฺฐานํ. สํโยชเน วชฺชม\-ปสฺสมานาติ อวิชฺชาย ปทฏฺฐานํ, น หิ ชาตุ สํโยชนสงฺค\-สตฺตา โอฆํ ตเรยฺยุํ วิปุลํ มหนฺตนฺติ อุปาทานสฺส ปทฏฺฐานํ. กาเมสุ สตฺตาติ กามา ทฺวิธา วตฺถุกามา จ กิเลสกามา จ, ตตฺถ กิเลสกามา กามตณฺหา กามตณฺหาย ยุตฺตา ภวนฺติ รูปตณฺหา ภวตณฺหา ลกฺขเณน หาเรน, สํโยชเน วชฺชม\-ปสฺสมานาติ สํโยชนสฺส, โย ตตฺถ ฉนฺทราโค ตสฺส กิํ ปทฏฺฐานํ, สุขา เวทนา เทฺว จ อินฺทฺริยานิ สุขินฺทฺริยญฺจ โสมนสฺสินฺทฺริยญฺจ, อิติ สุขาย เวทนาย คหิตาย ตโยปิ เวทนา คหิตา โหนฺติ. เวทนากฺขนฺเธ คหิเต สพฺเพ ปญฺจกฺขนฺธา คหิตา โหนฺติ. รูปสทฺทคนฺธรสโผฏฺฐพฺพา คหิตา, \csromanfolio{270}วตฺถุกาเมสุ คหิเตสุ สพฺพานิ ฉ พาหิรานิ อายตนานิ คหิตานิ โหนฺติ. อชฺฌตฺติก\-พาหิเรสุ อายตเนสุ โย สโต, อยํ วุจฺจเต ลกฺขโณ หาโร, ตตฺถ โย โอฬาริกมฺหิ กิเลเส อชฺฌาวสิโต สพฺพกิเลเสสุ โย น ตโต สุขุมตเรสุ น วีตราโค ภวติ. ตตฺถ พาหิร\-สํโยชนํ มมนฺติ อชฺฌตฺต\-สํโยชนํ อหนฺติ. ตตฺถ ภควโต โก อธิปฺปาโย, เย โอฆํ ตริตุกามา เต สํโยชนีเยสุ ธมฺเมสุ อาทีนวา\-นุปสฺสิโน วิหริสฺสนฺตีติ อยเมตฺถ ภควโต อธิปฺปาโย. กาเมสุ สตฺตาติ เยสุ จ สตฺตา เยน จ สตฺตา เยสญฺจ สตฺตา อยํ จตุพฺพิโธ อากาโร สพฺเพสํ หารภาคิโย.}

\proseitem{79}{ตตฺถ กตมานิ ตีณิ วิปลฺลาสานิ ปทฏฺฐานานิ จ. จิตฺต\-วิปลฺลาสสฺส ทิฏฺฐิ\-วิปลฺลาสสฺส สญฺญา\-วิปลฺลาสสฺส ตโย วิปลฺลาสา ตีณิ อกุสลมูลานิ ปทฏฺฐานํ. ตีณิ อกุสลมูลานิ หีนปฺปณีต\-การิย\-กมฺมสฺส ปทฏฺฐานํ. จตุนฺนญฺจ อุปาทานานํ โทโส อกุสลมูลํ ทิสฺสติ. หีนปฺปณีต\-การิย\-กมฺมสฺส ปทฏฺฐานํ. ยถา มาตุยา วา ปิตุโน วา อญฺญตรสฺส วา ปุน อุฬารสฺส ภิกฺขุโน อภยํ เทติ. ตตฺถ อญฺโญ มิจฺฉา ปฏิปชฺเชยฺย กาเยน วา วาจาย วา. ตตฺถ โส พฺยาปาทมุ\-ปาทาย เตสํ อุฬารานํ รกฺขา\-วรณ\-คุตฺติยา อนุปาลยนฺโต โย อุฬารานํ อภยํ เทติ. เตสํ อภเย ทินฺเน โย ตตฺถ มิจฺฉา ปฏิปชฺเชยฺย. ตตฺถ โส พฺยาปาทํ อุปาทายนฺโต โทสชํ กมฺมํ กโรติ. โย ตตฺถ อสาธุ อินฺทฺริยา นีวรณํ ยํ เตสํ อภยํ ทกฺขิณโต สญฺญํ อิทํ ปณีตํ การณํ มยา ปุน ตตฺถ มิจฺฉา\-ปฏิปตฺติ อยํ พฺยาปาโท หีนคมิ\-ว\-กมฺมํ โลโภ โมโห จ อิมานิ นีวรณานิ วจนานิ ตานิ จตฺตาริ อุปาทานานิ เตหิ จตูหิ อุปาทาเนหิ โย โส อุปาทาโน อิตฺถี วา ปุริโส วา เตสํ ปญฺจ\-กฺขนฺธานํ เตเยว อุปาทาโน สมุทโย อิทํ ทุกฺขญฺจ สมุทโย จ โสเยว เทสนาหาโร.}

\prose{ตตฺถ กาเมสุ เย น ปชฺชนฺติ, เต อาทีนวา\-นุปสฺสนาย ปชฺชนฺติ. อิติสฺสา กามธาตุยา นิกฺขมิตุ\-กามตา, อยํ วุจฺจติ เนกฺขมฺม\-จฺฉนฺโท. โย ตตฺถ อนภิสงฺขารานํ กิญฺจิ วิโสเธติ ตสฺส ธาวรา วา, อยํ อพฺยาปาท\-จฺฉนฺโท. กิญฺจิ วิหิํสติ, อยํ วิหิํสาฉนฺโท. อิติ เนกฺขมฺมา\-ภินีหตา ตโย ฉนฺทา. เนกฺขมฺม\-จฺฉนฺโท อพฺยาปาท\-จฺฉนฺโท อวิหิํสาฉนฺโท. \csromanfolio{271}ตตฺถ เนกฺขมฺม\-จฺฉนฺโท อโลโภ. อพฺยาปาท\-จฺฉนฺโท อโทโส. อวิหิํสาฉนฺโท อโมโห. อิมานิ ตีณิ กุสลมูลานิ อฏฺฐสุ สมฺปตฺเตสุ ปรหิตานิ, เตสํเยว จตุนฺนํ อุปาทานานํ นิโรธาย สํวตฺตนฺติ. สเจ วา ปุน กมฺมํ กเรยฺย กณฺหํ วา สุกฺกํ วา. ตสฺส วิปากหานาย สํวตฺตนฺติ. อิทํ กมฺมํ อกณฺหํ อสุกฺกํ กมฺมกฺขยาย สํวตฺตติ. ตตฺถ โย ติณฺณํ อกุสลมูลานํ นิโรโธ, อยํ นิโรโธ. โสเยว มคฺโค ตตฺถ ปฏิปทานิ อิมานิ เทฺว สจฺจานิ อิมานิ จตฺตาริ สจฺจานิ อาวฏฺโฏ หาโร.}

\prose{กาเมสุ สตฺตาติ เย เสกฺขา, เต เอเกเนวากาเรน สตฺตา. เย ปุถุชฺชนา, เต ทฺวีหากาเรหิ สตฺตา, ตสฺสายํ ปโญฺห วิภชฺช\-พฺยากรณีโย วตฺตพฺโพ. กิญฺจาปิ โสตาปนฺโน ปฏิเสวนาย, โน จ โข อภินิเวเส สตฺโต โย หิ อปจยาย ปทหติ, น อุปจยาย. เสกฺโข หิ กิเลสวเสน กาเม ปฏิเสวติ. ปุถุชฺชโน ปน กิเลส\-สมุฏฺฐานาย กาเม ปฏิเสวติ. ตตฺถ กาเมสุ สตฺตานํ จตุ- โอฆํ ตริสฺสตีติ วิภชฺช\-พฺยากรณีโย, อยํ วิภตฺติ.}

\proseitem{80}{ปริวตฺตโนติ กาเม เย เนว สชฺชนฺติ น จ สํโยชเนหิ สํยุตฺตา, เต โอฆํ ตริสฺสนฺติ วิปุลํ มหนฺตนฺติ. อยํ สุตฺตสฺส ปฏิปกฺโข.}

\prose{เววจนนฺติ โย กาเมสุ สตฺโต โย จ ตตฺถ กามานํ คุโณ, ตตฺถ วิโส สตฺโต. เยปิ กามานํ อาหารา ธมฺมา, ตตฺถ วิโส สตฺโต. ตตฺถิมํ กามานํ เววจนํ ปาโก รโช สลฺลํ คณฺโฑ อีติ อุปทฺทโวติ. ยานิ วา ปน อญฺญานิ เววจนานิ ตตฺถ วิโส สตฺโตติ เววจนํ. สตฺโต พนฺโธ มุจฺฉิโต คธิโต อชฺโฌสิโต กาเม อชฺฌาปนฺนา ปริมุตฺโต ตพฺพหุล\-วิหารีติ. ยานิ วา ปน อญฺญานิ เววจนานิ, อยํเววจโน นาม.  กาม\-ปฺปจาร\-ปญฺญตฺติยา กิเลสโค จรปญฺญตฺติยา ปญฺญตฺตา จิตฺตนฺติ เววจนํ. สตฺโต ตพฺพหุล\-วิหารีติ ยานิ วา ปน อญฺญานิ. อิเม กาม\-ปฺปจาร\-ปญฺญตฺติยา กิเลส\-โคจร\-ปญฺญตฺติยา ปญฺญตฺตา, พีชปญฺญตฺติยา ปญฺญตฺตา, สงฺขารา สํโยชน\-ปญฺญตฺติยา ปญฺญตฺตา, อุปาทานํ เหตุปญฺญตฺติยา ปญฺญตฺตํ, ปุคฺคโล ปุถุปญฺญตฺติยา ปญฺญตฺโต.}

\prose{โอตรโณติ อิมาย ปฏิจฺจ\-สมุปฺปาโท ทุกฺขญฺจ สมุทโย จ. เย กิเลสา เย สงฺขารา สํโยชนานิ จ ปญฺจสุ ขนฺเธสุ สงฺขาร\-กฺขนฺโธ \csromanfolio{272}ธมฺมายตเนสุ อกุสลา ธมฺมายตนานิ อินฺทฺริเยสุ สุขินฺทฺริยญฺจ, โสมนสฺสินฺทฺริยญฺจ, อยํ อินฺทฺริโยตรโณ.}

\prose{โสธโนติ เอตฺตโก. เอเสว อารมฺโภ นิทฺทิสิตพฺโพ สุตฺตตฺโถ.}

\prose{อธิฏฺฐาโนติ อิเม ธมฺมา อตฺถิ เอกตฺตตาย ปญฺญตฺตา อตฺถิ เวมตฺตตาย. เย สญฺญา พาหิโร กาเม, เต เวมตฺตตาย ปญฺญตฺตา. ปญฺจสุ กามคุเณสุ สตฺตาติ ปริยุฏฺฐาน\-วิปลฺลาสา เวมตฺตตาย ปญฺญตฺตา โอฆํ ตเรยฺยุํ. วิปุลํ มหนฺตนฺติ อวิชฺชา เอกตฺตตาย ปญฺญตฺตา.}

\prose{ปริกฺขาโรติ ตสฺส โก เหตุ โก ปจฺจโย. อารมฺมณ\-ปจฺจยตาย ปจฺจโย. อโยนิโส จ มนสิกาโร สนฺนิสฺสยสฺส ปจฺจยตาย ปจฺจโย. อวิชฺชา สมนนฺตร\-ปจฺจยตาย ปจฺจโย. ราคานุสโย เหตุ\-ปจฺจยตาย ปจฺจโย. อยํ เหตุ, อยํ ปจฺจโย.}

\prose{สมาโรปโน ปจฺจโยติ เย กาเมสุ สตฺตา สุคตา สุรูปาติ อยํ กามธาตุยา ฉนฺโท ราโค เต อปุญฺญมยา สงฺขารา. เต กิํ ปจฺจยา, อวิชฺชา ปจฺจยา. เต กิสฺส ปจฺจยา, วิญฺญาณสฺส ปจฺจยา. อิติ อวิชฺชาปจฺจยา สงฺขารา. สงฺขาร\-ปจฺจยา วิญฺญาณํ ยาว ชรามรณํ เอวเมตสฺส เกวลสฺส มหโต ทุกฺข\-กฺขนฺธสฺส สมุทโย โหติ เอกํ สุตฺตํ คตํ.  ปญฺจนีวรณิกํ สุตฺตํ กาตพฺพํ.}

\proseitem{81}{ตตฺถ กตโม เทสนาหาโร นาม ยา จ อภิชฺฌา โย จ พฺยาปาโท ยญฺจ อุทฺธจฺจํ, อยํ ตณฺหา. ยญฺจ ถินมิทฺธํ, ยญฺจ กุกฺกุจฺจํ ยา จ วิจิกิจฺฉา, อยํ ทิฏฺฐิ. ยา ปน กายสฺส อกมฺมนิยตา กิญฺจาปิ ตํ มิทฺธํ โน ตุ สภาว\-กิเลส\-ตาย กิเลโส, อิติ ยา จ จิตฺตสลฺลียนา ยา จ กายา\-กมฺม\-นิยตา, อยํ ปกฺโขปกิเลโส น ตุ สภาวกิเลโส. ตตฺถ อตฺตสญฺญา\-นุป\-จิตฺตํ กิลมโถ กุกฺกุจฺจา\-นุป\-จิตฺตํ ถินํ ยา จิตฺตสฺส ลียนา, อิติ อิเม ปญฺจ นีวรณา จตฺตาริ นีวรณานิ สภาวกิเลสา ถินมิทฺธํ นีวรณ\-ปกฺโข\-ปกิเลโส. ยถา จตฺตาโร อาสวา สภาว- อาสวตาย อาสวา โน ตุ จิตฺต\-สาส\-วตาย อาสวา. สภาวตาย อาสวา. ปกฺเข อาสวตาย อาสวา. อถ ปนาห สุตฺตนฺตํ เยน เต สมฺปยุตฺตา วา วิปฺปยุตฺตา วา อาสวา, เตเยว เอเต วตฺตพฺพา สาสวา วา อนาสวา วา. (1)}

\prose{\csromanfolio{273}ตตฺถ กตโม วิจโย. อภิชฺฌา กามตณฺหา รูปตณฺหา ภวตณฺหา. ยํ วา ปน กิญฺจิ อชฺโฌสานคตํ สาสวา อภิชฺฌิตสฺส เมตฺตานุปสฺสิย โย อนตฺถํ จรติ. ตตฺถ โย พฺยาปาทํ อุปฺปาเทติ, อจริ จริสฺสตีติ. เอวํ นว อาฆาตวตฺถูนิ กตฺตพฺพานิ, ตสฺเสวํ พฺยาปาทา\-นุปสฺสิสฺส กิเลโส โย ปริทาโห กายกิลมโถ อกมฺมนิยตา มิทฺธํ. จิตฺตา\-นุปสฺสิสฺส ปฏิฆาเตน ขิยนา, อิทํ ถินมิทฺธํ. ตตฺถ อธิกรณ\-อวูปสโม, อิทํ อุทฺธจฺจํ. ยํ กิํ กสถมีติ\csromansharedfootnote{ddfb1f21c41da3d8}{กรถมีติ (อิ, ก)} อิทํ กุกฺกุจฺจํ. ยํ ยถา อิทํ สนฺตีรณํ, อยํ วิจิกิจฺฉา. ตตฺถ อวิชฺชา จ ตณฺหา จ อตฺถิ, อิทํ ปริยุฏฺฐานํ. อาวรณํ นีวรณํ ฉทนํ อุปกฺกิเลโส จ อตฺถิ, อิทํ กามจฺฉนฺโท กามราค\-ปริยุฏฺฐานสฺส ปทฏฺฐานํ. พฺยาปาโท พฺยาปาท\-ปริยุฏฺฐานสฺส ปทฏฺฐานํ. ถินมิทฺธํ ถินมิทฺธ\-ปริยุฏฺฐานสฺส ปทฏฺฐานํ. อุทฺธจฺจ\-กุกฺกุจฺจํ อวิชฺชา\-ปริยุฏฺฐานสฺส ปทฏฺฐานํ. วิจิกิจฺฉา วิจิกิจฺฉา\-ปริยุฏฺฐานสฺส ปทฏฺฐานํ. กามราค\-ปริยุฏฺฐานํ อนุสย\-สํโยชนสฺส ปทฏฺฐานํ. พฺยาปาท\-ปริยุฏฺฐานํ ปฏิฆ\-สํโยชนสฺส ปทฏฺฐานํ. ถินมิทฺธ\-ปริยุฏฺฐานํ มาน\-สํโยชนสฺส ปทฏฺฐานํ. อวิชฺชา\-ปริยุฏฺฐานญฺจ วิจิกิจฺฉาปริยุฏฺฐานญฺจ ทิฏฺฐิ\-สํโยชนสฺส ปทฏฺฐานํ. (2)}

\prose{ตตฺถ กตโม ลกฺขโณ หาโร. กามราค\-ปริยุฏฺฐาเน วุตฺเต สพฺพานิ ปริยุฏฺฐานานิ วุตฺตานิ โหนฺตีติ. สํโยชเนสุ วุตฺเตสุ สพฺพ\-สํโยชนานิ วุตฺตานิ โหนฺติ. อยํ ลกฺขโณ หาโร. (3)}

\proseitem{82}{ตตฺถ กตโม จตุพฺยูโห หาโร. เย อิเม ปญฺจ นีวรณา ฌาน\-ปฏิปกฺโข. โส ทุกฺขสมุทโย. ยํ ผลํ, อิทํ ทุกฺขํ. ตตฺถ กามจฺฉนฺทสฺส เนกฺขมฺม\-วิตกฺโก ปฏิปกฺโข. พฺยาปาทสฺส อพฺยาปาทวิตกฺโก ปฏิปกฺโข. ติณฺณํ นีวรณานํ อวิหิํสา\-วิตกฺโก ปฏิปกฺโข, อิติ อิเม ตโย วิตกฺกา. เนกฺขมฺม\-วิตกฺโก สมาธิ\-กฺขนฺธํ ภชติ. อพฺยาปาทวิตกฺโก สีลกฺขนฺธํ ภชติ. อวิหิํสา\-วิตกฺโก ปญฺญากฺขนฺธํ ภชติ. อิเม ตโย ขนฺธา. อริโย อฏฺฐงฺคิโก มคฺโค นีวรณ\-ปฺปหานาย สํวตฺตติ. ยํ นีวรณ\-ปฺปหานํ, อยํ นิโรโธ. อิมานิ จตฺตาริ สจฺจานิ. อยํ จตุพฺยูโห หาโร. (4)}

\prose{ตตฺถ กตโม อาวฏฺโฏ หาโร. ปญฺจ นีวรณานิ ทส ภวนฺติ. ยทปิ อชฺฌตฺตํ สารชฺชติ, ตทปิ นีวรณํ. ยทปิ พหิทฺธา สารชฺชติ, ตทปิ นีวรณํ, \csromanfolio{274}เอวํ ยาว วิจิกิจฺฉา อิเม ทส นีวรณา. อชฺฌตฺต\-พหิทฺธา กิเลสา อิมานิ เทฺว สํโยชนานิ อชฺฌตฺตสํโยชนญฺจ พหิทฺธาสํโยชนญฺจ. ตตฺถ อหนฺติ อชฺฌตฺตํ, มมนฺติ พหิทฺธา. สกฺกายทิฏฺฐิ อชฺฌตฺตํ, เอกสฏฺฐิ ทิฏฺฐิคตานิ พหิทฺธา. โย อชฺฌตฺตํ ฉนฺทราโค รูเปสุ อวีตราโค ภวติ อวีตจฺฉนฺโท. เอวํ ยาว วิญฺญาเณ, อยํ อชฺฌตฺตา ตณฺหา. ยํ ฉสุ พาหิเรสุ อายตเนสุ ตีสุ จ ภเวสุ อชฺโฌสานํ, อยํ พหิทฺธา ตณฺหา. อิมานิ เทฺว สจฺจานิ สํโยชนานิ สํโยชนียา จ ธมฺมา. ตตฺถ สํโยชเนสุ ธมฺเมสุ ยา นิพฺพิทา\-นุปสฺสนา จ, อยํ มคฺโค. ยํ สํโยชน\-ปฺปหานํ, อยํ นิโรโธ. อยํ อาวฏฺโฏ หาโร. (5)}

\prose{ตตฺถ กตโม วิภตฺติหาโร. สํโยชนนฺติ น เอตํ เอกํเสน. มานสํโยชนํ ทิฏฺฐิภาคิยนฺติ น ตํ เอกํเสน อทิฏฺฐมานํ นิสฺสายมานํ น ปชหติ. โย ปญฺจ อุทฺธมฺภาคิโย มาโน กิญฺจาปิ โส ทิฏฺฐิปกฺเข สิยา. น ตุ โอรมฺภาคิยํ สํโยชนํ ตสฺส ปหานาย สํวตฺตตีติ. โย จ อหํกาโร น ปวิทฺโธยํ ปนสฺส เอวํ โหติ. กทาสุ นามาหํ ตํ สนฺตํ อายตนํ สจฺฉิกตฺวา อุปสมฺปชฺช วิหริสฺสามิ, ยํ อริยา สนฺตํ อายตนํ อุปสมฺปชฺช วิหริสฺสนฺตีติ, อยํ อภิชฺฌา น จ ตํ นีวรณํ. อตฺถิ ปน อรหโต กาย\-กิเลส\-มิทฺธญฺจ โอกฺกมติ น จ ตํ นีวรณํ ตสฺส ถินมิทฺธํ นีวรณนฺติ. น เอกํเสน. อยํ วิภตฺติหาโร. (6)}

\prose{ปริวตฺตโนติ ปญฺจ นีวรณา ปญฺจงฺคิเกน ฌาเนน ปหานํ คจฺฉนฺติ. อยํ เตสํ ปฏิปกฺโข นีวรโณ อสุกสฺส ปหีนาติ น อญฺญา\-นุมินิตพฺพํ, ปรมตฺถ\-มชฺฌตฺตํ, อยํ ปริวตฺตนา. (7)}

\prose{ตตฺถ กตโม เววจโน. กามจฺฉนฺโท ฉนฺทราโค เปมํ นิกนฺตีติ เววจนํ. นีวรณํ ฉทนํ อุปกฺกิเลโส ปริยุฏฺฐานนฺติ เววจนํ. (8)}

\prose{ปญฺญตฺตีติ อวิชฺชาปจฺจยา กิจฺจ\-ปญฺญตฺติยา\csromansharedfootnote{715a34203f93c6ae}{ปจฺจาปญฺญตฺติยา (ก)} ปญฺญตฺติ, พฺยาปาโท วิกฺเขป\-ปญฺญตฺติยา ปญฺญตฺติ, ถินมิทฺธํ อสมุ\-คฺฆาต\-ปญฺญตฺติยา ปญฺญตฺติ. เอวํ สพฺเพปิ เอเต ปญฺจ นีวรณา อิมมฺหิ สุตฺเต วิกฺเขป\-ปญฺญตฺติยา ปญฺญตฺติ. (9)}

\prose{ตตฺถ กตโม โอตรโณ. อิเม ปญฺจ นีวรณา อวิชฺชา จ ตณฺหา จ ตตฺถ อวิชฺชามูลา นีวรณา. ยา ตณฺหา อิเม สงฺขารา, เต อวิชฺชาปจฺจยา \csromanfolio{275}อิเม เทฺว ธมฺมา ปญฺจสุ ขนฺเธสุ สงฺขารกฺขนฺธ\-ปริยาปนฺนา, อายตเนสุ ธมฺมายตนํ, ธาตูสุ ธมฺมธาตุ, อินฺทฺริเยสุ อิเมสํ ธมฺมานํ ปทฏฺฐานํ สุขินฺทฺริยสฺส จ โสมนสฺสิ\-นฺทฺริยสฺส จ อิตฺถินฺทฺริยสฺส จ ปุริสิ\-นฺทฺริยสฺส จ. (10)}

\prose{ตตฺถ กตโม โสธโน หาโร. อิทํ สุตฺตํ ยถา อารพฺภ นิกฺขิตฺตํ โส อตฺโถ ภาสิโต อิเมหิ ปญฺจหิ ปเทหิ.(11)}

\prose{ตตฺถ กามจฺฉนฺโท จ พฺยาปาโท จ วิจิกิจฺฉา จ น เอกตฺตตายา ปญฺญตฺตา, กามาติ น เอกตฺตตาย ปญฺญตฺตา, อถ ขลุ เวมตฺตตาย ปญฺญตฺตา. อยํ อธิฏฺฐาโน หาโร. (12)}

\prose{ตตฺถ กตโม ปริกฺขาโร. กามจฺฉนฺทสฺส อโยนิโส มนสิกาโร สุภารมฺมณ\-ปจฺจโย. สุภนิมิตฺตญฺจ เหตุ. พฺยาปาทสฺส อโยนิโส มนสิกาโร อาฆาตวตฺถูนิ จ ปจฺจโย. ปฏิฆานุสโย เหตุ. ถินมิทฺธสฺส ปฏิสํหาโร ปจฺจโย. ปวตฺติยา กิลมถา จลนา ตญฺจ เหตุ. อุทฺธจฺจ\-กุกฺกุจฺจสฺส รชนียํ อารมฺมณิยํ อสฺสาทิยากินฺทฺริยํ ตาว อปริปุณฺณญฺจ ญาณํ ปจฺจโย. กามสญฺญา จ ทิฏฺฐิอนุสโย จ เหตุ. วิจิกิจฺฉาย นว มานวิธา อารมฺมณํ มานานุสโย, โสว ปจฺจโย. วิจิกิจฺฉา\-นุสโย เหตุ. เอเต ปญฺจ ธมฺมา สเหตุ สปฺปจฺจยา อุปฺปชฺชนฺติ. (13)}

\prose{ตตฺถ กตโม สมาโรปโน หาโร. อิเม ปญฺจ นีวรณา จตฺตาโรปิ เอเต อาสวา คณฺฑาปิ\csromansharedfootnote{6d3e52e03f02a0a1}{ตณฺหาปิ (อิ)} เอเต สลฺลาปิ เอเต อุปาทานานิ เอเต. เตสุ เอว พาหิเรสุ ธมฺเมสุ สํกิเลส\-ภาคิยํ สุตฺตนฺติ ปญฺญตฺติํ คจฺฉติ. อยํ สมาโรปโน หาโร. (14) นิทฺทิฏฺฐํ สํกิเลสิก\-ภาคิยํ สุตฺตํ.}

\csromanheader{2}{83. มโนปุพฺพงฺคมา ธมฺมาติ คาถา.}
\prose{ตตฺถ กตโม เทสนา หาโร. อิมมฺหิ สุตฺเต โก อตฺโถ ขนฺธ\-ววตฺถาเนน วิญฺญาณ\-กฺขนฺธํ เทเสติ. ธาตุ\-ววตฺถาเนน มโนวิญฺญาณ\-ธาตุํ. อายตน\-ววตฺถาเนน มนายตนํ, อินฺทฺริยววตฺตาเนน มนินฺทฺริยํ. ตสฺส กิํ ปุพฺพงฺคมา ธมฺมา สํขิตฺเตน ฉ ธมฺมา ปุพฺพงฺคมา ธมฺมา กุสลมูลานิ จ อกุสลมูลานิ จ อนิมิตฺตํ อิมมฺหิ สุตฺเต กุสลมูลํ เทสิตํ.  ตตฺถ กตมา มโนปุพฺพงฺคมา ธมฺมา. มโน เตสํ ปุพฺพงฺคมํ, ยถาปิ พลสฺส ราชา ปุพฺพงฺคโม,}

\prose{\csromanfolio{276}เอวเมว ธมฺมานํ มโนปุพฺพงฺคมา. ตตฺถ ติวิธานํ ปุพฺพงฺคมานํ เนกฺขมฺม\-จฺฉนฺเทน อพฺยาปาท\-จฺฉนฺเทน อวิหิํสา\-ฉนฺเทน. อโลภสฺส เนกฺขมฺม\-จฺฉนฺเทน ปุพฺพงฺคมา. อโทสสฺส อพฺยาปาท\-จฺฉนฺเทน ปุพฺพงฺคมา. อโมหสฺส อวิหิํสา\-ฉนฺเทน ปุพฺพงฺคมา. ตตฺถ \textbf{มโนเสฏฺฐา}ติ มนสา อิเม ธมฺมา อุสฺสฏา มเนน วา นิมฺมิตา. มโนว อิเมสํ ธมฺมานํ เสฏฺโฐติ มโนว อิเมสํ ธมฺมานํ เสฏฺฐเชฏฺโฐติ มโนว อิเมสํ ธมฺมานํ อาธิปจฺจํ กโรตีติ มโนเสฏฺฐา. \textbf{มโนชวา}ติ ยตฺถ มโน คจฺฉติ. ตตฺถ อิเม ธมฺมา คจฺฉนฺตีติ มโนชวา. ยถา วาโต สีฆํ คจฺฉติ อญฺโญ วา โกจิ สีฆํ คามโก วุจฺจเต วาตชโวติ ปกฺขิคามิโกติ, เอวเมว อิเม ธมฺมา มเนน สมฺปชายมานา คจฺฉนฺติ, ตตฺถ อิเม ธมฺมา คจฺฉนฺตีติ มโนชวาติ. เต ติวิธา ฉนฺท\-สมุ\-ทานิตา อนาวิลตา จ สงฺกปฺโป. สตฺตวิธา จ กายิกํ สุจริตํ วาจสิกํ สุจริตํ, เต ทส กุสล\-กมฺม\-ปถา. ตตฺถ \textbf{มนสา เจ ปสนฺเนนา}ติ มโนกมฺมํ. \textbf{ภาสติ วา}ติ วจีกมฺมํ. \textbf{กโรติ วา}ติ กายกมฺมํ. อิเมหิ อิมสฺมิํ สุตฺเต ทส กุสล\-กมฺม\-ปถา ปรมาปิ สนฺตา สีลวตา ปรมา. โส ภวติ วิวตฺติยํ น โลกนิยฺยานาย วาสนาภาคิยํ สุตฺตํ ภวติ. อยํ เทสนา.}

\prose{ตตฺถ กตโม วิจโย หาโร. มโนปุพฺพงฺคมา ธมฺมาติ กุสลมูลานิ จ อฏฺฐงฺค\-สมฺมตฺตานิ. อิทํ สุตฺตํ.}

\prose{ยุตฺตีติ ทสนฺนํ กุสล\-กมฺม\-ปถานํ โย วิปาโก, โส สุขเวทนีโย อพฺยาปาท\-สฺสงฺค\-มาโน. ฉายาว อนปายินีติ อนุคจฺฉติ อตฺถิ เอสา ยุตฺติ.}

\prose{ปทฏฺฐานนฺติ อฏฺฐารสนฺนํ มโนปวิจารานํ ปทฏฺฐานํ. มโนปุพฺพงฺคมา ธมฺมาติ สพฺพ\-กุสล\-ปกฺขสฺส อิเม ธามฺมา ปทฏฺฐานํ. \textbf{มนสา เจ ปสนฺเนนา}ติ โย เจตโส ปสาโท, อิทํ สทฺธิ\-นฺทฺริยสฺส ปทฏฺฐานํ. \textbf{ภาสติ วา}ติ สมฺมาวาจาย. \textbf{กโรติ วา}ติ สมฺมา\-กมฺมนฺตสฺส จ สมฺมาวายามสฺส จ ปทฏฺฐานํ.}

\prose{ลกฺขโณติ อิติ ปุพฺพงฺคมา ธมฺมาติ เวทนา\-ปุพฺพงฺคมาปิ เอเต, สญฺญา\-ปุพฺพงฺคมาปิ เอเต, สงฺขาร\-ปุพฺพงฺคมาปิ เอเต. เย เกจิ ธมฺมา สหชาตา สพฺเพ ปุพฺพงฺคมา เอเตสํ ธมฺมานํ. \textbf{ตโต นํ สุขมเนฺวตี}ติ โสมนสฺสมปิ นํ อเนฺวติ ยํ สุสุขจฺฉายา ตทาปิ นํ สุขํ ตทป อเนฺวติ.}

\proseitem{84}{\csromanfolio{277}ตตฺถ กตโม จตุพฺยูโห หาโร. มโนปุพฺพงฺคมาติ น อิทํ เอกาทิวจนํ. กิํ การณา, สพฺเพ เยว อิเม ฉวิญฺญาณกายา, อิมมฺหิ ภควโต โก อธิปฺปาโย, เย สุเขน อตฺถิกา, เต มนํ ปสาเทนฺตีติ อยํ อิมมฺหิ สุตฺเต ภควโต อธิปฺปาโย. อตฺโถ ปุพฺเพเยว นิทฺทิฏฺโฐ.}

\prose{ยานิ หิ กุสลมูลานิ, ตานิ อฏฺฐานิสํส\-มตฺตา เหตุ, อยํ อฏฺฐงฺคิโก มคฺโค. ทส ฐานานิ เทสนาเหตูนิ เทสนาปจฺจยา นิทฺเทสนา จ. ตตฺถ ยํ มญฺเญ ทุกฺเขน สห นามรูปํ วิญฺญาณสจฺจนฺติ องฺเคน กุสลมูลํ ปหียติ, อยํ อปฺปหีน\-ภูมิยํ สมุทโย. ยํ เตสํ ปหานา, อยํ นิโรโธ. อิมานิ จตฺตาริ สจฺจานิ. อยํ อาวฏฺโฏ หาโร. วิภตฺตีติ\csromandash{}}

\setlength{\csromangathapagewidth}{0pt}
\setlength{\csromangathawakwidth}{0pt}
\csromangathameasuregroup{\textsuperscript{*}มโนปุพฺพงฺคมา ธมฺมา, \\ มนสา เจ ปสนฺเนน, \\ ตโต นํ สุขมเนฺวติ,}{\csromangathabat{\textsuperscript{*}มโนปุพฺพงฺคมา ธมฺมา,}{มโนเสฏฺฐา มโนชวา.} \\ \csromangathabat{มนสา เจ ปสนฺเนน,}{ภาสติ วา กโรติ วา.} \\ \csromangathabat{ตโต นํ สุขมเนฺวติ,}{ฉายาว อนปายินีติ.}}
\csromangathasetleft{\textsuperscript{*}มโนปุพฺพงฺคมา ธมฺมา, \\ มนสา เจ ปสนฺเนน, \\ ตโต นํ สุขมเนฺวติ,}
\csromangathagroup{\csromangathastanza{\csromangathabat{\csromansymbolfootnote{*}{เหฏฺฐา 112, 184 ปิฏฺเฐสุปิ โถกํ วิสทิสํ.}มโนปุพฺพงฺคมา ธมฺมา,}{มโนเสฏฺฐา มโนชวา.} \\ \csromangathabat{มนสา เจ ปสนฺเนน,}{ภาสติ วา กโรติ วา.} \\ \csromangathabat{ตโต นํ สุขมเนฺวติ,}{ฉายาว อนปายินีติ.}}}

\prose{ตํ น เอกํเสน สมณสฺส วา พฺราหฺมณสฺส วา ปน โหติ. ตสฺส วา มิจฺฉา\-ทิฏฺฐิกสฺส สกสตฺเถ จิตฺตํ ปสาเทติ, เตน จ ปสนฺเนน จิตฺเตน ภาสติ พฺยากโรติ น ตํ สุขมเนฺวติ น ฉายาว อนุคามินี, ทุกฺขเมว ตํ อเนฺวติ. ยถา วหนฺตํ จกฺกํ ปทมเนฺวติ, อิทํ ตํ วิภชฺช\-พฺยากรณียํ. มนสา เจ ปสนฺเนน กายกมฺมํ วจีกมฺมํ สุขเวทนียนฺติ สมคฺคเต สุขเวทนียํ มิจฺฉคฺคเต ทุกฺข\-เวทนียํ, อยํ วิภตฺติ.}

\prose{ตตฺถ กตโม ปริวตฺตโน หาโร. มโนปุพฺพงฺคมา ธมฺมาติ ยํ มนสา ปทุฏฺเฐน ภาสติ วา กโรติ วา ทุกฺขมสฺสา\-นุคามินี, เอตานิเยว เทฺว สุตฺตานิ ภาสิตานิ เอสเอว จ ปฏิปกฺโข.  เววจนนฺติ ยทิทํ มโน จิตฺตํ วิญฺญาณํ มนินฺทฺริยํ มโนวิญฺญาณ\-ธาตุ.}

\prose{ปญฺญตฺตีติ มโนปุพฺพงฺคมา ธมฺมาติ อยํ มโน กิญฺจิ ปญฺญตฺติยา ปญฺญตฺตํ. ธมฺมาติ กุสลกมฺมปถ\-ปญฺญตฺติยา ปญฺญตฺตํ. มโนเสฏฺฐาติ วิสิฏฺฐ\-ปญฺญตฺติยา ปญฺญตฺตํ. มโนชวาติ สหปญฺญตฺติยา ปญฺญตฺตํ. จิตฺตนฺติ เนกฺขมฺม\-ปญฺญตฺติยา ปญฺญตฺตํ. มนสา เจ ปสนฺเนนาติ สทฺธินฺทฺริย\-ปญฺญตฺติยา ปญฺญตฺตํ. มนสา เจ ปสนฺเนนาติ อนาวิลสงฺกปฺป\-ทุติยชฺฌาน\-ปญฺญตฺติยา ปญฺญตฺตํ. มนสา เจ \csromanfolio{278}ปสนฺเนนาติ อสฺสทฺธานํ ปฏิปกฺข\-ปญฺญตฺติยา ปญฺญตฺตํ. ภาสติ วาติ สมฺมาวาจา\-ปญฺญตฺติยา ปญฺญตฺตํ. กโรติ วาติ สมฺมากมฺมนฺต\-ปญฺญตฺติยา ปญฺญตฺตํ. ตโต นํ สุขมเนฺวตีติ ฌาน\-สมาธานํ. อินฺทฺริเยสุ มนินฺทฺริยํ. ปฏิจฺจ\-สมุปฺปาเท วิญฺญาณํ. มโนปุพฺพงฺคมา ธมฺมาติ เมตฺตา จ มุทุตา จ ฌาเนสุ ทุติยํ ฌานํ ตติยญฺจ. ขนฺเธสุ สงฺขารกฺขนฺธ\-ปริยาปนฺโน. ธาตูสุ ธมฺมธาตุ, อายตเนสุ ธมฺมายตนํ. ยํ กุสลํ อินฺทฺริเยสุ สุขินฺทฺริยญฺจ โสมนสฺสินฺทฺริยญฺจ ปทฏฺฐานํ. อิเมสํ ธมฺมานํ ปฏิจฺจ\-สมุปฺปนฺนานํ ผสฺสปจฺจยา สุขเวทนีโย ผสฺโส สุขเวทนา มโนปวิจาเรสุ โสมนสฺส\-วิจาโร ฉตฺติํเสสุ ปฐมปเทสุ ฉ โสมนสฺส\-เนกฺขมฺมสฺสิตา. อิติ อยํ โอตรโณ หาโร.}

\prose{ตตฺถ กตโม โสธโน หาโร. ยํ อตฺถํ อารพฺภ อิทํ สุตฺตํ ภาสิตํ. โส อตฺโถ นิยุตฺโต เอตมตฺถํ อารพฺภ สุตฺตํ. อยํ โสธโน หาโร.}

\proseitem{85}{ตตฺถ กตโม อธิฏฺฐาโน หาโร. มโนปุพฺพงฺคมา ธมฺมาติ เววจน\-ปญฺญตฺติ, น เอกตฺตปญฺญตฺติ. ธมฺมาติ เอกโต น เววจน\-ปญฺญตฺติ. มนสา เจ ปสนฺเนนาติ โส ปสาโท ทฺวิโธ อชฺฌตฺตญฺจ อพฺยาปาทา วิกฺขมฺภนพหิทฺธา จ โอกปฺปนโต. โส อชฺฌตฺตปสาโท ทฺวิโธ. สมุคฺฆาต\-ปสาโท จ วิกฺขมฺภน\-ปสาโท จ พฺยาปาท\-ปริยุฏฺฐานํ. วิฆาโต น มูลปสาโท ชาตมูลมฺปิ วา. ปสาโท สพฺยาปาทํ วิฆาเตน. ตโต นํ สุขมเนฺวตีติ สุขํ กายิกญฺจ เจตสิกญฺจ อปฺปิยวิปฺปโยโคปิ ปิยสมฺปโยโคปิ เนกฺขมฺม\-สุขมฺปิ ปุถุชฺชน\-สุขมฺปิ ปีติสมฺโพชฺฌงฺคมฺปิ เจตสิกํ สุขํ. ยมฺปิ ปสฺสทฺธกาโย สุขํ เวเทติ, ตมฺปิ กายิกํ สุขํ โพชฺฌงฺคา จ เจตสิกํ สุขํ. ยมฺปิ ปสฺสทฺธกาโย สุขํ เวเทสิ, ตมฺปิ ตญฺจ สุขปทฏฺฐานํ ปญฺญตฺติยา ยถาวุตฺตํ ตํ อปรามฏฺฐํ กุสลานํ ธมฺมานํ. อเนฺวตีติ อปฺปนา สนฺทิสฺสติ น จายํ วา ปตฺตภูโต อเนฺวติ. ตทิทํ สุตฺตํ ทฺวีหิ อากาเรหิ อธิฏฺฐาตพฺพํ. เหตุนา จ โย ปสนฺนมานโส วิปาเกน จ โย ทุกฺขเวทนีโย.}

\prose{ปริกฺขาโรติ ภควา ปญฺจสเตน ภิกฺขุ\-สํเฆน นครํ ปวิสติ ราชคหํ. ตตฺถ มนุสฺโส ปุคฺคโล ภควนฺตํ ปริวิสติ, ตสฺส ปสาโท อุปฺปนฺโน กุสลมูล\-ปุพฺพโยคา\-วจ\-โรปิ. โส อญฺเญสญฺจ อกฺขาติ, อิทํ วาจํ \csromanfolio{279}ภาสภิ ลาภา เตสํ, เยสํ นิเวสนํ ภควา ปวิสติ, อมฺหากมฺปิ ยทิ ภเวยฺย มยมฺปิ ภควโต สํปสาทํ ลจฺฉมฺหาติ. เยน ภควา เตนญฺชลิํ ปณาเมตฺวา “นโม ภควโต นโม ภควโต”ติ อพฺยาปาทมาโน เอกมนฺเต อฏฺฐาสิ. ตทนนฺตเร ภควา อิมํ สุตฺตํ อภาสิตฺถ “มโนปุพฺพงฺคมา ธมฺมา”ติ. สพฺพํ สุตฺตํ ตถา ยํ ปเรสํ ภาสติ อิทํ วาจากมฺมํ. ยํ อญฺชลิํ ปณาเมติ, อิทํ กายกมฺมํ. โย มโนปสาโท, อิทํ มโนกมฺมํ. ตตฺถ ยํ ปเรสํ ปกาเสติ ภาสติ วณฺณํ. เยสํ ภควา นิเวสนํ คจฺฉตีติ. สพฺพํ ตสฺส อโลโภ กุสลมูลํ. ยํ ภควติ เมตฺตจิตฺโต, ตสฺส อโทโส กุสลมูลํ. ยํ อญฺชลิํ ปณาเมติ มานญฺจ นิคฺคณฺหาติ, ตตฺถสฺส อโมโห กุสลมูลํ ปาตุภวติ. ยํ อุฬารปญฺญํ ปฏิลภติ, อิทมสฺส ทิฏฺฐิวิปลฺลาส\-ปฺปหานํ. ยํ ตถาเยว สํวโร โหติ, อิทมสฺส สญฺญาวิปลฺลาส\-ปฺปหานํ. ยํ มนสฺส ปสาทนํ, อิทมสฺส จิตฺตวิปลฺลาสปฺปหานนฺติ อกุสล\-วิปลฺลาสานํ วิกฺขมฺภนํ ปหานํ ปจฺจโย. ตีณิ กุสลมูลานิ โย อนาวิล\-จิตฺตสงฺกปฺโป, โส ตสฺส มนสิกาโรติ วุจฺจติ. ยํ กิเลเสหิ วิกฺขมฺภนํ อิติ วิปลฺลาสา จ อารมฺมณา สปฺปจฺจยตาย ปจฺจโย กุสลมูลานิ จ สนฺทิสฺสยตาย ปจฺจโย, โส จ มนสิกาโร เหตุนา อิมินา ปจฺจเยน จิตฺตํ อุปฺปนฺนํ. ตตฺถ ยํ สสตฺถา\-รมฺมณํ จิตฺตํ ปวตฺตํ อยํ พุทฺธานุสฺสติ. ยมฺปิ ภควโต คุเณ มนสิ กโรติ, อยมสฺส ธมฺมานุสฺสติ. ตตฺถ สติสมฺปชญฺญํ เหตุ อยญฺจ ปจฺจโย. วาจา ปญฺญา เหตุ วิตกฺกวิจารา ปจฺจโย. กายสงฺขารา กมฺมสฺส อภิสงฺขาโร นาม เหตุ วา อปฺปจฺจโย สุขเวทนียสฺส กมฺมสฺส อุปจโย เหตุกา กมฺมสฺส ปจฺจโย.}

\proseitem{86}{ตตฺถ กตโม สมาโรปโน หาโร. มนสาเยว ปสนฺเนน สโตเยเวตฺถ ปสนฺโน อปิ จ จิตฺตโวทานา สตฺตา วิมุจฺจนฺตีติ เตน สตฺตา จิตฺต\-ปุพฺพงฺคมา จิตฺเตน ปสนฺเนน เจตนาปิ ตตฺถ จิตฺตภูตา ภวนฺตีติ ปฏิฆา อยํ เจตนานํ ปสาเทน กาโย จสฺส ปสาโท, โส จ อารภติ ปสาเทน ปสนฺโน สญฺญานนฺติ จสฺส อวิปรีตา, โส ปญฺจวิโธ วิกฺขมฺภนา, กายปสฺสมฺภนาเยวา ปสาโท จิตฺตสิโต จิตฺตํ ปน ปุพฺพํเยว ปสนฺนํ. อยํ สมาโรปนา. เอวํ ปญฺจนฺนมฺปิ ปสาโท. ตโต นํ สุขมเนฺวตีติ กตมํ ภควา นิทฺทิสติ. น หิ อตฺตสจฺจํ ตสฺส กมฺมสฺส วิปาโก อเนฺวติ. ตสฺส อุปาโย อนุคจฺฉติ ยทา สิตปจฺจยา \csromanfolio{280}อุปฺปชฺชเต โสมนสฺสํ อวิปฺปฏิสาโรปิ อเนฺวติ. อยํ สมาโรปโน หาโร.}

\prose{มหานาม สกฺกสฺส สุตฺตํ\csromansharedfootnote{b12a3829d58d1607}{สํ 3. 322 ปิฏฺเฐ.}. ตสฺมิํ เจ สมเย อสฺสโต อสมฺปชาโน กาลํ กเรยฺย กา เม ภวติ. อสฺสโต อภิสมาหาโร โย มา ภายิ มหานาม ยํ ตํ จิตฺตํ ทีฆรตฺตํ สทฺธา\-ปริภาวิตํ สีลปริภาวิตํ สุตจาคปริภาวิตนฺติ วิตฺถาเรน กาตพฺพํ. จาเคน จ ปญฺญาย จ กิํ ทสฺเสติ, ยา สทฺธา สา เจตโส ปสาโท. ยา อนาวิล\-สงฺกปฺปิตา, สา สทฺธา, กิํ การณา, อนาวิล\-ลกฺขณา, อนาวิล\-ลกฺขณา หิ สทฺธา. อปเร อาหุ คุณปริสุทฺธิ\-นิฏฺฐา\-คมน\-ลกฺขณา, ยญฺจ อปเร วา วจน\-ปฏิคฺคห\-ลกฺขณา สทฺธา. อปโร ปริยาโย อตฺตานํ ยทิ เอวํ โอกปฺเปติ “นาหํ กิญฺจิ ชานามีติ เอสา อหํ ตตฺถ อนุญฺญตฺตา อนญฺญตา”ติ. อยํ สทฺธาติ. อปโร ปริยาโย เอกสฏฺฐิยา ทิฏฺฐิคตานํ อาทีนวา\-นุปสฺสนา อนิจฺจํ ทุกฺขม\-นตฺตาติ. เตน จ ปทิฏฺฐํ ภวติ ยถา คมฺภีเร อุทปาเน อุทกํ จกฺขุนา ปสฺสติ น จ กาเยน อภิสมฺภุนาติ. เอวมสฺส อริยา นิชฺฌาน\-กฺขนฺติยา ทิฏฺฐิ ภวติ, น จ สจฺฉิกตา. อยํ วุจฺจติ สทฺธา. สา จ โลกิกา. อปโร ปริยาโย ขมติ ปุถุชฺชน\-ภูตสฺส วีสติ จาติ โก สกฺกายาธีนา น นิเวโส. น เอตํ เอกนฺติ นยสญฺญา ยถาภูตํ ทิฏฺฐิยา ตุ ขลุ มุทูหิ ปญฺจหิ อินฺทฺริเยหิ ทสฺสนมคฺเคน ปหีนา ภวนฺติ. ทิฏฺเฐกฏฺฐา จ กิเลสา, อยํ สทฺธา.}

\prose{โสตา\-ปตฺตงฺคม\-ทุกฺขายํ ภูมิยํ ปริปุณฺณา วุจฺจติ. ตสฺมิํ เยว ภูมิยํ เสกฺขสีลํ อริยา ธารนฺติ วุจฺจติ. ตสฺมิํเยว ภูมิยํ มุทุปญฺญา ปญฺญินฺทฺริยนฺติ วุจฺจติ. ตสฺมิํ เยว ภูมิยํ ขนฺเธหิ อนตฺถิกตา, อยํ จาโค. ตสฺมา สทฺธา จาคา\-ธิฏฺฐาเนน นิทฺทิสิตพฺพา. ยติเกน\csromansharedfootnote{55011d6bf07c9fb1}{เตน (ก)} ภิยฺโย มเนน สา หิสฺส วิปรีตา ทิฏฺฐิกา อสฺสทฺธา, สา นยน- อุปธีสุ ปมตฺตา สมาทินฺนา. ตตฺถ สทฺธินฺทฺริยํ โย กามํ ปริวิสฺสนฺติ อิติ สนฺต\-ปาป\-ปฏินิสฺสคฺคา น จาคาธิฏฺฐานํ ปญฺญินฺทฺริเยน ปญฺญา\-ธิฏฺฐานํ, สีเลน อุปสมา\-ธิฏฺฐานํ. อิเม จตฺตาโร ธมฺมา สีลํ ปริภาวยนฺติ สทฺธา สีลํ จาโค จ ปญฺญาติ. ตตฺถ สทฺธาย โอฆํ ตรติ. ยํ สีลํ, อยํ อปฺปมาโท. โย จาโค, อิทํ ปญฺญาย กมฺมํ. ยา ปญฺญา, อิทํ ปญฺญินฺทฺริยํ, ตตฺถ ยํ สทฺธินฺทฺริยํ. ตํ ตีสุ อเวจฺจ\-ปฺปสาเทสุ. ยํ สีลํ, ตํ สทฺธินฺทฺริเยสุ. โย จาโค, โส \csromanfolio{281}จตูสุ ฌาเนสุ. ยา ปญฺญา, สา สจฺเจสุ, สติ สพฺพตฺถ\-คามินี. ตสฺส เสกฺขสฺส ภทฺทิกา ภติ, ภทฺทิโก อภิสมฺปราโย. ตสฺส สมฺมุฏฺฐสฺส\-ติกสฺส สีลํ กโรนฺตสฺส น กาย\-สมฺมุฏฺฐ\-สฺสติ\-ตาย ตานิ วา อินฺทฺริยานิ ตํ วา กุสลมูลํ กมฺมวิปากํ ภวติ. ตสฺส ติกสฺส อตฺถนิทฺเทโส. ตตฺถ สทฺธา สีลํ จาโค ปญฺญา จตฺตาโร ธมฺมา. ยา สทฺธา ยา จ ปญฺญา, อิทํ มโนสุจริตํ. ยํ สีลํ, อิทํ กายิกํ วาจสิกํ สุจริตํ. โย จาโค, อิทํ เจตสิกํ อโลโภ สุจริตํ. อิติ จิตฺเต คหิเต ปญฺจกฺขนฺธา คหิตา ภวนฺติ. อิเมหิ ธมฺเมหิ สุจริตํ อิทํ ทุกฺขญฺจ อริยสจฺจํ ปทฏฺฐานํ มคฺคสฺส.}

\proseitem{87}{ตตฺถ กตโม วิจโย หาโร. ยา จ สทฺธา ยญฺจ สีลํ, ตํ กิสฺส กโรติ. ยา สทฺธา ตาย ภควนฺตํ อนุสฺสรติ มตฺเตนปิ หตฺถินา สมาคตา, อสฺส โภ กุกฺกุรา สพฺพํ สีเลน นปฺปฏิปชฺชติ กาเยน วา วาจาย วา ฐานํ วิสารโท ภวตีติ อวิปฺปฏิสารี ปญฺญา ยสฺส ปญฺญตฺตํ อุปฏฺฐเปติ. ตสฺส อขณฺฑสฺส สีลํ ยํ น ปจฺฉิ ตสฺสํ โมหสฺส อกุสลจิตฺตํ อุปฺปชฺชติ มิจฺฉาทิฏฺฐิ\-สหคตํ วา, อยํ วิจโย หาโร.  ธมฺมวาทิโน ภทฺทิการาติ ภวิสฺสติ อตฺถิ เอสา ยุตฺติ.}

\prose{ตตฺถ กตโม ปทฏฺฐาโน หาโร. ยมิทํ จิตฺตํ ทีฆรตฺตํ ปริภาวิตํ สทฺธาย สีเลน จาเคน ปญฺญาย สมาธินา ปฐม\-ชฺฌานสฺส ปทฏฺฐานํ. ยา สทฺธา อสฺส อนาวิล\-สงฺกปฺโป, ตํ ทุติย\-ชฺฌานสฺส ปทฏฺฐานํ. ตีณิ จ อเวจฺจปฺปสาทา ยํ สีลํ, ตํ อริยกนฺตํ, ตํ สีลกฺขนฺธสฺส ปทฏฺฐานํ. ยา ปญฺญา, สา ปญฺญา\-กฺขนฺธสฺส ปทฏฺฐานํ. อิเม จ ธมฺมา อิทญฺจ จิตฺตํ เอโกทิภูต\-สมาธิสฺส ปทฏฺฐานํ. สทฺธา สทฺธิ\-นฺทฺริยสฺส ปทฏฺฐานํ. จาโค สมาธิ\-นฺทฺริยสฺส ปทฏฺฐานํ. ปญฺญา ปญฺญินฺทฺริยสฺส ปทฏฺฐานํ. สทฺธา จ ปญฺญา จ วิปสฺสนาย ปทฏฺฐานํ. สีลญฺจ จาโค จ สมถสฺส ปทฏฺฐานํ. สทฺธา จ ปญฺญา จ อวิชฺชา วิราคาย ปญฺญาวิมุตฺติยา ปทฏฺฐานํ. สีลญฺจ จาโค จ ราควิราคาย เจโตวิมุตฺติยา ปทฏฺฐานํ.}

\prose{ตตฺถ กตโม ลกฺขโณ หาโร. วิญฺญาเณ วุตฺเต สทฺธา\-สติ\-ภาวิเต สพฺเพ ปญฺจกฺขนฺธา วุตฺตา ภวนฺติ. สทฺธาย ภณิตาย สพฺพานิ สตฺต ธนานิ ภณิตานิ โหนฺติ สทฺธาธนํ ฯเปฯ. สีลกฺขนฺเธ วุตฺเต สมาธิ\-กฺขนฺโธ จ ปญฺญากฺขนฺโธ จ วุตฺตา ภวนฺติ. ยํ ตํ จิตฺตํ ทีฆรตฺตํ ปริภาวิตํ \csromanfolio{282}ปจฺฉิมเก กาเล น ตทนุปริวตฺติ ภวิสฺสตีติ เนตํ ฐานํ วิชฺชติ. ตตฺถ สญฺญาปิ ตทนุปริวตฺตินี ภวติ. เยปิ ตชฺชาติกา ธมฺมา, เตปิ ตทนุปริวตฺติโน ภวนฺติ. รูปสญฺญา รูปสญฺเจตนา\-นุปสฺสน\-มนสิกาโร เอวํ ฉนฺนํ อายตนานํ วิญฺญาณกาเย, อยํ ลกฺขโณ หาโร.}

\prose{ตตฺถ กตโม จตุพฺยูโห หาโร. อิธ สุตฺเต ภควโต โก อธิปฺปาโย. เย ภทฺทิกํ ภติํ อากงฺเขยฺย ภทฺทิกญฺจ อภิสมฺปรายํ, เต สทฺธํ สีลํ จาคํ ปญฺญญฺจ มนสิ กริสฺสนฺติ, อยํ อธิปฺปาโย. เย จญฺเญปิ สตฺตา ตถาคตสฺส สมฺมุขํ น ปฏิยุชฺฌนฺเต, อิมํ ธมฺมํ โสตา อวิปฺปฏิสารโต กาลํ กริสฺสนฺตีติ, อยํ อธิปฺปาโย.}

\proseitem{88}{ตตฺถ กตโม อาวฏฺโฏ หาโร. อิทมฺปิ จตฺตาโร ธมฺมา สทฺธา จ ปญฺญา จ อสฺสทฺธิยญฺจ อวิชฺชญฺจ หนนฺติ. สีลญฺจ จาโค จ ตณฺหา จ โทสญฺจ หนนฺติ. ตสฺส เทฺว มูลานิ ปหียนฺติ. ทุกฺขํ นิวตฺเตติ อปฺปหีน\-ภูมิยญฺจ ทฺวิมูลานิ ปญฺจกฺขนฺธา. เทฺว อริยสจฺจานิ สมโถ จ วิปสฺสนา จ. ทฺวินฺนํ มูลานํ ปหานํ. อิมานิ เทฺว สจฺจานิ นิโรโธ จ มคฺโค จ. อยํ อาวฏฺโฏ หาโร.}

\prose{ตตฺถ กตโม วิภตฺติ. ยํ ตํ จิตฺตํ สทฺธา\-ปริภาวิตํ ฯเปฯ. สเจ ปุถุชฺชนสฺส ตสฺสปิ ภทฺทิกา ภติ ภวิสฺสตีติ น เอกํเสน ตสฺส กมฺมํ ทิฏฺเฐเยว ธมฺเม วิปากนฺติ ปจฺเจสฺสติ, อปรมฺหิ วา ปริยาเย ภวิสฺสติ. ยํ วา อตีตํ วิปากาย ปจฺจุปฏฺฐิตํ, ตปฺปจฺจยานิ เจตานิ, เย ยถา มหา\-กมฺม\-วิภงฺเค “เตนายํ วิภชฺชพฺยากรณิโย นิทฺเทโส ธมฺมจาริโน ยา ภทฺทิกา ภตี”ติ.}

\prose{ตตฺถ กตมา ปริวตฺตนา. อสฺสทฺธิยํ ทุสฺสีลฺยํ ยํ มจฺเฉรํ ทุปฺปญฺญํ จ\csromansharedfootnote{b24441032f969ef7}{ทุปฺปญฺญิยํ (ก)} ยญฺจ ปฏิปกฺเขน ปหีนา ภวนฺติ, อยํ ปริวตฺตนา.}

\prose{ตตฺถ กตมํ เววจนํ. ยํ ตํ จิตฺตํ ทีฆรตฺตํ ปริภาวิตํ จิตฺตํ มโนวิญฺญาณํ ฯเปฯ ยํ สทฺธาพลํ สทฺธินฺทฺริยํ, ยํ สีลํ, ตํ สุจริตํ, สํยโมนิยโม ทโม ธนฺธตา อิมานิ ตสฺส เววจนานิ. โย จาโค, โส ปฏินิสฺสคฺโค, อโลโภ โวสคฺโค จาโค ยิฏฺฐานํ. ยา ปญฺญา, สา ปญฺญตฺตา ปญฺญปฺปภา ปญฺญินฺทฺริยํ ปญฺญาพลํ.}

\prose{\csromanfolio{283}ตตฺถ กตมา ปญฺญตฺติ. ยํ ตํ จิตฺตํ พีชํ ปญฺญตฺติยา ปญฺญตฺตํ. ปริภาวนา วาสนา ปญฺญตฺติยา ปญฺญตฺติ. สทฺธา ปสาท\-ปญฺญตฺติยา ปญฺญตฺตา. สีลํ สุจริต\-ปญฺญตฺติยา ปญฺญตฺตํ. จาโค ปุญฺญกิริย\-ปญฺญตฺติยา ปญฺญตฺโต. ปญฺญา วีมํส\-ปญฺญตฺติยา ปญฺญตฺตา. อิเม ตโย ธมฺมา สทฺธา สีลํ จาโค ปญฺญวโต ปาริสุทฺธิํ คจฺฉนฺติ.}

\prose{ตตฺถ กตโม โอตรโณ. ยํ จิตฺตํ, ตํ ขนฺเธสุ วิญฺญาณ\-กฺขนฺโธ. ธาตูสุ มโนวิญฺญาณ\-ธาตุ, อายตเนสุ มนายตนํ. เย จตฺตาโร ธมฺมา, เต ขนฺเธสุ สงฺขาร\-กฺขนฺเธ ปริยาปนฺนา ฯเปฯ ธาตูสุ อายตเนสุ.}

\prose{ตตฺถ กตโม โสธโน หาโร. อิทํ ภควโต ภาสิตํ มหานาเมน สกฺเกน ปุจฺฉิเตน สพฺพํ ตํ นิยุตฺตํ.}

\prose{ตตฺถ กตโม อธิฏฺฐาโน. อิทํ จิตฺตํ เวมตฺตตาย ปญฺญตฺตํ อกุสเลหิ จิตฺเตหิ อปริภาวิเตหิ ปริภาวิตนฺติ ยานิ ปุน ปริภาวิตานิ อญฺเญสมฺปิ ตตฺถ อุปาทาย ปญฺญตฺตํ สพฺเพปิเม จตฺตาโร ธมฺมา เอกตฺตตาย ปญฺญตฺตา. ภทฺทิกา ภตีติ กามโภคิโน เตว รูปธาตุ อรูปธาตุ มนุสฺสาติ สพฺพา ภทฺทิกา ภติ ตเทว กถาย ปญฺญตฺตํ, อยํ ปญฺญตฺติ.}

\prose{ตตฺถ กตโม ปริกฺขาโร. จิตฺตสฺส อินฺทฺริยานิ ปจฺจโย อาธิปเตยฺย\-ปจฺจยตาย มนสิกาโร. เหตุ\-ปจฺจยตาย ปจฺจโย. สทฺธาย โลกิกา ปญฺญา เหตุ\-ปจฺจยตาย ปจฺจโย. โยนิโส จ มนสิกาโร ปจฺจโย. สีลสฺส ปติรูป\-เทส\-วาโส ปจฺจโย. อตฺต\-สมฺมา\-ปณิธานญฺจ เหตุ. จาคสฺส อโลโภ เหตุ. อวิปฺปฏิสาโร จ เหตุปจฺจโย. ปญฺญา ปรโต จ โฆโส อชฺฌตฺตญฺจ โยนิโส มนสิกาโร เหตุปจฺจโย จ.}

\prose{ตตฺถ กตโม สมาโรปโน, ยํ ตํ จิตฺตํ ทีฆรตฺตํ ปริภาวิตนฺติ เจตสิกาปิ. เอตฺถ สพฺเพ ธมฺมา ปริภาวิตา ภทฺทิกา เต ภติ ภวิสฺสติ, ภทฺทิกา อุปปตฺติโก อภิสมฺปราโย. อิติ เย เกจิ มนุสฺสกา อุปโภค\-ปริโภคา สพฺเพ ภทฺทิกา ภติเยว, อยํ สมาโรปโน.}

\proseitem{89}{\csromanfolio{284}อุทฺธํ อโธ สพฺพธิ วีตราโคติ คาถา\csromansharedfootnote{d2c7105c6058f59c}{ขุ 1. 170 อุทาเน.}. ตตฺถ กิํ อุทฺธํ นาม, ยํ อิโต อุทฺธํ ภวิสฺสติ อนาคามี, อิทํ อุทฺธํ. อโธ นาม ยมติกฺกนฺตมตีตํ, อิทมโวจ อปทานตนฺติ อุทฺธํ. ตตฺถ อตีเตน สสฺสตทิฏฺฐิ ปุพฺพนฺตา\-กปฺปิกานํ อปรนฺตทิฏฺฐิ เกสญฺจิ, อุจฺเฉททิฏฺฐิํ ยํ\csromansharedfootnote{ce4d0bd09824260d}{อุจฺเฉททิฏฺฐิยํ (ก)} วุตฺต\-กปฺปิกานํ อิมา เจว ทิฏฺฐิโย อุจฺเฉททิฏฺฐิ จ สสฺสตทิฏฺฐิ จ. ตตฺถายํ สสฺสตทิฏฺฐิ อิมานิ ปนฺนรส ปทานิ สกฺกายทิฏฺฐิ สสฺสตํ ภชนฺติ. รูปวนฺตํ เม อตฺตา, อตฺตนิ เม รูปํ, รูปํ เม อตฺตาติ ยทุจฺจเต ปญฺญํ ปริทหนฺติ. ยา อุจฺเฉททิฏฺฐิ สา ปญฺจสตานิ อุจฺเฉทํ ภชนฺติ. เต “ตํ ชีวํ ตํ สรีรนฺ”ติ ปสฺสนฺติ, รูปํ เม อตฺตาติ ตถารูปา จตุพฺพิธา สกฺกายทิฏฺฐิ อุจฺเฉเทน จ สสฺสเตน จ. เอวํ ปญฺจสุ ขนฺเธสุ วีสติ\-วตฺถุกาย ทิฏฺฐิยา ปนฺนรส ปทานิ ปุพฺพนฺตํ ภชนฺติ. สสฺสต\-ทิฏฺฐิยา ปญฺจ ปทานิ อปรนฺตํ ภชนฺติ อุจฺเฉท\-ทิฏฺฐิยา. ตตฺถ “อยมหมสฺมี”ติ ปสฺสนฺตา รูปํ อตฺตโต สมนุปสฺสติ, โส อุจฺเฉทวาที รูปวนฺตญฺจ อตฺตานํ, อตฺตนิ จ รูปํ, รูปสฺมิํ วา อตฺตาติ โส ปสฺสติ จาติ อิติ อุจฺเฉททิฏฺฐิ จ, อตฺตโต ปฏิสฺสรติ สสฺสตทิฏฺฐิ ปุพฺพนฺตโต จ ปฏิสฺสรติ. “อยมหมสฺมี”ติ นสม\-นุ\-ปสฺสติ. ตสฺส ทิฏฺฐาสวา ปหานํ คจฺฉนฺติ. โย ตีสุ อทฺธาสุ ปุพฺพนฺเต จ อปรนฺเต จ เตน เตน นิทฺทิฏฺฐาเนน อุทฺธํ อโธ สพฺพธิ วีตราโค “อหมสฺมี”ติ น อนุปสฺสตีติ อิมินา ทฺวาเรน อิมินา ปโยเคน อิมินา อุปาเยน อิทํ ทสฺสนภูมิ จ โสตาปตฺติผลญฺจ โส อริโย ปโยโค อนภาวํ\-คเตน สํสาเรน อปุนพฺภวาติ โย โกจิ อริโย ปโยโค ปุนพฺภวาย มุทูนิ วา ปญฺจินฺทฺริยานิ มชฺฌานิ อธิมตฺตานิ วา สพฺพํ อปุนพฺภว\-ปฺปหานาย สํวตฺตนฺติ. อหนฺติ ทิฏฺโฐโฆ กาโมโฆ ภโวโฆ อวิชฺโชโฆ จ โอธิโส. ตตฺถ เทสนาหาเรน จตฺตาริ สจฺจานิ ปญฺจหิ อินฺทฺริเยหิ โสตาปตฺติ\-ผเลน จ เทฺว สจฺจานิ มคฺโค จ นิโรโธ จ. สกฺกาย\-สมุทเยน เทฺว สจฺจานิ ทุกฺขญฺจ สมุทโย จ, อยํ เทสนา หาโร.}

\prose{ตตฺถ กตโม วิจโย. “อยมหมสฺมี”ติ อสมนุปสฺสนฺโต ตีณิ ทสฺสน\-ปฺปหาตพฺพานิ สํโยชนานิ ปชหติ. อยํ วิจโย.}

\prose{ตตฺถ กตมา ยุตฺติ. ติวิธา ปุคฺคลา โกจิ อุคฺฆฏิตญฺญู\csromansymbolfootnote{*}{สํ 1. 452 ปิฏฺเฐ ปสฺสิตพฺพํ.}โกจิ วิปญฺจิตญฺญู โกจิ เนยฺโย. * อุคฺฆฏิตญฺญู ติกฺขินฺทฺริโย จ ตโต วิปญฺจิตญฺญู * มุทินฺทฺริโย ตโต มุทินฺทฺริเยหิ เนยฺโย. ตตฺถ อุคฺฆฏิตญฺญู ติกฺขินฺทฺริย\-ตาย \csromanfolio{285}ทสฺสน\-ภูมิมา\-คมฺม โสตาปตฺติผลํ ปาปุณาติ, เอกพีชโก ภวติ. อยํ ปฐโม โสตาปนฺโน. วิปญฺจิตญฺญู มุทูหิ อินฺทฺริเยหิ ทสฺสน\-ภูมิมา\-คมฺม โสตาปตฺติผลํ ปาปุณาติ, โกลํโกโล จ โหติ. อยํ ทุติโย โสตาปนฺโน. ตตฺถ เนยฺโย ทสฺสน\-ภูมิมา\-คมฺม โสตาปตฺติผลํ ปาปุณาติ, สตฺตกฺขตฺตุ\-ปรโม จ ภวติ. อยํ ตติโย โสตาปนฺโน.}

\prose{อตฺถิ เอสา ยุตฺติ มุทุมชฺฌา\-ธิมตฺเตหิ อินฺทฺริเยหิ มุทุมชฺฌา\-ธิมตฺตํ ภูมิํ สจฺฉิกเรยฺย สกฺกายทิฏฺฐิ\-ปฺปหาเนน วา ทิฏฺฐิคตานิ ปชหติ. อยํ ยุตฺติ.}

\prose{ตตฺถ กตโม ปทฏฺฐาโน, ตตฺถ สกฺกายทิฏฺฐิ สพฺพ\-มิจฺฉาทิฏฺฐิยา ปทฏฺฐานํ. สกฺกาโย นามรูปสฺส ปทฏฺฐานํ. นามรูปํ สกฺกาย\-ทิฏฺฐิยา ปทฏฺฐานํ. ปญฺจ อินฺทฺริยานิ รูปีนิ รูปราคสฺส ปทฏฺฐานํ. สฬายตนํ อหํการสฺส ปทฏฺฐานํ. ตตฺถ กตโม ลกฺขโณ, ทฺวีสุ ทิฏฺฐีสุ ปหีนาสุ ตตฺถ เอกา ทิฏฺฐิ ทิฏฺฐิคตานิ ปหานํ คจฺฉนฺติ. อุทฺธํ จ อโธ จ วีตราโค สพฺพ\-รชนีเยสุ วีตราโค โหติ. ตชฺชา ปรภูมิยํ, อิทํ ปจฺจยนฺติ ยถาภูตํ ปสฺสติ. โส สพฺพ\-ปฏิจฺจ\-สมุปฺปาทํ อามสติ. อยํ ลกฺขโณ หาโร.}

\proseitem{90}{ตตฺถ กตโม จตูพฺยูโห หาโร. อิมมฺหิ สุตฺเต ภควโต โก อธิปฺปาโย. เย สตฺตา เย นาภิรมิสฺสนฺติ, เต ทิฏฺฐิ\-ปฺปหานาย วายมิสฺสนฺติ. อยเมตฺถ ภควโต อธิปฺปาโย. อยํ จตุพฺยูโห หาโร.}

\prose{ตตฺถ กตโม อาวฏฺโฏ หาโร. ยานิมานิ มุทูนิ ปญฺจินฺทฺริยานิ ตานิ โอรมฺภาคิยานิ ปญฺจินฺทฺริยานิ. สพฺเพน สพฺพํ สมูหนนฺติ อภิชฺฌา\-พฺยาปาโท จ ภาวนากาเรน เสกฺขาย วิมุตฺติยา พลํ สทฺธา, อุทฺธมฺ\-ภาคิยานิ ทิฏฺฐิวเสน พลํ สทฺธา, วีริยินฺทฺริยํ อารภิตตฺตา สตินฺทฺริยํ ปคฺคหิตตฺตา อจฺจนฺตํ นิฏฺฐํ คจฺฉนฺติ. ตตฺถ ยานิ อินฺทฺริยานิ, อยํ มคฺโค สํกิเลส\-ปฺปหานํ. อยํ นิโรโธ อายติํ อนุปฺปาทธมฺโม, อิทํ ทุกฺขํ. อยํ อาวฏฺโฏ หาโร.}

\prose{ตตฺถ กตโม วิภตฺติ หาโร. “อยมหมสฺมี”ติ โย สมนุปสฺสติ, โส จ โข อธิมตฺเตน โลกิกา ยํ ภูมิยํ น ตุ อริเยน ปโยเคน โส สกฺกายทิฏฺฐิ ปชหติ. ยํ วุจฺจติ ตชฺชาย ภูมิยา อธิมตฺตาย. ตตฺถ ตชฺชาย ภูมิยํ ปญฺจหิ อากาเรหิ อธิมตฺตตํ ปฏิลภติ สีเลน วเตน พาหุสฺสจฺเจน สมาธินา เนกฺขมฺม\-สุเขน. ตตฺถ อปฺปตฺเต \csromanfolio{286}ปตฺตสญฺญี อธิมานํ คณฺหาติ. เอตสฺมิํ เยว วตฺถุ\-ปฺปตฺติยํ ภควา อิทํ สุตฺตํ ภาสติ. สีลวา วตมตฺเตนาติ. ตตฺถ โย อปฺปตฺเต ปตฺตสญฺญี ตสฺส โย สมาธิ, โส สามิโส กาปุริสเสวิโต ปน โส กาปุริสา วุจฺจนฺติ ปุถุชฺชนา. อามิสํ ยญฺจ อริยมคฺคมา\-คมฺม โลกิกา อนริยํ เตน สมาธิ โหติ อนริโย กาปุริสเสวิโต. โย ปน อริยากาเรน ยถาภูตํ น ชานาติ น ปสฺสติ\csromansharedfootnote{e54e851c257820bb}{ฌานาติ ปสฺสติ (อิ)}, โส อธิคมนํ ปชหติ โย อริเยน สมาธินา อกาปุริสเสวิเตน นิรามิเสน นียติ, ตตฺถ อกาปุริสา วุจฺจนฺติ อริยปุคฺคลา. โย เตหิ เสวิโต สมาธิ, โส อกาปุริสเสวิโต. ตสฺมา เอกํ วิภชฺช\-พฺยากรณียํ “อยมหมสฺมี”ติ อสมนุปสฺสนฺโต ตถา ปาเตติ.}

\prose{ตตฺถ กตมา ปริวตฺตนา, อิมาย ทสฺสนภูมิยา กิเลสา ปหาตพฺพา, เตหิ ปหียนฺติ อนิทฺทิฏฺฐาปิ ภควตา นิทฺทิสิตพฺพา โย.}

\prose{ตตฺถ กตมํ เววจนํ. ยา สกฺกาย\-ทิฏฺฐิยา อตฺตทิฏฺฐิยา. อยํ ภูมิ. เย กิเลสา ปหาตพฺพา. เต อปฺปหียนฺติ อนิทฺทิฏฺฐาปิ ภควตา สสฺสตทิฏฺฐิ จ อุจฺเฉททิฏฺฐิ จ, สา ปริยนฺต\-ทิฏฺฐิ จ. ยา อปริยนฺต\-ทิฏฺฐิ จ, สา สสฺสตทิฏฺฐิ จ. ยา อุจฺเฉททิฏฺฐิ, สา นตฺถิกา ทิฏฺฐิ. ยา สสฺสตทิฏฺฐิ, สา อกิริยทิฏฺฐิ. อิทํ เววจนํ.}

\prose{ตตฺถ กตมา ปญฺญตฺติ. ตณฺหา สํโยชน\-ปญฺญตฺติยา ปญฺญตฺตา. มคฺโค ปฏิลาภ\-ปญฺญตฺติยา ปญฺญตฺโต. อินฺทฺริยา ปฏิลาภ\-ปญฺญตฺติยา ปญฺญตฺตาติ.  ตตฺถ กตโม โอตรโณ. สกฺกาโย ทุกฺขํ ทสฺสน\-ปฺปหาตพฺโพ. สมุทโย มคฺโค. อินฺทฺริยานิ ตานิ จ นิทฺทิฏฺฐานิ ขนฺธ\-ธาตุ\-อายตเนสุ.}

\prose{ตตฺถ กตโม โสธโน หาโร. ยญฺหิ อารพฺภ ภควตา อิทํ สุตฺตํ ภาสิตํ, โส อารพฺภ นิทฺทิฏฺโฐ.  ตตฺถ กตโม ปริกฺขาโร. นามรูปสฺส เหตุปจฺจโยปิ วิญฺญาณํ เหตุ พีชํ. เตน อวิชฺชา จ สงฺขารา จ ปจฺจโย. นิวตฺตินโย น อปโร ปริยาโย สพฺพภโว, เย จ สพฺพภวสฺส เหตุ ปรภณฺฑ\-ปจฺจโย อิติ สมฺมาทิฏฺฐิ ปรโต จ โฆโส โยนิโส จ มนสิกาโร ปจฺจโย. ยา ปญฺญา อุปฺปาเทติ, เอสา เหตุ สมฺมาทิฏฺฐิยา สมฺมาสงฺกปฺโป ภวติ, ยา สมฺมาสมาธิ\csromansharedfootnote{306fb30a17b58e2f}{สมฺมาทิฏฺฐิ (อิ)}, อยํ ปริกฺขาโร.}

\prose{\csromanfolio{287}ตตฺถ กตโม สมาโรปโน. “อยมหมสฺมี”ติ อสมนุปสฺสี ทุกฺขโต โรคโต ฯเปฯ ปนฺนรส ปทานิ. สีลานิ ภควา กิมตฺถิยานิ กิมานิสํสานิ. สีลานิ อานนฺท อวิปฺปฏิสาร\-ตฺถานิ ยาว วิมุตฺติ. ตตฺถ ทุวิโธ อตฺโถ ปุริสตฺโถ จ วจนตฺโถ จ.}

\proseitem{91}{ตตฺถ กตโม ปุริสตฺโถ. ยายํ น ปจฺฉานุตาปิตา อยํ อวิปฺปฏิสาโร, อยํ ปุริสตฺโถ. ยถา โกจิ พฺรูหยติ อิมตฺถมาเสวติ โส ภเณยฺย, กิญฺจิ มเมตฺถ อธีนํ ตสฺสตฺถาย อิทํ กิริยํ อารภามีติ. อยํ ปุริสตฺโถ.}

\prose{ตตฺถ กตโม วจนตฺโถ. สีลานิ กายิกํ วา วาจสิกํ วา สุจริตํ อวิปฺปฏิสาโรติ. ตตฺถ สีลสฺส วตสฺส จ ภาโสเยว. อนญฺญา สุคตกมฺมตา สุจริตํ อยํ อวิปฺปฏิสาโร. เอวํ ยาว วิมุตฺตีติ เอกเมกสฺมิํ ปเท เทฺว อตฺถา ปุริสตฺโถ จ วจนตฺโถ จ, ยถา อิมมฺหิ สุตฺเต เอวํ สพฺเพสุ สุตฺเตสุ เทฺว เทฺว อตฺถา. อยํ หิ ปรมตฺโถ อุตฺตมตฺโถ จ. ยํ นิพฺพาน\-สจฺฉิกํ นิสฺสาย ยํ สกํ สจฺฉิกาตพฺพํ ภวติ, โส วุจฺจติ กตสฺส\csromansharedfootnote{65a82f0d03dbe14b}{สตสฺส (ก)} กตฺโถติ. อยํ ปุน เววจนํ สมฺปชานาติ. อิมินา นิยุตฺตตฺถมภิลพฺภนฺติ วจนตฺโถ. ตตฺถ ยํ อตฺถํ สาวโก อภิกงฺขติ. ตสฺส โย ปฏิลาโภ, อยํ ปุริสตฺโถ. ยํ ยํ ภควา ธมฺมํ เทเสติ, ตสฺส ตสฺส ธมฺมสฺส ยา อตฺถวิญฺญตฺติ. อยํ อตฺโถ, ตตฺถ สีลานํ อวิปฺปฏิสาโร อตฺโถปิ อานิสํโสปิ. เอโส จ อานิสํโส ยํ ทุคฺคติํ น คจฺฉติ. ยถา ตํ ภควตา เอสานิสํโส ธมฺเม สุจิณฺเณ น ทุคฺคติํ คจฺฉติ ธมฺมจารี, อยํ อตฺโถ.}

\prose{ยํ ปุริโส ภาวนา\-ภูมิยํ สีลานิ อารพฺภ สีเลน สํยุตฺโต โหติ เอวํ ยาว วิมุตฺติ ตถา สีลกฺขนฺโธ. ตตฺถ โย จ อวิปฺปฏิสาโร อนุสยวเสน นิทฺทิฏฺโฐ, ตญฺจ สีลํ อยํ สีลกฺขนฺโธ. ปาโมชฺช\-ปีติ\-ปสฺสทฺธีติ จ สมาธิ\-นฺทฺริเยน, อยํ สมาธิ\-กฺขนฺโธ. ยํ สมาหิโต ยถาภูตํ ปชานาติ, อยํ ปญฺญากฺขนฺโธ. อิเม ตโย ขนฺธา สีลํ สมาธิ ปญฺญา จ ตถา สีลํ ปริปูเรติ ยํ วีริยินฺทฺริยํ เตน การเณน โส สีลํ ปริปูเรติ, อนุปฺปนฺนสฺส จ อกุสลสฺส อนุปฺปาทาย วายมติ, อุปฺปนฺนสฺส จ ปหานาย อนุปฺปนฺนสฺส จ กุสลสฺส อุปฺปาทาย, อุปฺปนฺนสฺส จ กุสลสฺส \csromanfolio{288}ภิยฺโยภาวาย อิติ วีริยินฺทฺริยํ นิทฺทิฏฺฐํ. ตตฺถ โย สมาธิ\-กฺขนฺโธ, อิทํ สมาธินฺทฺริยํ. ปญฺญากฺขนฺโธ ปญฺญินฺทฺริยํ, ตํ จตูสุ สมฺม\-ปฺปธาเนสุ ทฏฺฐพฺพํ. ตถา โย อนุปฺปนฺนสฺส จ อกุสลสฺส อนุปฺปาทาย วายมติ, อิทํ ปฐมํ สมฺมปฺปธานํ. ยํ อุปฺปนฺนสฺส, อิทํ ทุติยํ. จตฺตาริ สมฺม\-ปฺปธานานิ จตูสุ ฌาเนสุ ปสฺสิตพฺพานิ. ตถา สีลกฺขนฺเธน เนกฺขมฺมธาตุ จ อธิกา\csromansharedfootnote{368411095038d7ed}{อาทิกา (อิ)}, ตโย จ วิตกฺกา เนกฺขมฺม\-วิตกฺโก อพฺยาปาทวิตกฺโก อวิหิํสา\-วิตกฺโก จ. สาธารณภูตา. ยา ปิยายมานสฺส ปาโมชฺเชน อิทํ กายิกํ สุขํ อานิตํ อนิยมีติเปเมน, อิทํ ทุกฺขํ. โย ตตฺถ อวิกฺเขโป, อยํ สมาธิ. อิทํ ปญฺจงฺคิกํ ปฐมํ ฌานํ. ยา เจตสิกา ปสฺสทฺธิ สวิตกฺกํ สวิจารํ วิโรธนํ, โย กิเลโส จ ปริทาโห, โส ปฐเม ฌาเน นิรุทฺโธ. ตถา ยา จ กิเลส\-ปสฺสทฺธิ ยา จ วิตกฺก\-วิจารานํ ปสฺสทฺธิ, อุภเยปิ เอเต ธมฺเม ปสฺสทฺธายํ. ตตฺถ กายสฺส จิตฺตสฺส จ สุขํ สุขายนา, อิทํ ปีติสุขิโน ปสฺสทฺธิ. โยปิ เอโกทิภาโว จิตฺตสฺส, เตน เอโกทิภาเวน ยํ จิตฺตสฺส อชฺฌตฺตํ สมฺปสาทนํ, อิทํ จตุตฺถํ ฌานงฺคํ. อิติ อชฺฌตฺตญฺจ สมฺปสาโท เจตโส จ เอโกทิภาโว ปีติ จ สุขญฺจ, อิทํ ทุติยํ ฌานํ จตุรงฺคิกํ. โย ปสฺสทฺธกาโย สุขํ เวเทติ, เตน อธิมตฺเตน สุเขน ผริตฺวา สุขํ เจตสิกํ ยํ, โส ปีติวีตราโค เอวํ ตสฺส ปีติวีตราค\-ตาย อุเปกฺขํ ปฏิลภติ. โส ปีติยา จ วิราคา อุเปกฺขํ ปฏิลภติ. สุขญฺจ ปฏิสํเวเทติ. สติ จ สมฺมา ปญฺญาย ปฏิลภติ. สเจ สติ เอกคฺคตา อิทํ ปญฺจงฺคิกํ ตติยํ ฌานํ. ยํ สุขิโน จิตฺตํ สมาธิยติ, อยํ เอกคฺคตาย ปรา\-วิธาน\-ภาคิยา, ปฐเม ฌาเน อตฺถิ จิตฺเตกคฺคตา โน จกฺขุสฺส เวทนา สพฺพํ ปาริปูริํ คจฺฉติ. ยถา จตุตฺเถ ฌาเน, ตถา ยา อุเปกฺขา ปสฺสมฺภยํ สติสมฺปชญฺญํ จิตฺเตกคฺคตา จ, อิทํ จตุตฺถํ ฌานํ.}

\proseitem{92}{ยถา สมาธิ ทสฺสยิตพฺพํ, ตถา ปญฺญินฺทฺริยํ ตํ จตูสุ อริยสจฺเจสุ ปสฺสิตพฺพํ. ยํ สมาหิโต ยถาภูตํ ปชานาติ, สา ปชานนา จตุพฺพิธา อสุภโต ทุกฺขโต อนตฺตโต จ ยทารมฺมณํ, ตํ ทุกฺขํ อริยสจฺจํ, ยํ ปชานนฺโต นิพฺพินฺทติ วิมุจฺจติ ตถา ยํ กามาสวสฺส ปหานํ ภวาสวสฺส ทิฏฺฐาสวสฺส อวิชฺชาสวสฺส, อยํ นิโรโธ \csromanfolio{289}อปฺปหีน\-ภูมิยํ อาสวสมุทโย. อิมานิ จตฺตาริ อริยสจฺจานิ ยถา ปญฺญินฺทฺริยํ ปสฺสิตพฺพํ. ยถายํ สมาหิโต ยถาภูตํ ปชานาติ, อยํ ทสฺสนภูมิ. โสตาปตฺติผลญฺจ ยถาภูตํ ปชานนฺโต นิพฺพินฺทตีติ. อิทํ ตนุกญฺจ. กามราค\-พฺยาปาทํ สกทาคามิ\-ผลญฺจ ยํ นิพฺพินฺทติ วิรชฺชติ, อยํ ปฐมชฺฌาน\-ภาวนาภูมิ จ ราควิราคา เจโตวิมุตฺติ อนาคามิผลญฺจ. ยํ วิมุตฺติ วิมุจฺจติ, อยํ อวิชฺชาวิราคา ปญฺญาวิมุตฺติ อรหตฺตญฺจ. อิเม อวิปฺปฏิสารา จ วีริยินฺทฺริยญฺจ จตฺตาโร สมฺมปฺปธานา อวิปฺปฏิสารา ตญฺจ อุปริ ยาว สมาธิ, เอวํ เต จตฺตาริ ฌานานิ สมาธินฺทฺริยญฺจ ยํ สมาหิโต ยถาภูตํ ปชานาติ. อิเม จตฺตาโร สติปฏฺฐานา สีลปาริปูริมุ\-ปาทาย จาคสํหิเตน จ นิพฺเพธิกานญฺจ นิมิตฺตานํ อนาวิลมนา, อิทํ สตินฺทฺริยํ จตฺตาโร สติปฏฺฐานา. ยํ ปุน อิมาย ธมฺมเทสนาย ตีสุ ฐาเนสุ ทิฏฺโฐคมนกินฺทฺริยํ กิเลสปหาเนน จ เสกฺขสีลํ, อิทํ สทฺธินฺทฺริยํ. จตฺตาริ จ โสตาปตฺติยงฺคานิ ผลานิ. สมาธิ\-นฺทฺริยานิ โสปนิยาหารีนิ สพฺพสุตฺเตสุ นิทฺทิสิตพฺพานิ. ยํ ฌานํ ปฏิลภนํ วิริยคหิตํเยว ญาณํ ปฏิสฺสรโต, อยํ สุตมยี ปญฺญา. โย สมาธิ ปุพฺพาปร\-นิมิตฺตา\-ภาโส อโนมคติตาย ยถากาโม, อยํ จินฺตามยี ปญฺญา, ยํ ตถาสมาหิโต ยถาภูตํ ปชานาติ, อยํ ภาวนามยี ปญฺญา. อยํ สตฺตนิทฺเทโส.}

\prose{อิมํ สุตฺตํ นิพฺเพธ\-ภาคิยํ พุชฺฌ\-การธิกํ พุชฺฌิตพฺพํ. เยหิ องฺเคหิ สมนฺนาคตํ ตํ พุชฺฌิสฺสนฺติ ตสฺส องฺคานิ พุชฺฌิสฺสนฺติ, เตน โพชฺฌงฺคา. ตถา อาทิโต ยาว สีลํ วตํ เจตนา กรณียา, กิสฺส สีลานิ ปาริปูเรติ. อนุปฺปนฺนสฺส จ อกุสลสฺส อนุปฺปาทาย อุปฺปนฺนสฺส จ อกุสลสฺส ปหานาย อนุปฺปนฺนสฺส กุสลสฺส อุปฺปาทาย อุปฺปนฺนสฺส จ กุสลสฺส ภิยฺโยภาวาย, อิทํ วีริยํ ตสฺส ตสฺส พุชฺฌิตสฺส องฺคนฺติ. อยํ วีริย\-สมฺโพชฺฌงฺโค. อิมินา วีริเยน เทฺว ธมฺมา อาทิโต อวิปฺปฏิสาโร ปาโมชฺชญฺจ ยา ปุน ปีติ อวิปฺปฏิสาร\-ปจฺจยา ปาโมชฺชปจฺจยา, อยํ ปีติสมฺโพชฺฌงฺโค. ยํ ปีติมนสฺส กาโย ปสฺสมฺภติ. อยํ ปสฺสทฺธิ\-สมฺโพชฺฌงฺโค. เตน กายิก\-สุข\-มานิตํ ยํ สุขิโน จิตฺตํ สมาธิยติ, อยํ สมาธิ\-สมฺโพชฺฌงฺโค. ยํ สมาหิโต ยถาภูตํ ปชานาติ, อยํ ธมฺมวิจย\-สมฺโพชฺฌงฺโค. ยา สีลมุปาทาย ปญฺจนฺนํ โพชฺฌงฺคานํ อุปาทา\-ยานุ\-โลมตา นิมิตฺตายนา ปีติภาคิยานญฺจ \csromanfolio{290}วิเสสภาคิยานญฺจ อปิลาปนตา สหคตา โหติ อนวมคฺโค, อยํ สติสมฺโพชฺฌงฺโค. ยํ ยถาภูตํ ปชานาติ, อจฺจารทฺธ\-วีริยํ กโรติ. อุทฺธจฺจ\-ภูมีติ กตา อภิปตฺถิตํ เปเสติ. โกสชฺชภูมีติ ครหิโต รหิเตหิ องฺเคหิ พุชฺฌติ ยํ จกฺขุ\-สมถปถํ, สา อุเปกฺขาติ. เตน สา อุเปกฺขา ตสฺส โพชฺฌงฺคสฺส องฺคนฺติ กริตฺวา อุเปกฺขา\-สมฺโพชฺฌงฺโคติ วุจฺจเต. เอโส สุตฺตนิทฺเทโส.}

\proseitem{93}{ตตฺถ กตมา เทสนา. อสฺมิํ สุตฺเต จตฺตาริ อริยสจฺจานิ เทสิตานิ.  ตตฺถ กตโม วิจโย, สีลวโต อวิปฺปฏิสาโร ยาว วิมุตฺติ อิมิสฺสาย ปุจฺฉาย มินิ กิมตฺถสฺสมีติ เทฺว ปทานิ ปุจฺฉา เทฺว ปทานิ วิสชฺชนานิ ทฺวีหิ ปเทหิ เทฺว อภิญฺญํ ทฺวีหิ เจว ปเทหิ วิสชฺชนา กิํ ปุจฺฉติ นิพฺพาธิกํ กายภูมิํ กมฺมสฺส ตถา หิ ปติฏฺฐา จ อเสกฺเข ธมฺเม อุปฺปาเทติ.  ตตฺถ กตมา ยุตฺติ, สีลวโต อวิปฺปฏิสาโร ภวติ กิํ นิจฺฉนฺทสฺส จ วิราโค อตฺถิ เอสา ยุตฺติ.  ตตฺถ กตมํ ปทฏฺฐานํ. วีริยํ วีริยิ\-นฺทฺริยสฺส ปทฏฺฐานํ. สมาธิ สมาธิ\-นฺทฺริยสฺส ปทฏฺฐานํ. ปญฺญา ปญฺญินฺทฺริยสฺส ปทฏฺฐานํ. วีริยํ อโทสสฺส ปทฏฺฐานํ. สมาธิ อโลภสฺส ปทฏฺฐานํ. ปญฺญา อโมหสฺส ปทฏฺฐานํ. วีริยินฺทฺริยํ ติณฺณํ มคฺคงฺคานํ ปทฏฺฐานํ, สมฺมาวาจาย สมฺมา\-กมฺมนฺตสฺส สมฺมาอาชีวสฺส. สมาธินฺทฺริยํ ติณฺณํ มคฺคงฺคานํ ปทฏฺฐานํ, สมฺมา\-สงฺกปฺปสฺส สมฺมาวาจาย สมฺมาสมาธิโน. ปญฺญินฺทฺริยํ ทฺวินฺนํ มคฺคงฺคานํ ปทฏฺฐานํ, สมฺมาสติยา สมฺมาทิฏฺฐิยา จ.}

\prose{ตตฺถ กตโม ลกฺขโณ, สีลกฺขนฺเธ วุตฺเต สพฺเพ ตโย ขนฺธา วุตฺตา ภวนฺติ, สีลเมว หิ เสโลปมตา ยถา เสโล สพฺพ\-ปจฺจตฺถิเกหิ อกรณีโย เอวํ ตํ จิตฺตํ สพฺพกิเลเสหิ น กมฺปตีติ, อยํ อโมโห. วิรตฺตํ\csromansharedfootnote{8ef00d88dd4a2c8b}{ขุ 1. 125 อุทานปาฬิยํ.} รชนีเยสูติ อยํ อโลโภ. โกปเนยฺเย น กุปฺปตีติ อยํ อโทโส. ตตฺถ ปญฺญา อโมโห กุสลมูลํ, อโลโภ อโลโภเยว, อโทโส อโทโสเยว. อิเมหิ ตีหิ กุสลมูเลหิ เสกฺขภูมิยํ ฐิโต อเสกฺขมคฺคํ อุปฺปาเทติ. เสกฺขภูมิ สมฺปตฺติ\-กมฺม\-ธมฺเม อุปฺปาเทติ, สา จ สมฺมาวิมุตฺติ, ยญฺจ วิมุตฺติ รสญาณทสฺสนํ อิเม ทส อเสกฺขานํ อรหตฺตํ ธมฺมา. ตตฺถ อฏฺฐงฺคิเกน มคฺเคน จตุพฺพิธา ภาวนาปิ ลพฺภติ. สีลภาวนา กายภาวนา \csromanfolio{291}จิตฺตภาวนา ปญฺญาภาวนา จ. ตตฺถ สมฺมา\-กมฺมนฺเตน สมฺมาอาชีเวน จ กาโย ภาวิโต. สมฺมาวาจาย สมฺมาวายาเมน จ สีลํ ภาวิต. สมฺมา\-สงฺกปฺเปน สมฺมาสมาธินา จ จิตฺตํ ภาวิตํ. สมฺมาทิฏฺฐิยา สมฺมาสติยา จ ปญฺญา ภาวิตา. อิมาย จตุพฺพิธาย ภาวนาย เทฺว ธมฺมา ภาวนา\-ปาริปูริํ คจฺฉนฺติ จิตฺตํ ปญฺญญฺจ. จิตฺตํ ภาวนาย สมโถ, ปญฺญา ภาวนาย วิปสฺสนา. ตตฺถ ปญฺญา อวิชฺชาปหาเนน จิตฺตํ อุปกฺกิเลเสหิ อมิสฺสีกตนฺติ. ปญฺญา ภาวนาย จิตฺต\-ภาวนํเยว ปริปูเรติ. เอวํ ยสฺส สุภาวิตํ จิตฺตํ กุโต ตํ ทุกฺขเมสฺสตีติ. อปิ จ โข ปน ตสฺส อายสฺมโต อพฺยาปาทธาตุ อธิมุตฺตา, น โส เปตํ สมาปนฺโน ตสฺส สงฺขาปหารํ เทติ, สงฺขา\-วิตกฺกิเต สรีเร ทุกฺขํ น เวทิยติ, อยํ สุตฺตตฺโถ.}

\proseitem{94}{ตตฺถ กตมา เทสนา. อิมมฺหิ สุตฺเต ทส อเสกฺขา อรหตฺตธมฺมา เทสิตา อปฺปมาณา จ สมฺมา วิภาวนา.  ตตฺถ กตโม วิจโย. เสโลปมตา เย เย ธมฺมา เวทนีย\-สุขทุกฺโข\-ปคตา, เต สพฺเพ นิรูปํ วานุปสฺสนฺตานํ วูปคตา กายโต เวทยิต\-ปริกฺขาโร อปฺปวตฺติโต ทุกฺขํ น เวทิยติ.  ตตฺถ กตมา ยุตฺติ, ยสฺเสวํ ภาวิตํ จิตฺตํ กุโต ตํ\csromansharedfootnote{b47fe170f88cd559}{นํ (ก)} ทุกฺขเมสฺสตีติ. ตีสุ ภาวนาสุ ทุกฺขํ นกฺขมติ จิตฺตํ จิตฺตภาวนาย จ. นิโรธ\-ภาวนาย จ อานนฺตริกา สมาธิ\-ภาวนาย จ.  อิติ ยสฺเสวํ ภาวิตํ จิตฺตนฺติ สมาธิ ผลสฺส ปทฏฺฐานํ.}

\prose{ตตฺถ กตโม ลกฺขโณ, ยสฺเสวํ ภาวิตํ\csromansharedfootnote{8ef00d88dd4a2c8b}{ขุ 1. 125 อุทานปาฬิยํ.} จิตฺตนฺติ จิตฺตานิ ภาวิตานิ ยถา ปฐมํ นิทฺทิฏฺฐานิ ปญฺญา สีลํ กาโย จิตฺตํ, สีลมฺปิ สุภาวิตํ กายิกเจตสิกญฺจ ฐิตตฺตา นานุ\-ป\-กมฺปตีติ เวทนาปิ ตถา สญฺญาปิ สงฺขาราปิ. กุโต ตํ ทุกฺขเมสฺสตีติ สุขมฺปิ นานุคจฺฉติ, อทุกฺขม\-สุขมฺปิ นาคตนฺติ.}

\prose{ตตฺถ กตโม จตุพฺยูโห หาโร, อิธ ภควโต โก อธิปฺปาโย. เย ทุกฺเขน อธิกา ภวิสฺสนฺติ, เต เอวรูปาหิ สมาปตฺตีหิ วิรหิสฺสนฺติ. อยเมตฺถ ภควโต อธิปฺปาโย. เย จ อปฺปสนฺนา, เต หิ ภวิสฺสนฺติ, \csromanfolio{292}ปสนฺนานญฺจ ปีติปาโมชฺชํ ภวิสฺสติ, อยํ ตตฺถ ภควโต อธิปฺปาโย.  อาวฏฺโฏติ นตฺถิ อาวฏฺฏนสฺส ภูมิ.}

\prose{วิภตฺตีติ ยสฺเสวํ ภาวิตํ จิตฺตํ กุโต ตํ ทุกฺขเมสฺสตีติ ทุวิโธ นิทฺเทโส, ทุกฺขเหตุ\-นิทฺเทโส จ ปฏิปกฺข\-นิทฺเทโส จ. โก โส ทุกฺขเหตุ, ยโต ทุกฺขํ อาคจฺฉติ ปฏิปกฺเข วุตฺเต เสสธมฺมานํ สีลํ เหตุ จ ปจฺจโย จ, เต สพฺเพ ธมฺมา วุตฺตา โหนฺติ. เอก\-โพธิปกฺขิเย ธมฺเม วุตฺเต สพฺเพ โพธคมนียา ธมฺมา วุตฺตา ภวนฺติ.}

\prose{ตตฺถ กตโม จตุพฺยูโห หาโร, อิมมฺหิ สุตฺเต ภควโต โก อธิปฺปาโย, เย อวิปฺปฏิสาเรน ฉนฺทิกา, เต สีลปาริปูรี ภวนฺติ ปาโมชฺช\-ฉนฺทิกา อวิปฺปฏิสารี\-ปาริปูรี, อยเมตฺถ ภควโต อธิปฺปาโย ฯเปฯ อยํ จตุพฺยูโห หาโร.}

\prose{ตตฺถ กตโม อาวฏฺโฏ, อิทํ สุตฺตํ นิพฺเพธ\-ภาคิยํ, โย นิพฺเพโธ, อยํ นิโรโธ. เยน นิพฺพิชฺฌติ, โส มคฺโค. ยํ นิพฺพิชฺฌติ, ตํ ทุกฺขํ. ยํ นิพฺเพธคามินา มคฺเคน ปหียติ, สมุทโยยํ วุตฺโต.}

\prose{ตตฺถ กตมา วิภตฺติ, สีลวโต อวิปฺปฏิสาโรติ วิภชฺช\-พฺยากรณียํ, ปรามสนฺตสฺส นตฺถิ อวิปฺปฏิสาโร ยาว โทสกตํ กาเยน วา วาจาย วา อกุสลํ อารภติ. กิญฺจิปิสฺส เอวํ โหติ “สุกตเมตํ สุจริตเมตํ โน จสฺส เตน อวิปฺปฏิสาเรน ปาโมชฺชํ ชายติ ยาว วิมุตฺติ, ตสฺส สีลวโต อวิปฺปฏิสาโร”ติ วิภชฺช\-พฺยากรณียํ, อยํ วิภตฺติหาโร.}

\prose{ตตฺถ กตมา ปริวตฺตนา, อิเมหิ สตฺตหิ อุปนิสา\-สมฺปตฺตีหิ เอกาทส อุปนิสา วิภตฺติยํ ปชหานํ ปชหนฺติ, อยํ ปริวตฺตนา.}

\prose{ตตฺถ กตมา เววจนา, อิเมสํ อริยธมฺมานํ พลโพชฺฌงฺค\-วิโมกฺข\-สมาธิ\-สมาปตฺตีนํ อิมานิ เววจนานิ.}

\prose{ตตฺถ กตมา ปญฺญตฺติ, สีลวโต อวิปฺปฏิสาโรติ สีลกฺขนฺเธ เนกฺขมฺม\-ปญฺญตฺติยา ปญฺญตฺตํ, นิสชฺช\-ปญฺญตฺติ จ เอวํ ทส องฺคานิ ทฺวีหิ ทฺวีหิ องฺเคหิ ปญฺญตฺตานิ.}

\prose{ตตฺถ กตโม โอตรโณ, อิทํ นิพฺเพธภาคิย\-สุตฺตํ ปญฺจสุ โอติณฺณํ ยถา ยํ ปฐมํ นิทฺทิฏฺฐํ เอวมิ\-นฺทฺริยา\-ทิ\-ขนฺธธาตุอายตเนสุ นิทฺทิสิตพฺพานิ.}

\prose{\csromanfolio{293}ตตฺถ กตโม โสธโน หาโร, สีลวโต อวิปฺปฏิสาโรติ น ตาว สุทฺโธ อารมฺโภ อวิปฺปฏิสาริโน ปาโมชฺชนฺติ น ตาว สุทฺโธ อารมฺโภ ยานิ เอกาทส ปทานิ เทสิตานิ ยทา ตทา สุทฺโธ อารมฺโภ, อยํ โสธโน.}

\prose{ตตฺถ กตโม อธิฏฺฐาโน, สีลเวมตฺตตาย ปญฺญตฺตํ เอวํ ทส ปทานิ สพฺพานิ สีลกฺขนฺธสฺส อานิสํโส, เต จ ปตีรูปเทสวาโส จ ปจฺจโย อตฺต\-สมฺมา\-ปณิธานญฺจ เหตุ, สมาธิ\-กฺขนฺธสฺส สุขํ เหตุ ปสฺสทฺธิ ปจฺจโย, เยน ฌานสหชาติ จ ฐานนฺติ ฌานงฺคา อปโร ปริยาโย กาเมสุ อาทีนวา\-นุปสฺสนา สมาธิโน ปจฺจโย เนกฺขมฺเม อานิสํส\-ทสฺสาวิตา เหตุ.}

\prose{ตตฺถ กตมา สมาโรปนา, ยํ วีริยินฺทฺริยํ, โส สีลกฺขนฺโธ. ยํ สีลํ, เต จตฺตาโร ธมฺมา ปธานา. ยํ ธมฺมานุธมฺม\-ปฏิปตฺติ, โส ปาติโมกฺข\-สํวโร.}

\proseitem{95}{ยสฺส เสโลปมํ จิตฺตนฺติ คาถา\csromansharedfootnote{b44e1b43556aaac4}{ขุ 1. 125 อุทาเน.}, เสโลปมนฺติ อุปมา ยถา เสโล วาเตน น กมฺปติ น อุเณฺหน น สีเตน สํกมฺปติ. ยถา อเนกา อเจตนา, เต อุเณฺหน มิลายนฺติ, สีเตน อวสุสฺสนฺติ, วาเตน ภชนฺติ. น เอวํ เสโล วิรตฺตํ รชนีเยสุ โทสนีเย น ทุสฺสตีติ การณํ โทสนีเย โทมนสฺสนฺตํ, น ทุฏฺเฐน วา กมฺปติ อุเณฺหน วา, โส มิลายติ สีเตน วา อวสุสฺสติ, เอวํ จิตฺตํ ราเคน นานุสฺสติ สีเตน กมฺปตีติ. กิํ การณํ วิรตฺตํ รชนีเยสุ โทสนีเย น ทุสฺสติ. กิํ การณํ, โทสนีเย ปนสฺสนฺติ น ทุสฺสติ, อทุฏฺฐํ ตํ น โกสิสฺสนฺติ, เตน กุปฺปนีเย น กุปฺปติ, ยสฺเสวํ ภาวิตํ จิตฺตํ กุโต ตํ ทุกฺขนิทฺเทโส จ กุโต เอวรูปสฺส ทุกฺขํ อาคมิสฺสตีติ นิทฺทิฏฺฐํ.}

\prose{ปริวตฺตนาติ กุโต ตํ ทุกฺขเมสฺสตีติ ยํ เจตสิกํ สุขํ อนุปาทิเสสา อยํ นตฺถิ โสปาทิเสสา อยํ อตฺถิ. ปุน เอวมาหํสุ ตํ ขณํ ตํ มุหุตฺตํ อุภยเมว อเวทยิตํ โสปาทิเสสํ ยญฺจ อนุปาทิเสสํ ยญฺจ ตํ ขณํ ตํ มุหุตฺตํ อนุปาทิเสสํ ยญฺจ โสปาทิเสสํ จ อเวทยิตํ. สุขมา\-ปนฺนสฺส อนาวตฺติกนฺติ อยเมตฺถ วิเสโส ปริวตฺตนา.}

\prose{\csromanfolio{294}ตตฺถ กตโม เววจโน. ยสฺเสวํ ภาวิตํ จิตฺตํ วา ภาวิตํ สุภาวิตํ อนุฏฺฐิตํ วตฺถุกตํ สุสมารทฺธํ. จิตฺตนฺติ มโน วิญฺญาณํ มนินฺทฺริยํ มโนวิญฺญาณ\-ธาตุ.}

\prose{ตตฺถ กตมา ปญฺญตฺติ. จิตฺตํ มโน สงฺขารา วูปสม\-ปญฺญตฺติยา ปญฺญตฺตํ. สมาธิ อเสกฺข\-ปญฺญตฺติยา ปญฺญตฺโต. ทุกฺขํ อุจฺฉินฺน\-ปญฺญตฺติยา ปญฺญตฺตํ.}

\prose{ตตฺถ กตโม โอตรโณ, จิตฺเต นิทฺทิฏฺเฐ ปญฺจกฺขนฺธา นิทฺทิฏฺฐา โหนฺติ, อยํ ขนฺเธสุ โอตรโณ, มโนวิญฺญาณ\-ธาตุยา นิทฺทิฏฺฐาย อฏฺฐารส ธาตุโย นิทฺทิฏฺฐา โหนฺติ, อยํ ธาตูสุ โอตรโณ. มนายตเน นิทฺทิฏฺเฐ สพฺพานิ อายตนานิ นิทฺทิฏฺฐานิ โหนฺติ. ตตฺถ มนายตนํ นามรูปสฺส ปทฏฺฐานํ. นามรูป\-ปจฺจยา สฬายตนํ. ตถา ปฏิจฺจ\-สมุปฺปาเท. อยํ โอตรโณ.  ตตฺถ กตโม โสธโน สุทฺโธเยว อารมฺโภ.}

\prose{ตตฺถ กตโม อธิฏฺฐาโน. ฉฬินฺทฺริยํ ภาวนา เอกตฺตายํ ปญฺญตฺติ ฉฏฺฐิเตน กาโย เอกตฺตาย ปญฺญตฺโต.}

\prose{ตตฺถ กตโม ปริกฺขาโร. จิตฺตสฺส ปุพฺพเหตุ สมุปฺปาทาย มนสิกาโร จ ตปฺโปณตา จ ยํ อสมาหิต\-ภูมิยํ จ วิเส\-สธมฺมานํ อภาวิตตฺตา จิตฺตสตตํ คจฺฉติ, สเจ สมาธิโน สุขํ เหตุ อวิปฺปฏิสาโร ปจฺจโย, อยํ เหตุ อยํ ปจฺจโย ปริกฺขาโร.}

\prose{ตตฺถ กตมา สมาโรปนา, ยสฺเสวํ ภาวิตนฺติ ตสฺส ธมฺมา สมา\-โรป\-ยิตพฺพา. กาโย สีลํ ปญฺญา ภาวิต\-จิตฺตนฺติ อนภิรตํ อนปณตํ อเนกํ อนุตํ อนาปชฺชาสตฺตํ อยํ สมญฺญายตนา น ตสฺส เสกฺขสฺส สมฺมาสมาธิ สพฺเพ อเสกฺขา ทส อรหนฺตธมฺมา นิทฺทิฏฺฐา โหนฺติ. อเสกฺข\-ภาคิยานิ สุตฺตานิ.}

\proseitem{96}{ยสฺส นูน ภนฺเต กายคตาสติ อภาวิตา, อยํ โส อญฺญตรํ สพฺรหฺมจาริํ\csromansharedfootnote{91342bdc0629b21e}{สพฺรหฺมจารีนํ (ก)} อาสชฺช สมาปชฺช อปฺปฏินิสชฺช ชนปท\-จาริกํ ปกฺกเมยฺย, โส อายสฺมา อิมสฺมิํ วิปฺปฏิชานาติ เทฺว ปชานิ ปฏิชานาติ จิตฺตภาวนายญฺจ ทิฏฺฐิยา ปหานํ, กายภาวนายญฺจ ทิฏฺฐิปฺปหานํ, กายภาวนายญฺจ ตณฺหาปหานํ, ยํ ปฐมํ อุปมํ กโรติ. อสุจินาปิ สุจินาปิ ปถวี เนว อฏฺฏิยติ น ชิคุจฺฉติ น ปีติปาโมชฺชํ ปฏิลภติ, \csromanfolio{295}เอวเมว หิ ปถวีสเมน โส เจตสา อนฺวเยน อปฺปเกน อเวเรน อพฺยาปชฺเชน วิหรามีติ. อิติ โส อายสฺมา กิํ ปฏิชานาติ, กายภาวนาย สุขินฺทฺริย\-ปหานํ ปฏิชานาติ, จิตฺตภาวนาย โสมนสฺสินฺทฺริย\-ปหานํ ปฏิชานาติ. กายิกา เวทนา ราคานุสยม\-นุคตานํ สุขินฺทฺริยํ ปฏิกฺขิปติ. น หิ เวทนา\-กฺขนฺธํ ยา เจตสิกา สุขเวทนา ตตฺถ อยํ ปฏิลาภ\-ปจฺจยา อุปฺปชฺชติ สุขํ โสมนสฺสํ. โสตํ ปฏิกฺขิปติ, น หิ มโน\-สมฺผสฺสชํ เวทนํ. ตตฺถ จตูสุ มหาภูเตสุ รูปกฺขนฺธสฺส อนุสย\-ปฏิฆ\-ปหานํ ภณติ. กาเม รูปญฺจ ตญฺจ อเสกฺข\-ภูมิยํ. กาเย กายานุปสฺสนา ทิฏฺฐ\-ธมฺม\-สุข\-วิหารญฺจ. พเลน จ อุสฺสาเหน จ สพฺพํ มนสิ กตตฺตานํ ปหานํ เม’ทํ กตาลิกาย จ ปุริเสน จ มณฺฑนก\-ชาติเกน จ, เอเตหิ อิมสฺส มาตาปิตุ\-สมฺภูตํ ปจฺจเวกฺขณํ, โส กาเยน จ กายา\-นุปสฺสนาย จ จิตฺเตน จ จิตฺตา\-นุปสฺสนาย จ เทฺว ธมฺเม ธาเรติ. กาย\-กิเลสวตฺถุํ จิตฺเตน จ จิตฺต\-สนฺนิสฺสเย จิตฺเตน สุภาวิเตน สตฺตนฺนํ จ สมาปตฺตีนํ วิหริตุํ ปฏิชานาติ.}

\prose{คหปติปุตฺโต\-ปม\-ตาย จ ยถา คหปติ\-ปุตฺตสฺส นานารงฺคานํ วตฺถ\-กรณฺฑโก ปุณฺโณ ภเวยฺย, โส ยํ ยเทว วตฺถยุคํ ปุพฺพณฺหสมเย อากงฺขติ, ปุพฺพณฺหสมเย นิพฺพาเปติ, เอวํ มชฺฌนฺหิก\-สมเย, สายนฺหสมเย, เอวเมว โส อายสฺมา จิตฺตสฺส สุภาวิตตฺตา ยถารูเปน วิหาเรน อากงฺขติ ปุพฺพณฺหสมยํ วิหริตุํ, ตถารูเปน\csromansharedfootnote{974cad16be724271}{ยถารูเปน (อิ, ก)} ปุพฺพณฺหสมยํ วิหรติ. มชฺฌนฺหิก\-สมเย. สายนฺหสมเย. เตน เวส อายสฺมตา อุปมาย เม อาสิตาย ปถวี วา อนุตฺตรา อินฺทฺริยภาวนา ภาวิตจิตฺเตน. เตน โส อายสฺมา อิทํ อฏฺฐวิธํ ภาวนํ ปฏิชานาติ จตูสุ มหาภูเตสุ, กายภาวนํ อุปกจณฺฑาลํ ปุริสเมตกํ ภวตลากาสุ จิตฺตภาวนํ, อิมาหิ ภาวนาหิ ตาย ภาวนาย จ สมถา ปาริปูริ\-มนฺเตหิ. อิเมหิ จตูหิ ปญฺญาปาริปูริ\-มนฺเตหิ.}

\proseitem{97}{กถํ อุปกจณฺฑาลํ ปฏิกูเลสุ ธมฺเมสุ อปฺปฏิกูลสญฺญี วิหรติ, กาโย ปกติยา อปฺปฏิกูลํ กาเย อุทฺธุมาตก\-สญฺญา สํขิตฺเตน นว สญฺญา อิเม ปฏิกูลา ธมฺมา เจโส อายสฺมา ปฏิกูลโต อชิคุจฺฉิโต กายคตาสติยา ภาวนา\-นุโยคม\-นุยุตฺโต วิหรติ, น หิ ตสฺส ชิคุจฺฉ\-ปฺปหาย จิตฺตํ ปฏิกูลติ.}

\prose{\csromanfolio{296}กถํ อปฺปฏิกูเลสุ ธมฺเมสุ ปฏิกูลสญฺญี วิหรตีติ กาโย สพฺพโลกสฺส อปฺปฏิกูโล ตํ โส อายสฺมา อสุภสญฺญาย วิหรติ. เอวํ อปฺปฏิกูเลสุ ธมฺเมสุ ปฏิกูลสญฺญี วิหรติ.}

\prose{กถํ ปฏิกูเลสุ จ อปฺปฏิกูเลสุ จ อปฺปฏิกูลสญฺญี วิหรตีติ อปิ สพฺโพยํ โลกสฺส ยมิทํ มุณฺโฑ ปตฺตปาณี กุเลสุ ปิณฺฑาย วิจรติ. เตน จ โส อายสฺมา สุวณฺณ\-ทุพฺพณฺเณน อปฺปฏิกูลสญฺญี จิตฺเตน จ กาเยน นิพฺพิทา\-สหคเตน อปฺปฏิกูลสญฺญี, เอวํ ปฏิกูเลสุ อปฺปฏิกูเลสุ จ ธมฺเมสุ อปฺปฏิกูลสญฺญี วิหรติ.}

\prose{กถํ ปฏิกูเลสุ จ ธมฺเมสุ อปฺปฏิกูลสญฺญี วิหรติ. ปฏิกูเลสุ จ ธมฺเมสุ สุภสญฺญิโน อิตฺถิรูเป ปฏิกูเลสุ จ ชิคุจฺฉิโน วินีลก\-วิปุพฺพเก ตตฺถ โส อายสฺมา ปฏิกูลสญฺญี วิหรติ.}

\prose{กถํ ปฏิกูเลสุ ธมฺเมสุ ตทุภยํ อภินิวชฺชยิตฺวา อุเปกฺขโก วิหรติ สโต จ สมฺปชาโน จ. อปฺปฏิกูเลลุ จ ธมฺเมสุ สุภสญฺญิโน อิตฺถิรูเป ปฏิกูเลสุ จ ชิคุจฺฉิโน วินีลก\-วิปุพฺพเก ตทุภยํ อภินิวชฺชยิตฺวา เนตํ มม เนโสหมสฺมิ เนโส เม อตฺตาติ วิหรติ. เอวํ ตทุภยํ อภินิวชฺชยิตฺวา อุเปกฺขโก วิหรติ สโต สมฺปชาโน.}

\prose{อปโร ปริยาโย. เตธาตุโก โลกสนฺนิวาโส สพฺพ\-พาล\-ปุถุชฺชนานํ อปฺปฏิกูล\-สญฺญา. ตตฺถ จ อายสฺมา สาริปุตฺโต อปฺปฏิกูลสญฺญี วิหรติ. เอวํ อปฺปฏิกูเลสุ ธมฺเมสุ ปฏิกูลสญฺญี วิหรติ.}

\prose{กถํ ปฏิกูเลสุ ธมฺเมสุ อปฺปฏิกูลสญฺญี วิหรติ. ปฏิกูล\-สญฺญิโน สพฺพเสกฺขา อิธ กา เตธาตุเก สพฺพโลเก. ตตฺถ กตโม ภูมิปฺปตฺโต สมาธิผเล สจฺฉิกโต อปฺปฏิกูลสญฺญี วิหรติ. กิํ การณํ, น หิ ตํ อตฺถิ ยสฺส โลกสฺส ปหานาย ปฏิกูลสญฺญี อุปฺปาเทยฺย.}

\prose{กถํ ปฏิกูเลสุ จ อปฺปฏิกูเลสุ จ ธมฺเมสุ ปฏิกูลสญฺญี วิหรติ. เตธาตุเก โลกสนฺนิวาเส ยาว กามโลก\-ภูมตา หิ ราคานํ วีตราคานํ ปฏิกูลสมตา รูปารูปธาตุํ อปฺปฏิกูล\-สมตา. ตตฺถ จ อายสฺมา สาริปุตฺโต ปฏิกูลสญฺญี วิหรติ. เอวํ ปฏิกูเลสุ จ อปฺปฏิกูเลสุ จ ธมฺเมสุ ปฏิกูลสญฺญี วิหรติ.}

\prose{\csromanfolio{297}กถํ ปฏิกูเลสุ จ อปฺปฏิกูเลสุ จ ธมฺเมสุ อปฺปฏิกูลสญฺญี วิหรติ, ยํ กิญฺจิ ปรโต ทุรุตฺตานํ ทุราคตานํ วจนปถานํ ตํ วจนํ อปฺปฏิกูลํ ยาวตา วาจโส อปฺปติรูปา ตถา ชนสฺส อปฺปฏิกูล\-สญฺญา. ตตฺถ อายสฺมา สาริปุตฺโต อภิญฺญาย สจฺฉิกโต อปฺปฏิกูลสญฺญี วิหรติ, เอวํ ปฏิกูเลสุ จ อปฺปฏิกูเลสุ จ ธมฺเมสุ อปฺปฏิกูลสญฺญี วิหรติ.}

\proseitem{98}{กถํ ปฏิกูเลสุ จ อปฺปฏิกูเลสุ จ ธมฺเมสุ ตทุภยํ อภินิวชฺชยิตฺวา อุเปกฺขโก จ วิหรติ สโต จ สมฺปชาโน. ยญฺจ เนสํ สมนุปสฺสติ เย ธมฺมา ทุจฺจริตา, เต ธมฺมา อปฺปฏิกูลา. ตตฺถ อายสฺมา สาริปุตฺโต อิติ ปฏิสญฺจิกฺขติ เย ธมฺมา ทุจฺจริตา, เต ธมฺมา อนิฏฺฐวิปากา. เย ธมฺมา สุจริตา, เต อาจยคามิโน. โส จ สุจริตํ อาจยคามินิํ กริตฺวา ทุจฺจริตํ อนิฏฺฐวิปากํ กริตฺวา ตทุภยํ อภินิวชฺชยิตฺวา อุเปกฺขโก วิหรติ.}

\prose{อถ ปฏิกูเลสุ จ ธมฺเมสุ อปฺปฏิกูเลสุ จ ปฏิกูลสญฺญี วิหรติ. ตณฺหา ปฏิกูลธมฺมา กิํ การณํ, ตณฺหาวเสน หิ สตฺตา ทฺวีหิ ธมฺเมหิ สตฺตา, กพฬีกาเร อาหาเร รสตณฺหาย สตฺตา, ผสฺเส สุขสญฺญาย สตฺตา. ตตฺถายสฺมา สาริปุตฺโต กพฬีกาเร จ อาหาเร ปฏิกูลสญฺญี วิหรติ, ผสฺเส จ ทุกฺขสญฺญี วิหรติ. เอวํ ปฏิกูเลสุ จ อปฺปฏิกูเลสุ จ ปฏิกูลสญฺญี วิหรติ.}

\prose{กถํ ปฏิกูเลสุ จ ธมฺเมสุ อปฺปฏิกูเลสุ จ ธมฺเมสุ อปฺปฏิกูลสญฺญี วิหรติ. ตณฺหากฺขยํ อนุตฺตรํ นิพฺพานํ ตถา พาล\-ปุถุชฺชนานํ ปฏิกูลสญฺญา ปหตสญฺญา จ. ตตฺถายสฺมโต สาริปุตฺตสฺส อปฺปฏิกูล\-สญฺญา อพฺยาปาทสญฺญา จ สามํ ปญฺญาย ปสฺสิตฺวา เอวํ ปฏิกูเลสุ จ ธมฺเมสุ อปฺปฏิกูลสญฺญี วิหรติ.}

\prose{กถํ ปฏิกูเลสุ จ อปฺปฏิกูเลสุ จ ธมฺเมสุ อปฺปฏิกูลสญฺญี วิหรติ, ตติเย จ นิพฺพาเน ปฏิกูล\-สญฺญิโน ยเสน จ กิตฺตินิ จ อปฺปฏิกูล\-สญฺญิโน. ตตฺถายสฺมา สาริปุตฺโต อสฺสาทญฺจ อาทีนวญฺจ นิสฺสรณญฺจ ยถาภูตํ สมฺมาปญฺญาย ปฏิชานนฺโต ปฏิกูลญฺจ อปฺปฏิกูลญฺจ ธมฺมํ ตทุภยํ อภินิวชฺชยิตฺวา อปฺปฏิกูลสญฺญี วิหรติ.}

\prose{\csromanfolio{298}กถํ ปฏิกูลํ อปฺปฏิกูลญฺจ ธมฺมํ ตทุภยํ อภินิวชฺชยิตฺวา อุเปกฺขโก วิหรติ สโต จ สมฺปชาโน จ, ยญฺจ สมนุปสฺสติ อนุนโย อปฺปฏิกูโล ธมฺโม ปฏิโฆ จ ปฏิกูโล ธมฺโม, ตตฺถายสฺมา สาริปุตฺโต อนุนยสฺส ปฏิฆ\-ปฺปหีนตฺตา อุเปกฺขโก วิหรติ สโต สมฺปชาโน จ. ยญฺจสฺส สมนุปสฺสติ อยํ ปญฺจวิธา อนุตฺตรา อินฺทฺริยภาวนา. อยํ สุตฺตนิทฺเทโส.}

\proseitem{99}{ตตฺถ กตโม เทสนาหาโร, อิมมฺหิ สุตฺเต กิํ เทสิตพฺพํ, ตตฺถ วุจฺจเต. อิมมฺหิ สุตฺเต ทิฏฺฐ\-ธมฺม\-สุข\-วิหาโร เทสิโต, ตถา วิมุตฺตํ จิตฺตํ ปจฺจเวกฺขณา จ อธิปญฺญา\-ธมฺมํ เทสิตํ.}

\prose{ตตฺถ กตโม วิจโย, เย กาเย กายานุปสฺสิโน วิหรนฺติ, เตสํ จิตฺตํ อนุนย\-ปฺปฏิเฆน น วิหรติ อนุนย\-ปฺปฏิเฆน จาภิรมมานสฺส จิตฺตํ สมคฺคตํ ภวิสฺสตีติ ภาวนาย พลเมตํ, อยํ วิจโย หาโร.}

\prose{ตตฺถ กตโม ยุตฺติหาโร, กายภาวนาย จ จิตฺตภาวนาย จ น กิญฺจิ สพฺรหฺมจารี อติมญฺญิสฺสตีติ. อตฺถิ เอสา ยุตฺติ, อยํ ยุตฺติหาโร.}

\prose{ตตฺถ กตโม ปทฏฺฐาโน หาโร. กายภาวนาย ปฐมสฺส สติ อุปฏฺฐานสฺส ปทฏฺฐานํ. ยา ปถวีสม\-จิตฺตตา, สา อนิจฺจา\-นุปสฺสนาย ปทฏฺฐานํ.}

\prose{ตตฺถ กตโม ลกฺขโณ, ยํ ปถวีสเมน เจตสา วิหรติ อตฺตานุปสฺสี ปถวีสเมน คิหี วิหรติ. โก อตฺโถ ปถวีสเมนาติ ยถา เย จ เสโลปมตาย อกมฺมยุตฺตา เอวเมว ปถวีสโม อยํ หิริยตาย. อยํ ลกฺขโณ.}

\prose{ตตฺถ กตโม จตุพฺยูโห หาโร, อิมมฺหิ พฺยากรเณ โก ตสฺส อายสฺมโต อธิปฺปาโย. เย เกจิ อรหนฺตา อินฺทฺริย\-ภาวนํ อากงฺขิยนฺติ, เต ปถวีสมตํ อุปฺปาทยิสฺสนฺตีติ. อยํ อธิปฺปาโย.}

\prose{ตตฺถ กตโม อาวฏฺโฏติ. นตฺถิ อาวฏฺฏสฺส ภูมิ.}

\prose{\csromanfolio{299}ตตฺถ กตโม วิภตฺติ, โย กายานุปสฺสี วิหรติ, โส ปถวีสม\-จิตฺตตํ ปฏิลภิสฺสตีติ น เอกํเสน. กิํ การณํ, เย ขณฺฑกาทิฉินฺนกาทิโน, น เต ปถวีสม\-จิตฺตตํ ปฏิลภนฺติ. สพฺพา กายคตาสติ เสกฺข\-ภาวนาย นิพฺพานํ ผลํ, อยํ วิภตฺติ.}

\prose{ตตฺถ กตโม ปริวตฺตโน หาโร, เย กายานุปสฺสิโน วิหริสฺสนฺติ, เตสํเยว กายปจฺจยา อุปฺปชฺเชยฺย อาสวา วิฆาต\-ปริฬาหา, อยํ ปริวตฺตโน หาโร.}

\prose{ตตฺถ กตโม โอตรโณ, ปญฺจกฺขนฺธา\csromansharedfootnote{bf3b6e966d364c8e}{สตฺเตสุ จ ปญฺจกฺขนฺธา (อิ)} อวิติณฺณา\csromansharedfootnote{62022b09f864980f}{อวติณฺณา (อิ)} พาวีสติ\-นฺทฺริยานิ, ตถา ยํ มนินฺทฺริยํ, ตํ มโนธาตุ มนายตนญฺจ. ยํ สมาธินฺทฺริยํ, ตํ ธมฺมธาตุ ธมฺมายตนญฺจ. อยํ โอตรโณ หาโร.}

\prose{ตตฺถ กตโม โสธโน หาโร. เย จ มนสา จตฺตาโร ภาเวตพฺพา, เต สพฺเพ ภาวิตา ยํ ตํ มเนน ปหีเน ปตฺตพฺพตํ สพฺพตฺถ เอตสฺส จ อตฺถาย อารมฺโภ, โส อตฺโถ สุทฺโธ. อยํ โสธโน หาโร.}

\prose{ตตฺถ กตโม อธิฏฺฐาโน, อยํ สมาธิ เอกตฺตตาย ปญฺญตฺโต, ฉ กายา เอกตฺตตาย ปญฺญตฺตา. ปญฺจินฺทฺริยานิ รูปีนิ รูปกาโย. ฉ เวทนากายา เวทนากาโย. ฉ สญฺญากายา สญฺญากาโย, ฉ เจตนากายา เจตนากาโย. ฉ วิญฺญาณกายา วิญฺญาณกาโย. สพฺเพปิ เอเต ธมฺมา ธมฺมกาโยติเยว สงฺขํ คจฺฉนฺติ. อยํ อธิฏฺฐาโน.}

\prose{ปริกฺขาโรติ สมาปตฺติ\-โกสลฺลญฺจ วีถิโกสลฺลญฺจ\csromansharedfootnote{3f78f0a38fd968a5}{มีติโกสลฺลญฺจ (อิ)} เหตุ. ยญฺจ โคจรโกสลฺลํ ยญฺจ กลฺลํ ตํ โกสลฺลํ ปจฺจโย. โวทานโกสลฺลํ เหตุ, กลฺลํ ปจฺจโย. สุขํ เหตุ, อพฺยาปชฺชํ ปจฺจโย. อยํ ปริกฺขาโร.}

\prose{ตตฺถ กตโม สมาโรปโนติ ยถา ปถวี สุจิมฺปิ นิกฺขีปนฺเต อสุจิมฺปิ นิกฺขิตฺเต ตาทิเสเยว เอวํ กาโย มนาปิเกหิปิ ผสฺเสหิ อมนาปิเกหิปิ ผสฺเสหิ ตาทิโสเยว ปฏิฆ\-สมฺผสฺเสน วา สุขาย เวทนาย ตาทิสํ โย จิตฺตํ. อิทํ สุตฺตํ วิภตฺตํ สโอปมฺมํ อุคฺฆฏิต\-ญฺญุสฺส ปุคฺคลสฺส วิภาเคน. ตตฺถ สมาโรปนาย อวกาโส นตฺถิ.}

\proseitem{100}{\csromanfolio{300}ตตฺถ กตมํ สุตฺตํ สํกิเลส\-ภาคิยํ, ยโต จ กุสเลหิ ธมฺเมหิ น วิโรธติ, น วฑฺฒติ, อิมํ อาทีนวํ ภควา เทเสติ, ตสฺมา ฉนฺนํ วิวเรยฺย, วิวฏํ นาติวสฺสติ, ตโต อาทีนวโต วิวเรยฺยาติ ตํ ตีหิ ธมฺเมหิ นาภิ\-ธํสิตาติ อสุภสญฺญาย ราเคน นาภิธํ\-สิยติ. เมตฺตาย โทเสน นาภิธํ\-สิยติ. วิปสฺสนาย โมเหน นาภิธํ\-สิยติ. เอวญฺจสฺส โย โย ธมฺโม ปฏิปกฺโข ตมฺหิ ตมฺหิ ธมฺเม ปริปูริสฺสติ. โย ตสฺส ธมฺมสฺส อกุสโล ธมฺโม ปฏิปกฺโข, เตน นาธิวาสิยติ.}

\prose{อปโร ปริยาโย เย อิเม ธมฺมา อตฺตนา น สกฺโกติ วุฏฺฐานํ, เต เอเต ธมฺมา เทสิตา. ฉนฺนม\-ติวสฺสตีติ เตหิ วิตกฺกํ เยน จ สกฺกา ปุน เทสิตํ จิตฺตํ วิภาเวตุํ ปริโยทาเปตุํ วิเวกนินฺนสฺส วิเวกโปณสฺส วิเวก\-ปพฺภารสฺส วุทฺธิํ วิรูฬฺหิํ เวปุลฺลตํ อาปชฺชติ กุสเลสุ ธมฺเมสุ, เสยฺยถาปิ นาม อุปฺปลํ วา กุมุทํ วา ปทุมํ วา อุทเก สุกฺกปกฺเข จนฺโท ยาวรตฺติ ยาวทิวโส อาคจฺฉติ, ตสฺส วุทฺธิเยว ปาฏิกงฺขิตพฺพา, น ปริหานิ, เอวํวิธํ ตํ จิตฺตํ นาภิธํ\-สิยติ. อปโรเปตฺถ โย อกูโฏ อสโฐ อมายาวี อุชุ ปุริโส ยถาภูตํ อตฺตานํ อาวิกโรติ. ตตฺถ โย ฉาเทติ ตสฺส อกุสลา ธมฺมา จิตฺตํ อนุธาวนฺติ. ฉนฺนม\-ติวสฺสตีติ โย ปน โหติ อสโฐ อกูโฏ อมายาวี อุชุ ปุริโส ยถาภูตํ อตฺตานํ อาวิกโรติ. ตสฺส จิตฺตํ อกุสเลหิ ธมฺเมหิ น วิทฺธํสิยติ, อยํ สุตฺตตฺโถ.}

\proseitem{101}{ตตฺถ กตมา เทสนา, อิธ เทสิตา ทส อกุสล\-กมฺมปถา อธิวสฺส\-นตาย ทส กุสล\-กมฺม\-ปถา อนธิ\-วสฺส\-นตาย อกุสเลหิ น วิสุชฺฌติ. ยถา วุตฺตํ ภควตา “จิตฺตสํกิเลสา ภิกฺขเว สตฺตา สํกิลิสฺสนฺตี”ติ.}

\prose{ตตฺถ กตโม วิจโย. ยสฺเสวํ จิตฺตํ อธิวาสิยติ, ตสฺส พุชฺฌิตสฺส ยํ ภเวยฺย กูเฏยฺย, ตํ อานนฺตริเยนปิ สตฺถริ วา คุณา\-นุกมฺป\-นตาย, อยํ วิจโย}

\prose{ตตฺถ กตมา ยุตฺตีติ เอวํ อนธิวสิ\-ยนฺตํ จิตฺตํ วุฏฺฐาติ. วุฏฺฐิตํ ปติฏฺฐหติ กุสเลสุ ธมฺเมสูติ อตฺถิ เอสา ยุตฺติ.}

\prose{\csromanfolio{301}ปทฏฺฐานนฺติ ฉนฺนม\-ติวสฺสตีติ ฉนฺนํ อสํวรานํ ปทฏฺฐานํ, วิวฏํ นาติวสฺสตีติ อฉนฺนํ สํวรณานํ. ตสฺมา ฉนฺนํ วิวเรยฺย วิวฏํ นาติวสฺสตีติ เทสนาย ปทฏฺฐานํ.}

\prose{ลกฺขโณติ ฉนฺนม\-ติวสฺสตีติ เย เกจิ วิจิตฺเตน ฉนฺเนน เอกลกฺขณา ธมฺมา สพฺเพ เต อวิทฺธํ\-สิยนฺติ. ตสฺมา ฉนฺนํ วิวเรยฺย. วิวฏํ นาติวสฺสตีติ เย เกจิ เตน อจฺฉนฺเนน เอกลกฺขณา ธมฺมา สพฺเพ เต นาติวสฺสนฺตีติ ลกฺขโณ หาโร.}

\prose{ตตฺถ กตโม จตุพฺยูโห หาโร, อิมมฺหิ สุตฺเต ภควโต โก อธิปฺปาโย, เยสํ เกสญฺจิ จิตฺตํ อกุสลา ธมฺมา อธิปฏิเทสิตา เต ยถาธมฺมํ ปฏิกริสฺสนฺตีติ อยํ ตตฺถ ภควโต อธิปฺปาโย. อยํ จตุพฺยูโห หาโร.}

\prose{อาวฏฺโฏติ ยํ ฉนฺนํ ตํ ทุวิธํ กมฺปมานํ สมุจฺฉิตพฺโพ. อานนฺตริย\-สมาธีนํ, ตตฺถ ปสฺสทฺธิยญฺจ มาโน อาสเว วฑฺเฒติ, อสฺสทฺธิเยน จ ปมาทํ คจฺฉติ, ปมาเทน โอนมติ, อุนฺนฬภาวํ คจฺฉติ. วุตฺตํ เจตํ ภควตา “อุนฺนฬานํ ปมตฺตานํ เตสํ วฑฺฒนฺติ อาสวา”ติ จตฺตาริ ตานิ อุปาทานานิ, ยานิ จตฺตาริ อุปาทานานิ, เต ปญฺจุ\-ปาทาน\-กฺขนฺธา ภวนฺติ. อิมานิ สจฺจานิ ทุกฺขญฺจ สมุทโย จ. ตสฺมา ฉนฺนํ วิวเรยฺยาติ เยน เหตุนา, เต อาสวา วฑฺฒนฺติ. เตสํ ปหีนตฺตา อาสวา ปหียนฺเต. ตตฺถ อปฺปมาเทน อสฺสทฺธิยํ ปหียติ อุทฺธจฺจกุกฺกุจฺจ\-ปฺปหาเนน โอฬาริกตา ตสฺส เทฺว ธมฺมา น สมโถ จ ภาวนา จ ปาริปูริํ คจฺฉนฺติ. โย เตสํ อาสวานํ ขโย, อยํ นิโรโธ. อิมานิ จตฺตาริ สจฺจานิ, อยํ อาวฏฺโฏ.}

\prose{ตตฺถ กตโม วิภตฺติ หาโร. ฉนฺนม\-ติวสฺสตีติ น เอกํโส. กิํ การณํ, ยสฺส อสฺสา นิวตฺตนา ยถาปิ เสกฺขานํ. ยถาวุตฺตํ ภควตา\csromandash{}}

\setlength{\csromangathapagewidth}{0pt}
\setlength{\csromangathawakwidth}{0pt}
\csromangathameasuregroup{}{“กิญฺจาปิ เสกฺโข ปกเรยฺย ปาปํ, \\ กาเยน วาจาย อุท เจตสา วา. \\ อภพฺโพ หิ ตสฺส ปริคุหนาย, \\ อภพฺพตา ทิฏฺฐปทสฺส โหตี”ติ.}
\csromangathameasurewak{“กิญฺจาปิ เสกฺโข ปกเรยฺย ปาปํ, \\ กาเยน วาจาย อุท เจตสา วา. \\ อภพฺโพ หิ ตสฺส ปริคุหนาย, \\ อภพฺพตา ทิฏฺฐปทสฺส โหตี”ติ.}
\setlength{\csromangathaleftcol}{0pt}
\csromangathagroup{\csromangathastanza{\csromangathawak{“กิญฺจาปิ เสกฺโข ปกเรยฺย ปาปํ, \\ กาเยน วาจาย อุท เจตสา วา. \\ อภพฺโพ หิ ตสฺส ปริคุหนาย, \\ อภพฺพตา ทิฏฺฐปทสฺส โหตี”ติ.}}}

\prose{\csromanfolio{302}กิญฺจาปิ เตสํ นิวารณํ จิตฺตํ โหติ. อปิ ตุ อปฺปจฺจยา สมาเย จ เต นิทฺทิสิตพฺพา, อยํ วิภตฺติหาโร.}

\prose{ตตฺถ กตโม ปริวตฺตโน หาโร, ฉนฺนม\-ติวสฺสตีติ ยสฺส เย ธมฺมา สพฺพํ อนวิวฏํ อติวสฺสิยติ, วิวฏํ นาติวสฺสติ, อวคุณนฺตํ นาติวสฺสติ. อยํ ปริวตฺตโน หาโร.}

\prose{ตตฺถ กตโม เววจโน หาโร, ฉนฺนนฺติ อาวุตํ นิวุตํ ปิหิตํ ปฏิกุชฺชิตํ สญฺฉนฺนํ ปโรธํ, วิวฏํ นาติวสฺสตีติ ยสฺส เต ธมฺมา ปพฺพชฺชิตา วิโนทํ นาธิวสฺสิตา วนฺติกตาติ, อยํ เววจโน หาโร.}

\prose{ตตฺถ กตโม ปญฺญตฺติ หาโร, ฉนฺนม\-ติวสฺสตีติ กิเลส\-ภาคิย\-ปญฺญตฺตํ วิวฏํ นาติวสฺสตีติ สธมฺมกิจฺจํ ยํ ปฏิปทา ปญฺญตฺติยา ปญฺญตฺตํ, ตสฺมา หิ ฉนฺนํ วิวเรยฺยาติ อนุสาสน\-ปญฺญตฺติยา ปญฺญตฺตํ, วิวฏํ นาติวสฺสตีติ นิทฺธานปญฺญตฺติยา ปญฺญตฺตํ, อยํ ปญฺญตฺติ หาโร.}

\prose{ตตฺถ กตโม โอตรโณ หาโร. ฉนฺนม\-ติวสฺสตีติ ตโย กิเลสา ราโค โทโส โมโห, เต ขนฺเธสุ สงฺขาร\-กฺขนฺโธ ฯเปฯ. เต ปุรา ยถา นิทฺทิฏฺฐํ ขนฺธ\-ธาตุ\-อายตเนสุ, อยํ โอตรโณ หาโร.}

\prose{ตตฺถ กตโม โสธโน หาโร, เยนารมฺเภน อิทํ สุตฺตํ ภาสติ โส อารมฺโภ นิยุตฺโต.}

\prose{อธิฏฺฐาโนติ ฉนฺนม\-ติวสฺสตีติ เอกตฺตตาย ปญฺญตฺตํ กิํ การณํ, อิทํ หิ อติวสฺสตีติ อิมสฺส จ อติวสฺสติ เอวญฺจ อติวสฺสตีติ อยํ เวมตฺตตาย ยา สุณสาธารเณหิ ลกฺขเณหิ ปญฺญาปิยติ, สา เอกตฺตปญฺญตฺติ.}

\prose{ตตฺถ กตโม ปริกฺขาโร, ยญฺจ ตํ อติวสฺสิยนฺติ, ตสฺส เทฺว เหตู เทฺว ปจฺจยา อกุสล\-ปสุเตว วาจ\-กตฺตาภิรติ จ. อิเม เทฺว อโยนิโส\-มนสิกาโร จ กุสลา ธมฺมา โวปสคฺคา จ, อิเม เทฺว ปจฺจยา.}

\prose{ตตฺถ กตโม สมาโรปโน, ฉนฺนม\-ติวสฺสตีติ เวมติ ปสฺสตีติ ฉนฺนํ ยํ ปริคฺคหิตุํ ยํ อเทสิตุํ อปฺปสฺสุตํ ยํ กถํกถา วิภูเตน อกุสลมูเลน ยํ ตณฺหาย จ เต วฑฺฒติ โทสาติ สนฺนิตฺวา เต อปฺปสกฺขเยน สงฺขารา. สงฺขาร\-ปจฺจยา วิญฺญาณํ ยาว ชรามรณํ, อยํ \csromanfolio{303}สมาโรปโน. ยํ ปุน ตถา เทสนา, ตสฺเสว อกุสลา ธมฺมา วุทฺธิํ วิรูฬฺหิํ เวปุลฺลตมา\-ปชฺชติ ตสฺส สงฺขารา นิโรธา, อยํ สมาโรปโน.}

\proseitem{102}{จตฺตาโร ปุคฺคลา\csromansharedfootnote{dea62cd07e8d3cf7}{อํ 1. 387 ปิฏฺเฐ.} ตโม\-ตมปรายโนติ ฯเปฯ. ตตฺถ กตโม วุจฺจเต ตโม นาม, โย ตโม อนฺธกาโร, ยถา วุตฺตํ ภควตา “ยถา อนฺธกาเร ตสฺมิํ ภยานเก สกมฺปิธาตุปุริโส น ปสฺสติ, เอวเมว อญฺญาณโต ตโม\-ปน\-นฺธกาโร ปาปก\-สกมฺม\-สวิปากํ น สทฺโธ โหติ. อิติ เอวํ ลกฺขณตา อญฺญาณํ ตโม อวิชฺชา โมโห, เยน สตฺตา ยถาภูตํ นปฺปชานนฺติ, อิติ วุจฺจติ ตโมติ. โส ติณฺณํ จกฺขูนํ ตโม มํสจกฺขุโน ทิพฺพจกฺขุโน ปญฺญาจกฺขุโน, อิเมสํ จกฺขูนํ อิธ ตโม นิทฺทิสิยติ อญฺญาณนฺติ. ตตฺถ กตมํ อญฺญาณํ อทสฺสนํ, อถ นิสฺสเย ยํ ปุพฺพนฺเต อญฺญาณํ อปรนฺเต อญฺญาณํ ปุพฺพนฺตา\-ปรนฺเต อญฺญาณํ เหตุมฺหิ อญฺญาณํ ปจฺจยมฺหิ อญฺญาณํ ตสฺส อญฺญาณิโน สมาธิ\-ภูตสฺส เอโส นิสฺสนฺโท. ยํ น ชานาติ อิทํ เสวิตพฺพํ อิทํ น มนสิ\-กาตพฺพนฺติ. โส เตน ตเมน นิทฺทิสิยติ ตโมปิ ยถา วุจฺจติ. มูโฬฺหติ เอวํ เจตนา. เตน ตเมน โส ปุคฺคโล วุจฺจติ. ตโมติ โส เตน ตเมน อสมูหเตน อสมุจฺฉินฺเนน ตปฺปรโม ภวติ ตปฺปรายโน, อยํ วุจฺจติ ปุคฺคโล ตโม ตมปรายโนติ. ปรายโนเยว ธมฺโม มนสิกาตพฺโพ โส ตโม ทหติ อญฺญจิตฺตํ อุปฏฺฐเปติ. เต จสฺส ธมฺมา นิชฺฌาน\-กฺขมนฺติ. โส สุตมยาย ปญฺญาย สมนุปสฺสติ. (1)}

\prose{ตตฺถ กตโม ตโม โชติปรายโน. โส เตน ปญฺญาวเสน อิริยติ เอวํ ตสฺเสว อิริยนฺตสฺส ปรายโน ภวติ. อยํ วุจฺจเต ปุคฺคโล ตโม โชติปรายโน. (2)}

\prose{ตตฺถ กตโม ปุคฺคโล โชติ โชติปรายโน\csromansharedfootnote{96698975a4fb5e17}{โชติปรายโน (อิ)}, ตตฺถ วุจฺจติ โชติ นาม ยํ ตสฺส เจ ตมสฺส ปฏิปกฺเขน เย จ ธมฺเม อนฺตมโส ญาณาโลโก, โส สุณธมฺโม ปุคฺคโล ตโม โชติปรายโน, ตตฺถ วุจฺจเต, โยยํ ปุคฺคโล ตโม โชติปรายโน, โส ยทิ ตถารูปํ กลฺยาณมิตฺตํ ปฏิลภติ, โย นํ อกุสลโต จ \csromanfolio{304}นิวาเรติ ภาวิต\-กุสลตาว ภาวี นิโยเชตีติ. เอวญฺจ สทฺธมฺมํ เทเสติ. อิเม ธมฺมา กุสลา, อิเม ธมฺมา อกุสลา. อิเม ธมฺมา สาวชฺชา, อิเม ธมฺมา อนวชฺชา. อิเม ธมฺมา เสวิตพฺพา, อิเม ธมฺมา น เสวิตพฺพา. อิเม ธมฺมา ภชิตพฺพา. อิเม ธมฺมา น ภชิตพฺพา. อิเม ธมฺมา อุปสมฺปชฺช วิหาตพฺพา, อิเม ธมฺมา น อุปสมฺปชฺช วิหาตพฺพา. อิเม ธมฺมา มนสิกาตพฺพา, อิเม ธมฺมา น มนสิกาตพฺพาติ. ปจฺจเต สญฺญาย ยถา สญฺญายติ สตินฺทฺริยานิ, โส เอวํ ปชานาติ. อิเม ธมฺมา กุสลา, อิเม ธมฺมา อกุสลา. อิเม ธมฺมา สาวชฺชา, อิเม ธมฺมา อนวชฺชา. อิเม ธมฺมา เสวิตพฺพา, อิเม ธมฺมา น เสวิตพฺพา. อิเม ธมฺมา ภาเวตพฺพา, อิเม ธมฺมา น ภาเวตพฺพา. อิเม ธมฺมา อุปสมฺปชฺช วิหาตพฺพา, อิเม ธมฺมา น อุปสมฺปชฺช วิหาตพฺพา. อิเม ธมฺมา มนสิกาตพฺพา, อิเม ธมฺมา น มนสิกาตพฺพาติ. โส เต ธมฺเม สุสุยฺยติ, โสตํ โอทหติ, อญฺญํ จิตฺตํ อุปฏฺฐเปติ, เต จสฺส ธมฺมา นิชฺฌาน\-กฺขมนฺติ, โส สุตมยาย ปญฺญาย สมนฺนาคโต โส เตน ปจฺจยวเสน อิริยติ เอวํ ตสฺเสว อิริยนฺติ ตปฺปรโม ภวติ ตปฺปรายโน. อยํ วุจฺจเต ปุคฺคโล ตโม ตมปรายโน. (3)}

\prose{ตตฺถ กตโม ปุคฺคโล โชติ ตมปรายโน, โชติ นาม ยา ตสฺเสว ตมสฺส ปฏิปกฺเขน เย ธมฺมา อนฺตมโส ญาณาโลโก, โส ปุน ธมฺโม. กตมา อุจฺจเต ปญฺญายโต ปณฺฑิโตติ วุจฺจเต, โส เอวํ ปชานาติ. อิเม ธมฺมา กุสลา, อิเม ธมฺมา อกุสลา. อิเม ธมฺมา สาวชฺชา, อิเม ธมฺมา อนวชฺชา. อิเม ธมฺมา เสวิตพฺพา, อิเม ธมฺมา น เสวิตพฺพา. อิเม ธมฺมา ภาวิตพฺพา, อิเม ธมฺมา น ภาวิตพฺพา. อิเม ธมฺมา อุปสมฺปชฺช วิหาตพฺพา, อิเม ธมฺมา น อุปสมฺปชฺช วิหาตพฺพา. อิเม ธมฺมา มนสิกาตพฺพา, อิเม ธมฺมา น มนสิกาตพฺพา. อิธ ปน ปาปมิตฺต\-สํเส\-วโน ปาปมิตฺต\-วสานุโค อกุสเล ธมฺเม อภิวฑฺเฒติ, กุสเล ธมฺเม ปชหติ. โส เตน ปมาเทน ปจฺจยสญฺญา อมนสิกตฺวา อสฺสติ\-อสมฺปชญฺญํ อาเสวติ. ตยา โย ปฏิปกฺโข ตโม, โส ปวฑฺเฒติ. โส ตมาภิภูโต ปรายโน ตมปรโม เจว ภวติ. อยํ วุจฺจติ ปุคฺคโล โชติ ตมปรายโน. (4)}

\proseitem{103}{ตตฺถ กตโม ปุคฺคโล โชติ โชติปรายโน. ตตฺถ วุจฺจเต โสยํ ปุคฺคโล กลฺยาณ\-มิตฺตสฺส สนฺนิสฺสิโต ภวติ สกฺกา สํโยคี \csromanfolio{305}กุสลํ คเวสี, โส กลฺยาณมิตฺเต อุปสงฺกมิตฺวา ปริปุจฺฉติ, ปริปญฺหยติ, กิํ กุสลํ, กิํ อกุสลํ. กิํ สาวชฺชํ, กิํ อนวชฺชํ. กิํ เสวิตพฺพํ, กิํ น เสวิตพฺพํ. กิํ ภาวิตพฺพํ, กิํ น ภาวิตพฺพํ. กิํ อุปสมฺปชฺช วิหาตพฺพํ, กิํ น อุปสมฺปชฺช วิหาตพฺพํ. กิํ มนสิกาตพฺพํ, กิํ น มนสิกาตพฺพํ. กถํ สํกิเลโส โหติ, กถํ โวทานํ โหติ. กถํ ปวตฺติ โหติ, กถํ นิวตฺติ โหติ. กถํ พนฺโธ โหติ, กถํ โมกฺโข โหติ. กถํ สกฺกาย\-สมุทโย โหติ, กถํ สกฺกายนิโรโธ โหติ. กถํ สกฺกาย\-สมุทโย โหติ, กถํ สกฺกายนิโรโธ โหติ. โส เอตฺถ เทสิตํ ยถา อุปฏฺฐิตํ ตถา สมฺปฏิปชฺชนฺโต โส เอวํ ปชานาติ. อิเม ธมฺมา กุสลา, อิเม ธมฺมา อกุสลา. เอวํ ฯเปฯ. ยาว กถํ สกฺกาย\-สมุทโย โหติ, กถํ สกฺกายนิโรโธ โหตีติ วิตฺถาเรน กาตพฺพํ. โส เต ธมฺเม อธิปาฏิกงฺขาติ เอวํ ลกฺขณํ ญาณํ วิชฺชา อาโลกํ วฑฺเฒติ. โส ปุคฺคโล ตปฺปรโม ภวติ ตปฺปรายโน, อยํ วุจฺจเต ปุคฺคโล โชติ โชติปรายโน.}

\prose{ตตฺถ กตโม ปุคฺคโล ตโม ตมปรายโน. โย อกุสลํ ธมฺมํ ทีเปติ. ตํ ภาวนาย หีนาสุ คตีสุ อุปปตฺติํ ทสฺเสติ, ตปฺปรโม ภวติ ตปฺปรายโน. อยํ วุจฺจเต ปุคฺคโล ตโม ตมปรายโน.}

\prose{ตตฺถ โย ปุคฺคโล ตโม โชติปรายโน. โส ตเมน อกุสลสฺส กมฺมสฺส วิปากํ ทสฺเสติ. ตเมติ ยํ จกฺขุ กลฺยาณ\-มิตฺตสฺส เยน อกุสเล ธมฺเม ปชหติ, กุสเล ธมฺเม อภิวฑฺฒติ.}

\prose{ตตฺถ โย จ ปณีตาสุ คตีสุ อุปปตฺติํ ทสฺเสติ, ตปฺปรโม เตน วุจฺจเต ตโม โชติปรายโน.}

\prose{ตตฺถ โย ปุคฺคโล โชติ ตมปรายโน. กุสลสฺส กมฺมวิปากํ ทสฺเสติ. ยํ จกฺขุ ปาป\-มิตฺตสํสคฺเคน ปาปมิตฺตุ\-ปเสเวน ปาปมิตฺต\-วสานุโค อกุสลํ ธมฺมํ อภิวฑฺฒติ, ตํ ภาวนาย หีนาสุ คตีสุ อุปปตฺติํ ทสฺเสติ. ตปฺปรโม เตน วุจฺจเต โชติ ตมปรายโน.}

\prose{ตตฺถ โย ปุคฺคโล โชติ โชติปรายโน. โส โชติตา ปภาตา\csromansharedfootnote{29137a15330652f7}{โชติตภาวตาย (อิ)} ยาว ปณีตาสุ คตีสุ อุปปตฺติํ ทสฺเสติ. ตปฺปรโม เตนาห โชติ โชติปรายโน.}

\prose{\csromanfolio{306}โชติ\-ตมปรายเนน ทส อกุสลานํ กมฺมานํ อุทยํ ทสฺเสติ. ตเมน ปุคฺคเลน อกุสลานํ กมฺมานํ วิปากํ ทสฺเสติ. น อกุสลานํ ธมฺมานํ วิปากํ ทสฺเสติ. ตเมน อฏฺฐ มิจฺฉตฺตานิ ทสฺเสติ. โชตินา อฏฺฐ สมฺมตฺตานิ ทสฺเสติ. โชตินา ตมปรายเนน ทส อกุสล\-กมฺมปเถ ทสฺเสติ. โชตินา ปณีตตฺตํ ทสฺเสติ. ตเมน โชติปรายเนน อตปนียํ ธมฺมํ ทสฺเสติ. โชตินา ตมปรายเนน ตปนียํ ธมฺมํ ทสฺเสติ. อยํ สุตฺตตฺโถ.}

\proseitem{104}{ตตฺถ กตโม เทสนา หาโร. อิมมฺหิ สุตฺเต กิํ เทสิตํ. ตตฺถ วุจฺจเต อิมมฺหิ สุตฺเต กุสลากุสลา ธมฺมา เทสิตา. กุสลากุสลานญฺจ ธมฺมานํ วิปาโก เทสิโต. หีนปฺปณีตานญฺจ สตฺตานํ คติ นานากรณํ เทสิตํ. อยํ เทสนา หาโร.}

\prose{ตตฺถ กตโม วิจโย หาโร. อกุสลสฺส กมฺมสฺส โย วิปากํ ปจฺจนุโภติ. ตตฺถ ฐิโต อกุสเล ธมฺเม อุปฺปาทิยติ วิจยนฺตํ ยุชฺชติ. กุสลสฺส กมฺมสฺส โย วิปากํ ปจฺจนุโภติ. ตตฺถ ฐิโต กุสเล ธมฺเม อุปฺปาทิยติ วิจยนฺตํ ยุชฺชติ. อยํ วิจโย ยุตฺติ จ.}

\prose{ตตฺถ กตโม ปทฏฺฐาโน หาโร. โย ปุคฺคโล โชติ, โส ปจฺจเวกฺขณาย ปทฏฺฐานํ. โย ปุคฺคโล ตโม, โส ตมาทินฺนํ วานุปสฺสนาย ปทฏฺฐานนฺติ ทสฺเสติ. ตเมน โชติปรายเนน อปฺปมาทสฺส ปทฏฺฐานํ ทสฺเสติ, ตโม อวิชฺชาย จ ทิฏฺฐิยา จ ปทฏฺฐานํ ทสฺเสติ. โชตินา ตมปรายเนน ปมาทสฺส จ ทิฏฺฐิยา จ ปทฏฺฐาน ทสฺเสติ. อยํ ปทฏฺฐาโน.}

\prose{ตตฺถ กตโม ลกฺขโณ หาโร. ตเมน ตมปรายเนน ตโมติ อวิชฺชาย นิทฺทิฏฺฐาย สพฺพกิเลส\-ธมฺมา นิทฺทิฏฺฐา โหนฺติ. ตเมน โชติปรายเนน โชติวิชฺชาย นิทฺทิฏฺฐาย สพฺเพ โพธิปกฺขิยธมฺมา นิทฺทิฏฺฐา โหนฺติ. โชติ\-ตมปรายเนน ปมาโท นิทฺทิฏฺโฐ โหติ. ตเมน โชติปรายเนน อปฺปมาโท นิทฺทิฏฺโฐ โหติ. อยํ ลกฺขโณ หาโร.}

\prose{ตตฺถ กตโม จตุพฺยูโห หาโร, อิมมฺหิ สุตฺเต ภควโต โก อธิปฺปาโย. เย สตฺตา นีจกุลิโน, น เต อิมํ สุตฺวา กุสเล ธมฺเม สมาทาย วตฺติสฺสนฺติ. เย สตฺตา อุจฺจกุลิโน, เต อิมํ ธมฺมเทสนํ \csromanfolio{307}สุตฺวา ภิยฺโยโส มตฺตาย กุสเล ธมฺเม สมาทาย วตฺติสฺสนฺตีติ. อยํ จตุพฺยูโห หาโร. ภูมิยํ อุปเทโส.}

\prose{ตตฺถ กตโม อาวฏฺโฏ หาโร, ยา อวิชฺชาโต ปภูติ ตณฺหา. อยํ สมุทโย. โย ตโม ตมปรายโน, อิทํ ทุกฺขํ. อิมานิ เทฺว สจฺจานิ ทุกฺขญฺจ สมุทโย จ โชติ เยน สุตฺเตน ธมฺเมน ปญฺญาปิยติ, โส ธมฺโม ปญฺญินฺทฺริยสฺส ปทฏฺฐานํ. เตน อโมเหน ตีณิ กุสลมูลานิ ปาริปูริํ คจฺฉนฺติ สคฺคสฺส ปทฏฺฐานํ.}

\prose{ตตฺถ กตมา วิภตฺติ. ตโม ตมปรายโนติ น เอกํ เสน, กิํ การณํ, อตฺถิ ตโม จ ภโว อปราปริย\-เวทนีเยน จ กุสเลน โชตินา ปุคฺคเลน สโหปตฺติภาเว. อตฺถิ โชติ จ ภโว อปราปริย\-เวทนีเยน จ อกุสเลน ตเมน ปุคฺคเลน สโหปตฺติภาเว ปริวตฺตนา ตเมสุ ปฏิปกฺโขติ โชตินา ตมปรายโน.}

\prose{ตตฺถ กตโม เววจโน, โย ตโม, โส เอวํ อตฺตพฺยาปาทาย ปฏิปนฺโน, โส อสฺสทฺธาย พาโล อกุสโล อพฺยตฺโต อนาทีนวทสฺสี. โย โชติ, โส อตฺตหิตาย ปฏิปนฺโน ปณฺฑิโต กุสโล พฺยตฺโต อาทีนวทสฺสี. อยํ เววจโน.}

\prose{ตตฺถ กตมา ปญฺญตฺติ, โส ปุคฺคโล วิปาก\-ปญฺญตฺติยา ปญฺญาปิยติ, อกุสเล ปริยาทินฺนตา ปญฺญาปิยติ. โชติ\-กุสลธมฺมุ\-ปปตฺติ\-ปญฺญตฺติยา ปญฺญาปิยติ กุสลธมฺม\-วิปาก\-ปญฺญตฺติยา จาติ.}

\prose{โอตรโณติ เย อวิชฺชาปจฺจยา สงฺขารา ยญฺจ ชรามรณํ ยา จ อวิชฺชา, ตํ ปทฏฺฐานํ, นิทฺเทเสน วิชฺชุปฺปาโท อวิชฺชานิโรโธ โย ยาว ชรามรณ\-นิโรโธ, อิเม เทฺว ธมฺมา สงฺขารกฺขนฺธ\-ปริยาปนฺนา. ธมฺมธาตุ ธมฺมายตนญฺจ ปทฏฺฐานํ นิทฺเทเสน ธาตูสุ.}

\prose{ตตฺถ กตโม โสธโน, อิมสฺส สุตฺตสฺส เทสิตสฺส อารมฺโภ. อธิฏฺฐาโนติ ตโมติ ภควา พฺรวีติ, น เอกํ ปุคฺคลํ เทเสติ. ยาวตา สตฺตานํ คติ, ตตฺถ เย ทุจฺจริต\-ธมฺเมน อุปปนฺนา, เต พหุลา\-ธิวจเนน ตโม นิทฺทิสติ. ยา โชติ สพฺพสตฺเตสุ กุสลธมฺโม\-ปปตฺติ สพฺพํ ตํ โชตีติ อภิลปติ อยเมกตา ปจฺจโย โยนิโสมนสิการ\-ปญฺญตฺติ จตุนฺนํ มหาภูตานํ ปุคฺคลานํ.}

\prose{\csromanfolio{308}ตตฺถ กตโม ปริกฺขาโร, อกุสลสฺส ปาปมิตฺตตา ปจฺจโย, อโยนิโส มนสิกาโร เหตุ. กุสลสฺส กลฺยาณมิตฺตตา ปจฺจโย, โยนิโส มนสิกาโร เหตุ.}

\prose{ตตฺถ กตมา สมาโรปนาติ, อิเธกจฺโจ นีเจ กุเล ปจฺจาชาโต โหตีติ นีเจ กุเล ปจฺจาชาโต รูเปสุ สทฺเทสุ คนฺเธสุ รเสสุ ผสฺเสสุ, โส อุปปนฺโน สพฺพมฺหิ มานุสฺสเก อุปโภค\-ปริโภเค. โชติ ปณีเตสุ กุสเลสุ อุปปนฺโน สพฺพมฺหิ มานุสฺสเก อุปโภค\-ปริโภเค อุปปนฺโนติ.}

\proseitem{105}{ตตฺถ กตมํ สํกิเลส\-ภาคิยํ นิพฺเพธ\-ภาคิยํ จ สุตฺตํ, น ตํ ทฬฺหํ พนฺธนมาหุ ธีราติ คาถา. เกน การเณน ตํ พนฺธนํ ทฬฺหํ จตูหิ การเณหิ อิสฺสริเยน สกฺกา โมเจตุํ ธเนน วา อญฺเญน วา ยาจนาย วา ปรายเนน วา. เยสุ จ อยํ ราโค มณิกุณฺฑเลสุ ปุตฺเตสุ ทาเรสุ จ ยา อเปกฺขา, อิทมสฺส เจตสิก\-พนฺธนํ. ตํ น สกฺกา อิสฺสริเยน วา ธเนน วา อญฺเญน วา ยาจนาย วา ปรายเนน วา โมเจตุํ. น จ ตตฺถ โกจิ อตฺถิ ปาฏิโภโค. อิมินา พนฺธนโต โมจยิตฺถาติ เทโว วา มนุสฺโส วา ตทิทํ พนฺธนํ ราคานุสเยน จ ฉสุ พาหิเรสุ จ อายตเนสุ พนฺธติ. รูเปสุ รูปตณฺหา พนฺธติ. ยาว ธมฺเมสุ ธมฺมตณฺหา. โย อิธ โลเก พนฺโธ ปรโลกสฺมิํ พนฺโธ นียติ, โส พนฺโธ ชายติ, พนฺโธ มียติ. พนฺโธ อสฺมา โลกา ปรํ โลกํ คจฺฉติ. น สกฺกา โมเจตุํ อญฺญตฺร อริยมคฺเคน อิมญฺจ พนฺธนํ. มรณภาวญฺจ อุปปตฺติ\-ภาวญฺจ ภยโต วิทิตฺวา ฉนฺทราคํ ปชหติ. โส อิมํ ฉนฺทราคํ ปชหิตฺวา อติกฺกมติ. อยญฺจ โลโก อิโต ปรํ ทุติโย.}

\prose{ตตฺถ ยํ พนฺธนา\-สงฺขารานํ ปหานํ อิทํ วุจฺจติ อุภเยสุ ฐาเนสุ วีริยํ, คนฺธปริวาโต\csromansharedfootnote{8854cf08128aeb3c}{คนฺถปริวโส (อิ), คนฺถปริวุโต (ก)} สุมุนิ โนปลิมฺปติ. ตเถว ปริคฺคเหสุ ปุตฺเตสุ ทาเรสุ จ อวูโฬฺห สลฺโลติ ตสฺเสว ตณฺหาย ปหานํ ทสฺเสติ. อยํ ตณฺหามูลสฺส ปหานา วเร\csromansharedfootnote{e5659b902bbed664}{อหนาวเร (อิ), อหนาวโร (ก)} อปฺปมตฺโตติ กาโม ปมาทวตฺตติ ปหานาย เนกฺขมฺมา\-ภิรโต อปฺปมาทวิหารี ภวติ. ตสฺส อาสยํ ปหานาย เนว อิมํ โลกํ อาสีสติ น ปรโลกํ. น อิธโลกํ นิสฺสิตํ, ปิยรูปํ สาตรูปํ อากงฺขติ. นาปิ ปรโลกํ นิสฺสิตํ ปิยรูปํ \csromanfolio{309}สาภรูปํ อากงฺขติ, เตน วุจฺจเต “นาสีสเต โลกมิมํ ปรํ โลกญฺจา”ติ. ยํ ตสฺส ปหานํ ตํ เฉทนํ อฏฺฐก\-วคฺคิ\-เยสุ มุนิ นิทฺทิฏฺโฐ. โส อิธ วิโรโธ อฏฺฐก\-วคฺคิ\-เยสุ นาสีสนํ อิธ อนาถา. ตถายํ ตณฺหาย ตสฺส ปริคฺคหสฺส วตฺถุกามสฺส เอกคาถาย เอเต สพฺเพ กามา ทสฺสิตา. เตน ภควา เทเสติ “เอตมฺปิ เฉตฺวาน ปริพฺพชนฺติ อนเปกฺขิโน สพฺพกาเม ปหายา”ติ. อิมิสฺสา คาถาย ทฺวิธา นิทฺเทโส สํสนฺทน\-นิทฺเทโส จ สมยนิทฺเทโส จ, ยถา อยํ คาถา สํกิเลสภาคิยญฺจ นิพฺเพธภาคิยญฺจ, เอวํ ตาย คาถาย สํกิเลสภาคิยญฺจ นิพฺเพธภาคิยญฺจ วิสชฺชนา. เอวํ คาถา สพฺพคาถาสุ พฺยากรเณสุ วา นิทฺทิฏฺฐํ สุตฺตํ.}

\proseitem{106}{ตตฺถ กตมา เทสนา, อิมํ สุตฺตํ เกนาธิปฺปาเยน เทสิตํ. เย ราคจริตา สตฺตา, เต กาเม ปชหิสฺสนฺตีติ อยํ ตตฺถ ภควโต อธิปฺปาโย.}

\prose{ตตฺถ กตโม วิจโย, ยสฺส ทสวตฺถุกา กิเลสา อุตฺติณฺณา วนฺตา วิทิตา. กตเม ทสวิธาติ, กิเลสกามา จ โอรมฺภาคิยอุทฺธมฺภาคิยา จ สญฺโญชนา ทสวตฺถุกานิ อายตนานิ, อยํ วิจโย.}

\prose{ตตฺถ กตมา ยุตฺติ. เย สารตฺตา เต คาฬฺห\-พนฺธเนน พนฺธนฺติ อตฺถิ เอสา ยุตฺติ.}

\prose{ตตฺถ กตโม ปทฏฺฐาโน, สารตฺโต มณิกุณฺฑเลสุ มมํการสฺส ปทฏฺฐานํ. อเปกฺขาติ อตีตวตฺถุสฺส สราคสฺส ปทฏฺฐานํ. เอตมฺปิ เฉตฺวาติ ภาวนาย ปทฏฺฐานํ.}

\prose{ตตฺถ กตโม ลกฺขโณ, สารตฺตจิตฺโต มณิกุณฺฑเลสุ โย อหํกาเร วิสตฺโต มมํกาเร วิสตฺโต, โย ปุตฺตทาเร สารตฺโต. เขตฺต\-วตฺถุสฺมิํ สารตฺโต. อยํ ลกฺขโณ หาโร.}

\prose{ตตฺถ กตโม จตุพฺยูโห หาโร, อิธ สุตฺเต ภควโต โก อธิปฺปาโย. เย นิพฺพาเนน ฉนฺทิกา ภวิสฺสนฺติ, เต ปุตฺตทาเร ตณฺหํ ปชหิสฺสนฺติ. อยํ ตตฺถ ภควโต อธิปฺปาโย. อิมานิ จตฺตาริ สจฺจานิ.}

\prose{ตตฺถ กตโม อาวฏฺโฏ, ยา ปุตฺตทาเร ตณฺหา, อยํ สมุทโย. เย อุปาทินฺน\-กฺขนฺธา, เต เย จ พาหิเรสุ รูเปสุ รูปปริคฺคโห, อิทํ ทุกฺขํ, \csromanfolio{310}ยํ ตตฺถ เฉทนียํ, อยํ นิโรโธ. เยน ภิชฺชติ, อยํ มคฺโค.  วิภตฺตีติ นตฺถิ วิภตฺติยา ภูมิ, ปริวตฺตโนติ ปฏิปกฺโข นิทฺทิฏฺโฐ.}

\prose{ตตฺถ กตโม เววจโน, นิทฺทิฏฺโฐ เววจโน.  ตตฺถ กตโม โอตรโณ. อตฺถิ ตณฺหา เอโก สตฺโต โอติณฺโณ ตปฺปจฺจยา วิญฺญาณํ ยาว ชารามรณํ. ยา ตตฺถ เวทนา, อยํ อวิชฺชา วิชฺชุปฺปาทา อวิชฺชานิโรโธ ยาว ชรามรณ\-นิโรโธ.}

\prose{ตตฺถ กตโม โสธโน, สุทฺโธ คาถาย อารมฺโภ.  ตตฺถ กตโม อธิฏฺฐาโน. น ตํ ทฬฺหํ พนฺธนมาหุ ธีราติ เอกตฺตตาย ปญฺญตฺตา, น เวมตฺตตาย. จตฺตาโร ราคา กามราโค รูปราโค ภวราโค ทิฏฺฐิราโค จาติ เอกตฺตตาย ปญฺญตฺตา.}

\prose{ตตฺถ กตโม ปริกฺขาโร, เยสํ ราโค มณิกุณฺฑเลสุ ตสฺส สุภสญฺญา เหตุ, อนุพฺยญฺชนโส จ นิมิตฺตคฺคาหิตา ปจฺจโย. ยาย เต ฉินฺนานิ ตสฺส อสุภสญฺญา เหตุ, นิมิตฺตคฺคหณอนุพฺยญฺชนคฺคหณวิโนทนํ ปจฺจโย.}

\prose{ตตฺถ กตโม สมาโรปโน, สารตฺโต มณิกุณฺฑเลสุ สมฺมูฬฺหวิโธ ทุฏฺฐาติปิ เอตมฺปิ\csromansharedfootnote{82aa07767baf9d48}{เอวมฺปิ (อิ, ก)} เฉตฺวาน ปริพฺพชนฺตีติ ตํ ปริญฺญาตตฺถํ ปริวชฺชิตตฺถํ ปชหิตา, อยํ สมาโรปโน.}

\proseitem{107}{ยํ เจตสิกํ ยํ ปกปฺปิตํ วิตฺถาเรน ปจฺจโย, ยํ วา เจตสิกํ กายิกํ เจตสิกํ กมฺมํ. กิํ การณา, เจตสิกา หิ เจตนา มโนกมฺมาติ วุจฺจเต, สา เจตนากมฺมํ, ยํ เจตสิกํ อิมํ กายิกญฺจ วาจสิกญฺจ อิมานิ ตีณิ กมฺมานิ นิทฺทิฏฺฐานิ. กายกมฺมํ วจีกมฺมญฺจ ตานิ กุสลานิ ปิยํ กาเยน จ วาจาย จ อารภติ ปรามสติ, อยํ วุจฺจติ สีลพฺพต\-ปรามาโส. สงฺกปฺปนา เต ติวิธา สงฺขารา ปุญฺญมยา อปุญฺญมยา อาเนญฺชมยา, ตปฺปจฺจยา วิญฺญาณํ เต อารมฺมณเมตํ โหติ วิญฺญาณสฺส ฐิติยา. ยา สุภสญฺญา สุขสญฺญา อตฺตสญฺญา จ. อิทํ เจตสิกํ. ยํ รูปูปคํ วิญฺญาณํ ติฏฺฐติ รูปารมฺมณํ รูปปติฏฺฐิตํ นนฺทูปเสจนํ วุทฺธิํ วุรูฬฺหิํ เวปุลฺลตํ คจฺฉติ, อยํ สงฺกปฺปนา, อิติ ยํ \csromanfolio{311}วิญฺญาณ\-ฏฺฐิตีสุ ฐิตํ ปฐมาภิ\-นิพฺพตฺติ\-อารมฺมณวเสน อุปาทานํ, อิทํ วุจฺจติ เจตสิกนฺติ.}

\prose{ตตฺถ ฐิตสฺส อรูปสฺส ยา นิกนฺติ อชฺโฌสานํ, อิทมฺปิ สกมฺปิตํ มนาปิเกสุ รูเปสุ ปิยรูป\-สาตรูเปสุ อาโภโค, อิทํ เจตสิกํ. ยํ เจเตติ สตฺเตสุ\csromansharedfootnote{679f374dfe3fdfbd}{สตฺตสุ (อิ)} มนาปิเกสุ อภิชฺฌา\-กายคนฺโถ ปฏิฆา\-นุสเยสุ พฺยาปาท\-กายคนฺโถ สพฺเพ จตฺตาโร คนฺถา, อยํ ปญฺจสุ กามคุเณสุ ปฐมาภิ\-นิปาโต จิตฺตสฺส ยา เจตนา ยสฺส ตตฺถ อสฺสาทา\-นุปสฺสิสฺส อเนกา ปาปกา อกุสลา ธมฺมา จิตฺตํ อรูปวติโย โหนฺติ. ปุคฺคโล ราคา\-นุพนฺธิ\-ภูโต เตหิ กิเลสกาเมหิ ยถา กามกรณีโย, อยํ วุจฺจเต กาเมสุ ปกปฺปนา. เอว สพฺเพ จตฺตาโร โอฆา. ยํ เตหิ กาเมหิ สํยุตฺโต วิหรติ ภาวิโต อชฺโฌสนฺโน, อยํ เจตนา. ยสฺส ตถายํ อวีตราคสฺส อธิคต\-เปมสฺส ตสฺส วิปริณาม\-ญฺญถาภาวา อุปฺปชฺชนฺติ โสก\-ปริเทว\-ทุกฺข\-โทมนสฺสุ\-ปายาสา ทุกฺขา\-นุ\-ปริวตฺติตํ วิญฺญาณํ โหติ สริตสฺส วยธมฺม\-สมุปฺปาโท จิตฺตํ ปริยาทิยติ, อิทํ วุจฺจติ ปกปฺปิตนฺติ.}

\prose{เอกเมกสฺส เจเตติ จ ปกปฺเปติ จ วิญฺญาณสฺส ฐิติ ยา โหติ, สา จ ฐิติ ทฺวิธา อารมฺมณฏฺฐิติ จ อาหารฏฺฐิติ จ, ตตฺถ ยา อารมฺมณฏฺฐิติ, อยํ นามรูปสฺส ปจฺจโย. ยา อาหารฏฺฐิติ ยา ปุนพฺภวา\-ภินิพฺพตฺติกา ฐิติ ยา จ โปโนภวิกา ฐิติ, อยํ วุจฺจติ อารมฺมณํ, ตํ โหติ วิญฺญาณสฺส ฐิติยา ตสฺส วิญฺญาณปจฺจยา นามรูปํ ยาว ชรามรณญฺจ เจเตติ, อถ จ ปุน ปตฺถยเต ยโต น โปโนภวิกา อนาคต\-วตฺถุมฺหิ, อยํ ปฏิปกฺโข นิทฺทิฏฺโฐ. น เจเตติ น ปตฺถยติ อถ จ ทูเสตีติ ทุวิโธ นิทฺเทโส. อสฺส ปุพฺเพ โหติ ตํ เจตสิกํ ตํ ปกปฺปิตํ อสมูหตํ ตปฺปจฺจยา, อยํ วิญฺญาณสฺส ฐิติ โหติ.}

\proseitem{108}{อถ วา ตสฺส อนุสยา อาวิภวนฺติ ตปฺปจฺจยา ตสฺส ปุนพฺภโว นิพฺพตฺตติ. อถ วา นํ สํกิยเต อปฺเปตุ อาคาเร วา, สุขุมา วา สนฺติ วา น สํกิยเต กาเม ตํ เอวํ นิจฺเจสุปิ อาคาเรสุ ชาโต โหติ. ตํ นยติ ยํ โน กปฺเปตุํ เอวํ สงฺขารา เจติตา ปกปฺปิตา จ \csromanfolio{312}อารมฺมณภูตา โหนฺติ, ยา จ เจตนา ยา จ ปกปฺปนา ยญฺจ วตฺถุ นิพฺพตฺตํ, อุโภปิ เอเต อารมฺมณํ วิญฺญาณสฺส ตถา เจตนาย จ สงฺกปฺปนาย จ ปตฺถนาย จ ภูตา สตฺตา เจเตติ จ สงฺกปฺเปติ จ. ยํ คเวสนา น จ เจเตติ น จ สงฺกปฺเปติ. กตเม จ สตฺตา ภูตา เย จ ตนุชาต- อณฺฑชาปิ อณฺฑกา อนุภินฺนา สํเสทชา น จ สมฺภินฺนา อิเม ภูตา, กตเม สมฺภเวสิโน คพฺภคตา อณฺฑคตา สํสรนฺโต อิเม น เจเตติ น ปตฺเถติ น จ สงฺกปฺเปติ. อนุสเย น จ ปุนพฺภโว นิพฺพตฺตีติ. เย ภูตา สตฺตา เย สมฺภเวสิโน, เต ถาวรา. เย วา สโต เจเตนฺติ ปตฺเถนฺติ จ เย ถาวรา. เต น จ เจเตนฺติ, น จ ปตฺเถนฺติ, น จ สงฺกปฺเปนฺติ, อนุสเยน จ สํสรนฺติ.}

\prose{อปโร ปริยาโย, เย อริยปุคฺคลา เสกฺขา, ตตฺถ เต น จ เจเตนฺติ, น จ สงฺกปฺเปนฺติ, อนุสเยน ปุน อุปฺปชฺชนฺติ.}

\prose{อปโร ปริยาโย, สุขุมา ปาณา ภูมิคตา อุทกคตา จกฺขุโน อาปาถํ นาคจฺฉนฺติ, เต น จ เจเตนฺติ, น จ สงฺกปฺเปนฺติ, อนุสเยน จ สํสรนฺติ.}

\prose{อปโร ปริยาโย, พาหิกา สพฺเพ ภิกฺขู อภิมานิกา, เต น จ เจเตนฺติ, น จ ปตฺถยนฺติ, อนุสเยน จ สํสรนฺติ, น จ เจเตนฺติ, น จ สงฺกปฺเปนฺติ, น จ อนุเสนฺติ. อารมฺมณมฺเปตํ น โหติ วิญฺญาณสฺส ฐิติยา.}

\prose{น จ เจเตตีติ ปริยุฏฺฐาน\-สมุคฺฆาตํ ทสฺเสติ. น จ อนุเสตีติ อนุสย\-สมุคฺฆาตํ ทสฺเสติ. น จ เจเตตีติ โอฬาริกานํ กิเลสานํ ปหานํ ทสฺเสติ. น จ อนุเสตีติ สุขุมานํ กิเลสานํ ปหานํ ทสฺเสติ. น จ เจเตตีติ เยน ภูมิ จ น จ ปตฺถยนฺตีติ สกทาคามี อนาคามี, น จ อนุเสตีติ อรหํ. น จ เจเตตีติ สีลกฺขนฺธสฺส ปฏิปกฺเขน ปหานํ ทสฺเสติ. น จ ปตฺถยตีติ สมาธิ\-กฺขนฺธสฺส ปฏิปกฺเขน ปหานํ ทสฺเสติ. น จ อนุสยตีติ ปญฺญา\-กฺขนฺธสฺส ปฏิปกฺเขน ปหานํ ทสฺเสติ. น จ เจเตตีติ อปุญฺญมยานํ สงฺขารานํ ปหานํ ทสฺเสติ. น จ ปตฺถยตีติ ปุญฺญมยานํ สงฺขารานํ ปหานํ ทสฺเสติ. น จ อนุเสตีติ อาเนญฺชมยานํ สงฺขารานํ ปหานํ ทสฺเสติ. น จ เจเตตีติ อนญฺญาต\-ญฺญสฺสามี\-ตินฺทฺริยํ, น จ ปตฺถยตีติ อญฺญินฺทฺริยํ, น จ อนุสยตีติ อญฺญาตาวิโน อินฺทฺริยํ. น จ เจเตตีติ มุทุกา อินฺทฺริยภาวนา, น จ}

\prose{\csromanfolio{313}ปตฺถยตีติ มชฺไฌ\-นฺทฺริยภาวนา, น จ อนุเสตีติ อธิมตฺตา อินฺทฺริย ภาวนา. อยํ สุตฺตตฺโถ.}

\proseitem{109}{ตตฺถ กตมา เทสนา. อิธ สุตฺเต จตฺตาริ สจฺจานิ เทสิตานิ. ยญฺจ เจตยิตํ ยญฺจ ปกปฺปิตํ อตฺถิ เอตํ อารมฺมณํ จิตฺตํ ปติฏฺฐติ วิจินติ\csromansharedfootnote{2a6083f30e9d643e}{วิจินยติ (อิ, ก)} ยุชฺชติ. น จ เจเตตีติ น จ ปตฺถยตีติ อตฺถิ เอวํ อารมฺมณํ อนุสเย วิญฺญาณมิติ วิจินิยติ ยุชฺชติ น จ เจเตติ น จ ปตฺถยติ. อนุสยปฺปหานา วิญฺญาณฏฺฐิติํ น คเวสนฺติ, วิจิยนฺตํ ยุชฺชติ. อยํ ยุตฺติวิจโย.}

\prose{ตตฺถ กตโม ปทฏฺฐาโน. เจตนา ปริยุฏฺฐานํ เจตนา\-ปริยุฏฺฐานสฺส ปทฏฺฐานํ. สงฺกปฺปนํ อุปาทานสฺส ปทฏฺฐานํ. อนุสโย ปริยุฏฺฐานสฺส ปทฏฺฐานํ. เตสํ ฉนฺทราค\-วินาสาย ภาวนา ภวราคสฺส ปหานํ.}

\prose{ตตฺถ กตโม ลกฺขโณ. ยํ เจตสิกนฺติ เวทยิตํ ปกปฺปิตํ อุคฺคหิตํ วิญฺญาตํ ตพฺพิญฺญาณํ อารมฺมณมฺปิ ปจฺจโยปิ.}

\prose{ตตฺถ กตโม จตุพฺยูโห. อิธ สุตฺเต ภควโต โก อธิปฺปาโย. เย ปุนพฺภวํ น อิจฺฉนฺติ, เต น เจตยิสฺสนฺติ น จ ปตฺถยิสฺสนฺตีติ, อยํ อธิปฺปาโย.}

\prose{อาวฏฺโฏติ ยา จ เจตนา ปตฺถนา จ อนุสโย จ วิญฺญาณฏฺฐิติ\-ปหานา จ, อิมานิ เทฺว สจฺจานิ.  วิภตฺตีติ นตฺถิ วิภตฺติยา ภูมิ.  ปริวตฺตนา ปน ปฏิปกฺขํ สุตฺตํ.}

\prose{ตตฺถ กตโม เววจโน. เจตนา รูปสญฺเจตนา ยาว\-ธมฺม\-สญฺเจตนา. โย อนุสโย, เต สตฺต อนุสยา.}

\prose{ปญฺญตฺตีติ เจตนา\-ปริยุฏฺฐานํ ปญฺญตฺติยา ปญฺญตฺตา. สงฺกปฺปนํ อุปาทาน\-ปญฺญตฺติยา ปญฺญตฺตํ. อนุสโย เหตุปญฺญตฺติยา ปญฺญตฺโต. วิญฺญาณฏฺฐิติ อุปปตฺติ\-เหตุ\-ปญฺญตฺติยา ปญฺญตฺตา. เจตนา สงฺกปฺปนา อนุสโย สมุจฺเฉโท ฉนฺทราควินย\-ปญฺญตฺติยา ปญฺญตฺโต. ปฐเม เกจิ ทฺวีหิ ปริวตฺตเกหิ ปฏิจฺจ\-สมุปฺปาโท อิทปฺปจฺจย\-ตาย มชฺฌปญฺญตฺติ.}

\prose{\csromanfolio{314}โอตรโณติ ทฺวีหิ ปริวตฺตเกหิ ทุกฺขญฺจ สมุทโย จ มชฺฌิมเกหิ มคฺโค จ นิโรโธ จ.  โสธโนติ สุตฺเต สุตฺตสฺส อารมฺโภ.}

\prose{อธิฏฺฐาโนติ ยญฺเจตยิตํ สพฺพํ อธิฏฺฐาเนน เอกตฺตาย ปญฺญตฺตํ. สงฺกปฺปิตนฺติ อุปาทาเนกตฺตาย ปญฺญตฺตํ. วิญฺญาณํ เอกตฺตาย ปญฺญตฺตํ.}

\prose{ปริกฺขาโรติ สุภญฺจ อารมฺมณํ อโยนิโส มนสิกาโร เจตนา เหตุ\-ปจฺจยตาย ปจฺจโย. วิญฺญาณสฺส ปติฏฺฐาโน ธมฺโม อารมฺมณ\-ปจฺจยตาย ปจฺจโย. ตสฺส มนสิกาโร เหตุ\-ปจฺจยตาย ปจฺจโย.}

\prose{ตตฺถ กตโม สมาโรปโน. อิทํ สุตฺตํ สญฺญิตํ ตตฺถ เจเตติ วิสชฺชนา อิติ นิทฺทิสิตพฺพา. ตสฺส ทิฏฺฐิยา วิญฺญาณปจฺจยา นามรูปํ ยาว ชรามรณํ, อยํ สมาโรปโน.  อารมฺมณเมตํ น โหติ วิญฺญาณสฺส ฐิติยา, วิญฺญาณนิโรธา นามรูป\-นิโรโธ, นามรูป\-นิโรธา ยาว ชรามรณ\-นิโรโธ.}

\proseitem{110}{ตตฺถ กตมํ สํกิเลสภาคิยญฺจ นิพฺเพธภาคิยญฺจ อเสกฺขภาคิยญฺจ สุตฺตํ. อยํ โลโก\csromansharedfootnote{ae09b214a326bfa2}{ขุ 1. 115 อุทาเน.} สนฺตาปชาโต ยาว เย หิ เกจิ สมณา วา พฺราหฺมณา วา ภเคน ภวสฺส วิปฺปโมกฺขมาหํสุ. สํกิเลส\-ภาคิยํ อุปธิํ หิ ปฏิจฺจ ทุกฺขมิทํ สมฺโภติ, ยา ตา ปน ตณฺหา ปหียนฺติ, ภวํ นาภินนฺทตีติ นิพฺเพธสฺส นิพฺพุตสฺส\csromansharedfootnote{fdeb1d74035382b2}{นิจฺจุตสฺส (อิ, ก)} ภิกฺขุโน อนุปาทาย ปุนพฺภโว น โหติ. อุปจฺจคา สพฺพภวานิ ตาทีติ อเสกฺข\-ภาคิยํ.}

\prose{ตตฺถ สนฺตาปชาโตติ ราคโช สนฺตาโป โทสโช โมหโชติ. เตสํ สตฺตานํ ฐานํ ทสฺเสติ. โลโก สนฺตาปชาโตติ ผสฺโส ติวิโธ สุขเวทนีโย ทุกฺขเวทนีโย อทุกฺขม\-สุข\-เวทนีโย. ตตฺถ สุขเวทนีโย ผสฺโส ราคสนฺตาโป, ทุกฺขเวทนีโย โทสสนฺตาโป, อทุกฺขม\-สุข\-เวทนีโย โมหสนฺตาโป. ยถา จ ภควา อาห ปฐมกสฺส วลาหกสฺส โคมคฺเค\csromansharedfootnote{ed44865329c0ff2b}{โกมคฺเค (อิ, ก) อํ 1. 135 ปิฏฺเฐ.} เยหิ คหปติปุตฺต ราคเชหิ โทสเชหิ โมหเชหิ สนฺตาเปหิ ทุกฺขํ สุปติ, เต มม สนฺตาปา น สนฺติ.}

\prose{\csromanfolio{315}โรคํ วทติ อตฺตโตติ เตหิ สนฺตาเปหิ สนฺตาปิโต ติวิธํ วิปลฺลาสํ ปฏิลภติ สญฺญา\-วิปลฺลาสํ จิตฺต\-วิปลฺลาสํ ทิฏฺฐิ\-วิปลฺลาสํ. ตตฺถ อสุเภ สุภนฺติ สญฺญาวิปลฺลาโส. ทุกฺเข สุขนฺติ จิตฺตวิปลฺลาโส. อนิจฺเจ นิจฺจนฺติ อนตฺตนิ อตฺตาติ ทิฏฺฐิ\-วิปลฺลาโส.}

\prose{ยถา จิตฺตสฺส วิปลฺลาโส สญฺญาทิฏฺฐิเต ติวิธา วิตกฺกา จิตฺตวิตกฺโก วิปลฺลาโส สญฺญาวิตกฺโก วิปลฺลาโส ทิฏฺฐิวิตกฺโก วิปลฺลาโสปิ. ตตฺถ อวิชฺชา วิปลฺลาโส โคจรา คติปเตยฺยภูมิ, ยถา หิ ตํ สญฺชานาติ ยถา วิชานาติ ยถา สญฺชานาติ จ วิชานาติ จ. ยถา ขนฺติ เจเตติ อิเม จตฺตาโร วิปลฺลาสา สตฺตา เยหิ จตุพฺพิธํ อตฺตภาว\-วตฺถุํ โรคภูตํ คณฺฑภูตํ “อตฺตา”ติ วทนฺติ. โรคํ วทติ อตฺตโตติ อยํ อาวฏฺโฏ. เยน เยน หิ มญฺญติ ตโต ตํ โหติ อญฺญถาติ สุภนฺติ มญฺญติ น ตถา โหติ. เอวํ สุขนฺติ นิจฺจํ อตฺตาติ โส อญฺญถา ภวเมว สนฺตํ อนาคตํ ภวํ ปตฺถยติ, เตน วุจฺจติ “ภวราโค”ติ. ภวเม\-วา\-ภินนฺทติ, ยํ อภินนฺทติ, ตํ ทุกฺขนฺติ ปญฺจกฺขนฺเธ นิทฺทิสิยติ. ยญฺจ ตปฺปจฺจยา โสกปริเทว\-ทุกฺขํ ตสฺส หิ ภาเวสฺสติ. เอตฺตาวตา สํกิเลโส โหติ. ปหานตฺถํ โข ปน พฺรหฺมจริยํ วุสฺสติ. ติณฺณํ สนฺตาปานํ ฉนฺทราค\-วินโย โหติ.}

\prose{อุปธิํ หิ ปฏิจฺจ ทุกฺขมิทํ ภวตีติ เย ภวเม\-วา\-ภินนฺทนฺติ ยสฺส ภาเวสฺสติ, ตํ ทุกฺขํ ตสฺส ทุกฺขสฺส ปหานมาห. สพฺพโส อุปาทานญฺจ ยํ นตฺถิ ทุกฺขสฺส สมฺภโวติ จตฺตาโร วิปลฺลาสา ยถา นิทฺทิฏฺฐอุปาทานมาห. ตสฺส ปฐโม วิปลฺลาโส กามุปาทานํ, ทุติยํ ทุฏฺฐุปาทานํ, ตติยํ สีลพฺพตุ\-ปาทานํ, จตุตฺถํ อตฺตวาทุ\-ปาทานํ, เตสํ โย ขโย นตฺถิ ทุกฺขสฺส สมฺภโว อุปธิ นิทานํ ทุกฺขนิโรธมาห. เอวเมตํ ยถาภูตํ สมฺมปฺปญฺญาย ปสฺสโต วิภวตณฺหา น โหติ. วิภวํ นาภินนฺทตีติ ทสฺสนภูมิํ มนฺเตติ สพฺพโส ตณฺหกฺขยํ นิพฺพานนฺติ เทฺว วิมุตฺติโย กเถติ ราควิราคญฺจ อวิชฺชาวิราคญฺจ. ตสฺส ภิกฺขุโนติ อนุปาทิเสส\-นิพฺพานธาตุํ มนฺเตติ. อยํ สุตฺตสฺส อตฺถนิทฺเทโส.}

\proseitem{111}{ตตฺถ กตโม วิจโย. ยสฺส ยตฺถ ปริฬาเหติ ตสฺส ปริฑยฺห\-นฺตสฺส โส ยถาภูตํ นตฺถิ นิพฺพินฺทติ จ, อยํ วิจโย จ ยุตฺติ จ.}

\prose{\csromanfolio{316}ปทฏฺฐาโน ราคโช ปริฬาโห สุขินฺทฺริยสฺส โทมนสฺสิ\-นฺทฺริยสฺส จ ปทฏฺฐานํ. โทสโช ปริฬาโห สุขินฺทฺริยสฺส โทมนสฺสิ\-นฺทฺริยสฺส จ ปทฏฺฐานํ. โมหโช ปริฬาโห อุเปกฺขิ\-นฺทฺริยสฺส โทมนสฺสิ\-นฺทฺริยสฺส จ ปทฏฺฐานํ.}

\prose{ตตฺถ กตโม ลกฺขโณ หาโร. ผสฺสปเรโต เวทนาปเรโต สญฺญาปเรโตปิ สงฺขารปเรโตปิ เยน เยน มญฺญติ ยทิ สุภนิมิตฺเตน ยทิ สุขินิมิตฺเตน ยทิ นิจฺจนิมิตฺเตน ยทิ อตฺตนิมิตฺเตน อสุเภ สุภนฺติ มญฺญติ, เอวํ สพฺพํ ราคเช ปริฬาเห วุตฺเต จตฺตาโร ปริฬาหา วุตฺตา ภวนฺติ. ราคโช โทสโช โมหโช ทิฏฺฐิโช จ ราคํ วทามีติ อตฺตโต วทติ. สพฺพานิ ปนฺนรส ปทานิ อนิจฺจํ ทุกฺขนฺติ.}

\prose{ตตฺถ กตโม จตุพฺยูโห. อิธ สุตฺเต ภควโต โก อธิปฺปาโย. เย ปริฬาเหน น อจฺฉนฺติ เต ภวํ นาภินนฺทนฺติ. เย ภวํ นาภินนฺทนฺติ, เต ปรินิพฺพายิสฺสนฺติ. อยํ อธิปฺปาโย.}

\prose{ตตฺถ กตโม อาวฏฺโฏ. สํกิเลส\-ภาคิเยน ทุกฺขญฺจ สมุทยญฺจ นิทฺทิสติ. นิพฺเพธ\-ภาคิเยน มคฺคญฺจ นิโรธญฺจ.}

\prose{ตตฺถ กตมา วิภตฺติ. สนฺตาปชาโต โรคชาโต โรคํ วทติ อตฺตโต ตํ น เอกํเสน โหติ อมนสิการา สนฺตาปชาโต โข น จ โรคํ อตฺตโต วทติ.}

\prose{ตตฺถ กตโม ปริวตฺตโน. ปกฺข\-ปฏิปกฺข\-นิทสฺสนตฺถํ ภูมิ ปริวตฺตนาย.}

\prose{ตตฺถ กตโม เววจโน หาโร. โรคญฺจ อตฺตโต วทติ สลฺลํ อตฺตโต วทติ. ปนฺนรส ปทานิ สพฺพานิ วตฺตพฺพานิ.}

\prose{ตตฺถ กตมา ปญฺญตฺติ. สนฺตาปชาโตติ โทมนสฺส\-ปทฏฺฐานํ. สพฺเพ วจน\-ปญฺญตฺติยา ปญฺญเปติ. โรคํ วทติ อตฺตโต วิปลฺลาโส สํกิเลส\-ปญฺญตฺติยา ปญฺญเปติ. ยํ นาภินนฺทติ, ตํ ทุกฺขนฺติ วิปลฺลาส\-นิกฺเขป\-ปญฺญตฺติยา ปญฺญตฺตา. เต อกตสตฺตา โลกา มชฺเฌน เวมตฺตตาย ปญฺญตฺตา.}

\prose{ตตฺถ กตโม โอตรโณ, สนฺตาปชาโตติ ตีณิ อกุสลมูลานิ, เต สงฺขารา สงฺขารกฺขนฺธ\-ปริยาปนฺนา, ธาตูสุ ธมฺมธาตุ, อายตเนสุ ธมฺมายตนํ. อินฺทฺริเยสุ อิตฺถินฺทฺริยํ ปุริสินฺทฺริยญฺจ ปทฏฺฐานํ.}

\prose{\csromanfolio{317}ตตฺถ กตโม โสธโน. สุทฺโธ สุตฺตสฺส อารมฺโภ.}

\prose{ตตฺถ กตโม อธิฏฺฐาโน หาโร. ปริฬาโหติ เย สตฺตา โลกา เอกตฺต\-ปญฺญตฺติยา ปญฺญตฺตา, เต อกตสตฺตา โลกา มชฺเฌน เวมตฺตตาย ปญฺญตฺตา.}

\prose{ตตฺถ กตโม ปริกฺขาโร. สนฺตาปชาโตติ อโยนิโส มนสิกาโร เหตุ, วิปลฺลาสญฺจ ปจฺจโย. ตตฺถ ทฺวีหิ ธมฺเมหิ อตฺตา อภินิวิฏฺฐา จิตฺตญฺจ เจตสิกญฺจ ธมฺเม อุภยานิ ตสฺส วิปรีเตน ปรามสโต. อปโร ปริยาโย, เจตสิเกหิ ธมฺเมหิ อตฺตสญฺญา อนตฺตสญฺญา สมุคฺฆาเตติ. อปโร ปริยาโย. อนิจฺจสญฺญา เจตสิเกสุ ธมฺเมสุ, น ตุ อตฺตสญฺญา. อิทํ วุจฺจติ จิตฺตนฺติ วา มโนติ วา วิญฺญาณนฺติ วา อิทํ ทีฆรตฺตํ อพฺภุคฺคตํ เอตํ มม, เอโสหมสฺมิ, เอโส เม อตฺตาติ. ตตฺถ เจตสิกา ธมฺมา\-นุปสฺสนา เอสาปิ ธมฺมสญฺญา. ตสฺส โก เหตุ, โก ปจฺจโย. อหํกาโร เหตุ, มมํกาโร ปจฺจโย.}

\prose{ตตฺถ กตโม สมาโรปโน. อยํ โลโก สนฺตาปชาโตติ อกุสลํ มนฺเตติ วิญฺญาณํ นามรูปสฺส ปจฺจโย ยาว ชรามรณนฺติ, อยํ สมาโรปโน.}

\proseitem{112}{เอวเมตํ ยถาภูตํ, สมฺมปฺปญฺญาย ปสฺสติ อกุสลมูลานํ ปหานํ. ตตฺถ อวิชฺชานิโรโธ อวิชฺชานิโรธา ยาว ชรามรณ\-นิโรโธ, อยํ สมาโรปโน.}

\prose{จตฺตาโร ปุคฺคลา\csromansharedfootnote{a4a66842f7e0ff2d}{อํ 1. 311 ปิฏฺเฐ.} อนุโสตคามี ปฏิโสตคามี ฐิตตฺโต, ติณฺโณ ปารงฺคโต ถเล ติฏฺฐติ พฺราหฺมโณติ.}

\prose{ตตฺถ โย อนุโสตคามี อยํ กาเม เสวติ. ปาปญฺจ กมฺมํ กโรติ ยาว กาเม ปฏิเสวติ. อิทํ โลโภ อกุสลมูลํ, โส เยว\csromansharedfootnote{cb63cef12ad01f03}{โสเยว สนฺธิ \csromandash{} ม.พ.ป.} ตณฺหา, โส เตหิ กาเมหิ วุยฺหติ อนุโสตคามีติ วุจฺจติ. โย ปุคฺคโล ตาหิ คมิโต ตปฺปจฺจยา ตสฺส เหตุ อกุสลกมฺมํ กโรติ กาเยน จ วาจาย จ, อยํ วุจฺจติ ปาปกมฺมํ กโรตีติ. ตสฺส ตีณิ โสตานิ สกฺกายทิฏฺฐิ วิจิกิจฺฉา สีลพฺพต\-ปรามาโส. อิเมหิ ตีหิ \csromanfolio{318}โสเตหิ ติวิธ\-ธาตุยํ อุปฺปชฺชติ กามธาตุยํ รูปธาตุยํ อรูปธาตุยํ. เตน ปฏิปกฺเขเน โย กาเม น ปฏิเสวติ. โย สีลวตํ น ปรามสติ. โย สกฺกาย\-ทิฏฺฐีนํ ปหานาย กาเมสุ ยถาภูตํ อาทีนวํ ปสฺสติ. เยน จ เต ธมฺเม ปฏิเสวติ. ยญฺจ ตปฺปจฺจยา ติฏฺฐติ พฺราหฺมโณติ อรหํ กิร. ตตฺถ อรหํ ตสฺส ปารงฺคโต โหติ, ปารงฺคตสฺส ถเล ติฏฺฐติ โสปาทิเสสา นิพฺพานธาตุ. อนุโสตคามีนีติ ทสฺสน\-ปฺปหาตพฺพานํ สํโยชนานํ อปฺปหานมาห. ปฏิโสตคามินีติ ผเล ทิฏฺเฐ\-กฏฺฐานญฺจ กิเลสานํ ปหานมาห, ฐิตตฺเตน ปญฺจนฺนํ โอรมฺภาคิยานํ สํโยชนานํ ปหานมาห. ตตฺถ อนุโสตคามินา มคฺครูปิมาห. ปฏิโสตคามินา ฐิตตฺเตน จ มคฺคมิติมาห. ปารงฺคเตน สาวกา อเสกฺขา จ สมฺมาสมฺพุทฺธา จ วุตฺตา. อนุโสตคามินา สกฺกาย\-สมุทย\-คามินิํ ปฏิปทมาห. ปฏิโสตคามินา ฐิตตฺเตน สกฺกาย\-นิโรธ\-คามินิํ ปฏิปทมาห. ปารงฺคเตน ทส อเสกฺขา อรหนฺตา ธมฺมา วุตฺตา. อยํ สุตฺตตฺโถ.}

\proseitem{113}{ตตฺถ กตมา เทสนา. อิมสฺมิํ หิ สุตฺเต จตฺตาริ อริยสจฺจานิ เทสิตานิ. เตธาตุก\-โลก\-สมติกฺกมนญฺจ.}

\prose{ตตฺถ กตโม วิจโย หาโร. โย กาเม ปฏิเสวติ ปาปํ\csromansharedfootnote{f2c0ef12660322e2}{ปาปกํ (อิ)} กเรยฺยาติ โย จ กาเม น ปฏิเสวติ โส ปาปกมฺมํ น กเรยฺยาติ โย จ อิเมหิ ทฺวีหิ ภูมีหิ อุตฺติณฺโณ ปารงฺคโตติ ยา วีมํสา อยํ วิจโย.}

\prose{ยุตฺตีติ ยุชฺชติ สุตฺเตสุ, นายุชฺชตีติ ยา วีมํสาย, อยํ ยุตฺติ.  ปทฏฺฐาโนติ อนุโสตคามินา สตฺตนฺนํ สํโยชนานํ ปทฏฺฐานํ. อกุสลสฺส กิริยา อกุสลสฺส มูลานํ ปทฏฺฐานํ. ปฏิโสตคามินา ยถาภูต\-ทสฺสนสฺส ปทฏฺฐานํ. ฐิตตฺเตน อสํหาริยาย\csromansharedfootnote{a6ef5640ad42113c}{อสหาริยาย (อิ)} ปทฏฺฐานํ. ปารงฺคโตติ กทาจิ ภูมิยา ปทฏฺฐานํ.}

\prose{ตตฺถ กตโม ลกฺขโณ หาโร. โย อนุโสตํ คจฺฉติ ตณฺหาวเสน. สพฺเพสมฺปิ กิเลสานํ วเสน คจฺฉติ. โย ปฏิโสตํ วายมติ. ตณฺหาย สพฺเพสมฺปิ โส กิเลสานํ วายมติ ปฏิโสตํ. \csromanfolio{319}โย อตฺตนา ฐิโต กาเยนปิ โส ฐิโต วาจาจิตฺเตนปิ โส ฐิโต. อยํ ลกฺขโณ หาโร.}

\prose{ตตฺถ กตโม จตุพฺยูโห. อิธ สุตฺเต ภควโต โก อธิปฺปาโย. เย อนุโสต\-คามินิยา ปฏิปทาย นาภิรมิสฺสนฺติ, เต ปฏิโสตํ วายมิสฺสนฺตีติ ยาว กทาจิ ภูมิยํ, อยํ อธิปฺปาโย.  อาวฏฺโฏติ อิธ สุตฺเต จตฺตาริ สุตฺตานิ เทสิตานิ.}

\prose{ตตฺถ กตโม วิภตฺติ หาโร. โย กาเม ปฏิเสวติ ปาปญฺจ กมฺมํ กโรติ. โส อนุโสตคามีติ น เอกํเสน โสตาปนฺโนปิ กาเม ปฏิเสวติ. ตํ ภาคิยญฺจ ปาปกมฺมํ กโรติ. กิญฺจาปิ เสกฺโขปิ กเรยฺย ปาปํ ยถา สุตฺเต นิทฺทิฏฺโฐ น จ โส อนุโสตคามี, อิทํ วิภชฺช\-พฺยากรณียํ. น จ กาเม ปฏิเสวติ น จ ปาปกมฺมํ กโรติ ปฏิโสตคามี น จ เอกํเสน สพฺเพ พาหิรโก กาเมสุ วีตราโค น จ กาเม ปฏิเสวติ, เตน จ ปาปกมฺมํ กโรติ อนุโสตคามี ปฏิโสตคามี, อยํ วิภตฺติ.}

\prose{ตตฺถ กตโม ปริวตฺตโน หาโร. นิทฺทิฏฺโฐ ปฏิปกฺโข. เววจโนติ กาเมสุ วตฺถุกามาปิ กิเลสกามาปิ รูปสทฺทคนฺธรสผสฺสปุตฺตทาร ทาสกมฺมกรโปริสญฺจ ปริคฺคหา.}

\prose{ปญฺญตฺตีติ สพฺเพ ปุถุชฺชนา เอกตฺตาย ปญฺญตฺตา. อนุโสตคามีติ กิเลส\-สมุทาจาร\-ปญฺญตฺติยา ปญฺญตฺตา. เย ปน เสกฺขา ปุคฺคลา, เต นิพฺพาน\-ปญฺญตฺติยา\csromansharedfootnote{c1bba330610b43c0}{นิฏฺฐานปญฺญตฺติยา (ก)} ปญฺญตฺตา. เย ปน อนาคามี, เต อสํหาริย ปญฺญตฺติยา ปญฺญตฺตา, อยํ ปญฺญตฺติ.}

\prose{โอตรโณติ โย อนุโสตคามี, โส ทุกฺขํ. เย ตสฺส ธมฺมา, เต ทุกฺขสฺส สมุทโย. ยํ รูปํ, อยํ รูปกฺขนฺโธ, เอวํ ปญฺจปิ ขนฺธา ปฏิจฺจ\-สมุปฺปาโท, เต กิเลสา สงฺขารกฺขนฺธ\-ปริยาปนฺนา ธมฺมายตนํ ธมฺมธาตุ อินฺทฺริเยสุ จ ปญฺญตฺตา.}

\prose{โสธโนติ เยนารมฺเภน อิทํ สุตฺตํ เทสิตํ, โส อารมฺโภ สพฺโพ สุทฺโธ.}

\prose{\csromanfolio{320}อธิฏฺฐาโนติ ปฏิโสตคามินา สพฺเพ โสตาปนฺนา เอกตฺเตน วา นิทฺทิฏฺฐา ราคานุสย\-ปฏิโสตคามิโน เสกฺขาว มคฺโค จ เสกฺโข จ ปุคฺคโล ฐิตตฺโตติ.}

\prose{วีตราโค เอกตฺตาย ปญฺญตฺโต. ปารงฺคโตติ สพฺเพ อรหนฺโต สพฺเพ ปจฺเจกพุทฺธา สมฺมาสมฺพุทฺธา จ เอกตฺตาย ปญฺญตฺตา.}

\prose{ปริกฺขาโรติ อนุโสตคามิโน ปาปมิตฺต\-ปจฺจโย กาม\-ปริยุฏฺฐานํ เหตุ. ปฏิโสตคามิโน เทฺว เหตู เทฺว ปจฺจยา จ ยาว สมฺมาทิฏฺฐิยา อุปฺปาทายทิฏฺฐิ\csromansharedfootnote{301c5ee3dd9a4d24}{อุปาทายทิฏฺฐิ (อิ)}, ตสฺส ปฏิลทฺธ\-มคฺโค เหตุ อารมฺโภ ปจฺจโย กายิโก เจตสิกสฺส โกฏฺฐาโส จ.  สมาโรปโนติ วิภตฺติ อิทํ สุตฺตํ นตฺถิ สมาโรปนาย ภูมิ.}

\proseitem{114}{ปญฺจานิสํสา โสตานุคตานํ ธมฺมานํ\csromansharedfootnote{aff3b73151d81464}{อํ 1. 504 ปิฏฺเฐ.} ยาว ทิฏฺฐิยา สุปฺปฏิวิทฺธานํ สุตฺตํ วิตฺถาเรน กาตพฺพํ. ยุญฺชโต ฆเฏนฺตสฺส วายมโต คิลาโน มรณกาเล เทวภูโต ปจฺเจกโพธิํ ปาปุณาติ. โสตานุคตาติ สทฺธมฺม\-สฺสวเนน กตํ โหติ. น จ อธิปญฺญา\-ธมฺม\-วิปสฺสนาย ตสฺส จิตฺตํ ตสิตํ โหติ, น จ อนิพฺพิทฺธตฺตํ, อิทํ จ สุตฺตํ ปญฺจนฺนํ ปุคฺคลานํ เทสิตํ, สทฺธานุสาริโน มุทินฺทฺริยสฺส ติกฺขิ\-นฺทฺริยสฺส จ ธมฺมานุสาริโน ติกฺขิ\-นฺทฺริยสฺส มุทินฺทฺริยสฺส จ. โย ปน โมหจริโต ปุคฺคโล น สกฺโกติ ยุญฺชิตุํ ฆฏิตุํ วายมิตุํ ยถาภูตํ ยถาสมาธิกา วิมุตฺติ ตํ ขณํ ตํ ลยํ ตํ มุหุตฺตํ ผลํ ทสฺเสติ. สาธุ ปริหายติ ปโร ตํ ทุยฺหติ, โน ตุ สุขอวิปากินี ภวติ. ตสฺส ทิฏฺเฐ เยว จ ธมฺเม อุปปชฺช\-อปราปริย\-เวทนียํ. ตตฺถ โย ปุคฺคโล ธมฺมานุสารี. ตสฺส ยทิ โสตานุคตา ธมฺมา โหนฺติ. โส ยุญฺชนฺโต ปาปุณาติ. โย ธมฺมานุสารี มุทินฺทฺริโย, โส คิลาโน ปาปุณาติ. โย สทฺธานุสารี ติกฺขินฺทฺริโย, โส มรณกาล\-สมเย ปาปุณาติ. โย มุทินฺทฺริโย, โส เทวภูโต ปาปุณาติ. ยทา เทวภูโต น ปาปุณาติ, น โส เตเนว ธมฺมราเคน ตาย ธมฺมนนฺทิยา ปจฺเจกโพธิํ ปาปุณาติ. โย โสตานุคเตสุ ยุญฺชติ ฆเฏติ วายมติ, โส ปุพฺพาปนฺเนน วิเสสํ สญฺชานาติ, สญฺชานนฺโต ปาปุณาติ. สเจ ปน คิลานสฺส มนสิกาโร โหติ, ตตฺถ ยุญฺชนฺโต ปาปุณาติ. สเจ ปนสฺส มรณกาเล สํวิคฺโค โหติ, \csromanfolio{321}ตตฺถ ยุญฺชนฺโต ปาปุณาติ. สเจ ปน น กตฺถจิ\csromansharedfootnote{70cfbd086c3fa6c7}{กตฺถ (อิ, ก), ตตฺถ (ก)} สํเวโค โหติ, ตสฺส เทวภูตสฺส สุขิโน ธมฺมภูตา ปาทา เอวํ อวิลปติ. โส เอวํ ชานาติ “อยํ โส ธมฺมวินโย ยตฺถ มยํ ปุพฺเพ มนุสฺสภูตา พฺรหฺมจริยํ จริมฺหา”ติ. อถ เทวภูโต ปาปุณาติ. ทิพฺเพสุ วา ปญฺจสุ กามคุเณสุ อชฺโฌสิโต โหติ ปมาทวิหารี, โส เตน กุสลมูเลน ปจฺเจกโพธิํ ปาปุณาติ.}

\prose{ยา ปรโตโฆเสน วจสา สุปริจิตา, อยํ สุตมยี ปญฺญา. เย ปน ธมฺมา โหนฺติ มนสา อนุเปกฺขิตา, อยํ จินฺตามยี ปญฺญา. ยํ ทิฏฺฐิยา สุปฺปฏิวิทฺธา, อยํ ภาวนามยี ปญฺญา. ยํ โสตานุคตา วจสา ปริจิตา โหนฺติ, โส จ ทิฏฺเฐ เยว ธมฺเม ปรินิพฺพายี, อยํ อรหํ ปุคฺคโล. โย อุปปชฺชติ เทวภูโต ปาปุณาติ, ตตฺถ จ ปรินิพฺพายติ, อยํ อนาคามี. โย เตน กุสลมูเลน ปจฺเจกโพธิํ ปาปุณาติ, อยํ ปุพฺพโยค\-สมฺภาร\-สมฺภูโต ปุคฺคโล.}

\prose{โสตานุคตา ธมฺมาติ ปฐมํ วิมุตฺตายตนํ, วจสา ปริจิตาติ ทุติยํ ตติยญฺจ วิมุตฺตายตนํ, มนสา อนุเปกฺขิตาติ จตุตฺถํ วิมุตฺตายตนํ, ทิฏฺฐิยา สุปฺปฏิวิทฺธาติ ปญฺจมํ วิมุตฺตายตนํ.}

\prose{โสตานุคตาย วิมุตฺติยา วจสา ยา วาจา สุปฺปฏิวิทฺธา อนุปุพฺพ\-ธมฺมสฺส โสเตน สุตฺวา สีลกฺขนฺเธ ปริปูเรติ, มนสา อนุเปกฺขิตา สมาธิ\-กฺขนฺธํ ปริปูเรติ, ทิฏฺฐิยา สุปฺปฏิวิทฺธา ปญฺญากฺขนฺธํ ปริปูเรติ.}

\prose{โสตานุคตา ธมฺมา พหุสฺสุตา โหนฺตีติ วิตฺถาเรน กาตพฺพํ. อิทํ ปฐมํ สทฺธาปทานํ มนสา อนุเปกฺขิตาติ ปฏิสลฺลาน\-พหุโล วิหรติ, วิตฺถาเรน กาตพฺพํ. อิทํ ทุติยํ สทฺธาปทานํ ทิฏฺฐิยา สุปฺปฏิวิทฺธาติ อนาสวา เจโตวิมุตฺติยา นาปรํ อิตฺถตฺตายาติ ปชานาตีติ. อิทํ ตติยํ สทฺธาปทานํ.}

\prose{โสตานุคตา ธมฺมาติ เสกฺขํ สตฺถา ทสฺเสติ. มนสา อนุเปกฺขิตาติ อรหตฺตํ สตฺถา ทสฺเสติ. ทิฏฺฐิยา สุปฺปฏิวิทฺธาติ ตถาคตํ อรหนฺตํ สมฺมา\-สมฺพุทฺธํ สตฺถา ทสฺเสติ.}

\prose{\csromanfolio{322}โสตานุคตา ธมฺมาติ กามานํ นิสฺสรณํ ทสฺเสติ. มนสา อนุเปกฺขิตาติ รูปธาตุยา นิสฺสรณํ ทสฺเสติ. ทิฏฺฐิยา สุปฺปฏิวิทฺธาติ เตธาตุกานํ นิสฺสรณํ ทสฺเสติ. อยํ สุตฺตตฺโถ.}

\proseitem{115}{ตตฺถ กตโม เทสนาหาโร. อิมมฺหิ สุตฺเต ตโย เอสนา เทสิตา โสตานุคเตหิ ธมฺเมหิ วจสา ปริจิเตหิ กาเมสนาย สมถมคฺโค. ทิฏฺฐิยา สุปฺปฏิวิทฺเธหิ พฺรหฺมจริเย\-สนาย สมถมคฺโค.}

\prose{วิจโยติ ยถา สุตฺตํ มนสิกโรนฺโต วิจินนฺโต สุตมยิปญฺญํ ปฏิลภติ. ยถา จ โส มนสิกโรตีติ ยถา สุตธมฺมา ตทา จินฺตามยิ\-ปญฺญํ ปฏิลภติ. ยถา ทิฏฺเฐว ธมฺเม มนสิกโรติ ตทา ภาวนามยิ\-ปญฺญํ ปฏิลภติ. อยํ วิจโย.}

\prose{สุเตน สุตมยิปญฺญํ ปฏิลภติ. จินฺตาย จินฺตามยิ\-ปญฺญํ ภาวนาย ภาวนามยิ\-ปญฺญํ ปฏิลภติ. อตฺถิ เอสา ยุตฺติ.}

\prose{ปทฏฺฐาโนติ โสตานุคตา ธมฺมาติ ธมฺม\-สฺสวนสฺส ปทฏฺฐานํ. วจสา ปริจิตาติ ยุญฺชนาย ปทฏฺฐานํ. มนสา อนุเปกฺขิตาติ ธมฺมา\-นุธมฺมาย วิปสฺสนาย ปทฏฺฐานํ. ทิฏฺฐิยา อนุเปกฺขิตาติ ปญฺญายปิ อนุเปกฺขิตา ทิฏฺฐิยาปิ อนุเปกฺขิตา.}

\prose{จตุพฺยูโหติ อิมมฺหิ สุตฺเต ภควโต โก อธิปฺปาโย. เย อิมาหิ ทฺวีหิ ปญฺญาหิ สมนฺนาคตา เตหิ...}

\prose{ส นิพฺพุโตติ มคฺคผลํ อนุปาทิเสสญฺจ นิพฺพานธาตุํ มนฺเตติ, ทาเนน โอฬาริกานํ กิเลสานํ ปหานํ มนฺเตติ. สีเลน มชฺฌิมานํ, ปญฺญาย สุขุม\-กิเลสานํ มนฺเตติ, ราคโทสโมห\-กฺขยา ส นิพฺพุโตติ กตา จ ภูมิ.}

\prose{\csromansymbolfootnote{*}{เหฏฺฐา 187 ปิฏฺเฐ ปสฺสิตพฺพํ.}ททโต ปุญฺญํ ปวฑฺฒติ,}

\prose{สํยมโต เวรํ น จียติ.}

\prose{กุสโล จ ชหาติ ปาปกนฺติ มคฺโค วุตฺโต. ราคโทสโมห\-กฺขยา ส นิพฺพุโตติ มคฺคผลมาห.}

\prose{\csromanfolio{323}ททโต ปุญฺญํ ปวฑฺฒติ, สํยมโตติ ตีหิ ปเทหิ โลกิกํ กุสลมูลํ วุตฺตํ. ราคโทสโมห\-กฺขยา ส นิพฺพุโตติ โลกุตฺตรํ กุสลมูลํ วุตฺตํ.}

\prose{ททโต ปุญฺญํ ปวฑฺฒติ, สํยมโต เวรํ น จียตีติ ปุถุชฺชนภูมิํ มนฺเตติ. กุสโล จ ชหาติ ปาปกนฺติ เสกฺขภูมิํ มนฺเตติ. ราคโทสโมห\-กฺขยา ส นิพฺพุโตติ อเสกฺขภูมิ วุตฺตา.}

\prose{ททโต ปุญฺญํ ปวฑฺฒติ, สํยมโต เวรํ น จียตีติ มคฺคนิยา ปฏิปทา วุตฺตา. กุสโล จ ชหาติ ปาปกนฺติ เสกฺขวิมุตฺติ. ราคโทสโมห\-กฺขยา ส นิพฺพุโตติ อเสกฺขวิมุตฺติ.}

\prose{ททโต ปุญฺญํ ปวฑฺฒติ, สํยมโต เวรํ น จียตีติ ทานกถํ สีลกถํ มคฺคกถํ โลกิกานํ ธมฺมานํ เทสนมาห. กุสโล จ ชหาติ ปาปกนฺติ โลเก อาทีนวา\-นุปสฺสนา. ราคโทสโมห\-กฺขยา ส นิพฺพุโตติ สามุกฺกํสิกาย ธมฺมเทสนายปิ ปฏิวิทฺธา.}

\prose{ททโต ปุญฺญํ ปวฑฺฒตีติ ปาณานํ อภยทาเนน ปาณาติปาตา เวรมณิ\-สตฺตานํ อภยํ เทติ. เอวํ สพฺพานิ สิกฺขาปทานิ กาตพฺพานิ. สํยมโต เวรํ น จียตีติ สีเล ปติฏฺฐาย จิตฺตํ สํยเมติ, ตสฺส สํยมโต ปาริปูริํ คจฺฉติ. ราคโทสโมห\-กฺขยา ส นิพฺพุโตติ เทฺว วิมุตฺติโย. อยํ สุตฺตนิทฺเทโส.}

\proseitem{116}{ตตฺถ กตมา เทสนา. อิมมฺหิ สุตฺเต กิํ เทสิตํ, เทฺว สุคติโย เทวา จ มนุสฺสา จ, ทิพฺพา จ ปญฺจกามคุณา, มานุสฺสกา จ. ทฺวีหิ ปเทหิ นิทฺเทโส. ททโต ปุญฺญํ ปวฑฺฒติ, สํยมโต เวรํ น จียติ, กุสโล จ ชหาติ ปาปกนฺติ มคฺโค วุตฺโต. ราคโทสโมห\-กฺขยา ส นิพฺพุโตติ เทฺว นิพฺพานธาตุโย เทสิตา โสปาทิเสสา จ อนุปาทิเสสา จ. อยํ เทสนา.}

\prose{วิจโยติ ททโต ปุญฺญํ ปวฑฺฒตีติ อิมินา ปฐเมน ปเทน ทาน\-มยิ\-กปุ\-ญฺญ\-กิริยวตฺถุ วุตฺตํ. เตนสฺส อานนฺตริยานํ กุสลานํ ธมฺมานํ. ทุติเยน ปเทน...}

\prose{ยนฺติ, นิยฺยานิกํ สาสนนฺติ, อยํ อธิปฺปาโย อสฺสวเนน จ อมนสิกาเรน จ อปฺปฏิเวเธน จ สกฺกาย\-สมุทย\-คามินี ปฏิปทา วุตฺตา. \csromanfolio{324}สวเนน จ มนสิกาเรน จ ปฏิเวเธน จ สกฺกาย\-นิโรธ\-คามินี ปฏิปทา วุตฺตา. อยํ อาวฏฺโฏ.}

\prose{วิภตฺตีติ เอกํส\-พฺยากรณีโย. นตฺถิ ตตฺถ วิภตฺติยา ภูมิ. ปริวตฺตนาติ เย ปญฺจานิสํสา, เต ปญฺจาทินา ปฏิปกฺเขน เตนว ทิฏฺเฐว ธมฺเม ปาปุณาติ, ตํ อุปปชฺชมานา อปโร ปริยาโย.}

\prose{เววจนนฺติ โสตานุคตา ธมฺมาติ ยํ สุตฺตํ ทิฏฺฐมฺปิ ปญฺญินฺทฺริยํ วิญฺญตฺตมฺปิ ทิฏฺฐิยา สุปฺปฏิวิทฺธมฺปิ วิภาวิตมฺปิ.}

\prose{ปญฺญตฺตีติ โสตานุคตา\-ธมฺมาติ เทสนา อวิชฺชา\-ปญฺญตฺติยา ปญฺญตฺตํ. มนสิกาโร ปาโมชฺช\-ปญฺญตฺติยา ปญฺญตฺโต, ทิฏฺฐธมฺมาปิ อานิสํส\-ปญฺญตฺติยา ปญฺญตฺตา.}

\prose{โอตรโณติ ติสฺโส ปญฺญา วจสา ปริจิเตสุ สุตมยีปญฺญา มนสา อนุเปกฺขิเตสุ จินฺตามยีปญฺญา ทิฏฺฐิยา สุปฺปฏิวิทฺธาสุ ภาวนามยี\-ปญฺญา. อิมานิ อริยสจฺจานิ อินฺทฺริยานิ วิชฺชุปฺปาทา อวิชฺชานิโรโธ ปฏิจฺจ\-สมุปฺปาโท อินฺทฺริเยสุ ตีณิ อินฺทฺริยานิ, อายตเนสุ ธมฺมายตน\-ปริยาปนฺนา ธาตูสุ ธมฺมธาตุ\-ปริยาปนฺนาติ.  โสธโนติ โย อารมฺโภ สุตฺตสฺส ปเวโส นิยุตฺโต.}

\prose{อธิฏฺฐาโนติ ปญฺจานิสํสาติ เวมตฺตตาย ปญฺญตฺตา อานิสํสา โสตา อนุคตาติ เวมตฺตตาย อริยโวหาโร ปญฺญตฺโต, ธมฺเม จ สวนนฺติ เอกตฺตตาย ปญฺญตฺตํ.}

\prose{ปริกฺขาโรติ ธมฺม\-สฺสวนสฺส ปยิรุปาสนา ปจฺจโย, สทฺธา เหตุ. มนสา อนุเปกฺขิตาติ อตฺถปฺปฏิสํเวทิตา ปจฺจโย, ธมฺมปฺปฏิสํเวทิตา เหตุ, ทิฏฺฐิยา สุปฺปฏิวิทฺธาติ สทฺธมฺมสฺสวนญฺจ มนสิกาโร จ ปจฺจโย, สุตมยี จินฺตามยี ปญฺญา เหตุ.  สมาโรปโนติ วิภตฺตํ สุตฺตํ อปโร ปริยาโย นิพฺพตฺติ พเล นตฺถิ. ตตฺถ สมาโรปนาย ภูมิ.}

\proseitem{117}{ตตฺถ กตมํ วาสนาภาคิยญฺจ นิพฺเพธภาคิยญฺจ สุตฺตํ. ททโต ปุญฺญํ ปวฑฺฒตีติ คาถา. ททโตติ ทาน\-มยิ\-กปุ\-ญฺญ\-กิริยวตฺถุ วุตฺตํ. สํยมโต เวรํ น จียตีติ สีลมยิ\-กปุ\-ญฺญ\-กิริยวตฺถุ วุตฺตํ. กุสโล จ ชหาติ ปาปกนฺติ โลภสฺส จ โมหสฺส จ พฺยาปาทสฺส จ ปหานมาห. ราคโทสโมห\-กฺขยา ส นิพฺพุโตติ โลภสฺส จ โมหสฺส จ \csromanfolio{325}พฺยาปาทสฺส จ ฉนฺทราควินยมา\-หาติ. ททโต ปุญฺญํ ปวฑฺฒตีติ คาถา อโลโภ กุสลมูลํ ภวติ. สํยมโต เวรํ น จียตีติ อโทโส กุสลมูลํ ภวติ. สํยมโต เวรํ น จียตีติ อเวรา อสปตฺตา อพฺยาปาทตาย สทา. กุสโล จ ชหาติ ปาปกนฺติ ญาณุปฺปาทา อญฺญาณนิโรโธ. จตุตฺถปเทน ราคโทสโมห\-กฺขเยน ราควิราคา เจโตวิมุตฺติ\-โมหกฺขเยน อวิชฺชาวิราคา ปญฺญาวิมุตฺติ, อยํ วิจโย.}

\prose{ยุตฺตีติ ทาเน ฐิโต อุภยํ หิ ปริปูเรติ. มจฺฉริยญฺจ ปชหติ. ปุญฺญญฺจ ปวฑฺฒติ. อตฺถิ เอสา ยุตฺติ.}

\prose{ปทฏฺฐานนฺติ ททโต ปุญฺญํ ปวฑฺฒตีติ จาคา\-ธิฏฺฐานสฺส ปทฏฺฐานํ. สํยมโต เวรํ น จียตีติ ปญฺญา\-ธิฏฺฐานสฺส ปทฏฺฐานํ กุสโล จ ชหาติ ปาปกนฺติ สจฺจา\-ธิฏฺฐานสฺส ปทฏฺฐานํ. ราคโทสโมห\-กฺขยา ส นิพฺพุโตติ อุปสมา\-ธิฏฺฐานสฺส ปทฏฺฐานํ. อยํ ปทฏฺฐาโน.}

\prose{ตตฺถ กตโม ลกฺขโณ. ททโต ปุญฺญํ ปวฑฺฒติ สํยมโต เวรํ น จียติ. ททโตปิ เวรํ น กริยาติ กุสโล จ ชหาติ ปาปกํ ราคโทสโมห\-กฺขยา ส นิพฺพุโต รูปกฺขยาปิ เวทนกฺขยาปิ, เยน รูเปน ทิฏฺฐํ, เตน ตถาคโต ปญฺญเปนฺโต ปญฺญเปยฺย รูปสฺส ขยา วิราคนิโรธาติ เอวํ ปญฺจกฺขนฺธา.}

\prose{จตุพฺยูโห อิธ ภควโต โก อธิปฺปาโย, เย มหาโภคานํ ปตฺถยิสฺสนฺติ. เต ทานํ ทสฺสนฺติ ปริสฺสย\-ปหานาย, เย อเวราภิ\-ฉนฺทกา, เต ปญฺจ เวรานิ ปชหิสฺสนฺติ, เย กุสลาภิ\-ฉนฺทกา, เต อฏฺฐงฺคิกํ มคฺคํ ภาเวสฺสนฺติ อฏฺฐนฺนํ มิจฺฉตฺตานํ ปหานาย. เย นิพฺพายิตุกามา, เต ราคโทสโมหํ ปชหิสฺสนฺตีติ อยํ ภควโต อธิปฺปาโย.}

\prose{อาวฏฺโฏติ ยญฺจ อททโต มจฺฉริยํ ยญฺจ สํยมโต เวรํ ยญฺจ อกุสลสฺส ปาปสฺส อปฺปหานํ, อยํ ทุกฺขนิทฺเทโส น สมุทโย. อโลเภน จ อโทเสน จ อโมเหน จ กุสเลน อิมานิ ตีณิ กุสลมูลานิ. เตสํ ปจฺจโย อฏฺฐ สมฺมตฺตานิ, อยํ มคฺโค. เตสํ ราค\-โทส\-โมหานํ ขยา, อยํ นิโรโธ.}

\prose{วิภตฺตีติ ททโต ปุญฺญํ ปวฑฺฒตีติ น เอกํเสน โย ราช\-ทณฺฑภเยน เทติ, โย จ อกปฺปิยสฺส ปริโภเคน สีลวนฺเตสุ เทติ, น ตสฺส \csromanfolio{326}ปุญฺญํ ปวฑฺฒตีติ โส เจตํ ทานํ อกุสเลน เทติ, ทณฺฑทานํ สตฺถทานํ อปุญฺญมยํ ปวฑฺฒติ, น ปุญฺญํ. สํยมโต เวรํ น จียตีติ น เอกํเสน กิํ การณํ ยญฺจ โย ปทํ ทิฏฺฐ\-ธมฺมิกํ ปสฺสติ ยทิ มม ราชาโน คเหตฺวา หตฺถํ วา ฉินฺเทยฺย ฯเปฯ น เตน สํยเมน เวรํ น กโรติ. โย ตุ เอวํ สมาทิยติ ปาณาติปาตสฺส ปาปโก วิปาโกติ, ทิฏฺเฐ เยว ธมฺเม อภิสมฺปราเย จ เอวํ สพฺพสฺส อกุสลสฺส เหตุโต อารติ. อิมินา สํยเมน เวรํ น จียติ.}

\proseitem{7}{ปริวตฺตนาติ ททโต ปุญฺญํ ปวฑฺฒตีติ อททโต ปุญฺญํ น ปวฑฺฒติ. ยํ ทานมยํ, ตํ สํยมโต เวรํ น จียติ, อสํยมโต เวรํ กรียติ. กุสโล จ ชหาติ ปาปกํ อกุสโล น ชหาติ. ราคโทสโมห\-กฺขยา สนิพฺพุโตติ ทูตํ เปเสตฺวา ปณีตํ เปเสตฺวาปิ น ปกฺโกสามิ, โส สยเมว ปน มหา\-ภิกฺขุสํฆ\-ปริวาโร อมฺหากํ วสนฏฺฐานํ สมฺปตฺโต อเมฺหหิ จ สนฺถาคารสาลา\csromansharedfootnote{71dee749c466dd18}{สนฺธาคารสาลา (ก)} การิตา, เอตฺถ มยํ ทสพลํ อาเนตฺวา มงฺคลํ ภณาเปมาติ จินฺเตตฺวา อุปสงฺกมิํสุ. เยน สนฺถาคารํ เตนุปสงฺกมิํสูติ ตํ ทิวสํ กิร สนฺถาคาเร จิตฺตกมฺมํ นิฏฺฐาเปตฺวา อฏฺฏกา มุตฺตมตฺตา โหนฺติ. พุทฺธา นาม อรญฺญชฺฌาสยา อรญฺญารามา อนฺโตคาเม วเสยฺยุํ วา โน วาติ ตสฺมา ภควโต มนํ ชานิตฺวาว ปฏิชคฺคิสฺสามาติ จินฺเตตฺวา เต ภควนฺตํ อุปสงฺกมิํสุ. อิทานิ ปน มนํ ลภิตฺวา ปฏิชคฺคิตุ\-กามา เยน สนฺถาคารํ, เตนุ\-ปสงฺกมิํสุ. สพฺพ\-สนฺถรินฺติ ยถา สพฺพํ สนฺถตํ โหติ เอวํ เยน ภควา เตนุปสงฺกมิํสูติ. เอตฺถ ปน เต มลฺลราชาโน สนฺถาคารํ ปฏิชคฺคิตฺวา นคร\-วีถิโยปิ สมฺมชฺชาเปตฺวา ธเช อุสฺสาเปตฺวา สุวณฺณ\-ฆฏิ\-กทลิโย จ ฐปาเปตฺวา สกลนครํ ทีปมาลาหิ วิปฺปกิณฺณ\-ตารกํ วิย กตฺวา ขีรปเก\csromansharedfootnote{2fa2fabe6787f4c4}{ขิรุปเก (อิ, ก)} ทารเก ขีรํ ปาเยถ, ทหเร กุมาเร ลหุํ ลหุํ โภชาเปตฺวา สยาเปถ, อุจฺจาสทฺทํ มากริ, อชฺช เอกรตฺติํ สตฺถา อนฺโตคาเมว วสิสฺสติ, พุทฺธา นาม อปฺปสทฺทกามา โหนฺตีติ เภริํ จราเปตฺวา สยํ ทณฺฑกทีปิกา อาทาย เยน ภควา เตนุ\-ปสงฺกมิํสุ. ภควนฺตํ เยว ปุรกฺขตฺวาติ ภควนฺตํ ปุรโต กตฺวา, ตตฺถ ภควา ภิกฺขุนญฺเจว อุปาสกานญฺจ มชฺเฌ นิสินฺโน อติวิย วิโรจติ. สมนฺต\-ปาสาทิโก สุวณฺณวณฺโณ อภิรูโป ทสฺสนีโย \csromanfolio{327}ปุรตฺถิม\-กายโต สุวณฺณวณฺณา รสฺมิ อุฏฺฐหิตฺวา คคนตเล อสีติหตฺถํ ฐานํ คณฺหาติ. ปจฺฉิมกายโต ทกฺขิณ\-หตฺถโต วามหตฺถโต สุวณฺณวณฺณา เหฏฺฐา ปาทตเลหิ ปวาฬ\-วณฺณ\-รสฺมิ อุฏฺฐหิตฺวา ฆนปถวิยํ อสีติหตฺถํ ฐานํ คณฺหาติ, เอวํ สมนฺตา อสีติ\-หตฺถ\-มตฺตํ ฐานํ ฉพฺพณฺณ\-พุทฺธรสฺมิโย วิชฺโชตมานา วิตณฺฑมานา วิธาวนฺติ, สพฺเพ ทิสาภาคา สุวณฺณ\-จมฺปกปุปฺเผหิ วิกิริยมานา วิย สุวณฺณฆฏโต นิกฺขนฺต\-สุวณฺณ\-รส\-ธาราหิ สิญฺจมานา วิย ปสาริต\-สุวณฺณ\-ปฏ\-ปริกฺขิตฺตา วิย เวรมฺภวาตสมุฏฺฐิตกิํ สุกกิํ สุการกณิการปุปฺผจุณฺณสโมกิณฺณา วิย วิปฺปกสนฺตํ อสีติ อนุพฺยญฺชนพฺยามปฺปภา ทฺวตฺติํส\-วรลกฺขณ\-สมุชฺชลํ สรีรํ สมุคฺคต\-ตารกํ วิย คคนตลํ วิกสิตมิว ปทุมวนํ สพฺพ\-ผาลิผุลฺโล วิย โยชนสติโก ปาริจฺฉตฺตโก ปฏิปาฏิยา ฐปิตานํ ทฺวตฺติํส\-จนฺทานํ ทฺวตฺติํส\-สูริยานํ ทฺวตฺติํส\-จกฺกวตฺตีนํ ทฺวตฺติํส\-เทวราชานํ ทฺวตฺติํส\-มหาพฺรหฺมานํ นิพฺพุโต อเสกฺขสฺส นตฺถิ นิพฺพุติ.}

\prose{เววจนนฺติ ททโต ปุญฺญํ ปวฑฺฒติ, อนุโมทโตปิ ปุญฺญํ ปวฑฺฒติ. จิตฺตสฺส สมาทหโตปิ เวยฺยาวจฺจ\-กิริยายปิ ปุญฺญํ ปวฑฺฒตีติ.}

\prose{ปญฺญตฺตีติ ททโต ปุญฺญํ ปวฑฺฒติ, อโลภสฺส ปฏินิสฺสยฆาตปญฺญตฺติยา ปญฺญตฺตํ. สํยมโต เวรํ น จียตีติ อโทสสฺส ปฏินิสฺสยฆาตปญฺญตฺติยา ปญฺญตฺตํ กุสโล จ ชหาติ ปาปกนฺติ อโมหสฺส ปฏินิสฺสยฆาตปญฺญตฺติยา ปญฺญตฺตํ.}

\prose{โอตรโณติ ปญฺจสุ อินฺทฺริเยสุ ททโต ปุญฺญํ ปวฑฺฒติ, สํยมโต เวรํ น จียติ สํยเมน สูลกฺขนฺโธ. โอติณฺโณ ฉสุ อินฺทฺริเยสุ สํวโร, อยํ สมาธิ\-กฺขนฺโธ, ยํ กุสโล จ ชหาติ ปาปกํ, อยํ ปญฺญากฺขนฺโธ, ราคโทสโมห\-กฺขยา ส นิพฺพุโตติ วิมุตฺติ\-กฺขนฺโธ. ธาตูสุ ธมฺมธาตุ, อายตเนสุ มนายตนํ.}

\prose{โสธโนติ เยนารมฺเภน อิทํ สุตฺตํ เทสิตํ โส อารมฺโภ สุทฺโธ.}

\csromanmatikaanchor{mtk.47}
\csromantocmarkref{section}{8. สุตฺตเวภงฺคิย}{mtk.47}
\bookmark[dest={mtk.47},level=1]{8. สุตฺตเวภงฺคิย}
\prose{อธิฏฺฐาโน ทานนฺติ เอกตฺตตาย ปญฺญตฺตํ. จาโค ปริจฺจาโค ธมฺมทานํ อามิสทานํ, อฏฺฐ ทานานิ วิตฺถาเรน กาตพฺพานิ, อยํ เวมตฺตตา. น จ \csromanfolio{328}ททโต เอกตฺต\-ปญฺญตฺติยา ปญฺญตฺตํ. ขนฺตี อนวชฺชนฺติ ปญฺญตฺติยา ปญฺญตฺตํ. ราคโทสโมห\-กฺขยา ส นิพฺพุโตติ โรธ\-วีริย\-ปญฺญตฺติยา\csromansharedfootnote{e9b5717847cd624b}{โยธ วีริยปญฺญตฺติยา (อิ, ก)} ปญฺญตฺตา.}

\prose{ปริกฺขาโรติ ทานสฺส ปาโมชฺชํ ปจฺจโย, อโลโภ เหตุ. สํยมโต โยนิโส มนสิกาโร เหตุ, ปริจฺจาโค ปจฺจโย. กุสโล จ ชหาติ ปาปกนฺติ ยถาภูต\-ทสฺสนํ ปจฺจโย, ญาณปฺปฏิลาโภ เหตุ. ราคโทสโมห\-กฺขยา ส นิพฺพุโตติ ปรโต จ โฆโส อชฺฌตฺตญฺจ โยนิโส มนสิกาโร มคฺโค จ เหตุ จ ปจฺจโย จ.}

\prose{สมาโรปโนติ ททโต ปุญฺญํ ปวฑฺฒตีติ คาถา ตสฺส สีลมฺปิ วฑฺฒติ. สํยโมปิ วฑฺฒติ. สํยมโต เวรํ น จียตีติ. อญฺเญปิ กิเลสา น จียนฺติ เยปิสฺส ตปฺปจฺจยา อุปฺปชฺเชยฺยุํ อาสวา วิฆาตา, เตปิสฺส น อุปฺปชฺชนฺติ. ราคโทสโมห\-กฺขยา ส นิพฺพุโตติ ราคโทสสฺสาปิ ขยา ราคานุสยสฺสปิ ขยา โทสสฺส โมหสฺสาปิ ส นิพฺพุโตติ โสปาทิเสสา นิพฺพานธาตุ อนุปาทิเสสาปิ. อยํ สมาโรปโน.}

\csromanheader{2}{เถรสฺส มหา\-กจฺจายนสฺส เปฏโกปเทเส}
\nitthitam{หารสฺส สมฺปาตภูมิ สมตฺตา.}
\proseitem{118}{ปุพฺพา โกฏิ น ปญฺญายติ\csromansharedfootnote{3b7d62f75bd77c34}{สํ 1. 387 ปิฏฺเฐ ปสฺสิตพฺพํ.} อวิชฺชาย จ ภวตณฺหาย จ. ตตฺถ อวิชฺชา\-นีวรณานํ ตณฺหา\-สํโยชนานํ สตฺตานํ ปุพฺพโกฏิ น ปญฺญายติ. ตตฺถ เย สตฺตา ตณฺหาสํโยชนา, เต อชฺโฌสาน\-พหุลา มนฺทวิปสฺสกา. เย ปน อุสฺสนฺน\-ทิฏฺฐิกา สตฺตา, เต วิปสฺสนา\-พหุลา มนฺทชฺโฌสานา.}

\prose{ตตฺถ ตณฺหาจริตา สตฺตา สตฺตสญฺญาภิ\-นิวิฏฺฐา อนุปฺปาท\-วย\-ทสฺสิโน. เต ปญฺจสุ ขนฺเธสุ อตฺตานํ สมนุปสฺสนฺติ “รูปวนฺตํ วา อตฺตานํ, อตฺตนิ วา รูปํ, รูปสฺมิํ วา อตฺตานนฺ”ติ. เอวํ ปญฺจกฺขนฺธา. อญฺเญหิ ขนฺเธหิ อตฺตานํ สมนุปสฺสนฺติ ตสฺส อุสฺสนฺน\-ทิฏฺฐิกา สตฺตา วิปสฺสมานา ขนฺเธ อุชุํ อตฺตโต สมนุปสฺสนฺติ. เต รูปํ อตฺตโต สมนุปสฺสนฺติ. ยํ รูปํ, โส อตฺตา. โย อหํ, ตํ รูปํ. โส รูปวินาสํ ปสฺสติ, อยํ \csromanfolio{329}อุจฺเฉทวาที. อิติ ปญฺจนฺนํ ขนฺธานํ ปฐมาภิ\-นิปาตา สกฺกาย\-ทิฏฺฐิโย ปญฺจ อุจฺเฉทํ ภชนฺติ “ตํ ชีวํ ตํ สรีรนฺ”ติ. เอกเมกมฺหิ ขนฺเธ ตีหิ ปเทหิ ปจฺฉิมเกหิ สสฺสตํ ภชติ “อญฺญํ ชีวํ อญฺญํ สรีรนฺ”ติ. อิโต พหิทฺธาเต ปพฺพชิตา ตณฺหาจริตา กามสุขลฺลิกา\-นุโยคมนุยุตฺตา วิหรนฺติ. เตน เย จ นิสฺสนฺเทน ทิฏฺฐิจริตา อตฺตกิลมถา\-นุโยคมนุยุตฺตา วิหรนฺติ. เตน เยว ทิฏฺฐิสุเขน เอตฺตาวตา พาหิรโก ปโยโค.}

\prose{ตตฺถ ทิฏฺฐิจริตา สตฺตา เย อริย\-ธมฺมวินยํ โอตรนฺติ, เต ธมฺมานุสาริโน โหนฺติ. เย ตณฺหาจริตา สตฺตา อริยํ ธมฺมวินยํ โอตรนฺติ, เต สทฺธานุสาริโน โหนฺติ}

\prose{ตตฺถ เย ทิฏฺฐิจริตา สตฺตา, เต กาเมสุ โทสทิฏฺฐี, น จ เย กาเมสุ อนุสยา สมูหตา, เต อตฺตกิลมถา\-นุโยคมนุยุตฺตา วิหรนฺติ. เตสํ สตฺถา ธมฺมํ เทเสติ. อญฺโญ วา สาวโก กาเมหิ นตฺถิ อตฺโถติ เต จ ปุพฺเพเยว กาเมหิ อนตฺถิกา อิติ กาเม อปฺปกสิเรน ปฏินิสฺสชฺชนฺติ. เต เจตสิเกน ทุกฺเขน อนชฺโฌสิตา. เตน วุจฺจติ “สุขา ปฏิปทา”ติ. เย ปน ตณฺหาจริตา สตฺตา, เต กาเมสุ อชฺโฌสิตา, เตสํ สตฺถา วา ธมฺมํ เทเสติ. อญฺญตโร วา ภิกฺขุ กาเมหิ นตฺถิ อตฺโถติ, เต ปิยรูปํ ทุกฺเขน ปฏินิสฺสชฺชนฺติ. เตน วุจฺจติ “ทุกฺขา ปฏิปทา”ติ. อิติ อิเม สพฺพสตฺตา ทฺวีสุ ปฏิปทาสุ สโมสรณํ คจฺฉนฺติ ทุกฺขายญฺจ สุขายญฺจ.}

\prose{ตตฺถ เย ทิฏฺฐิจริตา สตฺตา, เต ทฺวิธา มุทินฺทฺริยา จ ติกฺขินฺทฺริยา จ. ตตฺถ เย ทิฏฺฐิจริตา สตฺตา ติกฺขินฺทฺริยา สุเขน ปฏินิสฺสชฺชนฺติ, ขิปฺปญฺจ อภิสเมนฺติ, เตน วุจฺจติ “ขิปฺปาภิญฺญา สุขา ปฏิปทา”ติ. ตตฺถ เย ทิฏฺฐิจริตา สตฺตา มุทินฺทฺริยา ปฐมํ ติกฺขินฺทฺริยํ อุปาทาย ทนฺธตรํ อภิสเมนฺติ, เต สุเขน ปฏินิสฺสชฺชนฺติ, ทนฺธญฺจ อภิสเมนฺติ. เตน วุจฺจติ “สุขา ปฏิปทา ทนฺธาภิญฺญา”ติ. ตตฺถ ตณฺหาจริตา สตฺตา ทฺวิธา ติกฺขินฺทฺริยา จ มุทินฺทฺริยา จ. ตตฺถ เย ตณฺหาจริตา สตฺตา ติกฺขินฺทฺริยา ทุกฺเขน ปฏินิสฺสชฺชนฺติ, ขิปฺปญฺจ อภิสเมนฺติ. เตน วุจฺจติ “ทุกฺขา ปฏิปทา ขิปฺปาภิญฺญา”ติ. ตตฺถ เย ตณฺหาจริตา สตฺตา มุทินฺทฺริยา ปฐมํ ติกฺขินฺทฺริยํ อุปาทาย ทนฺธตรํ อภิสเมนฺติ, เต ทุกฺเขน ปฏินิสฺสชฺชนฺติ, ทนฺธญฺจ อภิสเมนฺติ. เตน วุจฺจติ “ทุกฺขา ปฏิปทา ทนฺธาภิญฺญา”ติ. อิมา จตสฺโส ปฏิปทาโย อปญฺจมา อฉฏฺฐา. เย หิ เกจิ นิพฺพุตา \csromanfolio{330}นิพฺพายิสฺสนฺติ วา อิมาหิ จตูหิ ปฏิปทาหิ อนญฺญาหิ อยํ ปฏิปทา\-จตุกฺเกน กิเลเส นิทฺทิสติ. ยา จตุกฺกมคฺเคน อริยธมฺเมสุ นิทฺทิสิตพฺพา, อยํ วุจฺจติ สีหวิกฺกีฬิโต นาม นโย.}

\proseitem{119}{ตตฺริเม จตฺตาโร อาหารา. จตฺตาโร วิปลฺลาสา อุปาทานา โยคา คนฺถา อาสวา โอฆา สลฺลา วิญฺญาณ\-ฏฺฐิติโย อคติคมนาติ, เอวํ อิมานิ สพฺพานิ ทส ปทานิ. อยํ สุตฺตสฺส สํสนฺทนา.}

\prose{จตฺตาโร อาหารา. ตตฺถ โย จ กพฬีกาโร อาหาโร โย จ ผสฺโส อาหาโร, อิเม ตณฺหาจริเตน ปหาตพฺพา. ตตฺถ โย จ มโนสญฺเจตนา\-หาโร โย จ วิญฺญาณาหาโร, อิเม ทิฏฺฐิจริเตน ปหาตพฺพา.}

\prose{ปฐโม อาหาโร ปฐโม วิปลฺลาโส, ทุติโย อาหาโร ทุติโย วิปลฺลาโส, ตติโย อาหาโร ตติโย วิปลฺลาโส, จตุตฺโถ อาหาโร จตุตฺโถ วิปลฺลาโส. อิเม จตฺตาโร วิปลฺลาสา อปญฺจมา อฉฏฺฐา. อิทญฺจ ปมาณา จตฺตาโร อาหารา.}

\prose{ตตฺถ ปฐเม วิปลฺลาเส ฐิโต กาเม อุปาทิยติ, อิทํ กามุปาทานํ. ทุติเย วิปลฺลาเส ฐิโต อนาคตํ ภวํ อุปาทิยติ, อิทํ สีลพฺพตุ\-ปาทานํ. ตติเย วิปลฺลาเส ฐิโต วิปรีโต ทิฏฺฐิํ อุปาทิยติ, อิทํ ทิฏฺฐุปาทานํ. จตุตฺเถ วิปลฺลาเส ฐิโต ขนฺเธ อตฺตโต อุปาทิยติ, อิทํ อตฺตวาทุ\-ปาทานํ.}

\prose{ตตฺถ กามุปาทาเน ฐิโต กาเม อภิชฺฌายติ คนฺถติ, อยํ อภิชฺฌา\-กายคนฺโถ. สีลพฺพตุ\-ปาทาเน ฐิโต พฺยาปาทํ คนฺถติ, อยํ พฺยาปาท\-กายคนฺโถ. ทิฏฺฐุปาทาเน ฐิโต ปรามาสํ คนฺถติ, อยํ ปรามาส\-กายคนฺโถ. อตฺตวาทุปาทาเน ฐิโต ปปญฺจนฺโต คนฺถติ, อยํ อิทํสจฺจา\-ภินิเวโส กายคนฺโถ.}

\prose{ตสฺส คนฺถิตา กิเลสา อาสวนฺติ. กิญฺจิ ปน วุจฺจติ วิปฺปฏิสาโร, เย วิปฺปฏิสารา\csromansharedfootnote{ceae2a6e2c9689bd}{โย วิปฺปฏิสาโร (อิ, ก)} เต อนุสยา. ตตฺถ อภิชฺฌา\-กายคนฺเถน กามาสโว, พฺยาปาท\-กายคนฺเถน ภวาสโว, ปรามาส\-กายคนฺเถน ทิฏฺฐาสโว, อิทํ สจฺจาภิ\-นิเวส\-กายคนฺเถน อวิชฺชาสโว.}

\prose{\csromanfolio{331}เต จตฺตาโร อาสวา เวปุลฺลภาวํ คตา โอฆา โหนฺติ, เตน วุจฺจนฺติ “โอฆา”ติ. ตตฺถ กามาสโว กาโมโฆ, ภวาสโว ภโวโฆ, อวิชฺชาสโว อวิชฺโชโฆ, ทิฏฺฐาสโว ทิฏฺโฐโฆ.}

\prose{เต จตฺตาโร โอฆา อาสยม\-นุปวิฏฺฐา อนุสย\-สหคตา วุจฺจนฺติ. สลฺลาติ หทยมาหจฺจ ติฏฺฐนฺตา. ตตฺถ กาโมโฆ ราคสลฺลํ, ภโวโฆ โทสสลฺลํ, อวิชฺโชโฆ โมหสลฺลํ, ทิฏฺโฐโฆ ทิฏฺฐิสลฺลํ.}

\prose{อิเมหิ จตูหิ สลฺเลหิ ปริยาทินฺนํ วิญฺญาณํ จตูสุ ธมฺเมสุ ติฏฺฐติ รูเป เวทนาย สญฺญาย สงฺขาเรสุ. อิมา จตสฺโส วิญฺญาณ\-ฏฺฐิติโย. ตตฺถ ราคสลฺเลน นนฺทูปเสจนํ รูปูปคํ วิญฺญาณํ ติฏฺฐติ. โทสสลฺเลน เวทนุปคํ. โมหสลฺเลน สญฺญูปคํ. ทิฏฺฐิสลฺเลน นนฺทูปเสจนํ สงฺขารูปคํ วิญฺญาณํ ติฏฺฐติ.}

\prose{จตูหิ วิญฺญาณ\-ฏฺฐิตีหิ จตุพฺพิธํ อคติํ คจฺฉนฺติ ฉนฺทา โทสา ภยา โมหา. ราเคน ฉนฺทา อคติํ คจฺฉติ, โทเสน โทสา อคติํ คจฺฉติ, โมเหน โมหา อคติํ คจฺฉติ, ทิฏฺฐิยา ภยา อคติํ คจฺฉติ. อิติ อิทญฺจ กมฺมํ อิเม จ กิเลสา. อยํ สํสารสฺส เหตุ.}

\proseitem{120}{ตตฺถิมา จตสฺโส ทิสา กพฬีการาหาโร “อสุเภ สุภนฺ”ติ วิปลฺลาโส กามุปาทานํ กามโยโค อภิชฺฌา\-กายคนฺโถ กามาสโว กาโมโฆ ราคสลฺลํ รูปูปคา วิญฺญาณฏฺฐิติ ฉนฺทา อคติคมนํ. อยํ ปฐมา ทิสา.}

\prose{ผสฺโส อาหาโร “ทุกฺเข สุขนฺ”ติ วิปลฺลาโส สีลพฺพตุ\-ปาทานํ ภวโยโค\-พฺยาปาโท กายคนฺโถ ภวาสโว ภโวโฆ โทสสลฺลํ เวทนุปคา วิญฺญาณฏฺฐิติ โทสา อคติคมนํ, อยํ ทุติยา ทิสา.}

\prose{มโนสญฺเจตนา\-หาโร “อนตฺตนิ อตฺตา”ติ วิปลฺลาโส ทิฏฺฐุปาทานํ ทิฏฺฐิโยโค ปรามาส\-กายคนฺโถ ทิฏฺฐาสโว ทิฏฺโฐโฆ ทิฏฺฐิสลฺลํ สญฺญูปคา วิญฺญาณฏฺฐิติ ภยา อคติคมนํ. อยํ ตติยา ทิสา.}

\prose{วิญฺญณหาโร “อนิจฺเจ นิจฺจนฺ”ติ วิปลฺลาโส อตฺตวาทุ\-ปาทานํ อวิชฺชาโยโค อิทํสจฺจา\-ภินิเวโส กายคนฺโถ อวิชฺชาสโว อวิชฺโชโฆ โมหสลฺลํ สงฺขารูปคา วิญฺญาณฏฺฐิติ โมหา อคติคมนํ, \csromanfolio{332}อยํ จตุตฺถี ทิสา. อิติ อิเมสํ ทสนฺนํ สุตฺตานํ ปฐเมน ปเทน ปฐมาย ทิสาย อาโลกนํ. อยํ วุจฺจติ ทิสาโลกนา.}

\prose{จตูหิ วิปลฺลาเสหิ อกุสลปกฺเข ทิสาวิโลกนา กิเลสํ สํโยเชตฺวา อยํ อกุสลปกฺเข ทิสาวิโลกนาย ภูมิ ปญฺจนฺนํ ทสนฺนํ สุตฺตานํ ยานิ ปฐมานิ ปทานิ อิเมสํ ธมฺมานํ โก อตฺโถ, เอโก อตฺโถ, พฺยญฺชนเมว นานํ. เอวํ ทุติยา เอวํ ตติยา เอวํ จตุตฺถี. อยํ ปฐมา สํสนฺทนา.}

\prose{อิมินา เปยฺยาเลน สพฺเพ กิเลสา จตูสุ ปเทสุ ปกฺขิปิตพฺพา. ตโต กุสลปกฺเข จตสฺโส ปฏิปทา\csromansharedfootnote{34241fcce8b5b73a}{สํ 1. 467 ปิฏฺเฐ ปสฺสิตพฺพํ.} จตฺตาริ ฌานานิ จตฺตาโร สติปฏฺฐานา จตฺตาโร วิหารา ทิพฺโพ พฺรหฺมา อริโย อาเนญฺโช จตฺตาโร สมฺมปฺปธานา จตฺตาโร อจฺฉริยา อพฺภุตธมฺมา จตฺตาโร อธิฏฺฐานา จตฺตาโร สมาธโย ฉนฺทสมาธิ วีริยสมาธิ จิตฺตสมาธิ วีมํสาสมาธิ. จตฺตาโร ธมฺมา สุขภาคิยา นาญฺญตฺร โพชฺฌงฺคา นาญฺญตฺร ตปสา นาญฺญติ\-นฺทฺริยสํวรา นาญฺญตฺร สพฺพนิสฺสคฺคา จตฺตาริ อปฺปมาณานิ.}

\prose{ตตฺถ ทุกฺขา ปฏิปทา ทนฺธาภิญฺญา ภาวิยมานา พหุลี\-กริยมานา ปฐมํ ฌานํ ปริปูเรติ, ปฐมํ ฌานํ ปริปุณฺณํ ปฐมํ สติปฏฺฐานํ ปริปูเรติ, ปฐมํ สติปฏฺฐานํ ปริปุณฺณํ ปฐมํ วิหารํ ปริปูเรติ, ปฐโม วิหาโร ปริปุณฺโณ ปฐมํ สมฺมปฺปธานํ ปริปูเรติ, ปฐมํ สมฺมปฺปธานํ ปริปุณฺณํ ปฐมํ อจฺฉริยํ อพฺภุตธมฺมํ ปริปูเรติ, ปฐโม อจฺฉริโย อพฺภุโต ธมฺโม ปริปุณฺโณ ปฐมํ อธิฏฺฐานํ ปริปูเรติ, ปฐมํ อธิฏฺฐานํ ปริปุณฺณํ ฉนฺทสมาธิํ ปริปูเรติ, ฉนฺทสมาธิ ปริปุณฺโณ อินฺทฺริย\-สํวรํ ปริปูเรติ, อินฺทฺริยสํวโร ปริปุณฺโณ ปฐมํ เมตฺตา- อปฺปมาณํ ปริปูเรติ. เอวํ ยาว สพฺพนิสฺสคฺโค จตุตฺถํ อปฺปมาณํ ปริปูเรติ.}

\prose{ตตฺถ ปฐมา จ ปฏิปทา ปฐมญฺจ ฌานํ ปฐมญฺจ สติปฏฺฐานํ ปฐโม จ วิหาโร ปฐมญฺจ สมฺมปฺปธานํ ปฐโม จ อจฺฉริโย อพฺภุโต ธมฺโม สจฺจาธิฏฺฐานญฺจ ฉนฺทสมาธิ จ อินฺทฺริยสํวโร จ เมตฺตา จ อปฺปมาณํ. อยํ ปฐมา ทิสา.}

\prose{ทุกฺขา จ\csromansharedfootnote{0dab868ddc91ee25}{ทุติยา จ (ก)} ปฏิปทา ขิปฺปาภิญฺญา ทุติยํ ฌานํ ทุติยญฺจ สติปฏฺฐานํ ทุติโย จ วิหาโร ทุติยญฺจ สมฺมปฺปธานํ ทุติโย จ อจฺฉริโย \csromanfolio{333}อพฺภุโต ธมฺโม จาคาธิฏฺฐานํ จิตฺตสมาธิ จตฺตาโร อิทฺธิปาทา กรุณา จ อปฺปมาณํ, อยํ ทุติยา ทิสา.}

\prose{สุขา จ\csromansharedfootnote{cbea2e174baa8858}{ตติยา จ (ก)} ปฏิปทา ทนฺธาภิญฺญา ตติยญฺจ ฌานํ ตติยญฺจ สติปฏฺฐานํ ตติโย จ วิหาโร ตติยญฺจ สมฺมปฺปธานํ ตติโย จ อจฺฉริโย อพฺภุโต ธมฺโม ปญฺญาธิฏฺฐานญฺจ วีริยสมาธิ จ โพชฺฌงฺคา จ มุทิตา จ อปฺปมาณํ. อยํ ตติยา ทิสา.}

\prose{สุขา จ\csromansharedfootnote{a1a90973aa3f26c5}{จตุตฺถี จ (ก)} ปฏิปทา ขิปฺปาภิญฺญา จตุตฺถํ ฌานํ จตุตฺถญฺจ สติปฏฺฐานํ จตุตฺโถ จ วิหาโร จตุตฺถญฺจ สมฺมปฺปธานํ จตุตฺโถ จ อจฺฉริโย อพฺภุโต ธมฺโม อุปสมาธิฏฺฐานญฺจ วีมํสาสมาธิ จ สพฺพนิสฺสคฺโค จ อุเปกฺขา อปฺปมาณญฺจ. อยํ จตุตฺถี ทิสา. อิมาสํ จตสฺสนฺนํ ทิสานํ อาโลกนา. อยํ วุจฺจติ ทิสาโลกโน นาม นโย.}

\prose{ตตฺถายํ โยชนา. จตฺตาโร จ อาหารา จตสฺโส จ ปฏิปทา, จตฺตาโร จ วิปลฺลาสา จตฺตาโร จ สติปฏฺฐานา, จตฺตาริ จ อุปาทานานิ จตฺตาริ จ ฌานานิ จตฺตาโร จ โยคา วิหารา จ, คนฺถา จ สมฺมปฺปธานา จ, อาสวา จ อจฺฉริยา อพฺภุตธมฺมา จ, โอฆา จ อธิฏฺฐานานิ จ, สลฺลา จ สมาธโย, วิญฺญาณ\-ฏฺฐิติโย จตฺตาโร จ สุขภาคิยา ธมฺมา, จตฺตาริ จ อคติคมนานิ จตฺตาริ จ อปฺปมาณานิ อิติ กุสลา\-กุสลานํ ปฏิปกฺข\-วเสน โยชนา, อยํ วุจฺจติ ทิสาโลกโน นโย.}

\prose{ตสฺส จตฺตาริ สามญฺญผลานิ ปริโยสานํ, โย จ ธมฺโม กุสลากุสล\-นิทฺเทเส ปฐโม ทิสานิทฺเทโส, อิมสฺส โสตาปตฺติผลํ ปริโยสานํ ทุติยํ สกทาคามิผลํ, ตติยํ อนาคามิผลํ, จตุตฺถํ อรหตฺตผลํ,}

\prose{ตตฺถ กตโม ติปุกฺขโล นโย, เย จ ทุกฺขาย ปฏิปทาย ทนฺธาภิญฺญาย ขิปฺปาภิญฺญาย จ นิยฺยนฺติ เทฺว ปุคฺคลา, เย จ สุขาย ปฏิปทาย ทนฺธาภิญฺญาย ขิปฺปาภิญฺญาย จ นิยฺยนฺติ เทฺว ปุคฺคลา.}

\prose{อิเมสํ จตุนฺนํ ปุคฺคลานํ\csromansymbolfootnote{*}{สํ 1. 467 ปิฏฺเฐ ปสฺสิตพฺพํ.}โย ปุคฺคโล สุขาย ปฏิปทาย ทนฺธาภิญฺญาย นิยฺยาติ, โย จ ปุคฺคโล ทุกฺขาย ปฏิปทาย ขิปฺปาภิญฺญาย นิยฺยาติ. อิเม เทฺว ปุคฺคลา ภวนฺติ. ตตฺถ โย สุขาย ปฏิปทาย ขิปฺปาภิญฺญาย \csromanfolio{334}นิยฺยาติ, อยํ อุคฺฆฏิตญฺญู. โย ปจฺฉิโม ปุคฺคโล สาธารโณ, อยํ วิปญฺจิตญฺญู. โย ปุคฺคโล ทนฺธาภิญฺญาย ทุกฺขาย ปฏิปทาย นิยฺยาติ, อยํ เนยฺโย. อิเม จตฺตาโร ภวิตฺวา ตีณิ โหนฺติ, ตตฺถ อุคฺฆฏิต\-ญฺญุสฺส สมถ\-ปุพฺพงฺคมา วิปสฺสนา, เนยฺยสฺส วิปสฺสนา\-ปุพฺพงฺคโม สมโถ, วิปญฺจิต\-ญฺญุสฺส สมถ\-วิปสฺสนา ยุคนทฺธา.  อุคฺฆฏิต\-ญฺญุสฺส มุทุกา เทสนา, เนยฺยสฺส ติกฺขา เทสนา, วิปญฺจิต\-ญฺญุสฺส ติกฺขมุทุกา เทสนา.}

\proseitem{8}{อุคฺฆฏิต\-ญฺญุสฺส อธิปญฺญา\-สิกฺขา, เนยฺยสฺส อธิจิตฺต\-สิกฺขา, วิปญฺจิต\-ญฺญุสฺส อธิสีลสิกฺขา. อิติ อิเมสํ ปุคฺคลานํ\csromansharedfootnote{e4f19bab0b970f76}{สํ 1. 452 ปิฏฺเฐ.} จตูหิ ปฏิปทาหิ\csromansharedfootnote{ceb6101ab8b6600f}{สํ 1. 467 ปิฏฺเฐ.} นิยฺยานํ.}

\prose{ตตฺถ อยํ สํกิเลโส, ตีณิ อกุสลมูลานิ ตโย ผสฺสา ติสฺโส เวทนา ตโย อุปวิจารา ตโย สํกิเลสา ตโย วิตกฺกา ตโย ปริฬาหา ตีณิ สงฺขต\-ลกฺขณานิ ติสฺโส ทุกฺขตาติ.}

\prose{ตีณิ อกุสลมูลานีติ โลโภ อกุสลมูลํ, โทโส อกุสลมูลํ, โมโห อกุสลมูลํ. ตโย ผสฺสาติ สุขเวทนีโย ผสฺโส, ทุกฺขเวทนีโย ผสฺโส, อทุกฺขม\-สุข\-เวทนีโย ผสฺโส. ติสฺโส เวทนาติ สุขา เวทนา ทุกฺขา เวทนา อทุกฺขมสุขา เวทนา. ตโย อุปวิจาราติ โสมนสฺโส\-ปวิจาโร โทมนสฺโส\-ปวิจาโร อุเปกฺโข\-ปวิจาโร. ตโย สํกิเลสาติ ราโค โทโส โมโห. ตโย วิตกฺกาติ กามวิตกฺโก พฺยาปาทวิตกฺโก วิหิํสาวิตกฺโก. ตโย ปริฬาหาติ ราคโช โทสโช โมหโช. ตีณิ สงฺขตลกฺขณานีติ อุปฺปาโท ฐิติ วโย. ติสฺโส ทุกฺขตาติ ทุกฺขทุกฺขตา วิปริณาม\-ทุกฺขตา สงฺขต\-ทุกฺขตา.}

\prose{ตตฺถ โลโภ อกุสลมูลํ กุโต สมุฏฺฐิตํ, ติวิธํ อารมฺมณํ มนาปิกํ อมนาปิกํ อุเปกฺขาฐานิยญฺจ. ตตฺถ มนาปิเกน อารมฺมเณน โลโภ อกุสลมูลํ สมุฏฺฐหติ. อิติ มนาปิกา อารมฺมณา สุขเวทนีโย ผสฺโส, สุขเวทนียํ ผสฺสํ ปฏิจฺจ อุปฺปชฺชเต สุขเวทนา, สุขเวทนํ ปฏิจฺจ อุปฺปชฺชเต โสมนสฺสู\-ปวิจาโร, โสมนสฺสู\-ปวิจารํ ปฏิจฺจ อุปฺปชฺชเต ราโค, ราคํ ปฏิจฺจ อุปฺปชฺชเต กามวิตกฺโก, กามวิตกฺกํ ปฏิจฺจ อุปฺปชฺชเต ราคโช ปริฬาโห, ราคชํ ปริฬาหํ ปฏิจฺจ อุปฺปชฺชเต อุปฺปาโท สงฺขต\-ลกฺขโณ, อุปฺปาทํ สงฺขต\-ลกฺขณํ ปฏิจฺจ อุปฺปชฺชเต วิปริณาม\-ทุกฺขตา.}

\prose{\csromanfolio{335}โทโส อกุสลมูลํ กุโต สมุฏฺฐิตํ, อมนาปิเกน อารมฺมเณน โทโส อกุสลมูลํ สมุฏฺฐิตํ. อิติ อมนาปิกา อารมฺมณา ทุกฺขเวทนีโย ผสฺโส, ทุกฺข\-เวทนียํ ผสฺสํ ปฏิจฺจ อุปฺปชฺชเต ทุกฺขเวทนา, ทุกฺขเวทนํ ปฏิจฺจ อุปฺปชฺชเต โทมนสฺสู\-ปวิจาโร, โทมนสฺสู\-ปวิจารํ ปฏิจฺจ อุปฺปชฺชเต โทโส, โทสํ ปฏิจฺจ อุปฺปชฺชเต พฺยาปาทวิตกฺโก, พฺยาปาท\-วิตกฺกํ ปฏิจฺจ อุปฺปชฺชเต โทสโช ปริฬาโห, โทสชํ ปริฬาหํ ปฏิจฺจ อุปฺปชฺชเต ฐิตสฺส อญฺญถตฺตํ สงฺขต\-ลกฺขณํ, ฐิตสฺส อญฺญถตฺตํ สงฺขต\-ลกฺขณํ ปฏิจฺจ อุปฺปชฺชเต ทุกฺขทุกฺขตา เวทนา.}

\prose{โมโห อกุสลมูลํ กุโต สมุฏฺฐิตํ, อุเปกฺขา\-ฐานิเยน อารมฺมเณน โมโห อกุสลมูลํ สมุฏฺฐิตํ. อิติ อุเปกฺขาฐานิยา อารมฺมณา อทุกฺขม\-สุข\-เวทนีโย ผสฺโส, อทุกฺขม\-สุข\-เวทนียํ ผสฺสํ ปฏิจฺจ อุปฺปชฺชเต อทุกฺขมสุขา เวทนา, อทุกฺขมสุข\-เวทนํ ปฏิจฺจ อุปฺปชฺชเต อุเปกฺขู\-ปวิจาโร, อุเปกฺขู\-ปวิจารํ ปฏิจฺจ อุปฺปชฺชเต โมโห, โมหํ ปฏิจฺจ อุปฺปชฺชเต วิหิํสาวิตกฺโก, วิหิํสา\-วิตกฺกํ ปฏิจฺจ อุปฺปชฺชเต โมหโช ปริฬาโห, โมหชํ ปริฬาหํ ปฏิจฺจ อุปฺปชฺชเต วโย สงฺขต\-ลกฺขณํ, วยํ สงฺขต\-ลกฺขณํ ปฏิจฺจ อุปฺปชฺชเต สงฺขต\-ทุกฺขตา, อิติ อยํ ติณฺณํ กิเลสานํ นิทฺเทโส, อยํ วุจฺจเต กุสลปกฺเข ติปุกฺขโล นโย.}

\prose{อิติ ตีณิ อกุสลมูลานิ\csromansharedfootnote{990f699f491d594c}{อํ 1. 202, 3 ปิฏฺเฐสุ.} น จตุตฺถานิ น ปญฺจมานิ, ตโย ผสฺสาติ ติสฺโส เวทนา ยาว สงฺขต\-ทุกฺขตาติ, โย โกจิ อกุสลปกฺโข, สพฺโพ โส ตีสุ อกุสลมูเลสุ สโมสรติ.}

\prose{ตตฺถ กตโม กุสลปกฺโข, ตีณิ กุสลมูลานิ\csromansharedfootnote{62aa02457b066b5c}{อํ 1. 237 ปิฏฺเฐ.}, ติสฺโส ปญฺญา สุตมยี ปญฺญา จินฺตามยี ปญฺญา ภาวนามยี ปญฺญา. ตโย สมาธี สวิตกฺก\-สวิจาโร ฯเปฯ. ติสฺโส สิกฺขา\csromansharedfootnote{990f699f491d594c}{อํ 1. 202, 3 ปิฏฺเฐสุ.} อธิสีลสิกฺขา ฯเปฯ สิกฺขา. ตีณิ นิมิตฺตานิ สมถ\-นิมิตฺตํ ปคฺคห\-นิมิตฺตํ อุเปกฺขา\-นิมิตฺตํ. ตโย วิตกฺกา เนกฺขมฺม\-วิตกฺโก ฯเปฯ อวิหิํสา\-วิตกฺโก. ตีณิ อินฺทฺริยานิ อนญฺญาต\-ญฺญสฺสามี\-ตินฺทฺริยนฺติ วิตฺถาโร. ตโย อุปวิจารา เนกฺขมฺมู\-ปวิจาโร อพฺยาปาทู\-ปวิจาโร อวิหิํสู\-ปวิจาโร. ติสฺโส เอสนา กาเมสนา ภเวสนา พฺรหฺมจริเย\-สนา. ตโย ขนฺธา สีลกฺขนฺโธ สมาธิ\-กฺขนฺโธ ปญฺญากฺขนฺโธ.}

\prose{\csromanfolio{336}ตตฺถ ยํ อโลโภ กุสลมูลํ, ตํ สุตมยิปญฺญํ ปริปูเรติ, สุตมยี ปญฺญา ปริปุณฺณา สวิตกฺกํ สวิจารํ สมาธิํ ปริปูเรติ, สวิตกฺโก สวิจาโร สมาธิ ปริปุณฺโณ อธิจิตฺต\-สิกฺขํ ปริปูเรติ, อธิจิตฺต\-สิกฺขา ปริปุณฺณา สมถ\-นิมิตฺตํ ปริปูเรติ, สมถ\-นิมิตฺตํ ปริปุณฺณํ เนกฺขมฺม\-วิตกฺกํ ปริปูเรติ, เนกฺขมฺม\-วิตกฺโก ปริปุณฺโณ อนญฺญาต\-ญฺญสฺสามี\-ตินฺทฺริยํ ปริปูเรติ, อนญฺญาต\-ญฺญสฺสามี\-ตินฺทฺริยํ ปริปุณฺณํ เนกฺขมฺมู\-ปวิจารํ ปริปูเรติ, เนกฺขมฺมู\-ปวิจาโร ปริปุณฺโณ กาเมสนํ ปชหติ, กาเมสน\-ปฺปหานํ สมาธิ\-กฺขนฺธํ ปริปูเรติ.}

\prose{อโทโส กุสลมูลํ จินฺตามยิ\-ปญฺญํ ปริปูเรติ, จินฺตามยี ปญฺญา ปริปุณฺณา อวิตกฺก\-วิจารมตฺตํ สมาธิํ ปริปูเรติ. อวิตกฺก\-วิจารมตฺโต สมาธิ ปริปุณฺโณ อธิสีล\-สิกฺขํ ปริปูเรติ, อธิสีลสิกฺขา ปริปุณฺณา อุเปกฺขา\-นิมิตฺตํ ปริปูเรติ, อุเปกฺขา\-นิมิตฺตํ ปริปุณฺณํ อพฺยาปาทวิตกฺกํ ปริปูเรติ, อพฺยาปาทวิตกฺโก ปริปุณฺโณ อญฺญินฺทฺริยํ ปริปูเรติ, อญฺญินฺทฺริยํ ปริปุณฺณํ อพฺยาปาทู\-ปวิจารํ ปริปูเรติ, อพฺยาปาทู\-ปวิจาโร ปริปุณฺโณ ภเวสนํ ปชหติ, ภเวสน\-ปฺปหานํ สีลกฺขนฺธํ ปริปูเรติ.}

\prose{อโมโห กุสลมูลํ ภาวนามยิ\-ปญฺญํ ปริปูเรติ, ภาวนามยี\-ปญฺญา ปริปุณฺณา อวิตกฺก\-อวิจารํ สมาธิํ ปริปูเรติ, อวิตกฺโก อวิจาโร สมาธิ ปริปุณฺโณ อธิปญฺญา\-สิกฺขํ ปริปูเรติ, อธิปญฺญา\-สิกฺขา ปริปุณฺณา ปคฺคห\-นิมิตฺตํ ปริปูเรติ, ปคฺคห\-นิมิตฺตํ ปริปุณฺณํ อญฺญาตาวิโน อินฺทฺริยํ ปริปูเรติ, อญฺญาตาวิโน อินฺทฺริยํ ปริปุณฺณํ อวิหํ สูปวิจารํ ปริปูเรติ, อวิหิํสู\-ปวิจาโร ปริปุณฺโณ พฺรหฺมจริเยสนํ ปริปูเรติ, พฺรหฺมจริเย\-สนา ปริปุณฺณา ปญฺญากฺขนฺธํ ปริปูเรติ.}

\prose{อิติ อิเม ตโย ธมฺมา กุสลปกฺขิกา สพฺเพ กุสลา ธมฺมา ตีหิ ติกนิทฺเทเสหิ นิทฺทิสิยนฺติ ตีณิ วิโมกฺขมุขานิ ตสฺส ปริโยสานํ. ตตฺถ ปฐเมน อปฺปณิหิตํ, ทุติเยน สุญฺญตํ, ตติเยน อนิมิตฺตํ. อยํ วุจฺจติ ทุติโย ติปุกฺขโล นาม นโย.}

\prose{ตตฺถ เย อิเม ตโย ปุคฺคลา อุคฺฆฏิตญฺญู วิปญฺจิตญฺญู เนยฺโยติ. อิเมสํ ติณฺณํ ปุคฺคลานํ เย จ ปุคฺคลา สุขาย ปฏิปทาย ขิปฺปาภิญฺญาย สุขาย ปฏิปทาย ทนฺธาภิญฺญาย จ นิยฺยนฺติ, เต เทฺว ปุคฺคลา. เย จ เทฺว ปุคฺคลา ทุกฺขาย ปฏิปทาย ขิปฺปาภิญฺญาย ทุกฺขาย ปฏิปทาย ทนฺธาภิญฺญาย จ \csromanfolio{337}นิยฺยนฺติ, อิเม จตฺตาโร เตน วิเสเสน เทฺว ภวนฺติ ทิฏฺฐิจริโต จ ตณฺหาจริโต จ. อิเม จตฺตาโร ภวิตฺวา ตโย ภวนฺติ, ตโย ภวิตฺวา เทฺว ภวนฺติ. อิเมสํ ทฺวินฺนํ ปุคฺคลานํ อยํ สํกิเลโส, อวิชฺชา จ ตณฺหา จ, อหิริกญฺจ อโนตฺตปฺปญฺจ, อสฺสติ จ อสมฺปชญฺญญฺจ, นีวรณานิ จ สํโยชนานิ จ, อชฺโฌสานญฺจ อภินิเวโส จ, อหํกาโร จ มมํกาโร จ, อสฺสทฺธิยญฺจ โทวจสฺสญฺจ, โกสชฺชญฺจ อโยนิโส จ มนสิกาโร, วิจิกิจฺฉา จ อภิชฺฌา จ, อสทฺธมฺมสฺสวนญฺจ อสมาปตฺติ จ.}

\proseitem{8}{ตตฺถ อวิชฺชา จ อหิริกญฺจ อสฺสติ จ นีวรณานิ จ อชฺโฌสานญฺจ อหํกาโร จ อสฺสทฺธิยญฺจ โกสชฺชญฺจ วิจิกิจฺฉา จ อสทฺธมฺมสฺสวนญฺจ, อยํ เอกา ทิสา.}

\prose{ตณฺหา จ อโนตฺตปฺปญฺจ อสมฺปชญฺญญฺจ สํโยชนานิ จ อภินิเวโส จ มมํกาโร จ โทวจสฺสตา จ อโยนิโส มนสิกาโร จ อภิชฺฌา จ อสมาปตฺติ จ, อยํ ทุติยา ทิสา. ทสนฺนํ ทุกานํ ทส ปทานิ ปฐมานิ กาตพฺพานิ. สํขิตฺเตน อตฺถํ ญาเปนฺติ ปฏิปกฺเข กณฺหปกฺขสฺส สพฺเพสํ ทุกานํ ทส ปทานิ ทุติยกานิ, อยํ ทุติยา ทิสา.}

\prose{อิติ อกุสลานํ ธมฺมานํ ทุกฺขนิทฺเทโส, อยํ สมุทโย. ยํ ตํ ธมฺมํ อชฺฌาวสติ นามญฺจ รูปญฺจ อิทํ ทุกฺขํ อิติ อยญฺจ สมุทโย, อิทญฺจ ทุกฺขํ, อิมานิ เทฺว สจฺจานิ ทุกฺขญฺจ สมุทโย จ นนฺทิยา\-วฏฺฏสฺส นยสฺส ปฐมนิทฺเทโส.}

\prose{ตตฺถ กตโม กุสลปกฺโข, สมโถ จ วิปสฺสนา จ, วิชฺชา จ จรณญฺจ, สติ จ สมฺปชญฺญญฺจ, หิรี จ โอตฺตปฺปญฺจ, อหํการ\-ปฺปหานญฺจ มมํการ\-ปฺปหานญฺจ, สมฺมาวายาโม จ โยนิโส จ มนสิกาโร, สมฺมาสติ จ สมฺมาสมาธิ จ, ปญฺญา จ นิพฺพิทา จ, สมาปตฺติ จ สทฺธมฺมสฺสวนญฺจ, โสมนสฺสญฺจ ธมฺมา\-นุธมฺม\-ปฺปฏิปตฺติ จ.}

\prose{ตตฺถ สมโถ จ วิชฺชา จ สติ จ หิรี จ อหํการ\-ปฺปหานญฺจ สมฺมาวายาโม จ สมฺมาสติ จ ปญฺญา จ สมาปตฺติ จ โสมนสฺสญฺจ, อิเม ธมฺมา เอกา ทิสา. วิปสฺสนา จ จรณญฺจ สมฺปชญฺญญฺจ โอตฺตปฺปญฺจ มมํการ\-ปฺปหานญฺจ โยนิโส มนสิกาโร จ สมฺมาสมาธิ จ นิพฺพิทา จ สทฺธมฺมสฺสวนญฺจ ธมฺมา\-นุธมฺม\-ปฺปฏิปตฺติ จ, อยํ ทุติยา ทิสา. อิติ กุสลปกฺเข จ อกุสลปกฺเข จ นนฺทิยา\-วฏฺฏสฺส ปน นยสฺส จตสฺโส ทิสา.}

\prose{\csromanfolio{338}ตาสุ ยานิ อกุสล\-ปกฺขสฺส ปฐมานิ ปทานิ อกุสลานิ กุสเลหิ ปหานํ คจฺฉนฺติ, ตานิ กุสลปกฺเข ทุติเยหิ ปเทหิ ปหานํ คจฺฉนฺติ. เตสํ ปหานา ราควิราคา เจโตวิมุตฺติ ยานิ อกุสล\-ปกฺขสฺส ทุติยานิ อกุสลปทานิ ปหานํ คจฺฉนฺติ, ตานิ กุสล\-ปกฺขสฺส ปฐเมหิ ปเทหิ ปหานํ คจฺฉนฺติ. เตสํ ปหานา อวิชฺชาวิราคา ปญฺญาวิมุตฺติ ปริโยสานํ. อิเมสํ ติณฺณํ นยานํ ปฐโม นโย สีหวิกฺกีฬิโต นาม. อฏฺฐ ปทานิ จตฺตาริ จ กุสลานิ จตฺตาริ จ อกุสลานิ อิมานิ อฏฺฐ ปทานิ มูลปทานิ, อตฺถนเยน ทุติโย ติปุกฺขโล. โส ฉหิ ธมฺเมหิ เนติ กุสลมูลานิ จ เนติ, อกุสลมูลานิ จ, อิติ อิมานิ ฉ ปทานิ ปุริมกานิ จ อฏฺฐ มูลปทานิ อิมานิ จุทฺทส ปทานิ อฏฺฐารสนฺนํ มูลปทานํ. ตตฺถ โย ปจฺฉิมโก นโย นนฺทิยาวฏฺโฏ, โส จตูหิ ธมฺเมหิ เนติ. อวิชฺชาย จ ตณฺหาย จ สมเถน จ วิปสฺสนาย จ, อิเม จตฺตาโร ธมฺมา อิมานิ อฏฺฐารส มูลปทานิ ตีสุ นเยสุ นิทฺทิฏฺฐานิ.}

\prose{ตตฺถ ยานิ นว ปทานิ กุสลานิ, ตตฺถ สพฺพํ กุสลํ สโมสรติ. เตสญฺจ นวนฺนํ มูลานํ จตฺตาริ ปทานิ สีหวิกฺกีฬิต\-นเย ตีณิ ติปุกฺขเล เทฺว นนฺทิยาวฏฺเฏ, อิจฺเจเต กุสลสฺส ปกฺขา. ตตฺถ ยานิ นว ปทานิ กุสลานิ, ตตฺถ สพฺพํ กุสลํ ยุชฺชติ. ตตฺถ สีหวิกฺกีฬิเต นเย จตฺตาริ ปทานิ ตีณิ ติปุกฺขเล เทฺว นนฺทิยาวฏฺเฏ อิมานิ นว ปทานิ กุสลานิ นิทฺทิฏฺฐานิ.}

\prose{ตตฺถ ยานิ นนฺทิยาวฏฺเฏ นเย จตฺตาริ ปทานิ, ตตฺถ อฏฺฐารส มูลปทานิ สโมสรนฺติ. ยถา กถํ, สมโถ จ อโลโภ จ อโทโส จ อสุภสญฺญา จ ทุกฺขสญฺญา จ อิมานิ กุสลปกฺเข ปญฺจ ปทานิ สมถํ ภชนฺติ. วิปสฺสนา จ อโมโห จ อนิจฺจสญฺญา จ อนตฺตสญฺญา จ อิมานิ จตฺตาริ ปทานิ วิปสฺสนํ ภชนฺติ. อิมานิ นว ปทานิ กุสลานิ ทฺวีสุ ปเทสุ โยชิตานิ, ตตฺถ อกุสลปกฺเข นวนฺนํ อกุสล\-มูลปทานํ ยา จ ตณฺหา โย จ โลโภ โย จ โทโส ยา จ สุภสญฺญา ยา จ สุขสญฺญา, อิมานิ ปญฺจ ปทานิ ตณฺหํ ภชนฺติ. ยา จ อวิชฺชา โย จ โมโห ยา จ นิจฺจสญฺญา ยา จ อตฺตสญฺญา, อิมานิ จตฺตาริ ปทานิ อวิชฺชํ ภชนฺติ. เอตานิ นว ปทานิ อกุสลานิ สุสํขิตฺตานิ. อิติ ตโย นยา เอกํ นยํ น ปวิฏฺฐา. เอวํ อฏฺฐารส มูลปทานิ นนฺทิยา\-วฏฺฏ\-นเย นิทฺทิสิตพฺพานิ.}

\prose{\csromanfolio{339}กถํ อฏฺฐารส มูลปทานิ, ติปุกฺขเล นเย ยุชฺชนฺติ, นวนฺนํ ปทานํ กุสลานํ วิปสฺสนา จ อโมโห จ อนิจฺจสญฺญา จ อนตฺตสญฺญา จ, อิมานิ จตฺตาริ ปทานิ. อโมโห จ สมโถ จ อโลโภ จ อสุภสญฺญา จ, อิมานิ จตฺตาริ ปทานิ. โลโภ จ โทโส จ, เอวํ อิมานิ นว ปทานิ ตีสุ กุสเลสุ โยเชตพฺพานิ. ตตฺถ นวนฺนํ ปทานํ อกุสลานํ ตณฺหา จ โลโภ จ สุภสญฺญา จ สุขสญฺญา จ, อิมานิ จตฺตาริ ปทานิ โลโภ อกุสลมูลํ. อวิชฺชา จ โมโห จ นิจฺจสญฺญา จ อตฺตสญฺญา จ อยํ โมโห อยํ โทโส, เย จ อิมานิ นว ปทานิ ตีสุ อกุสเลสุ โยชิตานิ. เอวํ อฏฺฐารส มูลปทานิ กุสลมูเลสุ จ โยเชตฺวา ติปุกฺขเลน นเยน นิทฺทิสิตพฺพานิ.}

\prose{กถํ อฏฺฐารส มูลปทานิ สีหวิกฺกีฬิเต นเย ยุชฺชนฺติ, ตณฺหา จ สุภสญฺญา จ, อยํ ปฐโม วิปลฺลาโส. โลโภ จ สุขสญฺญา จ, อยํ ทุติโย วิปลฺลาโส. อวิชฺชา จ นิจฺจสญฺญา จ, อยํ ตติโย วิปลฺลาโส. โมโห จ อตฺตสญฺญา จ, อยํ จ ตุตฺโถ วิปลฺลาโส. อิติ นว ปทานิ อกุสลมูลานิ จตูสุ ปเทสุ โยชิตานิ. ตตฺถ นวนฺนํ มูลปทานํ กุสลานํ สมโถ จ อสุภสญฺญา จ, อิทํ ปฐมํ สติปฏฺฐานํ. อโลโภ จ ทุกฺขสญฺญา จ, อิทํ ทุติยํ สติปฏฺฐานํ. วิปสฺสนา จ อนิจฺจสญฺญา จ, อิทํ ตติยํ สติปฏฺฐานํ. อโมโห จ อนตฺตสญฺญา จ, อิทํ จตุตฺถํ สติปฏฺฐานํ. อิมานิ อฏฺฐารส มูลปทานิ สีหวิกฺกีฬิตนยํ อนุปวิฏฺฐานิ. อิเมสํ ติณฺณํ นยานํ ยา ภูมิ จ โย ราโค จ โย โทโส จ เอกํ นยํ ปวิสติ. เอกสฺส นยสฺส อกุสเล วา ธมฺเม กุสเล วา ธมฺเม วิญฺญาเต ปฏิปกฺโข อเนฺวสิตพฺโพ, ปฏิปกฺเข อเนฺวสิตฺวา โส นโย นิทฺทิสิตพฺโพ, ตมฺหิ นเย นิทฺทิฏฺโฐ. ยถา เอกมฺหิ นเย สพฺเพ นยา ปวิฏฺฐา ตถา นิทฺทิสิตพฺพา. เอกมฺหิ จ นเย อฏฺฐารส มูลปทานิ ปวิฏฺฐานิ, ตมฺหิ ธมฺเม วิญฺญาเต สพฺเพ ธมฺมา วิญฺญาตา โหนฺติ. อิเมสํ ติณฺณํ นยานํ สีหวิกฺกีฬิต\-นยสฺส จตฺตาริ ผลานิ ปริโยสานํ. ปฐมาย ทิสาย ปฐมํ ผลํ, ทุติยาย ทิสาย ทุติยํ ผลํ, ตติยาย ทิสาย ตติยํ ผลํ, จตุตฺถาย ทิสาย จตุตฺถํ ผลํ. ติปุกฺขลสฺส นยสฺส ตีณิ วิโมกฺขมุขานิ ปริโยสานํ. ปฐมาย ทิสาย อปฺปณิหิตํ, ทุติยาย ทิสาย สุญฺญตํ, ตติยาย ทิสาย อนิมิตฺตํ. นนฺทิยา\-วฏฺฏสฺส นยสฺส ราควิราคา เจโตวิมุตฺติ อวิชฺชาวิราคา จ ปญฺญาวิมุตฺติ}

\prose{\csromanfolio{340}ปริโยสานํ. ปฐมาย ทิสาย ราควิราคา เจโตวิมุตฺติ, ทุติยาย ทิสาย อวิชฺชาวิราคา ปญฺญาวิมุตฺติ. อิเม ตโย นยา อิเมสํ ติณฺณํ นยานํ อฏฺฐารสนฺนํ มูลปทานํ อาโลกนา, อยํ วุจฺจติ ทิสาโลกโน\csromansharedfootnote{84fa83d40b40d8c8}{ทิสาโลจโน (ก)} นโย.  อาโลเกตฺวาน ชานาติ “อยํ ธมฺโม อิมํ ธมฺมํ ภชตี”ติ สมฺมา โยชนา. กุสลปกฺเข อกุสลปกฺเข จ อยํ นโย องฺกุโก นาม. อิเม ปญฺจ นยา.}

\csromanheader{2}{\textbf{ตตฺถิมา อุทฺทานคาถา}}
\setlength{\csromangathapagewidth}{0pt}
\setlength{\csromangathawakwidth}{0pt}
\csromangathameasuregroup{\textsuperscript{*}ตณฺหา จ อวิชฺชาปิ จ, \\ จตฺตาโร จ วิปลฺลาสา, \\ เย จ สติปฏฺฐานา, \\ เอตํ สพฺพํ กุสลํ, \\ สพฺพกุสลํ นวหิ ปเทหิ ยุชฺชติ, \\ เอเต เต มูลปทา, \\ ตณฺหา เจว อวิชฺชา จ, \\ โย เนติ สพฺเพสุ โยคยุตฺโต, \\ ยํ กุสลมูเลหิ, \\ ภูตํ ตถํ อวิตถํ, \\ โส เนติ วิปลฺลาเสหิ, \\ ธมฺเม ตํ นยํ วินยมาหุ, \\ เวยฺยากรเณ วุตฺเต, \\ ตโย อาโลกยติ, \\ ** โอโลเกตฺวา ทิสาโลจเนน, \\ สพฺเพ กุสลากุสเล,}{\csromangathabat{\textsuperscript{*}ตณฺหา จ อวิชฺชาปิ จ,}{โลโภ โทโส ตเถว โมโห จ.} \\ \csromangathabat{จตฺตาโร จ วิปลฺลาสา,}{กิเลสภูมี นว ปทานิ.} \\ \csromangathabat{เย จ สติปฏฺฐานา,}{สมโถ จ วิปสฺสนา กุสลมูลา.} \\ \csromangathabat{เอตํ สพฺพํ กุสลํ,}{อินฺทฺริยภูมี นว ปทานิ.} \\ \csromangathabat{สพฺพกุสลํ นวหิ ปเทหิ ยุชฺชติ,}{นวหิ เจว อกุสลํ.} \\ \csromangathabat{เอเต เต มูลปทา,}{อุภโต อฏฺฐารส ปทานิ.} \\ \csromangathabat{ตณฺหา เจว อวิชฺชา จ,}{สมโถ จ วิปสฺสนา.} \\ \csromangathabat{โย เนติ สพฺเพสุ โยคยุตฺโต,}{อยํ นโย นนฺทิยาวฏฺโฏ.} \\ \csromangathabat{ยํ กุสลมูเลหิ,}{นยติ กุสลอกุสลมูเลหิ.} \\ \csromangathabat{ภูตํ ตถํ อวิตถํ,}{ติปุกฺขลํ ตํ นยํ อาหุ.} \\ \csromangathabat{โส เนติ วิปลฺลาเสหิ,}{กิเลสอินฺทฺริเยหิ จ.} \\ \csromangathabat{ธมฺเม ตํ นยํ วินยมาหุ,}{สีหวิกฺกีฬิตํ นาม.} \\ \csromangathabat{เวยฺยากรเณ วุตฺเต,}{กุสลตาหิ อกุสลตาหิ จ.} \\ \csromangathabat{ตโย อาโลกยติ,}{อยํ นโย ทิสาโลจโน นาม.} \\ \csromangathabat{** โอโลเกตฺวา ทิสาโลจเนน,}{อุกฺขิปิย ยํ สมาเนติ.} \\ \csromangathabat{สพฺเพ กุสลากุสเล,}{อยํ นโย องฺกุโส นาม.} \\ นยสมุฏฺฐานํ.}
\csromangathameasurewak{นยสมุฏฺฐานํ.}
\csromangathasetleft{\textsuperscript{*}ตณฺหา จ อวิชฺชาปิ จ, \\ จตฺตาโร จ วิปลฺลาสา, \\ เย จ สติปฏฺฐานา, \\ เอตํ สพฺพํ กุสลํ, \\ สพฺพกุสลํ นวหิ ปเทหิ ยุชฺชติ, \\ เอเต เต มูลปทา, \\ ตณฺหา เจว อวิชฺชา จ, \\ โย เนติ สพฺเพสุ โยคยุตฺโต, \\ ยํ กุสลมูเลหิ, \\ ภูตํ ตถํ อวิตถํ, \\ โส เนติ วิปลฺลาเสหิ, \\ ธมฺเม ตํ นยํ วินยมาหุ, \\ เวยฺยากรเณ วุตฺเต, \\ ตโย อาโลกยติ, \\ ** โอโลเกตฺวา ทิสาโลจเนน, \\ สพฺเพ กุสลากุสเล,}
\csromangathagroupwithcloser{\csromangathastanza{\csromangathabat{\csromansymbolfootnote{*}{เหฏฺฐา 170 ปิฏฺเฐ ปสฺสิตพฺพํ. ** เหฏฺฐา 5 ปิฏฺเฐ.}ตณฺหา จ อวิชฺชาปิ จ,}{โลโภ โทโส ตเถว โมโห จ.} \\ \csromangathabat{จตฺตาโร จ วิปลฺลาสา,}{กิเลสภูมี นว ปทานิ.}}\csromangathastanza{\csromangathabat{เย จ สติปฏฺฐานา,}{สมโถ จ วิปสฺสนา กุสลมูลา.} \\ \csromangathabat{เอตํ สพฺพํ กุสลํ,}{อินฺทฺริยภูมี นว ปทานิ.}}\csromangathastanza{\csromangathabat{สพฺพกุสลํ นวหิ ปเทหิ ยุชฺชติ,}{นวหิ เจว อกุสลํ.} \\ \csromangathabat{เอเต เต มูลปทา,}{อุภโต อฏฺฐารส ปทานิ.}}\csromangathastanza{\csromangathabat{ตณฺหา เจว อวิชฺชา จ,}{สมโถ จ วิปสฺสนา.} \\ \csromangathabat{โย เนติ สพฺเพสุ โยคยุตฺโต,}{อยํ นโย นนฺทิยาวฏฺโฏ.}}\csromangathastanza{\csromangathabat{ยํ กุสลมูเลหิ,}{นยติ กุสลอกุสลมูเลหิ.} \\ \csromangathabat{ภูตํ ตถํ อวิตถํ,}{ติปุกฺขลํ ตํ นยํ อาหุ.}}\csromangathastanza{\csromangathabat{โส เนติ วิปลฺลาเสหิ,}{กิเลสอินฺทฺริเยหิ จ.} \\ \csromangathabat{ธมฺเม ตํ นยํ วินยมาหุ,}{สีหวิกฺกีฬิตํ นาม.}}\csromangathastanza{\csromangathabat{เวยฺยากรเณ วุตฺเต,}{กุสลตาหิ อกุสลตาหิ จ.} \\ \csromangathabat{ตโย อาโลกยติ,}{อยํ นโย ทิสาโลจโน นาม.}}\csromangathastanza{\csromangathabat{** โอโลเกตฺวา ทิสาโลจเนน,}{อุกฺขิปิย ยํ สมาเนติ.} \\ \csromangathabat{สพฺเพ กุสลากุสเล,}{อยํ นโย องฺกุโส นาม.}}\csromangathastanza{\csromangathawak{นยสมุฏฺฐานํ.}}}{เปฏโกปเทเส มหา\-กจฺจายนสฺส เถรสฺส สุตฺตวิภงฺคสฺส\csromansharedfootnote{338dd9d64fd7ee1e}{[มิสฺสินฺคฺ โนเต 0]} ทสฺสนํ สมตฺตํ.}
\prosewithcloser{\csromanfolio{341}ยานิ จตุกฺกานิ อกุสลานิ กุสลานิ จ สีหวิกฺกีฬิเต นเย นิทฺทิฏฺฐานิ, ติกานิ กุสลานิ จ อกุสลานิ จ ติปุกฺขเล นเย นิทฺทิฏฺฐานิ, ทุกานิ กุสลานิ จ อกุสลานิ จ นนฺทิยาวฏฺเฏ นเย นิทฺทิฏฺฐานิ. เยสุ ทฺวีสุ ธมฺเมสุ\csromansharedfootnote{e08e9761b27f701d}{วิสุทฺธีสุ (ก)} กุสเลสุ โส อตฺโถ ติเกสุ วิภชฺช\-มานสฺส ภวภูมิ, อถ จ สพฺโพ\csromansharedfootnote{23ed088967625ccc}{ปุพฺโพ (ก)} จ อตฺโถ ตีหิ พฺยญฺชเนหิ นิทฺทิสติ. ตตฺตกานิ วุจฺจติ. โย อตฺโถ จตูหิ ปเทหิ อฏฺฐวีส\-ติภาเคหิ นตฺถิภูมิ นิทฺทิสิตุํ, อวจรนฺโตว จตูหิ ปเทหิ นิทฺทิสติ. อิติ ยํ ยถา\-นิทฺทิฏฺฐสฺส อวิโกสนา อิทํ ปมาณํ. ยถา สพฺเพ สมาธโย ตีสุ สมาธีสุ ปริเยสิตพฺพา, สวิตกฺก\-สวิจาเร อวิตกฺก\-วิจารมตฺเต อวิตกฺก\-อวิจาเร อิทํ ปมาณํ, นตฺถิ จตุตฺโถ สมาธิ. ตถา ติสฺโส ปญฺญา จินฺตามยี สุตมยี ภาวนามยี สพฺพาสุ ปญฺญาสุ นิทฺทิสติ, นตฺถิ จตุตฺถี ปญฺญา น จินฺตามยี น สุตมยี น ภาวนามยี, ปญฺญา นาสฺส อตฺถิ อิเมสํ ธมฺมานํ ยา อวิกฺเขปนา, อิทํ วุจฺจติ ปมาณนฺติ.}{เถรสฺส มหา\-กจฺจายนสฺส ชมฺพุ\-วนวาสิโน เปฏโกปเทโส สมตฺโต.}
\nitthitam{\textbf{เปฏโกปเทส}\-\textbf{ปกรณํ นิฏฺฐิตํ.}}
\orphannote{ธมฺมจกฺกปฺปวตฺตน\-สุตฺตํ วิ 3. 14; สํ 3. 368 ปิฏฺฐาทีสุ ปสฺสิตพฺพํ.}

\orphannote{อํ 2. 252 ปิฏฺเฐ. ** สํ 2. 543 ปิฏฺเฐ.}

\orphannote{วิ 3. 15, 16; สํ 3. 369 ปิฏฺเฐสุ ปสฺสิตพฺพํ.}

\orphannote{ม 1. 103 ปิฏฺเฐปิ ปสฺสิตพฺพํ.}

\orphannote{อํ 1. 28, 31 ปิฏฺเฐสุปิ ปสฺสิตพฺพํ.}

\orphannote{สํ 1. 98 ปิฏฺเฐ; อุปริ 174 ปิฏฺเฐ ปน โถกํ วิสทิสํ.}

\orphannote{ที 3. 217; อํ 2. 338 ปิฏฺเฐสุปิ ปสฺสิตพฺพํ,}

\orphannote{** สํ 1. 224, 225 ปิฏฺเฐสุ.}

